%% Generated by Sphinx.
\def\sphinxdocclass{report}
\documentclass[letterpaper,10pt,english]{sphinxmanual}
\ifdefined\pdfpxdimen
   \let\sphinxpxdimen\pdfpxdimen\else\newdimen\sphinxpxdimen
\fi \sphinxpxdimen=.75bp\relax
\ifdefined\pdfimageresolution
    \pdfimageresolution= \numexpr \dimexpr1in\relax/\sphinxpxdimen\relax
\fi
%% let collapsible pdf bookmarks panel have high depth per default
\PassOptionsToPackage{bookmarksdepth=5}{hyperref}

\PassOptionsToPackage{booktabs}{sphinx}
\PassOptionsToPackage{colorrows}{sphinx}

\PassOptionsToPackage{warn}{textcomp}
\usepackage[utf8]{inputenc}
\ifdefined\DeclareUnicodeCharacter
% support both utf8 and utf8x syntaxes
  \ifdefined\DeclareUnicodeCharacterAsOptional
    \def\sphinxDUC#1{\DeclareUnicodeCharacter{"#1}}
  \else
    \let\sphinxDUC\DeclareUnicodeCharacter
  \fi
  \sphinxDUC{00A0}{\nobreakspace}
  \sphinxDUC{2500}{\sphinxunichar{2500}}
  \sphinxDUC{2502}{\sphinxunichar{2502}}
  \sphinxDUC{2514}{\sphinxunichar{2514}}
  \sphinxDUC{251C}{\sphinxunichar{251C}}
  \sphinxDUC{2572}{\textbackslash}
\fi
\usepackage{cmap}
\usepackage[T1]{fontenc}
\usepackage{amsmath,amssymb,amstext}
\usepackage{babel}



\usepackage{tgtermes}
\usepackage{tgheros}
\renewcommand{\ttdefault}{txtt}



\usepackage[Bjarne]{fncychap}
\usepackage[,numfigreset=1,mathnumfig,maxlistdepth=6]{sphinx}

\fvset{fontsize=auto}
\usepackage{geometry}


% Include hyperref last.
\usepackage{hyperref}
% Fix anchor placement for figures with captions.
\usepackage{hypcap}% it must be loaded after hyperref.
% Set up styles of URL: it should be placed after hyperref.
\urlstyle{same}

\addto\captionsenglish{\renewcommand{\contentsname}{Table of Contents}}

\usepackage{sphinxmessages}




\title{PSyclone User Guide}
\date{Mar 26, 2025}
\release{3.1.1\sphinxhyphen{}dev}
\author{Oakley Brunt, Andrew Coughtrie, Joshua Dendy, \\ Rupert Ford, Joerg Henrichs, Iva Kavcic, Andrew Porter, \\ Sergi Siso and Joseph Wallwork}
\newcommand{\sphinxlogo}{\vbox{}}
\renewcommand{\releasename}{Release}
\makeindex
\begin{document}

\ifdefined\shorthandoff
  \ifnum\catcode`\=\string=\active\shorthandoff{=}\fi
  \ifnum\catcode`\"=\active\shorthandoff{"}\fi
\fi

\pagestyle{empty}
\sphinxmaketitle
\pagestyle{plain}
\sphinxtableofcontents
\pagestyle{normal}
\phantomsection\label{\detokenize{index::doc}}


\sphinxAtStartPar
PSyclone is a source\sphinxhyphen{}to\sphinxhyphen{}source Fortran compiler designed to programmatically
optimise, parallelise and instrument HPC applications via user\sphinxhyphen{}provided
transformation scripts.

\sphinxAtStartPar
By encapsulating the performance\sphinxhyphen{}portability aspects (e.g. whether to
parallelise with OpenMP or OpenACC), these scripts enable a separation of
concerns between the scientific implementation and the optimisation choices.
This allows each aspect to be explored and developed largely independently.
Additionally, PSyclone supports the development of kernel\sphinxhyphen{}based, Fortran\sphinxhyphen{}embedded
DSLs following the PSyKAl model developed in the
\sphinxhref{https://www.metoffice.gov.uk/research/foundation/dynamics/next-generation}{GungHo project}.

\sphinxAtStartPar
PSyclone is currently used to support the
\sphinxhref{https://www.metoffice.gov.uk/research/modelling-systems/lfric/}{LFRic}
mixed finite\sphinxhyphen{}element PSyKAl DSL for the UK MetOffice’s next generation
modelling system and the
\sphinxhref{https://gtr.ukri.org/projects?ref=NE\%2FL01209X\%2F1}{GOcean}
finite\sphinxhyphen{}difference PSyKAl DSL for a prototype 2D ocean modelling system.
It is also used to insert GPU offloading directives into existing
directly\sphinxhyphen{}addressed MPI applications such as the
\sphinxhref{https://www.nemo-ocean.eu/}{NEMO ocean model}.

\sphinxAtStartPar
More detailed implementation information is available in the
\sphinxhref{https://psyclone-dev.readthedocs.io/}{Developer Guide}
and the
\sphinxhref{https://psyclone-ref.readthedocs.io/}{Reference Guide}.

\sphinxstepscope


\chapter{Getting Going}
\label{\detokenize{getting_going:getting-going}}\label{\detokenize{getting_going:id1}}\label{\detokenize{getting_going::doc}}

\section{Installation}
\label{\detokenize{getting_going:installation}}\label{\detokenize{getting_going:getting-going-download}}
\sphinxAtStartPar
The following instructions are intended for a PSyclone user who wants
to work with a released version of the code. If you are a developer or
wish to test a specific branch of PSyclone from the GitHub repository
please see the
\sphinxhref{https://psyclone-dev.readthedocs.io/en/latest/working\_practises.html\#dev-installation}{\DUrole{xref,std,std-ref}{Installation section in the Developer Guide}}.

\begin{sphinxuseclass}{sphinx-tabs}
\sphinxAtStartPar
From PyPI:

\sphinxAtStartPar
For a system\sphinxhyphen{}wide installation use:

\begin{sphinxVerbatim}[commandchars=\\\{\}]
pip\PYG{+w}{ }install\PYG{+w}{ }psyclone
\end{sphinxVerbatim}

\sphinxAtStartPar
For a user\sphinxhyphen{}local installation use:

\begin{sphinxVerbatim}[commandchars=\\\{\}]
pip\PYG{+w}{ }install\PYG{+w}{ }\PYGZhy{}\PYGZhy{}user\PYG{+w}{ }psyclone
\end{sphinxVerbatim}

\sphinxAtStartPar
For a specific release (where \sphinxcode{\sphinxupquote{X.Y.Z}} is the release version) use:

\begin{sphinxVerbatim}[commandchars=\\\{\}]
pip\PYG{+w}{ }install\PYG{+w}{ }\PYG{n+nv}{psyclone}\PYG{o}{=}\PYG{o}{=}X.Y.Z
\end{sphinxVerbatim}

\sphinxAtStartPar
For more information about using \sphinxcode{\sphinxupquote{pip}} or encapsulating the installation
in its own \sphinxcode{\sphinxupquote{virtual environment}} we recommend reading the
\sphinxhref{https://packaging.python.org/en/latest/tutorials/installing-packages/}{Python Packaging User Guide}.

\sphinxAtStartPar
From Spack:

\sphinxAtStartPar
To install psyclone to your loaded Spack installation use:

\begin{sphinxVerbatim}[commandchars=\\\{\}]
spack\PYG{+w}{ }install\PYG{+w}{ }psyclone
\end{sphinxVerbatim}

\sphinxAtStartPar
For more information about how to use Spack we recommend reading the
\sphinxhref{https://spack-tutorial.readthedocs.io/}{Spack documentation}.

\sphinxAtStartPar
From Source:

\sphinxAtStartPar
To download and install a specific PSyclone release (where \sphinxcode{\sphinxupquote{X.Y.Z}} is the release version)
from source, use:

\begin{sphinxVerbatim}[commandchars=\\\{\}]
wget\PYG{+w}{ }https://github.com/stfc/PSyclone/archive/X.Y.Z.tar.gz
tar\PYG{+w}{ }zxf\PYG{+w}{ }X.Y.Z.tar.gz
\PYG{n+nb}{cd}\PYG{+w}{ }PSyclone\PYGZhy{}X.Y.Z
pip\PYG{+w}{ }install\PYG{+w}{ }.
\end{sphinxVerbatim}

\end{sphinxuseclass}

\subsection{Installation location}
\label{\detokenize{getting_going:installation-location}}\label{\detokenize{getting_going:getting-going-install-loc}}
\sphinxAtStartPar
The PSyclone installation location will vary depending on the specific installation
method and options used. The \sphinxcode{\sphinxupquote{psyclone}} command will typically already be
prepended to your \sphinxcode{\sphinxupquote{PATH}} after following the instructions above, but sometimes
you will need to source the virtual environment or load the Spack module again
after restarting your terminal.

\sphinxAtStartPar
Once psyclone is in your \sphinxtitleref{PATH} you can execute \sphinxcode{\sphinxupquote{which psyclone}} to see
the installation directory. Some supporting files such as configuration,
examples and instrumentation libraries are installed under the \sphinxcode{\sphinxupquote{share/psyclone}}
directory relative to the psyclone installation. You can replace
\sphinxcode{\sphinxupquote{bin/psyclone}} in the string returned by \sphinxcode{\sphinxupquote{which psyclone}} with
\sphinxcode{\sphinxupquote{share/psyclone}} to find their location.


\section{Configuration}
\label{\detokenize{getting_going:configuration}}\label{\detokenize{getting_going:getting-going-configuration}}
\sphinxAtStartPar
Various aspects of PSyclone are controlled through a configuration
file, \sphinxcode{\sphinxupquote{psyclone.cfg}}.  The default version of this file is located in
the \sphinxcode{\sphinxupquote{share/psyclone/psyclone.cfg}} file in the {\hyperref[\detokenize{getting_going:getting-going-install-loc}]{\sphinxcrossref{\DUrole{std,std-ref}{Installation location}}}}.

\begin{sphinxadmonition}{warning}{Warning:}
\sphinxAtStartPar
If PSyclone is installed in ‘editable’ mode (\sphinxcode{\sphinxupquote{\sphinxhyphen{}e}} flag to \sphinxcode{\sphinxupquote{pip}}),
or in a non\sphinxhyphen{}standard location, then PSyclone will not be able to find the
default configuration file. There are two solutions to this:
\begin{enumerate}
\sphinxsetlistlabels{\arabic}{enumi}{enumii}{}{.}%
\item {} 
\sphinxAtStartPar
copy a configuration file to the location specified above.

\end{enumerate}

\sphinxAtStartPar
2. set the \sphinxcode{\sphinxupquote{PSYCLONE\_CONFIG}} environment variable (or the \sphinxcode{\sphinxupquote{\sphinxhyphen{}\sphinxhyphen{}config}}
flag) to the full\sphinxhyphen{}path to the configuration file, e.g.:

\begin{sphinxVerbatim}[commandchars=\\\{\}]
\PYG{n+nb}{export}\PYG{+w}{ }\PYG{n+nv}{PSYCLONE\PYGZus{}CONFIG}\PYG{o}{=}/some/path/PSyclone/config/psyclone.cfg
\end{sphinxVerbatim}
\end{sphinxadmonition}

\sphinxAtStartPar
See {\hyperref[\detokenize{configuration:configuration}]{\sphinxcrossref{\DUrole{std,std-ref}{Configuration}}}} for more details about the settings contained
within the config file.


\section{Running PSyclone}
\label{\detokenize{getting_going:running-psyclone}}\label{\detokenize{getting_going:getting-going-run}}
\sphinxAtStartPar
You are now ready to run PSyclone. One way of doing this is to use the \sphinxcode{\sphinxupquote{psyclone}}
command. To list the available options run: \sphinxcode{\sphinxupquote{psyclone \sphinxhyphen{}h}}, it should output:

\begin{sphinxVerbatim}[commandchars=\\\{\}]
\PYG{n}{usage}\PYG{p}{:} \PYG{n}{psyclone} \PYG{p}{[}\PYG{o}{\PYGZhy{}}\PYG{n}{h}\PYG{p}{]} \PYG{p}{[}\PYG{o}{\PYGZhy{}}\PYG{o}{\PYGZhy{}}\PYG{n}{version}\PYG{p}{]} \PYG{p}{[}\PYG{o}{\PYGZhy{}}\PYG{o}{\PYGZhy{}}\PYG{n}{config} \PYG{n}{CONFIG}\PYG{p}{]} \PYG{p}{[}\PYG{o}{\PYGZhy{}}\PYG{n}{s} \PYG{n}{SCRIPT}\PYG{p}{]} \PYG{p}{[}\PYG{o}{\PYGZhy{}}\PYG{n}{I} \PYG{n}{INCLUDE}\PYG{p}{]}
                \PYG{p}{[}\PYG{o}{\PYGZhy{}}\PYG{n}{l} \PYG{p}{\PYGZob{}}\PYG{n}{off}\PYG{p}{,}\PYG{n+nb}{all}\PYG{p}{,}\PYG{n}{output}\PYG{p}{\PYGZcb{}}\PYG{p}{]} \PYG{p}{[}\PYG{o}{\PYGZhy{}}\PYG{o}{\PYGZhy{}}\PYG{n}{profile} \PYG{p}{\PYGZob{}}\PYG{n}{invokes}\PYG{p}{,}\PYG{n}{routines}\PYG{p}{,}\PYG{n}{kernels}\PYG{p}{\PYGZcb{}}\PYG{p}{]}
                                \PYG{p}{[}\PYG{o}{\PYGZhy{}}\PYG{o}{\PYGZhy{}}\PYG{n}{backend} \PYG{p}{\PYGZob{}}\PYG{n}{enable}\PYG{o}{\PYGZhy{}}\PYG{n}{validation}\PYG{p}{,}\PYG{n}{disable}\PYG{o}{\PYGZhy{}}\PYG{n}{validation}\PYG{p}{\PYGZcb{}}\PYG{p}{]} \PYG{p}{[}\PYG{o}{\PYGZhy{}}\PYG{n}{o} \PYG{n}{OUTPUT\PYGZus{}FILE}\PYG{p}{]}
                                \PYG{p}{[}\PYG{o}{\PYGZhy{}}\PYG{n}{api} \PYG{n}{DSL}\PYG{p}{]} \PYG{p}{[}\PYG{o}{\PYGZhy{}}\PYG{n}{oalg} \PYG{n}{OUTPUT\PYGZus{}ALGORITHM\PYGZus{}FILE}\PYG{p}{]} \PYG{p}{[}\PYG{o}{\PYGZhy{}}\PYG{n}{opsy} \PYG{n}{OUTPUT\PYGZus{}PSY\PYGZus{}FILE}\PYG{p}{]}
                \PYG{p}{[}\PYG{o}{\PYGZhy{}}\PYG{n}{okern} \PYG{n}{OUTPUT\PYGZus{}KERNEL\PYGZus{}PATH}\PYG{p}{]} \PYG{p}{[}\PYG{o}{\PYGZhy{}}\PYG{n}{d} \PYG{n}{DIRECTORY}\PYG{p}{]} \PYG{p}{[}\PYG{o}{\PYGZhy{}}\PYG{n}{dm}\PYG{p}{]} \PYG{p}{[}\PYG{o}{\PYGZhy{}}\PYG{n}{nodm}\PYG{p}{]}
                \PYG{p}{[}\PYG{o}{\PYGZhy{}}\PYG{o}{\PYGZhy{}}\PYG{n}{kernel}\PYG{o}{\PYGZhy{}}\PYG{n}{renaming} \PYG{p}{\PYGZob{}}\PYG{n}{multiple}\PYG{p}{,}\PYG{n}{single}\PYG{p}{\PYGZcb{}}\PYG{p}{]}
                \PYG{n}{filename}

\PYG{n}{Transform} \PYG{n}{a} \PYG{n}{file} \PYG{n}{using} \PYG{n}{the} \PYG{n}{PSyclone} \PYG{n}{source}\PYG{o}{\PYGZhy{}}\PYG{n}{to}\PYG{o}{\PYGZhy{}}\PYG{n}{source} \PYG{n}{Fortran} \PYG{n}{compiler}

\PYG{n}{positional} \PYG{n}{arguments}\PYG{p}{:}
  \PYG{n}{filename}              \PYG{n+nb}{input} \PYG{n}{source} \PYG{n}{code}

\PYG{n}{options}\PYG{p}{:}
  \PYG{o}{\PYGZhy{}}\PYG{n}{h}\PYG{p}{,} \PYG{o}{\PYGZhy{}}\PYG{o}{\PYGZhy{}}\PYG{n}{help}            \PYG{n}{show} \PYG{n}{this} \PYG{n}{help} \PYG{n}{message} \PYG{o+ow}{and} \PYG{n}{exit}
  \PYG{o}{\PYGZhy{}}\PYG{o}{\PYGZhy{}}\PYG{n}{version}\PYG{p}{,} \PYG{o}{\PYGZhy{}}\PYG{n}{v}         \PYG{n}{display} \PYG{n}{version} \PYG{n}{information}
  \PYG{o}{\PYGZhy{}}\PYG{o}{\PYGZhy{}}\PYG{n}{config} \PYG{n}{CONFIG}\PYG{p}{,} \PYG{o}{\PYGZhy{}}\PYG{n}{c} \PYG{n}{CONFIG}
                        \PYG{n}{config} \PYG{n}{file} \PYG{k}{with} \PYG{n}{PSyclone} \PYG{n}{specific} \PYG{n}{options}
  \PYG{o}{\PYGZhy{}}\PYG{n}{s} \PYG{n}{SCRIPT}\PYG{p}{,} \PYG{o}{\PYGZhy{}}\PYG{o}{\PYGZhy{}}\PYG{n}{script} \PYG{n}{SCRIPT}
                        \PYG{n}{filename} \PYG{n}{of} \PYG{n}{a} \PYG{n}{PSyclone} \PYG{n}{optimisation} \PYG{n}{recipe}
  \PYG{o}{\PYGZhy{}}\PYG{n}{I} \PYG{n}{INCLUDE}\PYG{p}{,} \PYG{o}{\PYGZhy{}}\PYG{o}{\PYGZhy{}}\PYG{n}{include} \PYG{n}{INCLUDE}
                        \PYG{n}{path} \PYG{n}{to} \PYG{n}{Fortran} \PYG{n}{INCLUDE} \PYG{o+ow}{or} \PYG{n}{module} \PYG{n}{files}
  \PYG{o}{\PYGZhy{}}\PYG{n}{l} \PYG{p}{\PYGZob{}}\PYG{n}{off}\PYG{p}{,}\PYG{n+nb}{all}\PYG{p}{,}\PYG{n}{output}\PYG{p}{\PYGZcb{}}\PYG{p}{,} \PYG{o}{\PYGZhy{}}\PYG{o}{\PYGZhy{}}\PYG{n}{limit} \PYG{p}{\PYGZob{}}\PYG{n}{off}\PYG{p}{,}\PYG{n+nb}{all}\PYG{p}{,}\PYG{n}{output}\PYG{p}{\PYGZcb{}}
                        \PYG{n}{limit} \PYG{n}{the} \PYG{n}{Fortran} \PYG{n}{line} \PYG{n}{length} \PYG{n}{to} \PYG{l+m+mi}{132} \PYG{n}{characters} \PYG{p}{(}\PYG{n}{default} \PYG{l+s+s1}{\PYGZsq{}}\PYG{l+s+s1}{off}\PYG{l+s+s1}{\PYGZsq{}}\PYG{p}{)}\PYG{o}{.}
                        \PYG{n}{Use} \PYG{l+s+s1}{\PYGZsq{}}\PYG{l+s+s1}{all}\PYG{l+s+s1}{\PYGZsq{}} \PYG{n}{to} \PYG{n}{apply} \PYG{n}{limit} \PYG{n}{to} \PYG{n}{both} \PYG{n+nb}{input} \PYG{o+ow}{and} \PYG{n}{output} \PYG{n}{Fortran}\PYG{o}{.} \PYG{n}{Use}
                        \PYG{l+s+s1}{\PYGZsq{}}\PYG{l+s+s1}{output}\PYG{l+s+s1}{\PYGZsq{}} \PYG{n}{to} \PYG{n}{apply} \PYG{n}{line}\PYG{o}{\PYGZhy{}}\PYG{n}{length} \PYG{n}{limit} \PYG{n}{to} \PYG{n}{output} \PYG{n}{Fortran} \PYG{n}{only}\PYG{o}{.}
  \PYG{o}{\PYGZhy{}}\PYG{o}{\PYGZhy{}}\PYG{n}{profile} \PYG{p}{\PYGZob{}}\PYG{n}{invokes}\PYG{p}{,}\PYG{n}{routines}\PYG{p}{,}\PYG{n}{kernels}\PYG{p}{\PYGZcb{}}\PYG{p}{,} \PYG{o}{\PYGZhy{}}\PYG{n}{p} \PYG{p}{\PYGZob{}}\PYG{n}{invokes}\PYG{p}{,}\PYG{n}{routines}\PYG{p}{,}\PYG{n}{kernels}\PYG{p}{\PYGZcb{}}
                        \PYG{n}{add} \PYG{n}{profiling} \PYG{n}{hooks} \PYG{k}{for} \PYG{l+s+s1}{\PYGZsq{}}\PYG{l+s+s1}{kernels}\PYG{l+s+s1}{\PYGZsq{}}\PYG{p}{,} \PYG{l+s+s1}{\PYGZsq{}}\PYG{l+s+s1}{invokes}\PYG{l+s+s1}{\PYGZsq{}} \PYG{o+ow}{or} \PYG{l+s+s1}{\PYGZsq{}}\PYG{l+s+s1}{routines}\PYG{l+s+s1}{\PYGZsq{}}
  \PYG{o}{\PYGZhy{}}\PYG{o}{\PYGZhy{}}\PYG{n}{backend} \PYG{p}{\PYGZob{}}\PYG{n}{enable}\PYG{o}{\PYGZhy{}}\PYG{n}{validation}\PYG{p}{,}\PYG{n}{disable}\PYG{o}{\PYGZhy{}}\PYG{n}{validation}\PYG{p}{\PYGZcb{}}
                        \PYG{n}{options} \PYG{n}{to} \PYG{n}{control} \PYG{n}{the} \PYG{n}{PSyIR} \PYG{n}{backend} \PYG{n}{used} \PYG{k}{for} \PYG{n}{code} \PYG{n}{generation}\PYG{o}{.}
                        \PYG{n}{Use} \PYG{l+s+s1}{\PYGZsq{}}\PYG{l+s+s1}{disable\PYGZhy{}validation}\PYG{l+s+s1}{\PYGZsq{}} \PYG{n}{to} \PYG{n}{disable} \PYG{n}{the} \PYG{n}{validation} \PYG{n}{checks} \PYG{n}{that}
                        \PYG{n}{are} \PYG{n}{performed} \PYG{n}{by} \PYG{n}{default}\PYG{o}{.}
  \PYG{o}{\PYGZhy{}}\PYG{n}{o} \PYG{n}{OUTPUT\PYGZus{}FILE}        \PYG{p}{(}\PYG{n}{code}\PYG{o}{\PYGZhy{}}\PYG{n}{transformation} \PYG{n}{mode}\PYG{p}{)} \PYG{n}{output} \PYG{n}{file}
  \PYG{o}{\PYGZhy{}}\PYG{n}{api} \PYG{n}{DSL}\PYG{p}{,} \PYG{o}{\PYGZhy{}}\PYG{o}{\PYGZhy{}}\PYG{n}{psykal}\PYG{o}{\PYGZhy{}}\PYG{n}{dsl} \PYG{n}{DSL}
                        \PYG{n}{whether} \PYG{n}{to} \PYG{n}{use} \PYG{n}{a} \PYG{n}{PSyKAl} \PYG{n}{DSL} \PYG{p}{(}\PYG{n}{one} \PYG{n}{of} \PYG{p}{[}\PYG{l+s+s1}{\PYGZsq{}}\PYG{l+s+s1}{lfric}\PYG{l+s+s1}{\PYGZsq{}}\PYG{p}{,} \PYG{l+s+s1}{\PYGZsq{}}\PYG{l+s+s1}{gocean}\PYG{l+s+s1}{\PYGZsq{}}\PYG{p}{]}\PYG{p}{)}
  \PYG{o}{\PYGZhy{}}\PYG{n}{oalg} \PYG{n}{OUTPUT\PYGZus{}ALGORITHM\PYGZus{}FILE}
                        \PYG{p}{(}\PYG{n}{psykal} \PYG{n}{mode}\PYG{p}{)} \PYG{n}{filename} \PYG{n}{of} \PYG{n}{transformed} \PYG{n}{algorithm} \PYG{n}{code}
  \PYG{o}{\PYGZhy{}}\PYG{n}{opsy} \PYG{n}{OUTPUT\PYGZus{}PSY\PYGZus{}FILE}
                        \PYG{p}{(}\PYG{n}{psykal} \PYG{n}{mode}\PYG{p}{)} \PYG{n}{filename} \PYG{n}{of} \PYG{n}{generated} \PYG{n}{PSy}\PYG{o}{\PYGZhy{}}\PYG{n}{layer} \PYG{n}{code}
  \PYG{o}{\PYGZhy{}}\PYG{n}{okern} \PYG{n}{OUTPUT\PYGZus{}KERNEL\PYGZus{}PATH}
                        \PYG{p}{(}\PYG{n}{psykal} \PYG{n}{mode}\PYG{p}{)} \PYG{n}{directory} \PYG{o+ow}{in} \PYG{n}{which} \PYG{n}{to} \PYG{n}{put} \PYG{n}{transformed} \PYG{n}{kernels}\PYG{p}{,} \PYG{n}{default}
                        \PYG{o+ow}{is} \PYG{n}{the} \PYG{n}{current} \PYG{n}{working} \PYG{n}{directory}
  \PYG{o}{\PYGZhy{}}\PYG{n}{d} \PYG{n}{DIRECTORY}\PYG{p}{,} \PYG{o}{\PYGZhy{}}\PYG{o}{\PYGZhy{}}\PYG{n}{directory} \PYG{n}{DIRECTORY}
                        \PYG{p}{(}\PYG{n}{psykal} \PYG{n}{mode}\PYG{p}{)} \PYG{n}{path} \PYG{n}{to} \PYG{n}{a} \PYG{n}{root} \PYG{n}{directory} \PYG{n}{structure} \PYG{n}{containing} \PYG{n}{kernel}
                        \PYG{n}{source} \PYG{n}{code}\PYG{o}{.} \PYG{n}{Multiple} \PYG{n}{roots} \PYG{n}{can} \PYG{n}{be} \PYG{n}{specified} \PYG{n}{by} \PYG{n}{using} \PYG{n}{multiple} \PYG{o}{\PYGZhy{}}\PYG{n}{d}
                        \PYG{n}{arguments}\PYG{o}{.}
  \PYG{o}{\PYGZhy{}}\PYG{n}{dm}\PYG{p}{,} \PYG{o}{\PYGZhy{}}\PYG{o}{\PYGZhy{}}\PYG{n}{dist\PYGZus{}mem}       \PYG{p}{(}\PYG{n}{psykal} \PYG{n}{mode}\PYG{p}{)} \PYG{n}{generate} \PYG{n}{distributed} \PYG{n}{memory} \PYG{n}{code}
  \PYG{o}{\PYGZhy{}}\PYG{n}{nodm}\PYG{p}{,} \PYG{o}{\PYGZhy{}}\PYG{o}{\PYGZhy{}}\PYG{n}{no\PYGZus{}dist\PYGZus{}mem}  \PYG{p}{(}\PYG{n}{psykal} \PYG{n}{mode}\PYG{p}{)} \PYG{n}{do} \PYG{o+ow}{not} \PYG{n}{generate} \PYG{n}{distributed} \PYG{n}{memory} \PYG{n}{code}
  \PYG{o}{\PYGZhy{}}\PYG{o}{\PYGZhy{}}\PYG{n}{kernel}\PYG{o}{\PYGZhy{}}\PYG{n}{renaming} \PYG{p}{\PYGZob{}}\PYG{n}{multiple}\PYG{p}{,}\PYG{n}{single}\PYG{p}{\PYGZcb{}}
                        \PYG{p}{(}\PYG{n}{psykal} \PYG{n}{mode}\PYG{p}{)} \PYG{n}{naming} \PYG{n}{scheme} \PYG{n}{to} \PYG{n}{use} \PYG{n}{when} \PYG{n}{re}\PYG{o}{\PYGZhy{}}\PYG{n}{naming} \PYG{n}{transformed} \PYG{n}{kernels}
\end{sphinxVerbatim}

\sphinxAtStartPar
There is more detailed information about each flag in {\hyperref[\detokenize{psyclone_command:psyclone-command}]{\sphinxcrossref{\DUrole{std,std-ref}{The psyclone command}}}} section,
but the main parameters are the input source file that we aim to transform, and a transformation
recipe that is provided with the \sphinxcode{\sphinxupquote{\sphinxhyphen{}s}} flag.
In addition to these, note that psyclone can be used in two distinct modes:
the code\sphinxhyphen{}transformation mode (when no \sphinxcode{\sphinxupquote{\sphinxhyphen{}api}}/\sphinxcode{\sphinxupquote{\sphinxhyphen{}\sphinxhyphen{}psykal\sphinxhyphen{}dsl}} flags are provided) or the
PSyKAl DSL mode (when a \sphinxcode{\sphinxupquote{\sphinxhyphen{}api}}/\sphinxcode{\sphinxupquote{\sphinxhyphen{}\sphinxhyphen{}psykal\sphinxhyphen{}dsl}} flag is provided). The following sections provide
a brief introduction to each mode.


\subsection{PSyclone for Code Transformation}
\label{\detokenize{getting_going:psyclone-for-code-transformation}}
\sphinxAtStartPar
When using PSyclone for transforming existing Fortran files, only an
input source file is required:

\begin{sphinxVerbatim}[commandchars=\\\{\}]
\PYG{g+go}{psyclone input\PYGZus{}file.f90}
\end{sphinxVerbatim}

\sphinxAtStartPar
However, we usually want to redirect the output to a file so that we can later
compile it. We can do this using the \sphinxtitleref{\sphinxhyphen{}o} flag:

\begin{sphinxVerbatim}[commandchars=\\\{\}]
\PYG{g+go}{psyclone input\PYGZus{}file.f90 \PYGZhy{}o output.f90}
\end{sphinxVerbatim}

\sphinxAtStartPar
This should not transform the semantics of the code (only the syntax), and is
what we sometimes refer to as a “passthrough” run. This can be useful as an initial
correctness test when applying PSyclone to a new code.

\sphinxAtStartPar
However, PSyclone allows users to programatically change the source code of the
processed file. This is achieved using transformation recipes which are python scripts
with a \sphinxtitleref{trans} function defined. For example:

\begin{sphinxVerbatim}[commandchars=\\\{\}]
\PYG{k}{def}\PYG{+w}{ }\PYG{n+nf}{trans}\PYG{p}{(}\PYG{n}{psyir}\PYG{p}{)}\PYG{p}{:}
\PYG{+w}{    }\PYG{l+s+sd}{\PYGZsq{}\PYGZsq{}\PYGZsq{} Add OpenMP Parallel Loop directives.}

\PYG{l+s+sd}{    :param psyir: the PSyIR of the provided file.}
\PYG{l+s+sd}{    :type psyir: :py:class:`psyclone.psyir.nodes.FileContainer`}

\PYG{l+s+sd}{    \PYGZsq{}\PYGZsq{}\PYGZsq{}}
    \PYG{n}{omp\PYGZus{}trans} \PYG{o}{=} \PYG{n}{TransInfo}\PYG{p}{(}\PYG{p}{)}\PYG{o}{.}\PYG{n}{get\PYGZus{}trans\PYGZus{}name}\PYG{p}{(}\PYG{l+s+s1}{\PYGZsq{}}\PYG{l+s+s1}{OMPParallelLoopTrans}\PYG{l+s+s1}{\PYGZsq{}}\PYG{p}{)}

    \PYG{k}{for} \PYG{n}{loop} \PYG{o+ow}{in} \PYG{n}{psyir}\PYG{o}{.}\PYG{n}{walk}\PYG{p}{(}\PYG{n}{Loop}\PYG{p}{)}\PYG{p}{:}
        \PYG{k}{try}\PYG{p}{:}
            \PYG{n}{omp\PYGZus{}trans}\PYG{o}{.}\PYG{n}{apply}\PYG{p}{(}\PYG{n}{loop}\PYG{p}{)}
        \PYG{k}{except} \PYG{n}{TransformationError} \PYG{k}{as} \PYG{n}{err}\PYG{p}{:}
            \PYG{n+nb}{print}\PYG{p}{(}\PYG{l+s+sa}{f}\PYG{l+s+s2}{\PYGZdq{}}\PYG{l+s+s2}{Loop not paralellised because: }\PYG{l+s+si}{\PYGZob{}}\PYG{n}{err}\PYG{o}{.}\PYG{n}{value}\PYG{l+s+si}{\PYGZcb{}}\PYG{l+s+s2}{\PYGZdq{}}\PYG{p}{)}
\end{sphinxVerbatim}

\begin{sphinxadmonition}{warning}{Warning:}
\sphinxAtStartPar
Before PSyclone 3.0 the transformation scripts took a PSy object as argument:

\begin{sphinxVerbatim}[commandchars=\\\{\}]
\PYG{k}{def}\PYG{+w}{ }\PYG{n+nf}{trans}\PYG{p}{(}\PYG{n}{psy}\PYG{p}{)}\PYG{p}{:}
\PYG{+w}{    }\PYG{l+s+sd}{\PYGZsq{}\PYGZsq{}\PYGZsq{} Add OpenMP Parallel Loop directives.}

\PYG{l+s+sd}{    :param psy: the PSy object that PSyclone has constructed for the}
\PYG{l+s+sd}{                \PYGZsq{}invoke\PYGZsq{}(s) found in the Algorithm file.}
\PYG{l+s+sd}{    :type psy: :py:class:`psyclone.dynamo0p3.DynamoPSy`}

\PYG{l+s+sd}{    \PYGZsq{}\PYGZsq{}\PYGZsq{}}
    \PYG{k}{for} \PYG{n}{invoke} \PYG{o+ow}{in} \PYG{n}{psy}\PYG{o}{.}\PYG{n}{invokes}\PYG{o}{.}\PYG{n}{invoke\PYGZus{}list}\PYG{p}{:}
       \PYG{n}{invoke}\PYG{o}{.}\PYG{n}{schedule}
\end{sphinxVerbatim}

\sphinxAtStartPar
This is deprecated and will stop working in PSyclone releases post version 3.0
\end{sphinxadmonition}

\sphinxAtStartPar
And can be applied using the \sphinxtitleref{\sphinxhyphen{}s} flag:

\begin{sphinxVerbatim}[commandchars=\\\{\}]
\PYG{g+go}{psyclone input\PYGZus{}file.f90 \PYGZhy{}s trans\PYGZus{}script.py \PYGZhy{}o output.f90}
\end{sphinxVerbatim}

\sphinxAtStartPar
To see more complete examples of PSyclone for code transformation, see the
\sphinxcode{\sphinxupquote{examples/nemo}} folder in the PSyclone repository.


\subsection{PSyclone for PSyKAl DSLs}
\label{\detokenize{getting_going:psyclone-for-psykal-dsls}}
\sphinxAtStartPar
As indicated above, the \sphinxcode{\sphinxupquote{psyclone}} command can also be used to process PSyKAl
DSLs (\sphinxcode{\sphinxupquote{\sphinxhyphen{}\sphinxhyphen{}psykal\sphinxhyphen{}dsl}} flag). In this case the command takes as input the Fortran
source file containing the algorithm specification (in terms
of calls to \sphinxcode{\sphinxupquote{invoke()}}). It parses this, finds the necessary kernel
source files and produces two Fortran files. The first contains the
{\hyperref[\detokenize{introduction_to_psykal:psy-layer}]{\sphinxcrossref{\DUrole{std,std-ref}{middle, PSy\sphinxhyphen{}layer}}}} and the second a re\sphinxhyphen{}write of the
{\hyperref[\detokenize{introduction_to_psykal:algorithm-layer}]{\sphinxcrossref{\DUrole{std,std-ref}{algorithm code}}}} to use that layer. These files
are named according to the user\sphinxhyphen{}supplied arguments (options \sphinxcode{\sphinxupquote{\sphinxhyphen{}opsy}}
and \sphinxcode{\sphinxupquote{\sphinxhyphen{}oalg}} respectively). If those arguments are not supplied then the script
writes the re\sphinxhyphen{}written Fortran Algorithm layer to the terminal. For details of the other
command\sphinxhyphen{}line arguments please see the {\hyperref[\detokenize{psyclone_command:psyclone-command}]{\sphinxcrossref{\DUrole{std,std-ref}{The psyclone command}}}} Section.

\sphinxAtStartPar
Examples are provided in the \sphinxcode{\sphinxupquote{examples/lfric}} and \sphinxcode{\sphinxupquote{examples/gocean}} directories
of the PSyclone repository. Alternatively, if you have installed PSyclone using
\sphinxcode{\sphinxupquote{pip}} then they may be found in the \sphinxcode{\sphinxupquote{share/psyclone}} directory under your
PSyclone installation (see \sphinxtitleref{which psyclone} for the location of the
PSyclone installation).
In this case you should copy the whole \sphinxcode{\sphinxupquote{examples}} directory to some
convenient location before attempting to carry out the following instructions.

\sphinxAtStartPar
In this case we are going to use one of the LFRic examples:

\begin{sphinxVerbatim}[commandchars=\\\{\}]
\PYG{g+go}{cd \PYGZlt{}EGS\PYGZus{}HOME\PYGZgt{}/examples/lfric/eg1}
\PYG{g+go}{psyclone \PYGZhy{}\PYGZhy{}psykal\PYGZhy{}dsl lfric \PYGZhy{}d ../code \PYGZhy{}nodm \PYGZhy{}oalg alg.f90 \PYGZbs{}}
\PYG{g+go}{    \PYGZhy{}opsy psy.f90 ./single\PYGZus{}invoke.x90}
\end{sphinxVerbatim}

\sphinxAtStartPar
You should see two new files created, called \sphinxcode{\sphinxupquote{alg.f90}} (containing
the re\sphinxhyphen{}written algorithm layer) and \sphinxcode{\sphinxupquote{psy.f90}} (containing the
generated PSy\sphinxhyphen{} or middle\sphinxhyphen{}layer). Since this is an LFRic example the
Fortran source code has dependencies on the LFRic system and
therefore cannot be compiled stand\sphinxhyphen{}alone.

\sphinxAtStartPar
The PSy\sphinxhyphen{}layer that PSyclone creates is constructed using the PSyclone Internal
Representation ({\hyperref[\detokenize{psyir:psyir-ug}]{\sphinxcrossref{\DUrole{std,std-ref}{PSyIR}}}}). Accessing this is demonstrated
by the \sphinxcode{\sphinxupquote{print\_psyir\_trans.py}} script in the second LFRic example:

\begin{sphinxVerbatim}[commandchars=\\\{\}]
\PYG{g+go}{cd \PYGZlt{}EGS\PYGZus{}HOME\PYGZgt{}/examples/lfric/eg2}
\PYG{g+go}{psyclone \PYGZhy{}\PYGZhy{}psykal\PYGZhy{}dsl lfric \PYGZhy{}d ../code \PYGZhy{}s ./print\PYGZus{}psyir\PYGZus{}trans.py \PYGZbs{}}
\PYG{g+go}{    \PYGZhy{}opsy psy.f90 \PYGZhy{}oalg alg.f90 ./multi\PYGZus{}invoke\PYGZus{}mod.x90}
\end{sphinxVerbatim}

\sphinxAtStartPar
Take a look at the \sphinxcode{\sphinxupquote{print\_psyir\_trans.py}} script for more information. \sphinxstyleemphasis{Hint};
you can insert a single line in that script in order to break into the Python
interpreter during execution: \sphinxcode{\sphinxupquote{import pdb; pdb.set\_trace()}}. This then enables
interactive exploration of the PSyIR if you are interested.

\sphinxstepscope


\chapter{The psyclone command}
\label{\detokenize{psyclone_command:the-psyclone-command}}\label{\detokenize{psyclone_command:psyclone-command}}\label{\detokenize{psyclone_command::doc}}
\sphinxAtStartPar
The \sphinxcode{\sphinxupquote{psyclone}} command is an executable script designed to be run from the
command line, e.g.:

\begin{sphinxVerbatim}[commandchars=\\\{\}]
\PYG{o}{\PYGZgt{}} \PYG{n}{psyclone} \PYG{o}{\PYGZlt{}}\PYG{n}{args}\PYG{o}{\PYGZgt{}}
\end{sphinxVerbatim}

\sphinxAtStartPar
The optional \sphinxcode{\sphinxupquote{\sphinxhyphen{}h}} argument gives a description of the options provided
by the command:

\begin{sphinxVerbatim}[commandchars=\\\{\}]
\PYG{o}{\PYGZgt{}} \PYG{n}{psyclone} \PYG{o}{\PYGZhy{}}\PYG{n}{h}
 \PYG{n}{usage}\PYG{p}{:} \PYG{n}{psyclone} \PYG{p}{[}\PYG{o}{\PYGZhy{}}\PYG{n}{h}\PYG{p}{]} \PYG{p}{[}\PYG{o}{\PYGZhy{}}\PYG{o}{\PYGZhy{}}\PYG{n}{version}\PYG{p}{]} \PYG{p}{[}\PYG{o}{\PYGZhy{}}\PYG{o}{\PYGZhy{}}\PYG{n}{config} \PYG{n}{CONFIG}\PYG{p}{]} \PYG{p}{[}\PYG{o}{\PYGZhy{}}\PYG{n}{s} \PYG{n}{SCRIPT}\PYG{p}{]} \PYG{p}{[}\PYG{o}{\PYGZhy{}}\PYG{n}{I} \PYG{n}{INCLUDE}\PYG{p}{]}
                 \PYG{p}{[}\PYG{o}{\PYGZhy{}}\PYG{n}{l} \PYG{p}{\PYGZob{}}\PYG{n}{off}\PYG{p}{,}\PYG{n+nb}{all}\PYG{p}{,}\PYG{n}{output}\PYG{p}{\PYGZcb{}}\PYG{p}{]} \PYG{p}{[}\PYG{o}{\PYGZhy{}}\PYG{o}{\PYGZhy{}}\PYG{n}{profile} \PYG{p}{\PYGZob{}}\PYG{n}{invokes}\PYG{p}{,}\PYG{n}{routines}\PYG{p}{,}\PYG{n}{kernels}\PYG{p}{\PYGZcb{}}\PYG{p}{]}
                                 \PYG{p}{[}\PYG{o}{\PYGZhy{}}\PYG{o}{\PYGZhy{}}\PYG{n}{backend} \PYG{p}{\PYGZob{}}\PYG{n}{enable}\PYG{o}{\PYGZhy{}}\PYG{n}{validation}\PYG{p}{,}\PYG{n}{disable}\PYG{o}{\PYGZhy{}}\PYG{n}{validation}\PYG{p}{\PYGZcb{}}\PYG{p}{]} \PYG{p}{[}\PYG{o}{\PYGZhy{}}\PYG{n}{o} \PYG{n}{OUTPUT\PYGZus{}FILE}\PYG{p}{]}
                                 \PYG{p}{[}\PYG{o}{\PYGZhy{}}\PYG{n}{api} \PYG{n}{DSL}\PYG{p}{]} \PYG{p}{[}\PYG{o}{\PYGZhy{}}\PYG{n}{oalg} \PYG{n}{OUTPUT\PYGZus{}ALGORITHM\PYGZus{}FILE}\PYG{p}{]} \PYG{p}{[}\PYG{o}{\PYGZhy{}}\PYG{n}{opsy} \PYG{n}{OUTPUT\PYGZus{}PSY\PYGZus{}FILE}\PYG{p}{]}
                 \PYG{p}{[}\PYG{o}{\PYGZhy{}}\PYG{n}{okern} \PYG{n}{OUTPUT\PYGZus{}KERNEL\PYGZus{}PATH}\PYG{p}{]} \PYG{p}{[}\PYG{o}{\PYGZhy{}}\PYG{n}{d} \PYG{n}{DIRECTORY}\PYG{p}{]} \PYG{p}{[}\PYG{o}{\PYGZhy{}}\PYG{n}{dm}\PYG{p}{]} \PYG{p}{[}\PYG{o}{\PYGZhy{}}\PYG{n}{nodm}\PYG{p}{]}
                 \PYG{p}{[}\PYG{o}{\PYGZhy{}}\PYG{o}{\PYGZhy{}}\PYG{n}{kernel}\PYG{o}{\PYGZhy{}}\PYG{n}{renaming} \PYG{p}{\PYGZob{}}\PYG{n}{multiple}\PYG{p}{,}\PYG{n}{single}\PYG{p}{\PYGZcb{}}\PYG{p}{]}
                 \PYG{n}{filename}

 \PYG{n}{Transform} \PYG{n}{a} \PYG{n}{file} \PYG{n}{using} \PYG{n}{the} \PYG{n}{PSyclone} \PYG{n}{source}\PYG{o}{\PYGZhy{}}\PYG{n}{to}\PYG{o}{\PYGZhy{}}\PYG{n}{source} \PYG{n}{Fortran} \PYG{n}{compiler}

 \PYG{n}{positional} \PYG{n}{arguments}\PYG{p}{:}
   \PYG{n}{filename}              \PYG{n+nb}{input} \PYG{n}{source} \PYG{n}{code}

 \PYG{n}{options}\PYG{p}{:}
   \PYG{o}{\PYGZhy{}}\PYG{n}{h}\PYG{p}{,} \PYG{o}{\PYGZhy{}}\PYG{o}{\PYGZhy{}}\PYG{n}{help}            \PYG{n}{show} \PYG{n}{this} \PYG{n}{help} \PYG{n}{message} \PYG{o+ow}{and} \PYG{n}{exit}
   \PYG{o}{\PYGZhy{}}\PYG{o}{\PYGZhy{}}\PYG{n}{version}\PYG{p}{,} \PYG{o}{\PYGZhy{}}\PYG{n}{v}         \PYG{n}{display} \PYG{n}{version} \PYG{n}{information}
   \PYG{o}{\PYGZhy{}}\PYG{o}{\PYGZhy{}}\PYG{n}{config} \PYG{n}{CONFIG}\PYG{p}{,} \PYG{o}{\PYGZhy{}}\PYG{n}{c} \PYG{n}{CONFIG}
                         \PYG{n}{config} \PYG{n}{file} \PYG{k}{with} \PYG{n}{PSyclone} \PYG{n}{specific} \PYG{n}{options}
   \PYG{o}{\PYGZhy{}}\PYG{n}{s} \PYG{n}{SCRIPT}\PYG{p}{,} \PYG{o}{\PYGZhy{}}\PYG{o}{\PYGZhy{}}\PYG{n}{script} \PYG{n}{SCRIPT}
                         \PYG{n}{filename} \PYG{n}{of} \PYG{n}{a} \PYG{n}{PSyclone} \PYG{n}{optimisation} \PYG{n}{recipe}
   \PYG{o}{\PYGZhy{}}\PYG{n}{I} \PYG{n}{INCLUDE}\PYG{p}{,} \PYG{o}{\PYGZhy{}}\PYG{o}{\PYGZhy{}}\PYG{n}{include} \PYG{n}{INCLUDE}
                         \PYG{n}{path} \PYG{n}{to} \PYG{n}{Fortran} \PYG{n}{INCLUDE} \PYG{o+ow}{or} \PYG{n}{module} \PYG{n}{files}
   \PYG{o}{\PYGZhy{}}\PYG{n}{l} \PYG{p}{\PYGZob{}}\PYG{n}{off}\PYG{p}{,}\PYG{n+nb}{all}\PYG{p}{,}\PYG{n}{output}\PYG{p}{\PYGZcb{}}\PYG{p}{,} \PYG{o}{\PYGZhy{}}\PYG{o}{\PYGZhy{}}\PYG{n}{limit} \PYG{p}{\PYGZob{}}\PYG{n}{off}\PYG{p}{,}\PYG{n+nb}{all}\PYG{p}{,}\PYG{n}{output}\PYG{p}{\PYGZcb{}}
                         \PYG{n}{limit} \PYG{n}{the} \PYG{n}{Fortran} \PYG{n}{line} \PYG{n}{length} \PYG{n}{to} \PYG{l+m+mi}{132} \PYG{n}{characters} \PYG{p}{(}\PYG{n}{default} \PYG{l+s+s1}{\PYGZsq{}}\PYG{l+s+s1}{off}\PYG{l+s+s1}{\PYGZsq{}}\PYG{p}{)}\PYG{o}{.}
                         \PYG{n}{Use} \PYG{l+s+s1}{\PYGZsq{}}\PYG{l+s+s1}{all}\PYG{l+s+s1}{\PYGZsq{}} \PYG{n}{to} \PYG{n}{apply} \PYG{n}{limit} \PYG{n}{to} \PYG{n}{both} \PYG{n+nb}{input} \PYG{o+ow}{and} \PYG{n}{output} \PYG{n}{Fortran}\PYG{o}{.} \PYG{n}{Use}
                         \PYG{l+s+s1}{\PYGZsq{}}\PYG{l+s+s1}{output}\PYG{l+s+s1}{\PYGZsq{}} \PYG{n}{to} \PYG{n}{apply} \PYG{n}{line}\PYG{o}{\PYGZhy{}}\PYG{n}{length} \PYG{n}{limit} \PYG{n}{to} \PYG{n}{output} \PYG{n}{Fortran} \PYG{n}{only}\PYG{o}{.}
   \PYG{o}{\PYGZhy{}}\PYG{o}{\PYGZhy{}}\PYG{n}{profile} \PYG{p}{\PYGZob{}}\PYG{n}{invokes}\PYG{p}{,}\PYG{n}{routines}\PYG{p}{,}\PYG{n}{kernels}\PYG{p}{\PYGZcb{}}\PYG{p}{,} \PYG{o}{\PYGZhy{}}\PYG{n}{p} \PYG{p}{\PYGZob{}}\PYG{n}{invokes}\PYG{p}{,}\PYG{n}{routines}\PYG{p}{,}\PYG{n}{kernels}\PYG{p}{\PYGZcb{}}
                         \PYG{n}{add} \PYG{n}{profiling} \PYG{n}{hooks} \PYG{k}{for} \PYG{l+s+s1}{\PYGZsq{}}\PYG{l+s+s1}{kernels}\PYG{l+s+s1}{\PYGZsq{}}\PYG{p}{,} \PYG{l+s+s1}{\PYGZsq{}}\PYG{l+s+s1}{invokes}\PYG{l+s+s1}{\PYGZsq{}} \PYG{o+ow}{or} \PYG{l+s+s1}{\PYGZsq{}}\PYG{l+s+s1}{routines}\PYG{l+s+s1}{\PYGZsq{}}
   \PYG{o}{\PYGZhy{}}\PYG{o}{\PYGZhy{}}\PYG{n}{backend} \PYG{p}{\PYGZob{}}\PYG{n}{enable}\PYG{o}{\PYGZhy{}}\PYG{n}{validation}\PYG{p}{,}\PYG{n}{disable}\PYG{o}{\PYGZhy{}}\PYG{n}{validation}\PYG{p}{\PYGZcb{}}
                         \PYG{n}{options} \PYG{n}{to} \PYG{n}{control} \PYG{n}{the} \PYG{n}{PSyIR} \PYG{n}{backend} \PYG{n}{used} \PYG{k}{for} \PYG{n}{code} \PYG{n}{generation}\PYG{o}{.}
                         \PYG{n}{Use} \PYG{l+s+s1}{\PYGZsq{}}\PYG{l+s+s1}{disable\PYGZhy{}validation}\PYG{l+s+s1}{\PYGZsq{}} \PYG{n}{to} \PYG{n}{disable} \PYG{n}{the} \PYG{n}{validation} \PYG{n}{checks} \PYG{n}{that}
                         \PYG{n}{are} \PYG{n}{performed} \PYG{n}{by} \PYG{n}{default}\PYG{o}{.}
   \PYG{o}{\PYGZhy{}}\PYG{n}{o} \PYG{n}{OUTPUT\PYGZus{}FILE}        \PYG{p}{(}\PYG{n}{code}\PYG{o}{\PYGZhy{}}\PYG{n}{transformation} \PYG{n}{mode}\PYG{p}{)} \PYG{n}{output} \PYG{n}{file}
   \PYG{o}{\PYGZhy{}}\PYG{n}{api} \PYG{n}{DSL}\PYG{p}{,} \PYG{o}{\PYGZhy{}}\PYG{o}{\PYGZhy{}}\PYG{n}{psykal}\PYG{o}{\PYGZhy{}}\PYG{n}{dsl} \PYG{n}{DSL}
                         \PYG{n}{whether} \PYG{n}{to} \PYG{n}{use} \PYG{n}{a} \PYG{n}{PSyKAl} \PYG{n}{DSL} \PYG{p}{(}\PYG{n}{one} \PYG{n}{of} \PYG{p}{[}\PYG{l+s+s1}{\PYGZsq{}}\PYG{l+s+s1}{lfric}\PYG{l+s+s1}{\PYGZsq{}}\PYG{p}{,} \PYG{l+s+s1}{\PYGZsq{}}\PYG{l+s+s1}{gocean}\PYG{l+s+s1}{\PYGZsq{}}\PYG{p}{]}\PYG{p}{)}
   \PYG{o}{\PYGZhy{}}\PYG{n}{oalg} \PYG{n}{OUTPUT\PYGZus{}ALGORITHM\PYGZus{}FILE}
                         \PYG{p}{(}\PYG{n}{psykal} \PYG{n}{mode}\PYG{p}{)} \PYG{n}{filename} \PYG{n}{of} \PYG{n}{transformed} \PYG{n}{algorithm} \PYG{n}{code}
   \PYG{o}{\PYGZhy{}}\PYG{n}{opsy} \PYG{n}{OUTPUT\PYGZus{}PSY\PYGZus{}FILE}
                         \PYG{p}{(}\PYG{n}{psykal} \PYG{n}{mode}\PYG{p}{)} \PYG{n}{filename} \PYG{n}{of} \PYG{n}{generated} \PYG{n}{PSy}\PYG{o}{\PYGZhy{}}\PYG{n}{layer} \PYG{n}{code}
   \PYG{o}{\PYGZhy{}}\PYG{n}{okern} \PYG{n}{OUTPUT\PYGZus{}KERNEL\PYGZus{}PATH}
                         \PYG{p}{(}\PYG{n}{psykal} \PYG{n}{mode}\PYG{p}{)} \PYG{n}{directory} \PYG{o+ow}{in} \PYG{n}{which} \PYG{n}{to} \PYG{n}{put} \PYG{n}{transformed} \PYG{n}{kernels}\PYG{p}{,} \PYG{n}{default}
                         \PYG{o+ow}{is} \PYG{n}{the} \PYG{n}{current} \PYG{n}{working} \PYG{n}{directory}\PYG{o}{.}
   \PYG{o}{\PYGZhy{}}\PYG{n}{d} \PYG{n}{DIRECTORY}\PYG{p}{,} \PYG{o}{\PYGZhy{}}\PYG{o}{\PYGZhy{}}\PYG{n}{directory} \PYG{n}{DIRECTORY}
                         \PYG{p}{(}\PYG{n}{psykal} \PYG{n}{mode}\PYG{p}{)} \PYG{n}{path} \PYG{n}{to} \PYG{n}{a} \PYG{n}{root} \PYG{n}{directory} \PYG{n}{structure} \PYG{n}{containing} \PYG{n}{kernel}
                         \PYG{n}{source} \PYG{n}{code}\PYG{o}{.} \PYG{n}{Multiple} \PYG{n}{roots} \PYG{n}{can} \PYG{n}{be} \PYG{n}{specified} \PYG{n}{by} \PYG{n}{using} \PYG{n}{multiple} \PYG{o}{\PYGZhy{}}\PYG{n}{d}
                         \PYG{n}{arguments}\PYG{o}{.}
   \PYG{o}{\PYGZhy{}}\PYG{n}{dm}\PYG{p}{,} \PYG{o}{\PYGZhy{}}\PYG{o}{\PYGZhy{}}\PYG{n}{dist\PYGZus{}mem}       \PYG{p}{(}\PYG{n}{psykal} \PYG{n}{mode}\PYG{p}{)} \PYG{n}{generate} \PYG{n}{distributed} \PYG{n}{memory} \PYG{n}{code}
   \PYG{o}{\PYGZhy{}}\PYG{n}{nodm}\PYG{p}{,} \PYG{o}{\PYGZhy{}}\PYG{o}{\PYGZhy{}}\PYG{n}{no\PYGZus{}dist\PYGZus{}mem}  \PYG{p}{(}\PYG{n}{psykal} \PYG{n}{mode}\PYG{p}{)} \PYG{n}{do} \PYG{o+ow}{not} \PYG{n}{generate} \PYG{n}{distributed} \PYG{n}{memory} \PYG{n}{code}
   \PYG{o}{\PYGZhy{}}\PYG{o}{\PYGZhy{}}\PYG{n}{kernel}\PYG{o}{\PYGZhy{}}\PYG{n}{renaming} \PYG{p}{\PYGZob{}}\PYG{n}{multiple}\PYG{p}{,}\PYG{n}{single}\PYG{p}{\PYGZcb{}}
                         \PYG{p}{(}\PYG{n}{psykal} \PYG{n}{mode}\PYG{p}{)} \PYG{n}{naming} \PYG{n}{scheme} \PYG{n}{to} \PYG{n}{use} \PYG{n}{when} \PYG{n}{re}\PYG{o}{\PYGZhy{}}\PYG{n}{naming} \PYG{n}{transformed} \PYG{n}{kernels}
\end{sphinxVerbatim}


\section{Basic Use}
\label{\detokenize{psyclone_command:basic-use}}
\sphinxAtStartPar
The simplest way to use \sphinxcode{\sphinxupquote{psyclone}} is to provide a Fortran input source file:

\begin{sphinxVerbatim}[commandchars=\\\{\}]
\PYG{g+go}{psyclone input.f90}
\end{sphinxVerbatim}

\sphinxAtStartPar
If the input file is valid Fortran, PSyclone will print the output Fortran
(in this case the same unmodified code but with normalised syntax) to stdout.
Otherwise it will print the errors detected while parsing the Fortran file.

\sphinxAtStartPar
Usually we want to redirect the output to a file that we can later
compile. We can do this with the \sphinxtitleref{\sphinxhyphen{}o} flag:

\begin{sphinxVerbatim}[commandchars=\\\{\}]
\PYG{g+go}{psyclone input.f90 \PYGZhy{}o output.f90}
\end{sphinxVerbatim}


\section{Transformation script}
\label{\detokenize{psyclone_command:transformation-script}}
\sphinxAtStartPar
By default, the \sphinxcode{\sphinxupquote{psyclone}} command will not apply any transformation (other
than canonicalising the code and generating a normalised syntax). To apply
transformations to the code, a recipe needs to be specified with the \sphinxtitleref{\sphinxhyphen{}s} flag.
This option is discussed in more detail in the {\hyperref[\detokenize{transformations:sec-transformations-script}]{\sphinxcrossref{\DUrole{std,std-ref}{PSyclone User Scripts}}}}
section. With a transformation recipe the command looks like:

\begin{sphinxVerbatim}[commandchars=\\\{\}]
\PYG{g+go}{psyclone input.f90 \PYGZhy{}s transformation\PYGZus{}recipe.py}
\end{sphinxVerbatim}


\section{Fortran INCLUDE Files and Modules}
\label{\detokenize{psyclone_command:fortran-include-files-and-modules}}
\sphinxAtStartPar
If the source code to be processed by PSyclone
contains INCLUDE statements then the location of any INCLUDE’d files
\sphinxstyleemphasis{must} be supplied to PSyclone via the \sphinxcode{\sphinxupquote{\sphinxhyphen{}I}} or \sphinxcode{\sphinxupquote{\sphinxhyphen{}\sphinxhyphen{}include}}
option. (This is necessary because INCLUDE lines are a part of the
Fortran language and must therefore be parsed \sphinxhyphen{} they are not handled
by any pre\sphinxhyphen{}processing step.) Multiple locations may be specified by
using multiple \sphinxcode{\sphinxupquote{\sphinxhyphen{}I}} flags, e.g.:

\begin{sphinxVerbatim}[commandchars=\\\{\}]
\PYG{g+go}{psyclone \PYGZhy{}I /some/path \PYGZhy{}I /some/other/path input.f90}
\end{sphinxVerbatim}

\sphinxAtStartPar
If no include paths are specified then the directory containing the
source file currently being parsed is searched by default. If the
specified INCLUDE file is not found then PSyclone will abort with
an appropriate error. For example:

\begin{sphinxVerbatim}[commandchars=\\\{\}]
\PYG{g+go}{psyclone \PYGZhy{}I nonexisting test.f90}
\PYG{g+go}{PSyclone configuration error: Include path \PYGZsq{}nonexisting\PYGZsq{} does not exist}
\end{sphinxVerbatim}

\sphinxAtStartPar
Currently, the PSyKAl\sphinxhyphen{}based APIs (LFRic and GOcean \sphinxhyphen{} see below) will ignore
(but preserve) INCLUDE statements in algorithm\sphinxhyphen{}layer code. However, INCLUDE
statements in kernels will, in general, cause the kernel parsing to fail
unless the file(s) referenced in such statements are in the same directory
as the kernel file. Once kernel parsing has been re\sphinxhyphen{}implemented to use
fparser2 (issue \#239) and the PSyclone Intermediate Representation then the
behaviour will be the same as for generic code\sphinxhyphen{}transformations.

\sphinxAtStartPar
Since PSyclone does not attempt to be a full compiler, it does not require
that the code be available for any Fortran modules referred to by \sphinxcode{\sphinxupquote{use}}
statements. However, certain transformations \sphinxstyleemphasis{do} require that e.g. type
information be determined for all variables in the code being transformed.
In this case PSyclone \sphinxstyleemphasis{will} need to be able to find and process any
referenced modules. To do this it searches in the directories specified
by the \sphinxcode{\sphinxupquote{\sphinxhyphen{}I}}/\sphinxcode{\sphinxupquote{\sphinxhyphen{}\sphinxhyphen{}include}} flags.


\subsection{C Pre\sphinxhyphen{}processor \#include Files}
\label{\detokenize{psyclone_command:c-pre-processor-include-files}}
\sphinxAtStartPar
PSyclone currently only supports Fortran input. As such, if a file to
be processed contains CPP \sphinxcode{\sphinxupquote{\#include}} statements then it must first be
processed by a suitable pre\sphinxhyphen{}processor before being passed to PSyclone.
PSyclone will abort with an appropriate error if it encounters a
\sphinxcode{\sphinxupquote{\#include}} in any code being processed (whether or not a PSykAL DSL is
in use).


\section{Fortran line length}
\label{\detokenize{psyclone_command:fortran-line-length}}\label{\detokenize{psyclone_command:fort-line-length}}
\sphinxAtStartPar
By default the \sphinxcode{\sphinxupquote{psyclone}} command will generate Fortran code with no
consideration of Fortran line\sphinxhyphen{}length limits. As the line\sphinxhyphen{}length limit
for free\sphinxhyphen{}format Fortran is 132 characters, the code that is output may
be non\sphinxhyphen{}conformant.

\sphinxAtStartPar
Line length is not an issue for many compilers as they provide flags to
increase or disable Fortran standard line lengths limits. However this is
not the case for all compilers.

\sphinxAtStartPar
When either the \sphinxcode{\sphinxupquote{\sphinxhyphen{}l all}} or \sphinxcode{\sphinxupquote{\sphinxhyphen{}l output}} option is specified to
the \sphinxcode{\sphinxupquote{psyclone}} command, the output will be line wrapped so that the
output lines are always within the 132 character limit.

\sphinxAtStartPar
The \sphinxcode{\sphinxupquote{\sphinxhyphen{}l all}} additionally checks the input Fortran files for conformance
and raises an error if they do not conform.

\sphinxAtStartPar
Line wrapping is not performed by default. There are two reasons for
this. This first reason is that most compilers are able to cope with
long lines. The second reason is that the line wrapping implementation
could fail in certain pathological cases. The implementation and
limitations of line wrapping are discussed in the
{\hyperref[\detokenize{line_length:line-length-limitations}]{\sphinxcrossref{\DUrole{std,std-ref}{Limitations}}}} section.


\section{Backend Options}
\label{\detokenize{psyclone_command:backend-options}}\label{\detokenize{psyclone_command:id1}}
\sphinxAtStartPar
The final code generated by PSyclone is created by passing the PSyIR
tree to one of the ‘backends’ (see \sphinxhref{https://psyclone-dev.readthedocs.io/en/latest/psyir\_backends.html\#psyir-backends}{PSyIR Back\sphinxhyphen{}ends} in
the Developer Guide for more details). The \sphinxcode{\sphinxupquote{\sphinxhyphen{}\sphinxhyphen{}backend}} flag permits
a user to tune the behaviour of this code generation. Currently, the
only option is \sphinxcode{\sphinxupquote{\{en,dis\}able\sphinxhyphen{}validation}} which turns on/off the
validation checks performed when doing code generation. By default,
such validation is enabled as it is only at code\sphinxhyphen{}generation time that
certain constraints can be checked (since PSyclone does not mandate
the order in which code transformations are applied).  Occasionally,
these validation checks may raise false positives (due to incomplete
implementations), at which point it is useful to be able to disable
them.  The default behaviour may be changed by adding the
\sphinxcode{\sphinxupquote{BACKEND\_CHECKS\_ENABLED}} entry to the
{\hyperref[\detokenize{configuration:config-default-section}]{\sphinxcrossref{\DUrole{std,std-ref}{configuration file}}}}. Any
command\sphinxhyphen{}line setting always takes precendence though. It is
recommended that validation only be disabled as a last resort and for
as few input source files as possible.


\section{Automatic Profiling Instrumentation}
\label{\detokenize{psyclone_command:automatic-profiling-instrumentation}}
\sphinxAtStartPar
The \sphinxcode{\sphinxupquote{\sphinxhyphen{}\sphinxhyphen{}profile}} option allows the user to instruct PSyclone to automatically
insert profiling calls in addition to the code transformations specified in
the recipe.  This flag accepts the options: \sphinxcode{\sphinxupquote{routines}}, \sphinxcode{\sphinxupquote{invokes}} and
\sphinxcode{\sphinxupquote{kernels}}. PSyclone will insert profiling\sphinxhyphen{}start and \sphinxhyphen{}stop calls at the
beginning and end of each routine, PSy\sphinxhyphen{}layer invoke or PSy\sphinxhyphen{}layer kernel call,
respectively. The generated code must be linked against the PSyclone profiling
interface and the profiling tool itself. The application that calls the
PSyclone\sphinxhyphen{}generated code is responsible for initialising and finalising the
profiling library that is being used (if necessary). For more details on the use
of this profiling functionality please see the {\hyperref[\detokenize{profiling:profiling}]{\sphinxcrossref{\DUrole{std,std-ref}{Profiling}}}} section.


\section{Using PSyclone for PSyKAL DSLs}
\label{\detokenize{psyclone_command:using-psyclone-for-psykal-dsls}}
\sphinxAtStartPar
In addition to the default code\sphinxhyphen{}transformation mode, \sphinxcode{\sphinxupquote{psyclone}} can also
be used to process Fortran files that implement PSyKAL DSLs (see
{\hyperref[\detokenize{introduction_to_psykal:introduction-to-psykal}]{\sphinxcrossref{\DUrole{std,std-ref}{Introduction to PSyKAl}}}}). To do this you can choose a DSL API
with the \sphinxcode{\sphinxupquote{\sphinxhyphen{}api}} or \sphinxcode{\sphinxupquote{\sphinxhyphen{}\sphinxhyphen{}psykal\sphinxhyphen{}dsl}} flag.

\sphinxAtStartPar
The main difference is that, instead of providing a single file to process, for
PSyKAl DSLs PSyclone expects an algorithm\sphinxhyphen{}layer file that describes the high\sphinxhyphen{}level
view of an algorithm. PSyclone will use this algorithm file and the metadata of the
kernels that it calls to generate a PSy(Parallel System)\sphinxhyphen{}layer code that connects
the Algorithm layer to the Kernels. In this mode of operation, any supplied
transformation recipe is applied to the PSy\sphinxhyphen{}layer.

\sphinxAtStartPar
By default, the \sphinxcode{\sphinxupquote{psyclone}} command for PSyKAl APIs will generate distributed
memory (DM) code (unless otherwise specified in the {\hyperref[\detokenize{configuration:configuration}]{\sphinxcrossref{\DUrole{std,std-ref}{Configuration}}}} file).
Alternatively, whether or not to generate DM code can be specified as an
argument to the \sphinxcode{\sphinxupquote{psyclone}} command using the \sphinxcode{\sphinxupquote{\sphinxhyphen{}dm}}/\sphinxcode{\sphinxupquote{\sphinxhyphen{}\sphinxhyphen{}dist\_mem}} or
\sphinxcode{\sphinxupquote{\sphinxhyphen{}nodm}}/\sphinxcode{\sphinxupquote{\sphinxhyphen{}\sphinxhyphen{}no\_dist\_mem}} flags, respectively.
For exampe the following command will generate GOcean PSyKAl code with DM:

\begin{sphinxVerbatim}[commandchars=\\\{\}]
\PYG{g+go}{psyclone \PYGZhy{}api gocean \PYGZhy{}dm algorithm.f90}
\end{sphinxVerbatim}

\sphinxAtStartPar
See {\hyperref[\detokenize{introduction_to_psykal:psykal-usage}]{\sphinxcrossref{\DUrole{std,std-ref}{psyclone usage for PSyKAl}}}} section for more information
about how to use PSyKAl DSLs.


\subsection{PSyKAl file output}
\label{\detokenize{psyclone_command:psykal-file-output}}
\sphinxAtStartPar
By default the modified algorithm code and the generated PSy code are
output to the terminal. These can instead be output to files by using the
\sphinxcode{\sphinxupquote{\sphinxhyphen{}oalg \textless{}file\textgreater{}}} and \sphinxcode{\sphinxupquote{\sphinxhyphen{}opsy \textless{}file\textgreater{}}} options, respectively. For example, the
following will output the generated PSy code to the file ‘psy.f90’ but
the algorithm code will be output to the terminal:

\begin{sphinxVerbatim}[commandchars=\\\{\}]
\PYG{g+go}{psyclone \PYGZhy{}opsy psy.f90 algorithm.f90}
\end{sphinxVerbatim}

\sphinxAtStartPar
If PSyclone is being used to transform Kernels then the location to
write these to is specified using the \sphinxcode{\sphinxupquote{\sphinxhyphen{}okern \textless{}directory\textgreater{}}}
option. If this is not supplied then they are written to the current
working directory. By default, PSyclone will overwrite any kernel of
the same name in that directory. To change this behaviour, the user
can use the \sphinxcode{\sphinxupquote{\sphinxhyphen{}\sphinxhyphen{}no\_kernel\_clobber}} option. This causes PSyclone to
re\sphinxhyphen{}name any transformed kernel that would clash with any of those
already present in the output directory.


\subsection{Algorithm files with no invokes}
\label{\detokenize{psyclone_command:algorithm-files-with-no-invokes}}
\sphinxAtStartPar
If \sphinxcode{\sphinxupquote{psyclone}} is provided with a file that contains no
\sphinxcode{\sphinxupquote{invoke}} calls then the command outputs a warning to \sphinxcode{\sphinxupquote{stdout}} and
copies the input file to \sphinxcode{\sphinxupquote{stdout}}, or to the specified algorithm
file (if the \sphinxcode{\sphinxupquote{\sphinxhyphen{}oalg \textless{}file\textgreater{}}} option is used). No PSy code will be
output. If a file is specified using the \sphinxcode{\sphinxupquote{\sphinxhyphen{}opsy \textless{}file\textgreater{}}} option this file
will not be created.

\begin{sphinxVerbatim}[commandchars=\\\{\}]
\PYGZgt{}\PYG{+w}{ }psyclone\PYG{+w}{ }\PYGZhy{}opsy\PYG{+w}{ }psy.f90\PYG{+w}{ }\PYGZhy{}oalg\PYG{+w}{ }alg\PYGZus{}new.f90\PYG{+w}{ }empty\PYGZus{}alg.f90
Warning:\PYG{+w}{ }\PYG{l+s+s1}{\PYGZsq{}Algorithm Error: Algorithm file contains no invoke() calls: refusing to}
\PYG{l+s+s1}{generate empty PSy code\PYGZsq{}}
\end{sphinxVerbatim}


\subsection{Kernel search directory}
\label{\detokenize{psyclone_command:kernel-search-directory}}
\sphinxAtStartPar
When an algorithm file is parsed, the parser looks for the associated
kernel files. The way in which this is done requires that any
user\sphinxhyphen{}defined kernel routine (as opposed to {\hyperref[\detokenize{introduction_to_psykal:built-ins}]{\sphinxcrossref{\DUrole{std,std-ref}{Built\sphinxhyphen{}ins}}}}) called
within an invoke must have an explicit use statement. For example, the
following code gives an error:

\begin{sphinxVerbatim}[commandchars=\\\{\}]
\PYGZgt{}\PYG{+w}{ }cat\PYG{+w}{ }no\PYGZus{}use.f90
program\PYG{+w}{ }no\PYGZus{}use
\PYG{+w}{  }call\PYG{+w}{ }invoke\PYG{o}{(}testkern\PYGZus{}type\PYG{o}{(}a,b,c,d,e\PYG{o}{)}\PYG{o}{)}
end\PYG{+w}{ }program\PYG{+w}{ }no\PYGZus{}use
\PYGZgt{}\PYG{+w}{ }psyclone\PYG{+w}{ }\PYGZhy{}api\PYG{+w}{ }gocean\PYG{+w}{ }no\PYGZus{}use.f90
\PYG{l+s+s2}{\PYGZdq{}Parse Error: kernel call \PYGZsq{}testkern\PYGZus{}type\PYGZsq{} must either be named in a use statement or be a recognised built\PYGZhy{}in (one of \PYGZsq{}[]\PYGZsq{} for this API)\PYGZdq{}}
\end{sphinxVerbatim}

\sphinxAtStartPar
(If the chosen API has any {\hyperref[\detokenize{introduction_to_psykal:built-ins}]{\sphinxcrossref{\DUrole{std,std-ref}{Built\sphinxhyphen{}ins}}}} defined then
these will be listed within the \sphinxcode{\sphinxupquote{{[}{]}}} in the above error message.) If the
name of the kernel is provided in a use statement then the parser will
look for a file with the same name as the module in the use
statement. In the example below, the parser will look for a file
called “testkern.f90” or “testkern.F90”:

\begin{sphinxVerbatim}[commandchars=\\\{\}]
\PYGZgt{}\PYG{+w}{ }cat\PYG{+w}{ }use.f90
program\PYG{+w}{ }use
\PYG{+w}{  }use\PYG{+w}{ }testkern,\PYG{+w}{ }only\PYG{+w}{ }:\PYG{+w}{ }testkern\PYGZus{}type
\PYG{+w}{  }call\PYG{+w}{ }invoke\PYG{o}{(}testkern\PYGZus{}type\PYG{o}{(}a,b,c,d,e\PYG{o}{)}\PYG{o}{)}
end\PYG{+w}{ }program\PYG{+w}{ }use
\end{sphinxVerbatim}

\sphinxAtStartPar
Therefore, for PSyclone to find kernel files, the module name of a
kernel file must be the same as its filename. By default the parser
looks for the kernel file in the same directory as the algorithm
file. If this file is not found then an error is reported.

\begin{sphinxVerbatim}[commandchars=\\\{\}]
\PYGZgt{}\PYG{+w}{ }psyclone\PYG{+w}{ }use.f90
Kernel\PYG{+w}{ }file\PYG{+w}{ }\PYG{l+s+s1}{\PYGZsq{}testkern.[fF]90\PYGZsq{}}\PYG{+w}{ }not\PYG{+w}{ }found\PYG{+w}{ }\PYG{k}{in}\PYG{+w}{ }\PYGZlt{}location\PYGZgt{}
\end{sphinxVerbatim}

\sphinxAtStartPar
The \sphinxcode{\sphinxupquote{\sphinxhyphen{}d}} option can be used to tell \sphinxcode{\sphinxupquote{psyclone}} where to look for
kernel files by supplying it with a directory. The execution will recurse
from the specified directory path to look for the required file. There
must be only one instance of the specified file within (or below) the
specified directory:

\begin{sphinxVerbatim}[commandchars=\\\{\}]
\PYGZgt{}\PYG{+w}{ }\PYG{n+nb}{cd}\PYG{+w}{ }\PYGZlt{}PSYCLONEHOME\PYGZgt{}/src/psyclone
\PYGZgt{}\PYG{+w}{ }psyclone\PYG{+w}{ }\PYGZhy{}d\PYG{+w}{ }.\PYG{+w}{ }use.f90
More\PYG{+w}{ }than\PYG{+w}{ }one\PYG{+w}{ }match\PYG{+w}{ }\PYG{k}{for}\PYG{+w}{ }kernel\PYG{+w}{ }file\PYG{+w}{ }\PYG{l+s+s1}{\PYGZsq{}testkern.[fF]90\PYGZsq{}}\PYG{+w}{ }found!
\PYGZgt{}\PYG{+w}{ }psyclone\PYG{+w}{ }\PYGZhy{}d\PYG{+w}{ }tests/test\PYGZus{}files/dynamo0p3\PYG{+w}{ }\PYGZhy{}api\PYG{+w}{ }lfric\PYG{+w}{ }use.f90
\PYG{o}{[}code\PYG{+w}{ }output\PYG{o}{]}
\end{sphinxVerbatim}

\begin{sphinxadmonition}{note}{Note:}
\sphinxAtStartPar
The \sphinxcode{\sphinxupquote{\sphinxhyphen{}d}} option can be repeated to add as many search
directories as is required, with the constraint that there must be
only one instance of the specified file within (or below) the
specified directories.
\end{sphinxadmonition}


\subsection{Transforming PSyKAl Kernels}
\label{\detokenize{psyclone_command:transforming-psykal-kernels}}
\sphinxAtStartPar
When transforming kernels there are two use\sphinxhyphen{}cases to consider:
\begin{enumerate}
\sphinxsetlistlabels{\arabic}{enumi}{enumii}{}{.}%
\item {} 
\sphinxAtStartPar
a given kernel will be transformed only once and that version
then used from multiple, different Invokes and Algorithms;

\item {} 
\sphinxAtStartPar
a given kernel is used from multiple, different Invokes and
Algorithms and is transformed differently, depending on the
Invoke.

\end{enumerate}

\sphinxAtStartPar
Whenever PSyclone is used to transform a kernel, the new kernel must
be re\sphinxhyphen{}named in order to avoid clashing with other possible calls to
the original. By default (\sphinxcode{\sphinxupquote{\sphinxhyphen{}\sphinxhyphen{}kernel\sphinxhyphen{}renaming multiple}}), PSyclone
generates a new, unique name for each kernel that is
transformed. Since PSyclone is run on one Algorithm file at a time, it
uses the chosen kernel output directory (\sphinxcode{\sphinxupquote{\sphinxhyphen{}okern}}) to ensure that
names created by different invocations do not clash.  Therefore, when
building a single application, the same kernel output directory must
be used for each separate invocation of PSyclone.

\sphinxAtStartPar
Alternatively, in order to support use case 1, a user may specify
\sphinxcode{\sphinxupquote{\sphinxhyphen{}\sphinxhyphen{}kernel\sphinxhyphen{}renaming single}}: now, before transforming a kernel,
PSyclone will check the kernel output directory and if a transformed
version of that kernel is already present then that will be
used. Note, if the kernel file on disk does not match with what would
be generated then PSyclone will raise an exception.

\sphinxstepscope


\chapter{Configuration}
\label{\detokenize{configuration:configuration}}\label{\detokenize{configuration:id1}}\label{\detokenize{configuration::doc}}
\sphinxAtStartPar
PSyclone reads various run\sphinxhyphen{}time configuration options from
the \sphinxcode{\sphinxupquote{psyclone.cfg}} file. As described in
{\hyperref[\detokenize{getting_going:getting-going-configuration}]{\sphinxcrossref{\DUrole{std,std-ref}{Configuration}}}}, the default \sphinxcode{\sphinxupquote{psyclone.cfg}}
configuration file is installed in \sphinxcode{\sphinxupquote{\textless{}python\sphinxhyphen{}base\sphinxhyphen{}prefix\textgreater{}/share/psyclone/}}
during the installation process. The original version of this file
is in the \sphinxcode{\sphinxupquote{PSyclone/config}} directory of the PSyclone
distribution.

\sphinxAtStartPar
At execution\sphinxhyphen{}time, the user can specify a custom configuration file to
be used. This can either be done with the \sphinxcode{\sphinxupquote{\sphinxhyphen{}\sphinxhyphen{}config}} command\sphinxhyphen{}line
option, or by specifying the (full path to the) configuration file
to use via the \sphinxcode{\sphinxupquote{PSYCLONE\_CONFIG}} environment variable. If the specified
configuration file is not found then PSyclone will fall back to
searching in a list of default locations.

\sphinxAtStartPar
The ordering of these
locations depends upon whether PSyclone is being run within a Python
virtual environment (such as \sphinxcode{\sphinxupquote{venv}}). If no virtual environment is
detected then the locations searched, in order, are:
\begin{enumerate}
\sphinxsetlistlabels{\arabic}{enumi}{enumii}{}{.}%
\item {} 
\sphinxAtStartPar
\sphinxcode{\sphinxupquote{\$\{PWD\}/.psyclone/}}

\item {} 
\sphinxAtStartPar
\sphinxcode{\sphinxupquote{\$\{HOME\}/.local/share/psyclone/}}

\item {} 
\sphinxAtStartPar
\sphinxcode{\sphinxupquote{\textless{}python\sphinxhyphen{}base\sphinxhyphen{}dir\textgreater{}/share/psyclone/}}

\end{enumerate}

\sphinxAtStartPar
where \sphinxcode{\sphinxupquote{\textless{}python\sphinxhyphen{}base\sphinxhyphen{}dir\textgreater{}}} is the path stored in Python’s \sphinxcode{\sphinxupquote{sys.prefix}}.

\sphinxAtStartPar
If a virtual environment is detected then it is assumed that the
\sphinxcode{\sphinxupquote{share}} directory will be a part of that environment. In order to
maintain isolation of distinct virtual environments this directory is
then checked \sphinxstyleemphasis{before} the user’s home directory. i.e. the list of
locations searched is now:
\begin{enumerate}
\sphinxsetlistlabels{\arabic}{enumi}{enumii}{}{.}%
\item {} 
\sphinxAtStartPar
\sphinxcode{\sphinxupquote{\$\{PWD\}/.psyclone/}}

\item {} 
\sphinxAtStartPar
\sphinxcode{\sphinxupquote{\textless{}python\sphinxhyphen{}base\sphinxhyphen{}dir\textgreater{}/share/psyclone/}}

\item {} 
\sphinxAtStartPar
\sphinxcode{\sphinxupquote{\$\{HOME\}/.local/share/psyclone/}}

\end{enumerate}

\sphinxAtStartPar
As a last resort, the location
\sphinxcode{\sphinxupquote{\textless{}psyclone\sphinxhyphen{}src\sphinxhyphen{}base\textgreater{}/config/}}
is searched in case PSyclone was installed in editable mode.

\sphinxAtStartPar
Note that for developers a slightly different configuration handling
is implemented, see \sphinxhref{https://psyclone-dev.readthedocs.io/en/latest/modules.html\#dev-configuration}{Module: configuration} for details.


\section{Options}
\label{\detokenize{configuration:options}}
\sphinxAtStartPar
The configuration file is read by the Python ConfigParser class
(\sphinxurl{https://docs.python.org/3/library/configparser.html}) and must be
formatted accordingly. It currently consists of a \sphinxcode{\sphinxupquote{DEFAULT}}
section e.g.:

\begin{sphinxVerbatim}[commandchars=\\\{\}]
\PYG{p}{[}\PYG{n}{DEFAULT}\PYG{p}{]}
\PYG{n}{DISTRIBUTED\PYGZus{}MEMORY} \PYG{o}{=} \PYG{n}{true}
\PYG{n}{REPRODUCIBLE\PYGZus{}REDUCTIONS} \PYG{o}{=} \PYG{n}{false}
\PYG{n}{REPROD\PYGZus{}PAD\PYGZus{}SIZE} \PYG{o}{=} \PYG{l+m+mi}{8}
\PYG{n}{PSYIR\PYGZus{}ROOT\PYGZus{}NAME} \PYG{o}{=} \PYG{n}{psyir\PYGZus{}tmp}
\PYG{n}{VALID\PYGZus{}PSY\PYGZus{}DATA\PYGZus{}PREFIXES} \PYG{o}{=} \PYG{n}{profile}\PYG{p}{,} \PYG{n}{extract}
\PYG{n}{FORTRAN\PYGZus{}STANDARD} \PYG{o}{=} \PYG{n}{f2008}
\end{sphinxVerbatim}

\sphinxAtStartPar
and an optional API specific section, for example for the
\sphinxcode{\sphinxupquote{lfric}} section:

\begin{sphinxVerbatim}[commandchars=\\\{\}]
\PYG{p}{[}\PYG{n}{lfric}\PYG{p}{]}
\PYG{n}{access\PYGZus{}mapping} \PYG{o}{=} \PYG{n}{gh\PYGZus{}read}\PYG{p}{:} \PYG{n}{read}\PYG{p}{,} \PYG{n}{gh\PYGZus{}write}\PYG{p}{:} \PYG{n}{write}\PYG{p}{,} \PYG{n}{gh\PYGZus{}readwrite}\PYG{p}{:} \PYG{n}{readwrite}\PYG{p}{,}
                 \PYG{n}{gh\PYGZus{}inc}\PYG{p}{:} \PYG{n}{inc}\PYG{p}{,} \PYG{n}{gh\PYGZus{}readinc}\PYG{p}{:} \PYG{n}{readinc}\PYG{p}{,} \PYG{n}{gh\PYGZus{}sum}\PYG{p}{:} \PYG{n+nb}{sum}
\PYG{n}{COMPUTE\PYGZus{}ANNEXED\PYGZus{}DOFS} \PYG{o}{=} \PYG{n}{false}
\PYG{n}{supported\PYGZus{}fortran\PYGZus{}datatypes} \PYG{o}{=} \PYG{n}{real}\PYG{p}{,} \PYG{n}{integer}\PYG{p}{,} \PYG{n}{logical}
\PYG{n}{default\PYGZus{}kind} \PYG{o}{=} \PYG{n}{real}\PYG{p}{:} \PYG{n}{r\PYGZus{}def}\PYG{p}{,} \PYG{n}{integer}\PYG{p}{:} \PYG{n}{i\PYGZus{}def}\PYG{p}{,} \PYG{n}{logical}\PYG{p}{:} \PYG{n}{l\PYGZus{}def}
\PYG{n}{precision\PYGZus{}map} \PYG{o}{=} \PYG{n}{i\PYGZus{}def}\PYG{p}{:} \PYG{l+m+mi}{4}\PYG{p}{,}
                \PYG{n}{l\PYGZus{}def}\PYG{p}{:} \PYG{l+m+mi}{1}\PYG{p}{,}
                \PYG{n}{r\PYGZus{}def}\PYG{p}{:} \PYG{l+m+mi}{8}\PYG{p}{,}
                \PYG{n}{r\PYGZus{}double}\PYG{p}{:} \PYG{l+m+mi}{8}\PYG{p}{,}
                \PYG{n}{r\PYGZus{}ncdf}\PYG{p}{:} \PYG{l+m+mi}{8}\PYG{p}{,}
                \PYG{n}{r\PYGZus{}quad}\PYG{p}{:} \PYG{l+m+mi}{16}\PYG{p}{,}
                \PYG{n}{r\PYGZus{}second}\PYG{p}{:} \PYG{l+m+mi}{8}\PYG{p}{,}
                \PYG{n}{r\PYGZus{}single}\PYG{p}{:} \PYG{l+m+mi}{4}\PYG{p}{,}
                \PYG{n}{r\PYGZus{}solver}\PYG{p}{:} \PYG{l+m+mi}{4}\PYG{p}{,}
                \PYG{n}{r\PYGZus{}tran}\PYG{p}{:} \PYG{l+m+mi}{8}\PYG{p}{,}
                \PYG{n}{r\PYGZus{}bl}\PYG{p}{:} \PYG{l+m+mi}{8}\PYG{p}{,}
                \PYG{n}{r\PYGZus{}phys}\PYG{p}{:} \PYG{l+m+mi}{8}\PYG{p}{,}
                \PYG{n}{r\PYGZus{}um}\PYG{p}{:} \PYG{l+m+mi}{8}
\PYG{n}{RUN\PYGZus{}TIME\PYGZus{}CHECKS} \PYG{o}{=} \PYG{n}{false}
\PYG{n}{NUM\PYGZus{}ANY\PYGZus{}SPACE} \PYG{o}{=} \PYG{l+m+mi}{10}
\PYG{n}{NUM\PYGZus{}ANY\PYGZus{}DISCONTINUOUS\PYGZus{}SPACE} \PYG{o}{=} \PYG{l+m+mi}{10}
\end{sphinxVerbatim}

\sphinxAtStartPar
or for \sphinxcode{\sphinxupquote{gocean}}:

\begin{sphinxVerbatim}[commandchars=\\\{\}]
\PYG{p}{[}\PYG{n}{gocean}\PYG{p}{]}
\PYG{n}{access\PYGZus{}mapping} \PYG{o}{=} \PYG{n}{go\PYGZus{}read}\PYG{p}{:}\PYG{n}{read}\PYG{p}{,} \PYG{n}{go\PYGZus{}write}\PYG{p}{:}\PYG{n}{write}\PYG{p}{,} \PYG{n}{go\PYGZus{}readwrite}\PYG{p}{:}\PYG{n}{readwrite}
\PYG{n}{grid}\PYG{o}{\PYGZhy{}}\PYG{n}{properties} \PYG{o}{=} \PYG{n}{go\PYGZus{}grid\PYGZus{}xstop}\PYG{p}{:} \PYG{p}{\PYGZob{}}\PYG{l+m+mi}{0}\PYG{p}{\PYGZcb{}}\PYG{o}{\PYGZpc{}}\PYG{o}{\PYGZpc{}}\PYG{n}{grid}\PYG{o}{\PYGZpc{}}\PYG{o}{\PYGZpc{}}\PYG{n}{subdomain}\PYG{o}{\PYGZpc{}}\PYG{o}{\PYGZpc{}}\PYG{n}{internal}\PYG{o}{\PYGZpc{}}\PYG{o}{\PYGZpc{}}\PYG{n}{xstop}\PYG{p}{:} \PYG{n}{scalar}\PYG{p}{,}
               \PYG{n}{go\PYGZus{}grid\PYGZus{}ystop}\PYG{p}{:} \PYG{p}{\PYGZob{}}\PYG{l+m+mi}{0}\PYG{p}{\PYGZcb{}}\PYG{o}{\PYGZpc{}}\PYG{o}{\PYGZpc{}}\PYG{n}{grid}\PYG{o}{\PYGZpc{}}\PYG{o}{\PYGZpc{}}\PYG{n}{subdomain}\PYG{o}{\PYGZpc{}}\PYG{o}{\PYGZpc{}}\PYG{n}{internal}\PYG{o}{\PYGZpc{}}\PYG{o}{\PYGZpc{}}\PYG{n}{ystop}\PYG{p}{:} \PYG{n}{scalar}\PYG{p}{,}
               \PYG{n}{go\PYGZus{}grid\PYGZus{}data}\PYG{p}{:} \PYG{p}{\PYGZob{}}\PYG{l+m+mi}{0}\PYG{p}{\PYGZcb{}}\PYG{o}{\PYGZpc{}}\PYG{o}{\PYGZpc{}}\PYG{n}{data}\PYG{p}{:} \PYG{n}{array}\PYG{p}{,}
               \PYG{o}{.}\PYG{o}{.}\PYG{o}{.}
\end{sphinxVerbatim}

\sphinxAtStartPar
The meaning of the various entries is described in the following sub\sphinxhyphen{}sections.

\sphinxAtStartPar
Note that ConfigParser supports various forms of boolean entry
including “true/false”, “yes/no” and “1/0”. See
\sphinxurl{https://docs.python.org/3/library/configparser.html\#supported-datatypes}
for more details.


\subsection{\sphinxstyleliteralintitle{\sphinxupquote{DEFAULT}} Section}
\label{\detokenize{configuration:default-section}}\label{\detokenize{configuration:config-default-section}}
\sphinxAtStartPar
This section contains entries that are, in principle, applicable to all APIs
supported by PSyclone.


\begin{savenotes}\sphinxattablestart
\sphinxthistablewithglobalstyle
\sphinxthistablewithvlinesstyle
\centering
\begin{tabulary}{\linewidth}[t]{|l|L|l|}
\sphinxtoprule
\sphinxstyletheadfamily 
\sphinxAtStartPar
Entry
&\sphinxstyletheadfamily 
\sphinxAtStartPar
Description
&\sphinxstyletheadfamily 
\sphinxAtStartPar
Type
\\
\sphinxmidrule
\sphinxtableatstartofbodyhook
\sphinxAtStartPar
DISTRIBUTED\_MEMORY
&
\sphinxAtStartPar
Whether or not to generate code for distributed\sphinxhyphen{}memory
parallelism by default.  Note that this is currently
only supported for the LFRic (Dynamo 0.3) API.
&
\sphinxAtStartPar
bool
\\
\sphinxhline
\sphinxAtStartPar
REPRODUCIBLE\_REDUCTIONS
&
\sphinxAtStartPar
Whether or not to generate code for reproducible OpenMP
reductions (see {\hyperref[\detokenize{transformations:openmp-reductions}]{\sphinxcrossref{\DUrole{std,std-ref}{Reductions}}}}) by default.
&
\sphinxAtStartPar
bool
\\
\sphinxhline
\sphinxAtStartPar
REPROD\_PAD\_SIZE
&
\sphinxAtStartPar
If generating code for reproducible OpenMP reductions,
this setting controls the amount of padding used
between elements of the array in which each thread
accumulates its local reduction. (This prevents false
sharing of cache lines by different threads.)
&
\sphinxAtStartPar
int
\\
\sphinxhline
\sphinxAtStartPar
PSYIR\_ROOT\_NAME
&
\sphinxAtStartPar
The root for generated PSyIR symbol names if one is not
supplied when creating a symbol. Defaults to
“psyir\_tmp”.
&
\sphinxAtStartPar
str
\\
\sphinxhline
\sphinxAtStartPar
VALID\_PSY\_DATA\_PREFIXES
&
\sphinxAtStartPar
Which class prefixes are permitted in any
PSyData\sphinxhyphen{}related transformations. See {\hyperref[\detokenize{psy_data:psy-data}]{\sphinxcrossref{\DUrole{std,std-ref}{PSyData API}}}}
for details.
&
\sphinxAtStartPar
list of str
\\
\sphinxhline
\sphinxAtStartPar
BACKEND\_CHECKS\_ENABLED
&
\sphinxAtStartPar
Optional (defaults to True). Whether or not the PSyIR
backend should validate the tree that it is passed.
Can be overridden by the \sphinxcode{\sphinxupquote{\sphinxhyphen{}\sphinxhyphen{}backend}} command\sphinxhyphen{}line
flag (see {\hyperref[\detokenize{psyclone_command:backend-options}]{\sphinxcrossref{\DUrole{std,std-ref}{Backend Options}}}}).
&
\sphinxAtStartPar
bool
\\
\sphinxhline
\sphinxAtStartPar
FORTRAN\_STANDARD
&
\sphinxAtStartPar
Optional (defaults to f2008). The Fortran standard
that should be used by fparser. Valid values are
f2003 and f2008.
&
\sphinxAtStartPar
str
\\
\sphinxbottomrule
\end{tabulary}
\sphinxtableafterendhook\par
\sphinxattableend\end{savenotes}


\subsection{Common Sections}
\label{\detokenize{configuration:common-sections}}
\sphinxAtStartPar
The following entries must be defined for each API in order for PSyclone to
work as expected:


\begin{savenotes}\sphinxattablestart
\sphinxthistablewithglobalstyle
\sphinxthistablewithvlinesstyle
\centering
\begin{tabulary}{\linewidth}[t]{|l|L|}
\sphinxtoprule
\sphinxstyletheadfamily 
\sphinxAtStartPar
Entry
&\sphinxstyletheadfamily 
\sphinxAtStartPar
Description
\\
\sphinxmidrule
\sphinxtableatstartofbodyhook
\sphinxAtStartPar
access\_mapping
&
\sphinxAtStartPar
This field defines the strings that are used by a
particular API to indicate write, read, … access. Its
value is a comma separated list of access\sphinxhyphen{}string:access
pairs, e.g.:

\sphinxAtStartPar
\sphinxcode{\sphinxupquote{gh\_read: read, gh\_write: write, gh\_readwrite: readwrite,
gh\_inc: inc, gh\_readinc: gh\_sum: sum}}

\sphinxAtStartPar
At this stage these 6 types are defined for
read, write, read+write, increment,
read+increment and summation access by
PSyclone. Sum is a form of reduction. The
GOcean API does not support increment or sum,
so it only defines three mappings for read,
write, and readwrite.
\\
\sphinxbottomrule
\end{tabulary}
\sphinxtableafterendhook\par
\sphinxattableend\end{savenotes}


\subsection{\sphinxstyleliteralintitle{\sphinxupquote{lfric}} Section}
\label{\detokenize{configuration:lfric-section}}
\sphinxAtStartPar
This section contains configuration options that are only applicable when
using the LFRic (Dynamo 0.3) API.


\begin{savenotes}\sphinxattablestart
\sphinxthistablewithglobalstyle
\sphinxthistablewithvlinesstyle
\centering
\begin{tabulary}{\linewidth}[t]{|l|L|}
\sphinxtoprule
\sphinxstyletheadfamily 
\sphinxAtStartPar
Entry
&\sphinxstyletheadfamily 
\sphinxAtStartPar
Description
\\
\sphinxmidrule
\sphinxtableatstartofbodyhook
\sphinxAtStartPar
COMPUTE\_ANNEXED\_DOFS
&
\sphinxAtStartPar
Whether or not to perform redundant computation
over annexed dofs in order to reduce the number of
halo exchanges, see {\hyperref[\detokenize{dynamo0p3:lfric-annexed-dofs}]{\sphinxcrossref{\DUrole{std,std-ref}{Annexed DoFs}}}}.
\\
\sphinxhline
\sphinxAtStartPar
supported\_fortran\_datatypes
&
\sphinxAtStartPar
Captures the supported Fortran data types of LFRic
arguments, see {\hyperref[\detokenize{dynamo0p3:lfric-datatype-kind}]{\sphinxcrossref{\DUrole{std,std-ref}{Supported Data Types and Default Kind}}}}.
\\
\sphinxhline
\sphinxAtStartPar
default\_kind
&
\sphinxAtStartPar
Captures the default kinds (precisions) for the
supported Fortran data types in LFRic, see
{\hyperref[\detokenize{dynamo0p3:lfric-datatype-kind}]{\sphinxcrossref{\DUrole{std,std-ref}{Supported Data Types and Default Kind}}}}.
\\
\sphinxhline
\sphinxAtStartPar
precision\_map
&
\sphinxAtStartPar
Captures the value of the actual precisions in
bytes, see {\hyperref[\detokenize{dynamo0p3:lfric-precision-map}]{\sphinxcrossref{\DUrole{std,std-ref}{Precision Map}}}}
\\
\sphinxhline
\sphinxAtStartPar
RUN\_TIME\_CHECKS
&
\sphinxAtStartPar
Specifies whether to generate run\sphinxhyphen{}time validation
checks, see {\hyperref[\detokenize{dynamo0p3:lfric-run-time-checks}]{\sphinxcrossref{\DUrole{std,std-ref}{Run\sphinxhyphen{}time Checks}}}}.
\\
\sphinxhline
\sphinxAtStartPar
NUM\_ANY\_SPACE
&
\sphinxAtStartPar
Sets the number of \sphinxcode{\sphinxupquote{ANY\_SPACE}} function spaces
in LFRic, see {\hyperref[\detokenize{dynamo0p3:lfric-num-any-spaces}]{\sphinxcrossref{\DUrole{std,std-ref}{Number of Generalised ANY\_*\_SPACE Function Spaces}}}}.
\\
\sphinxhline
\sphinxAtStartPar
NUM\_ANY\_DISCONTINUOUS\_SPACE
&
\sphinxAtStartPar
Sets the number of \sphinxcode{\sphinxupquote{ANY\_DISCONTINUOUS\_SPACE}}
function spaces in LFRic, see
{\hyperref[\detokenize{dynamo0p3:lfric-num-any-spaces}]{\sphinxcrossref{\DUrole{std,std-ref}{Number of Generalised ANY\_*\_SPACE Function Spaces}}}}.
\\
\sphinxbottomrule
\end{tabulary}
\sphinxtableafterendhook\par
\sphinxattableend\end{savenotes}


\subsection{\sphinxstyleliteralintitle{\sphinxupquote{gocean}} Section}
\label{\detokenize{configuration:gocean-section}}
\sphinxAtStartPar
This section contains configuration options that are only applicable when
using the Gocean 1.0 API.


\begin{savenotes}\sphinxattablestart
\sphinxthistablewithglobalstyle
\sphinxthistablewithvlinesstyle
\centering
\begin{tabulary}{\linewidth}[t]{|l|L|}
\sphinxtoprule
\sphinxstyletheadfamily 
\sphinxAtStartPar
Entry
&\sphinxstyletheadfamily 
\sphinxAtStartPar
Description
\\
\sphinxmidrule
\sphinxtableatstartofbodyhook
\sphinxAtStartPar
iteration\sphinxhyphen{}spaces
&
\sphinxAtStartPar
This contains definitions of additional iteration spaces
used by PSyclone. A detailed description can be found
in the {\hyperref[\detokenize{gocean1p0:gocean-configuration-iteration-spaces}]{\sphinxcrossref{\DUrole{std,std-ref}{Iteration\sphinxhyphen{}spaces}}}}
section of the GOcean1.0 chapter.
\\
\sphinxhline
\sphinxAtStartPar
grid\sphinxhyphen{}properties
&
\sphinxAtStartPar
This key contains definitions to access various grid
properties. A detailed description can be found
in the {\hyperref[\detokenize{gocean1p0:gocean-configuration-grid-properties}]{\sphinxcrossref{\DUrole{std,std-ref}{Grid Properties}}}}
section of the GOcean1.0 chapter.
\\
\sphinxbottomrule
\end{tabulary}
\sphinxtableafterendhook\par
\sphinxattableend\end{savenotes}

\sphinxstepscope


\chapter{Tutorial and Examples}
\label{\detokenize{tutorials_and_examples:tutorial-and-examples}}\label{\detokenize{tutorials_and_examples:tutorial}}\label{\detokenize{tutorials_and_examples::doc}}

\section{Tutorial}
\label{\detokenize{tutorials_and_examples:id1}}
\sphinxAtStartPar
PSyclone provides a hands\sphinxhyphen{}on tutorial. The easiest way to follow it is reading
the \sphinxhref{https://github.com/stfc/PSyclone/tree/master/tutorial/practicals}{Readme files in github}.
The tutorial is divided into two sections, a first section that introduces
PSyclone and how to
\sphinxhref{https://github.com/stfc/PSyclone/tree/master/tutorial/practicals/generic}{use it to transform generic Fortran code}
(this is the recommended starting point for everybody).
And a second section about
\sphinxhref{https://github.com/stfc/PSyclone/tree/master/tutorial/practicals/LFRic}{the LFRic DSL}
(this is only recommended for people interested in PSyKAL DSLs and LFRic in particular).

\sphinxAtStartPar
To do the proposed hands\sphinxhyphen{}on you will need a linux shell with Python installed and to
download the hands\sphinxhyphen{}on directory with:

\begin{sphinxVerbatim}[commandchars=\\\{\}]
git\PYG{+w}{ }clone\PYG{+w}{ }\PYGZhy{}\PYGZhy{}recursive\PYG{+w}{ }git@github.com:stfc/PSyclone.git
\PYG{n+nb}{cd}\PYG{+w}{ }PSyclone
\PYG{c+c1}{\PYGZsh{} If psyclone isn\PYGZsq{}t already installed you can use \PYGZsq{}pip\PYGZsq{} in this folder to}
\PYG{c+c1}{\PYGZsh{} install a version that matches the downloaded tutorials}
pip\PYG{+w}{ }install\PYG{+w}{ }.
\PYG{n+nb}{cd}\PYG{+w}{ }tutorial/practicals
\end{sphinxVerbatim}


\section{Examples}
\label{\detokenize{tutorials_and_examples:examples}}\label{\detokenize{tutorials_and_examples:id2}}
\sphinxAtStartPar
Various examples of the use of PSyclone are provided under the
\sphinxcode{\sphinxupquote{examples}} directory in the Git repository. If you have installed
PSyclone using \sphinxcode{\sphinxupquote{pip}} then the examples may be found in
\sphinxcode{\sphinxupquote{share/psyclone/examples}} in psyclone  {\hyperref[\detokenize{getting_going:getting-going-install-loc}]{\sphinxcrossref{\DUrole{std,std-ref}{Installation location}}}}.

\sphinxAtStartPar
Running any of these examples requires that PSyclone be installed on
the host system, see Section {\hyperref[\detokenize{getting_going:getting-going}]{\sphinxcrossref{\DUrole{std,std-ref}{Getting Going}}}}.
This section is intended to provide an overview of the various examples
so that a user can find one that is appropriate to them. For details of
what each example does and how to run each example please see the
\sphinxcode{\sphinxupquote{README.md}} files in the associated directories.

\sphinxAtStartPar
For the purposes of correctness checking, the whole suite of examples
may be executed using Gnu \sphinxcode{\sphinxupquote{make}} (this functionality is used by GitHub
Actions alongside the test suite). The default target is \sphinxcode{\sphinxupquote{transform}} which
just performs the PSyclone code transformation steps for each
example. For those examples that support it, the \sphinxcode{\sphinxupquote{compile}} target
also requests that the generated code be compiled. The \sphinxcode{\sphinxupquote{notebook}}
target checks the various Jupyter notebooks using \sphinxcode{\sphinxupquote{nbconvert}}.

\begin{sphinxadmonition}{note}{Note:}
\sphinxAtStartPar
As outlined in the {\hyperref[\detokenize{getting_going:getting-going-run}]{\sphinxcrossref{\DUrole{std,std-ref}{Run}}}} section, if
working with the examples from a PSyclone installation, it is
advisable to copy the whole \sphinxcode{\sphinxupquote{examples}} directory to some
convenient location before running them. If you have copied the
\sphinxcode{\sphinxupquote{examples}} directory but still wish to use \sphinxcode{\sphinxupquote{make}} then you
will also have to set the \sphinxcode{\sphinxupquote{PSYCLONE\_CONFIG}} environment variable
to the full path to the PSyclone configuration file, e.g.
\sphinxcode{\sphinxupquote{PSYCLONE\_CONFIG=/some/path/psyclone.cfg make}}.
\end{sphinxadmonition}


\subsection{Compilation}
\label{\detokenize{tutorials_and_examples:compilation}}\label{\detokenize{tutorials_and_examples:examples-compilation}}
\sphinxAtStartPar
Some of the examples support compilation (and some even execution of
a compiled binary). Please consult the \sphinxcode{\sphinxupquote{README.md}} to check which ones
can be compiled and executed.

\sphinxAtStartPar
As mentioned above, by default each example will execute the
\sphinxcode{\sphinxupquote{transform}} target, which performs the PSyclone code transformation
steps. In order to compile the sources, use the target \sphinxcode{\sphinxupquote{compile}}:

\begin{sphinxVerbatim}[commandchars=\\\{\}]
make\PYG{+w}{ }compile
\end{sphinxVerbatim}

\sphinxAtStartPar
which will first perform the transformation steps before compiling
any created Fortan source files. If the example also supports running
a compiled and linked binary, use the target:

\begin{sphinxVerbatim}[commandchars=\\\{\}]
make\PYG{+w}{ }run
\end{sphinxVerbatim}

\sphinxAtStartPar
This will first trigger compilation using the \sphinxcode{\sphinxupquote{compile}} target, and
then execute the program with any parameters that might be required
(check the corresponding \sphinxcode{\sphinxupquote{README.md}} document for details).

\sphinxAtStartPar
All \sphinxcode{\sphinxupquote{Makefile}}s support the variables \sphinxcode{\sphinxupquote{F90}} and \sphinxcode{\sphinxupquote{F90FLAGS}} to specify
the compiler and compilation flags to use. By default, the Gnu Fortran
compiler (\sphinxcode{\sphinxupquote{gfortran}}) is used, and the compilation flags will be set
to debugging. If you want to change the compiler or flags, just define
these as environment variables:

\begin{sphinxVerbatim}[commandchars=\\\{\}]
\PYG{n+nv}{F90}\PYG{o}{=}ifort\PYG{+w}{ }\PYG{n+nv}{F90FLAGS}\PYG{o}{=}\PYG{l+s+s2}{\PYGZdq{}\PYGZhy{}g \PYGZhy{}check bounds\PYGZdq{}}\PYG{+w}{ }make\PYG{+w}{ }compile
\end{sphinxVerbatim}

\sphinxAtStartPar
To clean all compiled files (and potential output files from a run),
use:

\begin{sphinxVerbatim}[commandchars=\\\{\}]
make\PYG{+w}{ }clean
\end{sphinxVerbatim}

\sphinxAtStartPar
This will clean up in the \sphinxcode{\sphinxupquote{examples}} directory. If you want to change compilers
or compiler flags, you should run \sphinxcode{\sphinxupquote{make allclean}}, see the section
about {\hyperref[\detokenize{tutorials_and_examples:examples-dependencies}]{\sphinxcrossref{\DUrole{std,std-ref}{Dependencies}}}} for details.


\subsection{Supported Compilers}
\label{\detokenize{tutorials_and_examples:supported-compilers}}\label{\detokenize{tutorials_and_examples:id3}}
\sphinxAtStartPar
All examples have been tested with the following compilers.
Please let the developers know if you have problems using a compiler
that has been tested or if you are working with a different compiler
so it can be recorded in this table.


\begin{savenotes}\sphinxattablestart
\sphinxthistablewithglobalstyle
\sphinxthistablewithvlinesstyle
\centering
\begin{tabulary}{\linewidth}[t]{|l|L|}
\sphinxtoprule
\sphinxstyletheadfamily 
\sphinxAtStartPar
Compiler
&\sphinxstyletheadfamily 
\sphinxAtStartPar
Version
\\
\sphinxmidrule
\sphinxtableatstartofbodyhook
\sphinxAtStartPar
Gnu Fortran
&
\sphinxAtStartPar
9.3
\\
\sphinxhline
\sphinxAtStartPar
Intel Fortran
&
\sphinxAtStartPar
17, 21
\\
\sphinxhline
\sphinxAtStartPar
NVIDIA Fortran
&
\sphinxAtStartPar
23.5
\\
\sphinxbottomrule
\end{tabulary}
\sphinxtableafterendhook\par
\sphinxattableend\end{savenotes}


\subsection{Dependencies}
\label{\detokenize{tutorials_and_examples:dependencies}}\label{\detokenize{tutorials_and_examples:examples-dependencies}}
\sphinxAtStartPar
Any required library that is included in PSyclone (typically
the infrastructure libraries for the APIs, or {\hyperref[\detokenize{libraries:libraries}]{\sphinxcrossref{\DUrole{std,std-ref}{PSyData wrapper
libraries}}}}) will automatically be compiled with the same
compiler and compilation flags as the examples.

\begin{sphinxadmonition}{note}{Note:}
\sphinxAtStartPar
Once a dependent library is compiled, changing the
compilation flags will not trigger a recompilation
of this library. For example, if an example is first compiled
with debug options, and later the same or a different
example is compiled with optimisations, the dependent library
will not automatically be recompiled!
\end{sphinxadmonition}

\sphinxAtStartPar
All \sphinxcode{\sphinxupquote{Makefile}}s support an \sphinxcode{\sphinxupquote{allclean}} target, which will not only
clean the current directory, but also all libraries the current
example depends on.

\begin{sphinxadmonition}{important}{Important:}
\sphinxAtStartPar
Using \sphinxcode{\sphinxupquote{make allclean}} is especially important if
the compiler is changed. Typically, one compiler cannot
read module information from a different compiler, and
then compilation will fail.
\end{sphinxadmonition}


\subsubsection{NetCDF}
\label{\detokenize{tutorials_and_examples:netcdf}}
\sphinxAtStartPar
Some examples require NetCDF for compilation. Installation of NetCDF
is described in detail in
\sphinxhref{https://github.com/stfc/PSyclone/tree/master/tutorial/practicals\#user-content-netcdf-library-lfric-examples}{the hands\sphinxhyphen{}on practicals documentation}.


\section{PSyIR Examples}
\label{\detokenize{tutorials_and_examples:psyir-examples}}
\sphinxAtStartPar
Examples may all be found in the \sphinxcode{\sphinxupquote{examples/psyir}} directory. Read the
\sphinxcode{\sphinxupquote{README.md}} file in this directory for full details.


\subsection{Example 1: Constructing PSyIR and Generating Code}
\label{\detokenize{tutorials_and_examples:example-1-constructing-psyir-and-generating-code}}
\sphinxAtStartPar
\sphinxcode{\sphinxupquote{create.py}} is a Python script that demonstrates the use of the various
\sphinxcode{\sphinxupquote{create}} methods to build a PSyIR tree from scratch.


\subsection{Example 2: Creating PSyIR for Structure Types}
\label{\detokenize{tutorials_and_examples:example-2-creating-psyir-for-structure-types}}
\sphinxAtStartPar
\sphinxcode{\sphinxupquote{create\_structure\_types.py}} demonstrates the representation of
structure types (i.e. Fortran derived types or C structs) in the PSyIR.


\section{GOcean Examples}
\label{\detokenize{tutorials_and_examples:gocean-examples}}

\subsection{Example 1: Loop transformations}
\label{\detokenize{tutorials_and_examples:example-1-loop-transformations}}
\sphinxAtStartPar
Examples of applying various transformations (loop fusion, OpenMP,
OpenMP Taskloop, OpenACC, OpenCL) to the semi\sphinxhyphen{}PSyKAl’d version of
the Shallow benchmark. (“semi” because not all kernels are called
from within invoke()’s.) Also includes an example of generating a
DAG from an InvokeSchedule.


\subsection{Example 2: OpenACC}
\label{\detokenize{tutorials_and_examples:example-2-openacc}}
\sphinxAtStartPar
This is a simple but complete example of using PSyclone to enable an
application to run on a GPU by adding OpenACC directives. A \sphinxcode{\sphinxupquote{Makefile}}
is included which will use PSyclone to generate the PSy code and
transformed kernels and then compile the application. This compilation
requires that the \sphinxhref{https://github.com/stfc/dl\_esm\_inf}{dl\_esm\_inf library}
be installed/available \sphinxhyphen{} it is provided as a Git submodule of the PSyclone
project (see \sphinxhref{https://psyclone-dev.readthedocs.io/en/latest/working\_practises.html\#dev-installation}{Installation} in the Developers’ Guide
for details on working with submodules).

\sphinxAtStartPar
The supplied \sphinxcode{\sphinxupquote{Makefile}} also provides a second, \sphinxcode{\sphinxupquote{profile}} target which
performs the same OpenACC transformations but then encloses the whole
of the resulting PSy layer in a profiling region. By linking this with
the PSyclone NVTX profiling wrapper (and the NVTX library itself), the
resulting application can be profiled using NVIDIA’s \sphinxtitleref{nvprof} or
\sphinxtitleref{nvvp} tools.


\subsection{Example 3: OpenCL}
\label{\detokenize{tutorials_and_examples:example-3-opencl}}
\sphinxAtStartPar
Example of the use of PSyclone to generate an OpenCL driver version of
the PSy layer and OpenCL kernels. The \sphinxcode{\sphinxupquote{Makefile}} in this example provides
a target (\sphinxtitleref{make compile\sphinxhyphen{}ocl}) to compile the generated OpenCL code. This
requires an OpenCL implementation installed in the system. Read the README
provided in the example folder for more details about how to compile and
execute the generated OpenCL code.


\subsection{Example 4: Kernels containing use statements}
\label{\detokenize{tutorials_and_examples:example-4-kernels-containing-use-statements}}
\sphinxAtStartPar
Transforming kernels for use with either OpenACC or OpenCL requires
that we handle those that access data and/or routines via module
\sphinxcode{\sphinxupquote{use}} statements. This example shows the various forms for which
support is being implemented. Although there is support for converting
global\sphinxhyphen{}data accesses into kernel arguments, PSyclone does not yet support
nested \sphinxcode{\sphinxupquote{use}} of modules (i.e. data accessed via a module that in turn
imports that symbol from another module) and kernels that call other
kernels (Issue \#342).


\subsection{Example 5: PSyData}
\label{\detokenize{tutorials_and_examples:example-5-psydata}}\label{\detokenize{tutorials_and_examples:gocean-example-psydata}}
\sphinxAtStartPar
This directory contains all examples that use the
{\hyperref[\detokenize{psy_data:psy-data}]{\sphinxcrossref{\DUrole{std,std-ref}{PSyData API}}}}. At this stage there are three
runnable examples:


\subsubsection{Example 5.1: Kernel data extraction}
\label{\detokenize{tutorials_and_examples:example-5-1-kernel-data-extraction}}
\sphinxAtStartPar
This example shows the use of kernel data extraction in PSyclone.
It instruments each of the two invokes in the example program
with the PSyData\sphinxhyphen{}based kernel extraction code. Detailed compilation
instructions are in the \sphinxcode{\sphinxupquote{README.md}} file, including how to switch
from using the stand\sphinxhyphen{}alone extraction library to the NetCDF\sphinxhyphen{}based one
(see {\hyperref[\detokenize{psyke:extraction-libraries}]{\sphinxcrossref{\DUrole{std,std-ref}{Extraction Libraries}}}} for details).

\sphinxAtStartPar
The \sphinxcode{\sphinxupquote{Makefile}} in this example will create the binary that extracts
the data at run time, as well as two driver programs that can read in
the extracted data, call the kernel, and compare the results. These
driver programs are independent of the dl\_esm\_inf infrastructure library.
These drivers can only read the corresponding file format, i.e. a NetCDF
driver program cannot read in extraction data that is based on Fortran IO
and vice versa.

\begin{sphinxadmonition}{note}{Note:}
\sphinxAtStartPar
At this stage the driver program still needs the infrastructure
library when compiling the kernels, see \#1757.
\end{sphinxadmonition}


\subsubsection{Example 5.2: Profiling}
\label{\detokenize{tutorials_and_examples:example-5-2-profiling}}
\sphinxAtStartPar
This example shows how to use the profiling support in PSyclone.
It instruments two invoke statements and can link in with any
of the following profiling wrapper libraries: template,
simple\_timer, dl\_timer, TAU, Vernier, and DrHook (see
{\hyperref[\detokenize{profiling:profiling-third-party-tools}]{\sphinxcrossref{\DUrole{std,std-ref}{Interface to Third Party Profiling Tools}}}}). The \sphinxcode{\sphinxupquote{README.md}}
file contains detailed instructions on how to build the
different executables. By default (i.e. just using \sphinxcode{\sphinxupquote{make}}
without additional parameters) it links in with the
template profiling library included in PSyclone. This library just
prints out the name of the module and region before and after each
invoke is executed. This example can actually be executed to
test the behaviour of the various profiling wrappers, and is
also useful if you want to develop your own wrapper libraries.


\subsubsection{Example 5.3: Read\sphinxhyphen{}only\sphinxhyphen{}verification}
\label{\detokenize{tutorials_and_examples:example-5-3-read-only-verification}}\label{\detokenize{tutorials_and_examples:gocean-example-readonly}}
\sphinxAtStartPar
This example shows the use of read\sphinxhyphen{}only\sphinxhyphen{}verification with PSyclone.
It instruments each of the two invokes in the example program
with the PSyData\sphinxhyphen{}based read\sphinxhyphen{}only\sphinxhyphen{}verification code.
It uses the dl\_esm\_inf\sphinxhyphen{}specific read\sphinxhyphen{}only\sphinxhyphen{}verification library
(\sphinxcode{\sphinxupquote{lib/read\_only/dl\_esm\_inf/}}).

\begin{sphinxadmonition}{note}{Note:}
\sphinxAtStartPar
The \sphinxcode{\sphinxupquote{update\_field\_mod}} subroutine contains some very
buggy and non\sphinxhyphen{}standard code to change the value of some
read\sphinxhyphen{}only variables and fields, even though the variables
are all declared with
\sphinxcode{\sphinxupquote{intent(in)}}. It uses the addresses of variables and
then out\sphinxhyphen{}of\sphinxhyphen{}bound writes to a writeable array to
actually overwrite the read\sphinxhyphen{}only variables. Using
array bounds checking at runtime will be triggered by these
out\sphinxhyphen{}of\sphinxhyphen{}bound writes.
\end{sphinxadmonition}

\sphinxAtStartPar
The \sphinxcode{\sphinxupquote{Makefile}} in this example will link with the compiled
read\sphinxhyphen{}only\sphinxhyphen{}verification library. You can execute the created
binary and it will print two warnings about modified
read\sphinxhyphen{}only variables:

\begin{sphinxVerbatim}[commandchars=\\\{\}]
\PYGZhy{}\PYGZhy{}\PYGZhy{}\PYGZhy{}\PYGZhy{}\PYGZhy{}\PYGZhy{}\PYGZhy{}\PYGZhy{}\PYGZhy{}\PYGZhy{}\PYGZhy{}\PYGZhy{}\PYGZhy{}\PYGZhy{}\PYGZhy{}\PYGZhy{}\PYGZhy{}\PYGZhy{}\PYGZhy{}\PYGZhy{}\PYGZhy{}\PYGZhy{}\PYGZhy{}\PYGZhy{}\PYGZhy{}\PYGZhy{}\PYGZhy{}\PYGZhy{}\PYGZhy{}\PYGZhy{}\PYGZhy{}\PYGZhy{}\PYGZhy{}\PYGZhy{}\PYGZhy{}\PYGZhy{}\PYGZhy{}
Double precision field b\PYGZus{}fld has been modified in main : update
Original checksum:   4611686018427387904
New checksum:        4638355772470722560
\PYGZhy{}\PYGZhy{}\PYGZhy{}\PYGZhy{}\PYGZhy{}\PYGZhy{}\PYGZhy{}\PYGZhy{}\PYGZhy{}\PYGZhy{}\PYGZhy{}\PYGZhy{}\PYGZhy{}\PYGZhy{}\PYGZhy{}\PYGZhy{}\PYGZhy{}\PYGZhy{}\PYGZhy{}\PYGZhy{}\PYGZhy{}\PYGZhy{}\PYGZhy{}\PYGZhy{}\PYGZhy{}\PYGZhy{}\PYGZhy{}\PYGZhy{}\PYGZhy{}\PYGZhy{}\PYGZhy{}\PYGZhy{}\PYGZhy{}\PYGZhy{}\PYGZhy{}\PYGZhy{}\PYGZhy{}\PYGZhy{}
\PYGZhy{}\PYGZhy{}\PYGZhy{}\PYGZhy{}\PYGZhy{}\PYGZhy{}\PYGZhy{}\PYGZhy{}\PYGZhy{}\PYGZhy{}\PYGZhy{}\PYGZhy{}\PYGZhy{}\PYGZhy{}\PYGZhy{}\PYGZhy{}\PYGZhy{}\PYGZhy{}\PYGZhy{}\PYGZhy{}\PYGZhy{}\PYGZhy{}\PYGZhy{}\PYGZhy{}\PYGZhy{}\PYGZhy{}\PYGZhy{}\PYGZhy{}\PYGZhy{}\PYGZhy{}\PYGZhy{}\PYGZhy{}\PYGZhy{}\PYGZhy{}\PYGZhy{}\PYGZhy{}\PYGZhy{}\PYGZhy{}
Double precision variable z has been modified in main : update
Original value:    1.0000000000000000
New value:         123.00000000000000
\PYGZhy{}\PYGZhy{}\PYGZhy{}\PYGZhy{}\PYGZhy{}\PYGZhy{}\PYGZhy{}\PYGZhy{}\PYGZhy{}\PYGZhy{}\PYGZhy{}\PYGZhy{}\PYGZhy{}\PYGZhy{}\PYGZhy{}\PYGZhy{}\PYGZhy{}\PYGZhy{}\PYGZhy{}\PYGZhy{}\PYGZhy{}\PYGZhy{}\PYGZhy{}\PYGZhy{}\PYGZhy{}\PYGZhy{}\PYGZhy{}\PYGZhy{}\PYGZhy{}\PYGZhy{}\PYGZhy{}\PYGZhy{}\PYGZhy{}\PYGZhy{}\PYGZhy{}\PYGZhy{}\PYGZhy{}\PYGZhy{}
\end{sphinxVerbatim}


\subsubsection{Example 5.4: Value Range Check}
\label{\detokenize{tutorials_and_examples:example-5-4-value-range-check}}\label{\detokenize{tutorials_and_examples:gocean-example-value-range-check}}
\sphinxAtStartPar
This example shows the use of valid number verification with PSyclone.
It instruments each of the two invokes in the example program
with the PSyData\sphinxhyphen{}based Value\sphinxhyphen{}Range\sphinxhyphen{}Check code.
It uses the dl\_esm\_inf\sphinxhyphen{}specific value range check library
(\sphinxcode{\sphinxupquote{lib/value\_range\_check/dl\_esm\_inf/}}).

\begin{sphinxadmonition}{note}{Note:}
\sphinxAtStartPar
The \sphinxcode{\sphinxupquote{update\_field\_mod}} subroutine contains code
that will trigger a division by 0 to create NaNs. If
the compiler happens to add code that handles floating point
exceptions , this will take effect before the value testing
is done by the PSyData\sphinxhyphen{}based verification code.
\end{sphinxadmonition}

\sphinxAtStartPar
The \sphinxcode{\sphinxupquote{Makefile}} in this example will link with the compiled
value\_range\_check library. You can then execute the binary
and enable the value range check by setting environments
(see {\hyperref[\detokenize{psy_data:psydata-value-range-check}]{\sphinxcrossref{\DUrole{std,std-ref}{value range check}}}} for
details).

\begin{sphinxVerbatim}[commandchars=\\\{\}]
\PYG{n+nv}{PSYVERIFY\PYGZus{}\PYGZus{}main\PYGZus{}\PYGZus{}init\PYGZus{}\PYGZus{}b\PYGZus{}fld}\PYG{o}{=}\PYG{l+m}{2}:3\PYG{+w}{ }./value\PYGZus{}range\PYGZus{}check
...
PSyData:\PYG{+w}{ }Variable\PYG{+w}{ }b\PYGZus{}fld\PYG{+w}{ }has\PYG{+w}{ }the\PYG{+w}{ }value\PYG{+w}{ }\PYG{l+m}{0}.0000000000000000\PYG{+w}{ }at\PYG{+w}{ }index/indices\PYG{+w}{ }\PYG{l+m}{6}\PYG{+w}{ }\PYG{l+m}{1}\PYG{+w}{ }\PYG{k}{in}\PYG{+w}{ }module\PYG{+w}{ }\PYG{l+s+s1}{\PYGZsq{}main\PYGZsq{}},\PYG{+w}{ }region\PYG{+w}{ }\PYG{l+s+s1}{\PYGZsq{}init\PYGZsq{}},\PYG{+w}{ }which\PYG{+w}{ }is\PYG{+w}{ }not\PYG{+w}{ }between\PYG{+w}{ }\PYG{l+s+s1}{\PYGZsq{}2.0000000000000000\PYGZsq{}}\PYG{+w}{ }and\PYG{+w}{ }\PYG{l+s+s1}{\PYGZsq{}3.0000000000000000\PYGZsq{}}.
...
PSyData:\PYG{+w}{ }Variable\PYG{+w}{ }a\PYGZus{}fld\PYG{+w}{ }has\PYG{+w}{ }the\PYG{+w}{ }invalid\PYG{+w}{ }value\PYG{+w}{ }\PYG{l+s+s1}{\PYGZsq{}Inf\PYGZsq{}}\PYG{+w}{ }at\PYG{+w}{ }index/indices\PYG{+w}{ }\PYG{l+m}{1}\PYG{+w}{ }\PYG{l+m}{1}\PYG{+w}{ }\PYG{k}{in}\PYG{+w}{ }module\PYG{+w}{ }\PYG{l+s+s1}{\PYGZsq{}main\PYGZsq{}},\PYG{+w}{ }region\PYG{+w}{ }\PYG{l+s+s1}{\PYGZsq{}update\PYGZsq{}}.
\end{sphinxVerbatim}

\sphinxAtStartPar
As indicated in {\hyperref[\detokenize{psy_data:psydata-value-range-check}]{\sphinxcrossref{\DUrole{std,std-ref}{value range check}}}}, you can
also check a variable in all kernels of a module, or in any instrumented
code region (since the example has only one module, both settings below
will create the same warnings):

\begin{sphinxVerbatim}[commandchars=\\\{\}]
\PYG{n+nv}{PSYVERIFY\PYGZus{}\PYGZus{}main\PYGZus{}\PYGZus{}b\PYGZus{}fld}\PYG{o}{=}\PYG{l+m}{2}:3\PYG{+w}{ }./value\PYGZus{}range\PYGZus{}check
\PYG{n+nv}{PSYVERIFY\PYGZus{}\PYGZus{}b\PYGZus{}fld}\PYG{o}{=}\PYG{l+m}{2}:3\PYG{+w}{ }./value\PYGZus{}range\PYGZus{}check
...
PSyData:\PYG{+w}{ }Variable\PYG{+w}{ }b\PYGZus{}fld\PYG{+w}{ }has\PYG{+w}{ }the\PYG{+w}{ }value\PYG{+w}{ }\PYG{l+m}{0}.0000000000000000\PYG{+w}{ }at\PYG{+w}{ }index/indices\PYG{+w}{ }\PYG{l+m}{6}\PYG{+w}{ }\PYG{l+m}{1}\PYG{+w}{ }\PYG{k}{in}\PYG{+w}{ }module\PYG{+w}{ }\PYG{l+s+s1}{\PYGZsq{}main\PYGZsq{}},\PYG{+w}{ }region\PYG{+w}{ }\PYG{l+s+s1}{\PYGZsq{}init\PYGZsq{}},\PYG{+w}{ }which\PYG{+w}{ }is\PYG{+w}{ }not\PYG{+w}{ }between\PYG{+w}{ }\PYG{l+s+s1}{\PYGZsq{}2.0000000000000000\PYGZsq{}}\PYG{+w}{ }and\PYG{+w}{ }\PYG{l+s+s1}{\PYGZsq{}3.0000000000000000\PYGZsq{}}.
...
PSyData:\PYG{+w}{ }Variable\PYG{+w}{ }b\PYGZus{}fld\PYG{+w}{ }has\PYG{+w}{ }the\PYG{+w}{ }value\PYG{+w}{ }\PYG{l+m}{0}.0000000000000000\PYG{+w}{ }at\PYG{+w}{ }index/indices\PYG{+w}{ }\PYG{l+m}{6}\PYG{+w}{ }\PYG{l+m}{1}\PYG{+w}{ }\PYG{k}{in}\PYG{+w}{ }module\PYG{+w}{ }\PYG{l+s+s1}{\PYGZsq{}main\PYGZsq{}},\PYG{+w}{ }region\PYG{+w}{ }\PYG{l+s+s1}{\PYGZsq{}update\PYGZsq{}},\PYG{+w}{ }which\PYG{+w}{ }is\PYG{+w}{ }not\PYG{+w}{ }between\PYG{+w}{ }\PYG{l+s+s1}{\PYGZsq{}2.0000000000000000\PYGZsq{}}\PYG{+w}{ }and\PYG{+w}{ }\PYG{l+s+s1}{\PYGZsq{}3.0000000000000000\PYGZsq{}}.
\end{sphinxVerbatim}

\sphinxAtStartPar
Notice that now a warning is created for both kernels: \sphinxcode{\sphinxupquote{init}} and \sphinxcode{\sphinxupquote{update}}.

\sphinxAtStartPar
Support for checking arbitrary Fortran code is tracked as issue \#2741.


\subsection{Example 6: PSy\sphinxhyphen{}layer Code Creation using PSyIR}
\label{\detokenize{tutorials_and_examples:example-6-psy-layer-code-creation-using-psyir}}
\sphinxAtStartPar
This example informs the development of the code generation of PSy\sphinxhyphen{}layer
code using the PSyIR language backends.


\section{LFRic Examples}
\label{\detokenize{tutorials_and_examples:lfric-examples}}\label{\detokenize{tutorials_and_examples:examples-lfric}}
\sphinxAtStartPar
These examples illustrate the functionality of PSyclone for the LFRic
domain.


\subsection{Example 1: Basic Operation}
\label{\detokenize{tutorials_and_examples:example-1-basic-operation}}
\sphinxAtStartPar
Basic operation of PSyclone with an \sphinxcode{\sphinxupquote{invoke()}} containing two
kernels, one {\hyperref[\detokenize{dynamo0p3:lfric-kernel}]{\sphinxcrossref{\DUrole{std,std-ref}{user\sphinxhyphen{}supplied}}}}, the other a
{\hyperref[\detokenize{dynamo0p3:lfric-built-ins}]{\sphinxcrossref{\DUrole{std,std-ref}{Built\sphinxhyphen{}in}}}}. Code is generated both with and
without distributed\sphinxhyphen{}memory support. Also demonstrates the use of the
\sphinxcode{\sphinxupquote{\sphinxhyphen{}d}} flag to specify where to search for user\sphinxhyphen{}supplied kernel code
(see {\hyperref[\detokenize{psyclone_command:psyclone-command}]{\sphinxcrossref{\DUrole{std,std-ref}{The psyclone command}}}} section for more details).


\subsection{Example 2: Applying Transformations}
\label{\detokenize{tutorials_and_examples:example-2-applying-transformations}}
\sphinxAtStartPar
A more complex example showing the use of PSyclone
{\hyperref[\detokenize{dynamo0p3:lfric-api-transformations}]{\sphinxcrossref{\DUrole{std,std-ref}{transformations}}}} to
change the generated PSy\sphinxhyphen{}layer code. Provides examples of
kernel\sphinxhyphen{}inlining and loop\sphinxhyphen{}fusion transformations.


\subsection{Example 3: Distributed and Shared Memory}
\label{\detokenize{tutorials_and_examples:example-3-distributed-and-shared-memory}}
\sphinxAtStartPar
Shows the use of colouring and OpenMP for the Dynamo 0.3 API. Includes
multi\sphinxhyphen{}kernel, named invokes with both user\sphinxhyphen{}supplied and built\sphinxhyphen{}in
kernels. Also shows the use of \sphinxcode{\sphinxupquote{Wchi}} function space metadata for
coordinate fields in LFRic.


\subsection{Example 4: Multiple Built\sphinxhyphen{}ins, Named Invokes and Boundary Conditions}
\label{\detokenize{tutorials_and_examples:example-4-multiple-built-ins-named-invokes-and-boundary-conditions}}
\sphinxAtStartPar
Demonstrates the use of the special \sphinxcode{\sphinxupquote{enforce\_bc\_kernel}} which
PSyclone recognises as a boundary\sphinxhyphen{}condition kernel.


\subsection{Example 5: Stencils}
\label{\detokenize{tutorials_and_examples:example-5-stencils}}
\sphinxAtStartPar
Example of kernels which require stencil information.


\subsection{Example 6: Reductions}
\label{\detokenize{tutorials_and_examples:example-6-reductions}}
\sphinxAtStartPar
Example of applying OpenMP to an InvokeSchedule containing kernels
that perform reduction operations. Two scripts are provided, one of
which demonstrates how to request that PSyclone generate code for a
reproducible OpenMP reduction. (The default OpenMP reduction is not
guaranteed to be reproducible from one run to the next on the same
number of threads.)


\subsection{Example 7: Column\sphinxhyphen{}Matrix Assembly Operators}
\label{\detokenize{tutorials_and_examples:example-7-column-matrix-assembly-operators}}
\sphinxAtStartPar
Example of kernels requiring Column\sphinxhyphen{}Matrix Assembly operators.


\subsection{Example 8: Redundant Computation}
\label{\detokenize{tutorials_and_examples:example-8-redundant-computation}}
\sphinxAtStartPar
Example of the use of the redundant\sphinxhyphen{}computation and move
transformations to eliminate and re\sphinxhyphen{}order halo exchanges.


\subsection{Example 9: Writing to Discontinuous Fields}
\label{\detokenize{tutorials_and_examples:example-9-writing-to-discontinuous-fields}}
\sphinxAtStartPar
Demonstrates the behaviour of PSyclone for kernels that read and write
quantities on horizontally\sphinxhyphen{}discontinuous function spaces. In addition,
this example demonstrates how to write a PSyclone transformation script
that only colours loops over continuous spaces.


\subsection{Example 10: Inter\sphinxhyphen{}grid Kernels}
\label{\detokenize{tutorials_and_examples:example-10-inter-grid-kernels}}
\sphinxAtStartPar
Demonstrates the use of “inter\sphinxhyphen{}grid” kernels that prolong or restrict
fields (map between grids of different resolutions), as well as the
use of \sphinxcode{\sphinxupquote{ANY\_DISCONTINUOUS\_SPACE}} function space metadata.


\subsection{Example 11: Asynchronous Halo Exchanges}
\label{\detokenize{tutorials_and_examples:example-11-asynchronous-halo-exchanges}}
\sphinxAtStartPar
Example of the use of transformations to introduce redundant computation,
split synchronous halo exchanges into asynchronous exchanges (start and
stop) and move the starts of those exchanges in order to overlap them
with computation.


\subsection{Example 12: Code Extraction}
\label{\detokenize{tutorials_and_examples:example-12-code-extraction}}
\sphinxAtStartPar
Example of applying code extraction to Nodes in an Invoke Schedule:

\begin{sphinxVerbatim}[commandchars=\\\{\}]
\PYGZgt{}\PYG{+w}{ }psyclone\PYG{+w}{ }\PYGZhy{}nodm\PYG{+w}{ }\PYGZhy{}s\PYG{+w}{ }./extract\PYGZus{}nodes.py\PYG{+w}{ }\PYG{l+s+se}{\PYGZbs{}}
\PYG{+w}{    }gw\PYGZus{}mixed\PYGZus{}schur\PYGZus{}preconditioner\PYGZus{}alg\PYGZus{}mod.x90
\end{sphinxVerbatim}

\sphinxAtStartPar
or to a Kernel in an Invoke after applying transformations:

\begin{sphinxVerbatim}[commandchars=\\\{\}]
\PYGZgt{}\PYG{+w}{ }psyclone\PYG{+w}{ }\PYGZhy{}nodm\PYG{+w}{ }\PYGZhy{}s\PYG{+w}{ }./extract\PYGZus{}kernel\PYGZus{}with\PYGZus{}transformations.py\PYG{+w}{ }\PYG{l+s+se}{\PYGZbs{}}
\PYG{+w}{    }gw\PYGZus{}mixed\PYGZus{}schur\PYGZus{}preconditioner\PYGZus{}alg\PYGZus{}mod.x90
\end{sphinxVerbatim}

\sphinxAtStartPar
For now it only inserts comments in appropriate locations while the
the full support for code extraction is being developed.

\sphinxAtStartPar
This example also contains a Python helper script \sphinxcode{\sphinxupquote{find\_kernel.py}}
which displays the names and Schedules of Invokes containing call(s)
to the specified Kernel:

\begin{sphinxVerbatim}[commandchars=\\\{\}]
\PYGZgt{}\PYG{+w}{ }python\PYG{+w}{ }find\PYGZus{}kernel.py
\end{sphinxVerbatim}


\subsection{Example 13 : Kernel Transformation}
\label{\detokenize{tutorials_and_examples:example-13-kernel-transformation}}
\sphinxAtStartPar
Demonstrates how an LFRic kernel can be transformed. The example
transformation makes Kernel values constant where appropriate. For
example, the number of levels is usually passed into a kernel by
argument but the transformation allows a particular value to be
specified which the transformation then sets as a parameter in the
kernel. Hard\sphinxhyphen{}coding values in a kernel helps the compiler to do a
better job when optimising the code.


\subsection{Example 14: OpenACC}
\label{\detokenize{tutorials_and_examples:example-14-openacc}}
\sphinxAtStartPar
Example of adding OpenACC directives in the LFRic API.
A single transformation script (\sphinxcode{\sphinxupquote{acc\_parallel.py}}) is provided
which demonstrates how to add OpenACC Kernels and Enter Data
directives to the PSy\sphinxhyphen{}layer. It supports distributed memory being
switched on by placing an OpenACC Kernels directive around each
(parallelisable) loop, rather than having one for the whole invoke.
This approach avoids having halo exchanges within an OpenACC Parallel
region. The script also uses {\hyperref[\detokenize{transformations:available-kernel-trans}]{\sphinxcrossref{\DUrole{std,std-ref}{ACCRoutineTrans}}}}
to transform the one user\sphinxhyphen{}supplied kernel through
the addition of an \sphinxcode{\sphinxupquote{!\$acc routine}} directive. This ensures that the
compiler builds a version suitable for execution on the accelerator (GPU).

\sphinxAtStartPar
This script is used by the supplied Makefile. The invocation of PSyclone
within that Makefile also specifies the \sphinxcode{\sphinxupquote{\sphinxhyphen{}\sphinxhyphen{}profile invokes}} option so that
each \sphinxcode{\sphinxupquote{invoke}} is enclosed within profiling calipers (by default the
‘template’ profiling library supplied with PSyclone is used at the link
stage). Compilation of the example using the NVIDIA compiler may be performed
by e.g.:

\begin{sphinxVerbatim}[commandchars=\\\{\}]
\PYGZgt{}\PYG{+w}{ }\PYG{n+nv}{F90}\PYG{o}{=}nvfortran\PYG{+w}{ }\PYG{n+nv}{F90FLAGS}\PYG{o}{=}\PYG{l+s+s2}{\PYGZdq{}\PYGZhy{}acc \PYGZhy{}Minfo=all\PYGZdq{}}\PYG{+w}{ }make\PYG{+w}{ }compile
\end{sphinxVerbatim}

\sphinxAtStartPar
Launching the resulting binary with \sphinxcode{\sphinxupquote{NV\_ACC\_NOTIFY}} set will show details
of the kernel launches and data transfers:

\begin{sphinxVerbatim}[commandchars=\\\{\}]
\PYGZgt{}\PYG{+w}{ }\PYG{n+nv}{NV\PYGZus{}ACC\PYGZus{}NOTIFY}\PYG{o}{=}\PYG{l+m}{3}\PYG{+w}{ }./example\PYGZus{}openacc
...
\PYG{+w}{  }Step\PYG{+w}{             }\PYG{l+m}{5}\PYG{+w}{ }:\PYG{+w}{ }\PYG{n+nv}{chksm}\PYG{+w}{ }\PYG{o}{=}\PYG{+w}{    }\PYG{l+m}{2}.1098315506694516E\PYGZhy{}004
\PYG{+w}{  }PreStart\PYG{+w}{ }called\PYG{+w}{ }\PYG{k}{for}\PYG{+w}{ }module\PYG{+w}{ }\PYG{l+s+s1}{\PYGZsq{}main\PYGZus{}psy\PYGZsq{}}\PYG{+w}{ }region\PYG{+w}{ }\PYG{l+s+s1}{\PYGZsq{}invoke\PYGZus{}2\PYGZhy{}setval\PYGZus{}c\PYGZhy{}r2\PYGZsq{}}
\PYG{+w}{ }upload\PYG{+w}{ }CUDA\PYG{+w}{ }data\PYG{+w}{  }\PYG{n+nv}{file}\PYG{o}{=}PSyclone/examples/lfric/eg14/main\PYGZus{}psy.f90\PYG{+w}{ }\PYG{k}{function}\PYG{o}{=}invoke\PYGZus{}2\PYG{+w}{ }\PYG{n+nv}{line}\PYG{o}{=}\PYG{l+m}{183}\PYG{+w}{ }\PYG{n+nv}{device}\PYG{o}{=}\PYG{l+m}{0}\PYG{+w}{ }\PYG{n+nv}{threadid}\PYG{o}{=}\PYG{l+m}{1}\PYG{+w}{ }\PYG{n+nv}{variable}\PYG{o}{=}.attach.\PYG{+w}{ }\PYG{n+nv}{bytes}\PYG{o}{=}\PYG{l+m}{144}
\PYG{+w}{ }upload\PYG{+w}{ }CUDA\PYG{+w}{ }data\PYG{+w}{  }\PYG{n+nv}{file}\PYG{o}{=}PSyclone/examples/lfric/eg14/main\PYGZus{}psy.f90\PYG{+w}{ }\PYG{k}{function}\PYG{o}{=}invoke\PYGZus{}2\PYG{+w}{ }\PYG{n+nv}{line}\PYG{o}{=}\PYG{l+m}{183}\PYG{+w}{ }\PYG{n+nv}{device}\PYG{o}{=}\PYG{l+m}{0}\PYG{+w}{ }\PYG{n+nv}{threadid}\PYG{o}{=}\PYG{l+m}{1}\PYG{+w}{ }\PYG{n+nv}{variable}\PYG{o}{=}.attach.\PYG{+w}{ }\PYG{n+nv}{bytes}\PYG{o}{=}\PYG{l+m}{144}
\PYG{+w}{ }launch\PYG{+w}{ }CUDA\PYG{+w}{ }kernel\PYG{+w}{  }\PYG{n+nv}{file}\PYG{o}{=}PSyclone/examples/lfric/eg14/main\PYGZus{}psy.f90\PYG{+w}{ }\PYG{k}{function}\PYG{o}{=}invoke\PYGZus{}2\PYG{+w}{ }\PYG{n+nv}{line}\PYG{o}{=}\PYG{l+m}{186}\PYG{+w}{ }\PYG{n+nv}{device}\PYG{o}{=}\PYG{l+m}{0}\PYG{+w}{ }\PYG{n+nv}{threadid}\PYG{o}{=}\PYG{l+m}{1}\PYG{+w}{ }\PYG{n+nv}{num\PYGZus{}gangs}\PYG{o}{=}\PYG{l+m}{5}\PYG{+w}{ }\PYG{n+nv}{num\PYGZus{}workers}\PYG{o}{=}\PYG{l+m}{1}\PYG{+w}{ }\PYG{n+nv}{vector\PYGZus{}length}\PYG{o}{=}\PYG{l+m}{128}\PYG{+w}{ }\PYG{n+nv}{grid}\PYG{o}{=}\PYG{l+m}{5}\PYG{+w}{ }\PYG{n+nv}{block}\PYG{o}{=}\PYG{l+m}{128}
\PYG{+w}{  }PostEnd\PYG{+w}{ }called\PYG{+w}{ }\PYG{k}{for}\PYG{+w}{ }module\PYG{+w}{ }\PYG{l+s+s1}{\PYGZsq{}main\PYGZus{}psy\PYGZsq{}}\PYG{+w}{ }region\PYG{+w}{ }\PYG{l+s+s1}{\PYGZsq{}invoke\PYGZus{}2\PYGZhy{}setval\PYGZus{}c\PYGZhy{}r2\PYGZsq{}}
\PYG{+w}{ }download\PYG{+w}{ }CUDA\PYG{+w}{ }data\PYG{+w}{  }\PYG{n+nv}{file}\PYG{o}{=}PSyclone/src/psyclone/tests/test\PYGZus{}files/dynamo0p3/infrastructure//field/field\PYGZus{}r64\PYGZus{}mod.f90\PYG{+w}{ }\PYG{k}{function}\PYG{o}{=}log\PYGZus{}minmax\PYG{+w}{ }\PYG{n+nv}{line}\PYG{o}{=}\PYG{l+m}{756}\PYG{+w}{ }\PYG{n+nv}{device}\PYG{o}{=}\PYG{l+m}{0}\PYG{+w}{ }\PYG{n+nv}{threadid}\PYG{o}{=}\PYG{l+m}{1}\PYG{+w}{ }\PYG{n+nv}{variable}\PYG{o}{=}self\PYGZpc{}data\PYG{o}{(}:\PYG{o}{)}\PYG{+w}{ }\PYG{n+nv}{bytes}\PYG{o}{=}\PYG{l+m}{4312}
\PYG{+w}{ }\PYG{l+m}{20230807214504}.374+0100:INFO\PYG{+w}{ }:\PYG{+w}{ }Min/max\PYG{+w}{ }minmax\PYG{+w}{ }of\PYG{+w}{ }\PYG{n+nv}{field1}\PYG{+w}{ }\PYG{o}{=}\PYG{+w}{   }\PYG{l+m}{0}.30084014E+00\PYG{+w}{  }\PYG{l+m}{0}.17067212E+01
...
\end{sphinxVerbatim}

\sphinxAtStartPar
However, performance will be very poor as, with the limited
optimisations and directives currently applied, the NVIDIA compiler
refuses to run the user\sphinxhyphen{}supplied kernel in parallel.


\subsection{Example 15: CPU Optimisation of Matvec}
\label{\detokenize{tutorials_and_examples:example-15-cpu-optimisation-of-matvec}}
\sphinxAtStartPar
Example of optimising the LFRic matvec kernel for CPUs. This is work
in progress with the idea being that PSyclone transformations will be
able to reproduce hand\sphinxhyphen{}optimised code.

\sphinxAtStartPar
There is one script which, when run:

\begin{sphinxVerbatim}[commandchars=\\\{\}]
\PYGZgt{}\PYG{+w}{ }psyclone\PYG{+w}{ }./matvec\PYGZus{}opt.py\PYG{+w}{ }../code/gw\PYGZus{}mixed\PYGZus{}schur\PYGZus{}preconditioner\PYGZus{}alg\PYGZus{}mod.x90
\end{sphinxVerbatim}

\sphinxAtStartPar
will print out the modified matvec kernel code. At the moment no
transformations are included (as they are work\sphinxhyphen{}in\sphinxhyphen{}progress) so the
code that is output is the same as the original (but looks different
as it has been translated to PSyIR and then output by the PSyIR
Fortran back\sphinxhyphen{}end).


\subsection{Example 16: Generating LFRic Code Using LFRic\sphinxhyphen{}specific PSyIR}
\label{\detokenize{tutorials_and_examples:example-16-generating-lfric-code-using-lfric-specific-psyir}}
\sphinxAtStartPar
This example shows how LFRic\sphinxhyphen{}specific PSyIR can be used to create
LFRic kernel code. There is one Python script provided which when run:

\begin{sphinxVerbatim}[commandchars=\\\{\}]
\PYGZgt{}\PYG{+w}{ }python\PYG{+w}{ }create.py
\end{sphinxVerbatim}

\sphinxAtStartPar
will print out generated LFRic kernel code. The script makes use of
LFRic\sphinxhyphen{}specific data symbols to simplify code generation.


\subsection{Example 17: Runnable Simplified Examples}
\label{\detokenize{tutorials_and_examples:example-17-runnable-simplified-examples}}
\sphinxAtStartPar
This directory contains three simplified LFRic examples that can be
compiled and executed \sphinxhyphen{} of course, a suitable Fortran compiler is
required. The examples are using a subset of the LFRic infrastructure
library, which is contained in PSyclone and which has been slightly
modified to make it easier to create stand\sphinxhyphen{}alone, non\sphinxhyphen{}MPI LFRic codes.


\subsubsection{Example 17.1: A Simple Runnable Example}
\label{\detokenize{tutorials_and_examples:example-17-1-a-simple-runnable-example}}
\sphinxAtStartPar
The subdirectory \sphinxcode{\sphinxupquote{full\_example}} contains a very simple example code
that uses PSyclone to process two invokes. It uses unit\sphinxhyphen{}testing
code from various classes to create the required data structures like
initial grid etc. The code can be compiled with \sphinxcode{\sphinxupquote{make compile}}, and
the binary executed with either \sphinxcode{\sphinxupquote{make run}} or \sphinxcode{\sphinxupquote{./example}}.


\subsubsection{Example 17.2: A Simple Runnable Example With NetCDF}
\label{\detokenize{tutorials_and_examples:example-17-2-a-simple-runnable-example-with-netcdf}}
\sphinxAtStartPar
The subdirectory \sphinxcode{\sphinxupquote{full\_example\_netcdf}} contains code very similar
to the previous example, but uses NetCDF to read the initial grid
from the NetCDF file \sphinxcode{\sphinxupquote{mesh\_BiP128x16\sphinxhyphen{}400x100.nc}}.
Installation of NetCDF is described in
\sphinxhref{https://github.com/stfc/PSyclone/tree/master/tutorial/practicals\#user-content-netcdf-library-lfric-examples}{the hands\sphinxhyphen{}on practicals documentation}.
The code can be compiled with \sphinxcode{\sphinxupquote{make compile}}, and
the binary executed with either \sphinxcode{\sphinxupquote{make run}} or \sphinxcode{\sphinxupquote{./example}}.


\subsubsection{Example 17.3: Kernel Data Extraction}
\label{\detokenize{tutorials_and_examples:example-17-3-kernel-data-extraction}}
\sphinxAtStartPar
The example in the subdirectory \sphinxcode{\sphinxupquote{full\_example\_extract}} shows the
use of {\hyperref[\detokenize{psyke:psyke}]{\sphinxcrossref{\DUrole{std,std-ref}{kernel extraction}}}}. The code can be compiled with
\sphinxcode{\sphinxupquote{make compile}}, and the binary executed with either \sphinxcode{\sphinxupquote{make run}} or
\sphinxcode{\sphinxupquote{./extract.standalone}}. By default, it will be using
a stand\sphinxhyphen{}alone extraction library (see {\hyperref[\detokenize{psyke:extraction-libraries}]{\sphinxcrossref{\DUrole{std,std-ref}{Extraction Libraries}}}}).
If you want to use the NetCDF version, set the environment variable
\sphinxcode{\sphinxupquote{TYPE}} to be \sphinxcode{\sphinxupquote{netcdf}}:

\begin{sphinxVerbatim}[commandchars=\\\{\}]
\PYG{n+nv}{TYPE}\PYG{o}{=}netcdf\PYG{+w}{ }make\PYG{+w}{ }compile
\end{sphinxVerbatim}

\sphinxAtStartPar
This requires the installation of a NetCDF development environment
(see \sphinxhref{https://github.com/stfc/PSyclone/tree/master/tutorial/practicals\#user-content-netcdf-library-lfric-examples}{here}
for installing NetCDF). The binary will be called \sphinxcode{\sphinxupquote{extract.netcdf}},
and the output files will have the \sphinxcode{\sphinxupquote{.nc}} extension.

\sphinxAtStartPar
Running the compiled binary will create two Fortran binary files or
two NetCDF files if the NetCDF library was used. They contain
the input and output parameters for the two invokes in this example:

\begin{sphinxVerbatim}[commandchars=\\\{\}]
\PYG{n+nb}{cd}\PYG{+w}{ }full\PYGZus{}example\PYGZus{}extraction
\PYG{n+nv}{TYPE}\PYG{o}{=}netcdf\PYG{+w}{ }make\PYG{+w}{ }compile
./extract.netcdf
ncdump\PYG{+w}{ }./main\PYGZhy{}update.nc\PYG{+w}{ }\PYG{p}{|}\PYG{+w}{ }less
\end{sphinxVerbatim}


\subsection{Example 18: Special Accesses of Continuous Fields \sphinxhyphen{} Incrementing After Reading and Writing Before (Potentially) Reading}
\label{\detokenize{tutorials_and_examples:example-18-special-accesses-of-continuous-fields-incrementing-after-reading-and-writing-before-potentially-reading}}
\sphinxAtStartPar
Example containing one kernel with a \sphinxcode{\sphinxupquote{GH\_READINC}} access and one
with a \sphinxcode{\sphinxupquote{GH\_WRITE}} access, both for continuous fields. A kernel with
\sphinxcode{\sphinxupquote{GH\_READINC}} access first reads the field data and then increments
the field data. This contrasts with a \sphinxcode{\sphinxupquote{GH\_INC}} access which simply
increments the field data. As an increment is effectively a read
followed by a write, it may not be clear why we need to distinguish
between these cases. The reason for distinguishing is that the
\sphinxcode{\sphinxupquote{GH\_INC}} access is able to remove a halo exchange (or at least
reduce its depth by one) in certain circumstances, whereas a
\sphinxcode{\sphinxupquote{GH\_READINC}} is not able to take advantage of this optimisation.

\sphinxAtStartPar
A kernel with a \sphinxcode{\sphinxupquote{GH\_WRITE}} access for a continuous field must guarantee to
write the same value to a given shared DoF, independent of which cell
is being updated. As \sphinxhref{https://psyclone-dev.readthedocs.io/en/latest/APIs.html\#iterators-continuous}{\DUrole{xref,std,std-ref}{described}}
in the Developer Guide, this means that annexed DoFs are computed
correctly without the need to iterate into the L1 halo and thus can
remove the need for halo exchanges on those fields that are read.


\subsection{Example 19: Mixed Precision}
\label{\detokenize{tutorials_and_examples:example-19-mixed-precision}}
\sphinxAtStartPar
This example shows the use of the LFRic {\hyperref[\detokenize{dynamo0p3:lfric-mixed-precision}]{\sphinxcrossref{\DUrole{std,std-ref}{mixed\sphinxhyphen{}precision support}}}} to call a kernel with {\hyperref[\detokenize{dynamo0p3:lfric-mixed-precision-scalars}]{\sphinxcrossref{\DUrole{std,std-ref}{scalars}}}}, {\hyperref[\detokenize{dynamo0p3:lfric-mixed-precision-fields}]{\sphinxcrossref{\DUrole{std,std-ref}{fields}}}}
and {\hyperref[\detokenize{dynamo0p3:lfric-mixed-precision-lma-operators}]{\sphinxcrossref{\DUrole{std,std-ref}{operators}}}} of different
precision.


\subsection{Example 20: Algorithm Generation}
\label{\detokenize{tutorials_and_examples:example-20-algorithm-generation}}\label{\detokenize{tutorials_and_examples:lfric-alg-gen-example}}
\sphinxAtStartPar
Illustration of the use of the \sphinxcode{\sphinxupquote{psyclone\sphinxhyphen{}kern}} tool to create an
algorithm layer for a kernel. A makefile is provide that also
runs \sphinxcode{\sphinxupquote{psyclone}} to create an executable program from the generated
algorithm layer and original kernel code. To see the generated
algorithm layer run:

\begin{sphinxVerbatim}[commandchars=\\\{\}]
\PYG{n+nb}{cd}\PYG{+w}{ }eg20/
psyclone\PYGZhy{}kern\PYG{+w}{ }\PYGZhy{}gen\PYG{+w}{ }alg\PYG{+w}{ }../code/testkern\PYGZus{}mod.F90
\end{sphinxVerbatim}


\section{NEMO Examples}
\label{\detokenize{tutorials_and_examples:nemo-examples}}
\sphinxAtStartPar
These examples may all be found in the \sphinxcode{\sphinxupquote{examples/nemo}} directory.


\subsection{Example 1: OpenMP parallelisation of tra\_adv}
\label{\detokenize{tutorials_and_examples:example-1-openmp-parallelisation-of-tra-adv}}
\sphinxAtStartPar
Demonstrates the use of PSyclone to parallelise loops in a NEMO
tracer\sphinxhyphen{}advection benchmark using OpenMP for CPUs and for GPUs.


\subsection{Example 2: OpenMP parallelisation of traldf\_iso}
\label{\detokenize{tutorials_and_examples:example-2-openmp-parallelisation-of-traldf-iso}}
\sphinxAtStartPar
Demonstrates the use of PSyclone to parallelise in some NEMO
tracer\sphinxhyphen{}diffusion code using OpenMP for CPUs and for GPUs.


\subsection{Example 3: OpenACC parallelisation of tra\_adv}
\label{\detokenize{tutorials_and_examples:example-3-openacc-parallelisation-of-tra-adv}}
\sphinxAtStartPar
Demonstrates the introduction of simple OpenACC parallelisation (using the
\sphinxcode{\sphinxupquote{data}} and \sphinxcode{\sphinxupquote{kernels}} directives) for a NEMO tracer\sphinxhyphen{}advection benchmark.


\subsection{Example 4: Transforming Fortran code to the SIR}
\label{\detokenize{tutorials_and_examples:example-4-transforming-fortran-code-to-the-sir}}\label{\detokenize{tutorials_and_examples:nemo-eg4-sir}}
\sphinxAtStartPar
Demonstrates that simple Fortran code can be transformed to the Stencil
Intermediate Representation (SIR). The SIR is the front\sphinxhyphen{}end language to DAWN
(\sphinxurl{https://github.com/MeteoSwiss-APN/dawn}), a tool which generates
optimised cuda, or gridtools code. Thus various simple Fortran
examples and the computational part of the tracer\sphinxhyphen{}advection benchmark
can be transformed to optimised cuda and/or gridtools code by using
PSyclone and then DAWN.


\subsection{Example 5: Kernel Data Extraction}
\label{\detokenize{tutorials_and_examples:example-5-kernel-data-extraction}}
\sphinxAtStartPar
This example shows the use of kernel data extraction in PSyclone for
generic Fortran code. It instruments each kernel in the NEMO tracer\sphinxhyphen{}advection
benchmark with the PSyData\sphinxhyphen{}based kernel extraction code. Detailed
compilation instructions are in the \sphinxcode{\sphinxupquote{README.md}} file, including how
to switch from using the stand\sphinxhyphen{}alone extraction library to the NetCDF\sphinxhyphen{}based
one (see {\hyperref[\detokenize{psyke:extraction-libraries}]{\sphinxcrossref{\DUrole{std,std-ref}{Extraction Libraries}}}} for details).


\subsection{Example 6: Read\sphinxhyphen{}only Verification}
\label{\detokenize{tutorials_and_examples:example-6-read-only-verification}}
\sphinxAtStartPar
This example shows the use of read\sphinxhyphen{}only verification with PSyclone for
generic Fortran code. It instruments each kernel in a small Fortran
program with the PSyData\sphinxhyphen{}based read\sphinxhyphen{}only verification code. Detailed
compilation instructions are in the \sphinxcode{\sphinxupquote{README.md}} file.


\subsection{Scripts}
\label{\detokenize{tutorials_and_examples:scripts}}
\sphinxAtStartPar
This contains examples of two different scripts that aid the use of PSyclone
with the full NEMO model. The first, \sphinxtitleref{process\_nemo.py} is a simple wrapper
script that allows a user to control which source files are transformed, which
only have profiling instrumentation added and which are ignored altogether.
The second, \sphinxtitleref{kernels\_trans.py} is a PSyclone transformation script which
adds the largest possible OpenACC Kernels regions to the code being processed.

\sphinxAtStartPar
For more details see the \sphinxcode{\sphinxupquote{examples/nemo/README.md}} file.

\sphinxAtStartPar
Note that these scripts are here to support the ongoing development of PSyclone
to transform the NEMO source. They are \sphinxstyleemphasis{not} intended as ‘turn\sphinxhyphen{}key’ solutions
but as a starting point.

\sphinxstepscope


\chapter{Libraries}
\label{\detokenize{libraries:libraries}}\label{\detokenize{libraries:id1}}\label{\detokenize{libraries::doc}}
\sphinxAtStartPar
PSyclone provides {\hyperref[\detokenize{psy_data:psy-data}]{\sphinxcrossref{\DUrole{std,std-ref}{PSyData\sphinxhyphen{}API\sphinxhyphen{}based}}}} wrappers to
various external libraries. These wrapper libraries provide PSyclone
transformations that insert callbacks to an external library at runtime.
The callbacks then allow third\sphinxhyphen{}party libraries to access data structures
at specified locations in the code for different purposes, such as
profiling and extraction of argument values.

\sphinxAtStartPar
These wrapper libraries can be found under the \sphinxcode{\sphinxupquote{lib}} directory in the Git
repository. If you have installed PSyclone using \sphinxcode{\sphinxupquote{pip}} then the libraries
may be found in \sphinxcode{\sphinxupquote{share/psyclone/lib}} in PSyclone {\hyperref[\detokenize{getting_going:getting-going-install-loc}]{\sphinxcrossref{\DUrole{std,std-ref}{Installation location}}}}.

\begin{sphinxadmonition}{note}{Note:}
\sphinxAtStartPar
If working with wrapper libraries from a PSyclone installation,
it is advisable to copy the entire \sphinxcode{\sphinxupquote{lib}} directory to some
convenient location before building and using them. The provided
\sphinxcode{\sphinxupquote{Makefile}}s support the options to specify paths to the
libraries and their dependencies, see {\hyperref[\detokenize{libraries:libraries-compilation}]{\sphinxcrossref{\DUrole{std,std-ref}{compilation}}}} for more information.
\end{sphinxadmonition}


\section{Available libraries}
\label{\detokenize{libraries:available-libraries}}
\sphinxAtStartPar
An overview of the currently available functionality is below. For details
of what each library does and how to build and use it please see the related
sections in the User Guide and the specific \sphinxcode{\sphinxupquote{README.md}} files in the
associated directories.


\subsection{Profiling}
\label{\detokenize{libraries:profiling}}
\sphinxAtStartPar
PSyclone provides wrapper libraries for some common performance profiling
tools, such as dl\_timer, TAU, and Dr Hook. More information can be found in
the {\hyperref[\detokenize{profiling:profiling}]{\sphinxcrossref{\DUrole{std,std-ref}{Profiling}}}} section.

\sphinxAtStartPar
Profiling libraries are located in the \sphinxcode{\sphinxupquote{lib/profiling}} \sphinxhref{https://github.com/stfc/PSyclone/tree/master/lib/profiling}{directory}.
For detailed instructions on how to build and use them please refer
to their specific \sphinxcode{\sphinxupquote{README.md}} documentation.


\subsection{Kernel Data Extraction}
\label{\detokenize{libraries:kernel-data-extraction}}
\sphinxAtStartPar
These libraries enable PSyclone to add callbacks that provide access
to all input variables before, and output variables after a kernel
invocation. More information can be found in the
{\hyperref[\detokenize{psyke:psyke}]{\sphinxcrossref{\DUrole{std,std-ref}{PSy Kernel Extractor (PSyKE)}}}} section.

\sphinxAtStartPar
Example libraries that extract input and output data into a NetCDF file
for {\hyperref[\detokenize{dynamo0p3:lfric-api}]{\sphinxcrossref{\DUrole{std,std-ref}{LFRic}}}} and
{\hyperref[\detokenize{gocean1p0:gocean-api}]{\sphinxcrossref{\DUrole{std,std-ref}{GOcean}}}} APIs are included with PSyclone in the
\sphinxcode{\sphinxupquote{lib/extract/netcdf}} \sphinxhref{https://github.com/stfc/PSyclone/tree/master/lib/extract/netcdf}{directory}.
For detailed instructions on how to build and use these libraries
please refer to their specific \sphinxcode{\sphinxupquote{README.md}} documentation.


\subsection{Access Verification}
\label{\detokenize{libraries:access-verification}}
\sphinxAtStartPar
Read\sphinxhyphen{}only libraries check that a field declared as read\sphinxhyphen{}only is not
modified during a kernel call. More information can be found in the
{\hyperref[\detokenize{psy_data:psydata-read-verification}]{\sphinxcrossref{\DUrole{std,std-ref}{Read\sphinxhyphen{}Only Verification}}}} section.

\sphinxAtStartPar
The libraries for {\hyperref[\detokenize{dynamo0p3:lfric-api}]{\sphinxcrossref{\DUrole{std,std-ref}{LFRic}}}},
{\hyperref[\detokenize{gocean1p0:gocean-api}]{\sphinxcrossref{\DUrole{std,std-ref}{GOcean}}}} APIs and generic Fortran code
are included with PSyclone in
the \sphinxcode{\sphinxupquote{lib/read\_only}} \sphinxhref{https://github.com/stfc/PSyclone/tree/master/lib/read\_only}{directory}.
For detailed instructions on how to build and use these libraries
please refer to their specific \sphinxcode{\sphinxupquote{README.md}} documentation.


\subsection{Value Range Check}
\label{\detokenize{libraries:value-range-check}}
\sphinxAtStartPar
These libraries can test if user\sphinxhyphen{}defined variables are within a
specified range. Additionally, they also verify that they are
not \sphinxcode{\sphinxupquote{NaN}} or infinite.  More information can be
found in the {\hyperref[\detokenize{psy_data:psydata-value-range-check}]{\sphinxcrossref{\DUrole{std,std-ref}{Value Range Check}}}} section.

\sphinxAtStartPar
The libraries for {\hyperref[\detokenize{dynamo0p3:lfric-api}]{\sphinxcrossref{\DUrole{std,std-ref}{LFRic}}}} and
{\hyperref[\detokenize{gocean1p0:gocean-api}]{\sphinxcrossref{\DUrole{std,std-ref}{GOcean}}}} APIs are included with PSyclone in
the \sphinxcode{\sphinxupquote{lib/value\_range\_check}} \sphinxhref{https://github.com/stfc/PSyclone/tree/master/lib/value\_range\_check}{directory}.
For detailed instructions on how to build and use these libraries
please refer to their specific \sphinxcode{\sphinxupquote{README.md}} documentation.


\section{Dependencies}
\label{\detokenize{libraries:dependencies}}\label{\detokenize{libraries:libraries-dependencies}}
\sphinxAtStartPar
Building and using the wrapper libraries requires that PSyclone be installed
on the host system, see section {\hyperref[\detokenize{getting_going:getting-going}]{\sphinxcrossref{\DUrole{std,std-ref}{Getting Going}}}}. A
Fortran compiler (e.g. Gnu Fortran compiler, \sphinxcode{\sphinxupquote{gfortran}}, is free and easily
installed) and Gnu Make are also required.

\sphinxAtStartPar
The majority of wrapper libraries use \sphinxhref{https://pypi.org/project/Jinja/}{Jinja} templates to create PSyData\sphinxhyphen{}derived
classes (please refer to \sphinxhref{https://psyclone-dev.readthedocs.io/en/latest/psy\_data.html\#psy-data}{psy\_data} and \sphinxhref{https://psyclone-dev.readthedocs.io/en/latest/psy\_data.html\#jinja}{Jinja Support in the Base Class}
for full details about the PSyData API).

\sphinxAtStartPar
Compilation of \sphinxcode{\sphinxupquote{extract}}, \sphinxcode{\sphinxupquote{value\_range\_check}}, \sphinxcode{\sphinxupquote{read\_only}} and some of the
profiling wrapper libraries depends on infrastructure libraries relevant
to the API they are used for. The {\hyperref[\detokenize{dynamo0p3:lfric-api}]{\sphinxcrossref{\DUrole{std,std-ref}{LFRic API}}}} uses the
LFRic infrastructure and {\hyperref[\detokenize{gocean1p0:gocean-api}]{\sphinxcrossref{\DUrole{std,std-ref}{GOcean}}}} uses the
dl\_esm\_inf library. The LFRic infrastructure can be obtained from the
LFRic \sphinxhref{https://code.metoffice.gov.uk/trac/lfric/browser}{code repository},
however this requires access to the \sphinxhref{https://code.metoffice.gov.uk/trac/home}{Met Office Science Repository Service
(MOSRS)}. A useful contact for
LFRic\sphinxhyphen{}related questions (including access to MOSRS) is the \sphinxhref{mailto:lfric@cmpd1.metoffice.gov.uk}{“lfric” mailing
list} which gathers the Met Office and
external LFRic developers and users. The dl\_esm\_inf library is freely
available and can be downloaded from \sphinxurl{https://github.com/stfc/dl\_esm\_inf}.

\sphinxAtStartPar
Some libraries require NetCDF for compilation. Installation of NetCDF is
described in details in the \sphinxhref{https://github.com/stfc/PSyclone/tree/master/tutorial/practicals\#user-content-netcdf-library-lfric-examples}{hands\sphinxhyphen{}on practicals documentation}.

\sphinxAtStartPar
Profiling wrapper libraries that depend on external tools (e.g. dl\_timer)
require these tools be installed and configured beforehand.


\section{Compilation}
\label{\detokenize{libraries:compilation}}\label{\detokenize{libraries:libraries-compilation}}
\sphinxAtStartPar
Each library is compiled with \sphinxcode{\sphinxupquote{make}} using the provided \sphinxcode{\sphinxupquote{Makefile}} that
has configurable options for compiler flags and locations of dependencies.

\sphinxAtStartPar
As in case of {\hyperref[\detokenize{tutorials_and_examples:examples-compilation}]{\sphinxcrossref{\DUrole{std,std-ref}{examples}}}}, \sphinxcode{\sphinxupquote{F90}} and
\sphinxcode{\sphinxupquote{F90FLAGS}} specify the compiler and compilation flags to use. The default
value for \sphinxcode{\sphinxupquote{F90}} is \sphinxcode{\sphinxupquote{gfortran}}.

\sphinxAtStartPar
Locations of the top\sphinxhyphen{}level \sphinxcode{\sphinxupquote{lib}} directory and the required Jinja templates
are specified with the \sphinxcode{\sphinxupquote{PSYDATA\_LIB\_DIR}} and \sphinxcode{\sphinxupquote{LIB\_TMPLT\_DIR}} variables.
For testing purposes their default values are set to relative paths to the
respective directories in the PSyclone repository.

\sphinxAtStartPar
The locations of the infrastructure libraries for LFRic and GOcean
applications can be configured with the variables \sphinxcode{\sphinxupquote{LFRIC\_INF\_DIR}} and
\sphinxcode{\sphinxupquote{GOCEAN\_INF\_DIR}}, respectively. Their default values are set to relative
paths to the locations of these libraries in the PSyclone repository. The
dl\_esm\_inf library is provided as a Git submodule of the PSyclone
project (see \sphinxhref{https://psyclone-dev.readthedocs.io/en/latest/working\_practises.html\#dev-installation}{Installation} in the Developers’ Guide
for details on working with submodules) and a pared\sphinxhyphen{}down version of LFRic
infrastructure is also available in the PSyclone repository (please refer
to the \sphinxcode{\sphinxupquote{README.md}} documentation of relevant wrapper libraries). However,
the infrastructure libraries are not available in a PSyclone installation
and they need to be downloaded separately, see {\hyperref[\detokenize{libraries:libraries-dependencies}]{\sphinxcrossref{\DUrole{std,std-ref}{Dependencies}}}} for more information. In this case
\sphinxcode{\sphinxupquote{LFRIC\_INF\_DIR}} and \sphinxcode{\sphinxupquote{GOCEAN\_INF\_DIR}} \sphinxstylestrong{must be set} to the exact paths
to where the respective infrastructure source can be found. For instance,

\begin{sphinxVerbatim}[commandchars=\\\{\}]
\PYG{n+nv}{GOCEAN\PYGZus{}INF\PYGZus{}DIR}\PYG{o}{=}\PYG{n+nv}{\PYGZdl{}HOME}/dl\PYGZus{}esm\PYGZus{}inf/finite\PYGZus{}difference\PYG{+w}{ }make
\end{sphinxVerbatim}

\sphinxAtStartPar
Profiling wrapper libraries that depend on external tools have specific
variables that configure paths to where these libraries are located in a
user environment.

\sphinxAtStartPar
For more information on how to build and configure a specific library
please refer to its \sphinxcode{\sphinxupquote{README.md}} documentation.

\sphinxAtStartPar
Similar to compilation of the {\hyperref[\detokenize{tutorials_and_examples:examples-compilation}]{\sphinxcrossref{\DUrole{std,std-ref}{examples}}}}, the
compiled library can be removed by running \sphinxcode{\sphinxupquote{make clean}}. There is also
the \sphinxcode{\sphinxupquote{allclean}} target that removes the compiled wrapper library as well
as the compiled infrastructure library that the wrapper may depend on.

\sphinxAtStartPar
The compilation of wrapper libraries was tested with the Gnu and Intel
Fortran compilers, see {\hyperref[\detokenize{tutorials_and_examples:supported-compilers}]{\sphinxcrossref{\DUrole{std,std-ref}{here}}}} for the full list.
Please let the PSyclone developers know if you have problems using a
compiler that has been tested or if you are working with a different compiler.

\sphinxstepscope


\chapter{PSyIR: the PSyclone Intermediate Representation}
\label{\detokenize{psyir:psyir-the-psyclone-intermediate-representation}}\label{\detokenize{psyir:psyir-ug}}\label{\detokenize{psyir::doc}}
\sphinxAtStartPar
The PSyIR is at the heart of PSyclone, representing code for existing
code and PSyKAl DSLs (at both the PSy\sphinxhyphen{} and kernel\sphinxhyphen{}layer levels). A PSyIR
tree may be constructed from scratch (in Python) or by processing existing
source code using a frontend. Transformations act on the PSyIR and
ultimately the generated code is produced by one of the PSyIR’s
backends.


\section{PSyIR Nodes}
\label{\detokenize{psyir:psyir-nodes}}
\sphinxAtStartPar
The PSyIR consists of classes whose instances can be connected
together to form a tree which represent computation in a
syntax\sphinxhyphen{}independent way. These classes all inherit from the \sphinxcode{\sphinxupquote{Node}}
baseclass and, as a result, PSyIR instances are often referred to
collectively as ‘PSyIR nodes’.

\sphinxAtStartPar
At the present time PSyIR classes can be essentially split into two
types: language\sphinxhyphen{}level nodes, which are nodes that the PSyIR backends
support, and therefore they can be directly translated to code; and
higher\sphinxhyphen{}level nodes, which are additional nodes that each domain can
insert. These nodes must implement a \sphinxtitleref{lower\_to\_language\_level} method
in order to be converted to their equivalent representation using only
language\sphinxhyphen{}level nodes. This then permits code to be generated for them.

\sphinxAtStartPar
The rest of this document describes only the language\sphinxhyphen{}level nodes, but as
all nodes inherit from the same base classes, the methods described here
are applicable to all PSyIR nodes.


\section{Available language\sphinxhyphen{}level nodes}
\label{\detokenize{psyir:available-language-level-nodes}}\begin{itemize}
\item {} 
\sphinxAtStartPar
\sphinxhref{https://psyclone-ref.readthedocs.io/en/latest/autogenerated/psyclone.psyir.nodes.html\#psyclone.psyir.nodes.ArrayMember}{ArrayMember}

\item {} 
\sphinxAtStartPar
\sphinxhref{https://psyclone-ref.readthedocs.io/en/latest/autogenerated/psyclone.psyir.nodes.html\#psyclone.psyir.nodes.ArrayReference}{ArrayReference}

\item {} 
\sphinxAtStartPar
\sphinxhref{https://psyclone-ref.readthedocs.io/en/latest/autogenerated/psyclone.psyir.nodes.html\#psyclone.psyir.nodes.ArrayOfStructuresMember}{ArrayOfStructuresMember}

\item {} 
\sphinxAtStartPar
\sphinxhref{https://psyclone-ref.readthedocs.io/en/latest/autogenerated/psyclone.psyir.nodes.html\#psyclone.psyir.nodes.ArrayOfStructuresReference}{ArrayOfStructuresReference}

\item {} 
\sphinxAtStartPar
\sphinxhref{https://psyclone-ref.readthedocs.io/en/latest/autogenerated/psyclone.psyir.nodes.html\#psyclone.psyir.nodes.Assignment}{Assignment}

\item {} 
\sphinxAtStartPar
\sphinxhref{https://psyclone-ref.readthedocs.io/en/latest/autogenerated/psyclone.psyir.nodes.html\#psyclone.psyir.nodes.BinaryOperation}{BinaryOperation}

\item {} 
\sphinxAtStartPar
\sphinxhref{https://psyclone-ref.readthedocs.io/en/latest/autogenerated/psyclone.psyir.nodes.html\#psyclone.psyir.nodes.Call}{Call}

\item {} 
\sphinxAtStartPar
\sphinxhref{https://psyclone-ref.readthedocs.io/en/latest/autogenerated/psyclone.psyir.nodes.html\#psyclone.psyir.nodes.CodeBlock}{CodeBlock}

\item {} 
\sphinxAtStartPar
\sphinxhref{https://psyclone-ref.readthedocs.io/en/latest/autogenerated/psyclone.psyir.nodes.html\#psyclone.psyir.nodes.Container}{Container}

\item {} 
\sphinxAtStartPar
\sphinxhref{https://psyclone-ref.readthedocs.io/en/latest/autogenerated/psyclone.psyir.nodes.html\#psyclone.psyir.nodes.FileContainer}{FileContainer}

\item {} 
\sphinxAtStartPar
\sphinxhref{https://psyclone-ref.readthedocs.io/en/latest/autogenerated/psyclone.psyir.nodes.html\#psyclone.psyir.nodes.IfBlock}{IfBlock}

\item {} 
\sphinxAtStartPar
\sphinxhref{https://psyclone-ref.readthedocs.io/en/latest/autogenerated/psyclone.psyir.nodes.html\#psyclone.psyir.nodes.IntrinsicCall}{IntrinsicCall}

\item {} 
\sphinxAtStartPar
\sphinxhref{https://psyclone-ref.readthedocs.io/en/latest/autogenerated/psyclone.psyir.nodes.html\#psyclone.psyir.nodes.Literal}{Literal}

\item {} 
\sphinxAtStartPar
\sphinxhref{https://psyclone-ref.readthedocs.io/en/latest/autogenerated/psyclone.psyir.nodes.html\#psyclone.psyir.nodes.Loop}{Loop}

\item {} 
\sphinxAtStartPar
\sphinxhref{https://psyclone-ref.readthedocs.io/en/latest/autogenerated/psyclone.psyir.nodes.html\#psyclone.psyir.nodes.Member}{Member}

\item {} 
\sphinxAtStartPar
\sphinxhref{https://psyclone-ref.readthedocs.io/en/latest/autogenerated/psyclone.psyir.nodes.html\#psyclone.psyir.nodes.Node}{Node}

\item {} 
\sphinxAtStartPar
\sphinxhref{https://psyclone-ref.readthedocs.io/en/latest/autogenerated/psyclone.psyir.nodes.html\#psyclone.psyir.nodes.Range}{Range}

\item {} 
\sphinxAtStartPar
\sphinxhref{https://psyclone-ref.readthedocs.io/en/latest/autogenerated/psyclone.psyir.nodes.html\#psyclone.psyir.nodes.Reference}{Reference}

\item {} 
\sphinxAtStartPar
\sphinxhref{https://psyclone-ref.readthedocs.io/en/latest/autogenerated/psyclone.psyir.nodes.html\#psyclone.psyir.nodes.Return}{Return}

\item {} 
\sphinxAtStartPar
\sphinxhref{https://psyclone-ref.readthedocs.io/en/latest/autogenerated/psyclone.psyir.nodes.html\#psyclone.psyir.nodes.Routine}{Routine}

\item {} 
\sphinxAtStartPar
\sphinxhref{https://psyclone-ref.readthedocs.io/en/latest/autogenerated/psyclone.psyir.nodes.html\#psyclone.psyir.nodes.Schedule}{Schedule}

\item {} 
\sphinxAtStartPar
\sphinxhref{https://psyclone-ref.readthedocs.io/en/latest/autogenerated/psyclone.psyir.nodes.html\#psyclone.psyir.nodes.Statement}{Statement}

\item {} 
\sphinxAtStartPar
\sphinxhref{https://psyclone-ref.readthedocs.io/en/latest/autogenerated/psyclone.psyir.nodes.html\#psyclone.psyir.nodes.StructureMember}{StructureMember}

\item {} 
\sphinxAtStartPar
\sphinxhref{https://psyclone-ref.readthedocs.io/en/latest/autogenerated/psyclone.psyir.nodes.html\#psyclone.psyir.nodes.StructureReference}{StructureReference}

\item {} 
\sphinxAtStartPar
\sphinxhref{https://psyclone-ref.readthedocs.io/en/latest/autogenerated/psyclone.psyir.nodes.html\#psyclone.psyir.nodes.UnaryOperation}{UnaryOperation}

\item {} 
\sphinxAtStartPar
\sphinxhref{https://psyclone-ref.readthedocs.io/en/latest/autogenerated/psyclone.psyir.nodes.html\#psyclone.psyir.nodes.WhileLoop}{WhileLoop}

\end{itemize}


\section{Text Representation}
\label{\detokenize{psyir:text-representation}}
\sphinxAtStartPar
When developing a transformation script it is often necessary to examine
the structure of the PSyIR. All nodes in the PSyIR have the \sphinxcode{\sphinxupquote{view}} method
that provides a text\sphinxhyphen{}representation of that node and all of its descendants.
If the \sphinxcode{\sphinxupquote{termcolor}} package is installed (see {\hyperref[\detokenize{getting_going:getting-going}]{\sphinxcrossref{\DUrole{std,std-ref}{Getting Going}}}}) then
colour highlighting is used as part of the output string.
For instance, part of the Schedule constructed for the second NEMO
\sphinxhref{https://github.com/stfc/PSyclone/blob/master/examples/nemo/eg2/omp\_levels\_trans.py}{example} is rendered as:

\noindent\sphinxincludegraphics{{schedule_with_indices}.png}

\sphinxAtStartPar
Note that in this view, only those nodes which are children of
Schedules have their indices shown. This means that nodes representing
e.g. loop bounds or the conditional part of \sphinxcode{\sphinxupquote{if}} statements are not
indexed. For the example shown, the PSyIR node representing the
\sphinxcode{\sphinxupquote{if(l\_hst)}} code would be reached by
\sphinxcode{\sphinxupquote{schedule.children{[}14{]}.if\_body.children{[}1{]}}} or, using the shorthand
notation (see below), \sphinxcode{\sphinxupquote{schedule{[}14{]}.if\_body{[}1{]}}} where \sphinxcode{\sphinxupquote{schedule}} is
the overall parent Schedule node (omitted from the above image).

\sphinxAtStartPar
One problem with the \sphinxcode{\sphinxupquote{view}} method is that the output can become very
large for big ASTs and is not readable for users unfamiliar with the PSyIR.
An alternative to it is the \sphinxcode{\sphinxupquote{debug\_string}} method that generates a
text representation with Fortran\sphinxhyphen{}like syntax but on which the high abstraction
constructs have not yet been lowered to Fortran level and instead they will be
embedded as \sphinxtitleref{\textless{} node \textgreater{}} expressions.


\section{Tree Navigation}
\label{\detokenize{psyir:tree-navigation}}
\sphinxAtStartPar
Each PSyIR node provides several ways to navigate the AST. These can be
categorised as homogeneous naviation methods (available in all nodes), and
heterogenous or semantic navigation methods (different methods available
depending on the node type). The homogeneous methods must be used for generic
code navigation that should work regardless of its context. However, when
the context is known, we recommend using the semantic methods to increase
the code readability.

\sphinxAtStartPar
The homogeneous navigation methods are:
\begin{quote}
\index{children() (psyclone.psyir.nodes.Node method)@\spxentry{children()}\spxextra{psyclone.psyir.nodes.Node method}}

\begin{fulllineitems}
\phantomsection\label{\detokenize{psyir:psyclone.psyir.nodes.Node.children}}
\pysigstartsignatures
\pysiglinewithargsret{\sphinxcode{\sphinxupquote{Node.}}\sphinxbfcode{\sphinxupquote{children}}}{}{}
\pysigstopsignatures\begin{quote}\begin{description}
\sphinxlineitem{Returns}
\sphinxAtStartPar
the immediate children of this Node.

\sphinxlineitem{Return type}
\sphinxAtStartPar
List{[}\sphinxcode{\sphinxupquote{psyclone.psyir.nodes.Node}}{]}

\end{description}\end{quote}

\end{fulllineitems}

\index{siblings() (psyclone.psyir.nodes.Node method)@\spxentry{siblings()}\spxextra{psyclone.psyir.nodes.Node method}}

\begin{fulllineitems}
\phantomsection\label{\detokenize{psyir:psyclone.psyir.nodes.Node.siblings}}
\pysigstartsignatures
\pysiglinewithargsret{\sphinxcode{\sphinxupquote{Node.}}\sphinxbfcode{\sphinxupquote{siblings}}}{}{}
\pysigstopsignatures\begin{quote}\begin{description}
\sphinxlineitem{Returns}
\sphinxAtStartPar
list of sibling nodes, including self.

\sphinxlineitem{Return type}
\sphinxAtStartPar
List{[}\sphinxcode{\sphinxupquote{psyclone.psyir.nodes.Node}}{]}

\end{description}\end{quote}

\end{fulllineitems}

\index{parent() (psyclone.psyir.nodes.Node method)@\spxentry{parent()}\spxextra{psyclone.psyir.nodes.Node method}}

\begin{fulllineitems}
\phantomsection\label{\detokenize{psyir:psyclone.psyir.nodes.Node.parent}}
\pysigstartsignatures
\pysiglinewithargsret{\sphinxcode{\sphinxupquote{Node.}}\sphinxbfcode{\sphinxupquote{parent}}}{}{}
\pysigstopsignatures\begin{quote}\begin{description}
\sphinxlineitem{Returns}
\sphinxAtStartPar
the parent node.

\sphinxlineitem{Return type}
\sphinxAtStartPar
\sphinxcode{\sphinxupquote{psyclone.psyir.nodes.Node}} or NoneType

\end{description}\end{quote}

\end{fulllineitems}

\index{root() (psyclone.psyir.nodes.Node method)@\spxentry{root()}\spxextra{psyclone.psyir.nodes.Node method}}

\begin{fulllineitems}
\phantomsection\label{\detokenize{psyir:psyclone.psyir.nodes.Node.root}}
\pysigstartsignatures
\pysiglinewithargsret{\sphinxcode{\sphinxupquote{Node.}}\sphinxbfcode{\sphinxupquote{root}}}{}{}
\pysigstopsignatures\begin{quote}\begin{description}
\sphinxlineitem{Returns}
\sphinxAtStartPar
the root node of the PSyIR tree.

\sphinxlineitem{Return type}
\sphinxAtStartPar
\sphinxcode{\sphinxupquote{psyclone.psyir.nodes.Node}}

\end{description}\end{quote}

\end{fulllineitems}

\index{walk() (psyclone.psyir.nodes.Node method)@\spxentry{walk()}\spxextra{psyclone.psyir.nodes.Node method}}

\begin{fulllineitems}
\phantomsection\label{\detokenize{psyir:psyclone.psyir.nodes.Node.walk}}
\pysigstartsignatures
\pysiglinewithargsret{\sphinxcode{\sphinxupquote{Node.}}\sphinxbfcode{\sphinxupquote{walk}}}{}{}
\pysigstopsignatures
\sphinxAtStartPar
Recurse through the PSyIR tree and return all objects that are
an instance of ‘my\_type’, which is either a single class or a tuple
of classes. In the latter case all nodes are returned that are
instances of any classes in the tuple. The recursion into the tree
is stopped if an instance of ‘stop\_type’ (which is either a single
class or a tuple of classes) is found. This can be used to avoid
analysing e.g. inlined kernels, or as performance optimisation to
reduce the number of recursive calls. The recursion into the tree is
also stopped if the (optional) ‘depth’ level is reached.
\begin{quote}\begin{description}
\sphinxlineitem{Parameters}\begin{itemize}
\item {} 
\sphinxAtStartPar
\sphinxstyleliteralstrong{\sphinxupquote{my\_type}} (\sphinxhref{https://docs.python.org/3/library/functions.html\#type}{\sphinxstyleliteralemphasis{\sphinxupquote{type}}}\sphinxstyleliteralemphasis{\sphinxupquote{ | }}\sphinxstyleliteralemphasis{\sphinxupquote{Tuple}}\sphinxstyleliteralemphasis{\sphinxupquote{{[}}}\sphinxhref{https://docs.python.org/3/library/functions.html\#type}{\sphinxstyleliteralemphasis{\sphinxupquote{type}}}\sphinxstyleliteralemphasis{\sphinxupquote{, }}\sphinxstyleliteralemphasis{\sphinxupquote{...}}\sphinxstyleliteralemphasis{\sphinxupquote{{]}}}) \textendash{} the class(es) for which the instances are collected.

\item {} 
\sphinxAtStartPar
\sphinxstyleliteralstrong{\sphinxupquote{stop\_type}} (\sphinxstyleliteralemphasis{\sphinxupquote{Optional}}\sphinxstyleliteralemphasis{\sphinxupquote{{[}}}\sphinxhref{https://docs.python.org/3/library/functions.html\#type}{\sphinxstyleliteralemphasis{\sphinxupquote{type}}}\sphinxstyleliteralemphasis{\sphinxupquote{ | }}\sphinxstyleliteralemphasis{\sphinxupquote{Tuple}}\sphinxstyleliteralemphasis{\sphinxupquote{{[}}}\sphinxhref{https://docs.python.org/3/library/functions.html\#type}{\sphinxstyleliteralemphasis{\sphinxupquote{type}}}\sphinxstyleliteralemphasis{\sphinxupquote{, }}\sphinxstyleliteralemphasis{\sphinxupquote{...}}\sphinxstyleliteralemphasis{\sphinxupquote{{]}}}\sphinxstyleliteralemphasis{\sphinxupquote{{]}}}) \textendash{} class(es) at which recursion is halted (optional).

\item {} 
\sphinxAtStartPar
\sphinxstyleliteralstrong{\sphinxupquote{depth}} (\sphinxstyleliteralemphasis{\sphinxupquote{Optional}}\sphinxstyleliteralemphasis{\sphinxupquote{{[}}}\sphinxhref{https://docs.python.org/3/library/functions.html\#int}{\sphinxstyleliteralemphasis{\sphinxupquote{int}}}\sphinxstyleliteralemphasis{\sphinxupquote{{]}}}) \textendash{} the depth value the instances must have (optional).

\end{itemize}

\sphinxlineitem{Returns}
\sphinxAtStartPar
list with all nodes that are instances of my\_type                   starting at and including this node.

\sphinxlineitem{Return type}
\sphinxAtStartPar
List{[}\sphinxcode{\sphinxupquote{psyclone.psyir.nodes.Node}}{]}

\end{description}\end{quote}

\end{fulllineitems}

\index{get\_sibling\_lists() (psyclone.psyir.nodes.Node method)@\spxentry{get\_sibling\_lists()}\spxextra{psyclone.psyir.nodes.Node method}}

\begin{fulllineitems}
\phantomsection\label{\detokenize{psyir:psyclone.psyir.nodes.Node.get_sibling_lists}}
\pysigstartsignatures
\pysiglinewithargsret{\sphinxcode{\sphinxupquote{Node.}}\sphinxbfcode{\sphinxupquote{get\_sibling\_lists}}}{}{}
\pysigstopsignatures
\sphinxAtStartPar
Recurse through the PSyIR tree and return lists of Nodes that are
instances of ‘my\_type’ and are immediate siblings. Here ‘my\_type’ is
either a single class or a tuple of classes. In the latter case all
nodes are returned that are instances of any classes in the tuple. The
recursion into the tree is stopped if an instance of ‘stop\_type’ (which
is either a single class or a tuple of classes) is found.
\begin{quote}\begin{description}
\sphinxlineitem{Parameters}\begin{itemize}
\item {} 
\sphinxAtStartPar
\sphinxstyleliteralstrong{\sphinxupquote{my\_type}} (\sphinxhref{https://docs.python.org/3/library/functions.html\#type}{\sphinxstyleliteralemphasis{\sphinxupquote{type}}}\sphinxstyleliteralemphasis{\sphinxupquote{ | }}\sphinxstyleliteralemphasis{\sphinxupquote{Tuple}}\sphinxstyleliteralemphasis{\sphinxupquote{{[}}}\sphinxhref{https://docs.python.org/3/library/functions.html\#type}{\sphinxstyleliteralemphasis{\sphinxupquote{type}}}\sphinxstyleliteralemphasis{\sphinxupquote{, }}\sphinxstyleliteralemphasis{\sphinxupquote{...}}\sphinxstyleliteralemphasis{\sphinxupquote{{]}}}) \textendash{} the class(es) for which the instances are collected.

\item {} 
\sphinxAtStartPar
\sphinxstyleliteralstrong{\sphinxupquote{stop\_type}} (\sphinxstyleliteralemphasis{\sphinxupquote{Optional}}\sphinxstyleliteralemphasis{\sphinxupquote{{[}}}\sphinxhref{https://docs.python.org/3/library/functions.html\#type}{\sphinxstyleliteralemphasis{\sphinxupquote{type}}}\sphinxstyleliteralemphasis{\sphinxupquote{ | }}\sphinxstyleliteralemphasis{\sphinxupquote{Tuple}}\sphinxstyleliteralemphasis{\sphinxupquote{{[}}}\sphinxhref{https://docs.python.org/3/library/functions.html\#type}{\sphinxstyleliteralemphasis{\sphinxupquote{type}}}\sphinxstyleliteralemphasis{\sphinxupquote{, }}\sphinxstyleliteralemphasis{\sphinxupquote{...}}\sphinxstyleliteralemphasis{\sphinxupquote{{]}}}\sphinxstyleliteralemphasis{\sphinxupquote{{]}}}) \textendash{} class(es) at which recursion is halted (optional).

\end{itemize}

\sphinxlineitem{Returns}
\sphinxAtStartPar
list of lists, each of which containing nodes that are
instances of my\_type and are immediate siblings, starting at
and including this node.

\sphinxlineitem{Return type}
\sphinxAtStartPar
List{[}List{[}\sphinxcode{\sphinxupquote{psyclone.psyir.nodes.Node}}{]}{]}

\end{description}\end{quote}

\end{fulllineitems}

\index{ancestor() (psyclone.psyir.nodes.Node method)@\spxentry{ancestor()}\spxextra{psyclone.psyir.nodes.Node method}}

\begin{fulllineitems}
\phantomsection\label{\detokenize{psyir:psyclone.psyir.nodes.Node.ancestor}}
\pysigstartsignatures
\pysiglinewithargsret{\sphinxcode{\sphinxupquote{Node.}}\sphinxbfcode{\sphinxupquote{ancestor}}}{}{}
\pysigstopsignatures
\sphinxAtStartPar
Search back up the tree and check whether this node has an ancestor
that is an instance of the supplied type. If it does then we return
it otherwise we return None. An individual (or tuple of) (sub\sphinxhyphen{})
class(es) to ignore may be provided via the \sphinxtitleref{excluding} argument. If
\sphinxtitleref{include\_self} is True then the current node is included in the search.
If \sphinxtitleref{limit} is provided then the search ceases if/when the supplied
node is encountered.
If \sphinxtitleref{shared\_with} is provided, then the ancestor search will find an
ancestor of both this node and the node provided as \sphinxtitleref{shared\_with} if
such an ancestor exists.
\begin{quote}\begin{description}
\sphinxlineitem{Parameters}\begin{itemize}
\item {} 
\sphinxAtStartPar
\sphinxstyleliteralstrong{\sphinxupquote{my\_type}} (\sphinxhref{https://docs.python.org/3/library/functions.html\#type}{\sphinxstyleliteralemphasis{\sphinxupquote{type}}}\sphinxstyleliteralemphasis{\sphinxupquote{ | }}\sphinxstyleliteralemphasis{\sphinxupquote{Tuple}}\sphinxstyleliteralemphasis{\sphinxupquote{{[}}}\sphinxhref{https://docs.python.org/3/library/functions.html\#type}{\sphinxstyleliteralemphasis{\sphinxupquote{type}}}\sphinxstyleliteralemphasis{\sphinxupquote{, }}\sphinxstyleliteralemphasis{\sphinxupquote{...}}\sphinxstyleliteralemphasis{\sphinxupquote{{]}}}) \textendash{} class(es) to search for.

\item {} 
\sphinxAtStartPar
\sphinxstyleliteralstrong{\sphinxupquote{excluding}} (\sphinxstyleliteralemphasis{\sphinxupquote{Optional}}\sphinxstyleliteralemphasis{\sphinxupquote{{[}}}\sphinxhref{https://docs.python.org/3/library/functions.html\#type}{\sphinxstyleliteralemphasis{\sphinxupquote{type}}}\sphinxstyleliteralemphasis{\sphinxupquote{ | }}\sphinxstyleliteralemphasis{\sphinxupquote{Tuple}}\sphinxstyleliteralemphasis{\sphinxupquote{{[}}}\sphinxhref{https://docs.python.org/3/library/functions.html\#type}{\sphinxstyleliteralemphasis{\sphinxupquote{type}}}\sphinxstyleliteralemphasis{\sphinxupquote{, }}\sphinxstyleliteralemphasis{\sphinxupquote{...}}\sphinxstyleliteralemphasis{\sphinxupquote{{]}}}\sphinxstyleliteralemphasis{\sphinxupquote{{]}}}) \textendash{} (sub\sphinxhyphen{})class(es) to ignore or None.

\item {} 
\sphinxAtStartPar
\sphinxstyleliteralstrong{\sphinxupquote{include\_self}} (\sphinxhref{https://docs.python.org/3/library/functions.html\#bool}{\sphinxstyleliteralemphasis{\sphinxupquote{bool}}}) \textendash{} whether or not to include this node in the                                   search.

\item {} 
\sphinxAtStartPar
\sphinxstyleliteralstrong{\sphinxupquote{limit}} (Optional{[}\sphinxcode{\sphinxupquote{psyclone.psyir.nodes.Node}}{]}) \textendash{} an optional node at which to stop the search.

\item {} 
\sphinxAtStartPar
\sphinxstyleliteralstrong{\sphinxupquote{shared\_with}} (Optional{[}\sphinxcode{\sphinxupquote{psyclone.psyir.nodes.Node}}{]}) \textendash{} an optional node which must also have the
found node as an ancestor.

\end{itemize}

\sphinxlineitem{Returns}
\sphinxAtStartPar
First ancestor Node that is an instance of any of the                   requested classes or None if not found.

\sphinxlineitem{Return type}
\sphinxAtStartPar
Optional{[}\sphinxcode{\sphinxupquote{psyclone.psyir.nodes.Node}}{]}

\sphinxlineitem{Raises}\begin{itemize}
\item {} 
\sphinxAtStartPar
\sphinxhref{https://docs.python.org/3/library/exceptions.html\#TypeError}{\sphinxstyleliteralstrong{\sphinxupquote{TypeError}}} \textendash{} if \sphinxtitleref{excluding} is provided but is not a type or                            tuple of types.

\item {} 
\sphinxAtStartPar
\sphinxhref{https://docs.python.org/3/library/exceptions.html\#TypeError}{\sphinxstyleliteralstrong{\sphinxupquote{TypeError}}} \textendash{} if \sphinxtitleref{limit} is provided but is not an instance                            of Node.

\end{itemize}

\end{description}\end{quote}

\end{fulllineitems}

\index{scope() (psyclone.psyir.nodes.Node method)@\spxentry{scope()}\spxextra{psyclone.psyir.nodes.Node method}}

\begin{fulllineitems}
\phantomsection\label{\detokenize{psyir:psyclone.psyir.nodes.Node.scope}}
\pysigstartsignatures
\pysiglinewithargsret{\sphinxcode{\sphinxupquote{Node.}}\sphinxbfcode{\sphinxupquote{scope}}}{}{}
\pysigstopsignatures
\sphinxAtStartPar
Some nodes (e.g. Schedule and Container) allow symbols to be
scoped via an attached symbol table. This property returns the closest
ScopingNode node including self.
\begin{quote}\begin{description}
\sphinxlineitem{Returns}
\sphinxAtStartPar
the closest ancestor ScopingNode node.

\sphinxlineitem{Return type}
\sphinxAtStartPar
\sphinxcode{\sphinxupquote{psyclone.psyir.node.ScopingNode}}

\sphinxlineitem{Raises}
\sphinxAtStartPar
\sphinxstyleliteralstrong{\sphinxupquote{SymbolError}} \textendash{} if there is no ScopingNode ancestor.

\end{description}\end{quote}

\end{fulllineitems}

\index{path\_from() (psyclone.psyir.nodes.Node method)@\spxentry{path\_from()}\spxextra{psyclone.psyir.nodes.Node method}}

\begin{fulllineitems}
\phantomsection\label{\detokenize{psyir:psyclone.psyir.nodes.Node.path_from}}
\pysigstartsignatures
\pysiglinewithargsret{\sphinxcode{\sphinxupquote{Node.}}\sphinxbfcode{\sphinxupquote{path\_from}}}{}{}
\pysigstopsignatures
\sphinxAtStartPar
Find the path in the psyir tree between ancestor and node and
returns a list containing the path.

\sphinxAtStartPar
The result of this method can be used to find the node from its
ancestor for example by:

\begin{sphinxVerbatim}[commandchars=\\\{\}]
\PYG{g+gp}{\PYGZgt{}\PYGZgt{}\PYGZgt{} }\PYG{n}{index\PYGZus{}list} \PYG{o}{=} \PYG{n}{node}\PYG{o}{.}\PYG{n}{path\PYGZus{}from}\PYG{p}{(}\PYG{n}{ancestor}\PYG{p}{)}
\PYG{g+gp}{\PYGZgt{}\PYGZgt{}\PYGZgt{} }\PYG{n}{cursor} \PYG{o}{=} \PYG{n}{ancestor}
\PYG{g+gp}{\PYGZgt{}\PYGZgt{}\PYGZgt{} }\PYG{k}{for} \PYG{n}{index} \PYG{o+ow}{in} \PYG{n}{index\PYGZus{}list}\PYG{p}{:}
\PYG{g+gp}{\PYGZgt{}\PYGZgt{}\PYGZgt{} }   \PYG{n}{cursor} \PYG{o}{=} \PYG{n}{cursor}\PYG{o}{.}\PYG{n}{children}\PYG{p}{[}\PYG{n}{index}\PYG{p}{]}
\PYG{g+gp}{\PYGZgt{}\PYGZgt{}\PYGZgt{} }\PYG{k}{assert} \PYG{n}{cursor} \PYG{o+ow}{is} \PYG{n}{node}
\end{sphinxVerbatim}
\begin{quote}\begin{description}
\sphinxlineitem{Parameters}
\sphinxAtStartPar
\sphinxstyleliteralstrong{\sphinxupquote{ancestor}} (\sphinxcode{\sphinxupquote{psyclone.psyir.nodes.Node}}) \textendash{} an ancestor node of self to find the path from.

\sphinxlineitem{Raises}
\sphinxAtStartPar
\sphinxhref{https://docs.python.org/3/library/exceptions.html\#ValueError}{\sphinxstyleliteralstrong{\sphinxupquote{ValueError}}} \textendash{} if ancestor is not an ancestor of self.

\sphinxlineitem{Returns}
\sphinxAtStartPar
a list of child indices representing the path between
ancestor and self.

\sphinxlineitem{Return type}
\sphinxAtStartPar
List{[}\sphinxhref{https://docs.python.org/3/library/functions.html\#int}{int}{]}

\end{description}\end{quote}

\end{fulllineitems}

\end{quote}

\sphinxAtStartPar
In addition to the navigation methods, nodes also have homogeneous methods to
interrogate their location and surrounding nodes.
\begin{quote}
\index{immediately\_precedes() (psyclone.psyir.nodes.Node method)@\spxentry{immediately\_precedes()}\spxextra{psyclone.psyir.nodes.Node method}}

\begin{fulllineitems}
\phantomsection\label{\detokenize{psyir:psyclone.psyir.nodes.Node.immediately_precedes}}
\pysigstartsignatures
\pysiglinewithargsret{\sphinxcode{\sphinxupquote{Node.}}\sphinxbfcode{\sphinxupquote{immediately\_precedes}}}{}{}
\pysigstopsignatures\begin{quote}\begin{description}
\sphinxlineitem{Parameters}
\sphinxAtStartPar
\sphinxstyleliteralstrong{\sphinxupquote{node}} \textendash{} the node to compare it to.

\sphinxlineitem{Returns}
\sphinxAtStartPar
whether this node immediately precedes the given \sphinxtitleref{node}.

\sphinxlineitem{Return type}
\sphinxAtStartPar
\sphinxhref{https://docs.python.org/3/library/functions.html\#bool}{bool}

\end{description}\end{quote}

\end{fulllineitems}

\index{immediately\_follows() (psyclone.psyir.nodes.Node method)@\spxentry{immediately\_follows()}\spxextra{psyclone.psyir.nodes.Node method}}

\begin{fulllineitems}
\phantomsection\label{\detokenize{psyir:psyclone.psyir.nodes.Node.immediately_follows}}
\pysigstartsignatures
\pysiglinewithargsret{\sphinxcode{\sphinxupquote{Node.}}\sphinxbfcode{\sphinxupquote{immediately\_follows}}}{}{}
\pysigstopsignatures\begin{quote}\begin{description}
\sphinxlineitem{Parameters}
\sphinxAtStartPar
\sphinxstyleliteralstrong{\sphinxupquote{node}} \textendash{} the node to compare it to.

\sphinxlineitem{Returns}
\sphinxAtStartPar
whether this node immediately follows the given \sphinxtitleref{node}.

\sphinxlineitem{Return type}
\sphinxAtStartPar
\sphinxhref{https://docs.python.org/3/library/functions.html\#bool}{bool}

\end{description}\end{quote}

\end{fulllineitems}

\index{position() (psyclone.psyir.nodes.Node method)@\spxentry{position()}\spxextra{psyclone.psyir.nodes.Node method}}

\begin{fulllineitems}
\phantomsection\label{\detokenize{psyir:psyclone.psyir.nodes.Node.position}}
\pysigstartsignatures
\pysiglinewithargsret{\sphinxcode{\sphinxupquote{Node.}}\sphinxbfcode{\sphinxupquote{position}}}{}{}
\pysigstopsignatures
\sphinxAtStartPar
Find a Node’s position relative to its parent Node (starting
with 0 if it does not have a parent).
\begin{quote}\begin{description}
\sphinxlineitem{Returns}
\sphinxAtStartPar
relative position of a Node to its parent

\sphinxlineitem{Return type}
\sphinxAtStartPar
\sphinxhref{https://docs.python.org/3/library/functions.html\#int}{int}

\end{description}\end{quote}

\end{fulllineitems}

\index{abs\_position() (psyclone.psyir.nodes.Node method)@\spxentry{abs\_position()}\spxextra{psyclone.psyir.nodes.Node method}}

\begin{fulllineitems}
\phantomsection\label{\detokenize{psyir:psyclone.psyir.nodes.Node.abs_position}}
\pysigstartsignatures
\pysiglinewithargsret{\sphinxcode{\sphinxupquote{Node.}}\sphinxbfcode{\sphinxupquote{abs\_position}}}{}{}
\pysigstopsignatures
\sphinxAtStartPar
Find a Node’s absolute position in the tree (starting with 0 if
it is the root). Needs to be computed dynamically from the
starting position (0) as its position may change.
\begin{quote}\begin{description}
\sphinxlineitem{Returns}
\sphinxAtStartPar
absolute position of a Node in the tree.

\sphinxlineitem{Return type}
\sphinxAtStartPar
\sphinxhref{https://docs.python.org/3/library/functions.html\#int}{int}

\sphinxlineitem{Raises}
\sphinxAtStartPar
\sphinxstyleliteralstrong{\sphinxupquote{InternalError}} \textendash{} if the absolute position cannot be found.

\end{description}\end{quote}

\end{fulllineitems}

\index{sameParent() (psyclone.psyir.nodes.Node method)@\spxentry{sameParent()}\spxextra{psyclone.psyir.nodes.Node method}}

\begin{fulllineitems}
\phantomsection\label{\detokenize{psyir:psyclone.psyir.nodes.Node.sameParent}}
\pysigstartsignatures
\pysiglinewithargsret{\sphinxcode{\sphinxupquote{Node.}}\sphinxbfcode{\sphinxupquote{sameParent}}}{}{}
\pysigstopsignatures\begin{quote}\begin{description}
\sphinxlineitem{Returns}
\sphinxAtStartPar
True if \sphinxtitleref{node\_2} has the same parent as this node, False                   otherwise.

\sphinxlineitem{Return type}
\sphinxAtStartPar
\sphinxhref{https://docs.python.org/3/library/functions.html\#bool}{bool}

\end{description}\end{quote}

\end{fulllineitems}

\end{quote}

\sphinxAtStartPar
The semantic navigation methods are:
\begin{itemize}
\item {} \begin{description}
\sphinxlineitem{\sphinxcode{\sphinxupquote{Schedule}}:}
\sphinxAtStartPar
subscript operator for indexing the statements (children) inside the
Schedule, e.g. \sphinxcode{\sphinxupquote{sched{[}3{]}}} or \sphinxcode{\sphinxupquote{sched{[}2:4{]}}}.

\end{description}

\item {} \begin{description}
\sphinxlineitem{\sphinxcode{\sphinxupquote{Assignment}}:}\index{lhs() (psyclone.psyir.nodes.Assignment method)@\spxentry{lhs()}\spxextra{psyclone.psyir.nodes.Assignment method}}

\begin{fulllineitems}
\phantomsection\label{\detokenize{psyir:psyclone.psyir.nodes.Assignment.lhs}}
\pysigstartsignatures
\pysiglinewithargsret{\sphinxcode{\sphinxupquote{Assignment.}}\sphinxbfcode{\sphinxupquote{lhs}}}{}{}
\pysigstopsignatures\begin{quote}\begin{description}
\sphinxlineitem{Returns}
\sphinxAtStartPar
the child node representing the Left\sphinxhyphen{}Hand Side of the             assignment.

\sphinxlineitem{Return type}
\sphinxAtStartPar
\sphinxcode{\sphinxupquote{psyclone.psyir.nodes.Node}}

\sphinxlineitem{Raises}
\sphinxAtStartPar
\sphinxstyleliteralstrong{\sphinxupquote{InternalError}} \textendash{} Node has fewer children than expected.

\end{description}\end{quote}

\end{fulllineitems}

\index{rhs() (psyclone.psyir.nodes.Assignment method)@\spxentry{rhs()}\spxextra{psyclone.psyir.nodes.Assignment method}}

\begin{fulllineitems}
\phantomsection\label{\detokenize{psyir:psyclone.psyir.nodes.Assignment.rhs}}
\pysigstartsignatures
\pysiglinewithargsret{\sphinxcode{\sphinxupquote{Assignment.}}\sphinxbfcode{\sphinxupquote{rhs}}}{}{}
\pysigstopsignatures\begin{quote}\begin{description}
\sphinxlineitem{Returns}
\sphinxAtStartPar
the child node representing the Right\sphinxhyphen{}Hand Side of the             assignment.

\sphinxlineitem{Return type}
\sphinxAtStartPar
\sphinxcode{\sphinxupquote{psyclone.psyir.nodes.Node}}

\sphinxlineitem{Raises}
\sphinxAtStartPar
\sphinxstyleliteralstrong{\sphinxupquote{InternalError}} \textendash{} Node has fewer children than expected.

\end{description}\end{quote}

\end{fulllineitems}


\end{description}

\item {} \begin{description}
\sphinxlineitem{\sphinxcode{\sphinxupquote{IfBlock}}:}\index{condition() (psyclone.psyir.nodes.IfBlock method)@\spxentry{condition()}\spxextra{psyclone.psyir.nodes.IfBlock method}}

\begin{fulllineitems}
\phantomsection\label{\detokenize{psyir:psyclone.psyir.nodes.IfBlock.condition}}
\pysigstartsignatures
\pysiglinewithargsret{\sphinxcode{\sphinxupquote{IfBlock.}}\sphinxbfcode{\sphinxupquote{condition}}}{}{}
\pysigstopsignatures
\sphinxAtStartPar
Return the PSyIR Node representing the conditional expression
of this IfBlock.
\begin{quote}\begin{description}
\sphinxlineitem{Returns}
\sphinxAtStartPar
IfBlock conditional expression.

\sphinxlineitem{Return type}
\sphinxAtStartPar
\sphinxcode{\sphinxupquote{psyclone.psyir.nodes.Node}}

\sphinxlineitem{Raises}
\sphinxAtStartPar
\sphinxstyleliteralstrong{\sphinxupquote{InternalError}} \textendash{} If the IfBlock node does not have the correct             number of children.

\end{description}\end{quote}

\end{fulllineitems}

\index{if\_body() (psyclone.psyir.nodes.IfBlock method)@\spxentry{if\_body()}\spxextra{psyclone.psyir.nodes.IfBlock method}}

\begin{fulllineitems}
\phantomsection\label{\detokenize{psyir:psyclone.psyir.nodes.IfBlock.if_body}}
\pysigstartsignatures
\pysiglinewithargsret{\sphinxcode{\sphinxupquote{IfBlock.}}\sphinxbfcode{\sphinxupquote{if\_body}}}{}{}
\pysigstopsignatures
\sphinxAtStartPar
Return the Schedule executed when the IfBlock evaluates to True.
\begin{quote}\begin{description}
\sphinxlineitem{Returns}
\sphinxAtStartPar
Schedule to be executed when IfBlock evaluates to True.

\sphinxlineitem{Return type}
\sphinxAtStartPar
\sphinxcode{\sphinxupquote{psyclone.psyir.nodes.Schedule}}

\sphinxlineitem{Raises}
\sphinxAtStartPar
\sphinxstyleliteralstrong{\sphinxupquote{InternalError}} \textendash{} If the IfBlock node does not have the correct             number of children.

\end{description}\end{quote}

\end{fulllineitems}

\index{else\_body() (psyclone.psyir.nodes.IfBlock method)@\spxentry{else\_body()}\spxextra{psyclone.psyir.nodes.IfBlock method}}

\begin{fulllineitems}
\phantomsection\label{\detokenize{psyir:psyclone.psyir.nodes.IfBlock.else_body}}
\pysigstartsignatures
\pysiglinewithargsret{\sphinxcode{\sphinxupquote{IfBlock.}}\sphinxbfcode{\sphinxupquote{else\_body}}}{}{}
\pysigstopsignatures
\sphinxAtStartPar
If available return the Schedule executed when the IfBlock
evaluates to False, otherwise return None.
\begin{quote}\begin{description}
\sphinxlineitem{Returns}
\sphinxAtStartPar
Schedule to be executed when IfBlock evaluates             to False, if it doesn’t exist returns None.

\sphinxlineitem{Return type}
\sphinxAtStartPar
\sphinxcode{\sphinxupquote{psyclone.psyir.nodes.Schedule}} or NoneType

\end{description}\end{quote}

\end{fulllineitems}


\end{description}

\item {} \begin{description}
\sphinxlineitem{\sphinxcode{\sphinxupquote{Loop}}:}\index{loop\_body() (psyclone.psyir.nodes.Loop method)@\spxentry{loop\_body()}\spxextra{psyclone.psyir.nodes.Loop method}}

\begin{fulllineitems}
\phantomsection\label{\detokenize{psyir:psyclone.psyir.nodes.Loop.loop_body}}
\pysigstartsignatures
\pysiglinewithargsret{\sphinxcode{\sphinxupquote{Loop.}}\sphinxbfcode{\sphinxupquote{loop\_body}}}{}{}
\pysigstopsignatures\begin{quote}\begin{description}
\sphinxlineitem{Returns}
\sphinxAtStartPar
the PSyIR Schedule with the loop body statements.

\sphinxlineitem{Return type}
\sphinxAtStartPar
\sphinxcode{\sphinxupquote{psyclone.psyir.nodes.Schedule}}

\end{description}\end{quote}

\end{fulllineitems}


\end{description}

\item {} \begin{description}
\sphinxlineitem{\sphinxcode{\sphinxupquote{WhileLoop}}:}\index{condition() (psyclone.psyir.nodes.WhileLoop method)@\spxentry{condition()}\spxextra{psyclone.psyir.nodes.WhileLoop method}}

\begin{fulllineitems}
\phantomsection\label{\detokenize{psyir:psyclone.psyir.nodes.WhileLoop.condition}}
\pysigstartsignatures
\pysiglinewithargsret{\sphinxcode{\sphinxupquote{WhileLoop.}}\sphinxbfcode{\sphinxupquote{condition}}}{}{}
\pysigstopsignatures
\sphinxAtStartPar
Return the PSyIR Node representing the conditional expression
of this WhileLoop.
\begin{quote}\begin{description}
\sphinxlineitem{Returns}
\sphinxAtStartPar
WhileLoop conditional expression.

\sphinxlineitem{Return type}
\sphinxAtStartPar
\sphinxcode{\sphinxupquote{psyclone.psyir.nodes.Node}}

\sphinxlineitem{Raises}
\sphinxAtStartPar
\sphinxstyleliteralstrong{\sphinxupquote{InternalError}} \textendash{} If the WhileLoop node does not have the             correct number of children.

\end{description}\end{quote}

\end{fulllineitems}

\index{loop\_body() (psyclone.psyir.nodes.WhileLoop method)@\spxentry{loop\_body()}\spxextra{psyclone.psyir.nodes.WhileLoop method}}

\begin{fulllineitems}
\phantomsection\label{\detokenize{psyir:psyclone.psyir.nodes.WhileLoop.loop_body}}
\pysigstartsignatures
\pysiglinewithargsret{\sphinxcode{\sphinxupquote{WhileLoop.}}\sphinxbfcode{\sphinxupquote{loop\_body}}}{}{}
\pysigstopsignatures
\sphinxAtStartPar
Return the Schedule executed when the WhileLoop condition is True.
\begin{quote}\begin{description}
\sphinxlineitem{Returns}
\sphinxAtStartPar
Schedule to be executed when WhileLoop condition is True.

\sphinxlineitem{Return type}
\sphinxAtStartPar
\sphinxcode{\sphinxupquote{psyclone.psyir.nodes.Schedule}}

\sphinxlineitem{Raises}
\sphinxAtStartPar
\sphinxstyleliteralstrong{\sphinxupquote{InternalError}} \textendash{} If the WhileLoop node does not have the             correct number of children.

\end{description}\end{quote}

\end{fulllineitems}


\end{description}

\item {} \begin{description}
\sphinxlineitem{\sphinxcode{\sphinxupquote{Array}} nodes (e.g. \sphinxcode{\sphinxupquote{ArrayReference}}, \sphinxcode{\sphinxupquote{ArrayOfStructuresReference}}):}\index{indices() (psyclone.psyir.nodes.ArrayReference method)@\spxentry{indices()}\spxextra{psyclone.psyir.nodes.ArrayReference method}}

\begin{fulllineitems}
\phantomsection\label{\detokenize{psyir:psyclone.psyir.nodes.ArrayReference.indices}}
\pysigstartsignatures
\pysiglinewithargsret{\sphinxcode{\sphinxupquote{ArrayReference.}}\sphinxbfcode{\sphinxupquote{indices}}}{}{}
\pysigstopsignatures
\sphinxAtStartPar
Supports semantic\sphinxhyphen{}navigation by returning the list of nodes
representing the index expressions for this array reference.
\begin{quote}\begin{description}
\sphinxlineitem{Returns}
\sphinxAtStartPar
the PSyIR nodes representing the array\sphinxhyphen{}index expressions.

\sphinxlineitem{Return type}
\sphinxAtStartPar
list of \sphinxcode{\sphinxupquote{psyclone.psyir.nodes.Node}}

\sphinxlineitem{Raises}
\sphinxAtStartPar
\sphinxstyleliteralstrong{\sphinxupquote{InternalError}} \textendash{} if this node has no children or if they are                                not valid array\sphinxhyphen{}index expressions.

\end{description}\end{quote}

\end{fulllineitems}


\end{description}

\item {} \begin{description}
\sphinxlineitem{\sphinxcode{\sphinxupquote{RegionDirective}}:}\index{dir\_body() (psyclone.psyir.nodes.RegionDirective method)@\spxentry{dir\_body()}\spxextra{psyclone.psyir.nodes.RegionDirective method}}

\begin{fulllineitems}
\phantomsection\label{\detokenize{psyir:psyclone.psyir.nodes.RegionDirective.dir_body}}
\pysigstartsignatures
\pysiglinewithargsret{\sphinxcode{\sphinxupquote{RegionDirective.}}\sphinxbfcode{\sphinxupquote{dir\_body}}}{}{}
\pysigstopsignatures\begin{quote}\begin{description}
\sphinxlineitem{Returns}
\sphinxAtStartPar
the Schedule associated with this directive.

\sphinxlineitem{Return type}
\sphinxAtStartPar
\sphinxcode{\sphinxupquote{psyclone.psyir.nodes.Schedule}}

\sphinxlineitem{Raises}
\sphinxAtStartPar
\sphinxstyleliteralstrong{\sphinxupquote{InternalError}} \textendash{} if this node does not have a Schedule as                                its first child.

\end{description}\end{quote}

\end{fulllineitems}

\index{clauses() (psyclone.psyir.nodes.RegionDirective method)@\spxentry{clauses()}\spxextra{psyclone.psyir.nodes.RegionDirective method}}

\begin{fulllineitems}
\phantomsection\label{\detokenize{psyir:psyclone.psyir.nodes.RegionDirective.clauses}}
\pysigstartsignatures
\pysiglinewithargsret{\sphinxcode{\sphinxupquote{RegionDirective.}}\sphinxbfcode{\sphinxupquote{clauses}}}{}{}
\pysigstopsignatures\begin{quote}\begin{description}
\sphinxlineitem{Returns}
\sphinxAtStartPar
the Clauses associated with this directive.

\sphinxlineitem{Return type}
\sphinxAtStartPar
List of \sphinxcode{\sphinxupquote{psyclone.psyir.nodes.Clause}}

\end{description}\end{quote}

\end{fulllineitems}


\end{description}

\item {} \begin{description}
\sphinxlineitem{Nodes representing accesses of data within a structure (e.g. \sphinxcode{\sphinxupquote{StructureReference}}, \sphinxcode{\sphinxupquote{StructureMember}}):}\index{member() (psyclone.psyir.nodes.StructureReference method)@\spxentry{member()}\spxextra{psyclone.psyir.nodes.StructureReference method}}

\begin{fulllineitems}
\phantomsection\label{\detokenize{psyir:psyclone.psyir.nodes.StructureReference.member}}
\pysigstartsignatures
\pysiglinewithargsret{\sphinxcode{\sphinxupquote{StructureReference.}}\sphinxbfcode{\sphinxupquote{member}}}{}{}
\pysigstopsignatures\begin{quote}\begin{description}
\sphinxlineitem{Returns}
\sphinxAtStartPar
the PSyIR child representing the accessor component.

\sphinxlineitem{Return type}
\sphinxAtStartPar
\sphinxcode{\sphinxupquote{psyclone.psyir.nodes.Member}}

\sphinxlineitem{Raises}
\sphinxAtStartPar
\sphinxstyleliteralstrong{\sphinxupquote{InternalError}} \textendash{} if the first child of this node is not an
instance of Member.

\end{description}\end{quote}

\end{fulllineitems}


\end{description}

\end{itemize}


\section{DataTypes}
\label{\detokenize{psyir:datatypes}}
\sphinxAtStartPar
The PSyIR supports the following datatypes: \sphinxcode{\sphinxupquote{ScalarType}},
\sphinxcode{\sphinxupquote{ArrayType}}, \sphinxcode{\sphinxupquote{StructureType}}, \sphinxcode{\sphinxupquote{UnresolvedType}}, \sphinxcode{\sphinxupquote{UnsupportedType}}
and \sphinxcode{\sphinxupquote{NoType}}.  These datatypes are used when creating instances of
DataSymbol, RoutineSymbol and Literal (although note that \sphinxcode{\sphinxupquote{NoType}} may
only be used with a RoutineSymbol). \sphinxcode{\sphinxupquote{UnresolvedType}} and \sphinxcode{\sphinxupquote{UnsupportedType}}
are both used when processing existing code. The former is used
when a symbol is being imported from some other scope (e.g. via a USE
statement in Fortran) that hasn’t yet been resolved and the latter is
used when an unsupported form of declaration is encountered.

\sphinxAtStartPar
More information on each of these various datatypes is given in the
following subsections.


\subsection{Scalar DataType}
\label{\detokenize{psyir:scalar-datatype}}
\sphinxAtStartPar
A Scalar datatype consists of an intrinsic and a precision.

\sphinxAtStartPar
The intrinsic can be one of \sphinxcode{\sphinxupquote{INTEGER}}, \sphinxcode{\sphinxupquote{REAL}}, \sphinxcode{\sphinxupquote{BOOLEAN}} and
\sphinxcode{\sphinxupquote{CHARACTER}}.

\sphinxAtStartPar
The precision can be \sphinxcode{\sphinxupquote{UNDEFINED}}, \sphinxcode{\sphinxupquote{SINGLE}}, \sphinxcode{\sphinxupquote{DOUBLE}}, an integer
value specifying the precision in bytes, or a datasymbol (see Section
{\hyperref[\detokenize{psyir:symbol-label}]{\sphinxcrossref{\DUrole{std,std-ref}{Symbols and Symbol Tables}}}}) that contains precision information. Note that
\sphinxcode{\sphinxupquote{UNDEFINED}}, \sphinxcode{\sphinxupquote{SINGLE}} and \sphinxcode{\sphinxupquote{DOUBLE}} allow the precision to be set
by the system so may be different for different architectures. For
example:

\begin{sphinxVerbatim}[commandchars=\\\{\}]
\PYG{g+gp}{\PYGZgt{}\PYGZgt{}\PYGZgt{} }\PYG{n}{char\PYGZus{}type} \PYG{o}{=} \PYG{n}{ScalarType}\PYG{p}{(}\PYG{n}{ScalarType}\PYG{o}{.}\PYG{n}{Intrinsic}\PYG{o}{.}\PYG{n}{CHARACTER}\PYG{p}{,}
\PYG{g+gp}{... }                       \PYG{n}{ScalarType}\PYG{o}{.}\PYG{n}{Precision}\PYG{o}{.}\PYG{n}{UNDEFINED}\PYG{p}{)}
\PYG{g+gp}{\PYGZgt{}\PYGZgt{}\PYGZgt{} }\PYG{n}{int\PYGZus{}type} \PYG{o}{=} \PYG{n}{ScalarType}\PYG{p}{(}\PYG{n}{ScalarType}\PYG{o}{.}\PYG{n}{Intrinsic}\PYG{o}{.}\PYG{n}{INTEGER}\PYG{p}{,}
\PYG{g+gp}{... }                      \PYG{n}{ScalarType}\PYG{o}{.}\PYG{n}{Precision}\PYG{o}{.}\PYG{n}{SINGLE}\PYG{p}{)}
\PYG{g+gp}{\PYGZgt{}\PYGZgt{}\PYGZgt{} }\PYG{n}{bool\PYGZus{}type} \PYG{o}{=} \PYG{n}{ScalarType}\PYG{p}{(}\PYG{n}{ScalarType}\PYG{o}{.}\PYG{n}{Intrinsic}\PYG{o}{.}\PYG{n}{BOOLEAN}\PYG{p}{,} \PYG{l+m+mi}{4}\PYG{p}{)}
\PYG{g+gp}{\PYGZgt{}\PYGZgt{}\PYGZgt{} }\PYG{n}{symbol} \PYG{o}{=} \PYG{n}{DataSymbol}\PYG{p}{(}\PYG{l+s+s2}{\PYGZdq{}}\PYG{l+s+s2}{rdef}\PYG{l+s+s2}{\PYGZdq{}}\PYG{p}{,} \PYG{n}{int\PYGZus{}type}\PYG{p}{,} \PYG{n}{initial\PYGZus{}value}\PYG{o}{=}\PYG{l+m+mi}{4}\PYG{p}{)}
\PYG{g+gp}{\PYGZgt{}\PYGZgt{}\PYGZgt{} }\PYG{n}{scalar\PYGZus{}type} \PYG{o}{=} \PYG{n}{ScalarType}\PYG{p}{(}\PYG{n}{ScalarType}\PYG{o}{.}\PYG{n}{Intrinsic}\PYG{o}{.}\PYG{n}{REAL}\PYG{p}{,} \PYG{n}{symbol}\PYG{p}{)}
\end{sphinxVerbatim}

\sphinxAtStartPar
For convenience PSyclone predefines a number of scalar datatypes:

\sphinxAtStartPar
\sphinxcode{\sphinxupquote{REAL\_TYPE}}, \sphinxcode{\sphinxupquote{INTEGER\_TYPE}}, \sphinxcode{\sphinxupquote{BOOLEAN\_TYPE}} and
\sphinxcode{\sphinxupquote{CHARACTER\_TYPE}} all have precision set to \sphinxcode{\sphinxupquote{UNDEFINED}};

\sphinxAtStartPar
\sphinxcode{\sphinxupquote{REAL\_SINGLE\_TYPE}}, \sphinxcode{\sphinxupquote{REAL\_DOUBLE\_TYPE}}, \sphinxcode{\sphinxupquote{INTEGER\_SINGLE\_TYPE}}
and \sphinxcode{\sphinxupquote{INTEGER\_DOUBLE\_TYPE}};

\sphinxAtStartPar
\sphinxcode{\sphinxupquote{REAL4\_TYPE}}, \sphinxcode{\sphinxupquote{REAL8\_TYPE}}, \sphinxcode{\sphinxupquote{INTEGER4\_TYPE}} and
\sphinxcode{\sphinxupquote{INTEGER8\_TYPE}}.


\subsection{Array DataType}
\label{\detokenize{psyir:array-datatype}}
\sphinxAtStartPar
An Array datatype itself has another datatype (or \sphinxcode{\sphinxupquote{DataTypeSymbol}})
specifying the type of its elements and a shape. The shape can have an
arbitrary number of dimensions. Each dimension captures what is known
about its extent. It is necessary to distinguish between four cases:


\begin{savenotes}\sphinxattablestart
\sphinxthistablewithglobalstyle
\sphinxthistablewithvlinesstyle
\centering
\begin{tabulary}{\linewidth}[t]{|p{9cm}|L|}
\sphinxtoprule
\sphinxstyletheadfamily 
\sphinxAtStartPar
Description
&\sphinxstyletheadfamily 
\sphinxAtStartPar
Entry in \sphinxcode{\sphinxupquote{shape}} list
\\
\sphinxmidrule
\sphinxtableatstartofbodyhook
\sphinxAtStartPar
An array has a static extent known at
compile time.
&
\sphinxAtStartPar
\sphinxcode{\sphinxupquote{ArrayType.ArrayBounds}}
containing integer \sphinxcode{\sphinxupquote{Literal}}
values
\\
\sphinxhline
\sphinxAtStartPar
An array has an extent defined by another
symbol or (constant) PSyIR expression.
&
\sphinxAtStartPar
\sphinxcode{\sphinxupquote{ArrayType.ArrayBounds}}
containing \sphinxcode{\sphinxupquote{Reference}} or
\sphinxcode{\sphinxupquote{Operation}} nodes
\\
\sphinxhline
\sphinxAtStartPar
An array has a definite extent which is not
known at compile time but can be queried
at runtime.
&
\sphinxAtStartPar
\sphinxcode{\sphinxupquote{ArrayType.Extent.ATTRIBUTE}}
\\
\sphinxhline
\sphinxAtStartPar
It is not known whether an array has memory
allocated to it in the current scoping unit.
&
\sphinxAtStartPar
\sphinxcode{\sphinxupquote{ArrayType.Extent.DEFERRED}}
\\
\sphinxbottomrule
\end{tabulary}
\sphinxtableafterendhook\par
\sphinxattableend\end{savenotes}

\sphinxAtStartPar
where \sphinxcode{\sphinxupquote{ArrayType.ArrayBounds}} is a \sphinxcode{\sphinxupquote{namedtuple}} with \sphinxcode{\sphinxupquote{lower}} and
\sphinxcode{\sphinxupquote{upper}} members holding the lower\sphinxhyphen{} and upper\sphinxhyphen{}bounds of the extent of a
given array dimension.

\sphinxAtStartPar
The distinction between the last two cases is that in the former the
extents are known but are kept internally with the array (for example
an assumed shape array in Fortran) and in the latter the array has not
yet been allocated any memory (for example the declaration of an
allocatable array in Fortran) so the extents may have not been defined
yet.

\sphinxAtStartPar
For example:

\begin{sphinxVerbatim}[commandchars=\\\{\}]
\PYG{g+gp}{\PYGZgt{}\PYGZgt{}\PYGZgt{} }\PYG{n}{array\PYGZus{}type} \PYG{o}{=} \PYG{n}{ArrayType}\PYG{p}{(}\PYG{n}{REAL4\PYGZus{}TYPE}\PYG{p}{,} \PYG{p}{[}\PYG{l+m+mi}{5}\PYG{p}{,} \PYG{l+m+mi}{10}\PYG{p}{]}\PYG{p}{)}

\PYG{g+gp}{\PYGZgt{}\PYGZgt{}\PYGZgt{} }\PYG{n}{n\PYGZus{}var} \PYG{o}{=} \PYG{n}{DataSymbol}\PYG{p}{(}\PYG{l+s+s2}{\PYGZdq{}}\PYG{l+s+s2}{n}\PYG{l+s+s2}{\PYGZdq{}}\PYG{p}{,} \PYG{n}{INTEGER\PYGZus{}TYPE}\PYG{p}{)}
\PYG{g+gp}{\PYGZgt{}\PYGZgt{}\PYGZgt{} }\PYG{n}{array\PYGZus{}type} \PYG{o}{=} \PYG{n}{ArrayType}\PYG{p}{(}\PYG{n}{INTEGER\PYGZus{}TYPE}\PYG{p}{,} \PYG{p}{[}\PYG{n}{Reference}\PYG{p}{(}\PYG{n}{n\PYGZus{}var}\PYG{p}{)}\PYG{p}{,}
\PYG{g+gp}{... }                                      \PYG{n}{Reference}\PYG{p}{(}\PYG{n}{n\PYGZus{}var}\PYG{p}{)}\PYG{p}{]}\PYG{p}{)}

\PYG{g+gp}{\PYGZgt{}\PYGZgt{}\PYGZgt{} }\PYG{n}{array\PYGZus{}type} \PYG{o}{=} \PYG{n}{ArrayType}\PYG{p}{(}\PYG{n}{REAL8\PYGZus{}TYPE}\PYG{p}{,} \PYG{p}{[}\PYG{n}{ArrayType}\PYG{o}{.}\PYG{n}{Extent}\PYG{o}{.}\PYG{n}{ATTRIBUTE}\PYG{p}{,}
\PYG{g+gp}{... }                                    \PYG{n}{ArrayType}\PYG{o}{.}\PYG{n}{Extent}\PYG{o}{.}\PYG{n}{ATTRIBUTE}\PYG{p}{]}\PYG{p}{)}

\PYG{g+gp}{\PYGZgt{}\PYGZgt{}\PYGZgt{} }\PYG{n}{array\PYGZus{}type} \PYG{o}{=} \PYG{n}{ArrayType}\PYG{p}{(}\PYG{n}{BOOLEAN\PYGZus{}TYPE}\PYG{p}{,} \PYG{p}{[}\PYG{n}{ArrayType}\PYG{o}{.}\PYG{n}{Extent}\PYG{o}{.}\PYG{n}{DEFERRED}\PYG{p}{]}\PYG{p}{)}
\end{sphinxVerbatim}


\subsection{Structure Datatype}
\label{\detokenize{psyir:structure-datatype}}
\sphinxAtStartPar
A Structure datatype consists of a dictionary of components where the
name of each component is used as the corresponding key. Each component
is stored as a named tuple with \sphinxcode{\sphinxupquote{name}}, \sphinxcode{\sphinxupquote{datatype}} and \sphinxcode{\sphinxupquote{visibility}}
members.

\sphinxAtStartPar
For example:

\begin{sphinxVerbatim}[commandchars=\\\{\}]
\PYG{c+c1}{\PYGZsh{} Shorthand for a scalar type with REAL\PYGZus{}KIND precision}
\PYG{n}{SCALAR\PYGZus{}TYPE} \PYG{o}{=} \PYG{n}{ScalarType}\PYG{p}{(}\PYG{n}{ScalarType}\PYG{o}{.}\PYG{n}{Intrinsic}\PYG{o}{.}\PYG{n}{REAL}\PYG{p}{,} \PYG{n}{REAL\PYGZus{}KIND}\PYG{p}{)}

\PYG{c+c1}{\PYGZsh{} Structure\PYGZhy{}type definition}
\PYG{n}{GRID\PYGZus{}TYPE} \PYG{o}{=} \PYG{n}{StructureType}\PYG{o}{.}\PYG{n}{create}\PYG{p}{(}\PYG{p}{[}
    \PYG{p}{(}\PYG{l+s+s2}{\PYGZdq{}}\PYG{l+s+s2}{dx}\PYG{l+s+s2}{\PYGZdq{}}\PYG{p}{,} \PYG{n}{SCALAR\PYGZus{}TYPE}\PYG{p}{,} \PYG{n}{Symbol}\PYG{o}{.}\PYG{n}{Visibility}\PYG{o}{.}\PYG{n}{PUBLIC}\PYG{p}{)}\PYG{p}{,}
    \PYG{p}{(}\PYG{l+s+s2}{\PYGZdq{}}\PYG{l+s+s2}{dy}\PYG{l+s+s2}{\PYGZdq{}}\PYG{p}{,} \PYG{n}{SCALAR\PYGZus{}TYPE}\PYG{p}{,} \PYG{n}{Symbol}\PYG{o}{.}\PYG{n}{Visibility}\PYG{o}{.}\PYG{n}{PUBLIC}\PYG{p}{)}\PYG{p}{]}\PYG{p}{)}

\PYG{n}{GRID\PYGZus{}TYPE\PYGZus{}SYMBOL} \PYG{o}{=} \PYG{n}{DataTypeSymbol}\PYG{p}{(}\PYG{l+s+s2}{\PYGZdq{}}\PYG{l+s+s2}{grid\PYGZus{}type}\PYG{l+s+s2}{\PYGZdq{}}\PYG{p}{,} \PYG{n}{GRID\PYGZus{}TYPE}\PYG{p}{)}

\PYG{c+c1}{\PYGZsh{} A structure\PYGZhy{}type containing other structure types}
\PYG{n}{FIELD\PYGZus{}TYPE\PYGZus{}DEF} \PYG{o}{=} \PYG{n}{StructureType}\PYG{o}{.}\PYG{n}{create}\PYG{p}{(}
    \PYG{p}{[}\PYG{p}{(}\PYG{l+s+s2}{\PYGZdq{}}\PYG{l+s+s2}{data}\PYG{l+s+s2}{\PYGZdq{}}\PYG{p}{,} \PYG{n}{ArrayType}\PYG{p}{(}\PYG{n}{SCALAR\PYGZus{}TYPE}\PYG{p}{,} \PYG{p}{[}\PYG{l+m+mi}{10}\PYG{p}{]}\PYG{p}{)}\PYG{p}{,} \PYG{n}{Symbol}\PYG{o}{.}\PYG{n}{Visibility}\PYG{o}{.}\PYG{n}{PUBLIC}\PYG{p}{)}\PYG{p}{,}
     \PYG{p}{(}\PYG{l+s+s2}{\PYGZdq{}}\PYG{l+s+s2}{grid}\PYG{l+s+s2}{\PYGZdq{}}\PYG{p}{,} \PYG{n}{GRID\PYGZus{}TYPE\PYGZus{}SYMBOL}\PYG{p}{,} \PYG{n}{Symbol}\PYG{o}{.}\PYG{n}{Visibility}\PYG{o}{.}\PYG{n}{PUBLIC}\PYG{p}{)}\PYG{p}{,}
     \PYG{p}{(}\PYG{l+s+s2}{\PYGZdq{}}\PYG{l+s+s2}{sub\PYGZus{}meshes}\PYG{l+s+s2}{\PYGZdq{}}\PYG{p}{,} \PYG{n}{ArrayType}\PYG{p}{(}\PYG{n}{GRID\PYGZus{}TYPE\PYGZus{}SYMBOL}\PYG{p}{,} \PYG{p}{[}\PYG{l+m+mi}{3}\PYG{p}{]}\PYG{p}{)}\PYG{p}{,}
      \PYG{n}{Symbol}\PYG{o}{.}\PYG{n}{Visibility}\PYG{o}{.}\PYG{n}{PUBLIC}\PYG{p}{)}\PYG{p}{,}
     \PYG{p}{(}\PYG{l+s+s2}{\PYGZdq{}}\PYG{l+s+s2}{flag}\PYG{l+s+s2}{\PYGZdq{}}\PYG{p}{,} \PYG{n}{INTEGER4\PYGZus{}TYPE}\PYG{p}{,} \PYG{n}{Symbol}\PYG{o}{.}\PYG{n}{Visibility}\PYG{o}{.}\PYG{n}{PUBLIC}\PYG{p}{)}\PYG{p}{]}\PYG{p}{)}
\end{sphinxVerbatim}


\subsection{Unknown DataType}
\label{\detokenize{psyir:unknown-datatype}}
\sphinxAtStartPar
If a PSyIR frontend encounters an unsupported declaration then the
corresponding Symbol is given \sphinxhref{https://psyclone-ref.readthedocs.io/en/latest/autogenerated/psyclone.psyir.symbols.html\#psyclone.psyir.symbols.UnsupportedType}{UnsupportedType}.
The text of the original declaration is stored in the type object and is
available via the \sphinxcode{\sphinxupquote{declaration}} property.


\subsection{NoType}
\label{\detokenize{psyir:notype}}
\sphinxAtStartPar
\sphinxcode{\sphinxupquote{NoType}} represents the empty type, equivalent to \sphinxcode{\sphinxupquote{void}} in C. It
is currently only used to describe a RoutineSymbol that has no return
type (such as a Fortran subroutine).


\section{Symbols and Symbol Tables}
\label{\detokenize{psyir:symbols-and-symbol-tables}}\label{\detokenize{psyir:symbol-label}}
\sphinxAtStartPar
Some PSyIR nodes have an associated Symbol Table
(\sphinxtitleref{psyclone.psyir.symbols.SymbolTable}) which keeps a record of the
Symbols (\sphinxtitleref{psyclone.psyir.symbols.Symbol}) specified and used within them.

\sphinxAtStartPar
Symbol Tables can be nested (i.e. a node with an attached symbol table
can be an ancestor or descendent of a node with an attached symbol
table). If the same symbol name is used in a hierarchy of symbol tables
then the symbol within the symbol table attached to the closest
ancestor node is in scope. By default, symbol tables are aware of
other symbol tables and will return information about relevant symbols
from all symbol tables.

\sphinxAtStartPar
The \sphinxcode{\sphinxupquote{SymbolTable}} has the following interface:
\index{SymbolTable (class in psyclone.psyir.symbols)@\spxentry{SymbolTable}\spxextra{class in psyclone.psyir.symbols}}

\begin{fulllineitems}
\phantomsection\label{\detokenize{psyir:psyclone.psyir.symbols.SymbolTable}}
\pysigstartsignatures
\pysiglinewithargsret{\sphinxbfcode{\sphinxupquote{class\DUrole{w,w}{  }}}\sphinxcode{\sphinxupquote{psyclone.psyir.symbols.}}\sphinxbfcode{\sphinxupquote{SymbolTable}}}{\sphinxparam{\DUrole{n,n}{node}\DUrole{o,o}{=}\DUrole{default_value}{None}}\sphinxparamcomma \sphinxparam{\DUrole{n,n}{default\_visibility}\DUrole{o,o}{=}\DUrole{default_value}{Visibility.PUBLIC}}}{}
\pysigstopsignatures
\sphinxAtStartPar
Encapsulates the symbol table and provides methods to add new
symbols and look up existing symbols. Nested scopes are supported
and, by default, the add and lookup methods take any ancestor
symbol tables into consideration (ones attached to nodes that are
ancestors of the node that this symbol table is attached to). If the
default visibility is not specified then it defaults to
Symbol.Visbility.PUBLIC.
\begin{quote}\begin{description}
\sphinxlineitem{Parameters}\begin{itemize}
\item {} 
\sphinxAtStartPar
\sphinxstyleliteralstrong{\sphinxupquote{node}} (Optional{[}\sphinxcode{\sphinxupquote{psyclone.psyir.nodes.Schedule}} |
\sphinxcode{\sphinxupquote{psyclone.psyir.nodes.Container}}{]}) \textendash{} reference to the Schedule or Container to which this
symbol table belongs.

\item {} 
\sphinxAtStartPar
\sphinxstyleliteralstrong{\sphinxupquote{default\_visibility}} \textendash{} optional default visibility value for this
symbol table, if not provided it defaults to PUBLIC visibility.

\end{itemize}

\sphinxlineitem{Raises}
\sphinxAtStartPar
\sphinxhref{https://docs.python.org/3/library/exceptions.html\#TypeError}{\sphinxstyleliteralstrong{\sphinxupquote{TypeError}}} \textendash{} if node argument is not a Schedule or a Container.

\end{description}\end{quote}

\end{fulllineitems}


\sphinxAtStartPar
Where each element is a \sphinxcode{\sphinxupquote{Symbol}} with an immutable name:
\index{Symbol (class in psyclone.psyir.symbols)@\spxentry{Symbol}\spxextra{class in psyclone.psyir.symbols}}

\begin{fulllineitems}
\phantomsection\label{\detokenize{psyir:psyclone.psyir.symbols.Symbol}}
\pysigstartsignatures
\pysiglinewithargsret{\sphinxbfcode{\sphinxupquote{class\DUrole{w,w}{  }}}\sphinxcode{\sphinxupquote{psyclone.psyir.symbols.}}\sphinxbfcode{\sphinxupquote{Symbol}}}{\sphinxparam{\DUrole{n,n}{name}}\sphinxparamcomma \sphinxparam{\DUrole{n,n}{visibility}\DUrole{o,o}{=}\DUrole{default_value}{Visibility.PUBLIC}}\sphinxparamcomma \sphinxparam{\DUrole{n,n}{interface}\DUrole{o,o}{=}\DUrole{default_value}{None}}}{}
\pysigstopsignatures
\sphinxAtStartPar
Generic Symbol item for the Symbol Table and PSyIR References.
It has an immutable name label because it must always match with the
key in the SymbolTable. If the symbol is private then it is only visible
to those nodes that are descendants of the Node to which its containing
Symbol Table belongs.
\begin{quote}\begin{description}
\sphinxlineitem{Parameters}\begin{itemize}
\item {} 
\sphinxAtStartPar
\sphinxstyleliteralstrong{\sphinxupquote{name}} (\sphinxhref{https://docs.python.org/3/library/stdtypes.html\#str}{\sphinxstyleliteralemphasis{\sphinxupquote{str}}}) \textendash{} name of the symbol.

\item {} 
\sphinxAtStartPar
\sphinxstyleliteralstrong{\sphinxupquote{visibility}} (\sphinxcode{\sphinxupquote{psyclone.psyir.symbols.Symbol.Visibility}}) \textendash{} the visibility of the symbol.

\item {} 
\sphinxAtStartPar
\sphinxstyleliteralstrong{\sphinxupquote{interface}} (Optional{[}         \sphinxcode{\sphinxupquote{psyclone.psyir.symbols.symbol.SymbolInterface}}{]}) \textendash{} optional object describing the interface to this         symbol (i.e. whether it is passed as a routine argument or         accessed in some other way). Defaults to         {\hyperref[\detokenize{psyir:psyclone.psyir.symbols.AutomaticInterface}]{\sphinxcrossref{\sphinxcode{\sphinxupquote{psyclone.psyir.symbols.AutomaticInterface}}}}}

\end{itemize}

\sphinxlineitem{Raises}
\sphinxAtStartPar
\sphinxhref{https://docs.python.org/3/library/exceptions.html\#TypeError}{\sphinxstyleliteralstrong{\sphinxupquote{TypeError}}} \textendash{} if the name is not a str.

\end{description}\end{quote}

\end{fulllineitems}


\sphinxAtStartPar
There are several \sphinxcode{\sphinxupquote{Symbol}} sub\sphinxhyphen{}classes to represent different
labeled entities in the PSyIR. At the moment the available symbols
are:
\begin{itemize}
\item {} \index{ContainerSymbol (class in psyclone.psyir.symbols)@\spxentry{ContainerSymbol}\spxextra{class in psyclone.psyir.symbols}}

\begin{fulllineitems}
\phantomsection\label{\detokenize{psyir:psyclone.psyir.symbols.ContainerSymbol}}
\pysigstartsignatures
\pysiglinewithargsret{\sphinxbfcode{\sphinxupquote{class\DUrole{w,w}{  }}}\sphinxcode{\sphinxupquote{psyclone.psyir.symbols.}}\sphinxbfcode{\sphinxupquote{ContainerSymbol}}}{\sphinxparam{\DUrole{n,n}{name}}\sphinxparamcomma \sphinxparam{\DUrole{o,o}{**}\DUrole{n,n}{kwargs}}}{}
\pysigstopsignatures
\sphinxAtStartPar
Symbol that represents a reference to a Container. The reference
is lazy evaluated, this means that the Symbol will be created without
parsing and importing the referenced container, but this can be imported
when needed.
\begin{quote}\begin{description}
\sphinxlineitem{Parameters}\begin{itemize}
\item {} 
\sphinxAtStartPar
\sphinxstyleliteralstrong{\sphinxupquote{name}} (\sphinxhref{https://docs.python.org/3/library/stdtypes.html\#str}{\sphinxstyleliteralemphasis{\sphinxupquote{str}}}) \textendash{} name of the symbol.

\item {} 
\sphinxAtStartPar
\sphinxstyleliteralstrong{\sphinxupquote{wildcard\_import}} (\sphinxhref{https://docs.python.org/3/library/functions.html\#bool}{\sphinxstyleliteralemphasis{\sphinxupquote{bool}}}) \textendash{} if all public Symbols of the Container are
imported into the current scope. Defaults to False.

\item {} 
\sphinxAtStartPar
\sphinxstyleliteralstrong{\sphinxupquote{is\_intrinsic}} (\sphinxhref{https://docs.python.org/3/library/functions.html\#bool}{\sphinxstyleliteralemphasis{\sphinxupquote{bool}}}) \textendash{} if the module is an intrinsic import. Defauts
to False.

\item {} 
\sphinxAtStartPar
\sphinxstyleliteralstrong{\sphinxupquote{kwargs}} (\sphinxstyleliteralemphasis{\sphinxupquote{unwrapped dict.}}) \textendash{} additional keyword arguments provided by
{\hyperref[\detokenize{psyir:psyclone.psyir.symbols.Symbol}]{\sphinxcrossref{\sphinxcode{\sphinxupquote{psyclone.psyir.symbols.Symbol}}}}}.

\end{itemize}

\end{description}\end{quote}

\end{fulllineitems}


\item {} \index{DataSymbol (class in psyclone.psyir.symbols)@\spxentry{DataSymbol}\spxextra{class in psyclone.psyir.symbols}}

\begin{fulllineitems}
\phantomsection\label{\detokenize{psyir:psyclone.psyir.symbols.DataSymbol}}
\pysigstartsignatures
\pysiglinewithargsret{\sphinxbfcode{\sphinxupquote{class\DUrole{w,w}{  }}}\sphinxcode{\sphinxupquote{psyclone.psyir.symbols.}}\sphinxbfcode{\sphinxupquote{DataSymbol}}}{\sphinxparam{\DUrole{n,n}{name}}\sphinxparamcomma \sphinxparam{\DUrole{n,n}{datatype}}\sphinxparamcomma \sphinxparam{\DUrole{n,n}{is\_constant}\DUrole{o,o}{=}\DUrole{default_value}{False}}\sphinxparamcomma \sphinxparam{\DUrole{n,n}{initial\_value}\DUrole{o,o}{=}\DUrole{default_value}{None}}\sphinxparamcomma \sphinxparam{\DUrole{o,o}{**}\DUrole{n,n}{kwargs}}}{}
\pysigstopsignatures
\sphinxAtStartPar
Symbol identifying a data element. It contains information about:
the datatype, the shape (in column\sphinxhyphen{}major order) and the interface
to that symbol (i.e. Local, Global, Argument).
\begin{quote}\begin{description}
\sphinxlineitem{Parameters}\begin{itemize}
\item {} 
\sphinxAtStartPar
\sphinxstyleliteralstrong{\sphinxupquote{name}} (\sphinxhref{https://docs.python.org/3/library/stdtypes.html\#str}{\sphinxstyleliteralemphasis{\sphinxupquote{str}}}) \textendash{} name of the symbol.

\item {} 
\sphinxAtStartPar
\sphinxstyleliteralstrong{\sphinxupquote{datatype}} (\sphinxcode{\sphinxupquote{psyclone.psyir.symbols.DataType}}) \textendash{} data type of the symbol.

\item {} 
\sphinxAtStartPar
\sphinxstyleliteralstrong{\sphinxupquote{is\_constant}} (\sphinxhref{https://docs.python.org/3/library/functions.html\#bool}{\sphinxstyleliteralemphasis{\sphinxupquote{bool}}}) \textendash{} whether this DataSymbol is a compile\sphinxhyphen{}time
constant (default is False). If True then an \sphinxtitleref{initial\_value} must
also be provided.

\item {} 
\sphinxAtStartPar
\sphinxstyleliteralstrong{\sphinxupquote{initial\_value}} (Optional{[}item of TYPE\_MAP\_TO\_PYTHON |
\sphinxcode{\sphinxupquote{psyclone.psyir.nodes.Node}}{]}) \textendash{} sets a fixed known expression as an initial
value for this DataSymbol. If \sphinxtitleref{is\_constant} is True then this
Symbol will always have this value. If the value is None then this
symbol does not have an initial value (and cannot be a constant).
Otherwise it can receive PSyIR expressions or Python intrinsic types
available in the TYPE\_MAP\_TO\_PYTHON map. By default it is None.

\item {} 
\sphinxAtStartPar
\sphinxstyleliteralstrong{\sphinxupquote{kwargs}} (\sphinxstyleliteralemphasis{\sphinxupquote{unwrapped dict.}}) \textendash{} additional keyword arguments provided by
\sphinxcode{\sphinxupquote{psyclone.psyir.symbols.TypedSymbol}}

\end{itemize}

\end{description}\end{quote}

\end{fulllineitems}


\item {} \index{DataTypeSymbol (class in psyclone.psyir.symbols)@\spxentry{DataTypeSymbol}\spxextra{class in psyclone.psyir.symbols}}

\begin{fulllineitems}
\phantomsection\label{\detokenize{psyir:psyclone.psyir.symbols.DataTypeSymbol}}
\pysigstartsignatures
\pysiglinewithargsret{\sphinxbfcode{\sphinxupquote{class\DUrole{w,w}{  }}}\sphinxcode{\sphinxupquote{psyclone.psyir.symbols.}}\sphinxbfcode{\sphinxupquote{DataTypeSymbol}}}{\sphinxparam{\DUrole{n,n}{name}}\sphinxparamcomma \sphinxparam{\DUrole{n,n}{datatype}}\sphinxparamcomma \sphinxparam{\DUrole{n,n}{visibility}\DUrole{o,o}{=}\DUrole{default_value}{Visibility.PUBLIC}}\sphinxparamcomma \sphinxparam{\DUrole{n,n}{interface}\DUrole{o,o}{=}\DUrole{default_value}{None}}}{}
\pysigstopsignatures
\sphinxAtStartPar
Symbol identifying a user\sphinxhyphen{}defined type (e.g. a derived type in Fortran).
\begin{quote}\begin{description}
\sphinxlineitem{Parameters}\begin{itemize}
\item {} 
\sphinxAtStartPar
\sphinxstyleliteralstrong{\sphinxupquote{name}} (\sphinxhref{https://docs.python.org/3/library/stdtypes.html\#str}{\sphinxstyleliteralemphasis{\sphinxupquote{str}}}) \textendash{} the name of this symbol.

\item {} 
\sphinxAtStartPar
\sphinxstyleliteralstrong{\sphinxupquote{datatype}} (\sphinxcode{\sphinxupquote{psyclone.psyir.symbols.DataType}}) \textendash{} the type represented by this symbol.

\item {} 
\sphinxAtStartPar
\sphinxstyleliteralstrong{\sphinxupquote{visibility}} (\sphinxcode{\sphinxupquote{psyclone.psyir.symbols.Symbol.Visibility}}) \textendash{} the visibility of this symbol.

\item {} 
\sphinxAtStartPar
\sphinxstyleliteralstrong{\sphinxupquote{interface}} (\sphinxcode{\sphinxupquote{psyclone.psyir.symbols.SymbolInterface}}) \textendash{} the interface to this symbol.

\end{itemize}

\end{description}\end{quote}

\end{fulllineitems}


\item {} \index{IntrinsicSymbol (class in psyclone.psyir.symbols)@\spxentry{IntrinsicSymbol}\spxextra{class in psyclone.psyir.symbols}}

\begin{fulllineitems}
\phantomsection\label{\detokenize{psyir:psyclone.psyir.symbols.IntrinsicSymbol}}
\pysigstartsignatures
\pysiglinewithargsret{\sphinxbfcode{\sphinxupquote{class\DUrole{w,w}{  }}}\sphinxcode{\sphinxupquote{psyclone.psyir.symbols.}}\sphinxbfcode{\sphinxupquote{IntrinsicSymbol}}}{\sphinxparam{\DUrole{n,n}{name}}\sphinxparamcomma \sphinxparam{\DUrole{n,n}{intrinsic}}\sphinxparamcomma \sphinxparam{\DUrole{o,o}{**}\DUrole{n,n}{kwargs}}}{}
\pysigstopsignatures
\sphinxAtStartPar
Symbol identifying a callable intrinsic routine.
\begin{quote}\begin{description}
\sphinxlineitem{Parameters}\begin{itemize}
\item {} 
\sphinxAtStartPar
\sphinxstyleliteralstrong{\sphinxupquote{name}} (\sphinxhref{https://docs.python.org/3/library/stdtypes.html\#str}{\sphinxstyleliteralemphasis{\sphinxupquote{str}}}) \textendash{} name of the symbol.

\item {} 
\sphinxAtStartPar
\sphinxstyleliteralstrong{\sphinxupquote{intrinsic}} (\sphinxcode{\sphinxupquote{psyclone.psyir.nodes.IntrinsicCall.Intrinsic}}) \textendash{} the intrinsic enum describing this Symbol.

\item {} 
\sphinxAtStartPar
\sphinxstyleliteralstrong{\sphinxupquote{kwargs}} (\sphinxstyleliteralemphasis{\sphinxupquote{unwrapped dict.}}) \textendash{} additional keyword arguments provided by
\sphinxcode{\sphinxupquote{psyclone.psyir.symbols.TypedSymbol}}

\end{itemize}

\end{description}\end{quote}

\sphinxAtStartPar
\# TODO \#2541: Currently name and the intrinsic should match, we really
\# just need the name, and make all the Intrinsic singature information
\# live inside the IntrinsicSymbol class.

\end{fulllineitems}


\item {} \index{RoutineSymbol (class in psyclone.psyir.symbols)@\spxentry{RoutineSymbol}\spxextra{class in psyclone.psyir.symbols}}

\begin{fulllineitems}
\phantomsection\label{\detokenize{psyir:psyclone.psyir.symbols.RoutineSymbol}}
\pysigstartsignatures
\pysiglinewithargsret{\sphinxbfcode{\sphinxupquote{class\DUrole{w,w}{  }}}\sphinxcode{\sphinxupquote{psyclone.psyir.symbols.}}\sphinxbfcode{\sphinxupquote{RoutineSymbol}}}{\sphinxparam{\DUrole{n,n}{name}}\sphinxparamcomma \sphinxparam{\DUrole{n,n}{datatype}\DUrole{o,o}{=}\DUrole{default_value}{None}}\sphinxparamcomma \sphinxparam{\DUrole{o,o}{**}\DUrole{n,n}{kwargs}}}{}
\pysigstopsignatures
\sphinxAtStartPar
Symbol identifying a callable routine.
\begin{quote}\begin{description}
\sphinxlineitem{Parameters}\begin{itemize}
\item {} 
\sphinxAtStartPar
\sphinxstyleliteralstrong{\sphinxupquote{name}} (\sphinxhref{https://docs.python.org/3/library/stdtypes.html\#str}{\sphinxstyleliteralemphasis{\sphinxupquote{str}}}) \textendash{} name of the symbol.

\item {} 
\sphinxAtStartPar
\sphinxstyleliteralstrong{\sphinxupquote{datatype}} (\sphinxcode{\sphinxupquote{psyclone.psyir.symbols.DataType}}) \textendash{} data type of the symbol. Default to NoType().

\item {} 
\sphinxAtStartPar
\sphinxstyleliteralstrong{\sphinxupquote{kwargs}} (\sphinxstyleliteralemphasis{\sphinxupquote{unwrapped dict.}}) \textendash{} additional keyword arguments provided by                    \sphinxcode{\sphinxupquote{psyclone.psyir.symbols.TypedSymbol}}

\end{itemize}

\end{description}\end{quote}

\end{fulllineitems}


\item {} \index{GenericInterfaceSymbol (class in psyclone.psyir.symbols)@\spxentry{GenericInterfaceSymbol}\spxextra{class in psyclone.psyir.symbols}}

\begin{fulllineitems}
\phantomsection\label{\detokenize{psyir:psyclone.psyir.symbols.GenericInterfaceSymbol}}
\pysigstartsignatures
\pysiglinewithargsret{\sphinxbfcode{\sphinxupquote{class\DUrole{w,w}{  }}}\sphinxcode{\sphinxupquote{psyclone.psyir.symbols.}}\sphinxbfcode{\sphinxupquote{GenericInterfaceSymbol}}}{\sphinxparam{\DUrole{n,n}{name}}\sphinxparamcomma \sphinxparam{\DUrole{n,n}{routines}}\sphinxparamcomma \sphinxparam{\DUrole{o,o}{**}\DUrole{n,n}{kwargs}}}{}
\pysigstopsignatures
\sphinxAtStartPar
Symbol identifying a generic interface that maps to a number of
different callable routines.
\begin{quote}\begin{description}
\sphinxlineitem{Parameters}\begin{itemize}
\item {} 
\sphinxAtStartPar
\sphinxstyleliteralstrong{\sphinxupquote{name}} (\sphinxhref{https://docs.python.org/3/library/stdtypes.html\#str}{\sphinxstyleliteralemphasis{\sphinxupquote{str}}}) \textendash{} name of the interface.

\item {} 
\sphinxAtStartPar
\sphinxstyleliteralstrong{\sphinxupquote{routines}} (list{[}tuple{[}
{\hyperref[\detokenize{psyir:psyclone.psyir.symbols.RoutineSymbol}]{\sphinxcrossref{\sphinxcode{\sphinxupquote{psyclone.psyir.symbols.RoutineSymbol}}}}},
bool{]}{]}) \textendash{} the routines that this interface provides access
to and whether or not each of them is a module procedure.

\item {} 
\sphinxAtStartPar
\sphinxstyleliteralstrong{\sphinxupquote{kwargs}} (\sphinxstyleliteralemphasis{\sphinxupquote{unwrapped dict.}}) \textendash{} additional keyword arguments provided by
\sphinxcode{\sphinxupquote{psyclone.psyir.symbols.TypedSymbol}}

\end{itemize}

\end{description}\end{quote}

\end{fulllineitems}


\end{itemize}

\sphinxAtStartPar
See the reference guide for the full API documentation of the
\sphinxhref{https://psyclone-ref.readthedocs.io/en/latest/autogenerated/psyclone.psyir.symbols.html\#psyclone.psyir.symbols.SymbolTable}{SymbolTable}
and the \sphinxhref{https://psyclone-ref.readthedocs.io/en/latest/autogenerated/psyclone.psyir.symbols.html\#module-psyclone.psyir.symbols}{Symbol types}.


\subsection{Symbol Interfaces}
\label{\detokenize{psyir:symbol-interfaces}}
\sphinxAtStartPar
Each symbol has a Symbol Interface with the information about how the
variable data is provided into the local context. The currently available
Interfaces are:
\begin{itemize}
\item {} \index{AutomaticInterface (class in psyclone.psyir.symbols)@\spxentry{AutomaticInterface}\spxextra{class in psyclone.psyir.symbols}}

\begin{fulllineitems}
\phantomsection\label{\detokenize{psyir:psyclone.psyir.symbols.AutomaticInterface}}
\pysigstartsignatures
\pysigline{\sphinxbfcode{\sphinxupquote{class\DUrole{w,w}{  }}}\sphinxcode{\sphinxupquote{psyclone.psyir.symbols.}}\sphinxbfcode{\sphinxupquote{AutomaticInterface}}}
\pysigstopsignatures
\sphinxAtStartPar
The symbol is declared without attributes. Its data will live
during the local context.

\end{fulllineitems}


\item {} \index{DefaultModuleInterface (class in psyclone.psyir.symbols)@\spxentry{DefaultModuleInterface}\spxextra{class in psyclone.psyir.symbols}}

\begin{fulllineitems}
\phantomsection\label{\detokenize{psyir:psyclone.psyir.symbols.DefaultModuleInterface}}
\pysigstartsignatures
\pysigline{\sphinxbfcode{\sphinxupquote{class\DUrole{w,w}{  }}}\sphinxcode{\sphinxupquote{psyclone.psyir.symbols.}}\sphinxbfcode{\sphinxupquote{DefaultModuleInterface}}}
\pysigstopsignatures
\sphinxAtStartPar
The symbol contains data declared in a module scope without additional
attributes.

\end{fulllineitems}


\item {} \index{ImportInterface (class in psyclone.psyir.symbols)@\spxentry{ImportInterface}\spxextra{class in psyclone.psyir.symbols}}

\begin{fulllineitems}
\phantomsection\label{\detokenize{psyir:psyclone.psyir.symbols.ImportInterface}}
\pysigstartsignatures
\pysiglinewithargsret{\sphinxbfcode{\sphinxupquote{class\DUrole{w,w}{  }}}\sphinxcode{\sphinxupquote{psyclone.psyir.symbols.}}\sphinxbfcode{\sphinxupquote{ImportInterface}}}{\sphinxparam{\DUrole{n,n}{container\_symbol}}\sphinxparamcomma \sphinxparam{\DUrole{n,n}{orig\_name}\DUrole{o,o}{=}\DUrole{default_value}{None}}}{}
\pysigstopsignatures
\sphinxAtStartPar
Describes the interface to a Symbol that is imported from an
external PSyIR container. The symbol can be renamed on import and,
if so, its original name in the Container is specified using the
optional ‘orig\_name’ argument.
\begin{quote}\begin{description}
\sphinxlineitem{Parameters}\begin{itemize}
\item {} 
\sphinxAtStartPar
\sphinxstyleliteralstrong{\sphinxupquote{container\_symbol}} ({\hyperref[\detokenize{psyir:psyclone.psyir.symbols.ContainerSymbol}]{\sphinxcrossref{\sphinxcode{\sphinxupquote{psyclone.psyir.symbols.ContainerSymbol}}}}}) \textendash{} symbol representing the external container         from which the symbol is imported.

\item {} 
\sphinxAtStartPar
\sphinxstyleliteralstrong{\sphinxupquote{orig\_name}} (\sphinxstyleliteralemphasis{\sphinxupquote{Optional}}\sphinxstyleliteralemphasis{\sphinxupquote{{[}}}\sphinxhref{https://docs.python.org/3/library/stdtypes.html\#str}{\sphinxstyleliteralemphasis{\sphinxupquote{str}}}\sphinxstyleliteralemphasis{\sphinxupquote{{]}}}) \textendash{} the name of the symbol in the         external container before it is renamed, or None (the default) if         it is not renamed.

\end{itemize}

\sphinxlineitem{Raises}
\sphinxAtStartPar
\sphinxhref{https://docs.python.org/3/library/exceptions.html\#TypeError}{\sphinxstyleliteralstrong{\sphinxupquote{TypeError}}} \textendash{} if the orig\_name argument is an unexpected type.

\end{description}\end{quote}

\end{fulllineitems}


\item {} \index{ArgumentInterface (class in psyclone.psyir.symbols)@\spxentry{ArgumentInterface}\spxextra{class in psyclone.psyir.symbols}}

\begin{fulllineitems}
\phantomsection\label{\detokenize{psyir:psyclone.psyir.symbols.ArgumentInterface}}
\pysigstartsignatures
\pysiglinewithargsret{\sphinxbfcode{\sphinxupquote{class\DUrole{w,w}{  }}}\sphinxcode{\sphinxupquote{psyclone.psyir.symbols.}}\sphinxbfcode{\sphinxupquote{ArgumentInterface}}}{\sphinxparam{\DUrole{n,n}{access}\DUrole{o,o}{=}\DUrole{default_value}{None}}}{}
\pysigstopsignatures
\sphinxAtStartPar
Captures the interface to a Symbol that is accessed as a routine
argument.
\begin{quote}\begin{description}
\sphinxlineitem{Parameters}
\sphinxAtStartPar
\sphinxstyleliteralstrong{\sphinxupquote{access}} (\sphinxcode{\sphinxupquote{psyclone.psyir.symbols.ArgumentInterface.Access}}) \textendash{} specifies how the argument is used in the Schedule

\end{description}\end{quote}

\end{fulllineitems}


\item {} \index{StaticInterface (class in psyclone.psyir.symbols)@\spxentry{StaticInterface}\spxextra{class in psyclone.psyir.symbols}}

\begin{fulllineitems}
\phantomsection\label{\detokenize{psyir:psyclone.psyir.symbols.StaticInterface}}
\pysigstartsignatures
\pysigline{\sphinxbfcode{\sphinxupquote{class\DUrole{w,w}{  }}}\sphinxcode{\sphinxupquote{psyclone.psyir.symbols.}}\sphinxbfcode{\sphinxupquote{StaticInterface}}}
\pysigstopsignatures
\sphinxAtStartPar
The symbol contains data that is kept alive through the execution
of the program.

\end{fulllineitems}


\item {} \index{CommonBlockInterface (class in psyclone.psyir.symbols)@\spxentry{CommonBlockInterface}\spxextra{class in psyclone.psyir.symbols}}

\begin{fulllineitems}
\phantomsection\label{\detokenize{psyir:psyclone.psyir.symbols.CommonBlockInterface}}
\pysigstartsignatures
\pysigline{\sphinxbfcode{\sphinxupquote{class\DUrole{w,w}{  }}}\sphinxcode{\sphinxupquote{psyclone.psyir.symbols.}}\sphinxbfcode{\sphinxupquote{CommonBlockInterface}}}
\pysigstopsignatures
\sphinxAtStartPar
A symbol declared in the local scope but acts as a global that
can be accessed by any scope referencing the same CommonBlock name.

\end{fulllineitems}


\item {} \index{UnresolvedInterface (class in psyclone.psyir.symbols)@\spxentry{UnresolvedInterface}\spxextra{class in psyclone.psyir.symbols}}

\begin{fulllineitems}
\phantomsection\label{\detokenize{psyir:psyclone.psyir.symbols.UnresolvedInterface}}
\pysigstartsignatures
\pysigline{\sphinxbfcode{\sphinxupquote{class\DUrole{w,w}{  }}}\sphinxcode{\sphinxupquote{psyclone.psyir.symbols.}}\sphinxbfcode{\sphinxupquote{UnresolvedInterface}}}
\pysigstopsignatures
\sphinxAtStartPar
We have a symbol but we don’t know where it is declared.

\end{fulllineitems}


\item {} \index{UnknownInterface (class in psyclone.psyir.symbols)@\spxentry{UnknownInterface}\spxextra{class in psyclone.psyir.symbols}}

\begin{fulllineitems}
\phantomsection\label{\detokenize{psyir:psyclone.psyir.symbols.UnknownInterface}}
\pysigstartsignatures
\pysigline{\sphinxbfcode{\sphinxupquote{class\DUrole{w,w}{  }}}\sphinxcode{\sphinxupquote{psyclone.psyir.symbols.}}\sphinxbfcode{\sphinxupquote{UnknownInterface}}}
\pysigstopsignatures
\sphinxAtStartPar
We have a symbol with a declaration but PSyclone does not support its
attributes.

\end{fulllineitems}


\item {} \index{PreprocessorInterface (class in psyclone.psyir.symbols)@\spxentry{PreprocessorInterface}\spxextra{class in psyclone.psyir.symbols}}

\begin{fulllineitems}
\phantomsection\label{\detokenize{psyir:psyclone.psyir.symbols.PreprocessorInterface}}
\pysigstartsignatures
\pysigline{\sphinxbfcode{\sphinxupquote{class\DUrole{w,w}{  }}}\sphinxcode{\sphinxupquote{psyclone.psyir.symbols.}}\sphinxbfcode{\sphinxupquote{PreprocessorInterface}}}
\pysigstopsignatures
\sphinxAtStartPar
The symbol exists in the file through compiler macros or preprocessor
directives.

\sphinxAtStartPar
Note that this is different from UnresolvedInterface because the backend
will not check if is importing statements that could bring them into
scope.

\end{fulllineitems}


\end{itemize}


\section{Creating PSyIR}
\label{\detokenize{psyir:creating-psyir}}

\subsection{Symbol names}
\label{\detokenize{psyir:symbol-names}}
\sphinxAtStartPar
PSyIR symbol names can be specified by a user. For example:

\begin{sphinxVerbatim}[commandchars=\\\{\}]
\PYG{n}{var\PYGZus{}name} \PYG{o}{=} \PYG{l+s+s2}{\PYGZdq{}}\PYG{l+s+s2}{my\PYGZus{}name}\PYG{l+s+s2}{\PYGZdq{}}
\PYG{n}{symbol\PYGZus{}table} \PYG{o}{=} \PYG{n}{SymbolTable}\PYG{p}{(}\PYG{p}{)}
\PYG{n}{data} \PYG{o}{=} \PYG{n}{DataSymbol}\PYG{p}{(}\PYG{n}{var\PYGZus{}name}\PYG{p}{,} \PYG{n}{REAL\PYGZus{}TYPE}\PYG{p}{)}
\PYG{n}{symbol\PYGZus{}table}\PYG{o}{.}\PYG{n}{add}\PYG{p}{(}\PYG{n}{data}\PYG{p}{)}
\PYG{n}{reference} \PYG{o}{=} \PYG{n}{Reference}\PYG{p}{(}\PYG{n}{data}\PYG{p}{)}
\end{sphinxVerbatim}

\sphinxAtStartPar
However, the \sphinxcode{\sphinxupquote{SymbolTable}} \sphinxcode{\sphinxupquote{add()}} method will raise an exception if a
user tries to add a symbol with the same name as a symbol already existing
in the symbol table.

\sphinxAtStartPar
Alternatively, the \sphinxcode{\sphinxupquote{SymbolTable}} also provides the \sphinxcode{\sphinxupquote{new\_symbol()}} method
(see Section {\hyperref[\detokenize{psyir:symbol-label}]{\sphinxcrossref{\DUrole{std,std-ref}{Symbols and Symbol Tables}}}} for more details) that uses a new distinct
name from any existing names in the symbol table. By default the generated
name is the value \sphinxcode{\sphinxupquote{PSYIR\_ROOT\_NAME}} variable specified in the \sphinxcode{\sphinxupquote{DEFAULT}}
section of the PSyclone config file, followed by an optional “\_” and
an integer. For example, the following code:

\begin{sphinxVerbatim}[commandchars=\\\{\}]
\PYG{k+kn}{from}\PYG{+w}{ }\PYG{n+nn}{psyclone}\PYG{n+nn}{.}\PYG{n+nn}{psyir}\PYG{n+nn}{.}\PYG{n+nn}{symbols}\PYG{+w}{ }\PYG{k+kn}{import} \PYG{n}{SymbolTable}
\PYG{n}{symbol\PYGZus{}table} \PYG{o}{=} \PYG{n}{SymbolTable}\PYG{p}{(}\PYG{p}{)}
\PYG{k}{for} \PYG{n}{i} \PYG{o+ow}{in} \PYG{n+nb}{range}\PYG{p}{(}\PYG{l+m+mi}{0}\PYG{p}{,} \PYG{l+m+mi}{3}\PYG{p}{)}\PYG{p}{:}
    \PYG{n}{var\PYGZus{}name} \PYG{o}{=} \PYG{n}{symbol\PYGZus{}table}\PYG{o}{.}\PYG{n}{new\PYGZus{}symbol}\PYG{p}{(}\PYG{p}{)}\PYG{o}{.}\PYG{n}{name}
    \PYG{n+nb}{print}\PYG{p}{(}\PYG{n}{var\PYGZus{}name}\PYG{p}{)}
\end{sphinxVerbatim}

\sphinxAtStartPar
gives the following output:

\begin{sphinxVerbatim}[commandchars=\\\{\}]
\PYG{n}{psyir\PYGZus{}tmp}
\PYG{n}{psyir\PYGZus{}tmp\PYGZus{}0}
\PYG{n}{psyir\PYGZus{}tmp\PYGZus{}1}
\end{sphinxVerbatim}

\sphinxAtStartPar
As the root name (\sphinxcode{\sphinxupquote{psyir\_tmp}} in the example above) is specified in
PSyclone’s config file it can be set to whatever the user wants.

\begin{sphinxadmonition}{note}{Note:}
\sphinxAtStartPar
The particular format used to create a unique name is the
responsibility of the SymbolTable class and may change in the
future.
\end{sphinxadmonition}

\sphinxAtStartPar
A user might want to create a name that has some meaning in the
context in which it is used e.g. \sphinxcode{\sphinxupquote{idx}} for an index, \sphinxcode{\sphinxupquote{i}} for an
iterator, or \sphinxcode{\sphinxupquote{temp}} for a temperature field. To support more
readable names, the \sphinxcode{\sphinxupquote{new\_symbol()}} method allows the user to specify a
root name as an argument to the method which then takes the place of
the default root name. For example, the following code:

\begin{sphinxVerbatim}[commandchars=\\\{\}]
\PYG{k+kn}{from}\PYG{+w}{ }\PYG{n+nn}{psyclone}\PYG{n+nn}{.}\PYG{n+nn}{psyir}\PYG{n+nn}{.}\PYG{n+nn}{symbols}\PYG{+w}{ }\PYG{k+kn}{import} \PYG{n}{SymbolTable}
\PYG{n}{symbol\PYGZus{}table} \PYG{o}{=} \PYG{n}{SymbolTable}\PYG{p}{(}\PYG{p}{)}
\PYG{k}{for} \PYG{n}{i} \PYG{o+ow}{in} \PYG{n+nb}{range}\PYG{p}{(}\PYG{l+m+mi}{0}\PYG{p}{,} \PYG{l+m+mi}{3}\PYG{p}{)}\PYG{p}{:}
    \PYG{n}{var\PYGZus{}name} \PYG{o}{=} \PYG{n}{symbol\PYGZus{}table}\PYG{o}{.}\PYG{n}{new\PYGZus{}symbol}\PYG{p}{(}\PYG{n}{root\PYGZus{}name}\PYG{o}{=}\PYG{l+s+s2}{\PYGZdq{}}\PYG{l+s+s2}{something}\PYG{l+s+s2}{\PYGZdq{}}\PYG{p}{)}
    \PYG{n+nb}{print}\PYG{p}{(}\PYG{n}{var\PYGZus{}name}\PYG{p}{)}
\end{sphinxVerbatim}

\sphinxAtStartPar
gives the following output:

\begin{sphinxVerbatim}[commandchars=\\\{\}]
\PYG{n}{something}
\PYG{n}{something\PYGZus{}0}
\PYG{n}{something\PYGZus{}1}
\end{sphinxVerbatim}

\sphinxAtStartPar
By default, \sphinxcode{\sphinxupquote{new\_symbol()}} creates generic symbols, but often the user
will want to specify a Symbol subclass with some given parameters. The
\sphinxcode{\sphinxupquote{new\_symbol()}} method accepts a \sphinxcode{\sphinxupquote{symbol\_type}} parameter to specify the
subclass.  Arguments for the constructor of that subclass may be supplied
as keyword arguments. For example, the following code:

\begin{sphinxVerbatim}[commandchars=\\\{\}]
\PYG{k+kn}{from}\PYG{+w}{ }\PYG{n+nn}{psyclone}\PYG{n+nn}{.}\PYG{n+nn}{psyir}\PYG{n+nn}{.}\PYG{n+nn}{symbols}\PYG{+w}{ }\PYG{k+kn}{import} \PYG{n}{SymbolTable}\PYG{p}{,} \PYG{n}{DataSymbol}\PYG{p}{,} \PYG{n}{REAL\PYGZus{}TYPE}
\PYG{n}{symbol\PYGZus{}table} \PYG{o}{=} \PYG{n}{SymbolTable}\PYG{p}{(}\PYG{p}{)}
\PYG{n}{symbol\PYGZus{}table}\PYG{o}{.}\PYG{n}{new\PYGZus{}symbol}\PYG{p}{(}\PYG{n}{root\PYGZus{}name}\PYG{o}{=}\PYG{l+s+s2}{\PYGZdq{}}\PYG{l+s+s2}{something}\PYG{l+s+s2}{\PYGZdq{}}\PYG{p}{,}
                        \PYG{n}{symbol\PYGZus{}type}\PYG{o}{=}\PYG{n}{DataSymbol}\PYG{p}{,}
                        \PYG{n}{datatype}\PYG{o}{=}\PYG{n}{REAL\PYGZus{}TYPE}\PYG{p}{,}
                        \PYG{n}{is\PYGZus{}constant}\PYG{o}{=}\PYG{k+kc}{True}\PYG{p}{,}
                        \PYG{n}{initial\PYGZus{}value}\PYG{o}{=}\PYG{l+m+mi}{3}\PYG{p}{)}
\end{sphinxVerbatim}

\sphinxAtStartPar
declares a symbol named “something” of REAL\_TYPE datatype where the
\sphinxcode{\sphinxupquote{is\_constant}} and \sphinxcode{\sphinxupquote{initial\_value}} arguments will be passed to the
DataSymbol constructor.

\sphinxAtStartPar
An example of using the \sphinxcode{\sphinxupquote{new\_symbol()}} method can be found in the
PSyclone \sphinxcode{\sphinxupquote{examples/psyir}} directory.


\subsection{Nodes}
\label{\detokenize{psyir:nodes}}
\sphinxAtStartPar
PSyIR nodes are connected together via parent and child methods
provided by the \sphinxcode{\sphinxupquote{Node}} baseclass.

\sphinxAtStartPar
These nodes can be created in isolation and then connected
together. For example:

\begin{sphinxVerbatim}[commandchars=\\\{\}]
\PYG{n}{assignment} \PYG{o}{=} \PYG{n}{Assignment}\PYG{p}{(}\PYG{p}{)}
\PYG{n}{literal} \PYG{o}{=} \PYG{n}{Literal}\PYG{p}{(}\PYG{l+s+s2}{\PYGZdq{}}\PYG{l+s+s2}{0.0}\PYG{l+s+s2}{\PYGZdq{}}\PYG{p}{,} \PYG{n}{REAL\PYGZus{}TYPE}\PYG{p}{)}
\PYG{n}{reference} \PYG{o}{=} \PYG{n}{Reference}\PYG{p}{(}\PYG{n}{symbol}\PYG{p}{)}
\PYG{n}{assignment}\PYG{o}{.}\PYG{n}{children} \PYG{o}{=} \PYG{p}{[}\PYG{n}{reference}\PYG{p}{,} \PYG{n}{literal}\PYG{p}{]}
\end{sphinxVerbatim}

\sphinxAtStartPar
However, as connections get more complicated, creating the correct
connections can become difficult to manage and error prone. Further,
in some cases children must be collected together within a
\sphinxcode{\sphinxupquote{Schedule}} (e.g. for \sphinxcode{\sphinxupquote{IfBlock}}, \sphinxcode{\sphinxupquote{Loop}} and \sphinxcode{\sphinxupquote{WhileLoop}}).

\sphinxAtStartPar
To simplify this complexity, each of the Kernel\sphinxhyphen{}layer nodes which
contain other nodes have a static \sphinxcode{\sphinxupquote{create}} method which helps
construct the PSyIR using a bottom up approach. Using this method, the
above example then becomes:

\begin{sphinxVerbatim}[commandchars=\\\{\}]
\PYG{n}{literal} \PYG{o}{=} \PYG{n}{Literal}\PYG{p}{(}\PYG{l+s+s2}{\PYGZdq{}}\PYG{l+s+s2}{0.0}\PYG{l+s+s2}{\PYGZdq{}}\PYG{p}{,} \PYG{n}{REAL\PYGZus{}TYPE}\PYG{p}{)}
\PYG{n}{reference} \PYG{o}{=} \PYG{n}{Reference}\PYG{p}{(}\PYG{n}{symbol}\PYG{p}{)}
\PYG{n}{assignment} \PYG{o}{=} \PYG{n}{Assignment}\PYG{o}{.}\PYG{n}{create}\PYG{p}{(}\PYG{n}{reference}\PYG{p}{,} \PYG{n}{literal}\PYG{p}{)}
\end{sphinxVerbatim}

\sphinxAtStartPar
Creating the PSyIR to represent a complicated access of a member of a
structure is best performed using the \sphinxcode{\sphinxupquote{create()}} method of the
appropriate \sphinxcode{\sphinxupquote{Reference}} subclass. For a relatively straightforward
access such as (the Fortran) \sphinxcode{\sphinxupquote{field1\%region\%nx}}, this would be:

\begin{sphinxVerbatim}[commandchars=\\\{\}]
\PYG{k+kn}{from}\PYG{+w}{ }\PYG{n+nn}{psyclone}\PYG{n+nn}{.}\PYG{n+nn}{psyir}\PYG{n+nn}{.}\PYG{n+nn}{nodes}\PYG{+w}{ }\PYG{k+kn}{import} \PYG{n}{StructureReference}
\PYG{n}{fld\PYGZus{}sym} \PYG{o}{=} \PYG{n}{symbol\PYGZus{}table}\PYG{o}{.}\PYG{n}{lookup}\PYG{p}{(}\PYG{l+s+s2}{\PYGZdq{}}\PYG{l+s+s2}{field1}\PYG{l+s+s2}{\PYGZdq{}}\PYG{p}{)}
\PYG{n}{ref} \PYG{o}{=} \PYG{n}{StructureReference}\PYG{o}{.}\PYG{n}{create}\PYG{p}{(}\PYG{n}{fld\PYGZus{}sym}\PYG{p}{,} \PYG{p}{[}\PYG{l+s+s2}{\PYGZdq{}}\PYG{l+s+s2}{region}\PYG{l+s+s2}{\PYGZdq{}}\PYG{p}{,} \PYG{l+s+s2}{\PYGZdq{}}\PYG{l+s+s2}{nx}\PYG{l+s+s2}{\PYGZdq{}}\PYG{p}{]}\PYG{p}{)}
\end{sphinxVerbatim}

\sphinxAtStartPar
where \sphinxcode{\sphinxupquote{symbol\_table}} is assumed to be a pre\sphinxhyphen{}populated Symbol Table
containing an entry for “field1”.

\sphinxAtStartPar
A more complicated access involving arrays of structures such as
\sphinxcode{\sphinxupquote{field1\%sub\_grids(idx, 1)\%nx}} would be constructed as:

\begin{sphinxVerbatim}[commandchars=\\\{\}]
\PYG{k+kn}{from}\PYG{+w}{ }\PYG{n+nn}{psyclone}\PYG{n+nn}{.}\PYG{n+nn}{psyir}\PYG{n+nn}{.}\PYG{n+nn}{symbols}\PYG{+w}{ }\PYG{k+kn}{import} \PYG{n}{INTEGER\PYGZus{}TYPE}
\PYG{k+kn}{from}\PYG{+w}{ }\PYG{n+nn}{psyclone}\PYG{n+nn}{.}\PYG{n+nn}{psyir}\PYG{n+nn}{.}\PYG{n+nn}{nodes}\PYG{+w}{ }\PYG{k+kn}{import} \PYG{n}{StructureReference}\PYG{p}{,} \PYG{n}{Reference}\PYG{p}{,} \PYG{n}{Literal}
\PYG{n}{idx\PYGZus{}sym} \PYG{o}{=} \PYG{n}{symbol\PYGZus{}table}\PYG{o}{.}\PYG{n}{lookup}\PYG{p}{(}\PYG{l+s+s2}{\PYGZdq{}}\PYG{l+s+s2}{idx}\PYG{l+s+s2}{\PYGZdq{}}\PYG{p}{)}
\PYG{n}{fld\PYGZus{}sym} \PYG{o}{=} \PYG{n}{symbol\PYGZus{}table}\PYG{o}{.}\PYG{n}{lookup}\PYG{p}{(}\PYG{l+s+s2}{\PYGZdq{}}\PYG{l+s+s2}{field1}\PYG{l+s+s2}{\PYGZdq{}}\PYG{p}{)}
\PYG{n}{ref} \PYG{o}{=} \PYG{n}{StructureReference}\PYG{o}{.}\PYG{n}{create}\PYG{p}{(}\PYG{n}{fld\PYGZus{}sym}\PYG{p}{,}
    \PYG{p}{[}\PYG{p}{(}\PYG{l+s+s2}{\PYGZdq{}}\PYG{l+s+s2}{sub\PYGZus{}grids}\PYG{l+s+s2}{\PYGZdq{}}\PYG{p}{,} \PYG{p}{[}\PYG{n}{Reference}\PYG{p}{(}\PYG{n}{idx\PYGZus{}sym}\PYG{p}{)}\PYG{p}{,} \PYG{n}{Literal}\PYG{p}{(}\PYG{l+s+s2}{\PYGZdq{}}\PYG{l+s+s2}{1}\PYG{l+s+s2}{\PYGZdq{}}\PYG{p}{,} \PYG{n}{INTEGER\PYGZus{}TYPE}\PYG{p}{)}\PYG{p}{]}\PYG{p}{)}\PYG{p}{,}
     \PYG{l+s+s2}{\PYGZdq{}}\PYG{l+s+s2}{nx}\PYG{l+s+s2}{\PYGZdq{}}\PYG{p}{]}\PYG{p}{)}
\end{sphinxVerbatim}

\sphinxAtStartPar
Note that the list of quantities passed to the \sphinxcode{\sphinxupquote{create()}} method now
contains a 2\sphinxhyphen{}tuple in order to describe the array access.

\sphinxAtStartPar
More examples of using this approach can be found in the PSyclone
\sphinxcode{\sphinxupquote{examples/psyir}} directory.


\section{Comparing PSyIR nodes}
\label{\detokenize{psyir:comparing-psyir-nodes}}
\sphinxAtStartPar
The \sphinxcode{\sphinxupquote{==}} (equality) operator for PSyIR nodes performs a specialised equality check
to compare the value of each node. This is also useful when comparing entire
subtrees since the equality operator automatically recurses through the children
and compares each child with the appropriate equality semantics, e.g.

\begin{sphinxVerbatim}[commandchars=\\\{\}]
\PYG{c+c1}{\PYGZsh{} Is the loop upper bound expression exactly the same?}
\PYG{k}{if} \PYG{n}{loop1}\PYG{o}{.}\PYG{n}{stop\PYGZus{}expr} \PYG{o}{==} \PYG{n}{loop2}\PYG{o}{.}\PYG{n}{stop\PYGZus{}expr}\PYG{p}{:}
        \PYG{n+nb}{print}\PYG{p}{(}\PYG{l+s+s2}{\PYGZdq{}}\PYG{l+s+s2}{Same upper bound!}\PYG{l+s+s2}{\PYGZdq{}}\PYG{p}{)}
\end{sphinxVerbatim}

\sphinxAtStartPar
The equality operator will handle expressions like \sphinxcode{\sphinxupquote{my\_array\%my\_field(:3)}} with the
derived type fields and the range components automatically, but it cannot handle
symbolically equivalent fields, i.e. \sphinxcode{\sphinxupquote{my\_array\%my\_field(:3) != my\_array\%my\_field(:2+1)}}.

\sphinxAtStartPar
Annotations and code comments are ignored in the equality comparison since they don’t
alter the semantic meaning of the code. So these two statements compare to True:

\begin{sphinxVerbatim}[commandchars=\\\{\}]
\PYG{n}{a}\PYG{+w}{ }\PYG{o}{=}\PYG{+w}{ }\PYG{n}{a}\PYG{+w}{ }\PYG{o}{+}\PYG{+w}{ }\PYG{l+m+mi}{1}
\PYG{n}{a}\PYG{+w}{ }\PYG{o}{=}\PYG{+w}{ }\PYG{n}{a}\PYG{+w}{ }\PYG{o}{+}\PYG{+w}{ }\PYG{l+m+mi}{1}\PYG{+w}{ }\PYG{c}{!Increases a by 1}
\end{sphinxVerbatim}

\sphinxAtStartPar
Sometimes there are cases where one really means to check for the specific instance
of a node. In this case, Python provides the \sphinxcode{\sphinxupquote{is}} operator, e.g.

\begin{sphinxVerbatim}[commandchars=\\\{\}]
\PYG{c+c1}{\PYGZsh{} Is the self instance part of this routine?}
\PYG{n}{is\PYGZus{}here} \PYG{o}{=} \PYG{n+nb}{any}\PYG{p}{(}\PYG{n}{node} \PYG{o+ow}{is} \PYG{n+nb+bp}{self} \PYG{k}{for} \PYG{n}{node} \PYG{o+ow}{in} \PYG{n}{routine}\PYG{o}{.}\PYG{n}{walk}\PYG{p}{(}\PYG{n}{Node}\PYG{p}{)}\PYG{p}{)}
\end{sphinxVerbatim}

\sphinxAtStartPar
Additionally, PSyIR nodes cannot be used as map keys or similar. The easiest way
to do this is just use the id as the key:

\begin{sphinxVerbatim}[commandchars=\\\{\}]
\PYG{n}{node\PYGZus{}map} \PYG{o}{=} \PYG{p}{\PYGZob{}}\PYG{p}{\PYGZcb{}}
\PYG{n}{node\PYGZus{}map}\PYG{p}{[}\PYG{n+nb}{id}\PYG{p}{(}\PYG{n}{mynode}\PYG{p}{)}\PYG{p}{]} \PYG{o}{=} \PYG{l+s+s2}{\PYGZdq{}}\PYG{l+s+s2}{element}\PYG{l+s+s2}{\PYGZdq{}}
\end{sphinxVerbatim}


\section{Modifying the PSyIR}
\label{\detokenize{psyir:modifying-the-psyir}}
\sphinxAtStartPar
Once we have a complete PSyIR AST there are 2 ways to modify its contents
and/or structure: by applying transformations (see next section
{\hyperref[\detokenize{transformations:transformations}]{\sphinxcrossref{\DUrole{std,std-ref}{Transformations}}}}), or by direct PSyIR API methods. This section
describes some of the methods that the PSyIR classes provide to
modify the PSyIR AST in a consistent way (e.g. without breaking its many
internal references). Some complete examples of modifying the PSyIR can be found in the
PSyclone \sphinxcode{\sphinxupquote{examples/psyir/modify.py}} script.

\sphinxAtStartPar
The rest of this section introduces examples of the available direct PSyIR
modification methods.


\subsection{Renaming symbols}
\label{\detokenize{psyir:renaming-symbols}}
\sphinxAtStartPar
The symbol table provides the method \sphinxcode{\sphinxupquote{rename\_symbol()}} that given a symbol
and an unused name will rename the symbol. The symbol renaming will affect
all the references in the PSyIR AST to that symbol. For example, the PSyIR
representing the following Fortran code:

\begin{sphinxVerbatim}[commandchars=\\\{\}]
\PYG{k}{subroutine }\PYG{n}{work}\PYG{p}{(}\PYG{n}{psyir\PYGZus{}tmp}\PYG{p}{)}
\PYG{+w}{    }\PYG{k+kt}{real}\PYG{p}{,}\PYG{+w}{ }\PYG{k}{intent}\PYG{p}{(}\PYG{n}{inout}\PYG{p}{)}\PYG{+w}{ }\PYG{k+kd}{::}\PYG{+w}{ }\PYG{n}{psyir\PYGZus{}tmp}
\PYG{+w}{    }\PYG{n}{psyir\PYGZus{}tmp}\PYG{o}{=}\PYG{l+m+mf}{0.0}
\PYG{k}{end }\PYG{k}{subroutine}
\end{sphinxVerbatim}

\sphinxAtStartPar
could be modified by the following PSyIR statements:

\begin{sphinxVerbatim}[commandchars=\\\{\}]
\PYG{n}{symbol} \PYG{o}{=} \PYG{n}{symbol\PYGZus{}table}\PYG{o}{.}\PYG{n}{lookup}\PYG{p}{(}\PYG{l+s+s2}{\PYGZdq{}}\PYG{l+s+s2}{psyir\PYGZus{}tmp}\PYG{l+s+s2}{\PYGZdq{}}\PYG{p}{)}
\PYG{n}{symbol\PYGZus{}table}\PYG{o}{.}\PYG{n}{rename\PYGZus{}symbol}\PYG{p}{(}\PYG{n}{tmp\PYGZus{}symbol}\PYG{p}{,} \PYG{l+s+s2}{\PYGZdq{}}\PYG{l+s+s2}{new\PYGZus{}variable}\PYG{l+s+s2}{\PYGZdq{}}\PYG{p}{)}
\end{sphinxVerbatim}

\sphinxAtStartPar
which would result in the following Fortran output code:

\begin{sphinxVerbatim}[commandchars=\\\{\}]
\PYG{k}{subroutine }\PYG{n}{work}\PYG{p}{(}\PYG{n}{new\PYGZus{}variable}\PYG{p}{)}
\PYG{+w}{    }\PYG{k+kt}{real}\PYG{p}{,}\PYG{+w}{ }\PYG{k}{intent}\PYG{p}{(}\PYG{n}{inout}\PYG{p}{)}\PYG{+w}{ }\PYG{k+kd}{::}\PYG{+w}{ }\PYG{n}{new\PYGZus{}variable}
\PYG{+w}{    }\PYG{n}{new\PYGZus{}variable}\PYG{o}{=}\PYG{l+m+mf}{0.0}
\PYG{k}{end }\PYG{k}{subroutine}
\end{sphinxVerbatim}


\subsection{Specialising symbols}
\label{\detokenize{psyir:specialising-symbols}}
\sphinxAtStartPar
The Symbol class provides the method \sphinxcode{\sphinxupquote{specialise()}} that given a
subclass of Symbol will change the Symbol instance to the specified
subclass. If the subclass has any additional properties then these
would need to be set explicitly.

\begin{sphinxVerbatim}[commandchars=\\\{\}]
\PYG{n}{symbol} \PYG{o}{=} \PYG{n}{Symbol}\PYG{p}{(}\PYG{l+s+s2}{\PYGZdq{}}\PYG{l+s+s2}{name}\PYG{l+s+s2}{\PYGZdq{}}\PYG{p}{)}
\PYG{n}{symbol}\PYG{o}{.}\PYG{n}{specialise}\PYG{p}{(}\PYG{n}{RoutineSymbol}\PYG{p}{)}
\PYG{c+c1}{\PYGZsh{} Symbol is now a RoutineSymbol}
\end{sphinxVerbatim}

\sphinxAtStartPar
This method is useful as it allows the class of a symbol to be changed
without affecting any references to it.


\subsection{Replacing PSyIR nodes}
\label{\detokenize{psyir:replacing-psyir-nodes}}
\sphinxAtStartPar
In certain cases one might want to replace a node in a PSyIR tree with
another node. All nodes provide the \sphinxtitleref{replace\_with()} method to replace
the node and its descendants with another given node and its
descendants.

\begin{sphinxVerbatim}[commandchars=\\\{\}]
\PYG{n}{node}\PYG{o}{.}\PYG{n}{replace\PYGZus{}with}\PYG{p}{(}\PYG{n}{new\PYGZus{}node}\PYG{p}{)}
\end{sphinxVerbatim}

\sphinxAtStartPar
When the node being replaced is part of a named context (in Calls or
Operations) the name of the argument is conserved by default. For example

\begin{sphinxVerbatim}[commandchars=\\\{\}]
\PYG{k}{call }\PYG{n}{named\PYGZus{}subroutine}\PYG{p}{(}\PYG{n}{name1}\PYG{o}{=}\PYG{l+m+mi}{1}\PYG{p}{)}
\end{sphinxVerbatim}

\begin{sphinxVerbatim}[commandchars=\\\{\}]
\PYG{n}{call}\PYG{o}{.}\PYG{n}{children}\PYG{p}{[}\PYG{l+m+mi}{0}\PYG{p}{]}\PYG{o}{.}\PYG{n}{replace\PYGZus{}with}\PYG{p}{(}\PYG{n}{Literal}\PYG{p}{(}\PYG{l+s+s1}{\PYGZsq{}}\PYG{l+s+s1}{2}\PYG{l+s+s1}{\PYGZsq{}}\PYG{p}{,} \PYG{n}{INTEGER\PYGZus{}TYPE}\PYG{p}{)}\PYG{p}{)}
\end{sphinxVerbatim}

\sphinxAtStartPar
will become:

\begin{sphinxVerbatim}[commandchars=\\\{\}]
\PYG{k}{call }\PYG{n}{named\PYGZus{}subroutine}\PYG{p}{(}\PYG{n}{name1}\PYG{o}{=}\PYG{l+m+mi}{2}\PYG{p}{)}
\end{sphinxVerbatim}

\sphinxAtStartPar
This behaviour can be changed with the \sphinxtitleref{keep\_name\_in\_context} parameter.

\begin{sphinxVerbatim}[commandchars=\\\{\}]
\PYG{n}{call}\PYG{o}{.}\PYG{n}{children}\PYG{p}{[}\PYG{l+m+mi}{0}\PYG{p}{]}\PYG{o}{.}\PYG{n}{replace\PYGZus{}with}\PYG{p}{(}
    \PYG{n}{Literal}\PYG{p}{(}\PYG{l+s+s1}{\PYGZsq{}}\PYG{l+s+s1}{3}\PYG{l+s+s1}{\PYGZsq{}}\PYG{p}{,} \PYG{n}{INTEGER\PYGZus{}TYPE}\PYG{p}{)}\PYG{p}{,}
    \PYG{n}{keep\PYGZus{}name\PYGZus{}in\PYGZus{}context}\PYG{o}{=}\PYG{k+kc}{False}
\PYG{p}{)}
\end{sphinxVerbatim}

\sphinxAtStartPar
will become:

\begin{sphinxVerbatim}[commandchars=\\\{\}]
\PYG{k}{call }\PYG{n}{named\PYGZus{}subroutine}\PYG{p}{(}\PYG{l+m+mi}{3}\PYG{p}{)}
\end{sphinxVerbatim}


\subsection{Detaching PSyIR nodes}
\label{\detokenize{psyir:detaching-psyir-nodes}}
\sphinxAtStartPar
Sometimes we just may wish to detach a certain PSyIR subtree in order to remove
it from the root tree but we don’t want to delete it altogether, as it may
be re\sphinxhyphen{}inserted again in another location. To achieve this, all nodes
provide the detach method:

\begin{sphinxVerbatim}[commandchars=\\\{\}]
\PYG{n}{tmp} \PYG{o}{=} \PYG{n}{node}\PYG{o}{.}\PYG{n}{detach}\PYG{p}{(}\PYG{p}{)}
\end{sphinxVerbatim}


\subsection{Copying nodes}
\label{\detokenize{psyir:copying-nodes}}
\sphinxAtStartPar
Copying a PSyIR node and its children is often useful in order to avoid
repeating the creation of similar PSyIR subtrees. The result of the copy
allows the modification of the original and the copied subtrees independently,
without altering the other subtree. Note that this is not equivalent to the
Python \sphinxcode{\sphinxupquote{copy}} or \sphinxcode{\sphinxupquote{deepcopy}} functionality provided in the \sphinxcode{\sphinxupquote{copy}} library.
This method performs a bespoke copy operation where some components of the
tree, like children, are recursively copied, while others, like the top\sphinxhyphen{}level
parent reference are not.

\begin{sphinxVerbatim}[commandchars=\\\{\}]
\PYG{n}{new\PYGZus{}node} \PYG{o}{=} \PYG{n}{node}\PYG{o}{.}\PYG{n}{copy}\PYG{p}{(}\PYG{p}{)}
\end{sphinxVerbatim}


\subsection{Named arguments}
\label{\detokenize{psyir:named-arguments}}
\sphinxAtStartPar
The \sphinxtitleref{Call} node (and its sub\sphinxhyphen{}classes) support named arguments.

\sphinxAtStartPar
Named arguments can be set or modified via the \sphinxtitleref{create()},
\sphinxtitleref{append\_named\_arg()}, \sphinxtitleref{insert\_named\_arg()} or \sphinxtitleref{replace\_named\_arg()}
methods.

\sphinxAtStartPar
If an argument is inserted directly (via the children list) then it is
assumed that this is not a named argument. If the top node of an
argument is replaced by removing and inserting a new node then it is
assumed that this argument is no longer a named argument. If it is
replaced with the \sphinxtitleref{replace\_with} method, it has a \sphinxtitleref{keep\_name\_in\_context}
argument to choose the desired behaviour (defaults to True).
If arguments are re\sphinxhyphen{}ordered then the names follow the re\sphinxhyphen{}ordering.

\sphinxAtStartPar
The names of named arguments can be accessed via the \sphinxtitleref{argument\_names}
property. This list has an entry for each argument and either contains
a name or None (if this is not a named argument).

\sphinxAtStartPar
The PSyIR does not constrain which arguments are specified as being
named and what those names are. It is the developer’s responsibility
to make sure that these names are consistent with any intrinsics that
will be generated by the back\sphinxhyphen{}end. In the future, it is expected that
the PSyIR will know about the number and type of arguments expected by
Operation nodes, beyond simply being unary, binary or nary.

\sphinxAtStartPar
One restriction that Fortran has (but the PSyIR does not) is that all
named arguments should be at the end of the argument list. If this is
not the case then the Fortran backend writer will raise an exception.

\sphinxstepscope


\chapter{Transformations}
\label{\detokenize{transformations:transformations}}\label{\detokenize{transformations:id1}}\label{\detokenize{transformations::doc}}
\sphinxAtStartPar
As discussed in the previous section, transformations can be applied
to the PSyIR to modify it. Typically transformations will be used to
optimise the provided source file, or the PSy and/or Kernel layer(s)
in the PSyKAl DSLs, for a particular architecture. However, transformations
could be added for other reasons, such as to aid debugging or for
performance monitoring.


\section{PSyclone User Scripts}
\label{\detokenize{transformations:psyclone-user-scripts}}\label{\detokenize{transformations:sec-transformations-script}}
\sphinxAtStartPar
A convenient way to apply transformations to a codebase is through the
{\hyperref[\detokenize{psyclone_command:psyclone-command}]{\sphinxcrossref{\DUrole{std,std-ref}{The psyclone command}}}} tool, which has the optional \sphinxcode{\sphinxupquote{\sphinxhyphen{}s \textless{}SCRIPT\_NAME\textgreater{}}}
flag that allows users to specify a script file to programatically modify
input code:

\begin{sphinxVerbatim}[commandchars=\\\{\}]
\PYG{o}{\PYGZgt{}} \PYG{n}{psyclone} \PYG{o}{\PYGZhy{}}\PYG{n}{s} \PYG{n}{optimise}\PYG{o}{.}\PYG{n}{py} \PYG{n}{input\PYGZus{}source}\PYG{o}{.}\PYG{n}{f90}
\end{sphinxVerbatim}

\sphinxAtStartPar
In this case, the current directory is prepended to the Python search path
\sphinxstylestrong{PYTHONPATH} which will then be used to try to find the script file. Thus,
the search begins in the current directory and continues over any pre\sphinxhyphen{}existing
directories in the search path, failing if the file cannot be found.

\sphinxAtStartPar
Alternatively, script files may be specified with a path. In this case
the file must exist in the specified location. This location is then added to
the Python search path \sphinxstylestrong{PYTHONPATH} as before. For example:

\begin{sphinxVerbatim}[commandchars=\\\{\}]
\PYG{o}{\PYGZgt{}} \PYG{n}{psyclone} \PYG{o}{\PYGZhy{}}\PYG{n}{s} \PYG{o}{.}\PYG{o}{/}\PYG{n}{optimise}\PYG{o}{.}\PYG{n}{py} \PYG{n}{input\PYGZus{}source}\PYG{o}{.}\PYG{n}{f90}
\PYG{o}{\PYGZgt{}} \PYG{n}{psyclone} \PYG{o}{\PYGZhy{}}\PYG{n}{s} \PYG{o}{.}\PYG{o}{.}\PYG{o}{/}\PYG{n}{scripts}\PYG{o}{/}\PYG{n}{optimise}\PYG{o}{.}\PYG{n}{py} \PYG{n}{input\PYGZus{}source}\PYG{o}{.}\PYG{n}{f90}
\PYG{o}{\PYGZgt{}} \PYG{n}{psyclone} \PYG{o}{\PYGZhy{}}\PYG{n}{s} \PYG{o}{/}\PYG{n}{home}\PYG{o}{/}\PYG{n}{me}\PYG{o}{/}\PYG{n}{PSyclone}\PYG{o}{/}\PYG{n}{scripts}\PYG{o}{/}\PYG{n}{optimise}\PYG{o}{.}\PYG{n}{py} \PYG{n}{input\PYGZus{}source}\PYG{o}{.}\PYG{n}{f90}
\end{sphinxVerbatim}

\sphinxAtStartPar
A valid PSyclone user script file must contain a \sphinxstylestrong{trans} function which accepts
a {\hyperref[\detokenize{psyir:psyir-ug}]{\sphinxcrossref{\DUrole{std,std-ref}{PSyIR node}}}} representing the root of the psy\sphinxhyphen{}layer
code (as a FileConatainer):

\begin{sphinxVerbatim}[commandchars=\\\{\}]
\PYG{k}{def}\PYG{+w}{ }\PYG{n+nf}{trans}\PYG{p}{(}\PYG{n}{psyir}\PYG{p}{)}\PYG{p}{:}
    \PYG{c+c1}{\PYGZsh{} ...}
\end{sphinxVerbatim}

\sphinxAtStartPar
The example below adds an OpenMP directive to a specific PSyKAL kernel:

\begin{sphinxVerbatim}[commandchars=\\\{\}]
\PYG{k}{def}\PYG{+w}{ }\PYG{n+nf}{trans}\PYG{p}{(}\PYG{n}{psyir}\PYG{p}{)}\PYG{p}{:}
    \PYG{k+kn}{from}\PYG{+w}{ }\PYG{n+nn}{psyclone}\PYG{n+nn}{.}\PYG{n+nn}{transformations}\PYG{+w}{ }\PYG{k+kn}{import} \PYG{n}{OMPParallelLoopTrans}
    \PYG{k+kn}{from}\PYG{+w}{ }\PYG{n+nn}{psyclone}\PYG{n+nn}{.}\PYG{n+nn}{psyir}\PYG{n+nn}{.}\PYG{n+nn}{node}\PYG{+w}{ }\PYG{k+kn}{import} \PYG{n}{Routine}
    \PYG{k}{for} \PYG{n}{subroutine} \PYG{o+ow}{in} \PYG{n}{psyir}\PYG{o}{.}\PYG{n}{walk}\PYG{p}{(}\PYG{n}{Routine}\PYG{p}{)}\PYG{p}{:}
        \PYG{k}{if} \PYG{n}{subroutine}\PYG{o}{.}\PYG{n}{name} \PYG{o}{==} \PYG{l+s+s1}{\PYGZsq{}}\PYG{l+s+s1}{invoke\PYGZus{}0\PYGZus{}v3\PYGZus{}kernel\PYGZus{}type}\PYG{l+s+s1}{\PYGZsq{}}\PYG{p}{:}
            \PYG{n}{ol} \PYG{o}{=} \PYG{n}{OMPParallelLoopTrans}\PYG{p}{(}\PYG{p}{)}
            \PYG{n}{ol}\PYG{o}{.}\PYG{n}{apply}\PYG{p}{(}\PYG{n}{subroutine}\PYG{o}{.}\PYG{n}{children}\PYG{p}{[}\PYG{l+m+mi}{0}\PYG{p}{]}\PYG{p}{)}
\end{sphinxVerbatim}

\sphinxAtStartPar
The script may apply as many transformations as is required for the intended
optimisation, and may also apply transformations to all the routines (i.e. invokes
and/or kernels) contained within the provided tree.
The {\hyperref[\detokenize{tutorials_and_examples:examples}]{\sphinxcrossref{\DUrole{std,std-ref}{examples section}}}} provides a list of psyclone user scripts
and associated usage instructions for multiple applications.


\subsection{Script Global Variables}
\label{\detokenize{transformations:script-global-variables}}
\sphinxAtStartPar
In addition to the \sphinxcode{\sphinxupquote{trans}} function, there are special global variables that can be set
to control some of the behaviours of the front\sphinxhyphen{}end (before the optimisation function
is applied). These are:

\begin{sphinxVerbatim}[commandchars=\\\{\}]
\PYG{c+c1}{\PYGZsh{} List of all files that psyclone will skip processing}
\PYG{n}{FILES\PYGZus{}TO\PYGZus{}SKIP} \PYG{o}{=} \PYG{p}{[}\PYG{l+s+s2}{\PYGZdq{}}\PYG{l+s+s2}{boken\PYGZus{}file1.f90}\PYG{l+s+s2}{\PYGZdq{}}\PYG{p}{,} \PYG{l+s+s2}{\PYGZdq{}}\PYG{l+s+s2}{boken\PYGZus{}file2.f90}\PYG{l+s+s2}{\PYGZdq{}}\PYG{p}{]}

\PYG{c+c1}{\PYGZsh{} Boolean to decide whether PSyclone should chase external modules while}
\PYG{c+c1}{\PYGZsh{} creating a PSyIR tree in order to obtain better external symbol information.}
\PYG{c+c1}{\PYGZsh{} It can also be a list of module names for more precise control}
\PYG{n}{RESOLVE\PYGZus{}IMPORTS} \PYG{o}{=} \PYG{p}{[}\PYG{l+s+s2}{\PYGZdq{}}\PYG{l+s+s2}{relevant\PYGZus{}module1.f90}\PYG{l+s+s2}{\PYGZdq{}}\PYG{p}{,} \PYG{l+s+s2}{\PYGZdq{}}\PYG{l+s+s2}{relevant\PYGZus{}module2.f90}\PYG{l+s+s2}{\PYGZdq{}}\PYG{p}{]}

\PYG{k}{def}\PYG{+w}{ }\PYG{n+nf}{trans}\PYG{p}{(}\PYG{n}{psyir}\PYG{p}{)}\PYG{p}{:}
    \PYG{c+c1}{\PYGZsh{} ...}
\end{sphinxVerbatim}


\subsection{PSyKAl algorithm code transformations}
\label{\detokenize{transformations:psykal-algorithm-code-transformations}}
\sphinxAtStartPar
When using PSyKAl, the \sphinxcode{\sphinxupquote{trans}} functions is used to transform the PSy\sphinxhyphen{}layer (the
layer in charge of the Parallel\sphinxhyphen{}System and Loops traversal orders), however, a
second optional transformation entry point \sphinxcode{\sphinxupquote{trans\_alg}} can be provided to
directly transform the Algorithm\sphinxhyphen{}layer (this is currently only implemented for
GOcean, but in the future it will also affect the LFRic DSL).

\begin{sphinxVerbatim}[commandchars=\\\{\}]
\PYG{k}{def}\PYG{+w}{ }\PYG{n+nf}{trans\PYGZus{}alg}\PYG{p}{(}\PYG{n}{psyir}\PYG{p}{)}\PYG{p}{:}
    \PYG{c+c1}{\PYGZsh{} ...}
\end{sphinxVerbatim}

\sphinxAtStartPar
As with the \sphinxtitleref{trans()} function it is up to the script what it does with
the algorithm PSyIR. Note that the \sphinxtitleref{trans\_alg()} script is applied to
the algorithm layer before the PSy\sphinxhyphen{}layer is generated so any changes
applied to the algorithm layer will be reflected in the PSy\sphinxhyphen{}layer PSyIR tree
object that is passed to the \sphinxtitleref{trans()} function.

\sphinxAtStartPar
For example, if the \sphinxtitleref{trans\_alg()} function in the script merged two
\sphinxtitleref{invoke} calls into one then the PSyIR node passed to the
\sphinxtitleref{trans()} function of the script would only contain one Routine
associated with the merged invoke.

\sphinxAtStartPar
An example of the use of a script making use of the \sphinxtitleref{trans\_alg()}
function can be found in examples/gocean/eg7.


\section{Finding transformations}
\label{\detokenize{transformations:finding-transformations}}
\sphinxAtStartPar
Transformations can be imported directly, but the user needs to know
what transformations are available. A helper class \sphinxstylestrong{TransInfo} is
provided to show the available transformations

\begin{sphinxadmonition}{note}{Note:}
\sphinxAtStartPar
The directory layout of PSyclone is currently being restructured.
As a result of this some transformations are already in the new
locations, while others have not been moved yet. Transformations
in the new locations can at the moment not be found using the
\sphinxstylestrong{TransInfo} approach, and need to be imported directly from
the path indicated in the documentation.
\end{sphinxadmonition}
\index{TransInfo (class in psyclone.psyGen)@\spxentry{TransInfo}\spxextra{class in psyclone.psyGen}}

\begin{fulllineitems}
\phantomsection\label{\detokenize{transformations:psyclone.psyGen.TransInfo}}
\pysigstartsignatures
\pysiglinewithargsret{\sphinxbfcode{\sphinxupquote{class\DUrole{w,w}{  }}}\sphinxcode{\sphinxupquote{psyclone.psyGen.}}\sphinxbfcode{\sphinxupquote{TransInfo}}}{\sphinxparam{\DUrole{n,n}{module}\DUrole{o,o}{=}\DUrole{default_value}{None}}\sphinxparamcomma \sphinxparam{\DUrole{n,n}{base\_class}\DUrole{o,o}{=}\DUrole{default_value}{None}}}{}
\pysigstopsignatures
\sphinxAtStartPar
This class provides information about, and access, to the available
transformations in this implementation of PSyclone. New transformations
will be picked up automatically as long as they subclass the abstract
Transformation class.

\sphinxAtStartPar
For example:

\begin{sphinxVerbatim}[commandchars=\\\{\}]
\PYG{g+gp}{\PYGZgt{}\PYGZgt{}\PYGZgt{} }\PYG{k+kn}{from}\PYG{+w}{ }\PYG{n+nn}{psyclone}\PYG{n+nn}{.}\PYG{n+nn}{psyGen}\PYG{+w}{ }\PYG{k+kn}{import} \PYG{n}{TransInfo}
\PYG{g+gp}{\PYGZgt{}\PYGZgt{}\PYGZgt{} }\PYG{n}{t} \PYG{o}{=} \PYG{n}{TransInfo}\PYG{p}{(}\PYG{p}{)}
\PYG{g+gp}{\PYGZgt{}\PYGZgt{}\PYGZgt{} }\PYG{n+nb}{print}\PYG{p}{(}\PYG{n}{t}\PYG{o}{.}\PYG{n}{list}\PYG{p}{)}
\PYG{g+go}{There is 1 transformation available:}
\PYG{g+go}{  1: SwapTrans, A test transformation}
\PYG{g+gp}{\PYGZgt{}\PYGZgt{}\PYGZgt{} }\PYG{c+c1}{\PYGZsh{} accessing a transformation by index}
\PYG{g+gp}{\PYGZgt{}\PYGZgt{}\PYGZgt{} }\PYG{n}{trans} \PYG{o}{=} \PYG{n}{t}\PYG{o}{.}\PYG{n}{get\PYGZus{}trans\PYGZus{}num}\PYG{p}{(}\PYG{l+m+mi}{1}\PYG{p}{)}
\PYG{g+gp}{\PYGZgt{}\PYGZgt{}\PYGZgt{} }\PYG{c+c1}{\PYGZsh{} accessing a transformation by name}
\PYG{g+gp}{\PYGZgt{}\PYGZgt{}\PYGZgt{} }\PYG{n}{trans} \PYG{o}{=} \PYG{n}{t}\PYG{o}{.}\PYG{n}{get\PYGZus{}trans\PYGZus{}name}\PYG{p}{(}\PYG{l+s+s2}{\PYGZdq{}}\PYG{l+s+s2}{SwapTrans}\PYG{l+s+s2}{\PYGZdq{}}\PYG{p}{)}
\end{sphinxVerbatim}
\index{get\_trans\_name() (psyclone.psyGen.TransInfo method)@\spxentry{get\_trans\_name()}\spxextra{psyclone.psyGen.TransInfo method}}

\begin{fulllineitems}
\phantomsection\label{\detokenize{transformations:psyclone.psyGen.TransInfo.get_trans_name}}
\pysigstartsignatures
\pysiglinewithargsret{\sphinxbfcode{\sphinxupquote{get\_trans\_name}}}{\sphinxparam{\DUrole{n,n}{name}}}{}
\pysigstopsignatures
\sphinxAtStartPar
return the transformation with this name (use list() first to see
available transformations)

\end{fulllineitems}

\index{get\_trans\_num() (psyclone.psyGen.TransInfo method)@\spxentry{get\_trans\_num()}\spxextra{psyclone.psyGen.TransInfo method}}

\begin{fulllineitems}
\phantomsection\label{\detokenize{transformations:psyclone.psyGen.TransInfo.get_trans_num}}
\pysigstartsignatures
\pysiglinewithargsret{\sphinxbfcode{\sphinxupquote{get\_trans\_num}}}{\sphinxparam{\DUrole{n,n}{number}}}{}
\pysigstopsignatures
\sphinxAtStartPar
return the transformation with this number (use list() first to
see available transformations)

\end{fulllineitems}

\index{list (psyclone.psyGen.TransInfo property)@\spxentry{list}\spxextra{psyclone.psyGen.TransInfo property}}

\begin{fulllineitems}
\phantomsection\label{\detokenize{transformations:psyclone.psyGen.TransInfo.list}}
\pysigstartsignatures
\pysigline{\sphinxbfcode{\sphinxupquote{property\DUrole{w,w}{  }}}\sphinxbfcode{\sphinxupquote{list}}}
\pysigstopsignatures
\sphinxAtStartPar
return a string with a human readable list of the available
transformations

\end{fulllineitems}

\index{num\_trans (psyclone.psyGen.TransInfo property)@\spxentry{num\_trans}\spxextra{psyclone.psyGen.TransInfo property}}

\begin{fulllineitems}
\phantomsection\label{\detokenize{transformations:psyclone.psyGen.TransInfo.num_trans}}
\pysigstartsignatures
\pysigline{\sphinxbfcode{\sphinxupquote{property\DUrole{w,w}{  }}}\sphinxbfcode{\sphinxupquote{num\_trans}}}
\pysigstopsignatures
\sphinxAtStartPar
return the number of transformations available

\end{fulllineitems}


\end{fulllineitems}



\section{Validating and Applying transformations}
\label{\detokenize{transformations:validating-and-applying-transformations}}
\sphinxAtStartPar
Each transformation must provide at least two functions for the
user: one for validation, i.e. to verify that a certain transformation
can be applied, and one to actually apply the transformation. They are
described in detail in the
{\hyperref[\detokenize{transformations:sec-transformations-available}]{\sphinxcrossref{\DUrole{std,std-ref}{overview of all transformations}}}},
but the following general guidelines apply.


\subsection{Validation}
\label{\detokenize{transformations:validation}}
\sphinxAtStartPar
Each transformation provides a function \sphinxcode{\sphinxupquote{validate}}. This function
can be called by the user, and it will raise an exception if the
transformation can not be applied (and otherwise will return nothing).
Validation will always be called when a transformation is applied.
The parameters for \sphinxcode{\sphinxupquote{validate}} can change from transformation to
transformation, but each \sphinxcode{\sphinxupquote{validate}} function accepts a parameter
\sphinxcode{\sphinxupquote{options}}. This parameter is either \sphinxcode{\sphinxupquote{None}}, or a dictionary of
string keys, that will provide additional parameters to the validation
process. For example, some validation functions allow part of the
validation process to be disabled in order to allow the HPC expert
to apply a transformation that they know to be safe, even if the
more general validation process might reject it. Those parameters
are documented for each transformation, and will show up as
a parameter, e.g.: \sphinxcode{\sphinxupquote{options{[}"node\sphinxhyphen{}type\sphinxhyphen{}check"{]}}}. As a simple
example:

\begin{sphinxVerbatim}[commandchars=\\\{\}]
\PYG{c+c1}{\PYGZsh{} The validation might reject the application, but in this}
\PYG{c+c1}{\PYGZsh{} specific case it is safe to apply the transformation,}
\PYG{c+c1}{\PYGZsh{} so disable the node type check:}
\PYG{n}{my\PYGZus{}transform}\PYG{o}{.}\PYG{n}{validate}\PYG{p}{(}\PYG{n}{node}\PYG{p}{,} \PYG{p}{\PYGZob{}}\PYG{l+s+s2}{\PYGZdq{}}\PYG{l+s+s2}{node\PYGZhy{}type\PYGZhy{}check}\PYG{l+s+s2}{\PYGZdq{}}\PYG{p}{:} \PYG{k+kc}{False}\PYG{p}{\PYGZcb{}}\PYG{p}{)}
\end{sphinxVerbatim}


\subsection{Application}
\label{\detokenize{transformations:application}}\label{\detokenize{transformations:transformations-application}}
\sphinxAtStartPar
Each transformation provides a function \sphinxcode{\sphinxupquote{apply}} which will
apply the transformation. It will first validate the transform
by calling the \sphinxcode{\sphinxupquote{validate}} function. Each \sphinxcode{\sphinxupquote{apply}} function
takes the same \sphinxcode{\sphinxupquote{options}} parameter as the \sphinxcode{\sphinxupquote{validate}} function
described above. Besides potentially modifying the validation
process, optional parameters for the transformation are also
provided this way. A simple example:

\begin{sphinxVerbatim}[commandchars=\\\{\}]
\PYG{n}{kctrans} \PYG{o}{=} \PYG{n}{Dynamo0p3KernelConstTrans}\PYG{p}{(}\PYG{p}{)}
\PYG{n}{kctrans}\PYG{o}{.}\PYG{n}{apply}\PYG{p}{(}\PYG{n}{kernel}\PYG{p}{,} \PYG{p}{\PYGZob{}}\PYG{l+s+s2}{\PYGZdq{}}\PYG{l+s+s2}{element\PYGZus{}order\PYGZus{}h}\PYG{l+s+s2}{\PYGZdq{}}\PYG{p}{:} \PYG{l+m+mi}{0}\PYG{p}{,} \PYG{l+s+s2}{\PYGZdq{}}\PYG{l+s+s2}{element\PYGZus{}order\PYGZus{}v}\PYG{l+s+s2}{\PYGZdq{}}\PYG{p}{:} \PYG{l+m+mi}{0}\PYG{p}{,} \PYG{l+s+s2}{\PYGZdq{}}\PYG{l+s+s2}{quadrature}\PYG{l+s+s2}{\PYGZdq{}}\PYG{p}{:} \PYG{k+kc}{True}\PYG{p}{\PYGZcb{}}\PYG{p}{)}
\end{sphinxVerbatim}

\sphinxAtStartPar
The same \sphinxcode{\sphinxupquote{options}} dictionary will be used when calling \sphinxcode{\sphinxupquote{validate}}.


\section{Available transformations}
\label{\detokenize{transformations:available-transformations}}\label{\detokenize{transformations:sec-transformations-available}}
\sphinxAtStartPar
Some transformations are generic as the schedule structure is
independent of the API, however it often makes sense to specialise
these for a particular API by adding API\sphinxhyphen{}specific errors checks. Some
transformations are API\sphinxhyphen{}specific. Currently these different types of
transformation are indicated by their names.

\sphinxAtStartPar
The generic transformations currently available are listed in
alphabetical order below (a number of these have specialisations which
can be found in the API\sphinxhyphen{}specific sections).

\begin{sphinxadmonition}{note}{Note:}
\sphinxAtStartPar
PSyclone currently only supports OpenCL and
KernelImportsToArguments transformations for the GOcean 1.0
API, the OpenACC Data transformation is limited to
the generic code transformation and the GOcean 1.0 API and the
OpenACC Kernels transformation is limited to the generic code
transformation and the LFRic API.
\end{sphinxadmonition}

\begin{sphinxadmonition}{note}{Note:}
\sphinxAtStartPar
The directory layout of PSyclone is currently being restructured.
As a result of this some transformations are already in the new
locations, while others have not been moved yet.
\end{sphinxadmonition}


\bigskip\hrule\bigskip



\begin{fulllineitems}

\pysigstartsignatures
\pysigline{\sphinxbfcode{\sphinxupquote{class\DUrole{w,w}{  }}}\sphinxcode{\sphinxupquote{psyclone.psyir.transformations.}}\sphinxbfcode{\sphinxupquote{Abs2CodeTrans}}}
\pysigstopsignatures
\sphinxAtStartPar
Provides a transformation from a PSyIR ABS Operator node to
equivalent code in a PSyIR tree. Validity checks are also
performed.

\sphinxAtStartPar
The transformation replaces

\begin{sphinxVerbatim}[commandchars=\\\{\}]
\PYG{n}{R} \PYG{o}{=} \PYG{n}{ABS}\PYG{p}{(}\PYG{n}{X}\PYG{p}{)}
\end{sphinxVerbatim}

\sphinxAtStartPar
with the following logic:

\begin{sphinxVerbatim}[commandchars=\\\{\}]
\PYG{n}{IF} \PYG{n}{X} \PYG{o}{\PYGZlt{}} \PYG{l+m+mf}{0.0}\PYG{p}{:}
    \PYG{n}{R} \PYG{o}{=} \PYG{n}{X}\PYG{o}{*}\PYG{o}{\PYGZhy{}}\PYG{l+m+mf}{1.0}
\PYG{n}{ELSE}\PYG{p}{:}
    \PYG{n}{R} \PYG{o}{=} \PYG{n}{X}
\end{sphinxVerbatim}


\begin{fulllineitems}

\pysigstartsignatures
\pysiglinewithargsret{\sphinxbfcode{\sphinxupquote{apply}}}{\sphinxparam{\DUrole{n,n}{node}}\sphinxparamcomma \sphinxparam{\DUrole{n,n}{options}\DUrole{o,o}{=}\DUrole{default_value}{None}}}{}
\pysigstopsignatures
\sphinxAtStartPar
Apply the ABS intrinsic conversion transformation to the specified
node. This node must be an ABS UnaryOperation. The ABS
UnaryOperation is converted to equivalent inline code. This is
implemented as a PSyIR transform from:

\begin{sphinxVerbatim}[commandchars=\\\{\}]
\PYG{n}{R} \PYG{o}{=} \PYG{o}{.}\PYG{o}{.}\PYG{o}{.} \PYG{n}{ABS}\PYG{p}{(}\PYG{n}{X}\PYG{p}{)} \PYG{o}{.}\PYG{o}{.}\PYG{o}{.}
\end{sphinxVerbatim}

\sphinxAtStartPar
to:

\begin{sphinxVerbatim}[commandchars=\\\{\}]
\PYG{n}{tmp\PYGZus{}abs} \PYG{o}{=} \PYG{n}{X}
\PYG{k}{if} \PYG{n}{tmp\PYGZus{}abs} \PYG{o}{\PYGZlt{}} \PYG{l+m+mf}{0.0}\PYG{p}{:}
    \PYG{n}{res\PYGZus{}abs} \PYG{o}{=} \PYG{n}{tmp\PYGZus{}abs}\PYG{o}{*}\PYG{o}{\PYGZhy{}}\PYG{l+m+mf}{1.0}
\PYG{k}{else}\PYG{p}{:}
    \PYG{n}{res\PYGZus{}abs} \PYG{o}{=} \PYG{n}{tmp\PYGZus{}abs}
\PYG{n}{R} \PYG{o}{=} \PYG{o}{.}\PYG{o}{.}\PYG{o}{.} \PYG{n}{res\PYGZus{}abs} \PYG{o}{.}\PYG{o}{.}\PYG{o}{.}
\end{sphinxVerbatim}

\sphinxAtStartPar
where \sphinxcode{\sphinxupquote{X}} could be an arbitrarily complex PSyIR expression
and \sphinxcode{\sphinxupquote{...}} could be arbitrary PSyIR code.

\sphinxAtStartPar
This transformation requires the operation node to be a
descendent of an assignment and will raise an exception if
this is not the case.
\begin{quote}\begin{description}
\sphinxlineitem{Parameters}\begin{itemize}
\item {} 
\sphinxAtStartPar
\sphinxstyleliteralstrong{\sphinxupquote{node}} (\sphinxcode{\sphinxupquote{psyclone.psyir.nodes.UnaryOperation}}) \textendash{} an ABS UnaryOperation node.

\item {} 
\sphinxAtStartPar
\sphinxstyleliteralstrong{\sphinxupquote{options}} (\sphinxstyleliteralemphasis{\sphinxupquote{Optional}}\sphinxstyleliteralemphasis{\sphinxupquote{{[}}}\sphinxstyleliteralemphasis{\sphinxupquote{Dict}}\sphinxstyleliteralemphasis{\sphinxupquote{{[}}}\sphinxhref{https://docs.python.org/3/library/stdtypes.html\#str}{\sphinxstyleliteralemphasis{\sphinxupquote{str}}}\sphinxstyleliteralemphasis{\sphinxupquote{, }}\sphinxstyleliteralemphasis{\sphinxupquote{Any}}\sphinxstyleliteralemphasis{\sphinxupquote{{]}}}\sphinxstyleliteralemphasis{\sphinxupquote{{]}}}) \textendash{} a dictionary with options for transformations.

\end{itemize}

\end{description}\end{quote}

\end{fulllineitems}


\end{fulllineitems}



\bigskip\hrule\bigskip



\begin{fulllineitems}

\pysigstartsignatures
\pysigline{\sphinxbfcode{\sphinxupquote{class\DUrole{w,w}{  }}}\sphinxcode{\sphinxupquote{psyclone.transformations.}}\sphinxbfcode{\sphinxupquote{ACCDataTrans}}}
\pysigstopsignatures
\sphinxAtStartPar
Add an OpenACC data region around a list of nodes in the PSyIR.
COPYIN, COPYOUT and COPY clauses are added as required.

\sphinxAtStartPar
For example:

\begin{sphinxVerbatim}[commandchars=\\\{\}]
\PYG{g+gp}{\PYGZgt{}\PYGZgt{}\PYGZgt{} }\PYG{k+kn}{from}\PYG{+w}{ }\PYG{n+nn}{psyclone}\PYG{n+nn}{.}\PYG{n+nn}{psyir}\PYG{n+nn}{.}\PYG{n+nn}{frontend}\PYG{+w}{ }\PYG{k+kn}{import} \PYG{n}{FortranReader}
\PYG{g+gp}{\PYGZgt{}\PYGZgt{}\PYGZgt{} }\PYG{n}{psyir} \PYG{o}{=} \PYG{n}{FortranReader}\PYG{p}{(}\PYG{p}{)}\PYG{o}{.}\PYG{n}{psyir\PYGZus{}from\PYGZus{}source}\PYG{p}{(}\PYG{n}{NEMO\PYGZus{}SOURCE\PYGZus{}FILE}\PYG{p}{)}
\PYG{g+gp}{\PYGZgt{}\PYGZgt{}\PYGZgt{}}
\PYG{g+gp}{\PYGZgt{}\PYGZgt{}\PYGZgt{} }\PYG{k+kn}{from}\PYG{+w}{ }\PYG{n+nn}{psyclone}\PYG{n+nn}{.}\PYG{n+nn}{transformations}\PYG{+w}{ }\PYG{k+kn}{import} \PYG{n}{ACCDataTrans}
\PYG{g+gp}{\PYGZgt{}\PYGZgt{}\PYGZgt{} }\PYG{k+kn}{from}\PYG{+w}{ }\PYG{n+nn}{psyclone}\PYG{n+nn}{.}\PYG{n+nn}{psyir}\PYG{n+nn}{.}\PYG{n+nn}{transformations}\PYG{+w}{ }\PYG{k+kn}{import} \PYG{n}{ACCKernelsTrans}
\PYG{g+gp}{\PYGZgt{}\PYGZgt{}\PYGZgt{} }\PYG{n}{ktrans} \PYG{o}{=} \PYG{n}{ACCKernelsTrans}\PYG{p}{(}\PYG{p}{)}
\PYG{g+gp}{\PYGZgt{}\PYGZgt{}\PYGZgt{} }\PYG{n}{dtrans} \PYG{o}{=} \PYG{n}{ACCDataTrans}\PYG{p}{(}\PYG{p}{)}
\PYG{g+gp}{\PYGZgt{}\PYGZgt{}\PYGZgt{}}
\PYG{g+gp}{\PYGZgt{}\PYGZgt{}\PYGZgt{} }\PYG{n}{schedule} \PYG{o}{=} \PYG{n}{psyir}\PYG{o}{.}\PYG{n}{children}\PYG{p}{[}\PYG{l+m+mi}{0}\PYG{p}{]}
\PYG{g+gp}{\PYGZgt{}\PYGZgt{}\PYGZgt{} }\PYG{c+c1}{\PYGZsh{} Uncomment the following line to see a text view of the schedule}
\PYG{g+gp}{\PYGZgt{}\PYGZgt{}\PYGZgt{} }\PYG{c+c1}{\PYGZsh{} print(schedule.view())}
\PYG{g+gp}{\PYGZgt{}\PYGZgt{}\PYGZgt{}}
\PYG{g+gp}{\PYGZgt{}\PYGZgt{}\PYGZgt{} }\PYG{c+c1}{\PYGZsh{} Add a kernels construct for execution on the device}
\PYG{g+gp}{\PYGZgt{}\PYGZgt{}\PYGZgt{} }\PYG{n}{kernels} \PYG{o}{=} \PYG{n}{schedule}\PYG{o}{.}\PYG{n}{children}\PYG{p}{[}\PYG{l+m+mi}{9}\PYG{p}{]}
\PYG{g+gp}{\PYGZgt{}\PYGZgt{}\PYGZgt{} }\PYG{n}{ktrans}\PYG{o}{.}\PYG{n}{apply}\PYG{p}{(}\PYG{n}{kernels}\PYG{p}{)}
\PYG{g+gp}{\PYGZgt{}\PYGZgt{}\PYGZgt{}}
\PYG{g+gp}{\PYGZgt{}\PYGZgt{}\PYGZgt{} }\PYG{c+c1}{\PYGZsh{} Enclose the kernels in a data construct}
\PYG{g+gp}{\PYGZgt{}\PYGZgt{}\PYGZgt{} }\PYG{n}{kernels} \PYG{o}{=} \PYG{n}{schedule}\PYG{o}{.}\PYG{n}{children}\PYG{p}{[}\PYG{l+m+mi}{9}\PYG{p}{]}
\PYG{g+gp}{\PYGZgt{}\PYGZgt{}\PYGZgt{} }\PYG{n}{dtrans}\PYG{o}{.}\PYG{n}{apply}\PYG{p}{(}\PYG{n}{kernels}\PYG{p}{)}
\end{sphinxVerbatim}


\begin{fulllineitems}

\pysigstartsignatures
\pysiglinewithargsret{\sphinxbfcode{\sphinxupquote{apply}}}{\sphinxparam{\DUrole{n,n}{node}}\sphinxparamcomma \sphinxparam{\DUrole{n,n}{options}\DUrole{o,o}{=}\DUrole{default_value}{None}}}{}
\pysigstopsignatures
\sphinxAtStartPar
Put the supplied node or list of nodes within an OpenACC data region.
\begin{quote}\begin{description}
\sphinxlineitem{Parameters}\begin{itemize}
\item {} 
\sphinxAtStartPar
\sphinxstyleliteralstrong{\sphinxupquote{node}} ((list of) \sphinxcode{\sphinxupquote{psyclone.psyir.nodes.Node}}) \textendash{} the PSyIR node(s) to enclose in the data region.

\item {} 
\sphinxAtStartPar
\sphinxstyleliteralstrong{\sphinxupquote{options}} (\sphinxstyleliteralemphasis{\sphinxupquote{Optional}}\sphinxstyleliteralemphasis{\sphinxupquote{{[}}}\sphinxstyleliteralemphasis{\sphinxupquote{Dict}}\sphinxstyleliteralemphasis{\sphinxupquote{{[}}}\sphinxhref{https://docs.python.org/3/library/stdtypes.html\#str}{\sphinxstyleliteralemphasis{\sphinxupquote{str}}}\sphinxstyleliteralemphasis{\sphinxupquote{, }}\sphinxstyleliteralemphasis{\sphinxupquote{Any}}\sphinxstyleliteralemphasis{\sphinxupquote{{]}}}\sphinxstyleliteralemphasis{\sphinxupquote{{]}}}) \textendash{} a dictionary with options for transformations.

\end{itemize}

\end{description}\end{quote}

\end{fulllineitems}


\end{fulllineitems}



\bigskip\hrule\bigskip



\begin{fulllineitems}

\pysigstartsignatures
\pysigline{\sphinxbfcode{\sphinxupquote{class\DUrole{w,w}{  }}}\sphinxcode{\sphinxupquote{psyclone.transformations.}}\sphinxbfcode{\sphinxupquote{ACCEnterDataTrans}}}
\pysigstopsignatures
\sphinxAtStartPar
Adds an OpenACC “enter data” directive to a Schedule.
For example:

\begin{sphinxVerbatim}[commandchars=\\\{\}]
\PYG{g+gp}{\PYGZgt{}\PYGZgt{}\PYGZgt{} }\PYG{k+kn}{from}\PYG{+w}{ }\PYG{n+nn}{psyclone}\PYG{n+nn}{.}\PYG{n+nn}{parse}\PYG{n+nn}{.}\PYG{n+nn}{algorithm}\PYG{+w}{ }\PYG{k+kn}{import} \PYG{n}{parse}
\PYG{g+gp}{\PYGZgt{}\PYGZgt{}\PYGZgt{} }\PYG{k+kn}{from}\PYG{+w}{ }\PYG{n+nn}{psyclone}\PYG{n+nn}{.}\PYG{n+nn}{psyGen}\PYG{+w}{ }\PYG{k+kn}{import} \PYG{n}{PSyFactory}
\PYG{g+gp}{\PYGZgt{}\PYGZgt{}\PYGZgt{} }\PYG{n}{api} \PYG{o}{=} \PYG{l+s+s2}{\PYGZdq{}}\PYG{l+s+s2}{gocean}\PYG{l+s+s2}{\PYGZdq{}}
\PYG{g+gp}{\PYGZgt{}\PYGZgt{}\PYGZgt{} }\PYG{n}{ast}\PYG{p}{,} \PYG{n}{invokeInfo} \PYG{o}{=} \PYG{n}{parse}\PYG{p}{(}\PYG{n}{GOCEAN\PYGZus{}SOURCE\PYGZus{}FILE}\PYG{p}{,} \PYG{n}{api}\PYG{o}{=}\PYG{n}{api}\PYG{p}{)}
\PYG{g+gp}{\PYGZgt{}\PYGZgt{}\PYGZgt{} }\PYG{n}{psy} \PYG{o}{=} \PYG{n}{PSyFactory}\PYG{p}{(}\PYG{n}{api}\PYG{p}{)}\PYG{o}{.}\PYG{n}{create}\PYG{p}{(}\PYG{n}{invokeInfo}\PYG{p}{)}
\PYG{g+gp}{\PYGZgt{}\PYGZgt{}\PYGZgt{}}
\PYG{g+gp}{\PYGZgt{}\PYGZgt{}\PYGZgt{} }\PYG{k+kn}{from}\PYG{+w}{ }\PYG{n+nn}{psyclone}\PYG{n+nn}{.}\PYG{n+nn}{transformations}\PYG{+w}{ }\PYG{k+kn}{import}         \PYG{n}{ACCEnterDataTrans}\PYG{p}{,} \PYG{n}{ACCLoopTrans}\PYG{p}{,} \PYG{n}{ACCParallelTrans}
\PYG{g+gp}{\PYGZgt{}\PYGZgt{}\PYGZgt{} }\PYG{n}{dtrans} \PYG{o}{=} \PYG{n}{ACCEnterDataTrans}\PYG{p}{(}\PYG{p}{)}
\PYG{g+gp}{\PYGZgt{}\PYGZgt{}\PYGZgt{} }\PYG{n}{ltrans} \PYG{o}{=} \PYG{n}{ACCLoopTrans}\PYG{p}{(}\PYG{p}{)}
\PYG{g+gp}{\PYGZgt{}\PYGZgt{}\PYGZgt{} }\PYG{n}{ptrans} \PYG{o}{=} \PYG{n}{ACCParallelTrans}\PYG{p}{(}\PYG{p}{)}
\PYG{g+gp}{\PYGZgt{}\PYGZgt{}\PYGZgt{}}
\PYG{g+gp}{\PYGZgt{}\PYGZgt{}\PYGZgt{} }\PYG{n}{schedule} \PYG{o}{=} \PYG{n}{psy}\PYG{o}{.}\PYG{n}{invokes}\PYG{o}{.}\PYG{n}{get}\PYG{p}{(}\PYG{l+s+s1}{\PYGZsq{}}\PYG{l+s+s1}{invoke\PYGZus{}0}\PYG{l+s+s1}{\PYGZsq{}}\PYG{p}{)}\PYG{o}{.}\PYG{n}{schedule}
\PYG{g+gp}{\PYGZgt{}\PYGZgt{}\PYGZgt{} }\PYG{c+c1}{\PYGZsh{} Uncomment the following line to see a text view of the schedule}
\PYG{g+gp}{\PYGZgt{}\PYGZgt{}\PYGZgt{} }\PYG{c+c1}{\PYGZsh{} print(schedule.view())}
\PYG{g+gp}{\PYGZgt{}\PYGZgt{}\PYGZgt{}}
\PYG{g+gp}{\PYGZgt{}\PYGZgt{}\PYGZgt{} }\PYG{c+c1}{\PYGZsh{} Apply the OpenACC Loop transformation to *every* loop in the schedule}
\PYG{g+gp}{\PYGZgt{}\PYGZgt{}\PYGZgt{} }\PYG{k}{for} \PYG{n}{child} \PYG{o+ow}{in} \PYG{n}{schedule}\PYG{o}{.}\PYG{n}{children}\PYG{p}{[}\PYG{p}{:}\PYG{p}{]}\PYG{p}{:}
\PYG{g+gp}{... }    \PYG{n}{ltrans}\PYG{o}{.}\PYG{n}{apply}\PYG{p}{(}\PYG{n}{child}\PYG{p}{)}
\PYG{g+gp}{\PYGZgt{}\PYGZgt{}\PYGZgt{}}
\PYG{g+gp}{\PYGZgt{}\PYGZgt{}\PYGZgt{} }\PYG{c+c1}{\PYGZsh{} Enclose all of these loops within a single OpenACC parallel region}
\PYG{g+gp}{\PYGZgt{}\PYGZgt{}\PYGZgt{} }\PYG{n}{ptrans}\PYG{o}{.}\PYG{n}{apply}\PYG{p}{(}\PYG{n}{schedule}\PYG{p}{)}
\PYG{g+gp}{\PYGZgt{}\PYGZgt{}\PYGZgt{}}
\PYG{g+gp}{\PYGZgt{}\PYGZgt{}\PYGZgt{} }\PYG{c+c1}{\PYGZsh{} Add an enter data directive}
\PYG{g+gp}{\PYGZgt{}\PYGZgt{}\PYGZgt{} }\PYG{n}{dtrans}\PYG{o}{.}\PYG{n}{apply}\PYG{p}{(}\PYG{n}{schedule}\PYG{p}{)}
\PYG{g+gp}{\PYGZgt{}\PYGZgt{}\PYGZgt{}}
\PYG{g+gp}{\PYGZgt{}\PYGZgt{}\PYGZgt{} }\PYG{c+c1}{\PYGZsh{} Uncomment the following line to see a text view of the schedule}
\PYG{g+gp}{\PYGZgt{}\PYGZgt{}\PYGZgt{} }\PYG{c+c1}{\PYGZsh{} print(schedule.view())}
\end{sphinxVerbatim}


\begin{fulllineitems}

\pysigstartsignatures
\pysiglinewithargsret{\sphinxbfcode{\sphinxupquote{apply}}}{\sphinxparam{\DUrole{n,n}{node}}\sphinxparamcomma \sphinxparam{\DUrole{n,n}{options}\DUrole{o,o}{=}\DUrole{default_value}{\{\}}}}{}
\pysigstopsignatures
\sphinxAtStartPar
Adds an OpenACC “enter data” directive to the invoke associated
with the supplied Schedule. Any fields accessed by OpenACC kernels
within this schedule will be added to this data region in
order to ensure they remain on the target device.
\begin{quote}\begin{description}
\sphinxlineitem{Parameters}\begin{itemize}
\item {} 
\sphinxAtStartPar
\sphinxstyleliteralstrong{\sphinxupquote{node}} (\DUrole{sphinx_autodoc_typehints-type}{\sphinxcode{\sphinxupquote{Schedule}}}) \textendash{} schedule to which to add an “enter data” directive.

\item {} 
\sphinxAtStartPar
\sphinxstyleliteralstrong{\sphinxupquote{options}} (\DUrole{sphinx_autodoc_typehints-type}{\sphinxhref{https://docs.python.org/3/library/typing.html\#typing.Optional}{\sphinxcode{\sphinxupquote{Optional}}}{[}\sphinxhref{https://docs.python.org/3/library/typing.html\#typing.Dict}{\sphinxcode{\sphinxupquote{Dict}}}{[}\sphinxhref{https://docs.python.org/3/library/stdtypes.html\#str}{\sphinxcode{\sphinxupquote{str}}}, \sphinxhref{https://docs.python.org/3/library/typing.html\#typing.Any}{\sphinxcode{\sphinxupquote{Any}}}{]}{]}}) \textendash{} a dictionary with options for transformations.

\item {} 
\sphinxAtStartPar
\sphinxstyleliteralstrong{\sphinxupquote{options}}\sphinxstyleliteralstrong{\sphinxupquote{{[}}}\sphinxstyleliteralstrong{\sphinxupquote{"async\_queue"}}\sphinxstyleliteralstrong{\sphinxupquote{{]}}} (\sphinxstyleliteralemphasis{\sphinxupquote{Union}}\sphinxstyleliteralemphasis{\sphinxupquote{{[}}}\sphinxhref{https://docs.python.org/3/library/functions.html\#bool}{\sphinxstyleliteralemphasis{\sphinxupquote{bool}}}\sphinxstyleliteralemphasis{\sphinxupquote{, }}\sphinxhref{https://docs.python.org/3/library/functions.html\#int}{\sphinxstyleliteralemphasis{\sphinxupquote{int}}}\sphinxstyleliteralemphasis{\sphinxupquote{{]}}}) \textendash{} force the transformation to use the
specified async stream if not False.

\end{itemize}

\end{description}\end{quote}

\end{fulllineitems}


\end{fulllineitems}



\bigskip\hrule\bigskip



\begin{fulllineitems}

\pysigstartsignatures
\pysigline{\sphinxbfcode{\sphinxupquote{class\DUrole{w,w}{  }}}\sphinxcode{\sphinxupquote{psyclone.psyir.transformations.}}\sphinxbfcode{\sphinxupquote{ACCKernelsTrans}}}
\pysigstopsignatures
\sphinxAtStartPar
Enclose a sub\sphinxhyphen{}set of nodes from a Schedule within an OpenACC kernels
region (i.e. within “!\$acc kernels” … “!\$acc end kernels” directives).

\sphinxAtStartPar
For example:

\begin{sphinxVerbatim}[commandchars=\\\{\}]
\PYG{g+gp}{\PYGZgt{}\PYGZgt{}\PYGZgt{} }\PYG{k+kn}{from}\PYG{+w}{ }\PYG{n+nn}{psyclone}\PYG{n+nn}{.}\PYG{n+nn}{psyir}\PYG{n+nn}{.}\PYG{n+nn}{frontend}\PYG{+w}{ }\PYG{k+kn}{import} \PYG{n}{FortranReader}
\PYG{g+gp}{\PYGZgt{}\PYGZgt{}\PYGZgt{} }\PYG{n}{psyir} \PYG{o}{=} \PYG{n}{FortranReader}\PYG{p}{(}\PYG{p}{)}\PYG{o}{.}\PYG{n}{psyir\PYGZus{}from\PYGZus{}source}\PYG{p}{(}\PYG{n}{NEMO\PYGZus{}SOURCE\PYGZus{}FILE}\PYG{p}{)}
\PYG{g+gp}{\PYGZgt{}\PYGZgt{}\PYGZgt{}}
\PYG{g+gp}{\PYGZgt{}\PYGZgt{}\PYGZgt{} }\PYG{k+kn}{from}\PYG{+w}{ }\PYG{n+nn}{psyclone}\PYG{n+nn}{.}\PYG{n+nn}{psyir}\PYG{n+nn}{.}\PYG{n+nn}{transformations}\PYG{+w}{ }\PYG{k+kn}{import} \PYG{n}{ACCKernelsTrans}
\PYG{g+gp}{\PYGZgt{}\PYGZgt{}\PYGZgt{} }\PYG{n}{ktrans} \PYG{o}{=} \PYG{n}{ACCKernelsTrans}\PYG{p}{(}\PYG{p}{)}
\PYG{g+gp}{\PYGZgt{}\PYGZgt{}\PYGZgt{}}
\PYG{g+gp}{\PYGZgt{}\PYGZgt{}\PYGZgt{} }\PYG{n}{schedule} \PYG{o}{=} \PYG{n}{psyir}\PYG{o}{.}\PYG{n}{children}\PYG{p}{[}\PYG{l+m+mi}{0}\PYG{p}{]}
\PYG{g+gp}{\PYGZgt{}\PYGZgt{}\PYGZgt{} }\PYG{c+c1}{\PYGZsh{} Uncomment the following line to see a text view of the schedule}
\PYG{g+gp}{\PYGZgt{}\PYGZgt{}\PYGZgt{} }\PYG{c+c1}{\PYGZsh{} print(schedule.view())}
\PYG{g+gp}{\PYGZgt{}\PYGZgt{}\PYGZgt{} }\PYG{n}{kernels} \PYG{o}{=} \PYG{n}{schedule}\PYG{o}{.}\PYG{n}{children}\PYG{p}{[}\PYG{l+m+mi}{9}\PYG{p}{]}
\PYG{g+gp}{\PYGZgt{}\PYGZgt{}\PYGZgt{} }\PYG{c+c1}{\PYGZsh{} Transform the kernel}
\PYG{g+gp}{\PYGZgt{}\PYGZgt{}\PYGZgt{} }\PYG{n}{ktrans}\PYG{o}{.}\PYG{n}{apply}\PYG{p}{(}\PYG{n}{kernels}\PYG{p}{)}
\end{sphinxVerbatim}


\begin{fulllineitems}

\pysigstartsignatures
\pysiglinewithargsret{\sphinxbfcode{\sphinxupquote{apply}}}{\sphinxparam{\DUrole{n,n}{node}}\sphinxparamcomma \sphinxparam{\DUrole{n,n}{options}\DUrole{o,o}{=}\DUrole{default_value}{\{\}}}}{}
\pysigstopsignatures
\sphinxAtStartPar
Enclose the supplied list of PSyIR nodes within an OpenACC
Kernels region.
\begin{quote}\begin{description}
\sphinxlineitem{Parameters}\begin{itemize}
\item {} 
\sphinxAtStartPar
\sphinxstyleliteralstrong{\sphinxupquote{node}} (\DUrole{sphinx_autodoc_typehints-type}{\sphinxhref{https://docs.python.org/3/library/typing.html\#typing.Union}{\sphinxcode{\sphinxupquote{Union}}}{[}\sphinxcode{\sphinxupquote{Node}}, \sphinxhref{https://docs.python.org/3/library/typing.html\#typing.List}{\sphinxcode{\sphinxupquote{List}}}{[}\sphinxcode{\sphinxupquote{Node}}{]}{]}}) \textendash{} a node or list of nodes in the PSyIR to enclose.

\item {} 
\sphinxAtStartPar
\sphinxstyleliteralstrong{\sphinxupquote{options}} (\DUrole{sphinx_autodoc_typehints-type}{\sphinxhref{https://docs.python.org/3/library/typing.html\#typing.Dict}{\sphinxcode{\sphinxupquote{Dict}}}{[}\sphinxhref{https://docs.python.org/3/library/stdtypes.html\#str}{\sphinxcode{\sphinxupquote{str}}}, \sphinxhref{https://docs.python.org/3/library/typing.html\#typing.Any}{\sphinxcode{\sphinxupquote{Any}}}{]}}) \textendash{} a dictionary with options for transformations.

\item {} 
\sphinxAtStartPar
\sphinxstyleliteralstrong{\sphinxupquote{options}}\sphinxstyleliteralstrong{\sphinxupquote{{[}}}\sphinxstyleliteralstrong{\sphinxupquote{"default\_present"}}\sphinxstyleliteralstrong{\sphinxupquote{{]}}} (\sphinxhref{https://docs.python.org/3/library/functions.html\#bool}{\sphinxstyleliteralemphasis{\sphinxupquote{bool}}}) \textendash{} whether or not the kernels
region should have the ‘default present’ attribute (indicating
that data is already on the accelerator). When using managed
memory this option should be False.

\item {} 
\sphinxAtStartPar
\sphinxstyleliteralstrong{\sphinxupquote{options}}\sphinxstyleliteralstrong{\sphinxupquote{{[}}}\sphinxstyleliteralstrong{\sphinxupquote{"disable\_loop\_check"}}\sphinxstyleliteralstrong{\sphinxupquote{{]}}} (\sphinxhref{https://docs.python.org/3/library/functions.html\#bool}{\sphinxstyleliteralemphasis{\sphinxupquote{bool}}}) \textendash{} whether to disable the check
that the supplied region contains 1 or more loops. Default is False
(i.e. the check is enabled).

\item {} 
\sphinxAtStartPar
\sphinxstyleliteralstrong{\sphinxupquote{options}}\sphinxstyleliteralstrong{\sphinxupquote{{[}}}\sphinxstyleliteralstrong{\sphinxupquote{"async\_queue"}}\sphinxstyleliteralstrong{\sphinxupquote{{]}}} (Union{[}bool, \sphinxcode{\sphinxupquote{psyclone.psyir.nodes.DataNode}}{]}) \textendash{} whether or not to add the ‘async’ clause
to the new directive and if so, which queue to associate it with.
True to enable for the default queue or a queue value specified
with an int or PSyIR expression.

\item {} 
\sphinxAtStartPar
\sphinxstyleliteralstrong{\sphinxupquote{options}}\sphinxstyleliteralstrong{\sphinxupquote{{[}}}\sphinxstyleliteralstrong{\sphinxupquote{"allow\_string"}}\sphinxstyleliteralstrong{\sphinxupquote{{]}}} (\sphinxhref{https://docs.python.org/3/library/functions.html\#bool}{\sphinxstyleliteralemphasis{\sphinxupquote{bool}}}) \textendash{} whether to allow the
transformation on assignments involving character types. Defaults
to False.

\item {} 
\sphinxAtStartPar
\sphinxstyleliteralstrong{\sphinxupquote{options}}\sphinxstyleliteralstrong{\sphinxupquote{{[}}}\sphinxstyleliteralstrong{\sphinxupquote{"verbose"}}\sphinxstyleliteralstrong{\sphinxupquote{{]}}} (\sphinxhref{https://docs.python.org/3/library/functions.html\#bool}{\sphinxstyleliteralemphasis{\sphinxupquote{bool}}}) \textendash{} log the reason the validation failed,
at the moment with a comment in the provided PSyIR node.

\end{itemize}

\end{description}\end{quote}

\end{fulllineitems}


\end{fulllineitems}



\bigskip\hrule\bigskip



\begin{fulllineitems}

\pysigstartsignatures
\pysigline{\sphinxbfcode{\sphinxupquote{class\DUrole{w,w}{  }}}\sphinxcode{\sphinxupquote{psyclone.transformations.}}\sphinxbfcode{\sphinxupquote{ACCLoopTrans}}}
\pysigstopsignatures
\sphinxAtStartPar
Adds an OpenACC loop directive to a loop. This directive must be within
the scope of some OpenACC Parallel region (at code\sphinxhyphen{}generation time).

\sphinxAtStartPar
For example:

\begin{sphinxVerbatim}[commandchars=\\\{\}]
\PYG{g+gp}{\PYGZgt{}\PYGZgt{}\PYGZgt{} }\PYG{k+kn}{from}\PYG{+w}{ }\PYG{n+nn}{psyclone}\PYG{n+nn}{.}\PYG{n+nn}{parse}\PYG{n+nn}{.}\PYG{n+nn}{algorithm}\PYG{+w}{ }\PYG{k+kn}{import} \PYG{n}{parse}
\PYG{g+gp}{\PYGZgt{}\PYGZgt{}\PYGZgt{} }\PYG{k+kn}{from}\PYG{+w}{ }\PYG{n+nn}{psyclone}\PYG{n+nn}{.}\PYG{n+nn}{parse}\PYG{n+nn}{.}\PYG{n+nn}{utils}\PYG{+w}{ }\PYG{k+kn}{import} \PYG{n}{ParseError}
\PYG{g+gp}{\PYGZgt{}\PYGZgt{}\PYGZgt{} }\PYG{k+kn}{from}\PYG{+w}{ }\PYG{n+nn}{psyclone}\PYG{n+nn}{.}\PYG{n+nn}{psyGen}\PYG{+w}{ }\PYG{k+kn}{import} \PYG{n}{PSyFactory}
\PYG{g+gp}{\PYGZgt{}\PYGZgt{}\PYGZgt{} }\PYG{k+kn}{from}\PYG{+w}{ }\PYG{n+nn}{psyclone}\PYG{n+nn}{.}\PYG{n+nn}{errors}\PYG{+w}{ }\PYG{k+kn}{import} \PYG{n}{GenerationError}
\PYG{g+gp}{\PYGZgt{}\PYGZgt{}\PYGZgt{} }\PYG{n}{api} \PYG{o}{=} \PYG{l+s+s2}{\PYGZdq{}}\PYG{l+s+s2}{gocean}\PYG{l+s+s2}{\PYGZdq{}}
\PYG{g+gp}{\PYGZgt{}\PYGZgt{}\PYGZgt{} }\PYG{n}{ast}\PYG{p}{,} \PYG{n}{invokeInfo} \PYG{o}{=} \PYG{n}{parse}\PYG{p}{(}\PYG{n}{GOCEAN\PYGZus{}SOURCE\PYGZus{}FILE}\PYG{p}{,} \PYG{n}{api}\PYG{o}{=}\PYG{n}{api}\PYG{p}{)}
\PYG{g+gp}{\PYGZgt{}\PYGZgt{}\PYGZgt{} }\PYG{n}{psy} \PYG{o}{=} \PYG{n}{PSyFactory}\PYG{p}{(}\PYG{n}{api}\PYG{p}{)}\PYG{o}{.}\PYG{n}{create}\PYG{p}{(}\PYG{n}{invokeInfo}\PYG{p}{)}
\PYG{g+gp}{\PYGZgt{}\PYGZgt{}\PYGZgt{}}
\PYG{g+gp}{\PYGZgt{}\PYGZgt{}\PYGZgt{} }\PYG{k+kn}{from}\PYG{+w}{ }\PYG{n+nn}{psyclone}\PYG{n+nn}{.}\PYG{n+nn}{psyGen}\PYG{+w}{ }\PYG{k+kn}{import} \PYG{n}{TransInfo}
\PYG{g+gp}{\PYGZgt{}\PYGZgt{}\PYGZgt{} }\PYG{n}{t} \PYG{o}{=} \PYG{n}{TransInfo}\PYG{p}{(}\PYG{p}{)}
\PYG{g+gp}{\PYGZgt{}\PYGZgt{}\PYGZgt{} }\PYG{n}{ltrans} \PYG{o}{=} \PYG{n}{t}\PYG{o}{.}\PYG{n}{get\PYGZus{}trans\PYGZus{}name}\PYG{p}{(}\PYG{l+s+s1}{\PYGZsq{}}\PYG{l+s+s1}{ACCLoopTrans}\PYG{l+s+s1}{\PYGZsq{}}\PYG{p}{)}
\PYG{g+gp}{\PYGZgt{}\PYGZgt{}\PYGZgt{} }\PYG{n}{rtrans} \PYG{o}{=} \PYG{n}{t}\PYG{o}{.}\PYG{n}{get\PYGZus{}trans\PYGZus{}name}\PYG{p}{(}\PYG{l+s+s1}{\PYGZsq{}}\PYG{l+s+s1}{ACCParallelTrans}\PYG{l+s+s1}{\PYGZsq{}}\PYG{p}{)}
\PYG{g+gp}{\PYGZgt{}\PYGZgt{}\PYGZgt{}}
\PYG{g+gp}{\PYGZgt{}\PYGZgt{}\PYGZgt{} }\PYG{n}{schedule} \PYG{o}{=} \PYG{n}{psy}\PYG{o}{.}\PYG{n}{invokes}\PYG{o}{.}\PYG{n}{get}\PYG{p}{(}\PYG{l+s+s1}{\PYGZsq{}}\PYG{l+s+s1}{invoke\PYGZus{}0}\PYG{l+s+s1}{\PYGZsq{}}\PYG{p}{)}\PYG{o}{.}\PYG{n}{schedule}
\PYG{g+gp}{\PYGZgt{}\PYGZgt{}\PYGZgt{} }\PYG{c+c1}{\PYGZsh{} Uncomment the following line to see a text view of the schedule}
\PYG{g+gp}{\PYGZgt{}\PYGZgt{}\PYGZgt{} }\PYG{c+c1}{\PYGZsh{} print(schedule.view())}
\PYG{g+gp}{\PYGZgt{}\PYGZgt{}\PYGZgt{}}
\PYG{g+gp}{\PYGZgt{}\PYGZgt{}\PYGZgt{} }\PYG{c+c1}{\PYGZsh{} Apply the OpenACC Loop transformation to *every* loop in the schedule}
\PYG{g+gp}{\PYGZgt{}\PYGZgt{}\PYGZgt{} }\PYG{k}{for} \PYG{n}{child} \PYG{o+ow}{in} \PYG{n}{schedule}\PYG{o}{.}\PYG{n}{children}\PYG{p}{[}\PYG{p}{:}\PYG{p}{]}\PYG{p}{:}
\PYG{g+gp}{... }    \PYG{n}{ltrans}\PYG{o}{.}\PYG{n}{apply}\PYG{p}{(}\PYG{n}{child}\PYG{p}{)}
\PYG{g+gp}{\PYGZgt{}\PYGZgt{}\PYGZgt{}}
\PYG{g+gp}{\PYGZgt{}\PYGZgt{}\PYGZgt{} }\PYG{c+c1}{\PYGZsh{} Enclose all of these loops within a single OpenACC parallel region}
\PYG{g+gp}{\PYGZgt{}\PYGZgt{}\PYGZgt{} }\PYG{n}{rtrans}\PYG{o}{.}\PYG{n}{apply}\PYG{p}{(}\PYG{n}{schedule}\PYG{p}{)}
\PYG{g+gp}{\PYGZgt{}\PYGZgt{}\PYGZgt{}}
\end{sphinxVerbatim}


\begin{fulllineitems}

\pysigstartsignatures
\pysiglinewithargsret{\sphinxbfcode{\sphinxupquote{apply}}}{\sphinxparam{\DUrole{n,n}{node}}\sphinxparamcomma \sphinxparam{\DUrole{n,n}{options}\DUrole{o,o}{=}\DUrole{default_value}{None}}}{}
\pysigstopsignatures
\sphinxAtStartPar
Apply the ACCLoop transformation to the specified node. This node
must be a Loop since this transformation corresponds to
inserting a directive immediately before a loop, e.g.:

\begin{sphinxVerbatim}[commandchars=\\\{\}]
\PYG{c}{!\PYGZdl{}ACC LOOP}
\PYG{k}{do}\PYG{+w}{ }\PYG{p}{.}\PYG{p}{.}\PYG{p}{.}
\PYG{+w}{   }\PYG{p}{.}\PYG{p}{.}\PYG{p}{.}
\PYG{k}{end }\PYG{k}{do}
\end{sphinxVerbatim}

\sphinxAtStartPar
At code\sphinxhyphen{}generation time (when
\sphinxcode{\sphinxupquote{psyclone.psyir.nodes.ACCLoopDirective.gen\_code()}} is called),
this node must be within (i.e. a child of) a PARALLEL region.
\begin{quote}\begin{description}
\sphinxlineitem{Parameters}\begin{itemize}
\item {} 
\sphinxAtStartPar
\sphinxstyleliteralstrong{\sphinxupquote{node}} (\sphinxcode{\sphinxupquote{psyclone.psyir.nodes.Loop}}) \textendash{} the supplied node to which we will apply the
Loop transformation.

\item {} 
\sphinxAtStartPar
\sphinxstyleliteralstrong{\sphinxupquote{options}} (\sphinxstyleliteralemphasis{\sphinxupquote{Optional}}\sphinxstyleliteralemphasis{\sphinxupquote{{[}}}\sphinxstyleliteralemphasis{\sphinxupquote{Dict}}\sphinxstyleliteralemphasis{\sphinxupquote{{[}}}\sphinxhref{https://docs.python.org/3/library/stdtypes.html\#str}{\sphinxstyleliteralemphasis{\sphinxupquote{str}}}\sphinxstyleliteralemphasis{\sphinxupquote{, }}\sphinxstyleliteralemphasis{\sphinxupquote{Any}}\sphinxstyleliteralemphasis{\sphinxupquote{{]}}}\sphinxstyleliteralemphasis{\sphinxupquote{{]}}}) \textendash{} a dictionary with options for transformations.

\item {} 
\sphinxAtStartPar
\sphinxstyleliteralstrong{\sphinxupquote{options}}\sphinxstyleliteralstrong{\sphinxupquote{{[}}}\sphinxstyleliteralstrong{\sphinxupquote{"collapse"}}\sphinxstyleliteralstrong{\sphinxupquote{{]}}} (\sphinxhref{https://docs.python.org/3/library/functions.html\#int}{\sphinxstyleliteralemphasis{\sphinxupquote{int}}}) \textendash{} number of nested loops to collapse.

\item {} 
\sphinxAtStartPar
\sphinxstyleliteralstrong{\sphinxupquote{options}}\sphinxstyleliteralstrong{\sphinxupquote{{[}}}\sphinxstyleliteralstrong{\sphinxupquote{"independent"}}\sphinxstyleliteralstrong{\sphinxupquote{{]}}} (\sphinxhref{https://docs.python.org/3/library/functions.html\#bool}{\sphinxstyleliteralemphasis{\sphinxupquote{bool}}}) \textendash{} whether to add the “independent”
clause to the directive (not strictly necessary within
PARALLEL regions).

\item {} 
\sphinxAtStartPar
\sphinxstyleliteralstrong{\sphinxupquote{options}}\sphinxstyleliteralstrong{\sphinxupquote{{[}}}\sphinxstyleliteralstrong{\sphinxupquote{"sequential"}}\sphinxstyleliteralstrong{\sphinxupquote{{]}}} (\sphinxhref{https://docs.python.org/3/library/functions.html\#bool}{\sphinxstyleliteralemphasis{\sphinxupquote{bool}}}) \textendash{} whether to add the “seq” clause to
the directive.

\item {} 
\sphinxAtStartPar
\sphinxstyleliteralstrong{\sphinxupquote{options}}\sphinxstyleliteralstrong{\sphinxupquote{{[}}}\sphinxstyleliteralstrong{\sphinxupquote{"gang"}}\sphinxstyleliteralstrong{\sphinxupquote{{]}}} (\sphinxhref{https://docs.python.org/3/library/functions.html\#bool}{\sphinxstyleliteralemphasis{\sphinxupquote{bool}}}) \textendash{} whether to add the “gang” clause to the
directive.

\item {} 
\sphinxAtStartPar
\sphinxstyleliteralstrong{\sphinxupquote{options}}\sphinxstyleliteralstrong{\sphinxupquote{{[}}}\sphinxstyleliteralstrong{\sphinxupquote{"vector"}}\sphinxstyleliteralstrong{\sphinxupquote{{]}}} (\sphinxhref{https://docs.python.org/3/library/functions.html\#bool}{\sphinxstyleliteralemphasis{\sphinxupquote{bool}}}) \textendash{} whether to add the “vector” clause to
the directive.

\end{itemize}

\end{description}\end{quote}

\end{fulllineitems}


\end{fulllineitems}



\bigskip\hrule\bigskip



\begin{fulllineitems}

\pysigstartsignatures
\pysiglinewithargsret{\sphinxbfcode{\sphinxupquote{class\DUrole{w,w}{  }}}\sphinxcode{\sphinxupquote{psyclone.transformations.}}\sphinxbfcode{\sphinxupquote{ACCParallelTrans}}}{\sphinxparam{\DUrole{n,n}{default\_present}\DUrole{o,o}{=}\DUrole{default_value}{True}}}{}
\pysigstopsignatures
\sphinxAtStartPar
Create an OpenACC parallel region by inserting an ‘acc parallel’
directive.

\begin{sphinxVerbatim}[commandchars=\\\{\}]
\PYG{g+gp}{\PYGZgt{}\PYGZgt{}\PYGZgt{} }\PYG{k+kn}{from}\PYG{+w}{ }\PYG{n+nn}{psyclone}\PYG{n+nn}{.}\PYG{n+nn}{psyGen}\PYG{+w}{ }\PYG{k+kn}{import} \PYG{n}{TransInfo}
\PYG{g+gp}{\PYGZgt{}\PYGZgt{}\PYGZgt{} }\PYG{k+kn}{from}\PYG{+w}{ }\PYG{n+nn}{psyclone}\PYG{n+nn}{.}\PYG{n+nn}{psyir}\PYG{n+nn}{.}\PYG{n+nn}{frontend}\PYG{n+nn}{.}\PYG{n+nn}{fortran}\PYG{+w}{ }\PYG{k+kn}{import} \PYG{n}{FortranReader}
\PYG{g+gp}{\PYGZgt{}\PYGZgt{}\PYGZgt{} }\PYG{k+kn}{from}\PYG{+w}{ }\PYG{n+nn}{psyclone}\PYG{n+nn}{.}\PYG{n+nn}{psyir}\PYG{n+nn}{.}\PYG{n+nn}{backend}\PYG{n+nn}{.}\PYG{n+nn}{fortran}\PYG{+w}{ }\PYG{k+kn}{import} \PYG{n}{FortranWriter}
\PYG{g+gp}{\PYGZgt{}\PYGZgt{}\PYGZgt{} }\PYG{k+kn}{from}\PYG{+w}{ }\PYG{n+nn}{psyclone}\PYG{n+nn}{.}\PYG{n+nn}{psyir}\PYG{n+nn}{.}\PYG{n+nn}{nodes}\PYG{+w}{ }\PYG{k+kn}{import} \PYG{n}{Loop}
\PYG{g+gp}{\PYGZgt{}\PYGZgt{}\PYGZgt{} }\PYG{n}{psyir} \PYG{o}{=} \PYG{n}{FortranReader}\PYG{p}{(}\PYG{p}{)}\PYG{o}{.}\PYG{n}{psyir\PYGZus{}from\PYGZus{}source}\PYG{p}{(}\PYG{l+s+s2}{\PYGZdq{}\PYGZdq{}\PYGZdq{}}
\PYG{g+gp}{... }\PYG{l+s+s2}{program do\PYGZus{}loop}
\PYG{g+gp}{... }\PYG{l+s+s2}{    real, dimension(10) :: A}
\PYG{g+gp}{... }\PYG{l+s+s2}{    integer i}
\PYG{g+gp}{... }\PYG{l+s+s2}{    do i = 1, 10}
\PYG{g+gp}{... }\PYG{l+s+s2}{      A(i) = i}
\PYG{g+gp}{... }\PYG{l+s+s2}{    end do}
\PYG{g+gp}{... }\PYG{l+s+s2}{end program do\PYGZus{}loop}
\PYG{g+gp}{... }\PYG{l+s+s2}{\PYGZdq{}\PYGZdq{}\PYGZdq{}}\PYG{p}{)}
\PYG{g+gp}{\PYGZgt{}\PYGZgt{}\PYGZgt{} }\PYG{n}{ptrans} \PYG{o}{=} \PYG{n}{TransInfo}\PYG{p}{(}\PYG{p}{)}\PYG{o}{.}\PYG{n}{get\PYGZus{}trans\PYGZus{}name}\PYG{p}{(}\PYG{l+s+s1}{\PYGZsq{}}\PYG{l+s+s1}{ACCParallelTrans}\PYG{l+s+s1}{\PYGZsq{}}\PYG{p}{)}
\PYG{g+gp}{\PYGZgt{}\PYGZgt{}\PYGZgt{}}
\PYG{g+gp}{\PYGZgt{}\PYGZgt{}\PYGZgt{} }\PYG{c+c1}{\PYGZsh{} Enclose the loop within a OpenACC PARALLEL region}
\PYG{g+gp}{\PYGZgt{}\PYGZgt{}\PYGZgt{} }\PYG{n}{ptrans}\PYG{o}{.}\PYG{n}{apply}\PYG{p}{(}\PYG{n}{psyir}\PYG{o}{.}\PYG{n}{walk}\PYG{p}{(}\PYG{n}{Loop}\PYG{p}{)}\PYG{p}{)}
\PYG{g+gp}{\PYGZgt{}\PYGZgt{}\PYGZgt{} }\PYG{n+nb}{print}\PYG{p}{(}\PYG{n}{FortranWriter}\PYG{p}{(}\PYG{p}{)}\PYG{p}{(}\PYG{n}{psyir}\PYG{p}{)}\PYG{p}{)}
\PYG{g+go}{program do\PYGZus{}loop}
\PYG{g+go}{  real, dimension(10) :: a}
\PYG{g+go}{  integer :: i}

\PYG{g+go}{  !\PYGZdl{}acc parallel default(present)}
\PYG{g+go}{  do i = 1, 10, 1}
\PYG{g+go}{    a(i) = i}
\PYG{g+go}{  enddo}
\PYG{g+go}{  !\PYGZdl{}acc end parallel}

\PYG{g+go}{end program do\PYGZus{}loop}
\end{sphinxVerbatim}


\begin{fulllineitems}

\pysigstartsignatures
\pysiglinewithargsret{\sphinxbfcode{\sphinxupquote{apply}}}{\sphinxparam{\DUrole{n,n}{target\_nodes}}\sphinxparamcomma \sphinxparam{\DUrole{n,n}{options}\DUrole{o,o}{=}\DUrole{default_value}{None}}}{}
\pysigstopsignatures
\sphinxAtStartPar
Encapsulate given nodes with the ACCParallelDirective.
\begin{quote}\begin{description}
\sphinxlineitem{Parameters}\begin{itemize}
\item {} 
\sphinxAtStartPar
\sphinxstyleliteralstrong{\sphinxupquote{target\_nodes}} (\sphinxcode{\sphinxupquote{psyclone.psyir.nodes.Node}} |
List{[}\sphinxcode{\sphinxupquote{psyclone.psyir.nodes.Node}}{]}) \textendash{} a single Node or a list of Nodes.

\item {} 
\sphinxAtStartPar
\sphinxstyleliteralstrong{\sphinxupquote{options}} (\sphinxstyleliteralemphasis{\sphinxupquote{Optional}}\sphinxstyleliteralemphasis{\sphinxupquote{{[}}}\sphinxstyleliteralemphasis{\sphinxupquote{Dict}}\sphinxstyleliteralemphasis{\sphinxupquote{{[}}}\sphinxhref{https://docs.python.org/3/library/stdtypes.html\#str}{\sphinxstyleliteralemphasis{\sphinxupquote{str}}}\sphinxstyleliteralemphasis{\sphinxupquote{, }}\sphinxstyleliteralemphasis{\sphinxupquote{Any}}\sphinxstyleliteralemphasis{\sphinxupquote{{]}}}\sphinxstyleliteralemphasis{\sphinxupquote{{]}}}) \textendash{} a dictionary with options for transformations.

\item {} 
\sphinxAtStartPar
\sphinxstyleliteralstrong{\sphinxupquote{options}}\sphinxstyleliteralstrong{\sphinxupquote{{[}}}\sphinxstyleliteralstrong{\sphinxupquote{"node\sphinxhyphen{}type\sphinxhyphen{}check"}}\sphinxstyleliteralstrong{\sphinxupquote{{]}}} (\sphinxhref{https://docs.python.org/3/library/functions.html\#bool}{\sphinxstyleliteralemphasis{\sphinxupquote{bool}}}) \textendash{} this flag controls if the
type of the nodes enclosed in the region should be tested to
avoid using unsupported nodes inside a region.

\item {} 
\sphinxAtStartPar
\sphinxstyleliteralstrong{\sphinxupquote{options}}\sphinxstyleliteralstrong{\sphinxupquote{{[}}}\sphinxstyleliteralstrong{\sphinxupquote{"default\_present"}}\sphinxstyleliteralstrong{\sphinxupquote{{]}}} (\sphinxhref{https://docs.python.org/3/library/functions.html\#bool}{\sphinxstyleliteralemphasis{\sphinxupquote{bool}}}) \textendash{} this flag controls if the
inserted directive should include the default\_present clause.

\end{itemize}

\end{description}\end{quote}

\end{fulllineitems}


\end{fulllineitems}



\bigskip\hrule\bigskip



\begin{fulllineitems}

\pysigstartsignatures
\pysigline{\sphinxbfcode{\sphinxupquote{class\DUrole{w,w}{  }}}\sphinxcode{\sphinxupquote{psyclone.psyir.transformations.}}\sphinxbfcode{\sphinxupquote{AllArrayAccess2LoopTrans}}}
\pysigstopsignatures
\sphinxAtStartPar
Provides a transformation from a PSyIR Assignment containing
constant index accesses to an array into single\sphinxhyphen{}trip loops. For
example:

\begin{sphinxVerbatim}[commandchars=\\\{\}]
\PYG{g+gp}{\PYGZgt{}\PYGZgt{}\PYGZgt{} }\PYG{k+kn}{from}\PYG{+w}{ }\PYG{n+nn}{psyclone}\PYG{n+nn}{.}\PYG{n+nn}{psyir}\PYG{n+nn}{.}\PYG{n+nn}{transformations}\PYG{+w}{ }\PYG{k+kn}{import} \PYG{n}{AllArrayAccess2LoopTrans}
\PYG{g+gp}{\PYGZgt{}\PYGZgt{}\PYGZgt{} }\PYG{k+kn}{from}\PYG{+w}{ }\PYG{n+nn}{psyclone}\PYG{n+nn}{.}\PYG{n+nn}{psyir}\PYG{n+nn}{.}\PYG{n+nn}{backend}\PYG{n+nn}{.}\PYG{n+nn}{fortran}\PYG{+w}{ }\PYG{k+kn}{import} \PYG{n}{FortranWriter}
\PYG{g+gp}{\PYGZgt{}\PYGZgt{}\PYGZgt{} }\PYG{k+kn}{from}\PYG{+w}{ }\PYG{n+nn}{psyclone}\PYG{n+nn}{.}\PYG{n+nn}{psyir}\PYG{n+nn}{.}\PYG{n+nn}{frontend}\PYG{n+nn}{.}\PYG{n+nn}{fortran}\PYG{+w}{ }\PYG{k+kn}{import} \PYG{n}{FortranReader}
\PYG{g+gp}{\PYGZgt{}\PYGZgt{}\PYGZgt{} }\PYG{k+kn}{from}\PYG{+w}{ }\PYG{n+nn}{psyclone}\PYG{n+nn}{.}\PYG{n+nn}{psyir}\PYG{n+nn}{.}\PYG{n+nn}{nodes}\PYG{+w}{ }\PYG{k+kn}{import} \PYG{n}{Assignment}
\PYG{g+gp}{\PYGZgt{}\PYGZgt{}\PYGZgt{} }\PYG{n}{code} \PYG{o}{=} \PYG{p}{(}\PYG{l+s+s2}{\PYGZdq{}}\PYG{l+s+s2}{program example}\PYG{l+s+se}{\PYGZbs{}n}\PYG{l+s+s2}{\PYGZdq{}}
\PYG{g+gp}{... }        \PYG{l+s+s2}{\PYGZdq{}}\PYG{l+s+s2}{  real a(10,10), b(10,10)}\PYG{l+s+se}{\PYGZbs{}n}\PYG{l+s+s2}{\PYGZdq{}}
\PYG{g+gp}{... }        \PYG{l+s+s2}{\PYGZdq{}}\PYG{l+s+s2}{  integer :: n}\PYG{l+s+se}{\PYGZbs{}n}\PYG{l+s+s2}{\PYGZdq{}}
\PYG{g+gp}{... }        \PYG{l+s+s2}{\PYGZdq{}}\PYG{l+s+s2}{  a(1,n\PYGZhy{}1) = b(1,n\PYGZhy{}1)}\PYG{l+s+se}{\PYGZbs{}n}\PYG{l+s+s2}{\PYGZdq{}}
\PYG{g+gp}{... }        \PYG{l+s+s2}{\PYGZdq{}}\PYG{l+s+s2}{end program example}\PYG{l+s+se}{\PYGZbs{}n}\PYG{l+s+s2}{\PYGZdq{}}\PYG{p}{)}
\PYG{g+gp}{\PYGZgt{}\PYGZgt{}\PYGZgt{} }\PYG{n}{psyir} \PYG{o}{=} \PYG{n}{FortranReader}\PYG{p}{(}\PYG{p}{)}\PYG{o}{.}\PYG{n}{psyir\PYGZus{}from\PYGZus{}source}\PYG{p}{(}\PYG{n}{code}\PYG{p}{)}
\PYG{g+gp}{\PYGZgt{}\PYGZgt{}\PYGZgt{} }\PYG{n}{assignment} \PYG{o}{=} \PYG{n}{psyir}\PYG{o}{.}\PYG{n}{walk}\PYG{p}{(}\PYG{n}{Assignment}\PYG{p}{)}\PYG{p}{[}\PYG{l+m+mi}{0}\PYG{p}{]}
\PYG{g+gp}{\PYGZgt{}\PYGZgt{}\PYGZgt{} }\PYG{n}{AllArrayAccess2LoopTrans}\PYG{p}{(}\PYG{p}{)}\PYG{o}{.}\PYG{n}{apply}\PYG{p}{(}\PYG{n}{assignment}\PYG{p}{)}
\PYG{g+gp}{\PYGZgt{}\PYGZgt{}\PYGZgt{} }\PYG{n+nb}{print}\PYG{p}{(}\PYG{n}{FortranWriter}\PYG{p}{(}\PYG{p}{)}\PYG{p}{(}\PYG{n}{psyir}\PYG{p}{)}\PYG{p}{)}
\PYG{g+go}{program example}
\PYG{g+go}{  real, dimension(10,10) :: a}
\PYG{g+go}{  real, dimension(10,10) :: b}
\PYG{g+go}{  integer :: n}
\PYG{g+go}{  integer :: idx}
\PYG{g+go}{  integer :: idx\PYGZus{}1}

\PYG{g+go}{  do idx = 1, 1, 1}
\PYG{g+go}{    do idx\PYGZus{}1 = n \PYGZhy{} 1, n \PYGZhy{} 1, 1}
\PYG{g+go}{      a(idx,idx\PYGZus{}1) = b(idx,idx\PYGZus{}1)}
\PYG{g+go}{    enddo}
\PYG{g+go}{  enddo}

\PYG{g+go}{end program example}
\end{sphinxVerbatim}


\begin{fulllineitems}

\pysigstartsignatures
\pysiglinewithargsret{\sphinxbfcode{\sphinxupquote{apply}}}{\sphinxparam{\DUrole{n,n}{node}}\sphinxparamcomma \sphinxparam{\DUrole{n,n}{options}\DUrole{o,o}{=}\DUrole{default_value}{None}}}{}
\pysigstopsignatures
\sphinxAtStartPar
Apply the AllArrayAccess2Loop transformation if the supplied
node is an Assignment with an Array Reference on its
left\sphinxhyphen{}hand\sphinxhyphen{}side. Each constant array index access (i.e. one not
containing a loop iterator or a range) is then transformed into
an iterator and the assignment placed within a single\sphinxhyphen{}trip loop,
subject to any constraints in the ArrayAccess2Loop transformation.

\sphinxAtStartPar
If any of the AllArrayAccess2Loop constraints are not
satisfied for a loop index then this transformation does
nothing for that index and silently moves to the next.
\begin{quote}\begin{description}
\sphinxlineitem{Parameters}\begin{itemize}
\item {} 
\sphinxAtStartPar
\sphinxstyleliteralstrong{\sphinxupquote{node}} (\sphinxcode{\sphinxupquote{psyclone.psyir.nodes.Assignment}}) \textendash{} an assignment.

\item {} 
\sphinxAtStartPar
\sphinxstyleliteralstrong{\sphinxupquote{options}} (\sphinxstyleliteralemphasis{\sphinxupquote{Optional}}\sphinxstyleliteralemphasis{\sphinxupquote{{[}}}\sphinxstyleliteralemphasis{\sphinxupquote{Dict}}\sphinxstyleliteralemphasis{\sphinxupquote{{[}}}\sphinxhref{https://docs.python.org/3/library/stdtypes.html\#str}{\sphinxstyleliteralemphasis{\sphinxupquote{str}}}\sphinxstyleliteralemphasis{\sphinxupquote{, }}\sphinxstyleliteralemphasis{\sphinxupquote{Any}}\sphinxstyleliteralemphasis{\sphinxupquote{{]}}}\sphinxstyleliteralemphasis{\sphinxupquote{{]}}}) \textendash{} a dictionary with options for transformations.
This is an optional argument that defaults to None.

\end{itemize}

\end{description}\end{quote}

\end{fulllineitems}


\end{fulllineitems}



\bigskip\hrule\bigskip



\begin{fulllineitems}

\pysigstartsignatures
\pysigline{\sphinxbfcode{\sphinxupquote{class\DUrole{w,w}{  }}}\sphinxcode{\sphinxupquote{psyclone.psyir.transformations.}}\sphinxbfcode{\sphinxupquote{ArrayAccess2LoopTrans}}}
\pysigstopsignatures
\sphinxAtStartPar
Provides a transformation to transform a constant index access to
an array (i.e. one that does not contain a loop iterator) to a
single trip loop. For example:

\begin{sphinxVerbatim}[commandchars=\\\{\}]
\PYG{g+gp}{\PYGZgt{}\PYGZgt{}\PYGZgt{} }\PYG{k+kn}{from}\PYG{+w}{ }\PYG{n+nn}{psyclone}\PYG{n+nn}{.}\PYG{n+nn}{psyir}\PYG{n+nn}{.}\PYG{n+nn}{transformations}\PYG{+w}{ }\PYG{k+kn}{import} \PYG{n}{ArrayAccess2LoopTrans}
\PYG{g+gp}{\PYGZgt{}\PYGZgt{}\PYGZgt{} }\PYG{k+kn}{from}\PYG{+w}{ }\PYG{n+nn}{psyclone}\PYG{n+nn}{.}\PYG{n+nn}{psyir}\PYG{n+nn}{.}\PYG{n+nn}{backend}\PYG{n+nn}{.}\PYG{n+nn}{fortran}\PYG{+w}{ }\PYG{k+kn}{import} \PYG{n}{FortranWriter}
\PYG{g+gp}{\PYGZgt{}\PYGZgt{}\PYGZgt{} }\PYG{k+kn}{from}\PYG{+w}{ }\PYG{n+nn}{psyclone}\PYG{n+nn}{.}\PYG{n+nn}{psyir}\PYG{n+nn}{.}\PYG{n+nn}{frontend}\PYG{n+nn}{.}\PYG{n+nn}{fortran}\PYG{+w}{ }\PYG{k+kn}{import} \PYG{n}{FortranReader}
\PYG{g+gp}{\PYGZgt{}\PYGZgt{}\PYGZgt{} }\PYG{k+kn}{from}\PYG{+w}{ }\PYG{n+nn}{psyclone}\PYG{n+nn}{.}\PYG{n+nn}{psyir}\PYG{n+nn}{.}\PYG{n+nn}{nodes}\PYG{+w}{ }\PYG{k+kn}{import} \PYG{n}{Assignment}
\PYG{g+gp}{\PYGZgt{}\PYGZgt{}\PYGZgt{} }\PYG{n}{code} \PYG{o}{=} \PYG{p}{(}\PYG{l+s+s2}{\PYGZdq{}}\PYG{l+s+s2}{program example}\PYG{l+s+se}{\PYGZbs{}n}\PYG{l+s+s2}{\PYGZdq{}}
\PYG{g+gp}{... }        \PYG{l+s+s2}{\PYGZdq{}}\PYG{l+s+s2}{  real a(10)}\PYG{l+s+se}{\PYGZbs{}n}\PYG{l+s+s2}{\PYGZdq{}}
\PYG{g+gp}{... }        \PYG{l+s+s2}{\PYGZdq{}}\PYG{l+s+s2}{  a(1) = 0.0}\PYG{l+s+se}{\PYGZbs{}n}\PYG{l+s+s2}{\PYGZdq{}}
\PYG{g+gp}{... }        \PYG{l+s+s2}{\PYGZdq{}}\PYG{l+s+s2}{end program example}\PYG{l+s+se}{\PYGZbs{}n}\PYG{l+s+s2}{\PYGZdq{}}\PYG{p}{)}
\PYG{g+gp}{\PYGZgt{}\PYGZgt{}\PYGZgt{} }\PYG{n}{psyir} \PYG{o}{=} \PYG{n}{FortranReader}\PYG{p}{(}\PYG{p}{)}\PYG{o}{.}\PYG{n}{psyir\PYGZus{}from\PYGZus{}source}\PYG{p}{(}\PYG{n}{code}\PYG{p}{)}
\PYG{g+gp}{\PYGZgt{}\PYGZgt{}\PYGZgt{} }\PYG{n}{assignment} \PYG{o}{=} \PYG{n}{psyir}\PYG{o}{.}\PYG{n}{walk}\PYG{p}{(}\PYG{n}{Assignment}\PYG{p}{)}\PYG{p}{[}\PYG{l+m+mi}{0}\PYG{p}{]}
\PYG{g+gp}{\PYGZgt{}\PYGZgt{}\PYGZgt{} }\PYG{n}{ArrayAccess2LoopTrans}\PYG{p}{(}\PYG{p}{)}\PYG{o}{.}\PYG{n}{apply}\PYG{p}{(}\PYG{n}{assignment}\PYG{o}{.}\PYG{n}{lhs}\PYG{o}{.}\PYG{n}{children}\PYG{p}{[}\PYG{l+m+mi}{0}\PYG{p}{]}\PYG{p}{)}
\PYG{g+gp}{\PYGZgt{}\PYGZgt{}\PYGZgt{} }\PYG{n+nb}{print}\PYG{p}{(}\PYG{n}{FortranWriter}\PYG{p}{(}\PYG{p}{)}\PYG{p}{(}\PYG{n}{psyir}\PYG{p}{)}\PYG{p}{)}
\PYG{g+go}{program example}
\PYG{g+go}{  real, dimension(10) :: a}
\PYG{g+go}{  integer :: ji}

\PYG{g+go}{  do ji = 1, 1, 1}
\PYG{g+go}{    a(ji) = 0.0}
\PYG{g+go}{  enddo}

\PYG{g+go}{end program example}
\end{sphinxVerbatim}


\begin{fulllineitems}

\pysigstartsignatures
\pysiglinewithargsret{\sphinxbfcode{\sphinxupquote{apply}}}{\sphinxparam{\DUrole{n,n}{node}}\sphinxparamcomma \sphinxparam{\DUrole{n,n}{options}\DUrole{o,o}{=}\DUrole{default_value}{None}}}{}
\pysigstopsignatures
\sphinxAtStartPar
Apply the ArrayAccess2Loop transformation if the supplied node
is an access to an array index within an Array Reference that
is on the left\sphinxhyphen{}hand\sphinxhyphen{}side of an Assignment node. The access
must be a scalar (i.e. not a range) and must not include a
loop variable (as we are transforming a single access to a
loop).

\sphinxAtStartPar
If the constraints are satisfied then the array access is replaced
with a loop iterator and a single trip loop. The new loop will be
placed immediately around the assignment.
\begin{quote}\begin{description}
\sphinxlineitem{Parameters}\begin{itemize}
\item {} 
\sphinxAtStartPar
\sphinxstyleliteralstrong{\sphinxupquote{node}} (\sphinxcode{\sphinxupquote{psyclone.psyir.nodes.Node}}) \textendash{} an array index.

\item {} 
\sphinxAtStartPar
\sphinxstyleliteralstrong{\sphinxupquote{options}} (\sphinxstyleliteralemphasis{\sphinxupquote{Optional}}\sphinxstyleliteralemphasis{\sphinxupquote{{[}}}\sphinxstyleliteralemphasis{\sphinxupquote{Dict}}\sphinxstyleliteralemphasis{\sphinxupquote{{[}}}\sphinxhref{https://docs.python.org/3/library/stdtypes.html\#str}{\sphinxstyleliteralemphasis{\sphinxupquote{str}}}\sphinxstyleliteralemphasis{\sphinxupquote{, }}\sphinxstyleliteralemphasis{\sphinxupquote{Any}}\sphinxstyleliteralemphasis{\sphinxupquote{{]}}}\sphinxstyleliteralemphasis{\sphinxupquote{{]}}}) \textendash{} a dictionary with options for transformations.
This is an optional argument that defaults to None.

\end{itemize}

\end{description}\end{quote}

\end{fulllineitems}


\end{fulllineitems}



\bigskip\hrule\bigskip



\begin{fulllineitems}

\pysigstartsignatures
\pysigline{\sphinxbfcode{\sphinxupquote{class\DUrole{w,w}{  }}}\sphinxcode{\sphinxupquote{psyclone.psyir.transformations.}}\sphinxbfcode{\sphinxupquote{ArrayAssignment2LoopsTrans}}}
\pysigstopsignatures
\sphinxAtStartPar
Provides a transformation from a PSyIR Array Range to a PSyIR
Loop. For example:

\begin{sphinxVerbatim}[commandchars=\\\{\}]
\PYG{g+gp}{\PYGZgt{}\PYGZgt{}\PYGZgt{} }\PYG{k+kn}{from}\PYG{+w}{ }\PYG{n+nn}{psyclone}\PYG{n+nn}{.}\PYG{n+nn}{psyir}\PYG{n+nn}{.}\PYG{n+nn}{frontend}\PYG{n+nn}{.}\PYG{n+nn}{fortran}\PYG{+w}{ }\PYG{k+kn}{import} \PYG{n}{FortranReader}
\PYG{g+gp}{\PYGZgt{}\PYGZgt{}\PYGZgt{} }\PYG{k+kn}{from}\PYG{+w}{ }\PYG{n+nn}{psyclone}\PYG{n+nn}{.}\PYG{n+nn}{psyir}\PYG{n+nn}{.}\PYG{n+nn}{nodes}\PYG{+w}{ }\PYG{k+kn}{import} \PYG{n}{Assignment}
\PYG{g+gp}{\PYGZgt{}\PYGZgt{}\PYGZgt{} }\PYG{k+kn}{from}\PYG{+w}{ }\PYG{n+nn}{psyclone}\PYG{n+nn}{.}\PYG{n+nn}{psyir}\PYG{n+nn}{.}\PYG{n+nn}{transformations}\PYG{+w}{ }\PYG{k+kn}{import} \PYG{n}{ArrayAssignment2LoopsTrans}
\PYG{g+gp}{\PYGZgt{}\PYGZgt{}\PYGZgt{} }\PYG{n}{code} \PYG{o}{=} \PYG{l+s+s2}{\PYGZdq{}\PYGZdq{}\PYGZdq{}}
\PYG{g+gp}{... }\PYG{l+s+s2}{subroutine sub()}
\PYG{g+gp}{... }\PYG{l+s+s2}{  real :: tmp(10)}
\PYG{g+gp}{... }\PYG{l+s+s2}{  tmp(:) = tmp(:) + 3}
\PYG{g+gp}{... }\PYG{l+s+s2}{end subroutine sub}\PYG{l+s+s2}{\PYGZdq{}\PYGZdq{}\PYGZdq{}}
\PYG{g+gp}{\PYGZgt{}\PYGZgt{}\PYGZgt{} }\PYG{n}{psyir} \PYG{o}{=} \PYG{n}{FortranReader}\PYG{p}{(}\PYG{p}{)}\PYG{o}{.}\PYG{n}{psyir\PYGZus{}from\PYGZus{}source}\PYG{p}{(}\PYG{n}{code}\PYG{p}{)}
\PYG{g+gp}{\PYGZgt{}\PYGZgt{}\PYGZgt{} }\PYG{n}{assignment} \PYG{o}{=} \PYG{n}{psyir}\PYG{o}{.}\PYG{n}{walk}\PYG{p}{(}\PYG{n}{Assignment}\PYG{p}{)}\PYG{p}{[}\PYG{l+m+mi}{0}\PYG{p}{]}
\PYG{g+gp}{\PYGZgt{}\PYGZgt{}\PYGZgt{} }\PYG{n}{trans} \PYG{o}{=} \PYG{n}{ArrayAssignment2LoopsTrans}\PYG{p}{(}\PYG{p}{)}
\PYG{g+gp}{\PYGZgt{}\PYGZgt{}\PYGZgt{} }\PYG{n}{trans}\PYG{o}{.}\PYG{n}{apply}\PYG{p}{(}\PYG{n}{assignment}\PYG{p}{)}
\PYG{g+gp}{\PYGZgt{}\PYGZgt{}\PYGZgt{} }\PYG{n+nb}{print}\PYG{p}{(}\PYG{n}{psyir}\PYG{o}{.}\PYG{n}{debug\PYGZus{}string}\PYG{p}{(}\PYG{p}{)}\PYG{p}{)}
\PYG{g+go}{subroutine sub()}
\PYG{g+go}{  real, dimension(10) :: tmp}
\PYG{g+go}{  integer :: idx}

\PYG{g+go}{  do idx = LBOUND(tmp, dim=1), UBOUND(tmp, dim=1), 1}
\PYG{g+go}{    tmp(idx) = tmp(idx) + 3}
\PYG{g+go}{  enddo}

\PYG{g+go}{end subroutine sub}
\end{sphinxVerbatim}

\sphinxAtStartPar
By default the transformation will reject character arrays, though this
can be overriden by setting the ‘allow\_string’ option to True. Note that
PSyclone expresses syntax such as \sphinxtitleref{character(LEN=100)} as
UnsupportedFortranType, and this transformation will convert unknown or
unsupported types to loops.


\begin{fulllineitems}

\pysigstartsignatures
\pysiglinewithargsret{\sphinxbfcode{\sphinxupquote{apply}}}{\sphinxparam{\DUrole{n,n}{node}}\sphinxparamcomma \sphinxparam{\DUrole{n,n}{options}\DUrole{o,o}{=}\DUrole{default_value}{None}}}{}
\pysigstopsignatures
\sphinxAtStartPar
Apply the transformation to the specified array assignment node.
Each range node within the assignment is replaced with an explicit
loop. The bounds of the loop are determined from the bounds of the
array range on the left hand side of the assignment.
\begin{quote}\begin{description}
\sphinxlineitem{Parameters}\begin{itemize}
\item {} 
\sphinxAtStartPar
\sphinxstyleliteralstrong{\sphinxupquote{node}} (\sphinxcode{\sphinxupquote{psyclone.psyir.nodes.Assignment}}) \textendash{} an Assignment node.

\item {} 
\sphinxAtStartPar
\sphinxstyleliteralstrong{\sphinxupquote{options}}\sphinxstyleliteralstrong{\sphinxupquote{{[}}}\sphinxstyleliteralstrong{\sphinxupquote{"allow\_string"}}\sphinxstyleliteralstrong{\sphinxupquote{{]}}} (\sphinxhref{https://docs.python.org/3/library/functions.html\#bool}{\sphinxstyleliteralemphasis{\sphinxupquote{bool}}}) \textendash{} whether to allow the
transformation on a character type array range. Defaults to False.

\item {} 
\sphinxAtStartPar
\sphinxstyleliteralstrong{\sphinxupquote{options}}\sphinxstyleliteralstrong{\sphinxupquote{{[}}}\sphinxstyleliteralstrong{\sphinxupquote{"verbose"}}\sphinxstyleliteralstrong{\sphinxupquote{{]}}} (\sphinxhref{https://docs.python.org/3/library/functions.html\#bool}{\sphinxstyleliteralemphasis{\sphinxupquote{bool}}}) \textendash{} log the reason the validation failed,
at the moment with a comment in the provided PSyIR node.

\end{itemize}

\end{description}\end{quote}

\end{fulllineitems}


\end{fulllineitems}



\bigskip\hrule\bigskip



\begin{fulllineitems}

\pysigstartsignatures
\pysigline{\sphinxbfcode{\sphinxupquote{class\DUrole{w,w}{  }}}\sphinxcode{\sphinxupquote{psyclone.psyir.transformations.}}\sphinxbfcode{\sphinxupquote{ChunkLoopTrans}}}
\pysigstopsignatures
\sphinxAtStartPar
Apply a chunking transformation to a loop (in order to permit a
chunked parallelisation). For example:

\begin{sphinxVerbatim}[commandchars=\\\{\}]
\PYG{g+gp}{\PYGZgt{}\PYGZgt{}\PYGZgt{} }\PYG{k+kn}{from}\PYG{+w}{ }\PYG{n+nn}{psyclone}\PYG{n+nn}{.}\PYG{n+nn}{psyir}\PYG{n+nn}{.}\PYG{n+nn}{frontend}\PYG{n+nn}{.}\PYG{n+nn}{fortran}\PYG{+w}{ }\PYG{k+kn}{import} \PYG{n}{FortranReader}
\PYG{g+gp}{\PYGZgt{}\PYGZgt{}\PYGZgt{} }\PYG{k+kn}{from}\PYG{+w}{ }\PYG{n+nn}{psyclone}\PYG{n+nn}{.}\PYG{n+nn}{psyir}\PYG{n+nn}{.}\PYG{n+nn}{nodes}\PYG{+w}{ }\PYG{k+kn}{import} \PYG{n}{Loop}
\PYG{g+gp}{\PYGZgt{}\PYGZgt{}\PYGZgt{} }\PYG{k+kn}{from}\PYG{+w}{ }\PYG{n+nn}{psyclone}\PYG{n+nn}{.}\PYG{n+nn}{psyir}\PYG{n+nn}{.}\PYG{n+nn}{transformations}\PYG{+w}{ }\PYG{k+kn}{import} \PYG{n}{ChunkLoopTrans}
\PYG{g+gp}{\PYGZgt{}\PYGZgt{}\PYGZgt{} }\PYG{n}{psyir} \PYG{o}{=} \PYG{n}{FortranReader}\PYG{p}{(}\PYG{p}{)}\PYG{o}{.}\PYG{n}{psyir\PYGZus{}from\PYGZus{}source}\PYG{p}{(}\PYG{l+s+s2}{\PYGZdq{}\PYGZdq{}\PYGZdq{}}
\PYG{g+gp}{... }\PYG{l+s+s2}{subroutine sub()}
\PYG{g+gp}{... }\PYG{l+s+s2}{    integer :: ji, tmp(100)}
\PYG{g+gp}{... }\PYG{l+s+s2}{    do ji=1, 100}
\PYG{g+gp}{... }\PYG{l+s+s2}{        tmp(ji) = 2 * ji}
\PYG{g+gp}{... }\PYG{l+s+s2}{    enddo}
\PYG{g+gp}{... }\PYG{l+s+s2}{end subroutine sub}\PYG{l+s+s2}{\PYGZdq{}\PYGZdq{}\PYGZdq{}}\PYG{p}{)}
\PYG{g+gp}{\PYGZgt{}\PYGZgt{}\PYGZgt{} }\PYG{n}{loop} \PYG{o}{=} \PYG{n}{psyir}\PYG{o}{.}\PYG{n}{walk}\PYG{p}{(}\PYG{n}{Loop}\PYG{p}{)}\PYG{p}{[}\PYG{l+m+mi}{0}\PYG{p}{]}
\PYG{g+gp}{\PYGZgt{}\PYGZgt{}\PYGZgt{} }\PYG{n}{ChunkLoopTrans}\PYG{p}{(}\PYG{p}{)}\PYG{o}{.}\PYG{n}{apply}\PYG{p}{(}\PYG{n}{loop}\PYG{p}{)}
\end{sphinxVerbatim}

\sphinxAtStartPar
will generate:

\begin{sphinxVerbatim}[commandchars=\\\{\}]
\PYG{k}{subroutine }\PYG{n}{sub}\PYG{p}{(}\PYG{p}{)}
\PYG{+w}{    }\PYG{k+kt}{integer}\PYG{+w}{ }\PYG{k+kd}{::}\PYG{+w}{ }\PYG{n}{ji}
\PYG{+w}{    }\PYG{k+kt}{integer}\PYG{p}{,}\PYG{+w}{ }\PYG{k}{dimension}\PYG{p}{(}\PYG{l+m+mi}{100}\PYG{p}{)}\PYG{+w}{ }\PYG{k+kd}{::}\PYG{+w}{ }\PYG{n}{tmp}
\PYG{+w}{    }\PYG{k+kt}{integer}\PYG{+w}{ }\PYG{k+kd}{::}\PYG{+w}{ }\PYG{n}{ji\PYGZus{}el\PYGZus{}inner}
\PYG{+w}{    }\PYG{k+kt}{integer}\PYG{+w}{ }\PYG{k+kd}{::}\PYG{+w}{ }\PYG{n}{ji\PYGZus{}out\PYGZus{}var}
\PYG{+w}{    }\PYG{k}{do }\PYG{n}{ji\PYGZus{}out\PYGZus{}var}\PYG{+w}{ }\PYG{o}{=}\PYG{+w}{ }\PYG{l+m+mi}{1}\PYG{p}{,}\PYG{+w}{ }\PYG{l+m+mi}{100}\PYG{p}{,}\PYG{+w}{ }\PYG{l+m+mi}{32}
\PYG{+w}{        }\PYG{n}{ji\PYGZus{}el\PYGZus{}inner}\PYG{+w}{ }\PYG{o}{=}\PYG{+w}{ }\PYG{n+nb}{MIN}\PYG{p}{(}\PYG{n}{ji\PYGZus{}out\PYGZus{}var}\PYG{+w}{ }\PYG{o}{+}\PYG{+w}{ }\PYG{p}{(}\PYG{l+m+mi}{32}\PYG{+w}{ }\PYG{o}{\PYGZhy{}}\PYG{+w}{ }\PYG{l+m+mi}{1}\PYG{p}{)}\PYG{p}{,}\PYG{+w}{ }\PYG{l+m+mi}{100}\PYG{p}{)}
\PYG{+w}{        }\PYG{k}{do }\PYG{n}{ji}\PYG{+w}{ }\PYG{o}{=}\PYG{+w}{ }\PYG{n}{ji\PYGZus{}out\PYGZus{}var}\PYG{p}{,}\PYG{+w}{ }\PYG{n}{ji\PYGZus{}el\PYGZus{}inner}\PYG{p}{,}\PYG{+w}{ }\PYG{l+m+mi}{1}
\PYG{+w}{            }\PYG{n}{tmp}\PYG{p}{(}\PYG{n}{ji}\PYG{p}{)}\PYG{+w}{ }\PYG{o}{=}\PYG{+w}{ }\PYG{l+m+mi}{2}\PYG{+w}{ }\PYG{o}{*}\PYG{+w}{ }\PYG{n}{ji}
\PYG{+w}{        }\PYG{k}{enddo}
\PYG{k}{    }\PYG{k}{enddo}
\PYG{k}{end }\PYG{k}{subroutine }\PYG{n}{sub}
\end{sphinxVerbatim}


\begin{fulllineitems}

\pysigstartsignatures
\pysiglinewithargsret{\sphinxbfcode{\sphinxupquote{apply}}}{\sphinxparam{\DUrole{n,n}{node}}\sphinxparamcomma \sphinxparam{\DUrole{n,n}{options}\DUrole{o,o}{=}\DUrole{default_value}{None}}}{}
\pysigstopsignatures
\sphinxAtStartPar
Converts the given Loop node into a nested loop where the outer
loop is over chunks and the inner loop is over each individual element
of the chunk.
\begin{quote}\begin{description}
\sphinxlineitem{Parameters}\begin{itemize}
\item {} 
\sphinxAtStartPar
\sphinxstyleliteralstrong{\sphinxupquote{node}} (\sphinxcode{\sphinxupquote{psyclone.psyir.nodes.Loop}}) \textendash{} the loop to transform.

\item {} 
\sphinxAtStartPar
\sphinxstyleliteralstrong{\sphinxupquote{options}} (\sphinxstyleliteralemphasis{\sphinxupquote{Optional}}\sphinxstyleliteralemphasis{\sphinxupquote{{[}}}\sphinxstyleliteralemphasis{\sphinxupquote{Dict}}\sphinxstyleliteralemphasis{\sphinxupquote{{[}}}\sphinxhref{https://docs.python.org/3/library/stdtypes.html\#str}{\sphinxstyleliteralemphasis{\sphinxupquote{str}}}\sphinxstyleliteralemphasis{\sphinxupquote{, }}\sphinxstyleliteralemphasis{\sphinxupquote{Any}}\sphinxstyleliteralemphasis{\sphinxupquote{{]}}}\sphinxstyleliteralemphasis{\sphinxupquote{{]}}}) \textendash{} a dict with options for transformations.

\item {} 
\sphinxAtStartPar
\sphinxstyleliteralstrong{\sphinxupquote{options}}\sphinxstyleliteralstrong{\sphinxupquote{{[}}}\sphinxstyleliteralstrong{\sphinxupquote{"chunksize"}}\sphinxstyleliteralstrong{\sphinxupquote{{]}}} (\sphinxhref{https://docs.python.org/3/library/functions.html\#int}{\sphinxstyleliteralemphasis{\sphinxupquote{int}}}) \textendash{} The size to chunk over for this                 transformation. If not specified, the value 32 is used.

\end{itemize}

\end{description}\end{quote}

\end{fulllineitems}


\end{fulllineitems}



\bigskip\hrule\bigskip



\begin{fulllineitems}

\pysigstartsignatures
\pysigline{\sphinxbfcode{\sphinxupquote{class\DUrole{w,w}{  }}}\sphinxcode{\sphinxupquote{psyclone.transformations.}}\sphinxbfcode{\sphinxupquote{ColourTrans}}}
\pysigstopsignatures
\sphinxAtStartPar
Apply a colouring transformation to a loop (in order to permit a
subsequent parallelisation over colours). For example:

\begin{sphinxVerbatim}[commandchars=\\\{\}]
\PYG{g+gp}{\PYGZgt{}\PYGZgt{}\PYGZgt{} }\PYG{n}{invoke} \PYG{o}{=} \PYG{o}{.}\PYG{o}{.}\PYG{o}{.}
\PYG{g+gp}{\PYGZgt{}\PYGZgt{}\PYGZgt{} }\PYG{n}{schedule} \PYG{o}{=} \PYG{n}{invoke}\PYG{o}{.}\PYG{n}{schedule}
\PYG{g+gp}{\PYGZgt{}\PYGZgt{}\PYGZgt{}}
\PYG{g+gp}{\PYGZgt{}\PYGZgt{}\PYGZgt{} }\PYG{n}{ctrans} \PYG{o}{=} \PYG{n}{ColourTrans}\PYG{p}{(}\PYG{p}{)}
\PYG{g+gp}{\PYGZgt{}\PYGZgt{}\PYGZgt{}}
\PYG{g+gp}{\PYGZgt{}\PYGZgt{}\PYGZgt{} }\PYG{c+c1}{\PYGZsh{} Colour all of the loops}
\PYG{g+gp}{\PYGZgt{}\PYGZgt{}\PYGZgt{} }\PYG{k}{for} \PYG{n}{child} \PYG{o+ow}{in} \PYG{n}{schedule}\PYG{o}{.}\PYG{n}{children}\PYG{p}{:}
\PYG{g+gp}{\PYGZgt{}\PYGZgt{}\PYGZgt{} }    \PYG{n}{ctrans}\PYG{o}{.}\PYG{n}{apply}\PYG{p}{(}\PYG{n}{child}\PYG{p}{)}
\PYG{g+gp}{\PYGZgt{}\PYGZgt{}\PYGZgt{}}
\PYG{g+gp}{\PYGZgt{}\PYGZgt{}\PYGZgt{} }\PYG{c+c1}{\PYGZsh{} Uncomment the following line to see a text view of the schedule}
\PYG{g+gp}{\PYGZgt{}\PYGZgt{}\PYGZgt{} }\PYG{c+c1}{\PYGZsh{} print(schedule.view())}
\end{sphinxVerbatim}


\begin{fulllineitems}

\pysigstartsignatures
\pysiglinewithargsret{\sphinxbfcode{\sphinxupquote{apply}}}{\sphinxparam{\DUrole{n,n}{node}}\sphinxparamcomma \sphinxparam{\DUrole{n,n}{options}\DUrole{o,o}{=}\DUrole{default_value}{None}}}{}
\pysigstopsignatures
\sphinxAtStartPar
Converts the Loop represented by \sphinxcode{\sphinxupquote{node}} into a
nested loop where the outer loop is over colours and the inner
loop is over cells of that colour.
\begin{quote}\begin{description}
\sphinxlineitem{Parameters}\begin{itemize}
\item {} 
\sphinxAtStartPar
\sphinxstyleliteralstrong{\sphinxupquote{node}} (\sphinxcode{\sphinxupquote{psyclone.psyir.nodes.Loop}}) \textendash{} the loop to transform.

\item {} 
\sphinxAtStartPar
\sphinxstyleliteralstrong{\sphinxupquote{options}} (\sphinxstyleliteralemphasis{\sphinxupquote{Optional}}\sphinxstyleliteralemphasis{\sphinxupquote{{[}}}\sphinxstyleliteralemphasis{\sphinxupquote{Dict}}\sphinxstyleliteralemphasis{\sphinxupquote{{[}}}\sphinxhref{https://docs.python.org/3/library/stdtypes.html\#str}{\sphinxstyleliteralemphasis{\sphinxupquote{str}}}\sphinxstyleliteralemphasis{\sphinxupquote{, }}\sphinxstyleliteralemphasis{\sphinxupquote{Any}}\sphinxstyleliteralemphasis{\sphinxupquote{{]}}}\sphinxstyleliteralemphasis{\sphinxupquote{{]}}}) \textendash{} options for the transformation.

\end{itemize}

\end{description}\end{quote}

\end{fulllineitems}


\end{fulllineitems}



\bigskip\hrule\bigskip



\begin{fulllineitems}

\pysigstartsignatures
\pysigline{\sphinxbfcode{\sphinxupquote{class\DUrole{w,w}{  }}}\sphinxcode{\sphinxupquote{psyclone.psyir.transformations.}}\sphinxbfcode{\sphinxupquote{DotProduct2CodeTrans}}}
\pysigstopsignatures
\sphinxAtStartPar
Provides a transformation from a PSyIR DOT\_PRODUCT Operator node to
equivalent code in a PSyIR tree. Validity checks are also
performed.

\sphinxAtStartPar
If \sphinxcode{\sphinxupquote{R}} is a scalar and \sphinxcode{\sphinxupquote{A}}, and \sphinxcode{\sphinxupquote{B}} have dimension \sphinxcode{\sphinxupquote{N}},
The transformation replaces:

\begin{sphinxVerbatim}[commandchars=\\\{\}]
\PYG{n}{R}\PYG{+w}{ }\PYG{o}{=}\PYG{+w}{ }\PYG{p}{.}\PYG{p}{.}\PYG{p}{.}\PYG{+w}{ }\PYG{n+nb}{DOT\PYGZus{}PRODUCT}\PYG{p}{(}\PYG{n}{A}\PYG{p}{,}\PYG{n}{B}\PYG{p}{)}\PYG{+w}{ }\PYG{p}{.}\PYG{p}{.}\PYG{p}{.}
\end{sphinxVerbatim}

\sphinxAtStartPar
with the following code:

\begin{sphinxVerbatim}[commandchars=\\\{\}]
\PYG{n}{TMP}\PYG{+w}{ }\PYG{o}{=}\PYG{+w}{ }\PYG{l+m+mf}{0.0}
\PYG{k}{do }\PYG{n}{I}\PYG{o}{=}\PYG{l+m+mi}{1}\PYG{p}{,}\PYG{n}{N}
\PYG{+w}{    }\PYG{n}{TMP}\PYG{+w}{ }\PYG{o}{=}\PYG{+w}{ }\PYG{n}{TMP}\PYG{+w}{ }\PYG{o}{+}\PYG{+w}{ }\PYG{n}{A}\PYG{p}{(}\PYG{n}{i}\PYG{p}{)}\PYG{o}{*}\PYG{n}{B}\PYG{p}{(}\PYG{n}{i}\PYG{p}{)}
\PYG{n}{R}\PYG{+w}{ }\PYG{o}{=}\PYG{+w}{ }\PYG{p}{.}\PYG{p}{.}\PYG{p}{.}\PYG{+w}{ }\PYG{n}{TMP}\PYG{+w}{ }\PYG{p}{.}\PYG{p}{.}\PYG{p}{.}
\end{sphinxVerbatim}

\sphinxAtStartPar
For example:

\begin{sphinxVerbatim}[commandchars=\\\{\}]
\PYG{g+gp}{\PYGZgt{}\PYGZgt{}\PYGZgt{} }\PYG{k+kn}{from}\PYG{+w}{ }\PYG{n+nn}{psyclone}\PYG{n+nn}{.}\PYG{n+nn}{psyir}\PYG{n+nn}{.}\PYG{n+nn}{backend}\PYG{n+nn}{.}\PYG{n+nn}{fortran}\PYG{+w}{ }\PYG{k+kn}{import} \PYG{n}{FortranWriter}
\PYG{g+gp}{\PYGZgt{}\PYGZgt{}\PYGZgt{} }\PYG{k+kn}{from}\PYG{+w}{ }\PYG{n+nn}{psyclone}\PYG{n+nn}{.}\PYG{n+nn}{psyir}\PYG{n+nn}{.}\PYG{n+nn}{frontend}\PYG{n+nn}{.}\PYG{n+nn}{fortran}\PYG{+w}{ }\PYG{k+kn}{import} \PYG{n}{FortranReader}
\PYG{g+gp}{\PYGZgt{}\PYGZgt{}\PYGZgt{} }\PYG{k+kn}{from}\PYG{+w}{ }\PYG{n+nn}{psyclone}\PYG{n+nn}{.}\PYG{n+nn}{psyir}\PYG{n+nn}{.}\PYG{n+nn}{nodes}\PYG{+w}{ }\PYG{k+kn}{import} \PYG{n}{IntrinsicCall}
\PYG{g+gp}{\PYGZgt{}\PYGZgt{}\PYGZgt{} }\PYG{k+kn}{from}\PYG{+w}{ }\PYG{n+nn}{psyclone}\PYG{n+nn}{.}\PYG{n+nn}{psyir}\PYG{n+nn}{.}\PYG{n+nn}{transformations}\PYG{+w}{ }\PYG{k+kn}{import} \PYG{n}{DotProduct2CodeTrans}
\PYG{g+gp}{\PYGZgt{}\PYGZgt{}\PYGZgt{} }\PYG{n}{code} \PYG{o}{=} \PYG{p}{(}\PYG{l+s+s2}{\PYGZdq{}}\PYG{l+s+s2}{subroutine dot\PYGZus{}product\PYGZus{}test(v1,v2)}\PYG{l+s+se}{\PYGZbs{}n}\PYG{l+s+s2}{\PYGZdq{}}
\PYG{g+gp}{... }        \PYG{l+s+s2}{\PYGZdq{}}\PYG{l+s+s2}{real,intent(in) :: v1(10), v2(10)}\PYG{l+s+se}{\PYGZbs{}n}\PYG{l+s+s2}{\PYGZdq{}}
\PYG{g+gp}{... }        \PYG{l+s+s2}{\PYGZdq{}}\PYG{l+s+s2}{real :: result}\PYG{l+s+se}{\PYGZbs{}n}\PYG{l+s+s2}{\PYGZdq{}}
\PYG{g+gp}{... }        \PYG{l+s+s2}{\PYGZdq{}}\PYG{l+s+s2}{result = dot\PYGZus{}product(v1,v2)}\PYG{l+s+se}{\PYGZbs{}n}\PYG{l+s+s2}{\PYGZdq{}}
\PYG{g+gp}{... }        \PYG{l+s+s2}{\PYGZdq{}}\PYG{l+s+s2}{end subroutine}\PYG{l+s+se}{\PYGZbs{}n}\PYG{l+s+s2}{\PYGZdq{}}\PYG{p}{)}
\PYG{g+gp}{\PYGZgt{}\PYGZgt{}\PYGZgt{} }\PYG{n}{psyir} \PYG{o}{=} \PYG{n}{FortranReader}\PYG{p}{(}\PYG{p}{)}\PYG{o}{.}\PYG{n}{psyir\PYGZus{}from\PYGZus{}source}\PYG{p}{(}\PYG{n}{code}\PYG{p}{)}
\PYG{g+gp}{\PYGZgt{}\PYGZgt{}\PYGZgt{} }\PYG{n}{trans} \PYG{o}{=} \PYG{n}{DotProduct2CodeTrans}\PYG{p}{(}\PYG{p}{)}
\PYG{g+gp}{\PYGZgt{}\PYGZgt{}\PYGZgt{} }\PYG{n}{trans}\PYG{o}{.}\PYG{n}{apply}\PYG{p}{(}\PYG{n}{psyir}\PYG{o}{.}\PYG{n}{walk}\PYG{p}{(}\PYG{n}{IntrinsicCall}\PYG{p}{)}\PYG{p}{[}\PYG{l+m+mi}{0}\PYG{p}{]}\PYG{p}{)}
\PYG{g+gp}{\PYGZgt{}\PYGZgt{}\PYGZgt{} }\PYG{n+nb}{print}\PYG{p}{(}\PYG{n}{FortranWriter}\PYG{p}{(}\PYG{p}{)}\PYG{p}{(}\PYG{n}{psyir}\PYG{p}{)}\PYG{p}{)}
\PYG{g+go}{subroutine dot\PYGZus{}product\PYGZus{}test(v1, v2)}
\PYG{g+go}{  real, dimension(10), intent(in) :: v1}
\PYG{g+go}{  real, dimension(10), intent(in) :: v2}
\PYG{g+go}{  real :: result}
\PYG{g+go}{  integer :: i}
\PYG{g+go}{  real :: res\PYGZus{}dot\PYGZus{}product}

\PYG{g+go}{  res\PYGZus{}dot\PYGZus{}product = 0.0}
\PYG{g+go}{  do i = 1, 10, 1}
\PYG{g+go}{    res\PYGZus{}dot\PYGZus{}product = res\PYGZus{}dot\PYGZus{}product + v1(i) * v2(i)}
\PYG{g+go}{  enddo}
\PYG{g+go}{  result = res\PYGZus{}dot\PYGZus{}product}

\PYG{g+go}{end subroutine dot\PYGZus{}product\PYGZus{}test}
\end{sphinxVerbatim}


\begin{fulllineitems}

\pysigstartsignatures
\pysiglinewithargsret{\sphinxbfcode{\sphinxupquote{apply}}}{\sphinxparam{\DUrole{n,n}{node}}\sphinxparamcomma \sphinxparam{\DUrole{n,n}{options}\DUrole{o,o}{=}\DUrole{default_value}{None}}}{}
\pysigstopsignatures
\sphinxAtStartPar
Apply the DOT\_PRODUCT intrinsic conversion transformation to the
specified node. This node must be a DOT\_PRODUCT
BinaryOperation.  If the transformation is successful then an
assignment which includes a DOT\_PRODUCT BinaryOperation node
is converted to equivalent inline code.
\begin{quote}\begin{description}
\sphinxlineitem{Parameters}\begin{itemize}
\item {} 
\sphinxAtStartPar
\sphinxstyleliteralstrong{\sphinxupquote{node}} (\sphinxcode{\sphinxupquote{psyclone.psyir.nodes.BinaryOperation}}) \textendash{} a DOT\_PRODUCT Binary\sphinxhyphen{}Operation node.

\item {} 
\sphinxAtStartPar
\sphinxstyleliteralstrong{\sphinxupquote{options}} (\sphinxhref{https://docs.python.org/3/library/stdtypes.html\#dict}{\sphinxstyleliteralemphasis{\sphinxupquote{dict}}}\sphinxstyleliteralemphasis{\sphinxupquote{ of }}\sphinxstyleliteralemphasis{\sphinxupquote{str:str}}\sphinxstyleliteralemphasis{\sphinxupquote{ or }}\sphinxstyleliteralemphasis{\sphinxupquote{None}}) \textendash{} a dictionary with options for transformations.

\end{itemize}

\end{description}\end{quote}

\end{fulllineitems}


\end{fulllineitems}



\bigskip\hrule\bigskip



\begin{fulllineitems}

\pysigstartsignatures
\pysiglinewithargsret{\sphinxbfcode{\sphinxupquote{class\DUrole{w,w}{  }}}\sphinxcode{\sphinxupquote{psyclone.psyir.transformations.extract\_trans.}}\sphinxbfcode{\sphinxupquote{ExtractTrans}}}{\sphinxparam{\DUrole{n,n}{node\_class=\textless{}class \textquotesingle{}psyclone.psyir.nodes.extract\_node.ExtractNode\textquotesingle{}\textgreater{}}}}{}
\pysigstopsignatures
\sphinxAtStartPar
This transformation inserts an ExtractNode or a node derived
from ExtractNode into the PSyIR of a schedule. At code creation
time this node will use the PSyData API to create code that can
write the input and output parameters to a file. The node might
also create a stand\sphinxhyphen{}alone driver program that can read the created
file and then execute the instrumented region.
Examples are given in the derived classes DynamoExtractTrans and
GOceanExtractTrans.

\sphinxAtStartPar
After applying the transformation the Nodes marked for extraction are
children of the ExtractNode.
Nodes to extract can be individual constructs within an Invoke (e.g.
Loops containing a Kernel or BuiltIn call) or entire Invokes. This
functionality does not support distributed memory.
\begin{quote}\begin{description}
\sphinxlineitem{Parameters}
\sphinxAtStartPar
\sphinxstyleliteralstrong{\sphinxupquote{node\_class}} (\sphinxcode{\sphinxupquote{psyclone.psyir.nodes.ExtractNode}} or         derived class) \textendash{} The Node class of which an instance will be inserted         into the tree (defaults to ExtractNode), but can be any derived class.

\end{description}\end{quote}


\begin{fulllineitems}

\pysigstartsignatures
\pysiglinewithargsret{\sphinxbfcode{\sphinxupquote{apply}}}{\sphinxparam{\DUrole{n,n}{nodes}}\sphinxparamcomma \sphinxparam{\DUrole{n,n}{options}\DUrole{o,o}{=}\DUrole{default_value}{None}}}{}
\pysigstopsignatures
\sphinxAtStartPar
Apply this transformation to a subset of the nodes within a
schedule \sphinxhyphen{} i.e. enclose the specified Nodes in the
schedule within a single PSyData region.
\begin{quote}\begin{description}
\sphinxlineitem{Parameters}\begin{itemize}
\item {} 
\sphinxAtStartPar
\sphinxstyleliteralstrong{\sphinxupquote{nodes}} (\sphinxcode{\sphinxupquote{psyclone.psyir.nodes.Node}} or list of                      \sphinxcode{\sphinxupquote{psyclone.psyir.nodes.Node}}) \textendash{} can be a single node or a list of nodes.

\item {} 
\sphinxAtStartPar
\sphinxstyleliteralstrong{\sphinxupquote{options}} (\sphinxstyleliteralemphasis{\sphinxupquote{Optional}}\sphinxstyleliteralemphasis{\sphinxupquote{{[}}}\sphinxstyleliteralemphasis{\sphinxupquote{Dict}}\sphinxstyleliteralemphasis{\sphinxupquote{{[}}}\sphinxhref{https://docs.python.org/3/library/stdtypes.html\#str}{\sphinxstyleliteralemphasis{\sphinxupquote{str}}}\sphinxstyleliteralemphasis{\sphinxupquote{, }}\sphinxstyleliteralemphasis{\sphinxupquote{Any}}\sphinxstyleliteralemphasis{\sphinxupquote{{]}}}\sphinxstyleliteralemphasis{\sphinxupquote{{]}}}) \textendash{} a dictionary with options for transformations.

\item {} 
\sphinxAtStartPar
\sphinxstyleliteralstrong{\sphinxupquote{options}}\sphinxstyleliteralstrong{\sphinxupquote{{[}}}\sphinxstyleliteralstrong{\sphinxupquote{"prefix"}}\sphinxstyleliteralstrong{\sphinxupquote{{]}}} (\sphinxhref{https://docs.python.org/3/library/stdtypes.html\#str}{\sphinxstyleliteralemphasis{\sphinxupquote{str}}}) \textendash{} a prefix to use for the PSyData module             name (\sphinxcode{\sphinxupquote{PREFIX\_psy\_data\_mod}}) and the PSyDataType             (\sphinxcode{\sphinxupquote{PREFIX\_PSYDATATYPE}}) \sphinxhyphen{} a “\_” will be added automatically.             It defaults to “”.

\item {} 
\sphinxAtStartPar
\sphinxstyleliteralstrong{\sphinxupquote{options}}\sphinxstyleliteralstrong{\sphinxupquote{{[}}}\sphinxstyleliteralstrong{\sphinxupquote{"region\_name"}}\sphinxstyleliteralstrong{\sphinxupquote{{]}}} (\sphinxstyleliteralemphasis{\sphinxupquote{(}}\sphinxhref{https://docs.python.org/3/library/stdtypes.html\#str}{\sphinxstyleliteralemphasis{\sphinxupquote{str}}}\sphinxstyleliteralemphasis{\sphinxupquote{,}}\sphinxhref{https://docs.python.org/3/library/stdtypes.html\#str}{\sphinxstyleliteralemphasis{\sphinxupquote{str}}}\sphinxstyleliteralemphasis{\sphinxupquote{)}}) \textendash{} an optional name to             use for this PSyData area, provided as a 2\sphinxhyphen{}tuple containing a             location name followed by a local name. The pair of strings             should uniquely identify a region unless aggregate information             is required (and is supported by the runtime library).

\end{itemize}

\end{description}\end{quote}

\end{fulllineitems}


\end{fulllineitems}



\bigskip\hrule\bigskip



\begin{fulllineitems}

\pysigstartsignatures
\pysigline{\sphinxbfcode{\sphinxupquote{class\DUrole{w,w}{  }}}\sphinxcode{\sphinxupquote{psyclone.psyir.transformations.}}\sphinxbfcode{\sphinxupquote{HoistLocalArraysTrans}}}
\pysigstopsignatures
\sphinxAtStartPar
This transformation takes a Routine and promotes any local, ‘automatic’
arrays to Container scope:

\begin{sphinxVerbatim}[commandchars=\\\{\}]
\PYG{g+gp}{\PYGZgt{}\PYGZgt{}\PYGZgt{} }\PYG{k+kn}{from}\PYG{+w}{ }\PYG{n+nn}{psyclone}\PYG{n+nn}{.}\PYG{n+nn}{psyir}\PYG{n+nn}{.}\PYG{n+nn}{backend}\PYG{n+nn}{.}\PYG{n+nn}{fortran}\PYG{+w}{ }\PYG{k+kn}{import} \PYG{n}{FortranWriter}
\PYG{g+gp}{\PYGZgt{}\PYGZgt{}\PYGZgt{} }\PYG{k+kn}{from}\PYG{+w}{ }\PYG{n+nn}{psyclone}\PYG{n+nn}{.}\PYG{n+nn}{psyir}\PYG{n+nn}{.}\PYG{n+nn}{frontend}\PYG{n+nn}{.}\PYG{n+nn}{fortran}\PYG{+w}{ }\PYG{k+kn}{import} \PYG{n}{FortranReader}
\PYG{g+gp}{\PYGZgt{}\PYGZgt{}\PYGZgt{} }\PYG{k+kn}{from}\PYG{+w}{ }\PYG{n+nn}{psyclone}\PYG{n+nn}{.}\PYG{n+nn}{psyir}\PYG{n+nn}{.}\PYG{n+nn}{nodes}\PYG{+w}{ }\PYG{k+kn}{import} \PYG{n}{Assignment}
\PYG{g+gp}{\PYGZgt{}\PYGZgt{}\PYGZgt{} }\PYG{k+kn}{from}\PYG{+w}{ }\PYG{n+nn}{psyclone}\PYG{n+nn}{.}\PYG{n+nn}{psyir}\PYG{n+nn}{.}\PYG{n+nn}{transformations}\PYG{+w}{ }\PYG{k+kn}{import} \PYG{n}{HoistLocalArraysTrans}
\PYG{g+gp}{\PYGZgt{}\PYGZgt{}\PYGZgt{} }\PYG{n}{code} \PYG{o}{=} \PYG{p}{(}\PYG{l+s+s2}{\PYGZdq{}}\PYG{l+s+s2}{module test\PYGZus{}mod}\PYG{l+s+se}{\PYGZbs{}n}\PYG{l+s+s2}{\PYGZdq{}}
\PYG{g+gp}{... }        \PYG{l+s+s2}{\PYGZdq{}}\PYG{l+s+s2}{contains}\PYG{l+s+se}{\PYGZbs{}n}\PYG{l+s+s2}{\PYGZdq{}}
\PYG{g+gp}{... }        \PYG{l+s+s2}{\PYGZdq{}}\PYG{l+s+s2}{  subroutine test\PYGZus{}sub(n)}\PYG{l+s+se}{\PYGZbs{}n}\PYG{l+s+s2}{\PYGZdq{}}
\PYG{g+gp}{... }        \PYG{l+s+s2}{\PYGZdq{}}\PYG{l+s+s2}{  integer :: i,j,n}\PYG{l+s+se}{\PYGZbs{}n}\PYG{l+s+s2}{\PYGZdq{}}
\PYG{g+gp}{... }        \PYG{l+s+s2}{\PYGZdq{}}\PYG{l+s+s2}{  real :: a(n,n)}\PYG{l+s+se}{\PYGZbs{}n}\PYG{l+s+s2}{\PYGZdq{}}
\PYG{g+gp}{... }        \PYG{l+s+s2}{\PYGZdq{}}\PYG{l+s+s2}{  real :: value = 1.0}\PYG{l+s+se}{\PYGZbs{}n}\PYG{l+s+s2}{\PYGZdq{}}
\PYG{g+gp}{... }        \PYG{l+s+s2}{\PYGZdq{}}\PYG{l+s+s2}{  do i=1,n}\PYG{l+s+se}{\PYGZbs{}n}\PYG{l+s+s2}{\PYGZdq{}}
\PYG{g+gp}{... }        \PYG{l+s+s2}{\PYGZdq{}}\PYG{l+s+s2}{    do j=1,n}\PYG{l+s+se}{\PYGZbs{}n}\PYG{l+s+s2}{\PYGZdq{}}
\PYG{g+gp}{... }        \PYG{l+s+s2}{\PYGZdq{}}\PYG{l+s+s2}{      a(i,j) = value}\PYG{l+s+se}{\PYGZbs{}n}\PYG{l+s+s2}{\PYGZdq{}}
\PYG{g+gp}{... }        \PYG{l+s+s2}{\PYGZdq{}}\PYG{l+s+s2}{    end do}\PYG{l+s+se}{\PYGZbs{}n}\PYG{l+s+s2}{\PYGZdq{}}
\PYG{g+gp}{... }        \PYG{l+s+s2}{\PYGZdq{}}\PYG{l+s+s2}{  end do}\PYG{l+s+se}{\PYGZbs{}n}\PYG{l+s+s2}{\PYGZdq{}}
\PYG{g+gp}{... }        \PYG{l+s+s2}{\PYGZdq{}}\PYG{l+s+s2}{  end subroutine test\PYGZus{}sub}\PYG{l+s+se}{\PYGZbs{}n}\PYG{l+s+s2}{\PYGZdq{}}
\PYG{g+gp}{... }        \PYG{l+s+s2}{\PYGZdq{}}\PYG{l+s+s2}{end module test\PYGZus{}mod}\PYG{l+s+se}{\PYGZbs{}n}\PYG{l+s+s2}{\PYGZdq{}}\PYG{p}{)}
\PYG{g+gp}{\PYGZgt{}\PYGZgt{}\PYGZgt{} }\PYG{n}{psyir} \PYG{o}{=} \PYG{n}{FortranReader}\PYG{p}{(}\PYG{p}{)}\PYG{o}{.}\PYG{n}{psyir\PYGZus{}from\PYGZus{}source}\PYG{p}{(}\PYG{n}{code}\PYG{p}{)}
\PYG{g+gp}{\PYGZgt{}\PYGZgt{}\PYGZgt{} }\PYG{n}{hoist} \PYG{o}{=} \PYG{n}{HoistLocalArraysTrans}\PYG{p}{(}\PYG{p}{)}
\PYG{g+gp}{\PYGZgt{}\PYGZgt{}\PYGZgt{} }\PYG{n}{hoist}\PYG{o}{.}\PYG{n}{apply}\PYG{p}{(}\PYG{n}{psyir}\PYG{o}{.}\PYG{n}{walk}\PYG{p}{(}\PYG{n}{Routine}\PYG{p}{)}\PYG{p}{[}\PYG{l+m+mi}{0}\PYG{p}{]}\PYG{p}{)}
\PYG{g+gp}{\PYGZgt{}\PYGZgt{}\PYGZgt{} }\PYG{n+nb}{print}\PYG{p}{(}\PYG{n}{FortranWriter}\PYG{p}{(}\PYG{p}{)}\PYG{p}{(}\PYG{n}{psyir}\PYG{p}{)}\PYG{o}{.}\PYG{n}{lower}\PYG{p}{(}\PYG{p}{)}\PYG{p}{)}
\PYG{g+go}{module test\PYGZus{}mod}
\PYG{g+go}{  implicit none}
\PYG{g+go}{  real, allocatable, dimension(:,:), private :: a}
\PYG{g+go}{  public}

\PYG{g+go}{  public :: test\PYGZus{}sub}

\PYG{g+go}{  contains}
\PYG{g+go}{  subroutine test\PYGZus{}sub(n)}
\PYG{g+go}{    integer :: n}
\PYG{g+go}{    integer :: i}
\PYG{g+go}{    integer :: j}
\PYG{g+go}{    real :: value = 1.0}

\PYG{g+go}{    if (.not.allocated(a) .or. ubound(a, 1) /= n .or. ubound(a, 2) /= n) then}
\PYG{g+go}{      if (allocated(a)) then}
\PYG{g+go}{        deallocate(a)}
\PYG{g+go}{      end if}
\PYG{g+go}{      allocate(a(1 : n, 1 : n))}
\PYG{g+go}{    end if}
\PYG{g+go}{    do i = 1, n, 1}
\PYG{g+go}{      do j = 1, n, 1}
\PYG{g+go}{        a(i,j) = value}
\PYG{g+go}{      enddo}
\PYG{g+go}{    enddo}

\PYG{g+go}{  end subroutine test\PYGZus{}sub}

\PYG{g+go}{end module test\PYGZus{}mod}
\end{sphinxVerbatim}

\sphinxAtStartPar
By default, the target routine will be rejected if it is found to contain
an ACCRoutineDirective since this usually implies that the routine will be
launched in parallel on the OpenACC device. This check can be disabled
by setting ‘allow\_accroutine’ to True in the \sphinxtitleref{options} dictionary.


\begin{fulllineitems}

\pysigstartsignatures
\pysiglinewithargsret{\sphinxbfcode{\sphinxupquote{apply}}}{\sphinxparam{\DUrole{n,n}{node}}\sphinxparamcomma \sphinxparam{\DUrole{n,n}{options}\DUrole{o,o}{=}\DUrole{default_value}{None}}}{}
\pysigstopsignatures
\sphinxAtStartPar
Applies the transformation to the supplied Routine node,
moving any local arrays up to Container scope and adding
a suitable allocation when they are first accessed. If there
are no local arrays or the supplied Routine is a program then
this method does nothing.
\begin{quote}\begin{description}
\sphinxlineitem{Parameters}\begin{itemize}
\item {} 
\sphinxAtStartPar
\sphinxstyleliteralstrong{\sphinxupquote{node}} (\sphinxcode{\sphinxupquote{psyclone.psyir.nodes.Routine}}) \textendash{} target PSyIR node.

\item {} 
\sphinxAtStartPar
\sphinxstyleliteralstrong{\sphinxupquote{options}} (\sphinxstyleliteralemphasis{\sphinxupquote{Optional}}\sphinxstyleliteralemphasis{\sphinxupquote{{[}}}\sphinxstyleliteralemphasis{\sphinxupquote{Dict}}\sphinxstyleliteralemphasis{\sphinxupquote{{[}}}\sphinxhref{https://docs.python.org/3/library/stdtypes.html\#str}{\sphinxstyleliteralemphasis{\sphinxupquote{str}}}\sphinxstyleliteralemphasis{\sphinxupquote{, }}\sphinxstyleliteralemphasis{\sphinxupquote{Any}}\sphinxstyleliteralemphasis{\sphinxupquote{{]}}}\sphinxstyleliteralemphasis{\sphinxupquote{{]}}}) \textendash{} a dictionary with options for transformations.

\item {} 
\sphinxAtStartPar
\sphinxstyleliteralstrong{\sphinxupquote{options}}\sphinxstyleliteralstrong{\sphinxupquote{{[}}}\sphinxstyleliteralstrong{\sphinxupquote{"allow\_accroutine"}}\sphinxstyleliteralstrong{\sphinxupquote{{]}}} (\sphinxhref{https://docs.python.org/3/library/functions.html\#bool}{\sphinxstyleliteralemphasis{\sphinxupquote{bool}}}) \textendash{} permit the target routine             to contain an ACCRoutineDirective. These are forbidden by default             because their presence usually indicates that the routine will be             run in parallel on the OpenACC device.

\end{itemize}

\end{description}\end{quote}

\end{fulllineitems}


\end{fulllineitems}



\bigskip\hrule\bigskip



\begin{fulllineitems}

\pysigstartsignatures
\pysigline{\sphinxbfcode{\sphinxupquote{class\DUrole{w,w}{  }}}\sphinxcode{\sphinxupquote{psyclone.psyir.transformations.}}\sphinxbfcode{\sphinxupquote{HoistLoopBoundExprTrans}}}
\pysigstopsignatures
\sphinxAtStartPar
This transformation moves complex bounds expressions out of the loop
construct and places them in integer scalar assignments before the loop.

\begin{sphinxVerbatim}[commandchars=\\\{\}]
\PYG{g+gp}{\PYGZgt{}\PYGZgt{}\PYGZgt{} }\PYG{k+kn}{from}\PYG{+w}{ }\PYG{n+nn}{psyclone}\PYG{n+nn}{.}\PYG{n+nn}{psyir}\PYG{n+nn}{.}\PYG{n+nn}{backend}\PYG{n+nn}{.}\PYG{n+nn}{fortran}\PYG{+w}{ }\PYG{k+kn}{import} \PYG{n}{FortranWriter}
\PYG{g+gp}{\PYGZgt{}\PYGZgt{}\PYGZgt{} }\PYG{k+kn}{from}\PYG{+w}{ }\PYG{n+nn}{psyclone}\PYG{n+nn}{.}\PYG{n+nn}{psyir}\PYG{n+nn}{.}\PYG{n+nn}{frontend}\PYG{n+nn}{.}\PYG{n+nn}{fortran}\PYG{+w}{ }\PYG{k+kn}{import} \PYG{n}{FortranReader}
\PYG{g+gp}{\PYGZgt{}\PYGZgt{}\PYGZgt{} }\PYG{k+kn}{from}\PYG{+w}{ }\PYG{n+nn}{psyclone}\PYG{n+nn}{.}\PYG{n+nn}{psyir}\PYG{n+nn}{.}\PYG{n+nn}{nodes}\PYG{+w}{ }\PYG{k+kn}{import} \PYG{n}{Loop}
\PYG{g+gp}{\PYGZgt{}\PYGZgt{}\PYGZgt{} }\PYG{k+kn}{from}\PYG{+w}{ }\PYG{n+nn}{psyclone}\PYG{n+nn}{.}\PYG{n+nn}{psyir}\PYG{n+nn}{.}\PYG{n+nn}{transformations}\PYG{+w}{ }\PYG{k+kn}{import} \PYG{n}{HoistTrans}
\PYG{g+gp}{\PYGZgt{}\PYGZgt{}\PYGZgt{} }\PYG{n}{code} \PYG{o}{=} \PYG{p}{(}\PYG{l+s+s2}{\PYGZdq{}}\PYG{l+s+s2}{program test}\PYG{l+s+se}{\PYGZbs{}n}\PYG{l+s+s2}{\PYGZdq{}}
\PYG{g+gp}{... }        \PYG{l+s+s2}{\PYGZdq{}}\PYG{l+s+s2}{  use mymod, only: mytype}\PYG{l+s+se}{\PYGZbs{}n}\PYG{l+s+s2}{\PYGZdq{}}
\PYG{g+gp}{... }        \PYG{l+s+s2}{\PYGZdq{}}\PYG{l+s+s2}{  integer :: i,j,n}\PYG{l+s+se}{\PYGZbs{}n}\PYG{l+s+s2}{\PYGZdq{}}
\PYG{g+gp}{... }        \PYG{l+s+s2}{\PYGZdq{}}\PYG{l+s+s2}{  real :: a(n)}\PYG{l+s+se}{\PYGZbs{}n}\PYG{l+s+s2}{\PYGZdq{}}
\PYG{g+gp}{... }        \PYG{l+s+s2}{\PYGZdq{}}\PYG{l+s+s2}{  do i=mytype}\PYG{l+s+si}{\PYGZpc{}s}\PYG{l+s+s2}{tart, UBOUND(a,1)}\PYG{l+s+se}{\PYGZbs{}n}\PYG{l+s+s2}{\PYGZdq{}}
\PYG{g+gp}{... }        \PYG{l+s+s2}{\PYGZdq{}}\PYG{l+s+s2}{    a(i) = 1.0}\PYG{l+s+se}{\PYGZbs{}n}\PYG{l+s+s2}{\PYGZdq{}}
\PYG{g+gp}{... }        \PYG{l+s+s2}{\PYGZdq{}}\PYG{l+s+s2}{  end do}\PYG{l+s+se}{\PYGZbs{}n}\PYG{l+s+s2}{\PYGZdq{}}
\PYG{g+gp}{... }        \PYG{l+s+s2}{\PYGZdq{}}\PYG{l+s+s2}{end program}\PYG{l+s+se}{\PYGZbs{}n}\PYG{l+s+s2}{\PYGZdq{}}\PYG{p}{)}
\PYG{g+gp}{\PYGZgt{}\PYGZgt{}\PYGZgt{} }\PYG{n}{psyir} \PYG{o}{=} \PYG{n}{FortranReader}\PYG{p}{(}\PYG{p}{)}\PYG{o}{.}\PYG{n}{psyir\PYGZus{}from\PYGZus{}source}\PYG{p}{(}\PYG{n}{code}\PYG{p}{)}
\PYG{g+gp}{\PYGZgt{}\PYGZgt{}\PYGZgt{} }\PYG{n}{hoist} \PYG{o}{=} \PYG{n}{HoistLoopBoundExprTrans}\PYG{p}{(}\PYG{p}{)}
\PYG{g+gp}{\PYGZgt{}\PYGZgt{}\PYGZgt{} }\PYG{n}{hoist}\PYG{o}{.}\PYG{n}{apply}\PYG{p}{(}\PYG{n}{psyir}\PYG{o}{.}\PYG{n}{walk}\PYG{p}{(}\PYG{n}{Loop}\PYG{p}{)}\PYG{p}{[}\PYG{l+m+mi}{0}\PYG{p}{]}\PYG{p}{)}
\PYG{g+gp}{\PYGZgt{}\PYGZgt{}\PYGZgt{} }\PYG{n+nb}{print}\PYG{p}{(}\PYG{n}{FortranWriter}\PYG{p}{(}\PYG{p}{)}\PYG{p}{(}\PYG{n}{psyir}\PYG{p}{)}\PYG{p}{)}
\PYG{g+go}{program test}
\PYG{g+go}{  use mymod, only : mytype}
\PYG{g+go}{  integer :: i}
\PYG{g+go}{  integer :: j}
\PYG{g+go}{  integer :: n}
\PYG{g+go}{  real, dimension(n) :: a}
\PYG{g+go}{  integer :: loop\PYGZus{}bound}
\PYG{g+go}{  integer :: loop\PYGZus{}bound\PYGZus{}1}

\PYG{g+go}{  loop\PYGZus{}bound\PYGZus{}1 = UBOUND(a, 1)}
\PYG{g+go}{  loop\PYGZus{}bound = mytype\PYGZpc{}start}
\PYG{g+go}{  do i = loop\PYGZus{}bound, loop\PYGZus{}bound\PYGZus{}1, 1}
\PYG{g+go}{    a(i) = 1.0}
\PYG{g+go}{  enddo}

\PYG{g+go}{end program test}
\end{sphinxVerbatim}


\begin{fulllineitems}

\pysigstartsignatures
\pysiglinewithargsret{\sphinxbfcode{\sphinxupquote{apply}}}{\sphinxparam{\DUrole{n,n}{node}}\sphinxparamcomma \sphinxparam{\DUrole{n,n}{options}\DUrole{o,o}{=}\DUrole{default_value}{None}}}{}
\pysigstopsignatures
\sphinxAtStartPar
Move complex bounds expressions out of the given loop construct and
place them in integer scalar assignments before the loop.
\begin{quote}\begin{description}
\sphinxlineitem{Parameters}\begin{itemize}
\item {} 
\sphinxAtStartPar
\sphinxstyleliteralstrong{\sphinxupquote{node}} (\sphinxcode{\sphinxupquote{psyclone.psyir.nodes.Loop}}) \textendash{} target PSyIR loop.

\item {} 
\sphinxAtStartPar
\sphinxstyleliteralstrong{\sphinxupquote{options}} (\sphinxstyleliteralemphasis{\sphinxupquote{Dict}}\sphinxstyleliteralemphasis{\sphinxupquote{{[}}}\sphinxhref{https://docs.python.org/3/library/stdtypes.html\#str}{\sphinxstyleliteralemphasis{\sphinxupquote{str}}}\sphinxstyleliteralemphasis{\sphinxupquote{, }}\sphinxstyleliteralemphasis{\sphinxupquote{Any}}\sphinxstyleliteralemphasis{\sphinxupquote{{]}}}) \textendash{} a dictionary with options for transformations.

\end{itemize}

\end{description}\end{quote}

\end{fulllineitems}


\end{fulllineitems}



\bigskip\hrule\bigskip



\begin{fulllineitems}

\pysigstartsignatures
\pysigline{\sphinxbfcode{\sphinxupquote{class\DUrole{w,w}{  }}}\sphinxcode{\sphinxupquote{psyclone.psyir.transformations.}}\sphinxbfcode{\sphinxupquote{HoistTrans}}}
\pysigstopsignatures
\sphinxAtStartPar
This transformation takes an assignment and moves it outside of
its parent loop if it is valid to do so. If as a result the loop body
becomes empty, the loop will be removed altogether. For example:

\begin{sphinxVerbatim}[commandchars=\\\{\}]
\PYG{g+gp}{\PYGZgt{}\PYGZgt{}\PYGZgt{} }\PYG{k+kn}{from}\PYG{+w}{ }\PYG{n+nn}{psyclone}\PYG{n+nn}{.}\PYG{n+nn}{psyir}\PYG{n+nn}{.}\PYG{n+nn}{backend}\PYG{n+nn}{.}\PYG{n+nn}{fortran}\PYG{+w}{ }\PYG{k+kn}{import} \PYG{n}{FortranWriter}
\PYG{g+gp}{\PYGZgt{}\PYGZgt{}\PYGZgt{} }\PYG{k+kn}{from}\PYG{+w}{ }\PYG{n+nn}{psyclone}\PYG{n+nn}{.}\PYG{n+nn}{psyir}\PYG{n+nn}{.}\PYG{n+nn}{frontend}\PYG{n+nn}{.}\PYG{n+nn}{fortran}\PYG{+w}{ }\PYG{k+kn}{import} \PYG{n}{FortranReader}
\PYG{g+gp}{\PYGZgt{}\PYGZgt{}\PYGZgt{} }\PYG{k+kn}{from}\PYG{+w}{ }\PYG{n+nn}{psyclone}\PYG{n+nn}{.}\PYG{n+nn}{psyir}\PYG{n+nn}{.}\PYG{n+nn}{nodes}\PYG{+w}{ }\PYG{k+kn}{import} \PYG{n}{Assignment}
\PYG{g+gp}{\PYGZgt{}\PYGZgt{}\PYGZgt{} }\PYG{k+kn}{from}\PYG{+w}{ }\PYG{n+nn}{psyclone}\PYG{n+nn}{.}\PYG{n+nn}{psyir}\PYG{n+nn}{.}\PYG{n+nn}{transformations}\PYG{+w}{ }\PYG{k+kn}{import} \PYG{n}{HoistTrans}
\PYG{g+gp}{\PYGZgt{}\PYGZgt{}\PYGZgt{} }\PYG{n}{code} \PYG{o}{=} \PYG{p}{(}\PYG{l+s+s2}{\PYGZdq{}}\PYG{l+s+s2}{program test}\PYG{l+s+se}{\PYGZbs{}n}\PYG{l+s+s2}{\PYGZdq{}}
\PYG{g+gp}{... }        \PYG{l+s+s2}{\PYGZdq{}}\PYG{l+s+s2}{  integer :: i,j,n}\PYG{l+s+se}{\PYGZbs{}n}\PYG{l+s+s2}{\PYGZdq{}}
\PYG{g+gp}{... }        \PYG{l+s+s2}{\PYGZdq{}}\PYG{l+s+s2}{  real :: a(n,n)}\PYG{l+s+se}{\PYGZbs{}n}\PYG{l+s+s2}{\PYGZdq{}}
\PYG{g+gp}{... }        \PYG{l+s+s2}{\PYGZdq{}}\PYG{l+s+s2}{  real value}\PYG{l+s+se}{\PYGZbs{}n}\PYG{l+s+s2}{\PYGZdq{}}
\PYG{g+gp}{... }        \PYG{l+s+s2}{\PYGZdq{}}\PYG{l+s+s2}{  do i=1,n}\PYG{l+s+se}{\PYGZbs{}n}\PYG{l+s+s2}{\PYGZdq{}}
\PYG{g+gp}{... }        \PYG{l+s+s2}{\PYGZdq{}}\PYG{l+s+s2}{    value = 1.0}\PYG{l+s+se}{\PYGZbs{}n}\PYG{l+s+s2}{\PYGZdq{}}
\PYG{g+gp}{... }        \PYG{l+s+s2}{\PYGZdq{}}\PYG{l+s+s2}{    do j=1,n}\PYG{l+s+se}{\PYGZbs{}n}\PYG{l+s+s2}{\PYGZdq{}}
\PYG{g+gp}{... }        \PYG{l+s+s2}{\PYGZdq{}}\PYG{l+s+s2}{      a(i,j) = value}\PYG{l+s+se}{\PYGZbs{}n}\PYG{l+s+s2}{\PYGZdq{}}
\PYG{g+gp}{... }        \PYG{l+s+s2}{\PYGZdq{}}\PYG{l+s+s2}{    end do}\PYG{l+s+se}{\PYGZbs{}n}\PYG{l+s+s2}{\PYGZdq{}}
\PYG{g+gp}{... }        \PYG{l+s+s2}{\PYGZdq{}}\PYG{l+s+s2}{  end do}\PYG{l+s+se}{\PYGZbs{}n}\PYG{l+s+s2}{\PYGZdq{}}
\PYG{g+gp}{... }        \PYG{l+s+s2}{\PYGZdq{}}\PYG{l+s+s2}{end program}\PYG{l+s+se}{\PYGZbs{}n}\PYG{l+s+s2}{\PYGZdq{}}\PYG{p}{)}
\PYG{g+gp}{\PYGZgt{}\PYGZgt{}\PYGZgt{} }\PYG{n}{psyir} \PYG{o}{=} \PYG{n}{FortranReader}\PYG{p}{(}\PYG{p}{)}\PYG{o}{.}\PYG{n}{psyir\PYGZus{}from\PYGZus{}source}\PYG{p}{(}\PYG{n}{code}\PYG{p}{)}
\PYG{g+gp}{\PYGZgt{}\PYGZgt{}\PYGZgt{} }\PYG{n}{hoist} \PYG{o}{=} \PYG{n}{HoistTrans}\PYG{p}{(}\PYG{p}{)}
\PYG{g+gp}{\PYGZgt{}\PYGZgt{}\PYGZgt{} }\PYG{n}{hoist}\PYG{o}{.}\PYG{n}{apply}\PYG{p}{(}\PYG{n}{psyir}\PYG{o}{.}\PYG{n}{walk}\PYG{p}{(}\PYG{n}{Assignment}\PYG{p}{)}\PYG{p}{[}\PYG{l+m+mi}{0}\PYG{p}{]}\PYG{p}{)}
\PYG{g+gp}{\PYGZgt{}\PYGZgt{}\PYGZgt{} }\PYG{n+nb}{print}\PYG{p}{(}\PYG{n}{FortranWriter}\PYG{p}{(}\PYG{p}{)}\PYG{p}{(}\PYG{n}{psyir}\PYG{p}{)}\PYG{p}{)}
\PYG{g+go}{program test}
\PYG{g+go}{  integer :: i}
\PYG{g+go}{  integer :: j}
\PYG{g+go}{  integer :: n}
\PYG{g+go}{  real, dimension(n,n) :: a}
\PYG{g+go}{  real :: value}

\PYG{g+go}{  value = 1.0}
\PYG{g+go}{  do i = 1, n, 1}
\PYG{g+go}{    do j = 1, n, 1}
\PYG{g+go}{      a(i,j) = value}
\PYG{g+go}{    enddo}
\PYG{g+go}{  enddo}

\PYG{g+go}{end program test}
\end{sphinxVerbatim}


\begin{fulllineitems}

\pysigstartsignatures
\pysiglinewithargsret{\sphinxbfcode{\sphinxupquote{apply}}}{\sphinxparam{\DUrole{n,n}{node}}\sphinxparamcomma \sphinxparam{\DUrole{n,n}{options}\DUrole{o,o}{=}\DUrole{default_value}{None}}}{}
\pysigstopsignatures
\sphinxAtStartPar
Applies the hoist transformation to the supplied assignment node
within a loop, moving the assignment outside of the loop if it
is valid to do so. Issue \#1445 will also look to extend this
transformation to other types of node.
\begin{quote}\begin{description}
\sphinxlineitem{Parameters}\begin{itemize}
\item {} 
\sphinxAtStartPar
\sphinxstyleliteralstrong{\sphinxupquote{node}} (subclass of \sphinxcode{\sphinxupquote{psyclone.psyir.nodes.Assignment}}) \textendash{} target PSyIR node.

\item {} 
\sphinxAtStartPar
\sphinxstyleliteralstrong{\sphinxupquote{options}} (\sphinxstyleliteralemphasis{\sphinxupquote{Optional}}\sphinxstyleliteralemphasis{\sphinxupquote{{[}}}\sphinxstyleliteralemphasis{\sphinxupquote{Dict}}\sphinxstyleliteralemphasis{\sphinxupquote{{[}}}\sphinxhref{https://docs.python.org/3/library/stdtypes.html\#str}{\sphinxstyleliteralemphasis{\sphinxupquote{str}}}\sphinxstyleliteralemphasis{\sphinxupquote{, }}\sphinxstyleliteralemphasis{\sphinxupquote{Any}}\sphinxstyleliteralemphasis{\sphinxupquote{{]}}}\sphinxstyleliteralemphasis{\sphinxupquote{{]}}}) \textendash{} a dictionary with options for transformations.

\end{itemize}

\end{description}\end{quote}

\end{fulllineitems}


\end{fulllineitems}



\bigskip\hrule\bigskip



\begin{fulllineitems}

\pysigstartsignatures
\pysigline{\sphinxbfcode{\sphinxupquote{class\DUrole{w,w}{  }}}\sphinxcode{\sphinxupquote{psyclone.psyir.transformations.}}\sphinxbfcode{\sphinxupquote{InlineTrans}}}
\pysigstopsignatures
\sphinxAtStartPar
This transformation takes a Call (which may have a return value)
and replaces it with the body of the target routine. It is used as
follows:

\begin{sphinxVerbatim}[commandchars=\\\{\}]
\PYG{g+gp}{\PYGZgt{}\PYGZgt{}\PYGZgt{} }\PYG{k+kn}{from}\PYG{+w}{ }\PYG{n+nn}{psyclone}\PYG{n+nn}{.}\PYG{n+nn}{psyir}\PYG{n+nn}{.}\PYG{n+nn}{backend}\PYG{n+nn}{.}\PYG{n+nn}{fortran}\PYG{+w}{ }\PYG{k+kn}{import} \PYG{n}{FortranWriter}
\PYG{g+gp}{\PYGZgt{}\PYGZgt{}\PYGZgt{} }\PYG{k+kn}{from}\PYG{+w}{ }\PYG{n+nn}{psyclone}\PYG{n+nn}{.}\PYG{n+nn}{psyir}\PYG{n+nn}{.}\PYG{n+nn}{frontend}\PYG{n+nn}{.}\PYG{n+nn}{fortran}\PYG{+w}{ }\PYG{k+kn}{import} \PYG{n}{FortranReader}
\PYG{g+gp}{\PYGZgt{}\PYGZgt{}\PYGZgt{} }\PYG{k+kn}{from}\PYG{+w}{ }\PYG{n+nn}{psyclone}\PYG{n+nn}{.}\PYG{n+nn}{psyir}\PYG{n+nn}{.}\PYG{n+nn}{nodes}\PYG{+w}{ }\PYG{k+kn}{import} \PYG{n}{Call}\PYG{p}{,} \PYG{n}{Routine}
\PYG{g+gp}{\PYGZgt{}\PYGZgt{}\PYGZgt{} }\PYG{k+kn}{from}\PYG{+w}{ }\PYG{n+nn}{psyclone}\PYG{n+nn}{.}\PYG{n+nn}{psyir}\PYG{n+nn}{.}\PYG{n+nn}{transformations}\PYG{+w}{ }\PYG{k+kn}{import} \PYG{n}{InlineTrans}
\PYG{g+gp}{\PYGZgt{}\PYGZgt{}\PYGZgt{} }\PYG{n}{code} \PYG{o}{=} \PYG{l+s+s2}{\PYGZdq{}\PYGZdq{}\PYGZdq{}}
\PYG{g+gp}{... }\PYG{l+s+s2}{module test\PYGZus{}mod}
\PYG{g+gp}{... }\PYG{l+s+s2}{contains}
\PYG{g+gp}{... }\PYG{l+s+s2}{  subroutine run\PYGZus{}it()}
\PYG{g+gp}{... }\PYG{l+s+s2}{    integer :: i}
\PYG{g+gp}{... }\PYG{l+s+s2}{    real :: a(10)}
\PYG{g+gp}{... }\PYG{l+s+s2}{    do i=1,10}
\PYG{g+gp}{... }\PYG{l+s+s2}{      a(i) = 1.0}
\PYG{g+gp}{... }\PYG{l+s+s2}{      call sub(a(i))}
\PYG{g+gp}{... }\PYG{l+s+s2}{    end do}
\PYG{g+gp}{... }\PYG{l+s+s2}{  end subroutine run\PYGZus{}it}
\PYG{g+gp}{... }\PYG{l+s+s2}{  subroutine sub(x)}
\PYG{g+gp}{... }\PYG{l+s+s2}{    real, intent(inout) :: x}
\PYG{g+gp}{... }\PYG{l+s+s2}{    x = 2.0*x}
\PYG{g+gp}{... }\PYG{l+s+s2}{  end subroutine sub}
\PYG{g+gp}{... }\PYG{l+s+s2}{end module test\PYGZus{}mod}\PYG{l+s+s2}{\PYGZdq{}\PYGZdq{}\PYGZdq{}}
\PYG{g+gp}{\PYGZgt{}\PYGZgt{}\PYGZgt{} }\PYG{n}{psyir} \PYG{o}{=} \PYG{n}{FortranReader}\PYG{p}{(}\PYG{p}{)}\PYG{o}{.}\PYG{n}{psyir\PYGZus{}from\PYGZus{}source}\PYG{p}{(}\PYG{n}{code}\PYG{p}{)}
\PYG{g+gp}{\PYGZgt{}\PYGZgt{}\PYGZgt{} }\PYG{n}{call} \PYG{o}{=} \PYG{n}{psyir}\PYG{o}{.}\PYG{n}{walk}\PYG{p}{(}\PYG{n}{Call}\PYG{p}{)}\PYG{p}{[}\PYG{l+m+mi}{0}\PYG{p}{]}
\PYG{g+gp}{\PYGZgt{}\PYGZgt{}\PYGZgt{} }\PYG{n}{inline\PYGZus{}trans} \PYG{o}{=} \PYG{n}{InlineTrans}\PYG{p}{(}\PYG{p}{)}
\PYG{g+gp}{\PYGZgt{}\PYGZgt{}\PYGZgt{} }\PYG{n}{inline\PYGZus{}trans}\PYG{o}{.}\PYG{n}{apply}\PYG{p}{(}\PYG{n}{call}\PYG{p}{)}
\PYG{g+gp}{\PYGZgt{}\PYGZgt{}\PYGZgt{} }\PYG{c+c1}{\PYGZsh{} Uncomment the following line to see a text view of the schedule}
\PYG{g+gp}{\PYGZgt{}\PYGZgt{}\PYGZgt{} }\PYG{c+c1}{\PYGZsh{} print(psyir.walk(Routine)[0].view())}
\PYG{g+gp}{\PYGZgt{}\PYGZgt{}\PYGZgt{} }\PYG{n+nb}{print}\PYG{p}{(}\PYG{n}{FortranWriter}\PYG{p}{(}\PYG{p}{)}\PYG{p}{(}\PYG{n}{psyir}\PYG{o}{.}\PYG{n}{walk}\PYG{p}{(}\PYG{n}{Routine}\PYG{p}{)}\PYG{p}{[}\PYG{l+m+mi}{0}\PYG{p}{]}\PYG{p}{)}\PYG{p}{)}
\PYG{g+go}{subroutine run\PYGZus{}it()}
\PYG{g+go}{  integer :: i}
\PYG{g+go}{  real, dimension(10) :: a}

\PYG{g+go}{  do i = 1, 10, 1}
\PYG{g+go}{    a(i) = 1.0}
\PYG{g+go}{    a(i) = 2.0 * a(i)}
\PYG{g+go}{  enddo}

\PYG{g+go}{end subroutine run\PYGZus{}it}
\end{sphinxVerbatim}

\begin{sphinxadmonition}{warning}{Warning:}
\sphinxAtStartPar
Routines/calls with any of the following characteristics are not
supported and will result in a TransformationError:
\begin{itemize}
\item {} 
\sphinxAtStartPar
the routine is not in the same file as the call;

\item {} 
\sphinxAtStartPar
the routine contains an early Return statement;

\item {} 
\sphinxAtStartPar
the routine contains a variable with UnknownInterface;

\item {} 
\sphinxAtStartPar
the routine contains a variable with StaticInterface;

\item {} 
\sphinxAtStartPar
the routine contains an UnsupportedType variable with
ArgumentInterface;

\item {} 
\sphinxAtStartPar
the routine has a named argument;

\item {} 
\sphinxAtStartPar
the shape of any array arguments as declared inside the routine does
not match the shape of the arrays being passed as arguments;

\item {} 
\sphinxAtStartPar
the routine accesses an un\sphinxhyphen{}resolved symbol;

\item {} 
\sphinxAtStartPar
the routine accesses a symbol declared in the Container to which it
belongs.

\end{itemize}

\sphinxAtStartPar
Some of these restrictions will be lifted by \#924.
\end{sphinxadmonition}


\begin{fulllineitems}

\pysigstartsignatures
\pysiglinewithargsret{\sphinxbfcode{\sphinxupquote{apply}}}{\sphinxparam{\DUrole{n,n}{node}}\sphinxparamcomma \sphinxparam{\DUrole{n,n}{options}\DUrole{o,o}{=}\DUrole{default_value}{None}}}{}
\pysigstopsignatures
\sphinxAtStartPar
Takes the body of the routine that is the target of the supplied
call and replaces the call with it.
\begin{quote}\begin{description}
\sphinxlineitem{Parameters}\begin{itemize}
\item {} 
\sphinxAtStartPar
\sphinxstyleliteralstrong{\sphinxupquote{node}} (\sphinxcode{\sphinxupquote{psyclone.psyir.nodes.Routine}}) \textendash{} target PSyIR node.

\item {} 
\sphinxAtStartPar
\sphinxstyleliteralstrong{\sphinxupquote{options}} (\sphinxstyleliteralemphasis{\sphinxupquote{Optional}}\sphinxstyleliteralemphasis{\sphinxupquote{{[}}}\sphinxstyleliteralemphasis{\sphinxupquote{Dict}}\sphinxstyleliteralemphasis{\sphinxupquote{{[}}}\sphinxhref{https://docs.python.org/3/library/stdtypes.html\#str}{\sphinxstyleliteralemphasis{\sphinxupquote{str}}}\sphinxstyleliteralemphasis{\sphinxupquote{, }}\sphinxstyleliteralemphasis{\sphinxupquote{Any}}\sphinxstyleliteralemphasis{\sphinxupquote{{]}}}\sphinxstyleliteralemphasis{\sphinxupquote{{]}}}) \textendash{} a dictionary with options for transformations.

\item {} 
\sphinxAtStartPar
\sphinxstyleliteralstrong{\sphinxupquote{options}}\sphinxstyleliteralstrong{\sphinxupquote{{[}}}\sphinxstyleliteralstrong{\sphinxupquote{"force"}}\sphinxstyleliteralstrong{\sphinxupquote{{]}}} (\sphinxhref{https://docs.python.org/3/library/functions.html\#bool}{\sphinxstyleliteralemphasis{\sphinxupquote{bool}}}) \textendash{} whether or not to permit the inlining
of Routines containing CodeBlocks. Default is False.

\end{itemize}

\end{description}\end{quote}

\end{fulllineitems}


\end{fulllineitems}



\bigskip\hrule\bigskip



\begin{fulllineitems}

\pysigstartsignatures
\pysigline{\sphinxbfcode{\sphinxupquote{class\DUrole{w,w}{  }}}\sphinxcode{\sphinxupquote{psyclone.domain.common.transformations.}}\sphinxbfcode{\sphinxupquote{KernelModuleInlineTrans}}}
\pysigstopsignatures
\sphinxAtStartPar
Brings the routine being called into the same Container as the call
site. For example:

\begin{sphinxVerbatim}[commandchars=\\\{\}]
\PYG{k+kn}{from}\PYG{+w}{ }\PYG{n+nn}{psyclone}\PYG{n+nn}{.}\PYG{n+nn}{domain}\PYG{n+nn}{.}\PYG{n+nn}{common}\PYG{n+nn}{.}\PYG{n+nn}{transformations}\PYG{+w}{ }\PYG{k+kn}{import} \PYGZbs{}
        \PYG{n}{KernelModuleInlineTrans}

\PYG{n}{inline\PYGZus{}trans} \PYG{o}{=} \PYG{n}{KernelModuleInlineTrans}\PYG{p}{(}\PYG{p}{)}
\PYG{n}{inline\PYGZus{}trans}\PYG{o}{.}\PYG{n}{apply}\PYG{p}{(}\PYG{n}{schedule}\PYG{o}{.}\PYG{n}{walk}\PYG{p}{(}\PYG{n}{CodedKern}\PYG{p}{)}\PYG{p}{[}\PYG{l+m+mi}{0}\PYG{p}{]}\PYG{p}{)}

\PYG{n+nb}{print}\PYG{p}{(}\PYG{n}{schedule}\PYG{o}{.}\PYG{n}{parent}\PYG{o}{.}\PYG{n}{view}\PYG{p}{(}\PYG{p}{)}\PYG{p}{)}
\end{sphinxVerbatim}

\begin{sphinxadmonition}{warning}{Warning:}
\sphinxAtStartPar
Not all Routines can be moved. This transformation will reject
attempts to move routines that access private data in the
original Container.
\end{sphinxadmonition}


\begin{fulllineitems}

\pysigstartsignatures
\pysiglinewithargsret{\sphinxbfcode{\sphinxupquote{apply}}}{\sphinxparam{\DUrole{n,n}{node}}\sphinxparamcomma \sphinxparam{\DUrole{n,n}{options}\DUrole{o,o}{=}\DUrole{default_value}{None}}}{}
\pysigstopsignatures
\sphinxAtStartPar
Bring the kernel subroutine into this Container.
\begin{quote}\begin{description}
\sphinxlineitem{Parameters}\begin{itemize}
\item {} 
\sphinxAtStartPar
\sphinxstyleliteralstrong{\sphinxupquote{node}} (\sphinxcode{\sphinxupquote{psyclone.psyGen.CodedKern}}) \textendash{} the kernel to module\sphinxhyphen{}inline.

\item {} 
\sphinxAtStartPar
\sphinxstyleliteralstrong{\sphinxupquote{options}} (\sphinxstyleliteralemphasis{\sphinxupquote{Optional}}\sphinxstyleliteralemphasis{\sphinxupquote{{[}}}\sphinxstyleliteralemphasis{\sphinxupquote{Dict}}\sphinxstyleliteralemphasis{\sphinxupquote{{[}}}\sphinxhref{https://docs.python.org/3/library/stdtypes.html\#str}{\sphinxstyleliteralemphasis{\sphinxupquote{str}}}\sphinxstyleliteralemphasis{\sphinxupquote{, }}\sphinxstyleliteralemphasis{\sphinxupquote{Any}}\sphinxstyleliteralemphasis{\sphinxupquote{{]}}}\sphinxstyleliteralemphasis{\sphinxupquote{{]}}}) \textendash{} a dictionary with options for transformations.

\end{itemize}

\sphinxlineitem{Raises}\begin{itemize}
\item {} 
\sphinxAtStartPar
\sphinxstyleliteralstrong{\sphinxupquote{TransformationError}} \textendash{} if the called Routine cannot be brought
into this Container because of a name clash with another Routine.

\item {} 
\sphinxAtStartPar
\sphinxhref{https://docs.python.org/3/library/exceptions.html\#NotImplementedError}{\sphinxstyleliteralstrong{\sphinxupquote{NotImplementedError}}} \textendash{} if node is a Call (rather than a
CodedKern) and the name of the called routine does not match that
of the caller.

\end{itemize}

\end{description}\end{quote}

\end{fulllineitems}


\end{fulllineitems}



\bigskip\hrule\bigskip



\begin{fulllineitems}

\pysigstartsignatures
\pysigline{\sphinxbfcode{\sphinxupquote{class\DUrole{w,w}{  }}}\sphinxcode{\sphinxupquote{psyclone.psyir.transformations.}}\sphinxbfcode{\sphinxupquote{LoopFuseTrans}}}
\pysigstopsignatures
\sphinxAtStartPar
Provides a generic loop\sphinxhyphen{}fuse transformation to two Nodes in the
PSyIR of a Schedule after performing validity checks for the supplied
Nodes. Examples are given in the descriptions of any children classes.

\sphinxAtStartPar
If loops have different named loop variables, when possible the loop
variable of the second loop will be renamed to be the same as the first
loop. This has the side effect that the second loop’s variable will no
longer have its value modified, with the expectation that that value
isn’t used anymore.

\sphinxAtStartPar
Note that the validation of this transformation still has several
shortcomings, especially for domain API loops. Use at your own risk.


\begin{fulllineitems}

\pysigstartsignatures
\pysiglinewithargsret{\sphinxbfcode{\sphinxupquote{apply}}}{\sphinxparam{\DUrole{n,n}{node1}}\sphinxparamcomma \sphinxparam{\DUrole{n,n}{node2}}\sphinxparamcomma \sphinxparam{\DUrole{n,n}{options}\DUrole{o,o}{=}\DUrole{default_value}{None}}}{}
\pysigstopsignatures
\sphinxAtStartPar
Fuses two loops represented by \sphinxtitleref{psyclone.psyir.nodes.Node} objects
after performing validity checks.

\sphinxAtStartPar
If the two loops don’t have the same loop variable, the second loop’s
variable (and references to it inside the loop) will be changed to be
references to the first loop’s variable before merging. This has the
side effect that the second loop’s variable will no longer have its
value modified, with the expectation that that value isn’t used after.
\begin{quote}\begin{description}
\sphinxlineitem{Parameters}\begin{itemize}
\item {} 
\sphinxAtStartPar
\sphinxstyleliteralstrong{\sphinxupquote{node1}} (\sphinxcode{\sphinxupquote{psyclone.psyir.nodes.Node}}) \textendash{} the first Node that is being checked.

\item {} 
\sphinxAtStartPar
\sphinxstyleliteralstrong{\sphinxupquote{node2}} (\sphinxcode{\sphinxupquote{psyclone.psyir.nodes.Node}}) \textendash{} the second Node that is being checked.

\item {} 
\sphinxAtStartPar
\sphinxstyleliteralstrong{\sphinxupquote{options}} (\sphinxstyleliteralemphasis{\sphinxupquote{Optional}}\sphinxstyleliteralemphasis{\sphinxupquote{{[}}}\sphinxstyleliteralemphasis{\sphinxupquote{Dict}}\sphinxstyleliteralemphasis{\sphinxupquote{{[}}}\sphinxhref{https://docs.python.org/3/library/stdtypes.html\#str}{\sphinxstyleliteralemphasis{\sphinxupquote{str}}}\sphinxstyleliteralemphasis{\sphinxupquote{, }}\sphinxstyleliteralemphasis{\sphinxupquote{Any}}\sphinxstyleliteralemphasis{\sphinxupquote{{]}}}\sphinxstyleliteralemphasis{\sphinxupquote{{]}}}) \textendash{} a dictionary with options for transformations.

\end{itemize}

\end{description}\end{quote}

\end{fulllineitems}


\end{fulllineitems}



\bigskip\hrule\bigskip



\begin{fulllineitems}

\pysigstartsignatures
\pysigline{\sphinxbfcode{\sphinxupquote{class\DUrole{w,w}{  }}}\sphinxcode{\sphinxupquote{psyclone.psyir.transformations.}}\sphinxbfcode{\sphinxupquote{LoopSwapTrans}}}
\pysigstopsignatures
\sphinxAtStartPar
Provides a loop\sphinxhyphen{}swap transformation, e.g.:

\begin{sphinxVerbatim}[commandchars=\\\{\}]
\PYG{k}{DO }\PYG{n}{j}\PYG{o}{=}\PYG{l+m+mi}{1}\PYG{p}{,}\PYG{+w}{ }\PYG{n}{m}
\PYG{+w}{    }\PYG{k}{DO }\PYG{n}{i}\PYG{o}{=}\PYG{l+m+mi}{1}\PYG{p}{,}\PYG{+w}{ }\PYG{n}{n}
\end{sphinxVerbatim}

\sphinxAtStartPar
becomes:

\begin{sphinxVerbatim}[commandchars=\\\{\}]
\PYG{k}{DO }\PYG{n}{i}\PYG{o}{=}\PYG{l+m+mi}{1}\PYG{p}{,}\PYG{+w}{ }\PYG{n}{n}
\PYG{+w}{    }\PYG{k}{DO }\PYG{n}{j}\PYG{o}{=}\PYG{l+m+mi}{1}\PYG{p}{,}\PYG{+w}{ }\PYG{n}{m}
\end{sphinxVerbatim}

\sphinxAtStartPar
This transform is used as follows:

\begin{sphinxVerbatim}[commandchars=\\\{\}]
\PYG{g+gp}{\PYGZgt{}\PYGZgt{}\PYGZgt{} }\PYG{k+kn}{from}\PYG{+w}{ }\PYG{n+nn}{psyclone}\PYG{n+nn}{.}\PYG{n+nn}{parse}\PYG{n+nn}{.}\PYG{n+nn}{algorithm}\PYG{+w}{ }\PYG{k+kn}{import} \PYG{n}{parse}
\PYG{g+gp}{\PYGZgt{}\PYGZgt{}\PYGZgt{} }\PYG{k+kn}{from}\PYG{+w}{ }\PYG{n+nn}{psyclone}\PYG{n+nn}{.}\PYG{n+nn}{psyGen}\PYG{+w}{ }\PYG{k+kn}{import} \PYG{n}{PSyFactory}
\PYG{g+gp}{\PYGZgt{}\PYGZgt{}\PYGZgt{} }\PYG{n}{ast}\PYG{p}{,} \PYG{n}{invokeInfo} \PYG{o}{=} \PYG{n}{parse}\PYG{p}{(}\PYG{l+s+s2}{\PYGZdq{}}\PYG{l+s+s2}{shallow\PYGZus{}alg.f90}\PYG{l+s+s2}{\PYGZdq{}}\PYG{p}{)}
\PYG{g+gp}{\PYGZgt{}\PYGZgt{}\PYGZgt{} }\PYG{n}{psy} \PYG{o}{=} \PYG{n}{PSyFactory}\PYG{p}{(}\PYG{l+s+s2}{\PYGZdq{}}\PYG{l+s+s2}{gocean}\PYG{l+s+s2}{\PYGZdq{}}\PYG{p}{)}\PYG{o}{.}\PYG{n}{create}\PYG{p}{(}\PYG{n}{invokeInfo}\PYG{p}{)}
\PYG{g+gp}{\PYGZgt{}\PYGZgt{}\PYGZgt{} }\PYG{n}{schedule} \PYG{o}{=} \PYG{n}{psy}\PYG{o}{.}\PYG{n}{invokes}\PYG{o}{.}\PYG{n}{get}\PYG{p}{(}\PYG{l+s+s1}{\PYGZsq{}}\PYG{l+s+s1}{invoke\PYGZus{}0}\PYG{l+s+s1}{\PYGZsq{}}\PYG{p}{)}\PYG{o}{.}\PYG{n}{schedule}
\PYG{g+gp}{\PYGZgt{}\PYGZgt{}\PYGZgt{} }\PYG{c+c1}{\PYGZsh{} Uncomment the following line to see a text view of the schedule}
\PYG{g+gp}{\PYGZgt{}\PYGZgt{}\PYGZgt{} }\PYG{c+c1}{\PYGZsh{} print(schedule.view())}
\PYG{g+gp}{\PYGZgt{}\PYGZgt{}\PYGZgt{}}
\PYG{g+gp}{\PYGZgt{}\PYGZgt{}\PYGZgt{} }\PYG{k+kn}{from}\PYG{+w}{ }\PYG{n+nn}{psyclone}\PYG{n+nn}{.}\PYG{n+nn}{transformations}\PYG{+w}{ }\PYG{k+kn}{import} \PYG{n}{LoopSwapTrans}
\PYG{g+gp}{\PYGZgt{}\PYGZgt{}\PYGZgt{} }\PYG{n}{swap} \PYG{o}{=} \PYG{n}{LoopSwapTrans}\PYG{p}{(}\PYG{p}{)}
\PYG{g+gp}{\PYGZgt{}\PYGZgt{}\PYGZgt{} }\PYG{n}{swap}\PYG{o}{.}\PYG{n}{apply}\PYG{p}{(}\PYG{n}{schedule}\PYG{o}{.}\PYG{n}{children}\PYG{p}{[}\PYG{l+m+mi}{0}\PYG{p}{]}\PYG{p}{)}
\PYG{g+gp}{\PYGZgt{}\PYGZgt{}\PYGZgt{} }\PYG{c+c1}{\PYGZsh{} Uncomment the following line to see a text view of the schedule}
\PYG{g+gp}{\PYGZgt{}\PYGZgt{}\PYGZgt{} }\PYG{c+c1}{\PYGZsh{} print(schedule.view())}
\end{sphinxVerbatim}


\begin{fulllineitems}

\pysigstartsignatures
\pysiglinewithargsret{\sphinxbfcode{\sphinxupquote{apply}}}{\sphinxparam{\DUrole{n,n}{node}}\sphinxparamcomma \sphinxparam{\DUrole{n,n}{options}\DUrole{o,o}{=}\DUrole{default_value}{None}}}{}
\pysigstopsignatures
\sphinxAtStartPar
The argument \sphinxcode{\sphinxupquote{outer}} must be a loop which has exactly
one inner loop. This transform then swaps the outer and inner loop.
\begin{quote}\begin{description}
\sphinxlineitem{Parameters}\begin{itemize}
\item {} 
\sphinxAtStartPar
\sphinxstyleliteralstrong{\sphinxupquote{outer}} (\sphinxcode{\sphinxupquote{psyclone.psyir.nodes.Loop}}) \textendash{} the node representing the outer loop.

\item {} 
\sphinxAtStartPar
\sphinxstyleliteralstrong{\sphinxupquote{options}} (\sphinxstyleliteralemphasis{\sphinxupquote{Optional}}\sphinxstyleliteralemphasis{\sphinxupquote{{[}}}\sphinxstyleliteralemphasis{\sphinxupquote{Dict}}\sphinxstyleliteralemphasis{\sphinxupquote{{[}}}\sphinxhref{https://docs.python.org/3/library/stdtypes.html\#str}{\sphinxstyleliteralemphasis{\sphinxupquote{str}}}\sphinxstyleliteralemphasis{\sphinxupquote{, }}\sphinxstyleliteralemphasis{\sphinxupquote{Any}}\sphinxstyleliteralemphasis{\sphinxupquote{{]}}}\sphinxstyleliteralemphasis{\sphinxupquote{{]}}}) \textendash{} a dictionary with options for transformations.

\end{itemize}

\sphinxlineitem{Raises}
\sphinxAtStartPar
\sphinxstyleliteralstrong{\sphinxupquote{TransformationError}} \textendash{} if the supplied node does not                                      allow a loop swap to be done.

\end{description}\end{quote}

\end{fulllineitems}


\end{fulllineitems}



\bigskip\hrule\bigskip



\begin{fulllineitems}

\pysigstartsignatures
\pysigline{\sphinxbfcode{\sphinxupquote{class\DUrole{w,w}{  }}}\sphinxcode{\sphinxupquote{psyclone.psyir.transformations.}}\sphinxbfcode{\sphinxupquote{LoopTiling2DTrans}}}
\pysigstopsignatures
\sphinxAtStartPar
Apply a 2D loop tiling transformation to a loop. For example:

\begin{sphinxVerbatim}[commandchars=\\\{\}]
\PYG{g+gp}{\PYGZgt{}\PYGZgt{}\PYGZgt{} }\PYG{k+kn}{from}\PYG{+w}{ }\PYG{n+nn}{psyclone}\PYG{n+nn}{.}\PYG{n+nn}{psyir}\PYG{n+nn}{.}\PYG{n+nn}{frontend}\PYG{n+nn}{.}\PYG{n+nn}{fortran}\PYG{+w}{ }\PYG{k+kn}{import} \PYG{n}{FortranReader}
\PYG{g+gp}{\PYGZgt{}\PYGZgt{}\PYGZgt{} }\PYG{k+kn}{from}\PYG{+w}{ }\PYG{n+nn}{psyclone}\PYG{n+nn}{.}\PYG{n+nn}{psyir}\PYG{n+nn}{.}\PYG{n+nn}{nodes}\PYG{+w}{ }\PYG{k+kn}{import} \PYG{n}{Loop}
\PYG{g+gp}{\PYGZgt{}\PYGZgt{}\PYGZgt{} }\PYG{k+kn}{from}\PYG{+w}{ }\PYG{n+nn}{psyclone}\PYG{n+nn}{.}\PYG{n+nn}{psyir}\PYG{n+nn}{.}\PYG{n+nn}{transformations}\PYG{+w}{ }\PYG{k+kn}{import} \PYG{n}{LoopTiling2DTrans}
\PYG{g+gp}{\PYGZgt{}\PYGZgt{}\PYGZgt{} }\PYG{n}{psyir} \PYG{o}{=} \PYG{n}{FortranReader}\PYG{p}{(}\PYG{p}{)}\PYG{o}{.}\PYG{n}{psyir\PYGZus{}from\PYGZus{}source}\PYG{p}{(}\PYG{l+s+s2}{\PYGZdq{}\PYGZdq{}\PYGZdq{}}
\PYG{g+gp}{... }\PYG{l+s+s2}{subroutine sub()}
\PYG{g+gp}{... }\PYG{l+s+s2}{    integer :: ji, tmp(100)}
\PYG{g+gp}{... }\PYG{l+s+s2}{    do i=1, 100}
\PYG{g+gp}{... }\PYG{l+s+s2}{      do j=1, 100}
\PYG{g+gp}{... }\PYG{l+s+s2}{        tmp(i, j) = 2 * tmp(i, j)}
\PYG{g+gp}{... }\PYG{l+s+s2}{      enddo}
\PYG{g+gp}{... }\PYG{l+s+s2}{    enddo}
\PYG{g+gp}{... }\PYG{l+s+s2}{end subroutine sub}\PYG{l+s+s2}{\PYGZdq{}\PYGZdq{}\PYGZdq{}}\PYG{p}{)}
\PYG{g+gp}{\PYGZgt{}\PYGZgt{}\PYGZgt{} }\PYG{n}{loop} \PYG{o}{=} \PYG{n}{psyir}\PYG{o}{.}\PYG{n}{walk}\PYG{p}{(}\PYG{n}{Loop}\PYG{p}{)}\PYG{p}{[}\PYG{l+m+mi}{0}\PYG{p}{]}
\PYG{g+gp}{\PYGZgt{}\PYGZgt{}\PYGZgt{} }\PYG{n}{LoopTiling2DTrans}\PYG{p}{(}\PYG{p}{)}\PYG{o}{.}\PYG{n}{apply}\PYG{p}{(}\PYG{n}{loop}\PYG{p}{)}
\end{sphinxVerbatim}

\sphinxAtStartPar
will generate:

\begin{sphinxVerbatim}[commandchars=\\\{\}]
\PYG{k}{subroutine }\PYG{n}{sub}\PYG{p}{(}\PYG{p}{)}
\PYG{+w}{    }\PYG{k+kt}{integer}\PYG{+w}{ }\PYG{k+kd}{::}\PYG{+w}{ }\PYG{n}{ji}
\PYG{+w}{    }\PYG{k+kt}{integer}\PYG{p}{,}\PYG{+w}{ }\PYG{k}{dimension}\PYG{p}{(}\PYG{l+m+mi}{100}\PYG{p}{)}\PYG{+w}{ }\PYG{k+kd}{::}\PYG{+w}{ }\PYG{n}{tmp}
\PYG{+w}{    }\PYG{k+kt}{integer}\PYG{+w}{ }\PYG{k+kd}{::}\PYG{+w}{ }\PYG{n}{ji\PYGZus{}el\PYGZus{}inner}
\PYG{+w}{    }\PYG{k+kt}{integer}\PYG{+w}{ }\PYG{k+kd}{::}\PYG{+w}{ }\PYG{n}{ji\PYGZus{}out\PYGZus{}var}
\PYG{+w}{    }\PYG{k}{do }\PYG{n}{i\PYGZus{}out\PYGZus{}var}\PYG{+w}{ }\PYG{o}{=}\PYG{+w}{ }\PYG{l+m+mi}{1}\PYG{p}{,}\PYG{+w}{ }\PYG{l+m+mi}{100}\PYG{p}{,}\PYG{+w}{ }\PYG{l+m+mi}{32}
\PYG{+w}{      }\PYG{n}{i\PYGZus{}el\PYGZus{}inner}\PYG{+w}{ }\PYG{o}{=}\PYG{+w}{ }\PYG{n+nb}{MIN}\PYG{p}{(}\PYG{n}{i\PYGZus{}out\PYGZus{}var}\PYG{+w}{ }\PYG{o}{+}\PYG{+w}{ }\PYG{p}{(}\PYG{l+m+mi}{32}\PYG{+w}{ }\PYG{o}{\PYGZhy{}}\PYG{+w}{ }\PYG{l+m+mi}{1}\PYG{p}{)}\PYG{p}{,}\PYG{+w}{ }\PYG{l+m+mi}{100}\PYG{p}{)}
\PYG{+w}{      }\PYG{k}{do }\PYG{n}{j\PYGZus{}out\PYGZus{}var}\PYG{+w}{ }\PYG{o}{=}\PYG{+w}{ }\PYG{l+m+mi}{1}\PYG{p}{,}\PYG{+w}{ }\PYG{l+m+mi}{100}\PYG{p}{,}\PYG{+w}{ }\PYG{l+m+mi}{32}
\PYG{+w}{        }\PYG{k}{do }\PYG{n}{i}\PYG{+w}{ }\PYG{o}{=}\PYG{+w}{ }\PYG{n}{i\PYGZus{}out\PYGZus{}var}\PYG{p}{,}\PYG{+w}{ }\PYG{n}{i\PYGZus{}el\PYGZus{}inner}\PYG{p}{,}\PYG{+w}{ }\PYG{l+m+mi}{1}
\PYG{+w}{          }\PYG{n}{j\PYGZus{}el\PYGZus{}inner}\PYG{+w}{ }\PYG{o}{=}\PYG{+w}{ }\PYG{n+nb}{MIN}\PYG{p}{(}\PYG{n}{j\PYGZus{}out\PYGZus{}var}\PYG{+w}{ }\PYG{o}{+}\PYG{+w}{ }\PYG{p}{(}\PYG{l+m+mi}{32}\PYG{+w}{ }\PYG{o}{\PYGZhy{}}\PYG{+w}{ }\PYG{l+m+mi}{1}\PYG{p}{)}\PYG{p}{,}\PYG{+w}{ }\PYG{l+m+mi}{100}\PYG{p}{)}
\PYG{+w}{          }\PYG{k}{do }\PYG{n}{j}\PYG{+w}{ }\PYG{o}{=}\PYG{+w}{ }\PYG{n}{j\PYGZus{}out\PYGZus{}var}\PYG{p}{,}\PYG{+w}{ }\PYG{n}{j\PYGZus{}el\PYGZus{}inner}\PYG{p}{,}\PYG{+w}{ }\PYG{l+m+mi}{1}
\PYG{+w}{            }\PYG{n}{tmp}\PYG{p}{(}\PYG{n}{i}\PYG{p}{,}\PYG{+w}{ }\PYG{n}{j}\PYG{p}{)}\PYG{+w}{ }\PYG{o}{=}\PYG{+w}{ }\PYG{l+m+mi}{2}\PYG{+w}{ }\PYG{o}{*}\PYG{+w}{ }\PYG{n}{tmp}\PYG{p}{(}\PYG{n}{i}\PYG{p}{,}\PYG{+w}{ }\PYG{n}{j}\PYG{p}{)}
\PYG{+w}{          }\PYG{k}{enddo}
\PYG{k}{        }\PYG{k}{enddo}
\PYG{k}{      }\PYG{k}{enddo}
\PYG{k}{    }\PYG{k}{enddo}
\PYG{k}{end }\PYG{k}{subroutine }\PYG{n}{sub}
\end{sphinxVerbatim}


\begin{fulllineitems}

\pysigstartsignatures
\pysiglinewithargsret{\sphinxbfcode{\sphinxupquote{apply}}}{\sphinxparam{\DUrole{n,n}{node}}\sphinxparamcomma \sphinxparam{\DUrole{n,n}{options}\DUrole{o,o}{=}\DUrole{default_value}{None}}}{}
\pysigstopsignatures
\sphinxAtStartPar
Converts the given 2D Loop construct into a tiled version of the nested
loops.
\begin{quote}\begin{description}
\sphinxlineitem{Parameters}\begin{itemize}
\item {} 
\sphinxAtStartPar
\sphinxstyleliteralstrong{\sphinxupquote{node}} (\sphinxcode{\sphinxupquote{psyclone.psyir.nodes.Loop}}) \textendash{} the loop to transform.

\item {} 
\sphinxAtStartPar
\sphinxstyleliteralstrong{\sphinxupquote{options}} (\sphinxstyleliteralemphasis{\sphinxupquote{Optional}}\sphinxstyleliteralemphasis{\sphinxupquote{{[}}}\sphinxstyleliteralemphasis{\sphinxupquote{Dict}}\sphinxstyleliteralemphasis{\sphinxupquote{{[}}}\sphinxhref{https://docs.python.org/3/library/stdtypes.html\#str}{\sphinxstyleliteralemphasis{\sphinxupquote{str}}}\sphinxstyleliteralemphasis{\sphinxupquote{, }}\sphinxstyleliteralemphasis{\sphinxupquote{Any}}\sphinxstyleliteralemphasis{\sphinxupquote{{]}}}\sphinxstyleliteralemphasis{\sphinxupquote{{]}}}) \textendash{} a dict with options for transformations.

\item {} 
\sphinxAtStartPar
\sphinxstyleliteralstrong{\sphinxupquote{options}}\sphinxstyleliteralstrong{\sphinxupquote{{[}}}\sphinxstyleliteralstrong{\sphinxupquote{"tilesize"}}\sphinxstyleliteralstrong{\sphinxupquote{{]}}} (\sphinxhref{https://docs.python.org/3/library/functions.html\#int}{\sphinxstyleliteralemphasis{\sphinxupquote{int}}}) \textendash{} The size of the resulting tile,                 currently square tiles are always used. If not                 specified, the value 32 is used.

\end{itemize}

\end{description}\end{quote}

\end{fulllineitems}


\end{fulllineitems}



\bigskip\hrule\bigskip



\begin{fulllineitems}

\pysigstartsignatures
\pysigline{\sphinxbfcode{\sphinxupquote{class\DUrole{w,w}{  }}}\sphinxcode{\sphinxupquote{psyclone.psyir.transformations.}}\sphinxbfcode{\sphinxupquote{Matmul2CodeTrans}}}
\pysigstopsignatures
\sphinxAtStartPar
Provides a transformation from a PSyIR MATMUL Operator node to
equivalent code in a PSyIR tree. Validity checks are also
performed.

\sphinxAtStartPar
For a matrix\sphinxhyphen{}vector multiplication, if the dimensions of \sphinxcode{\sphinxupquote{R}}, \sphinxcode{\sphinxupquote{A}},
and \sphinxcode{\sphinxupquote{B}} are \sphinxcode{\sphinxupquote{R(N)}}, \sphinxcode{\sphinxupquote{A(N,M)}}, \sphinxcode{\sphinxupquote{B(M)}}, the transformation replaces:

\begin{sphinxVerbatim}[commandchars=\\\{\}]
\PYG{n}{R}\PYG{o}{=}\PYG{n+nb}{MATMUL}\PYG{p}{(}\PYG{n}{A}\PYG{p}{,}\PYG{n}{B}\PYG{p}{)}
\end{sphinxVerbatim}

\sphinxAtStartPar
with the following code:

\begin{sphinxVerbatim}[commandchars=\\\{\}]
\PYG{k}{do }\PYG{n}{i}\PYG{o}{=}\PYG{l+m+mi}{1}\PYG{p}{,}\PYG{n}{N}
\PYG{+w}{    }\PYG{n}{R}\PYG{p}{(}\PYG{n}{i}\PYG{p}{)}\PYG{+w}{ }\PYG{o}{=}\PYG{+w}{ }\PYG{l+m+mf}{0.0}
\PYG{+w}{    }\PYG{k}{do }\PYG{n}{j}\PYG{o}{=}\PYG{l+m+mi}{1}\PYG{p}{,}\PYG{n}{M}
\PYG{+w}{        }\PYG{n}{R}\PYG{p}{(}\PYG{n}{i}\PYG{p}{)}\PYG{+w}{ }\PYG{o}{=}\PYG{+w}{ }\PYG{n}{R}\PYG{p}{(}\PYG{n}{i}\PYG{p}{)}\PYG{+w}{ }\PYG{o}{+}\PYG{+w}{ }\PYG{n}{A}\PYG{p}{(}\PYG{n}{i}\PYG{p}{,}\PYG{n}{j}\PYG{p}{)}\PYG{+w}{ }\PYG{o}{*}\PYG{+w}{ }\PYG{n}{B}\PYG{p}{(}\PYG{n}{j}\PYG{p}{)}
\end{sphinxVerbatim}

\sphinxAtStartPar
For a matrix\sphinxhyphen{}matrix multiplication, if the dimensions of \sphinxcode{\sphinxupquote{R}}, \sphinxcode{\sphinxupquote{A}},
and \sphinxcode{\sphinxupquote{B}} are \sphinxcode{\sphinxupquote{R(P,M)}}, \sphinxcode{\sphinxupquote{A(P,N)}}, \sphinxcode{\sphinxupquote{B(N,M)}}, the MATMUL is replaced
with the following code:

\begin{sphinxVerbatim}[commandchars=\\\{\}]
\PYG{k}{do }\PYG{n}{j}\PYG{o}{=}\PYG{l+m+mi}{1}\PYG{p}{,}\PYG{n}{M}
\PYG{+w}{    }\PYG{k}{do }\PYG{n}{i}\PYG{o}{=}\PYG{l+m+mi}{1}\PYG{p}{,}\PYG{n}{P}
\PYG{+w}{        }\PYG{n}{R}\PYG{p}{(}\PYG{n}{i}\PYG{p}{,}\PYG{n}{j}\PYG{p}{)}\PYG{+w}{ }\PYG{o}{=}\PYG{+w}{ }\PYG{l+m+mf}{0.0}
\PYG{+w}{        }\PYG{k}{do }\PYG{n}{ii}\PYG{o}{=}\PYG{l+m+mi}{1}\PYG{p}{,}\PYG{n}{N}
\PYG{+w}{            }\PYG{n}{R}\PYG{p}{(}\PYG{n}{i}\PYG{p}{,}\PYG{n}{j}\PYG{p}{)}\PYG{+w}{ }\PYG{o}{=}\PYG{+w}{ }\PYG{n}{R}\PYG{p}{(}\PYG{n}{i}\PYG{p}{,}\PYG{n}{j}\PYG{p}{)}\PYG{+w}{ }\PYG{o}{+}\PYG{+w}{ }\PYG{n}{A}\PYG{p}{(}\PYG{n}{i}\PYG{p}{,}\PYG{n}{ii}\PYG{p}{)}\PYG{+w}{ }\PYG{o}{*}\PYG{+w}{ }\PYG{n}{B}\PYG{p}{(}\PYG{n}{ii}\PYG{p}{,}\PYG{n}{j}\PYG{p}{)}
\end{sphinxVerbatim}

\sphinxAtStartPar
Note that this transformation does \sphinxstyleemphasis{not} support the case where \sphinxcode{\sphinxupquote{A}} is
a rank\sphinxhyphen{}1 array.


\begin{fulllineitems}

\pysigstartsignatures
\pysiglinewithargsret{\sphinxbfcode{\sphinxupquote{apply}}}{\sphinxparam{\DUrole{n,n}{node}}\sphinxparamcomma \sphinxparam{\DUrole{n,n}{options}\DUrole{o,o}{=}\DUrole{default_value}{None}}}{}
\pysigstopsignatures
\sphinxAtStartPar
Apply the MATMUL intrinsic conversion transformation to the
specified node. This node must be a MATMUL IntrinsicCall. The first
argument must currently have two dimensions while the second must have
either one or two dimensions. Each argument is permitted to have
additional dimensions (i.e. more than 2) but in each case it is only
the first one or two which may be ranges. Further, the ranges must
currently be for the full index space for that dimension (i.e. array
subsections are not supported). If the transformation is
successful then an assignment which includes a MATMUL
IntrinsicCall node is converted to equivalent inline code.
\begin{quote}\begin{description}
\sphinxlineitem{Parameters}\begin{itemize}
\item {} 
\sphinxAtStartPar
\sphinxstyleliteralstrong{\sphinxupquote{node}} (\sphinxcode{\sphinxupquote{psyclone.psyir.nodes.IntrinsicCall}}) \textendash{} a MATMUL IntrinsicCall node.

\item {} 
\sphinxAtStartPar
\sphinxstyleliteralstrong{\sphinxupquote{options}} (\sphinxstyleliteralemphasis{\sphinxupquote{Optional}}\sphinxstyleliteralemphasis{\sphinxupquote{{[}}}\sphinxstyleliteralemphasis{\sphinxupquote{Dict}}\sphinxstyleliteralemphasis{\sphinxupquote{{[}}}\sphinxhref{https://docs.python.org/3/library/stdtypes.html\#str}{\sphinxstyleliteralemphasis{\sphinxupquote{str}}}\sphinxstyleliteralemphasis{\sphinxupquote{, }}\sphinxstyleliteralemphasis{\sphinxupquote{Any}}\sphinxstyleliteralemphasis{\sphinxupquote{{]}}}\sphinxstyleliteralemphasis{\sphinxupquote{{]}}}) \textendash{} options for the transformation.

\end{itemize}

\end{description}\end{quote}

\end{fulllineitems}


\end{fulllineitems}


\begin{sphinxadmonition}{note}{Note:}
\sphinxAtStartPar
This transformation is currently limited to translating the
matrix vector form of MATMUL to equivalent PSyIR code.
\end{sphinxadmonition}


\bigskip\hrule\bigskip



\begin{fulllineitems}

\pysigstartsignatures
\pysigline{\sphinxbfcode{\sphinxupquote{class\DUrole{w,w}{  }}}\sphinxcode{\sphinxupquote{psyclone.psyir.transformations.}}\sphinxbfcode{\sphinxupquote{Max2CodeTrans}}}
\pysigstopsignatures
\sphinxAtStartPar
Provides a transformation from a PSyIR MAX Intrinsic node to
equivalent code in a PSyIR tree. Validity checks are also
performed (by a parent class).

\sphinxAtStartPar
The transformation replaces

\begin{sphinxVerbatim}[commandchars=\\\{\}]
\PYG{n}{R} \PYG{o}{=} \PYG{n}{MAX}\PYG{p}{(}\PYG{n}{A}\PYG{p}{,} \PYG{n}{B}\PYG{p}{,} \PYG{n}{C} \PYG{o}{.}\PYG{o}{.}\PYG{o}{.}\PYG{p}{)}
\end{sphinxVerbatim}

\sphinxAtStartPar
with the following logic:

\begin{sphinxVerbatim}[commandchars=\\\{\}]
\PYG{n}{R} \PYG{o}{=} \PYG{n}{A}
\PYG{k}{if} \PYG{n}{B} \PYG{o}{\PYGZgt{}} \PYG{n}{R}\PYG{p}{:}
    \PYG{n}{R} \PYG{o}{=} \PYG{n}{B}
\PYG{k}{if} \PYG{n}{C} \PYG{o}{\PYGZgt{}} \PYG{n}{R}\PYG{p}{:}
    \PYG{n}{R} \PYG{o}{=} \PYG{n}{C}
\PYG{o}{.}\PYG{o}{.}\PYG{o}{.}
\end{sphinxVerbatim}


\begin{fulllineitems}

\pysigstartsignatures
\pysiglinewithargsret{\sphinxbfcode{\sphinxupquote{apply}}}{\sphinxparam{\DUrole{n,n}{node}}\sphinxparamcomma \sphinxparam{\DUrole{n,n}{options}\DUrole{o,o}{=}\DUrole{default_value}{None}}}{}
\pysigstopsignatures
\sphinxAtStartPar
Apply this utility transformation to the specified node. This node
must be a MIN or MAX IntrinsicCall. The intrinsic is converted to
equivalent inline code. This is implemented as a PSyIR transform from:

\begin{sphinxVerbatim}[commandchars=\\\{\}]
\PYG{n}{R} \PYG{o}{=} \PYG{o}{.}\PYG{o}{.}\PYG{o}{.} \PYG{p}{[}\PYG{n}{MIN} \PYG{o+ow}{or} \PYG{n}{MAX}\PYG{p}{]}\PYG{p}{(}\PYG{n}{A}\PYG{p}{,} \PYG{n}{B}\PYG{p}{,} \PYG{n}{C} \PYG{o}{.}\PYG{o}{.}\PYG{o}{.}\PYG{p}{)} \PYG{o}{.}\PYG{o}{.}\PYG{o}{.}
\end{sphinxVerbatim}

\sphinxAtStartPar
to:

\begin{sphinxVerbatim}[commandchars=\\\{\}]
\PYG{n}{res} \PYG{o}{=} \PYG{n}{A}
\PYG{n}{tmp} \PYG{o}{=} \PYG{n}{B}
\PYG{n}{IF} \PYG{n}{tmp} \PYG{p}{[}\PYG{o}{\PYGZlt{}} \PYG{o+ow}{or} \PYG{o}{\PYGZgt{}}\PYG{p}{]} \PYG{n}{res}\PYG{p}{:}
    \PYG{n}{res} \PYG{o}{=} \PYG{n}{tmp}
\PYG{n}{tmp} \PYG{o}{=} \PYG{n}{C}
\PYG{n}{IF} \PYG{n}{tmp} \PYG{p}{[}\PYG{o}{\PYGZlt{}} \PYG{o+ow}{or} \PYG{o}{\PYGZgt{}}\PYG{p}{]} \PYG{n}{res}\PYG{p}{:}
    \PYG{n}{res} \PYG{o}{=} \PYG{n}{tmp}
\PYG{o}{.}\PYG{o}{.}\PYG{o}{.}
\PYG{n}{R} \PYG{o}{=} \PYG{o}{.}\PYG{o}{.}\PYG{o}{.} \PYG{n}{res} \PYG{o}{.}\PYG{o}{.}\PYG{o}{.}
\end{sphinxVerbatim}

\sphinxAtStartPar
where \sphinxcode{\sphinxupquote{A}}, \sphinxcode{\sphinxupquote{B}}, \sphinxcode{\sphinxupquote{C}} … could be arbitrarily complex
PSyIR expressions and the \sphinxcode{\sphinxupquote{...}} before and after \sphinxcode{\sphinxupquote{{[}MIN or
MAX{]}(A, B, C ...)}} can be arbitrary PSyIR code.

\sphinxAtStartPar
This transformation requires the IntrinsicCall node to be a
descendent of an assignment and will raise an exception if
this is not the case.
\begin{quote}\begin{description}
\sphinxlineitem{Parameters}\begin{itemize}
\item {} 
\sphinxAtStartPar
\sphinxstyleliteralstrong{\sphinxupquote{node}} (\sphinxcode{\sphinxupquote{psyclone.psyir.nodes.IntrinsicCall}}) \textendash{} a MIN or MAX intrinsic.

\item {} 
\sphinxAtStartPar
\sphinxstyleliteralstrong{\sphinxupquote{options}} (\sphinxstyleliteralemphasis{\sphinxupquote{Optional}}\sphinxstyleliteralemphasis{\sphinxupquote{{[}}}\sphinxstyleliteralemphasis{\sphinxupquote{Dict}}\sphinxstyleliteralemphasis{\sphinxupquote{{[}}}\sphinxhref{https://docs.python.org/3/library/stdtypes.html\#str}{\sphinxstyleliteralemphasis{\sphinxupquote{str}}}\sphinxstyleliteralemphasis{\sphinxupquote{, }}\sphinxstyleliteralemphasis{\sphinxupquote{Any}}\sphinxstyleliteralemphasis{\sphinxupquote{{]}}}\sphinxstyleliteralemphasis{\sphinxupquote{{]}}}) \textendash{} a dictionary with options for transformations.

\end{itemize}

\end{description}\end{quote}

\end{fulllineitems}


\end{fulllineitems}


\begin{sphinxadmonition}{warning}{Warning:}
\sphinxAtStartPar
This transformation assumes that the MAX Intrinsic acts on
PSyIR Real scalar data and does not check that this is
not the case. Once issue \#658 is on master then this
limitation can be fixed.
\end{sphinxadmonition}


\bigskip\hrule\bigskip



\begin{fulllineitems}

\pysigstartsignatures
\pysigline{\sphinxbfcode{\sphinxupquote{class\DUrole{w,w}{  }}}\sphinxcode{\sphinxupquote{psyclone.psyir.transformations.}}\sphinxbfcode{\sphinxupquote{Maxval2LoopTrans}}}
\pysigstopsignatures
\sphinxAtStartPar
Provides a transformation from a PSyIR MAXVAL IntrinsicCall node to
an equivalent PSyIR loop structure that is suitable for running in
parallel on CPUs and GPUs. Validity checks are also performed.

\sphinxAtStartPar
If MAXVAL contains a single positional argument which is an array,
the maximum value of all of the elements in the array is returned
in the the scalar R.

\begin{sphinxVerbatim}[commandchars=\\\{\}]
\PYG{n}{R}\PYG{+w}{ }\PYG{o}{=}\PYG{+w}{ }\PYG{n+nb}{MAXVAL}\PYG{p}{(}\PYG{k}{ARRAY}\PYG{p}{)}
\end{sphinxVerbatim}

\sphinxAtStartPar
For example, if the array is two dimensional, the equivalent code
for real data is:

\begin{sphinxVerbatim}[commandchars=\\\{\}]
\PYG{n}{R}\PYG{+w}{ }\PYG{o}{=}\PYG{+w}{ }\PYG{o}{\PYGZhy{}}\PYG{n+nb}{HUGE}\PYG{p}{(}\PYG{n}{R}\PYG{p}{)}
\PYG{k}{DO }\PYG{n}{J}\PYG{o}{=}\PYG{n+nb}{LBOUND}\PYG{p}{(}\PYG{k}{ARRAY}\PYG{p}{,}\PYG{l+m+mi}{2}\PYG{p}{)}\PYG{p}{,}\PYG{n+nb}{UBOUND}\PYG{p}{(}\PYG{k}{ARRAY}\PYG{p}{,}\PYG{l+m+mi}{2}\PYG{p}{)}
\PYG{+w}{  }\PYG{k}{DO }\PYG{n}{I}\PYG{o}{=}\PYG{n+nb}{LBOUND}\PYG{p}{(}\PYG{k}{ARRAY}\PYG{p}{,}\PYG{l+m+mi}{1}\PYG{p}{)}\PYG{p}{,}\PYG{n+nb}{UBOUND}\PYG{p}{(}\PYG{k}{ARRAY}\PYG{p}{,}\PYG{l+m+mi}{1}\PYG{p}{)}
\PYG{+w}{    }\PYG{n}{R}\PYG{+w}{ }\PYG{o}{=}\PYG{+w}{ }\PYG{n+nb}{MAX}\PYG{p}{(}\PYG{n}{R}\PYG{p}{,}\PYG{+w}{ }\PYG{k}{ARRAY}\PYG{p}{(}\PYG{n}{I}\PYG{p}{,}\PYG{n}{J}\PYG{p}{)}\PYG{p}{)}
\end{sphinxVerbatim}

\sphinxAtStartPar
If the mask argument is provided then the mask is used to
determine whether the maxval is applied:

\begin{sphinxVerbatim}[commandchars=\\\{\}]
\PYG{n}{R}\PYG{+w}{ }\PYG{o}{=}\PYG{+w}{ }\PYG{n+nb}{MAXVAL}\PYG{p}{(}\PYG{k}{ARRAY}\PYG{p}{,}\PYG{+w}{ }\PYG{n}{mask}\PYG{o}{=}\PYG{n+nb}{MOD}\PYG{p}{(}\PYG{k}{ARRAY}\PYG{p}{,}\PYG{+w}{ }\PYG{l+m+mf}{2.0}\PYG{p}{)}\PYG{o}{==}\PYG{l+m+mi}{1}\PYG{p}{)}
\end{sphinxVerbatim}

\sphinxAtStartPar
If the array is two dimensional, the equivalent code
for real data is:

\begin{sphinxVerbatim}[commandchars=\\\{\}]
\PYG{n}{R}\PYG{+w}{ }\PYG{o}{=}\PYG{+w}{ }\PYG{o}{\PYGZhy{}}\PYG{n+nb}{HUGE}\PYG{p}{(}\PYG{n}{R}\PYG{p}{)}
\PYG{k}{DO }\PYG{n}{J}\PYG{o}{=}\PYG{n+nb}{LBOUND}\PYG{p}{(}\PYG{k}{ARRAY}\PYG{p}{,}\PYG{l+m+mi}{2}\PYG{p}{)}\PYG{p}{,}\PYG{n+nb}{UBOUND}\PYG{p}{(}\PYG{k}{ARRAY}\PYG{p}{,}\PYG{l+m+mi}{2}\PYG{p}{)}
\PYG{+w}{  }\PYG{k}{DO }\PYG{n}{I}\PYG{o}{=}\PYG{n+nb}{LBOUND}\PYG{p}{(}\PYG{k}{ARRAY}\PYG{p}{,}\PYG{l+m+mi}{1}\PYG{p}{)}\PYG{p}{,}\PYG{n+nb}{UBOUND}\PYG{p}{(}\PYG{k}{ARRAY}\PYG{p}{,}\PYG{l+m+mi}{1}\PYG{p}{)}
\PYG{+w}{    }\PYG{k}{IF}\PYG{+w}{ }\PYG{p}{(}\PYG{n+nb}{MOD}\PYG{p}{(}\PYG{k}{ARRAY}\PYG{p}{(}\PYG{n}{I}\PYG{p}{,}\PYG{n}{J}\PYG{p}{)}\PYG{p}{,}\PYG{+w}{ }\PYG{l+m+mf}{2.0}\PYG{p}{)}\PYG{o}{==}\PYG{l+m+mi}{1}\PYG{p}{)}\PYG{+w}{ }\PYG{k}{THEN}
\PYG{k}{      }\PYG{n}{R}\PYG{+w}{ }\PYG{o}{=}\PYG{+w}{ }\PYG{n+nb}{MAX}\PYG{p}{(}\PYG{n}{R}\PYG{p}{,}\PYG{+w}{ }\PYG{k}{ARRAY}\PYG{p}{(}\PYG{n}{I}\PYG{p}{,}\PYG{n}{J}\PYG{p}{)}\PYG{p}{)}
\end{sphinxVerbatim}

\sphinxAtStartPar
The dimension argument is currently not supported and will result
in a TransformationError exception being raised.

\begin{sphinxVerbatim}[commandchars=\\\{\}]
\PYG{n}{R}\PYG{+w}{ }\PYG{o}{=}\PYG{+w}{ }\PYG{n+nb}{MAXVAL}\PYG{p}{(}\PYG{k}{ARRAY}\PYG{p}{,}\PYG{+w}{ }\PYG{k}{dimension}\PYG{o}{=}\PYG{l+m+mi}{2}\PYG{p}{)}
\end{sphinxVerbatim}

\sphinxAtStartPar
The array passed to MAXVAL may use any combination of array
syntax, array notation, array sections and scalar bounds:

\begin{sphinxVerbatim}[commandchars=\\\{\}]
\PYG{n}{R}\PYG{+w}{ }\PYG{o}{=}\PYG{+w}{ }\PYG{n+nb}{MAXVAL}\PYG{p}{(}\PYG{k}{ARRAY}\PYG{p}{)}\PYG{+w}{ }\PYG{c}{! array syntax}
\PYG{n}{R}\PYG{+w}{ }\PYG{o}{=}\PYG{+w}{ }\PYG{n+nb}{MAXVAL}\PYG{p}{(}\PYG{k}{ARRAY}\PYG{p}{(}\PYG{p}{:}\PYG{p}{,}\PYG{p}{:}\PYG{p}{)}\PYG{p}{)}\PYG{+w}{ }\PYG{c}{! array notation}
\PYG{n}{R}\PYG{+w}{ }\PYG{o}{=}\PYG{+w}{ }\PYG{n+nb}{MAXVAL}\PYG{p}{(}\PYG{k}{ARRAY}\PYG{p}{(}\PYG{l+m+mi}{1}\PYG{p}{:}\PYG{l+m+mi}{10}\PYG{p}{,}\PYG{n}{lo}\PYG{p}{:}\PYG{n}{hi}\PYG{p}{)}\PYG{p}{)}\PYG{+w}{ }\PYG{c}{! array sections}
\PYG{n}{R}\PYG{+w}{ }\PYG{o}{=}\PYG{+w}{ }\PYG{n+nb}{MAXVAL}\PYG{p}{(}\PYG{k}{ARRAY}\PYG{p}{(}\PYG{l+m+mi}{1}\PYG{p}{:}\PYG{l+m+mi}{10}\PYG{p}{,}\PYG{p}{:}\PYG{p}{)}\PYG{p}{)}\PYG{+w}{ }\PYG{c}{! mix of array section and array notation}
\PYG{n}{R}\PYG{+w}{ }\PYG{o}{=}\PYG{+w}{ }\PYG{n+nb}{MAXVAL}\PYG{p}{(}\PYG{k}{ARRAY}\PYG{p}{(}\PYG{l+m+mi}{1}\PYG{p}{:}\PYG{l+m+mi}{10}\PYG{p}{,}\PYG{l+m+mi}{2}\PYG{p}{)}\PYG{p}{)}\PYG{+w}{ }\PYG{c}{! mix of array section and scalar bound}
\end{sphinxVerbatim}

\sphinxAtStartPar
An example use of this transformation is given below:

\begin{sphinxVerbatim}[commandchars=\\\{\}]
\PYG{g+gp}{\PYGZgt{}\PYGZgt{}\PYGZgt{} }\PYG{k+kn}{from}\PYG{+w}{ }\PYG{n+nn}{psyclone}\PYG{n+nn}{.}\PYG{n+nn}{psyir}\PYG{n+nn}{.}\PYG{n+nn}{backend}\PYG{n+nn}{.}\PYG{n+nn}{fortran}\PYG{+w}{ }\PYG{k+kn}{import} \PYG{n}{FortranWriter}
\PYG{g+gp}{\PYGZgt{}\PYGZgt{}\PYGZgt{} }\PYG{k+kn}{from}\PYG{+w}{ }\PYG{n+nn}{psyclone}\PYG{n+nn}{.}\PYG{n+nn}{psyir}\PYG{n+nn}{.}\PYG{n+nn}{frontend}\PYG{n+nn}{.}\PYG{n+nn}{fortran}\PYG{+w}{ }\PYG{k+kn}{import} \PYG{n}{FortranReader}
\PYG{g+gp}{\PYGZgt{}\PYGZgt{}\PYGZgt{} }\PYG{k+kn}{from}\PYG{+w}{ }\PYG{n+nn}{psyclone}\PYG{n+nn}{.}\PYG{n+nn}{psyir}\PYG{n+nn}{.}\PYG{n+nn}{transformations}\PYG{+w}{ }\PYG{k+kn}{import} \PYG{n}{Maxval2LoopTrans}
\PYG{g+gp}{\PYGZgt{}\PYGZgt{}\PYGZgt{} }\PYG{n}{code} \PYG{o}{=} \PYG{p}{(}\PYG{l+s+s2}{\PYGZdq{}}\PYG{l+s+s2}{subroutine maxval\PYGZus{}test(array)}\PYG{l+s+se}{\PYGZbs{}n}\PYG{l+s+s2}{\PYGZdq{}}
\PYG{g+gp}{... }        \PYG{l+s+s2}{\PYGZdq{}}\PYG{l+s+s2}{  real :: array(10,10)}\PYG{l+s+se}{\PYGZbs{}n}\PYG{l+s+s2}{\PYGZdq{}}
\PYG{g+gp}{... }        \PYG{l+s+s2}{\PYGZdq{}}\PYG{l+s+s2}{  real :: result}\PYG{l+s+se}{\PYGZbs{}n}\PYG{l+s+s2}{\PYGZdq{}}
\PYG{g+gp}{... }        \PYG{l+s+s2}{\PYGZdq{}}\PYG{l+s+s2}{  result = maxval(array)}\PYG{l+s+se}{\PYGZbs{}n}\PYG{l+s+s2}{\PYGZdq{}}
\PYG{g+gp}{... }        \PYG{l+s+s2}{\PYGZdq{}}\PYG{l+s+s2}{end subroutine}\PYG{l+s+se}{\PYGZbs{}n}\PYG{l+s+s2}{\PYGZdq{}}\PYG{p}{)}
\PYG{g+gp}{\PYGZgt{}\PYGZgt{}\PYGZgt{} }\PYG{n}{psyir} \PYG{o}{=} \PYG{n}{FortranReader}\PYG{p}{(}\PYG{p}{)}\PYG{o}{.}\PYG{n}{psyir\PYGZus{}from\PYGZus{}source}\PYG{p}{(}\PYG{n}{code}\PYG{p}{)}
\PYG{g+gp}{\PYGZgt{}\PYGZgt{}\PYGZgt{} }\PYG{n}{sum\PYGZus{}node} \PYG{o}{=} \PYG{n}{psyir}\PYG{o}{.}\PYG{n}{children}\PYG{p}{[}\PYG{l+m+mi}{0}\PYG{p}{]}\PYG{o}{.}\PYG{n}{children}\PYG{p}{[}\PYG{l+m+mi}{0}\PYG{p}{]}\PYG{o}{.}\PYG{n}{children}\PYG{p}{[}\PYG{l+m+mi}{1}\PYG{p}{]}
\PYG{g+gp}{\PYGZgt{}\PYGZgt{}\PYGZgt{} }\PYG{n}{Maxval2LoopTrans}\PYG{p}{(}\PYG{p}{)}\PYG{o}{.}\PYG{n}{apply}\PYG{p}{(}\PYG{n}{sum\PYGZus{}node}\PYG{p}{)}
\PYG{g+gp}{\PYGZgt{}\PYGZgt{}\PYGZgt{} }\PYG{n+nb}{print}\PYG{p}{(}\PYG{n}{FortranWriter}\PYG{p}{(}\PYG{p}{)}\PYG{p}{(}\PYG{n}{psyir}\PYG{p}{)}\PYG{p}{)}
\PYG{g+go}{subroutine maxval\PYGZus{}test(array)}
\PYG{g+go}{  real, dimension(10,10) :: array}
\PYG{g+go}{  real :: result}
\PYG{g+go}{  integer :: idx}
\PYG{g+go}{  integer :: idx\PYGZus{}1}

\PYG{g+go}{  result = \PYGZhy{}HUGE(result)}
\PYG{g+go}{  do idx = 1, 10, 1}
\PYG{g+go}{    do idx\PYGZus{}1 = 1, 10, 1}
\PYG{g+go}{      result = MAX(result, array(idx\PYGZus{}1,idx))}
\PYG{g+go}{    enddo}
\PYG{g+go}{  enddo}

\PYG{g+go}{end subroutine maxval\PYGZus{}test}
\end{sphinxVerbatim}


\begin{fulllineitems}

\pysigstartsignatures
\pysiglinewithargsret{\sphinxbfcode{\sphinxupquote{apply}}}{\sphinxparam{\DUrole{n,n}{node}}\sphinxparamcomma \sphinxparam{\DUrole{n,n}{options}\DUrole{o,o}{=}\DUrole{default_value}{None}}}{}
\pysigstopsignatures
\sphinxAtStartPar
Apply the array\sphinxhyphen{}reduction intrinsic conversion transformation to
the specified node. This node must be one of these intrinsic
operations which is converted to an equivalent loop structure.
\begin{quote}\begin{description}
\sphinxlineitem{Parameters}\begin{itemize}
\item {} 
\sphinxAtStartPar
\sphinxstyleliteralstrong{\sphinxupquote{node}} (\sphinxcode{\sphinxupquote{psyclone.psyir.nodes.IntrinsicCall}}) \textendash{} an array\sphinxhyphen{}reduction intrinsic.

\item {} 
\sphinxAtStartPar
\sphinxstyleliteralstrong{\sphinxupquote{options}} (\sphinxstyleliteralemphasis{\sphinxupquote{Optional}}\sphinxstyleliteralemphasis{\sphinxupquote{{[}}}\sphinxstyleliteralemphasis{\sphinxupquote{Dict}}\sphinxstyleliteralemphasis{\sphinxupquote{{[}}}\sphinxhref{https://docs.python.org/3/library/stdtypes.html\#str}{\sphinxstyleliteralemphasis{\sphinxupquote{str}}}\sphinxstyleliteralemphasis{\sphinxupquote{, }}\sphinxstyleliteralemphasis{\sphinxupquote{Any}}\sphinxstyleliteralemphasis{\sphinxupquote{{]}}}\sphinxstyleliteralemphasis{\sphinxupquote{{]}}}) \textendash{} options for the transformation.

\end{itemize}

\end{description}\end{quote}

\end{fulllineitems}


\end{fulllineitems}



\bigskip\hrule\bigskip



\begin{fulllineitems}

\pysigstartsignatures
\pysigline{\sphinxbfcode{\sphinxupquote{class\DUrole{w,w}{  }}}\sphinxcode{\sphinxupquote{psyclone.psyir.transformations.}}\sphinxbfcode{\sphinxupquote{Min2CodeTrans}}}
\pysigstopsignatures
\sphinxAtStartPar
Provides a transformation from a PSyIR MIN Intrinsic node to
equivalent code in a PSyIR tree. Validity checks are also
performed (by a parent class).

\sphinxAtStartPar
The transformation replaces

\begin{sphinxVerbatim}[commandchars=\\\{\}]
\PYG{n}{R} \PYG{o}{=} \PYG{n}{MIN}\PYG{p}{(}\PYG{n}{A}\PYG{p}{,} \PYG{n}{B}\PYG{p}{,} \PYG{n}{C} \PYG{o}{.}\PYG{o}{.}\PYG{o}{.}\PYG{p}{)}
\end{sphinxVerbatim}

\sphinxAtStartPar
with the following logic:

\begin{sphinxVerbatim}[commandchars=\\\{\}]
\PYG{n}{R} \PYG{o}{=} \PYG{n}{A}
\PYG{k}{if} \PYG{n}{B} \PYG{o}{\PYGZlt{}} \PYG{n}{R}\PYG{p}{:}
    \PYG{n}{R} \PYG{o}{=} \PYG{n}{B}
\PYG{k}{if} \PYG{n}{C} \PYG{o}{\PYGZlt{}} \PYG{n}{R}\PYG{p}{:}
    \PYG{n}{R} \PYG{o}{=} \PYG{n}{C}
\PYG{o}{.}\PYG{o}{.}\PYG{o}{.}
\end{sphinxVerbatim}


\begin{fulllineitems}

\pysigstartsignatures
\pysiglinewithargsret{\sphinxbfcode{\sphinxupquote{apply}}}{\sphinxparam{\DUrole{n,n}{node}}\sphinxparamcomma \sphinxparam{\DUrole{n,n}{options}\DUrole{o,o}{=}\DUrole{default_value}{None}}}{}
\pysigstopsignatures
\sphinxAtStartPar
Apply this utility transformation to the specified node. This node
must be a MIN or MAX IntrinsicCall. The intrinsic is converted to
equivalent inline code. This is implemented as a PSyIR transform from:

\begin{sphinxVerbatim}[commandchars=\\\{\}]
\PYG{n}{R} \PYG{o}{=} \PYG{o}{.}\PYG{o}{.}\PYG{o}{.} \PYG{p}{[}\PYG{n}{MIN} \PYG{o+ow}{or} \PYG{n}{MAX}\PYG{p}{]}\PYG{p}{(}\PYG{n}{A}\PYG{p}{,} \PYG{n}{B}\PYG{p}{,} \PYG{n}{C} \PYG{o}{.}\PYG{o}{.}\PYG{o}{.}\PYG{p}{)} \PYG{o}{.}\PYG{o}{.}\PYG{o}{.}
\end{sphinxVerbatim}

\sphinxAtStartPar
to:

\begin{sphinxVerbatim}[commandchars=\\\{\}]
\PYG{n}{res} \PYG{o}{=} \PYG{n}{A}
\PYG{n}{tmp} \PYG{o}{=} \PYG{n}{B}
\PYG{n}{IF} \PYG{n}{tmp} \PYG{p}{[}\PYG{o}{\PYGZlt{}} \PYG{o+ow}{or} \PYG{o}{\PYGZgt{}}\PYG{p}{]} \PYG{n}{res}\PYG{p}{:}
    \PYG{n}{res} \PYG{o}{=} \PYG{n}{tmp}
\PYG{n}{tmp} \PYG{o}{=} \PYG{n}{C}
\PYG{n}{IF} \PYG{n}{tmp} \PYG{p}{[}\PYG{o}{\PYGZlt{}} \PYG{o+ow}{or} \PYG{o}{\PYGZgt{}}\PYG{p}{]} \PYG{n}{res}\PYG{p}{:}
    \PYG{n}{res} \PYG{o}{=} \PYG{n}{tmp}
\PYG{o}{.}\PYG{o}{.}\PYG{o}{.}
\PYG{n}{R} \PYG{o}{=} \PYG{o}{.}\PYG{o}{.}\PYG{o}{.} \PYG{n}{res} \PYG{o}{.}\PYG{o}{.}\PYG{o}{.}
\end{sphinxVerbatim}

\sphinxAtStartPar
where \sphinxcode{\sphinxupquote{A}}, \sphinxcode{\sphinxupquote{B}}, \sphinxcode{\sphinxupquote{C}} … could be arbitrarily complex
PSyIR expressions and the \sphinxcode{\sphinxupquote{...}} before and after \sphinxcode{\sphinxupquote{{[}MIN or
MAX{]}(A, B, C ...)}} can be arbitrary PSyIR code.

\sphinxAtStartPar
This transformation requires the IntrinsicCall node to be a
descendent of an assignment and will raise an exception if
this is not the case.
\begin{quote}\begin{description}
\sphinxlineitem{Parameters}\begin{itemize}
\item {} 
\sphinxAtStartPar
\sphinxstyleliteralstrong{\sphinxupquote{node}} (\sphinxcode{\sphinxupquote{psyclone.psyir.nodes.IntrinsicCall}}) \textendash{} a MIN or MAX intrinsic.

\item {} 
\sphinxAtStartPar
\sphinxstyleliteralstrong{\sphinxupquote{options}} (\sphinxstyleliteralemphasis{\sphinxupquote{Optional}}\sphinxstyleliteralemphasis{\sphinxupquote{{[}}}\sphinxstyleliteralemphasis{\sphinxupquote{Dict}}\sphinxstyleliteralemphasis{\sphinxupquote{{[}}}\sphinxhref{https://docs.python.org/3/library/stdtypes.html\#str}{\sphinxstyleliteralemphasis{\sphinxupquote{str}}}\sphinxstyleliteralemphasis{\sphinxupquote{, }}\sphinxstyleliteralemphasis{\sphinxupquote{Any}}\sphinxstyleliteralemphasis{\sphinxupquote{{]}}}\sphinxstyleliteralemphasis{\sphinxupquote{{]}}}) \textendash{} a dictionary with options for transformations.

\end{itemize}

\end{description}\end{quote}

\end{fulllineitems}


\end{fulllineitems}


\begin{sphinxadmonition}{warning}{Warning:}
\sphinxAtStartPar
This transformation assumes that the MIN Intrinsic acts on
PSyIR Real scalar data and does not check that this is
not the case. Once issue \#658 is on master then this
limitation can be fixed.
\end{sphinxadmonition}


\bigskip\hrule\bigskip



\begin{fulllineitems}

\pysigstartsignatures
\pysigline{\sphinxbfcode{\sphinxupquote{class\DUrole{w,w}{  }}}\sphinxcode{\sphinxupquote{psyclone.psyir.transformations.}}\sphinxbfcode{\sphinxupquote{Minval2LoopTrans}}}
\pysigstopsignatures
\sphinxAtStartPar
Provides a transformation from a PSyIR MINVAL IntrinsicCall node to
an equivalent PSyIR loop structure that is suitable for running in
parallel on CPUs and GPUs. Validity checks are also performed.

\sphinxAtStartPar
If MINVAL contains a single positional argument which is an array,
the minimum value of all of the elements in the array is returned
in the the scalar R.

\begin{sphinxVerbatim}[commandchars=\\\{\}]
\PYG{n}{R}\PYG{+w}{ }\PYG{o}{=}\PYG{+w}{ }\PYG{n+nb}{MINVAL}\PYG{p}{(}\PYG{k}{ARRAY}\PYG{p}{)}
\end{sphinxVerbatim}

\sphinxAtStartPar
For example, if the array is two dimensional, the equivalent code
for real data is:

\begin{sphinxVerbatim}[commandchars=\\\{\}]
\PYG{n}{R}\PYG{+w}{ }\PYG{o}{=}\PYG{+w}{ }\PYG{n+nb}{HUGE}\PYG{p}{(}\PYG{n}{R}\PYG{p}{)}
\PYG{k}{DO }\PYG{n}{J}\PYG{o}{=}\PYG{n+nb}{LBOUND}\PYG{p}{(}\PYG{k}{ARRAY}\PYG{p}{,}\PYG{l+m+mi}{2}\PYG{p}{)}\PYG{p}{,}\PYG{n+nb}{UBOUND}\PYG{p}{(}\PYG{k}{ARRAY}\PYG{p}{,}\PYG{l+m+mi}{2}\PYG{p}{)}
\PYG{+w}{  }\PYG{k}{DO }\PYG{n}{I}\PYG{o}{=}\PYG{n+nb}{LBOUND}\PYG{p}{(}\PYG{k}{ARRAY}\PYG{p}{,}\PYG{l+m+mi}{1}\PYG{p}{)}\PYG{p}{,}\PYG{n+nb}{UBOUND}\PYG{p}{(}\PYG{k}{ARRAY}\PYG{p}{,}\PYG{l+m+mi}{1}\PYG{p}{)}
\PYG{+w}{    }\PYG{n}{R}\PYG{+w}{ }\PYG{o}{=}\PYG{+w}{ }\PYG{n+nb}{MIN}\PYG{p}{(}\PYG{n}{R}\PYG{p}{,}\PYG{+w}{ }\PYG{k}{ARRAY}\PYG{p}{(}\PYG{n}{I}\PYG{p}{,}\PYG{n}{J}\PYG{p}{)}\PYG{p}{)}
\end{sphinxVerbatim}

\sphinxAtStartPar
If the mask argument is provided then the mask is used to
determine whether the minval is applied:

\begin{sphinxVerbatim}[commandchars=\\\{\}]
\PYG{n}{R}\PYG{+w}{ }\PYG{o}{=}\PYG{+w}{ }\PYG{n+nb}{MINVAL}\PYG{p}{(}\PYG{k}{ARRAY}\PYG{p}{,}\PYG{+w}{ }\PYG{n}{mask}\PYG{o}{=}\PYG{n+nb}{MOD}\PYG{p}{(}\PYG{k}{ARRAY}\PYG{p}{,}\PYG{+w}{ }\PYG{l+m+mf}{2.0}\PYG{p}{)}\PYG{o}{==}\PYG{l+m+mi}{1}\PYG{p}{)}
\end{sphinxVerbatim}

\sphinxAtStartPar
If the array is two dimensional, the equivalent code
for real data is:

\begin{sphinxVerbatim}[commandchars=\\\{\}]
\PYG{n}{R}\PYG{+w}{ }\PYG{o}{=}\PYG{+w}{ }\PYG{n+nb}{HUGE}\PYG{p}{(}\PYG{n}{R}\PYG{p}{)}
\PYG{k}{DO }\PYG{n}{J}\PYG{o}{=}\PYG{n+nb}{LBOUND}\PYG{p}{(}\PYG{k}{ARRAY}\PYG{p}{,}\PYG{l+m+mi}{2}\PYG{p}{)}\PYG{p}{,}\PYG{n+nb}{UBOUND}\PYG{p}{(}\PYG{k}{ARRAY}\PYG{p}{,}\PYG{l+m+mi}{2}\PYG{p}{)}
\PYG{+w}{  }\PYG{k}{DO }\PYG{n}{I}\PYG{o}{=}\PYG{n+nb}{LBOUND}\PYG{p}{(}\PYG{k}{ARRAY}\PYG{p}{,}\PYG{l+m+mi}{1}\PYG{p}{)}\PYG{p}{,}\PYG{n+nb}{UBOUND}\PYG{p}{(}\PYG{k}{ARRAY}\PYG{p}{,}\PYG{l+m+mi}{1}\PYG{p}{)}
\PYG{+w}{    }\PYG{k}{IF}\PYG{+w}{ }\PYG{p}{(}\PYG{n+nb}{MOD}\PYG{p}{(}\PYG{k}{ARRAY}\PYG{p}{(}\PYG{n}{I}\PYG{p}{,}\PYG{n}{J}\PYG{p}{)}\PYG{p}{,}\PYG{+w}{ }\PYG{l+m+mf}{2.0}\PYG{p}{)}\PYG{o}{==}\PYG{l+m+mi}{1}\PYG{p}{)}\PYG{+w}{ }\PYG{k}{THEN}
\PYG{k}{      }\PYG{n}{R}\PYG{+w}{ }\PYG{o}{=}\PYG{+w}{ }\PYG{n+nb}{MIN}\PYG{p}{(}\PYG{n}{R}\PYG{p}{,}\PYG{+w}{ }\PYG{k}{ARRAY}\PYG{p}{(}\PYG{n}{I}\PYG{p}{,}\PYG{n}{J}\PYG{p}{)}\PYG{p}{)}
\end{sphinxVerbatim}

\sphinxAtStartPar
The dimension argument is currently not supported and will result
in a TransformationError exception being raised.

\begin{sphinxVerbatim}[commandchars=\\\{\}]
\PYG{n}{R}\PYG{+w}{ }\PYG{o}{=}\PYG{+w}{ }\PYG{n+nb}{MINVAL}\PYG{p}{(}\PYG{k}{ARRAY}\PYG{p}{,}\PYG{+w}{ }\PYG{k}{dimension}\PYG{o}{=}\PYG{l+m+mi}{2}\PYG{p}{)}
\end{sphinxVerbatim}

\sphinxAtStartPar
The array passed to MINVAL may use any combination of array
syntax, array notation, array sections and scalar bounds:

\begin{sphinxVerbatim}[commandchars=\\\{\}]
\PYG{n}{R}\PYG{+w}{ }\PYG{o}{=}\PYG{+w}{ }\PYG{n+nb}{MINVAL}\PYG{p}{(}\PYG{k}{ARRAY}\PYG{p}{)}\PYG{+w}{ }\PYG{c}{! array syntax}
\PYG{n}{R}\PYG{+w}{ }\PYG{o}{=}\PYG{+w}{ }\PYG{n+nb}{MINVAL}\PYG{p}{(}\PYG{k}{ARRAY}\PYG{p}{(}\PYG{p}{:}\PYG{p}{,}\PYG{p}{:}\PYG{p}{)}\PYG{p}{)}\PYG{+w}{ }\PYG{c}{! array notation}
\PYG{n}{R}\PYG{+w}{ }\PYG{o}{=}\PYG{+w}{ }\PYG{n+nb}{MINVAL}\PYG{p}{(}\PYG{k}{ARRAY}\PYG{p}{(}\PYG{l+m+mi}{1}\PYG{p}{:}\PYG{l+m+mi}{10}\PYG{p}{,}\PYG{n}{lo}\PYG{p}{:}\PYG{n}{hi}\PYG{p}{)}\PYG{p}{)}\PYG{+w}{ }\PYG{c}{! array sections}
\PYG{n}{R}\PYG{+w}{ }\PYG{o}{=}\PYG{+w}{ }\PYG{n+nb}{MINVAL}\PYG{p}{(}\PYG{k}{ARRAY}\PYG{p}{(}\PYG{l+m+mi}{1}\PYG{p}{:}\PYG{l+m+mi}{10}\PYG{p}{,}\PYG{p}{:}\PYG{p}{)}\PYG{p}{)}\PYG{+w}{ }\PYG{c}{! mix of array section and array notation}
\PYG{n}{R}\PYG{+w}{ }\PYG{o}{=}\PYG{+w}{ }\PYG{n+nb}{MINVAL}\PYG{p}{(}\PYG{k}{ARRAY}\PYG{p}{(}\PYG{l+m+mi}{1}\PYG{p}{:}\PYG{l+m+mi}{10}\PYG{p}{,}\PYG{l+m+mi}{2}\PYG{p}{)}\PYG{p}{)}\PYG{+w}{ }\PYG{c}{! mix of array section and scalar bound}
\end{sphinxVerbatim}

\sphinxAtStartPar
For example:

\begin{sphinxVerbatim}[commandchars=\\\{\}]
\PYG{g+gp}{\PYGZgt{}\PYGZgt{}\PYGZgt{} }\PYG{k+kn}{from}\PYG{+w}{ }\PYG{n+nn}{psyclone}\PYG{n+nn}{.}\PYG{n+nn}{psyir}\PYG{n+nn}{.}\PYG{n+nn}{backend}\PYG{n+nn}{.}\PYG{n+nn}{fortran}\PYG{+w}{ }\PYG{k+kn}{import} \PYG{n}{FortranWriter}
\PYG{g+gp}{\PYGZgt{}\PYGZgt{}\PYGZgt{} }\PYG{k+kn}{from}\PYG{+w}{ }\PYG{n+nn}{psyclone}\PYG{n+nn}{.}\PYG{n+nn}{psyir}\PYG{n+nn}{.}\PYG{n+nn}{frontend}\PYG{n+nn}{.}\PYG{n+nn}{fortran}\PYG{+w}{ }\PYG{k+kn}{import} \PYG{n}{FortranReader}
\PYG{g+gp}{\PYGZgt{}\PYGZgt{}\PYGZgt{} }\PYG{k+kn}{from}\PYG{+w}{ }\PYG{n+nn}{psyclone}\PYG{n+nn}{.}\PYG{n+nn}{psyir}\PYG{n+nn}{.}\PYG{n+nn}{transformations}\PYG{+w}{ }\PYG{k+kn}{import} \PYG{n}{Minval2LoopTrans}
\PYG{g+gp}{\PYGZgt{}\PYGZgt{}\PYGZgt{} }\PYG{n}{code} \PYG{o}{=} \PYG{p}{(}\PYG{l+s+s2}{\PYGZdq{}}\PYG{l+s+s2}{subroutine minval\PYGZus{}test(array)}\PYG{l+s+se}{\PYGZbs{}n}\PYG{l+s+s2}{\PYGZdq{}}
\PYG{g+gp}{... }        \PYG{l+s+s2}{\PYGZdq{}}\PYG{l+s+s2}{  real :: array(10,10)}\PYG{l+s+se}{\PYGZbs{}n}\PYG{l+s+s2}{\PYGZdq{}}
\PYG{g+gp}{... }        \PYG{l+s+s2}{\PYGZdq{}}\PYG{l+s+s2}{  real :: result}\PYG{l+s+se}{\PYGZbs{}n}\PYG{l+s+s2}{\PYGZdq{}}
\PYG{g+gp}{... }        \PYG{l+s+s2}{\PYGZdq{}}\PYG{l+s+s2}{  result = minval(array)}\PYG{l+s+se}{\PYGZbs{}n}\PYG{l+s+s2}{\PYGZdq{}}
\PYG{g+gp}{... }        \PYG{l+s+s2}{\PYGZdq{}}\PYG{l+s+s2}{end subroutine}\PYG{l+s+se}{\PYGZbs{}n}\PYG{l+s+s2}{\PYGZdq{}}\PYG{p}{)}
\PYG{g+gp}{\PYGZgt{}\PYGZgt{}\PYGZgt{} }\PYG{n}{psyir} \PYG{o}{=} \PYG{n}{FortranReader}\PYG{p}{(}\PYG{p}{)}\PYG{o}{.}\PYG{n}{psyir\PYGZus{}from\PYGZus{}source}\PYG{p}{(}\PYG{n}{code}\PYG{p}{)}
\PYG{g+gp}{\PYGZgt{}\PYGZgt{}\PYGZgt{} }\PYG{n}{sum\PYGZus{}node} \PYG{o}{=} \PYG{n}{psyir}\PYG{o}{.}\PYG{n}{children}\PYG{p}{[}\PYG{l+m+mi}{0}\PYG{p}{]}\PYG{o}{.}\PYG{n}{children}\PYG{p}{[}\PYG{l+m+mi}{0}\PYG{p}{]}\PYG{o}{.}\PYG{n}{children}\PYG{p}{[}\PYG{l+m+mi}{1}\PYG{p}{]}
\PYG{g+gp}{\PYGZgt{}\PYGZgt{}\PYGZgt{} }\PYG{n}{Minval2LoopTrans}\PYG{p}{(}\PYG{p}{)}\PYG{o}{.}\PYG{n}{apply}\PYG{p}{(}\PYG{n}{sum\PYGZus{}node}\PYG{p}{)}
\PYG{g+gp}{\PYGZgt{}\PYGZgt{}\PYGZgt{} }\PYG{n+nb}{print}\PYG{p}{(}\PYG{n}{FortranWriter}\PYG{p}{(}\PYG{p}{)}\PYG{p}{(}\PYG{n}{psyir}\PYG{p}{)}\PYG{p}{)}
\PYG{g+go}{subroutine minval\PYGZus{}test(array)}
\PYG{g+go}{  real, dimension(10,10) :: array}
\PYG{g+go}{  real :: result}
\PYG{g+go}{  integer :: idx}
\PYG{g+go}{  integer :: idx\PYGZus{}1}

\PYG{g+go}{  result = HUGE(result)}
\PYG{g+go}{  do idx = 1, 10, 1}
\PYG{g+go}{    do idx\PYGZus{}1 = 1, 10, 1}
\PYG{g+go}{      result = MIN(result, array(idx\PYGZus{}1,idx))}
\PYG{g+go}{    enddo}
\PYG{g+go}{  enddo}

\PYG{g+go}{end subroutine minval\PYGZus{}test}
\end{sphinxVerbatim}


\begin{fulllineitems}

\pysigstartsignatures
\pysiglinewithargsret{\sphinxbfcode{\sphinxupquote{apply}}}{\sphinxparam{\DUrole{n,n}{node}}\sphinxparamcomma \sphinxparam{\DUrole{n,n}{options}\DUrole{o,o}{=}\DUrole{default_value}{None}}}{}
\pysigstopsignatures
\sphinxAtStartPar
Apply the array\sphinxhyphen{}reduction intrinsic conversion transformation to
the specified node. This node must be one of these intrinsic
operations which is converted to an equivalent loop structure.
\begin{quote}\begin{description}
\sphinxlineitem{Parameters}\begin{itemize}
\item {} 
\sphinxAtStartPar
\sphinxstyleliteralstrong{\sphinxupquote{node}} (\sphinxcode{\sphinxupquote{psyclone.psyir.nodes.IntrinsicCall}}) \textendash{} an array\sphinxhyphen{}reduction intrinsic.

\item {} 
\sphinxAtStartPar
\sphinxstyleliteralstrong{\sphinxupquote{options}} (\sphinxstyleliteralemphasis{\sphinxupquote{Optional}}\sphinxstyleliteralemphasis{\sphinxupquote{{[}}}\sphinxstyleliteralemphasis{\sphinxupquote{Dict}}\sphinxstyleliteralemphasis{\sphinxupquote{{[}}}\sphinxhref{https://docs.python.org/3/library/stdtypes.html\#str}{\sphinxstyleliteralemphasis{\sphinxupquote{str}}}\sphinxstyleliteralemphasis{\sphinxupquote{, }}\sphinxstyleliteralemphasis{\sphinxupquote{Any}}\sphinxstyleliteralemphasis{\sphinxupquote{{]}}}\sphinxstyleliteralemphasis{\sphinxupquote{{]}}}) \textendash{} options for the transformation.

\end{itemize}

\end{description}\end{quote}

\end{fulllineitems}


\end{fulllineitems}



\bigskip\hrule\bigskip

\phantomsection\label{\detokenize{transformations:sec-move-trans}}

\begin{fulllineitems}

\pysigstartsignatures
\pysigline{\sphinxbfcode{\sphinxupquote{class\DUrole{w,w}{  }}}\sphinxcode{\sphinxupquote{psyclone.transformations.}}\sphinxbfcode{\sphinxupquote{MoveTrans}}}
\pysigstopsignatures
\sphinxAtStartPar
Provides a transformation to move a node in the tree. For
example:

\begin{sphinxVerbatim}[commandchars=\\\{\}]
\PYG{g+gp}{\PYGZgt{}\PYGZgt{}\PYGZgt{} }\PYG{k+kn}{from}\PYG{+w}{ }\PYG{n+nn}{psyclone}\PYG{n+nn}{.}\PYG{n+nn}{parse}\PYG{n+nn}{.}\PYG{n+nn}{algorithm}\PYG{+w}{ }\PYG{k+kn}{import} \PYG{n}{parse}
\PYG{g+gp}{\PYGZgt{}\PYGZgt{}\PYGZgt{} }\PYG{k+kn}{from}\PYG{+w}{ }\PYG{n+nn}{psyclone}\PYG{n+nn}{.}\PYG{n+nn}{psyGen}\PYG{+w}{ }\PYG{k+kn}{import} \PYG{n}{PSyFactory}
\PYG{g+gp}{\PYGZgt{}\PYGZgt{}\PYGZgt{} }\PYG{n}{ast}\PYG{p}{,}\PYG{n}{invokeInfo}\PYG{o}{=}\PYG{n}{parse}\PYG{p}{(}\PYG{l+s+s2}{\PYGZdq{}}\PYG{l+s+s2}{dynamo.F90}\PYG{l+s+s2}{\PYGZdq{}}\PYG{p}{)}
\PYG{g+gp}{\PYGZgt{}\PYGZgt{}\PYGZgt{} }\PYG{n}{psy}\PYG{o}{=}\PYG{n}{PSyFactory}\PYG{p}{(}\PYG{l+s+s2}{\PYGZdq{}}\PYG{l+s+s2}{lfric}\PYG{l+s+s2}{\PYGZdq{}}\PYG{p}{)}\PYG{o}{.}\PYG{n}{create}\PYG{p}{(}\PYG{n}{invokeInfo}\PYG{p}{)}
\PYG{g+gp}{\PYGZgt{}\PYGZgt{}\PYGZgt{} }\PYG{n}{schedule}\PYG{o}{=}\PYG{n}{psy}\PYG{o}{.}\PYG{n}{invokes}\PYG{o}{.}\PYG{n}{get}\PYG{p}{(}\PYG{l+s+s1}{\PYGZsq{}}\PYG{l+s+s1}{invoke\PYGZus{}v3\PYGZus{}kernel\PYGZus{}type}\PYG{l+s+s1}{\PYGZsq{}}\PYG{p}{)}\PYG{o}{.}\PYG{n}{schedule}
\PYG{g+gp}{\PYGZgt{}\PYGZgt{}\PYGZgt{} }\PYG{c+c1}{\PYGZsh{} Uncomment the following line to see a text view of the schedule}
\PYG{g+gp}{\PYGZgt{}\PYGZgt{}\PYGZgt{} }\PYG{c+c1}{\PYGZsh{} print(schedule.view())}
\PYG{g+gp}{\PYGZgt{}\PYGZgt{}\PYGZgt{}}
\PYG{g+gp}{\PYGZgt{}\PYGZgt{}\PYGZgt{} }\PYG{k+kn}{from}\PYG{+w}{ }\PYG{n+nn}{psyclone}\PYG{n+nn}{.}\PYG{n+nn}{transformations}\PYG{+w}{ }\PYG{k+kn}{import} \PYG{n}{MoveTrans}
\PYG{g+gp}{\PYGZgt{}\PYGZgt{}\PYGZgt{} }\PYG{n}{trans}\PYG{o}{=}\PYG{n}{MoveTrans}\PYG{p}{(}\PYG{p}{)}
\PYG{g+gp}{\PYGZgt{}\PYGZgt{}\PYGZgt{} }\PYG{n}{trans}\PYG{o}{.}\PYG{n}{apply}\PYG{p}{(}\PYG{n}{schedule}\PYG{o}{.}\PYG{n}{children}\PYG{p}{[}\PYG{l+m+mi}{0}\PYG{p}{]}\PYG{p}{,} \PYG{n}{schedule}\PYG{o}{.}\PYG{n}{children}\PYG{p}{[}\PYG{l+m+mi}{2}\PYG{p}{]}\PYG{p}{,}
\PYG{g+gp}{... }            \PYG{n}{options} \PYG{o}{=} \PYG{p}{\PYGZob{}}\PYG{l+s+s2}{\PYGZdq{}}\PYG{l+s+s2}{position}\PYG{l+s+s2}{\PYGZdq{}}\PYG{p}{:}\PYG{l+s+s2}{\PYGZdq{}}\PYG{l+s+s2}{after}\PYG{l+s+s2}{\PYGZdq{}}\PYG{p}{)}
\PYG{g+gp}{\PYGZgt{}\PYGZgt{}\PYGZgt{} }\PYG{c+c1}{\PYGZsh{} Uncomment the following line to see a text view of the schedule}
\PYG{g+gp}{\PYGZgt{}\PYGZgt{}\PYGZgt{} }\PYG{c+c1}{\PYGZsh{} print(schedule.view())}
\end{sphinxVerbatim}

\sphinxAtStartPar
Nodes may only be moved to a new location with the same parent
and must not break any dependencies otherwise an exception is
raised.


\begin{fulllineitems}

\pysigstartsignatures
\pysiglinewithargsret{\sphinxbfcode{\sphinxupquote{apply}}}{\sphinxparam{\DUrole{n,n}{node}}\sphinxparamcomma \sphinxparam{\DUrole{n,n}{location}}\sphinxparamcomma \sphinxparam{\DUrole{n,n}{options}\DUrole{o,o}{=}\DUrole{default_value}{None}}}{}
\pysigstopsignatures
\sphinxAtStartPar
Move the node represented by \sphinxcode{\sphinxupquote{node}} before location
\sphinxcode{\sphinxupquote{location}} (which is also a node) by default and after
if the optional \sphinxtitleref{position} argument is set to ‘after’.
\begin{quote}\begin{description}
\sphinxlineitem{Parameters}\begin{itemize}
\item {} 
\sphinxAtStartPar
\sphinxstyleliteralstrong{\sphinxupquote{node}} (\sphinxcode{\sphinxupquote{psyclone.psyir.nodes.Node}}) \textendash{} the node to be moved.

\item {} 
\sphinxAtStartPar
\sphinxstyleliteralstrong{\sphinxupquote{location}} (\sphinxcode{\sphinxupquote{psyclone.psyir.nodes.Node}}) \textendash{} node before or after which the given node            should be moved.

\item {} 
\sphinxAtStartPar
\sphinxstyleliteralstrong{\sphinxupquote{options}} (\sphinxstyleliteralemphasis{\sphinxupquote{Optional}}\sphinxstyleliteralemphasis{\sphinxupquote{{[}}}\sphinxstyleliteralemphasis{\sphinxupquote{Dict}}\sphinxstyleliteralemphasis{\sphinxupquote{{[}}}\sphinxhref{https://docs.python.org/3/library/stdtypes.html\#str}{\sphinxstyleliteralemphasis{\sphinxupquote{str}}}\sphinxstyleliteralemphasis{\sphinxupquote{, }}\sphinxstyleliteralemphasis{\sphinxupquote{Any}}\sphinxstyleliteralemphasis{\sphinxupquote{{]}}}\sphinxstyleliteralemphasis{\sphinxupquote{{]}}}) \textendash{} a dictionary with options for transformations.

\item {} 
\sphinxAtStartPar
\sphinxstyleliteralstrong{\sphinxupquote{options}}\sphinxstyleliteralstrong{\sphinxupquote{{[}}}\sphinxstyleliteralstrong{\sphinxupquote{"position"}}\sphinxstyleliteralstrong{\sphinxupquote{{]}}} (\sphinxhref{https://docs.python.org/3/library/stdtypes.html\#str}{\sphinxstyleliteralemphasis{\sphinxupquote{str}}}) \textendash{} either ‘before’ or ‘after’.

\end{itemize}

\sphinxlineitem{Raises}\begin{itemize}
\item {} 
\sphinxAtStartPar
\sphinxstyleliteralstrong{\sphinxupquote{TransformationError}} \textendash{} if the given node is not an instance             of \sphinxcode{\sphinxupquote{psyclone.psyir.nodes.Node}}

\item {} 
\sphinxAtStartPar
\sphinxstyleliteralstrong{\sphinxupquote{TransformationError}} \textendash{} if the location is not valid.

\end{itemize}

\end{description}\end{quote}

\end{fulllineitems}


\end{fulllineitems}



\bigskip\hrule\bigskip



\begin{fulllineitems}

\pysigstartsignatures
\pysigline{\sphinxbfcode{\sphinxupquote{class\DUrole{w,w}{  }}}\sphinxcode{\sphinxupquote{psyclone.domain.gocean.transformations.}}\sphinxbfcode{\sphinxupquote{GOOpenCLTrans}}}
\pysigstopsignatures
\sphinxAtStartPar
Switches on/off the generation of an OpenCL PSy layer for a given
InvokeSchedule. Additionally, it will generate OpenCL kernels for
each of the kernels referenced by the Invoke. For example:

\begin{sphinxVerbatim}[commandchars=\\\{\}]
\PYG{g+gp}{\PYGZgt{}\PYGZgt{}\PYGZgt{} }\PYG{k+kn}{from}\PYG{+w}{ }\PYG{n+nn}{psyclone}\PYG{n+nn}{.}\PYG{n+nn}{parse}\PYG{n+nn}{.}\PYG{n+nn}{algorithm}\PYG{+w}{ }\PYG{k+kn}{import} \PYG{n}{parse}
\PYG{g+gp}{\PYGZgt{}\PYGZgt{}\PYGZgt{} }\PYG{k+kn}{from}\PYG{+w}{ }\PYG{n+nn}{psyclone}\PYG{n+nn}{.}\PYG{n+nn}{psyGen}\PYG{+w}{ }\PYG{k+kn}{import} \PYG{n}{PSyFactory}
\PYG{g+gp}{\PYGZgt{}\PYGZgt{}\PYGZgt{} }\PYG{n}{API} \PYG{o}{=} \PYG{l+s+s2}{\PYGZdq{}}\PYG{l+s+s2}{gocean}\PYG{l+s+s2}{\PYGZdq{}}
\PYG{g+gp}{\PYGZgt{}\PYGZgt{}\PYGZgt{} }\PYG{n}{FILENAME} \PYG{o}{=} \PYG{l+s+s2}{\PYGZdq{}}\PYG{l+s+s2}{shallow\PYGZus{}alg.f90}\PYG{l+s+s2}{\PYGZdq{}} \PYG{c+c1}{\PYGZsh{} examples/gocean/eg1}
\PYG{g+gp}{\PYGZgt{}\PYGZgt{}\PYGZgt{} }\PYG{n}{ast}\PYG{p}{,} \PYG{n}{invoke\PYGZus{}info} \PYG{o}{=} \PYG{n}{parse}\PYG{p}{(}\PYG{n}{FILENAME}\PYG{p}{,} \PYG{n}{api}\PYG{o}{=}\PYG{n}{API}\PYG{p}{)}
\PYG{g+gp}{\PYGZgt{}\PYGZgt{}\PYGZgt{} }\PYG{n}{psy} \PYG{o}{=} \PYG{n}{PSyFactory}\PYG{p}{(}\PYG{n}{API}\PYG{p}{,} \PYG{n}{distributed\PYGZus{}memory}\PYG{o}{=}\PYG{k+kc}{False}\PYG{p}{)}\PYG{o}{.}\PYG{n}{create}\PYG{p}{(}\PYG{n}{invoke\PYGZus{}info}\PYG{p}{)}
\PYG{g+gp}{\PYGZgt{}\PYGZgt{}\PYGZgt{} }\PYG{n}{schedule} \PYG{o}{=} \PYG{n}{psy}\PYG{o}{.}\PYG{n}{invokes}\PYG{o}{.}\PYG{n}{get}\PYG{p}{(}\PYG{l+s+s1}{\PYGZsq{}}\PYG{l+s+s1}{invoke\PYGZus{}0}\PYG{l+s+s1}{\PYGZsq{}}\PYG{p}{)}\PYG{o}{.}\PYG{n}{schedule}
\PYG{g+gp}{\PYGZgt{}\PYGZgt{}\PYGZgt{} }\PYG{n}{ocl\PYGZus{}trans} \PYG{o}{=} \PYG{n}{GOOpenCLTrans}\PYG{p}{(}\PYG{p}{)}
\PYG{g+gp}{\PYGZgt{}\PYGZgt{}\PYGZgt{} }\PYG{n}{ocl\PYGZus{}trans}\PYG{o}{.}\PYG{n}{apply}\PYG{p}{(}\PYG{n}{schedule}\PYG{p}{)}
\PYG{g+gp}{\PYGZgt{}\PYGZgt{}\PYGZgt{} }\PYG{n+nb}{print}\PYG{p}{(}\PYG{n}{schedule}\PYG{o}{.}\PYG{n}{view}\PYG{p}{(}\PYG{p}{)}\PYG{p}{)}
\end{sphinxVerbatim}


\begin{fulllineitems}

\pysigstartsignatures
\pysiglinewithargsret{\sphinxbfcode{\sphinxupquote{apply}}}{\sphinxparam{\DUrole{n,n}{node}}\sphinxparamcomma \sphinxparam{\DUrole{n,n}{options}\DUrole{o,o}{=}\DUrole{default_value}{None}}}{}
\pysigstopsignatures
\sphinxAtStartPar
Apply the OpenCL transformation to the supplied GOInvokeSchedule. This
causes PSyclone to generate an OpenCL version of the corresponding
PSy\sphinxhyphen{}layer routine. The generated code makes use of the FortCL
library (\sphinxurl{https://github.com/stfc/FortCL}) in order to manage the
OpenCL device directly from Fortran.
\begin{quote}\begin{description}
\sphinxlineitem{Parameters}\begin{itemize}
\item {} 
\sphinxAtStartPar
\sphinxstyleliteralstrong{\sphinxupquote{node}} (\sphinxcode{\sphinxupquote{psyclone.psyGen.GOInvokeSchedule}}) \textendash{} the InvokeSchedule to transform.

\item {} 
\sphinxAtStartPar
\sphinxstyleliteralstrong{\sphinxupquote{options}} (\sphinxhref{https://docs.python.org/3/library/stdtypes.html\#dict}{\sphinxstyleliteralemphasis{\sphinxupquote{dict}}}\sphinxstyleliteralemphasis{\sphinxupquote{ of }}\sphinxstyleliteralemphasis{\sphinxupquote{str:value}}\sphinxstyleliteralemphasis{\sphinxupquote{ or }}\sphinxstyleliteralemphasis{\sphinxupquote{None}}) \textendash{} set of option to tune the OpenCL generation.

\item {} 
\sphinxAtStartPar
\sphinxstyleliteralstrong{\sphinxupquote{options}}\sphinxstyleliteralstrong{\sphinxupquote{{[}}}\sphinxstyleliteralstrong{\sphinxupquote{"enable\_profiling"}}\sphinxstyleliteralstrong{\sphinxupquote{{]}}} (\sphinxhref{https://docs.python.org/3/library/functions.html\#bool}{\sphinxstyleliteralemphasis{\sphinxupquote{bool}}}) \textendash{} whether or not to set up the                 OpenCL environment with the profiling option enabled.

\item {} 
\sphinxAtStartPar
\sphinxstyleliteralstrong{\sphinxupquote{options}}\sphinxstyleliteralstrong{\sphinxupquote{{[}}}\sphinxstyleliteralstrong{\sphinxupquote{"out\_of\_order"}}\sphinxstyleliteralstrong{\sphinxupquote{{]}}} (\sphinxhref{https://docs.python.org/3/library/functions.html\#bool}{\sphinxstyleliteralemphasis{\sphinxupquote{bool}}}) \textendash{} whether or not to set up the                 OpenCL environment with the out\_of\_order option enabled.

\item {} 
\sphinxAtStartPar
\sphinxstyleliteralstrong{\sphinxupquote{options}}\sphinxstyleliteralstrong{\sphinxupquote{{[}}}\sphinxstyleliteralstrong{\sphinxupquote{"end\_barrier"}}\sphinxstyleliteralstrong{\sphinxupquote{{]}}} (\sphinxhref{https://docs.python.org/3/library/functions.html\#bool}{\sphinxstyleliteralemphasis{\sphinxupquote{bool}}}) \textendash{} whether or not to add an OpenCL                 barrier at the end of the transformed invoke.

\end{itemize}

\end{description}\end{quote}

\end{fulllineitems}


\end{fulllineitems}



\bigskip\hrule\bigskip



\begin{fulllineitems}

\pysigstartsignatures
\pysigline{\sphinxbfcode{\sphinxupquote{class\DUrole{w,w}{  }}}\sphinxcode{\sphinxupquote{psyclone.transformations.}}\sphinxbfcode{\sphinxupquote{OMPDeclareTargetTrans}}}
\pysigstopsignatures
\sphinxAtStartPar
Adds an OpenMP declare target directive to the specified routine.

\sphinxAtStartPar
For example:

\begin{sphinxVerbatim}[commandchars=\\\{\}]
\PYG{g+gp}{\PYGZgt{}\PYGZgt{}\PYGZgt{} }\PYG{k+kn}{from}\PYG{+w}{ }\PYG{n+nn}{psyclone}\PYG{n+nn}{.}\PYG{n+nn}{psyir}\PYG{n+nn}{.}\PYG{n+nn}{frontend}\PYG{n+nn}{.}\PYG{n+nn}{fortran}\PYG{+w}{ }\PYG{k+kn}{import} \PYG{n}{FortranReader}
\PYG{g+gp}{\PYGZgt{}\PYGZgt{}\PYGZgt{} }\PYG{k+kn}{from}\PYG{+w}{ }\PYG{n+nn}{psyclone}\PYG{n+nn}{.}\PYG{n+nn}{psyir}\PYG{n+nn}{.}\PYG{n+nn}{nodes}\PYG{+w}{ }\PYG{k+kn}{import} \PYG{n}{Loop}
\PYG{g+gp}{\PYGZgt{}\PYGZgt{}\PYGZgt{} }\PYG{k+kn}{from}\PYG{+w}{ }\PYG{n+nn}{psyclone}\PYG{n+nn}{.}\PYG{n+nn}{transformations}\PYG{+w}{ }\PYG{k+kn}{import} \PYG{n}{OMPDeclareTargetTrans}
\PYG{g+gp}{\PYGZgt{}\PYGZgt{}\PYGZgt{}}
\PYG{g+gp}{\PYGZgt{}\PYGZgt{}\PYGZgt{} }\PYG{n}{tree} \PYG{o}{=} \PYG{n}{FortranReader}\PYG{p}{(}\PYG{p}{)}\PYG{o}{.}\PYG{n}{psyir\PYGZus{}from\PYGZus{}source}\PYG{p}{(}\PYG{l+s+s2}{\PYGZdq{}\PYGZdq{}\PYGZdq{}}
\PYG{g+gp}{... }\PYG{l+s+s2}{    subroutine my\PYGZus{}subroutine(A)}
\PYG{g+gp}{... }\PYG{l+s+s2}{        integer, dimension(10, 10), intent(inout) :: A}
\PYG{g+gp}{... }\PYG{l+s+s2}{        integer :: i}
\PYG{g+gp}{... }\PYG{l+s+s2}{        integer :: j}
\PYG{g+gp}{... }\PYG{l+s+s2}{        do i = 1, 10}
\PYG{g+gp}{... }\PYG{l+s+s2}{            do j = 1, 10}
\PYG{g+gp}{... }\PYG{l+s+s2}{                A(i, j) = 0}
\PYG{g+gp}{... }\PYG{l+s+s2}{            end do}
\PYG{g+gp}{... }\PYG{l+s+s2}{        end do}
\PYG{g+gp}{... }\PYG{l+s+s2}{    end subroutine}
\PYG{g+gp}{... }\PYG{l+s+s2}{    }\PYG{l+s+s2}{\PYGZdq{}\PYGZdq{}\PYGZdq{}}
\PYG{g+gp}{\PYGZgt{}\PYGZgt{}\PYGZgt{} }\PYG{n}{omptargettrans} \PYG{o}{=} \PYG{n}{OMPDeclareTargetTrans}\PYG{p}{(}\PYG{p}{)}
\PYG{g+gp}{\PYGZgt{}\PYGZgt{}\PYGZgt{} }\PYG{n}{omptargettrans}\PYG{o}{.}\PYG{n}{apply}\PYG{p}{(}\PYG{n}{tree}\PYG{o}{.}\PYG{n}{walk}\PYG{p}{(}\PYG{n}{Routine}\PYG{p}{)}\PYG{p}{[}\PYG{l+m+mi}{0}\PYG{p}{]}\PYG{p}{)}
\end{sphinxVerbatim}

\sphinxAtStartPar
will generate:

\begin{sphinxVerbatim}[commandchars=\\\{\}]
\PYG{k}{subroutine }\PYG{n}{my\PYGZus{}subroutine}\PYG{p}{(}\PYG{n}{A}\PYG{p}{)}
\PYG{+w}{    }\PYG{k+kt}{integer}\PYG{p}{,}\PYG{+w}{ }\PYG{k}{dimension}\PYG{p}{(}\PYG{l+m+mi}{10}\PYG{p}{,}\PYG{+w}{ }\PYG{l+m+mi}{10}\PYG{p}{)}\PYG{p}{,}\PYG{+w}{ }\PYG{k}{intent}\PYG{p}{(}\PYG{n}{inout}\PYG{p}{)}\PYG{+w}{ }\PYG{k+kd}{::}\PYG{+w}{ }\PYG{n}{A}
\PYG{+w}{    }\PYG{k+kt}{integer}\PYG{+w}{ }\PYG{k+kd}{::}\PYG{+w}{ }\PYG{n}{i}
\PYG{+w}{    }\PYG{k+kt}{integer}\PYG{+w}{ }\PYG{k+kd}{::}\PYG{+w}{ }\PYG{n}{j}
\PYG{+w}{    }\PYG{c}{!\PYGZdl{}omp declare target}
\PYG{+w}{    }\PYG{k}{do }\PYG{n}{i}\PYG{+w}{ }\PYG{o}{=}\PYG{+w}{ }\PYG{l+m+mi}{1}\PYG{p}{,}\PYG{+w}{ }\PYG{l+m+mi}{10}
\PYG{+w}{        }\PYG{k}{do }\PYG{n}{j}\PYG{+w}{ }\PYG{o}{=}\PYG{+w}{ }\PYG{l+m+mi}{1}\PYG{p}{,}\PYG{+w}{ }\PYG{l+m+mi}{10}
\PYG{+w}{            }\PYG{n}{A}\PYG{p}{(}\PYG{n}{i}\PYG{p}{,}\PYG{+w}{ }\PYG{n}{j}\PYG{p}{)}\PYG{+w}{ }\PYG{o}{=}\PYG{+w}{ }\PYG{l+m+mi}{0}
\PYG{+w}{        }\PYG{k}{end }\PYG{k}{do}
\PYG{k}{    }\PYG{k}{end }\PYG{k}{do}
\PYG{k}{end }\PYG{k}{subroutine}
\end{sphinxVerbatim}


\begin{fulllineitems}

\pysigstartsignatures
\pysiglinewithargsret{\sphinxbfcode{\sphinxupquote{apply}}}{\sphinxparam{\DUrole{n,n}{node}}\sphinxparamcomma \sphinxparam{\DUrole{n,n}{options}\DUrole{o,o}{=}\DUrole{default_value}{None}}}{}
\pysigstopsignatures
\sphinxAtStartPar
Insert an OMPDeclareTargetDirective inside the provided routine or
associated PSyKAl kernel.
\begin{quote}\begin{description}
\sphinxlineitem{Parameters}\begin{itemize}
\item {} 
\sphinxAtStartPar
\sphinxstyleliteralstrong{\sphinxupquote{node}} (\sphinxcode{\sphinxupquote{psyclone.psyir.nodes.Routine}} |
\sphinxcode{\sphinxupquote{psyclone.psyGen.Kern}}) \textendash{} the kernel or routine which is the target of this
transformation.

\item {} 
\sphinxAtStartPar
\sphinxstyleliteralstrong{\sphinxupquote{options}} (\sphinxstyleliteralemphasis{\sphinxupquote{Optional}}\sphinxstyleliteralemphasis{\sphinxupquote{{[}}}\sphinxstyleliteralemphasis{\sphinxupquote{Dict}}\sphinxstyleliteralemphasis{\sphinxupquote{{[}}}\sphinxhref{https://docs.python.org/3/library/stdtypes.html\#str}{\sphinxstyleliteralemphasis{\sphinxupquote{str}}}\sphinxstyleliteralemphasis{\sphinxupquote{, }}\sphinxstyleliteralemphasis{\sphinxupquote{Any}}\sphinxstyleliteralemphasis{\sphinxupquote{{]}}}\sphinxstyleliteralemphasis{\sphinxupquote{{]}}}) \textendash{} a dictionary with options for transformations.

\item {} 
\sphinxAtStartPar
\sphinxstyleliteralstrong{\sphinxupquote{options}}\sphinxstyleliteralstrong{\sphinxupquote{{[}}}\sphinxstyleliteralstrong{\sphinxupquote{"force"}}\sphinxstyleliteralstrong{\sphinxupquote{{]}}} (\sphinxhref{https://docs.python.org/3/library/functions.html\#bool}{\sphinxstyleliteralemphasis{\sphinxupquote{bool}}}) \textendash{} whether to allow routines with
CodeBlocks to run on the GPU.

\end{itemize}

\end{description}\end{quote}

\end{fulllineitems}


\end{fulllineitems}



\bigskip\hrule\bigskip



\begin{fulllineitems}

\pysigstartsignatures
\pysiglinewithargsret{\sphinxbfcode{\sphinxupquote{class\DUrole{w,w}{  }}}\sphinxcode{\sphinxupquote{psyclone.psyir.transformations.}}\sphinxbfcode{\sphinxupquote{OMPLoopTrans}}}{\sphinxparam{\DUrole{n,n}{omp\_directive}\DUrole{o,o}{=}\DUrole{default_value}{\textquotesingle{}do\textquotesingle{}}}\sphinxparamcomma \sphinxparam{\DUrole{n,n}{omp\_schedule}\DUrole{o,o}{=}\DUrole{default_value}{\textquotesingle{}auto\textquotesingle{}}}}{}
\pysigstopsignatures
\sphinxAtStartPar
Adds an OpenMP directive to parallelise this loop. It can insert different
directives such as “omp do/for”, “omp parallel do/for”, “omp teams
distribute parallel do/for” or “omp loop” depending on the provided
parameters.
The OpenMP schedule to use can also be specified, but this will be ignored
in case of the “omp loop” (as the ‘schedule’ clause is not valid for this
specific directive). The configuration\sphinxhyphen{}defined ‘reprod’ parameter
also specifies whether a manual reproducible reproduction is to be used.
Note, reproducible in this case means obtaining the same results with the
same number of OpenMP threads, not for different numbers of OpenMP threads.
\begin{quote}\begin{description}
\sphinxlineitem{Parameters}\begin{itemize}
\item {} 
\sphinxAtStartPar
\sphinxstyleliteralstrong{\sphinxupquote{omp\_schedule}} (\sphinxhref{https://docs.python.org/3/library/stdtypes.html\#str}{\sphinxstyleliteralemphasis{\sphinxupquote{str}}}) \textendash{} the OpenMP schedule to use. Defaults to ‘auto’.

\item {} 
\sphinxAtStartPar
\sphinxstyleliteralstrong{\sphinxupquote{omp\_directive}} (\sphinxhref{https://docs.python.org/3/library/stdtypes.html\#str}{\sphinxstyleliteralemphasis{\sphinxupquote{str}}}) \textendash{} choose which OpenMP loop directive to use.         Defaults to “omp do”

\end{itemize}

\end{description}\end{quote}

\sphinxAtStartPar
For example:

\begin{sphinxVerbatim}[commandchars=\\\{\}]
\PYG{g+gp}{\PYGZgt{}\PYGZgt{}\PYGZgt{} }\PYG{k+kn}{from}\PYG{+w}{ }\PYG{n+nn}{psyclone}\PYG{n+nn}{.}\PYG{n+nn}{psyir}\PYG{n+nn}{.}\PYG{n+nn}{frontend}\PYG{n+nn}{.}\PYG{n+nn}{fortran}\PYG{+w}{ }\PYG{k+kn}{import} \PYG{n}{FortranReader}
\PYG{g+gp}{\PYGZgt{}\PYGZgt{}\PYGZgt{} }\PYG{k+kn}{from}\PYG{+w}{ }\PYG{n+nn}{psyclone}\PYG{n+nn}{.}\PYG{n+nn}{psyir}\PYG{n+nn}{.}\PYG{n+nn}{backend}\PYG{n+nn}{.}\PYG{n+nn}{fortran}\PYG{+w}{ }\PYG{k+kn}{import} \PYG{n}{FortranWriter}
\PYG{g+gp}{\PYGZgt{}\PYGZgt{}\PYGZgt{} }\PYG{k+kn}{from}\PYG{+w}{ }\PYG{n+nn}{psyclone}\PYG{n+nn}{.}\PYG{n+nn}{psyir}\PYG{n+nn}{.}\PYG{n+nn}{nodes}\PYG{+w}{ }\PYG{k+kn}{import} \PYG{n}{Loop}
\PYG{g+gp}{\PYGZgt{}\PYGZgt{}\PYGZgt{} }\PYG{k+kn}{from}\PYG{+w}{ }\PYG{n+nn}{psyclone}\PYG{n+nn}{.}\PYG{n+nn}{transformations}\PYG{+w}{ }\PYG{k+kn}{import} \PYG{n}{OMPLoopTrans}\PYG{p}{,} \PYG{n}{OMPParallelTrans}
\PYG{g+gp}{\PYGZgt{}\PYGZgt{}\PYGZgt{}}
\PYG{g+gp}{\PYGZgt{}\PYGZgt{}\PYGZgt{} }\PYG{n}{psyir} \PYG{o}{=} \PYG{n}{FortranReader}\PYG{p}{(}\PYG{p}{)}\PYG{o}{.}\PYG{n}{psyir\PYGZus{}from\PYGZus{}source}\PYG{p}{(}\PYG{l+s+s2}{\PYGZdq{}\PYGZdq{}\PYGZdq{}}
\PYG{g+gp}{... }\PYG{l+s+s2}{    subroutine my\PYGZus{}subroutine()}
\PYG{g+gp}{... }\PYG{l+s+s2}{        integer, dimension(10, 10) :: A}
\PYG{g+gp}{... }\PYG{l+s+s2}{        integer :: i}
\PYG{g+gp}{... }\PYG{l+s+s2}{        integer :: j}
\PYG{g+gp}{... }\PYG{l+s+s2}{        do i = 1, 10}
\PYG{g+gp}{... }\PYG{l+s+s2}{            do j = 1, 10}
\PYG{g+gp}{... }\PYG{l+s+s2}{                A(i, j) = 0}
\PYG{g+gp}{... }\PYG{l+s+s2}{            end do}
\PYG{g+gp}{... }\PYG{l+s+s2}{        end do}
\PYG{g+gp}{... }\PYG{l+s+s2}{    end subroutine}
\PYG{g+gp}{... }\PYG{l+s+s2}{    }\PYG{l+s+s2}{\PYGZdq{}\PYGZdq{}\PYGZdq{}}\PYG{p}{)}
\PYG{g+gp}{\PYGZgt{}\PYGZgt{}\PYGZgt{} }\PYG{n}{loop} \PYG{o}{=} \PYG{n}{psyir}\PYG{o}{.}\PYG{n}{walk}\PYG{p}{(}\PYG{n}{Loop}\PYG{p}{)}\PYG{p}{[}\PYG{l+m+mi}{0}\PYG{p}{]}
\PYG{g+gp}{\PYGZgt{}\PYGZgt{}\PYGZgt{} }\PYG{n}{omplooptrans1} \PYG{o}{=} \PYG{n}{OMPLoopTrans}\PYG{p}{(}\PYG{n}{omp\PYGZus{}schedule}\PYG{o}{=}\PYG{l+s+s2}{\PYGZdq{}}\PYG{l+s+s2}{dynamic}\PYG{l+s+s2}{\PYGZdq{}}\PYG{p}{,}
\PYG{g+gp}{... }                             \PYG{n}{omp\PYGZus{}directive}\PYG{o}{=}\PYG{l+s+s2}{\PYGZdq{}}\PYG{l+s+s2}{paralleldo}\PYG{l+s+s2}{\PYGZdq{}}\PYG{p}{)}
\PYG{g+gp}{\PYGZgt{}\PYGZgt{}\PYGZgt{} }\PYG{n}{omplooptrans1}\PYG{o}{.}\PYG{n}{apply}\PYG{p}{(}\PYG{n}{loop}\PYG{p}{)}
\PYG{g+gp}{\PYGZgt{}\PYGZgt{}\PYGZgt{} }\PYG{n+nb}{print}\PYG{p}{(}\PYG{n}{FortranWriter}\PYG{p}{(}\PYG{p}{)}\PYG{p}{(}\PYG{n}{psyir}\PYG{p}{)}\PYG{p}{)}
\PYG{g+go}{subroutine my\PYGZus{}subroutine()}
\PYG{g+go}{  integer, dimension(10,10) :: a}
\PYG{g+go}{  integer :: i}
\PYG{g+go}{  integer :: j}

\PYG{g+go}{  !\PYGZdl{}omp parallel do default(shared), private(i,j), schedule(dynamic)}
\PYG{g+go}{  do i = 1, 10, 1}
\PYG{g+go}{    do j = 1, 10, 1}
\PYG{g+go}{      a(i,j) = 0}
\PYG{g+go}{    enddo}
\PYG{g+go}{  enddo}
\PYG{g+go}{  !\PYGZdl{}omp end parallel do}

\PYG{g+go}{end subroutine my\PYGZus{}subroutine}
\end{sphinxVerbatim}


\begin{fulllineitems}

\pysigstartsignatures
\pysiglinewithargsret{\sphinxbfcode{\sphinxupquote{apply}}}{\sphinxparam{\DUrole{n,n}{node}}\sphinxparamcomma \sphinxparam{\DUrole{n,n}{options}\DUrole{o,o}{=}\DUrole{default_value}{None}}}{}
\pysigstopsignatures
\sphinxAtStartPar
Apply the OMPLoopTrans transformation to the specified PSyIR Loop.
\begin{quote}\begin{description}
\sphinxlineitem{Parameters}\begin{itemize}
\item {} 
\sphinxAtStartPar
\sphinxstyleliteralstrong{\sphinxupquote{node}} (\sphinxcode{\sphinxupquote{psyclone.psyir.nodes.Node}}) \textendash{} the supplied node to which we will apply the                      OMPLoopTrans transformation

\item {} 
\sphinxAtStartPar
\sphinxstyleliteralstrong{\sphinxupquote{options}} (\sphinxstyleliteralemphasis{\sphinxupquote{Optional}}\sphinxstyleliteralemphasis{\sphinxupquote{{[}}}\sphinxstyleliteralemphasis{\sphinxupquote{Dict}}\sphinxstyleliteralemphasis{\sphinxupquote{{[}}}\sphinxhref{https://docs.python.org/3/library/stdtypes.html\#str}{\sphinxstyleliteralemphasis{\sphinxupquote{str}}}\sphinxstyleliteralemphasis{\sphinxupquote{, }}\sphinxstyleliteralemphasis{\sphinxupquote{Any}}\sphinxstyleliteralemphasis{\sphinxupquote{{]}}}\sphinxstyleliteralemphasis{\sphinxupquote{{]}}}) \textendash{} a dictionary with options for transformations                        and validation.

\item {} 
\sphinxAtStartPar
\sphinxstyleliteralstrong{\sphinxupquote{options}}\sphinxstyleliteralstrong{\sphinxupquote{{[}}}\sphinxstyleliteralstrong{\sphinxupquote{"reprod"}}\sphinxstyleliteralstrong{\sphinxupquote{{]}}} (\sphinxhref{https://docs.python.org/3/library/functions.html\#bool}{\sphinxstyleliteralemphasis{\sphinxupquote{bool}}}) \textendash{} indicating whether reproducible reductions should be used.                 By default the value from the config file will be used.

\end{itemize}

\end{description}\end{quote}

\end{fulllineitems}



\begin{fulllineitems}

\pysigstartsignatures
\pysigline{\sphinxbfcode{\sphinxupquote{property\DUrole{w,w}{  }}}\sphinxbfcode{\sphinxupquote{omp\_directive}}}
\pysigstopsignatures\begin{quote}\begin{description}
\sphinxlineitem{Returns}
\sphinxAtStartPar
the type of OMP directive that this transformation will             insert.

\sphinxlineitem{Return type}
\sphinxAtStartPar
\sphinxhref{https://docs.python.org/3/library/stdtypes.html\#str}{str}

\end{description}\end{quote}

\end{fulllineitems}



\begin{fulllineitems}

\pysigstartsignatures
\pysigline{\sphinxbfcode{\sphinxupquote{property\DUrole{w,w}{  }}}\sphinxbfcode{\sphinxupquote{omp\_schedule}}}
\pysigstopsignatures\begin{quote}\begin{description}
\sphinxlineitem{Returns}
\sphinxAtStartPar
the OpenMP schedule that will be specified by             this transformation.

\sphinxlineitem{Return type}
\sphinxAtStartPar
\sphinxhref{https://docs.python.org/3/library/stdtypes.html\#str}{str}

\end{description}\end{quote}

\end{fulllineitems}


\end{fulllineitems}



\bigskip\hrule\bigskip



\begin{fulllineitems}

\pysigstartsignatures
\pysigline{\sphinxbfcode{\sphinxupquote{class\DUrole{w,w}{  }}}\sphinxcode{\sphinxupquote{psyclone.transformations.}}\sphinxbfcode{\sphinxupquote{OMPMasterTrans}}}
\pysigstopsignatures
\sphinxAtStartPar
Create an OpenMP MASTER region by inserting directives. The most
likely use case for this transformation is to wrap around task\sphinxhyphen{}based
transformations. Note that adding this directive requires a parent
OpenMP parallel region (which can be inserted by OMPParallelTrans),
otherwise it will produce an error in generation\sphinxhyphen{}time.

\sphinxAtStartPar
For example:

\begin{sphinxVerbatim}[commandchars=\\\{\}]
\PYG{g+gp}{\PYGZgt{}\PYGZgt{}\PYGZgt{} }\PYG{k+kn}{from}\PYG{+w}{ }\PYG{n+nn}{psyclone}\PYG{n+nn}{.}\PYG{n+nn}{parse}\PYG{n+nn}{.}\PYG{n+nn}{algorithm}\PYG{+w}{ }\PYG{k+kn}{import} \PYG{n}{parse}
\PYG{g+gp}{\PYGZgt{}\PYGZgt{}\PYGZgt{} }\PYG{k+kn}{from}\PYG{+w}{ }\PYG{n+nn}{psyclone}\PYG{n+nn}{.}\PYG{n+nn}{psyGen}\PYG{+w}{ }\PYG{k+kn}{import} \PYG{n}{PSyFactory}
\PYG{g+gp}{\PYGZgt{}\PYGZgt{}\PYGZgt{} }\PYG{n}{api} \PYG{o}{=} \PYG{l+s+s2}{\PYGZdq{}}\PYG{l+s+s2}{gocean}\PYG{l+s+s2}{\PYGZdq{}}
\PYG{g+gp}{\PYGZgt{}\PYGZgt{}\PYGZgt{} }\PYG{n}{ast}\PYG{p}{,} \PYG{n}{invokeInfo} \PYG{o}{=} \PYG{n}{parse}\PYG{p}{(}\PYG{n}{GOCEAN\PYGZus{}SOURCE\PYGZus{}FILE}\PYG{p}{,} \PYG{n}{api}\PYG{o}{=}\PYG{n}{api}\PYG{p}{)}
\PYG{g+gp}{\PYGZgt{}\PYGZgt{}\PYGZgt{} }\PYG{n}{psy} \PYG{o}{=} \PYG{n}{PSyFactory}\PYG{p}{(}\PYG{n}{api}\PYG{p}{)}\PYG{o}{.}\PYG{n}{create}\PYG{p}{(}\PYG{n}{invokeInfo}\PYG{p}{)}
\PYG{g+gp}{\PYGZgt{}\PYGZgt{}\PYGZgt{}}
\PYG{g+gp}{\PYGZgt{}\PYGZgt{}\PYGZgt{} }\PYG{k+kn}{from}\PYG{+w}{ }\PYG{n+nn}{psyclone}\PYG{n+nn}{.}\PYG{n+nn}{transformations}\PYG{+w}{ }\PYG{k+kn}{import} \PYG{n}{OMPParallelTrans}\PYG{p}{,} \PYG{n}{OMPMasterTrans}
\PYG{g+gp}{\PYGZgt{}\PYGZgt{}\PYGZgt{} }\PYG{n}{mastertrans} \PYG{o}{=} \PYG{n}{OMPMasterTrans}\PYG{p}{(}\PYG{p}{)}
\PYG{g+gp}{\PYGZgt{}\PYGZgt{}\PYGZgt{} }\PYG{n}{paralleltrans} \PYG{o}{=} \PYG{n}{OMPParallelTrans}\PYG{p}{(}\PYG{p}{)}
\PYG{g+gp}{\PYGZgt{}\PYGZgt{}\PYGZgt{}}
\PYG{g+gp}{\PYGZgt{}\PYGZgt{}\PYGZgt{} }\PYG{n}{schedule} \PYG{o}{=} \PYG{n}{psy}\PYG{o}{.}\PYG{n}{invokes}\PYG{o}{.}\PYG{n}{get}\PYG{p}{(}\PYG{l+s+s1}{\PYGZsq{}}\PYG{l+s+s1}{invoke\PYGZus{}0}\PYG{l+s+s1}{\PYGZsq{}}\PYG{p}{)}\PYG{o}{.}\PYG{n}{schedule}
\PYG{g+gp}{\PYGZgt{}\PYGZgt{}\PYGZgt{} }\PYG{c+c1}{\PYGZsh{} Uncomment the following line to see a text view of the schedule}
\PYG{g+gp}{\PYGZgt{}\PYGZgt{}\PYGZgt{} }\PYG{c+c1}{\PYGZsh{} print(schedule.view())}
\PYG{g+gp}{\PYGZgt{}\PYGZgt{}\PYGZgt{}}
\PYG{g+gp}{\PYGZgt{}\PYGZgt{}\PYGZgt{} }\PYG{c+c1}{\PYGZsh{} Enclose all of these loops within a single OpenMP}
\PYG{g+gp}{\PYGZgt{}\PYGZgt{}\PYGZgt{} }\PYG{c+c1}{\PYGZsh{} MASTER region}
\PYG{g+gp}{\PYGZgt{}\PYGZgt{}\PYGZgt{} }\PYG{n}{mastertrans}\PYG{o}{.}\PYG{n}{apply}\PYG{p}{(}\PYG{n}{schedule}\PYG{o}{.}\PYG{n}{children}\PYG{p}{)}
\PYG{g+gp}{\PYGZgt{}\PYGZgt{}\PYGZgt{} }\PYG{c+c1}{\PYGZsh{} Enclose all of these loops within a single OpenMP}
\PYG{g+gp}{\PYGZgt{}\PYGZgt{}\PYGZgt{} }\PYG{c+c1}{\PYGZsh{} PARALLEL region}
\PYG{g+gp}{\PYGZgt{}\PYGZgt{}\PYGZgt{} }\PYG{n}{paralleltrans}\PYG{o}{.}\PYG{n}{apply}\PYG{p}{(}\PYG{n}{schedule}\PYG{o}{.}\PYG{n}{children}\PYG{p}{)}
\PYG{g+gp}{\PYGZgt{}\PYGZgt{}\PYGZgt{} }\PYG{c+c1}{\PYGZsh{} Uncomment the following line to see a text view of the schedule}
\PYG{g+gp}{\PYGZgt{}\PYGZgt{}\PYGZgt{} }\PYG{c+c1}{\PYGZsh{} print(schedule.view())}
\end{sphinxVerbatim}


\begin{fulllineitems}

\pysigstartsignatures
\pysiglinewithargsret{\sphinxbfcode{\sphinxupquote{apply}}}{\sphinxparam{\DUrole{n,n}{target\_nodes}}\sphinxparamcomma \sphinxparam{\DUrole{n,n}{options}\DUrole{o,o}{=}\DUrole{default_value}{None}}}{}
\pysigstopsignatures
\sphinxAtStartPar
Apply this transformation to a subset of the nodes within a
schedule \sphinxhyphen{} i.e. enclose the specified Loops in the
schedule within a single parallel region.
\begin{quote}\begin{description}
\sphinxlineitem{Parameters}\begin{itemize}
\item {} 
\sphinxAtStartPar
\sphinxstyleliteralstrong{\sphinxupquote{target\_nodes}} ((list of) \sphinxcode{\sphinxupquote{psyclone.psyir.nodes.Node}}) \textendash{} a single Node or a list of Nodes.

\item {} 
\sphinxAtStartPar
\sphinxstyleliteralstrong{\sphinxupquote{options}} (\sphinxstyleliteralemphasis{\sphinxupquote{Optional}}\sphinxstyleliteralemphasis{\sphinxupquote{{[}}}\sphinxstyleliteralemphasis{\sphinxupquote{Dict}}\sphinxstyleliteralemphasis{\sphinxupquote{{[}}}\sphinxhref{https://docs.python.org/3/library/stdtypes.html\#str}{\sphinxstyleliteralemphasis{\sphinxupquote{str}}}\sphinxstyleliteralemphasis{\sphinxupquote{, }}\sphinxstyleliteralemphasis{\sphinxupquote{Any}}\sphinxstyleliteralemphasis{\sphinxupquote{{]}}}\sphinxstyleliteralemphasis{\sphinxupquote{{]}}}) \textendash{} a dictionary with options for transformations.

\item {} 
\sphinxAtStartPar
\sphinxstyleliteralstrong{\sphinxupquote{options}}\sphinxstyleliteralstrong{\sphinxupquote{{[}}}\sphinxstyleliteralstrong{\sphinxupquote{"node\sphinxhyphen{}type\sphinxhyphen{}check"}}\sphinxstyleliteralstrong{\sphinxupquote{{]}}} (\sphinxhref{https://docs.python.org/3/library/functions.html\#bool}{\sphinxstyleliteralemphasis{\sphinxupquote{bool}}}) \textendash{} this flag controls if the                 type of the nodes enclosed in the region should be tested                 to avoid using unsupported nodes inside a region.

\end{itemize}

\end{description}\end{quote}

\end{fulllineitems}



\begin{fulllineitems}

\pysigstartsignatures
\pysiglinewithargsret{\sphinxbfcode{\sphinxupquote{get\_node\_list}}}{\sphinxparam{\DUrole{n,n}{nodes}}}{}
\pysigstopsignatures
\sphinxAtStartPar
This is a helper function for region based transformations.
The parameter for any of those transformations is either a single
node, a schedule, or a list of nodes. This function converts this
into a list of nodes according to the parameter type. This function
will always return a copy, to avoid issues e.g. if a child list
of a node should be provided, and a transformation changes the order
in this list (which would then also change the order of the
nodes in the tree).
\begin{quote}\begin{description}
\sphinxlineitem{Parameters}\begin{itemize}
\item {} 
\sphinxAtStartPar
\sphinxstyleliteralstrong{\sphinxupquote{nodes}} (Union{[}\sphinxcode{\sphinxupquote{psyclone.psyir.nodes.Node}},
\sphinxcode{\sphinxupquote{psyclone.psyir.nodes.Schedule}},
List{[}\sphinxcode{\sphinxupquote{psyclone.psyir.nodes.Node}}{]}) \textendash{} can be a single node, a schedule or a list of nodes.

\item {} 
\sphinxAtStartPar
\sphinxstyleliteralstrong{\sphinxupquote{options}} (\sphinxstyleliteralemphasis{\sphinxupquote{Optional}}\sphinxstyleliteralemphasis{\sphinxupquote{{[}}}\sphinxstyleliteralemphasis{\sphinxupquote{Dict}}\sphinxstyleliteralemphasis{\sphinxupquote{{[}}}\sphinxhref{https://docs.python.org/3/library/stdtypes.html\#str}{\sphinxstyleliteralemphasis{\sphinxupquote{str}}}\sphinxstyleliteralemphasis{\sphinxupquote{,}}\sphinxstyleliteralemphasis{\sphinxupquote{Any}}\sphinxstyleliteralemphasis{\sphinxupquote{{]}}}\sphinxstyleliteralemphasis{\sphinxupquote{{]}}}) \textendash{} a dictionary with options for transformations.

\end{itemize}

\sphinxlineitem{Returns}
\sphinxAtStartPar
a list of nodes.

\sphinxlineitem{Return type}
\sphinxAtStartPar
List{[}\sphinxcode{\sphinxupquote{psyclone.psyir.nodes.Node}}{]}

\sphinxlineitem{Raises}
\sphinxAtStartPar
\sphinxstyleliteralstrong{\sphinxupquote{TransformationError}} \textendash{} if the supplied parameter is neither a             single Node, nor a Schedule, nor a list of Nodes.

\end{description}\end{quote}

\end{fulllineitems}



\begin{fulllineitems}

\pysigstartsignatures
\pysiglinewithargsret{\sphinxbfcode{\sphinxupquote{validate}}}{\sphinxparam{\DUrole{n,n}{node\_list}}\sphinxparamcomma \sphinxparam{\DUrole{n,n}{options}\DUrole{o,o}{=}\DUrole{default_value}{None}}}{}
\pysigstopsignatures
\sphinxAtStartPar
Check that the supplied list of Nodes are eligible to be
put inside a parallel region.
\begin{quote}\begin{description}
\sphinxlineitem{Parameters}\begin{itemize}
\item {} 
\sphinxAtStartPar
\sphinxstyleliteralstrong{\sphinxupquote{node\_list}} (\sphinxhref{https://docs.python.org/3/library/stdtypes.html\#list}{\sphinxstyleliteralemphasis{\sphinxupquote{list}}}) \textendash{} list of nodes to put into a parallel region

\item {} 
\sphinxAtStartPar
\sphinxstyleliteralstrong{\sphinxupquote{options}} \textendash{} a dictionary with options for transformations.        :type options: Optional{[}Dict{[}str, Any{]}{]}

\item {} 
\sphinxAtStartPar
\sphinxstyleliteralstrong{\sphinxupquote{options}}\sphinxstyleliteralstrong{\sphinxupquote{{[}}}\sphinxstyleliteralstrong{\sphinxupquote{"node\sphinxhyphen{}type\sphinxhyphen{}check"}}\sphinxstyleliteralstrong{\sphinxupquote{{]}}} (\sphinxhref{https://docs.python.org/3/library/functions.html\#bool}{\sphinxstyleliteralemphasis{\sphinxupquote{bool}}}) \textendash{} this flag controls whether             or not the type of the nodes enclosed in the region should be             tested to avoid using unsupported nodes inside a region.

\end{itemize}

\sphinxlineitem{Raises}\begin{itemize}
\item {} 
\sphinxAtStartPar
\sphinxstyleliteralstrong{\sphinxupquote{TransformationError}} \textendash{} if the supplied node is an             InvokeSchedule rather than being within an InvokeSchedule.

\item {} 
\sphinxAtStartPar
\sphinxstyleliteralstrong{\sphinxupquote{TransformationError}} \textendash{} if the supplied nodes are not all             children of the same parent (siblings).

\end{itemize}

\end{description}\end{quote}

\end{fulllineitems}


\end{fulllineitems}


\begin{sphinxadmonition}{note}{Note:}
\sphinxAtStartPar
PSyclone does not support (distributed\sphinxhyphen{}memory) halo swaps or
global sums within OpenMP master regions.  Attempting to
create a master region for a set of nodes that includes
halo swaps or global sums will produce an error. In such
cases it may be possible to re\sphinxhyphen{}order the nodes in the
Schedule such that the halo swaps or global sums are
performed outside the single region. The
{\hyperref[\detokenize{transformations:sec-move-trans}]{\sphinxcrossref{\DUrole{std,std-ref}{MoveTrans}}}} transformation may be used
for this.
\end{sphinxadmonition}


\bigskip\hrule\bigskip



\begin{fulllineitems}

\pysigstartsignatures
\pysiglinewithargsret{\sphinxbfcode{\sphinxupquote{class\DUrole{w,w}{  }}}\sphinxcode{\sphinxupquote{psyclone.transformations.}}\sphinxbfcode{\sphinxupquote{OMPParallelLoopTrans}}}{\sphinxparam{\DUrole{n,n}{omp\_directive}\DUrole{o,o}{=}\DUrole{default_value}{\textquotesingle{}do\textquotesingle{}}}\sphinxparamcomma \sphinxparam{\DUrole{n,n}{omp\_schedule}\DUrole{o,o}{=}\DUrole{default_value}{\textquotesingle{}auto\textquotesingle{}}}}{}
\pysigstopsignatures
\sphinxAtStartPar
Adds an OpenMP PARALLEL DO directive to a loop.

\sphinxAtStartPar
For example:

\begin{sphinxVerbatim}[commandchars=\\\{\}]
\PYG{g+gp}{\PYGZgt{}\PYGZgt{}\PYGZgt{} }\PYG{k+kn}{from}\PYG{+w}{ }\PYG{n+nn}{psyclone}\PYG{n+nn}{.}\PYG{n+nn}{parse}\PYG{n+nn}{.}\PYG{n+nn}{algorithm}\PYG{+w}{ }\PYG{k+kn}{import} \PYG{n}{parse}
\PYG{g+gp}{\PYGZgt{}\PYGZgt{}\PYGZgt{} }\PYG{k+kn}{from}\PYG{+w}{ }\PYG{n+nn}{psyclone}\PYG{n+nn}{.}\PYG{n+nn}{psyGen}\PYG{+w}{ }\PYG{k+kn}{import} \PYG{n}{PSyFactory}
\PYG{g+gp}{\PYGZgt{}\PYGZgt{}\PYGZgt{} }\PYG{n}{ast}\PYG{p}{,} \PYG{n}{invokeInfo} \PYG{o}{=} \PYG{n}{parse}\PYG{p}{(}\PYG{l+s+s2}{\PYGZdq{}}\PYG{l+s+s2}{dynamo.F90}\PYG{l+s+s2}{\PYGZdq{}}\PYG{p}{)}
\PYG{g+gp}{\PYGZgt{}\PYGZgt{}\PYGZgt{} }\PYG{n}{psy} \PYG{o}{=} \PYG{n}{PSyFactory}\PYG{p}{(}\PYG{l+s+s2}{\PYGZdq{}}\PYG{l+s+s2}{lfric}\PYG{l+s+s2}{\PYGZdq{}}\PYG{p}{)}\PYG{o}{.}\PYG{n}{create}\PYG{p}{(}\PYG{n}{invokeInfo}\PYG{p}{)}
\PYG{g+gp}{\PYGZgt{}\PYGZgt{}\PYGZgt{} }\PYG{n}{schedule} \PYG{o}{=} \PYG{n}{psy}\PYG{o}{.}\PYG{n}{invokes}\PYG{o}{.}\PYG{n}{get}\PYG{p}{(}\PYG{l+s+s1}{\PYGZsq{}}\PYG{l+s+s1}{invoke\PYGZus{}v3\PYGZus{}kernel\PYGZus{}type}\PYG{l+s+s1}{\PYGZsq{}}\PYG{p}{)}\PYG{o}{.}\PYG{n}{schedule}
\PYG{g+gp}{\PYGZgt{}\PYGZgt{}\PYGZgt{} }\PYG{c+c1}{\PYGZsh{} Uncomment the following line to see a text view of the schedule}
\PYG{g+gp}{\PYGZgt{}\PYGZgt{}\PYGZgt{} }\PYG{c+c1}{\PYGZsh{} print(schedule.view())}
\PYG{g+gp}{\PYGZgt{}\PYGZgt{}\PYGZgt{}}
\PYG{g+gp}{\PYGZgt{}\PYGZgt{}\PYGZgt{} }\PYG{k+kn}{from}\PYG{+w}{ }\PYG{n+nn}{psyclone}\PYG{n+nn}{.}\PYG{n+nn}{transformations}\PYG{+w}{ }\PYG{k+kn}{import} \PYG{n}{OMPParallelLoopTrans}
\PYG{g+gp}{\PYGZgt{}\PYGZgt{}\PYGZgt{} }\PYG{n}{trans} \PYG{o}{=} \PYG{n}{OMPParallelLoopTrans}\PYG{p}{(}\PYG{p}{)}
\PYG{g+gp}{\PYGZgt{}\PYGZgt{}\PYGZgt{} }\PYG{n}{trans}\PYG{o}{.}\PYG{n}{apply}\PYG{p}{(}\PYG{n}{schedule}\PYG{o}{.}\PYG{n}{children}\PYG{p}{[}\PYG{l+m+mi}{0}\PYG{p}{]}\PYG{p}{)}
\PYG{g+gp}{\PYGZgt{}\PYGZgt{}\PYGZgt{} }\PYG{c+c1}{\PYGZsh{} Uncomment the following line to see a text view of the schedule}
\PYG{g+gp}{\PYGZgt{}\PYGZgt{}\PYGZgt{} }\PYG{c+c1}{\PYGZsh{} print(schedule.view())}
\end{sphinxVerbatim}


\begin{fulllineitems}

\pysigstartsignatures
\pysiglinewithargsret{\sphinxbfcode{\sphinxupquote{apply}}}{\sphinxparam{\DUrole{n,n}{node}}\sphinxparamcomma \sphinxparam{\DUrole{n,n}{options}\DUrole{o,o}{=}\DUrole{default_value}{None}}}{}
\pysigstopsignatures
\sphinxAtStartPar
Apply an OMPParallelLoop Transformation to the supplied node
(which must be a Loop). In the generated code this corresponds to
wrapping the Loop with directives:

\begin{sphinxVerbatim}[commandchars=\\\{\}]
\PYG{c}{!\PYGZdl{}OMP PARALLEL DO ...}
\PYG{k}{do}\PYG{+w}{ }\PYG{p}{.}\PYG{p}{.}\PYG{p}{.}
\PYG{+w}{  }\PYG{p}{.}\PYG{p}{.}\PYG{p}{.}
\PYG{k}{end }\PYG{k}{do}
\PYG{c}{!\PYGZdl{}OMP END PARALLEL DO}
\end{sphinxVerbatim}
\begin{quote}\begin{description}
\sphinxlineitem{Parameters}\begin{itemize}
\item {} 
\sphinxAtStartPar
\sphinxstyleliteralstrong{\sphinxupquote{node}} (\sphinxcode{\sphinxupquote{psyclone.f2pygen.DoGen}}) \textendash{} the node (loop) to which to apply the transformation.

\item {} 
\sphinxAtStartPar
\sphinxstyleliteralstrong{\sphinxupquote{options}} (\sphinxstyleliteralemphasis{\sphinxupquote{Optional}}\sphinxstyleliteralemphasis{\sphinxupquote{{[}}}\sphinxstyleliteralemphasis{\sphinxupquote{Dict}}\sphinxstyleliteralemphasis{\sphinxupquote{{[}}}\sphinxhref{https://docs.python.org/3/library/stdtypes.html\#str}{\sphinxstyleliteralemphasis{\sphinxupquote{str}}}\sphinxstyleliteralemphasis{\sphinxupquote{, }}\sphinxstyleliteralemphasis{\sphinxupquote{Any}}\sphinxstyleliteralemphasis{\sphinxupquote{{]}}}\sphinxstyleliteralemphasis{\sphinxupquote{{]}}}) \textendash{} a dictionary with options for transformations                        and validation.

\end{itemize}

\end{description}\end{quote}

\end{fulllineitems}


\end{fulllineitems}



\bigskip\hrule\bigskip



\begin{fulllineitems}

\pysigstartsignatures
\pysigline{\sphinxbfcode{\sphinxupquote{class\DUrole{w,w}{  }}}\sphinxcode{\sphinxupquote{psyclone.transformations.}}\sphinxbfcode{\sphinxupquote{OMPParallelTrans}}}
\pysigstopsignatures
\sphinxAtStartPar
Create an OpenMP PARALLEL region by inserting directives. For
example:

\begin{sphinxVerbatim}[commandchars=\\\{\}]
\PYG{g+gp}{\PYGZgt{}\PYGZgt{}\PYGZgt{} }\PYG{k+kn}{from}\PYG{+w}{ }\PYG{n+nn}{psyclone}\PYG{n+nn}{.}\PYG{n+nn}{parse}\PYG{n+nn}{.}\PYG{n+nn}{algorithm}\PYG{+w}{ }\PYG{k+kn}{import} \PYG{n}{parse}
\PYG{g+gp}{\PYGZgt{}\PYGZgt{}\PYGZgt{} }\PYG{k+kn}{from}\PYG{+w}{ }\PYG{n+nn}{psyclone}\PYG{n+nn}{.}\PYG{n+nn}{parse}\PYG{n+nn}{.}\PYG{n+nn}{utils}\PYG{+w}{ }\PYG{k+kn}{import} \PYG{n}{ParseError}
\PYG{g+gp}{\PYGZgt{}\PYGZgt{}\PYGZgt{} }\PYG{k+kn}{from}\PYG{+w}{ }\PYG{n+nn}{psyclone}\PYG{n+nn}{.}\PYG{n+nn}{psyGen}\PYG{+w}{ }\PYG{k+kn}{import} \PYG{n}{PSyFactory}
\PYG{g+gp}{\PYGZgt{}\PYGZgt{}\PYGZgt{} }\PYG{k+kn}{from}\PYG{+w}{ }\PYG{n+nn}{psyclone}\PYG{n+nn}{.}\PYG{n+nn}{errors}\PYG{+w}{ }\PYG{k+kn}{import} \PYG{n}{GenerationError}
\PYG{g+gp}{\PYGZgt{}\PYGZgt{}\PYGZgt{} }\PYG{n}{api} \PYG{o}{=} \PYG{l+s+s2}{\PYGZdq{}}\PYG{l+s+s2}{gocean}\PYG{l+s+s2}{\PYGZdq{}}
\PYG{g+gp}{\PYGZgt{}\PYGZgt{}\PYGZgt{} }\PYG{n}{ast}\PYG{p}{,} \PYG{n}{invokeInfo} \PYG{o}{=} \PYG{n}{parse}\PYG{p}{(}\PYG{n}{GOCEAN\PYGZus{}SOURCE\PYGZus{}FILE}\PYG{p}{,} \PYG{n}{api}\PYG{o}{=}\PYG{n}{api}\PYG{p}{)}
\PYG{g+gp}{\PYGZgt{}\PYGZgt{}\PYGZgt{} }\PYG{n}{psy} \PYG{o}{=} \PYG{n}{PSyFactory}\PYG{p}{(}\PYG{n}{api}\PYG{p}{)}\PYG{o}{.}\PYG{n}{create}\PYG{p}{(}\PYG{n}{invokeInfo}\PYG{p}{)}
\PYG{g+gp}{\PYGZgt{}\PYGZgt{}\PYGZgt{}}
\PYG{g+gp}{\PYGZgt{}\PYGZgt{}\PYGZgt{} }\PYG{k+kn}{from}\PYG{+w}{ }\PYG{n+nn}{psyclone}\PYG{n+nn}{.}\PYG{n+nn}{psyGen}\PYG{+w}{ }\PYG{k+kn}{import} \PYG{n}{TransInfo}
\PYG{g+gp}{\PYGZgt{}\PYGZgt{}\PYGZgt{} }\PYG{n}{t} \PYG{o}{=} \PYG{n}{TransInfo}\PYG{p}{(}\PYG{p}{)}
\PYG{g+gp}{\PYGZgt{}\PYGZgt{}\PYGZgt{} }\PYG{n}{ltrans} \PYG{o}{=} \PYG{n}{t}\PYG{o}{.}\PYG{n}{get\PYGZus{}trans\PYGZus{}name}\PYG{p}{(}\PYG{l+s+s1}{\PYGZsq{}}\PYG{l+s+s1}{GOceanOMPLoopTrans}\PYG{l+s+s1}{\PYGZsq{}}\PYG{p}{)}
\PYG{g+gp}{\PYGZgt{}\PYGZgt{}\PYGZgt{} }\PYG{n}{rtrans} \PYG{o}{=} \PYG{n}{t}\PYG{o}{.}\PYG{n}{get\PYGZus{}trans\PYGZus{}name}\PYG{p}{(}\PYG{l+s+s1}{\PYGZsq{}}\PYG{l+s+s1}{OMPParallelTrans}\PYG{l+s+s1}{\PYGZsq{}}\PYG{p}{)}
\PYG{g+gp}{\PYGZgt{}\PYGZgt{}\PYGZgt{}}
\PYG{g+gp}{\PYGZgt{}\PYGZgt{}\PYGZgt{} }\PYG{n}{schedule} \PYG{o}{=} \PYG{n}{psy}\PYG{o}{.}\PYG{n}{invokes}\PYG{o}{.}\PYG{n}{get}\PYG{p}{(}\PYG{l+s+s1}{\PYGZsq{}}\PYG{l+s+s1}{invoke\PYGZus{}0}\PYG{l+s+s1}{\PYGZsq{}}\PYG{p}{)}\PYG{o}{.}\PYG{n}{schedule}
\PYG{g+gp}{\PYGZgt{}\PYGZgt{}\PYGZgt{} }\PYG{c+c1}{\PYGZsh{} Uncomment the following line to see a text view of the schedule}
\PYG{g+gp}{\PYGZgt{}\PYGZgt{}\PYGZgt{} }\PYG{c+c1}{\PYGZsh{} print(schedule.view())}
\PYG{g+gp}{\PYGZgt{}\PYGZgt{}\PYGZgt{}}
\PYG{g+gp}{\PYGZgt{}\PYGZgt{}\PYGZgt{} }\PYG{c+c1}{\PYGZsh{} Apply the OpenMP Loop transformation to *every* loop}
\PYG{g+gp}{\PYGZgt{}\PYGZgt{}\PYGZgt{} }\PYG{c+c1}{\PYGZsh{} in the schedule}
\PYG{g+gp}{\PYGZgt{}\PYGZgt{}\PYGZgt{} }\PYG{k}{for} \PYG{n}{child} \PYG{o+ow}{in} \PYG{n}{schedule}\PYG{o}{.}\PYG{n}{children}\PYG{p}{:}
\PYG{g+gp}{\PYGZgt{}\PYGZgt{}\PYGZgt{} }    \PYG{n}{ltrans}\PYG{o}{.}\PYG{n}{apply}\PYG{p}{(}\PYG{n}{child}\PYG{p}{)}
\PYG{g+gp}{\PYGZgt{}\PYGZgt{}\PYGZgt{}}
\PYG{g+gp}{\PYGZgt{}\PYGZgt{}\PYGZgt{} }\PYG{c+c1}{\PYGZsh{} Enclose all of these loops within a single OpenMP}
\PYG{g+gp}{\PYGZgt{}\PYGZgt{}\PYGZgt{} }\PYG{c+c1}{\PYGZsh{} PARALLEL region}
\PYG{g+gp}{\PYGZgt{}\PYGZgt{}\PYGZgt{} }\PYG{n}{rtrans}\PYG{o}{.}\PYG{n}{apply}\PYG{p}{(}\PYG{n}{schedule}\PYG{o}{.}\PYG{n}{children}\PYG{p}{)}
\PYG{g+gp}{\PYGZgt{}\PYGZgt{}\PYGZgt{} }\PYG{c+c1}{\PYGZsh{} Uncomment the following line to see a text view of the schedule}
\PYG{g+gp}{\PYGZgt{}\PYGZgt{}\PYGZgt{} }\PYG{c+c1}{\PYGZsh{} print(schedule.view())}
\end{sphinxVerbatim}


\begin{fulllineitems}

\pysigstartsignatures
\pysiglinewithargsret{\sphinxbfcode{\sphinxupquote{apply}}}{\sphinxparam{\DUrole{n,n}{target\_nodes}}\sphinxparamcomma \sphinxparam{\DUrole{n,n}{options}\DUrole{o,o}{=}\DUrole{default_value}{None}}}{}
\pysigstopsignatures
\sphinxAtStartPar
Apply this transformation to a subset of the nodes within a
schedule \sphinxhyphen{} i.e. enclose the specified Loops in the
schedule within a single parallel region.
\begin{quote}\begin{description}
\sphinxlineitem{Parameters}\begin{itemize}
\item {} 
\sphinxAtStartPar
\sphinxstyleliteralstrong{\sphinxupquote{target\_nodes}} ((list of) \sphinxcode{\sphinxupquote{psyclone.psyir.nodes.Node}}) \textendash{} a single Node or a list of Nodes.

\item {} 
\sphinxAtStartPar
\sphinxstyleliteralstrong{\sphinxupquote{options}} (\sphinxstyleliteralemphasis{\sphinxupquote{Optional}}\sphinxstyleliteralemphasis{\sphinxupquote{{[}}}\sphinxstyleliteralemphasis{\sphinxupquote{Dict}}\sphinxstyleliteralemphasis{\sphinxupquote{{[}}}\sphinxhref{https://docs.python.org/3/library/stdtypes.html\#str}{\sphinxstyleliteralemphasis{\sphinxupquote{str}}}\sphinxstyleliteralemphasis{\sphinxupquote{, }}\sphinxstyleliteralemphasis{\sphinxupquote{Any}}\sphinxstyleliteralemphasis{\sphinxupquote{{]}}}\sphinxstyleliteralemphasis{\sphinxupquote{{]}}}) \textendash{} a dictionary with options for transformations.

\item {} 
\sphinxAtStartPar
\sphinxstyleliteralstrong{\sphinxupquote{options}}\sphinxstyleliteralstrong{\sphinxupquote{{[}}}\sphinxstyleliteralstrong{\sphinxupquote{"node\sphinxhyphen{}type\sphinxhyphen{}check"}}\sphinxstyleliteralstrong{\sphinxupquote{{]}}} (\sphinxhref{https://docs.python.org/3/library/functions.html\#bool}{\sphinxstyleliteralemphasis{\sphinxupquote{bool}}}) \textendash{} this flag controls if the                 type of the nodes enclosed in the region should be tested                 to avoid using unsupported nodes inside a region.

\end{itemize}

\end{description}\end{quote}

\end{fulllineitems}



\begin{fulllineitems}

\pysigstartsignatures
\pysiglinewithargsret{\sphinxbfcode{\sphinxupquote{get\_node\_list}}}{\sphinxparam{\DUrole{n,n}{nodes}}}{}
\pysigstopsignatures
\sphinxAtStartPar
This is a helper function for region based transformations.
The parameter for any of those transformations is either a single
node, a schedule, or a list of nodes. This function converts this
into a list of nodes according to the parameter type. This function
will always return a copy, to avoid issues e.g. if a child list
of a node should be provided, and a transformation changes the order
in this list (which would then also change the order of the
nodes in the tree).
\begin{quote}\begin{description}
\sphinxlineitem{Parameters}\begin{itemize}
\item {} 
\sphinxAtStartPar
\sphinxstyleliteralstrong{\sphinxupquote{nodes}} (Union{[}\sphinxcode{\sphinxupquote{psyclone.psyir.nodes.Node}},
\sphinxcode{\sphinxupquote{psyclone.psyir.nodes.Schedule}},
List{[}\sphinxcode{\sphinxupquote{psyclone.psyir.nodes.Node}}{]}) \textendash{} can be a single node, a schedule or a list of nodes.

\item {} 
\sphinxAtStartPar
\sphinxstyleliteralstrong{\sphinxupquote{options}} (\sphinxstyleliteralemphasis{\sphinxupquote{Optional}}\sphinxstyleliteralemphasis{\sphinxupquote{{[}}}\sphinxstyleliteralemphasis{\sphinxupquote{Dict}}\sphinxstyleliteralemphasis{\sphinxupquote{{[}}}\sphinxhref{https://docs.python.org/3/library/stdtypes.html\#str}{\sphinxstyleliteralemphasis{\sphinxupquote{str}}}\sphinxstyleliteralemphasis{\sphinxupquote{,}}\sphinxstyleliteralemphasis{\sphinxupquote{Any}}\sphinxstyleliteralemphasis{\sphinxupquote{{]}}}\sphinxstyleliteralemphasis{\sphinxupquote{{]}}}) \textendash{} a dictionary with options for transformations.

\end{itemize}

\sphinxlineitem{Returns}
\sphinxAtStartPar
a list of nodes.

\sphinxlineitem{Return type}
\sphinxAtStartPar
List{[}\sphinxcode{\sphinxupquote{psyclone.psyir.nodes.Node}}{]}

\sphinxlineitem{Raises}
\sphinxAtStartPar
\sphinxstyleliteralstrong{\sphinxupquote{TransformationError}} \textendash{} if the supplied parameter is neither a             single Node, nor a Schedule, nor a list of Nodes.

\end{description}\end{quote}

\end{fulllineitems}



\begin{fulllineitems}

\pysigstartsignatures
\pysiglinewithargsret{\sphinxbfcode{\sphinxupquote{validate}}}{\sphinxparam{\DUrole{n,n}{node\_list}}\sphinxparamcomma \sphinxparam{\DUrole{n,n}{options}\DUrole{o,o}{=}\DUrole{default_value}{None}}}{}
\pysigstopsignatures
\sphinxAtStartPar
Perform OpenMP\sphinxhyphen{}specific validation checks.
\begin{quote}\begin{description}
\sphinxlineitem{Parameters}\begin{itemize}
\item {} 
\sphinxAtStartPar
\sphinxstyleliteralstrong{\sphinxupquote{node\_list}} (list of \sphinxcode{\sphinxupquote{psyclone.psyir.nodes.Node}}) \textendash{} list of Nodes to put within parallel region.

\item {} 
\sphinxAtStartPar
\sphinxstyleliteralstrong{\sphinxupquote{options}} (\sphinxstyleliteralemphasis{\sphinxupquote{Optional}}\sphinxstyleliteralemphasis{\sphinxupquote{{[}}}\sphinxstyleliteralemphasis{\sphinxupquote{Dict}}\sphinxstyleliteralemphasis{\sphinxupquote{{[}}}\sphinxhref{https://docs.python.org/3/library/stdtypes.html\#str}{\sphinxstyleliteralemphasis{\sphinxupquote{str}}}\sphinxstyleliteralemphasis{\sphinxupquote{, }}\sphinxstyleliteralemphasis{\sphinxupquote{Any}}\sphinxstyleliteralemphasis{\sphinxupquote{{]}}}\sphinxstyleliteralemphasis{\sphinxupquote{{]}}}) \textendash{} a dictionary with options for transformations.

\item {} 
\sphinxAtStartPar
\sphinxstyleliteralstrong{\sphinxupquote{options}}\sphinxstyleliteralstrong{\sphinxupquote{{[}}}\sphinxstyleliteralstrong{\sphinxupquote{"node\sphinxhyphen{}type\sphinxhyphen{}check"}}\sphinxstyleliteralstrong{\sphinxupquote{{]}}} (\sphinxhref{https://docs.python.org/3/library/functions.html\#bool}{\sphinxstyleliteralemphasis{\sphinxupquote{bool}}}) \textendash{} this flag controls if the                 type of the nodes enclosed in the region should be tested                 to avoid using unsupported nodes inside a region.

\end{itemize}

\sphinxlineitem{Raises}
\sphinxAtStartPar
\sphinxstyleliteralstrong{\sphinxupquote{TransformationError}} \textendash{} if the target Nodes are already within                                      some OMP parallel region.

\end{description}\end{quote}

\end{fulllineitems}


\end{fulllineitems}


\begin{sphinxadmonition}{note}{Note:}
\sphinxAtStartPar
PSyclone does not support (distributed\sphinxhyphen{}memory) halo swaps or
global sums within OpenMP parallel regions.  Attempting to
create a parallel region for a set of nodes that includes
halo swaps or global sums will produce an error. In such
cases it may be possible to re\sphinxhyphen{}order the nodes in the
Schedule such that the halo swaps or global sums are
performed outside the parallel region. The
{\hyperref[\detokenize{transformations:sec-move-trans}]{\sphinxcrossref{\DUrole{std,std-ref}{MoveTrans}}}} transformation may be used
for this.
\end{sphinxadmonition}


\bigskip\hrule\bigskip



\begin{fulllineitems}

\pysigstartsignatures
\pysiglinewithargsret{\sphinxbfcode{\sphinxupquote{class\DUrole{w,w}{  }}}\sphinxcode{\sphinxupquote{psyclone.transformations.}}\sphinxbfcode{\sphinxupquote{OMPSingleTrans}}}{\sphinxparam{\DUrole{n,n}{nowait}\DUrole{o,o}{=}\DUrole{default_value}{False}}}{}
\pysigstopsignatures
\sphinxAtStartPar
Create an OpenMP SINGLE region by inserting directives. The most
likely use case for this transformation is to wrap around task\sphinxhyphen{}based
transformations. The parent region for this should usually also be
a OMPParallelTrans.
\begin{quote}\begin{description}
\sphinxlineitem{Parameters}
\sphinxAtStartPar
\sphinxstyleliteralstrong{\sphinxupquote{nowait}} (\sphinxhref{https://docs.python.org/3/library/functions.html\#bool}{\sphinxstyleliteralemphasis{\sphinxupquote{bool}}}) \textendash{} whether to apply a nowait clause to this                        transformation. The default value is False

\end{description}\end{quote}

\sphinxAtStartPar
For example:

\begin{sphinxVerbatim}[commandchars=\\\{\}]
\PYG{g+gp}{\PYGZgt{}\PYGZgt{}\PYGZgt{} }\PYG{k+kn}{from}\PYG{+w}{ }\PYG{n+nn}{psyclone}\PYG{n+nn}{.}\PYG{n+nn}{parse}\PYG{n+nn}{.}\PYG{n+nn}{algorithm}\PYG{+w}{ }\PYG{k+kn}{import} \PYG{n}{parse}
\PYG{g+gp}{\PYGZgt{}\PYGZgt{}\PYGZgt{} }\PYG{k+kn}{from}\PYG{+w}{ }\PYG{n+nn}{psyclone}\PYG{n+nn}{.}\PYG{n+nn}{psyGen}\PYG{+w}{ }\PYG{k+kn}{import} \PYG{n}{PSyFactory}
\PYG{g+gp}{\PYGZgt{}\PYGZgt{}\PYGZgt{} }\PYG{n}{api} \PYG{o}{=} \PYG{l+s+s2}{\PYGZdq{}}\PYG{l+s+s2}{gocean}\PYG{l+s+s2}{\PYGZdq{}}
\PYG{g+gp}{\PYGZgt{}\PYGZgt{}\PYGZgt{} }\PYG{n}{ast}\PYG{p}{,} \PYG{n}{invokeInfo} \PYG{o}{=} \PYG{n}{parse}\PYG{p}{(}\PYG{n}{GOCEAN\PYGZus{}SOURCE\PYGZus{}FILE}\PYG{p}{,} \PYG{n}{api}\PYG{o}{=}\PYG{n}{api}\PYG{p}{)}
\PYG{g+gp}{\PYGZgt{}\PYGZgt{}\PYGZgt{} }\PYG{n}{psy} \PYG{o}{=} \PYG{n}{PSyFactory}\PYG{p}{(}\PYG{n}{api}\PYG{p}{)}\PYG{o}{.}\PYG{n}{create}\PYG{p}{(}\PYG{n}{invokeInfo}\PYG{p}{)}
\PYG{g+gp}{\PYGZgt{}\PYGZgt{}\PYGZgt{}}
\PYG{g+gp}{\PYGZgt{}\PYGZgt{}\PYGZgt{} }\PYG{k+kn}{from}\PYG{+w}{ }\PYG{n+nn}{psyclone}\PYG{n+nn}{.}\PYG{n+nn}{transformations}\PYG{+w}{ }\PYG{k+kn}{import} \PYG{n}{OMPParallelTrans}\PYG{p}{,} \PYG{n}{OMPSingleTrans}
\PYG{g+gp}{\PYGZgt{}\PYGZgt{}\PYGZgt{} }\PYG{n}{singletrans} \PYG{o}{=} \PYG{n}{OMPSingleTrans}\PYG{p}{(}\PYG{p}{)}
\PYG{g+gp}{\PYGZgt{}\PYGZgt{}\PYGZgt{} }\PYG{n}{paralleltrans} \PYG{o}{=} \PYG{n}{OMPParallelTrans}\PYG{p}{(}\PYG{p}{)}
\PYG{g+gp}{\PYGZgt{}\PYGZgt{}\PYGZgt{}}
\PYG{g+gp}{\PYGZgt{}\PYGZgt{}\PYGZgt{} }\PYG{n}{schedule} \PYG{o}{=} \PYG{n}{psy}\PYG{o}{.}\PYG{n}{invokes}\PYG{o}{.}\PYG{n}{get}\PYG{p}{(}\PYG{l+s+s1}{\PYGZsq{}}\PYG{l+s+s1}{invoke\PYGZus{}0}\PYG{l+s+s1}{\PYGZsq{}}\PYG{p}{)}\PYG{o}{.}\PYG{n}{schedule}
\PYG{g+gp}{\PYGZgt{}\PYGZgt{}\PYGZgt{} }\PYG{c+c1}{\PYGZsh{} Uncomment the following line to see a text view of the schedule}
\PYG{g+gp}{\PYGZgt{}\PYGZgt{}\PYGZgt{} }\PYG{c+c1}{\PYGZsh{} print(schedule.view())}
\PYG{g+gp}{\PYGZgt{}\PYGZgt{}\PYGZgt{}}
\PYG{g+gp}{\PYGZgt{}\PYGZgt{}\PYGZgt{} }\PYG{c+c1}{\PYGZsh{} Enclose all of these loops within a single OpenMP}
\PYG{g+gp}{\PYGZgt{}\PYGZgt{}\PYGZgt{} }\PYG{c+c1}{\PYGZsh{} SINGLE region}
\PYG{g+gp}{\PYGZgt{}\PYGZgt{}\PYGZgt{} }\PYG{n}{singletrans}\PYG{o}{.}\PYG{n}{apply}\PYG{p}{(}\PYG{n}{schedule}\PYG{o}{.}\PYG{n}{children}\PYG{p}{)}
\PYG{g+gp}{\PYGZgt{}\PYGZgt{}\PYGZgt{} }\PYG{c+c1}{\PYGZsh{} Enclose all of these loops within a single OpenMP}
\PYG{g+gp}{\PYGZgt{}\PYGZgt{}\PYGZgt{} }\PYG{c+c1}{\PYGZsh{} PARALLEL region}
\PYG{g+gp}{\PYGZgt{}\PYGZgt{}\PYGZgt{} }\PYG{n}{paralleltrans}\PYG{o}{.}\PYG{n}{apply}\PYG{p}{(}\PYG{n}{schedule}\PYG{o}{.}\PYG{n}{children}\PYG{p}{)}
\PYG{g+gp}{\PYGZgt{}\PYGZgt{}\PYGZgt{} }\PYG{c+c1}{\PYGZsh{} Uncomment the following line to see a text view of the schedule}
\PYG{g+gp}{\PYGZgt{}\PYGZgt{}\PYGZgt{} }\PYG{c+c1}{\PYGZsh{} print(schedule.view())}
\end{sphinxVerbatim}


\begin{fulllineitems}

\pysigstartsignatures
\pysiglinewithargsret{\sphinxbfcode{\sphinxupquote{apply}}}{\sphinxparam{\DUrole{n,n}{node\_list}}\sphinxparamcomma \sphinxparam{\DUrole{n,n}{options}\DUrole{o,o}{=}\DUrole{default_value}{None}}}{}
\pysigstopsignatures
\sphinxAtStartPar
Apply the OMPSingleTrans transformation to the specified node in a
Schedule.

\sphinxAtStartPar
At code\sphinxhyphen{}generation time this node must be within (i.e. a child of)
an OpenMP PARALLEL region. Code generation happens when
\sphinxcode{\sphinxupquote{OMPLoopDirective.gen\_code()}} is called, or when the PSyIR
tree is given to a backend.

\sphinxAtStartPar
If the keyword “nowait” is specified in the options, it will cause a
nowait clause to be added if it is set to True, otherwise no clause
will be added.
\begin{quote}\begin{description}
\sphinxlineitem{Parameters}\begin{itemize}
\item {} 
\sphinxAtStartPar
\sphinxstyleliteralstrong{\sphinxupquote{node\_list}} ((a list of) \sphinxcode{\sphinxupquote{psyclone.psyir.nodes.Node}}) \textendash{} the supplied node or node list to which we will                           apply the OMPSingleTrans transformation

\item {} 
\sphinxAtStartPar
\sphinxstyleliteralstrong{\sphinxupquote{options}} (\sphinxstyleliteralemphasis{\sphinxupquote{Optional}}\sphinxstyleliteralemphasis{\sphinxupquote{{[}}}\sphinxstyleliteralemphasis{\sphinxupquote{Dict}}\sphinxstyleliteralemphasis{\sphinxupquote{{[}}}\sphinxhref{https://docs.python.org/3/library/stdtypes.html\#str}{\sphinxstyleliteralemphasis{\sphinxupquote{str}}}\sphinxstyleliteralemphasis{\sphinxupquote{, }}\sphinxstyleliteralemphasis{\sphinxupquote{Any}}\sphinxstyleliteralemphasis{\sphinxupquote{{]}}}\sphinxstyleliteralemphasis{\sphinxupquote{{]}}}) \textendash{} a list with options for transformations                         and validation.

\item {} 
\sphinxAtStartPar
\sphinxstyleliteralstrong{\sphinxupquote{options}}\sphinxstyleliteralstrong{\sphinxupquote{{[}}}\sphinxstyleliteralstrong{\sphinxupquote{"nowait"}}\sphinxstyleliteralstrong{\sphinxupquote{{]}}} (\sphinxhref{https://docs.python.org/3/library/functions.html\#bool}{\sphinxstyleliteralemphasis{\sphinxupquote{bool}}}) \textendash{} indicating whether or not to use a nowait clause on this                 single region.

\end{itemize}

\end{description}\end{quote}

\end{fulllineitems}



\begin{fulllineitems}

\pysigstartsignatures
\pysiglinewithargsret{\sphinxbfcode{\sphinxupquote{get\_node\_list}}}{\sphinxparam{\DUrole{n,n}{nodes}}}{}
\pysigstopsignatures
\sphinxAtStartPar
This is a helper function for region based transformations.
The parameter for any of those transformations is either a single
node, a schedule, or a list of nodes. This function converts this
into a list of nodes according to the parameter type. This function
will always return a copy, to avoid issues e.g. if a child list
of a node should be provided, and a transformation changes the order
in this list (which would then also change the order of the
nodes in the tree).
\begin{quote}\begin{description}
\sphinxlineitem{Parameters}\begin{itemize}
\item {} 
\sphinxAtStartPar
\sphinxstyleliteralstrong{\sphinxupquote{nodes}} (Union{[}\sphinxcode{\sphinxupquote{psyclone.psyir.nodes.Node}},
\sphinxcode{\sphinxupquote{psyclone.psyir.nodes.Schedule}},
List{[}\sphinxcode{\sphinxupquote{psyclone.psyir.nodes.Node}}{]}) \textendash{} can be a single node, a schedule or a list of nodes.

\item {} 
\sphinxAtStartPar
\sphinxstyleliteralstrong{\sphinxupquote{options}} (\sphinxstyleliteralemphasis{\sphinxupquote{Optional}}\sphinxstyleliteralemphasis{\sphinxupquote{{[}}}\sphinxstyleliteralemphasis{\sphinxupquote{Dict}}\sphinxstyleliteralemphasis{\sphinxupquote{{[}}}\sphinxhref{https://docs.python.org/3/library/stdtypes.html\#str}{\sphinxstyleliteralemphasis{\sphinxupquote{str}}}\sphinxstyleliteralemphasis{\sphinxupquote{,}}\sphinxstyleliteralemphasis{\sphinxupquote{Any}}\sphinxstyleliteralemphasis{\sphinxupquote{{]}}}\sphinxstyleliteralemphasis{\sphinxupquote{{]}}}) \textendash{} a dictionary with options for transformations.

\end{itemize}

\sphinxlineitem{Returns}
\sphinxAtStartPar
a list of nodes.

\sphinxlineitem{Return type}
\sphinxAtStartPar
List{[}\sphinxcode{\sphinxupquote{psyclone.psyir.nodes.Node}}{]}

\sphinxlineitem{Raises}
\sphinxAtStartPar
\sphinxstyleliteralstrong{\sphinxupquote{TransformationError}} \textendash{} if the supplied parameter is neither a             single Node, nor a Schedule, nor a list of Nodes.

\end{description}\end{quote}

\end{fulllineitems}



\begin{fulllineitems}

\pysigstartsignatures
\pysigline{\sphinxbfcode{\sphinxupquote{property\DUrole{w,w}{  }}}\sphinxbfcode{\sphinxupquote{omp\_nowait}}}
\pysigstopsignatures\begin{quote}\begin{description}
\sphinxlineitem{Returns}
\sphinxAtStartPar
whether or not this Single region uses a nowait                       clause to remove the end barrier.

\sphinxlineitem{Return type}
\sphinxAtStartPar
\sphinxhref{https://docs.python.org/3/library/functions.html\#bool}{bool}

\end{description}\end{quote}

\end{fulllineitems}



\begin{fulllineitems}

\pysigstartsignatures
\pysiglinewithargsret{\sphinxbfcode{\sphinxupquote{validate}}}{\sphinxparam{\DUrole{n,n}{node\_list}}\sphinxparamcomma \sphinxparam{\DUrole{n,n}{options}\DUrole{o,o}{=}\DUrole{default_value}{None}}}{}
\pysigstopsignatures
\sphinxAtStartPar
Check that the supplied list of Nodes are eligible to be
put inside a parallel region.
\begin{quote}\begin{description}
\sphinxlineitem{Parameters}\begin{itemize}
\item {} 
\sphinxAtStartPar
\sphinxstyleliteralstrong{\sphinxupquote{node\_list}} (\sphinxhref{https://docs.python.org/3/library/stdtypes.html\#list}{\sphinxstyleliteralemphasis{\sphinxupquote{list}}}) \textendash{} list of nodes to put into a parallel region

\item {} 
\sphinxAtStartPar
\sphinxstyleliteralstrong{\sphinxupquote{options}} \textendash{} a dictionary with options for transformations.        :type options: Optional{[}Dict{[}str, Any{]}{]}

\item {} 
\sphinxAtStartPar
\sphinxstyleliteralstrong{\sphinxupquote{options}}\sphinxstyleliteralstrong{\sphinxupquote{{[}}}\sphinxstyleliteralstrong{\sphinxupquote{"node\sphinxhyphen{}type\sphinxhyphen{}check"}}\sphinxstyleliteralstrong{\sphinxupquote{{]}}} (\sphinxhref{https://docs.python.org/3/library/functions.html\#bool}{\sphinxstyleliteralemphasis{\sphinxupquote{bool}}}) \textendash{} this flag controls whether             or not the type of the nodes enclosed in the region should be             tested to avoid using unsupported nodes inside a region.

\end{itemize}

\sphinxlineitem{Raises}\begin{itemize}
\item {} 
\sphinxAtStartPar
\sphinxstyleliteralstrong{\sphinxupquote{TransformationError}} \textendash{} if the supplied node is an             InvokeSchedule rather than being within an InvokeSchedule.

\item {} 
\sphinxAtStartPar
\sphinxstyleliteralstrong{\sphinxupquote{TransformationError}} \textendash{} if the supplied nodes are not all             children of the same parent (siblings).

\end{itemize}

\end{description}\end{quote}

\end{fulllineitems}


\end{fulllineitems}


\begin{sphinxadmonition}{note}{Note:}
\sphinxAtStartPar
PSyclone does not support (distributed\sphinxhyphen{}memory) halo swaps or
global sums within OpenMP single regions.  Attempting to
create a single region for a set of nodes that includes
halo swaps or global sums will produce an error. In such
cases it may be possible to re\sphinxhyphen{}order the nodes in the
Schedule such that the halo swaps or global sums are
performed outside the single region. The
{\hyperref[\detokenize{transformations:sec-move-trans}]{\sphinxcrossref{\DUrole{std,std-ref}{MoveTrans}}}} transformation may be used
for this.
\end{sphinxadmonition}


\bigskip\hrule\bigskip



\begin{fulllineitems}

\pysigstartsignatures
\pysigline{\sphinxbfcode{\sphinxupquote{class\DUrole{w,w}{  }}}\sphinxcode{\sphinxupquote{psyclone.psyir.transformations.}}\sphinxbfcode{\sphinxupquote{OMPTargetTrans}}}
\pysigstopsignatures
\sphinxAtStartPar
Adds an OpenMP target directive to a region of code.

\sphinxAtStartPar
For example:

\begin{sphinxVerbatim}[commandchars=\\\{\}]
\PYG{g+gp}{\PYGZgt{}\PYGZgt{}\PYGZgt{} }\PYG{k+kn}{from}\PYG{+w}{ }\PYG{n+nn}{psyclone}\PYG{n+nn}{.}\PYG{n+nn}{psyir}\PYG{n+nn}{.}\PYG{n+nn}{frontend}\PYG{n+nn}{.}\PYG{n+nn}{fortran}\PYG{+w}{ }\PYG{k+kn}{import} \PYG{n}{FortranReader}
\PYG{g+gp}{\PYGZgt{}\PYGZgt{}\PYGZgt{} }\PYG{k+kn}{from}\PYG{+w}{ }\PYG{n+nn}{psyclone}\PYG{n+nn}{.}\PYG{n+nn}{psyir}\PYG{n+nn}{.}\PYG{n+nn}{backend}\PYG{n+nn}{.}\PYG{n+nn}{fortran}\PYG{+w}{ }\PYG{k+kn}{import} \PYG{n}{FortranWriter}
\PYG{g+gp}{\PYGZgt{}\PYGZgt{}\PYGZgt{} }\PYG{k+kn}{from}\PYG{+w}{ }\PYG{n+nn}{psyclone}\PYG{n+nn}{.}\PYG{n+nn}{psyir}\PYG{n+nn}{.}\PYG{n+nn}{nodes}\PYG{+w}{ }\PYG{k+kn}{import} \PYG{n}{Loop}
\PYG{g+gp}{\PYGZgt{}\PYGZgt{}\PYGZgt{} }\PYG{k+kn}{from}\PYG{+w}{ }\PYG{n+nn}{psyclone}\PYG{n+nn}{.}\PYG{n+nn}{psyir}\PYG{n+nn}{.}\PYG{n+nn}{transformations}\PYG{+w}{ }\PYG{k+kn}{import} \PYG{n}{OMPTargetTrans}
\PYG{g+gp}{\PYGZgt{}\PYGZgt{}\PYGZgt{}}
\PYG{g+gp}{\PYGZgt{}\PYGZgt{}\PYGZgt{} }\PYG{n}{tree} \PYG{o}{=} \PYG{n}{FortranReader}\PYG{p}{(}\PYG{p}{)}\PYG{o}{.}\PYG{n}{psyir\PYGZus{}from\PYGZus{}source}\PYG{p}{(}\PYG{l+s+s2}{\PYGZdq{}\PYGZdq{}\PYGZdq{}}
\PYG{g+gp}{... }\PYG{l+s+s2}{    subroutine my\PYGZus{}subroutine()}
\PYG{g+gp}{... }\PYG{l+s+s2}{        integer, dimension(10, 10) :: A}
\PYG{g+gp}{... }\PYG{l+s+s2}{        integer :: i}
\PYG{g+gp}{... }\PYG{l+s+s2}{        integer :: j}
\PYG{g+gp}{... }\PYG{l+s+s2}{        do i = 1, 10}
\PYG{g+gp}{... }\PYG{l+s+s2}{            do j = 1, 10}
\PYG{g+gp}{... }\PYG{l+s+s2}{                A(i, j) = 0}
\PYG{g+gp}{... }\PYG{l+s+s2}{            end do}
\PYG{g+gp}{... }\PYG{l+s+s2}{        end do}
\PYG{g+gp}{... }\PYG{l+s+s2}{    end subroutine}
\PYG{g+gp}{... }\PYG{l+s+s2}{    }\PYG{l+s+s2}{\PYGZdq{}\PYGZdq{}\PYGZdq{}}\PYG{p}{)}
\PYG{g+gp}{\PYGZgt{}\PYGZgt{}\PYGZgt{} }\PYG{n}{OMPTargetTrans}\PYG{p}{(}\PYG{p}{)}\PYG{o}{.}\PYG{n}{apply}\PYG{p}{(}\PYG{n}{tree}\PYG{o}{.}\PYG{n}{walk}\PYG{p}{(}\PYG{n}{Loop}\PYG{p}{)}\PYG{p}{[}\PYG{l+m+mi}{0}\PYG{p}{]}\PYG{p}{)}
\PYG{g+gp}{\PYGZgt{}\PYGZgt{}\PYGZgt{} }\PYG{n+nb}{print}\PYG{p}{(}\PYG{n}{FortranWriter}\PYG{p}{(}\PYG{p}{)}\PYG{p}{(}\PYG{n}{tree}\PYG{p}{)}\PYG{p}{)}
\PYG{g+go}{subroutine my\PYGZus{}subroutine()}
\PYG{g+go}{  integer, dimension(10,10) :: a}
\PYG{g+go}{  integer :: i}
\PYG{g+go}{  integer :: j}

\PYG{g+go}{  !\PYGZdl{}omp target}
\PYG{g+go}{  do i = 1, 10, 1}
\PYG{g+go}{    do j = 1, 10, 1}
\PYG{g+go}{      a(i,j) = 0}
\PYG{g+go}{    enddo}
\PYG{g+go}{  enddo}
\PYG{g+go}{  !\PYGZdl{}omp end target}

\PYG{g+go}{end subroutine my\PYGZus{}subroutine}
\end{sphinxVerbatim}


\begin{fulllineitems}

\pysigstartsignatures
\pysiglinewithargsret{\sphinxbfcode{\sphinxupquote{apply}}}{\sphinxparam{\DUrole{n,n}{node}}\sphinxparamcomma \sphinxparam{\DUrole{n,n}{options}\DUrole{o,o}{=}\DUrole{default_value}{None}}}{}
\pysigstopsignatures
\sphinxAtStartPar
Insert an OMPTargetDirective before the provided node or list
of nodes.
\begin{quote}\begin{description}
\sphinxlineitem{Parameters}\begin{itemize}
\item {} 
\sphinxAtStartPar
\sphinxstyleliteralstrong{\sphinxupquote{node}} (List{[}\sphinxcode{\sphinxupquote{psyclone.psyir.nodes.Node}}{]}) \textendash{} the PSyIR node or nodes to enclose in the OpenMP
target region.

\item {} 
\sphinxAtStartPar
\sphinxstyleliteralstrong{\sphinxupquote{options}} (\sphinxstyleliteralemphasis{\sphinxupquote{Optional}}\sphinxstyleliteralemphasis{\sphinxupquote{{[}}}\sphinxstyleliteralemphasis{\sphinxupquote{Dict}}\sphinxstyleliteralemphasis{\sphinxupquote{{[}}}\sphinxhref{https://docs.python.org/3/library/stdtypes.html\#str}{\sphinxstyleliteralemphasis{\sphinxupquote{str}}}\sphinxstyleliteralemphasis{\sphinxupquote{,}}\sphinxstyleliteralemphasis{\sphinxupquote{Any}}\sphinxstyleliteralemphasis{\sphinxupquote{{]}}}\sphinxstyleliteralemphasis{\sphinxupquote{{]}}}) \textendash{} a dictionary with options for transformations.

\end{itemize}

\end{description}\end{quote}

\end{fulllineitems}


\end{fulllineitems}



\bigskip\hrule\bigskip



\begin{fulllineitems}

\pysigstartsignatures
\pysiglinewithargsret{\sphinxbfcode{\sphinxupquote{class\DUrole{w,w}{  }}}\sphinxcode{\sphinxupquote{psyclone.transformations.}}\sphinxbfcode{\sphinxupquote{OMPTaskloopTrans}}}{\sphinxparam{\DUrole{n,n}{grainsize}\DUrole{o,o}{=}\DUrole{default_value}{None}}\sphinxparamcomma \sphinxparam{\DUrole{n,n}{num\_tasks}\DUrole{o,o}{=}\DUrole{default_value}{None}}\sphinxparamcomma \sphinxparam{\DUrole{n,n}{nogroup}\DUrole{o,o}{=}\DUrole{default_value}{False}}}{}
\pysigstopsignatures
\sphinxAtStartPar
Adds an OpenMP taskloop directive to a loop. Only one of grainsize or
num\_tasks must be specified.

\sphinxAtStartPar
TODO: \#1364 Taskloops do not yet support reduction clauses.
\begin{quote}\begin{description}
\sphinxlineitem{Parameters}\begin{itemize}
\item {} 
\sphinxAtStartPar
\sphinxstyleliteralstrong{\sphinxupquote{grainsize}} (\sphinxhref{https://docs.python.org/3/library/functions.html\#int}{\sphinxstyleliteralemphasis{\sphinxupquote{int}}}\sphinxstyleliteralemphasis{\sphinxupquote{ or }}\sphinxstyleliteralemphasis{\sphinxupquote{None}}) \textendash{} the grainsize to use in for this transformation.

\item {} 
\sphinxAtStartPar
\sphinxstyleliteralstrong{\sphinxupquote{num\_tasks}} (\sphinxhref{https://docs.python.org/3/library/functions.html\#int}{\sphinxstyleliteralemphasis{\sphinxupquote{int}}}\sphinxstyleliteralemphasis{\sphinxupquote{ or }}\sphinxstyleliteralemphasis{\sphinxupquote{None}}) \textendash{} the num\_tasks to use for this transformation.

\item {} 
\sphinxAtStartPar
\sphinxstyleliteralstrong{\sphinxupquote{nogroup}} (\sphinxhref{https://docs.python.org/3/library/functions.html\#bool}{\sphinxstyleliteralemphasis{\sphinxupquote{bool}}}) \textendash{} whether or not to use a nogroup clause for this
transformation. Default is False.

\end{itemize}

\end{description}\end{quote}

\sphinxAtStartPar
For example:

\begin{sphinxVerbatim}[commandchars=\\\{\}]
\PYG{g+gp}{\PYGZgt{}\PYGZgt{}\PYGZgt{} }\PYG{k+kn}{from}\PYG{+w}{ }\PYG{n+nn}{pysclone}\PYG{n+nn}{.}\PYG{n+nn}{parse}\PYG{n+nn}{.}\PYG{n+nn}{algorithm}\PYG{+w}{ }\PYG{k+kn}{import} \PYG{n}{parse}
\PYG{g+gp}{\PYGZgt{}\PYGZgt{}\PYGZgt{} }\PYG{k+kn}{from}\PYG{+w}{ }\PYG{n+nn}{psyclone}\PYG{n+nn}{.}\PYG{n+nn}{psyGen}\PYG{+w}{ }\PYG{k+kn}{import} \PYG{n}{PSyFactory}
\PYG{g+gp}{\PYGZgt{}\PYGZgt{}\PYGZgt{} }\PYG{n}{api} \PYG{o}{=} \PYG{l+s+s2}{\PYGZdq{}}\PYG{l+s+s2}{gocean}\PYG{l+s+s2}{\PYGZdq{}}
\PYG{g+gp}{\PYGZgt{}\PYGZgt{}\PYGZgt{} }\PYG{n}{ast}\PYG{p}{,} \PYG{n}{invokeInfo} \PYG{o}{=} \PYG{n}{parse}\PYG{p}{(}\PYG{n}{GOCEAN\PYGZus{}SOURCE\PYGZus{}FILE}\PYG{p}{,} \PYG{n}{api}\PYG{o}{=}\PYG{n}{api}\PYG{p}{)}
\PYG{g+gp}{\PYGZgt{}\PYGZgt{}\PYGZgt{} }\PYG{n}{psy} \PYG{o}{=} \PYG{n}{PSyFactory}\PYG{p}{(}\PYG{n}{api}\PYG{p}{)}\PYG{o}{.}\PYG{n}{create}\PYG{p}{(}\PYG{n}{invokeInfo}\PYG{p}{)}
\PYG{g+gp}{\PYGZgt{}\PYGZgt{}\PYGZgt{}}
\PYG{g+gp}{\PYGZgt{}\PYGZgt{}\PYGZgt{} }\PYG{k+kn}{from}\PYG{+w}{ }\PYG{n+nn}{psyclone}\PYG{n+nn}{.}\PYG{n+nn}{transformations}\PYG{+w}{ }\PYG{k+kn}{import} \PYG{n}{OMPParallelTrans}\PYG{p}{,} \PYG{n}{OMPSingleTrans}
\PYG{g+gp}{\PYGZgt{}\PYGZgt{}\PYGZgt{} }\PYG{k+kn}{from}\PYG{+w}{ }\PYG{n+nn}{psyclone}\PYG{n+nn}{.}\PYG{n+nn}{transformations}\PYG{+w}{ }\PYG{k+kn}{import} \PYG{n}{OMPTaskloopTrans}
\PYG{g+gp}{\PYGZgt{}\PYGZgt{}\PYGZgt{} }\PYG{k+kn}{from}\PYG{+w}{ }\PYG{n+nn}{psyclone}\PYG{n+nn}{.}\PYG{n+nn}{psyir}\PYG{n+nn}{.}\PYG{n+nn}{transformations}\PYG{+w}{ }\PYG{k+kn}{import} \PYG{n}{OMPTaskwaitTrans}
\PYG{g+gp}{\PYGZgt{}\PYGZgt{}\PYGZgt{} }\PYG{n}{singletrans} \PYG{o}{=} \PYG{n}{OMPSingleTrans}\PYG{p}{(}\PYG{p}{)}
\PYG{g+gp}{\PYGZgt{}\PYGZgt{}\PYGZgt{} }\PYG{n}{paralleltrans} \PYG{o}{=} \PYG{n}{OMPParallelTrans}\PYG{p}{(}\PYG{p}{)}
\PYG{g+gp}{\PYGZgt{}\PYGZgt{}\PYGZgt{} }\PYG{n}{tasklooptrans} \PYG{o}{=} \PYG{n}{OMPTaskloopTrans}\PYG{p}{(}\PYG{p}{)}
\PYG{g+gp}{\PYGZgt{}\PYGZgt{}\PYGZgt{} }\PYG{n}{taskwaittrans} \PYG{o}{=} \PYG{n}{OMPTaskwaitTrans}\PYG{p}{(}\PYG{p}{)}
\PYG{g+gp}{\PYGZgt{}\PYGZgt{}\PYGZgt{}}
\PYG{g+gp}{\PYGZgt{}\PYGZgt{}\PYGZgt{} }\PYG{n}{schedule} \PYG{o}{=} \PYG{n}{psy}\PYG{o}{.}\PYG{n}{invokes}\PYG{o}{.}\PYG{n}{get}\PYG{p}{(}\PYG{l+s+s1}{\PYGZsq{}}\PYG{l+s+s1}{invoke\PYGZus{}0}\PYG{l+s+s1}{\PYGZsq{}}\PYG{p}{)}\PYG{o}{.}\PYG{n}{schedule}
\PYG{g+gp}{\PYGZgt{}\PYGZgt{}\PYGZgt{} }\PYG{c+c1}{\PYGZsh{} Uncomment the following line to see a text view of the schedule}
\PYG{g+gp}{\PYGZgt{}\PYGZgt{}\PYGZgt{} }\PYG{c+c1}{\PYGZsh{} print(schedule.view())}
\PYG{g+gp}{\PYGZgt{}\PYGZgt{}\PYGZgt{}}
\PYG{g+gp}{\PYGZgt{}\PYGZgt{}\PYGZgt{} }\PYG{c+c1}{\PYGZsh{} Apply the OpenMP Taskloop transformation to *every* loop}
\PYG{g+gp}{\PYGZgt{}\PYGZgt{}\PYGZgt{} }\PYG{c+c1}{\PYGZsh{} in the schedule.}
\PYG{g+gp}{\PYGZgt{}\PYGZgt{}\PYGZgt{} }\PYG{c+c1}{\PYGZsh{} This ignores loop dependencies. These can be handled}
\PYG{g+gp}{\PYGZgt{}\PYGZgt{}\PYGZgt{} }\PYG{c+c1}{\PYGZsh{} by the OMPTaskwaitTrans}
\PYG{g+gp}{\PYGZgt{}\PYGZgt{}\PYGZgt{} }\PYG{k}{for} \PYG{n}{child} \PYG{o+ow}{in} \PYG{n}{schedule}\PYG{o}{.}\PYG{n}{children}\PYG{p}{:}
\PYG{g+gp}{\PYGZgt{}\PYGZgt{}\PYGZgt{} }    \PYG{n}{tasklooptrans}\PYG{o}{.}\PYG{n}{apply}\PYG{p}{(}\PYG{n}{child}\PYG{p}{)}
\PYG{g+gp}{\PYGZgt{}\PYGZgt{}\PYGZgt{} }\PYG{c+c1}{\PYGZsh{} Enclose all of these loops within a single OpenMP}
\PYG{g+gp}{\PYGZgt{}\PYGZgt{}\PYGZgt{} }\PYG{c+c1}{\PYGZsh{} SINGLE region}
\PYG{g+gp}{\PYGZgt{}\PYGZgt{}\PYGZgt{} }\PYG{n}{singletrans}\PYG{o}{.}\PYG{n}{apply}\PYG{p}{(}\PYG{n}{schedule}\PYG{o}{.}\PYG{n}{children}\PYG{p}{)}
\PYG{g+gp}{\PYGZgt{}\PYGZgt{}\PYGZgt{} }\PYG{c+c1}{\PYGZsh{} Enclose all of these loops within a single OpenMP}
\PYG{g+gp}{\PYGZgt{}\PYGZgt{}\PYGZgt{} }\PYG{c+c1}{\PYGZsh{} PARALLEL region}
\PYG{g+gp}{\PYGZgt{}\PYGZgt{}\PYGZgt{} }\PYG{n}{paralleltrans}\PYG{o}{.}\PYG{n}{apply}\PYG{p}{(}\PYG{n}{schedule}\PYG{o}{.}\PYG{n}{children}\PYG{p}{)}
\PYG{g+gp}{\PYGZgt{}\PYGZgt{}\PYGZgt{} }\PYG{c+c1}{\PYGZsh{} Ensure loop dependencies are satisfied}
\PYG{g+gp}{\PYGZgt{}\PYGZgt{}\PYGZgt{} }\PYG{n}{taskwaittrans}\PYG{o}{.}\PYG{n}{apply}\PYG{p}{(}\PYG{n}{schedule}\PYG{o}{.}\PYG{n}{children}\PYG{p}{)}
\PYG{g+gp}{\PYGZgt{}\PYGZgt{}\PYGZgt{} }\PYG{c+c1}{\PYGZsh{} Uncomment the following line to see a text view of the schedule}
\PYG{g+gp}{\PYGZgt{}\PYGZgt{}\PYGZgt{} }\PYG{c+c1}{\PYGZsh{} print(schedule.view())}
\end{sphinxVerbatim}


\begin{fulllineitems}

\pysigstartsignatures
\pysiglinewithargsret{\sphinxbfcode{\sphinxupquote{apply}}}{\sphinxparam{\DUrole{n,n}{node}}\sphinxparamcomma \sphinxparam{\DUrole{n,n}{options}\DUrole{o,o}{=}\DUrole{default_value}{None}}}{}
\pysigstopsignatures
\sphinxAtStartPar
Apply the OMPTaskloopTrans transformation to the specified node in
a Schedule. This node must be a Loop since this transformation
corresponds to wrapping the generated code with directives like so:

\begin{sphinxVerbatim}[commandchars=\\\{\}]
\PYG{c}{!\PYGZdl{}OMP TASKLOOP}
\PYG{k}{do}\PYG{+w}{ }\PYG{p}{.}\PYG{p}{.}\PYG{p}{.}
\PYG{+w}{   }\PYG{p}{.}\PYG{p}{.}\PYG{p}{.}
\PYG{k}{end }\PYG{k}{do}
\PYG{c}{!\PYGZdl{}OMP END TASKLOOP}
\end{sphinxVerbatim}

\sphinxAtStartPar
At code\sphinxhyphen{}generation time (when
\sphinxcode{\sphinxupquote{OMPTaskloopDirective.gen\_code()}} is called), this node must be
within (i.e. a child of) an OpenMP SERIAL region.

\sphinxAtStartPar
If the keyword “nogroup” is specified in the options, it will cause a
nogroup clause be generated if it is set to True. This will override
the value supplied to the constructor, but will only apply to the
apply call to which the value is supplied.
\begin{quote}\begin{description}
\sphinxlineitem{Parameters}\begin{itemize}
\item {} 
\sphinxAtStartPar
\sphinxstyleliteralstrong{\sphinxupquote{node}} (\sphinxcode{\sphinxupquote{psyclone.psyir.nodes.Node}}) \textendash{} the supplied node to which we will apply the                      OMPTaskloopTrans transformation

\item {} 
\sphinxAtStartPar
\sphinxstyleliteralstrong{\sphinxupquote{options}} (\sphinxstyleliteralemphasis{\sphinxupquote{Optional}}\sphinxstyleliteralemphasis{\sphinxupquote{{[}}}\sphinxstyleliteralemphasis{\sphinxupquote{Dict}}\sphinxstyleliteralemphasis{\sphinxupquote{{[}}}\sphinxhref{https://docs.python.org/3/library/stdtypes.html\#str}{\sphinxstyleliteralemphasis{\sphinxupquote{str}}}\sphinxstyleliteralemphasis{\sphinxupquote{, }}\sphinxstyleliteralemphasis{\sphinxupquote{Any}}\sphinxstyleliteralemphasis{\sphinxupquote{{]}}}\sphinxstyleliteralemphasis{\sphinxupquote{{]}}}) \textendash{} a dictionary with options for transformations                        and validation.

\item {} 
\sphinxAtStartPar
\sphinxstyleliteralstrong{\sphinxupquote{options}}\sphinxstyleliteralstrong{\sphinxupquote{{[}}}\sphinxstyleliteralstrong{\sphinxupquote{"nogroup"}}\sphinxstyleliteralstrong{\sphinxupquote{{]}}} (\sphinxhref{https://docs.python.org/3/library/functions.html\#bool}{\sphinxstyleliteralemphasis{\sphinxupquote{bool}}}) \textendash{} indicating whether a nogroup clause should be applied to
this taskloop.

\end{itemize}

\end{description}\end{quote}

\end{fulllineitems}



\begin{fulllineitems}

\pysigstartsignatures
\pysigline{\sphinxbfcode{\sphinxupquote{property\DUrole{w,w}{  }}}\sphinxbfcode{\sphinxupquote{omp\_grainsize}}}
\pysigstopsignatures
\sphinxAtStartPar
Returns the grainsize that will be specified by
this transformation. By default the grainsize
clause is not applied, so grainsize is None.
\begin{quote}\begin{description}
\sphinxlineitem{Returns}
\sphinxAtStartPar
The grainsize specified by this transformation.

\sphinxlineitem{Return type}
\sphinxAtStartPar
\sphinxhref{https://docs.python.org/3/library/functions.html\#int}{int} or None

\end{description}\end{quote}

\end{fulllineitems}



\begin{fulllineitems}

\pysigstartsignatures
\pysigline{\sphinxbfcode{\sphinxupquote{property\DUrole{w,w}{  }}}\sphinxbfcode{\sphinxupquote{omp\_num\_tasks}}}
\pysigstopsignatures
\sphinxAtStartPar
Returns the num\_tasks that will be specified
by this transformation. By default the num\_tasks
clause is not applied so num\_tasks is None.
\begin{quote}\begin{description}
\sphinxlineitem{Returns}
\sphinxAtStartPar
The grainsize specified by this transformation.

\sphinxlineitem{Return type}
\sphinxAtStartPar
\sphinxhref{https://docs.python.org/3/library/functions.html\#int}{int} or None

\end{description}\end{quote}

\end{fulllineitems}


\end{fulllineitems}



\bigskip\hrule\bigskip



\begin{fulllineitems}

\pysigstartsignatures
\pysigline{\sphinxbfcode{\sphinxupquote{class\DUrole{w,w}{  }}}\sphinxcode{\sphinxupquote{psyclone.psyir.transformations.}}\sphinxbfcode{\sphinxupquote{OMPTaskTrans}}}
\pysigstopsignatures
\sphinxAtStartPar
Apply an OpenMP Task Transformation to a Loop. The Loop must
be within an OpenMP Serial region (Single or Master) at codegen time.
Once lowering begins, no more modifications to the tree should occur
as the task directives do not recompute dependencies after lowering.
In the future it may be possible to do this through an \_update\_node
implementation.


\begin{fulllineitems}

\pysigstartsignatures
\pysiglinewithargsret{\sphinxbfcode{\sphinxupquote{apply}}}{\sphinxparam{\DUrole{n,n}{node}}\sphinxparamcomma \sphinxparam{\DUrole{n,n}{options}\DUrole{o,o}{=}\DUrole{default_value}{None}}}{}
\pysigstopsignatures
\sphinxAtStartPar
Apply the OMPTaskTrans to the specified node in a Schedule.

\sphinxAtStartPar
Can only be applied to a Loop.

\sphinxAtStartPar
The specified node is wrapped by directives during code generation
like so:

\begin{sphinxVerbatim}[commandchars=\\\{\}]
\PYG{c}{!\PYGZdl{}OMP TASK}
\PYG{p}{.}\PYG{p}{.}\PYG{p}{.}
\PYG{c}{!\PYGZdl{}OMP END TASK}
\end{sphinxVerbatim}

\sphinxAtStartPar
At code\sphinxhyphen{}generation time, this node must be
within (i.e. a child of) an OpenMP Serial region (OpenMP Single or
OpenMP Master)

\sphinxAtStartPar
Any kernels or Calls will be inlined into the region before the task
transformation is applied.
\begin{quote}\begin{description}
\sphinxlineitem{Parameters}\begin{itemize}
\item {} 
\sphinxAtStartPar
\sphinxstyleliteralstrong{\sphinxupquote{node}} (\sphinxcode{\sphinxupquote{psyclone.psyir.nodes.Loop}}) \textendash{} the supplied node to which we will apply the                      OMPTaskTrans transformation

\item {} 
\sphinxAtStartPar
\sphinxstyleliteralstrong{\sphinxupquote{options}} (\sphinxstyleliteralemphasis{\sphinxupquote{dictionary}}\sphinxstyleliteralemphasis{\sphinxupquote{ of }}\sphinxstyleliteralemphasis{\sphinxupquote{string:values}}\sphinxstyleliteralemphasis{\sphinxupquote{ or }}\sphinxstyleliteralemphasis{\sphinxupquote{None}}) \textendash{} a dictionary with options for transformations                        and validation.

\end{itemize}

\end{description}\end{quote}

\end{fulllineitems}


\end{fulllineitems}



\bigskip\hrule\bigskip



\begin{fulllineitems}

\pysigstartsignatures
\pysigline{\sphinxbfcode{\sphinxupquote{class\DUrole{w,w}{  }}}\sphinxcode{\sphinxupquote{psyclone.psyir.transformations.}}\sphinxbfcode{\sphinxupquote{OMPTaskwaitTrans}}}
\pysigstopsignatures
\sphinxAtStartPar
Adds zero or more OpenMP Taskwait directives to an OMP parallel region.
This transformation will add directives to satisfy dependencies between
Taskloop directives without an associated taskgroup (i.e. no nogroup
clause). It also tries to minimise the number added to maximise available
parallelism.

\sphinxAtStartPar
For example:

\begin{sphinxVerbatim}[commandchars=\\\{\}]
\PYG{g+gp}{\PYGZgt{}\PYGZgt{}\PYGZgt{} }\PYG{k+kn}{from}\PYG{+w}{ }\PYG{n+nn}{pysclone}\PYG{n+nn}{.}\PYG{n+nn}{parse}\PYG{n+nn}{.}\PYG{n+nn}{algorithm}\PYG{+w}{ }\PYG{k+kn}{import} \PYG{n}{parse}
\PYG{g+gp}{\PYGZgt{}\PYGZgt{}\PYGZgt{} }\PYG{k+kn}{from}\PYG{+w}{ }\PYG{n+nn}{psyclone}\PYG{n+nn}{.}\PYG{n+nn}{psyGen}\PYG{+w}{ }\PYG{k+kn}{import} \PYG{n}{PSyFactory}
\PYG{g+gp}{\PYGZgt{}\PYGZgt{}\PYGZgt{} }\PYG{n}{api} \PYG{o}{=} \PYG{l+s+s2}{\PYGZdq{}}\PYG{l+s+s2}{gocean}\PYG{l+s+s2}{\PYGZdq{}}
\PYG{g+gp}{\PYGZgt{}\PYGZgt{}\PYGZgt{} }\PYG{n}{filename} \PYG{o}{=} \PYG{l+s+s2}{\PYGZdq{}}\PYG{l+s+s2}{nemolite2d\PYGZus{}alg.f90}\PYG{l+s+s2}{\PYGZdq{}}
\PYG{g+gp}{\PYGZgt{}\PYGZgt{}\PYGZgt{} }\PYG{n}{ast}\PYG{p}{,} \PYG{n}{invokeInfo} \PYG{o}{=} \PYG{n}{parse}\PYG{p}{(}\PYG{n}{filename}\PYG{p}{,} \PYG{n}{api}\PYG{o}{=}\PYG{n}{api}\PYG{p}{,} \PYG{n}{invoke\PYGZus{}name}\PYG{o}{=}\PYG{l+s+s2}{\PYGZdq{}}\PYG{l+s+s2}{invoke}\PYG{l+s+s2}{\PYGZdq{}}\PYG{p}{)}
\PYG{g+gp}{\PYGZgt{}\PYGZgt{}\PYGZgt{} }\PYG{n}{psy} \PYG{o}{=} \PYG{n}{PSyFactory}\PYG{p}{(}\PYG{n}{api}\PYG{p}{)}\PYG{o}{.}\PYG{n}{create}\PYG{p}{(}\PYG{n}{invokeInfo}\PYG{p}{)}
\PYG{g+gp}{\PYGZgt{}\PYGZgt{}\PYGZgt{}}
\PYG{g+gp}{\PYGZgt{}\PYGZgt{}\PYGZgt{} }\PYG{k+kn}{from}\PYG{+w}{ }\PYG{n+nn}{psyclone}\PYG{n+nn}{.}\PYG{n+nn}{transformations}\PYG{+w}{ }\PYG{k+kn}{import} \PYG{n}{OMPParallelTrans}\PYG{p}{,} \PYG{n}{OMPSingleTrans}
\PYG{g+gp}{\PYGZgt{}\PYGZgt{}\PYGZgt{} }\PYG{k+kn}{from}\PYG{+w}{ }\PYG{n+nn}{psyclone}\PYG{n+nn}{.}\PYG{n+nn}{transformations}\PYG{+w}{ }\PYG{k+kn}{import} \PYG{n}{OMPTaskloopTrans}
\PYG{g+gp}{\PYGZgt{}\PYGZgt{}\PYGZgt{} }\PYG{k+kn}{from}\PYG{+w}{ }\PYG{n+nn}{psyclone}\PYG{n+nn}{.}\PYG{n+nn}{psyir}\PYG{n+nn}{.}\PYG{n+nn}{transformations}\PYG{+w}{ }\PYG{k+kn}{import} \PYG{n}{OMPTaskwaitTrans}
\PYG{g+gp}{\PYGZgt{}\PYGZgt{}\PYGZgt{} }\PYG{n}{singletrans} \PYG{o}{=} \PYG{n}{OMPSingleTrans}\PYG{p}{(}\PYG{p}{)}
\PYG{g+gp}{\PYGZgt{}\PYGZgt{}\PYGZgt{} }\PYG{n}{paralleltrans} \PYG{o}{=} \PYG{n}{OMPParallelTrans}\PYG{p}{(}\PYG{p}{)}
\PYG{g+gp}{\PYGZgt{}\PYGZgt{}\PYGZgt{} }\PYG{n}{tasklooptrans} \PYG{o}{=} \PYG{n}{OMPTaskloopTrans}\PYG{p}{(}\PYG{p}{)}
\PYG{g+gp}{\PYGZgt{}\PYGZgt{}\PYGZgt{} }\PYG{n}{taskwaittrans} \PYG{o}{=} \PYG{n}{OMPTaskwaitTrans}\PYG{p}{(}\PYG{p}{)}
\PYG{g+gp}{\PYGZgt{}\PYGZgt{}\PYGZgt{}}
\PYG{g+gp}{\PYGZgt{}\PYGZgt{}\PYGZgt{} }\PYG{n}{schedule} \PYG{o}{=} \PYG{n}{psy}\PYG{o}{.}\PYG{n}{invokes}\PYG{o}{.}\PYG{n}{get}\PYG{p}{(}\PYG{l+s+s1}{\PYGZsq{}}\PYG{l+s+s1}{invoke\PYGZus{}0}\PYG{l+s+s1}{\PYGZsq{}}\PYG{p}{)}\PYG{o}{.}\PYG{n}{schedule}
\PYG{g+gp}{\PYGZgt{}\PYGZgt{}\PYGZgt{} }\PYG{n+nb}{print}\PYG{p}{(}\PYG{n}{schedule}\PYG{o}{.}\PYG{n}{view}\PYG{p}{(}\PYG{p}{)}\PYG{p}{)}
\PYG{g+gp}{\PYGZgt{}\PYGZgt{}\PYGZgt{}}
\PYG{g+gp}{\PYGZgt{}\PYGZgt{}\PYGZgt{} }\PYG{c+c1}{\PYGZsh{} Apply the OpenMP Taskloop transformation to *every* loop}
\PYG{g+gp}{\PYGZgt{}\PYGZgt{}\PYGZgt{} }\PYG{c+c1}{\PYGZsh{} in the schedule.}
\PYG{g+gp}{\PYGZgt{}\PYGZgt{}\PYGZgt{} }\PYG{c+c1}{\PYGZsh{} This ignores loop dependencies. These are handled by the}
\PYG{g+gp}{\PYGZgt{}\PYGZgt{}\PYGZgt{} }\PYG{c+c1}{\PYGZsh{} taskwait transformation.}
\PYG{g+gp}{\PYGZgt{}\PYGZgt{}\PYGZgt{} }\PYG{k}{for} \PYG{n}{child} \PYG{o+ow}{in} \PYG{n}{schedule}\PYG{o}{.}\PYG{n}{children}\PYG{p}{:}
\PYG{g+gp}{\PYGZgt{}\PYGZgt{}\PYGZgt{} }    \PYG{n}{tasklooptrans}\PYG{o}{.}\PYG{n}{apply}\PYG{p}{(}\PYG{n}{child}\PYG{p}{,} \PYG{n}{nogroup} \PYG{o}{=} \PYG{n}{true}\PYG{p}{)}
\PYG{g+gp}{\PYGZgt{}\PYGZgt{}\PYGZgt{} }\PYG{c+c1}{\PYGZsh{} Enclose all of these loops within a single OpenMP}
\PYG{g+gp}{\PYGZgt{}\PYGZgt{}\PYGZgt{} }\PYG{c+c1}{\PYGZsh{} SINGLE region}
\PYG{g+gp}{\PYGZgt{}\PYGZgt{}\PYGZgt{} }\PYG{n}{singletrans}\PYG{o}{.}\PYG{n}{apply}\PYG{p}{(}\PYG{n}{schedule}\PYG{o}{.}\PYG{n}{children}\PYG{p}{)}
\PYG{g+gp}{\PYGZgt{}\PYGZgt{}\PYGZgt{} }\PYG{c+c1}{\PYGZsh{} Enclose all of these loops within a single OpenMP}
\PYG{g+gp}{\PYGZgt{}\PYGZgt{}\PYGZgt{} }\PYG{c+c1}{\PYGZsh{} PARALLEL region}
\PYG{g+gp}{\PYGZgt{}\PYGZgt{}\PYGZgt{} }\PYG{n}{paralleltrans}\PYG{o}{.}\PYG{n}{apply}\PYG{p}{(}\PYG{n}{schedule}\PYG{o}{.}\PYG{n}{children}\PYG{p}{)}
\PYG{g+gp}{\PYGZgt{}\PYGZgt{}\PYGZgt{} }\PYG{n}{taskwaittrans}\PYG{o}{.}\PYG{n}{apply}\PYG{p}{(}\PYG{n}{schedule}\PYG{o}{.}\PYG{n}{children}\PYG{p}{)}
\PYG{g+gp}{\PYGZgt{}\PYGZgt{}\PYGZgt{} }\PYG{n+nb}{print}\PYG{p}{(}\PYG{n}{schedule}\PYG{o}{.}\PYG{n}{view}\PYG{p}{(}\PYG{p}{)}\PYG{p}{)}
\end{sphinxVerbatim}


\begin{fulllineitems}

\pysigstartsignatures
\pysiglinewithargsret{\sphinxbfcode{\sphinxupquote{apply}}}{\sphinxparam{\DUrole{n,n}{node}}\sphinxparamcomma \sphinxparam{\DUrole{n,n}{options}\DUrole{o,o}{=}\DUrole{default_value}{None}}}{}
\pysigstopsignatures
\sphinxAtStartPar
Apply an OMPTaskwait Transformation to the supplied node
(which must be an OMPParallelDirective). In the generated code this
corresponds to adding zero or more OMPTaskwaitDirectives as
appropriate:

\begin{sphinxVerbatim}[commandchars=\\\{\}]
\PYG{c}{!\PYGZdl{}OMP PARALLEL}
\PYG{+w}{  }\PYG{p}{.}\PYG{p}{.}\PYG{p}{.}
\PYG{+w}{  }\PYG{c}{!\PYGZdl{}OMP TASKWAIT}
\PYG{+w}{  }\PYG{p}{.}\PYG{p}{.}\PYG{p}{.}
\PYG{+w}{  }\PYG{c}{!\PYGZdl{}OMP TASKWAIT}
\PYG{+w}{  }\PYG{p}{.}\PYG{p}{.}\PYG{p}{.}
\PYG{c}{!\PYGZdl{}OMP END PARALLEL}
\end{sphinxVerbatim}
\begin{quote}\begin{description}
\sphinxlineitem{Parameters}\begin{itemize}
\item {} 
\sphinxAtStartPar
\sphinxstyleliteralstrong{\sphinxupquote{node}} (\sphinxcode{\sphinxupquote{psyclone.psyir.nodes.OMPParallelDirective}}) \textendash{} the node to which to apply the transformation.

\item {} 
\sphinxAtStartPar
\sphinxstyleliteralstrong{\sphinxupquote{options}} (\sphinxstyleliteralemphasis{\sphinxupquote{Optional}}\sphinxstyleliteralemphasis{\sphinxupquote{{[}}}\sphinxstyleliteralemphasis{\sphinxupquote{Dict}}\sphinxstyleliteralemphasis{\sphinxupquote{{[}}}\sphinxhref{https://docs.python.org/3/library/stdtypes.html\#str}{\sphinxstyleliteralemphasis{\sphinxupquote{str}}}\sphinxstyleliteralemphasis{\sphinxupquote{, }}\sphinxstyleliteralemphasis{\sphinxupquote{Any}}\sphinxstyleliteralemphasis{\sphinxupquote{{]}}}\sphinxstyleliteralemphasis{\sphinxupquote{{]}}}) \textendash{} a dictionary with options for transformations                        and validation.

\item {} 
\sphinxAtStartPar
\sphinxstyleliteralstrong{\sphinxupquote{options}}\sphinxstyleliteralstrong{\sphinxupquote{{[}}}\sphinxstyleliteralstrong{\sphinxupquote{"fail\_on\_no\_taskloop"}}\sphinxstyleliteralstrong{\sphinxupquote{{]}}} (\sphinxhref{https://docs.python.org/3/library/functions.html\#bool}{\sphinxstyleliteralemphasis{\sphinxupquote{bool}}}) \textendash{} indicating whether this should throw an error if no                 OMPTaskloop nodes are found in this tree. This can be                 safely disabled as if there are no Taskloop nodes the                 result of this transformation is valid OpenMP code. Default                 is True

\end{itemize}

\end{description}\end{quote}

\end{fulllineitems}


\end{fulllineitems}



\bigskip\hrule\bigskip



\begin{fulllineitems}

\pysigstartsignatures
\pysigline{\sphinxbfcode{\sphinxupquote{class\DUrole{w,w}{  }}}\sphinxcode{\sphinxupquote{psyclone.psyir.transformations.}}\sphinxbfcode{\sphinxupquote{Product2LoopTrans}}}
\pysigstopsignatures
\sphinxAtStartPar
Provides a transformation from a PSyIR PRODUCT IntrinsicCall node to
an equivalent PSyIR loop structure that is suitable for running in
parallel on CPUs and GPUs. Validity checks are also performed.

\sphinxAtStartPar
If PRODUCT contains a single positional argument which is an array,
the maximum value of all of the elements in the array is returned
in the the scalar R.

\begin{sphinxVerbatim}[commandchars=\\\{\}]
\PYG{n}{R}\PYG{+w}{ }\PYG{o}{=}\PYG{+w}{ }\PYG{n+nb}{PRODUCT}\PYG{p}{(}\PYG{k}{ARRAY}\PYG{p}{)}
\end{sphinxVerbatim}

\sphinxAtStartPar
For example, if the array is two dimensional, the equivalent code
for real data is:

\begin{sphinxVerbatim}[commandchars=\\\{\}]
\PYG{n}{R}\PYG{+w}{ }\PYG{o}{=}\PYG{+w}{ }\PYG{l+m+mf}{1.0}
\PYG{k}{DO }\PYG{n}{J}\PYG{o}{=}\PYG{n+nb}{LBOUND}\PYG{p}{(}\PYG{k}{ARRAY}\PYG{p}{,}\PYG{l+m+mi}{2}\PYG{p}{)}\PYG{p}{,}\PYG{n+nb}{UBOUND}\PYG{p}{(}\PYG{k}{ARRAY}\PYG{p}{,}\PYG{l+m+mi}{2}\PYG{p}{)}
\PYG{+w}{  }\PYG{k}{DO }\PYG{n}{I}\PYG{o}{=}\PYG{n+nb}{LBOUND}\PYG{p}{(}\PYG{k}{ARRAY}\PYG{p}{,}\PYG{l+m+mi}{1}\PYG{p}{)}\PYG{p}{,}\PYG{n+nb}{UBOUND}\PYG{p}{(}\PYG{k}{ARRAY}\PYG{p}{,}\PYG{l+m+mi}{1}\PYG{p}{)}
\PYG{+w}{    }\PYG{n}{R}\PYG{+w}{ }\PYG{o}{=}\PYG{+w}{ }\PYG{n}{R}\PYG{+w}{ }\PYG{o}{*}\PYG{+w}{ }\PYG{k}{ARRAY}\PYG{p}{(}\PYG{n}{I}\PYG{p}{,}\PYG{n}{J}\PYG{p}{)}
\end{sphinxVerbatim}

\sphinxAtStartPar
If the mask argument is provided then the mask is used to
determine whether the product is applied:

\begin{sphinxVerbatim}[commandchars=\\\{\}]
\PYG{n}{R}\PYG{+w}{ }\PYG{o}{=}\PYG{+w}{ }\PYG{n+nb}{PRODUCT}\PYG{p}{(}\PYG{k}{ARRAY}\PYG{p}{,}\PYG{+w}{ }\PYG{n}{mask}\PYG{o}{=}\PYG{n+nb}{MOD}\PYG{p}{(}\PYG{k}{ARRAY}\PYG{p}{,}\PYG{+w}{ }\PYG{l+m+mf}{2.0}\PYG{p}{)}\PYG{o}{==}\PYG{l+m+mi}{1}\PYG{p}{)}
\end{sphinxVerbatim}

\sphinxAtStartPar
If the array is two dimensional, the equivalent code
for real data is:

\begin{sphinxVerbatim}[commandchars=\\\{\}]
\PYG{n}{R}\PYG{+w}{ }\PYG{o}{=}\PYG{+w}{ }\PYG{l+m+mf}{1.0}
\PYG{k}{DO }\PYG{n}{J}\PYG{o}{=}\PYG{n+nb}{LBOUND}\PYG{p}{(}\PYG{k}{ARRAY}\PYG{p}{,}\PYG{l+m+mi}{2}\PYG{p}{)}\PYG{p}{,}\PYG{n+nb}{UBOUND}\PYG{p}{(}\PYG{k}{ARRAY}\PYG{p}{,}\PYG{l+m+mi}{2}\PYG{p}{)}
\PYG{+w}{  }\PYG{k}{DO }\PYG{n}{I}\PYG{o}{=}\PYG{n+nb}{LBOUND}\PYG{p}{(}\PYG{k}{ARRAY}\PYG{p}{,}\PYG{l+m+mi}{1}\PYG{p}{)}\PYG{p}{,}\PYG{n+nb}{UBOUND}\PYG{p}{(}\PYG{k}{ARRAY}\PYG{p}{,}\PYG{l+m+mi}{1}\PYG{p}{)}
\PYG{+w}{    }\PYG{k}{IF}\PYG{+w}{ }\PYG{p}{(}\PYG{n+nb}{MOD}\PYG{p}{(}\PYG{k}{ARRAY}\PYG{p}{(}\PYG{n}{I}\PYG{p}{,}\PYG{n}{J}\PYG{p}{)}\PYG{p}{,}\PYG{+w}{ }\PYG{l+m+mf}{2.0}\PYG{p}{)}\PYG{o}{==}\PYG{l+m+mi}{1}\PYG{p}{)}\PYG{+w}{ }\PYG{k}{THEN}
\PYG{k}{      }\PYG{n}{R}\PYG{+w}{ }\PYG{o}{=}\PYG{+w}{ }\PYG{n}{R}\PYG{+w}{ }\PYG{o}{*}\PYG{+w}{ }\PYG{k}{ARRAY}\PYG{p}{(}\PYG{n}{I}\PYG{p}{,}\PYG{n}{J}\PYG{p}{)}
\end{sphinxVerbatim}

\sphinxAtStartPar
The dimension argument is currently not supported and will result
in a TransformationError exception being raised.

\begin{sphinxVerbatim}[commandchars=\\\{\}]
\PYG{n}{R}\PYG{+w}{ }\PYG{o}{=}\PYG{+w}{ }\PYG{n+nb}{PRODUCT}\PYG{p}{(}\PYG{k}{ARRAY}\PYG{p}{,}\PYG{+w}{ }\PYG{k}{dimension}\PYG{o}{=}\PYG{l+m+mi}{2}\PYG{p}{)}
\end{sphinxVerbatim}

\sphinxAtStartPar
The array passed to PRODUCT may use any combination of array
syntax, array notation, array sections and scalar bounds:

\begin{sphinxVerbatim}[commandchars=\\\{\}]
\PYG{n}{R}\PYG{+w}{ }\PYG{o}{=}\PYG{+w}{ }\PYG{n+nb}{PRODUCT}\PYG{p}{(}\PYG{k}{ARRAY}\PYG{p}{)}\PYG{+w}{ }\PYG{c}{! array syntax}
\PYG{n}{R}\PYG{+w}{ }\PYG{o}{=}\PYG{+w}{ }\PYG{n+nb}{PRODUCT}\PYG{p}{(}\PYG{k}{ARRAY}\PYG{p}{(}\PYG{p}{:}\PYG{p}{,}\PYG{p}{:}\PYG{p}{)}\PYG{p}{)}\PYG{+w}{ }\PYG{c}{! array notation}
\PYG{n}{R}\PYG{+w}{ }\PYG{o}{=}\PYG{+w}{ }\PYG{n+nb}{PRODUCT}\PYG{p}{(}\PYG{k}{ARRAY}\PYG{p}{(}\PYG{l+m+mi}{1}\PYG{p}{:}\PYG{l+m+mi}{10}\PYG{p}{,}\PYG{n}{lo}\PYG{p}{:}\PYG{n}{hi}\PYG{p}{)}\PYG{p}{)}\PYG{+w}{ }\PYG{c}{! array sections}
\PYG{n}{R}\PYG{+w}{ }\PYG{o}{=}\PYG{+w}{ }\PYG{n+nb}{PRODUCT}\PYG{p}{(}\PYG{k}{ARRAY}\PYG{p}{(}\PYG{l+m+mi}{1}\PYG{p}{:}\PYG{l+m+mi}{10}\PYG{p}{,}\PYG{p}{:}\PYG{p}{)}\PYG{p}{)}\PYG{+w}{ }\PYG{c}{! mix of array section and array notation}
\PYG{n}{R}\PYG{+w}{ }\PYG{o}{=}\PYG{+w}{ }\PYG{n+nb}{PRODUCT}\PYG{p}{(}\PYG{k}{ARRAY}\PYG{p}{(}\PYG{l+m+mi}{1}\PYG{p}{:}\PYG{l+m+mi}{10}\PYG{p}{,}\PYG{l+m+mi}{2}\PYG{p}{)}\PYG{p}{)}\PYG{+w}{ }\PYG{c}{! mix of array section and scalar bound}
\end{sphinxVerbatim}

\sphinxAtStartPar
An example use of this transformation is given below:

\begin{sphinxVerbatim}[commandchars=\\\{\}]
\PYG{g+gp}{\PYGZgt{}\PYGZgt{}\PYGZgt{} }\PYG{k+kn}{from}\PYG{+w}{ }\PYG{n+nn}{psyclone}\PYG{n+nn}{.}\PYG{n+nn}{psyir}\PYG{n+nn}{.}\PYG{n+nn}{backend}\PYG{n+nn}{.}\PYG{n+nn}{fortran}\PYG{+w}{ }\PYG{k+kn}{import} \PYG{n}{FortranWriter}
\PYG{g+gp}{\PYGZgt{}\PYGZgt{}\PYGZgt{} }\PYG{k+kn}{from}\PYG{+w}{ }\PYG{n+nn}{psyclone}\PYG{n+nn}{.}\PYG{n+nn}{psyir}\PYG{n+nn}{.}\PYG{n+nn}{frontend}\PYG{n+nn}{.}\PYG{n+nn}{fortran}\PYG{+w}{ }\PYG{k+kn}{import} \PYG{n}{FortranReader}
\PYG{g+gp}{\PYGZgt{}\PYGZgt{}\PYGZgt{} }\PYG{k+kn}{from}\PYG{+w}{ }\PYG{n+nn}{psyclone}\PYG{n+nn}{.}\PYG{n+nn}{psyir}\PYG{n+nn}{.}\PYG{n+nn}{transformations}\PYG{+w}{ }\PYG{k+kn}{import} \PYG{n}{Product2LoopTrans}
\PYG{g+gp}{\PYGZgt{}\PYGZgt{}\PYGZgt{} }\PYG{n}{code} \PYG{o}{=} \PYG{p}{(}\PYG{l+s+s2}{\PYGZdq{}}\PYG{l+s+s2}{subroutine product\PYGZus{}test(array)}\PYG{l+s+se}{\PYGZbs{}n}\PYG{l+s+s2}{\PYGZdq{}}
\PYG{g+gp}{... }        \PYG{l+s+s2}{\PYGZdq{}}\PYG{l+s+s2}{  real :: array(10,10)}\PYG{l+s+se}{\PYGZbs{}n}\PYG{l+s+s2}{\PYGZdq{}}
\PYG{g+gp}{... }        \PYG{l+s+s2}{\PYGZdq{}}\PYG{l+s+s2}{  real :: result}\PYG{l+s+se}{\PYGZbs{}n}\PYG{l+s+s2}{\PYGZdq{}}
\PYG{g+gp}{... }        \PYG{l+s+s2}{\PYGZdq{}}\PYG{l+s+s2}{  result = product(array)}\PYG{l+s+se}{\PYGZbs{}n}\PYG{l+s+s2}{\PYGZdq{}}
\PYG{g+gp}{... }        \PYG{l+s+s2}{\PYGZdq{}}\PYG{l+s+s2}{end subroutine}\PYG{l+s+se}{\PYGZbs{}n}\PYG{l+s+s2}{\PYGZdq{}}\PYG{p}{)}
\PYG{g+gp}{\PYGZgt{}\PYGZgt{}\PYGZgt{} }\PYG{n}{psyir} \PYG{o}{=} \PYG{n}{FortranReader}\PYG{p}{(}\PYG{p}{)}\PYG{o}{.}\PYG{n}{psyir\PYGZus{}from\PYGZus{}source}\PYG{p}{(}\PYG{n}{code}\PYG{p}{)}
\PYG{g+gp}{\PYGZgt{}\PYGZgt{}\PYGZgt{} }\PYG{n}{product\PYGZus{}node} \PYG{o}{=} \PYG{n}{psyir}\PYG{o}{.}\PYG{n}{children}\PYG{p}{[}\PYG{l+m+mi}{0}\PYG{p}{]}\PYG{o}{.}\PYG{n}{children}\PYG{p}{[}\PYG{l+m+mi}{0}\PYG{p}{]}\PYG{o}{.}\PYG{n}{children}\PYG{p}{[}\PYG{l+m+mi}{1}\PYG{p}{]}
\PYG{g+gp}{\PYGZgt{}\PYGZgt{}\PYGZgt{} }\PYG{n}{Product2LoopTrans}\PYG{p}{(}\PYG{p}{)}\PYG{o}{.}\PYG{n}{apply}\PYG{p}{(}\PYG{n}{product\PYGZus{}node}\PYG{p}{)}
\PYG{g+gp}{\PYGZgt{}\PYGZgt{}\PYGZgt{} }\PYG{n+nb}{print}\PYG{p}{(}\PYG{n}{FortranWriter}\PYG{p}{(}\PYG{p}{)}\PYG{p}{(}\PYG{n}{psyir}\PYG{p}{)}\PYG{p}{)}
\PYG{g+go}{subroutine product\PYGZus{}test(array)}
\PYG{g+go}{  real, dimension(10,10) :: array}
\PYG{g+go}{  real :: result}
\PYG{g+go}{  integer :: idx}
\PYG{g+go}{  integer :: idx\PYGZus{}1}

\PYG{g+go}{  result = 1.0}
\PYG{g+go}{  do idx = 1, 10, 1}
\PYG{g+go}{    do idx\PYGZus{}1 = 1, 10, 1}
\PYG{g+go}{      result = result * array(idx\PYGZus{}1,idx)}
\PYG{g+go}{    enddo}
\PYG{g+go}{  enddo}

\PYG{g+go}{end subroutine product\PYGZus{}test}
\end{sphinxVerbatim}


\begin{fulllineitems}

\pysigstartsignatures
\pysiglinewithargsret{\sphinxbfcode{\sphinxupquote{apply}}}{\sphinxparam{\DUrole{n,n}{node}}\sphinxparamcomma \sphinxparam{\DUrole{n,n}{options}\DUrole{o,o}{=}\DUrole{default_value}{None}}}{}
\pysigstopsignatures
\sphinxAtStartPar
Apply the array\sphinxhyphen{}reduction intrinsic conversion transformation to
the specified node. This node must be one of these intrinsic
operations which is converted to an equivalent loop structure.
\begin{quote}\begin{description}
\sphinxlineitem{Parameters}\begin{itemize}
\item {} 
\sphinxAtStartPar
\sphinxstyleliteralstrong{\sphinxupquote{node}} (\sphinxcode{\sphinxupquote{psyclone.psyir.nodes.IntrinsicCall}}) \textendash{} an array\sphinxhyphen{}reduction intrinsic.

\item {} 
\sphinxAtStartPar
\sphinxstyleliteralstrong{\sphinxupquote{options}} (\sphinxstyleliteralemphasis{\sphinxupquote{Optional}}\sphinxstyleliteralemphasis{\sphinxupquote{{[}}}\sphinxstyleliteralemphasis{\sphinxupquote{Dict}}\sphinxstyleliteralemphasis{\sphinxupquote{{[}}}\sphinxhref{https://docs.python.org/3/library/stdtypes.html\#str}{\sphinxstyleliteralemphasis{\sphinxupquote{str}}}\sphinxstyleliteralemphasis{\sphinxupquote{, }}\sphinxstyleliteralemphasis{\sphinxupquote{Any}}\sphinxstyleliteralemphasis{\sphinxupquote{{]}}}\sphinxstyleliteralemphasis{\sphinxupquote{{]}}}) \textendash{} options for the transformation.

\end{itemize}

\end{description}\end{quote}

\end{fulllineitems}


\end{fulllineitems}



\bigskip\hrule\bigskip



\begin{fulllineitems}

\pysigstartsignatures
\pysigline{\sphinxbfcode{\sphinxupquote{class\DUrole{w,w}{  }}}\sphinxcode{\sphinxupquote{psyclone.psyir.transformations.}}\sphinxbfcode{\sphinxupquote{ProfileTrans}}}
\pysigstopsignatures
\sphinxAtStartPar
Create a profile region around a list of statements. For
example:

\begin{sphinxVerbatim}[commandchars=\\\{\}]
\PYG{g+gp}{\PYGZgt{}\PYGZgt{}\PYGZgt{} }\PYG{k+kn}{from}\PYG{+w}{ }\PYG{n+nn}{psyclone}\PYG{n+nn}{.}\PYG{n+nn}{parse}\PYG{n+nn}{.}\PYG{n+nn}{algorithm}\PYG{+w}{ }\PYG{k+kn}{import} \PYG{n}{parse}
\PYG{g+gp}{\PYGZgt{}\PYGZgt{}\PYGZgt{} }\PYG{k+kn}{from}\PYG{+w}{ }\PYG{n+nn}{psyclone}\PYG{n+nn}{.}\PYG{n+nn}{parse}\PYG{n+nn}{.}\PYG{n+nn}{utils}\PYG{+w}{ }\PYG{k+kn}{import} \PYG{n}{ParseError}
\PYG{g+gp}{\PYGZgt{}\PYGZgt{}\PYGZgt{} }\PYG{k+kn}{from}\PYG{+w}{ }\PYG{n+nn}{psyclone}\PYG{n+nn}{.}\PYG{n+nn}{psyGen}\PYG{+w}{ }\PYG{k+kn}{import} \PYG{n}{PSyFactory}\PYG{p}{,} \PYG{n}{GenerationError}
\PYG{g+gp}{\PYGZgt{}\PYGZgt{}\PYGZgt{} }\PYG{k+kn}{from}\PYG{+w}{ }\PYG{n+nn}{psyclone}\PYG{n+nn}{.}\PYG{n+nn}{psyir}\PYG{n+nn}{.}\PYG{n+nn}{transformations}\PYG{+w}{ }\PYG{k+kn}{import} \PYG{n}{ProfileTrans}
\PYG{g+gp}{\PYGZgt{}\PYGZgt{}\PYGZgt{} }\PYG{n}{api} \PYG{o}{=} \PYG{l+s+s2}{\PYGZdq{}}\PYG{l+s+s2}{gocean}\PYG{l+s+s2}{\PYGZdq{}}
\PYG{g+gp}{\PYGZgt{}\PYGZgt{}\PYGZgt{} }\PYG{n}{filename} \PYG{o}{=} \PYG{l+s+s2}{\PYGZdq{}}\PYG{l+s+s2}{nemolite2d\PYGZus{}alg.f90}\PYG{l+s+s2}{\PYGZdq{}}
\PYG{g+gp}{\PYGZgt{}\PYGZgt{}\PYGZgt{} }\PYG{n}{ast}\PYG{p}{,} \PYG{n}{invokeInfo} \PYG{o}{=} \PYG{n}{parse}\PYG{p}{(}\PYG{n}{filename}\PYG{p}{,} \PYG{n}{api}\PYG{o}{=}\PYG{n}{api}\PYG{p}{,} \PYG{n}{invoke\PYGZus{}name}\PYG{o}{=}\PYG{l+s+s2}{\PYGZdq{}}\PYG{l+s+s2}{invoke}\PYG{l+s+s2}{\PYGZdq{}}\PYG{p}{)}
\PYG{g+gp}{\PYGZgt{}\PYGZgt{}\PYGZgt{} }\PYG{n}{psy} \PYG{o}{=} \PYG{n}{PSyFactory}\PYG{p}{(}\PYG{n}{api}\PYG{p}{)}\PYG{o}{.}\PYG{n}{create}\PYG{p}{(}\PYG{n}{invokeInfo}\PYG{p}{)}
\PYG{g+gp}{\PYGZgt{}\PYGZgt{}\PYGZgt{}}
\PYG{g+gp}{\PYGZgt{}\PYGZgt{}\PYGZgt{} }\PYG{n}{p\PYGZus{}trans} \PYG{o}{=} \PYG{n}{ProfileTrans}\PYG{p}{(}\PYG{p}{)}
\PYG{g+gp}{\PYGZgt{}\PYGZgt{}\PYGZgt{}}
\PYG{g+gp}{\PYGZgt{}\PYGZgt{}\PYGZgt{} }\PYG{n}{schedule} \PYG{o}{=} \PYG{n}{psy}\PYG{o}{.}\PYG{n}{invokes}\PYG{o}{.}\PYG{n}{get}\PYG{p}{(}\PYG{l+s+s1}{\PYGZsq{}}\PYG{l+s+s1}{invoke\PYGZus{}0}\PYG{l+s+s1}{\PYGZsq{}}\PYG{p}{)}\PYG{o}{.}\PYG{n}{schedule}
\PYG{g+gp}{\PYGZgt{}\PYGZgt{}\PYGZgt{} }\PYG{n+nb}{print}\PYG{p}{(}\PYG{n}{schedule}\PYG{o}{.}\PYG{n}{view}\PYG{p}{(}\PYG{p}{)}\PYG{p}{)}
\PYG{g+gp}{\PYGZgt{}\PYGZgt{}\PYGZgt{}}
\PYG{g+gp}{\PYGZgt{}\PYGZgt{}\PYGZgt{} }\PYG{c+c1}{\PYGZsh{} Enclose all children within a single profile region}
\PYG{g+gp}{\PYGZgt{}\PYGZgt{}\PYGZgt{} }\PYG{n}{p\PYGZus{}trans}\PYG{o}{.}\PYG{n}{apply}\PYG{p}{(}\PYG{n}{schedule}\PYG{o}{.}\PYG{n}{children}\PYG{p}{)}
\PYG{g+gp}{\PYGZgt{}\PYGZgt{}\PYGZgt{} }\PYG{n+nb}{print}\PYG{p}{(}\PYG{n}{schedule}\PYG{o}{.}\PYG{n}{view}\PYG{p}{(}\PYG{p}{)}\PYG{p}{)}
\end{sphinxVerbatim}

\sphinxAtStartPar
This implementation relies completely on the base class PSyDataTrans
for the actual work, it only adjusts the name etc, and the list
of valid nodes.


\begin{fulllineitems}

\pysigstartsignatures
\pysiglinewithargsret{\sphinxbfcode{\sphinxupquote{apply}}}{\sphinxparam{\DUrole{n,n}{nodes}}\sphinxparamcomma \sphinxparam{\DUrole{n,n}{options}\DUrole{o,o}{=}\DUrole{default_value}{None}}}{}
\pysigstopsignatures
\sphinxAtStartPar
Apply this transformation to a subset of the nodes within a
schedule \sphinxhyphen{} i.e. enclose the specified Nodes in the
schedule within a single PSyData region.
\begin{quote}\begin{description}
\sphinxlineitem{Parameters}\begin{itemize}
\item {} 
\sphinxAtStartPar
\sphinxstyleliteralstrong{\sphinxupquote{nodes}} (\sphinxcode{\sphinxupquote{psyclone.psyir.nodes.Node}} or list of                      \sphinxcode{\sphinxupquote{psyclone.psyir.nodes.Node}}) \textendash{} can be a single node or a list of nodes.

\item {} 
\sphinxAtStartPar
\sphinxstyleliteralstrong{\sphinxupquote{options}} (\sphinxstyleliteralemphasis{\sphinxupquote{Optional}}\sphinxstyleliteralemphasis{\sphinxupquote{{[}}}\sphinxstyleliteralemphasis{\sphinxupquote{Dict}}\sphinxstyleliteralemphasis{\sphinxupquote{{[}}}\sphinxhref{https://docs.python.org/3/library/stdtypes.html\#str}{\sphinxstyleliteralemphasis{\sphinxupquote{str}}}\sphinxstyleliteralemphasis{\sphinxupquote{, }}\sphinxstyleliteralemphasis{\sphinxupquote{Any}}\sphinxstyleliteralemphasis{\sphinxupquote{{]}}}\sphinxstyleliteralemphasis{\sphinxupquote{{]}}}) \textendash{} a dictionary with options for transformations.

\item {} 
\sphinxAtStartPar
\sphinxstyleliteralstrong{\sphinxupquote{options}}\sphinxstyleliteralstrong{\sphinxupquote{{[}}}\sphinxstyleliteralstrong{\sphinxupquote{"prefix"}}\sphinxstyleliteralstrong{\sphinxupquote{{]}}} (\sphinxhref{https://docs.python.org/3/library/stdtypes.html\#str}{\sphinxstyleliteralemphasis{\sphinxupquote{str}}}) \textendash{} a prefix to use for the PSyData module             name (\sphinxcode{\sphinxupquote{PREFIX\_psy\_data\_mod}}) and the PSyDataType             (\sphinxcode{\sphinxupquote{PREFIX\_PSYDATATYPE}}) \sphinxhyphen{} a “\_” will be added automatically.             It defaults to “”.

\item {} 
\sphinxAtStartPar
\sphinxstyleliteralstrong{\sphinxupquote{options}}\sphinxstyleliteralstrong{\sphinxupquote{{[}}}\sphinxstyleliteralstrong{\sphinxupquote{"region\_name"}}\sphinxstyleliteralstrong{\sphinxupquote{{]}}} (\sphinxstyleliteralemphasis{\sphinxupquote{(}}\sphinxhref{https://docs.python.org/3/library/stdtypes.html\#str}{\sphinxstyleliteralemphasis{\sphinxupquote{str}}}\sphinxstyleliteralemphasis{\sphinxupquote{,}}\sphinxhref{https://docs.python.org/3/library/stdtypes.html\#str}{\sphinxstyleliteralemphasis{\sphinxupquote{str}}}\sphinxstyleliteralemphasis{\sphinxupquote{)}}) \textendash{} an optional name to             use for this PSyData area, provided as a 2\sphinxhyphen{}tuple containing a             location name followed by a local name. The pair of strings             should uniquely identify a region unless aggregate information             is required (and is supported by the runtime library).

\end{itemize}

\end{description}\end{quote}

\end{fulllineitems}


\end{fulllineitems}



\bigskip\hrule\bigskip



\begin{fulllineitems}

\pysigstartsignatures
\pysiglinewithargsret{\sphinxbfcode{\sphinxupquote{class\DUrole{w,w}{  }}}\sphinxcode{\sphinxupquote{psyclone.psyir.transformations.}}\sphinxbfcode{\sphinxupquote{ReadOnlyVerifyTrans}}}{\sphinxparam{\DUrole{n,n}{node\_class=\textless{}class \textquotesingle{}psyclone.psyir.nodes.read\_only\_verify\_node.ReadOnlyVerifyNode\textquotesingle{}\textgreater{}}}}{}
\pysigstopsignatures
\sphinxAtStartPar
This transformation inserts a ReadOnlyVerifyNode or a node derived
from ReadOnlyVerifyNode into the PSyIR of a schedule. At code creation
time this node will use the PSyData API to create code that will
verify that read\sphinxhyphen{}only quantities are not modified.

\sphinxAtStartPar
After applying the transformation the Nodes marked for verification are
children of the ReadOnlyVerifyNode.
Nodes to verify can be individual constructs within an Invoke (e.g.
Loops containing a Kernel or BuiltIn call) or entire Invokes.
\begin{quote}\begin{description}
\sphinxlineitem{Parameters}
\sphinxAtStartPar
\sphinxstyleliteralstrong{\sphinxupquote{node\_class}} (\sphinxcode{\sphinxupquote{psyclone.psyir.nodes.ReadOnlyVerifyNode}} or         derived class) \textendash{} The class of Node which will be inserted         into the tree (defaults to ReadOnlyVerifyNode), but can be any         derived class.

\end{description}\end{quote}


\begin{fulllineitems}

\pysigstartsignatures
\pysiglinewithargsret{\sphinxbfcode{\sphinxupquote{apply}}}{\sphinxparam{\DUrole{n,n}{nodes}}\sphinxparamcomma \sphinxparam{\DUrole{n,n}{options}\DUrole{o,o}{=}\DUrole{default_value}{None}}}{}
\pysigstopsignatures
\sphinxAtStartPar
Apply this transformation to a subset of the nodes within a
schedule \sphinxhyphen{} i.e. enclose the specified Nodes in the
schedule within a single PSyData region.
\begin{quote}\begin{description}
\sphinxlineitem{Parameters}\begin{itemize}
\item {} 
\sphinxAtStartPar
\sphinxstyleliteralstrong{\sphinxupquote{nodes}} (\sphinxcode{\sphinxupquote{psyclone.psyir.nodes.Node}} or list of                      \sphinxcode{\sphinxupquote{psyclone.psyir.nodes.Node}}) \textendash{} can be a single node or a list of nodes.

\item {} 
\sphinxAtStartPar
\sphinxstyleliteralstrong{\sphinxupquote{options}} (\sphinxstyleliteralemphasis{\sphinxupquote{Optional}}\sphinxstyleliteralemphasis{\sphinxupquote{{[}}}\sphinxstyleliteralemphasis{\sphinxupquote{Dict}}\sphinxstyleliteralemphasis{\sphinxupquote{{[}}}\sphinxhref{https://docs.python.org/3/library/stdtypes.html\#str}{\sphinxstyleliteralemphasis{\sphinxupquote{str}}}\sphinxstyleliteralemphasis{\sphinxupquote{, }}\sphinxstyleliteralemphasis{\sphinxupquote{Any}}\sphinxstyleliteralemphasis{\sphinxupquote{{]}}}\sphinxstyleliteralemphasis{\sphinxupquote{{]}}}) \textendash{} a dictionary with options for transformations.

\item {} 
\sphinxAtStartPar
\sphinxstyleliteralstrong{\sphinxupquote{options}}\sphinxstyleliteralstrong{\sphinxupquote{{[}}}\sphinxstyleliteralstrong{\sphinxupquote{"prefix"}}\sphinxstyleliteralstrong{\sphinxupquote{{]}}} (\sphinxhref{https://docs.python.org/3/library/stdtypes.html\#str}{\sphinxstyleliteralemphasis{\sphinxupquote{str}}}) \textendash{} a prefix to use for the PSyData module             name (\sphinxcode{\sphinxupquote{PREFIX\_psy\_data\_mod}}) and the PSyDataType             (\sphinxcode{\sphinxupquote{PREFIX\_PSYDATATYPE}}) \sphinxhyphen{} a “\_” will be added automatically.             It defaults to “”.

\item {} 
\sphinxAtStartPar
\sphinxstyleliteralstrong{\sphinxupquote{options}}\sphinxstyleliteralstrong{\sphinxupquote{{[}}}\sphinxstyleliteralstrong{\sphinxupquote{"region\_name"}}\sphinxstyleliteralstrong{\sphinxupquote{{]}}} (\sphinxstyleliteralemphasis{\sphinxupquote{(}}\sphinxhref{https://docs.python.org/3/library/stdtypes.html\#str}{\sphinxstyleliteralemphasis{\sphinxupquote{str}}}\sphinxstyleliteralemphasis{\sphinxupquote{,}}\sphinxhref{https://docs.python.org/3/library/stdtypes.html\#str}{\sphinxstyleliteralemphasis{\sphinxupquote{str}}}\sphinxstyleliteralemphasis{\sphinxupquote{)}}) \textendash{} an optional name to             use for this PSyData area, provided as a 2\sphinxhyphen{}tuple containing a             location name followed by a local name. The pair of strings             should uniquely identify a region unless aggregate information             is required (and is supported by the runtime library).

\end{itemize}

\end{description}\end{quote}

\end{fulllineitems}


\end{fulllineitems}



\bigskip\hrule\bigskip



\begin{fulllineitems}

\pysigstartsignatures
\pysigline{\sphinxbfcode{\sphinxupquote{class\DUrole{w,w}{  }}}\sphinxcode{\sphinxupquote{psyclone.psyir.transformations.}}\sphinxbfcode{\sphinxupquote{Reference2ArrayRangeTrans}}}
\pysigstopsignatures
\sphinxAtStartPar
Provides a transformation from PSyIR Array Notation (a reference to
an Array) to a PSyIR Range. For example:

\begin{sphinxVerbatim}[commandchars=\\\{\}]
\PYG{g+gp}{\PYGZgt{}\PYGZgt{}\PYGZgt{} }\PYG{k+kn}{from}\PYG{+w}{ }\PYG{n+nn}{psyclone}\PYG{n+nn}{.}\PYG{n+nn}{psyir}\PYG{n+nn}{.}\PYG{n+nn}{backend}\PYG{n+nn}{.}\PYG{n+nn}{fortran}\PYG{+w}{ }\PYG{k+kn}{import} \PYG{n}{FortranWriter}
\PYG{g+gp}{\PYGZgt{}\PYGZgt{}\PYGZgt{} }\PYG{k+kn}{from}\PYG{+w}{ }\PYG{n+nn}{psyclone}\PYG{n+nn}{.}\PYG{n+nn}{psyir}\PYG{n+nn}{.}\PYG{n+nn}{frontend}\PYG{n+nn}{.}\PYG{n+nn}{fortran}\PYG{+w}{ }\PYG{k+kn}{import} \PYG{n}{FortranReader}
\PYG{g+gp}{\PYGZgt{}\PYGZgt{}\PYGZgt{} }\PYG{k+kn}{from}\PYG{+w}{ }\PYG{n+nn}{psyclone}\PYG{n+nn}{.}\PYG{n+nn}{psyir}\PYG{n+nn}{.}\PYG{n+nn}{nodes}\PYG{+w}{ }\PYG{k+kn}{import} \PYG{n}{Reference}
\PYG{g+gp}{\PYGZgt{}\PYGZgt{}\PYGZgt{} }\PYG{k+kn}{from}\PYG{+w}{ }\PYG{n+nn}{psyclone}\PYG{n+nn}{.}\PYG{n+nn}{psyir}\PYG{n+nn}{.}\PYG{n+nn}{transformations}\PYG{+w}{ }\PYG{k+kn}{import} \PYG{n}{TransformationError}
\PYG{g+gp}{\PYGZgt{}\PYGZgt{}\PYGZgt{} }\PYG{n}{CODE} \PYG{o}{=} \PYG{p}{(}\PYG{l+s+s2}{\PYGZdq{}}\PYG{l+s+s2}{program example}\PYG{l+s+se}{\PYGZbs{}n}\PYG{l+s+s2}{\PYGZdq{}}
\PYG{g+gp}{... }        \PYG{l+s+s2}{\PYGZdq{}}\PYG{l+s+s2}{real :: a(:)}\PYG{l+s+se}{\PYGZbs{}n}\PYG{l+s+s2}{\PYGZdq{}}
\PYG{g+gp}{... }        \PYG{l+s+s2}{\PYGZdq{}}\PYG{l+s+s2}{a = 0.0}\PYG{l+s+se}{\PYGZbs{}n}\PYG{l+s+s2}{\PYGZdq{}}
\PYG{g+gp}{... }        \PYG{l+s+s2}{\PYGZdq{}}\PYG{l+s+s2}{end program}\PYG{l+s+se}{\PYGZbs{}n}\PYG{l+s+s2}{\PYGZdq{}}\PYG{p}{)}
\PYG{g+gp}{\PYGZgt{}\PYGZgt{}\PYGZgt{} }\PYG{n}{trans} \PYG{o}{=} \PYG{n}{Reference2ArrayRangeTrans}\PYG{p}{(}\PYG{p}{)}
\PYG{g+gp}{\PYGZgt{}\PYGZgt{}\PYGZgt{} }\PYG{n}{psyir} \PYG{o}{=} \PYG{n}{FortranReader}\PYG{p}{(}\PYG{p}{)}\PYG{o}{.}\PYG{n}{psyir\PYGZus{}from\PYGZus{}source}\PYG{p}{(}\PYG{n}{CODE}\PYG{p}{)}
\PYG{g+gp}{\PYGZgt{}\PYGZgt{}\PYGZgt{} }\PYG{k}{for} \PYG{n}{reference} \PYG{o+ow}{in} \PYG{n}{psyir}\PYG{o}{.}\PYG{n}{walk}\PYG{p}{(}\PYG{n}{Reference}\PYG{p}{)}\PYG{p}{:}
\PYG{g+gp}{... }   \PYG{k}{try}\PYG{p}{:}
\PYG{g+gp}{... }       \PYG{n}{trans}\PYG{o}{.}\PYG{n}{apply}\PYG{p}{(}\PYG{n}{reference}\PYG{p}{)}
\PYG{g+gp}{... }   \PYG{k}{except} \PYG{n}{TransformationError}\PYG{p}{:}
\PYG{g+gp}{... }       \PYG{k}{pass}
\PYG{g+gp}{\PYGZgt{}\PYGZgt{}\PYGZgt{} }\PYG{n+nb}{print}\PYG{p}{(}\PYG{n}{FortranWriter}\PYG{p}{(}\PYG{p}{)}\PYG{p}{(}\PYG{n}{psyir}\PYG{p}{)}\PYG{p}{)}
\PYG{g+go}{program example}
\PYG{g+go}{  real, dimension(:) :: a}

\PYG{g+go}{  a(:) = 0.0}

\PYG{g+go}{end program example}
\end{sphinxVerbatim}

\sphinxAtStartPar
This transformation does not currently support arrays within
structures, see issue \#1858.


\begin{fulllineitems}

\pysigstartsignatures
\pysiglinewithargsret{\sphinxbfcode{\sphinxupquote{apply}}}{\sphinxparam{\DUrole{n,n}{node}}\sphinxparamcomma \sphinxparam{\DUrole{n,n}{options}\DUrole{o,o}{=}\DUrole{default_value}{None}}}{}
\pysigstopsignatures
\sphinxAtStartPar
Apply the Reference2ArrayRangeTrans transformation to the specified
node. The node must be a Reference to an array. The Reference
is replaced by an ArrayReference with appropriate explicit
range nodes (termed colon notation in Fortran).
\begin{quote}\begin{description}
\sphinxlineitem{Parameters}\begin{itemize}
\item {} 
\sphinxAtStartPar
\sphinxstyleliteralstrong{\sphinxupquote{node}} (\sphinxcode{\sphinxupquote{psyclone.psyir.nodes.Reference}}) \textendash{} a Reference node.

\item {} 
\sphinxAtStartPar
\sphinxstyleliteralstrong{\sphinxupquote{options}} (\sphinxstyleliteralemphasis{\sphinxupquote{Optional}}\sphinxstyleliteralemphasis{\sphinxupquote{{[}}}\sphinxstyleliteralemphasis{\sphinxupquote{Dict}}\sphinxstyleliteralemphasis{\sphinxupquote{{[}}}\sphinxhref{https://docs.python.org/3/library/stdtypes.html\#str}{\sphinxstyleliteralemphasis{\sphinxupquote{str}}}\sphinxstyleliteralemphasis{\sphinxupquote{, }}\sphinxstyleliteralemphasis{\sphinxupquote{Any}}\sphinxstyleliteralemphasis{\sphinxupquote{{]}}}\sphinxstyleliteralemphasis{\sphinxupquote{{]}}}) \textendash{} a dict with options for transformations.

\end{itemize}

\end{description}\end{quote}

\end{fulllineitems}


\end{fulllineitems}



\bigskip\hrule\bigskip

\phantomsection\label{\detokenize{transformations:replace-induction-variable-trans}}

\begin{fulllineitems}

\pysigstartsignatures
\pysigline{\sphinxbfcode{\sphinxupquote{class\DUrole{w,w}{  }}}\sphinxcode{\sphinxupquote{psyclone.psyir.transformations.}}\sphinxbfcode{\sphinxupquote{ReplaceInductionVariablesTrans}}}
\pysigstopsignatures
\sphinxAtStartPar
Move all supported induction variables out of the loop, and replace
their usage inside the loop. For example:

\begin{sphinxVerbatim}[commandchars=\\\{\}]
\PYG{g+gp}{\PYGZgt{}\PYGZgt{}\PYGZgt{} }\PYG{k+kn}{from}\PYG{+w}{ }\PYG{n+nn}{psyclone}\PYG{n+nn}{.}\PYG{n+nn}{psyir}\PYG{n+nn}{.}\PYG{n+nn}{frontend}\PYG{n+nn}{.}\PYG{n+nn}{fortran}\PYG{+w}{ }\PYG{k+kn}{import} \PYG{n}{FortranReader}
\PYG{g+gp}{\PYGZgt{}\PYGZgt{}\PYGZgt{} }\PYG{k+kn}{from}\PYG{+w}{ }\PYG{n+nn}{psyclone}\PYG{n+nn}{.}\PYG{n+nn}{psyir}\PYG{n+nn}{.}\PYG{n+nn}{nodes}\PYG{+w}{ }\PYG{k+kn}{import} \PYG{n}{Loop}
\PYG{g+gp}{\PYGZgt{}\PYGZgt{}\PYGZgt{} }\PYG{k+kn}{from}\PYG{+w}{ }\PYG{n+nn}{psyclone}\PYG{n+nn}{.}\PYG{n+nn}{psyir}\PYG{n+nn}{.}\PYG{n+nn}{transformations}\PYG{+w}{ }\PYG{k+kn}{import}             \PYG{n}{ReplaceInductionVariablesTrans}
\PYG{g+gp}{\PYGZgt{}\PYGZgt{}\PYGZgt{} }\PYG{k+kn}{from}\PYG{+w}{ }\PYG{n+nn}{psyclone}\PYG{n+nn}{.}\PYG{n+nn}{psyir}\PYG{n+nn}{.}\PYG{n+nn}{backend}\PYG{n+nn}{.}\PYG{n+nn}{fortran}\PYG{+w}{ }\PYG{k+kn}{import} \PYG{n}{FortranWriter}
\PYG{g+gp}{\PYGZgt{}\PYGZgt{}\PYGZgt{} }\PYG{n}{psyir} \PYG{o}{=} \PYG{n}{FortranReader}\PYG{p}{(}\PYG{p}{)}\PYG{o}{.}\PYG{n}{psyir\PYGZus{}from\PYGZus{}source}\PYG{p}{(}\PYG{l+s+s2}{\PYGZdq{}\PYGZdq{}\PYGZdq{}}
\PYG{g+gp}{... }\PYG{l+s+s2}{subroutine sub()}
\PYG{g+gp}{... }\PYG{l+s+s2}{    integer :: i, im, ic, tmp(100)}
\PYG{g+gp}{... }\PYG{l+s+s2}{    do i=1, 100}
\PYG{g+gp}{... }\PYG{l+s+s2}{        im = i \PYGZhy{} 1}
\PYG{g+gp}{... }\PYG{l+s+s2}{        ic = 2}
\PYG{g+gp}{... }\PYG{l+s+s2}{        tmp(i) = ic * im}
\PYG{g+gp}{... }\PYG{l+s+s2}{    enddo}
\PYG{g+gp}{... }\PYG{l+s+s2}{end subroutine sub}\PYG{l+s+s2}{\PYGZdq{}\PYGZdq{}\PYGZdq{}}\PYG{p}{)}
\PYG{g+gp}{\PYGZgt{}\PYGZgt{}\PYGZgt{} }\PYG{n}{loop} \PYG{o}{=} \PYG{n}{psyir}\PYG{o}{.}\PYG{n}{walk}\PYG{p}{(}\PYG{n}{Loop}\PYG{p}{)}\PYG{p}{[}\PYG{l+m+mi}{0}\PYG{p}{]}
\PYG{g+gp}{\PYGZgt{}\PYGZgt{}\PYGZgt{} }\PYG{n}{ReplaceInductionVariablesTrans}\PYG{p}{(}\PYG{p}{)}\PYG{o}{.}\PYG{n}{apply}\PYG{p}{(}\PYG{n}{loop}\PYG{p}{)}
\PYG{g+gp}{\PYGZgt{}\PYGZgt{}\PYGZgt{} }\PYG{n+nb}{print}\PYG{p}{(}\PYG{n}{FortranWriter}\PYG{p}{(}\PYG{p}{)}\PYG{p}{(}\PYG{n}{psyir}\PYG{p}{)}\PYG{p}{)}
\PYG{g+go}{subroutine sub()}
\PYG{g+go}{  integer :: i}
\PYG{g+go}{  integer :: im}
\PYG{g+go}{  integer :: ic}
\PYG{g+go}{  integer, dimension(100) :: tmp}

\PYG{g+go}{  do i = 1, 100, 1}
\PYG{g+go}{    tmp(i) = 2 * (i \PYGZhy{} 1)}
\PYG{g+go}{  enddo}
\PYG{g+go}{  ic = 2}
\PYG{g+go}{  im = i \PYGZhy{} 1 \PYGZhy{} 1}

\PYG{g+go}{end subroutine sub}
\end{sphinxVerbatim}

\sphinxAtStartPar
The replaced induction variables assignments are added after the loop,
so these variables will have the correct value if they are used elsewhere.

\sphinxAtStartPar
The following restrictions apply for the assignment to an induction
variable:
\begin{itemize}
\item {} 
\sphinxAtStartPar
the variable must be a scalar (i.e. no array access at all,
not even a constant like a(3) or a\%b(3)\%c)

\item {} 
\sphinxAtStartPar
none of variables on the right\sphinxhyphen{}hand side can be written in the loop body
(the loop variable is written in the Loop statement, not in
the body, so it can be used).

\item {} 
\sphinxAtStartPar
Only intrinsic function calls are allowed on the RHS (since they
are known to be elemental)

\item {} 
\sphinxAtStartPar
the assigned variable must not be read before the assignment.

\item {} 
\sphinxAtStartPar
the assigned variable cannot occur on the right\sphinxhyphen{}hand side
(e.g. \sphinxtitleref{k = k + 3}).

\item {} 
\sphinxAtStartPar
there must be only one assignment to this induction variable.

\end{itemize}


\begin{fulllineitems}

\pysigstartsignatures
\pysiglinewithargsret{\sphinxbfcode{\sphinxupquote{apply}}}{\sphinxparam{\DUrole{n,n}{node}}\sphinxparamcomma \sphinxparam{\DUrole{n,n}{options}\DUrole{o,o}{=}\DUrole{default_value}{None}}}{}
\pysigstopsignatures
\sphinxAtStartPar
Apply the ReplaceInductionVariablesTrans transformation to the
specified node. The node must be a loop. In case of nested
loops, the transformation might need to be applied several
times, from the inner\sphinxhyphen{}most loop outwards.
\begin{quote}\begin{description}
\sphinxlineitem{Parameters}
\sphinxAtStartPar
\sphinxstyleliteralstrong{\sphinxupquote{node}} (\sphinxcode{\sphinxupquote{psyclone.psyir.nodes.Loop}}) \textendash{} a Loop node.

\end{description}\end{quote}

\end{fulllineitems}


\end{fulllineitems}



\bigskip\hrule\bigskip



\begin{fulllineitems}

\pysigstartsignatures
\pysigline{\sphinxbfcode{\sphinxupquote{class\DUrole{w,w}{  }}}\sphinxcode{\sphinxupquote{psyclone.psyir.transformations.}}\sphinxbfcode{\sphinxupquote{Sign2CodeTrans}}}
\pysigstopsignatures
\sphinxAtStartPar
Provides a transformation from a PSyIR SIGN intrinsic node to
equivalent code in a PSyIR tree. Validity checks are also
performed.

\sphinxAtStartPar
The transformation replaces

\begin{sphinxVerbatim}[commandchars=\\\{\}]
\PYG{n}{R} \PYG{o}{=} \PYG{n}{SIGN}\PYG{p}{(}\PYG{n}{A}\PYG{p}{,} \PYG{n}{B}\PYG{p}{)}
\end{sphinxVerbatim}

\sphinxAtStartPar
with the following logic:

\begin{sphinxVerbatim}[commandchars=\\\{\}]
\PYG{n}{R} \PYG{o}{=} \PYG{n}{ABS}\PYG{p}{(}\PYG{n}{A}\PYG{p}{)}
\PYG{k}{if} \PYG{n}{B} \PYG{o}{\PYGZlt{}} \PYG{l+m+mf}{0.0}\PYG{p}{:}
    \PYG{n}{R} \PYG{o}{=} \PYG{n}{R}\PYG{o}{*}\PYG{o}{\PYGZhy{}}\PYG{l+m+mf}{1.0}
\end{sphinxVerbatim}

\sphinxAtStartPar
i.e. the value of \sphinxcode{\sphinxupquote{A}} with the sign of \sphinxcode{\sphinxupquote{B}}


\begin{fulllineitems}

\pysigstartsignatures
\pysiglinewithargsret{\sphinxbfcode{\sphinxupquote{apply}}}{\sphinxparam{\DUrole{n,n}{node}}\sphinxparamcomma \sphinxparam{\DUrole{n,n}{options}\DUrole{o,o}{=}\DUrole{default_value}{None}}}{}
\pysigstopsignatures
\sphinxAtStartPar
Apply the SIGN intrinsic conversion transformation to the specified
node. This node must be a SIGN IntrinsicCall. The SIGN
IntrinsicCall is converted to equivalent inline code. This
is implemented as a PSyIR transform from:

\begin{sphinxVerbatim}[commandchars=\\\{\}]
\PYG{n}{R} \PYG{o}{=} \PYG{o}{.}\PYG{o}{.}\PYG{o}{.} \PYG{n}{SIGN}\PYG{p}{(}\PYG{n}{A}\PYG{p}{,} \PYG{n}{B}\PYG{p}{)} \PYG{o}{.}\PYG{o}{.}\PYG{o}{.}
\end{sphinxVerbatim}

\sphinxAtStartPar
to:

\begin{sphinxVerbatim}[commandchars=\\\{\}]
\PYG{n}{tmp\PYGZus{}abs} \PYG{o}{=} \PYG{n}{A}
\PYG{k}{if} \PYG{n}{tmp\PYGZus{}abs} \PYG{o}{\PYGZlt{}} \PYG{l+m+mf}{0.0}\PYG{p}{:}
    \PYG{n}{res\PYGZus{}abs} \PYG{o}{=} \PYG{n}{tmp\PYGZus{}abs}\PYG{o}{*}\PYG{o}{\PYGZhy{}}\PYG{l+m+mf}{1.0}
\PYG{k}{else}\PYG{p}{:}
    \PYG{n}{res\PYGZus{}abs} \PYG{o}{=} \PYG{n}{tmp\PYGZus{}abs}
\PYG{n}{res\PYGZus{}sign} \PYG{o}{=} \PYG{n}{res\PYGZus{}abs}
\PYG{n}{tmp\PYGZus{}sign} \PYG{o}{=} \PYG{n}{B}
\PYG{k}{if} \PYG{n}{tmp\PYGZus{}sign} \PYG{o}{\PYGZlt{}} \PYG{l+m+mf}{0.0}\PYG{p}{:}
    \PYG{n}{res\PYGZus{}sign} \PYG{o}{=} \PYG{n}{res\PYGZus{}sign}\PYG{o}{*}\PYG{o}{\PYGZhy{}}\PYG{l+m+mf}{1.0}
\PYG{n}{R} \PYG{o}{=} \PYG{o}{.}\PYG{o}{.}\PYG{o}{.} \PYG{n}{res\PYGZus{}sign} \PYG{o}{.}\PYG{o}{.}\PYG{o}{.}
\end{sphinxVerbatim}

\sphinxAtStartPar
where \sphinxcode{\sphinxupquote{A}} and \sphinxcode{\sphinxupquote{B}} could be arbitrarily complex PSyIR
expressions, \sphinxcode{\sphinxupquote{...}} could be arbitrary PSyIR code and where
\sphinxcode{\sphinxupquote{ABS}} has been replaced with inline code by the NemoAbsTrans
transformation.

\sphinxAtStartPar
This transformation requires the IntrinsicCall node to be a child of
an assignment and will raise an exception if this is not the case.
\begin{quote}\begin{description}
\sphinxlineitem{Parameters}\begin{itemize}
\item {} 
\sphinxAtStartPar
\sphinxstyleliteralstrong{\sphinxupquote{node}} (\sphinxcode{\sphinxupquote{psyclone.psyir.nodes.IntrinsicCall}}) \textendash{} a SIGN IntrinsicCall node.

\item {} 
\sphinxAtStartPar
\sphinxstyleliteralstrong{\sphinxupquote{symbol\_table}} ({\hyperref[\detokenize{psyir:psyclone.psyir.symbols.SymbolTable}]{\sphinxcrossref{\sphinxcode{\sphinxupquote{psyclone.psyir.symbols.SymbolTable}}}}}) \textendash{} the symbol table.

\item {} 
\sphinxAtStartPar
\sphinxstyleliteralstrong{\sphinxupquote{options}} (\sphinxstyleliteralemphasis{\sphinxupquote{Optional}}\sphinxstyleliteralemphasis{\sphinxupquote{{[}}}\sphinxstyleliteralemphasis{\sphinxupquote{Dict}}\sphinxstyleliteralemphasis{\sphinxupquote{{[}}}\sphinxhref{https://docs.python.org/3/library/stdtypes.html\#str}{\sphinxstyleliteralemphasis{\sphinxupquote{str}}}\sphinxstyleliteralemphasis{\sphinxupquote{, }}\sphinxstyleliteralemphasis{\sphinxupquote{Any}}\sphinxstyleliteralemphasis{\sphinxupquote{{]}}}\sphinxstyleliteralemphasis{\sphinxupquote{{]}}}) \textendash{} a dictionary with options for transformations.

\end{itemize}

\end{description}\end{quote}

\end{fulllineitems}


\end{fulllineitems}



\bigskip\hrule\bigskip



\begin{fulllineitems}

\pysigstartsignatures
\pysigline{\sphinxbfcode{\sphinxupquote{class\DUrole{w,w}{  }}}\sphinxcode{\sphinxupquote{psyclone.psyir.transformations.}}\sphinxbfcode{\sphinxupquote{Sum2LoopTrans}}}
\pysigstopsignatures
\sphinxAtStartPar
Provides a transformation from a PSyIR SUM IntrinsicCall node to an
equivalent PSyIR loop structure that is suitable for running in
parallel on CPUs and GPUs. Validity checks are also performed.

\sphinxAtStartPar
If SUM contains a single positional argument which is an array,
all elements of that array are summed and the result returned in
the scalar R.

\begin{sphinxVerbatim}[commandchars=\\\{\}]
\PYG{n}{R}\PYG{+w}{ }\PYG{o}{=}\PYG{+w}{ }\PYG{n+nb}{SUM}\PYG{p}{(}\PYG{k}{ARRAY}\PYG{p}{)}
\end{sphinxVerbatim}

\sphinxAtStartPar
For example, if the array is two dimensional, the equivalent code
for real data is:

\begin{sphinxVerbatim}[commandchars=\\\{\}]
\PYG{n}{R}\PYG{+w}{ }\PYG{o}{=}\PYG{+w}{ }\PYG{l+m+mf}{0.0}
\PYG{k}{DO }\PYG{n}{J}\PYG{o}{=}\PYG{n+nb}{LBOUND}\PYG{p}{(}\PYG{k}{ARRAY}\PYG{p}{,}\PYG{l+m+mi}{2}\PYG{p}{)}\PYG{p}{,}\PYG{n+nb}{UBOUND}\PYG{p}{(}\PYG{k}{ARRAY}\PYG{p}{,}\PYG{l+m+mi}{2}\PYG{p}{)}
\PYG{+w}{  }\PYG{k}{DO }\PYG{n}{I}\PYG{o}{=}\PYG{n+nb}{LBOUND}\PYG{p}{(}\PYG{k}{ARRAY}\PYG{p}{,}\PYG{l+m+mi}{1}\PYG{p}{)}\PYG{p}{,}\PYG{n+nb}{UBOUND}\PYG{p}{(}\PYG{k}{ARRAY}\PYG{p}{,}\PYG{l+m+mi}{1}\PYG{p}{)}
\PYG{+w}{    }\PYG{n}{R}\PYG{+w}{ }\PYG{o}{=}\PYG{+w}{ }\PYG{n}{R}\PYG{+w}{ }\PYG{o}{+}\PYG{+w}{ }\PYG{k}{ARRAY}\PYG{p}{(}\PYG{n}{I}\PYG{p}{,}\PYG{n}{J}\PYG{p}{)}
\end{sphinxVerbatim}

\sphinxAtStartPar
If the mask argument is provided then the mask is used to
determine whether the sum is applied:

\begin{sphinxVerbatim}[commandchars=\\\{\}]
\PYG{n}{R}\PYG{+w}{ }\PYG{o}{=}\PYG{+w}{ }\PYG{n+nb}{SUM}\PYG{p}{(}\PYG{k}{ARRAY}\PYG{p}{,}\PYG{+w}{ }\PYG{n}{mask}\PYG{o}{=}\PYG{n+nb}{MOD}\PYG{p}{(}\PYG{k}{ARRAY}\PYG{p}{,}\PYG{+w}{ }\PYG{l+m+mf}{2.0}\PYG{p}{)}\PYG{o}{==}\PYG{l+m+mi}{1}\PYG{p}{)}
\end{sphinxVerbatim}

\sphinxAtStartPar
If the array is two dimensional, the equivalent code
for real data is:

\begin{sphinxVerbatim}[commandchars=\\\{\}]
\PYG{n}{R}\PYG{+w}{ }\PYG{o}{=}\PYG{+w}{ }\PYG{l+m+mf}{0.0}
\PYG{k}{DO }\PYG{n}{J}\PYG{o}{=}\PYG{n+nb}{LBOUND}\PYG{p}{(}\PYG{k}{ARRAY}\PYG{p}{,}\PYG{l+m+mi}{2}\PYG{p}{)}\PYG{p}{,}\PYG{n+nb}{UBOUND}\PYG{p}{(}\PYG{k}{ARRAY}\PYG{p}{,}\PYG{l+m+mi}{2}\PYG{p}{)}
\PYG{+w}{  }\PYG{k}{DO }\PYG{n}{I}\PYG{o}{=}\PYG{n+nb}{LBOUND}\PYG{p}{(}\PYG{k}{ARRAY}\PYG{p}{,}\PYG{l+m+mi}{1}\PYG{p}{)}\PYG{p}{,}\PYG{n+nb}{UBOUND}\PYG{p}{(}\PYG{k}{ARRAY}\PYG{p}{,}\PYG{l+m+mi}{1}\PYG{p}{)}
\PYG{+w}{    }\PYG{k}{IF}\PYG{+w}{ }\PYG{p}{(}\PYG{n+nb}{MOD}\PYG{p}{(}\PYG{k}{ARRAY}\PYG{p}{(}\PYG{n}{I}\PYG{p}{,}\PYG{n}{J}\PYG{p}{)}\PYG{p}{,}\PYG{+w}{ }\PYG{l+m+mf}{2.0}\PYG{p}{)}\PYG{o}{==}\PYG{l+m+mi}{1}\PYG{p}{)}\PYG{+w}{ }\PYG{k}{THEN}
\PYG{k}{      }\PYG{n}{R}\PYG{+w}{ }\PYG{o}{=}\PYG{+w}{ }\PYG{n}{R}\PYG{+w}{ }\PYG{o}{+}\PYG{+w}{ }\PYG{k}{ARRAY}\PYG{p}{(}\PYG{n}{I}\PYG{p}{,}\PYG{n}{J}\PYG{p}{)}
\end{sphinxVerbatim}

\sphinxAtStartPar
The dimension argument is currently not supported and will result
in a TransformationError exception being raised.

\begin{sphinxVerbatim}[commandchars=\\\{\}]
\PYG{n}{R}\PYG{+w}{ }\PYG{o}{=}\PYG{+w}{ }\PYG{n+nb}{SUM}\PYG{p}{(}\PYG{k}{ARRAY}\PYG{p}{,}\PYG{+w}{ }\PYG{k}{dimension}\PYG{o}{=}\PYG{l+m+mi}{2}\PYG{p}{)}
\end{sphinxVerbatim}

\sphinxAtStartPar
The array passed to MAXVAL may use any combination of array
syntax, array notation, array sections and scalar bounds:

\begin{sphinxVerbatim}[commandchars=\\\{\}]
\PYG{n}{R}\PYG{+w}{ }\PYG{o}{=}\PYG{+w}{ }\PYG{n+nb}{SUM}\PYG{p}{(}\PYG{k}{ARRAY}\PYG{p}{)}\PYG{+w}{ }\PYG{c}{! array syntax}
\PYG{n}{R}\PYG{+w}{ }\PYG{o}{=}\PYG{+w}{ }\PYG{n+nb}{SUM}\PYG{p}{(}\PYG{k}{ARRAY}\PYG{p}{(}\PYG{p}{:}\PYG{p}{,}\PYG{p}{:}\PYG{p}{)}\PYG{p}{)}\PYG{+w}{ }\PYG{c}{! array notation}
\PYG{n}{R}\PYG{+w}{ }\PYG{o}{=}\PYG{+w}{ }\PYG{n+nb}{SUM}\PYG{p}{(}\PYG{k}{ARRAY}\PYG{p}{(}\PYG{l+m+mi}{1}\PYG{p}{:}\PYG{l+m+mi}{10}\PYG{p}{,}\PYG{n}{lo}\PYG{p}{:}\PYG{n}{hi}\PYG{p}{)}\PYG{p}{)}\PYG{+w}{ }\PYG{c}{! array sections}
\PYG{n}{R}\PYG{+w}{ }\PYG{o}{=}\PYG{+w}{ }\PYG{n+nb}{SUM}\PYG{p}{(}\PYG{k}{ARRAY}\PYG{p}{(}\PYG{l+m+mi}{1}\PYG{p}{:}\PYG{l+m+mi}{10}\PYG{p}{,}\PYG{p}{:}\PYG{p}{)}\PYG{p}{)}\PYG{+w}{ }\PYG{c}{! mix of array section and array notation}
\PYG{n}{R}\PYG{+w}{ }\PYG{o}{=}\PYG{+w}{ }\PYG{n+nb}{SUM}\PYG{p}{(}\PYG{k}{ARRAY}\PYG{p}{(}\PYG{l+m+mi}{1}\PYG{p}{:}\PYG{l+m+mi}{10}\PYG{p}{,}\PYG{l+m+mi}{2}\PYG{p}{)}\PYG{p}{)}\PYG{+w}{ }\PYG{c}{! mix of array section and scalar bound}
\end{sphinxVerbatim}

\sphinxAtStartPar
For example:

\begin{sphinxVerbatim}[commandchars=\\\{\}]
\PYG{g+gp}{\PYGZgt{}\PYGZgt{}\PYGZgt{} }\PYG{k+kn}{from}\PYG{+w}{ }\PYG{n+nn}{psyclone}\PYG{n+nn}{.}\PYG{n+nn}{psyir}\PYG{n+nn}{.}\PYG{n+nn}{backend}\PYG{n+nn}{.}\PYG{n+nn}{fortran}\PYG{+w}{ }\PYG{k+kn}{import} \PYG{n}{FortranWriter}
\PYG{g+gp}{\PYGZgt{}\PYGZgt{}\PYGZgt{} }\PYG{k+kn}{from}\PYG{+w}{ }\PYG{n+nn}{psyclone}\PYG{n+nn}{.}\PYG{n+nn}{psyir}\PYG{n+nn}{.}\PYG{n+nn}{frontend}\PYG{n+nn}{.}\PYG{n+nn}{fortran}\PYG{+w}{ }\PYG{k+kn}{import} \PYG{n}{FortranReader}
\PYG{g+gp}{\PYGZgt{}\PYGZgt{}\PYGZgt{} }\PYG{k+kn}{from}\PYG{+w}{ }\PYG{n+nn}{psyclone}\PYG{n+nn}{.}\PYG{n+nn}{psyir}\PYG{n+nn}{.}\PYG{n+nn}{transformations}\PYG{+w}{ }\PYG{k+kn}{import} \PYG{n}{Sum2LoopTrans}
\PYG{g+gp}{\PYGZgt{}\PYGZgt{}\PYGZgt{} }\PYG{n}{code} \PYG{o}{=} \PYG{p}{(}\PYG{l+s+s2}{\PYGZdq{}}\PYG{l+s+s2}{subroutine sum\PYGZus{}test(array,n,m)}\PYG{l+s+se}{\PYGZbs{}n}\PYG{l+s+s2}{\PYGZdq{}}
\PYG{g+gp}{... }        \PYG{l+s+s2}{\PYGZdq{}}\PYG{l+s+s2}{  integer :: n, m}\PYG{l+s+se}{\PYGZbs{}n}\PYG{l+s+s2}{\PYGZdq{}}
\PYG{g+gp}{... }        \PYG{l+s+s2}{\PYGZdq{}}\PYG{l+s+s2}{  real :: array(10,10)}\PYG{l+s+se}{\PYGZbs{}n}\PYG{l+s+s2}{\PYGZdq{}}
\PYG{g+gp}{... }        \PYG{l+s+s2}{\PYGZdq{}}\PYG{l+s+s2}{  real :: result}\PYG{l+s+se}{\PYGZbs{}n}\PYG{l+s+s2}{\PYGZdq{}}
\PYG{g+gp}{... }        \PYG{l+s+s2}{\PYGZdq{}}\PYG{l+s+s2}{  result = sum(array)}\PYG{l+s+se}{\PYGZbs{}n}\PYG{l+s+s2}{\PYGZdq{}}
\PYG{g+gp}{... }        \PYG{l+s+s2}{\PYGZdq{}}\PYG{l+s+s2}{end subroutine}\PYG{l+s+se}{\PYGZbs{}n}\PYG{l+s+s2}{\PYGZdq{}}\PYG{p}{)}
\PYG{g+gp}{\PYGZgt{}\PYGZgt{}\PYGZgt{} }\PYG{n}{psyir} \PYG{o}{=} \PYG{n}{FortranReader}\PYG{p}{(}\PYG{p}{)}\PYG{o}{.}\PYG{n}{psyir\PYGZus{}from\PYGZus{}source}\PYG{p}{(}\PYG{n}{code}\PYG{p}{)}
\PYG{g+gp}{\PYGZgt{}\PYGZgt{}\PYGZgt{} }\PYG{n}{sum\PYGZus{}node} \PYG{o}{=} \PYG{n}{psyir}\PYG{o}{.}\PYG{n}{children}\PYG{p}{[}\PYG{l+m+mi}{0}\PYG{p}{]}\PYG{o}{.}\PYG{n}{children}\PYG{p}{[}\PYG{l+m+mi}{0}\PYG{p}{]}\PYG{o}{.}\PYG{n}{children}\PYG{p}{[}\PYG{l+m+mi}{1}\PYG{p}{]}
\PYG{g+gp}{\PYGZgt{}\PYGZgt{}\PYGZgt{} }\PYG{n}{Sum2LoopTrans}\PYG{p}{(}\PYG{p}{)}\PYG{o}{.}\PYG{n}{apply}\PYG{p}{(}\PYG{n}{sum\PYGZus{}node}\PYG{p}{)}
\PYG{g+gp}{\PYGZgt{}\PYGZgt{}\PYGZgt{} }\PYG{n+nb}{print}\PYG{p}{(}\PYG{n}{FortranWriter}\PYG{p}{(}\PYG{p}{)}\PYG{p}{(}\PYG{n}{psyir}\PYG{p}{)}\PYG{p}{)}
\PYG{g+go}{subroutine sum\PYGZus{}test(array, n, m)}
\PYG{g+go}{  integer :: n}
\PYG{g+go}{  integer :: m}
\PYG{g+go}{  real, dimension(10,10) :: array}
\PYG{g+go}{  real :: result}
\PYG{g+go}{  integer :: idx}
\PYG{g+go}{  integer :: idx\PYGZus{}1}

\PYG{g+go}{  result = 0.0}
\PYG{g+go}{  do idx = 1, 10, 1}
\PYG{g+go}{    do idx\PYGZus{}1 = 1, 10, 1}
\PYG{g+go}{      result = result + array(idx\PYGZus{}1,idx)}
\PYG{g+go}{    enddo}
\PYG{g+go}{  enddo}

\PYG{g+go}{end subroutine sum\PYGZus{}test}
\end{sphinxVerbatim}


\begin{fulllineitems}

\pysigstartsignatures
\pysiglinewithargsret{\sphinxbfcode{\sphinxupquote{apply}}}{\sphinxparam{\DUrole{n,n}{node}}\sphinxparamcomma \sphinxparam{\DUrole{n,n}{options}\DUrole{o,o}{=}\DUrole{default_value}{None}}}{}
\pysigstopsignatures
\sphinxAtStartPar
Apply the array\sphinxhyphen{}reduction intrinsic conversion transformation to
the specified node. This node must be one of these intrinsic
operations which is converted to an equivalent loop structure.
\begin{quote}\begin{description}
\sphinxlineitem{Parameters}\begin{itemize}
\item {} 
\sphinxAtStartPar
\sphinxstyleliteralstrong{\sphinxupquote{node}} (\sphinxcode{\sphinxupquote{psyclone.psyir.nodes.IntrinsicCall}}) \textendash{} an array\sphinxhyphen{}reduction intrinsic.

\item {} 
\sphinxAtStartPar
\sphinxstyleliteralstrong{\sphinxupquote{options}} (\sphinxstyleliteralemphasis{\sphinxupquote{Optional}}\sphinxstyleliteralemphasis{\sphinxupquote{{[}}}\sphinxstyleliteralemphasis{\sphinxupquote{Dict}}\sphinxstyleliteralemphasis{\sphinxupquote{{[}}}\sphinxhref{https://docs.python.org/3/library/stdtypes.html\#str}{\sphinxstyleliteralemphasis{\sphinxupquote{str}}}\sphinxstyleliteralemphasis{\sphinxupquote{, }}\sphinxstyleliteralemphasis{\sphinxupquote{Any}}\sphinxstyleliteralemphasis{\sphinxupquote{{]}}}\sphinxstyleliteralemphasis{\sphinxupquote{{]}}}) \textendash{} options for the transformation.

\end{itemize}

\end{description}\end{quote}

\end{fulllineitems}


\end{fulllineitems}



\bigskip\hrule\bigskip



\begin{fulllineitems}

\pysigstartsignatures
\pysigline{\sphinxbfcode{\sphinxupquote{class\DUrole{w,w}{  }}}\sphinxcode{\sphinxupquote{psyclone.psyir.transformations.}}\sphinxbfcode{\sphinxupquote{ScalarisationTrans}}}
\pysigstopsignatures
\sphinxAtStartPar
This transformation takes a Loop and converts any array accesses
to scalar if the results of the loop are unused, and the initial value
is unused. For example in the following snippet the value of a(i)
is only used inside the loop, so can be turned into a scalar, wheras
the values of b(i) are used in the following loop so are kept as an array:

\begin{sphinxVerbatim}[commandchars=\\\{\}]
\PYG{g+gp}{\PYGZgt{}\PYGZgt{}\PYGZgt{} }\PYG{k+kn}{from}\PYG{+w}{ }\PYG{n+nn}{psyclone}\PYG{n+nn}{.}\PYG{n+nn}{psyir}\PYG{n+nn}{.}\PYG{n+nn}{backend}\PYG{n+nn}{.}\PYG{n+nn}{fortran}\PYG{+w}{ }\PYG{k+kn}{import} \PYG{n}{FortranWriter}
\PYG{g+gp}{\PYGZgt{}\PYGZgt{}\PYGZgt{} }\PYG{k+kn}{from}\PYG{+w}{ }\PYG{n+nn}{psyclone}\PYG{n+nn}{.}\PYG{n+nn}{psyir}\PYG{n+nn}{.}\PYG{n+nn}{frontend}\PYG{n+nn}{.}\PYG{n+nn}{fortran}\PYG{+w}{ }\PYG{k+kn}{import} \PYG{n}{FortranReader}
\PYG{g+gp}{\PYGZgt{}\PYGZgt{}\PYGZgt{} }\PYG{k+kn}{from}\PYG{+w}{ }\PYG{n+nn}{psyclone}\PYG{n+nn}{.}\PYG{n+nn}{psyir}\PYG{n+nn}{.}\PYG{n+nn}{transformations}\PYG{+w}{ }\PYG{k+kn}{import} \PYG{n}{ScalarisationTrans}
\PYG{g+gp}{\PYGZgt{}\PYGZgt{}\PYGZgt{} }\PYG{k+kn}{from}\PYG{+w}{ }\PYG{n+nn}{psyclone}\PYG{n+nn}{.}\PYG{n+nn}{psyir}\PYG{n+nn}{.}\PYG{n+nn}{nodes}\PYG{+w}{ }\PYG{k+kn}{import} \PYG{n}{Loop}
\PYG{g+gp}{\PYGZgt{}\PYGZgt{}\PYGZgt{} }\PYG{n}{code} \PYG{o}{=} \PYG{l+s+s2}{\PYGZdq{}\PYGZdq{}\PYGZdq{}}\PYG{l+s+s2}{program test}
\PYG{g+gp}{... }\PYG{l+s+s2}{integer :: i,j}
\PYG{g+gp}{... }\PYG{l+s+s2}{real :: a(100), b(100)}
\PYG{g+gp}{... }\PYG{l+s+s2}{do i = 1,100}
\PYG{g+gp}{... }\PYG{l+s+s2}{  a(i) = i}
\PYG{g+gp}{... }\PYG{l+s+s2}{  b(i) = a(i) * a(i)}
\PYG{g+gp}{... }\PYG{l+s+s2}{end do}
\PYG{g+gp}{... }\PYG{l+s+s2}{do j = 1, 100}
\PYG{g+gp}{... }\PYG{l+s+s2}{ if(b(i) \PYGZgt{} 200) then}
\PYG{g+gp}{... }\PYG{l+s+s2}{   print *, b(i)}
\PYG{g+gp}{... }\PYG{l+s+s2}{ end if}
\PYG{g+gp}{... }\PYG{l+s+s2}{end do}
\PYG{g+gp}{... }\PYG{l+s+s2}{end program}\PYG{l+s+s2}{\PYGZdq{}\PYGZdq{}\PYGZdq{}}
\PYG{g+gp}{\PYGZgt{}\PYGZgt{}\PYGZgt{} }\PYG{n}{psyir} \PYG{o}{=} \PYG{n}{FortranReader}\PYG{p}{(}\PYG{p}{)}\PYG{o}{.}\PYG{n}{psyir\PYGZus{}from\PYGZus{}source}\PYG{p}{(}\PYG{n}{code}\PYG{p}{)}
\PYG{g+gp}{\PYGZgt{}\PYGZgt{}\PYGZgt{} }\PYG{n}{scalarise} \PYG{o}{=} \PYG{n}{ScalarisationTrans}\PYG{p}{(}\PYG{p}{)}
\PYG{g+gp}{\PYGZgt{}\PYGZgt{}\PYGZgt{} }\PYG{n}{scalarise}\PYG{o}{.}\PYG{n}{apply}\PYG{p}{(}\PYG{n}{psyir}\PYG{o}{.}\PYG{n}{walk}\PYG{p}{(}\PYG{n}{Loop}\PYG{p}{)}\PYG{p}{[}\PYG{l+m+mi}{0}\PYG{p}{]}\PYG{p}{)}
\PYG{g+gp}{\PYGZgt{}\PYGZgt{}\PYGZgt{} }\PYG{n+nb}{print}\PYG{p}{(}\PYG{n}{FortranWriter}\PYG{p}{(}\PYG{p}{)}\PYG{p}{(}\PYG{n}{psyir}\PYG{p}{)}\PYG{p}{)}
\PYG{g+go}{program test}
\PYG{g+go}{  integer :: i}
\PYG{g+go}{  integer :: j}
\PYG{g+go}{  real, dimension(100) :: a}
\PYG{g+go}{  real, dimension(100) :: b}
\PYG{g+go}{  real :: a\PYGZus{}scalar}

\PYG{g+go}{  do i = 1, 100, 1}
\PYG{g+go}{    a\PYGZus{}scalar = i}
\PYG{g+go}{    b(i) = a\PYGZus{}scalar * a\PYGZus{}scalar}
\PYG{g+go}{  enddo}
\PYG{g+go}{  do j = 1, 100, 1}
\PYG{g+go}{    if (b(i) \PYGZgt{} 200) then}
\PYG{g+go}{      ! PSyclone CodeBlock (unsupported code) reason:}
\PYG{g+go}{      !  \PYGZhy{} Unsupported statement: Print\PYGZus{}Stmt}
\PYG{g+go}{      PRINT *, b(i)}
\PYG{g+go}{    end if}
\PYG{g+go}{  enddo}

\PYG{g+go}{end program test}
\end{sphinxVerbatim}


\begin{fulllineitems}

\pysigstartsignatures
\pysiglinewithargsret{\sphinxbfcode{\sphinxupquote{apply}}}{\sphinxparam{\DUrole{n,n}{node}}\sphinxparamcomma \sphinxparam{\DUrole{n,n}{options}\DUrole{o,o}{=}\DUrole{default_value}{None}}}{}
\pysigstopsignatures
\sphinxAtStartPar
Apply the scalarisation transformation to a loop.
All of the array accesses that are identified as being able to be
scalarised will be transformed by this transformation.

\sphinxAtStartPar
An array access will be scalarised if:
1. All accesses to the array use the same indexing statement.
2. All References contained in the indexing statement are not modified
inside of the loop (loop variables are ok).
3. The array symbol is either not accessed again or is written to
as its next access. If the next access is inside a conditional
that is not an ancestor of the input loop, then PSyclone will
assume that we cannot scalarise that value instead of attempting to
understand the control flow.
4. The array symbol is a local variable.
\begin{quote}\begin{description}
\sphinxlineitem{Parameters}\begin{itemize}
\item {} 
\sphinxAtStartPar
\sphinxstyleliteralstrong{\sphinxupquote{node}} (\DUrole{sphinx_autodoc_typehints-type}{\sphinxcode{\sphinxupquote{Loop}}}) \textendash{} the supplied loop to apply scalarisation to.

\item {} 
\sphinxAtStartPar
\sphinxstyleliteralstrong{\sphinxupquote{options}} (\DUrole{sphinx_autodoc_typehints-type}{\sphinxhref{https://docs.python.org/3/library/typing.html\#typing.Optional}{\sphinxcode{\sphinxupquote{Optional}}}{[}\sphinxhref{https://docs.python.org/3/library/typing.html\#typing.Dict}{\sphinxcode{\sphinxupquote{Dict}}}{[}\sphinxhref{https://docs.python.org/3/library/stdtypes.html\#str}{\sphinxcode{\sphinxupquote{str}}}, \sphinxhref{https://docs.python.org/3/library/typing.html\#typing.Any}{\sphinxcode{\sphinxupquote{Any}}}{]}{]}}) \textendash{} a dictionary with options for transformations.

\end{itemize}

\sphinxlineitem{Return type}
\sphinxAtStartPar
\DUrole{sphinx_autodoc_typehints-type}{\sphinxhref{https://docs.python.org/3/library/constants.html\#None}{\sphinxcode{\sphinxupquote{None}}}}

\end{description}\end{quote}

\end{fulllineitems}


\end{fulllineitems}



\bigskip\hrule\bigskip



\section{Algorithm\sphinxhyphen{}layer}
\label{\detokenize{transformations:algorithm-layer}}
\sphinxAtStartPar
The gocean API supports the transformation of the algorithm
layer. In the future the LFRic API will also support this.
The ability to transformation the algorithm layer is new and
at this time no relevant transformations have been developed.


\section{Kernels}
\label{\detokenize{transformations:kernels}}
\sphinxAtStartPar
PSyclone supports the transformation of Kernels as well as PSy\sphinxhyphen{}layer
code. However, the transformation of kernels to produce new kernels
brings with it additional considerations, especially regarding the
naming of the resulting kernels. PSyclone supports two use cases:
\begin{enumerate}
\sphinxsetlistlabels{\arabic}{enumi}{enumii}{}{.}%
\item {} 
\sphinxAtStartPar
the HPC expert wishes to optimise the same kernel in different ways,
depending on where/how it is called;

\item {} 
\sphinxAtStartPar
the HPC expert wishes to transform the kernel just once and have the
new version used throughout the Algorithm file.

\end{enumerate}

\sphinxAtStartPar
The second case is really an optimisation of the first for the case
where the same set of transformations is applied to every instance of
a given kernel.

\sphinxAtStartPar
Since PSyclone is run separately for each Algorithm in a given
application, ensuring that there are no name clashes for kernels in
the application as a whole requires that some state is maintained
between PSyclone invocations. This is achieved by requiring that the
same kernel output directory is used for every invocation of PSyclone
when building a given application. However, this is under the control
of the user and therefore it is possible to use the same output
directory for a subset of algorithms that require the same kernel
transformation and then a different directory for another subset
requiring a different transformation. Of course, such use would
require care when building and linking the application since the
differently\sphinxhyphen{}optimised kernels would have the same names.

\sphinxAtStartPar
By default, transformed kernels are written to the current working
directory. Alternatively, the user may specify the location to which
to write the modified code via the \sphinxcode{\sphinxupquote{\sphinxhyphen{}okern}} command\sphinxhyphen{}line flag.

\sphinxAtStartPar
In order to support the two use cases given above, PSyclone supports
two different kernel\sphinxhyphen{}renaming schemes: “multiple” and “single”
(specified via the \sphinxcode{\sphinxupquote{\sphinxhyphen{}\sphinxhyphen{}kernel\sphinxhyphen{}renaming}} command\sphinxhyphen{}line flag). In the
default, “multiple” scheme, PSyclone ensures that each transformed
kernel is given a unique name (with reference to the contents of the
kernel output directory). In the “single” scheme, it is assumed that
any given kernel that is transformed is always transformed in the same
way (or left unchanged) and thus just one transformed version of it is
created. This assumption is checked by examining the Fortran code for
any pre\sphinxhyphen{}existing transformed version of that kernel. If another
transformed version of that kernel exists and does not match that
created by the current transformation then PSyclone will raise an
exception.


\subsection{Rules}
\label{\detokenize{transformations:rules}}
\sphinxAtStartPar
Kernel code that is to be transformed is subject to certain
restrictions. These rules are intended to make kernel transformations
as robust as possible, in particular by limiting the amount of
code that must be parsed by PSyclone (via fparser). The rules are
as follows:
\begin{enumerate}
\sphinxsetlistlabels{\arabic}{enumi}{enumii}{}{)}%
\item {} 
\sphinxAtStartPar
Any variable or procedure accessed by a kernel must either be explicitly
declared or named in the \sphinxcode{\sphinxupquote{only}} clause of a module \sphinxcode{\sphinxupquote{use}} statement
within the scope of the subroutine containing the kernel implementation.
This means that:
\begin{enumerate}
\sphinxsetlistlabels{\arabic}{enumii}{enumiii}{}{)}%
\item {} 
\sphinxAtStartPar
Kernel subroutines are forbidden from accessing data using COMMON
blocks;

\item {} 
\sphinxAtStartPar
Kernel subroutines are forbidden from calling procedures declared via
the EXTERN statement;

\item {} 
\sphinxAtStartPar
Kernel subroutines must not access data or procedures made available
via their parent (containing) module.

\end{enumerate}

\item {} 
\sphinxAtStartPar
The full Fortran source of a kernel must be available to PSyclone.
This includes the source of any modules from which it accesses
either routines or data. (However, kernel routines are permitted to make
use of Fortran intrinsic routines.)

\end{enumerate}

\sphinxAtStartPar
For instance, consider the following Fortran module containing the
\sphinxcode{\sphinxupquote{bc\_ssh\_code}} kernel:

\begin{sphinxVerbatim}[commandchars=\\\{\}]
\PYG{k}{module }\PYG{n}{boundary\PYGZus{}conditions\PYGZus{}mod}
\PYG{+w}{  }\PYG{k+kt}{real}\PYG{+w}{ }\PYG{k+kd}{::}\PYG{+w}{ }\PYG{n}{forbidden\PYGZus{}var}

\PYG{k}{contains}

\PYG{k}{  }\PYG{k}{subroutine }\PYG{n}{bc\PYGZus{}ssh\PYGZus{}code}\PYG{p}{(}\PYG{n}{ji}\PYG{p}{,}\PYG{+w}{ }\PYG{n}{jj}\PYG{p}{,}\PYG{+w}{ }\PYG{n}{istep}\PYG{p}{,}\PYG{+w}{ }\PYG{n}{ssha}\PYG{p}{)}
\PYG{+w}{    }\PYG{k}{use }\PYG{n}{kind\PYGZus{}params\PYGZus{}mod}\PYG{p}{,}\PYG{+w}{ }\PYG{k}{only}\PYG{p}{:}\PYG{+w}{ }\PYG{n}{go\PYGZus{}wp}
\PYG{+w}{    }\PYG{k}{use }\PYG{n}{model\PYGZus{}mod}\PYG{p}{,}\PYG{+w}{ }\PYG{k}{only}\PYG{p}{:}\PYG{+w}{ }\PYG{n}{rdt}
\PYG{+w}{    }\PYG{k+kt}{integer}\PYG{p}{,}\PYG{+w}{                     }\PYG{k}{intent}\PYG{p}{(}\PYG{n}{in}\PYG{p}{)}\PYG{+w}{    }\PYG{k+kd}{::}\PYG{+w}{ }\PYG{n}{ji}\PYG{p}{,}\PYG{+w}{ }\PYG{n}{jj}\PYG{p}{,}\PYG{+w}{ }\PYG{n}{istep}
\PYG{+w}{    }\PYG{k+kt}{real}\PYG{p}{(}\PYG{n}{go\PYGZus{}wp}\PYG{p}{)}\PYG{p}{,}\PYG{+w}{ }\PYG{k}{dimension}\PYG{p}{(}\PYG{p}{:}\PYG{p}{,}\PYG{p}{:}\PYG{p}{)}\PYG{p}{,}\PYG{+w}{ }\PYG{k}{intent}\PYG{p}{(}\PYG{n}{inout}\PYG{p}{)}\PYG{+w}{ }\PYG{k+kd}{::}\PYG{+w}{ }\PYG{n}{ssha}
\PYG{+w}{    }\PYG{k+kt}{real}\PYG{p}{(}\PYG{n}{go\PYGZus{}wp}\PYG{p}{)}\PYG{+w}{ }\PYG{k+kd}{::}\PYG{+w}{ }\PYG{n}{rtime}

\PYG{+w}{    }\PYG{n}{rtime}\PYG{+w}{ }\PYG{o}{=}\PYG{+w}{ }\PYG{k+kt}{real}\PYG{p}{(}\PYG{n}{istep}\PYG{p}{,}\PYG{+w}{ }\PYG{n}{go\PYGZus{}wp}\PYG{p}{)}\PYG{+w}{ }\PYG{o}{*}\PYG{+w}{ }\PYG{n}{rdt}
\PYG{+w}{    }\PYG{p}{.}\PYG{p}{.}\PYG{p}{.}
\PYG{+w}{  }\PYG{k}{end }\PYG{k}{subroutine }\PYG{n}{bc\PYGZus{}ssh\PYGZus{}code}

\PYG{k}{end }\PYG{k}{module }\PYG{n}{boundary\PYGZus{}conditions\PYGZus{}mod}
\end{sphinxVerbatim}

\sphinxAtStartPar
Since the kernel subroutine accesses data (the \sphinxcode{\sphinxupquote{rdt}} variable) from
the \sphinxcode{\sphinxupquote{model\_mod}} module, the source of that module must be available
to PSyclone if a transformation is applied to this kernel. Should
\sphinxcode{\sphinxupquote{rdt}} not actually be defined in \sphinxcode{\sphinxupquote{model\_mod}} (i.e. \sphinxcode{\sphinxupquote{model\_mod}}
itself imports it from another module) then the source containing its
definition must also be available to PSyclone. Note that the rules
forbid the \sphinxcode{\sphinxupquote{bc\_ssh\_code}} kernel from accessing the \sphinxcode{\sphinxupquote{forbidden\_var}}
variable that is available to it from the enclosing module scope.

\begin{sphinxadmonition}{note}{Note:}
\sphinxAtStartPar
these rules \sphinxstyleemphasis{only} apply to kernels that are the target of
PSyclone kernel transformations.
\end{sphinxadmonition}


\subsection{Available Kernel Transformations}
\label{\detokenize{transformations:available-kernel-transformations}}\label{\detokenize{transformations:available-kernel-trans}}
\sphinxAtStartPar
The transformations listed below have to be applied specifically to a PSyclone
kernel. There are a number of transformations not listed here that can be
applied to either or both the PSy\sphinxhyphen{}layer and Kernel\sphinxhyphen{}layer PSyIR.

\begin{sphinxadmonition}{note}{Note:}
\sphinxAtStartPar
Some of these transformations modify the PSyIR tree of both the
InvokeSchedule where the transformed CodedKernel is located and its
associated KernelSchedule.
\end{sphinxadmonition}


\bigskip\hrule\bigskip



\begin{fulllineitems}

\pysigstartsignatures
\pysigline{\sphinxbfcode{\sphinxupquote{class\DUrole{w,w}{  }}}\sphinxcode{\sphinxupquote{psyclone.transformations.}}\sphinxbfcode{\sphinxupquote{ACCRoutineTrans}}}
\pysigstopsignatures
\sphinxAtStartPar
Transform a kernel or routine by adding a “!\$acc routine” directive
(causing it to be compiled for the OpenACC accelerator device).
For example:

\begin{sphinxVerbatim}[commandchars=\\\{\}]
\PYG{g+gp}{\PYGZgt{}\PYGZgt{}\PYGZgt{} }\PYG{k+kn}{from}\PYG{+w}{ }\PYG{n+nn}{psyclone}\PYG{n+nn}{.}\PYG{n+nn}{parse}\PYG{n+nn}{.}\PYG{n+nn}{algorithm}\PYG{+w}{ }\PYG{k+kn}{import} \PYG{n}{parse}
\PYG{g+gp}{\PYGZgt{}\PYGZgt{}\PYGZgt{} }\PYG{k+kn}{from}\PYG{+w}{ }\PYG{n+nn}{psyclone}\PYG{n+nn}{.}\PYG{n+nn}{psyGen}\PYG{+w}{ }\PYG{k+kn}{import} \PYG{n}{PSyFactory}
\PYG{g+gp}{\PYGZgt{}\PYGZgt{}\PYGZgt{} }\PYG{n}{api} \PYG{o}{=} \PYG{l+s+s2}{\PYGZdq{}}\PYG{l+s+s2}{gocean}\PYG{l+s+s2}{\PYGZdq{}}
\PYG{g+gp}{\PYGZgt{}\PYGZgt{}\PYGZgt{} }\PYG{n}{ast}\PYG{p}{,} \PYG{n}{invokeInfo} \PYG{o}{=} \PYG{n}{parse}\PYG{p}{(}\PYG{n}{GOCEAN\PYGZus{}SOURCE\PYGZus{}FILE}\PYG{p}{,} \PYG{n}{api}\PYG{o}{=}\PYG{n}{api}\PYG{p}{)}
\PYG{g+gp}{\PYGZgt{}\PYGZgt{}\PYGZgt{} }\PYG{n}{psy} \PYG{o}{=} \PYG{n}{PSyFactory}\PYG{p}{(}\PYG{n}{api}\PYG{p}{)}\PYG{o}{.}\PYG{n}{create}\PYG{p}{(}\PYG{n}{invokeInfo}\PYG{p}{)}
\PYG{g+gp}{\PYGZgt{}\PYGZgt{}\PYGZgt{}}
\PYG{g+gp}{\PYGZgt{}\PYGZgt{}\PYGZgt{} }\PYG{k+kn}{from}\PYG{+w}{ }\PYG{n+nn}{psyclone}\PYG{n+nn}{.}\PYG{n+nn}{transformations}\PYG{+w}{ }\PYG{k+kn}{import} \PYG{n}{ACCRoutineTrans}
\PYG{g+gp}{\PYGZgt{}\PYGZgt{}\PYGZgt{} }\PYG{n}{rtrans} \PYG{o}{=} \PYG{n}{ACCRoutineTrans}\PYG{p}{(}\PYG{p}{)}
\PYG{g+gp}{\PYGZgt{}\PYGZgt{}\PYGZgt{}}
\PYG{g+gp}{\PYGZgt{}\PYGZgt{}\PYGZgt{} }\PYG{n}{schedule} \PYG{o}{=} \PYG{n}{psy}\PYG{o}{.}\PYG{n}{invokes}\PYG{o}{.}\PYG{n}{get}\PYG{p}{(}\PYG{l+s+s1}{\PYGZsq{}}\PYG{l+s+s1}{invoke\PYGZus{}0}\PYG{l+s+s1}{\PYGZsq{}}\PYG{p}{)}\PYG{o}{.}\PYG{n}{schedule}
\PYG{g+gp}{\PYGZgt{}\PYGZgt{}\PYGZgt{} }\PYG{c+c1}{\PYGZsh{} Uncomment the following line to see a text view of the schedule}
\PYG{g+gp}{\PYGZgt{}\PYGZgt{}\PYGZgt{} }\PYG{c+c1}{\PYGZsh{} print(schedule.view())}
\PYG{g+gp}{\PYGZgt{}\PYGZgt{}\PYGZgt{} }\PYG{n}{kern} \PYG{o}{=} \PYG{n}{schedule}\PYG{o}{.}\PYG{n}{children}\PYG{p}{[}\PYG{l+m+mi}{0}\PYG{p}{]}\PYG{o}{.}\PYG{n}{children}\PYG{p}{[}\PYG{l+m+mi}{0}\PYG{p}{]}\PYG{o}{.}\PYG{n}{children}\PYG{p}{[}\PYG{l+m+mi}{0}\PYG{p}{]}
\PYG{g+gp}{\PYGZgt{}\PYGZgt{}\PYGZgt{} }\PYG{c+c1}{\PYGZsh{} Transform the kernel}
\PYG{g+gp}{\PYGZgt{}\PYGZgt{}\PYGZgt{} }\PYG{n}{rtrans}\PYG{o}{.}\PYG{n}{apply}\PYG{p}{(}\PYG{n}{kern}\PYG{p}{)}
\end{sphinxVerbatim}


\begin{fulllineitems}

\pysigstartsignatures
\pysiglinewithargsret{\sphinxbfcode{\sphinxupquote{apply}}}{\sphinxparam{\DUrole{n,n}{node}}\sphinxparamcomma \sphinxparam{\DUrole{n,n}{options}\DUrole{o,o}{=}\DUrole{default_value}{None}}}{}
\pysigstopsignatures
\sphinxAtStartPar
Add the ‘!\$acc routine’ OpenACC directive into the code of the
supplied Kernel (in a PSyKAl API such as GOcean or LFRic) or directly
in the supplied Routine.
\begin{quote}\begin{description}
\sphinxlineitem{Parameters}\begin{itemize}
\item {} 
\sphinxAtStartPar
\sphinxstyleliteralstrong{\sphinxupquote{node}} (\sphinxcode{\sphinxupquote{psyclone.psyGen.Kern}} |
\sphinxcode{\sphinxupquote{psyclone.psyir.nodes.Routine}}) \textendash{} the kernel call or routine implementation to transform.

\item {} 
\sphinxAtStartPar
\sphinxstyleliteralstrong{\sphinxupquote{options}} (\sphinxstyleliteralemphasis{\sphinxupquote{Optional}}\sphinxstyleliteralemphasis{\sphinxupquote{{[}}}\sphinxstyleliteralemphasis{\sphinxupquote{Dict}}\sphinxstyleliteralemphasis{\sphinxupquote{{[}}}\sphinxhref{https://docs.python.org/3/library/stdtypes.html\#str}{\sphinxstyleliteralemphasis{\sphinxupquote{str}}}\sphinxstyleliteralemphasis{\sphinxupquote{, }}\sphinxstyleliteralemphasis{\sphinxupquote{Any}}\sphinxstyleliteralemphasis{\sphinxupquote{{]}}}\sphinxstyleliteralemphasis{\sphinxupquote{{]}}}) \textendash{} a dictionary with options for transformations.

\item {} 
\sphinxAtStartPar
\sphinxstyleliteralstrong{\sphinxupquote{options}}\sphinxstyleliteralstrong{\sphinxupquote{{[}}}\sphinxstyleliteralstrong{\sphinxupquote{"force"}}\sphinxstyleliteralstrong{\sphinxupquote{{]}}} (\sphinxhref{https://docs.python.org/3/library/functions.html\#bool}{\sphinxstyleliteralemphasis{\sphinxupquote{bool}}}) \textendash{} whether to allow routines with
CodeBlocks to run on the GPU.

\item {} 
\sphinxAtStartPar
\sphinxstyleliteralstrong{\sphinxupquote{options}}\sphinxstyleliteralstrong{\sphinxupquote{{[}}}\sphinxstyleliteralstrong{\sphinxupquote{"parallelism"}}\sphinxstyleliteralstrong{\sphinxupquote{{]}}} (\sphinxhref{https://docs.python.org/3/library/stdtypes.html\#str}{\sphinxstyleliteralemphasis{\sphinxupquote{str}}}) \textendash{} the level of parallelism that the
target routine (or a callee) exposes. One of “seq” (the default),
“vector”, “worker” or “gang”.

\end{itemize}

\end{description}\end{quote}

\end{fulllineitems}



\begin{fulllineitems}

\pysigstartsignatures
\pysiglinewithargsret{\sphinxbfcode{\sphinxupquote{validate}}}{\sphinxparam{\DUrole{n,n}{node}}\sphinxparamcomma \sphinxparam{\DUrole{n,n}{options}\DUrole{o,o}{=}\DUrole{default_value}{None}}}{}
\pysigstopsignatures
\sphinxAtStartPar
Perform checks that the supplied kernel or routine can be transformed.
\begin{quote}\begin{description}
\sphinxlineitem{Parameters}\begin{itemize}
\item {} 
\sphinxAtStartPar
\sphinxstyleliteralstrong{\sphinxupquote{node}} (\sphinxcode{\sphinxupquote{psyclone.psyGen.Kern}} |
\sphinxcode{\sphinxupquote{psyclone.psyir.nodes.Routine}}) \textendash{} the kernel or routine which is the target of this
transformation.

\item {} 
\sphinxAtStartPar
\sphinxstyleliteralstrong{\sphinxupquote{options}} (\sphinxstyleliteralemphasis{\sphinxupquote{Optional}}\sphinxstyleliteralemphasis{\sphinxupquote{{[}}}\sphinxstyleliteralemphasis{\sphinxupquote{Dict}}\sphinxstyleliteralemphasis{\sphinxupquote{{[}}}\sphinxhref{https://docs.python.org/3/library/stdtypes.html\#str}{\sphinxstyleliteralemphasis{\sphinxupquote{str}}}\sphinxstyleliteralemphasis{\sphinxupquote{, }}\sphinxstyleliteralemphasis{\sphinxupquote{Any}}\sphinxstyleliteralemphasis{\sphinxupquote{{]}}}\sphinxstyleliteralemphasis{\sphinxupquote{{]}}}) \textendash{} a dictionary with options for transformations.

\item {} 
\sphinxAtStartPar
\sphinxstyleliteralstrong{\sphinxupquote{options}}\sphinxstyleliteralstrong{\sphinxupquote{{[}}}\sphinxstyleliteralstrong{\sphinxupquote{"force"}}\sphinxstyleliteralstrong{\sphinxupquote{{]}}} (\sphinxhref{https://docs.python.org/3/library/functions.html\#bool}{\sphinxstyleliteralemphasis{\sphinxupquote{bool}}}) \textendash{} whether to allow routines with
CodeBlocks to run on the GPU.

\end{itemize}

\sphinxlineitem{Raises}\begin{itemize}
\item {} 
\sphinxAtStartPar
\sphinxstyleliteralstrong{\sphinxupquote{TransformationError}} \textendash{} if the node is not a kernel or a routine.

\item {} 
\sphinxAtStartPar
\sphinxstyleliteralstrong{\sphinxupquote{TransformationError}} \textendash{} if the target is a built\sphinxhyphen{}in kernel.

\item {} 
\sphinxAtStartPar
\sphinxstyleliteralstrong{\sphinxupquote{TransformationError}} \textendash{} if it is a kernel but without an
associated PSyIR.

\item {} 
\sphinxAtStartPar
\sphinxstyleliteralstrong{\sphinxupquote{TransformationError}} \textendash{} if any of the symbols in the kernel are
accessed via a module use statement.

\item {} 
\sphinxAtStartPar
\sphinxstyleliteralstrong{\sphinxupquote{TransformationError}} \textendash{} if the kernel contains any calls to other
routines.

\item {} 
\sphinxAtStartPar
\sphinxstyleliteralstrong{\sphinxupquote{TransformationError}} \textendash{} if the ‘parallelism’ option is supplied
but is not a recognised level of parallelism.

\end{itemize}

\end{description}\end{quote}

\end{fulllineitems}


\end{fulllineitems}



\bigskip\hrule\bigskip



\begin{fulllineitems}

\pysigstartsignatures
\pysigline{\sphinxbfcode{\sphinxupquote{class\DUrole{w,w}{  }}}\sphinxcode{\sphinxupquote{psyclone.psyir.transformations.}}\sphinxbfcode{\sphinxupquote{FoldConditionalReturnExpressionsTrans}}}
\pysigstopsignatures
\sphinxAtStartPar
Provides a transformation that folds conditional expressions with only
a return statement inside so that the Return statement is moved to the end
of the Routine and therefore it can be safely removed. This simplifies the
control flow of the kernel to facilitate other transformations like kernel
fusions. For example, the following code:

\begin{sphinxVerbatim}[commandchars=\\\{\}]
\PYG{k}{subroutine }\PYG{n}{test}\PYG{p}{(}\PYG{n}{i}\PYG{p}{)}
\PYG{+w}{    }\PYG{k}{if}\PYG{+w}{ }\PYG{p}{(}\PYG{n}{i}\PYG{+w}{ }\PYG{o}{\PYGZlt{}}\PYG{+w}{ }\PYG{l+m+mi}{5}\PYG{p}{)}\PYG{+w}{ }\PYG{k}{then}
\PYG{k}{        }\PYG{k}{return}
\PYG{k}{    }\PYG{k}{endif}
\PYG{k}{    }\PYG{k}{if}\PYG{+w}{ }\PYG{p}{(}\PYG{n}{i}\PYG{+w}{ }\PYG{o}{\PYGZgt{}}\PYG{+w}{ }\PYG{l+m+mi}{10}\PYG{p}{)}\PYG{+w}{ }\PYG{k}{then}
\PYG{k}{        }\PYG{k}{return}
\PYG{k}{    }\PYG{k}{endif}
\PYG{+w}{    }\PYG{c}{! CODE}
\PYG{k}{end }\PYG{k}{subroutine}
\end{sphinxVerbatim}

\sphinxAtStartPar
will be transformed to:

\begin{sphinxVerbatim}[commandchars=\\\{\}]
\PYG{k}{subroutine }\PYG{n}{test}\PYG{p}{(}\PYG{n}{i}\PYG{p}{)}
\PYG{+w}{    }\PYG{k}{if}\PYG{+w}{ }\PYG{p}{(}\PYG{p}{.}\PYG{n+nb}{not}\PYG{p}{.}\PYG{p}{(}\PYG{n}{i}\PYG{+w}{ }\PYG{o}{\PYGZlt{}}\PYG{+w}{ }\PYG{l+m+mi}{5}\PYG{p}{)}\PYG{p}{)}\PYG{+w}{ }\PYG{k}{then}
\PYG{k}{        }\PYG{k}{if}\PYG{+w}{ }\PYG{p}{(}\PYG{p}{.}\PYG{n+nb}{not}\PYG{p}{.}\PYG{p}{(}\PYG{n}{i}\PYG{+w}{ }\PYG{o}{\PYGZgt{}}\PYG{+w}{ }\PYG{l+m+mi}{10}\PYG{p}{)}\PYG{p}{)}\PYG{+w}{ }\PYG{k}{then}
\PYG{+w}{            }\PYG{c}{! CODE}
\PYG{+w}{        }\PYG{k}{endif}
\PYG{k}{    }\PYG{k}{endif}
\PYG{k}{end }\PYG{k}{subroutine}
\end{sphinxVerbatim}


\begin{fulllineitems}

\pysigstartsignatures
\pysiglinewithargsret{\sphinxbfcode{\sphinxupquote{apply}}}{\sphinxparam{\DUrole{n,n}{node}}\sphinxparamcomma \sphinxparam{\DUrole{n,n}{options}\DUrole{o,o}{=}\DUrole{default_value}{None}}}{}
\pysigstopsignatures
\sphinxAtStartPar
Apply this transformation to the supplied node.
\begin{quote}\begin{description}
\sphinxlineitem{Parameters}\begin{itemize}
\item {} 
\sphinxAtStartPar
\sphinxstyleliteralstrong{\sphinxupquote{node}} (\sphinxcode{\sphinxupquote{psyclone.psyir.nodes.Routine}}) \textendash{} the node to transform.

\item {} 
\sphinxAtStartPar
\sphinxstyleliteralstrong{\sphinxupquote{options}} (\sphinxstyleliteralemphasis{\sphinxupquote{Optional}}\sphinxstyleliteralemphasis{\sphinxupquote{{[}}}\sphinxstyleliteralemphasis{\sphinxupquote{Dict}}\sphinxstyleliteralemphasis{\sphinxupquote{{[}}}\sphinxhref{https://docs.python.org/3/library/stdtypes.html\#str}{\sphinxstyleliteralemphasis{\sphinxupquote{str}}}\sphinxstyleliteralemphasis{\sphinxupquote{, }}\sphinxstyleliteralemphasis{\sphinxupquote{Any}}\sphinxstyleliteralemphasis{\sphinxupquote{{]}}}\sphinxstyleliteralemphasis{\sphinxupquote{{]}}}) \textendash{} a dictionary with options for transformations.

\end{itemize}

\end{description}\end{quote}

\end{fulllineitems}



\begin{fulllineitems}

\pysigstartsignatures
\pysigline{\sphinxbfcode{\sphinxupquote{property\DUrole{w,w}{  }}}\sphinxbfcode{\sphinxupquote{name}}}
\pysigstopsignatures
\sphinxAtStartPar
Returns the name of this transformation as a string.

\end{fulllineitems}



\begin{fulllineitems}

\pysigstartsignatures
\pysiglinewithargsret{\sphinxbfcode{\sphinxupquote{validate}}}{\sphinxparam{\DUrole{n,n}{node}}\sphinxparamcomma \sphinxparam{\DUrole{n,n}{options}\DUrole{o,o}{=}\DUrole{default_value}{None}}}{}
\pysigstopsignatures
\sphinxAtStartPar
Ensure that it is valid to apply this transformation to the
supplied node.
\begin{quote}\begin{description}
\sphinxlineitem{Parameters}\begin{itemize}
\item {} 
\sphinxAtStartPar
\sphinxstyleliteralstrong{\sphinxupquote{node}} (\sphinxcode{\sphinxupquote{psyclone.psyir.nodes.Routine}}) \textendash{} the node to validate.

\item {} 
\sphinxAtStartPar
\sphinxstyleliteralstrong{\sphinxupquote{options}} (\sphinxstyleliteralemphasis{\sphinxupquote{Optional}}\sphinxstyleliteralemphasis{\sphinxupquote{{[}}}\sphinxstyleliteralemphasis{\sphinxupquote{Dict}}\sphinxstyleliteralemphasis{\sphinxupquote{{[}}}\sphinxhref{https://docs.python.org/3/library/stdtypes.html\#str}{\sphinxstyleliteralemphasis{\sphinxupquote{str}}}\sphinxstyleliteralemphasis{\sphinxupquote{, }}\sphinxstyleliteralemphasis{\sphinxupquote{Any}}\sphinxstyleliteralemphasis{\sphinxupquote{{]}}}\sphinxstyleliteralemphasis{\sphinxupquote{{]}}}) \textendash{} a dictionary with options for transformations.

\end{itemize}

\sphinxlineitem{Raises}
\sphinxAtStartPar
\sphinxstyleliteralstrong{\sphinxupquote{TransformationError}} \textendash{} if the node is not a Routine.

\end{description}\end{quote}

\end{fulllineitems}


\end{fulllineitems}



\bigskip\hrule\bigskip



\begin{fulllineitems}

\pysigstartsignatures
\pysigline{\sphinxbfcode{\sphinxupquote{class\DUrole{w,w}{  }}}\sphinxcode{\sphinxupquote{psyclone.transformations.}}\sphinxbfcode{\sphinxupquote{KernelImportsToArguments}}}
\pysigstopsignatures
\sphinxAtStartPar
Transformation that removes any accesses of imported data from the supplied
kernel and places them in the caller. The values/references are then passed
by argument into the kernel.


\begin{fulllineitems}

\pysigstartsignatures
\pysiglinewithargsret{\sphinxbfcode{\sphinxupquote{apply}}}{\sphinxparam{\DUrole{n,n}{node}}\sphinxparamcomma \sphinxparam{\DUrole{n,n}{options}\DUrole{o,o}{=}\DUrole{default_value}{None}}}{}
\pysigstopsignatures
\sphinxAtStartPar
Convert the imported variables used inside the kernel into arguments
and modify the InvokeSchedule to pass the same imported variables to
the kernel call.
\begin{quote}\begin{description}
\sphinxlineitem{Parameters}\begin{itemize}
\item {} 
\sphinxAtStartPar
\sphinxstyleliteralstrong{\sphinxupquote{node}} (\sphinxcode{\sphinxupquote{psyclone.psyGen.CodedKern}}) \textendash{} a kernel call.

\item {} 
\sphinxAtStartPar
\sphinxstyleliteralstrong{\sphinxupquote{options}} (\sphinxstyleliteralemphasis{\sphinxupquote{Optional}}\sphinxstyleliteralemphasis{\sphinxupquote{{[}}}\sphinxstyleliteralemphasis{\sphinxupquote{Dict}}\sphinxstyleliteralemphasis{\sphinxupquote{{[}}}\sphinxhref{https://docs.python.org/3/library/stdtypes.html\#str}{\sphinxstyleliteralemphasis{\sphinxupquote{str}}}\sphinxstyleliteralemphasis{\sphinxupquote{, }}\sphinxstyleliteralemphasis{\sphinxupquote{Any}}\sphinxstyleliteralemphasis{\sphinxupquote{{]}}}\sphinxstyleliteralemphasis{\sphinxupquote{{]}}}) \textendash{} a dictionary with options for transformations.

\end{itemize}

\end{description}\end{quote}

\end{fulllineitems}


\end{fulllineitems}


\begin{sphinxadmonition}{note}{Note:}
\sphinxAtStartPar
This transformation is only supported by the GOcean 1.0 API.
\end{sphinxadmonition}


\section{OpenMP}
\label{\detokenize{transformations:openmp}}
\sphinxAtStartPar
OpenMP is added to a code by using transformations. The OpenMP
transformations currently supported allow the addition of:
\begin{itemize}
\item {} 
\sphinxAtStartPar
an \sphinxstylestrong{OpenMP Parallel} directive

\item {} 
\sphinxAtStartPar
an \sphinxstylestrong{OpenMP Target} directive

\item {} 
\sphinxAtStartPar
an \sphinxstylestrong{OpenMP Declare Target} directive

\item {} 
\sphinxAtStartPar
an \sphinxstylestrong{OpenMP Do/For/Loop} directive

\item {} 
\sphinxAtStartPar
an \sphinxstylestrong{OpenMP Single} directive

\item {} 
\sphinxAtStartPar
an \sphinxstylestrong{OpenMP Master} directive

\item {} 
\sphinxAtStartPar
an \sphinxstylestrong{OpenMP Taskloop} directive

\item {} 
\sphinxAtStartPar
multiple \sphinxstylestrong{OpenMP Taskwait} directives; and

\item {} 
\sphinxAtStartPar
an \sphinxstylestrong{OpenMP Parallel Do} directive.

\end{itemize}

\sphinxAtStartPar
The generic versions of these transformations (i.e. ones that
theoretically work for all APIs) were given in the
{\hyperref[\detokenize{transformations:sec-transformations-available}]{\sphinxcrossref{\DUrole{std,std-ref}{Available transformations}}}} section. Examples of their use,
for both CPU and offload to GPU, may be found in the
\sphinxcode{\sphinxupquote{PSyclone/examples/nemo/scripts/omp\_?pu\_trans.py}} transformation scripts.

\sphinxAtStartPar
The API\sphinxhyphen{}specific versions of these transformations are described in the
API\sphinxhyphen{}specific sections of this document. Examples for the LFRic API may
be found in \sphinxcode{\sphinxupquote{PSyclone/examples/lfric/scripts}}.


\subsection{Reductions}
\label{\detokenize{transformations:reductions}}\label{\detokenize{transformations:openmp-reductions}}
\sphinxAtStartPar
PSyclone supports parallel scalar reductions.  If a scalar reduction is
specified in the Kernel metadata (see the API\sphinxhyphen{}specific sections for
details) then PSyclone ensures the appropriate reduction is performed.

\sphinxAtStartPar
In the case of distributed memory, PSyclone will add \sphinxstylestrong{GlobalSum’s}
at the appropriate locations. As can be inferred by the name, only
“summation” reductions are currently supported for distributed memory.

\sphinxAtStartPar
In the case of an OpenMP parallel loop the standard reduction support
will be used by default. For example

\begin{sphinxVerbatim}[commandchars=\\\{\}]
!\PYGZdl{}omp parallel do, reduction(+:x)
!loop
!\PYGZdl{}omp end parallel do
\end{sphinxVerbatim}

\sphinxAtStartPar
OpenMP reductions do not guarantee to give bit reproducible results
for different runs of the same problem even if the same problem is run
using the same resources. The reason for this is that the order in
which data is reduced is not mandated.

\sphinxAtStartPar
Therefore, an additional \sphinxstylestrong{reprod} option has been added to the
\sphinxstylestrong{OpenMP Do} transformation. If the reprod option is set to “True”
then the OpenMP reduction support is replaced with local per\sphinxhyphen{}thread
reductions which are reduced serially after the loop has
finished. This implementation guarantees to give bit\sphinxhyphen{}wise reproducible
results for different runs of the same problem using the same
resources, but will not bit\sphinxhyphen{}wise compare if the code is rerun with
different numbers of OpenMP threads.


\subsection{Restrictions}
\label{\detokenize{transformations:restrictions}}
\sphinxAtStartPar
If two reductions are used within an OpenMP region and the same
variable is used for both reductions then PSyclone will raise an
exception. In this case the solution is to use a different variable
for each reduction.

\sphinxAtStartPar
PSyclone does not support (distributed\sphinxhyphen{}memory) halo swaps or global
sums within OpenMP parallel regions.  Attempting to create a parallel
region for a set of nodes that includes halo swaps or global sums will
produce an error.  In such cases it may be possible to re\sphinxhyphen{}order the
nodes in the Schedule using the {\hyperref[\detokenize{transformations:sec-move-trans}]{\sphinxcrossref{\DUrole{std,std-ref}{MoveTrans}}}}
transformation.


\subsection{OpenMP Tasking}
\label{\detokenize{transformations:openmp-tasking}}
\sphinxAtStartPar
PSyclone supports OpenMP Tasking, through the \sphinxtitleref{OMPTaskloopTrans} and
\sphinxtitleref{OMPTaskwaitTrans} transformations. \sphinxtitleref{OMPTaskloopTrans}
transformations can be applied to loops, whilst the \sphinxtitleref{OMPTaskwaitTrans}
operator is applied to an OpenMP Parallel Region, and computes the dependencies
caused by Taskloops, and adds OpenMP Taskwait statements to satisfy those
dependencies. An example of using OpenMP tasking is available in
\sphinxtitleref{PSyclone/examples/nemo/eg1/openmp\_taskloop\_trans.py}.


\section{OpenCL}
\label{\detokenize{transformations:opencl}}
\sphinxAtStartPar
OpenCL is added to a code by using the \sphinxcode{\sphinxupquote{GOOpenCLTrans}} transformation (see the
{\hyperref[\detokenize{transformations:sec-transformations-available}]{\sphinxcrossref{\DUrole{std,std-ref}{Available transformations}}}} Section above).
Currently this transformation is only supported for the GOcean1.0 API and
is applied to the whole InvokeSchedule of an Invoke.
This transformation will add an OpenCL driver infrastructure to the PSy layer
and generate an OpenCL kernel for each of the Invoke kernels.
This means that all kernels in that Invoke will be executed on the OpenCL
device.
The PSy\sphinxhyphen{}layer OpenCL code generated by PSyclone is still Fortran and makes use
of the FortCL library (\sphinxurl{https://github.com/stfc/FortCL}) to access
OpenCL functionality. It also relies upon the device acceleration support
provided by the dl\_esm\_inf library (\sphinxurl{https://github.com/stfc/dl\_esm\_inf}).

\begin{sphinxadmonition}{note}{Note:}
\sphinxAtStartPar
The generated OpenCL kernels are written in a file called
opencl\_kernels\_\textless{}index\textgreater{}.cl where the index keeps increasing if the
file name already exist.
\end{sphinxadmonition}

\sphinxAtStartPar
The \sphinxcode{\sphinxupquote{GOOpenCLTrans}} transformation accepts an \sphinxtitleref{options} argument with a
map of optional parameters to tune the OpenCL host code in the PSy layer.
These options will be attached to the transformed InvokeSchedule.
The current available options are:


\begin{savenotes}\sphinxattablestart
\sphinxthistablewithglobalstyle
\centering
\begin{tabulary}{\linewidth}[t]{TTT}
\sphinxtoprule
\sphinxstyletheadfamily 
\sphinxAtStartPar
Option
&\sphinxstyletheadfamily 
\sphinxAtStartPar
Description
&\sphinxstyletheadfamily 
\sphinxAtStartPar
Default
\\
\sphinxmidrule
\sphinxtableatstartofbodyhook
\sphinxAtStartPar
end\_barrier
&
\sphinxAtStartPar
Whether a synchronization
barrier should be placed at the end of the
Invoke.
&
\sphinxAtStartPar
True
\\
\sphinxhline
\sphinxAtStartPar
enable\_profiling
&
\sphinxAtStartPar
Enables the profiling of OpenCL Kernels.
&
\sphinxAtStartPar
False
\\
\sphinxhline
\sphinxAtStartPar
out\_of\_order
&
\sphinxAtStartPar
Allows the OpenCL implementation to execute
the enqueued kernels out\sphinxhyphen{}of\sphinxhyphen{}order.
&
\sphinxAtStartPar
False
\\
\sphinxbottomrule
\end{tabulary}
\sphinxtableafterendhook\par
\sphinxattableend\end{savenotes}

\sphinxAtStartPar
Additionally, each individual kernel (inside the Invoke that is going to
be transformed) also accepts a map of options which
are provided by the \sphinxtitleref{set\_opencl\_options()} method of the \sphinxtitleref{Kern} object.
This can affect both the driver layer and/or the OpenCL kernels.
The current available options are:


\begin{savenotes}\sphinxattablestart
\sphinxthistablewithglobalstyle
\centering
\begin{tabulary}{\linewidth}[t]{TTT}
\sphinxtoprule
\sphinxstyletheadfamily 
\sphinxAtStartPar
Option
&\sphinxstyletheadfamily 
\sphinxAtStartPar
Description
&\sphinxstyletheadfamily 
\sphinxAtStartPar
Default
\\
\sphinxmidrule
\sphinxtableatstartofbodyhook
\sphinxAtStartPar
local\_size
&
\sphinxAtStartPar
Number of work\sphinxhyphen{}items to group together
in a work\sphinxhyphen{}group execution (kernel instances
executed at the same time).
&
\sphinxAtStartPar
64
\\
\sphinxhline
\sphinxAtStartPar
queue\_number
&
\sphinxAtStartPar
The identifier of the OpenCL command\_queue
to which the kernel should be submitted. If
the kernel has a dependency on another
kernel submitted to a different
command\_queue a barrier will be added to
guarantee the execution order.
&
\sphinxAtStartPar
1
\\
\sphinxbottomrule
\end{tabulary}
\sphinxtableafterendhook\par
\sphinxattableend\end{savenotes}

\sphinxAtStartPar
Below is an example of a PSyclone script that uses a \sphinxcode{\sphinxupquote{GOOpenCLTrans}} with
multiple InvokeSchedule and kernel\sphinxhyphen{}specific optimization options.

\begin{sphinxVerbatim}[commandchars=\\\{\},numbers=left,firstnumber=1,stepnumber=1]
\PYG{k}{def}\PYG{+w}{ }\PYG{n+nf}{trans}\PYG{p}{(}\PYG{n}{psyir}\PYG{p}{)}\PYG{p}{:}
\PYG{+w}{    }\PYG{l+s+sd}{\PYGZsq{}\PYGZsq{}\PYGZsq{}}
\PYG{l+s+sd}{    Applies OpenCL to the given PSy\PYGZhy{}layer.}

\PYG{l+s+sd}{    :param psyir: the PSyIR of the PSy\PYGZhy{}layer.}
\PYG{l+s+sd}{    :type psyir: :py:class:`psyclone.psyir.nodes.FileContainer`}

\PYG{l+s+sd}{    \PYGZsq{}\PYGZsq{}\PYGZsq{}}
    \PYG{n}{ocl\PYGZus{}trans} \PYG{o}{=} \PYG{n}{GOOpenCLTrans}\PYG{p}{(}\PYG{p}{)}
    \PYG{n}{fold\PYGZus{}trans} \PYG{o}{=} \PYG{n}{FoldConditionalReturnExpressionsTrans}\PYG{p}{(}\PYG{p}{)}
    \PYG{n}{move\PYGZus{}boundaries\PYGZus{}trans} \PYG{o}{=} \PYG{n}{GOMoveIterationBoundariesInsideKernelTrans}\PYG{p}{(}\PYG{p}{)}

    \PYG{c+c1}{\PYGZsh{} Provide kernel\PYGZhy{}specific OpenCL optimization options}
    \PYG{k}{for} \PYG{n}{idx}\PYG{p}{,} \PYG{n}{kern} \PYG{o+ow}{in} \PYG{n+nb}{enumerate}\PYG{p}{(}\PYG{n}{psyir}\PYG{o}{.}\PYG{n}{kernels}\PYG{p}{(}\PYG{p}{)}\PYG{p}{)}\PYG{p}{:}
        \PYG{c+c1}{\PYGZsh{} Move the PSy\PYGZhy{}layer loop boundaries inside the kernel as a kernel}
        \PYG{c+c1}{\PYGZsh{} mask, this allows to iterate through the whole domain}
        \PYG{n}{move\PYGZus{}boundaries\PYGZus{}trans}\PYG{o}{.}\PYG{n}{apply}\PYG{p}{(}\PYG{n}{kern}\PYG{p}{)}
        \PYG{c+c1}{\PYGZsh{} Change the syntax to remove the return statements introduced by the}
        \PYG{c+c1}{\PYGZsh{} previous transformation}
        \PYG{n}{fold\PYGZus{}trans}\PYG{o}{.}\PYG{n}{apply}\PYG{p}{(}\PYG{n}{kern}\PYG{o}{.}\PYG{n}{get\PYGZus{}kernel\PYGZus{}schedule}\PYG{p}{(}\PYG{p}{)}\PYG{p}{)}
        \PYG{c+c1}{\PYGZsh{} Specify the OpenCL queue and workgroup size of the kernel}
        \PYG{c+c1}{\PYGZsh{} In this case we dispatch each kernel in a different queue to check}
        \PYG{c+c1}{\PYGZsh{} that the output code has the necessary barriers to guarantee the}
        \PYG{c+c1}{\PYGZsh{} kernel execution order.}
        \PYG{n}{kern}\PYG{o}{.}\PYG{n}{set\PYGZus{}opencl\PYGZus{}options}\PYG{p}{(}\PYG{p}{\PYGZob{}}\PYG{l+s+s2}{\PYGZdq{}}\PYG{l+s+s2}{queue\PYGZus{}number}\PYG{l+s+s2}{\PYGZdq{}}\PYG{p}{:} \PYG{n}{idx}\PYG{o}{+}\PYG{l+m+mi}{1}\PYG{p}{,} \PYG{l+s+s1}{\PYGZsq{}}\PYG{l+s+s1}{local\PYGZus{}size}\PYG{l+s+s1}{\PYGZsq{}}\PYG{p}{:} \PYG{l+m+mi}{4}\PYG{p}{\PYGZcb{}}\PYG{p}{)}

    \PYG{c+c1}{\PYGZsh{} Transform the Schedule}
    \PYG{k}{for} \PYG{n}{schedule} \PYG{o+ow}{in} \PYG{n}{psyir}\PYG{o}{.}\PYG{n}{walk}\PYG{p}{(}\PYG{n}{InvokeSchedule}\PYG{p}{)}\PYG{p}{:}
        \PYG{n}{ocl\PYGZus{}trans}\PYG{o}{.}\PYG{n}{apply}\PYG{p}{(}\PYG{n}{schedule}\PYG{p}{,} \PYG{n}{options}\PYG{o}{=}\PYG{p}{\PYGZob{}}\PYG{l+s+s2}{\PYGZdq{}}\PYG{l+s+s2}{end\PYGZus{}barrier}\PYG{l+s+s2}{\PYGZdq{}}\PYG{p}{:} \PYG{k+kc}{True}\PYG{p}{\PYGZcb{}}\PYG{p}{)}
\end{sphinxVerbatim}

\sphinxAtStartPar
OpenCL delays the decision of which and where kernels will execute until
run\sphinxhyphen{}time, therefore it is important to use the environment variables provided
by FortCL and DL\_ESM\_INF to inform how things should execute. Specifically:
\begin{itemize}
\item {} 
\sphinxAtStartPar
\sphinxcode{\sphinxupquote{FORTCL\_KERNELS\_FILE}}: Point to the file containing the kernels to execute,
they can be compiled ahead\sphinxhyphen{}of\sphinxhyphen{}time or providing the source for JIT
compilation. To link more than a single kernel, one must merge all the
kernels generated by PSyclone in a single source file.

\item {} 
\sphinxAtStartPar
\sphinxcode{\sphinxupquote{FORTCL\_PLATFORM}}: If the system has more than 1 OpenCL platform. This
environment variable may be used to select which platform on which to execute
the kernels.

\item {} 
\sphinxAtStartPar
\sphinxcode{\sphinxupquote{DL\_ESM\_ALIGNMENT}}: When using OpenCL \textless{}= 1.2 the local\_size should be
exactly divisible by the total size. If this is not the case some
implementations fail silently. A way to solve this issue is to set the
\sphinxtitleref{DL\_ESM\_ALIGNMENT} variable to be equal to the local size.

\end{itemize}

\begin{sphinxadmonition}{note}{Note:}
\sphinxAtStartPar
The OpenCL generation can be combined with distributed memory
generation. In the case where there is more than one accelerator available
on each node, the PSyclone configuration file parameter
\sphinxcode{\sphinxupquote{OCL\_DEVICES\_PER\_NODE}} has to be set to the appropriate value and
the number of MPI\sphinxhyphen{}ranks\sphinxhyphen{}per\sphinxhyphen{}node set by the \sphinxtitleref{mpirun} command has to match
this value accordingly.

\sphinxAtStartPar
For instance if there are 2 accelerators per nodes, \sphinxtitleref{psyclone.cfg} should
have \sphinxcode{\sphinxupquote{OCL\_DEVICES\_PER\_NODE=2}} and the program must be executed with
\sphinxcode{\sphinxupquote{mpirun \sphinxhyphen{}n \textless{}total\_ranks\textgreater{} \sphinxhyphen{}ppn 2 ./application}} (Note: \sphinxtitleref{\sphinxhyphen{}ppn} is an
Intel MPI specific parameter, use equivalent configuration parameters for
other MPI implementations.)
\end{sphinxadmonition}

\sphinxAtStartPar
For example, an execution of a PSyclone generated OpenCL code using all the
mentioned run\sphinxhyphen{}time configuration options could look something like:

\begin{sphinxVerbatim}[commandchars=\\\{\}]
\PYG{n}{FORTCL\PYGZus{}PLATFORM}\PYG{o}{=}\PYG{l+m+mi}{3} \PYG{n}{FORTCL\PYGZus{}KERNELS\PYGZus{}FILE}\PYG{o}{=}\PYG{n}{allkernels}\PYG{o}{.}\PYG{n}{cl} \PYG{n}{DL\PYGZus{}ESM\PYGZus{}ALIGNMENT}\PYG{o}{=}\PYG{l+m+mi}{64} \PYGZbs{}
\PYG{n}{mpirun} \PYG{o}{\PYGZhy{}}\PYG{n}{n} \PYG{l+m+mi}{2} \PYG{o}{.}\PYG{o}{/}\PYG{n}{application}\PYG{o}{.}\PYG{n}{exe}
\end{sphinxVerbatim}


\section{OpenACC}
\label{\detokenize{transformations:openacc}}
\sphinxAtStartPar
PSyclone supports the generation of code targetting GPUs through the
addition of OpenACC directives. This is achieved by a user applying
various OpenACC transformations to the PSyIR before the final Fortran
code is generated. The steps to parallelisation are very similar to
those in OpenMP with the added complexity of managing the movement of
data to and from the GPU device. For the latter task PSyclone provides
the \sphinxcode{\sphinxupquote{ACCDataTrans}} and \sphinxcode{\sphinxupquote{ACCEnterDataTrans}} transformations, as
described in the {\hyperref[\detokenize{transformations:sec-transformations-available}]{\sphinxcrossref{\DUrole{std,std-ref}{Available transformations}}}} Section above.
These two transformations add statically\sphinxhyphen{} and dynamically\sphinxhyphen{}scoped data
regions, respectively. The former manages what data is on the remote
device for a specific section of code while the latter allows run\sphinxhyphen{}time
control of data movement. This second option is essential for
minimising data movement as, without it, PSyclone\sphinxhyphen{}generated code would
move data to and from the device upon every entry/exit of an
Invoke. The first option is mainly provided as an aid to incremental
porting and/or debugging of an OpenACC application as it provides
explicit control over what data is present on a device for a given
(part of an) Invoke routine.

\sphinxAtStartPar
The NVIDIA compiler compiler provides an alternative approach to controlling
data movement through its ‘managed memory’ option
(\sphinxcode{\sphinxupquote{\sphinxhyphen{}gpu=mem:managed}}). When this is enabled the compiler itself takes
on the task of ensuring that data is copied to/from the GPU when
required. (Note that this approach can struggle with Fortran code
containing derived types however.)

\sphinxAtStartPar
As well as ensuring the correct data is copied to and from the remote
device, OpenACC directives must also be added to a code in order to
tell the compiler how it should be parallelised. PSyclone provides the
\sphinxcode{\sphinxupquote{ACCKernelsTrans}}, \sphinxcode{\sphinxupquote{ACCParallelTrans}} and \sphinxcode{\sphinxupquote{ACCLoopTrans}}
transformations for this purpose. The simplest of these is \sphinxcode{\sphinxupquote{ACCKernelsTrans}}
(currently only supported for the generic code transformation
and LFRic API) which encloses the code represented by a sub\sphinxhyphen{}tree of
the PSyIR within an OpenACC \sphinxcode{\sphinxupquote{kernels}} region.  This essentially
gives free\sphinxhyphen{}reign to the compiler to automatically parallelise any
suitable loops within the specified region. An example of the use of
\sphinxcode{\sphinxupquote{ACCDataTrans}} and \sphinxcode{\sphinxupquote{ACCKernelsTrans}} may be found in
PSyclone/examples/nemo/eg3 and an example of \sphinxcode{\sphinxupquote{ACCKernelsTrans}} may
be found in PSyclone/examples/lfric/eg14.

\sphinxAtStartPar
However, as with any “automatic” approach, a more performant solution
can almost always be obtained by providing the compiler with more
explicit direction on how to parallelise the code.  The
\sphinxcode{\sphinxupquote{ACCParallelTrans}} and \sphinxcode{\sphinxupquote{ACCLoopTrans}} transformations allow the
user to define thread\sphinxhyphen{}parallel regions and, within those, define which
loops should be parallelised. For an example of their use please see
PSyclone/examples/gocean/eg2 or PSyclone/examples/lfric/eg14.

\sphinxAtStartPar
In order for a given section of code to be executed on a GPU, any
routines called from within that section must also have been compiled
for the GPU.  This then requires either that any such routines are
in\sphinxhyphen{}lined or that the OpenACC \sphinxcode{\sphinxupquote{routine}} directive be added to any
such routines.  This situation will occur routinely in those PSyclone
APIs that use the PSyKAl separation of concerns since the
user\sphinxhyphen{}supplied kernel routines are called from within
PSyclone\sphinxhyphen{}generated loops in the PSy layer. PSyclone therefore provides
the \sphinxcode{\sphinxupquote{ACCRoutineTrans}} transformation which, given a Kernel node in
the PSyIR, creates a new version of that kernel with the \sphinxcode{\sphinxupquote{routine}}
directive added. See either PSyclone/examples/gocean/eg2 or
PSyclone/examples/lfric/eg14 for an example.


\section{SIR}
\label{\detokenize{transformations:sir}}
\sphinxAtStartPar
It is currently not possible for PSyclone to output SIR code without
using a script. Three examples of such scripts are given in example 4
for the NEMO examples directory. The first \sphinxtitleref{sir\_trans.py} simply outputs SIR.
This will raise an exception if used with the \sphinxtitleref{tracer advection} example as
the example contains array\sphinxhyphen{}index notation which is not supported by
the SIR backend, but will generate code for the other examples. The
second, \sphinxtitleref{sir\_trans\_loop.py} includes transformations to hoist code out
of a loop, translate array\sphinxhyphen{}index notation into explicit loops and
translate a single access to an array dimension to a one\sphinxhyphen{}trip loop (to
make the code suitable for the SIR backend). This works with the
\sphinxtitleref{tracer\sphinxhyphen{}advection} example. The third script \sphinxtitleref{sir\_trans\_all.py}
additionally replaces any intrinsics with equivalent code and can also
be used with the \sphinxtitleref{tracer\sphinxhyphen{}advection} example (and the
\sphinxtitleref{intrinsic\_example.f90} example).

\sphinxstepscope


\chapter{Introduction to PSyKAl}
\label{\detokenize{introduction_to_psykal:introduction-to-psykal}}\label{\detokenize{introduction_to_psykal:id1}}\label{\detokenize{introduction_to_psykal::doc}}
\sphinxAtStartPar
PSyKAl is a kernel\sphinxhyphen{}based software architecture proposed in the \sphinxhref{https://www.metoffice.gov.uk/research/foundation/dynamics/next-generation}{GungHo project}
to design Fortran\sphinxhyphen{}embedded domain\sphinxhyphen{}specific languages that provide a clear
separations of concerns between the science code and the optimisation/parallelisation
details of an application.
The model distinguishes between three layers: the Algorithm layer, the
Kernel layer and the Parallelisation System (PSy) layer; which together give
the model its name.

\sphinxAtStartPar
The Algorithm layer is responsible for providing a high\sphinxhyphen{}level description of the
algorithm that the scientist wants to run. This layer operates on full fields
and includes calls to kernels and built\sphinxhyphen{}ins.

\sphinxAtStartPar
The Kernel layer contains the actual implementation of the kernels as functors.
Each functor implements a method that operates on an individual section of the fields.
Depending on the DSL, these can be a single element, a vertical column, or a set of
vertical columns. The kernels also specify some metadata that allow the data
dependencies between kernels to be determined. Built\sphinxhyphen{}ins are similar to kernels but
are provided by the PSyclone infrastructure itself.

\sphinxAtStartPar
The PSy layer acts as a bridge between the algorithm and kernel layers. It
contains the code that is responsible for performing concurrent execution
of kernels, including the way in which the iteration space is traversed,
while respecting kernel dependencies. This layer can be tuned for
specific platforms such as multi\sphinxhyphen{}node, multi\sphinxhyphen{}core and GPGPUs architectures
without affecting the user\sphinxhyphen{}supplied Algorithm and Kernel layers.

\sphinxAtStartPar
Rather than requiring that the PSy layer be written manually, PSyclone uses
the provided Algorithm and Kernel implementations to generate an inital PSyIR
for the PSy\sphinxhyphen{}layer, optionally with distributed\sphinxhyphen{}memory parallelism.
This then can be programatically optimised by applying PSyclone transformations
(e.g. kernel fusing, colouring, inlining, …) to better fit the target architecture.

\sphinxAtStartPar
The rest of this section describes how to use the \sphinxcode{\sphinxupquote{psyclone}} command to process PSyKAl
DSLs and how to implement each layer, providing examples for each of them.


\section{Usage}
\label{\detokenize{introduction_to_psykal:usage}}\label{\detokenize{introduction_to_psykal:psykal-usage}}
\sphinxAtStartPar
To use PSyclone to process a PSyKAl algorithm file, the \sphinxcode{\sphinxupquote{\sphinxhyphen{}api API\_NAME}} parameter
must be provided. In addition, distributed memory can be switched on or off by
using the \sphinxcode{\sphinxupquote{\sphinxhyphen{}dm}}/\sphinxcode{\sphinxupquote{\sphinxhyphen{}\sphinxhyphen{}dist\_mem}} or \sphinxcode{\sphinxupquote{\sphinxhyphen{}nodm}}/\sphinxcode{\sphinxupquote{\sphinxhyphen{}\sphinxhyphen{}no\_dist\_mem}} flags. For PSyKAl
DSLs, the optional transformation script provided by the \sphinxcode{\sphinxupquote{\sphinxhyphen{}s SCRIPT}}
parameter will be applied to the PSyIR of the PSy\sphinxhyphen{}layer (but the scripts can also
contain instructions to transform the code of the kernels used within it).

\sphinxAtStartPar
For example, the following command will process an LFRic PSyKAl algorithm file,
generate a PSy\sphinxhyphen{}layer containing distributed\sphinxhyphen{}memory parallelism and then transform
it using the \sphinxcode{\sphinxupquote{additional\_optimisations.py}} script:

\begin{sphinxVerbatim}[commandchars=\\\{\}]
psyclone\PYG{+w}{ }\PYGZhy{}api\PYG{+w}{ }lfric\PYG{+w}{ }\PYGZhy{}dm\PYG{+w}{ }\PYGZhy{}s\PYG{+w}{ }additional\PYGZus{}optimisations.py\PYG{+w}{ }algorithm.f90\PYG{+w}{ }\PYG{l+s+se}{\PYGZbs{}}
\PYG{+w}{   }\PYGZhy{}oalg\PYG{+w}{ }algorithm\PYGZus{}output.f90\PYG{+w}{ }\PYGZhy{}opsy\PYG{+w}{ }psy\PYGZus{}layer\PYGZus{}output.f90
\end{sphinxVerbatim}

\sphinxAtStartPar
To date, there are two PSyKAl DSL implementations: the
{\hyperref[\detokenize{dynamo0p3:lfric-api}]{\sphinxcrossref{\DUrole{std,std-ref}{LFRic PSyKAl API}}}}, a mixed
finite\sphinxhyphen{}element DSL used to implement the next\sphinxhyphen{}generation UK Met Office
atmospheric model dynamical core; and the
{\hyperref[\detokenize{gocean1p0:gocean-api}]{\sphinxcrossref{\DUrole{std,std-ref}{GOcean PSyKAl API}}}}, a finite difference ocean model
benchmark. For more details on these see the corresponding sections of
this User Guide.


\section{Algorithm layer}
\label{\detokenize{introduction_to_psykal:algorithm-layer}}\label{\detokenize{introduction_to_psykal:id2}}
\sphinxAtStartPar
As mentioned in the Introduction, the Algorithm layer provides a high\sphinxhyphen{}level
description of the algorithm that the scientist would like to run, in terms
of invocations to Kernel and Built\sphinxhyphen{}in operations. It operates on full fields
and therefore it is not allowed to call the individual kernel executors nor
include any parallelisation calls or directives. Instead, the algorithm layer
uses the \sphinxcode{\sphinxupquote{invoke}} subroutine, which takes as arguments one or more kernel
functor constructors (as a Fortran type constructor) and, optionally, a name
argument. The kernel functors, in turn, take as argument the full
fields they operate on and any other quantities specified in their metadata.

\sphinxAtStartPar
For example:

\begin{sphinxVerbatim}[commandchars=\\\{\}]
\PYG{p}{.}\PYG{p}{.}\PYG{p}{.}
\PYG{k}{call }\PYG{n}{invoke}\PYG{p}{(}\PYG{n}{kernel1}\PYG{p}{(}\PYG{n}{arg1}\PYG{p}{,}\PYG{n}{arg2}\PYG{p}{)}\PYG{p}{,}\PYG{+w}{ }\PYG{n}{kernel2}\PYG{p}{(}\PYG{n}{arg1}\PYG{p}{,}\PYG{+w}{ }\PYG{l+m+mi}{3}\PYG{p}{)}\PYG{p}{,}\PYG{+w}{ }\PYG{n}{name}\PYG{o}{=}\PYG{l+s+s2}{\PYGZdq{}Example\PYGZus{}Invoke\PYGZdq{}}\PYG{p}{)}
\PYG{p}{.}\PYG{p}{.}\PYG{p}{.}
\end{sphinxVerbatim}

\sphinxAtStartPar
A complete application can consist of many algorithm files, each of them
containing as many \sphinxcode{\sphinxupquote{invoke()}} calls as required. PSyclone is applied to
each individual algorithm layer file and must therefore be run multiple
times if multiple algorithm files exist in a project.

\sphinxAtStartPar
The algorithm developer is also able to reference more than one
Kernel/Built\sphinxhyphen{}in within an invoke call (as indicated in the previous example).
In fact this feature is encouraged for performance reasons.
\sphinxstylestrong{As a general guideline the developer should aim to use as few invokes
as possible, each with as many Kernel functors in them
as is possible}. The reason for this is that it allows for greater
freedom for optimisations in the PSy\sphinxhyphen{}layer as these
are limited to the contents of individual invokes \sphinxhyphen{} PSyclone
currently does not attempt to optimise the PSy layer over multiple
invoke calls.

\sphinxAtStartPar
As well as generating the PSy\sphinxhyphen{}layer code, PSyclone modifies the
Algorithm layer code, replacing \sphinxcode{\sphinxupquote{invoke}} calls with calls to the
generated PSy\sphinxhyphen{}layer subroutine(s) (plus the associated \sphinxcode{\sphinxupquote{use}} statements)
so that the algorithm code is compilable and linkable. For example,
the invoke above is translated into something like the following:

\begin{sphinxVerbatim}[commandchars=\\\{\}]
\PYG{p}{.}\PYG{p}{.}\PYG{p}{.}
\PYG{k}{use }\PYG{n}{psy}\PYG{p}{,}\PYG{+w}{ }\PYG{k}{only}\PYG{+w}{ }\PYG{p}{:}\PYG{+w}{ }\PYG{n}{invoke\PYGZus{}example\PYGZus{}invoke}
\PYG{p}{.}\PYG{p}{.}\PYG{p}{.}
\PYG{k}{call }\PYG{n}{invoke\PYGZus{}example\PYGZus{}invoke}\PYG{p}{(}\PYG{n}{arg1}\PYG{p}{,}\PYG{+w}{ }\PYG{n}{arg2}\PYG{p}{,}\PYG{+w}{ }\PYG{l+m+mi}{3}\PYG{p}{)}
\PYG{p}{.}\PYG{p}{.}\PYG{p}{.}
\end{sphinxVerbatim}

\sphinxAtStartPar
The \sphinxcode{\sphinxupquote{name}} argument in the invoke call is optional. If supplied it
must be a string literal. Labels are not case\sphinxhyphen{}sensitive and must be
valid Fortran names (e.g. \sphinxcode{\sphinxupquote{name="compute(1)"}} is invalid).
The label is used to name the corresponding PSy\sphinxhyphen{}layer routine generated
by PSyclone in order to make debugging and profiling outputs more readable.
So, for the above example, the generated PSy\sphinxhyphen{}layer subroutine
will be named “invoke\_example\_invoke”, otherwise a numeric index is given
(e.g. “invoke\_0”).
Each invoke label must be unique within an Algorithm source file.


\subsection{Limitations}
\label{\detokenize{introduction_to_psykal:limitations}}
\sphinxAtStartPar
There are limitations in the Fortran expressions that can be used inside
the invoke call kernel arguments. The current list of known restrictions
on the form of kernel arguments within an invoke is:
\begin{itemize}
\item {} 
\sphinxAtStartPar
No arithmetic expressions (e.g. \sphinxcode{\sphinxupquote{kernel\_type(a+b)}} or \sphinxcode{\sphinxupquote{kernel\_type(\sphinxhyphen{}a)}})

\item {} 
\sphinxAtStartPar
No named (optional) arguments (e.g. \sphinxcode{\sphinxupquote{kernel\_type(fn(my\_arg=a))}})

\end{itemize}

\sphinxAtStartPar
If you encounter any other limitations (or have a burning desire to use one
of the above forms) then please contact the PSyclone developers.


\section{Kernel layer}
\label{\detokenize{introduction_to_psykal:kernel-layer}}\label{\detokenize{introduction_to_psykal:id3}}
\sphinxAtStartPar
In the PSyKAl model, the Kernel code operates on an individual element of
a field (such as a column of cells). The reason for doing this is that it
gives the PSy layer flexibility in choosing the iteration order and exploiting
the spatial domain parallelisation. The Kernel layer is not allowed to include
any calls or directives related to parallelisation and works on
raw Fortran arrays (to allow the compiler to optimise the code).
Since a Kernel is called over the spatial domain (by the PSy layer) it
must take at least one field or operator as an argument.

\sphinxAtStartPar
Kernels are implemented as Fortran Functors. Functors are objects that can be
treated as if they are functions. As such they have two main interfaces for
calling and providing arguments to them: the object constructor (used by the
Algorithm layer) and a method that executes the code of the functor (used by
the PSy\sphinxhyphen{}layer).

\sphinxAtStartPar
PSyKal applications accept one or more modules providing kernels, each of which
can contain one or more kernel functors. Each kernel functor provides a set of
meta\sphinxhyphen{}data attributes and a method with its implementation.

\sphinxAtStartPar
In the example below the module \sphinxcode{\sphinxupquote{w3\_solver\_kernel\_mod}} contains one kernel
named \sphinxcode{\sphinxupquote{w3\_solver\_kernel\_type}} and its individual element execution method in
subroutine \sphinxcode{\sphinxupquote{w3\_solver\_code}}.

\sphinxAtStartPar
The metadata is API\sphinxhyphen{}specific and describes the kernel iteration space and
dependencies, so that PSyclone can generate correct PSy\sphinxhyphen{}layer code (including
any necessary halo\sphinxhyphen{}exchanges if DM is enabled). The
metadata is provided by the kernel developer, who must guarantee its
correctness. The example below shows meta\sphinxhyphen{}data for the \sphinxcode{\sphinxupquote{LFRic}} API:

\begin{sphinxVerbatim}[commandchars=\\\{\}]
\PYG{k}{module }\PYG{n}{w3\PYGZus{}solver\PYGZus{}kernel\PYGZus{}mod}

\PYG{p}{.}\PYG{p}{.}\PYG{p}{.}

\PYG{+w}{  }\PYG{k}{type}\PYG{p}{,}\PYG{+w}{ }\PYG{k}{public}\PYG{p}{,}\PYG{+w}{ }\PYG{k}{extends}\PYG{p}{(}\PYG{n}{kernel\PYGZus{}type}\PYG{p}{)}\PYG{+w}{ }\PYG{k+kd}{::}\PYG{+w}{ }\PYG{n}{w3\PYGZus{}solver\PYGZus{}kernel\PYGZus{}type}
\PYG{+w}{    }\PYG{k}{private}
\PYG{k}{    }\PYG{k}{type}\PYG{p}{(}\PYG{n}{arg\PYGZus{}type}\PYG{p}{)}\PYG{+w}{ }\PYG{k+kd}{::}\PYG{+w}{ }\PYG{n}{meta\PYGZus{}args}\PYG{p}{(}\PYG{l+m+mi}{4}\PYG{p}{)}\PYG{+w}{ }\PYG{o}{=}\PYG{+w}{ }\PYG{p}{(}\PYG{o}{/}\PYG{+w}{                 }\PYG{p}{\PYGZam{}}
\PYG{+w}{         }\PYG{n}{arg\PYGZus{}type}\PYG{p}{(}\PYG{n}{GH\PYGZus{}FIELD}\PYG{p}{,}\PYG{+w}{   }\PYG{n}{GH\PYGZus{}REAL}\PYG{p}{,}\PYG{+w}{ }\PYG{n}{GH\PYGZus{}WRITE}\PYG{p}{,}\PYG{+w}{ }\PYG{n}{W3}\PYG{p}{)}\PYG{p}{,}\PYG{+w}{   }\PYG{p}{\PYGZam{}}
\PYG{+w}{         }\PYG{n}{arg\PYGZus{}type}\PYG{p}{(}\PYG{n}{GH\PYGZus{}FIELD}\PYG{p}{,}\PYG{+w}{   }\PYG{n}{GH\PYGZus{}REAL}\PYG{p}{,}\PYG{+w}{ }\PYG{n}{GH\PYGZus{}READ}\PYG{p}{,}\PYG{+w}{  }\PYG{n}{W3}\PYG{p}{)}\PYG{p}{,}\PYG{+w}{   }\PYG{p}{\PYGZam{}}
\PYG{+w}{         }\PYG{n}{arg\PYGZus{}type}\PYG{p}{(}\PYG{n}{GH\PYGZus{}FIELD}\PYG{o}{*}\PYG{l+m+mi}{3}\PYG{p}{,}\PYG{+w}{ }\PYG{n}{GH\PYGZus{}REAL}\PYG{p}{,}\PYG{+w}{ }\PYG{n}{GH\PYGZus{}READ}\PYG{p}{,}\PYG{+w}{  }\PYG{n}{Wchi}\PYG{p}{)}\PYG{p}{,}\PYG{+w}{ }\PYG{p}{\PYGZam{}}
\PYG{+w}{         }\PYG{n}{arg\PYGZus{}type}\PYG{p}{(}\PYG{n}{GH\PYGZus{}SCALAR}\PYG{p}{,}\PYG{+w}{  }\PYG{n}{GH\PYGZus{}REAL}\PYG{p}{,}\PYG{+w}{ }\PYG{n}{GH\PYGZus{}READ}\PYG{p}{)}\PYG{+w}{         }\PYG{p}{\PYGZam{}}
\PYG{+w}{         }\PYG{o}{/}\PYG{p}{)}
\PYG{+w}{    }\PYG{k}{type}\PYG{p}{(}\PYG{n}{func\PYGZus{}type}\PYG{p}{)}\PYG{+w}{ }\PYG{k+kd}{::}\PYG{+w}{ }\PYG{n}{meta\PYGZus{}funcs}\PYG{p}{(}\PYG{l+m+mi}{2}\PYG{p}{)}\PYG{+w}{ }\PYG{o}{=}\PYG{+w}{ }\PYG{p}{(}\PYG{o}{/}\PYG{+w}{               }\PYG{p}{\PYGZam{}}
\PYG{+w}{         }\PYG{n}{func\PYGZus{}type}\PYG{p}{(}\PYG{n}{W3}\PYG{p}{,}\PYG{+w}{   }\PYG{n}{GH\PYGZus{}BASIS}\PYG{p}{)}\PYG{p}{,}\PYG{+w}{                     }\PYG{p}{\PYGZam{}}
\PYG{+w}{         }\PYG{n}{func\PYGZus{}type}\PYG{p}{(}\PYG{n}{Wchi}\PYG{p}{,}\PYG{+w}{ }\PYG{n}{GH\PYGZus{}DIFF\PYGZus{}BASIS}\PYG{p}{)}\PYG{+w}{                 }\PYG{p}{\PYGZam{}}
\PYG{+w}{         }\PYG{o}{/}\PYG{p}{)}
\PYG{+w}{    }\PYG{k+kt}{integer}\PYG{+w}{ }\PYG{k+kd}{::}\PYG{+w}{ }\PYG{n}{gh\PYGZus{}shape}\PYG{+w}{ }\PYG{o}{=}\PYG{+w}{ }\PYG{n}{GH\PYGZus{}QUADRATURE\PYGZus{}XYoZ}
\PYG{+w}{    }\PYG{k+kt}{integer}\PYG{+w}{ }\PYG{k+kd}{::}\PYG{+w}{ }\PYG{n}{operates\PYGZus{}on}\PYG{+w}{ }\PYG{o}{=}\PYG{+w}{ }\PYG{n}{CELL\PYGZus{}COLUMN}
\PYG{+w}{  }\PYG{k}{contains}
\PYG{k}{    }\PYG{k}{procedure}\PYG{p}{,}\PYG{+w}{ }\PYG{k}{nopass}\PYG{+w}{ }\PYG{k+kd}{::}\PYG{+w}{ }\PYG{n}{solver\PYGZus{}w3\PYGZus{}code}
\PYG{+w}{  }\PYG{k}{end }\PYG{k}{type}

\PYG{k}{contains}

\PYG{k}{  }\PYG{k}{subroutine }\PYG{n}{solver\PYGZus{}w3\PYGZus{}code}\PYG{p}{(}\PYG{n}{nlayers}\PYG{p}{,}\PYG{+w}{                                         }\PYG{p}{\PYGZam{}}
\PYG{+w}{                            }\PYG{n}{x}\PYG{p}{,}\PYG{+w}{ }\PYG{n}{rhs}\PYG{p}{,}\PYG{+w}{                                          }\PYG{p}{\PYGZam{}}
\PYG{+w}{                            }\PYG{n}{chi\PYGZus{}1}\PYG{p}{,}\PYG{+w}{ }\PYG{n}{chi\PYGZus{}2}\PYG{p}{,}\PYG{+w}{ }\PYG{n}{chi\PYGZus{}3}\PYG{p}{,}\PYG{+w}{ }\PYG{n}{ascalar}\PYG{p}{,}\PYG{+w}{                    }\PYG{p}{\PYGZam{}}
\PYG{+w}{                            }\PYG{n}{ndf\PYGZus{}w3}\PYG{p}{,}\PYG{+w}{ }\PYG{n}{undf\PYGZus{}w3}\PYG{p}{,}\PYG{+w}{ }\PYG{n}{map\PYGZus{}w3}\PYG{p}{,}\PYG{+w}{ }\PYG{n}{w3\PYGZus{}basis}\PYG{p}{,}\PYG{+w}{               }\PYG{p}{\PYGZam{}}
\PYG{+w}{                            }\PYG{n}{ndf\PYGZus{}wchi}\PYG{p}{,}\PYG{+w}{ }\PYG{n}{undf\PYGZus{}wchi}\PYG{p}{,}\PYG{+w}{ }\PYG{n}{map\PYGZus{}wchi}\PYG{p}{,}\PYG{+w}{ }\PYG{n}{wchi\PYGZus{}diff\PYGZus{}basis}\PYG{p}{,}\PYG{+w}{  }\PYG{p}{\PYGZam{}}
\PYG{+w}{                            }\PYG{n}{nqp\PYGZus{}h}\PYG{p}{,}\PYG{+w}{ }\PYG{n}{nqp\PYGZus{}v}\PYG{p}{,}\PYG{+w}{ }\PYG{n}{wqp\PYGZus{}h}\PYG{p}{,}\PYG{+w}{ }\PYG{n}{wqp\PYGZus{}v}\PYG{p}{)}
\PYG{+w}{    }\PYG{p}{.}\PYG{p}{.}\PYG{p}{.}
\PYG{+w}{  }\PYG{k}{end }\PYG{k}{subroutine }\PYG{n}{solver\PYGZus{}w3\PYGZus{}code}

\PYG{k}{end }\PYG{k}{module }\PYG{n}{w3\PYGZus{}solver\PYGZus{}kernel\PYGZus{}mod}
\end{sphinxVerbatim}

\sphinxAtStartPar
Note that the executor method can also be declared as a module procedure
interface to provide alternative implementations (e.g. different precisions)
of the kernel code. These are selected as appropriate by the Fortran compiler,
depending on the precision of the fields being passed to them:

\begin{sphinxVerbatim}[commandchars=\\\{\}]
\PYG{k}{type}\PYG{p}{,}\PYG{+w}{ }\PYG{k}{extends}\PYG{p}{(}\PYG{n}{kernel\PYGZus{}type}\PYG{p}{)}\PYG{+w}{ }\PYG{k+kd}{::}\PYG{+w}{ }\PYG{n}{kernel1}
\PYG{+w}{  }\PYG{p}{.}\PYG{p}{.}\PYG{p}{.}
\PYG{+w}{  }\PYG{k}{type}\PYG{p}{(}\PYG{p}{.}\PYG{p}{.}\PYG{p}{.}\PYG{p}{)}\PYG{+w}{ }\PYG{k+kd}{::}\PYG{+w}{ }\PYG{n}{meta\PYGZus{}args}\PYG{p}{(}\PYG{p}{.}\PYG{p}{.}\PYG{p}{.}\PYG{p}{)}\PYG{+w}{ }\PYG{o}{=}\PYG{+w}{ }\PYG{p}{(}\PYG{o}{/}\PYG{+w}{ }\PYG{p}{.}\PYG{p}{.}\PYG{p}{.}\PYG{+w}{ }\PYG{o}{/}\PYG{p}{)}
\PYG{+w}{  }\PYG{p}{.}\PYG{p}{.}\PYG{p}{.}
\PYG{+w}{  }\PYG{k+kt}{integer}\PYG{+w}{ }\PYG{k+kd}{::}\PYG{+w}{ }\PYG{n}{operates\PYGZus{}on}\PYG{+w}{ }\PYG{o}{=}\PYG{+w}{ }\PYG{p}{.}\PYG{p}{.}\PYG{p}{.}
\PYG{+w}{  }\PYG{p}{.}\PYG{p}{.}\PYG{p}{.}
\PYG{k}{end }\PYG{k}{type }\PYG{n}{kernel1}

\PYG{k}{interface}\PYG{+w}{ }\PYG{p}{.}\PYG{p}{.}\PYG{p}{.}
\PYG{+w}{  }\PYG{k}{module }\PYG{k}{procedure}\PYG{+w}{ }\PYG{p}{.}\PYG{p}{.}\PYG{p}{.}
\PYG{k}{end }\PYG{k}{interface}
\end{sphinxVerbatim}


\section{Built\sphinxhyphen{}ins}
\label{\detokenize{introduction_to_psykal:built-ins}}\label{\detokenize{introduction_to_psykal:id4}}
\sphinxAtStartPar
Built\sphinxhyphen{}ins are operations which can be specified within an invoke call in
the algorithm layer but do not require an associated kernel implementation
because they are provided by the infrastructure.

\sphinxAtStartPar
These are useful for commonly\sphinxhyphen{}used operations, as they reduce the amount of
code that the PSyKAl project has to maintain. In addition, they offer
potential performance advantages as their implementation can completely
change for different architectures and the PSy layer is free to implement
these operations in whatever way it chooses.

\sphinxAtStartPar
The list of supported Built\sphinxhyphen{}ins is API\sphinxhyphen{}specific and therefore these are
described under the documentation of each API. In general, PSyclone will
need to know the types of the arguments being passed to any Built\sphinxhyphen{}ins.
Each API provides a Fortran file that contains the metadata for all
Built\sphinxhyphen{}in operations supported for that API.

\begin{sphinxadmonition}{note}{Note:}
\sphinxAtStartPar
When a particular Built\sphinxhyphen{}in is used, the name of this
Built\sphinxhyphen{}in should not be used for anything else within the
same scope. For example, it is not valid to make use of a
Built\sphinxhyphen{}in called \sphinxcode{\sphinxupquote{setval\_c}} and for its parent subroutine
to also be called \sphinxcode{\sphinxupquote{setval\_c}}. In this case PSyclone will
raise an exception.
\end{sphinxadmonition}


\subsection{Example}
\label{\detokenize{introduction_to_psykal:example}}
\sphinxAtStartPar
In the following algorithm\sphinxhyphen{}layer example from the LFRic API, the invoke call
includes a call to two Built\sphinxhyphen{}ins (\sphinxcode{\sphinxupquote{setval\_c}} and \sphinxcode{\sphinxupquote{X\_divideby\_Y}}) and a
user\sphinxhyphen{}supplied kernel that operates on cell columns
(\sphinxcode{\sphinxupquote{matrix\_vector\_kernel\_mm\_type}}).
The \sphinxcode{\sphinxupquote{setval\_c}} Built\sphinxhyphen{}in sets all values in the field \sphinxcode{\sphinxupquote{Ax}} to \sphinxcode{\sphinxupquote{1.0}} and
the \sphinxcode{\sphinxupquote{X\_divideby\_Y}} Built\sphinxhyphen{}in divides values in the field \sphinxcode{\sphinxupquote{rhs}} by their
equivalent (per degree of freedom) values in the field \sphinxcode{\sphinxupquote{lumped\_weight}}
(see {\hyperref[\detokenize{dynamo0p3:lfric-built-ins}]{\sphinxcrossref{\DUrole{std,std-ref}{supported LFRic API Built\sphinxhyphen{}ins}}}}).
Notice that, unlike the kernel call, no \sphinxcode{\sphinxupquote{use}} association is required for
the Built\sphinxhyphen{}ins since they are provided as part of the environment.

\begin{sphinxVerbatim}[commandchars=\\\{\}]
\PYG{k}{module }\PYG{n}{solver\PYGZus{}mod}
\PYG{+w}{  }\PYG{p}{.}\PYG{p}{.}\PYG{p}{.}
\PYG{+w}{  }\PYG{k}{use }\PYG{n}{matrix\PYGZus{}vector\PYGZus{}mm\PYGZus{}kernel\PYGZus{}mod}\PYG{p}{,}\PYG{+w}{ }\PYG{k}{only}\PYG{p}{:}\PYG{+w}{ }\PYG{n}{matrix\PYGZus{}vector\PYGZus{}kernel\PYGZus{}mm\PYGZus{}type}
\PYG{+w}{  }\PYG{p}{.}\PYG{p}{.}\PYG{p}{.}

\PYG{+w}{  }\PYG{k}{subroutine }\PYG{n}{jacobi\PYGZus{}solver\PYGZus{}algorithm}\PYG{p}{(}\PYG{n}{lhs}\PYG{p}{,}\PYG{+w}{ }\PYG{n}{rhs}\PYG{p}{,}\PYG{+w}{ }\PYG{n}{mm}\PYG{p}{,}\PYG{+w}{ }\PYG{n}{mesh}\PYG{p}{,}\PYG{+w}{ }\PYG{n}{n\PYGZus{}iter}\PYG{p}{)}

\PYG{+w}{    }\PYG{k+kt}{integer}\PYG{p}{(}\PYG{n+nb}{kind}\PYG{o}{=}\PYG{n}{i\PYGZus{}def}\PYG{p}{)}\PYG{p}{,}\PYG{+w}{ }\PYG{k}{intent}\PYG{p}{(}\PYG{n}{in}\PYG{p}{)}\PYG{+w}{    }\PYG{k+kd}{::}\PYG{+w}{ }\PYG{n}{n\PYGZus{}iter}
\PYG{+w}{    }\PYG{k}{type}\PYG{p}{(}\PYG{n}{field\PYGZus{}type}\PYG{p}{)}\PYG{p}{,}\PYG{+w}{    }\PYG{k}{intent}\PYG{p}{(}\PYG{n}{inout}\PYG{p}{)}\PYG{+w}{ }\PYG{k+kd}{::}\PYG{+w}{ }\PYG{n}{lhs}
\PYG{+w}{    }\PYG{k}{type}\PYG{p}{(}\PYG{n}{field\PYGZus{}type}\PYG{p}{)}\PYG{p}{,}\PYG{+w}{    }\PYG{k}{intent}\PYG{p}{(}\PYG{n}{in}\PYG{p}{)}\PYG{+w}{    }\PYG{k+kd}{::}\PYG{+w}{ }\PYG{n}{rhs}
\PYG{+w}{    }\PYG{k}{type}\PYG{p}{(}\PYG{n}{operator\PYGZus{}type}\PYG{p}{)}\PYG{p}{,}\PYG{+w}{ }\PYG{k}{intent}\PYG{p}{(}\PYG{n}{in}\PYG{p}{)}\PYG{+w}{    }\PYG{k+kd}{::}\PYG{+w}{ }\PYG{n}{mm}
\PYG{+w}{    }\PYG{k}{type}\PYG{p}{(}\PYG{n}{mesh\PYGZus{}type}\PYG{p}{)}\PYG{p}{,}\PYG{+w}{     }\PYG{k}{intent}\PYG{p}{(}\PYG{n}{in}\PYG{p}{)}\PYG{+w}{    }\PYG{k+kd}{::}\PYG{+w}{ }\PYG{n}{mesh}
\PYG{+w}{    }\PYG{k}{type}\PYG{p}{(}\PYG{n}{field\PYGZus{}type}\PYG{p}{)}\PYG{+w}{                   }\PYG{k+kd}{::}\PYG{+w}{ }\PYG{n}{Ax}\PYG{p}{,}\PYG{+w}{ }\PYG{n}{lumped\PYGZus{}weight}

\PYG{+w}{    }\PYG{p}{.}\PYG{p}{.}\PYG{p}{.}

\PYG{+w}{    }\PYG{c}{! Compute mass lump}
\PYG{+w}{    }\PYG{k}{call }\PYG{n}{invoke}\PYG{p}{(}\PYG{+w}{ }\PYG{n}{name}\PYG{+w}{ }\PYG{o}{=}\PYG{+w}{ }\PYG{l+s+s2}{\PYGZdq{}Jacobi\PYGZus{}mass\PYGZus{}lump\PYGZdq{}}\PYG{p}{,}\PYG{+w}{                           }\PYG{p}{\PYGZam{}}
\PYG{+w}{                 }\PYG{n}{setval\PYGZus{}c}\PYG{p}{(}\PYG{n}{Ax}\PYG{p}{,}\PYG{+w}{ }\PYG{l+m+mf}{1.0\PYGZus{}r\PYGZus{}def}\PYG{p}{)}\PYG{p}{,}\PYG{+w}{                             }\PYG{p}{\PYGZam{}}
\PYG{+w}{                 }\PYG{n}{matrix\PYGZus{}vector\PYGZus{}kernel\PYGZus{}mm\PYGZus{}type}\PYG{p}{(}\PYG{n}{lumped\PYGZus{}weight}\PYG{p}{,}\PYG{+w}{ }\PYG{n}{Ax}\PYG{p}{,}\PYG{+w}{ }\PYG{n}{mm}\PYG{p}{)}\PYG{p}{,}\PYG{+w}{ }\PYG{p}{\PYGZam{}}
\PYG{+w}{                 }\PYG{n}{X\PYGZus{}divideby\PYGZus{}Y}\PYG{p}{(}\PYG{n}{lhs}\PYG{p}{,}\PYG{+w}{ }\PYG{n}{rhs}\PYG{p}{,}\PYG{+w}{ }\PYG{n}{lumped\PYGZus{}weight}\PYG{p}{)}\PYG{+w}{ }\PYG{p}{)}

\PYG{+w}{  }\PYG{k}{end }\PYG{k}{subroutine }\PYG{n}{jacobi\PYGZus{}solver\PYGZus{}algorithm}
\PYG{+w}{  }\PYG{p}{.}\PYG{p}{.}\PYG{p}{.}

\PYG{k}{end }\PYG{k}{module }\PYG{n}{solver\PYGZus{}mod}
\end{sphinxVerbatim}


\section{PSy layer}
\label{\detokenize{introduction_to_psykal:psy-layer}}\label{\detokenize{introduction_to_psykal:id5}}
\sphinxAtStartPar
In the PSyKAl model, the PSy layer is the bridge between the Algorithm
full\sphinxhyphen{}field operations and the Kernel/Built\sphinxhyphen{}Ins individual element
operations. As such, it is responsible for:
\begin{enumerate}
\sphinxsetlistlabels{\arabic}{enumi}{enumii}{}{.}%
\item {} 
\sphinxAtStartPar
calling any Kernel and expanding any Buit\sphinxhyphen{}In so that they iterate over
their specified interation space;

\item {} 
\sphinxAtStartPar
map the Kernel and Built\sphinxhyphen{}In arguments supplied by an Algorithm \sphinxcode{\sphinxupquote{invoke}}
call to the arguments required by a Built\sphinxhyphen{}in or Kernel method;

\item {} 
\sphinxAtStartPar
include any required distributed\sphinxhyphen{}memory operations such as halo swaps
and reductions to guarantee the correctness of the code;

\item {} 
\sphinxAtStartPar
providing an entry point for the optimisation expert to provide
additional optimisations for the target architecture.

\end{enumerate}

\sphinxAtStartPar
The PSy layer can be written manually but this is error prone and
potentially complex to optimise. Therefore, the PSyclone code\sphinxhyphen{}generation
system automatically generates an initial version of the PSy layer by
parsing the associated Algorithm file and each of the kernels used in it.
This initial version of the PSy\sphinxhyphen{}layer can be further tuned with the PSyclone
code\sphinxhyphen{}transformation capabilities by providing a transformation script.
Each API comes with a set of specialised transformations designed for that API.

\sphinxAtStartPar
Also, in addition to all the functionality available for PSyIR Nodes, the
PSy\sphinxhyphen{}layer nodes have a \sphinxtitleref{dag()} method (standing for directed acyclic graph),
which outputs the PSyIR nodes and their data dependencies. By default a file in
dot format is output with the name \sphinxcode{\sphinxupquote{dag}} and a file in svg format is
output with the name \sphinxcode{\sphinxupquote{dag.svg}}. The file name can be changed using
the \sphinxcode{\sphinxupquote{file\_name}} optional argument and the output file format can be
changed using the \sphinxcode{\sphinxupquote{file\_format}} optional argument. The file\_format
value is simply passed on to graphviz so the graphviz documentation
should be consulted for valid formats if svg is not required.

\begin{sphinxVerbatim}[commandchars=\\\{\}]
\PYG{o}{\PYGZgt{}}\PYG{o}{\PYGZgt{}}\PYG{o}{\PYGZgt{}}\PYG{+w}{ }\PYG{n}{schedule}\PYG{p}{.}\PYG{n}{dag}\PYG{p}{(}\PYG{n}{file\PYGZus{}name}\PYG{o}{=}\PYG{l+s+s2}{\PYGZdq{}lovely\PYGZdq{}}\PYG{p}{,}\PYG{+w}{ }\PYG{n}{file\PYGZus{}format}\PYG{o}{=}\PYG{l+s+s2}{\PYGZdq{}png\PYGZdq{}}\PYG{p}{)}
\end{sphinxVerbatim}

\begin{sphinxadmonition}{note}{Note:}
\sphinxAtStartPar
The dag method can be called from any node and will
output the dag for that node and all of its children.
\end{sphinxadmonition}

\sphinxAtStartPar
If we were to look at the LFRic eg6 example we would see the
following image:

\noindent{\hspace*{\fill}\sphinxincludegraphics[width=256\sphinxpxdimen]{{dag}.png}\hspace*{\fill}}

\sphinxAtStartPar
In the image, all PSyIR nodes with children are split into a start
vertex and an end vertex (for example the InvokeSchedule node has
both \sphinxtitleref{schedule\_start} and \sphinxtitleref{schedule\_end} vertices).
Blue arrows indicate that there is a parent to child relationship (from
a start node) or a child to parent relationship (to an end node).
Green arrows indicate that a Node depends on another Node later in the
schedule (which we call a forward dependence). Therefore the OMP parallel
loop must complete before the globalsum is performed.
Red arrows indicate that a Node depends on
another Node that is earlier in the schedule (which we call a backward
dependence). However the direction of the red arrows are reversed to
improve the flow of the dag layout. In this example the forward and
backward dependence is the same, however this is not always the
case. The two built\sphinxhyphen{}ins do not depend on each other, so they have no
associated green or red arrows.

\sphinxAtStartPar
The dependence graph output gives an indication of whether nodes can
be moved within the InvokeSchedule. In this case it is valid to run the
built\sphinxhyphen{}ins in either order. The underlying dependence analysis used to
create this graph is used to determine whether a transformation of a
Schedule is valid from the perspective of data dependencies.

\sphinxstepscope


\chapter{LFRic API}
\label{\detokenize{dynamo0p3:lfric-api}}\label{\detokenize{dynamo0p3:id1}}\label{\detokenize{dynamo0p3::doc}}
\sphinxAtStartPar
This section describes the LFRic application programming
interface (API). This API explains what a user needs to write in order
to make use of the LFRic API in PSyclone.

\sphinxAtStartPar
As with the majority of PSyclone APIs, the LFRic API specifies
how a user needs to write the algorithm layer and the kernel layer to
allow PSyclone to generate the PSy layer. These algorithm and kernel
APIs are discussed separately in the following sections.

\sphinxAtStartPar
The LFRic API supports the Met Office’s finite element (hereafter FEM)
based GungHo dynamical core (see \sphinxhref{https://docs.python.org/3/reference/introduction.html\#introduction}{Introduction}).
This dynamical core with atmospheric physics parameterisation
schemes is a part of the Met Office LFRic modelling system {[}\hyperlink{cite.zz_bibliography:id5}{AFH+19}{]},
currently being developed in preparation for exascale computing in the 2020s.
The LFRic repository and the associated wiki are hosted at the \sphinxhref{https://code.metoffice.gov.uk/trac/home}{Met Office
Science Repository Service}.
The code is BSD\sphinxhyphen{}licensed, however browsing the \sphinxhref{https://code.metoffice.gov.uk/trac/lfric/wiki}{LFRic wiki} and
\sphinxhref{https://code.metoffice.gov.uk/trac/lfric/browser}{code repository}
requires login access to MOSRS. For more technical details on the
implementation of LFRic, please see the \sphinxhref{https://code.metoffice.gov.uk/trac/lfric/attachment/wiki/LFRicDocumentationPapers/lfric\_documentation.pdf}{LFRic documentation}.


\section{Algorithm}
\label{\detokenize{dynamo0p3:algorithm}}\label{\detokenize{dynamo0p3:lfric-api-algorithm}}
\sphinxAtStartPar
The general requirements for the structure of an Algorithm are explained
in the {\hyperref[\detokenize{introduction_to_psykal:algorithm-layer}]{\sphinxcrossref{\DUrole{std,std-ref}{Algorithm layer}}}} section. This section explains the
LFRic\sphinxhyphen{}API\sphinxhyphen{}specific specialisations and extensions.

\sphinxAtStartPar
The LFRic API defines a set of objects, with specific meanings and
data\sphinxhyphen{}structures, that can be provided as arguments to Kernels within
invoke calls. These are: {\hyperref[\detokenize{dynamo0p3:lfric-scalar}]{\sphinxcrossref{\DUrole{std,std-ref}{scalar}}}},
{\hyperref[\detokenize{dynamo0p3:lfric-field}]{\sphinxcrossref{\DUrole{std,std-ref}{field}}}}, {\hyperref[\detokenize{dynamo0p3:lfric-field-vector}]{\sphinxcrossref{\DUrole{std,std-ref}{field vector}}}},
{\hyperref[\detokenize{dynamo0p3:lfric-operator}]{\sphinxcrossref{\DUrole{std,std-ref}{operator}}}},
{\hyperref[\detokenize{dynamo0p3:lfric-cma-operator}]{\sphinxcrossref{\DUrole{std,std-ref}{column\sphinxhyphen{}wise operator}}}},
{\hyperref[\detokenize{dynamo0p3:lfric-quadrature}]{\sphinxcrossref{\DUrole{std,std-ref}{Quadrature}}}},
{\hyperref[\detokenize{dynamo0p3:lfric-halo-depth}]{\sphinxcrossref{\DUrole{std,std-ref}{Halo Depth}}}} and
{\hyperref[\detokenize{dynamo0p3:lfric-alg-stencil}]{\sphinxcrossref{\DUrole{std,std-ref}{Stencil Extents}}}}. The example below showcases the
use of each of these arguments:

\begin{sphinxVerbatim}[commandchars=\\\{\}]
\PYG{k+kt}{real}\PYG{p}{(}\PYG{n+nb}{kind}\PYG{o}{=}\PYG{n}{r\PYGZus{}def}\PYG{p}{)}\PYG{+w}{           }\PYG{k+kd}{::}\PYG{+w}{ }\PYG{n}{rscalar}
\PYG{k+kt}{integer}\PYG{p}{(}\PYG{n+nb}{kind}\PYG{o}{=}\PYG{n}{i\PYGZus{}def}\PYG{p}{)}\PYG{+w}{        }\PYG{k+kd}{::}\PYG{+w}{ }\PYG{n}{iscalar}\PYG{p}{,}\PYG{+w}{ }\PYG{n}{halo\PYGZus{}depth}
\PYG{k+kt}{logical}\PYG{p}{(}\PYG{n+nb}{kind}\PYG{o}{=}\PYG{n}{l\PYGZus{}def}\PYG{p}{)}\PYG{+w}{        }\PYG{k+kd}{::}\PYG{+w}{ }\PYG{n}{lscalar}
\PYG{k+kt}{integer}\PYG{p}{(}\PYG{n+nb}{kind}\PYG{o}{=}\PYG{n}{i\PYGZus{}def}\PYG{p}{)}\PYG{+w}{        }\PYG{k+kd}{::}\PYG{+w}{ }\PYG{n}{stencil\PYGZus{}extent}
\PYG{k}{type}\PYG{p}{(}\PYG{n}{field\PYGZus{}type}\PYG{p}{)}\PYG{+w}{           }\PYG{k+kd}{::}\PYG{+w}{ }\PYG{n}{field1}\PYG{p}{,}\PYG{+w}{ }\PYG{n}{field2}\PYG{p}{,}\PYG{+w}{ }\PYG{n}{field3}
\PYG{k}{type}\PYG{p}{(}\PYG{n}{field\PYGZus{}type}\PYG{p}{)}\PYG{+w}{           }\PYG{k+kd}{::}\PYG{+w}{ }\PYG{n}{field5}\PYG{p}{(}\PYG{l+m+mi}{3}\PYG{p}{)}\PYG{p}{,}\PYG{+w}{ }\PYG{n}{field6}\PYG{p}{(}\PYG{l+m+mi}{3}\PYG{p}{)}
\PYG{k}{type}\PYG{p}{(}\PYG{n}{integer\PYGZus{}field\PYGZus{}type}\PYG{p}{)}\PYG{+w}{   }\PYG{k+kd}{::}\PYG{+w}{ }\PYG{n}{field7}
\PYG{k}{type}\PYG{p}{(}\PYG{n}{quadrature\PYGZus{}type}\PYG{p}{)}\PYG{+w}{      }\PYG{k+kd}{::}\PYG{+w}{ }\PYG{n}{qr}
\PYG{k}{type}\PYG{p}{(}\PYG{n}{operator\PYGZus{}type}\PYG{p}{)}\PYG{+w}{        }\PYG{k+kd}{::}\PYG{+w}{ }\PYG{n}{operator1}
\PYG{k}{type}\PYG{p}{(}\PYG{n}{columnwise\PYGZus{}operator\PYGZus{}type}\PYG{p}{)}\PYG{+w}{ }\PYG{k+kd}{::}\PYG{+w}{ }\PYG{n}{cma\PYGZus{}op1}
\PYG{p}{.}\PYG{p}{.}\PYG{p}{.}
\PYG{k}{call }\PYG{n}{invoke}\PYG{p}{(}\PYG{+w}{ }\PYG{n}{kernel1}\PYG{p}{(}\PYG{n}{field1}\PYG{p}{,}\PYG{+w}{ }\PYG{n}{field2}\PYG{p}{,}\PYG{+w}{ }\PYG{n}{operator1}\PYG{p}{,}\PYG{+w}{ }\PYG{n}{qr}\PYG{p}{)}\PYG{p}{,}\PYG{+w}{           }\PYG{p}{\PYGZam{}}
\PYG{+w}{             }\PYG{n}{builtin1}\PYG{p}{(}\PYG{n}{rscalar}\PYG{p}{,}\PYG{+w}{ }\PYG{n}{field2}\PYG{p}{,}\PYG{+w}{ }\PYG{n}{field3}\PYG{p}{)}\PYG{p}{,}\PYG{+w}{                }\PYG{p}{\PYGZam{}}
\PYG{+w}{             }\PYG{n}{int\PYGZus{}builtin2}\PYG{p}{(}\PYG{n}{iscalar}\PYG{p}{,}\PYG{+w}{ }\PYG{n}{field7}\PYG{p}{)}\PYG{p}{,}\PYG{+w}{                    }\PYG{p}{\PYGZam{}}
\PYG{+w}{             }\PYG{n}{kernel2}\PYG{p}{(}\PYG{n}{field1}\PYG{p}{,}\PYG{+w}{ }\PYG{n}{stencil\PYGZus{}extent}\PYG{p}{,}\PYG{+w}{ }\PYG{n}{field3}\PYG{p}{,}\PYG{+w}{ }\PYG{n}{lscalar}\PYG{p}{)}\PYG{p}{,}\PYG{+w}{ }\PYG{p}{\PYGZam{}}
\PYG{+w}{             }\PYG{n}{kernel3}\PYG{p}{(}\PYG{n}{field1}\PYG{p}{,}\PYG{+w}{ }\PYG{n}{halo\PYGZus{}depth}\PYG{p}{)}\PYG{p}{,}\PYG{+w}{                      }\PYG{p}{\PYGZam{}}
\PYG{+w}{             }\PYG{n}{assembly\PYGZus{}kernel}\PYG{p}{(}\PYG{n}{cma\PYGZus{}op1}\PYG{p}{,}\PYG{+w}{ }\PYG{n}{operator1}\PYG{p}{)}\PYG{p}{,}\PYG{+w}{              }\PYG{p}{\PYGZam{}}
\PYG{+w}{             }\PYG{n}{name}\PYG{o}{=}\PYG{l+s+s2}{\PYGZdq{}some\PYGZus{}calculation\PYGZdq{}}\PYG{p}{)}
\end{sphinxVerbatim}

\sphinxAtStartPar
Each of these argument types is described in more detail in
the next {\hyperref[\detokenize{dynamo0p3:lfric-alg-arg-types}]{\sphinxcrossref{\DUrole{std,std-ref}{section}}}}.

\sphinxAtStartPar
The LFRic API has support for inter\sphinxhyphen{}grid kernels (those that
map fields between grids of different resolution). At the Algorithm
layer, an \sphinxcode{\sphinxupquote{invoke}} of such kernels looks much like an
\sphinxcode{\sphinxupquote{invoke}} containing general\sphinxhyphen{}purpose kernels. The only restrictions to be
aware of are that inter\sphinxhyphen{}grid kernels accept only field or field\sphinxhyphen{}vectors
as arguments and that an \sphinxcode{\sphinxupquote{invoke}} may not mix inter\sphinxhyphen{}grid kernels with
any other kernel type.


\section{Algorithm Argument Types}
\label{\detokenize{dynamo0p3:algorithm-argument-types}}\label{\detokenize{dynamo0p3:lfric-alg-arg-types}}

\subsection{Scalar}
\label{\detokenize{dynamo0p3:scalar}}\label{\detokenize{dynamo0p3:lfric-scalar}}
\sphinxAtStartPar
In the LFRic API a scalar is a single\sphinxhyphen{}valued argument that is identified
with \sphinxcode{\sphinxupquote{GH\_SCALAR}} metadata. Scalar arguments can have \sphinxcode{\sphinxupquote{real}},
\sphinxcode{\sphinxupquote{integer}} or \sphinxcode{\sphinxupquote{logical}} data type in {\hyperref[\detokenize{dynamo0p3:lfric-kernel-valid-data-type}]{\sphinxcrossref{\DUrole{std,std-ref}{user\sphinxhyphen{}defined Kernels}}}} (\sphinxcode{\sphinxupquote{logical}} data type is not supported
in the {\hyperref[\detokenize{dynamo0p3:lfric-built-ins-dtype-access}]{\sphinxcrossref{\DUrole{std,std-ref}{LFRic Built\sphinxhyphen{}ins}}}}).


\subsection{Field}
\label{\detokenize{dynamo0p3:field}}\label{\detokenize{dynamo0p3:lfric-field}}
\sphinxAtStartPar
LFRic API fields, identified with \sphinxcode{\sphinxupquote{GH\_FIELD}} metadata, represent
FEM discretisations of various dynamical core prognostic and diagnostic
variables. In FEM, variables are discretised by placing them into a
function space (see {\hyperref[\detokenize{dynamo0p3:lfric-function-space}]{\sphinxcrossref{\DUrole{std,std-ref}{Supported Function Spaces}}}}) from which they
inherit a polynomial expansion via the basis functions of that space.
Field values at points within a cell are evaluated as the sum of a set
of basis functions multiplied by coefficients which are the data points.
Points of evaluation are determined by a quadrature object
({\hyperref[\detokenize{dynamo0p3:lfric-quadrature}]{\sphinxcrossref{\DUrole{std,std-ref}{Quadrature}}}}) and are independent of the function space
the field is on. Placement of field data points, also called degrees of
freedom (hereafter “DoFs”), is determined by the function space the field
is on.
LFRic fields passed as arguments to any {\hyperref[\detokenize{dynamo0p3:lfric-kernel-valid-data-type}]{\sphinxcrossref{\DUrole{std,std-ref}{LFRic kernel}}}} can be of \sphinxcode{\sphinxupquote{real}} or \sphinxcode{\sphinxupquote{integer}}
primitive type. In the LFRic infrastructure, these fields are
represented by instances of the \sphinxcode{\sphinxupquote{field\_type}} and \sphinxcode{\sphinxupquote{integer\_field\_type}}
classes, respectively.


\subsection{Field Vector}
\label{\detokenize{dynamo0p3:field-vector}}\label{\detokenize{dynamo0p3:lfric-field-vector}}
\sphinxAtStartPar
Depending on the function space a field lives on, the field data value at
a point can be a scalar or a vector (see {\hyperref[\detokenize{dynamo0p3:lfric-function-space}]{\sphinxcrossref{\DUrole{std,std-ref}{Supported Function Spaces}}}}
for the list of scalar and vector function spaces). There is an
additional option, called a \sphinxstyleemphasis{field vector}, to represent a bundle of
either scalar\sphinxhyphen{} or vector\sphinxhyphen{}valued fields.
Field vectors are represented as \sphinxcode{\sphinxupquote{GH\_FIELD*N}} where \sphinxcode{\sphinxupquote{N}} is the
size of the vector. The 3D coordinate field, for example, has
\sphinxcode{\sphinxupquote{(x, y, z)}} scalar values at the nodes and therefore has a
vector size of 3.


\subsection{Operator}
\label{\detokenize{dynamo0p3:operator}}\label{\detokenize{dynamo0p3:lfric-operator}}
\sphinxAtStartPar
Represents a matrix constructed on a per\sphinxhyphen{}cell basis using Local
Matrix Assembly (LMA) and is identified with \sphinxcode{\sphinxupquote{GH\_OPERATOR}}
metadata. In the LFRic infrastructure, operators are represented by
instances of the \sphinxcode{\sphinxupquote{operator\_type}} class. LFRic operators can only
have \sphinxcode{\sphinxupquote{real}}\sphinxhyphen{}valued data in {\hyperref[\detokenize{dynamo0p3:lfric-kernel-valid-data-type}]{\sphinxcrossref{\DUrole{std,std-ref}{user\sphinxhyphen{}defined Kernels}}}} ({\hyperref[\detokenize{dynamo0p3:lfric-built-ins}]{\sphinxcrossref{\DUrole{std,std-ref}{LFRic Built\sphinxhyphen{}ins}}}}
do not currently support operators).


\subsection{Column\sphinxhyphen{}wise Operator}
\label{\detokenize{dynamo0p3:column-wise-operator}}\label{\detokenize{dynamo0p3:lfric-cma-operator}}
\sphinxAtStartPar
The LFRic API has support for the construction and use of
column\sphinxhyphen{}wise/Column Matrix Assembly (CMA) operators whose metadata
identifier is \sphinxcode{\sphinxupquote{GH\_COLUMNWISE\_OPERATOR}}. In the LFRic
infrastructure, column\sphinxhyphen{}wise operators are represented by instances
of the \sphinxcode{\sphinxupquote{columnwise\_operator\_type}} class. As for the LMA operators
above, LFRic column\sphinxhyphen{}wise operators can only have \sphinxcode{\sphinxupquote{real}}\sphinxhyphen{}valued
{\hyperref[\detokenize{dynamo0p3:lfric-kernel-valid-data-type}]{\sphinxcrossref{\DUrole{std,std-ref}{data}}}}.

\sphinxAtStartPar
As the name suggests, these are operators constructed for a whole
column of the mesh. These are themselves constructed from the
Local Matrix Assembly (LMA) operators of each cell in the column.
The rules governing Kernels that have CMA operators as arguments
are given in the {\hyperref[\detokenize{dynamo0p3:lfric-kernel}]{\sphinxcrossref{\DUrole{std,std-ref}{Kernel}}}} section below.

\sphinxAtStartPar
There are three recognised Kernel types involving CMA operations;
construction, application (including inverse application) and
matrix\sphinxhyphen{}matrix. The following example sketches\sphinxhyphen{}out what the use
of such kernels might look like in the Algorithm layer:

\begin{sphinxVerbatim}[commandchars=\\\{\}]
\PYG{k}{use }\PYG{n}{field\PYGZus{}mod}\PYG{p}{,}\PYG{+w}{ }\PYG{k}{only}\PYG{p}{:}\PYG{+w}{ }\PYG{n}{field\PYGZus{}type}
\PYG{k}{use }\PYG{n}{operator\PYGZus{}mod}\PYG{p}{,}\PYG{+w}{ }\PYG{k}{only}\PYG{+w}{ }\PYG{p}{:}\PYG{+w}{ }\PYG{n}{operator\PYGZus{}type}
\PYG{k}{use }\PYG{n}{columnwise\PYGZus{}operator\PYGZus{}mod}\PYG{p}{,}\PYG{+w}{ }\PYG{k}{only}\PYG{+w}{ }\PYG{p}{:}\PYG{+w}{ }\PYG{n}{columnwise\PYGZus{}operator\PYGZus{}type}
\PYG{k}{type}\PYG{p}{(}\PYG{n}{field\PYGZus{}type}\PYG{p}{)}\PYG{+w}{ }\PYG{k+kd}{::}\PYG{+w}{ }\PYG{n}{field1}\PYG{p}{,}\PYG{+w}{ }\PYG{n}{field2}\PYG{p}{,}\PYG{+w}{ }\PYG{n}{field3}
\PYG{k}{type}\PYG{p}{(}\PYG{n}{operator\PYGZus{}type}\PYG{p}{)}\PYG{+w}{ }\PYG{k+kd}{::}\PYG{+w}{ }\PYG{n}{lma\PYGZus{}op1}\PYG{p}{,}\PYG{+w}{ }\PYG{n}{lma\PYGZus{}op2}
\PYG{k}{type}\PYG{p}{(}\PYG{n}{columnwise\PYGZus{}operator\PYGZus{}type}\PYG{p}{)}\PYG{+w}{ }\PYG{k+kd}{::}\PYG{+w}{ }\PYG{n}{cma\PYGZus{}op1}\PYG{p}{,}\PYG{+w}{ }\PYG{n}{cma\PYGZus{}op2}\PYG{p}{,}\PYG{+w}{ }\PYG{n}{cma\PYGZus{}op3}
\PYG{k+kt}{real}\PYG{p}{(}\PYG{n+nb}{kind}\PYG{o}{=}\PYG{n}{r\PYGZus{}def}\PYG{p}{)}\PYG{+w}{ }\PYG{k+kd}{::}\PYG{+w}{ }\PYG{n}{alpha}
\PYG{p}{.}\PYG{p}{.}\PYG{p}{.}
\PYG{k}{call }\PYG{n}{invoke}\PYG{p}{(}\PYG{+w}{                                                    }\PYG{p}{\PYGZam{}}
\PYG{+w}{        }\PYG{n}{assembly\PYGZus{}kernel}\PYG{p}{(}\PYG{n}{cma\PYGZus{}op1}\PYG{p}{,}\PYG{+w}{ }\PYG{n}{lma\PYGZus{}op1}\PYG{p}{,}\PYG{+w}{ }\PYG{n}{lma\PYGZus{}op2}\PYG{p}{)}\PYG{p}{,}\PYG{+w}{             }\PYG{p}{\PYGZam{}}
\PYG{+w}{        }\PYG{n}{assembly\PYGZus{}kernel2}\PYG{p}{(}\PYG{n}{cma\PYGZus{}op2}\PYG{p}{,}\PYG{+w}{ }\PYG{n}{lma\PYGZus{}op1}\PYG{p}{,}\PYG{+w}{ }\PYG{n}{lma\PYGZus{}op2}\PYG{p}{,}\PYG{+w}{ }\PYG{n}{field3}\PYG{p}{)}\PYG{p}{,}\PYG{+w}{    }\PYG{p}{\PYGZam{}}
\PYG{+w}{        }\PYG{n}{apply\PYGZus{}kernel}\PYG{p}{(}\PYG{n}{field1}\PYG{p}{,}\PYG{+w}{ }\PYG{n}{field2}\PYG{p}{,}\PYG{+w}{ }\PYG{n}{cma\PYGZus{}op1}\PYG{p}{)}\PYG{p}{,}\PYG{+w}{                  }\PYG{p}{\PYGZam{}}
\PYG{+w}{        }\PYG{n}{matrix\PYGZus{}matrix\PYGZus{}kernel}\PYG{p}{(}\PYG{n}{cma\PYGZus{}op3}\PYG{p}{,}\PYG{+w}{ }\PYG{n}{cma\PYGZus{}op1}\PYG{p}{,}\PYG{+w}{ }\PYG{n}{alpha}\PYG{p}{,}\PYG{+w}{ }\PYG{n}{cma\PYGZus{}op2}\PYG{p}{)}\PYG{p}{,}\PYG{+w}{ }\PYG{p}{\PYGZam{}}
\PYG{+w}{        }\PYG{n}{apply\PYGZus{}kernel}\PYG{p}{(}\PYG{n}{field3}\PYG{p}{,}\PYG{+w}{ }\PYG{n}{field1}\PYG{p}{,}\PYG{+w}{ }\PYG{n}{cma\PYGZus{}op3}\PYG{p}{)}\PYG{p}{,}\PYG{+w}{                  }\PYG{p}{\PYGZam{}}
\PYG{+w}{        }\PYG{n}{name}\PYG{o}{=}\PYG{l+s+s2}{\PYGZdq{}cma\PYGZus{}example\PYGZdq{}}\PYG{p}{)}
\end{sphinxVerbatim}

\sphinxAtStartPar
The above invoke uses two LMA operators to construct the CMA operator
\sphinxcode{\sphinxupquote{cma\_op1}}.  A second CMA operator, \sphinxcode{\sphinxupquote{cma\_op2}}, is assembled from
the same two LMA operators but also uses a field. The first of these
CMA operators is then applied to \sphinxcode{\sphinxupquote{field2}} and the result stored in
\sphinxcode{\sphinxupquote{field1}} (assuming that the metadata for \sphinxcode{\sphinxupquote{apply\_kernel}} specifies
that it is the first field argument that is written to). The two CMA
operators are then combined to produce a third, \sphinxcode{\sphinxupquote{cma\_op3}}. This is
then applied to \sphinxcode{\sphinxupquote{field1}} and the result stored in \sphinxcode{\sphinxupquote{field3}}.

\sphinxAtStartPar
Note that PSyclone identifies the type of kernels performing
column\sphinxhyphen{}wise operations based on their arguments as described in
metadata (see {\hyperref[\detokenize{dynamo0p3:lfric-cma-mdata-rules}]{\sphinxcrossref{\DUrole{std,std-ref}{Rules for Kernels that work with CMA Operators}}}} below). The names of the
kernels in the above example are purely illustrative and are not used
by PSyclone when determining kernel type.

\sphinxAtStartPar
A full example of CMA operator construction is available in
\sphinxcode{\sphinxupquote{examples/lfric/eg7}}.


\subsection{Quadrature}
\label{\detokenize{dynamo0p3:quadrature}}\label{\detokenize{dynamo0p3:lfric-quadrature}}
\sphinxAtStartPar
Kernels conforming to the LFRic API may require quadrature
information (specified using e.g. \sphinxcode{\sphinxupquote{gh\_shape = gh\_quadrature\_XYoZ}} in
the kernel metadata \sphinxhyphen{} see Section {\hyperref[\detokenize{dynamo0p3:lfric-gh-shape}]{\sphinxcrossref{\DUrole{std,std-ref}{gh\_shape and gh\_evaluator\_targets}}}}). This
information must be passed to the kernel from the Algorithm layer in
the form of one or more \sphinxcode{\sphinxupquote{quadrature\_type}} objects. These must be the
last arguments passed to the kernel (with the exception of \sphinxcode{\sphinxupquote{halo\_depth}}
\sphinxhyphen{} {\hyperref[\detokenize{dynamo0p3:lfric-halo-depth}]{\sphinxcrossref{\DUrole{std,std-ref}{Halo Depth}}}} \sphinxhyphen{} if the kernel requires it) and must be
provided in the same
order that they are specified in the kernel metadata, e.g. if the
metadata for kernel \sphinxcode{\sphinxupquote{pressure\_gradient\_kernel\_type}} specified
\sphinxcode{\sphinxupquote{gh\_shape = gh\_quadrature\_XYoZ}} and that for kernel
\sphinxcode{\sphinxupquote{geopotential\_gradient\_kernel}} had \sphinxcode{\sphinxupquote{gh\_shape(2) = (\textbackslash{}
gh\_quadrature\_XYoZ, gh\_quadrature\_face \textbackslash{})}} then the corresponding
invoke would look something like:

\begin{sphinxVerbatim}[commandchars=\\\{\}]
\PYG{p}{.}\PYG{p}{.}\PYG{p}{.}
\PYG{n}{qr\PYGZus{}xyoz}\PYG{+w}{ }\PYG{o}{=}\PYG{+w}{ }\PYG{n}{quadrature\PYGZus{}xyoz\PYGZus{}type}\PYG{p}{(}\PYG{n}{nqp\PYGZus{}h\PYGZus{}exact}\PYG{p}{,}\PYG{+w}{ }\PYG{n}{nqp\PYGZus{}h\PYGZus{}exact}\PYG{p}{,}\PYG{+w}{ }\PYG{n}{nqp\PYGZus{}v\PYGZus{}exact}\PYG{p}{,}\PYG{+w}{ }\PYG{n}{rule}\PYG{p}{)}
\PYG{n}{qr\PYGZus{}face}\PYG{+w}{ }\PYG{o}{=}\PYG{+w}{ }\PYG{n}{quadrature\PYGZus{}face\PYGZus{}type}\PYG{p}{(}\PYG{n}{nqp\PYGZus{}h\PYGZus{}exact}\PYG{p}{,}\PYG{+w}{ }\PYG{n}{nqp\PYGZus{}v\PYGZus{}exact}\PYG{p}{,}\PYG{+w}{ }\PYG{p}{.}\PYG{p}{.}\PYG{p}{.}\PYG{p}{,}\PYG{+w}{ }\PYG{n}{rule}\PYG{p}{)}
\PYG{k}{call }\PYG{n}{invoke}\PYG{p}{(}\PYG{n}{pressure\PYGZus{}gradient\PYGZus{}kernel\PYGZus{}type}\PYG{p}{(}\PYG{n}{rhs\PYGZus{}tmp}\PYG{p}{(}\PYG{n}{igh\PYGZus{}u}\PYG{p}{)}\PYG{p}{,}\PYG{+w}{ }\PYG{n}{rho}\PYG{p}{,}\PYG{+w}{ }\PYG{n}{theta}\PYG{p}{,}\PYG{+w}{ }\PYG{n}{qr\PYGZus{}xyoz}\PYG{p}{)}\PYG{p}{,}\PYG{+w}{ }\PYG{p}{\PYGZam{}}
\PYG{+w}{            }\PYG{n}{geopotential\PYGZus{}gradient\PYGZus{}kernel\PYGZus{}type}\PYG{p}{(}\PYG{n}{rhs\PYGZus{}tmp}\PYG{p}{(}\PYG{n}{igh\PYGZus{}u}\PYG{p}{)}\PYG{p}{,}\PYG{+w}{ }\PYG{n}{geopotential}\PYG{p}{,}\PYG{+w}{ }\PYG{p}{\PYGZam{}}
\PYG{+w}{                                              }\PYG{n}{qr\PYGZus{}xyoz}\PYG{p}{,}\PYG{+w}{ }\PYG{n}{qr\PYGZus{}face}\PYG{p}{)}\PYG{p}{)}
\end{sphinxVerbatim}

\sphinxAtStartPar
These quadrature objects specify the set(s) of points at which the
basis/differential\sphinxhyphen{}basis functions required by the kernel are to be evaluated.

\sphinxAtStartPar
The arguments \sphinxcode{\sphinxupquote{nqp\_h\_exact}} and \sphinxcode{\sphinxupquote{nqp\_v\_exact}} are namelist options from the
finite\sphinxhyphen{}element configuration file, generated as a function of the
finite\sphinxhyphen{}element model (FEM) order in the horizontal and vertical directions (see
{\hyperref[\detokenize{dynamo0p3:lfric-function-space}]{\sphinxcrossref{\DUrole{std,std-ref}{Supported Function Spaces}}}}). They specify the number of quadrature
points in the horizontal and vertical directions, respectively.

\sphinxAtStartPar
Here, \sphinxcode{\sphinxupquote{qr\_face}} has been passed the horizontal and vertical numbers of
quadrature points, meaning that the face in question is normal to a horizontal
plane. In general, \sphinxcode{\sphinxupquote{qr\_xyoz}} takes the number of quadrature points in each
direction (\sphinxstyleemphasis{x}, \sphinxstyleemphasis{y}, and \sphinxstyleemphasis{z}, in that order), whereas \sphinxcode{\sphinxupquote{qr\_face}} takes the
number of quadrature points in each direction tangential to the face in
question.


\subsection{Halo Depth}
\label{\detokenize{dynamo0p3:halo-depth}}\label{\detokenize{dynamo0p3:lfric-halo-depth}}
\sphinxAtStartPar
If a Kernel is written to iterate into the halo (has an \sphinxcode{\sphinxupquote{OPERATES\_ON}} of
\sphinxcode{\sphinxupquote{HALO\_CELL\_COLUMN}} or \sphinxcode{\sphinxupquote{OWNED\_AND\_HALO\_CELL\_COLUMN}}) then the halo depth
must be passed as a final, \sphinxcode{\sphinxupquote{integer}} argument to the Kernel.


\subsection{Stencil Extent}
\label{\detokenize{dynamo0p3:stencil-extent}}\label{\detokenize{dynamo0p3:lfric-alg-stencil}}
\sphinxAtStartPar
The metadata for a Kernel which operates on a cell\sphinxhyphen{}column may specify
that a Kernel performs a stencil operation on a field. Any such
metadata must provide a stencil type. See the
{\hyperref[\detokenize{dynamo0p3:lfric-api-meta-args}]{\sphinxcrossref{\DUrole{std,std-ref}{meta\_args}}}} section for more details. The supported
stencil types are \sphinxcode{\sphinxupquote{X1D}}, \sphinxcode{\sphinxupquote{Y1D}}, \sphinxcode{\sphinxupquote{XORY1D}}, \sphinxcode{\sphinxupquote{CROSS}}, \sphinxcode{\sphinxupquote{CROSS2D}} or
\sphinxcode{\sphinxupquote{REGION}}.

\sphinxAtStartPar
If a stencil operation is specified by the Kernel metadata, the
Algorithm layer must provide the \sphinxcode{\sphinxupquote{extent}} of the stencil (the
maximum distance from the central cell that the stencil extends). The
LFRic API expects this information to be added as an additional
\sphinxcode{\sphinxupquote{integer}} argument immediately after the relevant field when specifying
the Kernel via an \sphinxcode{\sphinxupquote{invoke}}.

\sphinxAtStartPar
For example:

\begin{sphinxVerbatim}[commandchars=\\\{\}]
\PYG{k+kt}{integer}\PYG{p}{(}\PYG{n+nb}{kind}\PYG{o}{=}\PYG{n}{i\PYGZus{}def}\PYG{p}{)}\PYG{+w}{ }\PYG{k+kd}{::}\PYG{+w}{ }\PYG{n}{extent}\PYG{+w}{ }\PYG{o}{=}\PYG{+w}{ }\PYG{l+m+mi}{2}
\PYG{k}{call }\PYG{n}{invoke}\PYG{p}{(}\PYG{n}{kernel}\PYG{p}{(}\PYG{n}{field1}\PYG{p}{,}\PYG{+w}{ }\PYG{n}{field2}\PYG{p}{,}\PYG{+w}{ }\PYG{n}{extent}\PYG{p}{)}\PYG{p}{)}
\end{sphinxVerbatim}

\sphinxAtStartPar
where \sphinxcode{\sphinxupquote{field2}} has kernel metadata specifying that it has a stencil
access.

\sphinxAtStartPar
\sphinxcode{\sphinxupquote{extent}}  may also be passed as a literal. For example:

\begin{sphinxVerbatim}[commandchars=\\\{\}]
\PYG{k}{call }\PYG{n}{invoke}\PYG{p}{(}\PYG{n}{kernel}\PYG{p}{(}\PYG{n}{field1}\PYG{p}{,}\PYG{+w}{ }\PYG{n}{field2}\PYG{p}{,}\PYG{+w}{ }\PYG{l+m+mi}{2}\PYG{p}{)}\PYG{p}{)}
\end{sphinxVerbatim}

\sphinxAtStartPar
where, again, \sphinxcode{\sphinxupquote{field2}} has kernel metadata specifying that it has a
stencil access.

\begin{sphinxadmonition}{note}{Note:}
\sphinxAtStartPar
The stencil extent specified in the Algorithm layer is not the
same as the stencil size passed in to the Kernel. The latter
contains the number of cells in the stencil which is dependent
on both the stencil type and extent.
\end{sphinxadmonition}

\sphinxAtStartPar
If the Kernel metadata specifies that the stencil is of type
\sphinxcode{\sphinxupquote{XORY1D}} (which means \sphinxcode{\sphinxupquote{X1D}} or \sphinxcode{\sphinxupquote{Y1D}}) then the algorithm layer
must specify whether the stencil is \sphinxcode{\sphinxupquote{X1D}} or \sphinxcode{\sphinxupquote{Y1D}} for that
particular kernel call. The LFRic API expects this information to
be added as an additional argument immediately after the relevant
stencil extent argument. The argument should be an \sphinxcode{\sphinxupquote{integer}} with
valid values being \sphinxcode{\sphinxupquote{x\_direction}} or \sphinxcode{\sphinxupquote{y\_direction}}, both being
supplied by the \sphinxcode{\sphinxupquote{LFRic}} infrastructure via the
\sphinxcode{\sphinxupquote{flux\_direction\_mod}} fortran module

\sphinxAtStartPar
For example:

\begin{sphinxVerbatim}[commandchars=\\\{\}]
\PYG{k}{use }\PYG{n}{flux\PYGZus{}direction\PYGZus{}mod}\PYG{p}{,}\PYG{+w}{ }\PYG{k}{only}\PYG{+w}{ }\PYG{p}{:}\PYG{+w}{ }\PYG{n}{x\PYGZus{}direction}
\PYG{k+kt}{integer}\PYG{p}{(}\PYG{n+nb}{kind}\PYG{o}{=}\PYG{n}{i\PYGZus{}def}\PYG{p}{)}\PYG{+w}{ }\PYG{k+kd}{::}\PYG{+w}{ }\PYG{n}{direction}\PYG{+w}{ }\PYG{o}{=}\PYG{+w}{ }\PYG{n}{x\PYGZus{}direction}
\PYG{k+kt}{integer}\PYG{p}{(}\PYG{n+nb}{kind}\PYG{o}{=}\PYG{n}{i\PYGZus{}def}\PYG{p}{)}\PYG{+w}{ }\PYG{k+kd}{::}\PYG{+w}{ }\PYG{n}{extent}\PYG{+w}{ }\PYG{o}{=}\PYG{+w}{ }\PYG{l+m+mi}{2}
\PYG{c}{! ...}
\PYG{k}{call }\PYG{n}{invoke}\PYG{p}{(}\PYG{n}{kernel}\PYG{p}{(}\PYG{n}{field1}\PYG{p}{,}\PYG{+w}{ }\PYG{n}{field2}\PYG{p}{,}\PYG{+w}{ }\PYG{n}{extent}\PYG{p}{,}\PYG{+w}{ }\PYG{n}{direction}\PYG{p}{)}\PYG{p}{)}
\end{sphinxVerbatim}

\sphinxAtStartPar
\sphinxcode{\sphinxupquote{direction}} may also be passed as a literal. For example:

\begin{sphinxVerbatim}[commandchars=\\\{\}]
\PYG{k}{use }\PYG{n}{flux\PYGZus{}direction\PYGZus{}mod}\PYG{p}{,}\PYG{+w}{ }\PYG{k}{only}\PYG{+w}{ }\PYG{p}{:}\PYG{+w}{ }\PYG{n}{x\PYGZus{}direction}
\PYG{k+kt}{integer}\PYG{p}{(}\PYG{n+nb}{kind}\PYG{o}{=}\PYG{n}{i\PYGZus{}def}\PYG{p}{)}\PYG{+w}{ }\PYG{k+kd}{::}\PYG{+w}{ }\PYG{n}{extent}\PYG{+w}{ }\PYG{o}{=}\PYG{+w}{ }\PYG{l+m+mi}{2}
\PYG{c}{! ...}
\PYG{k}{call }\PYG{n}{invoke}\PYG{p}{(}\PYG{n}{kernel}\PYG{p}{(}\PYG{n}{field1}\PYG{p}{,}\PYG{+w}{ }\PYG{n}{field2}\PYG{p}{,}\PYG{+w}{ }\PYG{n}{extent}\PYG{p}{,}\PYG{+w}{ }\PYG{n}{x\PYGZus{}direction}\PYG{p}{)}\PYG{p}{)}
\end{sphinxVerbatim}

\sphinxAtStartPar
If the stencil is of type \sphinxcode{\sphinxupquote{CROSS2D}} then the arrays passed to the kernel are
of different dimensions to those of other stencils. The \sphinxcode{\sphinxupquote{CROSS2D}} stencil is
designed for use when it is necessary for a kernel to know where the stencil
cells are, relative to the current cell. For this reason, the \sphinxcode{\sphinxupquote{stencil\_size}}
passed to the kernel is an array of length 4 containing sizes for each branch
of the stencil. The \sphinxcode{\sphinxupquote{stencil\_size}} array is always ordered: West, South,
East, North. This branch dimension is also part of the \sphinxcode{\sphinxupquote{stencil\_dofmap}} array
making it possible to loop over each branch of the stencil individually. The
invoke call for the \sphinxcode{\sphinxupquote{CROSS2D}} stencil remains of the same form as for other
stencils.

\sphinxAtStartPar
If certain fields use the same value of extent and/or direction then
the same variable, or literal value can be provided.

\sphinxAtStartPar
For example:

\begin{sphinxVerbatim}[commandchars=\\\{\}]
\PYG{k}{call }\PYG{n}{invoke}\PYG{p}{(}\PYG{n}{kernel1}\PYG{p}{(}\PYG{n}{field1}\PYG{p}{,}\PYG{+w}{ }\PYG{n}{field2}\PYG{p}{,}\PYG{+w}{ }\PYG{n}{extent}\PYG{p}{,}\PYG{+w}{  }\PYG{n}{field3}\PYG{p}{,}\PYG{+w}{ }\PYG{n}{extent}\PYG{p}{,}\PYG{+w}{ }\PYG{n}{direction}\PYG{p}{)}\PYG{p}{,}\PYG{+w}{ }\PYG{p}{\PYGZam{}}
\PYG{+w}{            }\PYG{n}{kernel2}\PYG{p}{(}\PYG{n}{field1}\PYG{p}{,}\PYG{+w}{ }\PYG{n}{field2}\PYG{p}{,}\PYG{+w}{ }\PYG{n}{extent2}\PYG{p}{,}\PYG{+w}{ }\PYG{n}{field4}\PYG{p}{,}\PYG{+w}{ }\PYG{n}{extent}\PYG{p}{,}\PYG{+w}{ }\PYG{n}{direction}\PYG{p}{)}\PYG{p}{)}
\end{sphinxVerbatim}

\sphinxAtStartPar
In the above example \sphinxcode{\sphinxupquote{field2}} and \sphinxcode{\sphinxupquote{field3}} in \sphinxcode{\sphinxupquote{kernel1}} and
\sphinxcode{\sphinxupquote{field4}} in \sphinxcode{\sphinxupquote{kernel2}} will have the same \sphinxcode{\sphinxupquote{extent}} value but
\sphinxcode{\sphinxupquote{field2}} in \sphinxcode{\sphinxupquote{kernel2}} may have a different value. Similarly,
\sphinxcode{\sphinxupquote{field3}} in \sphinxcode{\sphinxupquote{kernel1}} and \sphinxcode{\sphinxupquote{field4}} in \sphinxcode{\sphinxupquote{kernel2}} will have the
same \sphinxcode{\sphinxupquote{direction}} value.

\sphinxAtStartPar
An example of the use of stencils is available in \sphinxcode{\sphinxupquote{examples/lfric/eg5}}.

\sphinxAtStartPar
There is currently no attempt to perform type checking in PSyclone so
any errors in the type and/or position of arguments will not be picked
up until compile time. However, PSyclone does check for the correct
number of algorithm arguments. If the wrong number of arguments is
provided then an exception is raised.

\sphinxAtStartPar
For example, running test 19.2 from the LFRic API test suite gives:

\begin{sphinxVerbatim}[commandchars=\\\{\}]
\PYG{n+nb}{cd}\PYG{+w}{ }\PYGZlt{}PSYCLONEHOME\PYGZgt{}/src/psyclone/tests
psyclone\PYG{+w}{ }test\PYGZus{}files/dynamo0p3/19.2\PYGZus{}single\PYGZus{}stencil\PYGZus{}broken.f90
\PYG{l+s+s2}{\PYGZdq{}Generation Error: error: expected \PYGZsq{}5\PYGZsq{} arguments in the algorithm layer but found \PYGZsq{}4\PYGZsq{}.}
\PYG{l+s+s2}{Expected \PYGZsq{}4\PYGZsq{} standard arguments, \PYGZsq{}1\PYGZsq{} stencil arguments and \PYGZsq{}0\PYGZsq{} qr\PYGZus{}arguments\PYGZsq{}\PYGZdq{}}
\end{sphinxVerbatim}


\section{Mixed Precision}
\label{\detokenize{dynamo0p3:mixed-precision}}\label{\detokenize{dynamo0p3:lfric-mixed-precision}}
\sphinxAtStartPar
The LFRic API supports the ability to specify the precision required
by the model via precision variables. To make use of this, the code
developer must declare scalars, fields and operators in the algorithm
layer with the required LFRic\sphinxhyphen{}supported precision. In the current
implementation there are two supported precisions for \sphinxcode{\sphinxupquote{REAL}} data and
one each for \sphinxcode{\sphinxupquote{INTEGER}} and \sphinxcode{\sphinxupquote{LOGICAL}} data. The actual precision used in
the code can be set in a configuration file. For example, \sphinxcode{\sphinxupquote{INTEGER}} data
could be set to be 32\sphinxhyphen{}bit precision. As \sphinxcode{\sphinxupquote{REAL}} data has more than one
supported precision, different parts of the code can be configured to
have different precision.

\sphinxAtStartPar
The table below gives the currently supported datatypes, their
associated kernel metadata description and their precision:


\begin{savenotes}\sphinxattablestart
\sphinxthistablewithglobalstyle
\sphinxthistablewithvlinesstyle
\centering
\begin{tabulary}{\linewidth}[t]{|l|l|l|}
\sphinxtoprule
\sphinxstyletheadfamily 
\sphinxAtStartPar
Data Type
&\sphinxstyletheadfamily 
\sphinxAtStartPar
Kernel Metadata
&\sphinxstyletheadfamily 
\sphinxAtStartPar
Precision
\\
\sphinxmidrule
\sphinxtableatstartofbodyhook
\sphinxAtStartPar
REAL(R\_DEF)
&
\sphinxAtStartPar
GH\_SCALAR, GH\_REAL
&
\sphinxAtStartPar
R\_DEF
\\
\sphinxhline
\sphinxAtStartPar
REAL(R\_BL)
&
\sphinxAtStartPar
GH\_SCALAR, GH\_REAL
&
\sphinxAtStartPar
R\_BL
\\
\sphinxhline
\sphinxAtStartPar
REAL(R\_PHYS)
&
\sphinxAtStartPar
GH\_SCALAR, GH\_REAL
&
\sphinxAtStartPar
R\_PHYS
\\
\sphinxhline
\sphinxAtStartPar
REAL(R\_SOLVER)
&
\sphinxAtStartPar
GH\_SCALAR, GH\_REAL
&
\sphinxAtStartPar
R\_SOLVER
\\
\sphinxhline
\sphinxAtStartPar
REAL(R\_TRAN)
&
\sphinxAtStartPar
GH\_SCALAR, GH\_REAL
&
\sphinxAtStartPar
R\_TRAN
\\
\sphinxhline
\sphinxAtStartPar
INTEGER(I\_DEF)
&
\sphinxAtStartPar
GH\_SCALAR, GH\_INTEGER
&
\sphinxAtStartPar
I\_DEF
\\
\sphinxhline
\sphinxAtStartPar
LOGICAL(L\_DEF)
&
\sphinxAtStartPar
GH\_SCALAR, GH\_LOGICAL
&
\sphinxAtStartPar
L\_DEF
\\
\sphinxhline
\sphinxAtStartPar
FIELD\_TYPE
&
\sphinxAtStartPar
GH\_FIELD, GH\_REAL
&
\sphinxAtStartPar
R\_DEF
\\
\sphinxhline
\sphinxAtStartPar
R\_BL\_FIELD\_TYPE
&
\sphinxAtStartPar
GH\_FIELD, GH\_REAL
&
\sphinxAtStartPar
R\_BL
\\
\sphinxhline
\sphinxAtStartPar
R\_PHYS\_FIELD\_TYPE
&
\sphinxAtStartPar
GH\_FIELD, GH\_REAL
&
\sphinxAtStartPar
R\_PHYS
\\
\sphinxhline
\sphinxAtStartPar
R\_SOLVER\_FIELD\_TYPE
&
\sphinxAtStartPar
GH\_FIELD, GH\_REAL
&
\sphinxAtStartPar
R\_SOLVER
\\
\sphinxhline
\sphinxAtStartPar
R\_TRAN\_FIELD\_TYPE
&
\sphinxAtStartPar
GH\_FIELD, GH\_REAL
&
\sphinxAtStartPar
R\_TRAN
\\
\sphinxhline
\sphinxAtStartPar
INTEGER\_FIELD\_TYPE
&
\sphinxAtStartPar
GH\_FIELD, GH\_INTEGER
&
\sphinxAtStartPar
I\_DEF
\\
\sphinxhline
\sphinxAtStartPar
OPERATOR\_TYPE
&
\sphinxAtStartPar
GH\_OPERATOR, GH\_REAL
&
\sphinxAtStartPar
R\_DEF
\\
\sphinxhline
\sphinxAtStartPar
R\_SOLVER\_OPERATOR\_TYPE
&
\sphinxAtStartPar
GH\_OPERATOR, GH\_REAL
&
\sphinxAtStartPar
R\_SOLVER
\\
\sphinxhline
\sphinxAtStartPar
R\_TRAN\_OPERATOR\_TYPE
&
\sphinxAtStartPar
GH\_OPERATOR, GH\_REAL
&
\sphinxAtStartPar
R\_TRAN
\\
\sphinxhline
\sphinxAtStartPar
COLUMNWISE\_OPERATOR\_TYPE
&
\sphinxAtStartPar
GH\_COLUMNWISE\_OPERATOR, GH\_REAL
&
\sphinxAtStartPar
R\_SOLVER
\\
\sphinxbottomrule
\end{tabulary}
\sphinxtableafterendhook\par
\sphinxattableend\end{savenotes}

\sphinxAtStartPar
As can be seen from the above table, the kernel metadata does not
capture all of the precision options. For example, from the metadata
it is not possible to determine whether a \sphinxcode{\sphinxupquote{REAL}} scalar, \sphinxcode{\sphinxupquote{REAL}} field
or \sphinxcode{\sphinxupquote{REAL}} operator has precision \sphinxcode{\sphinxupquote{R\_DEF}}, \sphinxcode{\sphinxupquote{R\_SOLVER}} or \sphinxcode{\sphinxupquote{R\_TRAN}}.

\sphinxAtStartPar
If a scalar, field, or operator is specified with a particular
precision in the algorithm layer then any associated kernels that it
is passed to must have been written so that they support this
precision. If a kernel needs to support data that can be stored with
different precisions then appropriate precision\sphinxhyphen{}specific subroutines
should be written. These precision\sphinxhyphen{}specific subroutine should be
called via a generic interface (which lets Fortran choose the
appropriate subroutine based on the precision of its argument(s)).

\sphinxAtStartPar
Below is a simple example of an algorithm code calling the same
generic kernel twice with potentially different precision. The
implementation of the generic kernel such that it supports both 32\sphinxhyphen{}
and 64\sphinxhyphen{}bit precision is also shown. The use of LFRic names for
precision in the algorithm code allows precision to be controlled in a
simple way. For example, \sphinxcode{\sphinxupquote{r\_solver}} could be set to be 32\sphinxhyphen{}bits in
one configuration and 64\sphinxhyphen{}bits in another:

\begin{sphinxVerbatim}[commandchars=\\\{\}]
\PYG{k}{program }\PYG{n}{test}

\PYG{+w}{  }\PYG{k}{use }\PYG{n}{constants\PYGZus{}mod}\PYG{p}{,}\PYG{+w}{      }\PYG{k}{only}\PYG{+w}{ }\PYG{p}{:}\PYG{+w}{ }\PYG{n}{r\PYGZus{}def}\PYG{p}{,}\PYG{+w}{ }\PYG{n}{r\PYGZus{}solver}
\PYG{+w}{  }\PYG{k}{use }\PYG{n}{field\PYGZus{}mod}\PYG{p}{,}\PYG{+w}{          }\PYG{k}{only}\PYG{+w}{ }\PYG{p}{:}\PYG{+w}{ }\PYG{n}{field\PYGZus{}type}
\PYG{+w}{  }\PYG{k}{use }\PYG{n}{r\PYGZus{}solver\PYGZus{}field\PYGZus{}mod}\PYG{p}{,}\PYG{+w}{ }\PYG{k}{only}\PYG{+w}{ }\PYG{p}{:}\PYG{+w}{ }\PYG{n}{r\PYGZus{}solver\PYGZus{}field\PYGZus{}type}
\PYG{+w}{  }\PYG{k}{use }\PYG{n}{example\PYGZus{}mod}\PYG{p}{,}\PYG{+w}{        }\PYG{k}{only}\PYG{+w}{ }\PYG{p}{:}\PYG{+w}{ }\PYG{n}{example\PYGZus{}type}

\PYG{+w}{  }\PYG{k}{type}\PYG{p}{(}\PYG{n}{field\PYGZus{}type}\PYG{p}{)}\PYG{+w}{          }\PYG{k+kd}{::}\PYG{+w}{ }\PYG{n}{field\PYGZus{}r\PYGZus{}def}
\PYG{+w}{  }\PYG{k}{type}\PYG{p}{(}\PYG{n}{r\PYGZus{}solver\PYGZus{}field\PYGZus{}type}\PYG{p}{)}\PYG{+w}{ }\PYG{k+kd}{::}\PYG{+w}{ }\PYG{n}{field\PYGZus{}r\PYGZus{}solver}
\PYG{+w}{  }\PYG{k+kt}{real}\PYG{p}{(}\PYG{n+nb}{kind}\PYG{o}{=}\PYG{n}{r\PYGZus{}def}\PYG{p}{)}\PYG{+w}{          }\PYG{k+kd}{::}\PYG{+w}{ }\PYG{n}{x\PYGZus{}r\PYGZus{}def}
\PYG{+w}{  }\PYG{k+kt}{real}\PYG{p}{(}\PYG{n+nb}{kind}\PYG{o}{=}\PYG{n}{r\PYGZus{}solver}\PYG{p}{)}\PYG{+w}{       }\PYG{k+kd}{::}\PYG{+w}{ }\PYG{n}{x\PYGZus{}r\PYGZus{}solver}

\PYG{+w}{  }\PYG{k}{call }\PYG{n}{invoke}\PYG{p}{(}\PYG{+w}{ }\PYG{n}{example\PYGZus{}type}\PYG{p}{(}\PYG{n}{field\PYGZus{}r\PYGZus{}def}\PYG{p}{,}\PYG{+w}{ }\PYG{n}{x\PYGZus{}r\PYGZus{}def}\PYG{p}{)}\PYG{p}{,}\PYG{+w}{ }\PYG{p}{\PYGZam{}}
\PYG{+w}{               }\PYG{n}{example\PYGZus{}type}\PYG{p}{(}\PYG{n}{field\PYGZus{}r\PYGZus{}solver}\PYG{p}{,}\PYG{+w}{ }\PYG{n}{x\PYGZus{}r\PYGZus{}solver}\PYG{p}{)}\PYG{p}{)}

\PYG{k}{end }\PYG{k}{program }\PYG{n}{test}

\PYG{k}{module }\PYG{n}{example\PYGZus{}mod}

\PYG{+w}{  }\PYG{k}{use }\PYG{n}{argument\PYGZus{}mod}
\PYG{+w}{  }\PYG{k}{use }\PYG{n}{kernel\PYGZus{}mod}

\PYG{+w}{  }\PYG{k}{implicit }\PYG{k}{none}

\PYG{k}{  }\PYG{k}{type}\PYG{p}{,}\PYG{+w}{ }\PYG{k}{extends}\PYG{p}{(}\PYG{n}{kernel\PYGZus{}type}\PYG{p}{)}\PYG{+w}{ }\PYG{k+kd}{::}\PYG{+w}{ }\PYG{n}{example\PYGZus{}type}
\PYG{+w}{    }\PYG{k}{type}\PYG{p}{(}\PYG{n}{arg\PYGZus{}type}\PYG{p}{)}\PYG{p}{,}\PYG{+w}{ }\PYG{k}{dimension}\PYG{p}{(}\PYG{l+m+mi}{2}\PYG{p}{)}\PYG{+w}{ }\PYG{k+kd}{::}\PYG{+w}{ }\PYG{n}{meta\PYGZus{}args}\PYG{+w}{ }\PYG{o}{=}\PYG{+w}{ }\PYG{p}{(}\PYG{o}{/}\PYG{+w}{       }\PYG{p}{\PYGZam{}}
\PYG{+w}{         }\PYG{n}{arg\PYGZus{}type}\PYG{p}{(}\PYG{n}{gh\PYGZus{}field}\PYG{p}{,}\PYG{+w}{  }\PYG{n}{gh\PYGZus{}real}\PYG{p}{,}\PYG{+w}{ }\PYG{n}{gh\PYGZus{}readwrite}\PYG{p}{,}\PYG{+w}{ }\PYG{n}{w3}\PYG{p}{)}\PYG{p}{,}\PYG{+w}{ }\PYG{p}{\PYGZam{}}
\PYG{+w}{         }\PYG{n}{arg\PYGZus{}type}\PYG{p}{(}\PYG{n}{gh\PYGZus{}scalar}\PYG{p}{,}\PYG{+w}{ }\PYG{n}{gh\PYGZus{}real}\PYG{p}{,}\PYG{+w}{ }\PYG{n}{gh\PYGZus{}read}\PYG{+w}{ }\PYG{p}{)}\PYG{+w}{          }\PYG{p}{\PYGZam{}}
\PYG{+w}{         }\PYG{o}{/}\PYG{p}{)}
\PYG{+w}{     }\PYG{k+kt}{integer}\PYG{+w}{ }\PYG{k+kd}{::}\PYG{+w}{ }\PYG{n}{operates\PYGZus{}on}\PYG{+w}{ }\PYG{o}{=}\PYG{+w}{ }\PYG{n}{cell\PYGZus{}column}
\PYG{+w}{   }\PYG{k}{contains}
\PYG{k}{     }\PYG{k}{procedure}\PYG{p}{,}\PYG{+w}{ }\PYG{k}{nopass}\PYG{+w}{ }\PYG{k+kd}{::}\PYG{+w}{ }\PYG{n}{code}\PYG{+w}{ }\PYG{o}{=}\PYG{o}{\PYGZgt{}}\PYG{+w}{ }\PYG{n}{example\PYGZus{}code}
\PYG{+w}{  }\PYG{k}{end }\PYG{k}{type }\PYG{n}{example\PYGZus{}type}

\PYG{+w}{  }\PYG{k}{private}
\PYG{k}{  }\PYG{k}{public}\PYG{+w}{ }\PYG{k+kd}{::}\PYG{+w}{ }\PYG{n}{example\PYGZus{}code}

\PYG{+w}{  }\PYG{k}{interface }\PYG{n}{example\PYGZus{}code}
\PYG{+w}{    }\PYG{k}{module }\PYG{k}{procedure }\PYG{n}{example\PYGZus{}code\PYGZus{}32}
\PYG{+w}{    }\PYG{k}{module }\PYG{k}{procedure }\PYG{n}{example\PYGZus{}code\PYGZus{}64}
\PYG{+w}{  }\PYG{k}{end }\PYG{k}{interface }\PYG{n}{example\PYGZus{}code}

\PYG{k}{contains}

\PYG{k}{  }\PYG{k}{subroutine }\PYG{n}{example\PYGZus{}code\PYGZus{}32}\PYG{p}{(}\PYG{p}{.}\PYG{p}{.}\PYG{p}{.}\PYG{p}{,}\PYG{+w}{ }\PYG{n}{field1}\PYG{p}{,}\PYG{+w}{ }\PYG{n}{x}\PYG{p}{,}\PYG{+w}{ }\PYG{p}{.}\PYG{p}{.}\PYG{p}{.}\PYG{p}{)}
\PYG{+w}{    }\PYG{k+kt}{real}\PYG{o}{*}\PYG{l+m+mi}{4}\PYG{p}{,}\PYG{+w}{ }\PYG{k}{dimension}\PYG{p}{(}\PYG{p}{.}\PYG{p}{.}\PYG{p}{.}\PYG{p}{)}\PYG{p}{,}\PYG{+w}{ }\PYG{k}{intent}\PYG{p}{(}\PYG{n}{inout}\PYG{p}{)}\PYG{+w}{ }\PYG{k+kd}{::}\PYG{+w}{ }\PYG{n}{field1}
\PYG{+w}{    }\PYG{k+kt}{real}\PYG{o}{*}\PYG{l+m+mi}{4}\PYG{p}{,}\PYG{+w}{ }\PYG{k}{intent}\PYG{p}{(}\PYG{n}{in}\PYG{p}{)}\PYG{+w}{ }\PYG{k+kd}{::}\PYG{+w}{ }\PYG{n}{x}
\PYG{+w}{    }\PYG{k}{print}\PYG{+w}{ }\PYG{o}{*}\PYG{p}{,}\PYG{+w}{ }\PYG{l+s+s2}{\PYGZdq{}32\PYGZhy{}bit example called\PYGZdq{}}
\PYG{+w}{  }\PYG{k}{end }\PYG{k}{subroutine }\PYG{n}{example\PYGZus{}code\PYGZus{}32}

\PYG{+w}{  }\PYG{k}{subroutine }\PYG{n}{example\PYGZus{}code\PYGZus{}64}\PYG{p}{(}\PYG{p}{.}\PYG{p}{.}\PYG{p}{.}\PYG{p}{,}\PYG{+w}{ }\PYG{n}{field1}\PYG{p}{,}\PYG{+w}{ }\PYG{n}{x}\PYG{p}{,}\PYG{+w}{ }\PYG{p}{.}\PYG{p}{.}\PYG{p}{.}\PYG{p}{)}
\PYG{+w}{    }\PYG{k+kt}{real}\PYG{o}{*}\PYG{l+m+mi}{8}\PYG{p}{,}\PYG{+w}{ }\PYG{k}{dimension}\PYG{p}{(}\PYG{p}{.}\PYG{p}{.}\PYG{p}{.}\PYG{p}{)}\PYG{p}{,}\PYG{+w}{ }\PYG{k}{intent}\PYG{p}{(}\PYG{n}{inout}\PYG{p}{)}\PYG{+w}{ }\PYG{k+kd}{::}\PYG{+w}{ }\PYG{n}{field1}
\PYG{+w}{    }\PYG{k+kt}{real}\PYG{o}{*}\PYG{l+m+mi}{8}\PYG{p}{,}\PYG{+w}{ }\PYG{k}{intent}\PYG{p}{(}\PYG{n}{in}\PYG{p}{)}\PYG{+w}{ }\PYG{k+kd}{::}\PYG{+w}{ }\PYG{n}{x}
\PYG{+w}{    }\PYG{k}{print}\PYG{+w}{ }\PYG{o}{*}\PYG{p}{,}\PYG{+w}{ }\PYG{l+s+s2}{\PYGZdq{}64\PYGZhy{}bit example called\PYGZdq{}}
\PYG{+w}{  }\PYG{k}{end }\PYG{k}{subroutine }\PYG{n}{example\PYGZus{}code\PYGZus{}64}

\PYG{k}{end }\PYG{k}{module }\PYG{n}{example\PYGZus{}mod}
\end{sphinxVerbatim}

\sphinxAtStartPar
In order to support mixed precision, PSyclone needs to know the
precision (as specified in the algorithm layer) of any kernel
arguments that are of a type that supports different precisions (e.g.
\sphinxcode{\sphinxupquote{GH\_FIELD}}). The reason for this is that PSyclone needs to be able
to declare data with the correct precision information within the
PSy\sphinxhyphen{}layer to ensure that the correct flavour of kernels are called.

\sphinxAtStartPar
PSyclone must therefore determine this information from the algorithm
layer. The rules for whether PSyclone requires information for
particular LFRic datatypes and what it does with or without this
information are given below:


\subsection{Fields}
\label{\detokenize{dynamo0p3:fields}}\label{\detokenize{dynamo0p3:lfric-mixed-precision-fields}}
\sphinxAtStartPar
PSyclone must be able to determine the datatype of a field from the
algorithm layer declarations. If it is not able to do this, PSyclone
will abort with a message that indicates the problem.

\sphinxAtStartPar
Supported field types, their Fortran datatype and precisions are
outlined in the table below:


\begin{savenotes}\sphinxattablestart
\sphinxthistablewithglobalstyle
\sphinxthistablewithvlinesstyle
\centering
\begin{tabulary}{\linewidth}[t]{|l|l|l|}
\sphinxtoprule
\sphinxstyletheadfamily 
\sphinxAtStartPar
Field Type
&\sphinxstyletheadfamily 
\sphinxAtStartPar
Fortran Datatype
&\sphinxstyletheadfamily 
\sphinxAtStartPar
Precision
\\
\sphinxmidrule
\sphinxtableatstartofbodyhook
\sphinxAtStartPar
\sphinxcode{\sphinxupquote{field\_type}}
&
\sphinxAtStartPar
\sphinxcode{\sphinxupquote{real}}
&
\sphinxAtStartPar
\sphinxcode{\sphinxupquote{r\_def}}
\\
\sphinxhline
\sphinxAtStartPar
\sphinxcode{\sphinxupquote{r\_bl\_field\_type}}
&
\sphinxAtStartPar
\sphinxcode{\sphinxupquote{real}}
&
\sphinxAtStartPar
\sphinxcode{\sphinxupquote{r\_bl}}
\\
\sphinxhline
\sphinxAtStartPar
\sphinxcode{\sphinxupquote{r\_phys\_field\_type}}
&
\sphinxAtStartPar
\sphinxcode{\sphinxupquote{real}}
&
\sphinxAtStartPar
\sphinxcode{\sphinxupquote{r\_phys}}
\\
\sphinxhline
\sphinxAtStartPar
\sphinxcode{\sphinxupquote{r\_solver\_field\_type}}
&
\sphinxAtStartPar
\sphinxcode{\sphinxupquote{real}}
&
\sphinxAtStartPar
\sphinxcode{\sphinxupquote{r\_solver}}
\\
\sphinxhline
\sphinxAtStartPar
\sphinxcode{\sphinxupquote{r\_tran\_field\_type}}
&
\sphinxAtStartPar
\sphinxcode{\sphinxupquote{real}}
&
\sphinxAtStartPar
\sphinxcode{\sphinxupquote{r\_tran}}
\\
\sphinxhline
\sphinxAtStartPar
\sphinxcode{\sphinxupquote{integer\_field\_type}}
&
\sphinxAtStartPar
\sphinxcode{\sphinxupquote{integer}}
&
\sphinxAtStartPar
\sphinxcode{\sphinxupquote{i\_def}}
\\
\sphinxbottomrule
\end{tabulary}
\sphinxtableafterendhook\par
\sphinxattableend\end{savenotes}


\subsection{Field Vectors}
\label{\detokenize{dynamo0p3:field-vectors}}\label{\detokenize{dynamo0p3:lfric-mixed-precision-field-vectors}}
\sphinxAtStartPar
In addition to fields, LFRic supports an abstract vector type for fields,
used in the LFRic solver API. Please note that these structures are
different from the {\hyperref[\detokenize{dynamo0p3:lfric-field-vector}]{\sphinxcrossref{\DUrole{std,std-ref}{field vector}}}} implementation
of field bundles in the PSyclone LFRic API interface.

\sphinxAtStartPar
The LFRic abstract vector type has precision\sphinxhyphen{}specific implementations.
If PSyclone finds such a specifically declared field vector argument in the
algorithm layer, e.g. \sphinxcode{\sphinxupquote{r\_solver\_field\_vector\_type}}, it will assume that
the actual field being referenced is of the same datatype and precision
(see {\hyperref[\detokenize{dynamo0p3:lfric-mixed-precision-fields}]{\sphinxcrossref{\DUrole{std,std-ref}{above}}}} for details).
The correspondence between the available field types and their vector
implementations is given in the table below (note that only
\sphinxcode{\sphinxupquote{real}}\sphinxhyphen{}valued fields have abstract vector implementations for now):


\begin{savenotes}\sphinxattablestart
\sphinxthistablewithglobalstyle
\sphinxthistablewithvlinesstyle
\centering
\begin{tabulary}{\linewidth}[t]{|l|l|}
\sphinxtoprule
\sphinxstyletheadfamily 
\sphinxAtStartPar
Field Type
&\sphinxstyletheadfamily 
\sphinxAtStartPar
Field Vector Type
\\
\sphinxmidrule
\sphinxtableatstartofbodyhook
\sphinxAtStartPar
\sphinxcode{\sphinxupquote{field\_type}}
&
\sphinxAtStartPar
\sphinxcode{\sphinxupquote{field\_vector\_type}}
\\
\sphinxhline
\sphinxAtStartPar
\sphinxcode{\sphinxupquote{r\_bl\_field\_type}}
&
\sphinxAtStartPar
\sphinxcode{\sphinxupquote{r\_bl\_field\_vector\_type}}
\\
\sphinxhline
\sphinxAtStartPar
\sphinxcode{\sphinxupquote{r\_phys\_field\_type}}
&
\sphinxAtStartPar
\sphinxcode{\sphinxupquote{r\_phys\_field\_vector\_type}}
\\
\sphinxhline
\sphinxAtStartPar
\sphinxcode{\sphinxupquote{r\_solver\_field\_type}}
&
\sphinxAtStartPar
\sphinxcode{\sphinxupquote{r\_solver\_field\_vector\_type}}
\\
\sphinxhline
\sphinxAtStartPar
\sphinxcode{\sphinxupquote{r\_tran\_field\_type}}
&
\sphinxAtStartPar
\sphinxcode{\sphinxupquote{r\_tran\_field\_vector\_type}}
\\
\sphinxbottomrule
\end{tabulary}
\sphinxtableafterendhook\par
\sphinxattableend\end{savenotes}

\sphinxAtStartPar
If PSyclone finds an argument that is declared as an
\sphinxcode{\sphinxupquote{abstract\_field\_type}} then it will not know the actual type of the
argument. For instance, the following algorithm layer code will cause
PSyclone to raise an exception:

\begin{sphinxVerbatim}[commandchars=\\\{\}]
\PYG{c}{! ...}
\PYG{k}{class}\PYG{+w}{ }\PYG{p}{(}\PYG{n}{abstract\PYGZus{}vector\PYGZus{}type}\PYG{p}{)}\PYG{p}{,}\PYG{+w}{ }\PYG{k}{intent}\PYG{p}{(}\PYG{n}{inout}\PYG{p}{)}\PYG{+w}{ }\PYG{k+kd}{::}\PYG{+w}{ }\PYG{n}{x}
\PYG{c}{! ...}
\PYG{k}{select }\PYG{k}{type}\PYG{+w}{ }\PYG{p}{(}\PYG{n}{x}\PYG{p}{)}
\PYG{k}{type }\PYG{k}{is}\PYG{+w}{ }\PYG{p}{(}\PYG{n}{field\PYGZus{}vector\PYGZus{}type}\PYG{p}{)}
\PYG{+w}{  }\PYG{k}{call }\PYG{n}{invoke}\PYG{p}{(}\PYG{n}{testkern\PYGZus{}type}\PYG{p}{(}\PYG{n}{x}\PYG{p}{\PYGZpc{}}\PYG{n}{vector}\PYG{p}{(}\PYG{l+m+mi}{1}\PYG{p}{)}\PYG{p}{)}\PYG{p}{)}
\PYG{k}{class }\PYG{n}{default}
\PYG{+w}{  }\PYG{k}{print}\PYG{+w}{ }\PYG{o}{*}\PYG{p}{,}\PYG{l+s+s2}{\PYGZdq{}Error\PYGZdq{}}
\PYG{k}{end }\PYG{k}{select}
\PYG{c}{! ...}
\end{sphinxVerbatim}

\sphinxAtStartPar
The suggested solution to this is to add a pointer variable to the
code that is of the required type. This pointer can then be associated
with the argument and passed into the routine:

\begin{sphinxVerbatim}[commandchars=\\\{\}]
\PYG{c}{! ...}
\PYG{k}{class}\PYG{+w}{ }\PYG{p}{(}\PYG{n}{abstract\PYGZus{}vector\PYGZus{}type}\PYG{p}{)}\PYG{p}{,}\PYG{+w}{ }\PYG{k}{target}\PYG{p}{,}\PYG{+w}{ }\PYG{k}{intent}\PYG{p}{(}\PYG{n}{inout}\PYG{p}{)}\PYG{+w}{ }\PYG{k+kd}{::}\PYG{+w}{ }\PYG{n}{x}
\PYG{k}{type}\PYG{p}{(}\PYG{n}{field\PYGZus{}vector\PYGZus{}type}\PYG{p}{)}\PYG{p}{,}\PYG{+w}{ }\PYG{k}{pointer}\PYG{+w}{ }\PYG{k+kd}{::}\PYG{+w}{ }\PYG{n}{x\PYGZus{}ptr}
\PYG{c}{! ...}
\PYG{k}{select }\PYG{k}{type}\PYG{+w}{ }\PYG{p}{(}\PYG{n}{x}\PYG{p}{)}
\PYG{k}{type }\PYG{k}{is}\PYG{+w}{ }\PYG{p}{(}\PYG{n}{field\PYGZus{}vector\PYGZus{}type}\PYG{p}{)}
\PYG{+w}{  }\PYG{n}{x\PYGZus{}ptr}\PYG{+w}{ }\PYG{o}{=}\PYG{o}{\PYGZgt{}}\PYG{+w}{ }\PYG{n}{x}
\PYG{+w}{  }\PYG{k}{call }\PYG{n}{invoke}\PYG{p}{(}\PYG{n}{testkern\PYGZus{}type}\PYG{p}{(}\PYG{n}{x\PYGZus{}ptr}\PYG{p}{\PYGZpc{}}\PYG{n}{vector}\PYG{p}{(}\PYG{l+m+mi}{1}\PYG{p}{)}\PYG{p}{)}\PYG{p}{)}
\PYG{k}{class }\PYG{n}{default}
\PYG{+w}{  }\PYG{k}{print}\PYG{+w}{ }\PYG{o}{*}\PYG{p}{,}\PYG{l+s+s2}{\PYGZdq{}Error\PYGZdq{}}
\PYG{k}{end }\PYG{k}{select}
\PYG{c}{! ...}
\end{sphinxVerbatim}


\subsection{Scalars}
\label{\detokenize{dynamo0p3:scalars}}\label{\detokenize{dynamo0p3:lfric-mixed-precision-scalars}}
\sphinxAtStartPar
It is not mandatory for PSyclone to be able to determine the datatype
of a scalar from the algorithm layer. This constraint was considered
to be too restrictive as PSyclone currently only examines the
declarations in the same source file as the \sphinxcode{\sphinxupquote{invoke}} when
determining datatype. This means that if scalars are imported from
other modules (as is often the case) then their datatype cannot be
determined.

\sphinxAtStartPar
If the precision information for a scalar is found by PSyclone then
this is used. If the scalar declaration is found and it contains no
precision information then PSyclone will abort with a message that
indicates the problem (since this violates LFRic coding standards). If
no declaration information is found then default precision values are
used, as specified in the PSyclone config file (\sphinxcode{\sphinxupquote{r\_def}} for \sphinxcode{\sphinxupquote{real}},
\sphinxcode{\sphinxupquote{i\_def}} for \sphinxcode{\sphinxupquote{integer}} and \sphinxcode{\sphinxupquote{l\_def}} for \sphinxcode{\sphinxupquote{logical}}).

\sphinxAtStartPar
Supported precisions for scalars are outlined in the table below. If
an unsupported scalar precision is found then PSyclone will abort with
a message that indicates the problem.


\begin{savenotes}\sphinxattablestart
\sphinxthistablewithglobalstyle
\sphinxthistablewithvlinesstyle
\centering
\begin{tabulary}{\linewidth}[t]{|l|l|}
\sphinxtoprule
\sphinxstyletheadfamily 
\sphinxAtStartPar
Fortran Datatype
&\sphinxstyletheadfamily 
\sphinxAtStartPar
Supported Precision
\\
\sphinxmidrule
\sphinxtableatstartofbodyhook
\sphinxAtStartPar
\sphinxcode{\sphinxupquote{real}}
&
\sphinxAtStartPar
\sphinxcode{\sphinxupquote{r\_def}}, \sphinxcode{\sphinxupquote{r\_bl}}, \sphinxcode{\sphinxupquote{r\_phys}},
\sphinxcode{\sphinxupquote{r\_solver}}, \sphinxcode{\sphinxupquote{r\_tran}}
\\
\sphinxhline
\sphinxAtStartPar
\sphinxcode{\sphinxupquote{integer}}
&
\sphinxAtStartPar
\sphinxcode{\sphinxupquote{i\_def}}
\\
\sphinxhline
\sphinxAtStartPar
\sphinxcode{\sphinxupquote{logical}}
&
\sphinxAtStartPar
\sphinxcode{\sphinxupquote{l\_def}}
\\
\sphinxbottomrule
\end{tabulary}
\sphinxtableafterendhook\par
\sphinxattableend\end{savenotes}


\subsection{LMA Operators}
\label{\detokenize{dynamo0p3:lma-operators}}\label{\detokenize{dynamo0p3:lfric-mixed-precision-lma-operators}}
\sphinxAtStartPar
PSyclone must be able to determine the datatype of an LMA operator.
If it is not able to do this, PSyclone will abort with a message that
indicates the problem.

\sphinxAtStartPar
Supported LMA operator types, their Fortran datatype and precisions are
outlined in the table below:


\begin{savenotes}\sphinxattablestart
\sphinxthistablewithglobalstyle
\sphinxthistablewithvlinesstyle
\centering
\begin{tabulary}{\linewidth}[t]{|l|l|l|}
\sphinxtoprule
\sphinxstyletheadfamily 
\sphinxAtStartPar
Operator Type
&\sphinxstyletheadfamily 
\sphinxAtStartPar
Fortran Datatype
&\sphinxstyletheadfamily 
\sphinxAtStartPar
Precision
\\
\sphinxmidrule
\sphinxtableatstartofbodyhook
\sphinxAtStartPar
\sphinxcode{\sphinxupquote{operator\_type}}
&
\sphinxAtStartPar
\sphinxcode{\sphinxupquote{real}}
&
\sphinxAtStartPar
\sphinxcode{\sphinxupquote{r\_def}}
\\
\sphinxhline
\sphinxAtStartPar
\sphinxcode{\sphinxupquote{r\_solver\_operator\_type}}
&
\sphinxAtStartPar
\sphinxcode{\sphinxupquote{real}}
&
\sphinxAtStartPar
\sphinxcode{\sphinxupquote{r\_solver}}
\\
\sphinxhline
\sphinxAtStartPar
\sphinxcode{\sphinxupquote{r\_tran\_operator\_type}}
&
\sphinxAtStartPar
\sphinxcode{\sphinxupquote{real}}
&
\sphinxAtStartPar
\sphinxcode{\sphinxupquote{r\_tran}}
\\
\sphinxbottomrule
\end{tabulary}
\sphinxtableafterendhook\par
\sphinxattableend\end{savenotes}


\subsection{Column\sphinxhyphen{}wise Operators}
\label{\detokenize{dynamo0p3:column-wise-operators}}\label{\detokenize{dynamo0p3:lfric-mixed-precision-cma-operators}}
\sphinxAtStartPar
It is not mandatory for PSyclone to be able to determine the datatype
of a column\sphinxhyphen{}wise (CMA) operator. The reason for this is that only one
datatype is supported, a \sphinxcode{\sphinxupquote{columnwise\_operator\_type}} which contains
\sphinxcode{\sphinxupquote{real}}\sphinxhyphen{}valued data with precision \sphinxcode{\sphinxupquote{r\_solver}}. PSyclone can therefore
simply add this datatype in the PSy\sphinxhyphen{}layer. However, if the datatype
information is found in the algorithm layer and it is not of the
expected type then PSyclone will abort with a message that indicates
the problem.


\subsection{Consistency}
\label{\detokenize{dynamo0p3:consistency}}\label{\detokenize{dynamo0p3:lfric-mixed-precision-consistency}}
\sphinxAtStartPar
If PSyclone is able to determine the datatype of an LFRic datatype
then PSyclone also checks that this datatype is consistent with the
associated kernel metadata. If it is not consistent then PSyclone will
abort with a message that indicates the problem.


\section{PSy\sphinxhyphen{}layer}
\label{\detokenize{dynamo0p3:psy-layer}}\label{\detokenize{dynamo0p3:lfric-psy}}
\sphinxAtStartPar
The general details of the PSy\sphinxhyphen{}layer are explained in the
{\hyperref[\detokenize{introduction_to_psykal:psy-layer}]{\sphinxcrossref{\DUrole{std,std-ref}{PSy layer}}}} section. This section describes any LFRic\sphinxhyphen{}specific
issues.


\subsection{Module name}
\label{\detokenize{dynamo0p3:module-name}}
\sphinxAtStartPar
The PSy\sphinxhyphen{}layer code is contained within a Fortran module. The name of
the module is determined from the algorithm\sphinxhyphen{}layer name with “\_psy”
appended. The algorithm\sphinxhyphen{}layer name is the algorithm’s module name if it
is a module, its subroutine name if it is a subroutine that is not
within a module, or the program name if it is a program.

\sphinxAtStartPar
So, for example, if the algorithm code is contained within a module
called “fred” then the PSy\sphinxhyphen{}layer module name will be “fred\_psy”.


\subsubsection{Argument Intents}
\label{\detokenize{dynamo0p3:argument-intents}}\label{\detokenize{dynamo0p3:lfric-psy-arg-intents}}
\sphinxAtStartPar
LFRic {\hyperref[\detokenize{dynamo0p3:lfric-field}]{\sphinxcrossref{\DUrole{std,std-ref}{fields}}}}, {\hyperref[\detokenize{dynamo0p3:lfric-field-vector}]{\sphinxcrossref{\DUrole{std,std-ref}{field vectors}}}}, {\hyperref[\detokenize{dynamo0p3:lfric-operator}]{\sphinxcrossref{\DUrole{std,std-ref}{operators}}}} and
{\hyperref[\detokenize{dynamo0p3:lfric-cma-operator}]{\sphinxcrossref{\DUrole{std,std-ref}{column\sphinxhyphen{}wise operators}}}} are objects that
contain pointers to data rather than data. The data are accessed by proxies
of these objects and modified in {\hyperref[\detokenize{dynamo0p3:lfric-kernel}]{\sphinxcrossref{\DUrole{std,std-ref}{kernels}}}}.
As the objects themselves are not modified in the PSy layer, their Fortran
intents there are always \sphinxcode{\sphinxupquote{intent(in)}}.

\sphinxAtStartPar
The Fortran intent of {\hyperref[\detokenize{dynamo0p3:lfric-scalar}]{\sphinxcrossref{\DUrole{std,std-ref}{scalars}}}} is still defined
by their {\hyperref[\detokenize{dynamo0p3:lfric-kernel-valid-access}]{\sphinxcrossref{\DUrole{std,std-ref}{access metadata}}}} as they are
actual data. This means \sphinxcode{\sphinxupquote{intent(in)}} for \sphinxcode{\sphinxupquote{GH\_READ}} and \sphinxcode{\sphinxupquote{intent(out)}}
for \sphinxcode{\sphinxupquote{GH\_SUM}} (more details in {\hyperref[\detokenize{dynamo0p3:lfric-api-meta-args}]{\sphinxcrossref{\DUrole{std,std-ref}{meta\_args}}}}
section below).

\sphinxAtStartPar
The intent of other data structures is mandated by the relevant
LFRic API rules described in sections below.


\section{Kernel}
\label{\detokenize{dynamo0p3:kernel}}\label{\detokenize{dynamo0p3:lfric-kernel}}
\sphinxAtStartPar
The general requirements for the structure of a Kernel are explained
in the {\hyperref[\detokenize{introduction_to_psykal:kernel-layer}]{\sphinxcrossref{\DUrole{std,std-ref}{Kernel layer}}}} section. In the LFRic API there are six
different Kernel types; general purpose, CMA, inter\sphinxhyphen{}grid, domain, dof and
{\hyperref[\detokenize{dynamo0p3:lfric-built-ins}]{\sphinxcrossref{\DUrole{std,std-ref}{Built\sphinxhyphen{}ins}}}}. In the case of built\sphinxhyphen{}ins, PSyclone generates
the source of the kernels.  This section explains the rules for the
other five user\sphinxhyphen{}supplied kernel types and then goes on to describe
their metadata and subroutine arguments.

\sphinxAtStartPar
Domain kernels are distinct from the other four user\sphinxhyphen{}supplied kernel
types because they must be passed data for the whole domain rather
than a single cell\sphinxhyphen{}column or dof. This permits the use of kernels that have
not been written to conform to the single\sphinxhyphen{}column/dof approach which
simplifies the integration with existing code. Obviously, any
parallelisation in the ‘domain’ kernel must be consistent with that
in the rest of the application. The motivation for such kernels in
LFRic is that they allow existing, “i\sphinxhyphen{}first” physics code to be called from
the PSy layer. Since those routines currently contain their own,
i\sphinxhyphen{}first looping structure (and associated OpenMP parallelisation), the
most efficient way to use them is to avoid enclosing them within a
loop in the PSy layer. This is a temporary measure and these
kernels will ultimately be replaced once the LFRic infrastructure has
support for i\sphinxhyphen{}first kernels
(\sphinxurl{https://code.metoffice.gov.uk/trac/lfric/ticket/2154}). At that
point the looping (and associated parallelisation) will be put
back into the PSy layer.


\subsection{Rules for all User\sphinxhyphen{}Supplied Kernels that Operate on Cell\sphinxhyphen{}Columns}
\label{\detokenize{dynamo0p3:rules-for-all-user-supplied-kernels-that-operate-on-cell-columns}}\label{\detokenize{dynamo0p3:lfric-user-kernel-rules}}
\sphinxAtStartPar
In the following, ‘operator’ refers to both LMA and CMA operator
types.
\begin{enumerate}
\sphinxsetlistlabels{\arabic}{enumi}{enumii}{}{)}%
\item {} 
\sphinxAtStartPar
A Kernel must have at least one argument that is a field, field
vector, or operator. This rule reflects the fact that a Kernel
operates on some subset of the whole domain (e.g. a cell\sphinxhyphen{}column)
and is therefore designed to be called from within a loop that
iterates over those subsets of the domain.

\item {} 
\sphinxAtStartPar
The continuity of the iteration space of the Kernel is determined
from the function space of the modified argument (see Section
{\hyperref[\detokenize{dynamo0p3:lfric-function-space}]{\sphinxcrossref{\DUrole{std,std-ref}{Supported Function Spaces}}}} below).
If more than one argument is modified then the iteration space is taken
to be the largest required by any of those arguments. E.g. if a Kernel
writes to two fields, the first on \sphinxcode{\sphinxupquote{W3}} (discontinuous) and the
second on \sphinxcode{\sphinxupquote{W1}} (continuous), then the iteration space of that Kernel
will be determined by the field on the continuous space.

\item {} 
\sphinxAtStartPar
If any of the modified arguments are declared with the generic
function space metadata (e.g. \sphinxcode{\sphinxupquote{ANY\_SPACE\_\textless{}n\textgreater{}}}, see
{\hyperref[\detokenize{dynamo0p3:lfric-function-space}]{\sphinxcrossref{\DUrole{std,std-ref}{Supported Function Spaces}}}})
and their actual space cannot be determined statically then the
iteration space is assumed to be
\begin{enumerate}
\sphinxsetlistlabels{\arabic}{enumii}{enumiii}{}{)}%
\item {} 
\sphinxAtStartPar
discontinuous for \sphinxcode{\sphinxupquote{ANY\_DISCONTINUOUS\_SPACE\_\textless{}n\textgreater{}}};

\item {} 
\sphinxAtStartPar
continuous for \sphinxcode{\sphinxupquote{ANY\_SPACE\_\textless{}n\textgreater{}}} and \sphinxcode{\sphinxupquote{ANY\_W2}}.  This assumption
is always safe but leads to additional computation if the quantities
being updated are actually on discontinuous function spaces.

\end{enumerate}

\item {} 
\sphinxAtStartPar
Operators do not have halo operations operating on them as they
are either cell\sphinxhyphen{} (LMA) or column\sphinxhyphen{}based (CMA) and therefore act
like discontinuous fields.

\item {} 
\sphinxAtStartPar
Any Kernel that writes to an operator will have its iteration
space expanded such that valid values for the operator are
computed in the level\sphinxhyphen{}1 halo.

\item {} 
\sphinxAtStartPar
Any Kernel that reads from an operator must not access halos
beyond level 1. In this case PSyclone will check that the Kernel
does not require values beyond the level\sphinxhyphen{}1 halo. If it does then
PSyclone will abort.

\item {} 
\sphinxAtStartPar
Any Kernel that takes an operator argument must not also take
an \sphinxcode{\sphinxupquote{integer}}\sphinxhyphen{}valued field as an argument.

\end{enumerate}


\subsection{Rules specific to General\sphinxhyphen{}Purpose Kernels without CMA Operators}
\label{\detokenize{dynamo0p3:rules-specific-to-general-purpose-kernels-without-cma-operators}}\label{\detokenize{dynamo0p3:lfric-no-cma-mdata-rules}}\begin{enumerate}
\sphinxsetlistlabels{\arabic}{enumi}{enumii}{}{)}%
\item {} 
\sphinxAtStartPar
General\sphinxhyphen{}purpose kernels with \sphinxcode{\sphinxupquote{operates\_on = CELL\_COLUMN}} accept
arguments of any of the following types: field, field vector, LMA
operator, scalar (\sphinxcode{\sphinxupquote{real}}, \sphinxcode{\sphinxupquote{integer}} or \sphinxcode{\sphinxupquote{logical}}).

\item {} 
\sphinxAtStartPar
A Kernel is permitted to write to more than one
quantity (field or operator) and these quantities may be on the
same or different function spaces.

\item {} 
\sphinxAtStartPar
A Kernel may not write to a scalar argument. (Only
{\hyperref[\detokenize{dynamo0p3:lfric-built-ins}]{\sphinxcrossref{\DUrole{std,std-ref}{built\sphinxhyphen{}ins}}}} are permitted to do this.) Any
scalar arguments must therefore be declared in the metadata as
\sphinxcode{\sphinxupquote{GH\_READ}} \sphinxhyphen{} see {\hyperref[\detokenize{dynamo0p3:lfric-kernel-valid-access}]{\sphinxcrossref{\DUrole{std,std-ref}{below}}}}.

\end{enumerate}


\subsection{Rules for Kernels that work with CMA Operators}
\label{\detokenize{dynamo0p3:rules-for-kernels-that-work-with-cma-operators}}\label{\detokenize{dynamo0p3:lfric-cma-mdata-rules}}
\sphinxAtStartPar
The LFRic API has support for kernels that assemble, apply (or
inverse\sphinxhyphen{}apply) column\sphinxhyphen{}wise/Column Matrix Assembly (CMA) operators.
Such operators may also be used by matrix\sphinxhyphen{}matrix kernels. There are
thus three types of CMA\sphinxhyphen{}related kernels.  Since, by definition, CMA
operators only act on data within a column, they have no horizontal
dependencies. Therefore, kernels that write to them may be
parallelised without colouring.

\sphinxAtStartPar
All three CMA\sphinxhyphen{}related kernel types must obey the following rules:
\begin{enumerate}
\sphinxsetlistlabels{\arabic}{enumi}{enumii}{}{)}%
\item {} 
\sphinxAtStartPar
Since a CMA operator only acts within a single column of data,
stencil operations are not permitted.

\item {} 
\sphinxAtStartPar
No vector quantities (e.g. \sphinxcode{\sphinxupquote{GH\_FIELD*3}} \sphinxhyphen{} see below) are
permitted as arguments.

\item {} 
\sphinxAtStartPar
The kernel must operate on cell\sphinxhyphen{}columns.

\end{enumerate}

\sphinxAtStartPar
There are then additional rules specific to each of the three
CMA kernel types. These are described below.


\subsubsection{Assembly}
\label{\detokenize{dynamo0p3:assembly}}
\sphinxAtStartPar
CMA operators are themselves constructed from Local\sphinxhyphen{}Matrix\sphinxhyphen{}Assembly
(LMA) operators. Therefore, any kernel which assembles a CMA
operator must obey the following rules:
\begin{enumerate}
\sphinxsetlistlabels{\arabic}{enumi}{enumii}{}{)}%
\item {} 
\sphinxAtStartPar
Have one or more LMA operators as read\sphinxhyphen{}only arguments.

\item {} 
\sphinxAtStartPar
Have exactly one CMA operator argument which must have write access.

\item {} 
\sphinxAtStartPar
Other types of argument (e.g. scalars or fields) are permitted but
must be read\sphinxhyphen{}only.

\end{enumerate}


\subsubsection{Application and Inverse Application}
\label{\detokenize{dynamo0p3:application-and-inverse-application}}
\sphinxAtStartPar
Column\sphinxhyphen{}wise operators can only be applied to fields. CMA\sphinxhyphen{}Application
kernels must therefore:
\begin{enumerate}
\sphinxsetlistlabels{\arabic}{enumi}{enumii}{}{)}%
\item {} 
\sphinxAtStartPar
Have a single CMA operator as a read\sphinxhyphen{}only argument.

\item {} 
\sphinxAtStartPar
Have exactly two field arguments, one read\sphinxhyphen{}only and one that is written to.

\item {} 
\sphinxAtStartPar
The function spaces of the read and written fields must match the
from and to spaces, respectively, of the supplied CMA operator.

\end{enumerate}


\subsubsection{Matrix\sphinxhyphen{}Matrix}
\label{\detokenize{dynamo0p3:matrix-matrix}}
\sphinxAtStartPar
A kernel that has just column\sphinxhyphen{}wise operators as arguments and zero or
more read\sphinxhyphen{}only scalars is identified as performing a matrix\sphinxhyphen{}matrix
operation. In this case:
\begin{enumerate}
\sphinxsetlistlabels{\arabic}{enumi}{enumii}{}{)}%
\item {} 
\sphinxAtStartPar
Arguments must be CMA operators and, optionally, one or more scalars.

\item {} 
\sphinxAtStartPar
Exactly one of the CMA arguments must be written to while all other
arguments must be read\sphinxhyphen{}only.

\end{enumerate}


\subsection{Rules for Inter\sphinxhyphen{}Grid Kernels}
\label{\detokenize{dynamo0p3:rules-for-inter-grid-kernels}}\begin{enumerate}
\sphinxsetlistlabels{\arabic}{enumi}{enumii}{}{)}%
\item {} 
\sphinxAtStartPar
An inter\sphinxhyphen{}grid kernel is identified by the presence of a field or
field\sphinxhyphen{}vector argument with the optional \sphinxcode{\sphinxupquote{mesh\_arg}} metadata element (see
{\hyperref[\detokenize{dynamo0p3:lfric-intergrid-mdata}]{\sphinxcrossref{\DUrole{std,std-ref}{Inter\sphinxhyphen{}Grid Metadata}}}}).

\item {} 
\sphinxAtStartPar
An invoke that contains one or more inter\sphinxhyphen{}grid kernels must not contain
any other kernel types. (This restriction is an implementation decision
and could be lifted in future if there is a need.)

\item {} 
\sphinxAtStartPar
An inter\sphinxhyphen{}grid kernel is only permitted to have field or field\sphinxhyphen{}vector
arguments.

\item {} 
\sphinxAtStartPar
All inter\sphinxhyphen{}grid kernel arguments must have the \sphinxcode{\sphinxupquote{mesh\_arg}} metadata entry.

\item {} 
\sphinxAtStartPar
An inter\sphinxhyphen{}grid kernel (and metadata) must have at least one field on
each of the fine and coarse meshes. Specifying all fields as coarse or
fine is forbidden.

\item {} 
\sphinxAtStartPar
Fields on different meshes must always live on different function spaces.

\item {} 
\sphinxAtStartPar
All fields on a given mesh must be on the same function space.

\item {} 
\sphinxAtStartPar
An inter\sphinxhyphen{}grid kernel must operate on cell\sphinxhyphen{}columns.

\end{enumerate}

\sphinxAtStartPar
A consequence of Rules 5\sphinxhyphen{}7 is that an inter\sphinxhyphen{}grid kernel will
only involve two function spaces.


\subsection{Rules for User\sphinxhyphen{}Supplied Kernels that Operate on the Domain}
\label{\detokenize{dynamo0p3:rules-for-user-supplied-kernels-that-operate-on-the-domain}}
\sphinxAtStartPar
The rules for kernels that have \sphinxcode{\sphinxupquote{operates\_on = DOMAIN}} are a subset
of {\hyperref[\detokenize{dynamo0p3:lfric-no-cma-mdata-rules}]{\sphinxcrossref{\DUrole{std,std-ref}{those}}}} for kernels that operate
on a \sphinxcode{\sphinxupquote{CELL\_COLUMN}} without CMA Operators. Specifically:
\begin{enumerate}
\sphinxsetlistlabels{\arabic}{enumi}{enumii}{}{)}%
\item {} 
\sphinxAtStartPar
Only {\hyperref[\detokenize{dynamo0p3:lfric-scalar}]{\sphinxcrossref{\DUrole{std,std-ref}{scalar}}}}, {\hyperref[\detokenize{dynamo0p3:lfric-field}]{\sphinxcrossref{\DUrole{std,std-ref}{field}}}} and
{\hyperref[\detokenize{dynamo0p3:lfric-field-vector}]{\sphinxcrossref{\DUrole{std,std-ref}{field vector}}}} arguments are permitted.

\item {} 
\sphinxAtStartPar
All fields must be on discontinuous function spaces.

\item {} 
\sphinxAtStartPar
Stencil accesses are not permitted.

\end{enumerate}


\subsection{Rules for all User\sphinxhyphen{}Supplied Kernels that Operate on DoFs (DoF Kernels)}
\label{\detokenize{dynamo0p3:rules-for-all-user-supplied-kernels-that-operate-on-dofs-dof-kernels}}\label{\detokenize{dynamo0p3:lfric-dof-kernel-rules}}
\sphinxAtStartPar
Kernels that have \sphinxcode{\sphinxupquote{operates\_on = DOF}} and
{\hyperref[\detokenize{dynamo0p3:lfric-built-ins}]{\sphinxcrossref{\DUrole{std,std-ref}{LFRic Built\sphinxhyphen{}ins}}}} overlap significantly in their
scope, and the conventions that DoF Kernels must follow are influenced
by those for built\sphinxhyphen{}ins as a result. This includes {\hyperref[\detokenize{dynamo0p3:lfric-api-built-ins-metadata}]{\sphinxcrossref{\DUrole{std,std-ref}{metadata arguments}}}} and {\hyperref[\detokenize{dynamo0p3:lfric-built-ins-dtype-access}]{\sphinxcrossref{\DUrole{std,std-ref}{valid data types
and access modes}}}}. Naming conventions for DoF
Kernels should follow those for General\sphinxhyphen{}Purpose Kernels.

\sphinxAtStartPar
The list of rules for DoF Kernels is as follows:
\begin{enumerate}
\sphinxsetlistlabels{\arabic}{enumi}{enumii}{}{)}%
\item {} 
\sphinxAtStartPar
A DoF Kernel must have at least one argument that is a field. This rule
reflects that a Kernel operates on some subset of the whole domain
and is therefore designed to be called from within a loop that iterates
over those subsets of the domain. Fields (as opposed to e.g. operators)
are accepted for DoF Kernels because only they have a single value at
each DoF. Field vectors can be represented in a DoF kernel as a
collection of field arguments, each one corresponding to an index in the
field vector.

\item {} 
\sphinxAtStartPar
All Kernel arguments must be either fields or scalars (\sphinxtitleref{real\sphinxhyphen{}} and/or
\sphinxtitleref{integer}\sphinxhyphen{}valued). DoF Kernels cannot accept operators.

\item {} 
\sphinxAtStartPar
All field arguments to a given DoF Kernel must be on the same function
space so they have the same number of DoFs.

\item {} 
\sphinxAtStartPar
They must have at least one modified (i.e. written to) field argument. Unlike
built\sphinxhyphen{}ins, this is not limited and more than one modified argument is
allowed.

\item {} 
\sphinxAtStartPar
A Kernel may not write to a scalar argument. (Only built\sphinxhyphen{}ins are permitted
to do this.) Any scalar arguments must therefore be declared in the metadata
as \sphinxtitleref{GH\_READ} \sphinxhyphen{} see {\hyperref[\detokenize{dynamo0p3:lfric-kernel-valid-access}]{\sphinxcrossref{\DUrole{std,std-ref}{below}}}}

\item {} 
\sphinxAtStartPar
Kernels must be written to operate on a single DoF, such that field values
at the same dof location/index can be provided to the Kernel within a loop
over the DoFs of the function space of the field that is being updated.

\end{enumerate}


\subsection{Metadata}
\label{\detokenize{dynamo0p3:metadata}}\label{\detokenize{dynamo0p3:lfric-api-kernel-metadata}}
\sphinxAtStartPar
The code below outlines the elements of the LFRic API Kernel
metadata, 1) ‘{\hyperref[\detokenize{dynamo0p3:meta-args}]{\sphinxcrossref{meta\_args}}}’, 2) ‘{\hyperref[\detokenize{dynamo0p3:meta-funcs}]{\sphinxcrossref{meta\_funcs}}}’, 3) ‘{\hyperref[\detokenize{dynamo0p3:meta-reference-element}]{\sphinxcrossref{meta\_reference\_element}}}’,
4) ‘{\hyperref[\detokenize{dynamo0p3:meta-mesh}]{\sphinxcrossref{meta\_mesh}}}’, 5) ‘gh\_shape’ ({\hyperref[\detokenize{dynamo0p3:gh-shape-and-gh-evaluator-targets}]{\sphinxcrossref{gh\_shape and gh\_evaluator\_targets}}}),
6) ‘{\hyperref[\detokenize{dynamo0p3:operates-on}]{\sphinxcrossref{operates\_on}}}’ and 7) ‘{\hyperref[\detokenize{dynamo0p3:procedure}]{\sphinxcrossref{procedure}}}’:

\begin{sphinxVerbatim}[commandchars=\\\{\}]
\PYG{k}{type}\PYG{p}{,}\PYG{+w}{ }\PYG{k}{public}\PYG{p}{,}\PYG{+w}{ }\PYG{k}{extends}\PYG{p}{(}\PYG{n}{kernel\PYGZus{}type}\PYG{p}{)}\PYG{+w}{ }\PYG{k+kd}{::}\PYG{+w}{ }\PYG{n}{my\PYGZus{}kernel\PYGZus{}type}
\PYG{+w}{  }\PYG{k}{type}\PYG{p}{(}\PYG{n}{arg\PYGZus{}type}\PYG{p}{)}\PYG{+w}{ }\PYG{k+kd}{::}\PYG{+w}{ }\PYG{n}{meta\PYGZus{}args}\PYG{p}{(}\PYG{p}{.}\PYG{p}{.}\PYG{p}{.}\PYG{p}{)}\PYG{+w}{ }\PYG{o}{=}\PYG{+w}{ }\PYG{p}{(}\PYG{o}{/}\PYG{+w}{ }\PYG{p}{.}\PYG{p}{.}\PYG{p}{.}\PYG{+w}{ }\PYG{o}{/}\PYG{p}{)}
\PYG{+w}{  }\PYG{k}{type}\PYG{p}{(}\PYG{n}{func\PYGZus{}type}\PYG{p}{)}\PYG{+w}{ }\PYG{k+kd}{::}\PYG{+w}{ }\PYG{n}{meta\PYGZus{}funcs}\PYG{p}{(}\PYG{p}{.}\PYG{p}{.}\PYG{p}{.}\PYG{p}{)}\PYG{+w}{ }\PYG{o}{=}\PYG{+w}{ }\PYG{p}{(}\PYG{o}{/}\PYG{+w}{ }\PYG{p}{.}\PYG{p}{.}\PYG{p}{.}\PYG{+w}{ }\PYG{o}{/}\PYG{p}{)}
\PYG{+w}{  }\PYG{k}{type}\PYG{p}{(}\PYG{n}{reference\PYGZus{}element\PYGZus{}data\PYGZus{}type}\PYG{p}{)}\PYG{+w}{ }\PYG{k+kd}{::}\PYG{+w}{ }\PYG{n}{meta\PYGZus{}reference\PYGZus{}element}\PYG{p}{(}\PYG{p}{.}\PYG{p}{.}\PYG{p}{.}\PYG{p}{)}\PYG{+w}{ }\PYG{o}{=}\PYG{+w}{ }\PYG{p}{(}\PYG{o}{/}\PYG{+w}{ }\PYG{p}{.}\PYG{p}{.}\PYG{p}{.}\PYG{+w}{ }\PYG{o}{/}\PYG{p}{)}
\PYG{+w}{  }\PYG{k}{type}\PYG{p}{(}\PYG{n}{mesh\PYGZus{}data\PYGZus{}type}\PYG{p}{)}\PYG{+w}{ }\PYG{k+kd}{::}\PYG{+w}{ }\PYG{n}{meta\PYGZus{}mesh}\PYG{p}{(}\PYG{p}{.}\PYG{p}{.}\PYG{p}{.}\PYG{p}{)}\PYG{+w}{ }\PYG{o}{=}\PYG{+w}{ }\PYG{p}{(}\PYG{o}{/}\PYG{+w}{ }\PYG{p}{.}\PYG{p}{.}\PYG{p}{.}\PYG{+w}{ }\PYG{o}{/}\PYG{p}{)}
\PYG{+w}{  }\PYG{k+kt}{integer}\PYG{+w}{ }\PYG{k+kd}{::}\PYG{+w}{ }\PYG{n}{gh\PYGZus{}shape}\PYG{+w}{ }\PYG{o}{=}\PYG{+w}{ }\PYG{n}{gh\PYGZus{}quadrature\PYGZus{}XYoZ}
\PYG{+w}{  }\PYG{k+kt}{integer}\PYG{+w}{ }\PYG{k+kd}{::}\PYG{+w}{ }\PYG{n}{operates\PYGZus{}on}\PYG{+w}{ }\PYG{o}{=}\PYG{+w}{ }\PYG{n}{cell\PYGZus{}column}
\PYG{k}{contains}
\PYG{k}{  }\PYG{k}{procedure}\PYG{p}{,}\PYG{+w}{ }\PYG{k}{nopass}\PYG{+w}{ }\PYG{k+kd}{::}\PYG{+w}{ }\PYG{n}{my\PYGZus{}kernel\PYGZus{}code}
\PYG{k}{end }\PYG{k}{type}
\end{sphinxVerbatim}

\sphinxAtStartPar
These various metadata elements are discussed in order in the following
sections.


\subsubsection{meta\_args}
\label{\detokenize{dynamo0p3:meta-args}}\label{\detokenize{dynamo0p3:lfric-api-meta-args}}
\sphinxAtStartPar
The \sphinxcode{\sphinxupquote{meta\_args}} array specifies information about data that the
kernel code expects to be passed to it via its argument list. There is
one entry in the \sphinxcode{\sphinxupquote{meta\_args}} array for each \sphinxstylestrong{scalar}, \sphinxstylestrong{field},
or \sphinxstylestrong{operator} passed into the Kernel and the order that these occur
in the \sphinxcode{\sphinxupquote{meta\_args}} array must be the same as they are expected in
the kernel code argument list. The entry must be of \sphinxcode{\sphinxupquote{arg\_type}} which
itself contains metadata about the associated argument. The size of
the \sphinxcode{\sphinxupquote{meta\_args}} array must correspond to the number of \sphinxstylestrong{scalars},
\sphinxstylestrong{fields} and \sphinxstylestrong{operators} passed into the Kernel.

\begin{sphinxadmonition}{note}{Note:}
\sphinxAtStartPar
It makes no sense for a Kernel to have only \sphinxstylestrong{scalar} arguments
(because the PSy layer will call a Kernel for each point in the
spatial domain) and PSyclone will reject such Kernels.
\end{sphinxadmonition}

\sphinxAtStartPar
For example, if there are a total of 2 \sphinxstylestrong{scalar} / \sphinxstylestrong{field} /
\sphinxstylestrong{operator} entities being passed to the Kernel then the \sphinxcode{\sphinxupquote{meta\_args}}
array will be of size 2 and there will be two \sphinxcode{\sphinxupquote{arg\_type}} entries:

\begin{sphinxVerbatim}[commandchars=\\\{\}]
\PYG{k}{type}\PYG{p}{(}\PYG{n}{arg\PYGZus{}type}\PYG{p}{)}\PYG{+w}{ }\PYG{k+kd}{::}\PYG{+w}{ }\PYG{n}{meta\PYGZus{}args}\PYG{p}{(}\PYG{l+m+mi}{2}\PYG{p}{)}\PYG{+w}{ }\PYG{o}{=}\PYG{+w}{ }\PYG{p}{(}\PYG{o}{/}\PYG{+w}{                                  }\PYG{p}{\PYGZam{}}
\PYG{+w}{     }\PYG{n}{arg\PYGZus{}type}\PYG{p}{(}\PYG{+w}{ }\PYG{p}{.}\PYG{p}{.}\PYG{p}{.}\PYG{+w}{ }\PYG{p}{)}\PYG{p}{,}\PYG{+w}{                                                }\PYG{p}{\PYGZam{}}
\PYG{+w}{     }\PYG{n}{arg\PYGZus{}type}\PYG{p}{(}\PYG{+w}{ }\PYG{p}{.}\PYG{p}{.}\PYG{p}{.}\PYG{+w}{ }\PYG{p}{)}\PYG{+w}{                                                 }\PYG{p}{\PYGZam{}}
\PYG{+w}{     }\PYG{o}{/}\PYG{p}{)}
\end{sphinxVerbatim}

\sphinxAtStartPar
Argument metadata (information contained within the brackets of an
\sphinxcode{\sphinxupquote{arg\_type}} entry), describes either a \sphinxstylestrong{scalar}, a \sphinxstylestrong{field} or an
\sphinxstylestrong{operator} (either LMA or CMA).

\sphinxAtStartPar
The first argument\sphinxhyphen{}metadata entry describes whether the data that is
being passed is for a scalar (\sphinxcode{\sphinxupquote{GH\_SCALAR}}), a field (\sphinxcode{\sphinxupquote{GH\_FIELD}}) or
an operator (either \sphinxcode{\sphinxupquote{GH\_OPERATOR}} for LMA or \sphinxcode{\sphinxupquote{GH\_COLUMNWISE\_OPERATOR}}
for CMA). This information is mandatory.

\sphinxAtStartPar
Additionally, argument metadata can be used to describe a vector of
fields (see the {\hyperref[\detokenize{dynamo0p3:lfric-field-vector}]{\sphinxcrossref{\DUrole{std,std-ref}{Field Vector}}}} section for more
details).

\sphinxAtStartPar
As an example, the following \sphinxcode{\sphinxupquote{meta\_args}} metadata describes 4
entries, the first is a scalar, the next two are fields and the
fourth is an operator. The third entry is a field vector of size 3.

\begin{sphinxVerbatim}[commandchars=\\\{\}]
\PYG{k}{type}\PYG{p}{(}\PYG{n}{arg\PYGZus{}type}\PYG{p}{)}\PYG{+w}{ }\PYG{k+kd}{::}\PYG{+w}{ }\PYG{n}{meta\PYGZus{}args}\PYG{p}{(}\PYG{l+m+mi}{4}\PYG{p}{)}\PYG{+w}{ }\PYG{o}{=}\PYG{+w}{ }\PYG{p}{(}\PYG{o}{/}\PYG{+w}{                                  }\PYG{p}{\PYGZam{}}
\PYG{+w}{     }\PYG{n}{arg\PYGZus{}type}\PYG{p}{(}\PYG{n}{GH\PYGZus{}SCALAR}\PYG{p}{,}\PYG{+w}{ }\PYG{n}{GH\PYGZus{}REAL}\PYG{p}{,}\PYG{+w}{ }\PYG{p}{.}\PYG{p}{.}\PYG{p}{.}\PYG{p}{)}\PYG{p}{,}\PYG{+w}{                              }\PYG{p}{\PYGZam{}}
\PYG{+w}{     }\PYG{n}{arg\PYGZus{}type}\PYG{p}{(}\PYG{n}{GH\PYGZus{}FIELD}\PYG{p}{,}\PYG{+w}{ }\PYG{n}{GH\PYGZus{}INTEGER}\PYG{p}{,}\PYG{+w}{ }\PYG{p}{.}\PYG{p}{.}\PYG{p}{.}\PYG{+w}{ }\PYG{p}{)}\PYG{p}{,}\PYG{+w}{                           }\PYG{p}{\PYGZam{}}
\PYG{+w}{     }\PYG{n}{arg\PYGZus{}type}\PYG{p}{(}\PYG{n}{GH\PYGZus{}FIELD}\PYG{o}{*}\PYG{l+m+mi}{3}\PYG{p}{,}\PYG{+w}{ }\PYG{n}{GH\PYGZus{}REAL}\PYG{p}{,}\PYG{+w}{ }\PYG{p}{.}\PYG{p}{.}\PYG{p}{.}\PYG{+w}{ }\PYG{p}{)}\PYG{p}{,}\PYG{+w}{                            }\PYG{p}{\PYGZam{}}
\PYG{+w}{     }\PYG{n}{arg\PYGZus{}type}\PYG{p}{(}\PYG{n}{GH\PYGZus{}OPERATOR}\PYG{p}{,}\PYG{+w}{ }\PYG{n}{GH\PYGZus{}REAL}\PYG{p}{,}\PYG{+w}{ }\PYG{p}{.}\PYG{p}{.}\PYG{p}{.}\PYG{p}{)}\PYG{+w}{                             }\PYG{p}{\PYGZam{}}
\PYG{+w}{     }\PYG{o}{/}\PYG{p}{)}
\end{sphinxVerbatim}

\sphinxAtStartPar
The second item in a metadata entry describes the Fortran primitive
(intrinsic) type of the data of a kernel argument. The currently supported
values are \sphinxcode{\sphinxupquote{GH\_REAL}}, \sphinxcode{\sphinxupquote{GH\_INTEGER}} and \sphinxcode{\sphinxupquote{GH\_LOGICAL}} for \sphinxcode{\sphinxupquote{real}},
\sphinxcode{\sphinxupquote{integer}} and \sphinxcode{\sphinxupquote{logical}} data, respectively. This information is
mandatory. Valid data types for each LFRic API argument type are specified
later in this section (see {\hyperref[\detokenize{dynamo0p3:lfric-kernel-valid-data-type}]{\sphinxcrossref{\DUrole{std,std-ref}{Valid Data Types}}}}).

\sphinxAtStartPar
The third component of argument metadata describes how the Kernel
makes use of the data being passed into it (the way it is accessed
within a Kernel). This information is mandatory. There are currently 6
possible values of this metadata \sphinxcode{\sphinxupquote{GH\_READ}}, \sphinxcode{\sphinxupquote{GH\_WRITE}},
\sphinxcode{\sphinxupquote{GH\_READWRITE}}, \sphinxcode{\sphinxupquote{GH\_INC}}, \sphinxcode{\sphinxupquote{GH\_READINC}} and \sphinxcode{\sphinxupquote{GH\_SUM}}. However,
not all combinations of metadata entries are valid and PSyclone will
raise an exception if an invalid combination is specified. Valid
combinations are specified later in this section (see
{\hyperref[\detokenize{dynamo0p3:lfric-kernel-valid-access}]{\sphinxcrossref{\DUrole{std,std-ref}{Valid Access Modes}}}}).
\begin{itemize}
\item {} 
\sphinxAtStartPar
\sphinxcode{\sphinxupquote{GH\_READ}} indicates that the data is read and is unmodified.

\item {} 
\sphinxAtStartPar
\sphinxcode{\sphinxupquote{GH\_WRITE}} indicates the data is modified in the Kernel before
(optionally) being read. If any shared DoFs are written to then
different iterations of the Kernel must write the same value.

\item {} 
\sphinxAtStartPar
\sphinxcode{\sphinxupquote{GH\_READWRITE}} indicates that different iterations of a Kernel
update quantities which do not share DoFs, such as operators and
fields over discontinuous function spaces. If a Kernel modifies only
discontinuous fields and/or operators there is no need for
synchronisation or colouring when running such Kernels in parallel.
However, modifying another field with a \sphinxcode{\sphinxupquote{GH\_INC}} access in a
Kernel means that synchronisation or colouring is required for
parallel runs.

\item {} 
\sphinxAtStartPar
\sphinxcode{\sphinxupquote{GH\_INC}} indicates that different iterations of a Kernel make
contributions to shared values. For example, values at cell faces
may receive contributions from cells on either side of the
face. This means that such a Kernel needs appropriate
synchronisation (or colouring) to run in parallel.

\item {} 
\sphinxAtStartPar
\sphinxcode{\sphinxupquote{GH\_READINC}} indicates that the data is first read and then
subsequently incremented. Therefore this is equivalent to a
\sphinxcode{\sphinxupquote{GH\_READ}} followed by a \sphinxcode{\sphinxupquote{GH\_INC}}.

\item {} 
\sphinxAtStartPar
\sphinxcode{\sphinxupquote{GH\_SUM}} is an example of a reduction and is the only reduction
currently supported in PSyclone. This metadata indicates that values
are summed over calls to Kernel code.

\end{itemize}

\sphinxAtStartPar
For example:

\begin{sphinxVerbatim}[commandchars=\\\{\}]
\PYG{k}{type}\PYG{p}{(}\PYG{n}{arg\PYGZus{}type}\PYG{p}{)}\PYG{+w}{ }\PYG{k+kd}{::}\PYG{+w}{ }\PYG{n}{meta\PYGZus{}args}\PYG{p}{(}\PYG{l+m+mi}{6}\PYG{p}{)}\PYG{+w}{ }\PYG{o}{=}\PYG{+w}{ }\PYG{p}{(}\PYG{o}{/}\PYG{+w}{                            }\PYG{p}{\PYGZam{}}
\PYG{+w}{     }\PYG{n}{arg\PYGZus{}type}\PYG{p}{(}\PYG{n}{GH\PYGZus{}OPERATOR}\PYG{p}{,}\PYG{+w}{ }\PYG{n}{GH\PYGZus{}REAL}\PYG{p}{,}\PYG{+w}{    }\PYG{n}{GH\PYGZus{}READ}\PYG{p}{,}\PYG{+w}{      }\PYG{p}{.}\PYG{p}{.}\PYG{p}{.}\PYG{+w}{ }\PYG{p}{)}\PYG{p}{,}\PYG{+w}{    }\PYG{p}{\PYGZam{}}
\PYG{+w}{     }\PYG{n}{arg\PYGZus{}type}\PYG{p}{(}\PYG{n}{GH\PYGZus{}FIELD}\PYG{o}{*}\PYG{l+m+mi}{3}\PYG{p}{,}\PYG{+w}{  }\PYG{n}{GH\PYGZus{}REAL}\PYG{p}{,}\PYG{+w}{    }\PYG{n}{GH\PYGZus{}WRITE}\PYG{p}{,}\PYG{+w}{     }\PYG{p}{.}\PYG{p}{.}\PYG{p}{.}\PYG{+w}{ }\PYG{p}{)}\PYG{p}{,}\PYG{+w}{    }\PYG{p}{\PYGZam{}}
\PYG{+w}{     }\PYG{n}{arg\PYGZus{}type}\PYG{p}{(}\PYG{n}{GH\PYGZus{}FIELD}\PYG{p}{,}\PYG{+w}{    }\PYG{n}{GH\PYGZus{}REAL}\PYG{p}{,}\PYG{+w}{    }\PYG{n}{GH\PYGZus{}READWRITE}\PYG{p}{,}\PYG{+w}{ }\PYG{p}{.}\PYG{p}{.}\PYG{p}{.}\PYG{+w}{ }\PYG{p}{)}\PYG{p}{,}\PYG{+w}{    }\PYG{p}{\PYGZam{}}
\PYG{+w}{     }\PYG{n}{arg\PYGZus{}type}\PYG{p}{(}\PYG{n}{GH\PYGZus{}FIELD}\PYG{p}{,}\PYG{+w}{    }\PYG{n}{GH\PYGZus{}INTEGER}\PYG{p}{,}\PYG{+w}{ }\PYG{n}{GH\PYGZus{}INC}\PYG{p}{,}\PYG{+w}{       }\PYG{p}{.}\PYG{p}{.}\PYG{p}{.}\PYG{+w}{ }\PYG{p}{)}\PYG{p}{,}\PYG{+w}{    }\PYG{p}{\PYGZam{}}
\PYG{+w}{     }\PYG{n}{arg\PYGZus{}type}\PYG{p}{(}\PYG{n}{GH\PYGZus{}FIELD}\PYG{p}{,}\PYG{+w}{    }\PYG{n}{GH\PYGZus{}REAL}\PYG{p}{,}\PYG{+w}{    }\PYG{n}{GH\PYGZus{}READINC}\PYG{p}{,}\PYG{+w}{   }\PYG{p}{.}\PYG{p}{.}\PYG{p}{.}\PYG{+w}{ }\PYG{p}{)}\PYG{p}{,}\PYG{+w}{    }\PYG{p}{\PYGZam{}}
\PYG{+w}{     }\PYG{n}{arg\PYGZus{}type}\PYG{p}{(}\PYG{n}{GH\PYGZus{}SCALAR}\PYG{p}{,}\PYG{+w}{   }\PYG{n}{GH\PYGZus{}REAL}\PYG{p}{,}\PYG{+w}{    }\PYG{n}{GH\PYGZus{}SUM}\PYG{p}{)}\PYG{+w}{                 }\PYG{p}{\PYGZam{}}
\PYG{+w}{     }\PYG{o}{/}\PYG{p}{)}
\end{sphinxVerbatim}

\begin{sphinxadmonition}{warning}{Warning:}
\sphinxAtStartPar
It is important that \sphinxcode{\sphinxupquote{GH\_INC}} is not incorrectly used
in place of a \sphinxcode{\sphinxupquote{GH\_READINC}} access as it could result in
the reading of data from a dirty outermost halo when run
in parallel, giving incorrect results. The reason for
this is that PSyclone does not add a halo exchange for
the outermost modified halo level of a field before a
loop that contains a \sphinxcode{\sphinxupquote{GH\_INC}} access to that field,
i.e. a loop iterating to the level\sphinxhyphen{}\sphinxcode{\sphinxupquote{n}} halo will result
in a halo exchange to the level\sphinxhyphen{}(\sphinxcode{\sphinxupquote{n\sphinxhyphen{}1}}) halo being
added before the loop (which means no halo exchange is
added when \sphinxcode{\sphinxupquote{n==1}}). The reason this can be performed is
because any computation in the outermost halo will be
incorrect (will only compute partial sums) and PSyclone
therefore sets this halo level to dirty after the loop
has completed. There is, therefore, no reason to make the
values of the incremented field clean for the outermost
modified halo. However, this optimisation does require
that any (dirty) data in the outermost modified halo does
not result in exceptions. With some compilers an
exception can occur for a field that has not yet had its
outermost halo data written to, i.e. if the uninitialised
data is read. To avoid this potential problem in user
code it is recommended that a redundant computation
{\hyperref[\detokenize{dynamo0p3:lfric-api-transformations}]{\sphinxcrossref{\DUrole{std,std-ref}{transformation}}}}
is added to compute all \sphinxcode{\sphinxupquote{setval\_c}}, \sphinxcode{\sphinxupquote{setval\_x}} and
\sphinxcode{\sphinxupquote{setval\_random}} Built\sphinxhyphen{}in calls (see {\hyperref[\detokenize{dynamo0p3:lfric-built-ins}]{\sphinxcrossref{\DUrole{std,std-ref}{Built\sphinxhyphen{}ins}}}})
to the same halo depth as the associated \sphinxcode{\sphinxupquote{GH\_INC}}
access \sphinxhyphen{} which is level\sphinxhyphen{}1 without any redundant
computation transformations being applied to the
associated loops. This will guarantee that all data has
been initialised with a value before it is incremented
and avoid any potential exceptions.
\end{sphinxadmonition}

\begin{sphinxadmonition}{note}{Note:}
\sphinxAtStartPar
In the LFRic API only {\hyperref[\detokenize{dynamo0p3:lfric-built-ins}]{\sphinxcrossref{\DUrole{std,std-ref}{Built\sphinxhyphen{}ins}}}} are permitted
to write to scalar arguments (and hence perform reductions).
Furthermore, this permission is currently restricted to \sphinxcode{\sphinxupquote{real}}
scalars (\sphinxcode{\sphinxupquote{GH\_SCALAR, GH\_REAL}}) as the LFRic infrastructure
does not yet support \sphinxcode{\sphinxupquote{integer}} and \sphinxcode{\sphinxupquote{logical}} reductions.
\end{sphinxadmonition}

\sphinxAtStartPar
For a scalar, the argument metadata contains only these three entries.
However, fields and operators require further entries specifying
function\sphinxhyphen{}space information.
The meaning of these further entries differs depending on whether a
field or an operator is being described.

\sphinxAtStartPar
In the case of an operator, the fourth and fifth arguments describe
the \sphinxcode{\sphinxupquote{to}} and \sphinxcode{\sphinxupquote{from}} function spaces respectively. In the case of a
field the fourth argument specifies the function space that the field
lives on. More details about the supported function spaces are in
subsection {\hyperref[\detokenize{dynamo0p3:lfric-function-space}]{\sphinxcrossref{\DUrole{std,std-ref}{Supported Function Spaces}}}}.

\sphinxAtStartPar
For example, the metadata for a kernel that applies a column\sphinxhyphen{}wise
operator to a field might look like:

\begin{sphinxVerbatim}[commandchars=\\\{\}]
\PYG{k}{type}\PYG{p}{(}\PYG{n}{arg\PYGZus{}type}\PYG{p}{)}\PYG{+w}{ }\PYG{k+kd}{::}\PYG{+w}{ }\PYG{n}{meta\PYGZus{}args}\PYG{p}{(}\PYG{l+m+mi}{3}\PYG{p}{)}\PYG{+w}{ }\PYG{o}{=}\PYG{+w}{ }\PYG{p}{(}\PYG{o}{/}\PYG{+w}{                              }\PYG{p}{\PYGZam{}}
\PYG{+w}{     }\PYG{n}{arg\PYGZus{}type}\PYG{p}{(}\PYG{n}{GH\PYGZus{}FIELD}\PYG{p}{,}\PYG{+w}{ }\PYG{n}{GH\PYGZus{}REAL}\PYG{p}{,}\PYG{+w}{ }\PYG{n}{GH\PYGZus{}INC}\PYG{p}{,}\PYG{+w}{ }\PYG{n}{W1}\PYG{p}{)}\PYG{p}{,}\PYG{+w}{                    }\PYG{p}{\PYGZam{}}
\PYG{+w}{     }\PYG{n}{arg\PYGZus{}type}\PYG{p}{(}\PYG{n}{GH\PYGZus{}FIELD}\PYG{p}{,}\PYG{+w}{ }\PYG{n}{GH\PYGZus{}REAL}\PYG{p}{,}\PYG{+w}{ }\PYG{n}{GH\PYGZus{}READ}\PYG{p}{,}\PYG{+w}{ }\PYG{n}{W2H}\PYG{p}{)}\PYG{p}{,}\PYG{+w}{                  }\PYG{p}{\PYGZam{}}
\PYG{+w}{     }\PYG{n}{arg\PYGZus{}type}\PYG{p}{(}\PYG{n}{GH\PYGZus{}COLUMNWISE\PYGZus{}OPERATOR}\PYG{p}{,}\PYG{+w}{ }\PYG{n}{GH\PYGZus{}REAL}\PYG{p}{,}\PYG{+w}{ }\PYG{n}{GH\PYGZus{}READ}\PYG{p}{,}\PYG{+w}{ }\PYG{n}{W1}\PYG{p}{,}\PYG{+w}{ }\PYG{n}{W2H}\PYG{p}{)}\PYG{+w}{ }\PYG{p}{\PYGZam{}}
\PYG{+w}{     }\PYG{o}{/}\PYG{p}{)}
\end{sphinxVerbatim}

\sphinxAtStartPar
In some cases a Kernel may be written so that it works for fields and/or
operators from any type of a vector \sphinxcode{\sphinxupquote{W2*}} space (all \sphinxcode{\sphinxupquote{W2*}} spaces
except for the \sphinxcode{\sphinxupquote{W2*trace}} spaces, see Section
{\hyperref[\detokenize{dynamo0p3:lfric-function-space}]{\sphinxcrossref{\DUrole{std,std-ref}{Supported Function Spaces}}}} below).
In this case the metadata should be specified as being \sphinxcode{\sphinxupquote{ANY\_W2}}.

\begin{sphinxadmonition}{warning}{Warning:}
\sphinxAtStartPar
In the current implementation it is assumed that all
fields and/or operators specifying \sphinxcode{\sphinxupquote{ANY\_W2}} within a
kernel will use the \sphinxstylestrong{same} function space. It is up to
the user to ensure this is the case as otherwise invalid
code would be generated.
\end{sphinxadmonition}

\sphinxAtStartPar
It may be that a Kernel is written such that a field and/or operators
may be on/map\sphinxhyphen{}between any function space(s). In this case the metadata
should be specified as being one of \sphinxcode{\sphinxupquote{ANY\_SPACE\_1}}, …, \sphinxcode{\sphinxupquote{ANY\_SPACE\_\textless{}nmax\textgreater{}}}
(see {\hyperref[\detokenize{dynamo0p3:lfric-function-space}]{\sphinxcrossref{\DUrole{std,std-ref}{Supported Function Spaces}}}}), with the
number of spaces, \sphinxcode{\sphinxupquote{\textless{}nmax\textgreater{}}}, being set in the {\hyperref[\detokenize{configuration:configuration}]{\sphinxcrossref{\DUrole{std,std-ref}{PSyclone configuration
file}}}} (see {\hyperref[\detokenize{dynamo0p3:lfric-num-any-spaces}]{\sphinxcrossref{\DUrole{std,std-ref}{here}}}} for more
details on this option).

\sphinxAtStartPar
If the generic function spaces are known to be discontinuous the metadata
may be specified as being one of \sphinxcode{\sphinxupquote{ANY\_DISCONTINUOUS\_SPACE\_1}}, …,
\sphinxcode{\sphinxupquote{ANY\_DISCONTINUOUS\_SPACE\_\textless{}nmax\textgreater{}}} in order to avoid unnecessary computation
into the halos (see rules for
{\hyperref[\detokenize{dynamo0p3:lfric-user-kernel-rules}]{\sphinxcrossref{\DUrole{std,std-ref}{user\sphinxhyphen{}supplied kernels}}}} above).
The reason for having different names is that a Kernel might be written
to allow 2 or more arguments to be able to support any function space
but for a particular call the function spaces may have to be the same as
each other. Again, \sphinxcode{\sphinxupquote{\textless{}nmax\textgreater{}}} is the {\hyperref[\detokenize{dynamo0p3:lfric-num-any-spaces}]{\sphinxcrossref{\DUrole{std,std-ref}{configurable}}}} number of generalised discontinuous function spaces.

\sphinxAtStartPar
In the example below, the first field entry supports any function space but
it must be the same as the operator’s \sphinxcode{\sphinxupquote{to}} function space. Similarly,
the second field entry supports any function space but it must be the same
as the operator’s \sphinxcode{\sphinxupquote{from}} function space. Note, the metadata does not
forbid \sphinxcode{\sphinxupquote{ANY\_SPACE\_1}} and \sphinxcode{\sphinxupquote{ANY\_SPACE\_2}} from being the same.

\begin{sphinxVerbatim}[commandchars=\\\{\}]
\PYG{k}{type}\PYG{p}{(}\PYG{n}{arg\PYGZus{}type}\PYG{p}{)}\PYG{+w}{ }\PYG{k+kd}{::}\PYG{+w}{ }\PYG{n}{meta\PYGZus{}args}\PYG{p}{(}\PYG{l+m+mi}{3}\PYG{p}{)}\PYG{+w}{ }\PYG{o}{=}\PYG{+w}{ }\PYG{p}{(}\PYG{o}{/}\PYG{+w}{                                    }\PYG{p}{\PYGZam{}}
\PYG{+w}{     }\PYG{n}{arg\PYGZus{}type}\PYG{p}{(}\PYG{n}{GH\PYGZus{}FIELD}\PYG{p}{,}\PYG{+w}{    }\PYG{n}{GH\PYGZus{}REAL}\PYG{p}{,}\PYG{+w}{ }\PYG{n}{GH\PYGZus{}INC}\PYG{p}{,}\PYG{+w}{  }\PYG{n}{ANY\PYGZus{}SPACE\PYGZus{}1}\PYG{p}{)}\PYG{p}{,}\PYG{+w}{             }\PYG{p}{\PYGZam{}}
\PYG{+w}{     }\PYG{n}{arg\PYGZus{}type}\PYG{p}{(}\PYG{n}{GH\PYGZus{}FIELD}\PYG{o}{*}\PYG{l+m+mi}{3}\PYG{p}{,}\PYG{+w}{  }\PYG{n}{GH\PYGZus{}REAL}\PYG{p}{,}\PYG{+w}{ }\PYG{n}{GH\PYGZus{}INC}\PYG{p}{,}\PYG{+w}{  }\PYG{n}{ANY\PYGZus{}SPACE\PYGZus{}2}\PYG{p}{)}\PYG{p}{,}\PYG{+w}{             }\PYG{p}{\PYGZam{}}
\PYG{+w}{     }\PYG{n}{arg\PYGZus{}type}\PYG{p}{(}\PYG{n}{GH\PYGZus{}OPERATOR}\PYG{p}{,}\PYG{+w}{ }\PYG{n}{GH\PYGZus{}REAL}\PYG{p}{,}\PYG{+w}{ }\PYG{n}{GH\PYGZus{}READ}\PYG{p}{,}\PYG{+w}{ }\PYG{n}{ANY\PYGZus{}SPACE\PYGZus{}1}\PYG{p}{,}\PYG{+w}{ }\PYG{n}{ANY\PYGZus{}SPACE\PYGZus{}2}\PYG{p}{)}\PYG{+w}{ }\PYG{p}{\PYGZam{}}
\PYG{+w}{     }\PYG{o}{/}\PYG{p}{)}
\end{sphinxVerbatim}

\sphinxAtStartPar
Note also that the scope of this naming of any\sphinxhyphen{}space function spaces is
restricted to the argument list of individual kernels. I.e. if an
Invoke contains say, two kernel calls that each support arguments on
any function space, e.g. \sphinxcode{\sphinxupquote{ANY\_SPACE\_1}}, there is no requirement that
these two function spaces be the same. Put another way, if an Invoke
contained two calls of a kernel with arguments described by the above
metadata then the first field argument passed to each kernel call
need not be on the same space.


\paragraph{Valid Data Types}
\label{\detokenize{dynamo0p3:valid-data-types}}\label{\detokenize{dynamo0p3:lfric-kernel-valid-data-type}}
\sphinxAtStartPar
As mentioned earlier, the currently supported Fortran primitive
(intrinsic) types for kernel argument data are \sphinxcode{\sphinxupquote{real}}, \sphinxcode{\sphinxupquote{integer}}
and \sphinxcode{\sphinxupquote{logical}}, described by the \sphinxcode{\sphinxupquote{GH\_REAL}}, \sphinxcode{\sphinxupquote{GH\_INTEGER}} and
\sphinxcode{\sphinxupquote{GH\_LOGICAL}} metadata descriptors. Supported data types for each
argument type are given in the table below (please note that
{\hyperref[\detokenize{dynamo0p3:lfric-field-vector}]{\sphinxcrossref{\DUrole{std,std-ref}{field vectors}}}} follow the same rules as
the {\hyperref[\detokenize{dynamo0p3:lfric-field}]{\sphinxcrossref{\DUrole{std,std-ref}{LFRic fields}}}}):


\begin{savenotes}\sphinxattablestart
\sphinxthistablewithglobalstyle
\sphinxthistablewithvlinesstyle
\centering
\begin{tabulary}{\linewidth}[t]{|l|l|}
\sphinxtoprule
\sphinxstyletheadfamily 
\sphinxAtStartPar
Argument Type
&\sphinxstyletheadfamily 
\sphinxAtStartPar
Data Type
\\
\sphinxmidrule
\sphinxtableatstartofbodyhook
\sphinxAtStartPar
GH\_SCALAR
&
\sphinxAtStartPar
GH\_REAL, GH\_INTEGER, GH\_LOGICAL
\\
\sphinxhline
\sphinxAtStartPar
GH\_FIELD
&
\sphinxAtStartPar
GH\_REAL, GH\_INTEGER
\\
\sphinxhline
\sphinxAtStartPar
GH\_OPERATOR
&
\sphinxAtStartPar
GH\_REAL
\\
\sphinxhline
\sphinxAtStartPar
GH\_COLUMNWISE\_OPERATOR
&
\sphinxAtStartPar
GH\_REAL
\\
\sphinxbottomrule
\end{tabulary}
\sphinxtableafterendhook\par
\sphinxattableend\end{savenotes}


\paragraph{Valid Access Modes}
\label{\detokenize{dynamo0p3:valid-access-modes}}\label{\detokenize{dynamo0p3:lfric-kernel-valid-access}}
\sphinxAtStartPar
As mentioned earlier, not all combinations of metadata are
valid. Valid combinations for each argument type in
user\sphinxhyphen{}defined Kernels are summarised here. All argument types
(\sphinxcode{\sphinxupquote{GH\_SCALAR}}, \sphinxcode{\sphinxupquote{GH\_FIELD}}, \sphinxcode{\sphinxupquote{GH\_OPERATOR}} and
\sphinxcode{\sphinxupquote{GH\_COLUMNWISE\_OPERATOR}}) may be read within a Kernel and this
is specified in metadata using \sphinxcode{\sphinxupquote{GH\_READ}}. At least one kernel
argument must be listed as being modified. When data is \sphinxstyleemphasis{modified}
in a user\sphinxhyphen{}supplied Kernel (i.e. a Kernel that operates on a
\sphinxcode{\sphinxupquote{CELL\_COLUMN}}, see {\hyperref[\detokenize{dynamo0p3:lfric-operates-on}]{\sphinxcrossref{\DUrole{std,std-ref}{iteration space metadata}}}}) then the permitted access
modes depend upon the argument type and the function space it is on:


\begin{savenotes}\sphinxattablestart
\sphinxthistablewithglobalstyle
\sphinxthistablewithvlinesstyle
\centering
\begin{tabulary}{\linewidth}[t]{|l|l|l|}
\sphinxtoprule
\sphinxstyletheadfamily 
\sphinxAtStartPar
Argument Type
&\sphinxstyletheadfamily 
\sphinxAtStartPar
Function Space
&\sphinxstyletheadfamily 
\sphinxAtStartPar
Access Type
\\
\sphinxmidrule
\sphinxtableatstartofbodyhook
\sphinxAtStartPar
GH\_SCALAR
&
\sphinxAtStartPar
n/a
&
\sphinxAtStartPar
GH\_READ
\\
\sphinxhline
\sphinxAtStartPar
GH\_FIELD
&
\sphinxAtStartPar
Discontinuous
&
\sphinxAtStartPar
GH\_READ, GH\_WRITE,
GH\_READWRITE
\\
\sphinxhline
\sphinxAtStartPar
GH\_FIELD
&
\sphinxAtStartPar
Continuous
&
\sphinxAtStartPar
GH\_READ, GH\_WRITE,
GH\_INC, GH\_READINC
\\
\sphinxhline
\sphinxAtStartPar
GH\_OPERATOR
&
\sphinxAtStartPar
Any for both ‘to’ and ‘from’
&
\sphinxAtStartPar
GH\_READ, GH\_WRITE,
GH\_READWRITE
\\
\sphinxhline
\sphinxAtStartPar
GH\_COLUMNWISE\_OPERATOR
&
\sphinxAtStartPar
Any for both ‘to’ and ‘from’
&
\sphinxAtStartPar
GH\_READ, GH\_WRITE,
GH\_READWRITE
\\
\sphinxbottomrule
\end{tabulary}
\sphinxtableafterendhook\par
\sphinxattableend\end{savenotes}

\sphinxAtStartPar
Note that scalar arguments to user\sphinxhyphen{}defined Kernels must be read\sphinxhyphen{}only.
Only {\hyperref[\detokenize{dynamo0p3:lfric-built-ins}]{\sphinxcrossref{\DUrole{std,std-ref}{Built\sphinxhyphen{}ins}}}} are permitted to modify scalar
arguments. In practice this means that the only allowed access for the scalars
in user\sphinxhyphen{}defined Kernels is \sphinxcode{\sphinxupquote{GH\_READ}} (see the allowed accesses for arguments
in Built\sphinxhyphen{}ins in the {\hyperref[\detokenize{dynamo0p3:lfric-built-ins-dtype-access}]{\sphinxcrossref{\DUrole{std,std-ref}{section below}}}}).

\sphinxAtStartPar
Note also that a \sphinxcode{\sphinxupquote{GH\_FIELD}} argument that has \sphinxcode{\sphinxupquote{GH\_WRITE}} or
\sphinxcode{\sphinxupquote{GH\_READWRITE}} as its access pattern must typically (see below) be
on a horizontally\sphinxhyphen{}discontinuous function space (see
{\hyperref[\detokenize{dynamo0p3:lfric-function-space}]{\sphinxcrossref{\DUrole{std,std-ref}{Supported Function Spaces}}}} for the list of discontinuous function
spaces). Parallelisation of the loop over the horizontal domain for a
Kernel that updates such a field will not require colouring for either
of the above cases (since there are no shared entities).

\sphinxAtStartPar
There is however an exception to this \sphinxhyphen{} certain Kernels may write to
shared entities but each Kernel iteration is guaranteed to write the
\sphinxstyleemphasis{same value} to a given shared DoF. In this case, provided that the
first access to any such shared DoF is a write, the loop containing
such a Kernel may be parallelised without colouring. Therefore,
\sphinxcode{\sphinxupquote{GH\_WRITE}} access is permitted for \sphinxcode{\sphinxupquote{GH\_FIELD}} arguments on
continuous function spaces. Obviously, care must be taken to ensure
that the Kernel implementation satisfies the constraints just
described as PSyclone cannot currently check this.

\sphinxAtStartPar
If a field is described as being on \sphinxcode{\sphinxupquote{ANY\_SPACE\_*}}, there is currently no
way to determine its continuity from the metadata (unless we can statically
determine the space of the field being passed in). At the moment this type
of a user\sphinxhyphen{}supplied Kernel is always treated as if it is updating a field
that is on a function space that is continuous in the horizontal, even if
it is not (see rules for {\hyperref[\detokenize{dynamo0p3:lfric-user-kernel-rules}]{\sphinxcrossref{\DUrole{std,std-ref}{user\sphinxhyphen{}supplied kernels}}}} above).

\sphinxAtStartPar
There is no restriction on the number and function spaces of other
quantities that a general\sphinxhyphen{}purpose kernel can modify other than that it
must modify at least one. The rules for kernels involving CMA operators,
however, are stricter and only one argument may be modified (the CMA
operator itself for assembly, a field for CMA\sphinxhyphen{}application and a CMA
operator for matrix\sphinxhyphen{}matrix kernels). If a kernel writes to quantities
on different function spaces then PSyclone generates loop bounds
appropriate to the largest iteration space. This means that if a
single kernel updates one quantity on a continuous function space and
one on a discontinuous space then the resulting loop will include
cells in the level\sphinxhyphen{}1 halo since they are required for a quantity on a
continuous space. As a consequence, any quantities on a discontinuous
space will then be computed redundantly in the level\sphinxhyphen{}1 halo. Currently
PSyclone makes no attempt to take advantage of this (by e.g. setting
the appropriate level\sphinxhyphen{}1 halo to ‘clean’).

\sphinxAtStartPar
PSyclone ensures that both CMA and LMA operators are computed
(redundantly) out to the level\sphinxhyphen{}1 halo cells. This permits their use in
kernels which modify quantities on continuous function spaces and also
in subsequent redundant computation of other quantities on
discontinuous function spaces. In conjunction with this, PSyclone also
checks (when generating the PSy layer) that any kernels which read
operator values do not do so beyond the level\sphinxhyphen{}1 halo. If any such
accesses are found then PSyclone aborts.


\paragraph{Supported Function Spaces}
\label{\detokenize{dynamo0p3:supported-function-spaces}}\label{\detokenize{dynamo0p3:lfric-function-space}}
\sphinxAtStartPar
As mentioned in the {\hyperref[\detokenize{dynamo0p3:lfric-field}]{\sphinxcrossref{\DUrole{std,std-ref}{Field}}}} and {\hyperref[\detokenize{dynamo0p3:lfric-field-vector}]{\sphinxcrossref{\DUrole{std,std-ref}{Field Vector}}}}
sections, the function space of an argument specifies how it maps
onto the underlying topology and, additionally, whether the data at a
point is a vector. In LFRic API the dimension of the basis function
set for the scalar function spaces is 1 and for the vector function spaces
is 3 (see the table in {\hyperref[\detokenize{dynamo0p3:lfric-stub-generation-rules}]{\sphinxcrossref{\DUrole{std,std-ref}{Rules for General\sphinxhyphen{}Purpose Kernels}}}} for the
dimensions of the basis and differential basis functions).

\sphinxAtStartPar
Function spaces can share DoFs between cells in the horizontal, vertical
or both directions. Depending on the function space and FEM order,
the shared DoFs can lie on one or more cell entities (faces, edges
and vertices) in each direction. This property is referred to as the
\sphinxstylestrong{continuity} of a function space (horizontal, vertical or full).
Alternatively, if there are no shared DoFs a function space is described
as \sphinxstylestrong{discontinuous} (fully or in a particular direction).

\sphinxAtStartPar
The mixed FEM formulation is built on a foundation set of four function
spaces described below.
\begin{itemize}
\item {} 
\sphinxAtStartPar
\sphinxcode{\sphinxupquote{W0}} is the space of scalar functions with full continuity. The
shared DoFs lie on cell vertices in the lowest order FEM and on
all three entities in higher order FEM.

\item {} 
\sphinxAtStartPar
\sphinxcode{\sphinxupquote{W1}} is the space of vector functions with full continuity in the
tangential direction only. In the lowest order FEM the shared DoFs
lie on cell edges for each component, whereas in higher order they
also lie on cell faces.

\item {} 
\sphinxAtStartPar
\sphinxcode{\sphinxupquote{W2}} is the space of vector functions with full continuity in the
normal direction only. The shared DoFs lie on cell faces for each
component.

\item {} 
\sphinxAtStartPar
\sphinxcode{\sphinxupquote{W3}} is the space of scalar functions with full discontinuity.
All DoFs lie within the cell volume and are not shared across the
cell boundaries.

\end{itemize}

\sphinxAtStartPar
Other spaces required for representation of scalar or component\sphinxhyphen{}wise
vector variables are:
\begin{itemize}
\item {} 
\sphinxAtStartPar
\sphinxcode{\sphinxupquote{Wtheta}} is the space of scalar functions based on the vertical
part of \sphinxcode{\sphinxupquote{W2}}, discontinuous in the horizontal and continuous
in the vertical;

\item {} 
\sphinxAtStartPar
\sphinxcode{\sphinxupquote{W2H}} is the space of vector functions based on the horizontal
part of \sphinxcode{\sphinxupquote{W2}}, continuous in the horizontal and discontinuous
in the vertical;

\item {} 
\sphinxAtStartPar
\sphinxcode{\sphinxupquote{W2V}} is the space of vector functions based on the vertical
part of \sphinxcode{\sphinxupquote{W2}}, discontinuous in the horizontal and continuous
in the vertical;

\item {} 
\sphinxAtStartPar
\sphinxcode{\sphinxupquote{W2broken}} is the space of vector functions, locally identical
to the \sphinxcode{\sphinxupquote{W2}} space. However, DoFs are topologically discontinuous in
all directions despite their placement on cell faces;

\item {} 
\sphinxAtStartPar
\sphinxcode{\sphinxupquote{W2trace}} is the space of scalar functions defined only on cell faces,
resulting from taking the trace of a \sphinxcode{\sphinxupquote{W2}} space. DoFs are shared between
faces, hence making this space fully continuous;

\item {} 
\sphinxAtStartPar
\sphinxcode{\sphinxupquote{W2Htrace}} is the space of scalar functions defined only on cell faces
in the horizontal, resulting from taking the trace of a \sphinxcode{\sphinxupquote{W2H}} space.
DoFs are shared between horizontal faces, hence making this space
continuous in the horizontal and discontinuous in the vertical;

\item {} 
\sphinxAtStartPar
\sphinxcode{\sphinxupquote{W2Vtrace}} is the space of scalar functions defined only on cell faces
in the vertical, resulting from taking the trace of a \sphinxcode{\sphinxupquote{W2V}} space.
DoFs are shared between vertical faces, hence making this space
discontinuous in the horizontal and continuous in the vertical;

\item {} 
\sphinxAtStartPar
\sphinxcode{\sphinxupquote{Wchi}} is the space of scalar functions used to store coordinates
in LFRic. It is fully discontinuous except for the coordinate order
\sphinxcode{\sphinxupquote{0}} when it becomes the \sphinxcode{\sphinxupquote{W0}} space (i.e. fully continuous).
Please see the next section for more details on this function space.

\end{itemize}

\sphinxAtStartPar
The previously mentioned FEM order can be specified in the horizontal
and vertical directions independently, through the \sphinxcode{\sphinxupquote{element\_order\_h}} and
\sphinxcode{\sphinxupquote{element\_order\_v}} arguments in the function space initialisation. These
element orders dictate the polynomial order of the basis functions used to
represent a field. In most cases these element orders will be set to \sphinxstyleemphasis{0} in
LFRic, and increasing them will result in an increased number of DoFs per cell.
Increasing the element order in either direction will never change the
continuity of a space; however, it can increase the number of shared DoFs per
entity.

\sphinxAtStartPar
For example, the \sphinxcode{\sphinxupquote{W3}} space has a single DoF per cell at the lowest order,
situated in the cell volume. The basis functions are constant over the
cell, and the space is discontinuous. Increasing to \sphinxcode{\sphinxupquote{element\_order\_v=1}}
(the so\sphinxhyphen{}called ‘next\sphinxhyphen{}to\sphinxhyphen{}lowest order’ in the vertical direction) will result
in a second volume DoF per cell, yielding linear basis functions vertically,
and constant basis functions horizontally. The space remains discontinuous
in both directions.

\sphinxAtStartPar
In addition to the specific function space metadata, there are also
three generic function space metadata descriptors mentioned in
sections above:
\begin{itemize}
\item {} 
\sphinxAtStartPar
\sphinxcode{\sphinxupquote{ANY\_SPACE\_\textless{}n\textgreater{}}}, \sphinxstyleemphasis{n = 1, 2, … nmax}, for when the function
space of the argument(s) cannot be determined and/or for when
a Kernel has been written so that it works with fields on any
of the available spaces (as mentioned in the
{\hyperref[\detokenize{dynamo0p3:lfric-api-meta-args}]{\sphinxcrossref{\DUrole{std,std-ref}{meta\_args section}}}}, the number of
spaces, \sphinxcode{\sphinxupquote{\textless{}nmax\textgreater{}}}, is {\hyperref[\detokenize{dynamo0p3:lfric-num-any-spaces}]{\sphinxcrossref{\DUrole{std,std-ref}{configurable}}}});

\item {} 
\sphinxAtStartPar
\sphinxcode{\sphinxupquote{ANY\_DISCONTINUOUS\_SPACE\_\textless{}n\textgreater{}}}, \sphinxstyleemphasis{n = 1, 2, … nmax}, for when
the function space of the argument(s) cannot be determined
but is known to be discontinuous and/or for when a Kernel
has been written so that it works with fields on any of the
discontinuous spaces (again, the number of spaces, \sphinxcode{\sphinxupquote{\textless{}nmax\textgreater{}}},
is {\hyperref[\detokenize{dynamo0p3:lfric-num-any-spaces}]{\sphinxcrossref{\DUrole{std,std-ref}{configurable}}}});

\item {} 
\sphinxAtStartPar
\sphinxcode{\sphinxupquote{ANY\_W2}} for any type of a vector \sphinxcode{\sphinxupquote{W2*}} function space, i.e. \sphinxcode{\sphinxupquote{W2}},
\sphinxcode{\sphinxupquote{W2H}}, \sphinxcode{\sphinxupquote{W2V}} and \sphinxcode{\sphinxupquote{W2broken}} but not \sphinxcode{\sphinxupquote{W2*trace}} spaces.

\end{itemize}

\sphinxAtStartPar
As mentioned {\hyperref[\detokenize{dynamo0p3:lfric-user-kernel-rules}]{\sphinxcrossref{\DUrole{std,std-ref}{previously}}}} ,
\sphinxcode{\sphinxupquote{ANY\_SPACE\_\textless{}n\textgreater{}}} and \sphinxcode{\sphinxupquote{ANY\_W2}} function space types are treated as
continuous while \sphinxcode{\sphinxupquote{ANY\_DISCONTINUOUS\_SPACE\_\textless{}n\textgreater{}}} spaces are treated
as discontinuous.

\begin{sphinxadmonition}{note}{Note:}
\sphinxAtStartPar
The name and use of \sphinxcode{\sphinxupquote{ANY\_W2}} metadata (e.g. continuity and
vector or/and scalar basis of \sphinxcode{\sphinxupquote{W2*}} spaces the metadata
can represent) are being reviewed in PSyclone issue \#540.
\end{sphinxadmonition}

\sphinxAtStartPar
Since the LFRic API operates on columns of data, function spaces
are categorised as continuous or discontinuous with regard to their
\sphinxstylestrong{continuity in the horizontal}. For example, a \sphinxcode{\sphinxupquote{GH\_FIELD}} that
specifies \sphinxcode{\sphinxupquote{GH\_INC}} as its access pattern (see
:ref:lfric\sphinxhyphen{}kernel\sphinxhyphen{}valid\sphinxhyphen{}access: above) may be continuous in the vertical
(and discontinuous in the horizontal), continuous in the horizontal
(and discontinuous in the vertical), or continuous in both. In each
case the code is the same. This principle of horizontal continuity also
applies to the three generic \sphinxcode{\sphinxupquote{ANY\_*\_*}} function space identifiers
above. The valid metadata values for continuous and discontinuous
function spaces are summarised in the table below.


\begin{savenotes}\sphinxattablestart
\sphinxthistablewithglobalstyle
\sphinxthistablewithvlinesstyle
\centering
\begin{tabulary}{\linewidth}[t]{|l|l|}
\sphinxtoprule
\sphinxstyletheadfamily 
\sphinxAtStartPar
Function Space Continuity
&\sphinxstyletheadfamily 
\sphinxAtStartPar
Function Space Name
\\
\sphinxmidrule
\sphinxtableatstartofbodyhook
\sphinxAtStartPar
Continuous
&
\sphinxAtStartPar
W0, W1, W2, W2H, W2trace, W2Htrace,
ANY\_W2, ANY\_SPACE\_\textless{}n\textgreater{}
\\
\sphinxhline
\sphinxAtStartPar
Discontinuous
&
\sphinxAtStartPar
W2broken, W2V, W2Vtrace, W3, Wtheta,
ANY\_DISCONTINUOUS\_SPACE\_\textless{}n\textgreater{}
\\
\sphinxbottomrule
\end{tabulary}
\sphinxtableafterendhook\par
\sphinxattableend\end{savenotes}

\sphinxAtStartPar
Horizontally discontinuous function spaces and fields over them will not
need colouring so PSyclone does not perform it. If such attempt is made,
PSyclone will raise a \sphinxcode{\sphinxupquote{Generation Error}} in the \sphinxstylestrong{Dynamo0p3ColourTrans}
transformation (see {\hyperref[\detokenize{dynamo0p3:lfric-api-transformations}]{\sphinxcrossref{\DUrole{std,std-ref}{Transformations}}}} for more details
on transformations). An example of fields iterating over a discontinuous
function space \sphinxcode{\sphinxupquote{Wtheta}} is given in \sphinxcode{\sphinxupquote{examples/lfric/eg9}}, with the
\sphinxcode{\sphinxupquote{GH\_READWRITE}} access descriptor denoting an update to the relevant
fields. This example also demonstrates how to only colour loops over
continuous function spaces when transformations are applied.


\paragraph{Read\sphinxhyphen{}Only Function Spaces}
\label{\detokenize{dynamo0p3:read-only-function-spaces}}\label{\detokenize{dynamo0p3:lfric-ro-function-space}}
\sphinxAtStartPar
LFRic supports the concept of a \sphinxstylestrong{read\sphinxhyphen{}only function space}. A field
on such a function space must not be modified by any kernels contained
within \sphinxcode{\sphinxupquote{invoke}} calls (i.e. within any code that PSyclone is
responsible for). Further, a field on a read\sphinxhyphen{}only function space must
contain clean halos in order to avoid any halo exchanges that would
occur if the field is read within a kernel where redundant
computation is performed.

\sphinxAtStartPar
The primary reason for including a read\sphinxhyphen{}only function space is that it
does not need any halo\sphinxhyphen{}exchange support e.g. it does not require a
routing table, which can reduce the memory footprint.

\sphinxAtStartPar
Currently \sphinxcode{\sphinxupquote{Wchi}} is the only read\sphinxhyphen{}only function space in LFRic.

\sphinxAtStartPar
As a read\sphinxhyphen{}only function space is not modified, it does not matter
whether it is classified as continuous or discontinuous. LFRic
therefore treats read\sphinxhyphen{}only as a third category of function space.


\paragraph{Optional Field Metadata}
\label{\detokenize{dynamo0p3:optional-field-metadata}}
\sphinxAtStartPar
A field entry in the meta\_args array may have an optional fifth element.
This element describes either a stencil access or, for inter\sphinxhyphen{}grid kernels,
which mesh the field is on. Since an inter\sphinxhyphen{}grid kernel is not permitted
to have stencil accesses, these two options are mutually exclusive.
The metadata for each case is described in the following sections.


\subparagraph{Stencil Metadata}
\label{\detokenize{dynamo0p3:stencil-metadata}}
\sphinxAtStartPar
Stencil metadata specifies that the corresponding field argument is accessed
as a stencil operation within the Kernel.  Stencil metadata only makes sense
if the associated field is read within a Kernel i.e. it only makes
sense to specify stencil metadata if the first entry is \sphinxcode{\sphinxupquote{GH\_FIELD}}
and the second entry is \sphinxcode{\sphinxupquote{GH\_READ}}.

\sphinxAtStartPar
Stencil metadata is written in the following format:

\begin{sphinxVerbatim}[commandchars=\\\{\}]
\PYG{n}{STENCIL}\PYG{p}{(}\PYG{k}{type}\PYG{p}{)}
\end{sphinxVerbatim}

\sphinxAtStartPar
where \sphinxcode{\sphinxupquote{type}} may be one of \sphinxcode{\sphinxupquote{X1D}}, \sphinxcode{\sphinxupquote{Y1D}}, \sphinxcode{\sphinxupquote{XORY1D}}, \sphinxcode{\sphinxupquote{CROSS}},
\sphinxcode{\sphinxupquote{CROSS2D}} or \sphinxcode{\sphinxupquote{REGION}}.  As the stencil \sphinxcode{\sphinxupquote{extent}} (the maximum distance from
the central cell that the stencil extends) is not provided in the metadata,
it is expected to be provided by the algorithm writer as part of the
\sphinxcode{\sphinxupquote{invoke}} call (see Section {\hyperref[\detokenize{dynamo0p3:lfric-alg-stencil}]{\sphinxcrossref{\DUrole{std,std-ref}{Stencil Extent}}}}). As there
is currently no way to specify a fixed extent value for stencils in the
Kernel metadata, Kernels must therefore be written to support
different values of extent (i.e. stencils with a variable number of
cells).

\sphinxAtStartPar
The \sphinxcode{\sphinxupquote{XORY1D}} stencil type indicates that the Kernel can accept
either \sphinxcode{\sphinxupquote{X1D}} or \sphinxcode{\sphinxupquote{Y1D}} stencils. In this case it is up to the
algorithm developer to specify which of these it is from the algorithm
layer as part of the \sphinxcode{\sphinxupquote{invoke}} call (see Section
{\hyperref[\detokenize{dynamo0p3:lfric-alg-stencil}]{\sphinxcrossref{\DUrole{std,std-ref}{Stencil Extent}}}}).

\sphinxAtStartPar
For example, the following stencil (with \sphinxcode{\sphinxupquote{extent=2}}):

\begin{sphinxVerbatim}[commandchars=\\\{\}]
| 3 | 2 | 1 | 4 | 5 |
\end{sphinxVerbatim}

\sphinxAtStartPar
would be declared as:

\begin{sphinxVerbatim}[commandchars=\\\{\}]
\PYG{n}{STENCIL}\PYG{p}{(}\PYG{n}{X1D}\PYG{p}{)}
\end{sphinxVerbatim}

\sphinxAtStartPar
and the following stencil (with \sphinxcode{\sphinxupquote{extent=2}}):

\begin{sphinxVerbatim}[commandchars=\\\{\}]
|   |   | 9 |   |   |
|   |   | 8 |   |   |
| 3 | 2 | 1 | 6 | 7 |
|   |   | 4 |   |   |
|   |   | 5 |   |   |
\end{sphinxVerbatim}

\sphinxAtStartPar
would be declared as:

\begin{sphinxVerbatim}[commandchars=\\\{\}]
\PYG{n}{STENCIL}\PYG{p}{(}\PYG{n}{CROSS}\PYG{p}{)}
\end{sphinxVerbatim}

\sphinxAtStartPar
The \sphinxcode{\sphinxupquote{REGION}} stencil references a block of cells:

\begin{sphinxVerbatim}[commandchars=\\\{\}]
| 9 | 8 | 7 |
| 2 | 1 | 6 |
| 3 | 4 | 5 |
\end{sphinxVerbatim}

\sphinxAtStartPar
and would be declared as:

\begin{sphinxVerbatim}[commandchars=\\\{\}]
\PYG{n}{STENCIL}\PYG{p}{(}\PYG{n}{REGION}\PYG{p}{)}
\end{sphinxVerbatim}

\sphinxAtStartPar
Below is an example of stencil information within the full kernel metadata:

\begin{sphinxVerbatim}[commandchars=\\\{\}]
\PYG{k}{type}\PYG{p}{(}\PYG{n}{arg\PYGZus{}type}\PYG{p}{)}\PYG{+w}{ }\PYG{k+kd}{::}\PYG{+w}{ }\PYG{n}{meta\PYGZus{}args}\PYG{p}{(}\PYG{l+m+mi}{3}\PYG{p}{)}\PYG{+w}{ }\PYG{o}{=}\PYG{+w}{ }\PYG{p}{(}\PYG{o}{/}\PYG{+w}{                                  }\PYG{p}{\PYGZam{}}
\PYG{+w}{     }\PYG{n}{arg\PYGZus{}type}\PYG{p}{(}\PYG{n}{GH\PYGZus{}FIELD}\PYG{p}{,}\PYG{+w}{    }\PYG{n}{GH\PYGZus{}REAL}\PYG{p}{,}\PYG{+w}{ }\PYG{n}{GH\PYGZus{}INC}\PYG{p}{,}\PYG{+w}{  }\PYG{n}{W1}\PYG{p}{)}\PYG{p}{,}\PYG{+w}{                    }\PYG{p}{\PYGZam{}}
\PYG{+w}{     }\PYG{n}{arg\PYGZus{}type}\PYG{p}{(}\PYG{n}{GH\PYGZus{}FIELD}\PYG{p}{,}\PYG{+w}{    }\PYG{n}{GH\PYGZus{}REAL}\PYG{p}{,}\PYG{+w}{ }\PYG{n}{GH\PYGZus{}READ}\PYG{p}{,}\PYG{+w}{ }\PYG{n}{W2H}\PYG{p}{,}\PYG{+w}{ }\PYG{n}{STENCIL}\PYG{p}{(}\PYG{n}{CROSS}\PYG{p}{)}\PYG{p}{)}\PYG{p}{,}\PYG{+w}{   }\PYG{p}{\PYGZam{}}
\PYG{+w}{     }\PYG{n}{arg\PYGZus{}type}\PYG{p}{(}\PYG{n}{GH\PYGZus{}OPERATOR}\PYG{p}{,}\PYG{+w}{ }\PYG{n}{GH\PYGZus{}REAL}\PYG{p}{,}\PYG{+w}{ }\PYG{n}{GH\PYGZus{}READ}\PYG{p}{,}\PYG{+w}{ }\PYG{n}{W1}\PYG{p}{,}\PYG{+w}{ }\PYG{n}{W2H}\PYG{p}{)}\PYG{+w}{                }\PYG{p}{\PYGZam{}}
\PYG{+w}{     }\PYG{o}{/}\PYG{p}{)}
\end{sphinxVerbatim}

\sphinxAtStartPar
There is a full example of this distributed with PSyclone. It may
be found in \sphinxcode{\sphinxupquote{examples/lfric/eg5}}.


\subparagraph{Inter\sphinxhyphen{}Grid Metadata}
\label{\detokenize{dynamo0p3:inter-grid-metadata}}\label{\detokenize{dynamo0p3:lfric-intergrid-mdata}}
\sphinxAtStartPar
The alternative form of the optional fifth metadata argument for a
field specifies which mesh the associated field is on.  This is
required for inter\sphinxhyphen{}grid kernels which perform prolongation or
restriction operations on fields (or field vectors) existing on grids
of different resolutions.

\sphinxAtStartPar
Mesh metadata is written in the following format:

\begin{sphinxVerbatim}[commandchars=\\\{\}]
\PYG{n}{mesh\PYGZus{}arg}\PYG{o}{=}\PYG{k}{type}
\end{sphinxVerbatim}

\sphinxAtStartPar
where \sphinxcode{\sphinxupquote{type}} may be one of \sphinxcode{\sphinxupquote{GH\_COARSE}} or \sphinxcode{\sphinxupquote{GH\_FINE}}. Any kernel
having a field argument with this metadata is assumed to be an
inter\sphinxhyphen{}grid kernel and, as such, all of its other arguments (which
must also be fields) must have it specified too. An example of the
metadata for such a kernel is given below:

\begin{sphinxVerbatim}[commandchars=\\\{\}]
\PYG{k}{type}\PYG{p}{(}\PYG{n}{arg\PYGZus{}type}\PYG{p}{)}\PYG{+w}{ }\PYG{k+kd}{::}\PYG{+w}{ }\PYG{n}{meta\PYGZus{}args}\PYG{p}{(}\PYG{l+m+mi}{2}\PYG{p}{)}\PYG{+w}{ }\PYG{o}{=}\PYG{+w}{ }\PYG{p}{(}\PYG{o}{/}\PYG{+w}{                                      }\PYG{p}{\PYGZam{}}
\PYG{+w}{    }\PYG{n}{arg\PYGZus{}type}\PYG{p}{(}\PYG{n}{GH\PYGZus{}FIELD}\PYG{p}{,}\PYG{+w}{ }\PYG{n}{GH\PYGZus{}REAL}\PYG{p}{,}\PYG{+w}{ }\PYG{n}{GH\PYGZus{}READWRITE}\PYG{p}{,}\PYG{+w}{ }\PYG{n}{ANY\PYGZus{}DISCONTINUOUS\PYGZus{}SPACE\PYGZus{}1}\PYG{p}{,}\PYG{+w}{ }\PYG{p}{\PYGZam{}}
\PYG{+w}{                                                    }\PYG{n}{mesh\PYGZus{}arg}\PYG{o}{=}\PYG{n}{GH\PYGZus{}COARSE}\PYG{p}{)}\PYG{p}{,}\PYG{+w}{ }\PYG{p}{\PYGZam{}}
\PYG{+w}{    }\PYG{n}{arg\PYGZus{}type}\PYG{p}{(}\PYG{n}{GH\PYGZus{}FIELD}\PYG{p}{,}\PYG{+w}{ }\PYG{n}{GH\PYGZus{}REAL}\PYG{p}{,}\PYG{+w}{ }\PYG{n}{GH\PYGZus{}READ}\PYG{p}{,}\PYG{+w}{      }\PYG{n}{ANY\PYGZus{}DISCONTINUOUS\PYGZus{}SPACE\PYGZus{}2}\PYG{p}{,}\PYG{+w}{ }\PYG{p}{\PYGZam{}}
\PYG{+w}{                                                    }\PYG{n}{mesh\PYGZus{}arg}\PYG{o}{=}\PYG{n}{GH\PYGZus{}FINE}\PYG{+w}{  }\PYG{p}{)}\PYG{+w}{  }\PYG{p}{\PYGZam{}}
\PYG{+w}{    }\PYG{o}{/}\PYG{p}{)}
\end{sphinxVerbatim}

\sphinxAtStartPar
Note that an inter\sphinxhyphen{}grid kernel must have at least one field (or field\sphinxhyphen{}
vector) argument on each mesh type. Fields that are on different
meshes cannot be on the same function space while those on the same
mesh must also be on the same function space.


\paragraph{Column\sphinxhyphen{}wise Operators (CMA)}
\label{\detokenize{dynamo0p3:column-wise-operators-cma}}
\sphinxAtStartPar
In this section we provide example metadata for each of the three
recognised kernel types involving CMA operators.

\sphinxAtStartPar
Column\sphinxhyphen{}wise operators are constructed from cell\sphinxhyphen{}wise (local) operators.
Therefore, in order to \sphinxstylestrong{assemble} a CMA operator, a kernel must have at
least one read\sphinxhyphen{}only LMA operator, e.g.:

\begin{sphinxVerbatim}[commandchars=\\\{\}]
\PYG{k}{type}\PYG{p}{(}\PYG{n}{arg\PYGZus{}type}\PYG{p}{)}\PYG{+w}{ }\PYG{k+kd}{::}\PYG{+w}{ }\PYG{n}{meta\PYGZus{}args}\PYG{p}{(}\PYG{l+m+mi}{2}\PYG{p}{)}\PYG{+w}{ }\PYG{o}{=}\PYG{+w}{ }\PYG{p}{(}\PYG{o}{/}\PYG{+w}{                                                 }\PYG{p}{\PYGZam{}}
\PYG{+w}{     }\PYG{n}{arg\PYGZus{}type}\PYG{p}{(}\PYG{n}{GH\PYGZus{}OPERATOR}\PYG{p}{,}\PYG{+w}{            }\PYG{n}{GH\PYGZus{}REAL}\PYG{p}{,}\PYG{+w}{ }\PYG{n}{GH\PYGZus{}READ}\PYG{p}{,}\PYG{+w}{  }\PYG{n}{ANY\PYGZus{}SPACE\PYGZus{}1}\PYG{p}{,}\PYG{+w}{ }\PYG{n}{ANY\PYGZus{}SPACE\PYGZus{}2}\PYG{p}{)}\PYG{p}{,}\PYG{+w}{ }\PYG{p}{\PYGZam{}}
\PYG{+w}{     }\PYG{n}{arg\PYGZus{}type}\PYG{p}{(}\PYG{n}{GH\PYGZus{}COLUMNWISE\PYGZus{}OPERATOR}\PYG{p}{,}\PYG{+w}{ }\PYG{n}{GH\PYGZus{}REAL}\PYG{p}{,}\PYG{+w}{ }\PYG{n}{GH\PYGZus{}WRITE}\PYG{p}{,}\PYG{+w}{ }\PYG{n}{ANY\PYGZus{}SPACE\PYGZus{}1}\PYG{p}{,}\PYG{+w}{ }\PYG{n}{ANY\PYGZus{}SPACE\PYGZus{}2}\PYG{p}{)}\PYG{+w}{  }\PYG{p}{\PYGZam{}}
\PYG{+w}{     }\PYG{o}{/}\PYG{p}{)}
\end{sphinxVerbatim}

\sphinxAtStartPar
CMA operators (and their inverse) are \sphinxstylestrong{applied} to fields. Therefore any
kernel of this type must have one read\sphinxhyphen{}only CMA operator, one read\sphinxhyphen{}only
field and a field that is updated, e.g.:

\begin{sphinxVerbatim}[commandchars=\\\{\}]
\PYG{k}{type}\PYG{p}{(}\PYG{n}{arg\PYGZus{}type}\PYG{p}{)}\PYG{+w}{ }\PYG{k+kd}{::}\PYG{+w}{ }\PYG{n}{meta\PYGZus{}args}\PYG{p}{(}\PYG{l+m+mi}{3}\PYG{p}{)}\PYG{+w}{ }\PYG{o}{=}\PYG{+w}{ }\PYG{p}{(}\PYG{o}{/}\PYG{+w}{                                               }\PYG{p}{\PYGZam{}}
\PYG{+w}{     }\PYG{n}{arg\PYGZus{}type}\PYG{p}{(}\PYG{n}{GH\PYGZus{}FIELD}\PYG{p}{,}\PYG{+w}{               }\PYG{n}{GH\PYGZus{}REAL}\PYG{p}{,}\PYG{+w}{ }\PYG{n}{GH\PYGZus{}INC}\PYG{p}{,}\PYG{+w}{  }\PYG{n}{ANY\PYGZus{}SPACE\PYGZus{}1}\PYG{p}{)}\PYG{p}{,}\PYG{+w}{             }\PYG{p}{\PYGZam{}}
\PYG{+w}{     }\PYG{n}{arg\PYGZus{}type}\PYG{p}{(}\PYG{n}{GH\PYGZus{}FIELD}\PYG{p}{,}\PYG{+w}{               }\PYG{n}{GH\PYGZus{}REAL}\PYG{p}{,}\PYG{+w}{ }\PYG{n}{GH\PYGZus{}READ}\PYG{p}{,}\PYG{+w}{ }\PYG{n}{ANY\PYGZus{}SPACE\PYGZus{}2}\PYG{p}{)}\PYG{p}{,}\PYG{+w}{             }\PYG{p}{\PYGZam{}}
\PYG{+w}{     }\PYG{n}{arg\PYGZus{}type}\PYG{p}{(}\PYG{n}{GH\PYGZus{}COLUMNWISE\PYGZus{}OPERATOR}\PYG{p}{,}\PYG{+w}{ }\PYG{n}{GH\PYGZus{}REAL}\PYG{p}{,}\PYG{+w}{ }\PYG{n}{GH\PYGZus{}READ}\PYG{p}{,}\PYG{+w}{ }\PYG{n}{ANY\PYGZus{}SPACE\PYGZus{}1}\PYG{p}{,}\PYG{+w}{ }\PYG{n}{ANY\PYGZus{}SPACE\PYGZus{}2}\PYG{p}{)}\PYG{+w}{ }\PYG{p}{\PYGZam{}}
\PYG{+w}{     }\PYG{o}{/}\PYG{p}{)}
\end{sphinxVerbatim}

\sphinxAtStartPar
\sphinxstylestrong{Matrix\sphinxhyphen{}matrix} kernels compute the product/linear combination of CMA
operators. They must therefore have one such operator that is updated while
the rest are read\sphinxhyphen{}only. They may also have read\sphinxhyphen{}only scalar arguments, e.g.:

\begin{sphinxVerbatim}[commandchars=\\\{\}]
\PYG{k}{type}\PYG{p}{(}\PYG{n}{arg\PYGZus{}type}\PYG{p}{)}\PYG{+w}{ }\PYG{k+kd}{::}\PYG{+w}{ }\PYG{n}{meta\PYGZus{}args}\PYG{p}{(}\PYG{l+m+mi}{3}\PYG{p}{)}\PYG{+w}{ }\PYG{o}{=}\PYG{+w}{ }\PYG{p}{(}\PYG{o}{/}\PYG{+w}{                                                 }\PYG{p}{\PYGZam{}}
\PYG{+w}{     }\PYG{n}{arg\PYGZus{}type}\PYG{p}{(}\PYG{n}{GH\PYGZus{}COLUMNWISE\PYGZus{}OPERATOR}\PYG{p}{,}\PYG{+w}{ }\PYG{n}{GH\PYGZus{}REAL}\PYG{p}{,}\PYG{+w}{ }\PYG{n}{GH\PYGZus{}WRITE}\PYG{p}{,}\PYG{+w}{ }\PYG{n}{ANY\PYGZus{}SPACE\PYGZus{}1}\PYG{p}{,}\PYG{+w}{ }\PYG{n}{ANY\PYGZus{}SPACE\PYGZus{}2}\PYG{p}{)}\PYG{p}{,}\PYG{+w}{ }\PYG{p}{\PYGZam{}}
\PYG{+w}{     }\PYG{n}{arg\PYGZus{}type}\PYG{p}{(}\PYG{n}{GH\PYGZus{}COLUMNWISE\PYGZus{}OPERATOR}\PYG{p}{,}\PYG{+w}{ }\PYG{n}{GH\PYGZus{}REAL}\PYG{p}{,}\PYG{+w}{ }\PYG{n}{GH\PYGZus{}READ}\PYG{p}{,}\PYG{+w}{  }\PYG{n}{ANY\PYGZus{}SPACE\PYGZus{}1}\PYG{p}{,}\PYG{+w}{ }\PYG{n}{ANY\PYGZus{}SPACE\PYGZus{}2}\PYG{p}{)}\PYG{p}{,}\PYG{+w}{ }\PYG{p}{\PYGZam{}}
\PYG{+w}{     }\PYG{n}{arg\PYGZus{}type}\PYG{p}{(}\PYG{n}{GH\PYGZus{}COLUMNWISE\PYGZus{}OPERATOR}\PYG{p}{,}\PYG{+w}{ }\PYG{n}{GH\PYGZus{}REAL}\PYG{p}{,}\PYG{+w}{ }\PYG{n}{GH\PYGZus{}READ}\PYG{p}{,}\PYG{+w}{  }\PYG{n}{ANY\PYGZus{}SPACE\PYGZus{}1}\PYG{p}{,}\PYG{+w}{ }\PYG{n}{ANY\PYGZus{}SPACE\PYGZus{}2}\PYG{p}{)}\PYG{p}{,}\PYG{+w}{ }\PYG{p}{\PYGZam{}}
\PYG{+w}{     }\PYG{n}{arg\PYGZus{}type}\PYG{p}{(}\PYG{n}{GH\PYGZus{}SCALAR}\PYG{p}{,}\PYG{+w}{              }\PYG{n}{GH\PYGZus{}REAL}\PYG{p}{,}\PYG{+w}{ }\PYG{n}{GH\PYGZus{}READ}\PYG{p}{)}\PYG{+w}{ }\PYG{o}{/}\PYG{p}{)}
\end{sphinxVerbatim}

\begin{sphinxadmonition}{note}{Note:}
\sphinxAtStartPar
The order with which arguments are specified in metadata for CMA
kernels does not affect the process of identifying the type of
kernel (whether it is assembly, matrix\sphinxhyphen{}matrix etc.)
\end{sphinxadmonition}


\subsubsection{meta\_funcs}
\label{\detokenize{dynamo0p3:meta-funcs}}\label{\detokenize{dynamo0p3:lfric-meta-funcs}}
\sphinxAtStartPar
The (optional) second component of kernel metadata specifies
whether any quadrature or evaluator data is required for a given
function space. (If no quadrature or evaluator data is required then
this metadata should be omitted.) Consider the
following kernel metadata:

\begin{sphinxVerbatim}[commandchars=\\\{\}]
\PYG{k}{type}\PYG{p}{,}\PYG{+w}{ }\PYG{k}{extends}\PYG{p}{(}\PYG{n}{kernel\PYGZus{}type}\PYG{p}{)}\PYG{+w}{ }\PYG{k+kd}{::}\PYG{+w}{ }\PYG{n}{testkern\PYGZus{}operator\PYGZus{}type}
\PYG{+w}{  }\PYG{k}{type}\PYG{p}{(}\PYG{n}{arg\PYGZus{}type}\PYG{p}{)}\PYG{p}{,}\PYG{+w}{ }\PYG{k}{dimension}\PYG{p}{(}\PYG{l+m+mi}{3}\PYG{p}{)}\PYG{+w}{ }\PYG{k+kd}{::}\PYG{+w}{ }\PYG{n}{meta\PYGZus{}args}\PYG{+w}{ }\PYG{o}{=}\PYG{+w}{                 }\PYG{p}{\PYGZam{}}
\PYG{+w}{      }\PYG{p}{(}\PYG{o}{/}\PYG{+w}{ }\PYG{n}{arg\PYGZus{}type}\PYG{p}{(}\PYG{n}{gh\PYGZus{}operator}\PYG{p}{,}\PYG{+w}{ }\PYG{n}{gh\PYGZus{}real}\PYG{p}{,}\PYG{+w}{    }\PYG{n}{gh\PYGZus{}write}\PYG{p}{,}\PYG{+w}{ }\PYG{n}{w0}\PYG{p}{,}\PYG{+w}{ }\PYG{n}{w0}\PYG{p}{)}\PYG{p}{,}\PYG{+w}{ }\PYG{p}{\PYGZam{}}
\PYG{+w}{         }\PYG{n}{arg\PYGZus{}type}\PYG{p}{(}\PYG{n}{gh\PYGZus{}field}\PYG{o}{*}\PYG{l+m+mi}{3}\PYG{p}{,}\PYG{+w}{  }\PYG{n}{gh\PYGZus{}real}\PYG{p}{,}\PYG{+w}{    }\PYG{n}{gh\PYGZus{}read}\PYG{p}{,}\PYG{+w}{  }\PYG{n}{w1}\PYG{p}{)}\PYG{p}{,}\PYG{+w}{     }\PYG{p}{\PYGZam{}}
\PYG{+w}{         }\PYG{n}{arg\PYGZus{}type}\PYG{p}{(}\PYG{n}{gh\PYGZus{}scalar}\PYG{p}{,}\PYG{+w}{   }\PYG{n}{gh\PYGZus{}integer}\PYG{p}{,}\PYG{+w}{ }\PYG{n}{gh\PYGZus{}read}\PYG{p}{)}\PYG{+w}{           }\PYG{p}{\PYGZam{}}
\PYG{+w}{      }\PYG{o}{/}\PYG{p}{)}
\PYG{+w}{  }\PYG{k}{type}\PYG{p}{(}\PYG{n}{func\PYGZus{}type}\PYG{p}{)}\PYG{+w}{ }\PYG{k+kd}{::}\PYG{+w}{ }\PYG{n}{meta\PYGZus{}funcs}\PYG{p}{(}\PYG{l+m+mi}{2}\PYG{p}{)}\PYG{+w}{ }\PYG{o}{=}\PYG{+w}{                          }\PYG{p}{\PYGZam{}}
\PYG{+w}{      }\PYG{p}{(}\PYG{o}{/}\PYG{+w}{ }\PYG{n}{func\PYGZus{}type}\PYG{p}{(}\PYG{n}{w0}\PYG{p}{,}\PYG{+w}{ }\PYG{n}{gh\PYGZus{}basis}\PYG{p}{,}\PYG{+w}{ }\PYG{n}{gh\PYGZus{}diff\PYGZus{}basis}\PYG{p}{)}\PYG{+w}{               }\PYG{p}{\PYGZam{}}
\PYG{+w}{         }\PYG{n}{func\PYGZus{}type}\PYG{p}{(}\PYG{n}{w1}\PYG{p}{,}\PYG{+w}{ }\PYG{n}{gh\PYGZus{}basis}\PYG{p}{)}\PYG{+w}{                              }\PYG{p}{\PYGZam{}}
\PYG{+w}{      }\PYG{o}{/}\PYG{p}{)}
\PYG{+w}{  }\PYG{k+kt}{integer}\PYG{+w}{ }\PYG{k+kd}{::}\PYG{+w}{ }\PYG{n}{gh\PYGZus{}shape}\PYG{+w}{ }\PYG{o}{=}\PYG{+w}{ }\PYG{n}{gh\PYGZus{}quadrature\PYGZus{}XYoZ}
\PYG{+w}{  }\PYG{k+kt}{integer}\PYG{+w}{ }\PYG{k+kd}{::}\PYG{+w}{ }\PYG{n}{operates\PYGZus{}on}\PYG{+w}{ }\PYG{o}{=}\PYG{+w}{ }\PYG{n}{cell\PYGZus{}column}
\PYG{k}{contains}
\PYG{k}{  }\PYG{k}{procedure}\PYG{p}{,}\PYG{+w}{ }\PYG{k}{nopass}\PYG{+w}{ }\PYG{k+kd}{::}\PYG{+w}{ }\PYG{n}{code}\PYG{+w}{ }\PYG{o}{=}\PYG{o}{\PYGZgt{}}\PYG{+w}{ }\PYG{n}{testkern\PYGZus{}operator\PYGZus{}code}
\PYG{k}{end }\PYG{k}{type }\PYG{n}{testkern\PYGZus{}operator\PYGZus{}type}
\end{sphinxVerbatim}

\sphinxAtStartPar
The \sphinxcode{\sphinxupquote{arg\_type}} component of this metadata describes a kernel that
takes three arguments (an operator, a field and an \sphinxcode{\sphinxupquote{integer}}
scalar). Following the \sphinxcode{\sphinxupquote{meta\_args}} array we now have a
\sphinxcode{\sphinxupquote{meta\_funcs}} array. This allows the user to specify that the kernel
requires basis functions (\sphinxcode{\sphinxupquote{gh\_basis}}) and/or the differential of the
basis functions (\sphinxcode{\sphinxupquote{gh\_diff\_basis}}) on one or more of the function
spaces associated with the arguments listed in \sphinxcode{\sphinxupquote{meta\_args}}.  In this
case we require both for the W0 function space but only basis
functions for W1.

\begin{sphinxadmonition}{note}{Note:}
\sphinxAtStartPar
Basis and differential basis functions for both \sphinxcode{\sphinxupquote{real}}\sphinxhyphen{} and
\sphinxcode{\sphinxupquote{integer}}\sphinxhyphen{}valued field arguments have \sphinxcode{\sphinxupquote{real}} values on the
points on which these functions are {\hyperref[\detokenize{dynamo0p3:lfric-gh-shape}]{\sphinxcrossref{\DUrole{std,std-ref}{required}}}}.
\end{sphinxadmonition}


\subsubsection{meta\_reference\_element}
\label{\detokenize{dynamo0p3:meta-reference-element}}
\sphinxAtStartPar
A kernel that requires properties of the reference element in LFRic
specifies those properties through the \sphinxcode{\sphinxupquote{meta\_reference\_element}}
metadata entry.  (If no reference element properties are required then
this metadata should be omitted.)  Consider the following example
kernel metadata:

\begin{sphinxVerbatim}[commandchars=\\\{\}]
\PYG{k}{type}\PYG{p}{,}\PYG{+w}{ }\PYG{k}{extends}\PYG{p}{(}\PYG{n}{kernel\PYGZus{}type}\PYG{p}{)}\PYG{+w}{ }\PYG{k+kd}{::}\PYG{+w}{ }\PYG{n}{testkern\PYGZus{}type}
\PYG{+w}{  }\PYG{k}{type}\PYG{p}{(}\PYG{n}{arg\PYGZus{}type}\PYG{p}{)}\PYG{p}{,}\PYG{+w}{ }\PYG{k}{dimension}\PYG{p}{(}\PYG{l+m+mi}{2}\PYG{p}{)}\PYG{+w}{ }\PYG{k+kd}{::}\PYG{+w}{ }\PYG{n}{meta\PYGZus{}args}\PYG{+w}{ }\PYG{o}{=}\PYG{+w}{      }\PYG{p}{\PYGZam{}}
\PYG{+w}{      }\PYG{p}{(}\PYG{o}{/}\PYG{+w}{ }\PYG{n}{arg\PYGZus{}type}\PYG{p}{(}\PYG{n}{gh\PYGZus{}field}\PYG{p}{,}\PYG{+w}{ }\PYG{n}{gh\PYGZus{}real}\PYG{p}{,}\PYG{+w}{ }\PYG{n}{gh\PYGZus{}read}\PYG{p}{,}\PYG{+w}{ }\PYG{n}{w1}\PYG{p}{)}\PYG{p}{,}\PYG{+w}{ }\PYG{p}{\PYGZam{}}
\PYG{+w}{         }\PYG{n}{arg\PYGZus{}type}\PYG{p}{(}\PYG{n}{gh\PYGZus{}field}\PYG{p}{,}\PYG{+w}{ }\PYG{n}{gh\PYGZus{}real}\PYG{p}{,}\PYG{+w}{ }\PYG{n}{gh\PYGZus{}inc}\PYG{p}{,}\PYG{+w}{  }\PYG{n}{w0}\PYG{p}{)}\PYG{+w}{ }\PYG{o}{/}\PYG{p}{)}
\PYG{+w}{  }\PYG{k}{type}\PYG{p}{(}\PYG{n}{reference\PYGZus{}element\PYGZus{}data\PYGZus{}type}\PYG{p}{)}\PYG{p}{,}\PYG{+w}{ }\PYG{k}{dimension}\PYG{p}{(}\PYG{l+m+mi}{2}\PYG{p}{)}\PYG{+w}{ }\PYG{k+kd}{::}\PYG{+w}{               }\PYG{p}{\PYGZam{}}
\PYG{+w}{    }\PYG{n}{meta\PYGZus{}reference\PYGZus{}element}\PYG{+w}{ }\PYG{o}{=}\PYG{+w}{                                       }\PYG{p}{\PYGZam{}}
\PYG{+w}{      }\PYG{p}{(}\PYG{o}{/}\PYG{+w}{ }\PYG{n}{reference\PYGZus{}element\PYGZus{}data\PYGZus{}type}\PYG{p}{(}\PYG{n}{normals\PYGZus{}to\PYGZus{}horizontal\PYGZus{}faces}\PYG{p}{)}\PYG{p}{,}\PYG{+w}{ }\PYG{p}{\PYGZam{}}
\PYG{+w}{         }\PYG{n}{reference\PYGZus{}element\PYGZus{}data\PYGZus{}type}\PYG{p}{(}\PYG{n}{normals\PYGZus{}to\PYGZus{}vertical\PYGZus{}faces}\PYG{p}{)}\PYG{+w}{ }\PYG{o}{/}\PYG{p}{)}
\PYG{k}{contains}
\PYG{k}{  }\PYG{k}{procedure}\PYG{p}{,}\PYG{+w}{ }\PYG{k}{nopass}\PYG{+w}{ }\PYG{k+kd}{::}\PYG{+w}{ }\PYG{n}{code}\PYG{+w}{ }\PYG{o}{=}\PYG{o}{\PYGZgt{}}\PYG{+w}{ }\PYG{n}{testkern\PYGZus{}code}
\PYG{k}{end }\PYG{k}{type }\PYG{n}{testkern\PYGZus{}type}
\end{sphinxVerbatim}

\sphinxAtStartPar
This metadata specifies that the \sphinxcode{\sphinxupquote{testkern\_type}} kernel requires two
properties of the reference element. The supported properties are
listed below:


\begin{savenotes}\sphinxattablestart
\sphinxthistablewithglobalstyle
\sphinxthistablewithvlinesstyle
\centering
\begin{tabulary}{\linewidth}[t]{|p{5.5cm}|p{8.5cm}|}
\sphinxtoprule
\sphinxstyletheadfamily 
\sphinxAtStartPar
Name
&\sphinxstyletheadfamily 
\sphinxAtStartPar
Description
\\
\sphinxmidrule
\sphinxtableatstartofbodyhook
\sphinxAtStartPar
normals\_to\_horizontal\_faces
&
\sphinxAtStartPar
Array of normals pointing in the positive
(x, y, z) axis direction for each
horizontal face indexed as (component,
face).
\\
\sphinxhline
\sphinxAtStartPar
normals\_to\_vertical\_faces
&
\sphinxAtStartPar
Array of normals pointing in the positive
(x, y, z) axis direction for each vertical
face indexed as (component, face).
\\
\sphinxhline
\sphinxAtStartPar
normals\_to\_faces
&
\sphinxAtStartPar
Array of normals pointing in the positive
(x, y, z) axis direction for each face
indexed as (component, face).
\\
\sphinxhline
\sphinxAtStartPar
outward\_normals\_to\_horizontal\_faces
&
\sphinxAtStartPar
Array of outward\sphinxhyphen{}pointing normals for each
horizontal face indexed as (component,
face).
\\
\sphinxhline
\sphinxAtStartPar
outward\_normals\_to\_vertical\_faces
&
\sphinxAtStartPar
Array of outward\sphinxhyphen{}pointing normals for each
vertical face indexed as (component, face).
\\
\sphinxhline
\sphinxAtStartPar
outward\_normals\_to\_faces
&
\sphinxAtStartPar
Array of outward\sphinxhyphen{}pointing normals for each
face indexed as (component, face).
\\
\sphinxbottomrule
\end{tabulary}
\sphinxtableafterendhook\par
\sphinxattableend\end{savenotes}


\subsubsection{meta\_mesh}
\label{\detokenize{dynamo0p3:meta-mesh}}
\sphinxAtStartPar
A kernel that requires properties of the LFRic mesh object specifies
those properties through the \sphinxcode{\sphinxupquote{meta\_mesh}} metadata entry. (If no
mesh properties are required then this metadata should be omitted.)
Consider the following example kernel metadata:

\begin{sphinxVerbatim}[commandchars=\\\{\}]
\PYG{k}{type}\PYG{p}{,}\PYG{+w}{ }\PYG{k}{extends}\PYG{p}{(}\PYG{n}{kernel\PYGZus{}type}\PYG{p}{)}\PYG{+w}{ }\PYG{k+kd}{::}\PYG{+w}{ }\PYG{n}{testkern\PYGZus{}type}
\PYG{+w}{  }\PYG{k}{type}\PYG{p}{(}\PYG{n}{arg\PYGZus{}type}\PYG{p}{)}\PYG{p}{,}\PYG{+w}{ }\PYG{k}{dimension}\PYG{p}{(}\PYG{l+m+mi}{2}\PYG{p}{)}\PYG{+w}{ }\PYG{k+kd}{::}\PYG{+w}{ }\PYG{n}{meta\PYGZus{}args}\PYG{+w}{ }\PYG{o}{=}\PYG{+w}{      }\PYG{p}{\PYGZam{}}
\PYG{+w}{      }\PYG{p}{(}\PYG{o}{/}\PYG{+w}{ }\PYG{n}{arg\PYGZus{}type}\PYG{p}{(}\PYG{n}{gh\PYGZus{}field}\PYG{p}{,}\PYG{+w}{ }\PYG{n}{gh\PYGZus{}real}\PYG{p}{,}\PYG{+w}{ }\PYG{n}{gh\PYGZus{}read}\PYG{p}{,}\PYG{+w}{ }\PYG{n}{w1}\PYG{p}{)}\PYG{p}{,}\PYG{+w}{ }\PYG{p}{\PYGZam{}}
\PYG{+w}{         }\PYG{n}{arg\PYGZus{}type}\PYG{p}{(}\PYG{n}{gh\PYGZus{}field}\PYG{p}{,}\PYG{+w}{ }\PYG{n}{gh\PYGZus{}real}\PYG{p}{,}\PYG{+w}{ }\PYG{n}{gh\PYGZus{}inc}\PYG{p}{,}\PYG{+w}{  }\PYG{n}{w0}\PYG{p}{)}\PYG{+w}{ }\PYG{o}{/}\PYG{p}{)}
\PYG{+w}{  }\PYG{k}{type}\PYG{p}{(}\PYG{n}{mesh\PYGZus{}data\PYGZus{}type}\PYG{p}{)}\PYG{p}{,}\PYG{+w}{ }\PYG{k}{dimension}\PYG{p}{(}\PYG{l+m+mi}{1}\PYG{p}{)}\PYG{+w}{ }\PYG{k+kd}{::}\PYG{+w}{            }\PYG{p}{\PYGZam{}}
\PYG{+w}{    }\PYG{n}{meta\PYGZus{}mesh}\PYG{+w}{ }\PYG{o}{=}\PYG{+w}{                                    }\PYG{p}{\PYGZam{}}
\PYG{+w}{      }\PYG{p}{(}\PYG{o}{/}\PYG{+w}{ }\PYG{n}{mesh\PYGZus{}data\PYGZus{}type}\PYG{p}{(}\PYG{n}{adjacent\PYGZus{}face}\PYG{p}{)}\PYG{+w}{ }\PYG{o}{/}\PYG{p}{)}
\PYG{k}{contains}
\PYG{k}{  }\PYG{k}{procedure}\PYG{p}{,}\PYG{+w}{ }\PYG{k}{nopass}\PYG{+w}{ }\PYG{k+kd}{::}\PYG{+w}{ }\PYG{n}{code}\PYG{+w}{ }\PYG{o}{=}\PYG{o}{\PYGZgt{}}\PYG{+w}{ }\PYG{n}{testkern\PYGZus{}code}
\PYG{k}{end }\PYG{k}{type }\PYG{n}{testkern\PYGZus{}type}
\end{sphinxVerbatim}

\sphinxAtStartPar
This metadata specifies that the \sphinxcode{\sphinxupquote{testkern\_type}} kernel requires one
property of the mesh. There is currently one supported property:


\begin{savenotes}\sphinxattablestart
\sphinxthistablewithglobalstyle
\centering
\begin{tabulary}{\linewidth}[t]{TT}
\sphinxtoprule
\sphinxstyletheadfamily 
\sphinxAtStartPar
Name
&\sphinxstyletheadfamily 
\sphinxAtStartPar
Description
\\
\sphinxmidrule
\sphinxtableatstartofbodyhook
\sphinxAtStartPar
adjacent\_face
&
\sphinxAtStartPar
Local ID of a neighbouring face in each
horizontally\sphinxhyphen{}adjacent cell indexed as (face).
\\
\sphinxbottomrule
\end{tabulary}
\sphinxtableafterendhook\par
\sphinxattableend\end{savenotes}


\subsubsection{gh\_shape and gh\_evaluator\_targets}
\label{\detokenize{dynamo0p3:gh-shape-and-gh-evaluator-targets}}\label{\detokenize{dynamo0p3:lfric-gh-shape}}
\sphinxAtStartPar
If a kernel requires basis or differential\sphinxhyphen{}basis functions then the
metadata must also specify the set of points on which these functions
are required. This information is provided by the \sphinxcode{\sphinxupquote{gh\_shape}}
component of the metadata. Currently PSyclone supports four shapes;
\sphinxcode{\sphinxupquote{gh\_quadrature\_XYoZ}} for Gaussian quadrature points,
\sphinxcode{\sphinxupquote{gh\_quadrature\_face}} for quadrature points on cell faces,
\sphinxcode{\sphinxupquote{gh\_quadrature\_edge}} for quadrature points on cell edges and
\sphinxcode{\sphinxupquote{gh\_evaluator}} for evaluation at nodal points. If a kernel requires
just one of these then \sphinxcode{\sphinxupquote{gh\_shape}} is an \sphinxcode{\sphinxupquote{integer}} scalar. However, if
more than one is required then \sphinxcode{\sphinxupquote{gh\_shape}} becomes a one\sphinxhyphen{}dimensional,
\sphinxcode{\sphinxupquote{integer}} array, e.g.:

\begin{sphinxVerbatim}[commandchars=\\\{\}]
\PYG{k+kt}{integer}\PYG{+w}{ }\PYG{k+kd}{::}\PYG{+w}{ }\PYG{n}{gh\PYGZus{}shape}\PYG{p}{(}\PYG{l+m+mi}{2}\PYG{p}{)}\PYG{+w}{ }\PYG{o}{=}\PYG{+w}{ }\PYG{p}{(}\PYG{o}{/}\PYG{+w}{ }\PYG{n}{gh\PYGZus{}quadrature\PYGZus{}face}\PYG{p}{,}\PYG{+w}{ }\PYG{n}{gh\PYGZus{}quadrature\PYGZus{}edge}\PYG{+w}{ }\PYG{o}{/}\PYG{p}{)}
\end{sphinxVerbatim}

\sphinxAtStartPar
If a kernel requires an evaluator then there are two options: if an
evaluator is required for multiple function spaces then these can be
specified using the additional \sphinxcode{\sphinxupquote{gh\_evaluator\_targets}} metadata
entry. This entry is a one\sphinxhyphen{}dimensional, \sphinxcode{\sphinxupquote{integer}} array containing the
desired function spaces. For example, to request
basis/differential\sphinxhyphen{}basis functions evaluated on both W0 and W1, the
metadata would be:

\begin{sphinxVerbatim}[commandchars=\\\{\}]
\PYG{k+kt}{integer}\PYG{+w}{ }\PYG{k+kd}{::}\PYG{+w}{ }\PYG{n}{gh\PYGZus{}shape}\PYG{+w}{ }\PYG{o}{=}\PYG{+w}{ }\PYG{n}{gh\PYGZus{}evaluator}
\PYG{k+kt}{integer}\PYG{+w}{ }\PYG{k+kd}{::}\PYG{+w}{ }\PYG{n}{gh\PYGZus{}evaluator\PYGZus{}targets}\PYG{p}{(}\PYG{l+m+mi}{2}\PYG{p}{)}\PYG{+w}{ }\PYG{o}{=}\PYG{+w}{ }\PYG{p}{(}\PYG{o}{/}\PYG{n}{W0}\PYG{p}{,}\PYG{+w}{ }\PYG{n}{W1}\PYG{o}{/}\PYG{p}{)}
\end{sphinxVerbatim}

\sphinxAtStartPar
The kernel must have an argument (field or operator) on each of the
function spaces listed in \sphinxcode{\sphinxupquote{gh\_evaluator\_targets}}.
The default behaviour if \sphinxcode{\sphinxupquote{gh\_evaluator\_targets}} is not specified is
to provide evaluators for each function space associated with the
quantities that the kernel is updating. All necessary data is
extracted in the PSy layer and passed to the kernel(s) as required \sphinxhyphen{}
nothing is required from the Algorithm layer. If a kernel requires
quadrature on the other hand, the Algorithm writer must supply a
\sphinxcode{\sphinxupquote{quadrature\_type}} object for each specified quadrature as the last
argument(s) to the kernel (see Section {\hyperref[\detokenize{dynamo0p3:lfric-quadrature}]{\sphinxcrossref{\DUrole{std,std-ref}{Quadrature}}}}).

\sphinxAtStartPar
Note that it is an error for kernel metadata to specify a value for
\sphinxcode{\sphinxupquote{gh\_shape}} if no basis or differential\sphinxhyphen{}basis functions are required.
It is also an error to specify \sphinxcode{\sphinxupquote{gh\_evaluator\_targets}} if the kernel
does not require an evaluator (i.e. \sphinxcode{\sphinxupquote{gh\_shape != gh\_evaluator}}).


\subsubsection{operates\_on}
\label{\detokenize{dynamo0p3:operates-on}}\label{\detokenize{dynamo0p3:lfric-operates-on}}
\sphinxAtStartPar
The fourth type of metadata provided is \sphinxcode{\sphinxupquote{OPERATES\_ON}}. This
specifies that the Kernel has been written with the assumption that it
is supplied with the specified data for each field/operator argument.
The possible values for \sphinxcode{\sphinxupquote{OPERATES\_ON}} and their interpretation are
summarised in the following table:


\begin{savenotes}\sphinxattablestart
\sphinxthistablewithglobalstyle
\centering
\begin{tabulary}{\linewidth}[t]{TT}
\sphinxtoprule
\sphinxstyletheadfamily 
\sphinxAtStartPar
operates\_on
&\sphinxstyletheadfamily 
\sphinxAtStartPar
Data passed for each field/operator argument
\\
\sphinxmidrule
\sphinxtableatstartofbodyhook
\sphinxAtStartPar
\sphinxcode{\sphinxupquote{cell\_column}}
&
\sphinxAtStartPar
Single column of cells from the owned region (except
when performing an INC operation on continuous fields
when it will include one level of halo cells).
\\
\sphinxhline
\sphinxAtStartPar
\sphinxcode{\sphinxupquote{halo\_cell\_column}}
&
\sphinxAtStartPar
Single column of cells exclusively from halo region.
\\
\sphinxhline
\sphinxAtStartPar
\sphinxcode{\sphinxupquote{owned\_and\_halo\_cell\_column}}
&
\sphinxAtStartPar
Single column of cells but iteration space will include
both owned and halo regions.
\\
\sphinxhline
\sphinxAtStartPar
\sphinxcode{\sphinxupquote{dof}}
&
\sphinxAtStartPar
Single DoF .
\\
\sphinxhline
\sphinxAtStartPar
\sphinxcode{\sphinxupquote{domain}}
&
\sphinxAtStartPar
All columns of cells in the (sub\sphinxhyphen{})domain.
\\
\sphinxbottomrule
\end{tabulary}
\sphinxtableafterendhook\par
\sphinxattableend\end{savenotes}

\sphinxAtStartPar
(For a description of the concepts of ‘owned’ and ‘halo’ cells please see the
\sphinxhref{https://psyclone-dev.readthedocs.io/en/latest/APIs.html\#lfric-developers}{LFRic}.)


\subsubsection{procedure}
\label{\detokenize{dynamo0p3:procedure}}
\sphinxAtStartPar
The fifth and final type of metadata is \sphinxcode{\sphinxupquote{procedure}} metadata. This
specifies the name of the Kernel subroutine that this metadata
describes.

\sphinxAtStartPar
For example:

\begin{sphinxVerbatim}[commandchars=\\\{\}]
\PYG{k}{procedure}\PYG{p}{,}\PYG{+w}{ }\PYG{k}{nopass}\PYG{+w}{ }\PYG{k+kd}{::}\PYG{+w}{ }\PYG{n}{my\PYGZus{}kernel\PYGZus{}subroutine}
\end{sphinxVerbatim}


\subsection{Subroutine}
\label{\detokenize{dynamo0p3:subroutine}}\label{\detokenize{dynamo0p3:lfric-kern-subroutine}}

\subsubsection{Rules for General\sphinxhyphen{}Purpose Kernels}
\label{\detokenize{dynamo0p3:rules-for-general-purpose-kernels}}\label{\detokenize{dynamo0p3:lfric-stub-generation-rules}}
\sphinxAtStartPar
The arguments to general\sphinxhyphen{}purpose kernels (those that do not involve
either CMA operators or prolongation/restriction operations) that
operate on cell\sphinxhyphen{}columns follow a set of rules
which have been specified for the LFRic API. These rules are encoded
in the \sphinxcode{\sphinxupquote{generate()}} method within the \sphinxcode{\sphinxupquote{ArgOrdering}} abstract class
in the \sphinxcode{\sphinxupquote{dynamo0p3.py}} file. The rules, along with PSyclone’s naming
conventions, are:
\begin{enumerate}
\sphinxsetlistlabels{\arabic}{enumi}{enumii}{}{)}%
\item {} 
\sphinxAtStartPar
If an LMA operator is passed then include the \sphinxcode{\sphinxupquote{cells}} argument.
\sphinxcode{\sphinxupquote{cells}} is an \sphinxcode{\sphinxupquote{integer}} of kind \sphinxcode{\sphinxupquote{i\_def}} and has intent \sphinxcode{\sphinxupquote{in}}.

\item {} 
\sphinxAtStartPar
Include \sphinxcode{\sphinxupquote{nlayers}}, the number of layers in a column. \sphinxcode{\sphinxupquote{nlayers}}
is an \sphinxcode{\sphinxupquote{integer}} of kind \sphinxcode{\sphinxupquote{i\_def}} and has intent \sphinxcode{\sphinxupquote{in}}. PSyclone
will obtain the value of \sphinxcode{\sphinxupquote{nlayers}} to use for a particular kernel
from the first field or operator in the argument list.

\item {} 
\sphinxAtStartPar
For each scalar/field/vector\_field/operator in the order specified by
the meta\_args metadata:
\begin{enumerate}
\sphinxsetlistlabels{\arabic}{enumii}{enumiii}{}{)}%
\item {} 
\sphinxAtStartPar
If the current entry is a scalar quantity then include the Fortran
variable in the argument list. The intent is determined from the
metadata (see {\hyperref[\detokenize{dynamo0p3:lfric-api-meta-args}]{\sphinxcrossref{\DUrole{std,std-ref}{meta\_args}}}} for an explanation).

\item {} 
\sphinxAtStartPar
If the current entry is a field then include the field
array. The field array name is currently specified as being
\sphinxcode{\sphinxupquote{"field\_"\textless{}argument\_position\textgreater{}"\_"\textless{}field\_function\_space\textgreater{}}}. A
field array is a rank\sphinxhyphen{}1, \sphinxcode{\sphinxupquote{real}} array with extent equal to the
number of unique degrees of freedom for the space that the field
is on. Its precision (kind) depends on how it is defined in the
algorithm layer, see the {\hyperref[\detokenize{dynamo0p3:lfric-mixed-precision}]{\sphinxcrossref{\DUrole{std,std-ref}{Mixed Precision}}}} section
for more details. This value is passed in separately. Again, the
intent is determined from the metadata (see
{\hyperref[\detokenize{dynamo0p3:lfric-api-meta-args}]{\sphinxcrossref{\DUrole{std,std-ref}{meta\_args}}}}).
\begin{enumerate}
\sphinxsetlistlabels{\arabic}{enumiii}{enumiv}{}{)}%
\item {} 
\sphinxAtStartPar
If the field entry has a stencil access then add an \sphinxcode{\sphinxupquote{integer}} (or
if the stencil is of type \sphinxcode{\sphinxupquote{CROSS2D}}, an \sphinxcode{\sphinxupquote{integer}} rank\sphinxhyphen{}1 array of
extent 4 and kind \sphinxcode{\sphinxupquote{i\_def}}) stencil\sphinxhyphen{}size argument with intent \sphinxcode{\sphinxupquote{in}}.
This will supply the number of cells in the stencil or, in the case
of the \sphinxcode{\sphinxupquote{CROSS2D}} stencil, the number of cells in each branch of
the stencil.

\item {} 
\sphinxAtStartPar
If the stencil is of type \sphinxcode{\sphinxupquote{CROSS2D}} then an \sphinxcode{\sphinxupquote{integer}} of kind
\sphinxcode{\sphinxupquote{i\_def}} and intent \sphinxcode{\sphinxupquote{in}} for the max branch length is needed.
This is used in defining the dimensions of the stencil dofmap array
and is required due to the varying length of the branches of the
stencil when used on planar meshes.

\item {} 
\sphinxAtStartPar
Also needed is a stencil dofmap array of type \sphinxcode{\sphinxupquote{integer}}, kind
\sphinxcode{\sphinxupquote{i\_def}} and intent \sphinxcode{\sphinxupquote{in}} in either 2 or 3 dimensions. For a
\sphinxcode{\sphinxupquote{CROSS2D}} stencil the array needs dimensions of
(number\sphinxhyphen{}of\sphinxhyphen{}dofs\sphinxhyphen{}in\sphinxhyphen{}cell, max\sphinxhyphen{}branch\sphinxhyphen{}length, 4).
All other stencils need dimensions of (number\sphinxhyphen{}of\sphinxhyphen{}dofs\sphinxhyphen{}in\sphinxhyphen{}cell,
stencil\sphinxhyphen{}size).

\item {} 
\sphinxAtStartPar
If the field entry stencil access is of type \sphinxcode{\sphinxupquote{XORY1D}} then
add an additional \sphinxcode{\sphinxupquote{integer}} direction argument of kind
\sphinxcode{\sphinxupquote{i\_def}} and with intent \sphinxcode{\sphinxupquote{in}}.

\end{enumerate}

\item {} 
\sphinxAtStartPar
If the current entry is a field vector then for each dimension
of the vector, include a field array. The field array name is
specified as being using
\sphinxcode{\sphinxupquote{"field\_"\textless{}argument\_position\textgreater{}"\_"\textless{}field\_function\_space\textgreater{}"\_v"\textless{}vector\_position\textgreater{}}}.
A field array in a field vector is declared in the same way as a
field array (described in the previous step).

\item {} 
\sphinxAtStartPar
If the current entry is an operator then first include an
\sphinxcode{\sphinxupquote{integer}} extent of kind \sphinxcode{\sphinxupquote{i\_def}}. The name of this extent is
\sphinxcode{\sphinxupquote{\textless{}operator\_name\textgreater{}"\_ncell\_3d"}}. Next include the operator.  This
is a rank\sphinxhyphen{}3, \sphinxcode{\sphinxupquote{real}} array. Its precision (kind) depends on how
it is defined in the algorithm layer, see the
{\hyperref[\detokenize{dynamo0p3:lfric-mixed-precision}]{\sphinxcrossref{\DUrole{std,std-ref}{Mixed Precision}}}} section for more details. The
extent of the first dimension is \sphinxcode{\sphinxupquote{\textless{}operator\_name\textgreater{}"\_ncell\_3d"}},
and of the second and third dimension are the local degrees of
freedom for the \sphinxcode{\sphinxupquote{to}} and \sphinxcode{\sphinxupquote{from}} function spaces,
respectively. Again the intent is determined
from the metadata (see {\hyperref[\detokenize{dynamo0p3:lfric-api-meta-args}]{\sphinxcrossref{\DUrole{std,std-ref}{meta\_args}}}}).

\end{enumerate}

\item {} 
\sphinxAtStartPar
For each function space in the order they appear in the metadata arguments
(the \sphinxcode{\sphinxupquote{to}} function space of an operator is considered to be before the
\sphinxcode{\sphinxupquote{from}} function space of the same operator as it appears first in
lexicographic order)
\begin{enumerate}
\sphinxsetlistlabels{\arabic}{enumii}{enumiii}{}{)}%
\item {} 
\sphinxAtStartPar
Include the number of local degrees of freedom (i.e. number per\sphinxhyphen{}cell)
for the function space. This is an \sphinxcode{\sphinxupquote{integer}} of kind \sphinxcode{\sphinxupquote{i\_def}} and
has intent \sphinxcode{\sphinxupquote{in}}. The name of this argument is
\sphinxcode{\sphinxupquote{"ndf\_"\textless{}field\_function\_space\textgreater{}}}.

\item {} 
\sphinxAtStartPar
If there is a field on this space
\begin{enumerate}
\sphinxsetlistlabels{\arabic}{enumiii}{enumiv}{}{)}%
\item {} 
\sphinxAtStartPar
Include the unique number of degrees of freedom for the function
space. This is an \sphinxcode{\sphinxupquote{integer}} of kind \sphinxcode{\sphinxupquote{i\_def}} and has intent \sphinxcode{\sphinxupquote{in}}.
The name of this argument is \sphinxcode{\sphinxupquote{"undf\_"\textless{}field\_function\_space\textgreater{}}}.

\item {} 
\sphinxAtStartPar
Include the \sphinxstylestrong{dofmap} for this function space. This is an \sphinxcode{\sphinxupquote{integer}}
array of kind \sphinxcode{\sphinxupquote{i\_def}} with intent \sphinxcode{\sphinxupquote{in}}. It has one dimension
sized by the local degrees of freedom for the function space.

\end{enumerate}

\item {} 
\sphinxAtStartPar
For each operation on the function space (\sphinxcode{\sphinxupquote{basis}}, \sphinxcode{\sphinxupquote{diff\_basis}}),
in the order specified in the metadata, pass \sphinxcode{\sphinxupquote{real}} arrays of kind
\sphinxcode{\sphinxupquote{r\_def}} with intent \sphinxcode{\sphinxupquote{in}}. For each shape specified in the
\sphinxcode{\sphinxupquote{gh\_shape}} metadata entry:
\begin{enumerate}
\sphinxsetlistlabels{\arabic}{enumiii}{enumiv}{}{)}%
\item {} 
\sphinxAtStartPar
If shape is \sphinxcode{\sphinxupquote{gh\_quadrature\_*}} then the arrays are of rank four
and are named
\sphinxcode{\sphinxupquote{"basis\_"\textless{}field\_function\_space\textgreater{}\_\textless{}quadrature\_arg\_name\textgreater{}}} or
\sphinxcode{\sphinxupquote{"diff\_basis\_"\textless{}field\_function\_space\textgreater{}\_\textless{}quadrature\_arg\_name\textgreater{}}},
as appropriate:
\begin{enumerate}
\sphinxsetlistlabels{\arabic}{enumiv}{enumv}{}{)}%
\item {} 
\sphinxAtStartPar
If shape is \sphinxcode{\sphinxupquote{gh\_quadrature\_xyoz}} then the arrays have extent
(\sphinxcode{\sphinxupquote{dimension}}, \sphinxcode{\sphinxupquote{number\_of\_dofs}}, \sphinxcode{\sphinxupquote{np\_xy}}, \sphinxcode{\sphinxupquote{np\_z}}).

\item {} 
\sphinxAtStartPar
If shape is \sphinxcode{\sphinxupquote{gh\_quadrature\_face}} or \sphinxcode{\sphinxupquote{gh\_quadrature\_edge}}
then the  arrays have extent
(\sphinxcode{\sphinxupquote{dimension}}, \sphinxcode{\sphinxupquote{number\_of\_dofs}}, \sphinxcode{\sphinxupquote{np\_xyz}}, \sphinxcode{\sphinxupquote{nfaces}} or
\sphinxcode{\sphinxupquote{nedges}}).

\end{enumerate}

\item {} 
\sphinxAtStartPar
If shape is \sphinxcode{\sphinxupquote{gh\_evaluator}} then we pass one array for
each target function space (i.e. as specified by
\sphinxcode{\sphinxupquote{gh\_evaluator\_targets}}). Each of these arrays are of rank three
with extent (\sphinxcode{\sphinxupquote{dimension}}, \sphinxcode{\sphinxupquote{number\_of\_dofs}},
\sphinxcode{\sphinxupquote{ndf\_\textless{}target\_function\_space\textgreater{}}}). The name of the argument is
\sphinxcode{\sphinxupquote{"basis\_"\textless{}field\_function\_space\textgreater{}"\_on\_"\textless{}target\_function\_space\textgreater{}}} or
\sphinxcode{\sphinxupquote{"diff\_basis\_"\textless{}field\_function\_space\textgreater{}"\_on\_"\textless{}target\_function\_space\textgreater{}}},
as appropriate.

\end{enumerate}

\sphinxAtStartPar
Here \sphinxcode{\sphinxupquote{\textless{}quadrature\_arg\_name\textgreater{}}} is the name of the corresponding
quadrature object being passed to the Invoke.
\sphinxcode{\sphinxupquote{dimension}} is 1 or 3 and depends upon the function space
(see {\hyperref[\detokenize{dynamo0p3:lfric-function-space}]{\sphinxcrossref{\DUrole{std,std-ref}{Supported Function Spaces}}}} above for more information) and
whether or not it is a basis or a differential basis function (see
the table below). \sphinxcode{\sphinxupquote{number\_of\_dofs}} is the number of degrees of
freedom (DoFs) associated with the function space and \sphinxcode{\sphinxupquote{np\_*}} are
the number of points to be evaluated: i) \sphinxcode{\sphinxupquote{*\_xyz}} in
all directions (3D); ii) \sphinxcode{\sphinxupquote{*\_xy}} in the horizontal plane (2D);
iii) \sphinxcode{\sphinxupquote{*\_x, *\_y}} in the horizontal (1D); and iv) \sphinxcode{\sphinxupquote{*\_z}} in the
vertical (1D). \sphinxcode{\sphinxupquote{nfaces}} and \sphinxcode{\sphinxupquote{nedges}} are the number of horizontal
faces/edges obtained from the appropriate quadrature object supplied
to the Invoke.


\begin{savenotes}\sphinxattablestart
\sphinxthistablewithglobalstyle
\sphinxthistablewithvlinesstyle
\centering
\begin{tabulary}{\linewidth}[t]{|l|c|l|}
\sphinxtoprule
\sphinxstyletheadfamily 
\sphinxAtStartPar
Function Type
&\sphinxstyletheadfamily 
\sphinxAtStartPar
Dimension
&\sphinxstyletheadfamily 
\sphinxAtStartPar
Function Space Name
\\
\sphinxmidrule
\sphinxtableatstartofbodyhook\sphinxmultirow{2}{4}{%
\begin{varwidth}[t]{\sphinxcolwidth{1}{3}}
\sphinxAtStartPar
Basis
\par
\vskip-\baselineskip\vbox{\hbox{\strut}}\end{varwidth}%
}%
&
\sphinxAtStartPar
1
&
\sphinxAtStartPar
W0, W2trace, W2Htrace, W2Vtrace,
W3, Wtheta, Wchi
\\
\sphinxcline{2-3}\sphinxfixclines{3}\sphinxtablestrut{4}&
\sphinxAtStartPar
3
&
\sphinxAtStartPar
W1, W2, W2H, W2V, W2broken, ANY\_W2
\\
\sphinxhline\sphinxmultirow{2}{9}{%
\begin{varwidth}[t]{\sphinxcolwidth{1}{3}}
\sphinxAtStartPar
Differential
Basis
\par
\vskip-\baselineskip\vbox{\hbox{\strut}}\end{varwidth}%
}%
&
\sphinxAtStartPar
1
&
\sphinxAtStartPar
W2, W2H, W2V, W2broken, ANY\_W2
\\
\sphinxcline{2-3}\sphinxfixclines{3}\sphinxtablestrut{9}&
\sphinxAtStartPar
3
&
\sphinxAtStartPar
W0, W1, W2trace, W2Htrace,
W2Vtrace, W3, Wtheta, Wchi
\\
\sphinxbottomrule
\end{tabulary}
\sphinxtableafterendhook\par
\sphinxattableend\end{savenotes}

\end{enumerate}

\item {} 
\sphinxAtStartPar
If either the \sphinxcode{\sphinxupquote{normals\_to\_horizontal\_faces}} or
\sphinxcode{\sphinxupquote{outward\_normals\_to\_horizontal\_faces}} properties of the reference
element are required then pass the number of horizontal faces of the
reference element (\sphinxcode{\sphinxupquote{nfaces\_re\_h}}). Similarly, if either the
\sphinxcode{\sphinxupquote{normals\_to\_vertical\_faces}} or \sphinxcode{\sphinxupquote{outward\_normals\_to\_vertical\_faces}} are
required then pass the number of vertical faces (\sphinxcode{\sphinxupquote{nfaces\_re\_v}}). This
also holds for the \sphinxcode{\sphinxupquote{normals\_to\_faces}} and \sphinxcode{\sphinxupquote{outward\_normals\_to\_faces}}
where the number of all faces of the reference element (\sphinxcode{\sphinxupquote{nfaces\_re}})
is passed to the kernel. (All of these quantities are integers of kind
\sphinxcode{\sphinxupquote{i\_def}}.) Then, in the order specified in the
\sphinxcode{\sphinxupquote{meta\_reference\_element}} metadata:
\begin{enumerate}
\sphinxsetlistlabels{\arabic}{enumii}{enumiii}{}{)}%
\item {} 
\sphinxAtStartPar
For the \sphinxcode{\sphinxupquote{normals\_to\_horizontal/vertical\_faces}}, pass a rank\sphinxhyphen{}2
\sphinxcode{\sphinxupquote{integer}} array of kind \sphinxcode{\sphinxupquote{i\_def}} with dimensions
\sphinxcode{\sphinxupquote{(3, nfaces\_re\_h/v)}}.

\item {} 
\sphinxAtStartPar
For the \sphinxcode{\sphinxupquote{outward\_normals\_to\_horizontal/vertical\_faces}}, pass a rank\sphinxhyphen{}2
\sphinxcode{\sphinxupquote{integer}} array of kind \sphinxcode{\sphinxupquote{i\_def}} with dimensions
\sphinxcode{\sphinxupquote{(3, nfaces\_re\_h/v)}}.

\item {} 
\sphinxAtStartPar
For \sphinxcode{\sphinxupquote{normals\_to\_faces}} or \sphinxcode{\sphinxupquote{outward\_normals\_to\_faces}} pass
a rank\sphinxhyphen{}2 \sphinxcode{\sphinxupquote{integer}} array of kind \sphinxcode{\sphinxupquote{i\_def}} with dimensions
\sphinxcode{\sphinxupquote{(3, nfaces\_re)}}.

\end{enumerate}

\item {} 
\sphinxAtStartPar
If the \sphinxcode{\sphinxupquote{adjacent\_face}} mesh property is required then:
\begin{enumerate}
\sphinxsetlistlabels{\arabic}{enumii}{enumiii}{}{)}%
\item {} 
\sphinxAtStartPar
If the number of horizontal cell faces obtained from the reference
element (\sphinxcode{\sphinxupquote{nfaces\_re\_h}}) is not already being passed to the kernel (due
to rule 5 above) then supply it here. This is an \sphinxcode{\sphinxupquote{integer}} of kind
\sphinxcode{\sphinxupquote{i\_def}}.

\item {} 
\sphinxAtStartPar
Pass a rank\sphinxhyphen{}1, \sphinxcode{\sphinxupquote{integer}} array of kind \sphinxcode{\sphinxupquote{i\_def}} and extent
\sphinxcode{\sphinxupquote{nfaces\_re\_h}}.

\end{enumerate}

\item {} 
\sphinxAtStartPar
If Quadrature is required (\sphinxcode{\sphinxupquote{gh\_shape = gh\_quadrature\_*}}) then, for
each shape in the order specified in the \sphinxcode{\sphinxupquote{gh\_shape}} metadata:
\begin{enumerate}
\sphinxsetlistlabels{\arabic}{enumii}{enumiii}{}{)}%
\item {} 
\sphinxAtStartPar
Include \sphinxcode{\sphinxupquote{integer}}, scalar arguments of kind \sphinxcode{\sphinxupquote{i\_def}} with intent
\sphinxcode{\sphinxupquote{in}} that specify the extent of the basis/diff\sphinxhyphen{}basis arrays:
\begin{enumerate}
\sphinxsetlistlabels{\arabic}{enumiii}{enumiv}{}{)}%
\item {} 
\sphinxAtStartPar
If \sphinxcode{\sphinxupquote{gh\_shape}} is \sphinxcode{\sphinxupquote{gh\_quadrature\_XYoZ}} then pass
\sphinxcode{\sphinxupquote{np\_xy\_\textless{}quadrature\_arg\_name\textgreater{}}} and \sphinxcode{\sphinxupquote{np\_z\_\textless{}quadrature\_arg\_name\textgreater{}}}.

\item {} 
\sphinxAtStartPar
If \sphinxcode{\sphinxupquote{gh\_shape}} is \sphinxcode{\sphinxupquote{gh\_quadrature\_face}}/\sphinxcode{\sphinxupquote{\_edge}} then pass
\sphinxcode{\sphinxupquote{nfaces}}/\sphinxcode{\sphinxupquote{nedges\_\textless{}quadrature\_arg\_name\textgreater{}}} and
\sphinxcode{\sphinxupquote{np\_xyz\_\textless{}quadrature\_arg\_name\textgreater{}}}.

\end{enumerate}

\item {} 
\sphinxAtStartPar
Include weights which are \sphinxcode{\sphinxupquote{real}} arrays of kind \sphinxcode{\sphinxupquote{r\_def}}:
\begin{enumerate}
\sphinxsetlistlabels{\arabic}{enumiii}{enumiv}{}{)}%
\item {} 
\sphinxAtStartPar
If \sphinxcode{\sphinxupquote{gh\_quadrature\_XYoZ}} pass in
\sphinxcode{\sphinxupquote{weights\_xz\_\textless{}quadrature\_arg\_name\textgreater{}}} (rank one, extent
\sphinxcode{\sphinxupquote{np\_xy\_\textless{}quadrature\_arg\_name\textgreater{}}})
and \sphinxcode{\sphinxupquote{weights\_z\_\textless{}quadrature\_arg\_name\textgreater{}}} (rank one, extent
\sphinxcode{\sphinxupquote{np\_z\_\textless{}quadrature\_arg\_name\textgreater{}}}).

\item {} 
\sphinxAtStartPar
If \sphinxcode{\sphinxupquote{gh\_quadrature\_face}}/\sphinxcode{\sphinxupquote{\_edge}} pass in
\sphinxcode{\sphinxupquote{weights\_xyz\_\textless{}quadrature\_arg\_name\textgreater{}}} (rank two with extents
{[}\sphinxcode{\sphinxupquote{np\_xyz\_\textless{}quadrature\_arg\_name\textgreater{}}},
\sphinxcode{\sphinxupquote{nfaces/nedges\_\textless{}quadrature\_arg\_name\textgreater{}}}{]}).

\end{enumerate}

\end{enumerate}

\end{enumerate}


\paragraph{Examples}
\label{\detokenize{dynamo0p3:examples}}
\sphinxAtStartPar
For instance, if a kernel has only one written argument and requires an
evaluator then its metadata might be:

\begin{sphinxVerbatim}[commandchars=\\\{\}]
\PYG{k}{type}\PYG{p}{,}\PYG{+w}{ }\PYG{k}{extends}\PYG{p}{(}\PYG{n}{kernel\PYGZus{}type}\PYG{p}{)}\PYG{+w}{ }\PYG{k+kd}{::}\PYG{+w}{ }\PYG{n}{testkern\PYGZus{}operator\PYGZus{}type}
\PYG{+w}{   }\PYG{k}{type}\PYG{p}{(}\PYG{n}{arg\PYGZus{}type}\PYG{p}{)}\PYG{p}{,}\PYG{+w}{ }\PYG{k}{dimension}\PYG{p}{(}\PYG{l+m+mi}{2}\PYG{p}{)}\PYG{+w}{ }\PYG{k+kd}{::}\PYG{+w}{ }\PYG{n}{meta\PYGZus{}args}\PYG{+w}{ }\PYG{o}{=}\PYG{+w}{               }\PYG{p}{\PYGZam{}}
\PYG{+w}{        }\PYG{p}{(}\PYG{o}{/}\PYG{+w}{ }\PYG{n}{arg\PYGZus{}type}\PYG{p}{(}\PYG{n}{gh\PYGZus{}operator}\PYG{p}{,}\PYG{+w}{ }\PYG{n}{gh\PYGZus{}real}\PYG{p}{,}\PYG{+w}{ }\PYG{n}{gh\PYGZus{}write}\PYG{p}{,}\PYG{+w}{ }\PYG{n}{w0}\PYG{p}{,}\PYG{+w}{ }\PYG{n}{w1}\PYG{p}{)}\PYG{p}{,}\PYG{+w}{ }\PYG{p}{\PYGZam{}}
\PYG{+w}{           }\PYG{n}{arg\PYGZus{}type}\PYG{p}{(}\PYG{n}{gh\PYGZus{}field}\PYG{o}{*}\PYG{l+m+mi}{3}\PYG{p}{,}\PYG{+w}{  }\PYG{n}{gh\PYGZus{}real}\PYG{p}{,}\PYG{+w}{ }\PYG{n}{gh\PYGZus{}read}\PYG{p}{,}\PYG{+w}{  }\PYG{n}{w0}\PYG{p}{)}\PYG{+w}{ }\PYG{o}{/}\PYG{p}{)}
\PYG{+w}{   }\PYG{k}{type}\PYG{p}{(}\PYG{n}{func\PYGZus{}type}\PYG{p}{)}\PYG{+w}{ }\PYG{k+kd}{::}\PYG{+w}{ }\PYG{n}{meta\PYGZus{}funcs}\PYG{p}{(}\PYG{l+m+mi}{1}\PYG{p}{)}\PYG{+w}{ }\PYG{o}{=}\PYG{+w}{                        }\PYG{p}{\PYGZam{}}
\PYG{+w}{        }\PYG{p}{(}\PYG{o}{/}\PYG{+w}{ }\PYG{n}{func\PYGZus{}type}\PYG{p}{(}\PYG{n}{w0}\PYG{p}{,}\PYG{+w}{ }\PYG{n}{gh\PYGZus{}basis}\PYG{p}{)}\PYG{+w}{ }\PYG{o}{/}\PYG{p}{)}
\PYG{+w}{   }\PYG{k+kt}{integer}\PYG{+w}{ }\PYG{k+kd}{::}\PYG{+w}{ }\PYG{n}{operates\PYGZus{}on}\PYG{+w}{ }\PYG{o}{=}\PYG{+w}{ }\PYG{n}{cell\PYGZus{}column}
\PYG{+w}{   }\PYG{k+kt}{integer}\PYG{+w}{ }\PYG{k+kd}{::}\PYG{+w}{ }\PYG{n}{gh\PYGZus{}shape}\PYG{+w}{ }\PYG{o}{=}\PYG{+w}{ }\PYG{n}{gh\PYGZus{}evaluator}
\PYG{+w}{ }\PYG{k}{contains}
\PYG{k}{   }\PYG{k}{procedure}\PYG{p}{,}\PYG{+w}{ }\PYG{k}{nopass}\PYG{+w}{ }\PYG{k+kd}{::}\PYG{+w}{ }\PYG{n}{code}\PYG{+w}{ }\PYG{o}{=}\PYG{o}{\PYGZgt{}}\PYG{+w}{ }\PYG{n}{testkern\PYGZus{}operator\PYGZus{}code}
\PYG{k}{end }\PYG{k}{type }\PYG{n}{testkern\PYGZus{}operator\PYGZus{}type}
\end{sphinxVerbatim}

\sphinxAtStartPar
then we only pass the basis functions evaluated on \sphinxcode{\sphinxupquote{W0}} (the space of
the written kernel argument). The subroutine arguments will therefore
be:

\begin{sphinxVerbatim}[commandchars=\\\{\}]
\PYG{k}{subroutine }\PYG{n}{testkern\PYGZus{}operator\PYGZus{}code}\PYG{p}{(}\PYG{n}{cell}\PYG{p}{,}\PYG{+w}{ }\PYG{n}{nlayers}\PYG{p}{,}\PYG{+w}{ }\PYG{n}{ncell\PYGZus{}3d}\PYG{p}{,}\PYG{+w}{        }\PYG{p}{\PYGZam{}}
\PYG{+w}{     }\PYG{n}{local\PYGZus{}stencil}\PYG{p}{,}\PYG{+w}{ }\PYG{n}{xdata}\PYG{p}{,}\PYG{+w}{ }\PYG{n}{ydata}\PYG{p}{,}\PYG{+w}{ }\PYG{n}{zdata}\PYG{p}{,}\PYG{+w}{ }\PYG{n}{ndf\PYGZus{}w0}\PYG{p}{,}\PYG{+w}{ }\PYG{n}{undf\PYGZus{}w0}\PYG{p}{,}\PYG{+w}{ }\PYG{n}{map\PYGZus{}w0}\PYG{p}{,}\PYG{+w}{ }\PYG{p}{\PYGZam{}}
\PYG{+w}{     }\PYG{n}{basis\PYGZus{}w0\PYGZus{}on\PYGZus{}w0}\PYG{p}{,}\PYG{+w}{ }\PYG{n}{ndf\PYGZus{}w1}\PYG{p}{)}
\end{sphinxVerbatim}

\sphinxAtStartPar
where \sphinxcode{\sphinxupquote{local\_stencil}} is the operator, \sphinxcode{\sphinxupquote{xdata}}, \sphinxcode{\sphinxupquote{ydata}}
etc. are the three components of the field vector and \sphinxcode{\sphinxupquote{map\_w0}} is
the dofmap for the \sphinxcode{\sphinxupquote{W0}} function space.

\sphinxAtStartPar
If instead, \sphinxcode{\sphinxupquote{gh\_evaluator\_targets}} is specified in the metadata:

\begin{sphinxVerbatim}[commandchars=\\\{\}]
\PYG{k}{type}\PYG{p}{,}\PYG{+w}{ }\PYG{k}{extends}\PYG{p}{(}\PYG{n}{kernel\PYGZus{}type}\PYG{p}{)}\PYG{+w}{ }\PYG{k+kd}{::}\PYG{+w}{ }\PYG{n}{testkern\PYGZus{}operator\PYGZus{}type}
\PYG{+w}{   }\PYG{k}{type}\PYG{p}{(}\PYG{n}{arg\PYGZus{}type}\PYG{p}{)}\PYG{p}{,}\PYG{+w}{ }\PYG{k}{dimension}\PYG{p}{(}\PYG{l+m+mi}{2}\PYG{p}{)}\PYG{+w}{ }\PYG{k+kd}{::}\PYG{+w}{ }\PYG{n}{meta\PYGZus{}args}\PYG{+w}{ }\PYG{o}{=}\PYG{+w}{               }\PYG{p}{\PYGZam{}}
\PYG{+w}{        }\PYG{p}{(}\PYG{o}{/}\PYG{+w}{ }\PYG{n}{arg\PYGZus{}type}\PYG{p}{(}\PYG{n}{gh\PYGZus{}operator}\PYG{p}{,}\PYG{+w}{ }\PYG{n}{gh\PYGZus{}real}\PYG{p}{,}\PYG{+w}{ }\PYG{n}{gh\PYGZus{}write}\PYG{p}{,}\PYG{+w}{ }\PYG{n}{w0}\PYG{p}{,}\PYG{+w}{ }\PYG{n}{w1}\PYG{p}{)}\PYG{p}{,}\PYG{+w}{ }\PYG{p}{\PYGZam{}}
\PYG{+w}{           }\PYG{n}{arg\PYGZus{}type}\PYG{p}{(}\PYG{n}{gh\PYGZus{}field}\PYG{o}{*}\PYG{l+m+mi}{3}\PYG{p}{,}\PYG{+w}{  }\PYG{n}{gh\PYGZus{}real}\PYG{p}{,}\PYG{+w}{ }\PYG{n}{gh\PYGZus{}read}\PYG{p}{,}\PYG{+w}{  }\PYG{n}{w0}\PYG{p}{)}\PYG{+w}{ }\PYG{o}{/}\PYG{p}{)}
\PYG{+w}{   }\PYG{k}{type}\PYG{p}{(}\PYG{n}{func\PYGZus{}type}\PYG{p}{)}\PYG{+w}{ }\PYG{k+kd}{::}\PYG{+w}{ }\PYG{n}{meta\PYGZus{}funcs}\PYG{p}{(}\PYG{l+m+mi}{1}\PYG{p}{)}\PYG{+w}{ }\PYG{o}{=}\PYG{+w}{               }\PYG{p}{\PYGZam{}}
\PYG{+w}{        }\PYG{p}{(}\PYG{o}{/}\PYG{+w}{ }\PYG{n}{func\PYGZus{}type}\PYG{p}{(}\PYG{n}{w0}\PYG{p}{,}\PYG{+w}{ }\PYG{n}{gh\PYGZus{}basis}\PYG{p}{)}\PYG{+w}{ }\PYG{o}{/}\PYG{p}{)}
\PYG{+w}{   }\PYG{k+kt}{integer}\PYG{+w}{ }\PYG{k+kd}{::}\PYG{+w}{ }\PYG{n}{operates\PYGZus{}on}\PYG{+w}{ }\PYG{o}{=}\PYG{+w}{ }\PYG{n}{cell\PYGZus{}column}
\PYG{+w}{   }\PYG{k+kt}{integer}\PYG{+w}{ }\PYG{k+kd}{::}\PYG{+w}{ }\PYG{n}{gh\PYGZus{}shape}\PYG{+w}{ }\PYG{o}{=}\PYG{+w}{ }\PYG{n}{gh\PYGZus{}evaluator}
\PYG{+w}{   }\PYG{k+kt}{integer}\PYG{+w}{ }\PYG{k+kd}{::}\PYG{+w}{ }\PYG{n}{gh\PYGZus{}evaluator\PYGZus{}targets}\PYG{p}{(}\PYG{l+m+mi}{2}\PYG{p}{)}\PYG{+w}{ }\PYG{o}{=}\PYG{+w}{ }\PYG{p}{(}\PYG{o}{/}\PYG{n}{W0}\PYG{p}{,}\PYG{+w}{ }\PYG{n}{W1}\PYG{o}{/}\PYG{p}{)}
\PYG{+w}{ }\PYG{k}{contains}
\PYG{k}{   }\PYG{k}{procedure}\PYG{p}{,}\PYG{+w}{ }\PYG{k}{nopass}\PYG{+w}{ }\PYG{k+kd}{::}\PYG{+w}{ }\PYG{n}{code}\PYG{+w}{ }\PYG{o}{=}\PYG{o}{\PYGZgt{}}\PYG{+w}{ }\PYG{n}{testkern\PYGZus{}operator\PYGZus{}code}
\PYG{k}{end }\PYG{k}{type }\PYG{n}{testkern\PYGZus{}operator\PYGZus{}type}
\end{sphinxVerbatim}

\sphinxAtStartPar
then we will need to pass two sets of basis functions (evaluated at \sphinxcode{\sphinxupquote{W0}}
and at \sphinxcode{\sphinxupquote{W1}}):

\begin{sphinxVerbatim}[commandchars=\\\{\}]
\PYG{k}{subroutine }\PYG{n}{testkern\PYGZus{}operator\PYGZus{}code}\PYG{p}{(}\PYG{n}{cell}\PYG{p}{,}\PYG{+w}{ }\PYG{n}{nlayers}\PYG{p}{,}\PYG{+w}{ }\PYG{n}{ncell\PYGZus{}3d}\PYG{p}{,}\PYG{+w}{        }\PYG{p}{\PYGZam{}}
\PYG{+w}{     }\PYG{n}{local\PYGZus{}stencil}\PYG{p}{,}\PYG{+w}{ }\PYG{n}{xdata}\PYG{p}{,}\PYG{+w}{ }\PYG{n}{ydata}\PYG{p}{,}\PYG{+w}{ }\PYG{n}{zdata}\PYG{p}{,}\PYG{+w}{ }\PYG{n}{ndf\PYGZus{}w0}\PYG{p}{,}\PYG{+w}{ }\PYG{n}{undf\PYGZus{}w0}\PYG{p}{,}\PYG{+w}{ }\PYG{n}{map\PYGZus{}w0}\PYG{p}{,}\PYG{+w}{ }\PYG{p}{\PYGZam{}}
\PYG{+w}{     }\PYG{n}{basis\PYGZus{}w0\PYGZus{}on\PYGZus{}w0}\PYG{p}{,}\PYG{+w}{ }\PYG{n}{basis\PYGZus{}w0\PYGZus{}on\PYGZus{}w1}\PYG{p}{,}\PYG{+w}{ }\PYG{n}{ndf\PYGZus{}w1}\PYG{p}{)}
\end{sphinxVerbatim}

\sphinxAtStartPar
If the metadata specifies that a kernel requires both an evaluator
and quadrature:

\begin{sphinxVerbatim}[commandchars=\\\{\}]
\PYG{k}{type}\PYG{p}{,}\PYG{+w}{ }\PYG{k}{extends}\PYG{p}{(}\PYG{n}{kernel\PYGZus{}type}\PYG{p}{)}\PYG{+w}{ }\PYG{k+kd}{::}\PYG{+w}{ }\PYG{n}{testkern\PYGZus{}operator\PYGZus{}type}
\PYG{+w}{   }\PYG{k}{type}\PYG{p}{(}\PYG{n}{arg\PYGZus{}type}\PYG{p}{)}\PYG{p}{,}\PYG{+w}{ }\PYG{k}{dimension}\PYG{p}{(}\PYG{l+m+mi}{2}\PYG{p}{)}\PYG{+w}{ }\PYG{k+kd}{::}\PYG{+w}{ }\PYG{n}{meta\PYGZus{}args}\PYG{+w}{ }\PYG{o}{=}\PYG{+w}{               }\PYG{p}{\PYGZam{}}
\PYG{+w}{        }\PYG{p}{(}\PYG{o}{/}\PYG{+w}{ }\PYG{n}{arg\PYGZus{}type}\PYG{p}{(}\PYG{n}{gh\PYGZus{}operator}\PYG{p}{,}\PYG{+w}{ }\PYG{n}{gh\PYGZus{}real}\PYG{p}{,}\PYG{+w}{ }\PYG{n}{gh\PYGZus{}write}\PYG{p}{,}\PYG{+w}{ }\PYG{n}{w0}\PYG{p}{,}\PYG{+w}{ }\PYG{n}{w1}\PYG{p}{)}\PYG{p}{,}\PYG{+w}{ }\PYG{p}{\PYGZam{}}
\PYG{+w}{           }\PYG{n}{arg\PYGZus{}type}\PYG{p}{(}\PYG{n}{gh\PYGZus{}field}\PYG{o}{*}\PYG{l+m+mi}{3}\PYG{p}{,}\PYG{+w}{  }\PYG{n}{gh\PYGZus{}real}\PYG{p}{,}\PYG{+w}{ }\PYG{n}{gh\PYGZus{}read}\PYG{p}{,}\PYG{+w}{  }\PYG{n}{w0}\PYG{p}{)}\PYG{+w}{ }\PYG{o}{/}\PYG{p}{)}
\PYG{+w}{   }\PYG{k}{type}\PYG{p}{(}\PYG{n}{func\PYGZus{}type}\PYG{p}{)}\PYG{+w}{ }\PYG{k+kd}{::}\PYG{+w}{ }\PYG{n}{meta\PYGZus{}funcs}\PYG{p}{(}\PYG{l+m+mi}{1}\PYG{p}{)}\PYG{+w}{ }\PYG{o}{=}\PYG{+w}{                        }\PYG{p}{\PYGZam{}}
\PYG{+w}{        }\PYG{p}{(}\PYG{o}{/}\PYG{+w}{ }\PYG{n}{func\PYGZus{}type}\PYG{p}{(}\PYG{n}{w0}\PYG{p}{,}\PYG{+w}{ }\PYG{n}{gh\PYGZus{}basis}\PYG{p}{)}\PYG{+w}{ }\PYG{o}{/}\PYG{p}{)}
\PYG{+w}{   }\PYG{k+kt}{integer}\PYG{+w}{ }\PYG{k+kd}{::}\PYG{+w}{ }\PYG{n}{operates\PYGZus{}on}\PYG{+w}{ }\PYG{o}{=}\PYG{+w}{ }\PYG{n}{cell\PYGZus{}column}
\PYG{+w}{   }\PYG{k+kt}{integer}\PYG{+w}{ }\PYG{k+kd}{::}\PYG{+w}{ }\PYG{n}{gh\PYGZus{}shape}\PYG{p}{(}\PYG{l+m+mi}{2}\PYG{p}{)}\PYG{+w}{ }\PYG{o}{=}\PYG{+w}{ }\PYG{p}{(}\PYG{o}{/}\PYG{+w}{ }\PYG{n}{gh\PYGZus{}evaluator}\PYG{p}{,}\PYG{+w}{ }\PYG{n}{gh\PYGZus{}quadrature\PYGZus{}face}\PYG{+w}{ }\PYG{o}{/}\PYG{p}{)}
\PYG{+w}{ }\PYG{k}{contains}
\PYG{k}{   }\PYG{k}{procedure}\PYG{p}{,}\PYG{+w}{ }\PYG{k}{nopass}\PYG{+w}{ }\PYG{k+kd}{::}\PYG{+w}{ }\PYG{n}{code}\PYG{+w}{ }\PYG{o}{=}\PYG{o}{\PYGZgt{}}\PYG{+w}{ }\PYG{n}{testkern\PYGZus{}operator\PYGZus{}code}
\PYG{k}{end }\PYG{k}{type }\PYG{n}{testkern\PYGZus{}operator\PYGZus{}type}
\end{sphinxVerbatim}

\sphinxAtStartPar
then we will need to pass basis functions for both the evaluator and the
quadrature (where \sphinxcode{\sphinxupquote{qr\_face}} is the name of the face\sphinxhyphen{}quadrature object
passed to the Invoke):

\begin{sphinxVerbatim}[commandchars=\\\{\}]
\PYG{k}{subroutine }\PYG{n}{testkern\PYGZus{}operator\PYGZus{}code}\PYG{p}{(}\PYG{n}{cell}\PYG{p}{,}\PYG{+w}{ }\PYG{n}{nlayers}\PYG{p}{,}\PYG{+w}{ }\PYG{n}{ncell\PYGZus{}3d}\PYG{p}{,}\PYG{+w}{              }\PYG{p}{\PYGZam{}}
\PYG{+w}{     }\PYG{n}{local\PYGZus{}stencil}\PYG{p}{,}\PYG{+w}{ }\PYG{n}{xdata}\PYG{p}{,}\PYG{+w}{ }\PYG{n}{ydata}\PYG{p}{,}\PYG{+w}{ }\PYG{n}{zdata}\PYG{p}{,}\PYG{+w}{ }\PYG{n}{ndf\PYGZus{}w0}\PYG{p}{,}\PYG{+w}{ }\PYG{n}{undf\PYGZus{}w0}\PYG{p}{,}\PYG{+w}{ }\PYG{n}{map\PYGZus{}w0}\PYG{p}{,}\PYG{+w}{       }\PYG{p}{\PYGZam{}}
\PYG{+w}{     }\PYG{n}{basis\PYGZus{}w0\PYGZus{}on\PYGZus{}w0}\PYG{p}{,}\PYG{+w}{ }\PYG{n}{basis\PYGZus{}w0\PYGZus{}qr\PYGZus{}face}\PYG{p}{,}\PYG{+w}{ }\PYG{n}{ndf\PYGZus{}w1}\PYG{p}{,}\PYG{+w}{                          }\PYG{p}{\PYGZam{}}
\PYG{+w}{     }\PYG{n}{np\PYGZus{}xyz\PYGZus{}qr\PYGZus{}face}\PYG{p}{,}\PYG{+w}{ }\PYG{n}{nfaces\PYGZus{}qr\PYGZus{}face}\PYG{p}{,}\PYG{+w}{ }\PYG{n}{weights\PYGZus{}xyz\PYGZus{}qr\PYGZus{}face}\PYG{p}{)}
\end{sphinxVerbatim}

\sphinxAtStartPar
If the metadata specifies that the kernel requires a property of the
reference element:

\begin{sphinxVerbatim}[commandchars=\\\{\}]
\PYG{k}{type}\PYG{p}{,}\PYG{+w}{ }\PYG{k}{extends}\PYG{p}{(}\PYG{n}{kernel\PYGZus{}type}\PYG{p}{)}\PYG{+w}{ }\PYG{k+kd}{::}\PYG{+w}{ }\PYG{n}{testkern\PYGZus{}operator\PYGZus{}type}
\PYG{+w}{   }\PYG{k}{type}\PYG{p}{(}\PYG{n}{arg\PYGZus{}type}\PYG{p}{)}\PYG{p}{,}\PYG{+w}{ }\PYG{k}{dimension}\PYG{p}{(}\PYG{l+m+mi}{2}\PYG{p}{)}\PYG{+w}{ }\PYG{k+kd}{::}\PYG{+w}{ }\PYG{n}{meta\PYGZus{}args}\PYG{+w}{ }\PYG{o}{=}\PYG{+w}{               }\PYG{p}{\PYGZam{}}
\PYG{+w}{        }\PYG{p}{(}\PYG{o}{/}\PYG{+w}{ }\PYG{n}{arg\PYGZus{}type}\PYG{p}{(}\PYG{n}{gh\PYGZus{}operator}\PYG{p}{,}\PYG{+w}{ }\PYG{n}{gh\PYGZus{}real}\PYG{p}{,}\PYG{+w}{ }\PYG{n}{gh\PYGZus{}write}\PYG{p}{,}\PYG{+w}{ }\PYG{n}{w0}\PYG{p}{,}\PYG{+w}{ }\PYG{n}{w1}\PYG{p}{)}\PYG{p}{,}\PYG{+w}{ }\PYG{p}{\PYGZam{}}
\PYG{+w}{           }\PYG{n}{arg\PYGZus{}type}\PYG{p}{(}\PYG{n}{gh\PYGZus{}field}\PYG{o}{*}\PYG{l+m+mi}{3}\PYG{p}{,}\PYG{+w}{  }\PYG{n}{gh\PYGZus{}real}\PYG{p}{,}\PYG{+w}{ }\PYG{n}{gh\PYGZus{}read}\PYG{p}{,}\PYG{+w}{  }\PYG{n}{w0}\PYG{p}{)}\PYG{+w}{ }\PYG{o}{/}\PYG{p}{)}
\PYG{+w}{   }\PYG{k}{type}\PYG{p}{(}\PYG{n}{reference\PYGZus{}element\PYGZus{}data\PYGZus{}type}\PYG{p}{)}\PYG{+w}{ }\PYG{k+kd}{::}\PYG{+w}{ }\PYG{n}{meta\PYGZus{}reference\PYGZus{}element}\PYG{p}{(}\PYG{l+m+mi}{1}\PYG{p}{)}\PYG{+w}{ }\PYG{o}{=}\PYG{+w}{  }\PYG{p}{\PYGZam{}}
\PYG{+w}{        }\PYG{p}{(}\PYG{o}{/}\PYG{+w}{ }\PYG{n}{reference\PYGZus{}element\PYGZus{}data\PYGZus{}type}\PYG{p}{(}\PYG{n}{normals\PYGZus{}to\PYGZus{}horizontal\PYGZus{}faces}\PYG{p}{)}\PYG{+w}{ }\PYG{o}{/}\PYG{p}{)}
\PYG{+w}{   }\PYG{k+kt}{integer}\PYG{+w}{ }\PYG{k+kd}{::}\PYG{+w}{ }\PYG{n}{operates\PYGZus{}on}\PYG{+w}{ }\PYG{o}{=}\PYG{+w}{ }\PYG{n}{cell\PYGZus{}column}
\PYG{+w}{ }\PYG{k}{contains}
\PYG{k}{   }\PYG{k}{procedure}\PYG{p}{,}\PYG{+w}{ }\PYG{k}{nopass}\PYG{+w}{ }\PYG{k+kd}{::}\PYG{+w}{ }\PYG{n}{code}\PYG{+w}{ }\PYG{o}{=}\PYG{o}{\PYGZgt{}}\PYG{+w}{ }\PYG{n}{testkern\PYGZus{}operator\PYGZus{}code}
\PYG{k}{end }\PYG{k}{type }\PYG{n}{testkern\PYGZus{}operator\PYGZus{}type}
\end{sphinxVerbatim}

\sphinxAtStartPar
then the kernel must be passed the number of faces of the reference element
and the array of face normals in the specified direction (here horizontal):

\begin{sphinxVerbatim}[commandchars=\\\{\}]
\PYG{k}{subroutine }\PYG{n}{testkern\PYGZus{}operator\PYGZus{}code}\PYG{p}{(}\PYG{n}{cell}\PYG{p}{,}\PYG{+w}{ }\PYG{n}{nlayers}\PYG{p}{,}\PYG{+w}{ }\PYG{n}{ncell\PYGZus{}3d}\PYG{p}{,}\PYG{+w}{        }\PYG{p}{\PYGZam{}}
\PYG{+w}{     }\PYG{n}{local\PYGZus{}stencil}\PYG{p}{,}\PYG{+w}{ }\PYG{n}{xdata}\PYG{p}{,}\PYG{+w}{ }\PYG{n}{ydata}\PYG{p}{,}\PYG{+w}{ }\PYG{n}{zdata}\PYG{p}{,}\PYG{+w}{ }\PYG{n}{ndf\PYGZus{}w0}\PYG{p}{,}\PYG{+w}{ }\PYG{n}{undf\PYGZus{}w0}\PYG{p}{,}\PYG{+w}{ }\PYG{n}{map\PYGZus{}w0}\PYG{p}{,}\PYG{+w}{ }\PYG{p}{\PYGZam{}}
\PYG{+w}{     }\PYG{n}{nfaces\PYGZus{}re\PYGZus{}h}\PYG{p}{,}\PYG{+w}{ }\PYG{n}{normals\PYGZus{}face\PYGZus{}h}\PYG{p}{)}
\end{sphinxVerbatim}


\subsubsection{Rules for CMA Kernels}
\label{\detokenize{dynamo0p3:rules-for-cma-kernels}}
\sphinxAtStartPar
Kernels involving CMA operators are restricted to just three types;
assembly, application/inverse\sphinxhyphen{}application and matrix\sphinxhyphen{}matrix.
We give the rules for each of these in the sections below.


\paragraph{Assembly}
\label{\detokenize{dynamo0p3:id3}}
\sphinxAtStartPar
An assembly kernel requires the column\sphinxhyphen{}banded dofmap for both the to\sphinxhyphen{}
and from\sphinxhyphen{}function spaces of the CMA operator being assembled as well
as the number of DoFs for each of the dofmaps. The full set of rules is:
\begin{enumerate}
\sphinxsetlistlabels{\arabic}{enumi}{enumii}{}{)}%
\item {} 
\sphinxAtStartPar
Include the \sphinxcode{\sphinxupquote{cell}} argument. \sphinxcode{\sphinxupquote{cell}} is an \sphinxcode{\sphinxupquote{integer}} of kind
\sphinxcode{\sphinxupquote{i\_def\textasciigrave{}\textasciigrave{}and has intent \textasciigrave{}\textasciigrave{}in}}.

\item {} 
\sphinxAtStartPar
Include \sphinxcode{\sphinxupquote{nlayers}}, the number of layers in a column. \sphinxcode{\sphinxupquote{nlayers}}
is an \sphinxcode{\sphinxupquote{integer}} of kind \sphinxcode{\sphinxupquote{i\_def}} and has intent \sphinxcode{\sphinxupquote{in}}.

\item {} 
\sphinxAtStartPar
Include the total number of cells in the 2D mesh (including halos),
\sphinxcode{\sphinxupquote{ncell\_2d}}, which is an \sphinxcode{\sphinxupquote{integer}} of kind \sphinxcode{\sphinxupquote{i\_def}} with
intent \sphinxcode{\sphinxupquote{in}}.

\item {} 
\sphinxAtStartPar
Include the total number of cells, \sphinxcode{\sphinxupquote{ncell\_3d}}, which is an \sphinxcode{\sphinxupquote{integer}}
of kind \sphinxcode{\sphinxupquote{i\_def}} with intent \sphinxcode{\sphinxupquote{in}}.

\item {} 
\sphinxAtStartPar
For each argument in the \sphinxcode{\sphinxupquote{meta\_args}} metadata array:
\begin{enumerate}
\sphinxsetlistlabels{\arabic}{enumii}{enumiii}{}{)}%
\item {} 
\sphinxAtStartPar
If it is a LMA operator, include a \sphinxcode{\sphinxupquote{real}}, 3\sphinxhyphen{}dimensional
array. The first dimension is \sphinxcode{\sphinxupquote{ncell\_3d}}. The second and third
dimension are the local degrees of freedom for the \sphinxcode{\sphinxupquote{to}} and
\sphinxcode{\sphinxupquote{from}} spaces, respectively.  The precision of the array depends
on how it is defined in the algorithm layer, see the
{\hyperref[\detokenize{dynamo0p3:lfric-mixed-precision}]{\sphinxcrossref{\DUrole{std,std-ref}{Mixed Precision}}}} section for more details;

\item {} 
\sphinxAtStartPar
If it is a CMA operator, include a \sphinxcode{\sphinxupquote{real}}, 3\sphinxhyphen{}dimensional array
of kind \sphinxcode{\sphinxupquote{r\_solver}}. The first dimension is
\sphinxcode{\sphinxupquote{"bandwidth\_"\textless{}operator\_name\textgreater{}}}, the second is
\sphinxcode{\sphinxupquote{"nrow\_"\textless{}operator\_name\textgreater{}}}, and the third is \sphinxcode{\sphinxupquote{ncell\_2d}}.
\begin{enumerate}
\sphinxsetlistlabels{\arabic}{enumiii}{enumiv}{}{)}%
\item {} 
\sphinxAtStartPar
Include the number of rows in the banded matrix.  This is
an \sphinxcode{\sphinxupquote{integer}} of kind \sphinxcode{\sphinxupquote{i\_def}} with intent \sphinxcode{\sphinxupquote{in}} and is named
as \sphinxcode{\sphinxupquote{"nrow\_"\textless{}operator\_name\textgreater{}}}.

\item {} 
\sphinxAtStartPar
If the from\sphinxhyphen{}space of the operator is \sphinxstyleemphasis{not} the same as the
to\sphinxhyphen{}space then include the number of columns in the banded
matrix.  This is an \sphinxcode{\sphinxupquote{integer}} of kind \sphinxcode{\sphinxupquote{i\_def}} with intent
\sphinxcode{\sphinxupquote{in}} and is named as \sphinxcode{\sphinxupquote{"ncol\_"\textless{}operator\_name\textgreater{}}}.

\item {} 
\sphinxAtStartPar
Include the bandwidth of the banded matrix. This is an
\sphinxcode{\sphinxupquote{integer}} of kind \sphinxcode{\sphinxupquote{i\_def}} with intent \sphinxcode{\sphinxupquote{in}} and is named as
\sphinxcode{\sphinxupquote{"bandwidth\_"\textless{}operator\_name\textgreater{}}}.

\item {} 
\sphinxAtStartPar
Include banded\sphinxhyphen{}matrix parameter \sphinxcode{\sphinxupquote{alpha}}. This is an \sphinxcode{\sphinxupquote{integer}}
of kind \sphinxcode{\sphinxupquote{i\_def}} with intent \sphinxcode{\sphinxupquote{in}} and is named as
\sphinxcode{\sphinxupquote{"alpha\_"\textless{}operator\_name\textgreater{}}}.

\item {} 
\sphinxAtStartPar
Include banded\sphinxhyphen{}matrix parameter \sphinxcode{\sphinxupquote{beta}}. This is an \sphinxcode{\sphinxupquote{integer}}
of kind \sphinxcode{\sphinxupquote{i\_def}} with intent \sphinxcode{\sphinxupquote{in}} and is named as
\sphinxcode{\sphinxupquote{"beta\_"\textless{}operator\_name\textgreater{}}}.

\item {} 
\sphinxAtStartPar
Include banded\sphinxhyphen{}matrix parameter \sphinxcode{\sphinxupquote{gamma\_m}}. This is an \sphinxcode{\sphinxupquote{integer}}
of kind \sphinxcode{\sphinxupquote{i\_def}} with intent \sphinxcode{\sphinxupquote{in}} and is named as
\sphinxcode{\sphinxupquote{"gamma\_m\_"\textless{}operator\_name\textgreater{}}}.

\item {} 
\sphinxAtStartPar
Include banded\sphinxhyphen{}matrix parameter \sphinxcode{\sphinxupquote{gamma\_p}}. This is an \sphinxcode{\sphinxupquote{integer}}
of kind \sphinxcode{\sphinxupquote{i\_def}} with intent \sphinxcode{\sphinxupquote{in}} and is named as
\sphinxcode{\sphinxupquote{"gamma\_p\_"\textless{}operator\_name\textgreater{}}}.

\end{enumerate}

\item {} 
\sphinxAtStartPar
If it is a field or scalar argument then include arguments following
the same rules as for general\sphinxhyphen{}purpose kernels.

\end{enumerate}

\item {} 
\sphinxAtStartPar
For each unique function space in the order they appear in the
metadata arguments (the \sphinxcode{\sphinxupquote{to}} function space of an operator is
considered to be before the \sphinxcode{\sphinxupquote{from}} function space of the same
operator as it appears first in lexicographic order):
\begin{enumerate}
\sphinxsetlistlabels{\arabic}{enumii}{enumiii}{}{)}%
\item {} 
\sphinxAtStartPar
Include the number of degrees of freedom per cell for the space.
This is an \sphinxcode{\sphinxupquote{integer}} of kind \sphinxcode{\sphinxupquote{i\_def}} with intent \sphinxcode{\sphinxupquote{in}}. The name
of this argument is \sphinxcode{\sphinxupquote{"ndf\_"\textless{}arg\_function\_space\textgreater{}}}.

\item {} 
\sphinxAtStartPar
If there is a field on this space then:
\begin{enumerate}
\sphinxsetlistlabels{\arabic}{enumiii}{enumiv}{}{)}%
\item {} 
\sphinxAtStartPar
Include the unique number of degrees of freedom for the
function space. This is an \sphinxcode{\sphinxupquote{integer}} of kind \sphinxcode{\sphinxupquote{i\_def}} and has
intent \sphinxcode{\sphinxupquote{in}}. The name of this argument is
\sphinxcode{\sphinxupquote{"undf\_"\textless{}field\_function\_space\textgreater{}}}.

\item {} 
\sphinxAtStartPar
Include the dofmap for this space. This is an \sphinxcode{\sphinxupquote{integer}} array
of kind \sphinxcode{\sphinxupquote{i\_def}} with intent \sphinxcode{\sphinxupquote{in}}. It has one dimension
sized by the local degrees of freedom for the function space.

\end{enumerate}

\item {} 
\sphinxAtStartPar
If the CMA operator has this space as its to/from space then
include the column\sphinxhyphen{}banded dofmap, the list of offsets for the
to/from\sphinxhyphen{}space. This is an \sphinxcode{\sphinxupquote{integer}} array of rank 2 and kind
\sphinxcode{\sphinxupquote{i\_def}}. The first dimension is \sphinxcode{\sphinxupquote{"ndf\_"\textless{}arg\_function\_space\textgreater{}}}
and the second is \sphinxcode{\sphinxupquote{nlayers}}.

\end{enumerate}

\end{enumerate}


\paragraph{Application/Inverse\sphinxhyphen{}Application}
\label{\detokenize{dynamo0p3:application-inverse-application}}
\sphinxAtStartPar
A kernel applying a CMA operator requires the column\sphinxhyphen{}indirection
dofmap for both the to\sphinxhyphen{} and from\sphinxhyphen{}function spaces of the CMA
operator. Since it does not have any LMA operator arguments it does
not require the \sphinxcode{\sphinxupquote{ncell\_3d}} and \sphinxcode{\sphinxupquote{nlayers}} scalar arguments. (Since a
column\sphinxhyphen{}wise operator is, by definition, assembled for a whole column,
there is no loop over levels when applying it.)
The full set of rules is then:
\begin{enumerate}
\sphinxsetlistlabels{\arabic}{enumi}{enumii}{}{)}%
\item {} 
\sphinxAtStartPar
Include the \sphinxcode{\sphinxupquote{cell}} argument. \sphinxcode{\sphinxupquote{cell}} is an \sphinxcode{\sphinxupquote{integer}} of kind
\sphinxcode{\sphinxupquote{i\_def}} and has intent \sphinxcode{\sphinxupquote{in}}.

\item {} 
\sphinxAtStartPar
Include the total number of cells in the 2D mesh (including halos),
\sphinxcode{\sphinxupquote{ncell\_2d}}, which is an \sphinxcode{\sphinxupquote{integer}} of kind \sphinxcode{\sphinxupquote{i\_def}} with
intent \sphinxcode{\sphinxupquote{in}}.

\item {} 
\sphinxAtStartPar
For each argument in the \sphinxcode{\sphinxupquote{meta\_args}} metadata array:
\begin{enumerate}
\sphinxsetlistlabels{\arabic}{enumii}{enumiii}{}{)}%
\item {} 
\sphinxAtStartPar
If it is a field, include the field array. This is a \sphinxcode{\sphinxupquote{real}}
array of rank 1. Its precision (kind) depends on how it is
defined in the algorithm layer, see the
{\hyperref[\detokenize{dynamo0p3:lfric-mixed-precision}]{\sphinxcrossref{\DUrole{std,std-ref}{Mixed Precision}}}}. The field array name is
currently specified as being
\sphinxcode{\sphinxupquote{"field\_"\textless{}argument\_position\textgreater{}"\_"\textless{}field\_function\_space\textgreater{}}}. The
extent of the array is the number of unique degrees of freedom
for the function space that the field is on.  This value is
passed in separately. The intent of the argument is determined
from the metadata (see {\hyperref[\detokenize{dynamo0p3:lfric-api-meta-args}]{\sphinxcrossref{\DUrole{std,std-ref}{meta\_args}}}});

\item {} 
\sphinxAtStartPar
If it is a CMA operator, include it and its associated
parameters (see Rule 5 of CMA Assembly kernels).

\end{enumerate}

\item {} 
\sphinxAtStartPar
For each of the unique function spaces encountered in the
metadata arguments (the \sphinxcode{\sphinxupquote{to}} function space of an operator
is considered to be before the \sphinxcode{\sphinxupquote{from}} function space of the
same operator as it appears first in lexicographic order):
\begin{enumerate}
\sphinxsetlistlabels{\arabic}{enumii}{enumiii}{}{)}%
\item {} 
\sphinxAtStartPar
Include the number of degrees of freedom per cell for the associated
function space. This is an \sphinxcode{\sphinxupquote{integer}} of kind \sphinxcode{\sphinxupquote{i\_def}} with intent
\sphinxcode{\sphinxupquote{in}}. The name of this argument is \sphinxcode{\sphinxupquote{"ndf\_"\textless{}field\_function\_space\textgreater{}}};

\item {} 
\sphinxAtStartPar
Include the number of unique degrees of freedom for the associated
function space. This is an \sphinxcode{\sphinxupquote{integer}} of kind \sphinxcode{\sphinxupquote{i\_def}} with intent
\sphinxcode{\sphinxupquote{in}}. The name of this argument is \sphinxcode{\sphinxupquote{"undf\_"\textless{}field\_function\_space\textgreater{}}};

\item {} 
\sphinxAtStartPar
Include the dofmap for this function space. This is a rank\sphinxhyphen{}1 \sphinxcode{\sphinxupquote{integer}}
array of kind \sphinxcode{\sphinxupquote{i\_def}} with extent equal to the number of degrees of
freedom of the space (\sphinxcode{\sphinxupquote{"ndf\_"\textless{}field\_function\_space\textgreater{}}}).

\end{enumerate}

\item {} 
\sphinxAtStartPar
Include the indirection map for the to\sphinxhyphen{}space of the CMA operator.
This is a rank\sphinxhyphen{}1 \sphinxcode{\sphinxupquote{integer}} array of kind \sphinxcode{\sphinxupquote{i\_def}} with extent \sphinxcode{\sphinxupquote{nrow}}.

\item {} 
\sphinxAtStartPar
If the from\sphinxhyphen{}space of the operator is \sphinxstyleemphasis{not} the same as the to\sphinxhyphen{}space
then include the indirection map for the from\sphinxhyphen{}space of the CMA operator.
This is a rank\sphinxhyphen{}1 \sphinxcode{\sphinxupquote{integer}} array of kind \sphinxcode{\sphinxupquote{i\_def}} with extent \sphinxcode{\sphinxupquote{ncol}}.

\end{enumerate}


\paragraph{Matrix\sphinxhyphen{}Matrix}
\label{\detokenize{dynamo0p3:id4}}
\sphinxAtStartPar
Does not require any dofmaps and also does not require the \sphinxcode{\sphinxupquote{nlayers}}
and \sphinxcode{\sphinxupquote{ncell\_3d}} scalar arguments. The full set of rules are then:
\begin{enumerate}
\sphinxsetlistlabels{\arabic}{enumi}{enumii}{}{)}%
\item {} 
\sphinxAtStartPar
Include the \sphinxcode{\sphinxupquote{cell}} argument. \sphinxcode{\sphinxupquote{cell}} is an \sphinxcode{\sphinxupquote{integer}} of kind
\sphinxcode{\sphinxupquote{i\_def}} and has intent \sphinxcode{\sphinxupquote{in}}.

\item {} 
\sphinxAtStartPar
Include the total number of cells in the 2D mesh (including halos),
\sphinxcode{\sphinxupquote{ncell\_2d}}, which is an \sphinxcode{\sphinxupquote{integer}} of kind \sphinxcode{\sphinxupquote{i\_def}} with intent \sphinxcode{\sphinxupquote{in}}.

\item {} 
\sphinxAtStartPar
For each CMA operator or scalar argument specified in metadata:
\begin{enumerate}
\sphinxsetlistlabels{\arabic}{enumii}{enumiii}{}{)}%
\item {} 
\sphinxAtStartPar
If it is a CMA operator, include it and its associated
parameters (see Rule 5 of CMA Assembly kernels);

\item {} 
\sphinxAtStartPar
If it is a scalar argument include the corresponding Fortran
variable in the argument list with intent \sphinxcode{\sphinxupquote{in}}.

\end{enumerate}

\end{enumerate}


\subsubsection{Rules for Inter\sphinxhyphen{}Grid Kernels}
\label{\detokenize{dynamo0p3:id5}}
\sphinxAtStartPar
As already specified, inter\sphinxhyphen{}grid kernels are only permitted to take
fields and/or field\sphinxhyphen{}vectors as arguments. Fields (and field\sphinxhyphen{}vectors)
that are on different meshes must be on different function
spaces. Fields on the same mesh must also be on the same function
space.

\sphinxAtStartPar
Argument ordering follows the general pattern used for ‘normal’
kernels with field data being followed by dofmap data. The rules for
arguments to inter\sphinxhyphen{}grid kernels are as follows:
\begin{enumerate}
\sphinxsetlistlabels{\arabic}{enumi}{enumii}{}{)}%
\item {} 
\sphinxAtStartPar
Include \sphinxcode{\sphinxupquote{nlayers}}, the number of layers in a column. \sphinxcode{\sphinxupquote{nlayers}}
is an \sphinxcode{\sphinxupquote{integer}} of kind \sphinxcode{\sphinxupquote{i\_def}} and has intent \sphinxcode{\sphinxupquote{in}}.

\item {} 
\sphinxAtStartPar
Include the \sphinxcode{\sphinxupquote{cell\_map}} for the current cell (column). This is
an \sphinxcode{\sphinxupquote{integer}} array of rank two, kind \sphinxcode{\sphinxupquote{i\_def}} and intent \sphinxcode{\sphinxupquote{in}}
which provides the mapping from the coarse to the fine mesh. It
has extent \sphinxcode{\sphinxupquote{(ncell\_f\_per\_c\_x, ncell\_f\_per\_c\_y)}}.

\item {} 
\sphinxAtStartPar
Include \sphinxcode{\sphinxupquote{ncell\_f\_per\_c\_x}}, and \sphinxcode{\sphinxupquote{ncell\_f\_per\_c\_y}}, the numbers of
fine cells per coarse cell in the \sphinxcode{\sphinxupquote{x}} and \sphinxcode{\sphinxupquote{y}} directions,
respectively. These are integers of kind \sphinxcode{\sphinxupquote{i\_def}} and have intent
\sphinxcode{\sphinxupquote{in}}.

\item {} 
\sphinxAtStartPar
Include \sphinxcode{\sphinxupquote{ncell\_f}}, the number of cells (columns) in the fine mesh.
This is an \sphinxcode{\sphinxupquote{integer}} of kind \sphinxcode{\sphinxupquote{i\_def}} and has intent \sphinxcode{\sphinxupquote{in}}.

\item {} 
\sphinxAtStartPar
For each argument in the \sphinxcode{\sphinxupquote{meta\_args}} metadata array (which must be
a field or field\sphinxhyphen{}vector):
\begin{enumerate}
\sphinxsetlistlabels{\arabic}{enumii}{enumiii}{}{)}%
\item {} 
\sphinxAtStartPar
Pass in field data as done for a regular kernel.

\end{enumerate}

\item {} 
\sphinxAtStartPar
For each unique function space (of which there will currently be two)
in the order in which they are encountered in the \sphinxcode{\sphinxupquote{meta\_args}}
metadata array, include dofmap information:

\sphinxAtStartPar
If the dofmap is associated with an argument on the fine mesh:
\begin{enumerate}
\sphinxsetlistlabels{\arabic}{enumii}{enumiii}{}{)}%
\item {} 
\sphinxAtStartPar
Include \sphinxcode{\sphinxupquote{ndf\_fine}}, the number of DoFs per cell for the FS of
the field on the fine mesh;

\item {} 
\sphinxAtStartPar
Include \sphinxcode{\sphinxupquote{undf\_fine}}, the number of unique DoFs per cell for the FS
of the field on the fine mesh;

\item {} 
\sphinxAtStartPar
Include \sphinxcode{\sphinxupquote{dofmap\_fine}}, the \sphinxstyleemphasis{whole} dofmap for the fine mesh. This
is an \sphinxcode{\sphinxupquote{integer}} array of rank two and kind \sphinxcode{\sphinxupquote{i\_def}} with intent
\sphinxcode{\sphinxupquote{in}}. The extent of the first dimension is \sphinxcode{\sphinxupquote{ndf\_fine}} and that of
the second is \sphinxcode{\sphinxupquote{ncell\_f}}.

\end{enumerate}

\sphinxAtStartPar
else, the dofmap is associated with an argument on the coarse mesh:
\begin{enumerate}
\sphinxsetlistlabels{\arabic}{enumii}{enumiii}{}{)}%
\item {} 
\sphinxAtStartPar
Include \sphinxcode{\sphinxupquote{undf\_coarse}}, the number of unique DoFs for the coarse
field. This is an \sphinxcode{\sphinxupquote{integer}} of kind \sphinxcode{\sphinxupquote{i\_def}} with intent \sphinxcode{\sphinxupquote{in}};

\item {} 
\sphinxAtStartPar
Include \sphinxcode{\sphinxupquote{dofmap\_coarse}}, the dofmap for the current cell (column)
in the coarse mesh. This is an \sphinxcode{\sphinxupquote{integer}} array of rank one, kind
\sphinxcode{\sphinxupquote{i\_def\textasciigrave{}\textasciigrave{}and has intent \textasciigrave{}\textasciigrave{}in}}.

\end{enumerate}

\end{enumerate}


\subsubsection{Rules for Domain Kernels}
\label{\detokenize{dynamo0p3:rules-for-domain-kernels}}
\sphinxAtStartPar
The rules for kernels that have \sphinxcode{\sphinxupquote{operates\_on = DOMAIN}} are almost
identical to those for general\sphinxhyphen{}purpose kernels (described {\hyperref[\detokenize{dynamo0p3:lfric-stub-generation-rules}]{\sphinxcrossref{\DUrole{std,std-ref}{above}}}}), allowing for the fact that they
are not permitted any type of operator argument or any argument with a
stencil access. The only difference is that, since the kernel operates
on the whole domain, the number of columns in the mesh excluding those
in the halo (\sphinxcode{\sphinxupquote{ncell\_2d\_no\_halos}}), must be passed in. This is provided
as the second argument to the kernel (after \sphinxcode{\sphinxupquote{nlayers}}).
\sphinxcode{\sphinxupquote{ncell\_2d\_no\_halos}} is an \sphinxcode{\sphinxupquote{integer}} of kind \sphinxcode{\sphinxupquote{i\_def}} with intent \sphinxcode{\sphinxupquote{in}}.


\subsubsection{Rules for DoF Kernels}
\label{\detokenize{dynamo0p3:rules-for-dof-kernels}}
\sphinxAtStartPar
The rules for kernels that have \sphinxcode{\sphinxupquote{operates\_on = DOF}} are similar to those for
general\sphinxhyphen{}purpose kernels but, due to the restriction that only fields and
scalars can be passed to them, are much fewer. The full set of rules, along
with PSyclone’s naming conventions, are:
\begin{enumerate}
\sphinxsetlistlabels{\arabic}{enumi}{enumii}{}{)}%
\item {} 
\sphinxAtStartPar
Include \sphinxtitleref{df}, the index of the single dof to be operated on. This is an
\sphinxcode{\sphinxupquote{integer}} of of kind \sphinxcode{\sphinxupquote{i\_def}} with intent \sphinxcode{\sphinxupquote{in}}.

\item {} 
\sphinxAtStartPar
For each scalar/field in the order specified by the meta\_args metadata:
\begin{enumerate}
\sphinxsetlistlabels{\arabic}{enumii}{enumiii}{}{)}%
\item {} 
\sphinxAtStartPar
If the current entry is a scalar quantity then include the Fortran
variable in the argument list. The intent is determined from the
metadata (see meta\_args for an explanation).

\item {} 
\sphinxAtStartPar
If the current entry is a field then include the field array. The
field array name is currently specified as being \sphinxcode{\sphinxupquote{"field\_"
\textless{}argument\_position\textgreater{}}}. A field array is a rank\sphinxhyphen{}1, real array with
extent equal to the number of unique degrees of freedom for the space
that the field is on. Its precision (kind) depends on how it is
defined in the algorithm layer, see the {\hyperref[\detokenize{dynamo0p3:lfric-mixed-precision}]{\sphinxcrossref{\DUrole{std,std-ref}{Mixed Precision}}}} section for more details. This value is
passed in separately. Again, the intent is determined from the
metadata (see {\hyperref[\detokenize{dynamo0p3:lfric-api-meta-args}]{\sphinxcrossref{\DUrole{std,std-ref}{meta\_args}}}}).

\end{enumerate}

\item {} 
\sphinxAtStartPar
For each field vector in the order specified by the meta\_args metadata,
there needs to be an equivalent number of arguments in the kernel as
the dimension of the field vector. The dimension is specified in the
metadata. The arguments must be ordered following the indexing of the
field vector.

\end{enumerate}


\subsubsection{Argument Intents}
\label{\detokenize{dynamo0p3:lfric-kernel-arg-intents}}\label{\detokenize{dynamo0p3:id6}}
\sphinxAtStartPar
As described {\hyperref[\detokenize{dynamo0p3:lfric-psy-arg-intents}]{\sphinxcrossref{\DUrole{std,std-ref}{above}}}}, LFRic kernels read
and/or update the data pointed to by objects such as
{\hyperref[\detokenize{dynamo0p3:lfric-field}]{\sphinxcrossref{\DUrole{std,std-ref}{fields}}}} or {\hyperref[\detokenize{dynamo0p3:lfric-operator}]{\sphinxcrossref{\DUrole{std,std-ref}{operators}}}}.
This data is passed to the kernels as {\hyperref[\detokenize{dynamo0p3:lfric-kern-subroutine}]{\sphinxcrossref{\DUrole{std,std-ref}{subroutine arguments}}}} and their Fortran intents usually follow the
logic determined by their {\hyperref[\detokenize{dynamo0p3:lfric-kernel-valid-access}]{\sphinxcrossref{\DUrole{std,std-ref}{access modes}}}}.
\begin{itemize}
\item {} 
\sphinxAtStartPar
\sphinxcode{\sphinxupquote{GH\_READ}} indicates \sphinxcode{\sphinxupquote{intent(in)}} as the argument is only ever read from.

\item {} 
\sphinxAtStartPar
\sphinxcode{\sphinxupquote{GH\_WRITE}} (for discontinuous function spaces) indicates that the argument
is only written to in a kernel. The field and operator arguments’ data in
LFRic are always defined outside of a kernel so the argument intent for
this access type is \sphinxcode{\sphinxupquote{intent(inout)}}.

\item {} 
\sphinxAtStartPar
\sphinxcode{\sphinxupquote{GH\_INC}}, \sphinxcode{\sphinxupquote{GH\_READINC}} and \sphinxcode{\sphinxupquote{GH\_READWRITE}} indicate
\sphinxcode{\sphinxupquote{intent(inout)}} as the arguments are updated (albeit in a
different way due to different access to DoFs, see
{\hyperref[\detokenize{dynamo0p3:lfric-api-meta-args}]{\sphinxcrossref{\DUrole{std,std-ref}{meta\_args}}}} for more details).

\end{itemize}


\subsection{Kernel Naming Conventions}
\label{\detokenize{dynamo0p3:kernel-naming-conventions}}
\sphinxAtStartPar
LFRic development uses strict naming conventions related to kernels.
While they are not a requirement for PSyclone itself, any LFRic
development should follow these conventions (see e.g.
{\hyperref[\detokenize{tutorials_and_examples:examples-lfric}]{\sphinxcrossref{\DUrole{std,std-ref}{LFRic examples}}}} in PSyclone):
\begin{description}
\sphinxlineitem{Module name:}
\sphinxAtStartPar
\sphinxcode{\sphinxupquote{\textless{}base\_name\textgreater{}\_kernel\_mod}}

\sphinxlineitem{Kernel type name:}
\sphinxAtStartPar
\sphinxcode{\sphinxupquote{\textless{}base\_name\textgreater{}\_kernel\_type}}

\sphinxlineitem{Subroutine name:}
\sphinxAtStartPar
\sphinxcode{\sphinxupquote{\textless{}base\_name\textgreater{}\_code}}

\end{description}

\sphinxAtStartPar
The latest version of the LFRic coding style guidelines are availabe in this
\sphinxhref{https://code.metoffice.gov.uk/trac/lfric/wiki/LFRicTechnical/FortranCodingStandards}{LFRic wiki page}
(requires login access to MOSRS, see the above {\hyperref[\detokenize{dynamo0p3:lfric-api}]{\sphinxcrossref{\DUrole{std,std-ref}{introduction}}}}
to the LFRic API).


\section{Built\sphinxhyphen{}ins}
\label{\detokenize{dynamo0p3:built-ins}}\label{\detokenize{dynamo0p3:lfric-built-ins}}
\sphinxAtStartPar
The basic concept of a PSyclone Built\sphinxhyphen{}in is described in the
{\hyperref[\detokenize{introduction_to_psykal:built-ins}]{\sphinxcrossref{\DUrole{std,std-ref}{Built\sphinxhyphen{}ins}}}} section.  In the LFRic API, calls to
Built\sphinxhyphen{}ins generally follow a convention that the field/scalar written
to comes first in the argument list. LFRic Built\sphinxhyphen{}ins must conform to the
following rules:
\begin{enumerate}
\sphinxsetlistlabels{\arabic}{enumi}{enumii}{}{)}%
\item {} 
\sphinxAtStartPar
They must have one and only one modified (i.e. written to) argument.

\item {} 
\sphinxAtStartPar
They must operate on a DoF (\sphinxcode{\sphinxupquote{operates\_on = DOF}} metadata).

\item {} 
\sphinxAtStartPar
There must be at least one field in the argument list. This is so
that we know the number of DoFs to iterate over in the PSy layer.

\item {} 
\sphinxAtStartPar
Kernel arguments must be either fields or scalars (\sphinxcode{\sphinxupquote{real}}\sphinxhyphen{} and/or
\sphinxcode{\sphinxupquote{integer}}\sphinxhyphen{}valued).

\item {} 
\sphinxAtStartPar
All field arguments to a given Built\sphinxhyphen{}in must be on the same
function space. This is because all current Built\sphinxhyphen{}ins operate on
DoFs and therefore all fields should have the same number. It also
means that we can determine the number of DoFs uniquely when a
scalar is written to;

\item {} 
\sphinxAtStartPar
Built\sphinxhyphen{}ins that update \sphinxcode{\sphinxupquote{real}}\sphinxhyphen{}valued fields can, in general, only
read from other \sphinxcode{\sphinxupquote{real}}\sphinxhyphen{}valued fields, but they can take both \sphinxcode{\sphinxupquote{real}}
and \sphinxcode{\sphinxupquote{integer}} scalar arguments (see rule 8 for exceptions);

\item {} 
\sphinxAtStartPar
Built\sphinxhyphen{}ins that update \sphinxcode{\sphinxupquote{integer}}\sphinxhyphen{}valued fields can, in general, only
read from other \sphinxcode{\sphinxupquote{integer}}\sphinxhyphen{}valued fields and take \sphinxcode{\sphinxupquote{integer}} scalar
arguments (see rule 8 for exceptions);

\item {} 
\sphinxAtStartPar
The only two exceptions from the rules 6) and 7) above regarding the
same data type of “write” and “read” field arguments are Built\sphinxhyphen{}ins
that convert field data from \sphinxcode{\sphinxupquote{real}} to \sphinxcode{\sphinxupquote{integer}}, \sphinxcode{\sphinxupquote{real\_to\_int\_X}},
and from \sphinxcode{\sphinxupquote{integer}} to \sphinxcode{\sphinxupquote{real}}, \sphinxcode{\sphinxupquote{int\_to\_real\_X}}.

\end{enumerate}

\sphinxAtStartPar
The Built\sphinxhyphen{}ins supported for the LFRic API are listed in the related
subsections, grouped first by the data type of fields they operate on
({\hyperref[\detokenize{dynamo0p3:lfric-built-ins-real}]{\sphinxcrossref{\DUrole{std,std-ref}{real\sphinxhyphen{}valued}}}} and
{\hyperref[\detokenize{dynamo0p3:lfric-built-ins-int}]{\sphinxcrossref{\DUrole{std,std-ref}{integer\sphinxhyphen{}valued}}}}) and then by the mathematical
operation they perform.

\sphinxAtStartPar
The field arguments in Built\sphinxhyphen{}ins are the derived types that represent the
{\hyperref[\detokenize{dynamo0p3:lfric-field}]{\sphinxcrossref{\DUrole{std,std-ref}{LFRic fields}}}}, however mathematical operations are
actually performed on the data of the \sphinxstyleemphasis{field proxies} (e.g.
\sphinxcode{\sphinxupquote{field1\_proxy\%data(:)}}). For instance, \sphinxcode{\sphinxupquote{X\_plus\_Y}} Built\sphinxhyphen{}in adds the
values of two fields accessed via their proxies in a loop over DoFs:

\begin{sphinxVerbatim}[commandchars=\\\{\}]
\PYG{k}{DO }\PYG{n}{df}\PYG{o}{=}\PYG{n}{loop0\PYGZus{}start}\PYG{p}{,}\PYG{n}{loop0\PYGZus{}stop}
\PYG{+w}{   }\PYG{n}{field3\PYGZus{}proxy}\PYG{p}{\PYGZpc{}}\PYG{k}{data}\PYG{p}{(}\PYG{n}{df}\PYG{p}{)}\PYG{+w}{ }\PYG{o}{=}\PYG{+w}{ }\PYG{n}{field1\PYGZus{}proxy}\PYG{p}{\PYGZpc{}}\PYG{k}{data}\PYG{p}{(}\PYG{n}{df}\PYG{p}{)}\PYG{+w}{ }\PYG{o}{+}\PYG{+w}{ }\PYG{n}{field2\PYGZus{}proxy}\PYG{p}{\PYGZpc{}}\PYG{k}{data}\PYG{p}{(}\PYG{n}{df}\PYG{p}{)}
\end{sphinxVerbatim}

\sphinxAtStartPar
where the precise values of the loop limits depend on the use of
{\hyperref[\detokenize{introduction_to_psykal:psykal-usage}]{\sphinxcrossref{\DUrole{std,std-ref}{distributed memory}}}},
{\hyperref[\detokenize{dynamo0p3:lfric-annexed-dofs}]{\sphinxcrossref{\DUrole{std,std-ref}{annexed DoFs}}}} or both.

\sphinxAtStartPar
As described in the PSy\sphinxhyphen{}layer {\hyperref[\detokenize{dynamo0p3:lfric-psy-arg-intents}]{\sphinxcrossref{\DUrole{std,std-ref}{Argument Intents}}}} section, the Fortran intent of LFRic
{\hyperref[\detokenize{dynamo0p3:lfric-field}]{\sphinxcrossref{\DUrole{std,std-ref}{field}}}} objects is always \sphinxcode{\sphinxupquote{in}} (because it is only
the data pointed to from within the object that is modified). The field
or scalar that has its data modified by a Built\sphinxhyphen{}in is marked in \sphinxstylestrong{bold}.

\sphinxAtStartPar
For clarity, the calculation performed by each Built\sphinxhyphen{}in is described using
Fortran array syntax without the details about field proxies. The actual
implementation of the Built\sphinxhyphen{}in may change in future (\sphinxstyleemphasis{e.g.} it could be
implemented by PSyclone generating a call to an optimised Maths library).


\subsection{Metadata}
\label{\detokenize{dynamo0p3:lfric-api-built-ins-metadata}}\label{\detokenize{dynamo0p3:id7}}
\sphinxAtStartPar
The code below outlines the elements of the LFRic API Built\sphinxhyphen{}in
metadata for the Built\sphinxhyphen{}ins that update a \sphinxcode{\sphinxupquote{real}}\sphinxhyphen{}valued field,
1) ‘meta\_args’, 2) ‘operates\_on’ and 3) ‘procedure’:

\begin{sphinxVerbatim}[commandchars=\\\{\}]
\PYG{k}{type}\PYG{p}{,}\PYG{+w}{ }\PYG{k}{public}\PYG{p}{,}\PYG{+w}{ }\PYG{k}{extends}\PYG{p}{(}\PYG{n}{kernel\PYGZus{}type}\PYG{p}{)}\PYG{+w}{ }\PYG{k+kd}{::}\PYG{+w}{ }\PYG{n}{aX\PYGZus{}plus\PYGZus{}bY}
\PYG{+w}{   }\PYG{k}{private}
\PYG{k}{   }\PYG{k}{type}\PYG{p}{(}\PYG{n}{arg\PYGZus{}type}\PYG{p}{)}\PYG{+w}{ }\PYG{k+kd}{::}\PYG{+w}{ }\PYG{n}{meta\PYGZus{}args}\PYG{p}{(}\PYG{l+m+mi}{5}\PYG{p}{)}\PYG{+w}{ }\PYG{o}{=}\PYG{+w}{ }\PYG{p}{(}\PYG{o}{/}\PYG{+w}{                              }\PYG{p}{\PYGZam{}}
\PYG{+w}{        }\PYG{n}{arg\PYGZus{}type}\PYG{p}{(}\PYG{n}{GH\PYGZus{}FIELD}\PYG{p}{,}\PYG{+w}{  }\PYG{n}{GH\PYGZus{}REAL}\PYG{p}{,}\PYG{+w}{ }\PYG{n}{GH\PYGZus{}WRITE}\PYG{p}{,}\PYG{+w}{ }\PYG{n}{ANY\PYGZus{}SPACE\PYGZus{}1}\PYG{p}{)}\PYG{p}{,}\PYG{+w}{        }\PYG{p}{\PYGZam{}}
\PYG{+w}{        }\PYG{n}{arg\PYGZus{}type}\PYG{p}{(}\PYG{n}{GH\PYGZus{}SCALAR}\PYG{p}{,}\PYG{+w}{ }\PYG{n}{GH\PYGZus{}REAL}\PYG{p}{,}\PYG{+w}{ }\PYG{n}{GH\PYGZus{}READ}\PYG{+w}{              }\PYG{p}{)}\PYG{p}{,}\PYG{+w}{        }\PYG{p}{\PYGZam{}}
\PYG{+w}{        }\PYG{n}{arg\PYGZus{}type}\PYG{p}{(}\PYG{n}{GH\PYGZus{}FIELD}\PYG{p}{,}\PYG{+w}{  }\PYG{n}{GH\PYGZus{}REAL}\PYG{p}{,}\PYG{+w}{ }\PYG{n}{GH\PYGZus{}READ}\PYG{p}{,}\PYG{+w}{  }\PYG{n}{ANY\PYGZus{}SPACE\PYGZus{}1}\PYG{p}{)}\PYG{p}{,}\PYG{+w}{        }\PYG{p}{\PYGZam{}}
\PYG{+w}{        }\PYG{n}{arg\PYGZus{}type}\PYG{p}{(}\PYG{n}{GH\PYGZus{}SCALAR}\PYG{p}{,}\PYG{+w}{ }\PYG{n}{GH\PYGZus{}REAL}\PYG{p}{,}\PYG{+w}{ }\PYG{n}{GH\PYGZus{}READ}\PYG{+w}{              }\PYG{p}{)}\PYG{p}{,}\PYG{+w}{        }\PYG{p}{\PYGZam{}}
\PYG{+w}{        }\PYG{n}{arg\PYGZus{}type}\PYG{p}{(}\PYG{n}{GH\PYGZus{}FIELD}\PYG{p}{,}\PYG{+w}{  }\PYG{n}{GH\PYGZus{}REAL}\PYG{p}{,}\PYG{+w}{ }\PYG{n}{GH\PYGZus{}READ}\PYG{p}{,}\PYG{+w}{  }\PYG{n}{ANY\PYGZus{}SPACE\PYGZus{}1}\PYG{p}{)}\PYG{+w}{         }\PYG{p}{\PYGZam{}}
\PYG{+w}{        }\PYG{o}{/}\PYG{p}{)}
\PYG{+w}{   }\PYG{k+kt}{integer}\PYG{+w}{ }\PYG{k+kd}{::}\PYG{+w}{ }\PYG{n}{operates\PYGZus{}on}\PYG{+w}{ }\PYG{o}{=}\PYG{+w}{ }\PYG{n}{DOF}
\PYG{+w}{ }\PYG{k}{contains}
\PYG{k}{   }\PYG{k}{procedure}\PYG{p}{,}\PYG{+w}{ }\PYG{k}{nopass}\PYG{+w}{ }\PYG{k+kd}{::}\PYG{+w}{ }\PYG{n}{aX\PYGZus{}plus\PYGZus{}bY\PYGZus{}code}
\PYG{k}{end }\PYG{k}{type }\PYG{n}{aX\PYGZus{}plus\PYGZus{}bY}
\end{sphinxVerbatim}

\sphinxAtStartPar
As can be seen, the metadata for a Built\sphinxhyphen{}in kernel is a subset of that
for a {\hyperref[\detokenize{dynamo0p3:lfric-api-kernel-metadata}]{\sphinxcrossref{\DUrole{std,std-ref}{user\sphinxhyphen{}defined Kernel}}}} with the
exception that \sphinxcode{\sphinxupquote{operates\_on}} must be \sphinxcode{\sphinxupquote{DOF}} instead of \sphinxcode{\sphinxupquote{CELL\_COLUMN}}.

\sphinxAtStartPar
The metadata for the LFRic Built\sphinxhyphen{}ins that update an \sphinxcode{\sphinxupquote{integer}}\sphinxhyphen{}valued
field is similar:

\begin{sphinxVerbatim}[commandchars=\\\{\}]
\PYG{c}{!\PYGZgt{} ifield3 = ifield1 + ifield2}
\PYG{k}{type}\PYG{p}{,}\PYG{+w}{ }\PYG{k}{public}\PYG{p}{,}\PYG{+w}{ }\PYG{k}{extends}\PYG{p}{(}\PYG{n}{kernel\PYGZus{}type}\PYG{p}{)}\PYG{+w}{ }\PYG{k+kd}{::}\PYG{+w}{ }\PYG{n}{int\PYGZus{}X\PYGZus{}plus\PYGZus{}Y}
\PYG{+w}{   }\PYG{k}{private}
\PYG{k}{   }\PYG{k}{type}\PYG{p}{(}\PYG{n}{arg\PYGZus{}type}\PYG{p}{)}\PYG{+w}{ }\PYG{k+kd}{::}\PYG{+w}{ }\PYG{n}{meta\PYGZus{}args}\PYG{p}{(}\PYG{l+m+mi}{3}\PYG{p}{)}\PYG{+w}{ }\PYG{o}{=}\PYG{+w}{ }\PYG{p}{(}\PYG{o}{/}\PYG{+w}{                              }\PYG{p}{\PYGZam{}}
\PYG{+w}{        }\PYG{n}{arg\PYGZus{}type}\PYG{p}{(}\PYG{n}{GH\PYGZus{}FIELD}\PYG{p}{,}\PYG{+w}{ }\PYG{n}{GH\PYGZus{}INTEGER}\PYG{p}{,}\PYG{+w}{ }\PYG{n}{GH\PYGZus{}WRITE}\PYG{p}{,}\PYG{+w}{ }\PYG{n}{ANY\PYGZus{}SPACE\PYGZus{}1}\PYG{p}{)}\PYG{p}{,}\PYG{+w}{      }\PYG{p}{\PYGZam{}}
\PYG{+w}{        }\PYG{n}{arg\PYGZus{}type}\PYG{p}{(}\PYG{n}{GH\PYGZus{}FIELD}\PYG{p}{,}\PYG{+w}{ }\PYG{n}{GH\PYGZus{}INTEGER}\PYG{p}{,}\PYG{+w}{ }\PYG{n}{GH\PYGZus{}READ}\PYG{p}{,}\PYG{+w}{  }\PYG{n}{ANY\PYGZus{}SPACE\PYGZus{}1}\PYG{p}{)}\PYG{p}{,}\PYG{+w}{      }\PYG{p}{\PYGZam{}}
\PYG{+w}{        }\PYG{n}{arg\PYGZus{}type}\PYG{p}{(}\PYG{n}{GH\PYGZus{}FIELD}\PYG{p}{,}\PYG{+w}{ }\PYG{n}{GH\PYGZus{}INTEGER}\PYG{p}{,}\PYG{+w}{ }\PYG{n}{GH\PYGZus{}READ}\PYG{p}{,}\PYG{+w}{  }\PYG{n}{ANY\PYGZus{}SPACE\PYGZus{}1}\PYG{p}{)}\PYG{+w}{       }\PYG{p}{\PYGZam{}}
\PYG{+w}{        }\PYG{o}{/}\PYG{p}{)}
\PYG{+w}{   }\PYG{k+kt}{integer}\PYG{+w}{ }\PYG{k+kd}{::}\PYG{+w}{ }\PYG{n}{operates\PYGZus{}on}\PYG{+w}{ }\PYG{o}{=}\PYG{+w}{ }\PYG{n}{DOF}
\PYG{+w}{ }\PYG{k}{contains}
\PYG{k}{   }\PYG{k}{procedure}\PYG{p}{,}\PYG{+w}{ }\PYG{k}{nopass}\PYG{+w}{ }\PYG{k+kd}{::}\PYG{+w}{ }\PYG{n}{int\PYGZus{}X\PYGZus{}plus\PYGZus{}Y\PYGZus{}code}
\PYG{k}{end }\PYG{k}{type }\PYG{n}{int\PYGZus{}X\PYGZus{}plus\PYGZus{}Y}
\end{sphinxVerbatim}


\subsubsection{Valid Data Types and Access Modes}
\label{\detokenize{dynamo0p3:valid-data-types-and-access-modes}}\label{\detokenize{dynamo0p3:lfric-built-ins-dtype-access}}
\sphinxAtStartPar
The allowed data types and accesses for arguments in LFRic Built\sphinxhyphen{}in
kernels are a bit different than for the
{\hyperref[\detokenize{dynamo0p3:lfric-kernel-valid-access}]{\sphinxcrossref{\DUrole{std,std-ref}{user\sphinxhyphen{}defined Kernels}}}} and
are listed in the table below.


\begin{savenotes}\sphinxattablestart
\sphinxthistablewithglobalstyle
\sphinxthistablewithvlinesstyle
\centering
\begin{tabulary}{\linewidth}[t]{|l|l|l|l|}
\sphinxtoprule
\sphinxstyletheadfamily 
\sphinxAtStartPar
Argument Type
&\sphinxstyletheadfamily 
\sphinxAtStartPar
Data Type
&\sphinxstyletheadfamily 
\sphinxAtStartPar
Function Space
&\sphinxstyletheadfamily 
\sphinxAtStartPar
Access Type
\\
\sphinxmidrule
\sphinxtableatstartofbodyhook
\sphinxAtStartPar
GH\_SCALAR
&
\sphinxAtStartPar
GH\_INTEGER
&
\sphinxAtStartPar
n/a
&
\sphinxAtStartPar
GH\_READ
\\
\sphinxhline
\sphinxAtStartPar
GH\_SCALAR
&
\sphinxAtStartPar
GH\_REAL
&
\sphinxAtStartPar
n/a
&
\sphinxAtStartPar
GH\_READ, GH\_SUM
\\
\sphinxhline
\sphinxAtStartPar
GH\_FIELD
&
\sphinxAtStartPar
GH\_REAL, GH\_INTEGER
&
\sphinxAtStartPar
ANY\_SPACE\_\textless{}n\textgreater{}
&
\sphinxAtStartPar
GH\_READ, GH\_WRITE,
GH\_READWRITE
\\
\sphinxbottomrule
\end{tabulary}
\sphinxtableafterendhook\par
\sphinxattableend\end{savenotes}

\begin{sphinxadmonition}{note}{Note:}
\sphinxAtStartPar
Since the LFRic infrastructure does not currently support
\sphinxcode{\sphinxupquote{integer}} reductions, \sphinxcode{\sphinxupquote{integer}} scalar arguments in Built\sphinxhyphen{}ins
are restricted to having read\sphinxhyphen{}only access. Also, \sphinxcode{\sphinxupquote{logical}}
scalar arguments are not permitted.
\end{sphinxadmonition}


\subsection{Naming scheme}
\label{\detokenize{dynamo0p3:naming-scheme}}\label{\detokenize{dynamo0p3:lfric-built-ins-names}}
\sphinxAtStartPar
The supported Built\sphinxhyphen{}ins in the LFRic API are named according to the
scheme presented below. Any new Built\sphinxhyphen{}in needs to comply with these rules.
\begin{enumerate}
\sphinxsetlistlabels{\arabic}{enumi}{enumii}{}{)}%
\item {} 
\sphinxAtStartPar
Ordering of arguments in Built\sphinxhyphen{}ins calls follows
\sphinxstyleemphasis{LHS (result) \textless{}\sphinxhyphen{} RHS (operation on arguments)}
direction, except where a Built\sphinxhyphen{}in returns the \sphinxstyleemphasis{LHS} result to one of
the \sphinxstyleemphasis{RHS} arguments. In that case ordering of arguments remains as in
the \sphinxstyleemphasis{RHS} expression, with the returning \sphinxstyleemphasis{RHS} argument written as close
to the \sphinxstyleemphasis{LHS} as it can be without affecting the mathematical expression.

\item {} 
\sphinxAtStartPar
Field names begin with upper case in short form (e.g. \sphinxstylestrong{X}, \sphinxstylestrong{Y},
\sphinxstylestrong{Z}) and any case in long form (e.g. \sphinxstylestrong{Field1}, \sphinxstylestrong{field}).

\item {} 
\sphinxAtStartPar
Scalar names begin with lower case:  e.g. \sphinxstylestrong{a}, \sphinxstylestrong{b}, are \sphinxstylestrong{scalar1},
\sphinxstylestrong{scalar2}. Special names for scalars are: \sphinxstylestrong{constant} (or \sphinxstylestrong{c}),
\sphinxstylestrong{innprod} (inner/scalar product of two fields) and \sphinxstylestrong{sumfld}
(sum of a field).

\item {} 
\sphinxAtStartPar
Arguments in Built\sphinxhyphen{}ins variable declarations and constructs (PSyclone
Fortran and Python definitions):
\begin{enumerate}
\sphinxsetlistlabels{\arabic}{enumii}{enumiii}{}{)}%
\item {} 
\sphinxAtStartPar
Are always  written in long form and lower case (e.g. \sphinxstylestrong{field1},
\sphinxstylestrong{field2}, \sphinxstylestrong{scalar1}, \sphinxstylestrong{scalar2});

\item {} 
\sphinxAtStartPar
\sphinxstyleemphasis{LHS} result arguments are always listed first;

\item {} 
\sphinxAtStartPar
\sphinxstyleemphasis{RHS} arguments are listed in order of appearance in the mathematical
expression, except when one of them is the \sphinxstyleemphasis{LHS} result.

\end{enumerate}

\item {} 
\sphinxAtStartPar
Built\sphinxhyphen{}ins names in Fortran consist of:
\begin{enumerate}
\sphinxsetlistlabels{\arabic}{enumii}{enumiii}{}{)}%
\item {} 
\sphinxAtStartPar
\sphinxstyleemphasis{RHS} arguments in short form (e.g. \sphinxstylestrong{X}, \sphinxstylestrong{Y}, \sphinxstylestrong{a}, \sphinxstylestrong{b}) only;

\item {} 
\sphinxAtStartPar
Descriptive name of mathematical operation on \sphinxstyleemphasis{RHS} arguments in the
form  \sphinxcode{\sphinxupquote{\textless{}operationname\textgreater{}\_\textless{}RHSargs\textgreater{}}} or
\sphinxcode{\sphinxupquote{\textless{}RHSargs\textgreater{}\_\textless{}operationname\textgreater{}\_\textless{}RHSargs\textgreater{}}};

\item {} 
\sphinxAtStartPar
Prefix \sphinxcode{\sphinxupquote{"inc\_"}} where the result is returned to one of the \sphinxstyleemphasis{RHS}
arguments (i.e. \sphinxcode{\sphinxupquote{"inc\_"\textless{}RHSargs\textgreater{}\_\textless{}operationname\textgreater{}\_\textless{}RHSargs\textgreater{}}});

\item {} 
\sphinxAtStartPar
Prefix \sphinxcode{\sphinxupquote{"int\_"}} for the Built\sphinxhyphen{}in operations on the \sphinxcode{\sphinxupquote{integer}}\sphinxhyphen{}valued
field arguments (i.e. \sphinxcode{\sphinxupquote{"int\_inc\_"\textless{}RHSargs\textgreater{}\_\textless{}operationname\textgreater{}\_\textless{}RHSargs\textgreater{}}}).

\end{enumerate}

\item {} 
\sphinxAtStartPar
Built\sphinxhyphen{}ins names in Python definitions are similar to their Fortran
counterparts, with a few differences:
\begin{enumerate}
\sphinxsetlistlabels{\arabic}{enumii}{enumiii}{}{)}%
\item {} 
\sphinxAtStartPar
Operators and \sphinxstyleemphasis{RHS} arguments are all in upper case (e.g. \sphinxstylestrong{X},
\sphinxstylestrong{Y}, \sphinxstylestrong{A}, \sphinxstylestrong{B}, \sphinxstylestrong{Plus}, \sphinxstylestrong{Minus});

\item {} 
\sphinxAtStartPar
There are no underscores;

\end{enumerate}
\begin{enumerate}
\sphinxsetlistlabels{\arabic}{enumii}{enumiii}{}{)}%
\setcounter{enumii}{3}
\item {} 
\sphinxAtStartPar
Common suffix is \sphinxcode{\sphinxupquote{"Kern"}};

\end{enumerate}
\begin{enumerate}
\sphinxsetlistlabels{\arabic}{enumii}{enumiii}{}{)}%
\setcounter{enumii}{2}
\item {} 
\sphinxAtStartPar
Common prefix is \sphinxcode{\sphinxupquote{"LFRic"}} for the Built\sphinxhyphen{}in operations on the
\sphinxcode{\sphinxupquote{real}}\sphinxhyphen{}valued arguments and \sphinxcode{\sphinxupquote{"LFRicInt"}} for the Built\sphinxhyphen{}in
operations on the \sphinxcode{\sphinxupquote{integer}}\sphinxhyphen{}valued fields.

\end{enumerate}

\end{enumerate}


\subsection{Built\sphinxhyphen{}in operations on \sphinxstyleliteralintitle{\sphinxupquote{real}}\sphinxhyphen{}valued fields}
\label{\detokenize{dynamo0p3:built-in-operations-on-real-valued-fields}}\label{\detokenize{dynamo0p3:lfric-built-ins-real}}
\sphinxAtStartPar
As described {\hyperref[\detokenize{dynamo0p3:lfric-built-ins-dtype-access}]{\sphinxcrossref{\DUrole{std,std-ref}{above}}}}, Built\sphinxhyphen{}ins that
operate on \sphinxcode{\sphinxupquote{real}}\sphinxhyphen{}valued fields mandate \sphinxcode{\sphinxupquote{GH\_REAL}} as the kernel
metadata for fields and scalars.

\sphinxAtStartPar
The precision of fields and scalars, however, is determined by the
algorithm layer via precision variables as described in the {\hyperref[\detokenize{dynamo0p3:lfric-mixed-precision}]{\sphinxcrossref{\DUrole{std,std-ref}{Mixed
Precision}}}} section (see subsections on
{\hyperref[\detokenize{dynamo0p3:lfric-mixed-precision-fields}]{\sphinxcrossref{\DUrole{std,std-ref}{fields}}}} and
{\hyperref[\detokenize{dynamo0p3:lfric-mixed-precision-scalars}]{\sphinxcrossref{\DUrole{std,std-ref}{scalars}}}}).

\sphinxAtStartPar
For instance, field and scalar declarations for the \sphinxcode{\sphinxupquote{aX\_plus\_Y}}
Built\sphinxhyphen{}in that operates on \sphinxcode{\sphinxupquote{r\_solver\_field\_type}} and uses \sphinxcode{\sphinxupquote{r\_solver}}
scalar will be:

\begin{sphinxVerbatim}[commandchars=\\\{\}]
\PYG{k+kt}{real}\PYG{p}{(}\PYG{n+nb}{kind}\PYG{o}{=}\PYG{n}{r\PYGZus{}solver}\PYG{p}{)}\PYG{p}{,}\PYG{+w}{ }\PYG{k}{intent}\PYG{p}{(}\PYG{n}{in}\PYG{p}{)}\PYG{+w}{ }\PYG{k+kd}{::}\PYG{+w}{ }\PYG{n}{ascalar}
\PYG{k}{type}\PYG{p}{(}\PYG{n}{r\PYGZus{}solver\PYGZus{}field\PYGZus{}type}\PYG{p}{)}\PYG{p}{,}\PYG{+w}{ }\PYG{k}{intent}\PYG{p}{(}\PYG{n}{in}\PYG{p}{)}\PYG{+w}{ }\PYG{k+kd}{::}\PYG{+w}{ }\PYG{n}{zfield}\PYG{p}{,}\PYG{+w}{ }\PYG{n}{xfield}\PYG{p}{,}\PYG{+w}{ }\PYG{n}{yfield}
\end{sphinxVerbatim}

\sphinxAtStartPar
Mixing precisions is not explicitly forbidden, so we may have e.g.
\sphinxcode{\sphinxupquote{X\_divideby\_a}} Built\sphinxhyphen{}in where:

\begin{sphinxVerbatim}[commandchars=\\\{\}]
\PYG{k+kt}{real}\PYG{p}{(}\PYG{n+nb}{kind}\PYG{o}{=}\PYG{n}{r\PYGZus{}def}\PYG{p}{)}\PYG{p}{,}\PYG{+w}{ }\PYG{k}{intent}\PYG{p}{(}\PYG{n}{in}\PYG{p}{)}\PYG{+w}{ }\PYG{k+kd}{::}\PYG{+w}{ }\PYG{n}{ascalar}
\PYG{k}{type}\PYG{p}{(}\PYG{n}{r\PYGZus{}tran\PYGZus{}field\PYGZus{}type}\PYG{p}{)}\PYG{p}{,}\PYG{+w}{ }\PYG{k}{intent}\PYG{p}{(}\PYG{n}{in}\PYG{p}{)}\PYG{+w}{ }\PYG{k+kd}{::}\PYG{+w}{ }\PYG{n}{yfield}\PYG{p}{,}\PYG{+w}{ }\PYG{n}{xfield}
\end{sphinxVerbatim}

\sphinxAtStartPar
Certain Built\sphinxhyphen{}ins are currently restricted in the precision of the
arguments that they accept. Those that calculate the inner product
and sum of a field are restricted to \sphinxcode{\sphinxupquote{r\_def}} precision because the
scalar global reductions in the LFRic infrastructure are currently
only able to support \sphinxcode{\sphinxupquote{field\_type}} and hence have \sphinxcode{\sphinxupquote{r\_def}} precision.
In addition, all integer arguments to Built\sphinxhyphen{}ins are currently restricted
to \sphinxcode{\sphinxupquote{i\_def}} precision.


\subsubsection{Addition}
\label{\detokenize{dynamo0p3:addition}}
\sphinxAtStartPar
Built\sphinxhyphen{}ins that add (scaled) \sphinxcode{\sphinxupquote{real}}\sphinxhyphen{}valued fields and return the result
as a \sphinxcode{\sphinxupquote{real}}\sphinxhyphen{}valued field are denoted with the keyword \sphinxstylestrong{plus}.


\paragraph{X\_plus\_Y}
\label{\detokenize{dynamo0p3:x-plus-y}}
\sphinxAtStartPar
\sphinxstylestrong{X\_plus\_Y} (\sphinxstylestrong{field3}, \sphinxstyleemphasis{field1}, \sphinxstyleemphasis{field2})

\sphinxAtStartPar
Sums two fields and stores the result in the third field (\sphinxcode{\sphinxupquote{Z = X + Y}}):

\begin{sphinxVerbatim}[commandchars=\\\{\}]
\PYG{n}{field3}\PYG{p}{(}\PYG{p}{:}\PYG{p}{)}\PYG{+w}{ }\PYG{o}{=}\PYG{+w}{ }\PYG{n}{field1}\PYG{p}{(}\PYG{p}{:}\PYG{p}{)}\PYG{+w}{ }\PYG{o}{+}\PYG{+w}{ }\PYG{n}{field2}\PYG{p}{(}\PYG{p}{:}\PYG{p}{)}
\end{sphinxVerbatim}


\paragraph{inc\_X\_plus\_Y}
\label{\detokenize{dynamo0p3:inc-x-plus-y}}
\sphinxAtStartPar
\sphinxstylestrong{inc\_X\_plus\_Y} (\sphinxstylestrong{field1}, \sphinxstyleemphasis{field2})

\sphinxAtStartPar
Adds the second field to the first and returns it (\sphinxcode{\sphinxupquote{X = X + Y}}):

\begin{sphinxVerbatim}[commandchars=\\\{\}]
\PYG{n}{field1}\PYG{p}{(}\PYG{p}{:}\PYG{p}{)}\PYG{+w}{ }\PYG{o}{=}\PYG{+w}{ }\PYG{n}{field1}\PYG{p}{(}\PYG{p}{:}\PYG{p}{)}\PYG{+w}{ }\PYG{o}{+}\PYG{+w}{ }\PYG{n}{field2}\PYG{p}{(}\PYG{p}{:}\PYG{p}{)}
\end{sphinxVerbatim}


\paragraph{a\_plus\_X}
\label{\detokenize{dynamo0p3:a-plus-x}}
\sphinxAtStartPar
\sphinxstylestrong{a\_plus\_X} (\sphinxstylestrong{field2}, \sphinxstyleemphasis{rscalar}, \sphinxstyleemphasis{field1})

\sphinxAtStartPar
Adds a \sphinxcode{\sphinxupquote{real}} scalar value to all elements of a field and stores
the result in another field (\sphinxcode{\sphinxupquote{Y = a + X}}):

\begin{sphinxVerbatim}[commandchars=\\\{\}]
\PYG{n}{field2}\PYG{p}{(}\PYG{p}{:}\PYG{p}{)}\PYG{+w}{ }\PYG{o}{=}\PYG{+w}{ }\PYG{n}{rscalar}\PYG{+w}{ }\PYG{o}{+}\PYG{+w}{ }\PYG{n}{field1}\PYG{p}{(}\PYG{p}{:}\PYG{p}{)}
\end{sphinxVerbatim}


\paragraph{inc\_a\_plus\_X}
\label{\detokenize{dynamo0p3:inc-a-plus-x}}
\sphinxAtStartPar
\sphinxstylestrong{inc\_a\_plus\_X} (\sphinxstyleemphasis{rscalar}, \sphinxstylestrong{field})

\sphinxAtStartPar
Adds a \sphinxcode{\sphinxupquote{real}} scalar value to all elements of a field and returns
the field (\sphinxcode{\sphinxupquote{X = a + X}}):

\begin{sphinxVerbatim}[commandchars=\\\{\}]
\PYG{n}{field}\PYG{p}{(}\PYG{p}{:}\PYG{p}{)}\PYG{+w}{ }\PYG{o}{=}\PYG{+w}{ }\PYG{n}{rscalar}\PYG{+w}{ }\PYG{o}{+}\PYG{+w}{ }\PYG{n}{field}\PYG{p}{(}\PYG{p}{:}\PYG{p}{)}
\end{sphinxVerbatim}


\paragraph{aX\_plus\_Y}
\label{\detokenize{dynamo0p3:ax-plus-y}}
\sphinxAtStartPar
\sphinxstylestrong{aX\_plus\_Y} (\sphinxstylestrong{field3}, \sphinxstyleemphasis{rscalar}, \sphinxstyleemphasis{field1}, \sphinxstyleemphasis{field2})

\sphinxAtStartPar
Performs \sphinxcode{\sphinxupquote{Z = aX + Y}}:

\begin{sphinxVerbatim}[commandchars=\\\{\}]
\PYG{n}{field3}\PYG{p}{(}\PYG{p}{:}\PYG{p}{)}\PYG{+w}{ }\PYG{o}{=}\PYG{+w}{ }\PYG{n}{rscalar}\PYG{o}{*}\PYG{n}{field1}\PYG{p}{(}\PYG{p}{:}\PYG{p}{)}\PYG{+w}{ }\PYG{o}{+}\PYG{+w}{ }\PYG{n}{field2}\PYG{p}{(}\PYG{p}{:}\PYG{p}{)}
\end{sphinxVerbatim}


\paragraph{inc\_aX\_plus\_Y}
\label{\detokenize{dynamo0p3:inc-ax-plus-y}}
\sphinxAtStartPar
\sphinxstylestrong{inc\_aX\_plus\_Y} (\sphinxstyleemphasis{rscalar}, \sphinxstylestrong{field1}, \sphinxstyleemphasis{field2})

\sphinxAtStartPar
Performs \sphinxcode{\sphinxupquote{X = aX + Y}} (increments the first field):

\begin{sphinxVerbatim}[commandchars=\\\{\}]
\PYG{n}{field1}\PYG{p}{(}\PYG{p}{:}\PYG{p}{)}\PYG{+w}{ }\PYG{o}{=}\PYG{+w}{ }\PYG{n}{rscalar}\PYG{o}{*}\PYG{n}{field1}\PYG{p}{(}\PYG{p}{:}\PYG{p}{)}\PYG{+w}{ }\PYG{o}{+}\PYG{+w}{ }\PYG{n}{field2}\PYG{p}{(}\PYG{p}{:}\PYG{p}{)}
\end{sphinxVerbatim}


\paragraph{inc\_X\_plus\_bY}
\label{\detokenize{dynamo0p3:inc-x-plus-by}}
\sphinxAtStartPar
\sphinxstylestrong{inc\_X\_plus\_bY} (\sphinxstylestrong{field1}, \sphinxstyleemphasis{rscalar}, \sphinxstyleemphasis{field2})

\sphinxAtStartPar
Performs \sphinxcode{\sphinxupquote{X = X + bY}} (increments the first field):

\begin{sphinxVerbatim}[commandchars=\\\{\}]
\PYG{n}{field1}\PYG{p}{(}\PYG{p}{:}\PYG{p}{)}\PYG{+w}{ }\PYG{o}{=}\PYG{+w}{ }\PYG{n}{field1}\PYG{p}{(}\PYG{p}{:}\PYG{p}{)}\PYG{+w}{ }\PYG{o}{+}\PYG{+w}{ }\PYG{n}{rscalar}\PYG{o}{*}\PYG{n}{field2}\PYG{p}{(}\PYG{p}{:}\PYG{p}{)}
\end{sphinxVerbatim}


\paragraph{aX\_plus\_bY}
\label{\detokenize{dynamo0p3:ax-plus-by}}
\sphinxAtStartPar
\sphinxstylestrong{aX\_plus\_bY} (\sphinxstylestrong{field3}, \sphinxstyleemphasis{rscalar1}, \sphinxstyleemphasis{field1}, \sphinxstyleemphasis{rscalar2}, \sphinxstyleemphasis{field2})

\sphinxAtStartPar
Performs \sphinxcode{\sphinxupquote{Z = aX + bY}}:

\begin{sphinxVerbatim}[commandchars=\\\{\}]
\PYG{n}{field3}\PYG{p}{(}\PYG{p}{:}\PYG{p}{)}\PYG{+w}{ }\PYG{o}{=}\PYG{+w}{ }\PYG{n}{rscalar1}\PYG{o}{*}\PYG{n}{field1}\PYG{p}{(}\PYG{p}{:}\PYG{p}{)}\PYG{+w}{ }\PYG{o}{+}\PYG{+w}{ }\PYG{n}{rscalar2}\PYG{o}{*}\PYG{n}{field2}\PYG{p}{(}\PYG{p}{:}\PYG{p}{)}
\end{sphinxVerbatim}


\paragraph{inc\_aX\_plus\_bY}
\label{\detokenize{dynamo0p3:inc-ax-plus-by}}
\sphinxAtStartPar
\sphinxstylestrong{inc\_aX\_plus\_bY} (\sphinxstyleemphasis{rscalar1}, \sphinxstylestrong{field1}, \sphinxstyleemphasis{rscalar2}, \sphinxstyleemphasis{field2})

\sphinxAtStartPar
Performs \sphinxcode{\sphinxupquote{X = aX + bY}} (increments the first field):

\begin{sphinxVerbatim}[commandchars=\\\{\}]
\PYG{n}{field1}\PYG{p}{(}\PYG{p}{:}\PYG{p}{)}\PYG{+w}{ }\PYG{o}{=}\PYG{+w}{ }\PYG{n}{rscalar1}\PYG{o}{*}\PYG{n}{field1}\PYG{p}{(}\PYG{p}{:}\PYG{p}{)}\PYG{+w}{ }\PYG{o}{+}\PYG{+w}{ }\PYG{n}{rscalar2}\PYG{o}{*}\PYG{n}{field2}\PYG{p}{(}\PYG{p}{:}\PYG{p}{)}
\end{sphinxVerbatim}


\paragraph{aX\_plus\_aY}
\label{\detokenize{dynamo0p3:ax-plus-ay}}
\sphinxAtStartPar
\sphinxstylestrong{aX\_plus\_aY} (\sphinxstylestrong{field3}, \sphinxstyleemphasis{rscalar}, \sphinxstyleemphasis{field1}, \sphinxstyleemphasis{field2})

\sphinxAtStartPar
Performs \sphinxcode{\sphinxupquote{Z = aX + aY = a(X + Y)}}:

\begin{sphinxVerbatim}[commandchars=\\\{\}]
\PYG{n}{field3}\PYG{p}{(}\PYG{p}{:}\PYG{p}{)}\PYG{+w}{ }\PYG{o}{=}\PYG{+w}{ }\PYG{n}{rscalar}\PYG{o}{*}\PYG{p}{(}\PYG{n}{field1}\PYG{p}{(}\PYG{p}{:}\PYG{p}{)}\PYG{+w}{ }\PYG{o}{+}\PYG{+w}{ }\PYG{n}{field2}\PYG{p}{(}\PYG{p}{:}\PYG{p}{)}\PYG{p}{)}
\end{sphinxVerbatim}


\subsubsection{Subtraction}
\label{\detokenize{dynamo0p3:subtraction}}
\sphinxAtStartPar
Built\sphinxhyphen{}ins which subtract (scaled) \sphinxcode{\sphinxupquote{real}}\sphinxhyphen{}valued  fields and return the
result as a \sphinxcode{\sphinxupquote{real}}\sphinxhyphen{}valued field are denoted with the keyword \sphinxstylestrong{minus}.


\paragraph{X\_minus\_Y}
\label{\detokenize{dynamo0p3:x-minus-y}}
\sphinxAtStartPar
\sphinxstylestrong{X\_minus\_Y} (\sphinxstylestrong{field3}, \sphinxstyleemphasis{field1}, \sphinxstyleemphasis{field2})

\sphinxAtStartPar
Subtracts the second field from the first and returns the result in the
third field (\sphinxcode{\sphinxupquote{Z = X \sphinxhyphen{} Y}}):

\begin{sphinxVerbatim}[commandchars=\\\{\}]
\PYG{n}{field3}\PYG{p}{(}\PYG{p}{:}\PYG{p}{)}\PYG{+w}{ }\PYG{o}{=}\PYG{+w}{ }\PYG{n}{field1}\PYG{p}{(}\PYG{p}{:}\PYG{p}{)}\PYG{+w}{ }\PYG{o}{\PYGZhy{}}\PYG{+w}{ }\PYG{n}{field2}\PYG{p}{(}\PYG{p}{:}\PYG{p}{)}
\end{sphinxVerbatim}


\paragraph{inc\_X\_minus\_Y}
\label{\detokenize{dynamo0p3:inc-x-minus-y}}
\sphinxAtStartPar
\sphinxstylestrong{inc\_X\_minus\_Y} (\sphinxstylestrong{field1}, \sphinxstyleemphasis{field2})

\sphinxAtStartPar
Subtracts the second field from the first and returns it (\sphinxcode{\sphinxupquote{X = X \sphinxhyphen{} Y}}):

\begin{sphinxVerbatim}[commandchars=\\\{\}]
\PYG{n}{field1}\PYG{p}{(}\PYG{p}{:}\PYG{p}{)}\PYG{+w}{ }\PYG{o}{=}\PYG{+w}{ }\PYG{n}{field1}\PYG{p}{(}\PYG{p}{:}\PYG{p}{)}\PYG{+w}{ }\PYG{o}{\PYGZhy{}}\PYG{+w}{ }\PYG{n}{field2}\PYG{p}{(}\PYG{p}{:}\PYG{p}{)}
\end{sphinxVerbatim}


\paragraph{a\_minus\_X}
\label{\detokenize{dynamo0p3:a-minus-x}}
\sphinxAtStartPar
\sphinxstylestrong{a\_minus\_X} (\sphinxstylestrong{field2}, \sphinxstyleemphasis{rscalar}, \sphinxstyleemphasis{field1})

\sphinxAtStartPar
Subtracts all elements of a field from a \sphinxcode{\sphinxupquote{real}} scalar value and
stores the result in another field (\sphinxcode{\sphinxupquote{Y = a \sphinxhyphen{} X}}):

\begin{sphinxVerbatim}[commandchars=\\\{\}]
\PYG{n}{field2}\PYG{p}{(}\PYG{p}{:}\PYG{p}{)}\PYG{+w}{ }\PYG{o}{=}\PYG{+w}{ }\PYG{n}{rscalar}\PYG{+w}{ }\PYG{o}{\PYGZhy{}}\PYG{+w}{ }\PYG{n}{field1}\PYG{p}{(}\PYG{p}{:}\PYG{p}{)}
\end{sphinxVerbatim}


\paragraph{inc\_a\_minus\_X}
\label{\detokenize{dynamo0p3:inc-a-minus-x}}
\sphinxAtStartPar
\sphinxstylestrong{inc\_a\_minus\_X} (\sphinxstyleemphasis{rscalar}, \sphinxstylestrong{field})

\sphinxAtStartPar
Subtracts all elements of a field from a \sphinxcode{\sphinxupquote{real}} scalar value and
returns the field (\sphinxcode{\sphinxupquote{X = a \sphinxhyphen{} X}}):

\begin{sphinxVerbatim}[commandchars=\\\{\}]
\PYG{n}{field}\PYG{p}{(}\PYG{p}{:}\PYG{p}{)}\PYG{+w}{ }\PYG{o}{=}\PYG{+w}{ }\PYG{n}{rscalar}\PYG{+w}{ }\PYG{o}{\PYGZhy{}}\PYG{+w}{ }\PYG{n}{field}\PYG{p}{(}\PYG{p}{:}\PYG{p}{)}
\end{sphinxVerbatim}


\paragraph{X\_minus\_a}
\label{\detokenize{dynamo0p3:x-minus-a}}
\sphinxAtStartPar
\sphinxstylestrong{X\_minus\_a} (\sphinxstylestrong{field2}, \sphinxstyleemphasis{field1}, \sphinxstyleemphasis{rscalar})

\sphinxAtStartPar
Subtracts a \sphinxcode{\sphinxupquote{real}} scalar value from all elements of a field and
stores the result in another field (\sphinxcode{\sphinxupquote{Y = X \sphinxhyphen{} a}}):

\begin{sphinxVerbatim}[commandchars=\\\{\}]
\PYG{n}{field2}\PYG{p}{(}\PYG{p}{:}\PYG{p}{)}\PYG{+w}{ }\PYG{o}{=}\PYG{+w}{ }\PYG{n}{field1}\PYG{p}{(}\PYG{p}{:}\PYG{p}{)}\PYG{+w}{ }\PYG{o}{\PYGZhy{}}\PYG{+w}{ }\PYG{n}{rscalar}
\end{sphinxVerbatim}


\paragraph{inc\_X\_minus\_a}
\label{\detokenize{dynamo0p3:inc-x-minus-a}}
\sphinxAtStartPar
\sphinxstylestrong{inc\_X\_minus\_a} (\sphinxstylestrong{field}, \sphinxstyleemphasis{rscalar})

\sphinxAtStartPar
Subtracts a \sphinxcode{\sphinxupquote{real}} scalar value from all elements of a field and
returns the field (\sphinxcode{\sphinxupquote{X = X \sphinxhyphen{} a}}):

\begin{sphinxVerbatim}[commandchars=\\\{\}]
\PYG{n}{field}\PYG{p}{(}\PYG{p}{:}\PYG{p}{)}\PYG{+w}{ }\PYG{o}{=}\PYG{+w}{ }\PYG{n}{field}\PYG{p}{(}\PYG{p}{:}\PYG{p}{)}\PYG{+w}{ }\PYG{o}{\PYGZhy{}}\PYG{+w}{ }\PYG{n}{rscalar}
\end{sphinxVerbatim}


\paragraph{aX\_minus\_Y}
\label{\detokenize{dynamo0p3:ax-minus-y}}
\sphinxAtStartPar
\sphinxstylestrong{aX\_minus\_Y} (\sphinxstylestrong{field3}, \sphinxstyleemphasis{rscalar}, \sphinxstyleemphasis{field1}, \sphinxstyleemphasis{field2})

\sphinxAtStartPar
Performs \sphinxcode{\sphinxupquote{Z = aX \sphinxhyphen{} Y}}:

\begin{sphinxVerbatim}[commandchars=\\\{\}]
\PYG{n}{field3}\PYG{p}{(}\PYG{p}{:}\PYG{p}{)}\PYG{+w}{ }\PYG{o}{=}\PYG{+w}{ }\PYG{n}{rscalar}\PYG{o}{*}\PYG{n}{field1}\PYG{p}{(}\PYG{p}{:}\PYG{p}{)}\PYG{+w}{ }\PYG{o}{\PYGZhy{}}\PYG{+w}{ }\PYG{n}{field2}\PYG{p}{(}\PYG{p}{:}\PYG{p}{)}
\end{sphinxVerbatim}


\paragraph{X\_minus\_bY}
\label{\detokenize{dynamo0p3:x-minus-by}}
\sphinxAtStartPar
\sphinxstylestrong{X\_minus\_bY} (\sphinxstylestrong{field3}, \sphinxstyleemphasis{field1}, \sphinxstyleemphasis{rscalar}, \sphinxstyleemphasis{field2})

\sphinxAtStartPar
Performs \sphinxcode{\sphinxupquote{Z = X \sphinxhyphen{} bY}}:

\begin{sphinxVerbatim}[commandchars=\\\{\}]
\PYG{n}{field3}\PYG{p}{(}\PYG{p}{:}\PYG{p}{)}\PYG{+w}{ }\PYG{o}{=}\PYG{+w}{ }\PYG{n}{field1}\PYG{p}{(}\PYG{p}{:}\PYG{p}{)}\PYG{+w}{ }\PYG{o}{\PYGZhy{}}\PYG{+w}{ }\PYG{n}{rscalar}\PYG{o}{*}\PYG{n}{field2}\PYG{p}{(}\PYG{p}{:}\PYG{p}{)}
\end{sphinxVerbatim}


\paragraph{inc\_X\_minus\_bY}
\label{\detokenize{dynamo0p3:inc-x-minus-by}}
\sphinxAtStartPar
\sphinxstylestrong{inc\_X\_minus\_bY} (\sphinxstylestrong{field1}, \sphinxstyleemphasis{rscalar}, \sphinxstyleemphasis{field2})

\sphinxAtStartPar
Performs \sphinxcode{\sphinxupquote{X = X \sphinxhyphen{} bY}} (decrements the first field):

\begin{sphinxVerbatim}[commandchars=\\\{\}]
\PYG{n}{field1}\PYG{p}{(}\PYG{p}{:}\PYG{p}{)}\PYG{+w}{ }\PYG{o}{=}\PYG{+w}{ }\PYG{n}{field1}\PYG{p}{(}\PYG{p}{:}\PYG{p}{)}\PYG{+w}{ }\PYG{o}{\PYGZhy{}}\PYG{+w}{ }\PYG{n}{rscalar}\PYG{o}{*}\PYG{n}{field2}\PYG{p}{(}\PYG{p}{:}\PYG{p}{)}
\end{sphinxVerbatim}


\paragraph{aX\_minus\_bY}
\label{\detokenize{dynamo0p3:ax-minus-by}}
\sphinxAtStartPar
\sphinxstylestrong{aX\_minus\_bY} (\sphinxstylestrong{field3}, \sphinxstyleemphasis{rscalar1}, \sphinxstyleemphasis{field1}, \sphinxstyleemphasis{rscalar2}, \sphinxstyleemphasis{field2})

\sphinxAtStartPar
Performs \sphinxcode{\sphinxupquote{Z = aX \sphinxhyphen{} bY}}:

\begin{sphinxVerbatim}[commandchars=\\\{\}]
\PYG{n}{field3}\PYG{p}{(}\PYG{p}{:}\PYG{p}{)}\PYG{+w}{ }\PYG{o}{=}\PYG{+w}{ }\PYG{n}{rscalar1}\PYG{o}{*}\PYG{n}{field1}\PYG{p}{(}\PYG{p}{:}\PYG{p}{)}\PYG{+w}{ }\PYG{o}{\PYGZhy{}}\PYG{+w}{ }\PYG{n}{rscalar2}\PYG{o}{*}\PYG{n}{field2}\PYG{p}{(}\PYG{p}{:}\PYG{p}{)}
\end{sphinxVerbatim}


\subsubsection{Multiplication}
\label{\detokenize{dynamo0p3:multiplication}}
\sphinxAtStartPar
Built\sphinxhyphen{}ins which multiply (scaled) \sphinxcode{\sphinxupquote{real}}\sphinxhyphen{}valued fields and return the
result as a \sphinxcode{\sphinxupquote{real}}\sphinxhyphen{}valued field are denoted with the keyword \sphinxstylestrong{times}.


\paragraph{X\_times\_Y}
\label{\detokenize{dynamo0p3:x-times-y}}
\sphinxAtStartPar
\sphinxstylestrong{X\_times\_Y} (\sphinxstylestrong{field3}, \sphinxstyleemphasis{field1}, \sphinxstyleemphasis{field2})

\sphinxAtStartPar
Multiplies two fields DoF by DoF and returns the result in a
third field (\sphinxcode{\sphinxupquote{Z = X*Y}}):

\begin{sphinxVerbatim}[commandchars=\\\{\}]
\PYG{n}{field3}\PYG{p}{(}\PYG{p}{:}\PYG{p}{)}\PYG{+w}{ }\PYG{o}{=}\PYG{+w}{ }\PYG{n}{field1}\PYG{p}{(}\PYG{p}{:}\PYG{p}{)}\PYG{o}{*}\PYG{n}{field2}\PYG{p}{(}\PYG{p}{:}\PYG{p}{)}
\end{sphinxVerbatim}


\paragraph{inc\_X\_times\_Y}
\label{\detokenize{dynamo0p3:inc-x-times-y}}
\sphinxAtStartPar
\sphinxstylestrong{inc\_X\_times\_Y} (\sphinxstylestrong{field1}, \sphinxstyleemphasis{field2})

\sphinxAtStartPar
Multiplies the first field by the second and returns it (\sphinxcode{\sphinxupquote{X = X*Y}}):

\begin{sphinxVerbatim}[commandchars=\\\{\}]
\PYG{n}{field1}\PYG{p}{(}\PYG{p}{:}\PYG{p}{)}\PYG{+w}{ }\PYG{o}{=}\PYG{+w}{ }\PYG{n}{field1}\PYG{p}{(}\PYG{p}{:}\PYG{p}{)}\PYG{o}{*}\PYG{n}{field2}\PYG{p}{(}\PYG{p}{:}\PYG{p}{)}
\end{sphinxVerbatim}


\paragraph{inc\_aX\_times\_Y}
\label{\detokenize{dynamo0p3:inc-ax-times-y}}
\sphinxAtStartPar
\sphinxstylestrong{inc\_aX\_times\_Y} (\sphinxstyleemphasis{rscalar}, \sphinxstylestrong{field1}, \sphinxstyleemphasis{field2})

\sphinxAtStartPar
Performs \sphinxcode{\sphinxupquote{X = a*X*Y}} (increments the first field):

\begin{sphinxVerbatim}[commandchars=\\\{\}]
\PYG{n}{field1}\PYG{p}{(}\PYG{p}{:}\PYG{p}{)}\PYG{+w}{ }\PYG{o}{=}\PYG{+w}{ }\PYG{n}{rscalar}\PYG{o}{*}\PYG{n}{field1}\PYG{p}{(}\PYG{p}{:}\PYG{p}{)}\PYG{o}{*}\PYG{n}{field2}\PYG{p}{(}\PYG{p}{:}\PYG{p}{)}
\end{sphinxVerbatim}


\subsubsection{Scaling}
\label{\detokenize{dynamo0p3:scaling}}
\sphinxAtStartPar
Built\sphinxhyphen{}ins which scale \sphinxcode{\sphinxupquote{real}}\sphinxhyphen{}valued fields are technically cases of
multiplying a \sphinxcode{\sphinxupquote{real}}\sphinxhyphen{}valued field by a \sphinxcode{\sphinxupquote{real}} scalar and are hence
also denoted with the keyword \sphinxstylestrong{times}.


\paragraph{a\_times\_X}
\label{\detokenize{dynamo0p3:a-times-x}}
\sphinxAtStartPar
\sphinxstylestrong{a\_times\_X} (\sphinxstylestrong{field2}, \sphinxstyleemphasis{rscalar}, \sphinxstyleemphasis{field1})

\sphinxAtStartPar
Multiplies a field by a \sphinxcode{\sphinxupquote{real}} scalar value and stores the result
in another field (\sphinxcode{\sphinxupquote{Y = a*X}}):

\begin{sphinxVerbatim}[commandchars=\\\{\}]
\PYG{n}{field2}\PYG{p}{(}\PYG{p}{:}\PYG{p}{)}\PYG{+w}{ }\PYG{o}{=}\PYG{+w}{ }\PYG{n}{rscalar}\PYG{o}{*}\PYG{n}{field1}\PYG{p}{(}\PYG{p}{:}\PYG{p}{)}
\end{sphinxVerbatim}


\paragraph{inc\_a\_times\_X}
\label{\detokenize{dynamo0p3:inc-a-times-x}}
\sphinxAtStartPar
\sphinxstylestrong{inc\_a\_times\_X} (\sphinxstyleemphasis{rscalar}, \sphinxstylestrong{field})

\sphinxAtStartPar
Multiplies a field by a \sphinxcode{\sphinxupquote{real}} scalar value and returns the
field (\sphinxcode{\sphinxupquote{X = a*X}}):

\begin{sphinxVerbatim}[commandchars=\\\{\}]
\PYG{n}{field}\PYG{p}{(}\PYG{p}{:}\PYG{p}{)}\PYG{+w}{ }\PYG{o}{=}\PYG{+w}{ }\PYG{n}{rscalar}\PYG{o}{*}\PYG{n}{field}\PYG{p}{(}\PYG{p}{:}\PYG{p}{)}
\end{sphinxVerbatim}


\subsubsection{Division}
\label{\detokenize{dynamo0p3:division}}
\sphinxAtStartPar
Built\sphinxhyphen{}ins which divide \sphinxcode{\sphinxupquote{real}}\sphinxhyphen{}valued fields and return the result
as a \sphinxcode{\sphinxupquote{real}}\sphinxhyphen{}valued field are denoted with the keyword \sphinxstylestrong{divideby}.


\paragraph{X\_divideby\_Y}
\label{\detokenize{dynamo0p3:x-divideby-y}}
\sphinxAtStartPar
\sphinxstylestrong{X\_divideby\_Y} (\sphinxstylestrong{field3}, \sphinxstyleemphasis{field1}, \sphinxstyleemphasis{field2})

\sphinxAtStartPar
Divides the first field by the second field, DoF by DoF, and stores the
result in the third field (\sphinxcode{\sphinxupquote{Z = X/Y}}):

\begin{sphinxVerbatim}[commandchars=\\\{\}]
\PYG{n}{field3}\PYG{p}{(}\PYG{p}{:}\PYG{p}{)}\PYG{+w}{ }\PYG{o}{=}\PYG{+w}{ }\PYG{n}{field1}\PYG{p}{(}\PYG{p}{:}\PYG{p}{)}\PYG{o}{/}\PYG{n}{field2}\PYG{p}{(}\PYG{p}{:}\PYG{p}{)}
\end{sphinxVerbatim}


\paragraph{inc\_X\_divideby\_Y}
\label{\detokenize{dynamo0p3:inc-x-divideby-y}}
\sphinxAtStartPar
\sphinxstylestrong{inc\_X\_divideby\_Y} (\sphinxstylestrong{field1}, \sphinxstyleemphasis{field2})

\sphinxAtStartPar
Divides the first field by the second and returns it (\sphinxcode{\sphinxupquote{X = X/Y}}):

\begin{sphinxVerbatim}[commandchars=\\\{\}]
\PYG{n}{field1}\PYG{p}{(}\PYG{p}{:}\PYG{p}{)}\PYG{+w}{ }\PYG{o}{=}\PYG{+w}{ }\PYG{n}{field1}\PYG{p}{(}\PYG{p}{:}\PYG{p}{)}\PYG{o}{/}\PYG{n}{field2}\PYG{p}{(}\PYG{p}{:}\PYG{p}{)}
\end{sphinxVerbatim}


\paragraph{X\_divideby\_a}
\label{\detokenize{dynamo0p3:x-divideby-a}}
\sphinxAtStartPar
\sphinxstylestrong{X\_divideby\_a} (\sphinxstylestrong{field2}, \sphinxstyleemphasis{field1}, \sphinxstyleemphasis{rscalar})

\sphinxAtStartPar
Divides each field element by a \sphinxcode{\sphinxupquote{real}} scalar value and stores
the result in another field (\sphinxcode{\sphinxupquote{Y = X/a}}):

\begin{sphinxVerbatim}[commandchars=\\\{\}]
\PYG{n}{field2}\PYG{p}{(}\PYG{p}{:}\PYG{p}{)}\PYG{+w}{ }\PYG{o}{=}\PYG{+w}{ }\PYG{n}{field1}\PYG{p}{(}\PYG{p}{:}\PYG{p}{)}\PYG{o}{/}\PYG{n}{rscalar}
\end{sphinxVerbatim}


\paragraph{inc\_X\_divideby\_a}
\label{\detokenize{dynamo0p3:inc-x-divideby-a}}
\sphinxAtStartPar
\sphinxstylestrong{inc\_X\_divideby\_a} (\sphinxstylestrong{field}, \sphinxstyleemphasis{rscalar})

\sphinxAtStartPar
Divides each field element by a \sphinxcode{\sphinxupquote{real}} scalar value and returns
the field (\sphinxcode{\sphinxupquote{X = X/a}}):

\begin{sphinxVerbatim}[commandchars=\\\{\}]
\PYG{n}{field}\PYG{p}{(}\PYG{p}{:}\PYG{p}{)}\PYG{+w}{ }\PYG{o}{=}\PYG{+w}{ }\PYG{n}{field}\PYG{p}{(}\PYG{p}{:}\PYG{p}{)}\PYG{o}{/}\PYG{n}{rscalar}
\end{sphinxVerbatim}


\subsubsection{Inverse scaling}
\label{\detokenize{dynamo0p3:inverse-scaling}}
\sphinxAtStartPar
Built\sphinxhyphen{}ins which perform inverse scaling of \sphinxcode{\sphinxupquote{real}}\sphinxhyphen{}valued fields are
also denoted with the keyword \sphinxstylestrong{divideby} as they divide a \sphinxcode{\sphinxupquote{real}}
scalar by elements of a \sphinxcode{\sphinxupquote{real}}\sphinxhyphen{}valued field.


\paragraph{a\_divideby\_X}
\label{\detokenize{dynamo0p3:a-divideby-x}}
\sphinxAtStartPar
\sphinxstylestrong{a\_divideby\_X} (\sphinxstylestrong{field2}, \sphinxstyleemphasis{rscalar}, \sphinxstyleemphasis{field1})

\sphinxAtStartPar
Divides a \sphinxcode{\sphinxupquote{real}} scalar value by each field element and stores the
result in another field (\sphinxcode{\sphinxupquote{Y = a/X}}):

\begin{sphinxVerbatim}[commandchars=\\\{\}]
\PYG{n}{field2}\PYG{p}{(}\PYG{p}{:}\PYG{p}{)}\PYG{+w}{ }\PYG{o}{=}\PYG{+w}{ }\PYG{n}{rscalar}\PYG{o}{/}\PYG{n}{field1}\PYG{p}{(}\PYG{p}{:}\PYG{p}{)}
\end{sphinxVerbatim}


\paragraph{inc\_a\_divideby\_X}
\label{\detokenize{dynamo0p3:inc-a-divideby-x}}
\sphinxAtStartPar
\sphinxstylestrong{inc\_a\_divideby\_X} (\sphinxstyleemphasis{rscalar}, \sphinxstylestrong{field})

\sphinxAtStartPar
Divides a \sphinxcode{\sphinxupquote{real}} scalar value by each field element and returns
the field (\sphinxcode{\sphinxupquote{X = a/X}}):

\begin{sphinxVerbatim}[commandchars=\\\{\}]
\PYG{n}{field}\PYG{p}{(}\PYG{p}{:}\PYG{p}{)}\PYG{+w}{ }\PYG{o}{=}\PYG{+w}{ }\PYG{n}{rscalar}\PYG{o}{/}\PYG{n}{field}\PYG{p}{(}\PYG{p}{:}\PYG{p}{)}
\end{sphinxVerbatim}


\subsubsection{Setting to a value}
\label{\detokenize{dynamo0p3:setting-to-a-value}}
\sphinxAtStartPar
Built\sphinxhyphen{}ins which set \sphinxcode{\sphinxupquote{real}}\sphinxhyphen{}valued field elements to some \sphinxcode{\sphinxupquote{real}}
value are denoted with the keyword \sphinxstylestrong{setval}.


\paragraph{setval\_c}
\label{\detokenize{dynamo0p3:setval-c}}
\sphinxAtStartPar
\sphinxstylestrong{setval\_c} (\sphinxstylestrong{field}, \sphinxstyleemphasis{constant})

\sphinxAtStartPar
Sets all elements of a field \sphinxstyleemphasis{field} to a \sphinxcode{\sphinxupquote{real}} scalar
\sphinxstyleemphasis{constant} (\sphinxcode{\sphinxupquote{X = c}}):

\begin{sphinxVerbatim}[commandchars=\\\{\}]
\PYG{n}{field}\PYG{p}{(}\PYG{p}{:}\PYG{p}{)}\PYG{+w}{ }\PYG{o}{=}\PYG{+w}{ }\PYG{n}{constant}
\end{sphinxVerbatim}


\paragraph{setval\_X}
\label{\detokenize{dynamo0p3:setval-x}}
\sphinxAtStartPar
\sphinxstylestrong{setval\_X} (\sphinxstylestrong{field2}, \sphinxstyleemphasis{field1})

\sphinxAtStartPar
Sets a field \sphinxstyleemphasis{field2} equal (DoF per DoF) to another field
\sphinxstyleemphasis{field1} (\sphinxcode{\sphinxupquote{Y = X}}):

\begin{sphinxVerbatim}[commandchars=\\\{\}]
\PYG{n}{field2}\PYG{p}{(}\PYG{p}{:}\PYG{p}{)}\PYG{+w}{ }\PYG{o}{=}\PYG{+w}{ }\PYG{n}{field1}\PYG{p}{(}\PYG{p}{:}\PYG{p}{)}
\end{sphinxVerbatim}


\paragraph{setval\_random}
\label{\detokenize{dynamo0p3:setval-random}}
\sphinxAtStartPar
\sphinxstylestrong{setval\_random} (\sphinxstylestrong{field})

\sphinxAtStartPar
Fills all elements of a field \sphinxstyleemphasis{field} using a sequence of \sphinxcode{\sphinxupquote{real}},
pseudo\sphinxhyphen{}random numbers in the interval \sphinxcode{\sphinxupquote{0 \textless{}= x \textless{} 1}}:

\begin{sphinxVerbatim}[commandchars=\\\{\}]
\PYG{k}{do }\PYG{n}{df}\PYG{+w}{ }\PYG{o}{=}\PYG{+w}{ }\PYG{l+m+mi}{1}\PYG{p}{,}\PYG{+w}{ }\PYG{n}{ndofs}
\PYG{+w}{  }\PYG{n}{field}\PYG{p}{(}\PYG{n}{df}\PYG{p}{)}\PYG{+w}{ }\PYG{o}{=}\PYG{+w}{ }\PYG{n+nb}{RAND}\PYG{p}{(}\PYG{p}{)}
\PYG{k}{end }\PYG{k}{do}
\end{sphinxVerbatim}

\sphinxAtStartPar
where \sphinxcode{\sphinxupquote{RAND()}} is some function that returns a new pseudo\sphinxhyphen{}random number
each time it is called.

\begin{sphinxadmonition}{warning}{Warning:}
\sphinxAtStartPar
This Built\sphinxhyphen{}in is implemented using the Fortran \sphinxcode{\sphinxupquote{random\_number}}
intrinsic. Therefore no guarantee is made as to the quality of
the sequence of pseudo\sphinxhyphen{}random numbers, especially when running
in parallel.
\end{sphinxadmonition}


\subsubsection{Raising to power}
\label{\detokenize{dynamo0p3:raising-to-power}}
\sphinxAtStartPar
Built\sphinxhyphen{}ins which raise \sphinxcode{\sphinxupquote{real}}\sphinxhyphen{}valued field elements to an exponent are
denoted with the keyword \sphinxstylestrong{powreal} for a \sphinxcode{\sphinxupquote{real}} exponent or \sphinxstylestrong{powint}
for an \sphinxcode{\sphinxupquote{integer}} exponent.


\paragraph{inc\_X\_powreal\_a}
\label{\detokenize{dynamo0p3:inc-x-powreal-a}}
\sphinxAtStartPar
\sphinxstylestrong{inc\_X\_powreal\_a} (\sphinxstylestrong{field}, \sphinxstyleemphasis{rscalar})

\sphinxAtStartPar
Raises a field to a \sphinxcode{\sphinxupquote{real}} scalar value and returns the
field (\sphinxcode{\sphinxupquote{X = X**a}}):

\begin{sphinxVerbatim}[commandchars=\\\{\}]
\PYG{n}{field}\PYG{p}{(}\PYG{p}{:}\PYG{p}{)}\PYG{+w}{ }\PYG{o}{=}\PYG{+w}{ }\PYG{n}{field}\PYG{p}{(}\PYG{p}{:}\PYG{p}{)}\PYG{o}{**}\PYG{n}{rscalar}
\end{sphinxVerbatim}


\paragraph{inc\_X\_powint\_n}
\label{\detokenize{dynamo0p3:inc-x-powint-n}}
\sphinxAtStartPar
\sphinxstylestrong{inc\_X\_powint\_n} (\sphinxstylestrong{field}, \sphinxstyleemphasis{iscalar})

\sphinxAtStartPar
Raises a field to an \sphinxcode{\sphinxupquote{integer}} scalar value and returns
the field (\sphinxcode{\sphinxupquote{X = X**n}}):

\begin{sphinxVerbatim}[commandchars=\\\{\}]
\PYG{n}{field}\PYG{p}{(}\PYG{p}{:}\PYG{p}{)}\PYG{+w}{ }\PYG{o}{=}\PYG{+w}{ }\PYG{n}{field}\PYG{p}{(}\PYG{p}{:}\PYG{p}{)}\PYG{o}{**}\PYG{n}{iscalar}
\end{sphinxVerbatim}

\sphinxAtStartPar
where \sphinxcode{\sphinxupquote{iscalar}} is an \sphinxcode{\sphinxupquote{integer}} scalar of \sphinxcode{\sphinxupquote{i\_def}} precision.


\subsubsection{Inner product}
\label{\detokenize{dynamo0p3:inner-product}}
\sphinxAtStartPar
Built\sphinxhyphen{}ins which calculate the inner product of two \sphinxcode{\sphinxupquote{real}}\sphinxhyphen{}valued fields
or of a \sphinxcode{\sphinxupquote{real}}\sphinxhyphen{}valued field with itself and return the result as a
\sphinxcode{\sphinxupquote{real}} scalar are denoted with the keyword \sphinxstylestrong{innerproduct}.

\begin{sphinxadmonition}{note}{Note:}
\sphinxAtStartPar
When used with distributed memory these Built\sphinxhyphen{}ins will
trigger the addition of a global sum which may affect the
performance and/or scalability of the code.
Also, whilst the fields in these Built\sphinxhyphen{}ins can be of any
supported \sphinxcode{\sphinxupquote{real}} {\hyperref[\detokenize{dynamo0p3:lfric-mixed-precision}]{\sphinxcrossref{\DUrole{std,std-ref}{precision}}}},
the only currently supported precision for the global
reductions in the LFRic infrastructure is \sphinxcode{\sphinxupquote{r\_def}}, hence
the result will be converted accordingly.
\end{sphinxadmonition}


\paragraph{X\_innerproduct\_Y}
\label{\detokenize{dynamo0p3:x-innerproduct-y}}
\sphinxAtStartPar
\sphinxstylestrong{X\_innerproduct\_Y} (\sphinxstylestrong{innprod}, \sphinxstyleemphasis{field1}, \sphinxstyleemphasis{field2})

\sphinxAtStartPar
Computes the inner product of two fields, \sphinxstyleemphasis{field1}
and \sphinxstyleemphasis{field2}, \sphinxstyleemphasis{i.e.}:

\begin{sphinxVerbatim}[commandchars=\\\{\}]
\PYG{n}{innprod}\PYG{+w}{ }\PYG{o}{=}\PYG{+w}{ }\PYG{n+nb}{SUM}\PYG{p}{(}\PYG{n}{field1}\PYG{p}{(}\PYG{p}{:}\PYG{p}{)}\PYG{o}{*}\PYG{n}{field2}\PYG{p}{(}\PYG{p}{:}\PYG{p}{)}\PYG{p}{)}
\end{sphinxVerbatim}

\sphinxAtStartPar
where \sphinxstylestrong{innprod} is a \sphinxcode{\sphinxupquote{real}} scalar of \sphinxcode{\sphinxupquote{r\_def}} precision.


\paragraph{X\_innerproduct\_X}
\label{\detokenize{dynamo0p3:x-innerproduct-x}}
\sphinxAtStartPar
\sphinxstylestrong{X\_innerproduct\_X} (\sphinxstylestrong{innprod}, \sphinxstyleemphasis{field})

\sphinxAtStartPar
Computes the inner product of the field \sphinxstyleemphasis{field1} by itself, \sphinxstyleemphasis{i.e.}:

\begin{sphinxVerbatim}[commandchars=\\\{\}]
\PYG{n}{innprod}\PYG{+w}{ }\PYG{o}{=}\PYG{+w}{ }\PYG{n+nb}{SUM}\PYG{p}{(}\PYG{n}{field}\PYG{p}{(}\PYG{p}{:}\PYG{p}{)}\PYG{o}{*}\PYG{n}{field}\PYG{p}{(}\PYG{p}{:}\PYG{p}{)}\PYG{p}{)}
\end{sphinxVerbatim}

\sphinxAtStartPar
where \sphinxstylestrong{innprod} is a \sphinxcode{\sphinxupquote{real}} scalar of \sphinxcode{\sphinxupquote{r\_def}} precision.


\subsubsection{Sum of elements}
\label{\detokenize{dynamo0p3:sum-of-elements}}
\sphinxAtStartPar
A Built\sphinxhyphen{}in which sums the elements of a \sphinxcode{\sphinxupquote{real}}\sphinxhyphen{}valued field and returns
the result as a \sphinxcode{\sphinxupquote{real}} scalar is denoted with the keyword \sphinxstylestrong{sum}.

\begin{sphinxadmonition}{note}{Note:}
\sphinxAtStartPar
When used with distributed memory this Built\sphinxhyphen{}in will trigger
the addition of a global sum which may affect the
performance and/or scalability of the code.
Also, whilst the fields in these Built\sphinxhyphen{}ins can be of any
supported \sphinxcode{\sphinxupquote{real}} {\hyperref[\detokenize{dynamo0p3:lfric-mixed-precision}]{\sphinxcrossref{\DUrole{std,std-ref}{precision}}}},
the only currently supported precision for the global
reductions in the LFRic infrastructure is \sphinxcode{\sphinxupquote{r\_def}}, hence
the result will be converted accordingly.
\end{sphinxadmonition}


\paragraph{sum\_X}
\label{\detokenize{dynamo0p3:sum-x}}
\sphinxAtStartPar
\sphinxstylestrong{sum\_X} (\sphinxstylestrong{sumfld}, \sphinxstyleemphasis{field})

\sphinxAtStartPar
Sums all of the elements of the field \sphinxstyleemphasis{field} and returns the result
in the \sphinxcode{\sphinxupquote{real}} scalar variable \sphinxstyleemphasis{sumfld}:

\begin{sphinxVerbatim}[commandchars=\\\{\}]
\PYG{n}{sumfld}\PYG{+w}{ }\PYG{o}{=}\PYG{+w}{ }\PYG{n+nb}{SUM}\PYG{p}{(}\PYG{n}{field}\PYG{p}{(}\PYG{p}{:}\PYG{p}{)}\PYG{p}{)}
\end{sphinxVerbatim}

\sphinxAtStartPar
where \sphinxstylestrong{sumfld} is a \sphinxcode{\sphinxupquote{real}} scalar of \sphinxcode{\sphinxupquote{r\_def}} precision.


\subsubsection{Sign of elements}
\label{\detokenize{dynamo0p3:sign-of-elements}}
\sphinxAtStartPar
A Built\sphinxhyphen{}in which returns the sign of a \sphinxcode{\sphinxupquote{real}}\sphinxhyphen{}valued field is denoted
with the keyword \sphinxstylestrong{sign}.


\paragraph{sign\_X}
\label{\detokenize{dynamo0p3:sign-x}}
\sphinxAtStartPar
\sphinxstylestrong{sign\_X} (\sphinxstylestrong{field2}, \sphinxstyleemphasis{rscalar}, \sphinxstyleemphasis{field1})

\sphinxAtStartPar
Returns the sign of a \sphinxcode{\sphinxupquote{real}}\sphinxhyphen{}valued field, e.g. in Fortran:
\sphinxcode{\sphinxupquote{Y = sign(a, X)}}. Here \sphinxcode{\sphinxupquote{a}} is a \sphinxcode{\sphinxupquote{real}} scalar and \sphinxcode{\sphinxupquote{Y}} and \sphinxcode{\sphinxupquote{X}}
are \sphinxcode{\sphinxupquote{real}}\sphinxhyphen{}valued fields. The results are \sphinxcode{\sphinxupquote{a}} for \sphinxcode{\sphinxupquote{X \textgreater{}= 0}} and
\sphinxcode{\sphinxupquote{\sphinxhyphen{}a}} for \sphinxcode{\sphinxupquote{X \textless{} 0}}:

\begin{sphinxVerbatim}[commandchars=\\\{\}]
\PYG{n}{field2}\PYG{p}{(}\PYG{p}{:}\PYG{p}{)}\PYG{+w}{ }\PYG{o}{=}\PYG{+w}{ }\PYG{n+nb}{SIGN}\PYG{p}{(}\PYG{n}{rscalar}\PYG{p}{,}\PYG{+w}{ }\PYG{n}{field1}\PYG{p}{(}\PYG{p}{:}\PYG{p}{)}\PYG{p}{)}
\end{sphinxVerbatim}


\subsubsection{DoF\sphinxhyphen{}wise maximum of elements}
\label{\detokenize{dynamo0p3:dof-wise-maximum-of-elements}}
\sphinxAtStartPar
Built\sphinxhyphen{}ins which return the DoF\sphinxhyphen{}wise maximum of a \sphinxcode{\sphinxupquote{real}} scalar and
a \sphinxcode{\sphinxupquote{real}}\sphinxhyphen{}valued field are denoted with the keyword \sphinxstylestrong{max}.


\paragraph{max\_aX}
\label{\detokenize{dynamo0p3:max-ax}}
\sphinxAtStartPar
\sphinxstylestrong{max\_aX} (\sphinxstylestrong{field2}, \sphinxstyleemphasis{rscalar}, \sphinxstyleemphasis{field1})

\sphinxAtStartPar
Returns maximum of \sphinxstyleemphasis{rscalar} and each element of the field \sphinxstyleemphasis{field1} as
the second field \sphinxstylestrong{field2} (\sphinxcode{\sphinxupquote{Y = max(a, X)}}):

\begin{sphinxVerbatim}[commandchars=\\\{\}]
\PYG{n}{field2}\PYG{p}{(}\PYG{p}{:}\PYG{p}{)}\PYG{+w}{ }\PYG{o}{=}\PYG{+w}{ }\PYG{n+nb}{MAX}\PYG{p}{(}\PYG{n}{rscalar}\PYG{p}{,}\PYG{+w}{ }\PYG{n}{field1}\PYG{p}{(}\PYG{p}{:}\PYG{p}{)}\PYG{p}{)}
\end{sphinxVerbatim}


\paragraph{inc\_max\_aX}
\label{\detokenize{dynamo0p3:inc-max-ax}}
\sphinxAtStartPar
\sphinxstylestrong{inc\_max\_aX} (\sphinxstyleemphasis{rscalar}, \sphinxstylestrong{field})

\sphinxAtStartPar
Returns maximum of \sphinxstyleemphasis{rscalar} and each element of the field \sphinxstyleemphasis{field} in
the same field (\sphinxcode{\sphinxupquote{X = max(a, X)}}):

\begin{sphinxVerbatim}[commandchars=\\\{\}]
\PYG{n}{field}\PYG{p}{(}\PYG{p}{:}\PYG{p}{)}\PYG{+w}{ }\PYG{o}{=}\PYG{+w}{ }\PYG{n+nb}{MAX}\PYG{p}{(}\PYG{n}{rscalar}\PYG{p}{,}\PYG{+w}{ }\PYG{n}{field}\PYG{p}{(}\PYG{p}{:}\PYG{p}{)}\PYG{p}{)}
\end{sphinxVerbatim}


\subsubsection{DoF\sphinxhyphen{}wise minimum of elements}
\label{\detokenize{dynamo0p3:dof-wise-minimum-of-elements}}
\sphinxAtStartPar
Built\sphinxhyphen{}ins which return the DoF\sphinxhyphen{}wise minimum of a \sphinxcode{\sphinxupquote{real}} scalar and
a \sphinxcode{\sphinxupquote{real}}\sphinxhyphen{}valued field are denoted with the keyword \sphinxstylestrong{min}.


\paragraph{min\_aX}
\label{\detokenize{dynamo0p3:min-ax}}
\sphinxAtStartPar
\sphinxstylestrong{min\_aX} (\sphinxstylestrong{field2}, \sphinxstyleemphasis{rscalar}, \sphinxstyleemphasis{field1})

\sphinxAtStartPar
Returns minimum of \sphinxstyleemphasis{rscalar} and each element of the field \sphinxstyleemphasis{field1} as
the second field \sphinxstylestrong{field2} (\sphinxcode{\sphinxupquote{Y = min(a, X)}}):

\begin{sphinxVerbatim}[commandchars=\\\{\}]
\PYG{n}{field2}\PYG{p}{(}\PYG{p}{:}\PYG{p}{)}\PYG{+w}{ }\PYG{o}{=}\PYG{+w}{ }\PYG{n+nb}{MIN}\PYG{p}{(}\PYG{n}{rscalar}\PYG{p}{,}\PYG{+w}{ }\PYG{n}{field1}\PYG{p}{(}\PYG{p}{:}\PYG{p}{)}\PYG{p}{)}
\end{sphinxVerbatim}


\paragraph{inc\_min\_aX}
\label{\detokenize{dynamo0p3:inc-min-ax}}
\sphinxAtStartPar
\sphinxstylestrong{inc\_min\_aX} (\sphinxstyleemphasis{rscalar}, \sphinxstylestrong{field})

\sphinxAtStartPar
Returns minimum of \sphinxstyleemphasis{rscalar} and each element of the field \sphinxstyleemphasis{field} in
the same field (\sphinxcode{\sphinxupquote{X = min(a, X)}}):

\begin{sphinxVerbatim}[commandchars=\\\{\}]
\PYG{n}{field}\PYG{p}{(}\PYG{p}{:}\PYG{p}{)}\PYG{+w}{ }\PYG{o}{=}\PYG{+w}{ }\PYG{n+nb}{MIN}\PYG{p}{(}\PYG{n}{rscalar}\PYG{p}{,}\PYG{+w}{ }\PYG{n}{field}\PYG{p}{(}\PYG{p}{:}\PYG{p}{)}\PYG{p}{)}
\end{sphinxVerbatim}


\subsubsection{Conversion of \sphinxstyleliteralintitle{\sphinxupquote{real}} field elements}
\label{\detokenize{dynamo0p3:conversion-of-real-field-elements}}
\sphinxAtStartPar
Built\sphinxhyphen{}ins which take a \sphinxcode{\sphinxupquote{real}} field for conversion to a field of
a different datatype or precision are denoted by the datatype that the
input \sphinxcode{\sphinxupquote{real}} field will be converted to. A Built\sphinxhyphen{}in that converts a
\sphinxcode{\sphinxupquote{real}} to an \sphinxcode{\sphinxupquote{integer}} field is denoted by the phrase \sphinxstylestrong{to\_int}.
Likewise, a Built\sphinxhyphen{}in that converts a \sphinxcode{\sphinxupquote{real}} to a \sphinxcode{\sphinxupquote{real}} field is
denoted by the phrase \sphinxstylestrong{to\_real}.


\paragraph{real\_to\_int\_X}
\label{\detokenize{dynamo0p3:real-to-int-x}}\label{\detokenize{dynamo0p3:real-to-int-built-in}}
\sphinxAtStartPar
\sphinxstylestrong{real\_to\_int\_X} (\sphinxstylestrong{ifield2}, \sphinxstyleemphasis{field1})

\sphinxAtStartPar
Converts \sphinxcode{\sphinxupquote{real}}\sphinxhyphen{}valued field elements to \sphinxcode{\sphinxupquote{integer}}\sphinxhyphen{}valued field
elements, e.g. in Fortran this would be: \sphinxcode{\sphinxupquote{Y = INT(X, kind=i\_\textless{}prec\textgreater{})}}.
Here \sphinxcode{\sphinxupquote{Y}} is an \sphinxcode{\sphinxupquote{integer}}\sphinxhyphen{}valued field and \sphinxcode{\sphinxupquote{X}} is the
\sphinxcode{\sphinxupquote{real}}\sphinxhyphen{}valued field being converted:

\begin{sphinxVerbatim}[commandchars=\\\{\}]
\PYG{n}{ifield2}\PYG{p}{(}\PYG{p}{:}\PYG{p}{)}\PYG{+w}{ }\PYG{o}{=}\PYG{+w}{ }\PYG{n+nb}{INT}\PYG{p}{(}\PYG{n}{field1}\PYG{p}{(}\PYG{p}{:}\PYG{p}{)}\PYG{p}{,}\PYG{+w}{ }\PYG{n+nb}{kind}\PYG{o}{=}\PYG{n}{i\PYGZus{}}\PYG{o}{\PYGZlt{}}\PYG{n}{prec}\PYG{o}{\PYGZgt{}}\PYG{p}{)}
\end{sphinxVerbatim}

\sphinxAtStartPar
where \sphinxstylestrong{ifield2} is currently the only supported \sphinxcode{\sphinxupquote{integer}}\sphinxhyphen{}valued field
type in LFRic (\sphinxcode{\sphinxupquote{integer\_field\_type}} of \sphinxcode{\sphinxupquote{i\_def}} precision) and a \sphinxcode{\sphinxupquote{real}}
\sphinxhyphen{}valued field \sphinxstyleemphasis{field1} can be of any {\hyperref[\detokenize{dynamo0p3:lfric-mixed-precision}]{\sphinxcrossref{\DUrole{std,std-ref}{supported precisions}}}} for \sphinxcode{\sphinxupquote{GH\_REAL}} fields (e.g. \sphinxcode{\sphinxupquote{r\_tran}} for
\sphinxcode{\sphinxupquote{r\_tran\_field\_type}}).


\paragraph{real\_to\_real\_X}
\label{\detokenize{dynamo0p3:real-to-real-x}}\label{\detokenize{dynamo0p3:real-to-real-built-in}}
\sphinxAtStartPar
\sphinxstylestrong{real\_to\_real\_X} (\sphinxstylestrong{field2}, \sphinxstyleemphasis{field1})

\sphinxAtStartPar
Converts \sphinxcode{\sphinxupquote{real}}\sphinxhyphen{}valued field elements from a precision \sphinxcode{\sphinxupquote{r\_\textless{}prec\textgreater{}}}
to \sphinxcode{\sphinxupquote{real}}\sphinxhyphen{}valued field elements of a differing precision \sphinxcode{\sphinxupquote{r\_\textless{}prec\textgreater{}}},
e.g. in Fortran this would be: \sphinxcode{\sphinxupquote{Y = REAL(X, kind=r\_\textless{}prec\textgreater{})}}. Here \sphinxcode{\sphinxupquote{Y}}
and \sphinxcode{\sphinxupquote{X}} are both \sphinxcode{\sphinxupquote{real}}\sphinxhyphen{}valued fields, with \sphinxcode{\sphinxupquote{X}} being converted
to the precision of \sphinxcode{\sphinxupquote{Y}}:

\begin{sphinxVerbatim}[commandchars=\\\{\}]
\PYG{n}{field2}\PYG{p}{(}\PYG{p}{:}\PYG{p}{)}\PYG{+w}{ }\PYG{o}{=}\PYG{+w}{ }\PYG{k+kt}{REAL}\PYG{p}{(}\PYG{n}{field1}\PYG{p}{(}\PYG{p}{:}\PYG{p}{)}\PYG{p}{,}\PYG{+w}{ }\PYG{n+nb}{kind}\PYG{o}{=}\PYG{n}{r\PYGZus{}}\PYG{o}{\PYGZlt{}}\PYG{n}{prec}\PYG{o}{\PYGZgt{}}\PYG{p}{)}
\end{sphinxVerbatim}

\sphinxAtStartPar
\sphinxstylestrong{field2} and \sphinxstyleemphasis{field1} are \sphinxcode{\sphinxupquote{real}}\sphinxhyphen{}valued fields of any {\hyperref[\detokenize{dynamo0p3:lfric-mixed-precision}]{\sphinxcrossref{\DUrole{std,std-ref}{supported
precisions}}}} for \sphinxcode{\sphinxupquote{GH\_REAL}} fields (e.g. \sphinxcode{\sphinxupquote{r\_tran}}
for \sphinxcode{\sphinxupquote{r\_tran\_field\_type}}).


\subsection{Built\sphinxhyphen{}in operations on \sphinxstyleliteralintitle{\sphinxupquote{integer}}\sphinxhyphen{}valued fields}
\label{\detokenize{dynamo0p3:built-in-operations-on-integer-valued-fields}}\label{\detokenize{dynamo0p3:lfric-built-ins-int}}
\sphinxAtStartPar
The number of supported Built\sphinxhyphen{}in operations on the \sphinxcode{\sphinxupquote{integer}}\sphinxhyphen{}valued
fields is not as large as for their \sphinxcode{\sphinxupquote{real}} counterparts as not all
mathematical operations on \sphinxcode{\sphinxupquote{integer}}\sphinxhyphen{}valued fields make sense.

\sphinxAtStartPar
As described {\hyperref[\detokenize{dynamo0p3:lfric-built-ins-dtype-access}]{\sphinxcrossref{\DUrole{std,std-ref}{above}}}}, Built\sphinxhyphen{}ins that
operate on \sphinxcode{\sphinxupquote{integer}}\sphinxhyphen{}valued fields mandate \sphinxcode{\sphinxupquote{GH\_INTEGER}} as the kernel
metadata for fields and scalars. Both \sphinxcode{\sphinxupquote{integer}} scalar arguments and
\sphinxcode{\sphinxupquote{integer}}\sphinxhyphen{}valued fields can only currently have \sphinxcode{\sphinxupquote{i\_def}} precision,
as described in the {\hyperref[\detokenize{dynamo0p3:lfric-mixed-precision}]{\sphinxcrossref{\DUrole{std,std-ref}{Mixed Precision}}}} section.

\sphinxAtStartPar
For instance, field and scalar declarations for the \sphinxcode{\sphinxupquote{X\_minus\_a}}
Built\sphinxhyphen{}in will be:

\begin{sphinxVerbatim}[commandchars=\\\{\}]
\PYG{k+kt}{integer}\PYG{p}{(}\PYG{n+nb}{kind}\PYG{o}{=}\PYG{n}{i\PYGZus{}def}\PYG{p}{)}\PYG{p}{,}\PYG{+w}{ }\PYG{k}{intent}\PYG{p}{(}\PYG{n}{in}\PYG{p}{)}\PYG{+w}{ }\PYG{k+kd}{::}\PYG{+w}{ }\PYG{n}{ascalar}
\PYG{k}{type}\PYG{p}{(}\PYG{n}{integer\PYGZus{}field\PYGZus{}type}\PYG{p}{)}\PYG{p}{,}\PYG{+w}{ }\PYG{k}{intent}\PYG{p}{(}\PYG{n}{in}\PYG{p}{)}\PYG{+w}{ }\PYG{k+kd}{::}\PYG{+w}{ }\PYG{n}{yfield}\PYG{p}{,}\PYG{+w}{ }\PYG{n}{xfield}
\end{sphinxVerbatim}


\subsubsection{Addition}
\label{\detokenize{dynamo0p3:id8}}
\sphinxAtStartPar
Built\sphinxhyphen{}ins that add \sphinxcode{\sphinxupquote{integer}}\sphinxhyphen{}valued fields and return the result as
an \sphinxcode{\sphinxupquote{integer}}\sphinxhyphen{}valued field are denoted with the keyword \sphinxstylestrong{plus} and
the prefix \sphinxstylestrong{int}.


\paragraph{int\_X\_plus\_Y}
\label{\detokenize{dynamo0p3:int-x-plus-y}}
\sphinxAtStartPar
\sphinxstylestrong{int\_X\_plus\_Y} (\sphinxstylestrong{ifield3}, \sphinxstyleemphasis{ifield1}, \sphinxstyleemphasis{ifield2})

\sphinxAtStartPar
Sums two fields and stores the result in the third field (\sphinxcode{\sphinxupquote{Z = X + Y}}):

\begin{sphinxVerbatim}[commandchars=\\\{\}]
\PYG{n}{ifield3}\PYG{p}{(}\PYG{p}{:}\PYG{p}{)}\PYG{+w}{ }\PYG{o}{=}\PYG{+w}{ }\PYG{n}{ifield1}\PYG{p}{(}\PYG{p}{:}\PYG{p}{)}\PYG{+w}{ }\PYG{o}{+}\PYG{+w}{ }\PYG{n}{ifield2}\PYG{p}{(}\PYG{p}{:}\PYG{p}{)}
\end{sphinxVerbatim}


\paragraph{int\_inc\_X\_plus\_Y}
\label{\detokenize{dynamo0p3:int-inc-x-plus-y}}
\sphinxAtStartPar
\sphinxstylestrong{int\_inc\_X\_plus\_Y} (\sphinxstylestrong{ifield1}, \sphinxstyleemphasis{ifield2})

\sphinxAtStartPar
Adds the second field to the first and returns it (\sphinxcode{\sphinxupquote{X = X + Y}}):

\begin{sphinxVerbatim}[commandchars=\\\{\}]
\PYG{n}{ifield1}\PYG{p}{(}\PYG{p}{:}\PYG{p}{)}\PYG{+w}{ }\PYG{o}{=}\PYG{+w}{ }\PYG{n}{ifield1}\PYG{p}{(}\PYG{p}{:}\PYG{p}{)}\PYG{+w}{ }\PYG{o}{+}\PYG{+w}{ }\PYG{n}{ifield2}\PYG{p}{(}\PYG{p}{:}\PYG{p}{)}
\end{sphinxVerbatim}


\paragraph{int\_a\_plus\_X}
\label{\detokenize{dynamo0p3:int-a-plus-x}}
\sphinxAtStartPar
\sphinxstylestrong{int\_a\_plus\_X} (\sphinxstylestrong{ifield2}, \sphinxstyleemphasis{iscalar}, \sphinxstyleemphasis{ifield1})

\sphinxAtStartPar
Adds an \sphinxcode{\sphinxupquote{integer}} scalar value to all elements of a field and stores
the result in another field (\sphinxcode{\sphinxupquote{Y = a + X}}):

\begin{sphinxVerbatim}[commandchars=\\\{\}]
\PYG{n}{ifield2}\PYG{p}{(}\PYG{p}{:}\PYG{p}{)}\PYG{+w}{ }\PYG{o}{=}\PYG{+w}{ }\PYG{n}{iscalar}\PYG{+w}{ }\PYG{o}{+}\PYG{+w}{ }\PYG{n}{ifield1}\PYG{p}{(}\PYG{p}{:}\PYG{p}{)}
\end{sphinxVerbatim}


\paragraph{int\_inc\_a\_plus\_X}
\label{\detokenize{dynamo0p3:int-inc-a-plus-x}}
\sphinxAtStartPar
\sphinxstylestrong{int\_inc\_a\_plus\_X} (\sphinxstyleemphasis{iscalar}, \sphinxstylestrong{ifield})

\sphinxAtStartPar
Adds an \sphinxcode{\sphinxupquote{integer}} scalar value to all elements of a field and returns
the field (\sphinxcode{\sphinxupquote{X = a + X}}):

\begin{sphinxVerbatim}[commandchars=\\\{\}]
\PYG{n}{ifield}\PYG{p}{(}\PYG{p}{:}\PYG{p}{)}\PYG{+w}{ }\PYG{o}{=}\PYG{+w}{ }\PYG{n}{iscalar}\PYG{+w}{ }\PYG{o}{+}\PYG{+w}{ }\PYG{n}{ifield}\PYG{p}{(}\PYG{p}{:}\PYG{p}{)}
\end{sphinxVerbatim}


\subsubsection{Subtraction}
\label{\detokenize{dynamo0p3:id9}}
\sphinxAtStartPar
Built\sphinxhyphen{}ins which subtract \sphinxcode{\sphinxupquote{integer}}\sphinxhyphen{}valued fields and return the result
as an \sphinxcode{\sphinxupquote{integer}}\sphinxhyphen{}valued field are denoted with the keyword \sphinxstylestrong{minus}
and the prefix \sphinxstylestrong{int}.


\paragraph{int\_X\_minus\_Y}
\label{\detokenize{dynamo0p3:int-x-minus-y}}
\sphinxAtStartPar
\sphinxstylestrong{int\_X\_minus\_Y} (\sphinxstylestrong{ifield3}, \sphinxstyleemphasis{ifield1}, \sphinxstyleemphasis{ifield2})

\sphinxAtStartPar
Subtracts the second field from the first and returns the result in the
third field (\sphinxcode{\sphinxupquote{Z = X \sphinxhyphen{} Y}}):

\begin{sphinxVerbatim}[commandchars=\\\{\}]
\PYG{n}{ifield3}\PYG{p}{(}\PYG{p}{:}\PYG{p}{)}\PYG{+w}{ }\PYG{o}{=}\PYG{+w}{ }\PYG{n}{ifield1}\PYG{p}{(}\PYG{p}{:}\PYG{p}{)}\PYG{+w}{ }\PYG{o}{\PYGZhy{}}\PYG{+w}{ }\PYG{n}{ifield2}\PYG{p}{(}\PYG{p}{:}\PYG{p}{)}
\end{sphinxVerbatim}


\paragraph{int\_inc\_X\_minus\_Y}
\label{\detokenize{dynamo0p3:int-inc-x-minus-y}}
\sphinxAtStartPar
\sphinxstylestrong{int\_inc\_X\_minus\_Y} (\sphinxstylestrong{ifield1}, \sphinxstyleemphasis{ifield2})

\sphinxAtStartPar
Subtracts the second field from the first and returns it (\sphinxcode{\sphinxupquote{X = X \sphinxhyphen{} Y}}):

\begin{sphinxVerbatim}[commandchars=\\\{\}]
\PYG{n}{ifield1}\PYG{p}{(}\PYG{p}{:}\PYG{p}{)}\PYG{+w}{ }\PYG{o}{=}\PYG{+w}{ }\PYG{n}{ifield1}\PYG{p}{(}\PYG{p}{:}\PYG{p}{)}\PYG{+w}{ }\PYG{o}{\PYGZhy{}}\PYG{+w}{ }\PYG{n}{ifield2}\PYG{p}{(}\PYG{p}{:}\PYG{p}{)}
\end{sphinxVerbatim}


\paragraph{int\_a\_minus\_X}
\label{\detokenize{dynamo0p3:int-a-minus-x}}
\sphinxAtStartPar
\sphinxstylestrong{int\_a\_minus\_X} (\sphinxstylestrong{ifield2}, \sphinxstyleemphasis{iscalar}, \sphinxstyleemphasis{ifield1})

\sphinxAtStartPar
Subtracts all elements of a field from an \sphinxcode{\sphinxupquote{integer}} scalar value and
stores the result in another field (\sphinxcode{\sphinxupquote{Y = a \sphinxhyphen{} X}}):

\begin{sphinxVerbatim}[commandchars=\\\{\}]
\PYG{n}{ifield2}\PYG{p}{(}\PYG{p}{:}\PYG{p}{)}\PYG{+w}{ }\PYG{o}{=}\PYG{+w}{ }\PYG{n}{iscalar}\PYG{+w}{ }\PYG{o}{\PYGZhy{}}\PYG{+w}{ }\PYG{n}{ifield1}\PYG{p}{(}\PYG{p}{:}\PYG{p}{)}
\end{sphinxVerbatim}


\paragraph{int\_inc\_a\_minus\_X}
\label{\detokenize{dynamo0p3:int-inc-a-minus-x}}
\sphinxAtStartPar
\sphinxstylestrong{int\_inc\_a\_minus\_X} (\sphinxstyleemphasis{iscalar}, \sphinxstylestrong{ifield})

\sphinxAtStartPar
Subtracts all elements of a field from an \sphinxcode{\sphinxupquote{integer}} scalar value and
returns the field (\sphinxcode{\sphinxupquote{X = a \sphinxhyphen{} X}}):

\begin{sphinxVerbatim}[commandchars=\\\{\}]
\PYG{n}{ifield}\PYG{p}{(}\PYG{p}{:}\PYG{p}{)}\PYG{+w}{ }\PYG{o}{=}\PYG{+w}{ }\PYG{n}{iscalar}\PYG{+w}{ }\PYG{o}{\PYGZhy{}}\PYG{+w}{ }\PYG{n}{ifield}\PYG{p}{(}\PYG{p}{:}\PYG{p}{)}
\end{sphinxVerbatim}


\paragraph{int\_X\_minus\_a}
\label{\detokenize{dynamo0p3:int-x-minus-a}}
\sphinxAtStartPar
\sphinxstylestrong{int\_X\_minus\_a} (\sphinxstylestrong{ifield2}, \sphinxstyleemphasis{ifield1}, \sphinxstyleemphasis{iscalar})

\sphinxAtStartPar
Subtracts an \sphinxcode{\sphinxupquote{integer}} scalar value from all elements of a field and
stores the result in another field (\sphinxcode{\sphinxupquote{Y = X \sphinxhyphen{} a}}):

\begin{sphinxVerbatim}[commandchars=\\\{\}]
\PYG{n}{ifield2}\PYG{p}{(}\PYG{p}{:}\PYG{p}{)}\PYG{+w}{ }\PYG{o}{=}\PYG{+w}{  }\PYG{n}{ifield1}\PYG{p}{(}\PYG{p}{:}\PYG{p}{)}\PYG{+w}{ }\PYG{o}{\PYGZhy{}}\PYG{+w}{ }\PYG{n}{iscalar}
\end{sphinxVerbatim}


\paragraph{int\_inc\_X\_minus\_a}
\label{\detokenize{dynamo0p3:int-inc-x-minus-a}}
\sphinxAtStartPar
\sphinxstylestrong{int\_inc\_X\_minus\_a} (\sphinxstylestrong{ifield}, \sphinxstyleemphasis{iscalar})

\sphinxAtStartPar
Subtracts an \sphinxcode{\sphinxupquote{integer}} scalar value from all elements of a field and
returns the field (\sphinxcode{\sphinxupquote{X = X \sphinxhyphen{} a}}):

\begin{sphinxVerbatim}[commandchars=\\\{\}]
\PYG{n}{ifield}\PYG{p}{(}\PYG{p}{:}\PYG{p}{)}\PYG{+w}{ }\PYG{o}{=}\PYG{+w}{  }\PYG{n}{ifield}\PYG{p}{(}\PYG{p}{:}\PYG{p}{)}\PYG{+w}{ }\PYG{o}{\PYGZhy{}}\PYG{+w}{ }\PYG{n}{iscalar}
\end{sphinxVerbatim}


\subsubsection{Multiplication}
\label{\detokenize{dynamo0p3:id10}}
\sphinxAtStartPar
Built\sphinxhyphen{}ins which multiply \sphinxcode{\sphinxupquote{integer}}\sphinxhyphen{}valued fields and return the result
as an \sphinxcode{\sphinxupquote{integer}}\sphinxhyphen{}valued field are denoted with the keyword \sphinxstylestrong{times}
and the prefix \sphinxstylestrong{int}.


\paragraph{int\_X\_times\_Y}
\label{\detokenize{dynamo0p3:int-x-times-y}}
\sphinxAtStartPar
\sphinxstylestrong{int\_X\_times\_Y} (\sphinxstylestrong{ifield3}, \sphinxstyleemphasis{ifield1}, \sphinxstyleemphasis{ifield2})

\sphinxAtStartPar
Multiplies two fields DoF by DoF and returns the result in a
third field (\sphinxcode{\sphinxupquote{Z = X*Y}}):

\begin{sphinxVerbatim}[commandchars=\\\{\}]
\PYG{n}{ifield3}\PYG{p}{(}\PYG{p}{:}\PYG{p}{)}\PYG{+w}{ }\PYG{o}{=}\PYG{+w}{ }\PYG{n}{ifield1}\PYG{p}{(}\PYG{p}{:}\PYG{p}{)}\PYG{o}{*}\PYG{n}{ifield2}\PYG{p}{(}\PYG{p}{:}\PYG{p}{)}
\end{sphinxVerbatim}


\paragraph{int\_inc\_X\_times\_Y}
\label{\detokenize{dynamo0p3:int-inc-x-times-y}}
\sphinxAtStartPar
\sphinxstylestrong{int\_inc\_X\_times\_Y} (\sphinxstylestrong{ifield1}, \sphinxstyleemphasis{ifield2})

\sphinxAtStartPar
Multiplies the first field by the second and returns it (\sphinxcode{\sphinxupquote{X = X*Y}}):

\begin{sphinxVerbatim}[commandchars=\\\{\}]
\PYG{n}{ifield1}\PYG{p}{(}\PYG{p}{:}\PYG{p}{)}\PYG{+w}{ }\PYG{o}{=}\PYG{+w}{ }\PYG{n}{ifield1}\PYG{p}{(}\PYG{p}{:}\PYG{p}{)}\PYG{o}{*}\PYG{n}{ifield2}\PYG{p}{(}\PYG{p}{:}\PYG{p}{)}
\end{sphinxVerbatim}


\subsubsection{Scaling}
\label{\detokenize{dynamo0p3:id11}}
\sphinxAtStartPar
Built\sphinxhyphen{}ins which scale \sphinxcode{\sphinxupquote{integer}}\sphinxhyphen{}valued fields are denoted with the keyword
\sphinxstylestrong{times} and prefixed by the keyword \sphinxstylestrong{int}.


\paragraph{int\_a\_times\_X}
\label{\detokenize{dynamo0p3:int-a-times-x}}
\sphinxAtStartPar
\sphinxstylestrong{int\_a\_times\_X} (\sphinxstylestrong{ifield2}, \sphinxstyleemphasis{iscalar}, \sphinxstyleemphasis{ifield1})

\sphinxAtStartPar
Multiplies a field by an \sphinxcode{\sphinxupquote{integer}} scalar and stores the result
in another field (\sphinxcode{\sphinxupquote{Y = a*X}}):

\begin{sphinxVerbatim}[commandchars=\\\{\}]
\PYG{n}{ifield2}\PYG{p}{(}\PYG{p}{:}\PYG{p}{)}\PYG{+w}{ }\PYG{o}{=}\PYG{+w}{ }\PYG{n}{iscalar}\PYG{o}{*}\PYG{n}{ifield1}\PYG{p}{(}\PYG{p}{:}\PYG{p}{)}
\end{sphinxVerbatim}


\paragraph{int\_inc\_a\_times\_X}
\label{\detokenize{dynamo0p3:int-inc-a-times-x}}
\sphinxAtStartPar
\sphinxstylestrong{int\_inc\_a\_times\_X} (\sphinxstyleemphasis{iscalar}, \sphinxstylestrong{ifield})

\sphinxAtStartPar
Multiplies a field by an \sphinxcode{\sphinxupquote{integer}} scalar value and returns the
field (\sphinxcode{\sphinxupquote{X = a*X}}):

\begin{sphinxVerbatim}[commandchars=\\\{\}]
\PYG{n}{ifield}\PYG{p}{(}\PYG{p}{:}\PYG{p}{)}\PYG{+w}{ }\PYG{o}{=}\PYG{+w}{ }\PYG{n}{iscalar}\PYG{o}{*}\PYG{n}{ifield}\PYG{p}{(}\PYG{p}{:}\PYG{p}{)}
\end{sphinxVerbatim}


\subsubsection{Setting to a value}
\label{\detokenize{dynamo0p3:id12}}
\sphinxAtStartPar
Built\sphinxhyphen{}ins which set \sphinxcode{\sphinxupquote{integer}}\sphinxhyphen{}valued field elements to some \sphinxcode{\sphinxupquote{integer}}
value are denoted with the keyword \sphinxstylestrong{setval} and the prefix \sphinxstylestrong{int}.


\paragraph{int\_setval\_c}
\label{\detokenize{dynamo0p3:int-setval-c}}
\sphinxAtStartPar
\sphinxstylestrong{int\_setval\_c} (\sphinxstylestrong{ifield}, \sphinxstyleemphasis{constant})

\sphinxAtStartPar
Sets all elements of a field \sphinxstyleemphasis{ifield} to an \sphinxcode{\sphinxupquote{integer}} scalar
\sphinxstyleemphasis{constant} (\sphinxcode{\sphinxupquote{X = c}}):

\begin{sphinxVerbatim}[commandchars=\\\{\}]
\PYG{n}{ifield}\PYG{p}{(}\PYG{p}{:}\PYG{p}{)}\PYG{+w}{ }\PYG{o}{=}\PYG{+w}{ }\PYG{n}{constant}
\end{sphinxVerbatim}


\paragraph{int\_setval\_X}
\label{\detokenize{dynamo0p3:int-setval-x}}
\sphinxAtStartPar
\sphinxstylestrong{int\_setval\_X} (\sphinxstylestrong{ifield2}, \sphinxstyleemphasis{ifield1})

\sphinxAtStartPar
Sets a field \sphinxstyleemphasis{ifield2} equal (DoF per DoF) to another field
\sphinxstyleemphasis{ifield1} (\sphinxcode{\sphinxupquote{Y = X}}):

\begin{sphinxVerbatim}[commandchars=\\\{\}]
\PYG{n}{ifield2}\PYG{p}{(}\PYG{p}{:}\PYG{p}{)}\PYG{+w}{ }\PYG{o}{=}\PYG{+w}{ }\PYG{n}{ifield1}\PYG{p}{(}\PYG{p}{:}\PYG{p}{)}
\end{sphinxVerbatim}


\subsubsection{Sign of elements}
\label{\detokenize{dynamo0p3:id13}}
\sphinxAtStartPar
A Built\sphinxhyphen{}in which returns the sign of an \sphinxcode{\sphinxupquote{integer}}\sphinxhyphen{}valued field
is denoted with the keyword \sphinxstylestrong{sign} and the prefix \sphinxstylestrong{int}.


\paragraph{int\_sign\_X}
\label{\detokenize{dynamo0p3:int-sign-x}}
\sphinxAtStartPar
\sphinxstylestrong{int\_sign\_X} (\sphinxstylestrong{ifield2}, \sphinxstyleemphasis{iscalar}, \sphinxstyleemphasis{ifield1})

\sphinxAtStartPar
Returns the sign of an \sphinxcode{\sphinxupquote{integer}}\sphinxhyphen{}valued field, e.g. in Fortran:
\sphinxcode{\sphinxupquote{Y = sign(a, X)}}. Here \sphinxcode{\sphinxupquote{a}} is an \sphinxcode{\sphinxupquote{integer}} scalar and \sphinxcode{\sphinxupquote{Y}}
and \sphinxcode{\sphinxupquote{X}} are \sphinxcode{\sphinxupquote{integer}}\sphinxhyphen{}valued fields.
The results are \sphinxcode{\sphinxupquote{a}} for \sphinxcode{\sphinxupquote{X \textgreater{}= 0}} and \sphinxcode{\sphinxupquote{\sphinxhyphen{}a}} for \sphinxcode{\sphinxupquote{a \textless{} 0}}:

\begin{sphinxVerbatim}[commandchars=\\\{\}]
\PYG{n}{ifield2}\PYG{p}{(}\PYG{p}{:}\PYG{p}{)}\PYG{+w}{ }\PYG{o}{=}\PYG{+w}{ }\PYG{n+nb}{SIGN}\PYG{p}{(}\PYG{n}{iscalar}\PYG{p}{,}\PYG{+w}{ }\PYG{n}{ifield1}\PYG{p}{(}\PYG{p}{:}\PYG{p}{)}\PYG{p}{)}
\end{sphinxVerbatim}


\subsubsection{DoF\sphinxhyphen{}wise maximum of elements}
\label{\detokenize{dynamo0p3:id14}}
\sphinxAtStartPar
Built\sphinxhyphen{}ins which return the DoF\sphinxhyphen{}wise maximum of an \sphinxcode{\sphinxupquote{integer}} scalar
and an \sphinxcode{\sphinxupquote{integer}}\sphinxhyphen{}valued field are denoted with the keyword \sphinxstylestrong{max}.


\paragraph{int\_max\_aX}
\label{\detokenize{dynamo0p3:int-max-ax}}
\sphinxAtStartPar
\sphinxstylestrong{int\_max\_aX} (\sphinxstylestrong{ifield2}, \sphinxstyleemphasis{iscalar}, \sphinxstyleemphasis{ifield1})

\sphinxAtStartPar
Returns maximum of \sphinxstyleemphasis{iscalar} and each element of the field \sphinxstyleemphasis{ifield1} as
the second field \sphinxstylestrong{ifield2} (\sphinxcode{\sphinxupquote{Y = max(a, X)}}):

\begin{sphinxVerbatim}[commandchars=\\\{\}]
\PYG{n}{ifield2}\PYG{p}{(}\PYG{p}{:}\PYG{p}{)}\PYG{+w}{ }\PYG{o}{=}\PYG{+w}{ }\PYG{n+nb}{MAX}\PYG{p}{(}\PYG{n}{iscalar}\PYG{p}{,}\PYG{+w}{ }\PYG{n}{ifield1}\PYG{p}{(}\PYG{p}{:}\PYG{p}{)}\PYG{p}{)}
\end{sphinxVerbatim}


\paragraph{int\_inc\_max\_aX}
\label{\detokenize{dynamo0p3:int-inc-max-ax}}
\sphinxAtStartPar
\sphinxstylestrong{int\_inc\_max\_aX} (\sphinxstyleemphasis{iscalar}, \sphinxstylestrong{ifield})

\sphinxAtStartPar
Returns maximum of \sphinxstyleemphasis{iscalar} and each element of the field \sphinxstyleemphasis{ifield} in
the same field (\sphinxcode{\sphinxupquote{X = max(a, X)}}):

\begin{sphinxVerbatim}[commandchars=\\\{\}]
\PYG{n}{ifield}\PYG{p}{(}\PYG{p}{:}\PYG{p}{)}\PYG{+w}{ }\PYG{o}{=}\PYG{+w}{ }\PYG{n+nb}{MAX}\PYG{p}{(}\PYG{n}{iscalar}\PYG{p}{,}\PYG{+w}{ }\PYG{n}{ifield}\PYG{p}{(}\PYG{p}{:}\PYG{p}{)}\PYG{p}{)}
\end{sphinxVerbatim}


\subsubsection{DoF\sphinxhyphen{}wise minimum of elements}
\label{\detokenize{dynamo0p3:id15}}
\sphinxAtStartPar
Built\sphinxhyphen{}ins which return the DoF\sphinxhyphen{}wise minimum of an \sphinxcode{\sphinxupquote{integer}} scalar
and an \sphinxcode{\sphinxupquote{integer}}\sphinxhyphen{}valued field are denoted with the keyword \sphinxstylestrong{min}.


\paragraph{int\_min\_aX}
\label{\detokenize{dynamo0p3:int-min-ax}}
\sphinxAtStartPar
\sphinxstylestrong{int\_min\_aX} (\sphinxstylestrong{ifield2}, \sphinxstyleemphasis{iscalar}, \sphinxstyleemphasis{ifield1})

\sphinxAtStartPar
Returns minimum of \sphinxstyleemphasis{iscalar} and each element of the field \sphinxstyleemphasis{ifield1} as
the second field \sphinxstylestrong{ifield2} (\sphinxcode{\sphinxupquote{Y = min(a, X)}}):

\begin{sphinxVerbatim}[commandchars=\\\{\}]
\PYG{n}{ifield2}\PYG{p}{(}\PYG{p}{:}\PYG{p}{)}\PYG{+w}{ }\PYG{o}{=}\PYG{+w}{ }\PYG{n+nb}{MIN}\PYG{p}{(}\PYG{n}{iscalar}\PYG{p}{,}\PYG{+w}{ }\PYG{n}{ifield1}\PYG{p}{(}\PYG{p}{:}\PYG{p}{)}\PYG{p}{)}
\end{sphinxVerbatim}


\paragraph{int\_inc\_min\_aX}
\label{\detokenize{dynamo0p3:int-inc-min-ax}}
\sphinxAtStartPar
\sphinxstylestrong{int\_inc\_min\_aX} (\sphinxstyleemphasis{iscalar}, \sphinxstylestrong{ifield})

\sphinxAtStartPar
Returns minimum of \sphinxstyleemphasis{iscalar} and each element of the field \sphinxstyleemphasis{ifield} in
the same field (\sphinxcode{\sphinxupquote{X = min(a, X)}}):

\begin{sphinxVerbatim}[commandchars=\\\{\}]
\PYG{n}{ifield}\PYG{p}{(}\PYG{p}{:}\PYG{p}{)}\PYG{+w}{ }\PYG{o}{=}\PYG{+w}{ }\PYG{n+nb}{MIN}\PYG{p}{(}\PYG{n}{iscalar}\PYG{p}{,}\PYG{+w}{ }\PYG{n}{ifield}\PYG{p}{(}\PYG{p}{:}\PYG{p}{)}\PYG{p}{)}
\end{sphinxVerbatim}


\subsubsection{Conversion of \sphinxstyleliteralintitle{\sphinxupquote{integer}} to \sphinxstyleliteralintitle{\sphinxupquote{real}} field elements}
\label{\detokenize{dynamo0p3:conversion-of-integer-to-real-field-elements}}
\sphinxAtStartPar
A Built\sphinxhyphen{}in which takes an \sphinxcode{\sphinxupquote{integer}} field and converts it to
a \sphinxcode{\sphinxupquote{real}} field is denoted by the phrase \sphinxstylestrong{to\_real}.


\paragraph{int\_to\_real\_X}
\label{\detokenize{dynamo0p3:int-to-real-x}}\label{\detokenize{dynamo0p3:int-to-real-built-in}}
\sphinxAtStartPar
\sphinxstylestrong{int\_to\_real\_X} (\sphinxstylestrong{field2}, \sphinxstyleemphasis{ifield1})

\sphinxAtStartPar
Converts \sphinxcode{\sphinxupquote{integer}}\sphinxhyphen{}valued field elements to \sphinxcode{\sphinxupquote{real}}\sphinxhyphen{}valued field
elements, e.g. in Fortran this would be \sphinxcode{\sphinxupquote{Y = REAL(X, kind=r\_\textless{}prec\textgreater{})}}.
Here \sphinxcode{\sphinxupquote{Y}} is a \sphinxcode{\sphinxupquote{real}}\sphinxhyphen{}valued field and \sphinxcode{\sphinxupquote{X}} is the
\sphinxcode{\sphinxupquote{integer}}\sphinxhyphen{}valued field being converted:

\begin{sphinxVerbatim}[commandchars=\\\{\}]
\PYG{n}{field2}\PYG{p}{(}\PYG{p}{:}\PYG{p}{)}\PYG{+w}{ }\PYG{o}{=}\PYG{+w}{ }\PYG{k+kt}{REAL}\PYG{p}{(}\PYG{n}{ifield1}\PYG{p}{(}\PYG{p}{:}\PYG{p}{)}\PYG{p}{,}\PYG{+w}{ }\PYG{n+nb}{kind}\PYG{o}{=}\PYG{n}{r\PYGZus{}}\PYG{o}{\PYGZlt{}}\PYG{n}{prec}\PYG{o}{\PYGZgt{}}\PYG{p}{)}
\end{sphinxVerbatim}

\sphinxAtStartPar
where \sphinxstylestrong{ifield1} is currently the only supported \sphinxcode{\sphinxupquote{integer}}\sphinxhyphen{}valued
field type in LFRic (\sphinxcode{\sphinxupquote{integer\_field\_type}} of \sphinxcode{\sphinxupquote{i\_def}} precision).
The \sphinxcode{\sphinxupquote{real}}\sphinxhyphen{}valued \sphinxstylestrong{field1} can be of any {\hyperref[\detokenize{dynamo0p3:lfric-mixed-precision}]{\sphinxcrossref{\DUrole{std,std-ref}{supported
precisions}}}} for \sphinxcode{\sphinxupquote{GH\_REAL}} fields, hence
\sphinxcode{\sphinxupquote{r\_\textless{}prec\textgreater{}}} is determined from the algorithm layer (e.g.
\sphinxcode{\sphinxupquote{r\_solver}} for \sphinxcode{\sphinxupquote{r\_solver\_field\_type}}).


\section{Boundary Conditions}
\label{\detokenize{dynamo0p3:boundary-conditions}}
\sphinxAtStartPar
In the LFRic API, boundary conditions for a field or LMA operator can
be enforced by the algorithm developer by calling the Kernels
\sphinxcode{\sphinxupquote{enforce\_bc\_type}} or \sphinxcode{\sphinxupquote{enforce\_operator\_bc\_type}},
respectively. These kernels take a field or operator as input and apply
boundary conditions. For example:

\begin{sphinxVerbatim}[commandchars=\\\{\}]
\PYG{k}{call }\PYG{n}{invoke}\PYG{p}{(}\PYG{+w}{ }\PYG{n}{kernel\PYGZus{}type}\PYG{p}{(}\PYG{n}{field1}\PYG{p}{,}\PYG{+w}{ }\PYG{n}{field2}\PYG{p}{)}\PYG{p}{,}\PYG{+w}{      }\PYG{p}{\PYGZam{}}
\PYG{+w}{             }\PYG{n}{enforce\PYGZus{}bc\PYGZus{}type}\PYG{p}{(}\PYG{n}{field1}\PYG{p}{)}\PYG{p}{,}\PYG{+w}{          }\PYG{p}{\PYGZam{}}
\PYG{+w}{             }\PYG{n}{kernel\PYGZus{}with\PYGZus{}op\PYGZus{}type}\PYG{p}{(}\PYG{n}{field1}\PYG{p}{,}\PYG{+w}{ }\PYG{n}{op1}\PYG{p}{)}\PYG{p}{,}\PYG{+w}{ }\PYG{p}{\PYGZam{}}
\PYG{+w}{             }\PYG{n}{enforce\PYGZus{}operator\PYGZus{}bc\PYGZus{}type}\PYG{p}{(}\PYG{n}{op1}\PYG{p}{)}\PYG{+w}{     }\PYG{p}{\PYGZam{}}
\PYG{+w}{           }\PYG{p}{)}
\end{sphinxVerbatim}

\sphinxAtStartPar
The particular boundary conditions that are applied are not known by
PSyclone, PSyclone simply recognises these kernels by their names and passes
pre\sphinxhyphen{}specified dofmap and boundary\_value arrays into the kernel
implementations, the contents of which are set by the LFRic
infrastructure.

\sphinxAtStartPar
Up to and including version 1.4.0 of PSyclone, boundary conditions
were applied automatically after a call to \sphinxcode{\sphinxupquote{matrix\_vector\_type}} if
the field arguments were on a vector function space (one of \sphinxcode{\sphinxupquote{W1}},
\sphinxcode{\sphinxupquote{W2}}, \sphinxcode{\sphinxupquote{W2H}}, \sphinxcode{\sphinxupquote{W2V}} or \sphinxcode{\sphinxupquote{W2broken}}). With the subsequent introduction
of the ability to apply boundary conditions to operators this functionality
is no longer required and has been removed.

\sphinxAtStartPar
Example \sphinxcode{\sphinxupquote{eg4}} in the \sphinxcode{\sphinxupquote{examples/lfric}} directory includes a call
to \sphinxcode{\sphinxupquote{enforce\_bc\_kernel\_type}} so can be used to see the boundary condition
code that is added by PSyclone. See the \sphinxcode{\sphinxupquote{README}} in the
\sphinxcode{\sphinxupquote{examples/lfric}} directory for instructions on how to run this
example.

\sphinxAtStartPar
An example of applying boundary conditions to an operator is the kernel
\sphinxcode{\sphinxupquote{enforce\_operator\_bc\_kernel\_mod.F90}} in the
\sphinxcode{\sphinxupquote{\textless{}PSYCLONEHOME\textgreater{}/src/psyclone/tests/test\_files/dynamo0p3}} directory.
Since operators are discontinuous quantities, updating their values can
be safely performed in parallel (see Section {\hyperref[\detokenize{dynamo0p3:lfric-kernel}]{\sphinxcrossref{\DUrole{std,std-ref}{Kernel}}}}).
The \sphinxcode{\sphinxupquote{GH\_READWRITE}} access is used for updating discontinuous operators
(see subsection {\hyperref[\detokenize{dynamo0p3:lfric-kernel-valid-access}]{\sphinxcrossref{\DUrole{std,std-ref}{Valid Access Modes}}}} for more details).


\section{Conventions}
\label{\detokenize{dynamo0p3:conventions}}
\sphinxAtStartPar
The naming of LFRic API kernels and associated entities (types,
subroutines and modules) follows the PSyclone Fortran naming
conventions (see {\hyperref[\detokenize{fortran_naming_conventions:fortran-naming}]{\sphinxcrossref{\DUrole{std,std-ref}{Fortran Naming Conventions}}}}). However, PSyclone does not need
this convention to be followed apart from the stub generator (see the
{\hyperref[\detokenize{psyclone_kern:stub-generation}]{\sphinxcrossref{\DUrole{std,std-ref}{Kernel\sphinxhyphen{}stub Generator}}}} Section ) where the name of the metadata to be
parsed is determined from the module name.

\sphinxAtStartPar
The contents of the metadata is also usually declared private but this
does not affect PSyclone.

\sphinxAtStartPar
Finally, the \sphinxcode{\sphinxupquote{procedure}} metadata (located within the kernel
metadata) usually has \sphinxcode{\sphinxupquote{nopass}} specified but again this is ignored
by PSyclone.


\section{Configuration}
\label{\detokenize{dynamo0p3:configuration}}\label{\detokenize{dynamo0p3:lfric-api-configuration}}
\sphinxAtStartPar
The general and the LFRic\sphinxhyphen{}API\sphinxhyphen{}specific configuration options are described
in the {\hyperref[\detokenize{configuration:configuration}]{\sphinxcrossref{\DUrole{std,std-ref}{Configuration}}}} section.


\subsection{Annexed DoFs}
\label{\detokenize{dynamo0p3:annexed-dofs}}\label{\detokenize{dynamo0p3:lfric-annexed-dofs}}
\sphinxAtStartPar
When a kernel operates on DoFs (rather than cell\sphinxhyphen{}columns) for a continuous
field using distributed memory, PSyclone need only ensure that DoFs owned by a
processor are computed. However, for continuous fields, shared DoFs at
the boundary between processors must be replicated (as different cells
share the same DoF). Only one processor can own a DoF, therefore
processors will have continuous fields which contain DoFs that the
processor does not own. These unowned DoFs are called \sphinxtitleref{annexed} in the
LFRic API and are a separate, but related, concept to field halos.

\sphinxAtStartPar
When a kernel that operates on a cell\sphinxhyphen{}column needs to read a
continuous field then the annexed DoFs must be up\sphinxhyphen{}to\sphinxhyphen{}date on all
processors. If they are not then a halo exchange must be
added. Currently PSyclone defaults, for kernels which iterate over
DoFs, to iterating over only owned DoFs. This behaviour can be changed
by setting \sphinxtitleref{COMPUTE\_ANNEXED\_DOFS} to \sphinxcode{\sphinxupquote{true}} in the \sphinxtitleref{lfric}
section of the configuration file (see the {\hyperref[\detokenize{configuration:configuration}]{\sphinxcrossref{\DUrole{std,std-ref}{Configuration}}}}
section). PSyclone will then generate code to iterate over both owned
and annexed DoFs, thereby reducing the number of halo exchanges
required (at the expense of redundantly computing annexed DoFs). For
more details please refer to the \sphinxhref{https://psyclone-dev.readthedocs.io/en/latest/APIs.html\#lfric-developers}{LFRic}
developers section.


\subsection{Run\sphinxhyphen{}time Checks}
\label{\detokenize{dynamo0p3:run-time-checks}}\label{\detokenize{dynamo0p3:lfric-run-time-checks}}
\sphinxAtStartPar
PSyclone performs static consistency checks where possible. When this
is not possible PSyclone can generate run\sphinxhyphen{}time checks. As there may be
performance costs associated with run\sphinxhyphen{}time checks they may be switched
on or off by the \sphinxtitleref{RUN\_TIME\_CHECKS} option in the configuration file.

\sphinxAtStartPar
Currently run\sphinxhyphen{}time checks can be generated to:
\begin{enumerate}
\sphinxsetlistlabels{\arabic}{enumi}{enumii}{}{)}%
\item {} 
\sphinxAtStartPar
Check that a field with a read\sphinxhyphen{}only function space (see section
{\hyperref[\detokenize{dynamo0p3:lfric-ro-function-space}]{\sphinxcrossref{\DUrole{std,std-ref}{Read\sphinxhyphen{}Only Function Spaces}}}}) is not modified by a kernel. This is
enforced by checking that all fields that are marked (in kernel
metadata) as being updated by a kernel are not on a read\sphinxhyphen{}only function
space. A second check that is required for fields on read\sphinxhyphen{}only
function spaces is to ensure that the halo is clean before it is accessed.
This check is currently implemented within the LFRic
infrastructure halo exchange call (that the PSyclone LFRic API places
at appropriate locations). If the halo is clean then the halo exchange
will not be called. However, if the halo is not clean then the
resulting halo exchange call will cause the infrastructure to raise an
error (because the field is on a read\sphinxhyphen{}only space).

\item {} 
\sphinxAtStartPar
Check that the function space of a field is consistent with the
kernel function space metadata that the field’s data is passed
into. For example, if kernel metadata specifies that a field is on
the \sphinxcode{\sphinxupquote{W2}} function space then a run\sphinxhyphen{}time check is added to ensure that
the field object passed into the PSy layer is indeed on that space.
For more general kernel function space metadata, such as
\sphinxtitleref{ANY\_DISCONTINUOUS\_SPACE\_*} then a run\sphinxhyphen{}time check is added to
ensure that the field is on one of the discontinuous function
spaces supported in the LFRic API.

\end{enumerate}


\subsection{Supported Data Types and Default Kind}
\label{\detokenize{dynamo0p3:supported-data-types-and-default-kind}}\label{\detokenize{dynamo0p3:lfric-datatype-kind}}
\sphinxAtStartPar
The LFRic API supports three Fortran primitive (intrinsic) data
types, \sphinxcode{\sphinxupquote{real}}, \sphinxcode{\sphinxupquote{integer}} and \sphinxcode{\sphinxupquote{logical}} (listed in the
\sphinxtitleref{supported\_fortran\_datatypes} section of the {\hyperref[\detokenize{configuration:configuration}]{\sphinxcrossref{\DUrole{std,std-ref}{PSyclone
configuration file}}}}). All three data types are used
for {\hyperref[\detokenize{dynamo0p3:lfric-scalar}]{\sphinxcrossref{\DUrole{std,std-ref}{scalars}}}}. {\hyperref[\detokenize{dynamo0p3:lfric-field}]{\sphinxcrossref{\DUrole{std,std-ref}{Fields}}}} and
{\hyperref[\detokenize{dynamo0p3:lfric-field-vector}]{\sphinxcrossref{\DUrole{std,std-ref}{field vectors}}}} are allowed to have \sphinxcode{\sphinxupquote{real}}
and \sphinxcode{\sphinxupquote{integer}} data. {\hyperref[\detokenize{dynamo0p3:lfric-operator}]{\sphinxcrossref{\DUrole{std,std-ref}{Operators}}}} and
{\hyperref[\detokenize{dynamo0p3:lfric-cma-operator}]{\sphinxcrossref{\DUrole{std,std-ref}{column\sphinxhyphen{}wise operators}}}} are only allowed to
have \sphinxcode{\sphinxupquote{real}} data. These supported primitive types are linked to the
respective {\hyperref[\detokenize{dynamo0p3:lfric-kernel-valid-data-type}]{\sphinxcrossref{\DUrole{std,std-ref}{kernel data type}}}}
metadata descriptors, \sphinxcode{\sphinxupquote{GH\_REAL}} and \sphinxcode{\sphinxupquote{GH\_INTEGER}}.

\sphinxAtStartPar
The default kind (precision) for these supported data types is
set to \sphinxcode{\sphinxupquote{r\_def}}, \sphinxcode{\sphinxupquote{i\_def}} and \sphinxcode{\sphinxupquote{l\_def}}, respectively, in the
\sphinxcode{\sphinxupquote{default\_kind}} dictionary in the configuration file. These default
values are defined in the LFRic infrastructure code.

\begin{sphinxadmonition}{note}{Note:}
\sphinxAtStartPar
Whilst the \sphinxcode{\sphinxupquote{logical}} Fortran primitive (intrinsic) data
type is supported in the LFRic API for scalar arguments, it is
not yet available for fields and operators. This will be added
as required in future releases.
\end{sphinxadmonition}


\subsection{Precision Map}
\label{\detokenize{dynamo0p3:precision-map}}\label{\detokenize{dynamo0p3:lfric-precision-map}}
\sphinxAtStartPar
This gives the amount of storage (in bytes) associated with a
particular LFRic precision. The values for ‘r\_tran’, ‘r\_solver’,
‘r\_def’, ‘r\_bl’ and ‘r\_phys’ are set within LFRic infrastructure
according to CPP ifdefs. The values given in the configuration file
are the defaults. ‘l\_def’ is included in the dictionary so that it
contains a complete record of the various precision symbols used in
LFRic.

\begin{sphinxadmonition}{note}{Note:}
\sphinxAtStartPar
Storing the precision map in the LFRic API within PSyclone is a
temporary measure which will yield to the LFRic infrastructure
as the single source of precisions, as discussed in PSyclone
issue \#1941.
\end{sphinxadmonition}


\subsection{Number of Generalised \sphinxstyleliteralintitle{\sphinxupquote{ANY\_*\_SPACE}} Function Spaces}
\label{\detokenize{dynamo0p3:number-of-generalised-any-space-function-spaces}}\label{\detokenize{dynamo0p3:lfric-num-any-spaces}}
\sphinxAtStartPar
As outlined in the {\hyperref[\detokenize{dynamo0p3:lfric-api-meta-args}]{\sphinxcrossref{\DUrole{std,std-ref}{meta\_args}}}} and the
{\hyperref[\detokenize{dynamo0p3:lfric-function-space}]{\sphinxcrossref{\DUrole{std,std-ref}{Supported Function Spaces}}}} sections
above, the number of generalised \sphinxcode{\sphinxupquote{ANY\_SPACE\_\textless{}n\textgreater{}}} and
\sphinxcode{\sphinxupquote{ANY\_DISCONTINUOUS\_SPACE\_\textless{}n\textgreater{}}} function spaces can be set in the
{\hyperref[\detokenize{configuration:configuration}]{\sphinxcrossref{\DUrole{std,std-ref}{PSyclone configuration file}}}}.

\sphinxAtStartPar
The relevant parameters are \sphinxcode{\sphinxupquote{NUM\_ANY\_SPACE}} and
\sphinxcode{\sphinxupquote{NUM\_ANY\_DISCONTINUOUS\_SPACE}}, respectively. Their default values in
the configuration file are 10 and their allowed values are positive
non\sphinxhyphen{}zero integers. PSyclone will raise a \sphinxcode{\sphinxupquote{ConfigurationError}} if a
supplied value is invalid.


\section{Transformations}
\label{\detokenize{dynamo0p3:transformations}}\label{\detokenize{dynamo0p3:lfric-api-transformations}}
\sphinxAtStartPar
This section describes the LFRic API\sphinxhyphen{}specific transformations. In
cases, excepting \sphinxstylestrong{Dynamo0p3RedundantComputationTrans},
\sphinxstylestrong{Dynamo0p3AsyncHaloExchangeTrans} and \sphinxstylestrong{Dynamo0p3KernelConstTrans},
these transformations are specialisations of generic transformations
described in the {\hyperref[\detokenize{transformations:transformations}]{\sphinxcrossref{\DUrole{std,std-ref}{Transformations}}}} section. The difference
between these transformations and the generic ones is that these
perform LFRic API\sphinxhyphen{}specific checks to make sure the transformations
are valid. In practice these transformations perform the required
checks then call the generic ones internally.

\sphinxAtStartPar
The use of the LFRic API\sphinxhyphen{}specific transformations is exactly the
same as the equivalent generic ones in all cases excepting
\sphinxstylestrong{LFRicLoopFuseTrans}. In this case an additional optional argument
\sphinxstylestrong{same\_space} can be set when applying the transformation.
The reason for this is to allow loop fusion when one or more of the
iteration spaces is determined by a function space that is unknown by
PSyclone at compile time. This is the case when the \sphinxcode{\sphinxupquote{ANY\_SPACE\_\textless{}n\textgreater{}}}
function space is specified in the Kernel metadata. Adding
\sphinxcode{\sphinxupquote{\{"same\_space": True\}}} as option when applying the transformation allows
the user to specify that the spaces are the same (see
{\hyperref[\detokenize{transformations:sec-transformations-available}]{\sphinxcrossref{\DUrole{std,std-ref}{Available transformations}}}} for using options in transformations).
This option should therefore be used with caution. PSyclone will
raise an error if \sphinxstylestrong{same\_space} is used when at least one of the function
spaces is not \sphinxcode{\sphinxupquote{ANY\_SPACE\_\textless{}n\textgreater{}}} or both spaces are not the same. In general,
PSyclone will not allow loop fusion if it does not know the spaces
are the same. The exception are loops over discontinuous spaces (see
{\hyperref[\detokenize{dynamo0p3:lfric-function-space}]{\sphinxcrossref{\DUrole{std,std-ref}{Supported Function Spaces}}}} for list of discontinuous function spaces)
for which loop fusion is allowed (unless the loop bounds become different
due to a prior transformation).

\sphinxAtStartPar
The \sphinxstylestrong{Dynamo0p3RedundantComputationTrans} and
\sphinxstylestrong{Dynamo0p3AsyncHaloExchange} transformations are only valid for the
LFRic API. This is because this API is currently the only one
that supports distributed memory.  An example of redundant computation
can be found in \sphinxcode{\sphinxupquote{examples/lfric/eg8}} and an example of asynchronous
halo exchanges can be found in \sphinxcode{\sphinxupquote{examples/lfric/eg11}}.

\sphinxAtStartPar
The \sphinxstylestrong{Dynamo0p3KernelConstTrans} transformation is only valid for the
LFRic API. This is because the properties that it makes constant
are API specific.

\sphinxAtStartPar
The LFRic API\sphinxhyphen{}specific transformations currently available
are given below. Early transformations include “Dynamo0p3” or “Dynamo”
in their name to indicate that these transformations are only valid
for this particular API. More recent transformations typically include
“LFRic” in their name to indicate the same restriction. However, more
importantly, transformations that are specific to LFRic reside in the
LFRic\sphinxhyphen{}specific “psyclone.domain/lfric/transformations”
directory. Note, the early LFRic API\sphinxhyphen{}specific
transformations have not yet been migrated to this directory.

\begin{sphinxadmonition}{note}{Note:}
\sphinxAtStartPar
Only the loop\sphinxhyphen{}colouring and OpenMP transformations are currently
supported for loops that contain inter\sphinxhyphen{}grid kernels. Attempting
to apply other transformation types will result in PSyclone raising
an error.
\end{sphinxadmonition}


\begin{fulllineitems}

\pysigstartsignatures
\pysigline{\sphinxbfcode{\sphinxupquote{class\DUrole{w,w}{  }}}\sphinxcode{\sphinxupquote{psyclone.domain.lfric.transformations.}}\sphinxbfcode{\sphinxupquote{LFRicExtractTrans}}}
\pysigstopsignatures
\sphinxAtStartPar
LFRic API application of ExtractTrans transformation
to extract code into a stand\sphinxhyphen{}alone program. For example:

\begin{sphinxVerbatim}[commandchars=\\\{\}]
\PYG{g+gp}{\PYGZgt{}\PYGZgt{}\PYGZgt{} }\PYG{k+kn}{from}\PYG{+w}{ }\PYG{n+nn}{psyclone}\PYG{n+nn}{.}\PYG{n+nn}{parse}\PYG{n+nn}{.}\PYG{n+nn}{algorithm}\PYG{+w}{ }\PYG{k+kn}{import} \PYG{n}{parse}
\PYG{g+gp}{\PYGZgt{}\PYGZgt{}\PYGZgt{} }\PYG{k+kn}{from}\PYG{+w}{ }\PYG{n+nn}{psyclone}\PYG{n+nn}{.}\PYG{n+nn}{psyGen}\PYG{+w}{ }\PYG{k+kn}{import} \PYG{n}{PSyFactory}
\PYG{g+gp}{\PYGZgt{}\PYGZgt{}\PYGZgt{}}
\PYG{g+gp}{\PYGZgt{}\PYGZgt{}\PYGZgt{} }\PYG{n}{API} \PYG{o}{=} \PYG{l+s+s2}{\PYGZdq{}}\PYG{l+s+s2}{lfric}\PYG{l+s+s2}{\PYGZdq{}}
\PYG{g+gp}{\PYGZgt{}\PYGZgt{}\PYGZgt{} }\PYG{n}{FILENAME} \PYG{o}{=} \PYG{l+s+s2}{\PYGZdq{}}\PYG{l+s+s2}{solver\PYGZus{}alg.x90}\PYG{l+s+s2}{\PYGZdq{}}
\PYG{g+gp}{\PYGZgt{}\PYGZgt{}\PYGZgt{} }\PYG{n}{ast}\PYG{p}{,} \PYG{n}{invokeInfo} \PYG{o}{=} \PYG{n}{parse}\PYG{p}{(}\PYG{n}{FILENAME}\PYG{p}{,} \PYG{n}{api}\PYG{o}{=}\PYG{n}{API}\PYG{p}{)}
\PYG{g+gp}{\PYGZgt{}\PYGZgt{}\PYGZgt{} }\PYG{n}{psy} \PYG{o}{=} \PYG{n}{PSyFactory}\PYG{p}{(}\PYG{n}{API}\PYG{p}{,} \PYG{n}{distributed\PYGZus{}memory}\PYG{o}{=}\PYG{k+kc}{False}\PYG{p}{)}\PYG{o}{.}\PYG{n}{create}\PYG{p}{(}\PYG{n}{invoke\PYGZus{}info}\PYG{p}{)}
\PYG{g+gp}{\PYGZgt{}\PYGZgt{}\PYGZgt{} }\PYG{n}{schedule} \PYG{o}{=} \PYG{n}{psy}\PYG{o}{.}\PYG{n}{invokes}\PYG{o}{.}\PYG{n}{get}\PYG{p}{(}\PYG{l+s+s1}{\PYGZsq{}}\PYG{l+s+s1}{invoke\PYGZus{}0}\PYG{l+s+s1}{\PYGZsq{}}\PYG{p}{)}\PYG{o}{.}\PYG{n}{schedule}
\PYG{g+gp}{\PYGZgt{}\PYGZgt{}\PYGZgt{}}
\PYG{g+gp}{\PYGZgt{}\PYGZgt{}\PYGZgt{} }\PYG{k+kn}{from}\PYG{+w}{ }\PYG{n+nn}{psyclone}\PYG{n+nn}{.}\PYG{n+nn}{domain}\PYG{n+nn}{.}\PYG{n+nn}{lfric}\PYG{n+nn}{.}\PYG{n+nn}{transformations}\PYG{+w}{ }\PYG{k+kn}{import} \PYG{n}{LFRicExtractTrans}
\PYG{g+gp}{\PYGZgt{}\PYGZgt{}\PYGZgt{} }\PYG{n}{etrans} \PYG{o}{=} \PYG{n}{LFRicExtractTrans}\PYG{p}{(}\PYG{p}{)}
\PYG{g+gp}{\PYGZgt{}\PYGZgt{}\PYGZgt{}}
\PYG{g+gp}{\PYGZgt{}\PYGZgt{}\PYGZgt{} }\PYG{c+c1}{\PYGZsh{} Apply LFRicExtractTrans transformation to selected Nodes}
\PYG{g+gp}{\PYGZgt{}\PYGZgt{}\PYGZgt{} }\PYG{n}{etrans}\PYG{o}{.}\PYG{n}{apply}\PYG{p}{(}\PYG{n}{schedule}\PYG{o}{.}\PYG{n}{children}\PYG{p}{[}\PYG{l+m+mi}{0}\PYG{p}{:}\PYG{l+m+mi}{3}\PYG{p}{]}\PYG{p}{)}
\PYG{g+gp}{\PYGZgt{}\PYGZgt{}\PYGZgt{} }\PYG{n+nb}{print}\PYG{p}{(}\PYG{n}{schedule}\PYG{o}{.}\PYG{n}{view}\PYG{p}{(}\PYG{p}{)}\PYG{p}{)}
\end{sphinxVerbatim}


\begin{fulllineitems}

\pysigstartsignatures
\pysiglinewithargsret{\sphinxbfcode{\sphinxupquote{apply}}}{\sphinxparam{\DUrole{n,n}{nodes}}\sphinxparamcomma \sphinxparam{\DUrole{n,n}{options}\DUrole{o,o}{=}\DUrole{default_value}{None}}}{}
\pysigstopsignatures
\sphinxAtStartPar
Apply this transformation to a subset of the nodes within a
schedule \sphinxhyphen{} i.e. enclose the specified Nodes in the schedule within
a single PSyData region. It first uses the CallTreeUtils to determine
input\sphinxhyphen{} and output\sphinxhyphen{}parameters. If requested, it will then call
the LFRicExtractDriverCreator to write the stand\sphinxhyphen{}alone driver
program. Then it will call apply of the base class.
\begin{quote}\begin{description}
\sphinxlineitem{Parameters}\begin{itemize}
\item {} 
\sphinxAtStartPar
\sphinxstyleliteralstrong{\sphinxupquote{nodes}} (\sphinxcode{\sphinxupquote{psyclone.psyir.nodes.Node}} or                      List{[}\sphinxcode{\sphinxupquote{psyclone.psyir.nodes.Node}}{]}) \textendash{} can be a single node or a list of nodes.

\item {} 
\sphinxAtStartPar
\sphinxstyleliteralstrong{\sphinxupquote{options}} (\sphinxstyleliteralemphasis{\sphinxupquote{Optional}}\sphinxstyleliteralemphasis{\sphinxupquote{{[}}}\sphinxstyleliteralemphasis{\sphinxupquote{Dict}}\sphinxstyleliteralemphasis{\sphinxupquote{{[}}}\sphinxhref{https://docs.python.org/3/library/stdtypes.html\#str}{\sphinxstyleliteralemphasis{\sphinxupquote{str}}}\sphinxstyleliteralemphasis{\sphinxupquote{, }}\sphinxstyleliteralemphasis{\sphinxupquote{Any}}\sphinxstyleliteralemphasis{\sphinxupquote{{]}}}\sphinxstyleliteralemphasis{\sphinxupquote{{]}}}) \textendash{} a dictionary with options for transformations.

\item {} 
\sphinxAtStartPar
\sphinxstyleliteralstrong{\sphinxupquote{options}}\sphinxstyleliteralstrong{\sphinxupquote{{[}}}\sphinxstyleliteralstrong{\sphinxupquote{"prefix"}}\sphinxstyleliteralstrong{\sphinxupquote{{]}}} (\sphinxhref{https://docs.python.org/3/library/stdtypes.html\#str}{\sphinxstyleliteralemphasis{\sphinxupquote{str}}}) \textendash{} a prefix to use for the PSyData module             name (\sphinxcode{\sphinxupquote{prefix\_psy\_data\_mod}}) and the PSyDataType             (\sphinxcode{\sphinxupquote{prefix\_PSyDataType}}) \sphinxhyphen{} a “\_” will be added automatically.             It defaults to “extract”, resulting in e.g.             \sphinxcode{\sphinxupquote{extract\_psy\_data\_mod}}.

\item {} 
\sphinxAtStartPar
\sphinxstyleliteralstrong{\sphinxupquote{options}}\sphinxstyleliteralstrong{\sphinxupquote{{[}}}\sphinxstyleliteralstrong{\sphinxupquote{"create\_driver"}}\sphinxstyleliteralstrong{\sphinxupquote{{]}}} (\sphinxhref{https://docs.python.org/3/library/functions.html\#bool}{\sphinxstyleliteralemphasis{\sphinxupquote{bool}}}) \textendash{} whether or not to create a             driver program at code\sphinxhyphen{}generation time. If set, the driver will             be created in the current working directory with the name             “driver\sphinxhyphen{}MODULE\sphinxhyphen{}REGION.f90” where MODULE and REGION will be the             corresponding values for this region. Defaults to False.

\item {} 
\sphinxAtStartPar
\sphinxstyleliteralstrong{\sphinxupquote{options}}\sphinxstyleliteralstrong{\sphinxupquote{{[}}}\sphinxstyleliteralstrong{\sphinxupquote{"region\_name"}}\sphinxstyleliteralstrong{\sphinxupquote{{]}}} (\sphinxstyleliteralemphasis{\sphinxupquote{Tuple}}\sphinxstyleliteralemphasis{\sphinxupquote{{[}}}\sphinxhref{https://docs.python.org/3/library/stdtypes.html\#str}{\sphinxstyleliteralemphasis{\sphinxupquote{str}}}\sphinxstyleliteralemphasis{\sphinxupquote{,}}\sphinxhref{https://docs.python.org/3/library/stdtypes.html\#str}{\sphinxstyleliteralemphasis{\sphinxupquote{str}}}\sphinxstyleliteralemphasis{\sphinxupquote{{]}}}) \textendash{} an optional name to             use for this PSyData area, provided as a 2\sphinxhyphen{}tuple containing a             location name followed by a local name. The pair of strings             should uniquely identify a region unless aggregate information             is required (and is supported by the runtime library).

\end{itemize}

\end{description}\end{quote}

\end{fulllineitems}



\begin{fulllineitems}

\pysigstartsignatures
\pysiglinewithargsret{\sphinxbfcode{\sphinxupquote{validate}}}{\sphinxparam{\DUrole{n,n}{node\_list}}\sphinxparamcomma \sphinxparam{\DUrole{n,n}{options}\DUrole{o,o}{=}\DUrole{default_value}{None}}}{}
\pysigstopsignatures
\sphinxAtStartPar
Perform Dynamo0.3 API specific validation checks before applying
the transformation.
\begin{quote}\begin{description}
\sphinxlineitem{Parameters}\begin{itemize}
\item {} 
\sphinxAtStartPar
\sphinxstyleliteralstrong{\sphinxupquote{node\_list}} (List{[}\sphinxcode{\sphinxupquote{psyclone.psyir.nodes.Node}}{]}) \textendash{} the list of Node(s) we are checking.

\item {} 
\sphinxAtStartPar
\sphinxstyleliteralstrong{\sphinxupquote{options}} (\sphinxstyleliteralemphasis{\sphinxupquote{Optional}}\sphinxstyleliteralemphasis{\sphinxupquote{{[}}}\sphinxstyleliteralemphasis{\sphinxupquote{Dict}}\sphinxstyleliteralemphasis{\sphinxupquote{{[}}}\sphinxhref{https://docs.python.org/3/library/stdtypes.html\#str}{\sphinxstyleliteralemphasis{\sphinxupquote{str}}}\sphinxstyleliteralemphasis{\sphinxupquote{, }}\sphinxstyleliteralemphasis{\sphinxupquote{Any}}\sphinxstyleliteralemphasis{\sphinxupquote{{]}}}\sphinxstyleliteralemphasis{\sphinxupquote{{]}}}) \textendash{} a dictionary with options for transformations.

\end{itemize}

\sphinxlineitem{Raises}
\sphinxAtStartPar
\sphinxstyleliteralstrong{\sphinxupquote{TransformationError}} \textendash{} if transformation is applied to a Loop                                      over cells in a colour without its                                      parent Loop over colours.

\end{description}\end{quote}

\end{fulllineitems}


\end{fulllineitems}



\begin{fulllineitems}

\pysigstartsignatures
\pysigline{\sphinxbfcode{\sphinxupquote{class\DUrole{w,w}{  }}}\sphinxcode{\sphinxupquote{psyclone.domain.lfric.transformations.}}\sphinxbfcode{\sphinxupquote{LFRicLoopFuseTrans}}}
\pysigstopsignatures
\sphinxAtStartPar
LFRic API specialisation of the
\sphinxcode{\sphinxupquote{base class}} in order to fuse two Dynamo
loops after performing validity checks. For example:

\begin{sphinxVerbatim}[commandchars=\\\{\}]
\PYG{g+gp}{\PYGZgt{}\PYGZgt{}\PYGZgt{} }\PYG{k+kn}{from}\PYG{+w}{ }\PYG{n+nn}{psyclone}\PYG{n+nn}{.}\PYG{n+nn}{parse}\PYG{n+nn}{.}\PYG{n+nn}{algorithm}\PYG{+w}{ }\PYG{k+kn}{import} \PYG{n}{parse}
\PYG{g+gp}{\PYGZgt{}\PYGZgt{}\PYGZgt{} }\PYG{k+kn}{from}\PYG{+w}{ }\PYG{n+nn}{psyclone}\PYG{n+nn}{.}\PYG{n+nn}{psyGen}\PYG{+w}{ }\PYG{k+kn}{import} \PYG{n}{PSyFactory}
\PYG{g+gp}{\PYGZgt{}\PYGZgt{}\PYGZgt{}}
\PYG{g+gp}{\PYGZgt{}\PYGZgt{}\PYGZgt{} }\PYG{n}{API} \PYG{o}{=} \PYG{l+s+s2}{\PYGZdq{}}\PYG{l+s+s2}{lfric}\PYG{l+s+s2}{\PYGZdq{}}
\PYG{g+gp}{\PYGZgt{}\PYGZgt{}\PYGZgt{} }\PYG{n}{FILENAME} \PYG{o}{=} \PYG{l+s+s2}{\PYGZdq{}}\PYG{l+s+s2}{alg.x90}\PYG{l+s+s2}{\PYGZdq{}}
\PYG{g+gp}{\PYGZgt{}\PYGZgt{}\PYGZgt{} }\PYG{n}{ast}\PYG{p}{,} \PYG{n}{invokeInfo} \PYG{o}{=} \PYG{n}{parse}\PYG{p}{(}\PYG{n}{FILENAME}\PYG{p}{,} \PYG{n}{api}\PYG{o}{=}\PYG{n}{API}\PYG{p}{)}
\PYG{g+gp}{\PYGZgt{}\PYGZgt{}\PYGZgt{} }\PYG{n}{psy} \PYG{o}{=} \PYG{n}{PSyFactory}\PYG{p}{(}\PYG{n}{API}\PYG{p}{,} \PYG{n}{distributed\PYGZus{}memory}\PYG{o}{=}\PYG{k+kc}{False}\PYG{p}{)}\PYG{o}{.}\PYG{n}{create}\PYG{p}{(}\PYG{n}{invoke\PYGZus{}info}\PYG{p}{)}
\PYG{g+gp}{\PYGZgt{}\PYGZgt{}\PYGZgt{} }\PYG{n}{schedule} \PYG{o}{=} \PYG{n}{psy}\PYG{o}{.}\PYG{n}{invokes}\PYG{o}{.}\PYG{n}{get}\PYG{p}{(}\PYG{l+s+s1}{\PYGZsq{}}\PYG{l+s+s1}{invoke\PYGZus{}0}\PYG{l+s+s1}{\PYGZsq{}}\PYG{p}{)}\PYG{o}{.}\PYG{n}{schedule}
\PYG{g+gp}{\PYGZgt{}\PYGZgt{}\PYGZgt{}}
\PYG{g+gp}{\PYGZgt{}\PYGZgt{}\PYGZgt{} }\PYG{k+kn}{from}\PYG{+w}{ }\PYG{n+nn}{psyclone}\PYG{n+nn}{.}\PYG{n+nn}{domain}\PYG{n+nn}{.}\PYG{n+nn}{lfric}\PYG{n+nn}{.}\PYG{n+nn}{transformations}\PYG{+w}{ }\PYG{k+kn}{import} \PYG{n}{LFRicLoopFuseTrans}
\PYG{g+gp}{\PYGZgt{}\PYGZgt{}\PYGZgt{} }\PYG{n}{ftrans} \PYG{o}{=}  \PYG{n}{LFRicLoopFuseTrans}\PYG{p}{(}\PYG{p}{)}
\PYG{g+gp}{\PYGZgt{}\PYGZgt{}\PYGZgt{}}
\PYG{g+gp}{\PYGZgt{}\PYGZgt{}\PYGZgt{} }\PYG{n}{ftrans}\PYG{o}{.}\PYG{n}{apply}\PYG{p}{(}\PYG{n}{schedule}\PYG{p}{[}\PYG{l+m+mi}{0}\PYG{p}{]}\PYG{p}{,} \PYG{n}{schedule}\PYG{p}{[}\PYG{l+m+mi}{1}\PYG{p}{]}\PYG{p}{)}
\PYG{g+gp}{\PYGZgt{}\PYGZgt{}\PYGZgt{} }\PYG{n+nb}{print}\PYG{p}{(}\PYG{n}{schedule}\PYG{o}{.}\PYG{n}{view}\PYG{p}{(}\PYG{p}{)}\PYG{p}{)}
\end{sphinxVerbatim}

\sphinxAtStartPar
The optional argument \sphinxtitleref{same\_space} can be set as

\begin{sphinxVerbatim}[commandchars=\\\{\}]
\PYG{g+gp}{\PYGZgt{}\PYGZgt{}\PYGZgt{} }\PYG{n}{ftrans}\PYG{o}{.}\PYG{n}{apply}\PYG{p}{(}\PYG{n}{schedule}\PYG{p}{[}\PYG{l+m+mi}{0}\PYG{p}{]}\PYG{p}{,} \PYG{n}{schedule}\PYG{p}{[}\PYG{l+m+mi}{1}\PYG{p}{]}\PYG{p}{,} \PYG{p}{\PYGZob{}}\PYG{l+s+s2}{\PYGZdq{}}\PYG{l+s+s2}{same\PYGZus{}space}\PYG{l+s+s2}{\PYGZdq{}}\PYG{p}{:} \PYG{k+kc}{True}\PYG{p}{\PYGZcb{}}\PYG{p}{)}
\end{sphinxVerbatim}

\sphinxAtStartPar
when applying the transformation.


\begin{fulllineitems}

\pysigstartsignatures
\pysiglinewithargsret{\sphinxbfcode{\sphinxupquote{validate}}}{\sphinxparam{\DUrole{n,n}{node1}}\sphinxparamcomma \sphinxparam{\DUrole{n,n}{node2}}\sphinxparamcomma \sphinxparam{\DUrole{n,n}{options}\DUrole{o,o}{=}\DUrole{default_value}{None}}}{}
\pysigstopsignatures
\sphinxAtStartPar
Performs various checks to ensure that it is valid to apply
the LFRicLoopFuseTrans transformation to the supplied loops.
\begin{quote}\begin{description}
\sphinxlineitem{Parameters}\begin{itemize}
\item {} 
\sphinxAtStartPar
\sphinxstyleliteralstrong{\sphinxupquote{node1}} (\sphinxcode{\sphinxupquote{psyclone.domain.lfric.LFRicLoop}}) \textendash{} the first Loop to fuse.

\item {} 
\sphinxAtStartPar
\sphinxstyleliteralstrong{\sphinxupquote{node2}} (\sphinxcode{\sphinxupquote{psyclone.domain.lfric.LFRicLoop}}) \textendash{} the second Loop to fuse.

\item {} 
\sphinxAtStartPar
\sphinxstyleliteralstrong{\sphinxupquote{options}} (\sphinxstyleliteralemphasis{\sphinxupquote{Optional}}\sphinxstyleliteralemphasis{\sphinxupquote{{[}}}\sphinxstyleliteralemphasis{\sphinxupquote{Dict}}\sphinxstyleliteralemphasis{\sphinxupquote{{[}}}\sphinxhref{https://docs.python.org/3/library/stdtypes.html\#str}{\sphinxstyleliteralemphasis{\sphinxupquote{str}}}\sphinxstyleliteralemphasis{\sphinxupquote{, }}\sphinxstyleliteralemphasis{\sphinxupquote{Any}}\sphinxstyleliteralemphasis{\sphinxupquote{{]}}}\sphinxstyleliteralemphasis{\sphinxupquote{{]}}}) \textendash{} a dictionary with options for transformations.

\item {} 
\sphinxAtStartPar
\sphinxstyleliteralstrong{\sphinxupquote{options}}\sphinxstyleliteralstrong{\sphinxupquote{{[}}}\sphinxstyleliteralstrong{\sphinxupquote{"same\_space"}}\sphinxstyleliteralstrong{\sphinxupquote{{]}}} (\sphinxhref{https://docs.python.org/3/library/functions.html\#bool}{\sphinxstyleliteralemphasis{\sphinxupquote{bool}}}) \textendash{} this optional flag, set to \sphinxtitleref{True},             asserts that an unknown iteration space (i.e. \sphinxtitleref{ANY\_SPACE})             matches the other iteration space. This is set at the user’s own             risk. If both iteration spaces are discontinuous the loops can be             fused without having to use the \sphinxtitleref{same\_space} flag.

\end{itemize}

\sphinxlineitem{Raises}\begin{itemize}
\item {} 
\sphinxAtStartPar
\sphinxstyleliteralstrong{\sphinxupquote{TransformationError}} \textendash{} if either of the supplied loops contains                                      an inter\sphinxhyphen{}grid kernel.

\item {} 
\sphinxAtStartPar
\sphinxstyleliteralstrong{\sphinxupquote{TransformationError}} \textendash{} if one or both function spaces have                                      invalid names.

\item {} 
\sphinxAtStartPar
\sphinxstyleliteralstrong{\sphinxupquote{TransformationError}} \textendash{} if the \sphinxtitleref{same\_space} flag was set, but                                      does not apply because neither field                                      is on \sphinxtitleref{ANY\_SPACE} or the spaces are not                                      the same.

\item {} 
\sphinxAtStartPar
\sphinxstyleliteralstrong{\sphinxupquote{TransformationError}} \textendash{} if one or more of the iteration spaces                                      is unknown (\sphinxtitleref{ANY\_SPACE}) and the                                      \sphinxtitleref{same\_space} flag is not set to \sphinxtitleref{True}.

\item {} 
\sphinxAtStartPar
\sphinxstyleliteralstrong{\sphinxupquote{TransformationError}} \textendash{} if the loops are over different spaces                                      that are not both discontinuous and                                      the loops both iterate over cells.

\item {} 
\sphinxAtStartPar
\sphinxstyleliteralstrong{\sphinxupquote{TransformationError}} \textendash{} if the loops’ upper bound names are                                      not the same.

\item {} 
\sphinxAtStartPar
\sphinxstyleliteralstrong{\sphinxupquote{TransformationError}} \textendash{} if the halo\sphinxhyphen{}depth indices of two loops                                      are not the same.

\item {} 
\sphinxAtStartPar
\sphinxstyleliteralstrong{\sphinxupquote{TransformationError}} \textendash{} if each loop already contains a reduction.

\item {} 
\sphinxAtStartPar
\sphinxstyleliteralstrong{\sphinxupquote{TransformationError}} \textendash{} if the first loop has a reduction and                                      the second loop reads the result of                                      the reduction.

\end{itemize}

\end{description}\end{quote}

\end{fulllineitems}


\end{fulllineitems}



\begin{fulllineitems}

\pysigstartsignatures
\pysigline{\sphinxbfcode{\sphinxupquote{class\DUrole{w,w}{  }}}\sphinxcode{\sphinxupquote{psyclone.domain.lfric.transformations.}}\sphinxbfcode{\sphinxupquote{RaisePSyIR2LFRicKernTrans}}}
\pysigstopsignatures
\sphinxAtStartPar
Raise a generic PSyIR representation of a kernel\sphinxhyphen{}layer routine and
metadata to an LFRic version with specialised domain\sphinxhyphen{}specific
nodes and symbols. This is currently limited to the specialisation
of kernel metadata.

\begin{sphinxVerbatim}[commandchars=\\\{\}]
\PYG{g+gp}{\PYGZgt{}\PYGZgt{}\PYGZgt{} }\PYG{k+kn}{from}\PYG{+w}{ }\PYG{n+nn}{psyclone}\PYG{n+nn}{.}\PYG{n+nn}{configuration}\PYG{+w}{ }\PYG{k+kn}{import} \PYG{n}{Config}
\PYG{g+gp}{\PYGZgt{}\PYGZgt{}\PYGZgt{} }\PYG{k+kn}{from}\PYG{+w}{ }\PYG{n+nn}{psyclone}\PYG{n+nn}{.}\PYG{n+nn}{domain}\PYG{n+nn}{.}\PYG{n+nn}{lfric}\PYG{n+nn}{.}\PYG{n+nn}{transformations}\PYG{+w}{ }\PYG{k+kn}{import}             \PYG{n}{RaisePSyIR2LFRicKernTrans}
\PYG{g+gp}{\PYGZgt{}\PYGZgt{}\PYGZgt{} }\PYG{k+kn}{from}\PYG{+w}{ }\PYG{n+nn}{psyclone}\PYG{n+nn}{.}\PYG{n+nn}{psyir}\PYG{n+nn}{.}\PYG{n+nn}{frontend}\PYG{n+nn}{.}\PYG{n+nn}{fortran}\PYG{+w}{ }\PYG{k+kn}{import} \PYG{n}{FortranReader}
\PYG{g+gp}{\PYGZgt{}\PYGZgt{}\PYGZgt{} }\PYG{n}{config} \PYG{o}{=} \PYG{n}{Config}\PYG{o}{.}\PYG{n}{get}\PYG{p}{(}\PYG{p}{)}\PYG{o}{.}\PYG{n}{api\PYGZus{}conf}\PYG{p}{(}\PYG{l+s+s2}{\PYGZdq{}}\PYG{l+s+s2}{lfric}\PYG{l+s+s2}{\PYGZdq{}}\PYG{p}{)}
\PYG{g+gp}{\PYGZgt{}\PYGZgt{}\PYGZgt{} }\PYG{n}{CODE} \PYG{o}{=} \PYG{p}{(}\PYG{l+s+s2}{\PYGZdq{}\PYGZdq{}\PYGZdq{}}
\PYG{g+gp}{... }\PYG{l+s+s2}{MODULE example}
\PYG{g+gp}{... }\PYG{l+s+s2}{TYPE, EXTENDS(kernel\PYGZus{}type) :: compute\PYGZus{}cu}
\PYG{g+gp}{... }\PYG{l+s+s2}{  TYPE(arg\PYGZus{}type), DIMENSION(4) :: meta\PYGZus{}args = (/     \PYGZam{}}
\PYG{g+gp}{... }\PYG{l+s+s2}{    arg\PYGZus{}type(GH\PYGZus{}FIELD, GH\PYGZus{}REAL, GH\PYGZus{}INC, W1),         \PYGZam{}}
\PYG{g+gp}{... }\PYG{l+s+s2}{    arg\PYGZus{}type(GH\PYGZus{}FIELD, GH\PYGZus{}REAL, GH\PYGZus{}READ, W3),        \PYGZam{}}
\PYG{g+gp}{... }\PYG{l+s+s2}{    arg\PYGZus{}type(GH\PYGZus{}FIELD, GH\PYGZus{}REAL, GH\PYGZus{}READ, W3),        \PYGZam{}}
\PYG{g+gp}{... }\PYG{l+s+s2}{    arg\PYGZus{}type(GH\PYGZus{}FIELD, GH\PYGZus{}REAL, GH\PYGZus{}READ, W3)/)}
\PYG{g+gp}{... }\PYG{l+s+s2}{  INTEGER :: OPERATES\PYGZus{}ON = CELL\PYGZus{}COLUMN}
\PYG{g+gp}{... }\PYG{l+s+s2}{CONTAINS}
\PYG{g+gp}{... }\PYG{l+s+s2}{  PROCEDURE, NOPASS :: code =\PYGZgt{} compute\PYGZus{}cu\PYGZus{}code}
\PYG{g+gp}{... }\PYG{l+s+s2}{END TYPE compute\PYGZus{}cu}
\PYG{g+gp}{... }\PYG{l+s+s2}{contains}
\PYG{g+gp}{... }\PYG{l+s+s2}{  subroutine compute\PYGZus{}cu\PYGZus{}code()}
\PYG{g+gp}{... }\PYG{l+s+s2}{  end subroutine}
\PYG{g+gp}{... }\PYG{l+s+s2}{end module}\PYG{l+s+s2}{\PYGZdq{}\PYGZdq{}\PYGZdq{}}\PYG{p}{)}
\PYG{g+gp}{\PYGZgt{}\PYGZgt{}\PYGZgt{} }\PYG{n}{fortran\PYGZus{}reader} \PYG{o}{=} \PYG{n}{FortranReader}\PYG{p}{(}\PYG{p}{)}
\PYG{g+gp}{\PYGZgt{}\PYGZgt{}\PYGZgt{} }\PYG{n}{kernel\PYGZus{}container} \PYG{o}{=} \PYG{n}{fortran\PYGZus{}reader}\PYG{o}{.}\PYG{n}{psyir\PYGZus{}from\PYGZus{}source}\PYG{p}{(}\PYG{n}{CODE}\PYG{p}{)}
\PYG{g+gp}{\PYGZgt{}\PYGZgt{}\PYGZgt{} }\PYG{n}{trans} \PYG{o}{=} \PYG{n}{RaisePSyIR2LFRicKernTrans}\PYG{p}{(}\PYG{p}{)}
\PYG{g+gp}{\PYGZgt{}\PYGZgt{}\PYGZgt{} }\PYG{n}{trans}\PYG{o}{.}\PYG{n}{apply}\PYG{p}{(}\PYG{n}{kernel\PYGZus{}container}\PYG{p}{,} \PYG{p}{\PYGZob{}}\PYG{l+s+s2}{\PYGZdq{}}\PYG{l+s+s2}{metadata\PYGZus{}name}\PYG{l+s+s2}{\PYGZdq{}}\PYG{p}{:} \PYG{l+s+s2}{\PYGZdq{}}\PYG{l+s+s2}{compute\PYGZus{}cu}\PYG{l+s+s2}{\PYGZdq{}}\PYG{p}{\PYGZcb{}}\PYG{p}{)}
\end{sphinxVerbatim}


\begin{fulllineitems}

\pysigstartsignatures
\pysiglinewithargsret{\sphinxbfcode{\sphinxupquote{apply}}}{\sphinxparam{\DUrole{n,n}{node}}\sphinxparamcomma \sphinxparam{\DUrole{n,n}{options}\DUrole{o,o}{=}\DUrole{default_value}{None}}}{}
\pysigstopsignatures
\sphinxAtStartPar
Raise the supplied language\sphinxhyphen{}level kernel to LFRic\sphinxhyphen{}specific kernel
PSyIR. Specialises the kernel container to an LFRic\sphinxhyphen{}specific
subclass, populates this subclass with the kernel metadata
extracted from the metadata symbol as specified in
metadata\_name (which is supplied via the options argument) and
removes the symbol from the symbol table.
\begin{quote}\begin{description}
\sphinxlineitem{Parameters}\begin{itemize}
\item {} 
\sphinxAtStartPar
\sphinxstyleliteralstrong{\sphinxupquote{node}} (\sphinxcode{\sphinxupquote{psyclone.psyir.node.Container}}) \textendash{} a kernel represented in generic PSyIR.

\item {} 
\sphinxAtStartPar
\sphinxstyleliteralstrong{\sphinxupquote{options}} (\sphinxstyleliteralemphasis{\sphinxupquote{Optional}}\sphinxstyleliteralemphasis{\sphinxupquote{{[}}}\sphinxstyleliteralemphasis{\sphinxupquote{Dict}}\sphinxstyleliteralemphasis{\sphinxupquote{{[}}}\sphinxstyleliteralemphasis{\sphinxupquote{str: str}}\sphinxstyleliteralemphasis{\sphinxupquote{{]}}}\sphinxstyleliteralemphasis{\sphinxupquote{{]}}}) \textendash{} a dictionary with options for transformations.             This is expected to contain the metadata\_name.

\end{itemize}

\end{description}\end{quote}

\end{fulllineitems}



\begin{fulllineitems}

\pysigstartsignatures
\pysiglinewithargsret{\sphinxbfcode{\sphinxupquote{validate}}}{\sphinxparam{\DUrole{n,n}{node}}\sphinxparamcomma \sphinxparam{\DUrole{n,n}{options}\DUrole{o,o}{=}\DUrole{default_value}{None}}}{}
\pysigstopsignatures
\sphinxAtStartPar
Validate the supplied PSyIR tree.
\begin{quote}\begin{description}
\sphinxlineitem{Parameters}\begin{itemize}
\item {} 
\sphinxAtStartPar
\sphinxstyleliteralstrong{\sphinxupquote{node}} (\sphinxcode{\sphinxupquote{psyclone.psyir.node.Container}}) \textendash{} a PSyIR node that is the root of a PSyIR tree.

\item {} 
\sphinxAtStartPar
\sphinxstyleliteralstrong{\sphinxupquote{options}} (\sphinxstyleliteralemphasis{\sphinxupquote{Optional}}\sphinxstyleliteralemphasis{\sphinxupquote{{[}}}\sphinxstyleliteralemphasis{\sphinxupquote{Dict}}\sphinxstyleliteralemphasis{\sphinxupquote{{[}}}\sphinxstyleliteralemphasis{\sphinxupquote{str: str}}\sphinxstyleliteralemphasis{\sphinxupquote{{]}}}\sphinxstyleliteralemphasis{\sphinxupquote{{]}}}) \textendash{} a dictionary with options for transformations.

\end{itemize}

\sphinxlineitem{Raises}\begin{itemize}
\item {} 
\sphinxAtStartPar
\sphinxstyleliteralstrong{\sphinxupquote{TransformationError}} \textendash{} if the supplied node is not a
Container.

\item {} 
\sphinxAtStartPar
\sphinxstyleliteralstrong{\sphinxupquote{TransformationError}} \textendash{} if the supplied node argument has
a parent.

\item {} 
\sphinxAtStartPar
\sphinxstyleliteralstrong{\sphinxupquote{TransformationError}} \textendash{} if the metadata name has not been
provided in the options argument.

\item {} 
\sphinxAtStartPar
\sphinxstyleliteralstrong{\sphinxupquote{TransformationError}} \textendash{} if the metadata name has not been
set or does not exist in the code.

\item {} 
\sphinxAtStartPar
\sphinxstyleliteralstrong{\sphinxupquote{TransformationError}} \textendash{} if the metadata symbol does not
reside in a Container (as opposed to a FileContainer).

\end{itemize}

\end{description}\end{quote}

\end{fulllineitems}


\end{fulllineitems}



\begin{fulllineitems}

\pysigstartsignatures
\pysiglinewithargsret{\sphinxbfcode{\sphinxupquote{class\DUrole{w,w}{  }}}\sphinxcode{\sphinxupquote{psyclone.transformations.}}\sphinxbfcode{\sphinxupquote{DynamoOMPParallelLoopTrans}}}{\sphinxparam{\DUrole{n,n}{omp\_directive}\DUrole{o,o}{=}\DUrole{default_value}{\textquotesingle{}do\textquotesingle{}}}\sphinxparamcomma \sphinxparam{\DUrole{n,n}{omp\_schedule}\DUrole{o,o}{=}\DUrole{default_value}{\textquotesingle{}static\textquotesingle{}}}}{}
\pysigstopsignatures
\sphinxAtStartPar
Dynamo\sphinxhyphen{}specific OpenMP loop transformation. Adds Dynamo specific
validity checks. Actual transformation is done by the
\sphinxcode{\sphinxupquote{base class}}.
\begin{quote}\begin{description}
\sphinxlineitem{Parameters}\begin{itemize}
\item {} 
\sphinxAtStartPar
\sphinxstyleliteralstrong{\sphinxupquote{omp\_directive}} (\sphinxhref{https://docs.python.org/3/library/stdtypes.html\#str}{\sphinxstyleliteralemphasis{\sphinxupquote{str}}}) \textendash{} choose which OpenMP loop directive to use.
Defaults to “do”.

\item {} 
\sphinxAtStartPar
\sphinxstyleliteralstrong{\sphinxupquote{omp\_schedule}} (\sphinxhref{https://docs.python.org/3/library/stdtypes.html\#str}{\sphinxstyleliteralemphasis{\sphinxupquote{str}}}) \textendash{} the OpenMP schedule to use. Must be one of
‘runtime’, ‘static’, ‘dynamic’, ‘guided’ or ‘auto’. Defaults to
‘static’.

\end{itemize}

\end{description}\end{quote}


\begin{fulllineitems}

\pysigstartsignatures
\pysiglinewithargsret{\sphinxbfcode{\sphinxupquote{validate}}}{\sphinxparam{\DUrole{n,n}{node}}\sphinxparamcomma \sphinxparam{\DUrole{n,n}{options}\DUrole{o,o}{=}\DUrole{default_value}{None}}}{}
\pysigstopsignatures
\sphinxAtStartPar
Perform LFRic\sphinxhyphen{}specific loop validity checks then call the \sphinxtitleref{validate}
method of the base class.
\begin{quote}\begin{description}
\sphinxlineitem{Parameters}\begin{itemize}
\item {} 
\sphinxAtStartPar
\sphinxstyleliteralstrong{\sphinxupquote{node}} (\sphinxcode{\sphinxupquote{psyclone.psyir.nodes.Node}}) \textendash{} the Node in the Schedule to check

\item {} 
\sphinxAtStartPar
\sphinxstyleliteralstrong{\sphinxupquote{options}} (\sphinxstyleliteralemphasis{\sphinxupquote{Optional}}\sphinxstyleliteralemphasis{\sphinxupquote{{[}}}\sphinxstyleliteralemphasis{\sphinxupquote{Dict}}\sphinxstyleliteralemphasis{\sphinxupquote{{[}}}\sphinxhref{https://docs.python.org/3/library/stdtypes.html\#str}{\sphinxstyleliteralemphasis{\sphinxupquote{str}}}\sphinxstyleliteralemphasis{\sphinxupquote{, }}\sphinxstyleliteralemphasis{\sphinxupquote{Any}}\sphinxstyleliteralemphasis{\sphinxupquote{{]}}}\sphinxstyleliteralemphasis{\sphinxupquote{{]}}}) \textendash{} a dictionary with options for transformations.

\end{itemize}

\sphinxlineitem{Raises}\begin{itemize}
\item {} 
\sphinxAtStartPar
\sphinxstyleliteralstrong{\sphinxupquote{TransformationError}} \textendash{} if the supplied Node is not a LFRicLoop.

\item {} 
\sphinxAtStartPar
\sphinxstyleliteralstrong{\sphinxupquote{TransformationError}} \textendash{} if the associated loop requires
colouring.

\end{itemize}

\end{description}\end{quote}

\end{fulllineitems}


\end{fulllineitems}



\begin{fulllineitems}

\pysigstartsignatures
\pysigline{\sphinxbfcode{\sphinxupquote{class\DUrole{w,w}{  }}}\sphinxcode{\sphinxupquote{psyclone.transformations.}}\sphinxbfcode{\sphinxupquote{Dynamo0p3AsyncHaloExchangeTrans}}}
\pysigstopsignatures
\sphinxAtStartPar
Splits a synchronous halo exchange into a halo exchange start and
halo exchange end. For example:

\begin{sphinxVerbatim}[commandchars=\\\{\}]
\PYG{g+gp}{\PYGZgt{}\PYGZgt{}\PYGZgt{} }\PYG{k+kn}{from}\PYG{+w}{ }\PYG{n+nn}{psyclone}\PYG{n+nn}{.}\PYG{n+nn}{parse}\PYG{n+nn}{.}\PYG{n+nn}{algorithm}\PYG{+w}{ }\PYG{k+kn}{import} \PYG{n}{parse}
\PYG{g+gp}{\PYGZgt{}\PYGZgt{}\PYGZgt{} }\PYG{k+kn}{from}\PYG{+w}{ }\PYG{n+nn}{psyclone}\PYG{n+nn}{.}\PYG{n+nn}{psyGen}\PYG{+w}{ }\PYG{k+kn}{import} \PYG{n}{PSyFactory}
\PYG{g+gp}{\PYGZgt{}\PYGZgt{}\PYGZgt{} }\PYG{n}{api} \PYG{o}{=} \PYG{l+s+s2}{\PYGZdq{}}\PYG{l+s+s2}{lfric}\PYG{l+s+s2}{\PYGZdq{}}
\PYG{g+gp}{\PYGZgt{}\PYGZgt{}\PYGZgt{} }\PYG{n}{ast}\PYG{p}{,} \PYG{n}{invokeInfo} \PYG{o}{=} \PYG{n}{parse}\PYG{p}{(}\PYG{l+s+s2}{\PYGZdq{}}\PYG{l+s+s2}{file.f90}\PYG{l+s+s2}{\PYGZdq{}}\PYG{p}{,} \PYG{n}{api}\PYG{o}{=}\PYG{n}{api}\PYG{p}{)}
\PYG{g+gp}{\PYGZgt{}\PYGZgt{}\PYGZgt{} }\PYG{n}{psy}\PYG{o}{=}\PYG{n}{PSyFactory}\PYG{p}{(}\PYG{n}{api}\PYG{p}{)}\PYG{o}{.}\PYG{n}{create}\PYG{p}{(}\PYG{n}{invokeInfo}\PYG{p}{)}
\PYG{g+gp}{\PYGZgt{}\PYGZgt{}\PYGZgt{} }\PYG{n}{schedule} \PYG{o}{=} \PYG{n}{psy}\PYG{o}{.}\PYG{n}{invokes}\PYG{o}{.}\PYG{n}{get}\PYG{p}{(}\PYG{l+s+s1}{\PYGZsq{}}\PYG{l+s+s1}{invoke\PYGZus{}0}\PYG{l+s+s1}{\PYGZsq{}}\PYG{p}{)}\PYG{o}{.}\PYG{n}{schedule}
\PYG{g+gp}{\PYGZgt{}\PYGZgt{}\PYGZgt{} }\PYG{c+c1}{\PYGZsh{} Uncomment the following line to see a text view of the schedule}
\PYG{g+gp}{\PYGZgt{}\PYGZgt{}\PYGZgt{} }\PYG{c+c1}{\PYGZsh{} print(schedule.view())}
\PYG{g+gp}{\PYGZgt{}\PYGZgt{}\PYGZgt{}}
\PYG{g+gp}{\PYGZgt{}\PYGZgt{}\PYGZgt{} }\PYG{k+kn}{from}\PYG{+w}{ }\PYG{n+nn}{psyclone}\PYG{n+nn}{.}\PYG{n+nn}{transformations}\PYG{+w}{ }\PYG{k+kn}{import} \PYG{n}{Dynamo0p3AsyncHaloExchangeTrans}
\PYG{g+gp}{\PYGZgt{}\PYGZgt{}\PYGZgt{} }\PYG{n}{trans} \PYG{o}{=} \PYG{n}{Dynamo0p3AsyncHaloExchangeTrans}\PYG{p}{(}\PYG{p}{)}
\PYG{g+gp}{\PYGZgt{}\PYGZgt{}\PYGZgt{} }\PYG{n}{trans}\PYG{o}{.}\PYG{n}{apply}\PYG{p}{(}\PYG{n}{schedule}\PYG{o}{.}\PYG{n}{children}\PYG{p}{[}\PYG{l+m+mi}{0}\PYG{p}{]}\PYG{p}{)}
\PYG{g+gp}{\PYGZgt{}\PYGZgt{}\PYGZgt{} }\PYG{c+c1}{\PYGZsh{} Uncomment the following line to see a text view of the schedule}
\PYG{g+gp}{\PYGZgt{}\PYGZgt{}\PYGZgt{} }\PYG{c+c1}{\PYGZsh{} print(schedule.view())}
\end{sphinxVerbatim}


\begin{fulllineitems}

\pysigstartsignatures
\pysiglinewithargsret{\sphinxbfcode{\sphinxupquote{apply}}}{\sphinxparam{\DUrole{n,n}{node}}\sphinxparamcomma \sphinxparam{\DUrole{n,n}{options}\DUrole{o,o}{=}\DUrole{default_value}{None}}}{}
\pysigstopsignatures
\sphinxAtStartPar
Transforms a synchronous halo exchange, represented by a
HaloExchange node, into an asynchronous halo exchange,
represented by HaloExchangeStart and HaloExchangeEnd nodes.
\begin{quote}\begin{description}
\sphinxlineitem{Parameters}\begin{itemize}
\item {} 
\sphinxAtStartPar
\sphinxstyleliteralstrong{\sphinxupquote{node}} (\sphinxcode{\sphinxupquote{psyclone.psygen.HaloExchange}}) \textendash{} a synchronous haloexchange node.

\item {} 
\sphinxAtStartPar
\sphinxstyleliteralstrong{\sphinxupquote{options}} (\sphinxstyleliteralemphasis{\sphinxupquote{Optional}}\sphinxstyleliteralemphasis{\sphinxupquote{{[}}}\sphinxstyleliteralemphasis{\sphinxupquote{Dict}}\sphinxstyleliteralemphasis{\sphinxupquote{{[}}}\sphinxhref{https://docs.python.org/3/library/stdtypes.html\#str}{\sphinxstyleliteralemphasis{\sphinxupquote{str}}}\sphinxstyleliteralemphasis{\sphinxupquote{, }}\sphinxstyleliteralemphasis{\sphinxupquote{Any}}\sphinxstyleliteralemphasis{\sphinxupquote{{]}}}\sphinxstyleliteralemphasis{\sphinxupquote{{]}}}) \textendash{} a dictionary with options for transformations.

\end{itemize}

\end{description}\end{quote}

\end{fulllineitems}



\begin{fulllineitems}

\pysigstartsignatures
\pysigline{\sphinxbfcode{\sphinxupquote{property\DUrole{w,w}{  }}}\sphinxbfcode{\sphinxupquote{name}}}
\pysigstopsignatures\begin{quote}\begin{description}
\sphinxlineitem{Returns}
\sphinxAtStartPar
the name of this transformation as a string.

\sphinxlineitem{Return type}
\sphinxAtStartPar
\sphinxhref{https://docs.python.org/3/library/stdtypes.html\#str}{str}

\end{description}\end{quote}

\end{fulllineitems}



\begin{fulllineitems}

\pysigstartsignatures
\pysiglinewithargsret{\sphinxbfcode{\sphinxupquote{validate}}}{\sphinxparam{\DUrole{n,n}{node}}\sphinxparamcomma \sphinxparam{\DUrole{n,n}{options}}}{}
\pysigstopsignatures
\sphinxAtStartPar
Internal method to check whether the node is valid for this
transformation.
\begin{quote}\begin{description}
\sphinxlineitem{Parameters}\begin{itemize}
\item {} 
\sphinxAtStartPar
\sphinxstyleliteralstrong{\sphinxupquote{node}} (\sphinxcode{\sphinxupquote{psyclone.psygen.HaloExchange}}) \textendash{} a synchronous Halo Exchange node

\item {} 
\sphinxAtStartPar
\sphinxstyleliteralstrong{\sphinxupquote{options}} (\sphinxstyleliteralemphasis{\sphinxupquote{Optional}}\sphinxstyleliteralemphasis{\sphinxupquote{{[}}}\sphinxstyleliteralemphasis{\sphinxupquote{Dict}}\sphinxstyleliteralemphasis{\sphinxupquote{{[}}}\sphinxhref{https://docs.python.org/3/library/stdtypes.html\#str}{\sphinxstyleliteralemphasis{\sphinxupquote{str}}}\sphinxstyleliteralemphasis{\sphinxupquote{, }}\sphinxstyleliteralemphasis{\sphinxupquote{Any}}\sphinxstyleliteralemphasis{\sphinxupquote{{]}}}\sphinxstyleliteralemphasis{\sphinxupquote{{]}}}) \textendash{} a dictionary with options for transformations.

\end{itemize}

\sphinxlineitem{Raises}
\sphinxAtStartPar
\sphinxstyleliteralstrong{\sphinxupquote{TransformationError}} \textendash{} if the node argument is not a
HaloExchange (or subclass thereof)

\end{description}\end{quote}

\end{fulllineitems}


\end{fulllineitems}



\begin{fulllineitems}

\pysigstartsignatures
\pysigline{\sphinxbfcode{\sphinxupquote{class\DUrole{w,w}{  }}}\sphinxcode{\sphinxupquote{psyclone.transformations.}}\sphinxbfcode{\sphinxupquote{Dynamo0p3ColourTrans}}}
\pysigstopsignatures
\sphinxAtStartPar
Split a Dynamo 0.3 loop over cells into colours so that it can be
parallelised. For example:

\begin{sphinxVerbatim}[commandchars=\\\{\}]
\PYG{g+gp}{\PYGZgt{}\PYGZgt{}\PYGZgt{} }\PYG{k+kn}{from}\PYG{+w}{ }\PYG{n+nn}{psyclone}\PYG{n+nn}{.}\PYG{n+nn}{parse}\PYG{n+nn}{.}\PYG{n+nn}{algorithm}\PYG{+w}{ }\PYG{k+kn}{import} \PYG{n}{parse}
\PYG{g+gp}{\PYGZgt{}\PYGZgt{}\PYGZgt{} }\PYG{k+kn}{from}\PYG{+w}{ }\PYG{n+nn}{psyclone}\PYG{n+nn}{.}\PYG{n+nn}{psyGen}\PYG{+w}{ }\PYG{k+kn}{import} \PYG{n}{PSyFactory}
\PYG{g+gp}{\PYGZgt{}\PYGZgt{}\PYGZgt{} }\PYG{k+kn}{import}\PYG{+w}{ }\PYG{n+nn}{transformations}
\PYG{g+gp}{\PYGZgt{}\PYGZgt{}\PYGZgt{} }\PYG{k+kn}{import}\PYG{+w}{ }\PYG{n+nn}{os}
\PYG{g+gp}{\PYGZgt{}\PYGZgt{}\PYGZgt{} }\PYG{k+kn}{import}\PYG{+w}{ }\PYG{n+nn}{pytest}
\PYG{g+gp}{\PYGZgt{}\PYGZgt{}\PYGZgt{}}
\PYG{g+gp}{\PYGZgt{}\PYGZgt{}\PYGZgt{} }\PYG{n}{TEST\PYGZus{}API} \PYG{o}{=} \PYG{l+s+s2}{\PYGZdq{}}\PYG{l+s+s2}{lfric}\PYG{l+s+s2}{\PYGZdq{}}
\PYG{g+gp}{\PYGZgt{}\PYGZgt{}\PYGZgt{} }\PYG{n}{\PYGZus{}}\PYG{p}{,}\PYG{n}{info}\PYG{o}{=}\PYG{n}{parse}\PYG{p}{(}\PYG{n}{os}\PYG{o}{.}\PYG{n}{path}\PYG{o}{.}\PYG{n}{join}\PYG{p}{(}\PYG{n}{os}\PYG{o}{.}\PYG{n}{path}\PYG{o}{.}\PYG{n}{dirname}\PYG{p}{(}\PYG{n}{os}\PYG{o}{.}\PYG{n}{path}\PYG{o}{.}\PYG{n}{abspath}\PYG{p}{(}\PYG{n+nv+vm}{\PYGZus{}\PYGZus{}file\PYGZus{}\PYGZus{}}\PYG{p}{)}\PYG{p}{)}\PYG{p}{,}
\PYG{g+gp}{\PYGZgt{}\PYGZgt{}\PYGZgt{} }             \PYG{l+s+s2}{\PYGZdq{}}\PYG{l+s+s2}{tests}\PYG{l+s+s2}{\PYGZdq{}}\PYG{p}{,} \PYG{l+s+s2}{\PYGZdq{}}\PYG{l+s+s2}{test\PYGZus{}files}\PYG{l+s+s2}{\PYGZdq{}}\PYG{p}{,} \PYG{l+s+s2}{\PYGZdq{}}\PYG{l+s+s2}{dynamo0p3}\PYG{l+s+s2}{\PYGZdq{}}\PYG{p}{,}
\PYG{g+gp}{\PYGZgt{}\PYGZgt{}\PYGZgt{} }             \PYG{l+s+s2}{\PYGZdq{}}\PYG{l+s+s2}{4.6\PYGZus{}multikernel\PYGZus{}invokes.f90}\PYG{l+s+s2}{\PYGZdq{}}\PYG{p}{)}\PYG{p}{,}
\PYG{g+gp}{\PYGZgt{}\PYGZgt{}\PYGZgt{} }             \PYG{n}{api}\PYG{o}{=}\PYG{n}{TEST\PYGZus{}API}\PYG{p}{)}
\PYG{g+gp}{\PYGZgt{}\PYGZgt{}\PYGZgt{} }\PYG{n}{psy} \PYG{o}{=} \PYG{n}{PSyFactory}\PYG{p}{(}\PYG{n}{TEST\PYGZus{}API}\PYG{p}{)}\PYG{o}{.}\PYG{n}{create}\PYG{p}{(}\PYG{n}{info}\PYG{p}{)}
\PYG{g+gp}{\PYGZgt{}\PYGZgt{}\PYGZgt{} }\PYG{n}{invoke} \PYG{o}{=} \PYG{n}{psy}\PYG{o}{.}\PYG{n}{invokes}\PYG{o}{.}\PYG{n}{get}\PYG{p}{(}\PYG{l+s+s1}{\PYGZsq{}}\PYG{l+s+s1}{invoke\PYGZus{}0}\PYG{l+s+s1}{\PYGZsq{}}\PYG{p}{)}
\PYG{g+gp}{\PYGZgt{}\PYGZgt{}\PYGZgt{} }\PYG{n}{schedule} \PYG{o}{=} \PYG{n}{invoke}\PYG{o}{.}\PYG{n}{schedule}
\PYG{g+gp}{\PYGZgt{}\PYGZgt{}\PYGZgt{}}
\PYG{g+gp}{\PYGZgt{}\PYGZgt{}\PYGZgt{} }\PYG{n}{ctrans} \PYG{o}{=} \PYG{n}{Dynamo0p3ColourTrans}\PYG{p}{(}\PYG{p}{)}
\PYG{g+gp}{\PYGZgt{}\PYGZgt{}\PYGZgt{} }\PYG{n}{otrans} \PYG{o}{=} \PYG{n}{DynamoOMPParallelLoopTrans}\PYG{p}{(}\PYG{p}{)}
\PYG{g+gp}{\PYGZgt{}\PYGZgt{}\PYGZgt{}}
\PYG{g+gp}{\PYGZgt{}\PYGZgt{}\PYGZgt{} }\PYG{c+c1}{\PYGZsh{} Colour all of the loops}
\PYG{g+gp}{\PYGZgt{}\PYGZgt{}\PYGZgt{} }\PYG{k}{for} \PYG{n}{child} \PYG{o+ow}{in} \PYG{n}{schedule}\PYG{o}{.}\PYG{n}{children}\PYG{p}{:}
\PYG{g+gp}{\PYGZgt{}\PYGZgt{}\PYGZgt{} }    \PYG{n}{ctrans}\PYG{o}{.}\PYG{n}{apply}\PYG{p}{(}\PYG{n}{child}\PYG{p}{)}
\PYG{g+gp}{\PYGZgt{}\PYGZgt{}\PYGZgt{}}
\PYG{g+gp}{\PYGZgt{}\PYGZgt{}\PYGZgt{} }\PYG{c+c1}{\PYGZsh{} Then apply OpenMP to each of the colour loops}
\PYG{g+gp}{\PYGZgt{}\PYGZgt{}\PYGZgt{} }\PYG{k}{for} \PYG{n}{child} \PYG{o+ow}{in} \PYG{n}{schedule}\PYG{o}{.}\PYG{n}{children}\PYG{p}{:}
\PYG{g+gp}{\PYGZgt{}\PYGZgt{}\PYGZgt{} }    \PYG{n}{otrans}\PYG{o}{.}\PYG{n}{apply}\PYG{p}{(}\PYG{n}{child}\PYG{o}{.}\PYG{n}{children}\PYG{p}{[}\PYG{l+m+mi}{0}\PYG{p}{]}\PYG{p}{)}
\PYG{g+gp}{\PYGZgt{}\PYGZgt{}\PYGZgt{}}
\PYG{g+gp}{\PYGZgt{}\PYGZgt{}\PYGZgt{} }\PYG{c+c1}{\PYGZsh{} Uncomment the following line to see a text view of the schedule}
\PYG{g+gp}{\PYGZgt{}\PYGZgt{}\PYGZgt{} }\PYG{c+c1}{\PYGZsh{} print(schedule.view())}
\end{sphinxVerbatim}

\sphinxAtStartPar
Colouring in the LFRic (Dynamo 0.3) API is subject to the following rules:
\begin{itemize}
\item {} 
\sphinxAtStartPar
Only kernels which operate on ‘CELL\_COLUMN’s and which increment a
field on a continuous function space require colouring. Kernels that
update a field on a discontinuous function space will cause this
transformation to raise an exception. Kernels that only write to a field
on a continuous function space also do not require colouring but are
permitted.

\item {} 
\sphinxAtStartPar
A kernel may have at most one field with ‘GH\_INC’ access.

\item {} 
\sphinxAtStartPar
A separate colour map will be required for each field that is coloured
(if an invoke contains \textgreater{}1 kernel call).

\end{itemize}


\begin{fulllineitems}

\pysigstartsignatures
\pysiglinewithargsret{\sphinxbfcode{\sphinxupquote{apply}}}{\sphinxparam{\DUrole{n,n}{node}}\sphinxparamcomma \sphinxparam{\DUrole{n,n}{options}\DUrole{o,o}{=}\DUrole{default_value}{None}}}{}
\pysigstopsignatures
\sphinxAtStartPar
Performs LFRic\sphinxhyphen{}specific error checking and then uses the parent
class to convert the Loop represented by \sphinxcode{\sphinxupquote{node}} into a
nested loop where the outer loop is over colours and the inner
loop is over cells of that colour.
\begin{quote}\begin{description}
\sphinxlineitem{Parameters}\begin{itemize}
\item {} 
\sphinxAtStartPar
\sphinxstyleliteralstrong{\sphinxupquote{node}} (\sphinxcode{\sphinxupquote{psyclone.domain.lfric.LFRicLoop}}) \textendash{} the loop to transform.

\item {} 
\sphinxAtStartPar
\sphinxstyleliteralstrong{\sphinxupquote{options}} \textendash{} a dictionary with options for transformations.        :type options: Optional{[}Dict{[}str, Any{]}{]}

\end{itemize}

\end{description}\end{quote}

\end{fulllineitems}


\end{fulllineitems}



\begin{fulllineitems}

\pysigstartsignatures
\pysigline{\sphinxbfcode{\sphinxupquote{class\DUrole{w,w}{  }}}\sphinxcode{\sphinxupquote{psyclone.transformations.}}\sphinxbfcode{\sphinxupquote{Dynamo0p3KernelConstTrans}}}
\pysigstopsignatures
\sphinxAtStartPar
Modifies a kernel so that the number of dofs, number of layers and
number of quadrature points are fixed in the kernel rather than
being passed in by argument.

\begin{sphinxVerbatim}[commandchars=\\\{\}]
\PYG{g+gp}{\PYGZgt{}\PYGZgt{}\PYGZgt{} }\PYG{k+kn}{from}\PYG{+w}{ }\PYG{n+nn}{psyclone}\PYG{n+nn}{.}\PYG{n+nn}{parse}\PYG{n+nn}{.}\PYG{n+nn}{algorithm}\PYG{+w}{ }\PYG{k+kn}{import} \PYG{n}{parse}
\PYG{g+gp}{\PYGZgt{}\PYGZgt{}\PYGZgt{} }\PYG{k+kn}{from}\PYG{+w}{ }\PYG{n+nn}{psyclone}\PYG{n+nn}{.}\PYG{n+nn}{psyGen}\PYG{+w}{ }\PYG{k+kn}{import} \PYG{n}{PSyFactory}
\PYG{g+gp}{\PYGZgt{}\PYGZgt{}\PYGZgt{} }\PYG{n}{api} \PYG{o}{=} \PYG{l+s+s2}{\PYGZdq{}}\PYG{l+s+s2}{lfric}\PYG{l+s+s2}{\PYGZdq{}}
\PYG{g+gp}{\PYGZgt{}\PYGZgt{}\PYGZgt{} }\PYG{n}{ast}\PYG{p}{,} \PYG{n}{invokeInfo} \PYG{o}{=} \PYG{n}{parse}\PYG{p}{(}\PYG{l+s+s2}{\PYGZdq{}}\PYG{l+s+s2}{file.f90}\PYG{l+s+s2}{\PYGZdq{}}\PYG{p}{,} \PYG{n}{api}\PYG{o}{=}\PYG{n}{api}\PYG{p}{)}
\PYG{g+gp}{\PYGZgt{}\PYGZgt{}\PYGZgt{} }\PYG{n}{psy}\PYG{o}{=}\PYG{n}{PSyFactory}\PYG{p}{(}\PYG{n}{api}\PYG{p}{)}\PYG{o}{.}\PYG{n}{create}\PYG{p}{(}\PYG{n}{invokeInfo}\PYG{p}{)}
\PYG{g+gp}{\PYGZgt{}\PYGZgt{}\PYGZgt{} }\PYG{n}{schedule} \PYG{o}{=} \PYG{n}{psy}\PYG{o}{.}\PYG{n}{invokes}\PYG{o}{.}\PYG{n}{get}\PYG{p}{(}\PYG{l+s+s1}{\PYGZsq{}}\PYG{l+s+s1}{invoke\PYGZus{}0}\PYG{l+s+s1}{\PYGZsq{}}\PYG{p}{)}\PYG{o}{.}\PYG{n}{schedule}
\PYG{g+gp}{\PYGZgt{}\PYGZgt{}\PYGZgt{} }\PYG{c+c1}{\PYGZsh{} Uncomment the following line to see a text view of the schedule}
\PYG{g+gp}{\PYGZgt{}\PYGZgt{}\PYGZgt{} }\PYG{c+c1}{\PYGZsh{} print(schedule.view())}
\PYG{g+gp}{\PYGZgt{}\PYGZgt{}\PYGZgt{}}
\PYG{g+gp}{\PYGZgt{}\PYGZgt{}\PYGZgt{} }\PYG{k+kn}{from}\PYG{+w}{ }\PYG{n+nn}{psyclone}\PYG{n+nn}{.}\PYG{n+nn}{transformations}\PYG{+w}{ }\PYG{k+kn}{import} \PYG{n}{Dynamo0p3KernelConstTrans}
\PYG{g+gp}{\PYGZgt{}\PYGZgt{}\PYGZgt{} }\PYG{n}{trans} \PYG{o}{=} \PYG{n}{Dynamo0p3KernelConstTrans}\PYG{p}{(}\PYG{p}{)}
\PYG{g+gp}{\PYGZgt{}\PYGZgt{}\PYGZgt{} }\PYG{k}{for} \PYG{n}{kernel} \PYG{o+ow}{in} \PYG{n}{schedule}\PYG{o}{.}\PYG{n}{coded\PYGZus{}kernels}\PYG{p}{(}\PYG{p}{)}\PYG{p}{:}
\PYG{g+gp}{\PYGZgt{}\PYGZgt{}\PYGZgt{} }    \PYG{n}{trans}\PYG{o}{.}\PYG{n}{apply}\PYG{p}{(}\PYG{n}{kernel}\PYG{p}{,} \PYG{n}{number\PYGZus{}of\PYGZus{}layers}\PYG{o}{=}\PYG{l+m+mi}{150}\PYG{p}{)}
\PYG{g+gp}{\PYGZgt{}\PYGZgt{}\PYGZgt{} }    \PYG{n}{kernel\PYGZus{}schedule} \PYG{o}{=} \PYG{n}{kernel}\PYG{o}{.}\PYG{n}{get\PYGZus{}kernel\PYGZus{}schedule}\PYG{p}{(}\PYG{p}{)}
\PYG{g+gp}{\PYGZgt{}\PYGZgt{}\PYGZgt{} }    \PYG{c+c1}{\PYGZsh{} Uncomment the following line to see a text view of the}
\PYG{g+gp}{\PYGZgt{}\PYGZgt{}\PYGZgt{} }    \PYG{c+c1}{\PYGZsh{} symbol table}
\PYG{g+gp}{\PYGZgt{}\PYGZgt{}\PYGZgt{} }    \PYG{c+c1}{\PYGZsh{} print(kernel\PYGZus{}schedule.symbol\PYGZus{}table.view())}
\end{sphinxVerbatim}


\begin{fulllineitems}

\pysigstartsignatures
\pysiglinewithargsret{\sphinxbfcode{\sphinxupquote{apply}}}{\sphinxparam{\DUrole{n,n}{node}}\sphinxparamcomma \sphinxparam{\DUrole{n,n}{options}\DUrole{o,o}{=}\DUrole{default_value}{None}}}{}
\pysigstopsignatures
\sphinxAtStartPar
Transforms a kernel so that the values for the number of degrees of
freedom (if valid values for the element\_order\_h and element\_order\_v
args are provided), the number of quadrature points (if the quadrature
arg is set to True) and the number of layers (if a valid value
for the number\_of\_layers arg is provided) are constant in a
kernel rather than being passed in by argument.

\sphinxAtStartPar
The “cellshape”, “element\_order\_h”, “element\_order\_v” and
“number\_of\_layers” arguments are provided to mirror the namelist values
that are input into an LFRic model when it is run.

\sphinxAtStartPar
Quadrature support is currently limited to XYoZ in ths
transformation. In the case of XYoZ the number of quadrature
points in the horizontal are set to element\_order\_h+3, and in the
vertical to element\_order\_v+3. These values are set in the LFRic
infrastructure, so their value is derived.
\begin{quote}\begin{description}
\sphinxlineitem{Parameters}\begin{itemize}
\item {} 
\sphinxAtStartPar
\sphinxstyleliteralstrong{\sphinxupquote{node}} (\sphinxcode{\sphinxupquote{psyclone.domain.lfric.LFRicKern}}) \textendash{} a kernel node.

\item {} 
\sphinxAtStartPar
\sphinxstyleliteralstrong{\sphinxupquote{options}} (\sphinxstyleliteralemphasis{\sphinxupquote{Optional}}\sphinxstyleliteralemphasis{\sphinxupquote{{[}}}\sphinxstyleliteralemphasis{\sphinxupquote{Dict}}\sphinxstyleliteralemphasis{\sphinxupquote{{[}}}\sphinxhref{https://docs.python.org/3/library/stdtypes.html\#str}{\sphinxstyleliteralemphasis{\sphinxupquote{str}}}\sphinxstyleliteralemphasis{\sphinxupquote{, }}\sphinxstyleliteralemphasis{\sphinxupquote{Any}}\sphinxstyleliteralemphasis{\sphinxupquote{{]}}}\sphinxstyleliteralemphasis{\sphinxupquote{{]}}}) \textendash{} a dictionary with options for transformations.

\item {} 
\sphinxAtStartPar
\sphinxstyleliteralstrong{\sphinxupquote{options}}\sphinxstyleliteralstrong{\sphinxupquote{{[}}}\sphinxstyleliteralstrong{\sphinxupquote{"cellshape"}}\sphinxstyleliteralstrong{\sphinxupquote{{]}}} (\sphinxhref{https://docs.python.org/3/library/stdtypes.html\#str}{\sphinxstyleliteralemphasis{\sphinxupquote{str}}}) \textendash{} the shape of the cells. This is
provided as it helps determine the number of dofs a field has
for a particular function space. Currently only “quadrilateral”
is supported which is also the default value.

\item {} 
\sphinxAtStartPar
\sphinxstyleliteralstrong{\sphinxupquote{options}}\sphinxstyleliteralstrong{\sphinxupquote{{[}}}\sphinxstyleliteralstrong{\sphinxupquote{"element\_order\_h"}}\sphinxstyleliteralstrong{\sphinxupquote{{]}}} (\sphinxhref{https://docs.python.org/3/library/functions.html\#int}{\sphinxstyleliteralemphasis{\sphinxupquote{int}}}) \textendash{} the polynomial order of the
cell in the horizontal. In combination with cellshape and
element\_order\_v, this determines the number of dofs a field has
for a particular function space. If it is set to None (the
default), then the dofs values are not set as constants in the
kernel, otherwise they are.

\item {} 
\sphinxAtStartPar
\sphinxstyleliteralstrong{\sphinxupquote{options}}\sphinxstyleliteralstrong{\sphinxupquote{{[}}}\sphinxstyleliteralstrong{\sphinxupquote{"element\_order\_v"}}\sphinxstyleliteralstrong{\sphinxupquote{{]}}} (\sphinxhref{https://docs.python.org/3/library/functions.html\#int}{\sphinxstyleliteralemphasis{\sphinxupquote{int}}}) \textendash{} the polynomial order of the
cell in the vertical. In combination with cellshape and
element\_order\_h, this determines the number of dofs a field has
for a particular function space. If it is set to None (the
default), then the dofs values are not set as constants in the
kernel, otherwise they are.

\item {} 
\sphinxAtStartPar
\sphinxstyleliteralstrong{\sphinxupquote{options}}\sphinxstyleliteralstrong{\sphinxupquote{{[}}}\sphinxstyleliteralstrong{\sphinxupquote{"number\_of\_layers"}}\sphinxstyleliteralstrong{\sphinxupquote{{]}}} (\sphinxhref{https://docs.python.org/3/library/functions.html\#int}{\sphinxstyleliteralemphasis{\sphinxupquote{int}}}) \textendash{} the number of vertical
layers in the LFRic model mesh used for this particular run. If
this is set to None (the default) then the nlayers value is not
set as a constant in the kernel, otherwise it is.

\item {} 
\sphinxAtStartPar
\sphinxstyleliteralstrong{\sphinxupquote{options}}\sphinxstyleliteralstrong{\sphinxupquote{{[}}}\sphinxstyleliteralstrong{\sphinxupquote{"quadrature"}}\sphinxstyleliteralstrong{\sphinxupquote{{]}}} (\sphinxhref{https://docs.python.org/3/library/functions.html\#bool}{\sphinxstyleliteralemphasis{\sphinxupquote{bool}}}) \textendash{} whether the number of quadrature
points values are set as constants in the kernel (True) or not
(False). The default is False.

\end{itemize}

\end{description}\end{quote}

\end{fulllineitems}



\begin{fulllineitems}

\pysigstartsignatures
\pysigline{\sphinxbfcode{\sphinxupquote{property\DUrole{w,w}{  }}}\sphinxbfcode{\sphinxupquote{name}}}
\pysigstopsignatures\begin{quote}\begin{description}
\sphinxlineitem{Returns}
\sphinxAtStartPar
the name of this transformation as a string.

\sphinxlineitem{Return type}
\sphinxAtStartPar
\sphinxhref{https://docs.python.org/3/library/stdtypes.html\#str}{str}

\end{description}\end{quote}

\end{fulllineitems}



\begin{fulllineitems}

\pysigstartsignatures
\pysiglinewithargsret{\sphinxbfcode{\sphinxupquote{validate}}}{\sphinxparam{\DUrole{n,n}{node}}\sphinxparamcomma \sphinxparam{\DUrole{n,n}{options}\DUrole{o,o}{=}\DUrole{default_value}{None}}}{}
\pysigstopsignatures
\sphinxAtStartPar
This method checks whether the input arguments are valid for
this transformation.
\begin{quote}\begin{description}
\sphinxlineitem{Parameters}\begin{itemize}
\item {} 
\sphinxAtStartPar
\sphinxstyleliteralstrong{\sphinxupquote{node}} (\sphinxcode{\sphinxupquote{psyclone.domain.lfric.LFRicKern}}) \textendash{} a dynamo 0.3 kernel node.

\item {} 
\sphinxAtStartPar
\sphinxstyleliteralstrong{\sphinxupquote{options}} (\sphinxstyleliteralemphasis{\sphinxupquote{Optional}}\sphinxstyleliteralemphasis{\sphinxupquote{{[}}}\sphinxstyleliteralemphasis{\sphinxupquote{Dict}}\sphinxstyleliteralemphasis{\sphinxupquote{{[}}}\sphinxhref{https://docs.python.org/3/library/stdtypes.html\#str}{\sphinxstyleliteralemphasis{\sphinxupquote{str}}}\sphinxstyleliteralemphasis{\sphinxupquote{, }}\sphinxstyleliteralemphasis{\sphinxupquote{Any}}\sphinxstyleliteralemphasis{\sphinxupquote{{]}}}\sphinxstyleliteralemphasis{\sphinxupquote{{]}}}) \textendash{} a dictionary with options for transformations.

\item {} 
\sphinxAtStartPar
\sphinxstyleliteralstrong{\sphinxupquote{options}}\sphinxstyleliteralstrong{\sphinxupquote{{[}}}\sphinxstyleliteralstrong{\sphinxupquote{"cellshape"}}\sphinxstyleliteralstrong{\sphinxupquote{{]}}} (\sphinxhref{https://docs.python.org/3/library/stdtypes.html\#str}{\sphinxstyleliteralemphasis{\sphinxupquote{str}}}) \textendash{} the shape of the elements/cells.

\item {} 
\sphinxAtStartPar
\sphinxstyleliteralstrong{\sphinxupquote{options}}\sphinxstyleliteralstrong{\sphinxupquote{{[}}}\sphinxstyleliteralstrong{\sphinxupquote{"element\_order\_h"}}\sphinxstyleliteralstrong{\sphinxupquote{{]}}} (\sphinxhref{https://docs.python.org/3/library/functions.html\#int}{\sphinxstyleliteralemphasis{\sphinxupquote{int}}}) \textendash{} the horizontal order of the               elements/cells.

\item {} 
\sphinxAtStartPar
\sphinxstyleliteralstrong{\sphinxupquote{options}}\sphinxstyleliteralstrong{\sphinxupquote{{[}}}\sphinxstyleliteralstrong{\sphinxupquote{"element\_order\_v"}}\sphinxstyleliteralstrong{\sphinxupquote{{]}}} (\sphinxhref{https://docs.python.org/3/library/functions.html\#int}{\sphinxstyleliteralemphasis{\sphinxupquote{int}}}) \textendash{} the vertical order of the               elements/cells.

\item {} 
\sphinxAtStartPar
\sphinxstyleliteralstrong{\sphinxupquote{options}}\sphinxstyleliteralstrong{\sphinxupquote{{[}}}\sphinxstyleliteralstrong{\sphinxupquote{"number\_of\_layers"}}\sphinxstyleliteralstrong{\sphinxupquote{{]}}} (\sphinxhref{https://docs.python.org/3/library/functions.html\#int}{\sphinxstyleliteralemphasis{\sphinxupquote{int}}}) \textendash{} the number of layers to use.

\item {} 
\sphinxAtStartPar
\sphinxstyleliteralstrong{\sphinxupquote{options}}\sphinxstyleliteralstrong{\sphinxupquote{{[}}}\sphinxstyleliteralstrong{\sphinxupquote{"quadrature"}}\sphinxstyleliteralstrong{\sphinxupquote{{]}}} (\sphinxhref{https://docs.python.org/3/library/functions.html\#bool}{\sphinxstyleliteralemphasis{\sphinxupquote{bool}}}) \textendash{} whether quadrature dimension sizes             should or shouldn’t be set as constants in a kernel.

\end{itemize}

\sphinxlineitem{Raises}
\sphinxAtStartPar
\sphinxstyleliteralstrong{\sphinxupquote{TransformationError}} \textendash{} if the node argument is not a             dynamo 0.3 kernel, the cellshape argument is not set to             “quadrilateral”, the element\_order\_h or element\_order\_v arguments            are not a 0 or a positive integer, the number of layers argument            is not a positive integer, the quadrature argument is not a            boolean, neither element orders nor number of layers arguments are            set (as the transformation would then do nothing), or the             quadrature argument is True but the element order is not             provided (as the former needs the latter).

\end{description}\end{quote}

\end{fulllineitems}


\end{fulllineitems}



\begin{fulllineitems}

\pysigstartsignatures
\pysiglinewithargsret{\sphinxbfcode{\sphinxupquote{class\DUrole{w,w}{  }}}\sphinxcode{\sphinxupquote{psyclone.transformations.}}\sphinxbfcode{\sphinxupquote{Dynamo0p3OMPLoopTrans}}}{\sphinxparam{\DUrole{n,n}{omp\_schedule}\DUrole{o,o}{=}\DUrole{default_value}{\textquotesingle{}static\textquotesingle{}}}}{}
\pysigstopsignatures
\sphinxAtStartPar
LFRic (Dynamo 0.3) specific orphan OpenMP loop transformation. Adds
Dynamo\sphinxhyphen{}specific validity checks.
\begin{quote}\begin{description}
\sphinxlineitem{Parameters}
\sphinxAtStartPar
\sphinxstyleliteralstrong{\sphinxupquote{omp\_schedule}} (\sphinxhref{https://docs.python.org/3/library/stdtypes.html\#str}{\sphinxstyleliteralemphasis{\sphinxupquote{str}}}) \textendash{} the OpenMP schedule to use. Must be one of         ‘runtime’, ‘static’, ‘dynamic’, ‘guided’ or ‘auto’. Defaults to         ‘static’.

\end{description}\end{quote}


\begin{fulllineitems}

\pysigstartsignatures
\pysiglinewithargsret{\sphinxbfcode{\sphinxupquote{apply}}}{\sphinxparam{\DUrole{n,n}{node}}\sphinxparamcomma \sphinxparam{\DUrole{n,n}{options}\DUrole{o,o}{=}\DUrole{default_value}{None}}}{}
\pysigstopsignatures
\sphinxAtStartPar
Apply LFRic (Dynamo 0.3) specific OMPLoopTrans.
\begin{quote}\begin{description}
\sphinxlineitem{Parameters}\begin{itemize}
\item {} 
\sphinxAtStartPar
\sphinxstyleliteralstrong{\sphinxupquote{node}} (\sphinxcode{\sphinxupquote{psyclone.psyir.nodes.Node}}) \textendash{} the Node in the Schedule to check.

\item {} 
\sphinxAtStartPar
\sphinxstyleliteralstrong{\sphinxupquote{options}} (\sphinxstyleliteralemphasis{\sphinxupquote{Optional}}\sphinxstyleliteralemphasis{\sphinxupquote{{[}}}\sphinxhref{https://docs.python.org/3/library/stdtypes.html\#dict}{\sphinxstyleliteralemphasis{\sphinxupquote{dict}}}\sphinxstyleliteralemphasis{\sphinxupquote{{[}}}\sphinxhref{https://docs.python.org/3/library/stdtypes.html\#str}{\sphinxstyleliteralemphasis{\sphinxupquote{str}}}\sphinxstyleliteralemphasis{\sphinxupquote{, }}\sphinxstyleliteralemphasis{\sphinxupquote{Any}}\sphinxstyleliteralemphasis{\sphinxupquote{{]}}}\sphinxstyleliteralemphasis{\sphinxupquote{{]}}}) \textendash{} a dictionary with options for transformations
and validation.

\item {} 
\sphinxAtStartPar
\sphinxstyleliteralstrong{\sphinxupquote{options}}\sphinxstyleliteralstrong{\sphinxupquote{{[}}}\sphinxstyleliteralstrong{\sphinxupquote{"reprod"}}\sphinxstyleliteralstrong{\sphinxupquote{{]}}} (\sphinxhref{https://docs.python.org/3/library/functions.html\#bool}{\sphinxstyleliteralemphasis{\sphinxupquote{bool}}}) \textendash{} indicating whether reproducible reductions should be used.
By default the value from the config file will be used.

\end{itemize}

\end{description}\end{quote}

\end{fulllineitems}



\begin{fulllineitems}

\pysigstartsignatures
\pysiglinewithargsret{\sphinxbfcode{\sphinxupquote{validate}}}{\sphinxparam{\DUrole{n,n}{node}}\sphinxparamcomma \sphinxparam{\DUrole{n,n}{options}\DUrole{o,o}{=}\DUrole{default_value}{None}}}{}
\pysigstopsignatures
\sphinxAtStartPar
Perform LFRic (Dynamo 0.3) specific loop validity checks for the
OMPLoopTrans.
\begin{quote}\begin{description}
\sphinxlineitem{Parameters}\begin{itemize}
\item {} 
\sphinxAtStartPar
\sphinxstyleliteralstrong{\sphinxupquote{node}} (\sphinxcode{\sphinxupquote{psyclone.psyir.nodes.Node}}) \textendash{} the Node in the Schedule to check

\item {} 
\sphinxAtStartPar
\sphinxstyleliteralstrong{\sphinxupquote{options}} (\sphinxstyleliteralemphasis{\sphinxupquote{Optional}}\sphinxstyleliteralemphasis{\sphinxupquote{{[}}}\sphinxstyleliteralemphasis{\sphinxupquote{Dict}}\sphinxstyleliteralemphasis{\sphinxupquote{{[}}}\sphinxhref{https://docs.python.org/3/library/stdtypes.html\#str}{\sphinxstyleliteralemphasis{\sphinxupquote{str}}}\sphinxstyleliteralemphasis{\sphinxupquote{, }}\sphinxstyleliteralemphasis{\sphinxupquote{Any}}\sphinxstyleliteralemphasis{\sphinxupquote{{]}}}\sphinxstyleliteralemphasis{\sphinxupquote{{]}}}) \textendash{} a dictionary with options for transformations                         and validation.

\item {} 
\sphinxAtStartPar
\sphinxstyleliteralstrong{\sphinxupquote{options}}\sphinxstyleliteralstrong{\sphinxupquote{{[}}}\sphinxstyleliteralstrong{\sphinxupquote{"reprod"}}\sphinxstyleliteralstrong{\sphinxupquote{{]}}} (\sphinxhref{https://docs.python.org/3/library/functions.html\#bool}{\sphinxstyleliteralemphasis{\sphinxupquote{bool}}}) \textendash{} indicating whether reproducible reductions should be used.             By default the value from the config file will be used.

\end{itemize}

\sphinxlineitem{Raises}
\sphinxAtStartPar
\sphinxstyleliteralstrong{\sphinxupquote{TransformationError}} \textendash{} if an OMP loop transform would create             incorrect code.

\end{description}\end{quote}

\end{fulllineitems}


\end{fulllineitems}



\begin{fulllineitems}

\pysigstartsignatures
\pysigline{\sphinxbfcode{\sphinxupquote{class\DUrole{w,w}{  }}}\sphinxcode{\sphinxupquote{psyclone.transformations.}}\sphinxbfcode{\sphinxupquote{Dynamo0p3RedundantComputationTrans}}}
\pysigstopsignatures
\sphinxAtStartPar
This transformation allows the user to modify a loop’s bounds so
that redundant computation will be performed. Redundant
computation can result in halo exchanges being modified, new halo
exchanges being added or existing halo exchanges being removed.
\begin{itemize}
\item {} 
\sphinxAtStartPar
This transformation should be performed before any
parallelisation transformations (e.g. for OpenMP) to the loop in
question and will raise an exception if this is not the case.

\item {} 
\sphinxAtStartPar
This transformation can not be applied to a loop containing a
reduction and will again raise an exception if this is the case.

\item {} 
\sphinxAtStartPar
This transformation can only be used to add redundant
computation to a loop, not to remove it.

\item {} 
\sphinxAtStartPar
This transformation allows a loop that is already performing
redundant computation to be modified, but only if the depth is
increased.

\end{itemize}


\begin{fulllineitems}

\pysigstartsignatures
\pysiglinewithargsret{\sphinxbfcode{\sphinxupquote{apply}}}{\sphinxparam{\DUrole{n,n}{loop}}\sphinxparamcomma \sphinxparam{\DUrole{n,n}{options}\DUrole{o,o}{=}\DUrole{default_value}{None}}}{}
\pysigstopsignatures
\sphinxAtStartPar
Apply the redundant computation transformation to the loop
\sphinxcode{\sphinxupquote{loop}}. This transformation can be applied to loops iterating
over ‘cells or ‘dofs’. if \sphinxcode{\sphinxupquote{depth}} is set to a value then the
value will be the depth of the field’s halo over which redundant
computation will be performed. If \sphinxcode{\sphinxupquote{depth}} is not set to a
value then redundant computation will be performed to the full depth
of the field’s halo.
\begin{quote}\begin{description}
\sphinxlineitem{Parameters}\begin{itemize}
\item {} 
\sphinxAtStartPar
\sphinxstyleliteralstrong{\sphinxupquote{loop}} (\sphinxcode{\sphinxupquote{psyclone.psyGen.LFRicLoop}}) \textendash{} the loop that we are transforming.

\item {} 
\sphinxAtStartPar
\sphinxstyleliteralstrong{\sphinxupquote{options}} (\sphinxstyleliteralemphasis{\sphinxupquote{Optional}}\sphinxstyleliteralemphasis{\sphinxupquote{{[}}}\sphinxstyleliteralemphasis{\sphinxupquote{Dict}}\sphinxstyleliteralemphasis{\sphinxupquote{{[}}}\sphinxhref{https://docs.python.org/3/library/stdtypes.html\#str}{\sphinxstyleliteralemphasis{\sphinxupquote{str}}}\sphinxstyleliteralemphasis{\sphinxupquote{, }}\sphinxstyleliteralemphasis{\sphinxupquote{Any}}\sphinxstyleliteralemphasis{\sphinxupquote{{]}}}\sphinxstyleliteralemphasis{\sphinxupquote{{]}}}) \textendash{} a dictionary with options for transformations.

\item {} 
\sphinxAtStartPar
\sphinxstyleliteralstrong{\sphinxupquote{options}}\sphinxstyleliteralstrong{\sphinxupquote{{[}}}\sphinxstyleliteralstrong{\sphinxupquote{"depth"}}\sphinxstyleliteralstrong{\sphinxupquote{{]}}} (\sphinxhref{https://docs.python.org/3/library/functions.html\#int}{\sphinxstyleliteralemphasis{\sphinxupquote{int}}}) \textendash{} the depth of the stencil. Defaults                 to None.

\end{itemize}

\end{description}\end{quote}

\end{fulllineitems}



\begin{fulllineitems}

\pysigstartsignatures
\pysiglinewithargsret{\sphinxbfcode{\sphinxupquote{validate}}}{\sphinxparam{\DUrole{n,n}{node}}\sphinxparamcomma \sphinxparam{\DUrole{n,n}{options}\DUrole{o,o}{=}\DUrole{default_value}{None}}}{}
\pysigstopsignatures
\sphinxAtStartPar
Perform various checks to ensure that it is valid to apply the
RedundantComputation transformation to the supplied node
\begin{quote}\begin{description}
\sphinxlineitem{Parameters}\begin{itemize}
\item {} 
\sphinxAtStartPar
\sphinxstyleliteralstrong{\sphinxupquote{node}} (\sphinxcode{\sphinxupquote{psyclone.psyir.nodes.Node}}) \textendash{} the supplied node on which we are performing                     validity checks

\item {} 
\sphinxAtStartPar
\sphinxstyleliteralstrong{\sphinxupquote{options}} (\sphinxstyleliteralemphasis{\sphinxupquote{Optional}}\sphinxstyleliteralemphasis{\sphinxupquote{{[}}}\sphinxstyleliteralemphasis{\sphinxupquote{Dict}}\sphinxstyleliteralemphasis{\sphinxupquote{{[}}}\sphinxhref{https://docs.python.org/3/library/stdtypes.html\#str}{\sphinxstyleliteralemphasis{\sphinxupquote{str}}}\sphinxstyleliteralemphasis{\sphinxupquote{, }}\sphinxstyleliteralemphasis{\sphinxupquote{Any}}\sphinxstyleliteralemphasis{\sphinxupquote{{]}}}\sphinxstyleliteralemphasis{\sphinxupquote{{]}}}) \textendash{} a dictionary with options for transformations.

\item {} 
\sphinxAtStartPar
\sphinxstyleliteralstrong{\sphinxupquote{options}}\sphinxstyleliteralstrong{\sphinxupquote{{[}}}\sphinxstyleliteralstrong{\sphinxupquote{"depth"}}\sphinxstyleliteralstrong{\sphinxupquote{{]}}} (\sphinxhref{https://docs.python.org/3/library/functions.html\#int}{\sphinxstyleliteralemphasis{\sphinxupquote{int}}}) \textendash{} the depth of the stencil if the value                      is provided and None if not.

\end{itemize}

\sphinxlineitem{Raises}\begin{itemize}
\item {} 
\sphinxAtStartPar
\sphinxstyleliteralstrong{\sphinxupquote{TransformationError}} \textendash{} if the parent of the loop is a            \sphinxcode{\sphinxupquote{psyclone.psyir.nodes.Directive}}.

\item {} 
\sphinxAtStartPar
\sphinxstyleliteralstrong{\sphinxupquote{TransformationError}} \textendash{} if the parent of the loop is not a            \sphinxcode{\sphinxupquote{psyclone.psyir.nodes.Loop}} or a            \sphinxcode{\sphinxupquote{psyclone.psyGen.LFRicInvokeSchedule}}.

\item {} 
\sphinxAtStartPar
\sphinxstyleliteralstrong{\sphinxupquote{TransformationError}} \textendash{} if the parent of the loop is a            \sphinxcode{\sphinxupquote{psyclone.psyir.nodes.Loop}} but the original loop does            not iterate over ‘colour’.

\item {} 
\sphinxAtStartPar
\sphinxstyleliteralstrong{\sphinxupquote{TransformationError}} \textendash{} if the parent of the loop is a            \sphinxcode{\sphinxupquote{psyclone.psyir.nodes.Loop}} but the parent does not
iterate over ‘colours’.

\item {} 
\sphinxAtStartPar
\sphinxstyleliteralstrong{\sphinxupquote{TransformationError}} \textendash{} if the parent of the loop is a            \sphinxcode{\sphinxupquote{psyclone.psyir.nodes.Loop}} but the parent’s parent is            not a \sphinxcode{\sphinxupquote{psyclone.psyGen.LFRicInvokeSchedule}}.

\item {} 
\sphinxAtStartPar
\sphinxstyleliteralstrong{\sphinxupquote{TransformationError}} \textendash{} if this transformation is applied            when distributed memory is not switched on.

\item {} 
\sphinxAtStartPar
\sphinxstyleliteralstrong{\sphinxupquote{TransformationError}} \textendash{} if the loop does not iterate over            cells, dofs or colour.

\item {} 
\sphinxAtStartPar
\sphinxstyleliteralstrong{\sphinxupquote{TransformationError}} \textendash{} if the loop contains a kernel that
operates on halo cells.

\item {} 
\sphinxAtStartPar
\sphinxstyleliteralstrong{\sphinxupquote{TransformationError}} \textendash{} if the transformation is setting the            loop to the maximum halo depth but the loop already computes            to the maximum halo depth.

\item {} 
\sphinxAtStartPar
\sphinxstyleliteralstrong{\sphinxupquote{TransformationError}} \textendash{} if the transformation is setting the            loop to the maximum halo depth but the loop contains a stencil            access (as this would result in the field being accessed            beyond the halo depth).

\item {} 
\sphinxAtStartPar
\sphinxstyleliteralstrong{\sphinxupquote{TransformationError}} \textendash{} if the supplied depth value is not an            integer.

\item {} 
\sphinxAtStartPar
\sphinxstyleliteralstrong{\sphinxupquote{TransformationError}} \textendash{} if the supplied depth value is less            than 1.

\item {} 
\sphinxAtStartPar
\sphinxstyleliteralstrong{\sphinxupquote{TransformationError}} \textendash{} if the supplied depth value is not            greater than 1 when a continuous loop is modified as this is            the minimum valid value.

\item {} 
\sphinxAtStartPar
\sphinxstyleliteralstrong{\sphinxupquote{TransformationError}} \textendash{} if the supplied depth value is not            greater than the existing depth value, as we should not need            to undo existing transformations.

\item {} 
\sphinxAtStartPar
\sphinxstyleliteralstrong{\sphinxupquote{TransformationError}} \textendash{} if a depth value has been supplied            but the loop has already been set to the maximum halo depth.

\end{itemize}

\end{description}\end{quote}

\end{fulllineitems}


\end{fulllineitems}


\sphinxstepscope


\chapter{GOcean1.0 API}
\label{\detokenize{gocean1p0:gocean1-0-api}}\label{\detokenize{gocean1p0:gocean-api}}\label{\detokenize{gocean1p0::doc}}

\section{Introduction}
\label{\detokenize{gocean1p0:introduction}}\label{\detokenize{gocean1p0:gocean-intro}}
\sphinxAtStartPar
The GOcean 1.0 application programming interface (API) was originally
designed to support ocean models that use the finite\sphinxhyphen{}difference scheme
for two\sphinxhyphen{}dimensional domains.  However, the approach is not specific to
ocean models and can potentially be applied to any finite\sphinxhyphen{}difference
code.

\sphinxAtStartPar
As with all PSyclone APIs, the GOcean 1.0 API specifies how a user
must write the Algorithm Layer and the Kernel Layer to allow PSyclone
to generate the PSy Layer. These Algorithm and Kernel APIs are
discussed separately in the sections below. Before these we
describe the functionality provided by the GOcean Library.


\section{The GOcean Infrastructure Library \sphinxhyphen{} dl\_esm\_inf}
\label{\detokenize{gocean1p0:the-gocean-infrastructure-library-dl-esm-inf}}\label{\detokenize{gocean1p0:gocean-library}}
\sphinxAtStartPar
The use of PSyclone and the GOcean 1.0 API implies the use of a
standard set of data types and associated infrastructure. This is
provided by the GOcean infrastructure library \sphinxhyphen{} dl\_esm\_inf.
Currently this library is distributed separately from PSyclone and
is available from \sphinxurl{https://github.com/stfc/dl\_esm\_inf}.


\subsection{Grid}
\label{\detokenize{gocean1p0:grid}}\label{\detokenize{gocean1p0:gocean-grid}}
\sphinxAtStartPar
The dl\_esm\_inf library contains a \sphinxcode{\sphinxupquote{grid\_mod}} module which defines a
\sphinxcode{\sphinxupquote{grid\_type}} and associated constructor:

\begin{sphinxVerbatim}[commandchars=\\\{\}]
\PYG{k}{use }\PYG{n}{grid\PYGZus{}mod}
\PYG{p}{.}\PYG{p}{.}\PYG{p}{.}
\PYG{c}{!\PYGZgt{} The grid on which our fields are defined}
\PYG{k}{type}\PYG{p}{(}\PYG{n}{grid\PYGZus{}type}\PYG{p}{)}\PYG{p}{,}\PYG{+w}{ }\PYG{k}{target}\PYG{+w}{ }\PYG{k+kd}{::}\PYG{+w}{ }\PYG{n}{model\PYGZus{}grid}
\PYG{p}{.}\PYG{p}{.}\PYG{p}{.}
\PYG{c}{! Create the model grid}
\PYG{n}{model\PYGZus{}grid}\PYG{+w}{ }\PYG{o}{=}\PYG{+w}{ }\PYG{n}{grid\PYGZus{}type}\PYG{p}{(}\PYG{n}{GO\PYGZus{}ARAKAWA\PYGZus{}C}\PYG{p}{,}\PYG{+w}{                                 }\PYG{p}{\PYGZam{}}
\PYG{+w}{                       }\PYG{p}{(}\PYG{o}{/}\PYG{n}{GO\PYGZus{}BC\PYGZus{}EXTERNAL}\PYG{p}{,}\PYG{n}{GO\PYGZus{}BC\PYGZus{}EXTERNAL}\PYG{p}{,}\PYG{n}{GO\PYGZus{}BC\PYGZus{}NONE}\PYG{o}{/}\PYG{p}{)}\PYG{p}{,}\PYG{+w}{ }\PYG{p}{\PYGZam{}}
\PYG{+w}{                       }\PYG{n}{GO\PYGZus{}OFFSET\PYGZus{}NE}\PYG{p}{)}
\end{sphinxVerbatim}

\begin{sphinxadmonition}{note}{Note:}
\sphinxAtStartPar
The grid object itself must be declared with the \sphinxcode{\sphinxupquote{target}}
attribute. This is because each field object will contain a pointer to
it.
\end{sphinxadmonition}

\sphinxAtStartPar
The \sphinxcode{\sphinxupquote{grid\_type}} constructor takes three arguments:
\begin{enumerate}
\sphinxsetlistlabels{\arabic}{enumi}{enumii}{}{.}%
\item {} 
\sphinxAtStartPar
The type of grid (only GO\_ARAKAWA\_C is currently supported)

\item {} 
\sphinxAtStartPar
The boundary conditions on the domain for the \sphinxstyleemphasis{x}, \sphinxstyleemphasis{y} and \sphinxstyleemphasis{z} dimensions (see below). The value for the \sphinxstyleemphasis{z} dimension is currently ignored.

\item {} 
\sphinxAtStartPar
The ‘index offset’ \sphinxhyphen{} the convention used for indexing into offset fields.

\end{enumerate}

\sphinxAtStartPar
Three types of boundary condition are currently supported:


\begin{savenotes}\sphinxattablestart
\sphinxthistablewithglobalstyle
\sphinxthistablewithvlinesstyle
\centering
\begin{tabulary}{\linewidth}[t]{|l|L|}
\sphinxtoprule
\sphinxstyletheadfamily 
\sphinxAtStartPar
Name
&\sphinxstyletheadfamily 
\sphinxAtStartPar
Description
\\
\sphinxmidrule
\sphinxtableatstartofbodyhook
\sphinxAtStartPar
GO\_BC\_NONE
&
\sphinxAtStartPar
No boundary conditions are applied.
\\
\sphinxhline
\sphinxAtStartPar
GO\_BC\_EXTERNAL
&
\sphinxAtStartPar
Some external forcing is applied. This must be implemented by a kernel. The domain must be defined with a T\sphinxhyphen{}point mask (see {\hyperref[\detokenize{gocean1p0:gocean-grid-init}]{\sphinxcrossref{\DUrole{std,std-ref}{The grid\_init Routine}}}}).
\\
\sphinxhline
\sphinxAtStartPar
GO\_BC\_PERIODIC
&
\sphinxAtStartPar
Periodic boundary conditions are applied.
\\
\sphinxbottomrule
\end{tabulary}
\sphinxtableafterendhook\par
\sphinxattableend\end{savenotes}

\sphinxAtStartPar
The infrastructure requires this information in order to determine the
extent of the model grid.

\sphinxAtStartPar
The index offset is required because a model (kernel) developer has
choice in how they actually implement the staggering of variables on a
grid. This comes down to a choice of which grid points in the vicinity
of a given T point have the same array (\sphinxstyleemphasis{i}, \sphinxstyleemphasis{j}) indices. In
the diagram below, the image on the left corresponds to choosing those
points to the South and West of a T point to have the same (\sphinxstyleemphasis{i}, \sphinxstyleemphasis{j})
index. That on the right corresponds to choosing those points to the
North and East of the T point (this is the offset scheme used in the
NEMO ocean model):

\noindent\sphinxincludegraphics{{grid_offset_choices}.png}

\sphinxAtStartPar
The GOcean 1.0 API supports these two different offset schemes, which
we term \sphinxcode{\sphinxupquote{GO\_OFFSET\_SW}} and \sphinxcode{\sphinxupquote{GO\_OFFSET\_NE}}.

\sphinxAtStartPar
Note that the constructor does not specify the extent of the model
grid. This is because this information is normally obtained by reading
a file (a namelist file, a netcdf file etc.) which is specific to an
application.  Once this information has been obtained, a second
routine, \sphinxcode{\sphinxupquote{grid\_init}}, is provided with which to ‘load’ a grid object
with state. This is discussed below.


\subsubsection{The \sphinxstyleliteralintitle{\sphinxupquote{grid\_init}} Routine}
\label{\detokenize{gocean1p0:the-grid-init-routine}}\label{\detokenize{gocean1p0:gocean-grid-init}}
\sphinxAtStartPar
Once an application has determined the details of the model
configuration, it must use this information to populate the grid
object. This is done via a call to the \sphinxcode{\sphinxupquote{grid\_init}} subroutine:

\begin{sphinxVerbatim}[commandchars=\\\{\}]
\PYG{k}{subroutine }\PYG{n}{grid\PYGZus{}init}\PYG{p}{(}\PYG{n}{grid}\PYG{p}{,}\PYG{+w}{ }\PYG{n}{m}\PYG{p}{,}\PYG{+w}{ }\PYG{n}{n}\PYG{p}{,}\PYG{+w}{ }\PYG{n}{dxarg}\PYG{p}{,}\PYG{+w}{ }\PYG{n}{dyarg}\PYG{p}{,}\PYG{+w}{ }\PYG{n}{tmask}\PYG{p}{)}
\PYG{+w}{  }\PYG{c}{!\PYGZgt{} The grid object to configure}
\PYG{+w}{  }\PYG{k}{type}\PYG{p}{(}\PYG{n}{grid\PYGZus{}type}\PYG{p}{)}\PYG{p}{,}\PYG{+w}{ }\PYG{k}{intent}\PYG{p}{(}\PYG{n}{inout}\PYG{p}{)}\PYG{+w}{ }\PYG{k+kd}{::}\PYG{+w}{ }\PYG{n}{grid}
\PYG{+w}{  }\PYG{c}{!\PYGZgt{} Dimensions of the model grid}
\PYG{+w}{  }\PYG{k+kt}{integer}\PYG{p}{,}\PYG{+w}{         }\PYG{k}{intent}\PYG{p}{(}\PYG{n}{in}\PYG{p}{)}\PYG{+w}{    }\PYG{k+kd}{::}\PYG{+w}{ }\PYG{n}{m}\PYG{p}{,}\PYG{+w}{ }\PYG{n}{n}
\PYG{+w}{  }\PYG{c}{!\PYGZgt{} The (constant) grid spacing in x and y (m)}
\PYG{+w}{  }\PYG{k+kt}{real}\PYG{p}{(}\PYG{n}{wp}\PYG{p}{)}\PYG{p}{,}\PYG{+w}{        }\PYG{k}{intent}\PYG{p}{(}\PYG{n}{in}\PYG{p}{)}\PYG{+w}{    }\PYG{k+kd}{::}\PYG{+w}{ }\PYG{n}{dxarg}\PYG{p}{,}\PYG{+w}{ }\PYG{n}{dyarg}
\PYG{+w}{  }\PYG{c}{!\PYGZgt{} Optional T\PYGZhy{}point mask specifying whether each grid point is}
\PYG{+w}{  }\PYG{c}{!! wet (1), dry (0) or external (\PYGZhy{}1).}
\PYG{+w}{  }\PYG{k+kt}{integer}\PYG{p}{,}\PYG{+w}{ }\PYG{k}{dimension}\PYG{p}{(}\PYG{n}{m}\PYG{p}{,}\PYG{n}{n}\PYG{p}{)}\PYG{p}{,}\PYG{+w}{ }\PYG{k}{intent}\PYG{p}{(}\PYG{n}{in}\PYG{p}{)}\PYG{p}{,}\PYG{+w}{ }\PYG{k}{optional}\PYG{+w}{ }\PYG{k+kd}{::}\PYG{+w}{ }\PYG{n}{tmask}
\end{sphinxVerbatim}

\sphinxAtStartPar
If no T\sphinxhyphen{}mask is supplied then this routine configures the grid
appropriately for an all\sphinxhyphen{}wet domain with periodic boundary conditions
in both the \sphinxstyleemphasis{x}\sphinxhyphen{} and \sphinxstyleemphasis{y}\sphinxhyphen{}dimensions. It should also be noted that
currently only grids with constant resolution in \sphinxstyleemphasis{x} and \sphinxstyleemphasis{y} are
supported by this routine.


\subsection{Fields}
\label{\detokenize{gocean1p0:fields}}\label{\detokenize{gocean1p0:gocean-fields}}
\sphinxAtStartPar
Once a model has a grid defined it will require one or more fields.
The dl\_esm\_inf library contains a \sphinxcode{\sphinxupquote{field\_mod}} module which defines an
\sphinxcode{\sphinxupquote{r2d\_field}} type (real, 2\sphinxhyphen{}dimensional field) and associated
constructor:

\begin{sphinxVerbatim}[commandchars=\\\{\}]
\PYG{k}{use }\PYG{n}{field\PYGZus{}mod}
\PYG{p}{.}\PYG{p}{.}\PYG{p}{.}
\PYG{c}{!\PYGZgt{} Current (\PYGZsq{}now\PYGZsq{}) sea\PYGZhy{}surface height at different grid points}
\PYG{k}{type}\PYG{p}{(}\PYG{n}{r2d\PYGZus{}field}\PYG{p}{)}\PYG{+w}{ }\PYG{k+kd}{::}\PYG{+w}{ }\PYG{n}{sshn\PYGZus{}u\PYGZus{}fld}\PYG{p}{,}\PYG{+w}{ }\PYG{n}{sshn\PYGZus{}v\PYGZus{}fld}\PYG{p}{,}\PYG{+w}{ }\PYG{n}{sshn\PYGZus{}t\PYGZus{}fld}
\PYG{p}{.}\PYG{p}{.}\PYG{p}{.}

\PYG{c}{! Sea\PYGZhy{}surface height now (current time step)}
\PYG{n}{sshn\PYGZus{}u}\PYG{+w}{ }\PYG{o}{=}\PYG{+w}{ }\PYG{n}{r2d\PYGZus{}field}\PYG{p}{(}\PYG{n}{model\PYGZus{}grid}\PYG{p}{,}\PYG{+w}{ }\PYG{n}{GO\PYGZus{}U\PYGZus{}POINTS}\PYG{p}{)}
\PYG{n}{sshn\PYGZus{}v}\PYG{+w}{ }\PYG{o}{=}\PYG{+w}{ }\PYG{n}{r2d\PYGZus{}field}\PYG{p}{(}\PYG{n}{model\PYGZus{}grid}\PYG{p}{,}\PYG{+w}{ }\PYG{n}{GO\PYGZus{}V\PYGZus{}POINTS}\PYG{p}{)}
\PYG{n}{sshn\PYGZus{}t}\PYG{+w}{ }\PYG{o}{=}\PYG{+w}{ }\PYG{n}{r2d\PYGZus{}field}\PYG{p}{(}\PYG{n}{model\PYGZus{}grid}\PYG{p}{,}\PYG{+w}{ }\PYG{n}{GO\PYGZus{}T\PYGZus{}POINTS}\PYG{p}{)}
\end{sphinxVerbatim}

\sphinxAtStartPar
The constructor takes two arguments:
\begin{enumerate}
\sphinxsetlistlabels{\arabic}{enumi}{enumii}{}{.}%
\item {} 
\sphinxAtStartPar
The grid on which the field exists

\item {} 
\sphinxAtStartPar
The type of grid point at which the field is defined
(\sphinxcode{\sphinxupquote{GO\_U\_POINTS}}, \sphinxcode{\sphinxupquote{GO\_V\_POINTS}}, \sphinxcode{\sphinxupquote{GO\_T\_POINTS}} or \sphinxcode{\sphinxupquote{GO\_F\_POINTS}})

\end{enumerate}

\sphinxAtStartPar
Note that the grid object need not have been fully configured (by a
call to \sphinxcode{\sphinxupquote{grid\_init}} for instance) before it is passed into this
constructor.


\subsection{Example}
\label{\detokenize{gocean1p0:example}}
\sphinxAtStartPar
PSyclone is distributed with a full example of the use of the
GOcean Library. See \sphinxcode{\sphinxupquote{\textless{}PSYCLONEHOME\textgreater{}/examples/gocean/shallow\_alg.f90}}.  In what
follows we will walk through a slightly cut\sphinxhyphen{}down example for a
different program.

\sphinxAtStartPar
The following code illustrates the use of dl\_esm\_inf for constructing an
application:

\begin{sphinxVerbatim}[commandchars=\\\{\}]
\PYG{k}{program }\PYG{n}{gocean2d}
\PYG{+w}{  }\PYG{k}{use }\PYG{n}{grid\PYGZus{}mod}\PYG{+w}{  }\PYG{c}{! From dl\PYGZus{}esm\PYGZus{}inf}
\PYG{+w}{  }\PYG{k}{use }\PYG{n}{field\PYGZus{}mod}\PYG{+w}{ }\PYG{c}{! From dl\PYGZus{}esm\PYGZus{}inf}
\PYG{+w}{  }\PYG{k}{use }\PYG{n}{model\PYGZus{}mod}
\PYG{+w}{  }\PYG{k}{use }\PYG{n}{boundary\PYGZus{}conditions\PYGZus{}mod}

\PYG{+w}{  }\PYG{c}{!\PYGZgt{} The grid on which our fields are defined. Must have the \PYGZsq{}target\PYGZsq{}}
\PYG{+w}{  }\PYG{c}{!! attribute because each field object contains a pointer to it.}
\PYG{+w}{  }\PYG{k}{type}\PYG{p}{(}\PYG{n}{grid\PYGZus{}type}\PYG{p}{)}\PYG{p}{,}\PYG{+w}{ }\PYG{k}{target}\PYG{+w}{ }\PYG{k+kd}{::}\PYG{+w}{ }\PYG{n}{model\PYGZus{}grid}

\PYG{+w}{  }\PYG{c}{!\PYGZgt{} Current (\PYGZsq{}now\PYGZsq{}) velocity component fields}
\PYG{+w}{  }\PYG{k}{type}\PYG{p}{(}\PYG{n}{r2d\PYGZus{}field}\PYG{p}{)}\PYG{+w}{ }\PYG{k+kd}{::}\PYG{+w}{ }\PYG{n}{un\PYGZus{}fld}\PYG{p}{,}\PYG{+w}{ }\PYG{n}{vn\PYGZus{}fld}
\PYG{+w}{  }\PYG{c}{!\PYGZgt{} \PYGZsq{}After\PYGZsq{} velocity component fields}
\PYG{+w}{  }\PYG{k}{type}\PYG{p}{(}\PYG{n}{r2d\PYGZus{}field}\PYG{p}{)}\PYG{+w}{ }\PYG{k+kd}{::}\PYG{+w}{ }\PYG{n}{ua\PYGZus{}fld}\PYG{p}{,}\PYG{+w}{ }\PYG{n}{va\PYGZus{}fld}
\PYG{+w}{  }\PYG{p}{.}\PYG{p}{.}\PYG{p}{.}

\PYG{+w}{  }\PYG{c}{! time stepping index}
\PYG{+w}{  }\PYG{k+kt}{integer}\PYG{+w}{ }\PYG{k+kd}{::}\PYG{+w}{ }\PYG{n}{istp}

\PYG{+w}{  }\PYG{c}{! Create the model grid. We use a NE offset (i.e. the U, V and F}
\PYG{+w}{  }\PYG{c}{! points immediately to the North and East of a T point all have the}
\PYG{+w}{  }\PYG{c}{! same i,j index).  This is the same offset scheme as used by NEMO.}
\PYG{+w}{  }\PYG{n}{model\PYGZus{}grid}\PYG{+w}{ }\PYG{o}{=}\PYG{+w}{ }\PYG{n}{grid\PYGZus{}type}\PYG{p}{(}\PYG{n}{GO\PYGZus{}ARAKAWA\PYGZus{}C}\PYG{p}{,}\PYG{+w}{                                }\PYG{p}{\PYGZam{}}
\PYG{+w}{                        }\PYG{p}{(}\PYG{o}{/}\PYG{n}{GO\PYGZus{}BC\PYGZus{}EXTERNAL}\PYG{p}{,}\PYG{n}{GO\PYGZus{}BC\PYGZus{}EXTERNAL}\PYG{p}{,}\PYG{n}{GO\PYGZus{}BC\PYGZus{}NONE}\PYG{o}{/}\PYG{p}{)}\PYG{p}{,}\PYG{+w}{ }\PYG{p}{\PYGZam{}}
\PYG{+w}{                         }\PYG{n}{GO\PYGZus{}OFFSET\PYGZus{}NE}\PYG{p}{)}

\PYG{+w}{  }\PYG{c}{!! read in model parameters and configure the model grid}
\PYG{+w}{  }\PYG{k}{CALL }\PYG{n}{model\PYGZus{}init}\PYG{p}{(}\PYG{n}{model\PYGZus{}grid}\PYG{p}{)}

\PYG{+w}{  }\PYG{c}{! Create fields on this grid}

\PYG{+w}{  }\PYG{c}{! Velocity components now (current time step)}
\PYG{+w}{  }\PYG{n}{un\PYGZus{}fld}\PYG{+w}{ }\PYG{o}{=}\PYG{+w}{ }\PYG{n}{r2d\PYGZus{}field}\PYG{p}{(}\PYG{n}{model\PYGZus{}grid}\PYG{p}{,}\PYG{+w}{ }\PYG{n}{GO\PYGZus{}U\PYGZus{}POINTS}\PYG{p}{)}
\PYG{+w}{  }\PYG{n}{vn\PYGZus{}fld}\PYG{+w}{ }\PYG{o}{=}\PYG{+w}{ }\PYG{n}{r2d\PYGZus{}field}\PYG{p}{(}\PYG{n}{model\PYGZus{}grid}\PYG{p}{,}\PYG{+w}{ }\PYG{n}{GO\PYGZus{}V\PYGZus{}POINTS}\PYG{p}{)}

\PYG{+w}{  }\PYG{c}{! Velocity components \PYGZsq{}after\PYGZsq{} (next time step)}
\PYG{+w}{  }\PYG{n}{ua\PYGZus{}fld}\PYG{+w}{ }\PYG{o}{=}\PYG{+w}{ }\PYG{n}{r2d\PYGZus{}field}\PYG{p}{(}\PYG{n}{model\PYGZus{}grid}\PYG{p}{,}\PYG{+w}{ }\PYG{n}{GO\PYGZus{}U\PYGZus{}POINTS}\PYG{p}{)}
\PYG{+w}{  }\PYG{n}{va\PYGZus{}fld}\PYG{+w}{ }\PYG{o}{=}\PYG{+w}{ }\PYG{n}{r2d\PYGZus{}field}\PYG{p}{(}\PYG{n}{model\PYGZus{}grid}\PYG{p}{,}\PYG{+w}{ }\PYG{n}{GO\PYGZus{}V\PYGZus{}POINTS}\PYG{p}{)}

\PYG{+w}{  }\PYG{p}{.}\PYG{p}{.}\PYG{p}{.}

\PYG{+w}{  }\PYG{c}{!! time stepping}
\PYG{+w}{  }\PYG{k}{do }\PYG{n}{istp}\PYG{+w}{ }\PYG{o}{=}\PYG{+w}{ }\PYG{n}{nit000}\PYG{p}{,}\PYG{+w}{ }\PYG{n}{nitend}\PYG{p}{,}\PYG{+w}{ }\PYG{l+m+mi}{1}

\PYG{+w}{    }\PYG{k}{call }\PYG{n}{step}\PYG{p}{(}\PYG{n}{istp}\PYG{p}{,}\PYG{+w}{                               }\PYG{p}{\PYGZam{}}
\PYG{+w}{              }\PYG{n}{ua\PYGZus{}fld}\PYG{p}{,}\PYG{+w}{ }\PYG{n}{va\PYGZus{}fld}\PYG{p}{,}\PYG{+w}{ }\PYG{n}{un\PYGZus{}fld}\PYG{p}{,}\PYG{+w}{ }\PYG{n}{vn\PYGZus{}fld}\PYG{p}{,}\PYG{+w}{     }\PYG{p}{\PYGZam{}}
\PYG{+w}{              }\PYG{p}{.}\PYG{p}{.}\PYG{p}{.}\PYG{p}{)}
\PYG{+w}{  }\PYG{k}{end }\PYG{k}{do}
\PYG{+w}{  }\PYG{p}{.}\PYG{p}{.}\PYG{p}{.}
\PYG{k}{end }\PYG{k}{program }\PYG{n}{gocean2d}
\end{sphinxVerbatim}

\sphinxAtStartPar
The \sphinxcode{\sphinxupquote{model\_init}} routine is application specific since it must
determine details of the model configuration being run, \sphinxstyleemphasis{e.g.} by
reading a namelist file. An example might look something like:

\begin{sphinxVerbatim}[commandchars=\\\{\}]
\PYG{k}{subroutine }\PYG{n}{model\PYGZus{}init}\PYG{p}{(}\PYG{n}{grid}\PYG{p}{)}
\PYG{+w}{  }\PYG{k}{type}\PYG{p}{(}\PYG{n}{grid\PYGZus{}type}\PYG{p}{)}\PYG{p}{,}\PYG{+w}{ }\PYG{k}{intent}\PYG{p}{(}\PYG{n}{inout}\PYG{p}{)}\PYG{+w}{ }\PYG{k+kd}{::}\PYG{+w}{ }\PYG{n}{grid}

\PYG{+w}{  }\PYG{c}{!\PYGZgt{} Problem size, read from namelist}
\PYG{+w}{  }\PYG{k+kt}{integer}\PYG{+w}{ }\PYG{k+kd}{::}\PYG{+w}{ }\PYG{n}{jpiglo}\PYG{p}{,}\PYG{+w}{ }\PYG{n}{jpjglo}
\PYG{+w}{  }\PYG{k+kt}{real}\PYG{p}{(}\PYG{n}{wp}\PYG{p}{)}\PYG{+w}{ }\PYG{k+kd}{::}\PYG{+w}{ }\PYG{n}{dx}\PYG{p}{,}\PYG{+w}{ }\PYG{n}{dy}
\PYG{+w}{  }\PYG{k+kt}{integer}\PYG{p}{,}\PYG{+w}{ }\PYG{k}{dimension}\PYG{p}{(}\PYG{p}{:}\PYG{p}{,}\PYG{p}{:}\PYG{p}{)}\PYG{p}{,}\PYG{+w}{ }\PYG{k}{allocatable}\PYG{+w}{ }\PYG{k+kd}{::}\PYG{+w}{ }\PYG{n}{tmask}

\PYG{+w}{  }\PYG{c}{! Read model configuration from namelist}
\PYG{+w}{  }\PYG{k}{call }\PYG{n}{read\PYGZus{}namelist}\PYG{p}{(}\PYG{n}{jpiglo}\PYG{p}{,}\PYG{+w}{ }\PYG{n}{jpjglo}\PYG{p}{,}\PYG{+w}{ }\PYG{n}{dx}\PYG{p}{,}\PYG{+w}{ }\PYG{n}{dy}\PYG{p}{,}\PYG{+w}{ }\PYG{p}{\PYGZam{}}
\PYG{+w}{                     }\PYG{n}{nit000}\PYG{p}{,}\PYG{+w}{ }\PYG{n}{nitend}\PYG{p}{,}\PYG{+w}{ }\PYG{n}{irecord}\PYG{p}{,}\PYG{+w}{ }\PYG{p}{\PYGZam{}}
\PYG{+w}{                     }\PYG{n}{jphgr\PYGZus{}msh}\PYG{p}{,}\PYG{+w}{ }\PYG{n}{dep\PYGZus{}const}\PYG{p}{,}\PYG{+w}{ }\PYG{n}{rdt}\PYG{p}{,}\PYG{+w}{ }\PYG{n}{cbfr}\PYG{p}{,}\PYG{+w}{ }\PYG{n}{visc}\PYG{p}{)}

\PYG{+w}{  }\PYG{c}{! Set\PYGZhy{}up the T mask. This defines the model domain.}
\PYG{+w}{  }\PYG{k}{allocate}\PYG{p}{(}\PYG{n}{tmask}\PYG{p}{(}\PYG{n}{jpiglo}\PYG{p}{,}\PYG{n}{jpjglo}\PYG{p}{)}\PYG{p}{)}

\PYG{+w}{  }\PYG{k}{call }\PYG{n}{setup\PYGZus{}tpoints\PYGZus{}mask}\PYG{p}{(}\PYG{n}{jpiglo}\PYG{p}{,}\PYG{+w}{ }\PYG{n}{jpjglo}\PYG{p}{,}\PYG{+w}{ }\PYG{n}{tmask}\PYG{p}{)}

\PYG{+w}{  }\PYG{c}{! Having specified the T points mask, we can set up mesh parameters}
\PYG{+w}{  }\PYG{k}{call }\PYG{n}{grid\PYGZus{}init}\PYG{p}{(}\PYG{n}{grid}\PYG{p}{,}\PYG{+w}{ }\PYG{n}{jpiglo}\PYG{p}{,}\PYG{+w}{ }\PYG{n}{jpjglo}\PYG{p}{,}\PYG{+w}{ }\PYG{n}{dx}\PYG{p}{,}\PYG{+w}{ }\PYG{n}{dy}\PYG{p}{,}\PYG{+w}{ }\PYG{n}{tmask}\PYG{p}{)}

\PYG{+w}{  }\PYG{c}{! Clean\PYGZhy{}up. T\PYGZhy{}mask has been copied into the grid object.}
\PYG{+w}{  }\PYG{k}{deallocate}\PYG{p}{(}\PYG{n}{tmask}\PYG{p}{)}

\PYG{k}{end }\PYG{k}{subroutine }\PYG{n}{model\PYGZus{}init}
\end{sphinxVerbatim}

\sphinxAtStartPar
Here, only \sphinxcode{\sphinxupquote{grid\_type}} and the \sphinxcode{\sphinxupquote{grid\_init}} routine come from dl\_esm\_inf.
The remaining code is all application specific.

\sphinxAtStartPar
Once the grid object is fully configured and all fields have been
constructed, a simulation will proceed by performing calculations with
those fields.  In the example program given above, this calculation is
performed in the time\sphinxhyphen{}stepping loop within the \sphinxcode{\sphinxupquote{step}}
subroutine. The way in which this routine uses Invoke calls is
described in the {\hyperref[\detokenize{gocean1p0:gocean-invokes}]{\sphinxcrossref{\DUrole{std,std-ref}{Invokes}}}} Section.


\section{Algorithm}
\label{\detokenize{gocean1p0:algorithm}}\label{\detokenize{gocean1p0:gocean-api-algorithm}}
\sphinxAtStartPar
The Algorithm is the top\sphinxhyphen{}level specification of the natural science
implemented in the software. Essentially it consists of mesh setup,
field declarations, initialisation of fields and (a series of) Kernel
calls. Infrastructure to support these tasks is provided in
version 1.0 of the GOcean library (see {\hyperref[\detokenize{gocean1p0:gocean-library}]{\sphinxcrossref{\DUrole{std,std-ref}{The GOcean Infrastructure Library \sphinxhyphen{} dl\_esm\_inf}}}}).


\subsection{Invokes}
\label{\detokenize{gocean1p0:invokes}}\label{\detokenize{gocean1p0:gocean-invokes}}
\sphinxAtStartPar
The Kernels to call are specified through the use of Invokes, e.g.:

\begin{sphinxVerbatim}[commandchars=\\\{\}]
\PYG{k}{call }\PYG{n}{invoke}\PYG{p}{(}\PYG{+w}{ }\PYG{n}{kernel1}\PYG{p}{(}\PYG{n}{field1}\PYG{p}{,}\PYG{+w}{ }\PYG{n}{field2}\PYG{p}{)}\PYG{p}{,}\PYG{+w}{                      }\PYG{p}{\PYGZam{}}
\PYG{+w}{             }\PYG{n}{kernel2}\PYG{p}{(}\PYG{n}{field1}\PYG{p}{,}\PYG{+w}{ }\PYG{n}{field3}\PYG{p}{)}\PYG{+w}{                       }\PYG{p}{\PYGZam{}}
\PYG{+w}{           }\PYG{p}{)}
\end{sphinxVerbatim}

\sphinxAtStartPar
The location and number of these \sphinxcode{\sphinxupquote{call invoke(...)}} statements
within the source code is entirely up to the user. The only
requirement is that PSyclone must be run on every source file that
contains one or more Invokes.
The body of each Invoke specifies the kernels to be called, the order
in which they are to be applied and the fields (and scalars) that they
work with.

\sphinxAtStartPar
Note that the kernel names specified in an Invoke are the names of the
corresponding kernel \sphinxstyleemphasis{types} defined in the kernel metadata (see the
{\hyperref[\detokenize{gocean1p0:gocean-kernels}]{\sphinxcrossref{\DUrole{std,std-ref}{Kernel}}}} Section). These are not the same as the names
of the Fortran subroutines which contain the actual kernel code.
The kernel arguments are typically field objects, as described in the
{\hyperref[\detokenize{gocean1p0:gocean-fields}]{\sphinxcrossref{\DUrole{std,std-ref}{Fields}}}} Section, but they may also be scalar
quantities (real or integer).

\sphinxAtStartPar
In the example \sphinxcode{\sphinxupquote{gocean2d}} program shown earlier, there is only one
Invoke call and it is contained within the \sphinxcode{\sphinxupquote{step}} subroutine:

\begin{sphinxVerbatim}[commandchars=\\\{\}]
\PYG{k}{subroutine }\PYG{n}{step}\PYG{p}{(}\PYG{n}{istp}\PYG{p}{,}\PYG{+w}{                   }\PYG{p}{\PYGZam{}}
\PYG{+w}{                }\PYG{n}{ua}\PYG{p}{,}\PYG{+w}{ }\PYG{n}{va}\PYG{p}{,}\PYG{+w}{ }\PYG{n}{un}\PYG{p}{,}\PYG{+w}{ }\PYG{n}{vn}\PYG{p}{,}\PYG{+w}{         }\PYG{p}{\PYGZam{}}
\PYG{+w}{                }\PYG{n}{sshn\PYGZus{}t}\PYG{p}{,}\PYG{+w}{ }\PYG{n}{sshn\PYGZus{}u}\PYG{p}{,}\PYG{+w}{ }\PYG{n}{sshn\PYGZus{}v}\PYG{p}{,}\PYG{+w}{ }\PYG{p}{\PYGZam{}}
\PYG{+w}{                }\PYG{n}{ssha\PYGZus{}t}\PYG{p}{,}\PYG{+w}{ }\PYG{n}{ssha\PYGZus{}u}\PYG{p}{,}\PYG{+w}{ }\PYG{n}{ssha\PYGZus{}v}\PYG{p}{,}\PYG{+w}{ }\PYG{p}{\PYGZam{}}
\PYG{+w}{                }\PYG{n}{hu}\PYG{p}{,}\PYG{+w}{ }\PYG{n}{hv}\PYG{p}{,}\PYG{+w}{ }\PYG{n}{ht}\PYG{p}{)}
\PYG{+w}{  }\PYG{k}{use }\PYG{n}{kind\PYGZus{}params\PYGZus{}mod}\PYG{+w}{  }\PYG{c}{! From dl\PYGZus{}esm\PYGZus{}inf}
\PYG{+w}{  }\PYG{k}{use }\PYG{n}{grid\PYGZus{}mod}\PYG{+w}{         }\PYG{c}{! From dl\PYGZus{}esm\PYGZus{}inf}
\PYG{+w}{  }\PYG{k}{use }\PYG{n}{field\PYGZus{}mod}\PYG{+w}{        }\PYG{c}{! From dl\PYGZus{}esm\PYGZus{}inf}
\PYG{+w}{  }\PYG{k}{use }\PYG{n}{model\PYGZus{}mod}\PYG{p}{,}\PYG{+w}{ }\PYG{k}{only}\PYG{p}{:}\PYG{+w}{ }\PYG{n}{rdt}\PYG{+w}{ }\PYG{c}{! The model time\PYGZhy{}step}
\PYG{+w}{  }\PYG{k}{use }\PYG{n}{continuity\PYGZus{}mod}\PYG{p}{,}\PYG{+w}{  }\PYG{k}{only}\PYG{p}{:}\PYG{+w}{ }\PYG{n}{continuity}
\PYG{+w}{  }\PYG{k}{use }\PYG{n}{momentum\PYGZus{}mod}\PYG{p}{,}\PYG{+w}{    }\PYG{k}{only}\PYG{p}{:}\PYG{+w}{ }\PYG{n}{momentum\PYGZus{}u}\PYG{p}{,}\PYG{+w}{ }\PYG{n}{momentum\PYGZus{}v}
\PYG{+w}{  }\PYG{k}{use }\PYG{n}{boundary\PYGZus{}conditions\PYGZus{}mod}\PYG{p}{,}\PYG{+w}{ }\PYG{k}{only}\PYG{p}{:}\PYG{+w}{ }\PYG{n}{bc\PYGZus{}ssh}\PYG{p}{,}\PYG{+w}{ }\PYG{n}{bc\PYGZus{}solid\PYGZus{}u}
\PYG{+w}{  }\PYG{c}{!\PYGZgt{} The current time step}
\PYG{+w}{  }\PYG{k+kt}{integer}\PYG{p}{,}\PYG{+w}{         }\PYG{k}{intent}\PYG{p}{(}\PYG{n}{inout}\PYG{p}{)}\PYG{+w}{ }\PYG{k+kd}{::}\PYG{+w}{ }\PYG{n}{istp}
\PYG{+w}{  }\PYG{k}{type}\PYG{p}{(}\PYG{n}{r2d\PYGZus{}field}\PYG{p}{)}\PYG{p}{,}\PYG{+w}{ }\PYG{k}{intent}\PYG{p}{(}\PYG{n}{inout}\PYG{p}{)}\PYG{+w}{ }\PYG{k+kd}{::}\PYG{+w}{ }\PYG{n}{un}\PYG{p}{,}\PYG{+w}{ }\PYG{n}{vn}\PYG{p}{,}\PYG{+w}{ }\PYG{n}{sshn\PYGZus{}t}\PYG{p}{,}\PYG{+w}{ }\PYG{n}{sshn\PYGZus{}u}\PYG{p}{,}\PYG{+w}{ }\PYG{n}{sshn\PYGZus{}v}
\PYG{+w}{  }\PYG{k}{type}\PYG{p}{(}\PYG{n}{r2d\PYGZus{}field}\PYG{p}{)}\PYG{p}{,}\PYG{+w}{ }\PYG{k}{intent}\PYG{p}{(}\PYG{n}{inout}\PYG{p}{)}\PYG{+w}{ }\PYG{k+kd}{::}\PYG{+w}{ }\PYG{n}{ua}\PYG{p}{,}\PYG{+w}{ }\PYG{n}{va}\PYG{p}{,}\PYG{+w}{ }\PYG{n}{ssha\PYGZus{}t}\PYG{p}{,}\PYG{+w}{ }\PYG{n}{ssha\PYGZus{}u}\PYG{p}{,}\PYG{+w}{ }\PYG{n}{ssha\PYGZus{}v}
\PYG{+w}{  }\PYG{k}{type}\PYG{p}{(}\PYG{n}{r2d\PYGZus{}field}\PYG{p}{)}\PYG{p}{,}\PYG{+w}{ }\PYG{k}{intent}\PYG{p}{(}\PYG{n}{inout}\PYG{p}{)}\PYG{+w}{ }\PYG{k+kd}{::}\PYG{+w}{ }\PYG{n}{hu}\PYG{p}{,}\PYG{+w}{ }\PYG{n}{hv}\PYG{p}{,}\PYG{+w}{ }\PYG{n}{ht}

\PYG{+w}{  }\PYG{k}{call }\PYG{n}{invoke}\PYG{p}{(}\PYG{+w}{                                               }\PYG{p}{\PYGZam{}}
\PYG{+w}{              }\PYG{n}{continuity}\PYG{p}{(}\PYG{n}{ssha\PYGZus{}t}\PYG{p}{,}\PYG{+w}{ }\PYG{n}{sshn\PYGZus{}t}\PYG{p}{,}\PYG{+w}{ }\PYG{n}{sshn\PYGZus{}u}\PYG{p}{,}\PYG{+w}{ }\PYG{n}{sshn\PYGZus{}v}\PYG{p}{,}\PYG{+w}{     }\PYG{p}{\PYGZam{}}
\PYG{+w}{                         }\PYG{n}{hu}\PYG{p}{,}\PYG{+w}{ }\PYG{n}{hv}\PYG{p}{,}\PYG{+w}{ }\PYG{n}{un}\PYG{p}{,}\PYG{+w}{ }\PYG{n}{vn}\PYG{p}{,}\PYG{+w}{ }\PYG{n}{rdt}\PYG{p}{)}\PYG{p}{,}\PYG{+w}{               }\PYG{p}{\PYGZam{}}
\PYG{+w}{              }\PYG{n}{momentum\PYGZus{}u}\PYG{p}{(}\PYG{n}{ua}\PYG{p}{,}\PYG{+w}{ }\PYG{n}{un}\PYG{p}{,}\PYG{+w}{ }\PYG{n}{vn}\PYG{p}{,}\PYG{+w}{ }\PYG{n}{hu}\PYG{p}{,}\PYG{+w}{ }\PYG{n}{hv}\PYG{p}{,}\PYG{+w}{ }\PYG{n}{ht}\PYG{p}{,}\PYG{+w}{             }\PYG{p}{\PYGZam{}}
\PYG{+w}{                         }\PYG{n}{ssha\PYGZus{}u}\PYG{p}{,}\PYG{+w}{ }\PYG{n}{sshn\PYGZus{}t}\PYG{p}{,}\PYG{+w}{ }\PYG{n}{sshn\PYGZus{}u}\PYG{p}{,}\PYG{+w}{ }\PYG{n}{sshn\PYGZus{}v}\PYG{p}{)}\PYG{p}{,}\PYG{+w}{    }\PYG{p}{\PYGZam{}}
\PYG{+w}{              }\PYG{n}{momentum\PYGZus{}v}\PYG{p}{(}\PYG{n}{va}\PYG{p}{,}\PYG{+w}{ }\PYG{n}{un}\PYG{p}{,}\PYG{+w}{ }\PYG{n}{vn}\PYG{p}{,}\PYG{+w}{ }\PYG{n}{hu}\PYG{p}{,}\PYG{+w}{ }\PYG{n}{hv}\PYG{p}{,}\PYG{+w}{ }\PYG{n}{ht}\PYG{p}{,}\PYG{+w}{             }\PYG{p}{\PYGZam{}}
\PYG{+w}{                         }\PYG{n}{ssha\PYGZus{}v}\PYG{p}{,}\PYG{+w}{ }\PYG{n}{sshn\PYGZus{}t}\PYG{p}{,}\PYG{+w}{ }\PYG{n}{sshn\PYGZus{}u}\PYG{p}{,}\PYG{+w}{ }\PYG{n}{sshn\PYGZus{}v}\PYG{p}{)}\PYG{p}{,}\PYG{+w}{    }\PYG{p}{\PYGZam{}}
\PYG{+w}{              }\PYG{n}{bc\PYGZus{}ssh}\PYG{p}{(}\PYG{n}{istp}\PYG{p}{,}\PYG{+w}{ }\PYG{n}{ssha\PYGZus{}t}\PYG{p}{)}\PYG{p}{,}\PYG{+w}{                          }\PYG{p}{\PYGZam{}}
\PYG{+w}{              }\PYG{n}{bc\PYGZus{}solid\PYGZus{}u}\PYG{p}{(}\PYG{n}{ua}\PYG{p}{)}\PYG{p}{,}\PYG{+w}{                                }\PYG{p}{\PYGZam{}}
\PYG{+w}{              }\PYG{p}{.}\PYG{p}{.}\PYG{p}{.}
\PYG{+w}{             }\PYG{p}{)}
\PYG{k}{end }\PYG{k}{subroutine }\PYG{n}{step}
\end{sphinxVerbatim}

\sphinxAtStartPar
Note that in this example the grid was constructed for a
model with ‘external’ boundary conditions. These boundary conditions
are applied through several user\sphinxhyphen{}supplied kernels, two of which
(\sphinxcode{\sphinxupquote{bc\_ssh}} and \sphinxcode{\sphinxupquote{bc\_solid\_u}}) are include in the above code
fragment.


\section{Kernel}
\label{\detokenize{gocean1p0:kernel}}\label{\detokenize{gocean1p0:gocean-kernels}}
\sphinxAtStartPar
The general requirements for the structure of a Kernel are explained
in the {\hyperref[\detokenize{introduction_to_psykal:kernel-layer}]{\sphinxcrossref{\DUrole{std,std-ref}{Kernel layer}}}} section. This section explains the metadata
and subroutine arguments that are specific to the GOcean 1.0 API.


\subsection{Metadata}
\label{\detokenize{gocean1p0:metadata}}
\sphinxAtStartPar
The metadata for a GOcean 1.0 API kernel has four components:
\begin{enumerate}
\sphinxsetlistlabels{\arabic}{enumi}{enumii}{}{)}%
\item {} 
\sphinxAtStartPar
‘meta\_args’,

\item {} 
\sphinxAtStartPar
‘iterates\_over’,

\item {} 
\sphinxAtStartPar
‘index\_offset’ and

\item {} 
\sphinxAtStartPar
‘procedure’:

\end{enumerate}

\sphinxAtStartPar
These are illustrated in the code below:

\begin{sphinxVerbatim}[commandchars=\\\{\}]
\PYG{k}{type}\PYG{p}{,}\PYG{+w}{ }\PYG{k}{extends}\PYG{p}{(}\PYG{n}{kernel\PYGZus{}type}\PYG{p}{)}\PYG{+w}{ }\PYG{k+kd}{::}\PYG{+w}{ }\PYG{n}{my\PYGZus{}kernel\PYGZus{}type}
\PYG{+w}{   }\PYG{k}{type}\PYG{p}{(}\PYG{n}{go\PYGZus{}arg}\PYG{p}{)}\PYG{p}{,}\PYG{+w}{ }\PYG{k}{dimension}\PYG{p}{(}\PYG{p}{.}\PYG{p}{.}\PYG{p}{.}\PYG{p}{)}\PYG{+w}{ }\PYG{k+kd}{::}\PYG{+w}{ }\PYG{n}{meta\PYGZus{}args}\PYG{+w}{ }\PYG{o}{=}\PYG{+w}{ }\PYG{p}{(}\PYG{o}{/}\PYG{+w}{ }\PYG{p}{.}\PYG{p}{.}\PYG{p}{.}\PYG{+w}{ }\PYG{o}{/}\PYG{p}{)}
\PYG{+w}{   }\PYG{k+kt}{integer}\PYG{+w}{ }\PYG{k+kd}{::}\PYG{+w}{ }\PYG{n}{iterates\PYGZus{}over}\PYG{+w}{ }\PYG{o}{=}\PYG{+w}{ }\PYG{p}{.}\PYG{p}{.}\PYG{p}{.}
\PYG{+w}{   }\PYG{k+kt}{integer}\PYG{+w}{ }\PYG{k+kd}{::}\PYG{+w}{ }\PYG{n}{index\PYGZus{}offset}\PYG{+w}{ }\PYG{o}{=}\PYG{+w}{ }\PYG{p}{.}\PYG{p}{.}\PYG{p}{.}
\PYG{k}{contains}
\PYG{k}{  }\PYG{k}{procedure}\PYG{p}{,}\PYG{+w}{ }\PYG{k}{nopass}\PYG{+w}{ }\PYG{k+kd}{::}\PYG{+w}{ }\PYG{n}{code}\PYG{+w}{ }\PYG{o}{=}\PYG{o}{\PYGZgt{}}\PYG{+w}{ }\PYG{n}{my\PYGZus{}kernel\PYGZus{}code}
\PYG{k}{end }\PYG{k}{type }\PYG{n}{my\PYGZus{}kernel\PYGZus{}type}
\end{sphinxVerbatim}

\sphinxAtStartPar
These four metadata elements are discussed in order in the following
sections.


\subsubsection{Argument Metadata: meta\_args}
\label{\detokenize{gocean1p0:argument-metadata-meta-args}}
\sphinxAtStartPar
The \sphinxcode{\sphinxupquote{meta\_args}} array specifies information about data that the
kernel code expects to be passed to it via its argument list. There is
one entry in the \sphinxcode{\sphinxupquote{meta\_args}} array for each \sphinxstylestrong{scalar}, \sphinxstylestrong{field},
or \sphinxstylestrong{grid\sphinxhyphen{}property} passed into the Kernel. Their ordering in the
\sphinxcode{\sphinxupquote{meta\_args}} array must be the same as that in the kernel code
argument list. The entry must be of type \sphinxcode{\sphinxupquote{go\_arg}} which itself contains
metadata about the associated argument. The size of the meta\_args
array must correspond to the total number of \sphinxstylestrong{scalars}, \sphinxstylestrong{fields}
and \sphinxstylestrong{grid properties} passed into the Kernel.

\sphinxAtStartPar
For example, if there are a total of two \sphinxstylestrong{field} entities being passed
to the Kernel then the meta\_args array will be of size 2 and there
will be two entries of type \sphinxcode{\sphinxupquote{GO\_arg}}:

\begin{sphinxVerbatim}[commandchars=\\\{\}]
\PYG{k}{type}\PYG{p}{(}\PYG{n}{GO\PYGZus{}arg}\PYG{p}{)}\PYG{+w}{ }\PYG{k+kd}{::}\PYG{+w}{ }\PYG{n}{meta\PYGZus{}args}\PYG{p}{(}\PYG{l+m+mi}{2}\PYG{p}{)}\PYG{+w}{ }\PYG{o}{=}\PYG{+w}{ }\PYG{p}{(}\PYG{o}{/}\PYG{+w}{                                  }\PYG{p}{\PYGZam{}}
\PYG{+w}{     }\PYG{n}{go\PYGZus{}arg}\PYG{p}{(}\PYG{+w}{ }\PYG{p}{.}\PYG{p}{.}\PYG{p}{.}\PYG{+w}{ }\PYG{p}{)}\PYG{p}{,}\PYG{+w}{                                                }\PYG{p}{\PYGZam{}}
\PYG{+w}{     }\PYG{n}{go\PYGZus{}arg}\PYG{p}{(}\PYG{+w}{ }\PYG{p}{.}\PYG{p}{.}\PYG{p}{.}\PYG{+w}{ }\PYG{p}{)}\PYG{+w}{                                                 }\PYG{p}{\PYGZam{}}
\PYG{+w}{     }\PYG{o}{/}\PYG{p}{)}
\end{sphinxVerbatim}

\sphinxAtStartPar
Argument\sphinxhyphen{}metadata (metadata contained within the brackets of an
\sphinxcode{\sphinxupquote{go\_arg}} entry), describes either a \sphinxstylestrong{scalar}, a \sphinxstylestrong{field} or a \sphinxstylestrong{grid
property}.

\sphinxAtStartPar
The first argument\sphinxhyphen{}metadata entry describes how the kernel will access
the corresponding argument. As an example, the following \sphinxcode{\sphinxupquote{meta\_args}}
metadata describes four entries, the first one is written to by the
kernel while the remaining three are only read:

\begin{sphinxVerbatim}[commandchars=\\\{\}]
\PYG{k}{type}\PYG{p}{(}\PYG{n}{go\PYGZus{}arg}\PYG{p}{)}\PYG{+w}{ }\PYG{k+kd}{::}\PYG{+w}{ }\PYG{n}{meta\PYGZus{}args}\PYG{p}{(}\PYG{l+m+mi}{4}\PYG{p}{)}\PYG{+w}{ }\PYG{o}{=}\PYG{+w}{ }\PYG{p}{(}\PYG{o}{/}\PYG{+w}{                               }\PYG{p}{\PYGZam{}}
\PYG{+w}{     }\PYG{n}{go\PYGZus{}arg}\PYG{p}{(}\PYG{n}{GO\PYGZus{}WRITE}\PYG{p}{,}\PYG{+w}{ }\PYG{p}{.}\PYG{p}{.}\PYG{p}{.}\PYG{+w}{ }\PYG{p}{)}\PYG{p}{,}\PYG{+w}{                                    }\PYG{p}{\PYGZam{}}
\PYG{+w}{     }\PYG{n}{go\PYGZus{}arg}\PYG{p}{(}\PYG{n}{GO\PYGZus{}READ}\PYG{p}{,}\PYG{+w}{ }\PYG{p}{.}\PYG{p}{.}\PYG{p}{.}\PYG{+w}{ }\PYG{p}{)}\PYG{p}{,}\PYG{+w}{                                     }\PYG{p}{\PYGZam{}}
\PYG{+w}{     }\PYG{n}{go\PYGZus{}arg}\PYG{p}{(}\PYG{n}{GO\PYGZus{}READ}\PYG{p}{,}\PYG{+w}{ }\PYG{p}{.}\PYG{p}{.}\PYG{p}{.}\PYG{+w}{ }\PYG{p}{)}\PYG{p}{,}\PYG{+w}{                                     }\PYG{p}{\PYGZam{}}
\PYG{+w}{     }\PYG{n}{go\PYGZus{}arg}\PYG{p}{(}\PYG{n}{GO\PYGZus{}READ}\PYG{p}{,}\PYG{+w}{ }\PYG{p}{.}\PYG{p}{.}\PYG{p}{.}\PYG{p}{)}\PYG{+w}{                                       }\PYG{p}{\PYGZam{}}
\PYG{+w}{     }\PYG{o}{/}\PYG{p}{)}
\end{sphinxVerbatim}
\phantomsection\label{\detokenize{gocean1p0:gocean-grid-point-type}}
\sphinxAtStartPar
The second entry to argument\sphinxhyphen{}metadata (information contained within
the brackets of an \sphinxcode{\sphinxupquote{go\_arg}} type) describes the type of data
represented by the argument. This type falls into three categories;
field data, scalar data and grid properties. For field data the
metadata entry consists of the type of grid\sphinxhyphen{}point that field values
are defined on. Since the GOcean API supports fields on an Arakawa C
grid, the possible grid\sphinxhyphen{}point types are \sphinxcode{\sphinxupquote{GO\_CU}}, \sphinxcode{\sphinxupquote{GO\_CV}}, \sphinxcode{\sphinxupquote{GO\_CF}} and
\sphinxcode{\sphinxupquote{GO\_CT}}. GOcean Kernels can also take scalar quantities as
arguments. Since these do not live on grid\sphinxhyphen{}points they are specified
as either \sphinxcode{\sphinxupquote{GO\_R\_SCALAR}} or \sphinxcode{\sphinxupquote{GO\_I\_SCALAR}} depending on whether the
corresponding Fortran variable is a real or integer quantity.
Finally, grid\sphinxhyphen{}property entries are used to specify any properties of
the grid required by the kernel (\sphinxstyleemphasis{e.g.} the area of cells at U points or
whether T points are wet or dry).

\sphinxAtStartPar
For example:

\begin{sphinxVerbatim}[commandchars=\\\{\}]
\PYG{k}{type}\PYG{p}{(}\PYG{n}{go\PYGZus{}arg}\PYG{p}{)}\PYG{+w}{ }\PYG{k+kd}{::}\PYG{+w}{ }\PYG{n}{meta\PYGZus{}args}\PYG{p}{(}\PYG{l+m+mi}{4}\PYG{p}{)}\PYG{+w}{ }\PYG{o}{=}\PYG{+w}{ }\PYG{p}{(}\PYG{o}{/}\PYG{+w}{                                  }\PYG{p}{\PYGZam{}}
\PYG{+w}{     }\PYG{n}{go\PYGZus{}arg}\PYG{p}{(}\PYG{n}{GO\PYGZus{}WRITE}\PYG{p}{,}\PYG{+w}{ }\PYG{n}{GO\PYGZus{}CT}\PYG{p}{,}\PYG{+w}{ }\PYG{p}{.}\PYG{p}{.}\PYG{p}{.}\PYG{+w}{ }\PYG{p}{)}\PYG{p}{,}\PYG{+w}{                                }\PYG{p}{\PYGZam{}}
\PYG{+w}{     }\PYG{n}{go\PYGZus{}arg}\PYG{p}{(}\PYG{n}{GO\PYGZus{}READ}\PYG{p}{,}\PYG{+w}{  }\PYG{n}{GO\PYGZus{}CU}\PYG{p}{,}\PYG{+w}{ }\PYG{p}{.}\PYG{p}{.}\PYG{p}{.}\PYG{+w}{ }\PYG{p}{)}\PYG{p}{,}\PYG{+w}{                                }\PYG{p}{\PYGZam{}}
\PYG{+w}{     }\PYG{n}{go\PYGZus{}arg}\PYG{p}{(}\PYG{n}{GO\PYGZus{}READ}\PYG{p}{,}\PYG{+w}{  }\PYG{n}{GO\PYGZus{}R\PYGZus{}SCALAR}\PYG{p}{,}\PYG{+w}{ }\PYG{p}{.}\PYG{p}{.}\PYG{p}{.}\PYG{+w}{ }\PYG{p}{)}\PYG{p}{,}\PYG{+w}{                          }\PYG{p}{\PYGZam{}}
\PYG{+w}{     }\PYG{n}{go\PYGZus{}arg}\PYG{p}{(}\PYG{n}{GO\PYGZus{}READ}\PYG{p}{,}\PYG{+w}{  }\PYG{n}{GO\PYGZus{}GRID\PYGZus{}AREA\PYGZus{}U}\PYG{p}{)}\PYG{+w}{                              }\PYG{p}{\PYGZam{}}
\PYG{+w}{     }\PYG{o}{/}\PYG{p}{)}
\end{sphinxVerbatim}

\sphinxAtStartPar
Here, the first argument is a field on T points, the second is a field
on U points, the fourth is a real scalar and the fifth is a property
of the grid (cell area at U points).

\sphinxAtStartPar
The full list of supported grid properties in the GOcean 1.0 API is:


\begin{savenotes}\sphinxattablestart
\sphinxthistablewithglobalstyle
\centering
\sphinxcapstartof{table}
\sphinxthecaptionisattop
\sphinxcaption{Grid Properties Table}\label{\detokenize{gocean1p0:id1}}\label{\detokenize{gocean1p0:gocean-grid-props}}
\sphinxaftertopcaption
\begin{tabulary}{\linewidth}[t]{TTT}
\sphinxtoprule
\sphinxstyletheadfamily 
\sphinxAtStartPar
Name
&\sphinxstyletheadfamily 
\sphinxAtStartPar
Description
&\sphinxstyletheadfamily 
\sphinxAtStartPar
Type
\\
\sphinxmidrule
\sphinxtableatstartofbodyhook
\sphinxAtStartPar
go\_grid\_area\_t
&
\sphinxAtStartPar
Cell area at T point
&
\sphinxAtStartPar
Real array, rank=2
\\
\sphinxhline
\sphinxAtStartPar
go\_grid\_area\_u
&
\sphinxAtStartPar
Cell area at U point
&
\sphinxAtStartPar
Real array, rank=2
\\
\sphinxhline
\sphinxAtStartPar
go\_grid\_area\_v
&
\sphinxAtStartPar
Cell area at V point
&
\sphinxAtStartPar
Real array, rank=2
\\
\sphinxhline
\sphinxAtStartPar
go\_grid\_mask\_t
&
\sphinxAtStartPar
T\sphinxhyphen{}point mask (1=wet, 0=dry)
&
\sphinxAtStartPar
Integer array, rank=2
\\
\sphinxhline
\sphinxAtStartPar
go\_grid\_dx\_t
&
\sphinxAtStartPar
Grid spacing in x at T points
&
\sphinxAtStartPar
Real array, rank=2
\\
\sphinxhline
\sphinxAtStartPar
go\_grid\_dx\_u
&
\sphinxAtStartPar
Grid spacing in x at U points
&
\sphinxAtStartPar
Real array, rank=2
\\
\sphinxhline
\sphinxAtStartPar
go\_grid\_dx\_v
&
\sphinxAtStartPar
Grid spacing in x at V points
&
\sphinxAtStartPar
Real array, rank=2
\\
\sphinxhline
\sphinxAtStartPar
go\_grid\_dy\_t
&
\sphinxAtStartPar
Grid spacing in y at T points
&
\sphinxAtStartPar
Real array, rank=2
\\
\sphinxhline
\sphinxAtStartPar
go\_grid\_dy\_u
&
\sphinxAtStartPar
Grid spacing in y at U points
&
\sphinxAtStartPar
Real array, rank=2
\\
\sphinxhline
\sphinxAtStartPar
go\_grid\_dy\_v
&
\sphinxAtStartPar
Grid spacing in y at V points
&
\sphinxAtStartPar
Real array, rank=2
\\
\sphinxhline
\sphinxAtStartPar
go\_grid\_lat\_u
&
\sphinxAtStartPar
Latitude of U points (gphiu)
&
\sphinxAtStartPar
Real array, rank=2
\\
\sphinxhline
\sphinxAtStartPar
go\_grid\_lat\_v
&
\sphinxAtStartPar
Latitude of V points (gphiv)
&
\sphinxAtStartPar
Real array, rank=2
\\
\sphinxhline
\sphinxAtStartPar
go\_grid\_dx\_const
&
\sphinxAtStartPar
Grid spacing in x if constant
&
\sphinxAtStartPar
Real, scalar
\\
\sphinxhline
\sphinxAtStartPar
go\_grid\_dy\_const
&
\sphinxAtStartPar
Grid spacing in y if constant
&
\sphinxAtStartPar
Real, scalar
\\
\sphinxhline
\sphinxAtStartPar
go\_grid\_x\_min\_index
&
\sphinxAtStartPar
Minimum X index
&
\sphinxAtStartPar
Integer, scalar
\\
\sphinxhline
\sphinxAtStartPar
go\_grid\_x\_max\_index
&
\sphinxAtStartPar
Maximum X index
&
\sphinxAtStartPar
Integer, scalar
\\
\sphinxhline
\sphinxAtStartPar
go\_grid\_y\_min\_index
&
\sphinxAtStartPar
Minimum Y index
&
\sphinxAtStartPar
Integer, scalar
\\
\sphinxhline
\sphinxAtStartPar
go\_grid\_y\_max\_index
&
\sphinxAtStartPar
Maximum Y index
&
\sphinxAtStartPar
Integer, scalar
\\
\sphinxbottomrule
\end{tabulary}
\sphinxtableafterendhook\par
\sphinxattableend\end{savenotes}

\sphinxAtStartPar
These are defined in the psyclone config file (see
{\hyperref[\detokenize{gocean1p0:gocean-configuration}]{\sphinxcrossref{\DUrole{std,std-ref}{Configuration}}}}), and the user or infrastructure
library developer can provide additional entries if required.
PSyclone will query PSyclone’s Configuration class to get the
properties required. All of the rank\sphinxhyphen{}two arrays have the
first rank as longitude (\sphinxstyleemphasis{x}) and the second as latitude (\sphinxstyleemphasis{y}).

\sphinxAtStartPar
Scalars and fields contain a third argument\sphinxhyphen{}metadata entry which
describes whether the kernel accesses the corresponding argument with
a stencil. The value \sphinxcode{\sphinxupquote{GO\_POINTWISE}} indicates that there is no stencil
access. Metadata for a scalar field is limited to this value.
Grid\sphinxhyphen{}property arguments have no third metadata argument. If there
are no stencil accesses then the full argument metadata for our
previous example will be:

\begin{sphinxVerbatim}[commandchars=\\\{\}]
\PYG{k}{type}\PYG{p}{(}\PYG{n}{go\PYGZus{}arg}\PYG{p}{)}\PYG{+w}{ }\PYG{k+kd}{::}\PYG{+w}{ }\PYG{n}{meta\PYGZus{}args}\PYG{p}{(}\PYG{l+m+mi}{4}\PYG{p}{)}\PYG{+w}{ }\PYG{o}{=}\PYG{+w}{ }\PYG{p}{(}\PYG{o}{/}\PYG{+w}{                            }\PYG{p}{\PYGZam{}}
\PYG{+w}{     }\PYG{n}{go\PYGZus{}arg}\PYG{p}{(}\PYG{n}{GO\PYGZus{}WRITE}\PYG{p}{,}\PYG{+w}{ }\PYG{n}{GO\PYGZus{}CT}\PYG{p}{,}\PYG{+w}{       }\PYG{n}{GO\PYGZus{}POINTWISE}\PYG{p}{)}\PYG{p}{,}\PYG{+w}{            }\PYG{p}{\PYGZam{}}
\PYG{+w}{     }\PYG{n}{go\PYGZus{}arg}\PYG{p}{(}\PYG{n}{GO\PYGZus{}READ}\PYG{p}{,}\PYG{+w}{  }\PYG{n}{GO\PYGZus{}CU}\PYG{p}{,}\PYG{+w}{       }\PYG{n}{GO\PYGZus{}POINTWISE}\PYG{p}{)}\PYG{p}{,}\PYG{+w}{            }\PYG{p}{\PYGZam{}}
\PYG{+w}{     }\PYG{n}{go\PYGZus{}arg}\PYG{p}{(}\PYG{n}{GO\PYGZus{}READ}\PYG{p}{,}\PYG{+w}{  }\PYG{n}{GO\PYGZus{}R\PYGZus{}SCALAR}\PYG{p}{,}\PYG{+w}{ }\PYG{n}{GO\PYGZus{}POINTWISE}\PYG{p}{)}\PYG{p}{,}\PYG{+w}{            }\PYG{p}{\PYGZam{}}
\PYG{+w}{     }\PYG{n}{go\PYGZus{}arg}\PYG{p}{(}\PYG{n}{GO\PYGZus{}READ}\PYG{p}{,}\PYG{+w}{  }\PYG{n}{GO\PYGZus{}GRID\PYGZus{}AREA\PYGZus{}U}\PYG{p}{)}\PYG{+w}{                        }\PYG{p}{\PYGZam{}}
\PYG{+w}{     }\PYG{o}{/}\PYG{p}{)}
\end{sphinxVerbatim}

\sphinxAtStartPar
If a kernel accesses a field using a stencil then the third argument
metadata entry should take the form \sphinxcode{\sphinxupquote{go\_stencil(...)}}. Note, a stencil
access is only allowed for a field that is \sphinxcode{\sphinxupquote{READ}} by a kernel.

\sphinxAtStartPar
In the GOcean API, fields are implemented as two\sphinxhyphen{}dimensional
arrays. In Fortran, a standard 5\sphinxhyphen{}point stencil would look something
like the following:

\begin{sphinxVerbatim}[commandchars=\\\{\}]
\PYG{n}{a}\PYG{p}{(}\PYG{n}{i}\PYG{p}{,}\PYG{n}{j}\PYG{p}{)}\PYG{+w}{ }\PYG{o}{+}\PYG{+w}{ }\PYG{n}{a}\PYG{p}{(}\PYG{n}{i}\PYG{o}{+}\PYG{l+m+mi}{1}\PYG{p}{,}\PYG{n}{j}\PYG{p}{)}\PYG{+w}{ }\PYG{o}{+}\PYG{+w}{ }\PYG{n}{a}\PYG{p}{(}\PYG{n}{i}\PYG{o}{\PYGZhy{}}\PYG{l+m+mi}{1}\PYG{p}{,}\PYG{n}{j}\PYG{p}{)}\PYG{+w}{ }\PYG{o}{+}\PYG{+w}{ }\PYG{n}{a}\PYG{p}{(}\PYG{n}{i}\PYG{p}{,}\PYG{n}{j}\PYG{o}{+}\PYG{l+m+mi}{1}\PYG{p}{)}\PYG{+w}{ }\PYG{o}{+}\PYG{+w}{ }\PYG{n}{a}\PYG{p}{(}\PYG{n}{i}\PYG{p}{,}\PYG{n}{j}\PYG{o}{\PYGZhy{}}\PYG{l+m+mi}{1}\PYG{p}{)}
\end{sphinxVerbatim}

\sphinxAtStartPar
If we view the above accesses as co\sphinxhyphen{}ordinates relative to the \sphinxcode{\sphinxupquote{a(i,j)}}
access we get \sphinxcode{\sphinxupquote{(0,0), (1,0), (\sphinxhyphen{}1,0), (0,1), (0,\sphinxhyphen{}1)}}. If we then view
these accesses in graphical form with \sphinxcode{\sphinxupquote{i}} being in the horizontal
direction and \sphinxcode{\sphinxupquote{j}} in the vertical and with a \sphinxcode{\sphinxupquote{1}} indicating a
(depth\sphinxhyphen{}1) access and a \sphinxcode{\sphinxupquote{0}} indicating there is no access we get the
following:

\begin{sphinxVerbatim}[commandchars=\\\{\}]
\PYG{l+m+mi}{010}
\PYG{l+m+mi}{111}
\PYG{l+m+mi}{010}
\end{sphinxVerbatim}

\sphinxAtStartPar
In the GOcean API a stencil access is captured as a triplet of
integers (one row at a time from top to bottom) using the above view
i.e.

\begin{sphinxVerbatim}[commandchars=\\\{\}]
\PYG{n}{go\PYGZus{}stencil}\PYG{p}{(}\PYG{l+m+mi}{010}\PYG{p}{,}\PYG{l+m+mi}{111}\PYG{p}{,}\PYG{l+m+mi}{010}\PYG{p}{)}
\end{sphinxVerbatim}

\sphinxAtStartPar
So far we have only considered depth\sphinxhyphen{}1 stencils. In our notation
the depth of access is captured by the integer value (\sphinxcode{\sphinxupquote{0}} for no
access, \sphinxcode{\sphinxupquote{1}} for depth 1, \sphinxcode{\sphinxupquote{2}} for depth 2 etc). For example:

\begin{sphinxVerbatim}[commandchars=\\\{\}]
\PYG{n}{a}\PYG{p}{(}\PYG{n}{i}\PYG{p}{,}\PYG{n}{j}\PYG{p}{)}\PYG{+w}{ }\PYG{o}{+}\PYG{+w}{ }\PYG{n}{a}\PYG{p}{(}\PYG{n}{i}\PYG{p}{,}\PYG{n}{j}\PYG{o}{+}\PYG{l+m+mi}{1}\PYG{p}{)}\PYG{+w}{ }\PYG{o}{+}\PYG{+w}{ }\PYG{n}{a}\PYG{p}{(}\PYG{n}{i}\PYG{p}{,}\PYG{n}{j}\PYG{o}{+}\PYG{l+m+mi}{2}\PYG{p}{)}
\end{sphinxVerbatim}

\sphinxAtStartPar
would be captured as:

\begin{sphinxVerbatim}[commandchars=\\\{\}]
\PYG{n}{go\PYGZus{}stencil}\PYG{p}{(}\PYG{l+m+mi}{020}\PYG{p}{,}\PYG{l+m+mi}{010}\PYG{p}{,}\PYG{l+m+mi}{000}\PYG{p}{)}
\end{sphinxVerbatim}

\sphinxAtStartPar
All forms of stencil can be \sphinxstylestrong{summarised} using this triplet notation
up to a depth of \sphinxcode{\sphinxupquote{9}} apart from the central \sphinxcode{\sphinxupquote{a(i,j)}} value which can
either be \sphinxcode{\sphinxupquote{0}} (not accessed) or \sphinxcode{\sphinxupquote{1}} (accessed). Note, the central
value is not currently used by PSyclone. The notation is a \sphinxstylestrong{summary}
in two ways
\begin{enumerate}
\sphinxsetlistlabels{\arabic}{enumi}{enumii}{}{)}%
\item {} 
\sphinxAtStartPar
it only captures the depth of the stencil in a particular
direction, not the actual accesses. Therefore, there is no way to
distinguish between the stencil \sphinxcode{\sphinxupquote{a(i+2,j)}} and the stencil
\sphinxcode{\sphinxupquote{a(i+1,j) + a(i+2,j)}}.

\item {} 
\sphinxAtStartPar
when there are offsets for both \sphinxcode{\sphinxupquote{i}} and \sphinxcode{\sphinxupquote{j}} e.g. \sphinxcode{\sphinxupquote{a(i+1,j+1)}} it
only captures whether there is an access in that direction at a
particular depth, not the details of the access. For example, there
is no way to distinguish between \sphinxcode{\sphinxupquote{a(i+2,j+2)}} and \sphinxcode{\sphinxupquote{a(i+2,j+2) +
a(i+1,j+2) + a(i+2,j+1)}}.

\end{enumerate}

\sphinxAtStartPar
Whilst the description is a summary, it is accurate enough for
PSyclone as this information is primarily used to determine which grid
partitions must communicate with which for the purposes of placing
halo exchange calls. In this case, it is the depth and direction
information that is most important.


\subsubsection{Iterates Over}
\label{\detokenize{gocean1p0:iterates-over}}\label{\detokenize{gocean1p0:gocean-iterates-over}}
\sphinxAtStartPar
The second element of kernel metadata is \sphinxcode{\sphinxupquote{ITERATES\_OVER}}. This
specifies that the Kernel has been written with the assumption that it
is iterating over grid points of the specified type. By default the supported
values are: \sphinxcode{\sphinxupquote{GO\_INTERNAL\_PTS}}, \sphinxcode{\sphinxupquote{GO\_EXTERNAL\_PTS}} and \sphinxcode{\sphinxupquote{GO\_ALL\_PTS}}. These
may be understood by considering the following diagram of an example
model configuration:

\noindent\sphinxincludegraphics{{grids_SW_stagger}.png}

\sphinxAtStartPar
\sphinxcode{\sphinxupquote{GO\_INTERNAL\_PTS}} are then those points that are within the Model
domain (fuscia box), \sphinxcode{\sphinxupquote{GO\_EXTERNAL\_PTS}} are those outside the domain and
\sphinxcode{\sphinxupquote{GO\_ALL\_PTS}} encompasses all grid points in the model. The chosen value
is specified in the kernel\sphinxhyphen{}meta data like so:

\begin{sphinxVerbatim}[commandchars=\\\{\}]
\PYG{k+kt}{integer}\PYG{+w}{ }\PYG{k+kd}{::}\PYG{+w}{ }\PYG{n}{iterates\PYGZus{}over}\PYG{+w}{ }\PYG{o}{=}\PYG{+w}{ }\PYG{n}{GO\PYGZus{}INTERNAL\PYGZus{}PTS}
\end{sphinxVerbatim}

\sphinxAtStartPar
A user can use a config file (see {\hyperref[\detokenize{configuration:configuration}]{\sphinxcrossref{\DUrole{std,std-ref}{Configuration}}}}) to add
additional iteration spaces to PSyclone.


\subsubsection{Index Offset}
\label{\detokenize{gocean1p0:index-offset}}\label{\detokenize{gocean1p0:gocean-index-offset}}
\sphinxAtStartPar
The third element of kernel metadata, \sphinxcode{\sphinxupquote{INDEX\_OFFSET}}, specifies the
index\sphinxhyphen{}offset that the kernel uses. This is the same quantity as
supplied to the grid constructor (see the {\hyperref[\detokenize{gocean1p0:gocean-grid}]{\sphinxcrossref{\DUrole{std,std-ref}{Grid}}}}
Section for a description).

\sphinxAtStartPar
The GOcean 1.0 API supports two different offset schemes;
\sphinxcode{\sphinxupquote{GO\_OFFSET\_NE}}, \sphinxcode{\sphinxupquote{GO\_OFFSET\_SW}}. The scheme used by a kernel is specified
in the metadata as, e.g.:

\begin{sphinxVerbatim}[commandchars=\\\{\}]
\PYG{k+kt}{integer}\PYG{+w}{ }\PYG{k+kd}{::}\PYG{+w}{ }\PYG{n}{index\PYGZus{}offset}\PYG{+w}{ }\PYG{o}{=}\PYG{+w}{ }\PYG{n}{GO\PYGZus{}OFFSET\PYGZus{}NE}
\end{sphinxVerbatim}

\sphinxAtStartPar
Currently all kernels used in an application must use the same offset
scheme which must also be the same as passed to the grid constructor.


\subsubsection{Procedure}
\label{\detokenize{gocean1p0:procedure}}
\sphinxAtStartPar
The fourth and final type of metadata is \sphinxcode{\sphinxupquote{procedure}} metadata. This
specifies the name of the Kernel Fortran subroutine that this metadata
describes.

\sphinxAtStartPar
For example:

\begin{sphinxVerbatim}[commandchars=\\\{\}]
\PYG{k}{procedure}\PYG{+w}{ }\PYG{k+kd}{::}\PYG{+w}{ }\PYG{n}{my\PYGZus{}kernel\PYGZus{}code}
\end{sphinxVerbatim}


\subsection{Subroutine}
\label{\detokenize{gocean1p0:subroutine}}

\subsubsection{Rules}
\label{\detokenize{gocean1p0:rules}}
\sphinxAtStartPar
Kernel arguments follow a set of rules which have been specified for
the GOcean 1.0 API. These rules are encoded in the \sphinxcode{\sphinxupquote{gen\_code()}}
method of the \sphinxcode{\sphinxupquote{GOKern}} class in the \sphinxcode{\sphinxupquote{gocean1p0.py}} file. The
rules, along with PSyclone’s naming conventions, are:
\begin{enumerate}
\sphinxsetlistlabels{\arabic}{enumi}{enumii}{}{)}%
\item {} 
\sphinxAtStartPar
Every kernel has the indices of the current grid point as the first two arguments, \sphinxcode{\sphinxupquote{i}} and \sphinxcode{\sphinxupquote{j}}. These are integers and have intent \sphinxcode{\sphinxupquote{in}}.

\item {} 
\sphinxAtStartPar
For each field/scalar/grid property in the order specified by the meta\_args metadata:
\begin{enumerate}
\sphinxsetlistlabels{\arabic}{enumii}{enumiii}{}{)}%
\item {} 
\sphinxAtStartPar
For a field; the field array itself. A field array is a real array of kind \sphinxcode{\sphinxupquote{go\_wp}} and rank two. The first rank is longitude (\sphinxstyleemphasis{x}) and the second latitude (\sphinxstyleemphasis{y}).

\item {} 
\sphinxAtStartPar
For a scalar; the variable itself. A real scalar is of kind \sphinxcode{\sphinxupquote{go\_wp}}.

\item {} 
\sphinxAtStartPar
For a grid property; the array or variable (see the earlier table) containing the specified property.

\end{enumerate}

\end{enumerate}

\begin{sphinxadmonition}{note}{Note:}
\sphinxAtStartPar
Grid properties are not passed from the Algorithm
Layer. PSyclone generates the necessary lookups in the PSy Layer
and includes the resulting references in the arguments passed to
the kernel.
\end{sphinxadmonition}

\sphinxAtStartPar
As an example, consider the \sphinxcode{\sphinxupquote{bc\_solid\_u}} kernel that is used in the
\sphinxcode{\sphinxupquote{gocean2d}} program shown earlier. The metadata for this kernel is:

\begin{sphinxVerbatim}[commandchars=\\\{\}]
\PYG{+w}{ }\PYG{k}{type}\PYG{p}{,}\PYG{+w}{ }\PYG{k}{extends}\PYG{p}{(}\PYG{n}{kernel\PYGZus{}type}\PYG{p}{)}\PYG{+w}{ }\PYG{k+kd}{::}\PYG{+w}{ }\PYG{n}{bc\PYGZus{}solid\PYGZus{}u}
\PYG{+w}{   }\PYG{k}{type}\PYG{p}{(}\PYG{n}{go\PYGZus{}arg}\PYG{p}{)}\PYG{p}{,}\PYG{+w}{ }\PYG{k}{dimension}\PYG{p}{(}\PYG{l+m+mi}{2}\PYG{p}{)}\PYG{+w}{ }\PYG{k+kd}{::}\PYG{+w}{ }\PYG{n}{meta\PYGZus{}args}\PYG{+w}{ }\PYG{o}{=}\PYG{+w}{  }\PYG{p}{\PYGZam{}}
\PYG{+w}{        }\PYG{p}{(}\PYG{o}{/}\PYG{+w}{ }\PYG{n}{go\PYGZus{}arg}\PYG{p}{(}\PYG{n}{GO\PYGZus{}WRITE}\PYG{p}{,}\PYG{+w}{ }\PYG{n}{GO\PYGZus{}CU}\PYG{p}{,}\PYG{+w}{ }\PYG{n}{GO\PYGZus{}POINTWISE}\PYG{p}{)}\PYG{p}{,}\PYG{+w}{      }\PYG{p}{\PYGZam{}}
\PYG{+w}{           }\PYG{n}{go\PYGZus{}arg}\PYG{p}{(}\PYG{n}{GO\PYGZus{}READ}\PYG{p}{,}\PYG{+w}{      }\PYG{n}{GO\PYGZus{}GRID\PYGZus{}MASK\PYGZus{}T}\PYG{p}{)}\PYG{+w}{     }\PYG{p}{\PYGZam{}}
\PYG{+w}{         }\PYG{o}{/}\PYG{p}{)}

\PYG{+w}{   }\PYG{c}{!\PYGZgt{} This is a boundary\PYGZhy{}conditions kernel and therefore}
\PYG{+w}{   }\PYG{c}{!! acts on all points of the domain rather than just}
\PYG{+w}{   }\PYG{c}{!! those that are internal}
\PYG{+w}{   }\PYG{k+kt}{integer}\PYG{+w}{ }\PYG{k+kd}{::}\PYG{+w}{ }\PYG{n}{ITERATES\PYGZus{}OVER}\PYG{+w}{ }\PYG{o}{=}\PYG{+w}{ }\PYG{n}{GO\PYGZus{}ALL\PYGZus{}PTS}

\PYG{+w}{   }\PYG{k+kt}{integer}\PYG{+w}{ }\PYG{k+kd}{::}\PYG{+w}{ }\PYG{n}{index\PYGZus{}offset}\PYG{+w}{ }\PYG{o}{=}\PYG{+w}{ }\PYG{n}{GO\PYGZus{}OFFSET\PYGZus{}NE}

\PYG{k}{contains}
\PYG{k}{  }\PYG{k}{procedure}\PYG{p}{,}\PYG{+w}{ }\PYG{k}{nopass}\PYG{+w}{ }\PYG{k+kd}{::}\PYG{+w}{ }\PYG{n}{code}\PYG{+w}{ }\PYG{o}{=}\PYG{o}{\PYGZgt{}}\PYG{+w}{ }\PYG{n}{bc\PYGZus{}solid\PYGZus{}u\PYGZus{}code}
\PYG{k}{end }\PYG{k}{type }\PYG{n}{bc\PYGZus{}solid\PYGZus{}u}
\end{sphinxVerbatim}

\sphinxAtStartPar
The interface to the subroutine containing the implementation of this
kernel is:

\begin{sphinxVerbatim}[commandchars=\\\{\}]
\PYG{k}{subroutine }\PYG{n}{bc\PYGZus{}solid\PYGZus{}u\PYGZus{}code}\PYG{p}{(}\PYG{n}{ji}\PYG{p}{,}\PYG{+w}{ }\PYG{n}{jj}\PYG{p}{,}\PYG{+w}{ }\PYG{n}{ua}\PYG{p}{,}\PYG{+w}{ }\PYG{n}{tmask}\PYG{p}{)}
\PYG{+w}{  }\PYG{k+kt}{integer}\PYG{p}{,}\PYG{+w}{                  }\PYG{k}{intent}\PYG{p}{(}\PYG{n}{in}\PYG{p}{)}\PYG{+w}{    }\PYG{k+kd}{::}\PYG{+w}{ }\PYG{n}{ji}\PYG{p}{,}\PYG{+w}{ }\PYG{n}{jj}
\PYG{+w}{  }\PYG{k+kt}{integer}\PYG{p}{,}\PYG{+w}{  }\PYG{k}{dimension}\PYG{p}{(}\PYG{p}{:}\PYG{p}{,}\PYG{p}{:}\PYG{p}{)}\PYG{p}{,}\PYG{+w}{ }\PYG{k}{intent}\PYG{p}{(}\PYG{n}{in}\PYG{p}{)}\PYG{+w}{    }\PYG{k+kd}{::}\PYG{+w}{ }\PYG{n}{tmask}
\PYG{+w}{  }\PYG{k+kt}{real}\PYG{p}{(}\PYG{n}{wp}\PYG{p}{)}\PYG{p}{,}\PYG{+w}{ }\PYG{k}{dimension}\PYG{p}{(}\PYG{p}{:}\PYG{p}{,}\PYG{p}{:}\PYG{p}{)}\PYG{p}{,}\PYG{+w}{ }\PYG{k}{intent}\PYG{p}{(}\PYG{n}{inout}\PYG{p}{)}\PYG{+w}{ }\PYG{k+kd}{::}\PYG{+w}{ }\PYG{n}{ua}
\end{sphinxVerbatim}

\sphinxAtStartPar
As described above, the first two arguments to this subroutine specify
the grid\sphinxhyphen{}point at which the computation is to be performed. The third
argument is the field that this kernel updates and the fourth argument
is the T\sphinxhyphen{}point mask. The latter is a property of the grid and is
provided to the kernel call from the PSy Layer.

\sphinxAtStartPar
Comparing this interface definition with the use of the kernel in the
Invoke call:

\begin{sphinxVerbatim}[commandchars=\\\{\}]
\PYG{k}{call }\PYG{n}{invoke}\PYG{+w}{ }\PYG{p}{(}\PYG{+w}{ }\PYG{p}{.}\PYG{p}{.}\PYG{p}{.}\PYG{p}{,}\PYG{+w}{            }\PYG{p}{\PYGZam{}}
\PYG{+w}{              }\PYG{n}{bc\PYGZus{}solid\PYGZus{}u}\PYG{p}{(}\PYG{n}{ua}\PYG{p}{)}\PYG{p}{,}\PYG{+w}{ }\PYG{p}{\PYGZam{}}
\PYG{+w}{              }\PYG{p}{.}\PYG{p}{.}\PYG{p}{.}\PYG{+w}{ }\PYG{p}{)}
\end{sphinxVerbatim}

\sphinxAtStartPar
we see that in the Algorithm Layer the user need only provide the
field(s) (and possibly scalars) that a kernel operates on. The index
of the grid point and any grid properties are provided in the
(generated) PSy Layer where the kernel subroutine proper is called.


\section{Built\sphinxhyphen{}ins}
\label{\detokenize{gocean1p0:built-ins}}
\sphinxAtStartPar
The GOcean 1.0 API does not support any built\sphinxhyphen{}in operations.


\section{Conventions}
\label{\detokenize{gocean1p0:conventions}}
\sphinxAtStartPar
The GOcean 1.0 API kernel code conforms to the PSyclone Fortran naming
conventions (see {\hyperref[\detokenize{fortran_naming_conventions:fortran-naming}]{\sphinxcrossref{\DUrole{std,std-ref}{Fortran Naming Conventions}}}}). However, PSyclone’s support
for the GOcean 1.0 API does not rely on this convention.

\sphinxAtStartPar
The contents of the kernel metadata is usually declared private but this
does not affect PSyclone.

\sphinxAtStartPar
Finally, the \sphinxcode{\sphinxupquote{procedure}} metadata (located within the kernel
metadata) usually has \sphinxcode{\sphinxupquote{nopass}} specified but again this is ignored
by PSyclone.


\section{Configuration}
\label{\detokenize{gocean1p0:configuration}}\label{\detokenize{gocean1p0:gocean-configuration}}
\sphinxAtStartPar
The configuration file (see {\hyperref[\detokenize{configuration:configuration}]{\sphinxcrossref{\DUrole{std,std-ref}{Configuration}}}}) used by PSyclone
can contain GOcean 1.0 specific options. For example, after the
default section the GOcean 1.0 specific section looks like this:

\begin{sphinxVerbatim}[commandchars=\\\{\}]
[gocean]
iteration\PYGZhy{}spaces=offset\PYGZus{}sw:ct:test\PYGZus{}only:1:2:3:4
                 offset\PYGZus{}sw:ct:internal\PYGZus{}ns\PYGZus{}halo:\PYGZob{}start\PYGZcb{}\PYGZhy{}1:\PYGZob{}stop\PYGZcb{}+1:\PYGZob{}start\PYGZcb{}:\PYGZob{}stop\PYGZcb{}
\end{sphinxVerbatim}

\sphinxAtStartPar
The supported keys are listed in the next section.


\subsection{Iteration\sphinxhyphen{}spaces}
\label{\detokenize{gocean1p0:iteration-spaces}}\label{\detokenize{gocean1p0:gocean-configuration-iteration-spaces}}
\sphinxAtStartPar
This section lists additional iteration spaces that can be used in a kernel
metadata declaration to allow PSyclone to create a loop with different
loop boundaries. Each line of the \sphinxcode{\sphinxupquote{iteration\sphinxhyphen{}spaces}} declaration
contains 7 values, separated by ‘:’. The fields are:


\begin{savenotes}\sphinxattablestart
\sphinxthistablewithglobalstyle
\centering
\begin{tabulary}{\linewidth}[t]{TTT}
\sphinxtoprule
\sphinxstyletheadfamily 
\sphinxAtStartPar
Field
&\sphinxstyletheadfamily 
\sphinxAtStartPar
Description
&\sphinxstyletheadfamily 
\sphinxAtStartPar
Details
\\
\sphinxmidrule
\sphinxtableatstartofbodyhook
\sphinxAtStartPar
1
&
\sphinxAtStartPar
Index Offset
&
\sphinxAtStartPar
See {\hyperref[\detokenize{gocean1p0:gocean-index-offset}]{\sphinxcrossref{\DUrole{std,std-ref}{Index Offset}}}}.
\\
\sphinxhline
\sphinxAtStartPar
2
&
\sphinxAtStartPar
grid\sphinxhyphen{}point types
&
\sphinxAtStartPar
See {\hyperref[\detokenize{gocean1p0:gocean-grid-point-type}]{\sphinxcrossref{\DUrole{std,std-ref}{Grid point types}}}}.
\\
\sphinxhline
\sphinxAtStartPar
3
&
\sphinxAtStartPar
Iterates Over
&
\sphinxAtStartPar
See {\hyperref[\detokenize{gocean1p0:gocean-iterates-over}]{\sphinxcrossref{\DUrole{std,std-ref}{Iterates Over}}}}.
\\
\sphinxhline
\sphinxAtStartPar
4
&
\sphinxAtStartPar
Start index of outer loop
&
\sphinxAtStartPar
Start index of North\sphinxhyphen{}South loop.
\\
\sphinxhline
\sphinxAtStartPar
5
&
\sphinxAtStartPar
End index of outer loop
&
\sphinxAtStartPar
End index of North\sphinxhyphen{}South loop.
\\
\sphinxhline
\sphinxAtStartPar
6
&
\sphinxAtStartPar
Start index of inner loop
&
\sphinxAtStartPar
Start index of East\sphinxhyphen{}West loop.
\\
\sphinxhline
\sphinxAtStartPar
7
&
\sphinxAtStartPar
End index of inner loop
&
\sphinxAtStartPar
End index of East\sphinxhyphen{}West loop.
\\
\sphinxbottomrule
\end{tabulary}
\sphinxtableafterendhook\par
\sphinxattableend\end{savenotes}

\sphinxAtStartPar
Two special variables can be used in an iteration space: \sphinxcode{\sphinxupquote{\{start\}}} and \sphinxcode{\sphinxupquote{\{stop\}}}.
These values will be replaced by PSyclone with the correct loop boundaries for the
inner points of a grid (i.e. the non\sphinxhyphen{}halo area). This means that the depth\sphinxhyphen{}1 halo
region can be specified using \sphinxcode{\sphinxupquote{\{start\}\sphinxhyphen{}1}} and \sphinxcode{\sphinxupquote{\{stop\}+1}}.

\sphinxAtStartPar
For example, given the iteration\sphinxhyphen{}spaces declaration above, a kernel declared with
\sphinxcode{\sphinxupquote{iterates\_over=internal\_ns\_halo}} for a field type \sphinxcode{\sphinxupquote{ct}} and index offset
\sphinxcode{\sphinxupquote{offset\_sw}} would create the following loop boundaries:

\begin{sphinxVerbatim}[commandchars=\\\{\}]
\PYG{k}{DO }\PYG{n}{j}\PYG{o}{=}\PYG{l+m+mi}{2}\PYG{o}{\PYGZhy{}}\PYG{l+m+mi}{1}\PYG{p}{,}\PYG{n}{jstop}\PYG{o}{+}\PYG{l+m+mi}{1}
\PYG{+w}{  }\PYG{k}{DO }\PYG{n}{i}\PYG{o}{=}\PYG{l+m+mi}{2}\PYG{p}{,}\PYG{n}{istop}
\PYG{+w}{    }\PYG{k}{CALL}\PYG{+w}{ }\PYG{p}{(}\PYG{n}{i}\PYG{p}{,}\PYG{+w}{ }\PYG{n}{j}\PYG{p}{,}\PYG{+w}{ }\PYG{p}{.}\PYG{p}{.}\PYG{p}{.}\PYG{p}{)}
\PYG{+w}{  }\PYG{k}{END }\PYG{k}{DO}
\PYG{k}{END }\PYG{k}{DO}
\end{sphinxVerbatim}

\begin{sphinxadmonition}{warning}{Warning:}
\sphinxAtStartPar
With user defined iteration spaces it is possible that PSyclone will create code
that does not compile: if you specify syntactically correct, but semantically
incorrect boundary definitions, the PSyclone internal tests will accept the new
iteration space, but the compiler will not. For example
if one of the loop boundaries contains the name of a variable that is not defined,
compilation will fail. It is the responsibility of the user to make sure that
valid loop boundaries are specified in a new iteration space definition.
\end{sphinxadmonition}


\subsection{Grid Properties}
\label{\detokenize{gocean1p0:grid-properties}}\label{\detokenize{gocean1p0:gocean-configuration-grid-properties}}
\sphinxAtStartPar
Various grid properties can be specified as parameters to a kernel.
The actual names and meaning of these properties depend on the
infrastructure library used. By default PSyclone provides settings
for the dl\_esm\_inf infrastructure library. But the user or a library
developer can change or add definitions to the configuration file as
required.

\sphinxAtStartPar
The grid properties are specified as values for the key
\sphinxcode{\sphinxupquote{grid\sphinxhyphen{}properties}}. They consist of three entries, separated by “:”.
\begin{itemize}
\item {} 
\sphinxAtStartPar
The first entry is the name of the property as used in kernel metadata.

\item {} 
\sphinxAtStartPar
The next entry is the way of dereferencing the corresponding value in
Fortran. The expression \sphinxcode{\sphinxupquote{\{0\}}} is replaced with the field name that
is used. Note that any \sphinxcode{\sphinxupquote{\%}} must be replaced with \sphinxcode{\sphinxupquote{\%\%}} (due to the
way Python reads in configuration files).

\item {} 
\sphinxAtStartPar
The last entry specifies whether the value is an \sphinxcode{\sphinxupquote{array}} or a \sphinxcode{\sphinxupquote{scalar}}.

\end{itemize}

\sphinxAtStartPar
Below an excerpt from the configuration file that is distributed
with PSyclone:

\begin{sphinxVerbatim}[commandchars=\\\{\}]
grid\PYGZhy{}properties = go\PYGZus{}grid\PYGZus{}xstop: \PYGZob{}0\PYGZcb{}\PYGZpc{}\PYGZpc{}grid\PYGZpc{}\PYGZpc{}subdomain\PYGZpc{}\PYGZpc{}internal\PYGZpc{}\PYGZpc{}xstop: scalar,
                  go\PYGZus{}grid\PYGZus{}ystop: \PYGZob{}0\PYGZcb{}\PYGZpc{}\PYGZpc{}grid\PYGZpc{}\PYGZpc{}subdomain\PYGZpc{}\PYGZpc{}internal\PYGZpc{}\PYGZpc{}ystop: scalar,
                  go\PYGZus{}grid\PYGZus{}data: \PYGZob{}0\PYGZcb{}\PYGZpc{}\PYGZpc{}data: array,
                  go\PYGZus{}grid\PYGZus{}internal\PYGZus{}inner\PYGZus{}stop: \PYGZob{}0\PYGZcb{}\PYGZpc{}\PYGZpc{}internal\PYGZpc{}\PYGZpc{}xstop: scalar,
                  go\PYGZus{}grid\PYGZus{}internal\PYGZus{}outer\PYGZus{}stop: \PYGZob{}0\PYGZcb{}\PYGZpc{}\PYGZpc{}internal\PYGZpc{}\PYGZpc{}ystop: scalar,
                  go\PYGZus{}grid\PYGZus{}whole\PYGZus{}inner\PYGZus{}stop: \PYGZob{}0\PYGZcb{}\PYGZpc{}\PYGZpc{}whole\PYGZpc{}\PYGZpc{}xstop: scalar,
                  go\PYGZus{}grid\PYGZus{}whole\PYGZus{}outer\PYGZus{}stop: \PYGZob{}0\PYGZcb{}\PYGZpc{}\PYGZpc{}whole\PYGZpc{}\PYGZpc{}ystop: scalar,
...
\end{sphinxVerbatim}

\sphinxAtStartPar
Most of the property names can be set arbitrarily by the user (to match whatever
infrastructure library is being used), but PSyclone relies on a
small number of properties that must be defined with the right name:


\begin{savenotes}\sphinxattablestart
\sphinxthistablewithglobalstyle
\centering
\begin{tabular}[t]{*{2}{\X{1}{2}}}
\sphinxtoprule
\sphinxtableatstartofbodyhook
\sphinxAtStartPar
Key
&
\sphinxAtStartPar
Description
\\
\sphinxhline
\sphinxAtStartPar
go\_grid\_data
&
\sphinxAtStartPar
This property gives access to the raw 2d\sphinxhyphen{}field.
\\
\sphinxhline
\sphinxAtStartPar
go\_grid\_xstop,
go\_grid\_ystop
&
\sphinxAtStartPar
These values specify the upper loop boundary
when computing the constant loop boundaries.
\\
\sphinxhline
\begin{DUlineblock}{0em}
\item[] go\_grid\_\{internal,whole\}
\item[] \_\{inner,outer\}\_\{start,stop\}
\end{DUlineblock}
&
\sphinxAtStartPar
These eight values are required to
specify the loop boundaries depending on the
field space.
\\
\sphinxhline
\sphinxAtStartPar
go\_grid\_nx,
go\_grid\_ny
&
\sphinxAtStartPar
These properties are only required when OpenCL
is enabled. They specify the overall array
size (including any padding that the
infrastructure library might implement).
\\
\sphinxbottomrule
\end{tabular}
\sphinxtableafterendhook\par
\sphinxattableend\end{savenotes}


\subsection{Debug Mode}
\label{\detokenize{gocean1p0:debug-mode}}
\sphinxAtStartPar
The GOcean configuration also includes a boolean parameter to enable or disable
the generation of additional code which may impact performance but is useful
for debugging the application. By default it is set to False, but it can be
changed by updating the following line in the configuration file:

\begin{sphinxVerbatim}[commandchars=\\\{\}]
[gocean]
DEBUG\PYGZus{}MODE = true
\end{sphinxVerbatim}

\sphinxAtStartPar
Currently, only the OpenCL Invokes generate additional debugging code.


\section{Transformations}
\label{\detokenize{gocean1p0:transformations}}
\sphinxAtStartPar
In this section we describe the transformations that are specific to
the GOcean 1.0 API. For an overview of transformations in general see
{\hyperref[\detokenize{transformations:transformations}]{\sphinxcrossref{\DUrole{std,std-ref}{Transformations}}}}.


\begin{fulllineitems}

\pysigstartsignatures
\pysigline{\sphinxbfcode{\sphinxupquote{class\DUrole{w,w}{  }}}\sphinxcode{\sphinxupquote{psyclone.domain.gocean.transformations.}}\sphinxbfcode{\sphinxupquote{GOceanExtractTrans}}}
\pysigstopsignatures
\sphinxAtStartPar
GOcean1.0 API application of ExtractTrans transformation     to extract code into a stand\sphinxhyphen{}alone program. For example:

\begin{sphinxVerbatim}[commandchars=\\\{\}]
\PYG{g+gp}{\PYGZgt{}\PYGZgt{}\PYGZgt{} }\PYG{k+kn}{from}\PYG{+w}{ }\PYG{n+nn}{psyclone}\PYG{n+nn}{.}\PYG{n+nn}{parse}\PYG{n+nn}{.}\PYG{n+nn}{algorithm}\PYG{+w}{ }\PYG{k+kn}{import} \PYG{n}{parse}
\PYG{g+gp}{\PYGZgt{}\PYGZgt{}\PYGZgt{} }\PYG{k+kn}{from}\PYG{+w}{ }\PYG{n+nn}{psyclone}\PYG{n+nn}{.}\PYG{n+nn}{psyGen}\PYG{+w}{ }\PYG{k+kn}{import} \PYG{n}{PSyFactory}
\PYG{g+gp}{\PYGZgt{}\PYGZgt{}\PYGZgt{}}
\PYG{g+gp}{\PYGZgt{}\PYGZgt{}\PYGZgt{} }\PYG{n}{API} \PYG{o}{=} \PYG{l+s+s2}{\PYGZdq{}}\PYG{l+s+s2}{gocean}\PYG{l+s+s2}{\PYGZdq{}}
\PYG{g+gp}{\PYGZgt{}\PYGZgt{}\PYGZgt{} }\PYG{n}{FILENAME} \PYG{o}{=} \PYG{l+s+s2}{\PYGZdq{}}\PYG{l+s+s2}{shallow\PYGZus{}alg.f90}\PYG{l+s+s2}{\PYGZdq{}}
\PYG{g+gp}{\PYGZgt{}\PYGZgt{}\PYGZgt{} }\PYG{n}{ast}\PYG{p}{,} \PYG{n}{invokeInfo} \PYG{o}{=} \PYG{n}{parse}\PYG{p}{(}\PYG{n}{FILENAME}\PYG{p}{,} \PYG{n}{api}\PYG{o}{=}\PYG{n}{API}\PYG{p}{)}
\PYG{g+gp}{\PYGZgt{}\PYGZgt{}\PYGZgt{} }\PYG{n}{psy} \PYG{o}{=} \PYG{n}{PSyFactory}\PYG{p}{(}\PYG{n}{API}\PYG{p}{,} \PYG{n}{distributed\PYGZus{}memory}\PYG{o}{=}\PYG{k+kc}{False}\PYG{p}{)}\PYG{o}{.}\PYG{n}{create}\PYG{p}{(}\PYG{n}{invoke\PYGZus{}info}\PYG{p}{)}
\PYG{g+gp}{\PYGZgt{}\PYGZgt{}\PYGZgt{} }\PYG{n}{schedule} \PYG{o}{=} \PYG{n}{psy}\PYG{o}{.}\PYG{n}{invokes}\PYG{o}{.}\PYG{n}{get}\PYG{p}{(}\PYG{l+s+s1}{\PYGZsq{}}\PYG{l+s+s1}{invoke\PYGZus{}0}\PYG{l+s+s1}{\PYGZsq{}}\PYG{p}{)}\PYG{o}{.}\PYG{n}{schedule}
\PYG{g+gp}{\PYGZgt{}\PYGZgt{}\PYGZgt{}}
\PYG{g+gp}{\PYGZgt{}\PYGZgt{}\PYGZgt{} }\PYG{k+kn}{from}\PYG{+w}{ }\PYG{n+nn}{psyclone}\PYG{n+nn}{.}\PYG{n+nn}{domain}\PYG{n+nn}{.}\PYG{n+nn}{gocean}\PYG{n+nn}{.}\PYG{n+nn}{transformations}\PYG{+w}{ }\PYG{k+kn}{import} \PYG{n}{GOceanExtractTrans}
\PYG{g+gp}{\PYGZgt{}\PYGZgt{}\PYGZgt{} }\PYG{n}{etrans} \PYG{o}{=} \PYG{n}{GOceanExtractTrans}\PYG{p}{(}\PYG{p}{)}
\PYG{g+gp}{\PYGZgt{}\PYGZgt{}\PYGZgt{}}
\PYG{g+gp}{\PYGZgt{}\PYGZgt{}\PYGZgt{} }\PYG{c+c1}{\PYGZsh{} Apply GOceanExtractTrans transformation to selected Nodes}
\PYG{g+gp}{\PYGZgt{}\PYGZgt{}\PYGZgt{} }\PYG{n}{etrans}\PYG{o}{.}\PYG{n}{apply}\PYG{p}{(}\PYG{n}{schedule}\PYG{o}{.}\PYG{n}{children}\PYG{p}{[}\PYG{l+m+mi}{0}\PYG{p}{]}\PYG{p}{)}
\PYG{g+gp}{\PYGZgt{}\PYGZgt{}\PYGZgt{} }\PYG{n+nb}{print}\PYG{p}{(}\PYG{n}{schedule}\PYG{o}{.}\PYG{n}{view}\PYG{p}{(}\PYG{p}{)}\PYG{p}{)}
\end{sphinxVerbatim}


\begin{fulllineitems}

\pysigstartsignatures
\pysiglinewithargsret{\sphinxbfcode{\sphinxupquote{apply}}}{\sphinxparam{\DUrole{n,n}{nodes}}\sphinxparamcomma \sphinxparam{\DUrole{n,n}{options}\DUrole{o,o}{=}\DUrole{default_value}{None}}}{}
\pysigstopsignatures
\sphinxAtStartPar
Apply this transformation to a subset of the nodes within a
schedule \sphinxhyphen{} i.e. enclose the specified Nodes in the schedule within
a single PSyData region. Note that this implementation just calls
the base class, it is only added here to provide the documentation
for this function, since it accepts different options
to the base class (e.g. create\_driver, which is passed to the
ExtractNode instance that will be inserted.).
\begin{quote}\begin{description}
\sphinxlineitem{Parameters}\begin{itemize}
\item {} 
\sphinxAtStartPar
\sphinxstyleliteralstrong{\sphinxupquote{nodes}} (\sphinxcode{\sphinxupquote{psyclone.psyir.nodes.Node}} or list of                      \sphinxcode{\sphinxupquote{psyclone.psyir.nodes.Node}}) \textendash{} can be a single node or a list of nodes.

\item {} 
\sphinxAtStartPar
\sphinxstyleliteralstrong{\sphinxupquote{options}} (\sphinxstyleliteralemphasis{\sphinxupquote{Optional}}\sphinxstyleliteralemphasis{\sphinxupquote{{[}}}\sphinxstyleliteralemphasis{\sphinxupquote{Dict}}\sphinxstyleliteralemphasis{\sphinxupquote{{[}}}\sphinxhref{https://docs.python.org/3/library/stdtypes.html\#str}{\sphinxstyleliteralemphasis{\sphinxupquote{str}}}\sphinxstyleliteralemphasis{\sphinxupquote{, }}\sphinxstyleliteralemphasis{\sphinxupquote{Any}}\sphinxstyleliteralemphasis{\sphinxupquote{{]}}}\sphinxstyleliteralemphasis{\sphinxupquote{{]}}}) \textendash{} a dictionary with options for transformations.

\item {} 
\sphinxAtStartPar
\sphinxstyleliteralstrong{\sphinxupquote{options}}\sphinxstyleliteralstrong{\sphinxupquote{{[}}}\sphinxstyleliteralstrong{\sphinxupquote{"prefix"}}\sphinxstyleliteralstrong{\sphinxupquote{{]}}} (\sphinxhref{https://docs.python.org/3/library/stdtypes.html\#str}{\sphinxstyleliteralemphasis{\sphinxupquote{str}}}) \textendash{} a prefix to use for the PSyData module             name (\sphinxcode{\sphinxupquote{prefix\_psy\_data\_mod}}) and the PSyDataType             (\sphinxcode{\sphinxupquote{prefix\_PSyDataType}}) \sphinxhyphen{} a “\_” will be added automatically.             It defaults to “extract”, resulting in e.g.             \sphinxcode{\sphinxupquote{extract\_psy\_data\_mod}}.

\item {} 
\sphinxAtStartPar
\sphinxstyleliteralstrong{\sphinxupquote{options}}\sphinxstyleliteralstrong{\sphinxupquote{{[}}}\sphinxstyleliteralstrong{\sphinxupquote{"create\_driver"}}\sphinxstyleliteralstrong{\sphinxupquote{{]}}} (\sphinxhref{https://docs.python.org/3/library/functions.html\#bool}{\sphinxstyleliteralemphasis{\sphinxupquote{bool}}}) \textendash{} whether or not to create a             driver program at code\sphinxhyphen{}generation time. If set, the driver will             be created in the current working directory with the name             “driver\sphinxhyphen{}MODULE\sphinxhyphen{}REGION.f90” where MODULE and REGION will be the             corresponding values for this region. Defaults to False.

\item {} 
\sphinxAtStartPar
\sphinxstyleliteralstrong{\sphinxupquote{options}}\sphinxstyleliteralstrong{\sphinxupquote{{[}}}\sphinxstyleliteralstrong{\sphinxupquote{"region\_name"}}\sphinxstyleliteralstrong{\sphinxupquote{{]}}} (\sphinxstyleliteralemphasis{\sphinxupquote{(}}\sphinxhref{https://docs.python.org/3/library/stdtypes.html\#str}{\sphinxstyleliteralemphasis{\sphinxupquote{str}}}\sphinxstyleliteralemphasis{\sphinxupquote{,}}\sphinxhref{https://docs.python.org/3/library/stdtypes.html\#str}{\sphinxstyleliteralemphasis{\sphinxupquote{str}}}\sphinxstyleliteralemphasis{\sphinxupquote{)}}) \textendash{} an optional name to             use for this PSyData area, provided as a 2\sphinxhyphen{}tuple containing a             location name followed by a local name. The pair of strings             should uniquely identify a region unless aggregate information             is required (and is supported by the runtime library).

\end{itemize}

\end{description}\end{quote}

\end{fulllineitems}



\begin{fulllineitems}

\pysigstartsignatures
\pysiglinewithargsret{\sphinxbfcode{\sphinxupquote{validate}}}{\sphinxparam{\DUrole{n,n}{node\_list}}\sphinxparamcomma \sphinxparam{\DUrole{n,n}{options}\DUrole{o,o}{=}\DUrole{default_value}{None}}}{}
\pysigstopsignatures
\sphinxAtStartPar
Perform GOcean1.0 API specific validation checks before applying
the transformation.
\begin{quote}\begin{description}
\sphinxlineitem{Parameters}\begin{itemize}
\item {} 
\sphinxAtStartPar
\sphinxstyleliteralstrong{\sphinxupquote{node\_list}} (list of \sphinxcode{\sphinxupquote{psyclone.psyir.nodes.Node}}) \textendash{} the list of Node(s) we are checking.

\item {} 
\sphinxAtStartPar
\sphinxstyleliteralstrong{\sphinxupquote{options}} (\sphinxstyleliteralemphasis{\sphinxupquote{Optional}}\sphinxstyleliteralemphasis{\sphinxupquote{{[}}}\sphinxstyleliteralemphasis{\sphinxupquote{Dict}}\sphinxstyleliteralemphasis{\sphinxupquote{{[}}}\sphinxhref{https://docs.python.org/3/library/stdtypes.html\#str}{\sphinxstyleliteralemphasis{\sphinxupquote{str}}}\sphinxstyleliteralemphasis{\sphinxupquote{, }}\sphinxstyleliteralemphasis{\sphinxupquote{Any}}\sphinxstyleliteralemphasis{\sphinxupquote{{]}}}\sphinxstyleliteralemphasis{\sphinxupquote{{]}}}) \textendash{} a dictionary with options for transformations.

\item {} 
\sphinxAtStartPar
\sphinxstyleliteralstrong{\sphinxupquote{options}}\sphinxstyleliteralstrong{\sphinxupquote{{[}}}\sphinxstyleliteralstrong{\sphinxupquote{"create\_driver"}}\sphinxstyleliteralstrong{\sphinxupquote{{]}}} (\sphinxhref{https://docs.python.org/3/library/functions.html\#bool}{\sphinxstyleliteralemphasis{\sphinxupquote{bool}}}) \textendash{} whether or not to create a             driver program at code\sphinxhyphen{}generation time. If set, the driver will             be created in the current working directory with the name             “driver\sphinxhyphen{}MODULE\sphinxhyphen{}REGION.f90” where MODULE and REGION will be the             corresponding values for this region. This flag is forwarded to             the ExtractNode. Its default value is False.

\item {} 
\sphinxAtStartPar
\sphinxstyleliteralstrong{\sphinxupquote{options}}\sphinxstyleliteralstrong{\sphinxupquote{{[}}}\sphinxstyleliteralstrong{\sphinxupquote{"region\_name"}}\sphinxstyleliteralstrong{\sphinxupquote{{]}}} (\sphinxstyleliteralemphasis{\sphinxupquote{(}}\sphinxhref{https://docs.python.org/3/library/stdtypes.html\#str}{\sphinxstyleliteralemphasis{\sphinxupquote{str}}}\sphinxstyleliteralemphasis{\sphinxupquote{,}}\sphinxhref{https://docs.python.org/3/library/stdtypes.html\#str}{\sphinxstyleliteralemphasis{\sphinxupquote{str}}}\sphinxstyleliteralemphasis{\sphinxupquote{)}}) \textendash{} an optional name to             use for this data\sphinxhyphen{}extraction region, provided as a 2\sphinxhyphen{}tuple             containing a module name followed by a local name. The pair of             strings should uniquely identify a region unless aggregate             information is required (and is supported by the runtime             library). This option is forwarded to the PSyDataNode (where it             changes the region names) and to the ExtractNode (where it             changes the name of the created output files and the name of the             driver program).

\end{itemize}

\sphinxlineitem{Raises}
\sphinxAtStartPar
\sphinxstyleliteralstrong{\sphinxupquote{TransformationError}} \textendash{} if transformation is applied to an             inner Loop without its parent outer Loop.

\end{description}\end{quote}

\end{fulllineitems}


\end{fulllineitems}



\bigskip\hrule\bigskip



\begin{fulllineitems}

\pysigstartsignatures
\pysigline{\sphinxbfcode{\sphinxupquote{class\DUrole{w,w}{  }}}\sphinxcode{\sphinxupquote{psyclone.domain.gocean.transformations.}}\sphinxbfcode{\sphinxupquote{GOceanLoopFuseTrans}}}
\pysigstopsignatures
\sphinxAtStartPar
GOcean API specialisation of the \sphinxcode{\sphinxupquote{base class}}
in order to fuse two GOcean loops after performing validity checks (e.g.
that the loops are over the same grid\sphinxhyphen{}point type). For example:

\begin{sphinxVerbatim}[commandchars=\\\{\}]
\PYG{g+gp}{\PYGZgt{}\PYGZgt{}\PYGZgt{} }\PYG{k+kn}{from}\PYG{+w}{ }\PYG{n+nn}{psyclone}\PYG{n+nn}{.}\PYG{n+nn}{parse}\PYG{n+nn}{.}\PYG{n+nn}{algorithm}\PYG{+w}{ }\PYG{k+kn}{import} \PYG{n}{parse}
\PYG{g+gp}{\PYGZgt{}\PYGZgt{}\PYGZgt{} }\PYG{k+kn}{from}\PYG{+w}{ }\PYG{n+nn}{psyclone}\PYG{n+nn}{.}\PYG{n+nn}{psyGen}\PYG{+w}{ }\PYG{k+kn}{import} \PYG{n}{PSyFactory}
\PYG{g+gp}{\PYGZgt{}\PYGZgt{}\PYGZgt{} }\PYG{n}{ast}\PYG{p}{,} \PYG{n}{invokeInfo} \PYG{o}{=} \PYG{n}{parse}\PYG{p}{(}\PYG{l+s+s2}{\PYGZdq{}}\PYG{l+s+s2}{shallow\PYGZus{}alg.f90}\PYG{l+s+s2}{\PYGZdq{}}\PYG{p}{)}
\PYG{g+gp}{\PYGZgt{}\PYGZgt{}\PYGZgt{} }\PYG{n}{psy} \PYG{o}{=} \PYG{n}{PSyFactory}\PYG{p}{(}\PYG{l+s+s2}{\PYGZdq{}}\PYG{l+s+s2}{gocean}\PYG{l+s+s2}{\PYGZdq{}}\PYG{p}{)}\PYG{o}{.}\PYG{n}{create}\PYG{p}{(}\PYG{n}{invokeInfo}\PYG{p}{)}
\PYG{g+gp}{\PYGZgt{}\PYGZgt{}\PYGZgt{} }\PYG{n}{schedule} \PYG{o}{=} \PYG{n}{psy}\PYG{o}{.}\PYG{n}{invokes}\PYG{o}{.}\PYG{n}{get}\PYG{p}{(}\PYG{l+s+s1}{\PYGZsq{}}\PYG{l+s+s1}{invoke\PYGZus{}0}\PYG{l+s+s1}{\PYGZsq{}}\PYG{p}{)}\PYG{o}{.}\PYG{n}{schedule}
\PYG{g+gp}{\PYGZgt{}\PYGZgt{}\PYGZgt{} }\PYG{n+nb}{print}\PYG{p}{(}\PYG{n}{schedule}\PYG{o}{.}\PYG{n}{view}\PYG{p}{(}\PYG{p}{)}\PYG{p}{)}
\PYG{g+gp}{\PYGZgt{}\PYGZgt{}\PYGZgt{}}
\PYG{g+gp}{\PYGZgt{}\PYGZgt{}\PYGZgt{} }\PYG{k+kn}{from}\PYG{+w}{ }\PYG{n+nn}{psyclone}\PYG{n+nn}{.}\PYG{n+nn}{transformations}\PYG{+w}{ }\PYG{k+kn}{import} \PYG{n}{GOceanLoopFuseTrans}
\PYG{g+gp}{\PYGZgt{}\PYGZgt{}\PYGZgt{} }\PYG{n}{ftrans} \PYG{o}{=} \PYG{n}{GOceanLoopFuseTrans}\PYG{p}{(}\PYG{p}{)}
\PYG{g+gp}{\PYGZgt{}\PYGZgt{}\PYGZgt{} }\PYG{n}{ftrans}\PYG{o}{.}\PYG{n}{apply}\PYG{p}{(}\PYG{n}{schedule}\PYG{p}{[}\PYG{l+m+mi}{0}\PYG{p}{]}\PYG{p}{,} \PYG{n}{schedule}\PYG{p}{[}\PYG{l+m+mi}{1}\PYG{p}{]}\PYG{p}{)}
\PYG{g+gp}{\PYGZgt{}\PYGZgt{}\PYGZgt{} }\PYG{n+nb}{print}\PYG{p}{(}\PYG{n}{schedule}\PYG{o}{.}\PYG{n}{view}\PYG{p}{(}\PYG{p}{)}\PYG{p}{)}
\end{sphinxVerbatim}


\begin{fulllineitems}

\pysigstartsignatures
\pysiglinewithargsret{\sphinxbfcode{\sphinxupquote{validate}}}{\sphinxparam{\DUrole{n,n}{node1}}\sphinxparamcomma \sphinxparam{\DUrole{n,n}{node2}}\sphinxparamcomma \sphinxparam{\DUrole{n,n}{options}\DUrole{o,o}{=}\DUrole{default_value}{None}}}{}
\pysigstopsignatures
\sphinxAtStartPar
Checks if it is valid to apply the GOceanLoopFuseTrans
transform. It ensures that the fused loops are over
the same grid\sphinxhyphen{}point types, before calling the normal
LoopFuseTrans validation function.
\begin{quote}\begin{description}
\sphinxlineitem{Parameters}\begin{itemize}
\item {} 
\sphinxAtStartPar
\sphinxstyleliteralstrong{\sphinxupquote{node1}} (\sphinxcode{\sphinxupquote{psyclone.gocean1p0.GOLoop}}) \textendash{} the first Node representing a GOLoop.

\item {} 
\sphinxAtStartPar
\sphinxstyleliteralstrong{\sphinxupquote{node2}} (\sphinxcode{\sphinxupquote{psyclone.gocean1p0.GOLoop}}) \textendash{} the second Node representing a GOLoop.

\item {} 
\sphinxAtStartPar
\sphinxstyleliteralstrong{\sphinxupquote{options}} (\sphinxstyleliteralemphasis{\sphinxupquote{Optional}}\sphinxstyleliteralemphasis{\sphinxupquote{{[}}}\sphinxstyleliteralemphasis{\sphinxupquote{Dict}}\sphinxstyleliteralemphasis{\sphinxupquote{{[}}}\sphinxhref{https://docs.python.org/3/library/stdtypes.html\#str}{\sphinxstyleliteralemphasis{\sphinxupquote{str}}}\sphinxstyleliteralemphasis{\sphinxupquote{, }}\sphinxstyleliteralemphasis{\sphinxupquote{Any}}\sphinxstyleliteralemphasis{\sphinxupquote{{]}}}\sphinxstyleliteralemphasis{\sphinxupquote{{]}}}) \textendash{} a dictionary with options for transformations.

\end{itemize}

\sphinxlineitem{Raises}\begin{itemize}
\item {} 
\sphinxAtStartPar
\sphinxstyleliteralstrong{\sphinxupquote{TransformationError}} \textendash{} if the supplied loops are over                                      different grid\sphinxhyphen{}point types.

\item {} 
\sphinxAtStartPar
\sphinxstyleliteralstrong{\sphinxupquote{TransformationError}} \textendash{} if invalid parameters are passed in.

\end{itemize}

\end{description}\end{quote}

\end{fulllineitems}


\end{fulllineitems}



\bigskip\hrule\bigskip



\begin{fulllineitems}

\pysigstartsignatures
\pysiglinewithargsret{\sphinxbfcode{\sphinxupquote{class\DUrole{w,w}{  }}}\sphinxcode{\sphinxupquote{psyclone.transformations.}}\sphinxbfcode{\sphinxupquote{GOceanOMPParallelLoopTrans}}}{\sphinxparam{\DUrole{n,n}{omp\_directive}\DUrole{o,o}{=}\DUrole{default_value}{\textquotesingle{}do\textquotesingle{}}}\sphinxparamcomma \sphinxparam{\DUrole{n,n}{omp\_schedule}\DUrole{o,o}{=}\DUrole{default_value}{\textquotesingle{}static\textquotesingle{}}}}{}
\pysigstopsignatures
\sphinxAtStartPar
GOcean specific OpenMP Do loop transformation. Adds GOcean
specific validity checks (that supplied Loop is an inner or outer
loop). Actual transformation is done by
\sphinxcode{\sphinxupquote{base class}}.
\begin{quote}
\begin{quote}\begin{description}
\sphinxlineitem{param str omp\_directive}
\sphinxAtStartPar
choose which OpenMP loop directive to use.             Defaults to “do”.

\sphinxlineitem{param str omp\_schedule}
\sphinxAtStartPar
the OpenMP schedule to use. Must be one of             ‘runtime’, ‘static’, ‘dynamic’, ‘guided’ or ‘auto’. Defaults to             ‘static’.

\end{description}\end{quote}
\end{quote}


\begin{fulllineitems}

\pysigstartsignatures
\pysiglinewithargsret{\sphinxbfcode{\sphinxupquote{apply}}}{\sphinxparam{\DUrole{n,n}{node}}\sphinxparamcomma \sphinxparam{\DUrole{n,n}{options}\DUrole{o,o}{=}\DUrole{default_value}{None}}}{}
\pysigstopsignatures
\sphinxAtStartPar
Perform GOcean\sphinxhyphen{}specific loop validity checks then call
\sphinxcode{\sphinxupquote{OMPParallelLoopTrans.apply()}}.
\begin{quote}\begin{description}
\sphinxlineitem{Parameters}\begin{itemize}
\item {} 
\sphinxAtStartPar
\sphinxstyleliteralstrong{\sphinxupquote{node}} (\sphinxcode{\sphinxupquote{psyclone.psyir.nodes.Loop}}) \textendash{} a Loop node from an AST.

\item {} 
\sphinxAtStartPar
\sphinxstyleliteralstrong{\sphinxupquote{options}} (\sphinxstyleliteralemphasis{\sphinxupquote{Optional}}\sphinxstyleliteralemphasis{\sphinxupquote{{[}}}\sphinxstyleliteralemphasis{\sphinxupquote{Dict}}\sphinxstyleliteralemphasis{\sphinxupquote{{[}}}\sphinxhref{https://docs.python.org/3/library/stdtypes.html\#str}{\sphinxstyleliteralemphasis{\sphinxupquote{str}}}\sphinxstyleliteralemphasis{\sphinxupquote{, }}\sphinxstyleliteralemphasis{\sphinxupquote{Any}}\sphinxstyleliteralemphasis{\sphinxupquote{{]}}}\sphinxstyleliteralemphasis{\sphinxupquote{{]}}}) \textendash{} a dictionary with options for transformations                        and validation.

\end{itemize}

\sphinxlineitem{Raises}
\sphinxAtStartPar
\sphinxstyleliteralstrong{\sphinxupquote{TransformationError}} \textendash{} if the supplied node is not an inner or            outer loop.

\end{description}\end{quote}

\end{fulllineitems}


\end{fulllineitems}



\bigskip\hrule\bigskip



\begin{fulllineitems}

\pysigstartsignatures
\pysiglinewithargsret{\sphinxbfcode{\sphinxupquote{class\DUrole{w,w}{  }}}\sphinxcode{\sphinxupquote{psyclone.transformations.}}\sphinxbfcode{\sphinxupquote{GOceanOMPLoopTrans}}}{\sphinxparam{\DUrole{n,n}{omp\_directive}\DUrole{o,o}{=}\DUrole{default_value}{\textquotesingle{}do\textquotesingle{}}}\sphinxparamcomma \sphinxparam{\DUrole{n,n}{omp\_schedule}\DUrole{o,o}{=}\DUrole{default_value}{\textquotesingle{}static\textquotesingle{}}}}{}
\pysigstopsignatures
\sphinxAtStartPar
GOcean\sphinxhyphen{}specific orphan OpenMP loop transformation. Adds GOcean
specific validity checks (that the node is either an inner or outer
Loop).
\begin{quote}\begin{description}
\sphinxlineitem{Parameters}\begin{itemize}
\item {} 
\sphinxAtStartPar
\sphinxstyleliteralstrong{\sphinxupquote{omp\_directive}} (\sphinxhref{https://docs.python.org/3/library/stdtypes.html\#str}{\sphinxstyleliteralemphasis{\sphinxupquote{str}}}) \textendash{} choose which OpenMP loop directive to use.             Defaults to “do”.

\item {} 
\sphinxAtStartPar
\sphinxstyleliteralstrong{\sphinxupquote{omp\_schedule}} (\sphinxhref{https://docs.python.org/3/library/stdtypes.html\#str}{\sphinxstyleliteralemphasis{\sphinxupquote{str}}}) \textendash{} the OpenMP schedule to use. Must be one of             ‘runtime’, ‘static’, ‘dynamic’, ‘guided’ or ‘auto’. Defaults to             ‘static’.

\end{itemize}

\end{description}\end{quote}


\begin{fulllineitems}

\pysigstartsignatures
\pysiglinewithargsret{\sphinxbfcode{\sphinxupquote{validate}}}{\sphinxparam{\DUrole{n,n}{node}}\sphinxparamcomma \sphinxparam{\DUrole{n,n}{options}\DUrole{o,o}{=}\DUrole{default_value}{None}}}{}
\pysigstopsignatures
\sphinxAtStartPar
Checks that the supplied node is a valid target for parallelisation
using OMP directives.
\begin{quote}\begin{description}
\sphinxlineitem{Parameters}\begin{itemize}
\item {} 
\sphinxAtStartPar
\sphinxstyleliteralstrong{\sphinxupquote{node}} (\sphinxcode{\sphinxupquote{psyclone.psyir.nodes.Loop}}) \textendash{} the candidate loop for parallelising using OMP Do.

\item {} 
\sphinxAtStartPar
\sphinxstyleliteralstrong{\sphinxupquote{options}} (\sphinxstyleliteralemphasis{\sphinxupquote{Optional}}\sphinxstyleliteralemphasis{\sphinxupquote{{[}}}\sphinxstyleliteralemphasis{\sphinxupquote{Dict}}\sphinxstyleliteralemphasis{\sphinxupquote{{[}}}\sphinxhref{https://docs.python.org/3/library/stdtypes.html\#str}{\sphinxstyleliteralemphasis{\sphinxupquote{str}}}\sphinxstyleliteralemphasis{\sphinxupquote{, }}\sphinxstyleliteralemphasis{\sphinxupquote{Any}}\sphinxstyleliteralemphasis{\sphinxupquote{{]}}}\sphinxstyleliteralemphasis{\sphinxupquote{{]}}}) \textendash{} a dictionary with options for transformations.

\end{itemize}

\sphinxlineitem{Raises}
\sphinxAtStartPar
\sphinxstyleliteralstrong{\sphinxupquote{TransformationError}} \textendash{} if the loop\_type of the supplied Loop is                                      not “inner” or “outer”.

\end{description}\end{quote}

\end{fulllineitems}


\end{fulllineitems}



\bigskip\hrule\bigskip



\begin{fulllineitems}

\pysigstartsignatures
\pysigline{\sphinxbfcode{\sphinxupquote{class\DUrole{w,w}{  }}}\sphinxcode{\sphinxupquote{psyclone.domain.gocean.transformations.}}\sphinxbfcode{\sphinxupquote{GOConstLoopBoundsTrans}}}
\pysigstopsignatures
\sphinxAtStartPar
Use of a common constant variable for each loop bound within
a GOInvokeSchedule. By deafault, PSyclone generates loops where
the bounds are obtained by de\sphinxhyphen{}referencing a field object, e.g.:

\begin{sphinxVerbatim}[commandchars=\\\{\}]
\PYG{k}{DO }\PYG{n}{j}\PYG{+w}{ }\PYG{o}{=}\PYG{+w}{ }\PYG{n}{my\PYGZus{}field}\PYG{p}{\PYGZpc{}}\PYG{n}{grid}\PYG{p}{\PYGZpc{}}\PYG{n}{internal}\PYG{p}{\PYGZpc{}}\PYG{n}{ystart}\PYG{p}{,}\PYG{+w}{ }\PYG{n}{my\PYGZus{}field}\PYG{p}{\PYGZpc{}}\PYG{n}{grid}\PYG{p}{\PYGZpc{}}\PYG{n}{internal}\PYG{p}{\PYGZpc{}}\PYG{n}{ystop}
\end{sphinxVerbatim}

\sphinxAtStartPar
Some compilers are able to produce more efficient code if they are
provided with information on the relative trip\sphinxhyphen{}counts of the loops
within an Invoke. With constant loop bounds, PSyclone generates code
like:

\begin{sphinxVerbatim}[commandchars=\\\{\}]
\PYG{n}{ny}\PYG{+w}{ }\PYG{o}{=}\PYG{+w}{ }\PYG{n}{my\PYGZus{}field}\PYG{p}{\PYGZpc{}}\PYG{n}{grid}\PYG{p}{\PYGZpc{}}\PYG{n}{subdomain}\PYG{p}{\PYGZpc{}}\PYG{n}{internal}\PYG{p}{\PYGZpc{}}\PYG{n}{ystop}
\PYG{p}{.}\PYG{p}{.}\PYG{p}{.}
\PYG{k}{DO }\PYG{n}{j}\PYG{+w}{ }\PYG{o}{=}\PYG{+w}{ }\PYG{l+m+mi}{1}\PYG{p}{,}\PYG{+w}{ }\PYG{n}{ny}\PYG{o}{\PYGZhy{}}\PYG{l+m+mi}{1}
\end{sphinxVerbatim}

\sphinxAtStartPar
In practice, the application of the constant loop bounds transformation
looks something like, e.g.:

\begin{sphinxVerbatim}[commandchars=\\\{\}]
\PYG{g+gp}{\PYGZgt{}\PYGZgt{}\PYGZgt{} }\PYG{k+kn}{from}\PYG{+w}{ }\PYG{n+nn}{psyclone}\PYG{n+nn}{.}\PYG{n+nn}{parse}\PYG{n+nn}{.}\PYG{n+nn}{algorithm}\PYG{+w}{ }\PYG{k+kn}{import} \PYG{n}{parse}
\PYG{g+gp}{\PYGZgt{}\PYGZgt{}\PYGZgt{} }\PYG{k+kn}{from}\PYG{+w}{ }\PYG{n+nn}{psyclone}\PYG{n+nn}{.}\PYG{n+nn}{psyGen}\PYG{+w}{ }\PYG{k+kn}{import} \PYG{n}{PSyFactory}
\PYG{g+gp}{\PYGZgt{}\PYGZgt{}\PYGZgt{} }\PYG{k+kn}{import}\PYG{+w}{ }\PYG{n+nn}{os}
\PYG{g+gp}{\PYGZgt{}\PYGZgt{}\PYGZgt{} }\PYG{n}{TEST\PYGZus{}API} \PYG{o}{=} \PYG{l+s+s2}{\PYGZdq{}}\PYG{l+s+s2}{gocean}\PYG{l+s+s2}{\PYGZdq{}}
\PYG{g+gp}{\PYGZgt{}\PYGZgt{}\PYGZgt{} }\PYG{n}{\PYGZus{}}\PYG{p}{,} \PYG{n}{info} \PYG{o}{=} \PYG{n}{parse}\PYG{p}{(}\PYG{n}{os}\PYG{o}{.}\PYG{n}{path}\PYG{o}{.}\PYG{n}{join}\PYG{p}{(}\PYG{l+s+s2}{\PYGZdq{}}\PYG{l+s+s2}{tests}\PYG{l+s+s2}{\PYGZdq{}}\PYG{p}{,} \PYG{l+s+s2}{\PYGZdq{}}\PYG{l+s+s2}{test\PYGZus{}files}\PYG{l+s+s2}{\PYGZdq{}}\PYG{p}{,} \PYG{l+s+s2}{\PYGZdq{}}\PYG{l+s+s2}{gocean1p0}\PYG{l+s+s2}{\PYGZdq{}}\PYG{p}{,}
\PYG{g+gp}{... }                             \PYG{l+s+s2}{\PYGZdq{}}\PYG{l+s+s2}{single\PYGZus{}invoke.f90}\PYG{l+s+s2}{\PYGZdq{}}\PYG{p}{)}\PYG{p}{,}
\PYG{g+gp}{... }                \PYG{n}{api}\PYG{o}{=}\PYG{n}{TEST\PYGZus{}API}\PYG{p}{)}
\PYG{g+gp}{\PYGZgt{}\PYGZgt{}\PYGZgt{} }\PYG{n}{psy} \PYG{o}{=} \PYG{n}{PSyFactory}\PYG{p}{(}\PYG{n}{TEST\PYGZus{}API}\PYG{p}{)}\PYG{o}{.}\PYG{n}{create}\PYG{p}{(}\PYG{n}{info}\PYG{p}{)}
\PYG{g+gp}{\PYGZgt{}\PYGZgt{}\PYGZgt{} }\PYG{n}{invoke} \PYG{o}{=} \PYG{n}{psy}\PYG{o}{.}\PYG{n}{invokes}\PYG{o}{.}\PYG{n}{get}\PYG{p}{(}\PYG{l+s+s1}{\PYGZsq{}}\PYG{l+s+s1}{invoke\PYGZus{}0\PYGZus{}compute\PYGZus{}cu}\PYG{l+s+s1}{\PYGZsq{}}\PYG{p}{)}
\PYG{g+gp}{\PYGZgt{}\PYGZgt{}\PYGZgt{} }\PYG{n}{schedule} \PYG{o}{=} \PYG{n}{invoke}\PYG{o}{.}\PYG{n}{schedule}
\PYG{g+gp}{\PYGZgt{}\PYGZgt{}\PYGZgt{}}
\PYG{g+gp}{\PYGZgt{}\PYGZgt{}\PYGZgt{} }\PYG{k+kn}{from}\PYG{+w}{ }\PYG{n+nn}{psyclone}\PYG{n+nn}{.}\PYG{n+nn}{transformations}\PYG{+w}{ }\PYG{k+kn}{import} \PYG{n}{GOConstLoopBoundsTrans}
\PYG{g+gp}{\PYGZgt{}\PYGZgt{}\PYGZgt{} }\PYG{n}{clbtrans} \PYG{o}{=} \PYG{n}{GOConstLoopBoundsTrans}\PYG{p}{(}\PYG{p}{)}
\PYG{g+gp}{\PYGZgt{}\PYGZgt{}\PYGZgt{}}
\PYG{g+gp}{\PYGZgt{}\PYGZgt{}\PYGZgt{} }\PYG{n}{clbtrans}\PYG{o}{.}\PYG{n}{apply}\PYG{p}{(}\PYG{n}{schedule}\PYG{p}{)}
\PYG{g+gp}{\PYGZgt{}\PYGZgt{}\PYGZgt{} }\PYG{n+nb}{print}\PYG{p}{(}\PYG{n}{schedule}\PYG{o}{.}\PYG{n}{view}\PYG{p}{(}\PYG{p}{)}\PYG{p}{)}
\end{sphinxVerbatim}


\begin{fulllineitems}

\pysigstartsignatures
\pysiglinewithargsret{\sphinxbfcode{\sphinxupquote{apply}}}{\sphinxparam{\DUrole{n,n}{node}}\sphinxparamcomma \sphinxparam{\DUrole{n,n}{options}\DUrole{o,o}{=}\DUrole{default_value}{None}}}{}
\pysigstopsignatures
\sphinxAtStartPar
Modify the GOcean kernel loops in a GOInvokeSchedule to use
common constant loop bound variables.
\begin{quote}\begin{description}
\sphinxlineitem{Parameters}\begin{itemize}
\item {} 
\sphinxAtStartPar
\sphinxstyleliteralstrong{\sphinxupquote{node}} (\sphinxcode{\sphinxupquote{psyclone.gocean1p0.GOInvokeSchedule}}) \textendash{} the GOInvokeSchedule of which all loops will get the             constant loop bounds.

\item {} 
\sphinxAtStartPar
\sphinxstyleliteralstrong{\sphinxupquote{options}} (\sphinxstyleliteralemphasis{\sphinxupquote{Optional}}\sphinxstyleliteralemphasis{\sphinxupquote{{[}}}\sphinxstyleliteralemphasis{\sphinxupquote{Dict}}\sphinxstyleliteralemphasis{\sphinxupquote{{[}}}\sphinxhref{https://docs.python.org/3/library/stdtypes.html\#str}{\sphinxstyleliteralemphasis{\sphinxupquote{str}}}\sphinxstyleliteralemphasis{\sphinxupquote{, }}\sphinxstyleliteralemphasis{\sphinxupquote{Any}}\sphinxstyleliteralemphasis{\sphinxupquote{{]}}}\sphinxstyleliteralemphasis{\sphinxupquote{{]}}}) \textendash{} a dictionary with options for transformations.

\end{itemize}

\end{description}\end{quote}

\end{fulllineitems}



\begin{fulllineitems}

\pysigstartsignatures
\pysigline{\sphinxbfcode{\sphinxupquote{property\DUrole{w,w}{  }}}\sphinxbfcode{\sphinxupquote{name}}}
\pysigstopsignatures\begin{quote}\begin{description}
\sphinxlineitem{Returns}
\sphinxAtStartPar
the name of the Transformation as a string.

\sphinxlineitem{Return type}
\sphinxAtStartPar
\sphinxhref{https://docs.python.org/3/library/stdtypes.html\#str}{str}

\end{description}\end{quote}

\end{fulllineitems}



\begin{fulllineitems}

\pysigstartsignatures
\pysiglinewithargsret{\sphinxbfcode{\sphinxupquote{validate}}}{\sphinxparam{\DUrole{n,n}{node}}\sphinxparamcomma \sphinxparam{\DUrole{n,n}{options}\DUrole{o,o}{=}\DUrole{default_value}{None}}}{}
\pysigstopsignatures
\sphinxAtStartPar
Checks if it is valid to apply the GOConstLoopBoundsTrans
transform.
\begin{quote}\begin{description}
\sphinxlineitem{Parameters}\begin{itemize}
\item {} 
\sphinxAtStartPar
\sphinxstyleliteralstrong{\sphinxupquote{node}} (\sphinxcode{\sphinxupquote{psyclone.gocean1p0.GOInvokeSchedule}}) \textendash{} the GOInvokeSchedule to transform.

\item {} 
\sphinxAtStartPar
\sphinxstyleliteralstrong{\sphinxupquote{options}} (\sphinxstyleliteralemphasis{\sphinxupquote{Optional}}\sphinxstyleliteralemphasis{\sphinxupquote{{[}}}\sphinxstyleliteralemphasis{\sphinxupquote{Dict}}\sphinxstyleliteralemphasis{\sphinxupquote{{[}}}\sphinxhref{https://docs.python.org/3/library/stdtypes.html\#str}{\sphinxstyleliteralemphasis{\sphinxupquote{str}}}\sphinxstyleliteralemphasis{\sphinxupquote{, }}\sphinxstyleliteralemphasis{\sphinxupquote{Any}}\sphinxstyleliteralemphasis{\sphinxupquote{{]}}}\sphinxstyleliteralemphasis{\sphinxupquote{{]}}}) \textendash{} a dictionary with options for transformations.

\end{itemize}

\sphinxlineitem{Raises}\begin{itemize}
\item {} 
\sphinxAtStartPar
\sphinxstyleliteralstrong{\sphinxupquote{TransformationError}} \textendash{} if the supplied node is not a             GOInvokeSchedule.

\item {} 
\sphinxAtStartPar
\sphinxstyleliteralstrong{\sphinxupquote{TransformationError}} \textendash{} if the supplied schedule has loops with             a loop with loop\_type different than ‘inner’ or ‘outer’.

\item {} 
\sphinxAtStartPar
\sphinxstyleliteralstrong{\sphinxupquote{TransformationError}} \textendash{} if the supplied schedule has loops with             attributes for index\_offsets, field\_space, iteration\_space and             loop\_type that don’t appear in the GOLoop.bounds\_lookup table.

\item {} 
\sphinxAtStartPar
\sphinxstyleliteralstrong{\sphinxupquote{TransformationError}} \textendash{} if the supplied schedule doesn’t have a             field argument.

\end{itemize}

\end{description}\end{quote}

\end{fulllineitems}


\end{fulllineitems}



\bigskip\hrule\bigskip



\begin{fulllineitems}

\pysigstartsignatures
\pysigline{\sphinxbfcode{\sphinxupquote{class\DUrole{w,w}{  }}}\sphinxcode{\sphinxupquote{psyclone.domain.gocean.transformations.}}\sphinxbfcode{\sphinxupquote{GOMoveIterationBoundariesInsideKernelTrans}}}
\pysigstopsignatures
\sphinxAtStartPar
Provides a transformation that moves iteration boundaries that are
encoded in the Loops lower\_bound() and upper\_bound() methods to a mask
inside the kernel with the boundaries passed as kernel arguments.

\sphinxAtStartPar
For example the following kernel call:

\begin{sphinxVerbatim}[commandchars=\\\{\}]
\PYG{k}{do }\PYG{n}{i}\PYG{+w}{ }\PYG{o}{=}\PYG{+w}{ }\PYG{l+m+mi}{2}\PYG{p}{,}\PYG{+w}{ }\PYG{n}{N}\PYG{+w}{ }\PYG{o}{\PYGZhy{}}\PYG{+w}{ }\PYG{l+m+mi}{1}
\PYG{+w}{    }\PYG{k}{do }\PYG{n}{j}\PYG{+w}{ }\PYG{o}{=}\PYG{+w}{ }\PYG{l+m+mi}{2}\PYG{p}{,}\PYG{+w}{ }\PYG{n}{N}\PYG{+w}{ }\PYG{o}{\PYGZhy{}}\PYG{+w}{ }\PYG{l+m+mi}{1}
\PYG{+w}{        }\PYG{n}{kernel}\PYG{p}{(}\PYG{n}{i}\PYG{p}{,}\PYG{+w}{ }\PYG{n}{j}\PYG{p}{,}\PYG{+w}{ }\PYG{n}{field}\PYG{p}{)}
\PYG{+w}{    }\PYG{k}{end }\PYG{k}{do}
\PYG{k}{end }\PYG{k}{do}
\end{sphinxVerbatim}

\sphinxAtStartPar
will be transformed to:

\begin{sphinxVerbatim}[commandchars=\\\{\}]
\PYG{n}{startx}\PYG{+w}{ }\PYG{o}{=}\PYG{+w}{ }\PYG{l+m+mi}{2}
\PYG{n}{stopx}\PYG{+w}{ }\PYG{o}{=}\PYG{+w}{ }\PYG{n}{N}\PYG{+w}{ }\PYG{o}{\PYGZhy{}}\PYG{+w}{ }\PYG{l+m+mi}{1}
\PYG{n}{starty}\PYG{+w}{ }\PYG{o}{=}\PYG{+w}{ }\PYG{l+m+mi}{2}
\PYG{n}{stopy}\PYG{+w}{ }\PYG{o}{=}\PYG{+w}{ }\PYG{n}{N}\PYG{+w}{ }\PYG{o}{\PYGZhy{}}\PYG{+w}{ }\PYG{l+m+mi}{1}
\PYG{k}{do }\PYG{n}{i}\PYG{+w}{ }\PYG{o}{=}\PYG{+w}{ }\PYG{l+m+mi}{1}\PYG{p}{,}\PYG{+w}{ }\PYG{n}{size}\PYG{p}{(}\PYG{n}{field}\PYG{p}{,}\PYG{+w}{ }\PYG{l+m+mi}{1}\PYG{p}{)}
\PYG{+w}{    }\PYG{k}{do }\PYG{n}{j}\PYG{+w}{ }\PYG{o}{=}\PYG{+w}{ }\PYG{l+m+mi}{1}\PYG{p}{,}\PYG{+w}{ }\PYG{n}{size}\PYG{p}{(}\PYG{n}{field}\PYG{p}{,}\PYG{+w}{ }\PYG{l+m+mi}{2}\PYG{p}{)}
\PYG{+w}{        }\PYG{n}{kernel}\PYG{p}{(}\PYG{n}{i}\PYG{p}{,}\PYG{+w}{ }\PYG{n}{j}\PYG{p}{,}\PYG{+w}{ }\PYG{n}{field}\PYG{p}{,}\PYG{+w}{ }\PYG{n}{startx}\PYG{p}{,}\PYG{+w}{ }\PYG{n}{stopx}\PYG{p}{,}\PYG{+w}{ }\PYG{n}{starty}\PYG{p}{,}\PYG{+w}{ }\PYG{n}{stopy}\PYG{p}{)}
\PYG{+w}{    }\PYG{k}{end }\PYG{k}{do}
\PYG{k}{end }\PYG{k}{do}
\end{sphinxVerbatim}

\sphinxAtStartPar
additionally a mask like the following one will be introduced in the
kernel code:

\begin{sphinxVerbatim}[commandchars=\\\{\}]
\PYG{k}{if}\PYG{+w}{ }\PYG{p}{(}\PYG{n}{i}\PYG{+w}{ }\PYG{o}{\PYGZlt{}}\PYG{+w}{ }\PYG{n}{startx}\PYG{+w}{ }\PYG{p}{.}\PYG{n+nb}{or}\PYG{p}{.}\PYG{+w}{ }\PYG{n}{i}\PYG{+w}{ }\PYG{o}{\PYGZgt{}}\PYG{+w}{ }\PYG{n}{stopx}\PYG{+w}{ }\PYG{p}{.}\PYG{n+nb}{or}\PYG{p}{.}\PYG{+w}{ }\PYG{n}{j}\PYG{+w}{ }\PYG{o}{\PYGZlt{}}\PYG{+w}{ }\PYG{n}{starty}\PYG{+w}{ }\PYG{p}{.}\PYG{n+nb}{or}\PYG{p}{.}\PYG{+w}{ }\PYG{n}{j}\PYG{+w}{ }\PYG{o}{\PYGZgt{}}\PYG{+w}{ }\PYG{n}{stopy}\PYG{p}{)}\PYG{+w}{ }\PYG{k}{then}
\PYG{k}{    }\PYG{k}{return}
\PYG{k}{end }\PYG{k}{if}
\end{sphinxVerbatim}


\begin{fulllineitems}

\pysigstartsignatures
\pysiglinewithargsret{\sphinxbfcode{\sphinxupquote{apply}}}{\sphinxparam{\DUrole{n,n}{node}}\sphinxparamcomma \sphinxparam{\DUrole{n,n}{options}\DUrole{o,o}{=}\DUrole{default_value}{None}}}{}
\pysigstopsignatures
\sphinxAtStartPar
Apply this transformation to the supplied node.
\begin{quote}\begin{description}
\sphinxlineitem{Parameters}\begin{itemize}
\item {} 
\sphinxAtStartPar
\sphinxstyleliteralstrong{\sphinxupquote{node}} (\sphinxcode{\sphinxupquote{psyclone.gocean1p0.GOKern}}) \textendash{} the node to transform.

\item {} 
\sphinxAtStartPar
\sphinxstyleliteralstrong{\sphinxupquote{options}} (\sphinxstyleliteralemphasis{\sphinxupquote{Optional}}\sphinxstyleliteralemphasis{\sphinxupquote{{[}}}\sphinxstyleliteralemphasis{\sphinxupquote{Dict}}\sphinxstyleliteralemphasis{\sphinxupquote{{[}}}\sphinxhref{https://docs.python.org/3/library/stdtypes.html\#str}{\sphinxstyleliteralemphasis{\sphinxupquote{str}}}\sphinxstyleliteralemphasis{\sphinxupquote{, }}\sphinxstyleliteralemphasis{\sphinxupquote{Any}}\sphinxstyleliteralemphasis{\sphinxupquote{{]}}}\sphinxstyleliteralemphasis{\sphinxupquote{{]}}}) \textendash{} a dictionary with options for transformations.

\end{itemize}

\end{description}\end{quote}

\end{fulllineitems}



\begin{fulllineitems}

\pysigstartsignatures
\pysigline{\sphinxbfcode{\sphinxupquote{property\DUrole{w,w}{  }}}\sphinxbfcode{\sphinxupquote{name}}}
\pysigstopsignatures
\sphinxAtStartPar
Returns the name of this transformation as a string.

\end{fulllineitems}



\begin{fulllineitems}

\pysigstartsignatures
\pysiglinewithargsret{\sphinxbfcode{\sphinxupquote{validate}}}{\sphinxparam{\DUrole{n,n}{node}}\sphinxparamcomma \sphinxparam{\DUrole{n,n}{options}\DUrole{o,o}{=}\DUrole{default_value}{None}}}{}
\pysigstopsignatures
\sphinxAtStartPar
Ensure that it is valid to apply this transformation to the
supplied node.
\begin{quote}\begin{description}
\sphinxlineitem{Parameters}\begin{itemize}
\item {} 
\sphinxAtStartPar
\sphinxstyleliteralstrong{\sphinxupquote{node}} (\sphinxcode{\sphinxupquote{psyclone.gocean1p0.GOKern}}) \textendash{} the node to validate.

\item {} 
\sphinxAtStartPar
\sphinxstyleliteralstrong{\sphinxupquote{options}} (\sphinxstyleliteralemphasis{\sphinxupquote{Optional}}\sphinxstyleliteralemphasis{\sphinxupquote{{[}}}\sphinxstyleliteralemphasis{\sphinxupquote{Dict}}\sphinxstyleliteralemphasis{\sphinxupquote{{[}}}\sphinxhref{https://docs.python.org/3/library/stdtypes.html\#str}{\sphinxstyleliteralemphasis{\sphinxupquote{str}}}\sphinxstyleliteralemphasis{\sphinxupquote{, }}\sphinxstyleliteralemphasis{\sphinxupquote{Any}}\sphinxstyleliteralemphasis{\sphinxupquote{{]}}}\sphinxstyleliteralemphasis{\sphinxupquote{{]}}}) \textendash{} a dictionary with options for transformations.

\end{itemize}

\sphinxlineitem{Raises}
\sphinxAtStartPar
\sphinxstyleliteralstrong{\sphinxupquote{TransformationError}} \textendash{} if the node is not a GOKern.

\end{description}\end{quote}

\end{fulllineitems}


\end{fulllineitems}


\sphinxstepscope


\chapter{PSyclone Kernel Tools}
\label{\detokenize{psyclone_kern:psyclone-kernel-tools}}\label{\detokenize{psyclone_kern::doc}}
\sphinxAtStartPar
In addition to the \sphinxcode{\sphinxupquote{psyclone}} command, the PSyclone package also
provides tools related to generating code purely from kernel
metadata. Currently there are two such tools:
\begin{enumerate}
\sphinxsetlistlabels{\arabic}{enumi}{enumii}{}{.}%
\item {} 
\sphinxAtStartPar
{\hyperref[\detokenize{psyclone_kern:stub-generation}]{\sphinxcrossref{\DUrole{std,std-ref}{Kernel\sphinxhyphen{}stub Generator}}}}

\item {} 
\sphinxAtStartPar
{\hyperref[\detokenize{psyclone_kern:alg-generation}]{\sphinxcrossref{\DUrole{std,std-ref}{Algorithm Generator}}}}

\end{enumerate}

\sphinxAtStartPar
The kernel\sphinxhyphen{}stub generator takes a file containing kernel metadata as
input and outputs the (Fortran) kernel subroutine arguments and
declarations. The word “stub” is used to indicate that it is only the
subroutine arguments and their declarations that are generated; the
subroutine has no content.

\sphinxAtStartPar
The algorithm generator also takes a file containing a kernel
implementation but this time generates an appropriate algorithm layer
subroutine. This algorithm layer plus the associated kernel metadata
may then be processed with PSyclone in the usual way to generate code
which executes the supplied kernel.

\sphinxAtStartPar
This functionality is provided to the user via the \sphinxcode{\sphinxupquote{psyclone\sphinxhyphen{}kern}}
command, described in more detail below.


\section{The psyclone\sphinxhyphen{}kern Command}
\label{\detokenize{psyclone_kern:the-psyclone-kern-command}}
\sphinxAtStartPar
Before using the \sphinxcode{\sphinxupquote{psyclone\sphinxhyphen{}kern}} tool, PSyclone must be installed. If you
have not already done so, please follow the instructions for setting
up PSyclone in Section {\hyperref[\detokenize{getting_going:getting-going}]{\sphinxcrossref{\DUrole{std,std-ref}{Getting Going}}}}.

\sphinxAtStartPar
PSyclone will be installed in a particular location on your machine,
which will be referred to as the \sphinxcode{\sphinxupquote{\textless{}PSYCLONEINSTALL\textgreater{}}} directory. The
\sphinxcode{\sphinxupquote{psyclone\sphinxhyphen{}kern}} script comes with the PSyclone
installation. A quick check \sphinxcode{\sphinxupquote{\textgreater{} which psyclone\sphinxhyphen{}kern}} should return
the location of the \sphinxcode{\sphinxupquote{\textless{}PSYCLONEINSTALL\textgreater{}/bin}} directory.

\sphinxAtStartPar
The \sphinxcode{\sphinxupquote{psyclone\sphinxhyphen{}kern}} command has the following arguments:

\begin{sphinxVerbatim}[commandchars=\\\{\}]
\PYGZgt{}\PYG{+w}{ }psyclone\PYGZhy{}kern\PYG{+w}{ }\PYGZhy{}h
usage:\PYG{+w}{ }psyclone\PYGZhy{}kern\PYG{+w}{ }\PYG{o}{[}\PYGZhy{}h\PYG{o}{]}\PYG{+w}{ }\PYG{o}{[}\PYGZhy{}gen\PYG{+w}{ }\PYG{o}{\PYGZob{}}alg,stub\PYG{o}{\PYGZcb{}}\PYG{o}{]}\PYG{+w}{ }\PYG{o}{[}\PYGZhy{}o\PYG{+w}{ }OUT\PYGZus{}FILE\PYG{o}{]}\PYG{+w}{ }\PYG{o}{[}\PYGZhy{}api\PYG{+w}{ }API\PYG{o}{]}
\PYG{+w}{                     }\PYG{o}{[}\PYGZhy{}I\PYG{+w}{ }INCLUDE\PYG{o}{]}\PYG{+w}{ }\PYG{o}{[}\PYGZhy{}l\PYG{+w}{ }\PYG{o}{\PYGZob{}}off,all,output\PYG{o}{\PYGZcb{}}\PYG{o}{]}
\PYG{+w}{                     }\PYG{o}{[}\PYGZhy{}\PYGZhy{}config\PYG{+w}{ }CONFIG\PYG{o}{]}\PYG{+w}{ }\PYG{o}{[}\PYGZhy{}v\PYG{o}{]}
\PYG{+w}{                     }filename

Run\PYG{+w}{ }the\PYG{+w}{ }PSyclone\PYG{+w}{ }kernel\PYG{+w}{ }generator\PYG{+w}{ }on\PYG{+w}{ }a\PYG{+w}{ }particular\PYG{+w}{ }file

positional\PYG{+w}{ }arguments:
\PYG{+w}{  }filename\PYG{+w}{              }file\PYG{+w}{ }containing\PYG{+w}{ }Kernel\PYG{+w}{ }metadata

optional\PYG{+w}{ }arguments:
\PYG{+w}{  }\PYGZhy{}h,\PYG{+w}{ }\PYGZhy{}\PYGZhy{}help\PYG{+w}{            }show\PYG{+w}{ }this\PYG{+w}{ }\PYG{n+nb}{help}\PYG{+w}{ }message\PYG{+w}{ }and\PYG{+w}{ }\PYG{n+nb}{exit}
\PYG{+w}{  }\PYGZhy{}gen\PYG{+w}{ }\PYG{o}{\PYGZob{}}alg,stub\PYG{o}{)}\PYG{+w}{       }what\PYG{+w}{ }to\PYG{+w}{ }generate\PYG{+w}{ }\PYG{k}{for}\PYG{+w}{ }the\PYG{+w}{ }supplied\PYG{+w}{ }kernel
\PYG{+w}{                        }\PYG{o}{(}\PYG{n+nv}{alg}\PYG{o}{=}algorithm\PYG{+w}{ }layer,\PYG{+w}{ }\PYG{n+nv}{stub}\PYG{o}{=}kernel\PYGZhy{}stub\PYG{+w}{ }subroutine\PYG{o}{)}.
\PYG{+w}{                        }Defaults\PYG{+w}{ }to\PYG{+w}{ }stub.
\PYG{+w}{  }\PYGZhy{}o\PYG{+w}{ }OUT\PYGZus{}FILE\PYG{+w}{           }filename\PYG{+w}{ }\PYG{k}{for}\PYG{+w}{ }created\PYG{+w}{ }code.
\PYG{+w}{  }\PYGZhy{}api\PYG{+w}{ }API\PYG{+w}{              }choose\PYG{+w}{ }a\PYG{+w}{ }particular\PYG{+w}{ }API\PYG{+w}{ }from\PYG{+w}{ }\PYG{o}{[}\PYG{l+s+s1}{\PYGZsq{}lfric\PYGZsq{}},
\PYG{+w}{                        }\PYG{l+s+s1}{\PYGZsq{}gocean\PYGZsq{}}\PYG{o}{]}.
\PYG{+w}{  }\PYGZhy{}I\PYG{+w}{ }INCLUDE,\PYG{+w}{ }\PYGZhy{}\PYGZhy{}include\PYG{+w}{ }INCLUDE
\PYG{+w}{                        }path\PYG{+w}{ }to\PYG{+w}{ }Fortran\PYG{+w}{ }INCLUDE\PYG{+w}{ }or\PYG{+w}{ }module\PYG{+w}{ }files
\PYG{+w}{  }\PYGZhy{}l\PYG{+w}{ }\PYG{o}{\PYGZob{}}off,all,output\PYG{o}{\PYGZcb{}},\PYG{+w}{ }\PYGZhy{}\PYGZhy{}limit\PYG{+w}{ }\PYG{o}{\PYGZob{}}off,all,output\PYG{o}{\PYGZcb{}}
\PYG{+w}{                        }limit\PYG{+w}{ }the\PYG{+w}{ }Fortran\PYG{+w}{ }line\PYG{+w}{ }length\PYG{+w}{ }to\PYG{+w}{ }\PYG{l+m}{132}
\PYG{+w}{                        }characters\PYG{+w}{ }\PYG{o}{(}default\PYG{+w}{ }\PYG{l+s+s1}{\PYGZsq{}off\PYGZsq{}}\PYG{o}{)}.\PYG{+w}{ }Use\PYG{+w}{ }\PYG{l+s+s1}{\PYGZsq{}all\PYGZsq{}}\PYG{+w}{ }to
\PYG{+w}{                        }apply\PYG{+w}{ }limit\PYG{+w}{ }to\PYG{+w}{ }both\PYG{+w}{ }input\PYG{+w}{ }and\PYG{+w}{ }output
\PYG{+w}{                        }Fortran.\PYG{+w}{ }Use\PYG{+w}{ }\PYG{l+s+s1}{\PYGZsq{}output\PYGZsq{}}\PYG{+w}{ }to\PYG{+w}{ }apply\PYG{+w}{ }line\PYGZhy{}length
\PYG{+w}{                        }limit\PYG{+w}{ }to\PYG{+w}{ }output\PYG{+w}{ }Fortran\PYG{+w}{ }only.
\PYG{+w}{  }\PYGZhy{}\PYGZhy{}config\PYG{+w}{ }CONFIG,\PYG{+w}{ }\PYGZhy{}c\PYG{+w}{ }CONFIG
\PYG{+w}{                        }config\PYG{+w}{ }file\PYG{+w}{ }with\PYG{+w}{ }PSyclone\PYG{+w}{ }specific\PYG{+w}{ }options.
\PYG{+w}{  }\PYGZhy{}v,\PYG{+w}{ }\PYGZhy{}\PYGZhy{}version\PYG{+w}{         }display\PYG{+w}{ }version\PYG{+w}{ }information\PYG{+w}{ }\PYG{o}{(}\PYG{l+s+se}{\PYGZbs{} }\PYG{p}{|}release\PYG{p}{|}\PYG{l+s+se}{\PYGZbs{} }\PYG{o}{)}
\end{sphinxVerbatim}

\sphinxAtStartPar
The \sphinxcode{\sphinxupquote{\sphinxhyphen{}o}} option allows the user to specify that the output should be
written to a particular file. If this is not specified then the Python
\sphinxcode{\sphinxupquote{print}} statement is used to write to stdout.  Typically this
results in the output being printed to the terminal.

\sphinxAtStartPar
The \sphinxcode{\sphinxupquote{\sphinxhyphen{}l}}, or \sphinxcode{\sphinxupquote{\sphinxhyphen{}\sphinxhyphen{}limit}} option utilises the PSyclone support for
wrapping of lines within the 132 character limit in the generated Fortran code
(please see the {\hyperref[\detokenize{line_length:line-length}]{\sphinxcrossref{\DUrole{std,std-ref}{Line Length}}}} chapter for more details).


\section{Kernel\sphinxhyphen{}stub Generator}
\label{\detokenize{psyclone_kern:kernel-stub-generator}}\label{\detokenize{psyclone_kern:stub-generation}}

\subsection{Quick Start}
\label{\detokenize{psyclone_kern:quick-start}}\begin{enumerate}
\sphinxsetlistlabels{\arabic}{enumi}{enumii}{}{)}%
\item {} 
\sphinxAtStartPar
Use an existing Kernel file or create a Kernel file containing a
Kernel module with the required metadata and an empty Kernel
subroutine with no arguments.

\item {} 
\sphinxAtStartPar
Run the following command

\begin{sphinxVerbatim}[commandchars=\\\{\}]
\PYG{o}{\PYGZgt{}} \PYG{n}{psyclone}\PYG{o}{\PYGZhy{}}\PYG{n}{kern} \PYG{o}{\PYGZhy{}}\PYG{n}{api} \PYG{n}{lfric} \PYG{o}{\PYGZhy{}}\PYG{n}{gen} \PYG{n}{stub} \PYG{o}{\PYGZlt{}}\PYG{n}{PATH}\PYG{o}{\PYGZgt{}}\PYG{o}{/}\PYG{n}{my\PYGZus{}file}\PYG{o}{.}\PYG{n}{f90}
\end{sphinxVerbatim}

\item {} 
\sphinxAtStartPar
To have the generated code written to file rather than stdout use the
\sphinxtitleref{\sphinxhyphen{}o} flag

\begin{sphinxVerbatim}[commandchars=\\\{\}]
\PYG{o}{\PYGZgt{}} \PYG{n}{psyclone}\PYG{o}{\PYGZhy{}}\PYG{n}{kern} \PYG{o}{\PYGZhy{}}\PYG{n}{api} \PYG{n}{lfric} \PYG{o}{\PYGZhy{}}\PYG{n}{gen} \PYG{n}{stub} \PYG{o}{\PYGZhy{}}\PYG{n}{o} \PYG{n}{my\PYGZus{}stub\PYGZus{}file}\PYG{o}{.}\PYG{n}{f90} \PYG{o}{.}\PYG{o}{/}\PYG{n}{my\PYGZus{}kernel\PYGZus{}mod}\PYG{o}{.}\PYG{n}{f90}
\end{sphinxVerbatim}

\end{enumerate}

\sphinxAtStartPar
(Since stub generation is the default, the \sphinxcode{\sphinxupquote{\sphinxhyphen{}gen stub}} may be omitted
if desired.)


\subsection{Introduction}
\label{\detokenize{psyclone_kern:introduction}}
\sphinxAtStartPar
PSyclone provides a kernel stub generator for the LFRic API.
The kernel stub generator takes a kernel file as input and outputs the
kernel subroutine arguments and declarations. The word “stub” is used
to indicate that it is only the subroutine arguments and their
declarations that are generated; the subroutine has no content.

\sphinxAtStartPar
The primary reason the stub generator is useful is that it generates
the correct Kernel subroutine arguments and declarations for the
LFRic API as specified by the Kernel metadata. As the number of
arguments to Kernel subroutines can become large and the arguments
have to follow a particular order, it can become burdensome, and
potentially error prone, for the user to have to work out the
appropriate argument list if written by hand.

\sphinxAtStartPar
The stub generator can be used when creating a new Kernel. A Kernel
can first be written to specify the required metadata and then the
generator can be used to create the appropriate (empty) Kernel
subroutine. The user can then fill in the content of the subroutine.

\sphinxAtStartPar
The stub generator can also be used to check whether the arguments for
an existing Kernel are correct i.e. whether the Kernel subroutine and
Kernel metadata are consistent. One example would be where a Kernel is
updated resulting in a change to the metadata and subroutine
arguments.

\sphinxAtStartPar
The LFRic API requires Kernels to conform to a set of rules which
determine the required arguments and types for a particular
Kernel. These rules are required as the generated PSy layer needs to
know exactly how to call a Kernel. These rules are outlined in Section
{\hyperref[\detokenize{dynamo0p3:lfric-stub-generation-rules}]{\sphinxcrossref{\DUrole{std,std-ref}{Rules}}}}.

\sphinxAtStartPar
Therefore PSyclone has been coded with the LFRic API rules which
are then applied when reading the Kernel metadata to produce the
required Kernel call and its arguments in the generated PSy
layer. These same rules are used by the Kernel stub generator to
produce Kernel subroutine stubs, thereby guaranteeing that Kernel
calls from the PSy layer and the associated Kernel subroutines are
consistent.


\subsection{Kernels}
\label{\detokenize{psyclone_kern:kernels}}\label{\detokenize{psyclone_kern:stub-generation-kernels}}
\sphinxAtStartPar
Any LFRic kernel can be used as input to the stub generator.
Example Kernels can be found in the \sphinxcode{\sphinxupquote{examples/lfric}} repository or,
for more simple cases, in the \sphinxcode{\sphinxupquote{tests/test\_files/dynamo0p3}} directory.
These directories are located in the \sphinxcode{\sphinxupquote{\textless{}PSYCLONEHOME\textgreater{}/src/psyclone}}
directory where \sphinxcode{\sphinxupquote{\textless{}PSYCLONEHOME\textgreater{}}} refers to the location where you
download or clone PSyclone ({\hyperref[\detokenize{getting_going:getting-going}]{\sphinxcrossref{\DUrole{std,std-ref}{Getting Going}}}}).

\sphinxAtStartPar
In the \sphinxcode{\sphinxupquote{tests/test\_files/dynamo0p3}} directory the majority of examples
start with \sphinxcode{\sphinxupquote{testkern}}. Amongst the exceptions are: \sphinxcode{\sphinxupquote{testkern\_simple\_mod.f90}},
\sphinxcode{\sphinxupquote{ru\_kernel\_mod.f90}} and \sphinxcode{\sphinxupquote{matrix\_vector\_kernel\_mod.F90}}. The following
test kernels can be used to generate kernel stub code (running stub
generation from the \sphinxcode{\sphinxupquote{\textless{}PSYCLONEHOME\textgreater{}/src/psyclone}} directory):

\begin{sphinxVerbatim}[commandchars=\\\{\}]
\PYG{n}{tests}\PYG{o}{/}\PYG{n}{test\PYGZus{}files}\PYG{o}{/}\PYG{n}{dynamo0p3}\PYG{o}{/}\PYG{n}{testkern\PYGZus{}chi\PYGZus{}read\PYGZus{}mod}\PYG{o}{.}\PYG{n}{F90}
\PYG{n}{tests}\PYG{o}{/}\PYG{n}{test\PYGZus{}files}\PYG{o}{/}\PYG{n}{dynamo0p3}\PYG{o}{/}\PYG{n}{testkern\PYGZus{}coord\PYGZus{}w0\PYGZus{}mod}\PYG{o}{.}\PYG{n}{F90}
\PYG{n}{tests}\PYG{o}{/}\PYG{n}{test\PYGZus{}files}\PYG{o}{/}\PYG{n}{dynamo0p3}\PYG{o}{/}\PYG{n}{testkern\PYGZus{}operator\PYGZus{}mod}\PYG{o}{.}\PYG{n}{f90}
\PYG{n}{tests}\PYG{o}{/}\PYG{n}{test\PYGZus{}files}\PYG{o}{/}\PYG{n}{dynamo0p3}\PYG{o}{/}\PYG{n}{testkern\PYGZus{}operator\PYGZus{}nofield\PYGZus{}mod}\PYG{o}{.}\PYG{n}{f90}
\PYG{n}{tests}\PYG{o}{/}\PYG{n}{test\PYGZus{}files}\PYG{o}{/}\PYG{n}{dynamo0p3}\PYG{o}{/}\PYG{n}{ru\PYGZus{}kernel\PYGZus{}mod}\PYG{o}{.}\PYG{n}{f90}
\PYG{n}{tests}\PYG{o}{/}\PYG{n}{test\PYGZus{}files}\PYG{o}{/}\PYG{n}{dynamo0p3}\PYG{o}{/}\PYG{n}{testkern\PYGZus{}simple\PYGZus{}mod}\PYG{o}{.}\PYG{n}{f90}
\end{sphinxVerbatim}


\subsection{Example}
\label{\detokenize{psyclone_kern:example}}\label{\detokenize{psyclone_kern:stub-generation-example}}
\sphinxAtStartPar
A simple, single field example of a kernel that can be used as input for the
stub generator is found in \sphinxcode{\sphinxupquote{tests/test\_files/dynamo0p3/testkern\_simple\_mod.f90}} and
is shown below:

\phantomsection\label{\detokenize{psyclone_kern:simple-metadata}}\begin{quote}

\begin{sphinxVerbatim}[commandchars=\\\{\}]
\PYG{k}{module }\PYG{n}{simple\PYGZus{}mod}

\PYG{+w}{  }\PYG{k}{use }\PYG{n}{argument\PYGZus{}mod}
\PYG{+w}{  }\PYG{k}{use }\PYG{n}{fs\PYGZus{}continuity\PYGZus{}mod}
\PYG{+w}{  }\PYG{k}{use }\PYG{n}{kernel\PYGZus{}mod}
\PYG{+w}{  }\PYG{k}{use }\PYG{n}{constants\PYGZus{}mod}

\PYG{+w}{  }\PYG{k}{implicit }\PYG{k}{none}

\PYG{k}{  }\PYG{k}{type}\PYG{p}{,}\PYG{+w}{ }\PYG{k}{extends}\PYG{p}{(}\PYG{n}{kernel\PYGZus{}type}\PYG{p}{)}\PYG{+w}{ }\PYG{k+kd}{::}\PYG{+w}{ }\PYG{n}{simple\PYGZus{}type}
\PYG{+w}{    }\PYG{k}{type}\PYG{p}{(}\PYG{n}{arg\PYGZus{}type}\PYG{p}{)}\PYG{p}{,}\PYG{+w}{ }\PYG{k}{dimension}\PYG{p}{(}\PYG{l+m+mi}{1}\PYG{p}{)}\PYG{+w}{ }\PYG{k+kd}{::}\PYG{+w}{ }\PYG{n}{meta\PYGZus{}args}\PYG{+w}{ }\PYG{o}{=}\PYG{+w}{ }\PYG{p}{\PYGZam{}}
\PYG{+w}{         }\PYG{p}{(}\PYG{o}{/}\PYG{+w}{ }\PYG{n}{arg\PYGZus{}type}\PYG{p}{(}\PYG{n}{gh\PYGZus{}field}\PYG{p}{,}\PYG{+w}{ }\PYG{n}{gh\PYGZus{}real}\PYG{p}{,}\PYG{+w}{ }\PYG{n}{gh\PYGZus{}inc}\PYG{p}{,}\PYG{+w}{ }\PYG{n}{w1}\PYG{p}{)}\PYG{+w}{ }\PYG{o}{/}\PYG{p}{)}
\PYG{+w}{    }\PYG{k+kt}{integer}\PYG{+w}{ }\PYG{k+kd}{::}\PYG{+w}{ }\PYG{n}{operates\PYGZus{}on}\PYG{+w}{ }\PYG{o}{=}\PYG{+w}{ }\PYG{n}{cell\PYGZus{}column}
\PYG{+w}{  }\PYG{k}{contains}
\PYG{k}{    }\PYG{k}{procedure}\PYG{p}{,}\PYG{+w}{ }\PYG{k}{nopass}\PYG{+w}{ }\PYG{k+kd}{::}\PYG{+w}{ }\PYG{n}{code}\PYG{+w}{ }\PYG{o}{=}\PYG{o}{\PYGZgt{}}\PYG{+w}{ }\PYG{n}{simple\PYGZus{}code}
\PYG{+w}{  }\PYG{k}{end }\PYG{k}{type }\PYG{n}{simple\PYGZus{}type}

\PYG{k}{contains}

\PYG{k}{  }\PYG{k}{subroutine }\PYG{n}{simple\PYGZus{}code}\PYG{p}{(}\PYG{p}{)}
\PYG{+w}{  }\PYG{k}{end }\PYG{k}{subroutine}

\PYG{k}{end }\PYG{k}{module }\PYG{n}{simple\PYGZus{}mod}
\end{sphinxVerbatim}
\end{quote}

\begin{sphinxadmonition}{note}{Note:}
\sphinxAtStartPar
The module name \sphinxcode{\sphinxupquote{simple\_mod}} and the type name \sphinxcode{\sphinxupquote{simple\_type}}
share the same root \sphinxcode{\sphinxupquote{simple}} and have the extensions \sphinxcode{\sphinxupquote{\_mod}}
and \sphinxcode{\sphinxupquote{\_type}} respectively. This is a convention in LFRic API
and is required by the kernel stub generator as it needs to
determine the name of the type containing the metadata and infers
this by reading the module name. If this rule is not followed the
kernel stub generator will return with an error message
(see Section {\hyperref[\detokenize{psyclone_kern:stub-generation-errors}]{\sphinxcrossref{\DUrole{std,std-ref}{Errors}}}}).
\end{sphinxadmonition}

\begin{sphinxadmonition}{note}{Note:}
\sphinxAtStartPar
Whilst strictly the kernel stub generator only requires the Kernel
metadata to generate the appropriate stub code, the parser that
the generator relies on currently requires a dummy kernel subroutine
to exist.
\end{sphinxadmonition}

\sphinxAtStartPar
If we run the kernel stub generator on the \sphinxcode{\sphinxupquote{testkern\_simple\_mod.f90}} example:

\begin{sphinxVerbatim}[commandchars=\\\{\}]
\PYG{o}{\PYGZgt{}} \PYG{n}{psyclone}\PYG{o}{\PYGZhy{}}\PYG{n}{kern} \PYG{o}{\PYGZhy{}}\PYG{n}{api} \PYG{n}{lfric} \PYG{o}{\PYGZhy{}}\PYG{n}{gen} \PYG{n}{stub} \PYG{n}{tests}\PYG{o}{/}\PYG{n}{test\PYGZus{}files}\PYG{o}{/}\PYG{n}{dynamo0p3}\PYG{o}{/}\PYG{n}{testkern\PYGZus{}simple\PYGZus{}mod}\PYG{o}{.}\PYG{n}{f90}
\end{sphinxVerbatim}

\sphinxAtStartPar
we get the following kernel stub output:
\begin{quote}

\begin{sphinxVerbatim}[commandchars=\\\{\}]
\PYG{k}{MODULE }\PYG{n}{simple\PYGZus{}mod}
\PYG{+w}{  }\PYG{k}{IMPLICIT }\PYG{k}{NONE}
\PYG{k}{  }\PYG{k}{CONTAINS}
\PYG{k}{  }\PYG{k}{SUBROUTINE }\PYG{n}{simple\PYGZus{}code}\PYG{p}{(}\PYG{n}{nlayers}\PYG{p}{,}\PYG{+w}{ }\PYG{n}{field\PYGZus{}1\PYGZus{}w1}\PYG{p}{,}\PYG{+w}{ }\PYG{n}{ndf\PYGZus{}w1}\PYG{p}{,}\PYG{+w}{ }\PYG{n}{undf\PYGZus{}w1}\PYG{p}{,}\PYG{+w}{ }\PYG{n}{map\PYGZus{}w1}\PYG{p}{)}
\PYG{+w}{    }\PYG{k}{USE }\PYG{n}{constants\PYGZus{}mod}\PYG{p}{,}\PYG{+w}{ }\PYG{k}{ONLY}\PYG{p}{:}\PYG{+w}{ }\PYG{n}{r\PYGZus{}def}\PYG{p}{,}\PYG{+w}{ }\PYG{n}{i\PYGZus{}def}
\PYG{+w}{    }\PYG{k}{IMPLICIT }\PYG{k}{NONE}
\PYG{k}{    }\PYG{k+kt}{INTEGER}\PYG{p}{(}\PYG{n+nb}{KIND}\PYG{o}{=}\PYG{n}{i\PYGZus{}def}\PYG{p}{)}\PYG{p}{,}\PYG{+w}{ }\PYG{k}{intent}\PYG{p}{(}\PYG{n}{in}\PYG{p}{)}\PYG{+w}{ }\PYG{k+kd}{::}\PYG{+w}{ }\PYG{n}{nlayers}
\PYG{+w}{    }\PYG{k+kt}{INTEGER}\PYG{p}{(}\PYG{n+nb}{KIND}\PYG{o}{=}\PYG{n}{i\PYGZus{}def}\PYG{p}{)}\PYG{p}{,}\PYG{+w}{ }\PYG{k}{intent}\PYG{p}{(}\PYG{n}{in}\PYG{p}{)}\PYG{+w}{ }\PYG{k+kd}{::}\PYG{+w}{ }\PYG{n}{ndf\PYGZus{}w1}
\PYG{+w}{    }\PYG{k+kt}{INTEGER}\PYG{p}{(}\PYG{n+nb}{KIND}\PYG{o}{=}\PYG{n}{i\PYGZus{}def}\PYG{p}{)}\PYG{p}{,}\PYG{+w}{ }\PYG{k}{intent}\PYG{p}{(}\PYG{n}{in}\PYG{p}{)}\PYG{p}{,}\PYG{+w}{ }\PYG{k}{dimension}\PYG{p}{(}\PYG{n}{ndf\PYGZus{}w1}\PYG{p}{)}\PYG{+w}{ }\PYG{k+kd}{::}\PYG{+w}{ }\PYG{n}{map\PYGZus{}w1}
\PYG{+w}{    }\PYG{k+kt}{INTEGER}\PYG{p}{(}\PYG{n+nb}{KIND}\PYG{o}{=}\PYG{n}{i\PYGZus{}def}\PYG{p}{)}\PYG{p}{,}\PYG{+w}{ }\PYG{k}{intent}\PYG{p}{(}\PYG{n}{in}\PYG{p}{)}\PYG{+w}{ }\PYG{k+kd}{::}\PYG{+w}{ }\PYG{n}{undf\PYGZus{}w1}
\PYG{+w}{    }\PYG{k+kt}{REAL}\PYG{p}{(}\PYG{n+nb}{KIND}\PYG{o}{=}\PYG{n}{r\PYGZus{}def}\PYG{p}{)}\PYG{p}{,}\PYG{+w}{ }\PYG{k}{intent}\PYG{p}{(}\PYG{n}{inout}\PYG{p}{)}\PYG{p}{,}\PYG{+w}{ }\PYG{k}{dimension}\PYG{p}{(}\PYG{n}{undf\PYGZus{}w1}\PYG{p}{)}\PYG{+w}{ }\PYG{k+kd}{::}\PYG{+w}{ }\PYG{n}{field\PYGZus{}1\PYGZus{}w1}
\PYG{+w}{  }\PYG{k}{END }\PYG{k}{SUBROUTINE }\PYG{n}{simple\PYGZus{}code}
\PYG{k}{END }\PYG{k}{MODULE }\PYG{n}{simple\PYGZus{}mod}
\end{sphinxVerbatim}
\end{quote}

\sphinxAtStartPar
The subroutine content can then be copied into the required module,
used as the basis for a new module, or checked with an existing
subroutine for correctness.

\begin{sphinxadmonition}{note}{Note:}
\sphinxAtStartPar
The output does not currently conform to Met Office coding
standards so must be modified accordingly.
\end{sphinxadmonition}

\begin{sphinxadmonition}{note}{Note:}
\sphinxAtStartPar
The code will not compile without a) providing the
\sphinxcode{\sphinxupquote{constants\_mod}}, \sphinxcode{\sphinxupquote{argument\_mod}} and \sphinxcode{\sphinxupquote{kernel\_mod}} modules
in the compiler include path and b) adding in code that writes
to any arguments declared as intent \sphinxcode{\sphinxupquote{out}} or \sphinxcode{\sphinxupquote{inout}}. For a
quick check, the \sphinxcode{\sphinxupquote{USE}} declaration and \sphinxcode{\sphinxupquote{KIND}} declarations
can be removed and the \sphinxcode{\sphinxupquote{field\_1\_w1}} array can be initialised
with some value in the subroutine. At this point the Kernel
should compile successfully.
\end{sphinxadmonition}

\begin{sphinxadmonition}{note}{Note:}
\sphinxAtStartPar
Whilst there is only one field declared in the metadata there
are 5 arguments to the Kernel. The first argument \sphinxcode{\sphinxupquote{nlayers}}
specifies the number of layers in a column for a field. The
second argument is the array associated with the field. The
field array is dimensioned as the \sphinxstyleemphasis{number of unique degrees
of freedom} (hereafter \sphinxcode{\sphinxupquote{undf}}) which is also passed into
the kernel (the fourth argument). The naming convention is to
call each field a \sphinxcode{\sphinxupquote{field}}, followed by its position in the
(algorithm) argument list (which is reflected in the metadata
ordering). The third argument is the number of degrees of freedom
for the particular column and is used to dimension the final
argument which is the \sphinxstyleemphasis{degrees of freedom map} (dofmap) which
indicates the location of the required values in the field array.
The naming convention for the \sphinxcode{\sphinxupquote{dofmap}}, \sphinxcode{\sphinxupquote{undf}} and \sphinxcode{\sphinxupquote{ndf}} is
to append the name with the space that it is associated with.
\end{sphinxadmonition}

\sphinxAtStartPar
We now take a look at a more complicated example. The metadata in this
example is the same as an actual LFRic kernel, however the
subroutine content and various comments have been removed. The metadata
specifies that there are four fields passed by the algorithm layer, the
fourth of which is a vector field of size three. All three of the spaces
require a basis function and the \sphinxcode{\sphinxupquote{W0}} and \sphinxcode{\sphinxupquote{W2}} function spaces
additionally require a differential basis function. The content of the
Kernel, excluding the subroutine body, is given below:
\begin{quote}

\begin{sphinxVerbatim}[commandchars=\\\{\}]
\PYG{k}{module }\PYG{n}{ru\PYGZus{}kernel\PYGZus{}mod}

\PYG{k}{use }\PYG{n}{argument\PYGZus{}mod}
\PYG{k}{use }\PYG{n}{fs\PYGZus{}continuity\PYGZus{}mod}
\PYG{k}{use }\PYG{n}{kernel\PYGZus{}mod}
\PYG{k}{use }\PYG{n}{constants\PYGZus{}mod}

\PYG{k}{implicit }\PYG{k}{none}

\PYG{k}{private}

\PYG{k}{type}\PYG{p}{,}\PYG{+w}{ }\PYG{k}{public}\PYG{p}{,}\PYG{+w}{ }\PYG{k}{extends}\PYG{p}{(}\PYG{n}{kernel\PYGZus{}type}\PYG{p}{)}\PYG{+w}{ }\PYG{k+kd}{::}\PYG{+w}{ }\PYG{n}{ru\PYGZus{}kernel\PYGZus{}type}
\PYG{+w}{  }\PYG{k}{private}
\PYG{k}{  }\PYG{k}{type}\PYG{p}{(}\PYG{n}{arg\PYGZus{}type}\PYG{p}{)}\PYG{+w}{ }\PYG{k+kd}{::}\PYG{+w}{ }\PYG{n}{meta\PYGZus{}args}\PYG{p}{(}\PYG{l+m+mi}{6}\PYG{p}{)}\PYG{+w}{ }\PYG{o}{=}\PYG{+w}{ }\PYG{p}{(}\PYG{o}{/}\PYG{+w}{                                  }\PYG{p}{\PYGZam{}}
\PYG{+w}{       }\PYG{n}{arg\PYGZus{}type}\PYG{p}{(}\PYG{n}{GH\PYGZus{}FIELD}\PYG{p}{,}\PYG{+w}{   }\PYG{n}{GH\PYGZus{}REAL}\PYG{p}{,}\PYG{+w}{    }\PYG{n}{GH\PYGZus{}INC}\PYG{p}{,}\PYG{+w}{  }\PYG{n}{W2}\PYG{p}{)}\PYG{p}{,}\PYG{+w}{                  }\PYG{p}{\PYGZam{}}
\PYG{+w}{       }\PYG{n}{arg\PYGZus{}type}\PYG{p}{(}\PYG{n}{GH\PYGZus{}FIELD}\PYG{p}{,}\PYG{+w}{   }\PYG{n}{GH\PYGZus{}REAL}\PYG{p}{,}\PYG{+w}{    }\PYG{n}{GH\PYGZus{}READ}\PYG{p}{,}\PYG{+w}{ }\PYG{n}{W3}\PYG{p}{)}\PYG{p}{,}\PYG{+w}{                  }\PYG{p}{\PYGZam{}}
\PYG{+w}{       }\PYG{n}{arg\PYGZus{}type}\PYG{p}{(}\PYG{n}{GH\PYGZus{}SCALAR}\PYG{p}{,}\PYG{+w}{  }\PYG{n}{GH\PYGZus{}INTEGER}\PYG{p}{,}\PYG{+w}{ }\PYG{n}{GH\PYGZus{}READ}\PYG{p}{)}\PYG{p}{,}\PYG{+w}{                      }\PYG{p}{\PYGZam{}}
\PYG{+w}{       }\PYG{n}{arg\PYGZus{}type}\PYG{p}{(}\PYG{n}{GH\PYGZus{}SCALAR}\PYG{p}{,}\PYG{+w}{  }\PYG{n}{GH\PYGZus{}REAL}\PYG{p}{,}\PYG{+w}{    }\PYG{n}{GH\PYGZus{}READ}\PYG{p}{)}\PYG{p}{,}\PYG{+w}{                      }\PYG{p}{\PYGZam{}}
\PYG{+w}{       }\PYG{n}{arg\PYGZus{}type}\PYG{p}{(}\PYG{n}{GH\PYGZus{}FIELD}\PYG{p}{,}\PYG{+w}{   }\PYG{n}{GH\PYGZus{}REAL}\PYG{p}{,}\PYG{+w}{    }\PYG{n}{GH\PYGZus{}READ}\PYG{p}{,}\PYG{+w}{ }\PYG{n}{W0}\PYG{p}{)}\PYG{p}{,}\PYG{+w}{                  }\PYG{p}{\PYGZam{}}
\PYG{+w}{       }\PYG{n}{arg\PYGZus{}type}\PYG{p}{(}\PYG{n}{GH\PYGZus{}FIELD}\PYG{o}{*}\PYG{l+m+mi}{3}\PYG{p}{,}\PYG{+w}{ }\PYG{n}{GH\PYGZus{}REAL}\PYG{p}{,}\PYG{+w}{    }\PYG{n}{GH\PYGZus{}READ}\PYG{p}{,}\PYG{+w}{ }\PYG{n}{W0}\PYG{p}{)}\PYG{+w}{                   }\PYG{p}{\PYGZam{}}
\PYG{+w}{       }\PYG{o}{/}\PYG{p}{)}
\PYG{+w}{  }\PYG{k}{type}\PYG{p}{(}\PYG{n}{func\PYGZus{}type}\PYG{p}{)}\PYG{+w}{ }\PYG{k+kd}{::}\PYG{+w}{ }\PYG{n}{meta\PYGZus{}funcs}\PYG{p}{(}\PYG{l+m+mi}{3}\PYG{p}{)}\PYG{+w}{ }\PYG{o}{=}\PYG{+w}{ }\PYG{p}{(}\PYG{o}{/}\PYG{+w}{                                }\PYG{p}{\PYGZam{}}
\PYG{+w}{       }\PYG{n}{func\PYGZus{}type}\PYG{p}{(}\PYG{n}{W2}\PYG{p}{,}\PYG{+w}{ }\PYG{n}{GH\PYGZus{}BASIS}\PYG{p}{,}\PYG{+w}{ }\PYG{n}{GH\PYGZus{}DIFF\PYGZus{}BASIS}\PYG{p}{)}\PYG{p}{,}\PYG{+w}{                         }\PYG{p}{\PYGZam{}}
\PYG{+w}{       }\PYG{n}{func\PYGZus{}type}\PYG{p}{(}\PYG{n}{W3}\PYG{p}{,}\PYG{+w}{ }\PYG{n}{GH\PYGZus{}BASIS}\PYG{p}{)}\PYG{p}{,}\PYG{+w}{                                        }\PYG{p}{\PYGZam{}}
\PYG{+w}{       }\PYG{n}{func\PYGZus{}type}\PYG{p}{(}\PYG{n}{W0}\PYG{p}{,}\PYG{+w}{ }\PYG{n}{GH\PYGZus{}BASIS}\PYG{p}{,}\PYG{+w}{ }\PYG{n}{GH\PYGZus{}DIFF\PYGZus{}BASIS}\PYG{p}{)}\PYG{+w}{                          }\PYG{p}{\PYGZam{}}
\PYG{+w}{       }\PYG{o}{/}\PYG{p}{)}
\PYG{+w}{  }\PYG{k+kt}{integer}\PYG{+w}{ }\PYG{k+kd}{::}\PYG{+w}{ }\PYG{n}{operates\PYGZus{}on}\PYG{+w}{ }\PYG{o}{=}\PYG{+w}{ }\PYG{n}{CELL\PYGZus{}COLUMN}
\PYG{+w}{  }\PYG{k+kt}{integer}\PYG{+w}{ }\PYG{k+kd}{::}\PYG{+w}{ }\PYG{n}{gh\PYGZus{}shape}\PYG{+w}{ }\PYG{o}{=}\PYG{+w}{ }\PYG{n}{gh\PYGZus{}quadrature\PYGZus{}XYoZ}
\PYG{k}{contains}
\PYG{k}{  }\PYG{k}{procedure}\PYG{p}{,}\PYG{+w}{ }\PYG{k}{nopass}\PYG{+w}{ }\PYG{k+kd}{::}\PYG{+w}{ }\PYG{n}{ru\PYGZus{}code}
\PYG{k}{end }\PYG{k}{type}

\PYG{k}{public }\PYG{n}{ru\PYGZus{}code}

\PYG{k}{contains}

\PYG{k}{  }\PYG{k}{subroutine }\PYG{n}{ru\PYGZus{}code}\PYG{p}{(}\PYG{p}{)}
\PYG{+w}{  }\PYG{k}{end }\PYG{k}{subroutine }\PYG{n}{ru\PYGZus{}code}

\PYG{k}{end }\PYG{k}{module }\PYG{n}{ru\PYGZus{}kernel\PYGZus{}mod}
\end{sphinxVerbatim}
\end{quote}

\sphinxAtStartPar
If we run the kernel stub generator on this example:

\begin{sphinxVerbatim}[commandchars=\\\{\}]
\PYG{o}{\PYGZgt{}} \PYG{n}{psyclone}\PYG{o}{\PYGZhy{}}\PYG{n}{kern} \PYG{o}{\PYGZhy{}}\PYG{n}{api} \PYG{n}{lfric} \PYG{o}{\PYGZhy{}}\PYG{n}{gen} \PYG{n}{stub} \PYG{n}{tests}\PYG{o}{/}\PYG{n}{test\PYGZus{}files}\PYG{o}{/}\PYG{n}{dynamo0p3}\PYG{o}{/}\PYG{n}{ru\PYGZus{}kernel\PYGZus{}mod}\PYG{o}{.}\PYG{n}{f90}
\end{sphinxVerbatim}

\sphinxAtStartPar
we obtain the following output:
\begin{quote}

\begin{sphinxVerbatim}[commandchars=\\\{\}]
\PYG{k}{MODULE }\PYG{n}{ru\PYGZus{}mod}
\PYG{+w}{  }\PYG{k}{IMPLICIT }\PYG{k}{NONE}
\PYG{k}{  }\PYG{k}{CONTAINS}
\PYG{k}{  }\PYG{k}{SUBROUTINE }\PYG{n}{ru\PYGZus{}code}\PYG{p}{(}\PYG{n}{nlayers}\PYG{p}{,}\PYG{+w}{ }\PYG{n}{field\PYGZus{}1\PYGZus{}w2}\PYG{p}{,}\PYG{+w}{ }\PYG{n}{field\PYGZus{}2\PYGZus{}w3}\PYG{p}{,}\PYG{+w}{ }\PYG{n}{iscalar\PYGZus{}3}\PYG{p}{,}\PYG{+w}{ }\PYG{n}{rscalar\PYGZus{}4}\PYG{p}{,}\PYG{+w}{ }\PYG{p}{\PYGZam{}}
\PYG{+w}{                     }\PYG{n}{field\PYGZus{}5\PYGZus{}w0}\PYG{p}{,}\PYG{+w}{ }\PYG{n}{field\PYGZus{}6\PYGZus{}w0\PYGZus{}v1}\PYG{p}{,}\PYG{+w}{ }\PYG{n}{field\PYGZus{}6\PYGZus{}w0\PYGZus{}v2}\PYG{p}{,}\PYG{+w}{ }\PYG{n}{field\PYGZus{}6\PYGZus{}w0\PYGZus{}v3}\PYG{p}{,}\PYG{+w}{ }\PYG{p}{\PYGZam{}}
\PYG{+w}{                     }\PYG{n}{ndf\PYGZus{}w2}\PYG{p}{,}\PYG{+w}{ }\PYG{n}{undf\PYGZus{}w2}\PYG{p}{,}\PYG{+w}{ }\PYG{n}{map\PYGZus{}w2}\PYG{p}{,}\PYG{+w}{ }\PYG{n}{basis\PYGZus{}w2\PYGZus{}qr\PYGZus{}xyoz}\PYG{p}{,}\PYG{+w}{ }\PYG{p}{\PYGZam{}}
\PYG{+w}{                     }\PYG{n}{diff\PYGZus{}basis\PYGZus{}w2\PYGZus{}qr\PYGZus{}xyoz}\PYG{p}{,}\PYG{+w}{ }\PYG{n}{ndf\PYGZus{}w3}\PYG{p}{,}\PYG{+w}{ }\PYG{n}{undf\PYGZus{}w3}\PYG{p}{,}\PYG{+w}{ }\PYG{n}{map\PYGZus{}w3}\PYG{p}{,}\PYG{+w}{ }\PYG{p}{\PYGZam{}}
\PYG{+w}{                     }\PYG{n}{basis\PYGZus{}w3\PYGZus{}qr\PYGZus{}xyoz}\PYG{p}{,}\PYG{+w}{ }\PYG{n}{ndf\PYGZus{}w0}\PYG{p}{,}\PYG{+w}{ }\PYG{n}{undf\PYGZus{}w0}\PYG{p}{,}\PYG{+w}{ }\PYG{n}{map\PYGZus{}w0}\PYG{p}{,}\PYG{+w}{ }\PYG{p}{\PYGZam{}}
\PYG{+w}{                     }\PYG{n}{basis\PYGZus{}w0\PYGZus{}qr\PYGZus{}xyoz}\PYG{p}{,}\PYG{+w}{ }\PYG{n}{diff\PYGZus{}basis\PYGZus{}w0\PYGZus{}qr\PYGZus{}xyoz}\PYG{p}{,}\PYG{+w}{ }\PYG{p}{\PYGZam{}}
\PYG{+w}{                     }\PYG{n}{np\PYGZus{}xy\PYGZus{}qr\PYGZus{}xyoz}\PYG{p}{,}\PYG{+w}{ }\PYG{n}{np\PYGZus{}z\PYGZus{}qr\PYGZus{}xyoz}\PYG{p}{,}\PYG{+w}{ }\PYG{n}{weights\PYGZus{}xy\PYGZus{}qr\PYGZus{}xyoz}\PYG{p}{,}\PYG{+w}{ }\PYG{n}{weights\PYGZus{}z\PYGZus{}qr\PYGZus{}xyoz}\PYG{p}{)}
\PYG{+w}{    }\PYG{k}{USE }\PYG{n}{constants\PYGZus{}mod}\PYG{p}{,}\PYG{+w}{ }\PYG{k}{ONLY}\PYG{p}{:}\PYG{+w}{ }\PYG{n}{r\PYGZus{}def}\PYG{p}{,}\PYG{+w}{ }\PYG{n}{i\PYGZus{}def}
\PYG{+w}{    }\PYG{k}{IMPLICIT }\PYG{k}{NONE}
\PYG{k}{    }\PYG{k+kt}{INTEGER}\PYG{p}{(}\PYG{n+nb}{KIND}\PYG{o}{=}\PYG{n}{i\PYGZus{}def}\PYG{p}{)}\PYG{p}{,}\PYG{+w}{ }\PYG{k}{intent}\PYG{p}{(}\PYG{n}{in}\PYG{p}{)}\PYG{+w}{ }\PYG{k+kd}{::}\PYG{+w}{ }\PYG{n}{nlayers}
\PYG{+w}{    }\PYG{k+kt}{INTEGER}\PYG{p}{(}\PYG{n+nb}{KIND}\PYG{o}{=}\PYG{n}{i\PYGZus{}def}\PYG{p}{)}\PYG{p}{,}\PYG{+w}{ }\PYG{k}{intent}\PYG{p}{(}\PYG{n}{in}\PYG{p}{)}\PYG{+w}{ }\PYG{k+kd}{::}\PYG{+w}{ }\PYG{n}{ndf\PYGZus{}w0}
\PYG{+w}{    }\PYG{k+kt}{INTEGER}\PYG{p}{(}\PYG{n+nb}{KIND}\PYG{o}{=}\PYG{n}{i\PYGZus{}def}\PYG{p}{)}\PYG{p}{,}\PYG{+w}{ }\PYG{k}{intent}\PYG{p}{(}\PYG{n}{in}\PYG{p}{)}\PYG{p}{,}\PYG{+w}{ }\PYG{k}{dimension}\PYG{p}{(}\PYG{n}{ndf\PYGZus{}w0}\PYG{p}{)}\PYG{+w}{ }\PYG{k+kd}{::}\PYG{+w}{ }\PYG{n}{map\PYGZus{}w0}
\PYG{+w}{    }\PYG{k+kt}{INTEGER}\PYG{p}{(}\PYG{n+nb}{KIND}\PYG{o}{=}\PYG{n}{i\PYGZus{}def}\PYG{p}{)}\PYG{p}{,}\PYG{+w}{ }\PYG{k}{intent}\PYG{p}{(}\PYG{n}{in}\PYG{p}{)}\PYG{+w}{ }\PYG{k+kd}{::}\PYG{+w}{ }\PYG{n}{ndf\PYGZus{}w2}
\PYG{+w}{    }\PYG{k+kt}{INTEGER}\PYG{p}{(}\PYG{n+nb}{KIND}\PYG{o}{=}\PYG{n}{i\PYGZus{}def}\PYG{p}{)}\PYG{p}{,}\PYG{+w}{ }\PYG{k}{intent}\PYG{p}{(}\PYG{n}{in}\PYG{p}{)}\PYG{p}{,}\PYG{+w}{ }\PYG{k}{dimension}\PYG{p}{(}\PYG{n}{ndf\PYGZus{}w2}\PYG{p}{)}\PYG{+w}{ }\PYG{k+kd}{::}\PYG{+w}{ }\PYG{n}{map\PYGZus{}w2}
\PYG{+w}{    }\PYG{k+kt}{INTEGER}\PYG{p}{(}\PYG{n+nb}{KIND}\PYG{o}{=}\PYG{n}{i\PYGZus{}def}\PYG{p}{)}\PYG{p}{,}\PYG{+w}{ }\PYG{k}{intent}\PYG{p}{(}\PYG{n}{in}\PYG{p}{)}\PYG{+w}{ }\PYG{k+kd}{::}\PYG{+w}{ }\PYG{n}{ndf\PYGZus{}w3}
\PYG{+w}{    }\PYG{k+kt}{INTEGER}\PYG{p}{(}\PYG{n+nb}{KIND}\PYG{o}{=}\PYG{n}{i\PYGZus{}def}\PYG{p}{)}\PYG{p}{,}\PYG{+w}{ }\PYG{k}{intent}\PYG{p}{(}\PYG{n}{in}\PYG{p}{)}\PYG{p}{,}\PYG{+w}{ }\PYG{k}{dimension}\PYG{p}{(}\PYG{n}{ndf\PYGZus{}w3}\PYG{p}{)}\PYG{+w}{ }\PYG{k+kd}{::}\PYG{+w}{ }\PYG{n}{map\PYGZus{}w3}
\PYG{+w}{    }\PYG{k+kt}{INTEGER}\PYG{p}{(}\PYG{n+nb}{KIND}\PYG{o}{=}\PYG{n}{i\PYGZus{}def}\PYG{p}{)}\PYG{p}{,}\PYG{+w}{ }\PYG{k}{intent}\PYG{p}{(}\PYG{n}{in}\PYG{p}{)}\PYG{+w}{ }\PYG{k+kd}{::}\PYG{+w}{ }\PYG{n}{undf\PYGZus{}w2}\PYG{p}{,}\PYG{+w}{ }\PYG{n}{undf\PYGZus{}w3}\PYG{p}{,}\PYG{+w}{ }\PYG{n}{undf\PYGZus{}w0}
\PYG{+w}{    }\PYG{k+kt}{REAL}\PYG{p}{(}\PYG{n+nb}{KIND}\PYG{o}{=}\PYG{n}{r\PYGZus{}def}\PYG{p}{)}\PYG{p}{,}\PYG{+w}{ }\PYG{k}{intent}\PYG{p}{(}\PYG{n}{in}\PYG{p}{)}\PYG{+w}{ }\PYG{k+kd}{::}\PYG{+w}{ }\PYG{n}{rscalar\PYGZus{}4}
\PYG{+w}{    }\PYG{k+kt}{INTEGER}\PYG{p}{(}\PYG{n+nb}{KIND}\PYG{o}{=}\PYG{n}{i\PYGZus{}def}\PYG{p}{)}\PYG{p}{,}\PYG{+w}{ }\PYG{k}{intent}\PYG{p}{(}\PYG{n}{in}\PYG{p}{)}\PYG{+w}{ }\PYG{k+kd}{::}\PYG{+w}{ }\PYG{n}{iscalar\PYGZus{}3}
\PYG{+w}{    }\PYG{k+kt}{REAL}\PYG{p}{(}\PYG{n+nb}{KIND}\PYG{o}{=}\PYG{n}{r\PYGZus{}def}\PYG{p}{)}\PYG{p}{,}\PYG{+w}{ }\PYG{k}{intent}\PYG{p}{(}\PYG{n}{inout}\PYG{p}{)}\PYG{p}{,}\PYG{+w}{ }\PYG{k}{dimension}\PYG{p}{(}\PYG{n}{undf\PYGZus{}w2}\PYG{p}{)}\PYG{+w}{ }\PYG{k+kd}{::}\PYG{+w}{ }\PYG{n}{field\PYGZus{}1\PYGZus{}w2}
\PYG{+w}{    }\PYG{k+kt}{REAL}\PYG{p}{(}\PYG{n+nb}{KIND}\PYG{o}{=}\PYG{n}{r\PYGZus{}def}\PYG{p}{)}\PYG{p}{,}\PYG{+w}{ }\PYG{k}{intent}\PYG{p}{(}\PYG{n}{in}\PYG{p}{)}\PYG{p}{,}\PYG{+w}{ }\PYG{k}{dimension}\PYG{p}{(}\PYG{n}{undf\PYGZus{}w3}\PYG{p}{)}\PYG{+w}{ }\PYG{k+kd}{::}\PYG{+w}{ }\PYG{n}{field\PYGZus{}2\PYGZus{}w3}
\PYG{+w}{    }\PYG{k+kt}{REAL}\PYG{p}{(}\PYG{n+nb}{KIND}\PYG{o}{=}\PYG{n}{r\PYGZus{}def}\PYG{p}{)}\PYG{p}{,}\PYG{+w}{ }\PYG{k}{intent}\PYG{p}{(}\PYG{n}{in}\PYG{p}{)}\PYG{p}{,}\PYG{+w}{ }\PYG{k}{dimension}\PYG{p}{(}\PYG{n}{undf\PYGZus{}w0}\PYG{p}{)}\PYG{+w}{ }\PYG{k+kd}{::}\PYG{+w}{ }\PYG{n}{field\PYGZus{}5\PYGZus{}w0}
\PYG{+w}{    }\PYG{k+kt}{REAL}\PYG{p}{(}\PYG{n+nb}{KIND}\PYG{o}{=}\PYG{n}{r\PYGZus{}def}\PYG{p}{)}\PYG{p}{,}\PYG{+w}{ }\PYG{k}{intent}\PYG{p}{(}\PYG{n}{in}\PYG{p}{)}\PYG{p}{,}\PYG{+w}{ }\PYG{k}{dimension}\PYG{p}{(}\PYG{n}{undf\PYGZus{}w0}\PYG{p}{)}\PYG{+w}{ }\PYG{k+kd}{::}\PYG{+w}{ }\PYG{n}{field\PYGZus{}6\PYGZus{}w0\PYGZus{}v1}
\PYG{+w}{    }\PYG{k+kt}{REAL}\PYG{p}{(}\PYG{n+nb}{KIND}\PYG{o}{=}\PYG{n}{r\PYGZus{}def}\PYG{p}{)}\PYG{p}{,}\PYG{+w}{ }\PYG{k}{intent}\PYG{p}{(}\PYG{n}{in}\PYG{p}{)}\PYG{p}{,}\PYG{+w}{ }\PYG{k}{dimension}\PYG{p}{(}\PYG{n}{undf\PYGZus{}w0}\PYG{p}{)}\PYG{+w}{ }\PYG{k+kd}{::}\PYG{+w}{ }\PYG{n}{field\PYGZus{}6\PYGZus{}w0\PYGZus{}v2}
\PYG{+w}{    }\PYG{k+kt}{REAL}\PYG{p}{(}\PYG{n+nb}{KIND}\PYG{o}{=}\PYG{n}{r\PYGZus{}def}\PYG{p}{)}\PYG{p}{,}\PYG{+w}{ }\PYG{k}{intent}\PYG{p}{(}\PYG{n}{in}\PYG{p}{)}\PYG{p}{,}\PYG{+w}{ }\PYG{k}{dimension}\PYG{p}{(}\PYG{n}{undf\PYGZus{}w0}\PYG{p}{)}\PYG{+w}{ }\PYG{k+kd}{::}\PYG{+w}{ }\PYG{n}{field\PYGZus{}6\PYGZus{}w0\PYGZus{}v3}
\PYG{+w}{    }\PYG{k+kt}{INTEGER}\PYG{p}{(}\PYG{n+nb}{KIND}\PYG{o}{=}\PYG{n}{i\PYGZus{}def}\PYG{p}{)}\PYG{p}{,}\PYG{+w}{ }\PYG{k}{intent}\PYG{p}{(}\PYG{n}{in}\PYG{p}{)}\PYG{+w}{ }\PYG{k+kd}{::}\PYG{+w}{ }\PYG{n}{np\PYGZus{}xy\PYGZus{}qr\PYGZus{}xyoz}\PYG{p}{,}\PYG{+w}{ }\PYG{n}{np\PYGZus{}z\PYGZus{}qr\PYGZus{}xyoz}
\PYG{+w}{    }\PYG{k+kt}{REAL}\PYG{p}{(}\PYG{n+nb}{KIND}\PYG{o}{=}\PYG{n}{r\PYGZus{}def}\PYG{p}{)}\PYG{p}{,}\PYG{+w}{ }\PYG{k}{intent}\PYG{p}{(}\PYG{n}{in}\PYG{p}{)}\PYG{p}{,}\PYG{+w}{ }\PYG{k}{dimension}\PYG{p}{(}\PYG{l+m+mi}{3}\PYG{p}{,}\PYG{n}{ndf\PYGZus{}w2}\PYG{p}{,}\PYG{n}{np\PYGZus{}xy\PYGZus{}qr\PYGZus{}xyoz}\PYG{p}{,}\PYG{n}{np\PYGZus{}z\PYGZus{}qr\PYGZus{}xyoz}\PYG{p}{)}\PYG{+w}{ }\PYG{k+kd}{::}\PYG{+w}{ }\PYG{n}{basis\PYGZus{}w2\PYGZus{}qr\PYGZus{}xyoz}
\PYG{+w}{    }\PYG{k+kt}{REAL}\PYG{p}{(}\PYG{n+nb}{KIND}\PYG{o}{=}\PYG{n}{r\PYGZus{}def}\PYG{p}{)}\PYG{p}{,}\PYG{+w}{ }\PYG{k}{intent}\PYG{p}{(}\PYG{n}{in}\PYG{p}{)}\PYG{p}{,}\PYG{+w}{ }\PYG{k}{dimension}\PYG{p}{(}\PYG{l+m+mi}{1}\PYG{p}{,}\PYG{n}{ndf\PYGZus{}w2}\PYG{p}{,}\PYG{n}{np\PYGZus{}xy\PYGZus{}qr\PYGZus{}xyoz}\PYG{p}{,}\PYG{n}{np\PYGZus{}z\PYGZus{}qr\PYGZus{}xyoz}\PYG{p}{)}\PYG{+w}{ }\PYG{k+kd}{::}\PYG{+w}{ }\PYG{n}{diff\PYGZus{}basis\PYGZus{}w2\PYGZus{}qr\PYGZus{}xyoz}
\PYG{+w}{    }\PYG{k+kt}{REAL}\PYG{p}{(}\PYG{n+nb}{KIND}\PYG{o}{=}\PYG{n}{r\PYGZus{}def}\PYG{p}{)}\PYG{p}{,}\PYG{+w}{ }\PYG{k}{intent}\PYG{p}{(}\PYG{n}{in}\PYG{p}{)}\PYG{p}{,}\PYG{+w}{ }\PYG{k}{dimension}\PYG{p}{(}\PYG{l+m+mi}{1}\PYG{p}{,}\PYG{n}{ndf\PYGZus{}w3}\PYG{p}{,}\PYG{n}{np\PYGZus{}xy\PYGZus{}qr\PYGZus{}xyoz}\PYG{p}{,}\PYG{n}{np\PYGZus{}z\PYGZus{}qr\PYGZus{}xyoz}\PYG{p}{)}\PYG{+w}{ }\PYG{k+kd}{::}\PYG{+w}{ }\PYG{n}{basis\PYGZus{}w3\PYGZus{}qr\PYGZus{}xyoz}
\PYG{+w}{    }\PYG{k+kt}{REAL}\PYG{p}{(}\PYG{n+nb}{KIND}\PYG{o}{=}\PYG{n}{r\PYGZus{}def}\PYG{p}{)}\PYG{p}{,}\PYG{+w}{ }\PYG{k}{intent}\PYG{p}{(}\PYG{n}{in}\PYG{p}{)}\PYG{p}{,}\PYG{+w}{ }\PYG{k}{dimension}\PYG{p}{(}\PYG{l+m+mi}{1}\PYG{p}{,}\PYG{n}{ndf\PYGZus{}w0}\PYG{p}{,}\PYG{n}{np\PYGZus{}xy\PYGZus{}qr\PYGZus{}xyoz}\PYG{p}{,}\PYG{n}{np\PYGZus{}z\PYGZus{}qr\PYGZus{}xyoz}\PYG{p}{)}\PYG{+w}{ }\PYG{k+kd}{::}\PYG{+w}{ }\PYG{n}{basis\PYGZus{}w0\PYGZus{}qr\PYGZus{}xyoz}
\PYG{+w}{    }\PYG{k+kt}{REAL}\PYG{p}{(}\PYG{n+nb}{KIND}\PYG{o}{=}\PYG{n}{r\PYGZus{}def}\PYG{p}{)}\PYG{p}{,}\PYG{+w}{ }\PYG{k}{intent}\PYG{p}{(}\PYG{n}{in}\PYG{p}{)}\PYG{p}{,}\PYG{+w}{ }\PYG{k}{dimension}\PYG{p}{(}\PYG{l+m+mi}{3}\PYG{p}{,}\PYG{n}{ndf\PYGZus{}w0}\PYG{p}{,}\PYG{n}{np\PYGZus{}xy\PYGZus{}qr\PYGZus{}xyoz}\PYG{p}{,}\PYG{n}{np\PYGZus{}z\PYGZus{}qr\PYGZus{}xyoz}\PYG{p}{)}\PYG{+w}{ }\PYG{k+kd}{::}\PYG{+w}{ }\PYG{n}{diff\PYGZus{}basis\PYGZus{}w0\PYGZus{}qr\PYGZus{}xyoz}
\PYG{+w}{    }\PYG{k+kt}{REAL}\PYG{p}{(}\PYG{n+nb}{KIND}\PYG{o}{=}\PYG{n}{r\PYGZus{}def}\PYG{p}{)}\PYG{p}{,}\PYG{+w}{ }\PYG{k}{intent}\PYG{p}{(}\PYG{n}{in}\PYG{p}{)}\PYG{p}{,}\PYG{+w}{ }\PYG{k}{dimension}\PYG{p}{(}\PYG{n}{np\PYGZus{}xy\PYGZus{}qr\PYGZus{}xyoz}\PYG{p}{)}\PYG{+w}{ }\PYG{k+kd}{::}\PYG{+w}{ }\PYG{n}{weights\PYGZus{}xy\PYGZus{}qr\PYGZus{}xyoz}
\PYG{+w}{    }\PYG{k+kt}{REAL}\PYG{p}{(}\PYG{n+nb}{KIND}\PYG{o}{=}\PYG{n}{r\PYGZus{}def}\PYG{p}{)}\PYG{p}{,}\PYG{+w}{ }\PYG{k}{intent}\PYG{p}{(}\PYG{n}{in}\PYG{p}{)}\PYG{p}{,}\PYG{+w}{ }\PYG{k}{dimension}\PYG{p}{(}\PYG{n}{np\PYGZus{}z\PYGZus{}qr\PYGZus{}xyoz}\PYG{p}{)}\PYG{+w}{ }\PYG{k+kd}{::}\PYG{+w}{ }\PYG{n}{weights\PYGZus{}z\PYGZus{}qr\PYGZus{}xyoz}
\PYG{+w}{  }\PYG{k}{END }\PYG{k}{SUBROUTINE }\PYG{n}{ru\PYGZus{}code}
\PYG{k}{END }\PYG{k}{MODULE }\PYG{n}{ru\PYGZus{}mod}
\end{sphinxVerbatim}
\end{quote}

\sphinxAtStartPar
The above example demonstrates that the argument list can get quite
complex. Rather than going through an explanation of each argument you
are referred to Section {\hyperref[\detokenize{dynamo0p3:lfric-stub-generation-rules}]{\sphinxcrossref{\DUrole{std,std-ref}{Rules}}}} for
more details on the rules for argument types and argument ordering.
Regarding naming conventions for arguments you can see that the arrays
associated with the fields are labelled as 1\sphinxhyphen{}6 depending on their
position in the metadata. For a vector field, each vector results in a
different array. These are distinguished by appending \sphinxcode{\sphinxupquote{\_vx}} where \sphinxcode{\sphinxupquote{x}} is
the number of the vector.

\sphinxAtStartPar
The introduction of stencil operations on field arguments further complicates
the argument list of a kernel. An example of the use of the stub generator
for a kernel that performs stencil operations is provided in
\sphinxcode{\sphinxupquote{examples/lfric/eg5}}:

\begin{sphinxVerbatim}[commandchars=\\\{\}]
\PYG{o}{\PYGZgt{}} \PYG{n}{psyclone}\PYG{o}{\PYGZhy{}}\PYG{n}{kern} \PYG{o}{\PYGZhy{}}\PYG{n}{api} \PYG{n}{lfric} \PYG{o}{\PYGZhy{}}\PYG{n}{gen} \PYG{n}{stub} \PYG{o}{.}\PYG{o}{.}\PYG{o}{/}\PYG{o}{.}\PYG{o}{.}\PYG{o}{/}\PYG{n}{examples}\PYG{o}{/}\PYG{n}{lfric}\PYG{o}{/}\PYG{n}{eg5}\PYG{o}{/}\PYG{n}{conservative\PYGZus{}flux\PYGZus{}kernel\PYGZus{}mod}\PYG{o}{.}\PYG{n}{F90}
\end{sphinxVerbatim}


\subsection{Errors}
\label{\detokenize{psyclone_kern:errors}}\label{\detokenize{psyclone_kern:stub-generation-errors}}
\sphinxAtStartPar
The stub generator has been written to provide useful errors if
mistakes are found. If you run the generator and it does not produce a
useful error \sphinxhyphen{} and in particular if it produces a stack trace \sphinxhyphen{} please
contact the PSyclone developers.

\sphinxAtStartPar
The following tests do not produce stub kernel code either because
they are invalid or because they contain functionality that is not
supported in the stub generator:

\begin{sphinxVerbatim}[commandchars=\\\{\}]
\PYG{n}{tests}\PYG{o}{/}\PYG{n}{test\PYGZus{}files}\PYG{o}{/}\PYG{n}{dynamo0p3}\PYG{o}{/}\PYG{n}{testkern\PYGZus{}any\PYGZus{}space\PYGZus{}1\PYGZus{}mod}\PYG{o}{.}\PYG{n}{f90}
\PYG{n}{tests}\PYG{o}{/}\PYG{n}{test\PYGZus{}files}\PYG{o}{/}\PYG{n}{dynamo0p3}\PYG{o}{/}\PYG{n}{testkern\PYGZus{}any\PYGZus{}space\PYGZus{}4\PYGZus{}mod}\PYG{o}{.}\PYG{n}{f90}
\PYG{n}{tests}\PYG{o}{/}\PYG{n}{test\PYGZus{}files}\PYG{o}{/}\PYG{n}{dynamo0p3}\PYG{o}{/}\PYG{n}{testkern\PYGZus{}any\PYGZus{}discontinuous\PYGZus{}space\PYGZus{}op\PYGZus{}2\PYGZus{}mod}\PYG{o}{.}\PYG{n}{f90}
\PYG{n}{tests}\PYG{o}{/}\PYG{n}{test\PYGZus{}files}\PYG{o}{/}\PYG{n}{dynamo0p3}\PYG{o}{/}\PYG{n}{testkern\PYGZus{}dofs\PYGZus{}mod}\PYG{o}{.}\PYG{n}{f90}
\PYG{n}{tests}\PYG{o}{/}\PYG{n}{test\PYGZus{}files}\PYG{o}{/}\PYG{n}{dynamo0p3}\PYG{o}{/}\PYG{n}{testkern\PYGZus{}invalid\PYGZus{}fortran\PYGZus{}mod}\PYG{o}{.}\PYG{n}{f90}
\PYG{n}{tests}\PYG{o}{/}\PYG{n}{test\PYGZus{}files}\PYG{o}{/}\PYG{n}{dynamo0p3}\PYG{o}{/}\PYG{n}{testkern\PYGZus{}short\PYGZus{}name\PYGZus{}mod}\PYG{o}{.}\PYG{n}{f90}
\PYG{n}{tests}\PYG{o}{/}\PYG{n}{test\PYGZus{}files}\PYG{o}{/}\PYG{n}{dynamo0p3}\PYG{o}{/}\PYG{n}{testkern\PYGZus{}no\PYGZus{}datatype\PYGZus{}mod}\PYG{o}{.}\PYG{n}{f90}
\PYG{n}{tests}\PYG{o}{/}\PYG{n}{test\PYGZus{}files}\PYG{o}{/}\PYG{n}{dynamo0p3}\PYG{o}{/}\PYG{n}{testkern\PYGZus{}wrong\PYGZus{}file\PYGZus{}name}\PYG{o}{.}\PYG{n}{F90}
\end{sphinxVerbatim}

\sphinxAtStartPar
\sphinxcode{\sphinxupquote{testkern\_invalid\_fortran\_mod.f90}}, \sphinxcode{\sphinxupquote{testkern\_no\_datatype\_mod.f90}},
\sphinxcode{\sphinxupquote{testkern\_short\_name\_mod.f90}} and \sphinxcode{\sphinxupquote{testkern\_wrong\_file\_name.F90}} are designed to be
invalid for PSyclone stub generation testing purposes and should produce
appropriate errors. Two examples are below:

\begin{sphinxVerbatim}[commandchars=\\\{\}]
\PYG{o}{\PYGZgt{}} \PYG{n}{psyclone}\PYG{o}{\PYGZhy{}}\PYG{n}{kern} \PYG{o}{\PYGZhy{}}\PYG{n}{api} \PYG{n}{lfric} \PYG{o}{\PYGZhy{}}\PYG{n}{gen} \PYG{n}{stub} \PYG{n}{tests}\PYG{o}{/}\PYG{n}{test\PYGZus{}files}\PYG{o}{/}\PYG{n}{dynamo0p3}\PYG{o}{/}\PYG{n}{testkern\PYGZus{}invalid\PYGZus{}fortran\PYGZus{}mod}\PYG{o}{.}\PYG{n}{f90}
\PYG{n}{Error}\PYG{p}{:} \PYG{l+s+s1}{\PYGZsq{}}\PYG{l+s+s1}{Parse Error: Code appears to be invalid Fortran}\PYG{l+s+s1}{\PYGZsq{}}

\PYG{o}{\PYGZgt{}} \PYG{n}{psyclone}\PYG{o}{\PYGZhy{}}\PYG{n}{kern} \PYG{o}{\PYGZhy{}}\PYG{n}{api} \PYG{n}{lfric} \PYG{o}{\PYGZhy{}}\PYG{n}{gen} \PYG{n}{stub} \PYG{n}{tests}\PYG{o}{/}\PYG{n}{test\PYGZus{}files}\PYG{o}{/}\PYG{n}{dynamo0p3}\PYG{o}{/}\PYG{n}{testkern\PYGZus{}no\PYGZus{}datatype\PYGZus{}mod}\PYG{o}{.}\PYG{n}{f90}
\PYG{n}{Error}\PYG{p}{:} \PYG{l+s+s1}{\PYGZsq{}}\PYG{l+s+s1}{Parse Error: Kernel type testkern\PYGZus{}type does not exist}\PYG{l+s+s1}{\PYGZsq{}}
\end{sphinxVerbatim}

\sphinxAtStartPar
\sphinxcode{\sphinxupquote{testkern\_dofs\_mod.f90}} is an example with an unsupported feature, as the
\sphinxcode{\sphinxupquote{operates\_on}} metadata specifies \sphinxcode{\sphinxupquote{dof}}. Currently only kernels with
\sphinxcode{\sphinxupquote{operates\_on=CELL\_COLUMN}} are supported by the stub generator.

\sphinxAtStartPar
Generic function space metadata \sphinxcode{\sphinxupquote{any\_space}} and \sphinxcode{\sphinxupquote{any\_discontinuous\_space}}
(see Section {\hyperref[\detokenize{dynamo0p3:lfric-function-space}]{\sphinxcrossref{\DUrole{std,std-ref}{Supported Function Spaces}}}}
for function\sphinxhyphen{}space identifiers) are currently only supported for
{\hyperref[\detokenize{dynamo0p3:lfric-field}]{\sphinxcrossref{\DUrole{std,std-ref}{LFRic fields}}}} in the stub generator. Basis
and differential basis functions on these generic function spaces, required
for {\hyperref[\detokenize{dynamo0p3:lfric-quadrature}]{\sphinxcrossref{\DUrole{std,std-ref}{quadrature}}}} and
{\hyperref[\detokenize{dynamo0p3:lfric-gh-shape}]{\sphinxcrossref{\DUrole{std,std-ref}{evaluators}}}}, are not supported. Hence,
\sphinxcode{\sphinxupquote{testkern\_any\_space\_1\_mod.f90}}, \sphinxcode{\sphinxupquote{testkern\_any\_space\_4\_mod.f90}} and
\sphinxcode{\sphinxupquote{testkern\_any\_discontinuous\_space\_op\_2\_mod.f90}} should fail with
appropriate warnings because of that. For example:

\begin{sphinxVerbatim}[commandchars=\\\{\}]
\PYG{o}{\PYGZgt{}} \PYG{n}{psyclone}\PYG{o}{\PYGZhy{}}\PYG{n}{kern} \PYG{o}{\PYGZhy{}}\PYG{n}{api} \PYG{n}{lfric} \PYG{o}{\PYGZhy{}}\PYG{n}{gen} \PYG{n}{stub} \PYG{n}{tests}\PYG{o}{/}\PYG{n}{test\PYGZus{}files}\PYG{o}{/}\PYG{n}{dynamo0p3}\PYG{o}{/}\PYG{n}{testkern\PYGZus{}any\PYGZus{}space\PYGZus{}1\PYGZus{}mod}\PYG{o}{.}\PYG{n}{f90}
\PYG{n}{Error}\PYG{p}{:} \PYG{l+s+s2}{\PYGZdq{}}\PYG{l+s+s2}{Generation Error: Unsupported space for basis function, expecting}
\PYG{n}{one} \PYG{n}{of} \PYG{p}{[}\PYG{l+s+s1}{\PYGZsq{}}\PYG{l+s+s1}{w3}\PYG{l+s+s1}{\PYGZsq{}}\PYG{p}{,} \PYG{l+s+s1}{\PYGZsq{}}\PYG{l+s+s1}{wtheta}\PYG{l+s+s1}{\PYGZsq{}}\PYG{p}{,} \PYG{l+s+s1}{\PYGZsq{}}\PYG{l+s+s1}{w2v}\PYG{l+s+s1}{\PYGZsq{}}\PYG{p}{,} \PYG{l+s+s1}{\PYGZsq{}}\PYG{l+s+s1}{w2vtrace}\PYG{l+s+s1}{\PYGZsq{}}\PYG{p}{,} \PYG{l+s+s1}{\PYGZsq{}}\PYG{l+s+s1}{w2broken}\PYG{l+s+s1}{\PYGZsq{}}\PYG{p}{,} \PYG{l+s+s1}{\PYGZsq{}}\PYG{l+s+s1}{w0}\PYG{l+s+s1}{\PYGZsq{}}\PYG{p}{,} \PYG{l+s+s1}{\PYGZsq{}}\PYG{l+s+s1}{w1}\PYG{l+s+s1}{\PYGZsq{}}\PYG{p}{,} \PYG{l+s+s1}{\PYGZsq{}}\PYG{l+s+s1}{w2}\PYG{l+s+s1}{\PYGZsq{}}\PYG{p}{,}
\PYG{l+s+s1}{\PYGZsq{}}\PYG{l+s+s1}{w2trace}\PYG{l+s+s1}{\PYGZsq{}}\PYG{p}{,} \PYG{l+s+s1}{\PYGZsq{}}\PYG{l+s+s1}{w2h}\PYG{l+s+s1}{\PYGZsq{}}\PYG{p}{,} \PYG{l+s+s1}{\PYGZsq{}}\PYG{l+s+s1}{w2htrace}\PYG{l+s+s1}{\PYGZsq{}}\PYG{p}{,} \PYG{l+s+s1}{\PYGZsq{}}\PYG{l+s+s1}{any\PYGZus{}w2}\PYG{l+s+s1}{\PYGZsq{}}\PYG{p}{,} \PYG{l+s+s1}{\PYGZsq{}}\PYG{l+s+s1}{wchi}\PYG{l+s+s1}{\PYGZsq{}}\PYG{p}{]} \PYG{n}{but} \PYG{n}{found} \PYG{l+s+s1}{\PYGZsq{}}\PYG{l+s+s1}{any\PYGZus{}space\PYGZus{}1}\PYG{l+s+s1}{\PYGZsq{}}\PYG{l+s+s2}{\PYGZdq{}}
\end{sphinxVerbatim}

\sphinxAtStartPar
As noted above, if the LFRic API naming convention for module and type
names is not followed, the stub generator will return with an error
message. For example:

\begin{sphinxVerbatim}[commandchars=\\\{\}]
\PYG{o}{\PYGZgt{}} \PYG{n}{psyclone}\PYG{o}{\PYGZhy{}}\PYG{n}{kern} \PYG{o}{\PYGZhy{}}\PYG{n}{api} \PYG{n}{lfric} \PYG{o}{\PYGZhy{}}\PYG{n}{gen} \PYG{n}{stub} \PYG{n}{tests}\PYG{o}{/}\PYG{n}{test\PYGZus{}files}\PYG{o}{/}\PYG{n}{dynamo0p3}\PYG{o}{/}\PYG{n}{testkern\PYGZus{}wrong\PYGZus{}file\PYGZus{}name}\PYG{o}{.}\PYG{n}{F90}
\PYG{n}{Error}\PYG{p}{:} \PYG{l+s+s2}{\PYGZdq{}}\PYG{l+s+s2}{Parse Error: Error, module name }\PYG{l+s+s2}{\PYGZsq{}}\PYG{l+s+s2}{testkern\PYGZus{}wrong\PYGZus{}file\PYGZus{}name}\PYG{l+s+s2}{\PYGZsq{}}\PYG{l+s+s2}{ does not have}
\PYG{l+s+s1}{\PYGZsq{}}\PYG{l+s+s1}{\PYGZus{}mod}\PYG{l+s+s1}{\PYGZsq{}} \PYG{k}{as} \PYG{n}{an} \PYG{n}{extension}\PYG{o}{.} \PYG{n}{This} \PYG{n}{convention} \PYG{o+ow}{is} \PYG{n}{assumed}\PYG{o}{.}\PYG{l+s+s2}{\PYGZdq{}}
\end{sphinxVerbatim}


\section{Algorithm Generator}
\label{\detokenize{psyclone_kern:algorithm-generator}}\label{\detokenize{psyclone_kern:alg-generation}}

\subsection{Quick Start}
\label{\detokenize{psyclone_kern:id1}}\begin{enumerate}
\sphinxsetlistlabels{\arabic}{enumi}{enumii}{}{)}%
\item {} 
\sphinxAtStartPar
Use an existing Kernel file containing a full LFRic kernel implementation.

\item {} 
\sphinxAtStartPar
Run the following command

\begin{sphinxVerbatim}[commandchars=\\\{\}]
\PYG{o}{\PYGZgt{}} \PYG{n}{psyclone}\PYG{o}{\PYGZhy{}}\PYG{n}{kern} \PYG{o}{\PYGZhy{}}\PYG{n}{api} \PYG{n}{lfric} \PYG{o}{\PYGZhy{}}\PYG{n}{gen} \PYG{n}{alg} \PYG{o}{\PYGZlt{}}\PYG{n}{PATH}\PYG{o}{\PYGZgt{}}\PYG{o}{/}\PYG{n}{my\PYGZus{}kern\PYGZus{}file\PYGZus{}mod}\PYG{o}{.}\PYG{n}{f90}
\end{sphinxVerbatim}

\item {} 
\sphinxAtStartPar
The generated Algorithm code will be output to stdout by default. To have
it written to a file use the \sphinxtitleref{\sphinxhyphen{}o} flag.

\end{enumerate}


\subsection{Introduction}
\label{\detokenize{psyclone_kern:id2}}
\sphinxAtStartPar
The ability to generate a valid LFRic Algorithm layer that calls a given kernel
is useful for a number of reasons:
\begin{enumerate}
\sphinxsetlistlabels{\arabic}{enumi}{enumii}{}{)}%
\item {} 
\sphinxAtStartPar
Starting point for creating a test for a kernel;

\item {} 
\sphinxAtStartPar
Benchmarking an individual kernel;

\item {} 
\sphinxAtStartPar
Constructing a test harness for the adjoint of a kernel produced by
\sphinxhref{https://psyclone-adjoint.readthedocs.io/en/latest/introduction.html\#introduction}{\DUrole{xref,std,std-ref}{PSyAD}}.

\end{enumerate}

\sphinxAtStartPar
Currently algorithm generation is only supported for the LFRic API but it
could be extended to the GOcean API if desired.


\subsubsection{Mapping of Function Spaces}
\label{\detokenize{psyclone_kern:mapping-of-function-spaces}}
\sphinxAtStartPar
Every field or operator argument to an LFRic kernel must have its
function space(s) specified in the metadata of the kernel. This information
is used by the algorithm generation to ensure that each kernel argument
is correctly constructed. However, the metadata permits the use of certain
‘generic’ function\sphinxhyphen{}space specifiers (see {\hyperref[\detokenize{dynamo0p3:lfric-function-space}]{\sphinxcrossref{\DUrole{std,std-ref}{Supported Function Spaces}}}}). If an argument is specified as being on one of
these spaces then the algorithm generator chooses an appropriate, specific
function space for that argument. e.g. an argument that is specified as
being on \sphinxcode{\sphinxupquote{ANY\_SPACE\_\textless{}n\textgreater{}}} will be constructed on \sphinxcode{\sphinxupquote{W0}} while one on
\sphinxcode{\sphinxupquote{ANY\_DISCONTINUOUS\_SPACE\_\textless{}n\textgreater{}}} will be constructed on \sphinxcode{\sphinxupquote{W3}}.


\subsection{Example}
\label{\detokenize{psyclone_kern:id3}}
\sphinxAtStartPar
If we take the same kernel used in the stub\sphinxhyphen{}generation
{\hyperref[\detokenize{psyclone_kern:stub-generation-example}]{\sphinxcrossref{\DUrole{std,std-ref}{example}}}} then running

\begin{sphinxVerbatim}[commandchars=\\\{\}]
\PYG{o}{\PYGZgt{}} \PYG{n}{psyclone}\PYG{o}{\PYGZhy{}}\PYG{n}{kern} \PYG{o}{\PYGZhy{}}\PYG{n}{api} \PYG{n}{lfric} \PYG{o}{\PYGZhy{}}\PYG{n}{gen} \PYG{n}{alg} \PYG{n}{tests}\PYG{o}{/}\PYG{n}{test\PYGZus{}files}\PYG{o}{/}\PYG{n}{dynamo0p3}\PYG{o}{/}\PYG{n}{testkern\PYGZus{}simple\PYGZus{}mod}\PYG{o}{.}\PYG{n}{f90}
\end{sphinxVerbatim}

\sphinxAtStartPar
gives the following algorithm layer code:
\begin{quote}

\begin{sphinxVerbatim}[commandchars=\\\{\}]
\PYG{k}{module }\PYG{n}{test\PYGZus{}alg\PYGZus{}mod}
\PYG{+w}{  }\PYG{k}{implicit }\PYG{k}{none}
\PYG{k}{  }\PYG{k}{public}

\PYG{k}{contains}
\PYG{k}{  }\PYG{k}{subroutine }\PYG{n}{test\PYGZus{}alg}\PYG{p}{(}\PYG{n}{mesh}\PYG{p}{,}\PYG{+w}{ }\PYG{n}{chi}\PYG{p}{,}\PYG{+w}{ }\PYG{n}{panel\PYGZus{}id}\PYG{p}{)}
\PYG{+w}{    }\PYG{k}{use }\PYG{n}{field\PYGZus{}mod}\PYG{p}{,}\PYG{+w}{ }\PYG{k}{only}\PYG{+w}{ }\PYG{p}{:}\PYG{+w}{ }\PYG{n}{field\PYGZus{}type}
\PYG{+w}{    }\PYG{k}{use }\PYG{n}{function\PYGZus{}space\PYGZus{}mod}\PYG{p}{,}\PYG{+w}{ }\PYG{k}{only}\PYG{+w}{ }\PYG{p}{:}\PYG{+w}{ }\PYG{n}{function\PYGZus{}space\PYGZus{}type}
\PYG{+w}{    }\PYG{k}{use }\PYG{n}{fs\PYGZus{}continuity\PYGZus{}mod}\PYG{p}{,}\PYG{+w}{ }\PYG{k}{only}\PYG{+w}{ }\PYG{p}{:}\PYG{+w}{ }\PYG{n}{w1}
\PYG{+w}{    }\PYG{k}{use }\PYG{n}{function\PYGZus{}space\PYGZus{}collection\PYGZus{}mod}\PYG{p}{,}\PYG{+w}{ }\PYG{k}{only}\PYG{+w}{ }\PYG{p}{:}\PYG{+w}{ }\PYG{n}{function\PYGZus{}space\PYGZus{}collection}
\PYG{+w}{    }\PYG{k}{use }\PYG{n}{mesh\PYGZus{}mod}\PYG{p}{,}\PYG{+w}{ }\PYG{k}{only}\PYG{+w}{ }\PYG{p}{:}\PYG{+w}{ }\PYG{n}{mesh\PYGZus{}type}
\PYG{+w}{    }\PYG{k}{use }\PYG{n}{simple\PYGZus{}mod}\PYG{p}{,}\PYG{+w}{ }\PYG{k}{only}\PYG{+w}{ }\PYG{p}{:}\PYG{+w}{ }\PYG{n}{simple\PYGZus{}type}
\PYG{+w}{    }\PYG{k}{use }\PYG{n}{constants\PYGZus{}mod}\PYG{p}{,}\PYG{+w}{ }\PYG{k}{only}\PYG{+w}{ }\PYG{p}{:}\PYG{+w}{ }\PYG{n}{i\PYGZus{}def}\PYG{p}{,}\PYG{+w}{ }\PYG{n}{r\PYGZus{}def}
\PYG{+w}{    }\PYG{k+kt}{integer}\PYG{p}{(}\PYG{n+nb}{kind}\PYG{o}{=}\PYG{n}{i\PYGZus{}def}\PYG{p}{)}\PYG{p}{,}\PYG{+w}{ }\PYG{k}{parameter}\PYG{+w}{ }\PYG{k+kd}{::}\PYG{+w}{ }\PYG{n}{element\PYGZus{}order\PYGZus{}h}\PYG{+w}{ }\PYG{o}{=}\PYG{+w}{ }\PYG{l+m+mi}{1\PYGZus{}i\PYGZus{}def}
\PYG{+w}{    }\PYG{k+kt}{integer}\PYG{p}{(}\PYG{n+nb}{kind}\PYG{o}{=}\PYG{n}{i\PYGZus{}def}\PYG{p}{)}\PYG{p}{,}\PYG{+w}{ }\PYG{k}{parameter}\PYG{+w}{ }\PYG{k+kd}{::}\PYG{+w}{ }\PYG{n}{element\PYGZus{}order\PYGZus{}v}\PYG{+w}{ }\PYG{o}{=}\PYG{+w}{ }\PYG{l+m+mi}{1\PYGZus{}i\PYGZus{}def}
\PYG{+w}{    }\PYG{k}{type}\PYG{p}{(}\PYG{n}{mesh\PYGZus{}type}\PYG{p}{)}\PYG{p}{,}\PYG{+w}{ }\PYG{k}{pointer}\PYG{p}{,}\PYG{+w}{ }\PYG{k}{intent}\PYG{p}{(}\PYG{n}{in}\PYG{p}{)}\PYG{+w}{ }\PYG{k+kd}{::}\PYG{+w}{ }\PYG{n}{mesh}
\PYG{+w}{    }\PYG{k}{type}\PYG{p}{(}\PYG{n}{field\PYGZus{}type}\PYG{p}{)}\PYG{p}{,}\PYG{+w}{ }\PYG{k}{dimension}\PYG{p}{(}\PYG{l+m+mi}{3}\PYG{p}{)}\PYG{p}{,}\PYG{+w}{ }\PYG{k}{intent}\PYG{p}{(}\PYG{n}{in}\PYG{p}{)}\PYG{p}{,}\PYG{+w}{ }\PYG{k}{optional}\PYG{+w}{ }\PYG{k+kd}{::}\PYG{+w}{ }\PYG{n}{chi}
\PYG{+w}{    }\PYG{k}{type}\PYG{p}{(}\PYG{n}{field\PYGZus{}type}\PYG{p}{)}\PYG{p}{,}\PYG{+w}{ }\PYG{k}{intent}\PYG{p}{(}\PYG{n}{in}\PYG{p}{)}\PYG{p}{,}\PYG{+w}{ }\PYG{k}{optional}\PYG{+w}{ }\PYG{k+kd}{::}\PYG{+w}{ }\PYG{n}{panel\PYGZus{}id}
\PYG{+w}{    }\PYG{k}{TYPE}\PYG{p}{(}\PYG{n}{function\PYGZus{}space\PYGZus{}type}\PYG{p}{)}\PYG{p}{,}\PYG{+w}{ }\PYG{k}{POINTER}\PYG{+w}{ }\PYG{k+kd}{::}\PYG{+w}{ }\PYG{n}{vector\PYGZus{}space\PYGZus{}w1\PYGZus{}ptr}
\PYG{+w}{    }\PYG{k}{type}\PYG{p}{(}\PYG{n}{field\PYGZus{}type}\PYG{p}{)}\PYG{+w}{ }\PYG{k+kd}{::}\PYG{+w}{ }\PYG{n}{field\PYGZus{}1}

\PYG{+w}{    }\PYG{n}{vector\PYGZus{}space\PYGZus{}w1\PYGZus{}ptr}\PYG{+w}{ }\PYG{o}{=}\PYG{o}{\PYGZgt{}}\PYG{+w}{ }\PYG{n}{function\PYGZus{}space\PYGZus{}collection}\PYG{+w}{ }\PYG{p}{\PYGZpc{}}\PYG{+w}{ }\PYG{n}{get\PYGZus{}fs}\PYG{p}{(}\PYG{n}{mesh}\PYG{p}{,}\PYG{+w}{ }\PYG{n}{element\PYGZus{}order\PYGZus{}h}\PYG{p}{,}\PYG{+w}{ }\PYG{n}{element\PYGZus{}order\PYGZus{}v}\PYG{p}{,}\PYG{+w}{ }\PYG{n}{w1}\PYG{p}{)}
\PYG{+w}{    }\PYG{k}{call }\PYG{n}{field\PYGZus{}1}\PYG{+w}{ }\PYG{p}{\PYGZpc{}}\PYG{+w}{ }\PYG{n}{initialise}\PYG{p}{(}\PYG{n}{vector\PYGZus{}space}\PYG{o}{=}\PYG{n}{vector\PYGZus{}space\PYGZus{}w1\PYGZus{}ptr}\PYG{p}{,}\PYG{+w}{ }\PYG{n}{name}\PYG{o}{=}\PYG{l+s+s1}{\PYGZsq{}field\PYGZus{}1\PYGZsq{}}\PYG{p}{)}
\PYG{+w}{    }\PYG{k}{call }\PYG{n}{invoke}\PYG{p}{(}\PYG{n}{setval\PYGZus{}c}\PYG{p}{(}\PYG{n}{field\PYGZus{}1}\PYG{p}{,}\PYG{+w}{ }\PYG{l+m+mf}{1.0\PYGZus{}r\PYGZus{}def}\PYG{p}{)}\PYG{p}{,}\PYG{+w}{ }\PYG{n}{simple\PYGZus{}type}\PYG{p}{(}\PYG{n}{field\PYGZus{}1}\PYG{p}{)}\PYG{p}{)}

\PYG{+w}{  }\PYG{k}{end }\PYG{k}{subroutine }\PYG{n}{test\PYGZus{}alg}

\PYG{k}{end }\PYG{k}{module }\PYG{n}{test\PYGZus{}alg\PYGZus{}mod}
\end{sphinxVerbatim}
\end{quote}

\sphinxAtStartPar
Note that the generated code implements an Algorithm subroutine that
is intended to be called from within an LFRic application that has
already setup data structures for the mesh (and, optionally, the \sphinxtitleref{chi}
coordinate field and panel ID mapping).  Since the {\hyperref[\detokenize{psyclone_kern:simple-metadata}]{\sphinxcrossref{\DUrole{std,std-ref}{metadata}}}} for the \sphinxtitleref{simple\_type} kernel specifies that the
field argument is on \sphinxtitleref{W1}, the generated code must ensure that the
appropriate function space is set up and used to initialise the field.
Once that’s done, the interesting part is the \sphinxtitleref{invoke} call:
\begin{quote}

\begin{sphinxVerbatim}[commandchars=\\\{\}]
\PYG{k}{call }\PYG{n}{invoke}\PYG{p}{(}\PYG{n}{setval\PYGZus{}c}\PYG{p}{(}\PYG{n}{field\PYGZus{}1}\PYG{p}{,}\PYG{+w}{ }\PYG{l+m+mf}{1.0\PYGZus{}r\PYGZus{}def}\PYG{p}{)}\PYG{p}{,}\PYG{+w}{ }\PYG{p}{\PYGZam{}}
\PYG{+w}{            }\PYG{n}{simple\PYGZus{}type}\PYG{p}{(}\PYG{n}{field\PYGZus{}1}\PYG{p}{)}\PYG{p}{)}
\end{sphinxVerbatim}
\end{quote}

\sphinxAtStartPar
(where a line\sphinxhyphen{}break has been added for clarity). In this example the \sphinxtitleref{invoke}
is for two kernels: the first is a {\hyperref[\detokenize{introduction_to_psykal:built-ins}]{\sphinxcrossref{\DUrole{std,std-ref}{Built\sphinxhyphen{}in}}}} that gives
\sphinxtitleref{field\_1} the value \sphinxtitleref{1.0} everywhere and the second is the ‘simple’ kernel
itself which is passed the now initialised \sphinxtitleref{field\_1}.

\sphinxAtStartPar
This Algorithm code can now be processed by PSyclone in the normal way in
order to generate a transformed version plus an associated PSy\sphinxhyphen{}layer routine.
See {\hyperref[\detokenize{tutorials_and_examples:lfric-alg-gen-example}]{\sphinxcrossref{\DUrole{std,std-ref}{Example 20: Algorithm Generation}}}} for a full example of doing this.


\subsection{Limitations}
\label{\detokenize{psyclone_kern:limitations}}\begin{itemize}
\item {} 
\sphinxAtStartPar
Algorithm generation is only currently supported for the LFRic API.

\item {} 
\sphinxAtStartPar
All fields are currently set to unity. Obviously the generated algorithm
code may be edited to change this.

\item {} 
\sphinxAtStartPar
The generator does not currently recognise ‘special’ fields that hold
geometry information (such as Chi or the face IDs) and these too will
all be initialised to unity. This is the subject of Issue \#1708 (although
note that the generated code already permits the caller to supply Chi and/or
face IDs).

\item {} 
\sphinxAtStartPar
Kernels with operator arguments are not yet supported.

\item {} 
\sphinxAtStartPar
Kernels with stencil accesses are not yet supported.

\end{itemize}

\sphinxstepscope


\chapter{Line length}
\label{\detokenize{line_length:line-length}}\label{\detokenize{line_length:id1}}\label{\detokenize{line_length::doc}}
\sphinxAtStartPar
By default PSyclone will generate Fortran code with no consideration
of Fortran line\sphinxhyphen{}length limits.  As the line\sphinxhyphen{}length limit for
free\sphinxhyphen{}form Fortran is 132 characters, the code that is output may be
non\sphinxhyphen{}conformant.

\sphinxAtStartPar
Line length is not an issue for many compilers as they allow compiler
flags to be set which allow lines longer than the Fortran
standard. However this is not the case for all compilers.

\sphinxAtStartPar
PSyclone therefore supports the wrapping of lines within the 132
character limit. The next two sections discuss how this is done when
scripting and when working interactively respectively.


\section{Script}
\label{\detokenize{line_length:script}}
\sphinxAtStartPar
The \sphinxtitleref{psyclone} script provides the \sphinxhyphen{}l option to wrap lines. Please
see the {\hyperref[\detokenize{psyclone_command:fort-line-length}]{\sphinxcrossref{\DUrole{std,std-ref}{Fortran line length}}}} section for more details.


\section{Interactive}
\label{\detokenize{line_length:interactive}}
\sphinxAtStartPar
When using PSyclone interactively the line lengths of the input
algorithm and Kernel files can be checked by setting the
\sphinxcode{\sphinxupquote{psyclone.parse.algorithm.parse()}} function’s \sphinxcode{\sphinxupquote{line\_length}}
argument to \sphinxcode{\sphinxupquote{True}}.

\begin{sphinxVerbatim}[commandchars=\\\{\}]
\PYG{g+gp}{\PYGZgt{}\PYGZgt{}\PYGZgt{} }\PYG{k+kn}{from}\PYG{+w}{ }\PYG{n+nn}{psyclone}\PYG{n+nn}{.}\PYG{n+nn}{parse}\PYG{n+nn}{.}\PYG{n+nn}{algorithm}\PYG{+w}{ }\PYG{k+kn}{import} \PYG{n}{parse}
\PYG{g+gp}{\PYGZgt{}\PYGZgt{}\PYGZgt{} }\PYG{n}{ast}\PYG{p}{,} \PYG{n}{info} \PYG{o}{=} \PYG{n}{parse}\PYG{p}{(}\PYG{l+s+s2}{\PYGZdq{}}\PYG{l+s+s2}{argspec.F90}\PYG{l+s+s2}{\PYGZdq{}}\PYG{p}{,} \PYG{n}{line\PYGZus{}length}\PYG{o}{=}\PYG{k+kc}{True}\PYG{p}{)}
\end{sphinxVerbatim}

\sphinxAtStartPar
Similarly the \sphinxcode{\sphinxupquote{line\_length}} argument can be set to \sphinxcode{\sphinxupquote{True}} if calling
the \sphinxcode{\sphinxupquote{generator.generate()}} function. This function simply passes this
argument on to the \sphinxcode{\sphinxupquote{psyclone.parse.algorithm.parse()}} function.

\begin{sphinxVerbatim}[commandchars=\\\{\}]
\PYG{g+gp}{\PYGZgt{}\PYGZgt{}\PYGZgt{} }\PYG{k+kn}{from}\PYG{+w}{ }\PYG{n+nn}{psyclone}\PYG{n+nn}{.}\PYG{n+nn}{generator}\PYG{+w}{ }\PYG{k+kn}{import} \PYG{n}{generate}
\PYG{g+gp}{\PYGZgt{}\PYGZgt{}\PYGZgt{} }\PYG{n}{alg}\PYG{p}{,} \PYG{n}{psy} \PYG{o}{=} \PYG{n}{generate}\PYG{p}{(}\PYG{l+s+s2}{\PYGZdq{}}\PYG{l+s+s2}{argspec.F90}\PYG{l+s+s2}{\PYGZdq{}}\PYG{p}{,} \PYG{n}{line\PYGZus{}length}\PYG{o}{=}\PYG{k+kc}{True}\PYG{p}{)}
\end{sphinxVerbatim}

\sphinxAtStartPar
Line wrapping is performed as a post\sphinxhyphen{}processing step, i.e. after the
code has been generated. This is done by an instance of the
\sphinxcode{\sphinxupquote{line\_length.FortLineLength}} class. For example:

\begin{sphinxVerbatim}[commandchars=\\\{\}]
\PYG{g+gp}{\PYGZgt{}\PYGZgt{}\PYGZgt{} }\PYG{k+kn}{from}\PYG{+w}{ }\PYG{n+nn}{psyclone}\PYG{n+nn}{.}\PYG{n+nn}{generator}\PYG{+w}{ }\PYG{k+kn}{import} \PYG{n}{generate}
\PYG{g+gp}{\PYGZgt{}\PYGZgt{}\PYGZgt{} }\PYG{k+kn}{from}\PYG{+w}{ }\PYG{n+nn}{psyclone}\PYG{n+nn}{.}\PYG{n+nn}{line\PYGZus{}length}\PYG{+w}{ }\PYG{k+kn}{import} \PYG{n}{FortLineLength}
\PYG{g+gp}{\PYGZgt{}\PYGZgt{}\PYGZgt{} }\PYG{n}{psy}\PYG{p}{,} \PYG{n}{alg} \PYG{o}{=} \PYG{n}{generate}\PYG{p}{(}\PYG{l+s+s2}{\PYGZdq{}}\PYG{l+s+s2}{algspec.f90}\PYG{l+s+s2}{\PYGZdq{}}\PYG{p}{,} \PYG{n}{line\PYGZus{}length}\PYG{o}{=}\PYG{k+kc}{True}\PYG{p}{)}
\PYG{g+gp}{\PYGZgt{}\PYGZgt{}\PYGZgt{} }\PYG{n}{line\PYGZus{}length} \PYG{o}{=} \PYG{n}{FortLineLength}\PYG{p}{(}\PYG{p}{)}
\PYG{g+gp}{\PYGZgt{}\PYGZgt{}\PYGZgt{} }\PYG{n}{psy\PYGZus{}str} \PYG{o}{=} \PYG{n}{line\PYGZus{}length}\PYG{o}{.}\PYG{n}{process}\PYG{p}{(}\PYG{n+nb}{str}\PYG{p}{(}\PYG{n}{psy}\PYG{p}{)}\PYG{p}{)}
\PYG{g+gp}{\PYGZgt{}\PYGZgt{}\PYGZgt{} }\PYG{n+nb}{print} \PYG{n}{psy\PYGZus{}str}
\PYG{g+gp}{\PYGZgt{}\PYGZgt{}\PYGZgt{} }\PYG{n}{alg\PYGZus{}str} \PYG{o}{=} \PYG{n}{line\PYGZus{}length}\PYG{o}{.}\PYG{n}{process}\PYG{p}{(}\PYG{n+nb}{str}\PYG{p}{(}\PYG{n}{alg}\PYG{p}{)}\PYG{p}{)}
\PYG{g+gp}{\PYGZgt{}\PYGZgt{}\PYGZgt{} }\PYG{n+nb}{print} \PYG{n}{alg\PYGZus{}str}
\end{sphinxVerbatim}


\section{Limitations}
\label{\detokenize{line_length:limitations}}\label{\detokenize{line_length:line-length-limitations}}
\sphinxAtStartPar
The \sphinxcode{\sphinxupquote{line\_length.FortLineLength}} class is only partially aware
of Fortran syntax. This awareness is required so that appropriate
continuation characters can be used (for example \sphinxcode{\sphinxupquote{\&}} at the end of a
line and \sphinxcode{\sphinxupquote{!\$omp\&}} at the start of a line for OpenMP directives, \sphinxcode{\sphinxupquote{\&}} at
the end of a line for statements and \sphinxcode{\sphinxupquote{\&}} at the end of a line and \sphinxcode{\sphinxupquote{\&}}
at the beginning of a line for strings).

\sphinxAtStartPar
Whilst statements only require an \sphinxcode{\sphinxupquote{\&}} at the end of the line when line
wrapping with free\sphinxhyphen{}form fortran they may optionally also have an \sphinxcode{\sphinxupquote{\&}}
at the beginning of the subsequent line. In contrast, when splitting a
string over multiple lines an \sphinxcode{\sphinxupquote{\&}} is required at both
locations. Therefore an instance of the
\sphinxcode{\sphinxupquote{line\_length.FortLineLength}} class will always add \sphinxcode{\sphinxupquote{\&}} at the
beginning of a continuation line for a statement, in case the line is
split within a string.

\sphinxAtStartPar
One known situation that could cause an instance of the
\sphinxcode{\sphinxupquote{line\_length.FortLineLength}} class to fail is when an inline
comment is used at the end of a line to make it longer than the 132
character limit. Whilst PSyclone does not generate such code for the
PSy\sphinxhyphen{}layer, this might occur in Algorithm\sphinxhyphen{}layer code, even if the
Algorithm\sphinxhyphen{}layer code conforms to the 132 line length limit. The reason
for this is that PSyclone’s internal parser concatenates lines
together, thus a long line correctly split with continuation characters
in the Algorithm\sphinxhyphen{}layer becomes a line that needs to be split by an
instance of the \sphinxcode{\sphinxupquote{line\_length.FortLineLength}} class.

\sphinxstepscope


\chapter{Fortran Naming Conventions}
\label{\detokenize{fortran_naming_conventions:fortran-naming-conventions}}\label{\detokenize{fortran_naming_conventions:fortran-naming}}\label{\detokenize{fortran_naming_conventions::doc}}
\sphinxAtStartPar
There is a convention in the kernel code for the Dynamo0.3 and
GOcean1.0 APIs that if the name of the operation being performed is
\sphinxcode{\sphinxupquote{\textless{}name\textgreater{}}} then a kernel file is \sphinxcode{\sphinxupquote{\textless{}name\textgreater{}\_mod.{[}fF90{]}}}, the name of
the module inside the kernel file is \sphinxcode{\sphinxupquote{\textless{}name\textgreater{}\_mod}}, the name of the
kernel metadata in the module is \sphinxcode{\sphinxupquote{\textless{}name\textgreater{}\_type}} and the name of the
kernel subroutine in the module is \sphinxcode{\sphinxupquote{\textless{}name\textgreater{}\_code}}.

\sphinxAtStartPar
PSyclone itself does not rely on this convention apart from in the
stub generator (see the {\hyperref[\detokenize{psyclone_kern:stub-generation}]{\sphinxcrossref{\DUrole{std,std-ref}{Kernel\sphinxhyphen{}stub Generator}}}} Section) where the name
of the metadata to be parsed is determined from the module name.

\sphinxstepscope


\chapter{PSyData API}
\label{\detokenize{psy_data:psydata-api}}\label{\detokenize{psy_data:psy-data}}\label{\detokenize{psy_data::doc}}
\sphinxAtStartPar
PSyclone provides transformations that will insert callbacks to
an external library at runtime. These callbacks allow third\sphinxhyphen{}party
libraries to access data structures at specified locations in the
code. The PSyclone {\hyperref[\detokenize{libraries:libraries}]{\sphinxcrossref{\DUrole{std,std-ref}{wrappers}}}} to external libraries
are provided in \sphinxcode{\sphinxupquote{share/psyclone/lib}} in PSyclone {\hyperref[\detokenize{getting_going:getting-going-install-loc}]{\sphinxcrossref{\DUrole{std,std-ref}{Installation location}}}}.
Some example use cases are:
\begin{description}
\sphinxlineitem{Profiling:}
\sphinxAtStartPar
By inserting callbacks before and after a region of code,
performance measurements can be added. PSyclone provides
wrapper libraries for some common performance profiling tools,
see {\hyperref[\detokenize{profiling:profiling}]{\sphinxcrossref{\DUrole{std,std-ref}{Profiling}}}} for details.

\sphinxlineitem{Kernel Data Extraction:}
\sphinxAtStartPar
PSyclone provides the ability to add callbacks that provide access
to all input variables before, and output variables after a kernel
invocation. This can be used to automatically create tests for
a kernel, or to write a stand\sphinxhyphen{}alone driver that just calls one
kernel, which can be used for performance tuning. Two example
libraries that extract input and output data into either a Fortran
binary or a NetCDF file are included with PSyclone (see
{\hyperref[\detokenize{psyke:extraction-libraries}]{\sphinxcrossref{\DUrole{std,std-ref}{Extraction Libraries}}}}).

\sphinxlineitem{Access Verification:}
\sphinxAtStartPar
The callbacks can be used to make sure a field declared as read\sphinxhyphen{}only
is not modified during a kernel call (either because of an incorrect
declaration, or because memory is overwritten). The implementation
included in PSyclone uses a simple 64\sphinxhyphen{}bit checksum to detect changes
to a field (and scalar values). See {\hyperref[\detokenize{psy_data:psydata-read-verification}]{\sphinxcrossref{\DUrole{std,std-ref}{Read\sphinxhyphen{}Only Verification}}}}
for details.

\sphinxlineitem{Value Range Check:}
\sphinxAtStartPar
The callbacks can be used to make sure that all floating point input
and output parameters of a kernel are within a user\sphinxhyphen{}specified range.
Additionally, it will also verify that the values are not a \sphinxcode{\sphinxupquote{NaN}}
(not\sphinxhyphen{}a\sphinxhyphen{}number) or infinite. See {\hyperref[\detokenize{psy_data:psydata-value-range-check}]{\sphinxcrossref{\DUrole{std,std-ref}{Value Range Check}}}} for
the full description.

\sphinxlineitem{In\sphinxhyphen{}situ Visualisation:}
\sphinxAtStartPar
By giving access to output fields of a kernel, an in\sphinxhyphen{}situ visualisation
library can be used to plot fields while a (PSyclone\sphinxhyphen{}processed)
application is running. There is no example library available at
this stage, but the API has been designed with this application in mind.

\end{description}

\sphinxAtStartPar
The PSyData API should be general enough to allow these and other
applications to be developed and used.

\sphinxAtStartPar
PSyclone provides transformations that will insert callbacks to
the PSyData API, for example \sphinxcode{\sphinxupquote{ProfileTrans}}, \sphinxcode{\sphinxupquote{GOceanExtractTrans}}
and \sphinxcode{\sphinxupquote{LFRicExtractTrans}}. A user can develop additional transformations
and corresponding runtime libraries for additional functionality.
Refer to \sphinxhref{https://psyclone-dev.readthedocs.io/en/latest/psy\_data.html\#psy-data}{psy\_data} for full details about the PSyData API.


\section{Read\sphinxhyphen{}Only Verification}
\label{\detokenize{psy_data:read-only-verification}}\label{\detokenize{psy_data:psydata-read-verification}}
\sphinxAtStartPar
The PSyData interface is being used to verify that read\sphinxhyphen{}only variables
in a kernel are not overwritten. The \sphinxcode{\sphinxupquote{ReadOnlyVerifyTrans}} (in
\sphinxcode{\sphinxupquote{psyir.transformations.read\_only\_verify\_trans}}, or the
\sphinxhref{https://psyclone-ref.readthedocs.io/en/latest/autogenerated/psyclone.psyir.transformations.html\#classes}{Transformation Reference Guide}) uses the dependency
analysis to determine all read\sphinxhyphen{}only variables (i.e. arguments declared
to be read\sphinxhyphen{}only in metadata, most implicit arguments in LFRic, grid
properties in GOcean). A simple 64\sphinxhyphen{}bit checksum is then computed for all
these arguments before a kernel call, and compared with the checksum
after the kernel call. Any change in the checksum causes a message to
be printed at runtime, e.g.:

\begin{sphinxVerbatim}[commandchars=\\\{\}]
\PYG{o}{\PYGZhy{}}\PYG{o}{\PYGZhy{}}\PYG{o}{\PYGZhy{}}\PYG{o}{\PYGZhy{}}\PYG{o}{\PYGZhy{}}\PYG{o}{\PYGZhy{}}\PYG{o}{\PYGZhy{}}\PYG{o}{\PYGZhy{}}\PYG{o}{\PYGZhy{}}\PYG{o}{\PYGZhy{}}\PYG{o}{\PYGZhy{}}\PYG{o}{\PYGZhy{}}\PYG{o}{\PYGZhy{}}\PYG{o}{\PYGZhy{}}\PYG{o}{\PYGZhy{}}\PYG{o}{\PYGZhy{}}\PYG{o}{\PYGZhy{}}\PYG{o}{\PYGZhy{}}\PYG{o}{\PYGZhy{}}\PYG{o}{\PYGZhy{}}\PYG{o}{\PYGZhy{}}\PYG{o}{\PYGZhy{}}\PYG{o}{\PYGZhy{}}\PYG{o}{\PYGZhy{}}\PYG{o}{\PYGZhy{}}\PYG{o}{\PYGZhy{}}\PYG{o}{\PYGZhy{}}\PYG{o}{\PYGZhy{}}\PYG{o}{\PYGZhy{}}\PYG{o}{\PYGZhy{}}\PYG{o}{\PYGZhy{}}\PYG{o}{\PYGZhy{}}\PYG{o}{\PYGZhy{}}\PYG{o}{\PYGZhy{}}\PYG{o}{\PYGZhy{}}\PYG{o}{\PYGZhy{}}\PYG{o}{\PYGZhy{}}\PYG{o}{\PYGZhy{}}
\PYG{n}{Double} \PYG{n}{precision} \PYG{n}{field} \PYG{n}{b\PYGZus{}fld} \PYG{n}{has} \PYG{n}{been} \PYG{n}{modified} \PYG{o+ow}{in} \PYG{n}{main} \PYG{p}{:} \PYG{n}{update}
\PYG{n}{Original} \PYG{n}{checksum}\PYG{p}{:}   \PYG{l+m+mi}{4611686018427387904}
\PYG{n}{New} \PYG{n}{checksum}\PYG{p}{:}        \PYG{l+m+mi}{4638355772470722560}
\PYG{o}{\PYGZhy{}}\PYG{o}{\PYGZhy{}}\PYG{o}{\PYGZhy{}}\PYG{o}{\PYGZhy{}}\PYG{o}{\PYGZhy{}}\PYG{o}{\PYGZhy{}}\PYG{o}{\PYGZhy{}}\PYG{o}{\PYGZhy{}}\PYG{o}{\PYGZhy{}}\PYG{o}{\PYGZhy{}}\PYG{o}{\PYGZhy{}}\PYG{o}{\PYGZhy{}}\PYG{o}{\PYGZhy{}}\PYG{o}{\PYGZhy{}}\PYG{o}{\PYGZhy{}}\PYG{o}{\PYGZhy{}}\PYG{o}{\PYGZhy{}}\PYG{o}{\PYGZhy{}}\PYG{o}{\PYGZhy{}}\PYG{o}{\PYGZhy{}}\PYG{o}{\PYGZhy{}}\PYG{o}{\PYGZhy{}}\PYG{o}{\PYGZhy{}}\PYG{o}{\PYGZhy{}}\PYG{o}{\PYGZhy{}}\PYG{o}{\PYGZhy{}}\PYG{o}{\PYGZhy{}}\PYG{o}{\PYGZhy{}}\PYG{o}{\PYGZhy{}}\PYG{o}{\PYGZhy{}}\PYG{o}{\PYGZhy{}}\PYG{o}{\PYGZhy{}}\PYG{o}{\PYGZhy{}}\PYG{o}{\PYGZhy{}}\PYG{o}{\PYGZhy{}}\PYG{o}{\PYGZhy{}}\PYG{o}{\PYGZhy{}}\PYG{o}{\PYGZhy{}}
\end{sphinxVerbatim}

\sphinxAtStartPar
The transformation that adds read\sphinxhyphen{}only\sphinxhyphen{}verification to an application
can be applied for both the {\hyperref[\detokenize{dynamo0p3:lfric-api}]{\sphinxcrossref{\DUrole{std,std-ref}{LFRic}}}} and
{\hyperref[\detokenize{gocean1p0:gocean-api}]{\sphinxcrossref{\DUrole{std,std-ref}{GOcean API}}}} \sphinxhyphen{} no API\sphinxhyphen{}specific transformations are required.
It can also be used to verify existing Fortran code.
Below is an example that searches for each loop in a PSyKAl invoke code (which
will always surround kernel calls) and applies the transformation to each one.
This code has been successfully used as a global transformation with the LFRic
Gravity Wave application (the executable is named \sphinxcode{\sphinxupquote{gravity\_wave}})

\begin{sphinxVerbatim}[commandchars=\\\{\}]
\PYG{n}{def}\PYG{+w}{ }\PYG{n}{trans}\PYG{p}{(}\PYG{n}{psyir}\PYG{p}{)}\PYG{p}{:}
\PYG{+w}{    }\PYG{n}{from}\PYG{+w}{ }\PYG{n}{psyclone}\PYG{p}{.}\PYG{n}{psyir}\PYG{p}{.}\PYG{n}{transformations}\PYG{+w}{ }\PYG{k}{import }\PYG{n}{ReadOnlyVerifyTrans}
\PYG{+w}{    }\PYG{n}{from}\PYG{+w}{ }\PYG{n}{psyclone}\PYG{p}{.}\PYG{n}{psyir}\PYG{p}{.}\PYG{n}{nodes}\PYG{+w}{ }\PYG{k}{import }\PYG{n}{Loop}
\PYG{+w}{    }\PYG{n}{read\PYGZus{}only\PYGZus{}verify}\PYG{+w}{ }\PYG{o}{=}\PYG{+w}{ }\PYG{n}{ReadOnlyVerifyTrans}\PYG{p}{(}\PYG{p}{)}

\PYG{+w}{    }\PYG{n}{for}\PYG{+w}{ }\PYG{n}{loop}\PYG{+w}{ }\PYG{n}{in}\PYG{+w}{ }\PYG{n}{psyir}\PYG{p}{.}\PYG{n}{walk}\PYG{p}{(}\PYG{n}{Loop}\PYG{p}{)}\PYG{p}{:}
\PYG{+w}{        }\PYG{n}{read\PYGZus{}only\PYGZus{}verify}\PYG{p}{.}\PYG{n}{apply}\PYG{p}{(}\PYG{n}{loop}\PYG{p}{)}
\end{sphinxVerbatim}

\sphinxAtStartPar
Besides the transformation, a library is required to do the actual
verification at runtime. There are three implementations of the
read\sphinxhyphen{}only\sphinxhyphen{}verification library included in PSyclone: one for LFRic,
one for GOcean, and one for generic Fortran code.
These libraries support the environment variable \sphinxcode{\sphinxupquote{PSYDATA\_VERBOSE}}.
This can be used to control how much output is generated
by the read\sphinxhyphen{}only\sphinxhyphen{}verification library at runtime. If the
variable is not specified or has the value ‘0’, warnings will only
be printed if checksums change. If it is set to ‘1’, a message will be
printed before and after each kernel call that is checked. If the
variable is set to ‘2’, it will additionally print the name of each
variable that is checked.


\subsection{Read\sphinxhyphen{}Only Verification Library for LFRic}
\label{\detokenize{psy_data:read-only-verification-library-for-lfric}}
\sphinxAtStartPar
This library is contained in \sphinxcode{\sphinxupquote{lib/read\_only/lfric}} and it must be compiled
before compiling any LFRic\sphinxhyphen{}based application that uses read\sphinxhyphen{}only verification.
Compiling this library requires access to the LFRic infrastructure library
(since it must implement a generic interface for e.g. the LFRic
{\hyperref[\detokenize{dynamo0p3:lfric-field}]{\sphinxcrossref{\DUrole{std,std-ref}{field}}}} class).

\sphinxAtStartPar
The \sphinxcode{\sphinxupquote{Makefile}} uses the variable \sphinxcode{\sphinxupquote{LFRIC\_INF\_DIR}} to point to the
location where LFRic’s \sphinxcode{\sphinxupquote{field\_mod}} and \sphinxcode{\sphinxupquote{integer\_field\_mod}} have been
compiled. It defaults to the path to location of the pared\sphinxhyphen{}down LFRic
infrastructure located in a clone of PSyclone repository,
\sphinxcode{\sphinxupquote{\textless{}PSYCLONEHOME\textgreater{}/src/psyclone/tests/test\_files/dynamo0p3/infrastructure}},
but this will certainly need to be changed for any user (for instance with
PSyclone installation). The LFRic infrastructure library is not used in
linking the verification library. The application which uses the
read\sphinxhyphen{}only\sphinxhyphen{}verification library needs to link in the infrastructure
library anyway.

\begin{sphinxadmonition}{note}{Note:}
\sphinxAtStartPar
It is the responsibility of the user to make sure that the infrastructure
files used during compilation of the read\sphinxhyphen{}only\sphinxhyphen{}verification library are
also used when linking the application. Otherwise strange and
non\sphinxhyphen{}reproducible crashes might happen.
\end{sphinxadmonition}

\sphinxAtStartPar
Compilation of the library is done by invoking \sphinxcode{\sphinxupquote{make}} and setting
the required variables:

\begin{sphinxVerbatim}[commandchars=\\\{\}]
make\PYG{+w}{ }\PYG{n+nv}{LFRIC\PYGZus{}INF\PYGZus{}DIR}\PYG{o}{=}some\PYGZus{}path\PYG{+w}{ }\PYG{n+nv}{F90}\PYG{o}{=}ifort\PYG{+w}{ }\PYG{n+nv}{F90FLAGS}\PYG{o}{=}\PYG{l+s+s2}{\PYGZdq{}\PYGZhy{}\PYGZhy{}some\PYGZhy{}flag\PYGZdq{}}
\end{sphinxVerbatim}

\sphinxAtStartPar
This will create a library called \sphinxcode{\sphinxupquote{lib\_read\_only.a}}.

\sphinxAtStartPar
An executable example for using the LFRic read\sphinxhyphen{}only\sphinxhyphen{}verification library is
included in \sphinxcode{\sphinxupquote{tutorial/practicals/LFRic/building\_code/4\_psydata}} directory,
see \sphinxhref{https://github.com/stfc/PSyclone/tree/master/tutorial/practicals/LFRic/building\_code/4\_psydata}{this link for more information}.


\subsection{Read\sphinxhyphen{}Only\sphinxhyphen{}Verification Library for GOcean}
\label{\detokenize{psy_data:read-only-verification-library-for-gocean}}
\sphinxAtStartPar
This library is contained in the \sphinxcode{\sphinxupquote{lib/read\_only/dl\_esm\_inf}} directory and
it must be compiled before linking any GOcean\sphinxhyphen{}based application that uses
read\sphinxhyphen{}only verification. Compiling this library requires access to the
GOcean infrastructure library (since it must implement a generic interface
for e.g. the dl\_esm\_inf \sphinxcode{\sphinxupquote{r2d\_field}} class).

\sphinxAtStartPar
The \sphinxcode{\sphinxupquote{Makefile}} uses the variable \sphinxcode{\sphinxupquote{GOCEAN\_INF\_DIR}} to point to the
location where dl\_esm\_inf’s \sphinxcode{\sphinxupquote{field\_mod}} has been compiled. It
defaults to the relative path to location of the dl\_esm\_inf version
included in PSyclone repository as a Git submodule,
\sphinxcode{\sphinxupquote{\textless{}PSYCLONEHOME\textgreater{}/external/dl\_esm\_inf/finite\_difference/src}}. It can be
changed to a user\sphinxhyphen{}specified location if required (for instance with the
PSyclone installation).

\sphinxAtStartPar
The dl\_esm\_inf library is not used in linking the verification library.
The application which uses the read\sphinxhyphen{}only\sphinxhyphen{}verification library needs to
link in the infrastructure library anyway.

\sphinxAtStartPar
Compilation of the library is done by invoking \sphinxcode{\sphinxupquote{make}} and setting
the required variables:

\begin{sphinxVerbatim}[commandchars=\\\{\}]
make\PYG{+w}{ }\PYG{n+nv}{GOCEAN\PYGZus{}INF\PYGZus{}DIR}\PYG{o}{=}some\PYGZus{}path\PYG{+w}{ }\PYG{n+nv}{F90}\PYG{o}{=}ifort\PYG{+w}{ }\PYG{n+nv}{F90FLAGS}\PYG{o}{=}\PYG{l+s+s2}{\PYGZdq{}\PYGZhy{}\PYGZhy{}some\PYGZhy{}flag\PYGZdq{}}
\end{sphinxVerbatim}

\sphinxAtStartPar
This will create a library called \sphinxcode{\sphinxupquote{lib\_read\_only.a}}.
An executable example for using the GOcean read\sphinxhyphen{}only\sphinxhyphen{}verification
library is included in \sphinxcode{\sphinxupquote{examples/gocean/eg5/readonly}}, see
{\hyperref[\detokenize{tutorials_and_examples:gocean-example-readonly}]{\sphinxcrossref{\DUrole{std,std-ref}{Example 5.3: Read\sphinxhyphen{}only\sphinxhyphen{}verification}}}}.


\subsection{Read\sphinxhyphen{}Only\sphinxhyphen{}Verification Library for Generic Fortran}
\label{\detokenize{psy_data:read-only-verification-library-for-generic-fortran}}
\sphinxAtStartPar
This library is contained in the \sphinxcode{\sphinxupquote{lib/read\_only/generic}} directory and
it must be compiled before linking any existing Fortran code that uses
read\sphinxhyphen{}only verification.

\sphinxAtStartPar
Compilation of the library is done by invoking \sphinxcode{\sphinxupquote{make}} and setting
the required variables:

\begin{sphinxVerbatim}[commandchars=\\\{\}]
make\PYG{+w}{ }\PYG{n+nv}{F90}\PYG{o}{=}ifort\PYG{+w}{ }\PYG{n+nv}{F90FLAGS}\PYG{o}{=}\PYG{l+s+s2}{\PYGZdq{}\PYGZhy{}\PYGZhy{}some\PYGZhy{}flag\PYGZdq{}}
\end{sphinxVerbatim}

\sphinxAtStartPar
This will create a library called \sphinxcode{\sphinxupquote{lib\_read\_only.a}}.
An executable example for using the generic read\sphinxhyphen{}only\sphinxhyphen{}verification
library is included in \sphinxcode{\sphinxupquote{examples/nemo/eg6/}}.


\section{Value Range Check}
\label{\detokenize{psy_data:value-range-check}}\label{\detokenize{psy_data:psydata-value-range-check}}
\sphinxAtStartPar
This transformation can be used for both LFRic and GOcean APIs. It will
test all input and output parameters of a kernel to make sure they are
within a user\sphinxhyphen{}specified range. Additionally, it will also verify that floating
point values are not \sphinxcode{\sphinxupquote{NaN}} or infinite.

\sphinxAtStartPar
At runtime, environment variables must be specified to indicate which variables
are within what expected range, and optionally also at which location.
The range is specified as a \sphinxcode{\sphinxupquote{:}} separated tuple:

\begin{sphinxVerbatim}[commandchars=\\\{\}]
\PYG{l+m+mf}{1.1}\PYG{p}{:}\PYG{l+m+mf}{3.3}   \PYG{n}{A} \PYG{n}{value} \PYG{n}{between} \PYG{l+m+mf}{1.1} \PYG{o+ow}{and} \PYG{l+m+mf}{3.3} \PYG{p}{(}\PYG{n}{inclusive}\PYG{p}{)}\PYG{o}{.}
\PYG{p}{:}\PYG{l+m+mf}{3.3}      \PYG{n}{A} \PYG{n}{value} \PYG{n}{less} \PYG{n}{than} \PYG{o+ow}{or} \PYG{n}{equal} \PYG{n}{to} \PYG{l+m+mf}{3.3}
\PYG{l+m+mf}{1.1}\PYG{p}{:}      \PYG{n}{A} \PYG{n}{value} \PYG{n}{greater} \PYG{n}{than} \PYG{o+ow}{or} \PYG{n}{equal} \PYG{n}{to} \PYG{l+m+mf}{1.1}
\end{sphinxVerbatim}

\sphinxAtStartPar
The syntax for the environment variable is one of:
\begin{description}
\sphinxlineitem{\sphinxcode{\sphinxupquote{PSYVERIFY\_\_module\_\_kernel\_\_variable}}}
\sphinxAtStartPar
The specified variable is tested when calling the specified kernel in the
specified module.

\sphinxlineitem{\sphinxcode{\sphinxupquote{PSYVERIFY\_\_module\_\_variable}}}
\sphinxAtStartPar
The specified variable name is tested in all kernel calls of the
specified module that are instrumented with the ValueRangeCheckTrans
transformation.

\sphinxlineitem{\sphinxcode{\sphinxupquote{PSYVERIFY\_\_variable}}}
\sphinxAtStartPar
The specified variable name is tested in any instrumented code region.

\end{description}

\sphinxAtStartPar
If the module name or kernel name contains a \sphinxtitleref{\sphinxhyphen{}} (which can be inserted
by PSyclone, e.g. \sphinxtitleref{invoke\_compute\sphinxhyphen{}r1}), it needs to be replaced with an
underscore character in the environment variable (\sphinxtitleref{\_})

\sphinxAtStartPar
An example taken from the LFric tutorial (note that values greater than
4000 are actually valid, the upper limit was just chosen to show
a few warnings raised by the value range checker):

\begin{sphinxVerbatim}[commandchars=\\\{\}]
\PYG{n}{PSYVERIFY\PYGZus{}\PYGZus{}time\PYGZus{}evolution\PYGZus{}\PYGZus{}invoke\PYGZus{}initialise\PYGZus{}perturbation\PYGZus{}\PYGZus{}perturbation\PYGZus{}data}\PYG{o}{=}\PYG{l+m+mf}{0.0}\PYG{p}{:}\PYG{l+m+mi}{4000}
\PYG{n}{PSYVERIFY\PYGZus{}\PYGZus{}time\PYGZus{}evolution\PYGZus{}\PYGZus{}perturbation\PYGZus{}data}\PYG{o}{=}\PYG{l+m+mf}{0.0}\PYG{p}{:}\PYG{l+m+mi}{4000}
\PYG{n}{PSYVERIFY\PYGZus{}\PYGZus{}perturbation\PYGZus{}data}\PYG{o}{=}\PYG{l+m+mf}{0.0}\PYG{p}{:}\PYG{l+m+mi}{4000}
\end{sphinxVerbatim}

\begin{sphinxadmonition}{warning}{Warning:}
\sphinxAtStartPar
Note that while the field variable is called \sphinxtitleref{perturbation}, PSyclone will
append \sphinxtitleref{\_data} when the LFRic domain is used, so the name becomes
\sphinxtitleref{perturbation\_data}. You have to use
this name in LFRic in order to trigger the value range check. To verify
that the tests are done as expected, set the environment variable
\sphinxtitleref{PSYDATA\_VERBOSE} to 1, which will print which data is taken from the
environment variables:

\begin{sphinxVerbatim}[commandchars=\\\{\}]
PSyData:\PYG{+w}{ }checking\PYG{+w}{ }\PYG{l+s+s1}{\PYGZsq{}time\PYGZus{}evolution\PYGZsq{}}\PYG{+w}{ }region\PYG{+w}{ }\PYG{l+s+s1}{\PYGZsq{}invoke\PYGZus{}initialise\PYGZus{}perturbation\PYGZsq{}}\PYG{+w}{ }:\PYG{+w}{   }\PYG{l+m}{0}.0000000000000000\PYG{+w}{       }\PYGZlt{}\PYG{o}{=}\PYG{+w}{ }perturbation\PYGZus{}data\PYG{+w}{ }\PYGZlt{}\PYG{o}{=}\PYG{+w}{    }\PYG{l+m}{4000}.0000000000000
\end{sphinxVerbatim}
\end{sphinxadmonition}

\sphinxAtStartPar
If values outside the specified range are found, appropriate warnings are printed,
but the program is not aborted:

\begin{sphinxVerbatim}[commandchars=\\\{\}]
\PYG{n}{PSyData}\PYG{p}{:} \PYG{n}{Variable} \PYG{l+s+s1}{\PYGZsq{}}\PYG{l+s+s1}{perturbation\PYGZus{}data}\PYG{l+s+s1}{\PYGZsq{}} \PYG{n}{has} \PYG{n}{the} \PYG{n}{value} \PYG{l+m+mf}{4227.3587826606408} \PYG{n}{at} \PYG{n}{index}\PYG{o}{/}\PYG{n}{indices} \PYG{l+m+mi}{27051} \PYG{o+ow}{in} \PYG{n}{module} \PYG{l+s+s1}{\PYGZsq{}}\PYG{l+s+s1}{time\PYGZus{}evolution}\PYG{l+s+s1}{\PYGZsq{}}\PYG{p}{,} \PYG{n}{region} \PYG{l+s+s1}{\PYGZsq{}}\PYG{l+s+s1}{invoke\PYGZus{}initialise\PYGZus{}perturbation}\PYG{l+s+s1}{\PYGZsq{}}\PYG{p}{,} \PYG{n}{which} \PYG{o+ow}{is} \PYG{o+ow}{not} \PYG{n}{between} \PYG{l+s+s1}{\PYGZsq{}}\PYG{l+s+s1}{0.0000000000000000}\PYG{l+s+s1}{\PYGZsq{}} \PYG{o+ow}{and} \PYG{l+s+s1}{\PYGZsq{}}\PYG{l+s+s1}{4000.0000000000000}\PYG{l+s+s1}{\PYGZsq{}}\PYG{o}{.}
\end{sphinxVerbatim}

\sphinxAtStartPar
The library uses the function \sphinxcode{\sphinxupquote{IEEE\_IS\_FINITE}} from the ieee\_arithmetic module
for additionally verifying that values are not \sphinxcode{\sphinxupquote{NAN}} or \sphinxcode{\sphinxupquote{infinity}}
for any floating point variable, even if no \sphinxcode{\sphinxupquote{PSY\_VERIFY...}} environment
variable is set for this variable. Integer numbers do not have a bit pattern
for ‘infinity’ or \sphinxcode{\sphinxupquote{NaN}}, so they will only be tested for valid range
if a corresponding environment variable is specified.

\sphinxAtStartPar
The runtime libraries for GOcean and LFRic are based on a jinja\sphinxhyphen{}template
contained in the directory \sphinxcode{\sphinxupquote{\textless{}PSYCLONEHOME\textgreater{}/lib/value\_range\_check}}.
The respective API\sphinxhyphen{}specific libraries map the internal field structures
to Fortran basic types and call the functions from the base class to
handle those.

\sphinxAtStartPar
The relevant libraries for the LFRic and GOcean APIs are contained in
the \sphinxcode{\sphinxupquote{lib/value\_range\_check/lfric}} and \sphinxcode{\sphinxupquote{lib/value\_range\_check/dl\_esm\_inf}} subdirectories,
respectively. For more information on how to build and link these libraries,
please refer to the relevant \sphinxcode{\sphinxupquote{README.md}} files.


\section{Integrating PSyData Libraries into the LFRic Build Environment}
\label{\detokenize{psy_data:integrating-psydata-libraries-into-the-lfric-build-environment}}\label{\detokenize{psy_data:integrating-psy-data-lfric}}
\sphinxAtStartPar
The easiest way of integrating any PSyData\sphinxhyphen{}based library into the LFRic
build environment is:
\begin{itemize}
\item {} 
\sphinxAtStartPar
In the LFRic source tree create a new directory under \sphinxcode{\sphinxupquote{infrastructure/source}},
e.g. \sphinxcode{\sphinxupquote{infrastructure/source/psydata}}.

\item {} 
\sphinxAtStartPar
Build the PSyData wrapper stand\sphinxhyphen{}alone in \sphinxcode{\sphinxupquote{lib/extract/netcdf/lfric}} (which
will use NetCDF as output format) or \sphinxcode{\sphinxupquote{lib/extract/standalone/lfric}} (which
uses standard Fortran binary output format) by executing \sphinxcode{\sphinxupquote{make}}. The compiled
files will actually not be used, but this step will create all source
files (some of which are created by jinja). Do not copy
the compiled files into your LFRic build tree, since these files might be
compiled with an outdated version of the infrastructure files and be
incompatible with files in a current LFRic version.

\item {} 
\sphinxAtStartPar
Copy all processed source files (\sphinxcode{\sphinxupquote{extract\_netcdf\_base.f90}},
\sphinxcode{\sphinxupquote{kernel\_data\_netcdf.f90}}, \sphinxcode{\sphinxupquote{psy\_data\_base.f90}},
\sphinxcode{\sphinxupquote{read\_kernel\_data\_mod.f90}}) into \sphinxcode{\sphinxupquote{infrastructure/source/psydata}}

\item {} 
\sphinxAtStartPar
Start the LFRic build process as normal. The LFRic build environment will
copy the PSyData source files into the working directory and compile
them.

\item {} 
\sphinxAtStartPar
If the PSyData library needs additional include paths (e.g. when using an
external profiling tool), add the required paths to \sphinxcode{\sphinxupquote{\$FFLAGS}}.

\item {} 
\sphinxAtStartPar
If additional libraries are required at link time, add the paths
and libraries to \sphinxcode{\sphinxupquote{\$LDFLAGS}}. Alternatively, when a compiler wrapper
script is provided by a third\sphinxhyphen{}party tool (e.g. the profiling tool
TAU provides a script \sphinxcode{\sphinxupquote{tau\_f90.sh}}), either set the environment variable
\sphinxcode{\sphinxupquote{\$FC}}, or if this is only required at link time, the variable \sphinxcode{\sphinxupquote{\$LDMPI}}
to this compiler wrapper.

\end{itemize}

\begin{sphinxadmonition}{warning}{Warning:}
\sphinxAtStartPar
Only one PSyData library can be integrated at a time. Otherwise there
will be potentially several modules with the same name (e.g.
\sphinxcode{\sphinxupquote{psy\_data\_base}}), resulting in errors at compile time.
\end{sphinxadmonition}

\begin{sphinxadmonition}{note}{Note:}
\sphinxAtStartPar
With the new build system FAB this process might change.
\end{sphinxadmonition}

\sphinxstepscope


\chapter{Profiling}
\label{\detokenize{profiling:profiling}}\label{\detokenize{profiling:id1}}\label{\detokenize{profiling::doc}}
\sphinxAtStartPar
PSyclone has the ability to define regions that can be profiled
with various performance measurement tools. The profiling can
be enabled automatically using command line parameters like:

\begin{sphinxVerbatim}[commandchars=\\\{\}]
\PYG{n}{psyclone} \PYG{o}{\PYGZhy{}}\PYG{o}{\PYGZhy{}}\PYG{n}{profile} \PYG{n}{kernels} \PYG{o}{.}\PYG{o}{.}\PYG{o}{.}
\end{sphinxVerbatim}

\sphinxAtStartPar
Or, for finer\sphinxhyphen{}grained control, it may be applied via a profiling
transformation within a transformation script.

\sphinxAtStartPar
PSyclone can be used with a variety of existing profiling tools.
It currently supports dl\_timer, TAU, Vernier, Dr Hook, the NVIDIA GPU
profiling tools and it comes with a simple stand\sphinxhyphen{}alone timer library.
The {\hyperref[\detokenize{psy_data:psy-data}]{\sphinxcrossref{\DUrole{std,std-ref}{PSyData API}}}} (see also the
\sphinxhref{https://psyclone-dev.readthedocs.io/en/latest/psy\_data.html\#psy-data}{\DUrole{xref,std,std-ref}{Developer Guide}})
is utilised to implement wrapper libraries that connect the PSyclone
application to the profiling libraries. Certain adjustments to
the application’s build environment are required:
\begin{itemize}
\item {} 
\sphinxAtStartPar
The compiler needs to be able to find the module files for the
wrapper of the selected profiling library.

\item {} 
\sphinxAtStartPar
The application needs to be linked with the wrapper library
that interfaces between the PSyclone API and the
tool\sphinxhyphen{}specific API.

\item {} 
\sphinxAtStartPar
The tool\sphinxhyphen{}specific library also needs to be linked in.

\end{itemize}

\sphinxAtStartPar
It is the responsibility of the user to supply the corresponding
compiler command line options when building
the application that incorporates the PSyclone\sphinxhyphen{}generated code.


\section{Interface to Third Party Profiling Tools}
\label{\detokenize{profiling:interface-to-third-party-profiling-tools}}\label{\detokenize{profiling:profiling-third-party-tools}}
\sphinxAtStartPar
PSyclone comes with {\hyperref[\detokenize{libraries:libraries}]{\sphinxcrossref{\DUrole{std,std-ref}{wrapper libraries}}}} to support
usage of TAU, Vernier, Dr Hook, dl\_timer, NVTX (NVIDIA Tools Extension
library), and a simple non\sphinxhyphen{}thread\sphinxhyphen{}safe timing library. Support for further
profiling libraries will be added in the future. To compile the
wrapper libraries, change into the directory \sphinxcode{\sphinxupquote{lib/profiling}}
of PSyclone and type \sphinxcode{\sphinxupquote{make}} to compile all wrappers. If only
some of the wrappers are required, you can either use
\sphinxcode{\sphinxupquote{make wrapper\sphinxhyphen{}name}} (e.g. \sphinxcode{\sphinxupquote{make drhook}}), or change
into the corresponding directory and use \sphinxcode{\sphinxupquote{make}}. The
corresponding \sphinxcode{\sphinxupquote{README.md}} files contain additional parameters
that can be set in order to find third party profiling tools.

\sphinxAtStartPar
Below are short descriptions of each of the various wrapper
libraries that come with PSyclone:
\begin{description}
\sphinxlineitem{\sphinxcode{\sphinxupquote{lib/profiling/template}}}
\sphinxAtStartPar
This is a simple library that just prints out the name
as regions are entered and exited. It could act as a
template to develop new wrapper libraries, hence its
name.

\sphinxlineitem{\sphinxcode{\sphinxupquote{lib/profiling/simple\_timing}}}
\sphinxAtStartPar
This is a simple, stand\sphinxhyphen{}alone library that uses Fortran
system calls to measure the execution time, and reports
average, minimum and maximum execution time for all regions.
It is not MPI aware (i.e. it will just report independently
for each MPI process), and not thread\sphinxhyphen{}safe.

\sphinxlineitem{\sphinxcode{\sphinxupquote{lib/profiling/dl\_timer}}}
\sphinxAtStartPar
This wrapper uses the apeg\sphinxhyphen{}dl\_timer library. In order to use
this wrapper, you must download and install the dl\_timer library
from \sphinxurl{https://bitbucket.org/apeg/dl\_timer}. This library has
various compile\sphinxhyphen{}time options and may be built with MPI or OpenMP
support. Additional link options might therefore be required
(e.g. enabling OpenMP, or linking with MPI).

\sphinxlineitem{\sphinxcode{\sphinxupquote{lib/profiling/tau}}}
\sphinxAtStartPar
This wrapper uses TAU profiling and tracing toolkit. It can be
downloaded from \sphinxurl{https://www.cs.uoregon.edu/research/tau}.

\sphinxlineitem{\sphinxcode{\sphinxupquote{lib/profiling/drhook}}}
\sphinxAtStartPar
This wrapper uses the Dr Hook library. You need to contact
ECMWF to obtain a copy of Dr Hook.

\sphinxlineitem{\sphinxcode{\sphinxupquote{lib/profiling/vernier}}}
\sphinxAtStartPar
This wrapper uses the UK Met Office’s Vernier library. It can be
downloaded from \sphinxurl{https://github.com/MetOffice/Vernier}. This
library writes its output to files \sphinxcode{\sphinxupquote{vernier\sphinxhyphen{}out\sphinxhyphen{}\textless{}RANK\textgreater{}}}, and
will overwrite existing output files.

\sphinxlineitem{\sphinxcode{\sphinxupquote{lib/profiling/nvidia}}}
\sphinxAtStartPar
This is a wrapper library that maps the PSyclone profiling API
to the NVIDIA Tools Extension library (NVTX). This library is
available from \sphinxurl{https://developer.nvidia.com/cuda-toolkit}.

\sphinxlineitem{\sphinxcode{\sphinxupquote{lib/profiling/lfric\_timer}}}
\sphinxAtStartPar
This profile wrapper uses the timer functionality provided by
LFRic, and it comes in two different versions:
\begin{itemize}
\item {} 
\sphinxAtStartPar
\sphinxcode{\sphinxupquote{libpsy\_lfric\_timer.a}}
This library just contains the PSyData wrapper, but not the
actual timer code. It must therefore be linked with the LFRic
infrastructure library. It is meant to be used by LFRic only.

\item {} 
\sphinxAtStartPar
\sphinxcode{\sphinxupquote{libpsy\_lfric\_timer\_standalone.a}}
This library contains the LFRic timer object and its dependencies.
It can be used standalone (i.e. without LFRic) with any program.
A runnable example using a GOcean code is included in
\sphinxcode{\sphinxupquote{examples/gocean/eg5/profile}}.

\end{itemize}

\sphinxAtStartPar
The LFRic timer writes its output to a file called \sphinxcode{\sphinxupquote{timer.txt}}
in the current directory, and will overwrite this file if it
should already exist.

\end{description}

\sphinxAtStartPar
Any user can create similar wrapper libraries for
other profiling tools by providing a corresponding Fortran
module. The functions that need to be implemented are described in
the developer’s guide (\sphinxhref{https://psyclone-dev.readthedocs.io/en/latest/psy\_data.html\#psy-data}{psy\_data}).

\sphinxAtStartPar
Most libraries in \sphinxcode{\sphinxupquote{lib/profiling}} need to be linked in
with the corresponding 3rd party profiling tool, or use a compiler
wrapper provided by the tool which will provide the required additional
compiler parameters. The exceptions are the template and simple\_timing
libraries, which are stand alone. The profiling example in
\sphinxcode{\sphinxupquote{examples/gocean/eg5/profile}} can be used with any of the
wrapper libraries (except \sphinxcode{\sphinxupquote{nvidia}}) to see how they work.


\section{Required Modifications to the Program}
\label{\detokenize{profiling:required-modifications-to-the-program}}\label{\detokenize{profiling:required-profiling-calls}}
\sphinxAtStartPar
In order to guarantee that any profiling library is properly
initialised, PSyclone’s profiling wrappers utilise two additional
function calls that the user must manually insert into the program:


\subsection{profile\_PSyDataInit()}
\label{\detokenize{profiling:profile-psydatainit}}
\sphinxAtStartPar
This method needs to be called once to initialise the profiling tool.
At this stage this call is not automatically inserted by PSyclone, so
it is the responsibility of the user to add the call to an appropriate
location in the application:
\sphinxSetupCaptionForVerbatim{Adding profile\_PSyDataInit.}
\def\sphinxLiteralBlockLabel{\label{\detokenize{profiling:id2}}}
\fvset{hllines={, 3,}}%
\begin{sphinxVerbatim}[commandchars=\\\{\}]
\PYG{n}{use} \PYG{n}{profile\PYGZus{}psy\PYGZus{}data\PYGZus{}mod}\PYG{p}{,} \PYG{n}{only} \PYG{p}{:} \PYG{n}{profile\PYGZus{}PSyDataInit}
\PYG{o}{.}\PYG{o}{.}\PYG{o}{.}
\PYG{n}{call} \PYG{n}{profile\PYGZus{}PSyDataInit}\PYG{p}{(}\PYG{p}{)}
\end{sphinxVerbatim}
\sphinxresetverbatimhllines

\sphinxAtStartPar
The “appropriate” location might depend on the profiling library used.
For example, it might be necessary to invoke this before or after
a call to \sphinxcode{\sphinxupquote{MPI\_Init()}}.


\subsection{profile\_PSyDataShutdown()}
\label{\detokenize{profiling:profile-psydatashutdown}}
\sphinxAtStartPar
At the end of the program the function \sphinxcode{\sphinxupquote{profile\_PSyDataShutdown()}}
must be called.
It will make sure that the measurements are printed, files are flushed,
and that the profiling tool is closed correctly. Again at
this stage it is necessary to manually insert the call at an appropriate
location:
\sphinxSetupCaptionForVerbatim{Adding profile\_PSyDataShutdown.}
\def\sphinxLiteralBlockLabel{\label{\detokenize{profiling:id3}}}
\fvset{hllines={, 3,}}%
\begin{sphinxVerbatim}[commandchars=\\\{\}]
\PYG{n}{use} \PYG{n}{profile\PYGZus{}psy\PYGZus{}data\PYGZus{}mod}\PYG{p}{,} \PYG{n}{only} \PYG{p}{:} \PYG{n}{profile\PYGZus{}PSyDataShutdown}
\PYG{o}{.}\PYG{o}{.}\PYG{o}{.}
\PYG{n}{call} \PYG{n}{profile\PYGZus{}PSyDataShutdown}\PYG{p}{(}\PYG{p}{)}
\end{sphinxVerbatim}
\sphinxresetverbatimhllines

\sphinxAtStartPar
And again the appropriate location might depend on the profiling library
used (e.g. before or after a call to \sphinxcode{\sphinxupquote{MPI\_Finalize()}}).


\section{Profiling Command\sphinxhyphen{}Line Options}
\label{\detokenize{profiling:profiling-command-line-options}}
\sphinxAtStartPar
PSyclone offers two command\sphinxhyphen{}line options to automatically instrument
code with profiling regions. It can create profile regions around
a full invoke routine (including all kernel calls in this invoke), and/or
around each individual kernel (for the PSyKAl APIs ‘lfric’ and
‘gocean’).

\sphinxAtStartPar
The option \sphinxcode{\sphinxupquote{\sphinxhyphen{}\sphinxhyphen{}profile invokes}} will automatically add calls to
start and end a profile region at the beginning and end of every
invoke subroutine created by PSyclone. All kernels called within
this invoke subroutine will be included in the profiled region.

\sphinxAtStartPar
The option \sphinxcode{\sphinxupquote{\sphinxhyphen{}\sphinxhyphen{}profile routines}} is a synonym for ‘invokes’ but is
provided as it is more intuitive for users who are transforming
existing code. (In this case, PSyclone will put a profiling region
around every routine that it processes.)

\sphinxAtStartPar
The option \sphinxcode{\sphinxupquote{\sphinxhyphen{}\sphinxhyphen{}profile kernels}} will surround each outer loop
created by PSyclone with start and end profiling calls. Note that this
option is only available if PSyclone was invoked with a \sphinxcode{\sphinxupquote{\sphinxhyphen{}api}}
parameter. If you are only transforming existing code, this option
cannot be used as there is no concept of \sphinxtitleref{kernels}.

\begin{sphinxadmonition}{note}{Note:}
\sphinxAtStartPar
In some APIs (for example {\hyperref[\detokenize{dynamo0p3:lfric-api}]{\sphinxcrossref{\DUrole{std,std-ref}{LFRic}}}}
when using distributed memory) additional minor code might
get included in a profiled kernel section, for example
\sphinxcode{\sphinxupquote{setDirty()}} calls (expensive calls like \sphinxcode{\sphinxupquote{HaloExchange}}
are excluded).
\end{sphinxadmonition}

\begin{sphinxadmonition}{note}{Note:}
\sphinxAtStartPar
If the \sphinxcode{\sphinxupquote{kernels}} option is used in combination with an
optimisation script that introduces OpenACC then profiling
calls are automatically excluded from within OpenACC
regions (since the PSyData wrappers are not compiled for
GPU execution).
\end{sphinxadmonition}

\begin{sphinxadmonition}{note}{Note:}
\sphinxAtStartPar
It is still the responsibility of the user to manually
add the calls to \sphinxcode{\sphinxupquote{profile\_PSyDataInit}} and
\sphinxcode{\sphinxupquote{profile\_PSyDataShutdown}} to the
code base (see {\hyperref[\detokenize{profiling:required-profiling-calls}]{\sphinxcrossref{\DUrole{std,std-ref}{Required Modifications to the Program}}}}).
\end{sphinxadmonition}

\sphinxAtStartPar
PSyclone will modify the schedule of each invoke to insert the
profiling regions. Below we show an example of a schedule created
when instrumenting invokes \sphinxhyphen{} all children of a Profile\sphinxhyphen{}Node will
be part of the profiling region, including all loops created by
PSyclone and all kernel calls (note that for brevity, the nodes
holding the loop bounds have been omitted for all but the first loop):
\sphinxSetupCaptionForVerbatim{Instrumenting invokes.}
\def\sphinxLiteralBlockLabel{\label{\detokenize{profiling:id4}}}
\fvset{hllines={, 2,}}%
\begin{sphinxVerbatim}[commandchars=\\\{\}]
\PYG{n}{GOInvokeSchedule}\PYG{p}{[}\PYG{n}{invoke}\PYG{o}{=}\PYG{l+s+s1}{\PYGZsq{}}\PYG{l+s+s1}{invoke\PYGZus{}1}\PYG{l+s+s1}{\PYGZsq{}}\PYG{p}{]}
    \PYG{l+m+mi}{0}\PYG{p}{:} \PYG{p}{[}\PYG{n}{Profile}\PYG{p}{]}
        \PYG{n}{Schedule}\PYG{p}{[}\PYG{p}{]}
            \PYG{l+m+mi}{0}\PYG{p}{:} \PYG{n}{Loop}\PYG{p}{[}\PYG{n+nb}{type}\PYG{o}{=}\PYG{l+s+s1}{\PYGZsq{}}\PYG{l+s+s1}{outer}\PYG{l+s+s1}{\PYGZsq{}}\PYG{p}{,}\PYG{n}{field\PYGZus{}space}\PYG{o}{=}\PYG{l+s+s1}{\PYGZsq{}}\PYG{l+s+s1}{go\PYGZus{}cu}\PYG{l+s+s1}{\PYGZsq{}}\PYG{p}{,}\PYG{n}{it\PYGZus{}space}\PYG{o}{=}\PYG{l+s+s1}{\PYGZsq{}}\PYG{l+s+s1}{go\PYGZus{}internal\PYGZus{}pts}\PYG{l+s+s1}{\PYGZsq{}}\PYG{p}{]}
                \PYG{n}{Literal}\PYG{p}{[}\PYG{n}{value}\PYG{p}{:}\PYG{l+s+s1}{\PYGZsq{}}\PYG{l+s+s1}{2}\PYG{l+s+s1}{\PYGZsq{}}\PYG{p}{]}
                \PYG{n}{Literal}\PYG{p}{[}\PYG{n}{value}\PYG{p}{:}\PYG{l+s+s1}{\PYGZsq{}}\PYG{l+s+s1}{jstop}\PYG{l+s+s1}{\PYGZsq{}}\PYG{p}{]}
                \PYG{n}{Literal}\PYG{p}{[}\PYG{n}{value}\PYG{p}{:}\PYG{l+s+s1}{\PYGZsq{}}\PYG{l+s+s1}{1}\PYG{l+s+s1}{\PYGZsq{}}\PYG{p}{]}
                \PYG{n}{Schedule}\PYG{p}{[}\PYG{p}{]}
                    \PYG{l+m+mi}{0}\PYG{p}{:} \PYG{n}{Loop}\PYG{p}{[}\PYG{n+nb}{type}\PYG{o}{=}\PYG{l+s+s1}{\PYGZsq{}}\PYG{l+s+s1}{inner}\PYG{l+s+s1}{\PYGZsq{}}\PYG{p}{,}\PYG{n}{field\PYGZus{}space}\PYG{o}{=}\PYG{l+s+s1}{\PYGZsq{}}\PYG{l+s+s1}{go\PYGZus{}cu}\PYG{l+s+s1}{\PYGZsq{}}\PYG{p}{,}
                            \PYG{n}{it\PYGZus{}space}\PYG{o}{=}\PYG{l+s+s1}{\PYGZsq{}}\PYG{l+s+s1}{go\PYGZus{}internal\PYGZus{}pts}\PYG{l+s+s1}{\PYGZsq{}}\PYG{p}{]}
                        \PYG{o}{.}\PYG{o}{.}\PYG{o}{.}
                        \PYG{n}{Schedule}\PYG{p}{[}\PYG{p}{]}
                            \PYG{l+m+mi}{0}\PYG{p}{:} \PYG{n}{CodedKern} \PYG{n}{compute\PYGZus{}unew\PYGZus{}code}\PYG{p}{(}\PYG{n}{unew\PYGZus{}fld}\PYG{p}{,}\PYG{n}{uold\PYGZus{}fld}\PYG{p}{,}\PYG{n}{z\PYGZus{}fld}\PYG{p}{,}
                                       \PYG{n}{cv\PYGZus{}fld}\PYG{p}{,}\PYG{n}{h\PYGZus{}fld}\PYG{p}{,}\PYG{n}{tdt}\PYG{p}{,}\PYG{n}{dy}\PYG{p}{)} \PYG{p}{[}\PYG{n}{module\PYGZus{}inline}\PYG{o}{=}\PYG{k+kc}{False}\PYG{p}{]}
            \PYG{l+m+mi}{1}\PYG{p}{:} \PYG{n}{Loop}\PYG{p}{[}\PYG{n+nb}{type}\PYG{o}{=}\PYG{l+s+s1}{\PYGZsq{}}\PYG{l+s+s1}{outer}\PYG{l+s+s1}{\PYGZsq{}}\PYG{p}{,}\PYG{n}{field\PYGZus{}space}\PYG{o}{=}\PYG{l+s+s1}{\PYGZsq{}}\PYG{l+s+s1}{cv}\PYG{l+s+s1}{\PYGZsq{}}\PYG{p}{,}\PYG{n}{it\PYGZus{}space}\PYG{o}{=}\PYG{l+s+s1}{\PYGZsq{}}\PYG{l+s+s1}{internal\PYGZus{}pts}\PYG{l+s+s1}{\PYGZsq{}}\PYG{p}{]}
                \PYG{o}{.}\PYG{o}{.}\PYG{o}{.}
                \PYG{n}{Schedule}\PYG{p}{[}\PYG{p}{]}
                    \PYG{l+m+mi}{0}\PYG{p}{:} \PYG{n}{Loop}\PYG{p}{[}\PYG{n+nb}{type}\PYG{o}{=}\PYG{l+s+s1}{\PYGZsq{}}\PYG{l+s+s1}{inner}\PYG{l+s+s1}{\PYGZsq{}}\PYG{p}{,}\PYG{n}{field\PYGZus{}space}\PYG{o}{=}\PYG{l+s+s1}{\PYGZsq{}}\PYG{l+s+s1}{cv}\PYG{l+s+s1}{\PYGZsq{}}\PYG{p}{,}\PYG{n}{it\PYGZus{}space}\PYG{o}{=}\PYG{l+s+s1}{\PYGZsq{}}\PYG{l+s+s1}{internal\PYGZus{}pts}\PYG{l+s+s1}{\PYGZsq{}}\PYG{p}{]}
                        \PYG{o}{.}\PYG{o}{.}\PYG{o}{.}
                        \PYG{n}{Schedule}\PYG{p}{[}\PYG{p}{]}
                            \PYG{l+m+mi}{0}\PYG{p}{:} \PYG{n}{CodedKern} \PYG{n}{compute\PYGZus{}vnew\PYGZus{}code}\PYG{p}{(}\PYG{n}{vnew\PYGZus{}fld}\PYG{p}{,}\PYG{n}{vold\PYGZus{}fld}\PYG{p}{,}\PYG{n}{z\PYGZus{}fld}\PYG{p}{,}
                                       \PYG{n}{cu\PYGZus{}fld}\PYG{p}{,}\PYG{n}{h\PYGZus{}fld}\PYG{p}{,}\PYG{n}{tdt}\PYG{p}{,}\PYG{n}{dy}\PYG{p}{)} \PYG{p}{[}\PYG{n}{module\PYGZus{}inline}\PYG{o}{=}\PYG{k+kc}{False}\PYG{p}{]}
            \PYG{l+m+mi}{2}\PYG{p}{:} \PYG{n}{Loop}\PYG{p}{[}\PYG{n+nb}{type}\PYG{o}{=}\PYG{l+s+s1}{\PYGZsq{}}\PYG{l+s+s1}{outer}\PYG{l+s+s1}{\PYGZsq{}}\PYG{p}{,}\PYG{n}{field\PYGZus{}space}\PYG{o}{=}\PYG{l+s+s1}{\PYGZsq{}}\PYG{l+s+s1}{ct}\PYG{l+s+s1}{\PYGZsq{}}\PYG{p}{,}\PYG{n}{it\PYGZus{}space}\PYG{o}{=}\PYG{l+s+s1}{\PYGZsq{}}\PYG{l+s+s1}{internal\PYGZus{}pts}\PYG{l+s+s1}{\PYGZsq{}}\PYG{p}{]}
                \PYG{o}{.}\PYG{o}{.}\PYG{o}{.}
                \PYG{n}{Schedule}\PYG{p}{[}\PYG{p}{]}
                    \PYG{l+m+mi}{0}\PYG{p}{:} \PYG{n}{Loop}\PYG{p}{[}\PYG{n+nb}{type}\PYG{o}{=}\PYG{l+s+s1}{\PYGZsq{}}\PYG{l+s+s1}{inner}\PYG{l+s+s1}{\PYGZsq{}}\PYG{p}{,}\PYG{n}{field\PYGZus{}space}\PYG{o}{=}\PYG{l+s+s1}{\PYGZsq{}}\PYG{l+s+s1}{ct}\PYG{l+s+s1}{\PYGZsq{}}\PYG{p}{,}\PYG{n}{it\PYGZus{}space}\PYG{o}{=}\PYG{l+s+s1}{\PYGZsq{}}\PYG{l+s+s1}{internal\PYGZus{}pts}\PYG{l+s+s1}{\PYGZsq{}}\PYG{p}{]}
                        \PYG{o}{.}\PYG{o}{.}\PYG{o}{.}
                        \PYG{n}{Schedule}\PYG{p}{[}\PYG{p}{]}
                            \PYG{l+m+mi}{0}\PYG{p}{:} \PYG{n}{CodedKern} \PYG{n}{compute\PYGZus{}pnew\PYGZus{}code}\PYG{p}{(}\PYG{n}{pnew\PYGZus{}fld}\PYG{p}{,}\PYG{n}{pold\PYGZus{}fld}\PYG{p}{,}\PYG{n}{cu\PYGZus{}fld}\PYG{p}{,}
                                       \PYG{n}{cv\PYGZus{}fld}\PYG{p}{,}\PYG{n}{tdt}\PYG{p}{,}\PYG{n}{dx}\PYG{p}{,}\PYG{n}{dy}\PYG{p}{)} \PYG{p}{[}\PYG{n}{module\PYGZus{}inline}\PYG{o}{=}\PYG{k+kc}{False}\PYG{p}{]}
\end{sphinxVerbatim}
\sphinxresetverbatimhllines

\sphinxAtStartPar
And now the same schedule when instrumenting kernels. In this case
each loop nest and kernel call will be contained in a separate
region:
\sphinxSetupCaptionForVerbatim{Instrumenting kernels.}
\def\sphinxLiteralBlockLabel{\label{\detokenize{profiling:id5}}}
\fvset{hllines={, 2, 13, 24,}}%
\begin{sphinxVerbatim}[commandchars=\\\{\}]
\PYG{n}{GOInvokeSchedule}\PYG{p}{[}\PYG{n}{invoke}\PYG{o}{=}\PYG{l+s+s1}{\PYGZsq{}}\PYG{l+s+s1}{invoke\PYGZus{}1}\PYG{l+s+s1}{\PYGZsq{}}\PYG{p}{]}
    \PYG{l+m+mi}{0}\PYG{p}{:} \PYG{p}{[}\PYG{n}{Profile}\PYG{p}{]}
        \PYG{n}{Schedule}\PYG{p}{[}\PYG{p}{]}
            \PYG{l+m+mi}{0}\PYG{p}{:} \PYG{n}{Loop}\PYG{p}{[}\PYG{n+nb}{type}\PYG{o}{=}\PYG{l+s+s1}{\PYGZsq{}}\PYG{l+s+s1}{outer}\PYG{l+s+s1}{\PYGZsq{}}\PYG{p}{,}\PYG{n}{field\PYGZus{}space}\PYG{o}{=}\PYG{l+s+s1}{\PYGZsq{}}\PYG{l+s+s1}{go\PYGZus{}cu}\PYG{l+s+s1}{\PYGZsq{}}\PYG{p}{,}\PYG{n}{it\PYGZus{}space}\PYG{o}{=}\PYG{l+s+s1}{\PYGZsq{}}\PYG{l+s+s1}{go\PYGZus{}internal\PYGZus{}pts}\PYG{l+s+s1}{\PYGZsq{}}\PYG{p}{]}
                \PYG{o}{.}\PYG{o}{.}\PYG{o}{.}
                \PYG{n}{Schedule}\PYG{p}{[}\PYG{p}{]}
                    \PYG{l+m+mi}{0}\PYG{p}{:} \PYG{n}{Loop}\PYG{p}{[}\PYG{n+nb}{type}\PYG{o}{=}\PYG{l+s+s1}{\PYGZsq{}}\PYG{l+s+s1}{inner}\PYG{l+s+s1}{\PYGZsq{}}\PYG{p}{,}\PYG{n}{field\PYGZus{}space}\PYG{o}{=}\PYG{l+s+s1}{\PYGZsq{}}\PYG{l+s+s1}{go\PYGZus{}cu}\PYG{l+s+s1}{\PYGZsq{}}\PYG{p}{,}
                            \PYG{n}{it\PYGZus{}space}\PYG{o}{=}\PYG{l+s+s1}{\PYGZsq{}}\PYG{l+s+s1}{go\PYGZus{}internal\PYGZus{}pts}\PYG{l+s+s1}{\PYGZsq{}}\PYG{p}{]}
                        \PYG{o}{.}\PYG{o}{.}\PYG{o}{.}
                        \PYG{n}{Schedule}\PYG{p}{[}\PYG{p}{]}
                            \PYG{l+m+mi}{0}\PYG{p}{:} \PYG{n}{CodedKern} \PYG{n}{compute\PYGZus{}unew\PYGZus{}code}\PYG{p}{(}\PYG{n}{unew\PYGZus{}fld}\PYG{p}{,}\PYG{n}{uold\PYGZus{}fld}\PYG{p}{,}\PYG{n}{z\PYGZus{}fld}\PYG{p}{,}
                                    \PYG{n}{cv\PYGZus{}fld}\PYG{p}{,}\PYG{n}{h\PYGZus{}fld}\PYG{p}{,}\PYG{n}{tdt}\PYG{p}{,}\PYG{n}{dy}\PYG{p}{)} \PYG{p}{[}\PYG{n}{module\PYGZus{}inline}\PYG{o}{=}\PYG{k+kc}{False}\PYG{p}{]}
    \PYG{l+m+mi}{1}\PYG{p}{:} \PYG{p}{[}\PYG{n}{Profile}\PYG{p}{]}
        \PYG{n}{Schedule}\PYG{p}{[}\PYG{p}{]}
            \PYG{l+m+mi}{0}\PYG{p}{:} \PYG{n}{Loop}\PYG{p}{[}\PYG{n+nb}{type}\PYG{o}{=}\PYG{l+s+s1}{\PYGZsq{}}\PYG{l+s+s1}{outer}\PYG{l+s+s1}{\PYGZsq{}}\PYG{p}{,}\PYG{n}{field\PYGZus{}space}\PYG{o}{=}\PYG{l+s+s1}{\PYGZsq{}}\PYG{l+s+s1}{go\PYGZus{}cv}\PYG{l+s+s1}{\PYGZsq{}}\PYG{p}{,}\PYG{n}{it\PYGZus{}space}\PYG{o}{=}\PYG{l+s+s1}{\PYGZsq{}}\PYG{l+s+s1}{go\PYGZus{}internal\PYGZus{}pts}\PYG{l+s+s1}{\PYGZsq{}}\PYG{p}{]}
                \PYG{o}{.}\PYG{o}{.}\PYG{o}{.}
                \PYG{n}{Schedule}\PYG{p}{[}\PYG{p}{]}
                        \PYG{l+m+mi}{0}\PYG{p}{:} \PYG{n}{Loop}\PYG{p}{[}\PYG{n+nb}{type}\PYG{o}{=}\PYG{l+s+s1}{\PYGZsq{}}\PYG{l+s+s1}{inner}\PYG{l+s+s1}{\PYGZsq{}}\PYG{p}{,}\PYG{n}{field\PYGZus{}space}\PYG{o}{=}\PYG{l+s+s1}{\PYGZsq{}}\PYG{l+s+s1}{go\PYGZus{}cv}\PYG{l+s+s1}{\PYGZsq{}}\PYG{p}{,}
                            \PYG{n}{it\PYGZus{}space}\PYG{o}{=}\PYG{l+s+s1}{\PYGZsq{}}\PYG{l+s+s1}{go\PYGZus{}internal\PYGZus{}pts}\PYG{l+s+s1}{\PYGZsq{}}\PYG{p}{]}
                            \PYG{o}{.}\PYG{o}{.}\PYG{o}{.}
                            \PYG{n}{Schedule}\PYG{p}{[}\PYG{p}{]}
                                \PYG{l+m+mi}{0}\PYG{p}{:} \PYG{n}{CodedKern} \PYG{n}{compute\PYGZus{}vnew\PYGZus{}code}\PYG{p}{(}\PYG{n}{vnew\PYGZus{}fld}\PYG{p}{,}\PYG{n}{vold\PYGZus{}fld}\PYG{p}{,}\PYG{n}{z\PYGZus{}fld}\PYG{p}{,}
                                    \PYG{n}{cu\PYGZus{}fld}\PYG{p}{,}\PYG{n}{h\PYGZus{}fld}\PYG{p}{,}\PYG{n}{tdt}\PYG{p}{,}\PYG{n}{dy}\PYG{p}{)} \PYG{p}{[}\PYG{n}{module\PYGZus{}inline}\PYG{o}{=}\PYG{k+kc}{False}\PYG{p}{]}
    \PYG{l+m+mi}{2}\PYG{p}{:} \PYG{p}{[}\PYG{n}{Profile}\PYG{p}{]}
        \PYG{n}{Schedule}\PYG{p}{[}\PYG{p}{]}
            \PYG{l+m+mi}{0}\PYG{p}{:} \PYG{n}{Loop}\PYG{p}{[}\PYG{n+nb}{type}\PYG{o}{=}\PYG{l+s+s1}{\PYGZsq{}}\PYG{l+s+s1}{outer}\PYG{l+s+s1}{\PYGZsq{}}\PYG{p}{,}\PYG{n}{field\PYGZus{}space}\PYG{o}{=}\PYG{l+s+s1}{\PYGZsq{}}\PYG{l+s+s1}{go\PYGZus{}ct}\PYG{l+s+s1}{\PYGZsq{}}\PYG{p}{,}\PYG{n}{it\PYGZus{}space}\PYG{o}{=}\PYG{l+s+s1}{\PYGZsq{}}\PYG{l+s+s1}{go\PYGZus{}internal\PYGZus{}pts}\PYG{l+s+s1}{\PYGZsq{}}\PYG{p}{]}
                \PYG{o}{.}\PYG{o}{.}\PYG{o}{.}
                \PYG{n}{Schedule}\PYG{p}{[}\PYG{p}{]}
                    \PYG{l+m+mi}{0}\PYG{p}{:} \PYG{n}{Loop}\PYG{p}{[}\PYG{n+nb}{type}\PYG{o}{=}\PYG{l+s+s1}{\PYGZsq{}}\PYG{l+s+s1}{inner}\PYG{l+s+s1}{\PYGZsq{}}\PYG{p}{,}\PYG{n}{field\PYGZus{}space}\PYG{o}{=}\PYG{l+s+s1}{\PYGZsq{}}\PYG{l+s+s1}{go\PYGZus{}ct}\PYG{l+s+s1}{\PYGZsq{}}\PYG{p}{,}
                            \PYG{n}{it\PYGZus{}space}\PYG{o}{=}\PYG{l+s+s1}{\PYGZsq{}}\PYG{l+s+s1}{go\PYGZus{}internal\PYGZus{}pts}\PYG{l+s+s1}{\PYGZsq{}}\PYG{p}{]}
                        \PYG{o}{.}\PYG{o}{.}\PYG{o}{.}
                        \PYG{n}{Schedule}\PYG{p}{[}\PYG{p}{]}
                            \PYG{l+m+mi}{0}\PYG{p}{:} \PYG{n}{CodedKern} \PYG{n}{compute\PYGZus{}pnew\PYGZus{}code}\PYG{p}{(}\PYG{n}{pnew\PYGZus{}fld}\PYG{p}{,}\PYG{n}{pold\PYGZus{}fld}\PYG{p}{,}
                                    \PYG{n}{cu\PYGZus{}fld}\PYG{p}{,}\PYG{n}{cv\PYGZus{}fld}\PYG{p}{,}\PYG{n}{tdt}\PYG{p}{,}\PYG{n}{dx}\PYG{p}{,}\PYG{n}{dy}\PYG{p}{)} \PYG{p}{[}\PYG{n}{module\PYGZus{}inline}\PYG{o}{=}\PYG{k+kc}{False}\PYG{p}{]}
\end{sphinxVerbatim}
\sphinxresetverbatimhllines

\sphinxAtStartPar
Both options can be specified at the same time:
\sphinxSetupCaptionForVerbatim{Instrumenting kernels and invokes.}
\def\sphinxLiteralBlockLabel{\label{\detokenize{profiling:id6}}}
\fvset{hllines={, 2, 4, 16, 28,}}%
\begin{sphinxVerbatim}[commandchars=\\\{\}]
\PYG{n}{GOInvokeSchedule}\PYG{p}{[}\PYG{n}{invoke}\PYG{o}{=}\PYG{l+s+s1}{\PYGZsq{}}\PYG{l+s+s1}{invoke\PYGZus{}1}\PYG{l+s+s1}{\PYGZsq{}}\PYG{p}{]}
    \PYG{l+m+mi}{0}\PYG{p}{:} \PYG{p}{[}\PYG{n}{Profile}\PYG{p}{]}
        \PYG{n}{Schedule}\PYG{p}{[}\PYG{p}{]}
            \PYG{l+m+mi}{0}\PYG{p}{:} \PYG{p}{[}\PYG{n}{Profile}\PYG{p}{]}
                \PYG{n}{Schedule}\PYG{p}{[}\PYG{p}{]}
                    \PYG{l+m+mi}{0}\PYG{p}{:} \PYG{n}{Loop}\PYG{p}{[}\PYG{n+nb}{type}\PYG{o}{=}\PYG{l+s+s1}{\PYGZsq{}}\PYG{l+s+s1}{outer}\PYG{l+s+s1}{\PYGZsq{}}\PYG{p}{,}\PYG{n}{field\PYGZus{}space}\PYG{o}{=}\PYG{l+s+s1}{\PYGZsq{}}\PYG{l+s+s1}{go\PYGZus{}cu}\PYG{l+s+s1}{\PYGZsq{}}\PYG{p}{,}
                            \PYG{n}{it\PYGZus{}space}\PYG{o}{=}\PYG{l+s+s1}{\PYGZsq{}}\PYG{l+s+s1}{go\PYGZus{}internal\PYGZus{}pts}\PYG{l+s+s1}{\PYGZsq{}}\PYG{p}{]}
                        \PYG{o}{.}\PYG{o}{.}\PYG{o}{.}
                        \PYG{n}{Schedule}\PYG{p}{[}\PYG{p}{]}
                            \PYG{l+m+mi}{0}\PYG{p}{:} \PYG{n}{Loop}\PYG{p}{[}\PYG{n+nb}{type}\PYG{o}{=}\PYG{l+s+s1}{\PYGZsq{}}\PYG{l+s+s1}{inner}\PYG{l+s+s1}{\PYGZsq{}}\PYG{p}{,}\PYG{n}{field\PYGZus{}space}\PYG{o}{=}\PYG{l+s+s1}{\PYGZsq{}}\PYG{l+s+s1}{go\PYGZus{}cu}\PYG{l+s+s1}{\PYGZsq{}}\PYG{p}{,}
                                    \PYG{n}{it\PYGZus{}space}\PYG{o}{=}\PYG{l+s+s1}{\PYGZsq{}}\PYG{l+s+s1}{go\PYGZus{}internal\PYGZus{}pts}\PYG{l+s+s1}{\PYGZsq{}}\PYG{p}{]}
                                \PYG{o}{.}\PYG{o}{.}\PYG{o}{.}
                                \PYG{n}{Schedule}\PYG{p}{[}\PYG{p}{]}
                                    \PYG{l+m+mi}{0}\PYG{p}{:} \PYG{n}{CodedKern} \PYG{n}{compute\PYGZus{}unew\PYGZus{}code}\PYG{p}{(}\PYG{n}{unew\PYGZus{}fld}\PYG{p}{,}\PYG{n}{uold\PYGZus{}fld}\PYG{p}{,}
                                            \PYG{o}{.}\PYG{o}{.}\PYG{o}{.}\PYG{p}{)} \PYG{p}{[}\PYG{n}{module\PYGZus{}inline}\PYG{o}{=}\PYG{k+kc}{False}\PYG{p}{]}
            \PYG{l+m+mi}{1}\PYG{p}{:} \PYG{p}{[}\PYG{n}{Profile}\PYG{p}{]}
                \PYG{n}{Schedule}\PYG{p}{[}\PYG{p}{]}
                    \PYG{l+m+mi}{0}\PYG{p}{:} \PYG{n}{Loop}\PYG{p}{[}\PYG{n+nb}{type}\PYG{o}{=}\PYG{l+s+s1}{\PYGZsq{}}\PYG{l+s+s1}{outer}\PYG{l+s+s1}{\PYGZsq{}}\PYG{p}{,}\PYG{n}{field\PYGZus{}space}\PYG{o}{=}\PYG{l+s+s1}{\PYGZsq{}}\PYG{l+s+s1}{go\PYGZus{}cv}\PYG{l+s+s1}{\PYGZsq{}}\PYG{p}{,}
                            \PYG{n}{it\PYGZus{}space}\PYG{o}{=}\PYG{l+s+s1}{\PYGZsq{}}\PYG{l+s+s1}{go\PYGZus{}internal\PYGZus{}pts}\PYG{l+s+s1}{\PYGZsq{}}\PYG{p}{]}
                        \PYG{o}{.}\PYG{o}{.}\PYG{o}{.}
                        \PYG{n}{Schedule}\PYG{p}{[}\PYG{p}{]}
                                \PYG{l+m+mi}{0}\PYG{p}{:} \PYG{n}{Loop}\PYG{p}{[}\PYG{n+nb}{type}\PYG{o}{=}\PYG{l+s+s1}{\PYGZsq{}}\PYG{l+s+s1}{inner}\PYG{l+s+s1}{\PYGZsq{}}\PYG{p}{,}\PYG{n}{field\PYGZus{}space}\PYG{o}{=}\PYG{l+s+s1}{\PYGZsq{}}\PYG{l+s+s1}{go\PYGZus{}cv}\PYG{l+s+s1}{\PYGZsq{}}\PYG{p}{,}
                                    \PYG{n}{it\PYGZus{}space}\PYG{o}{=}\PYG{l+s+s1}{\PYGZsq{}}\PYG{l+s+s1}{go\PYGZus{}internal\PYGZus{}pts}\PYG{l+s+s1}{\PYGZsq{}}\PYG{p}{]}
                                    \PYG{o}{.}\PYG{o}{.}\PYG{o}{.}
                                    \PYG{n}{Schedule}\PYG{p}{[}\PYG{p}{]}
                                        \PYG{l+m+mi}{0}\PYG{p}{:} \PYG{n}{CodedKern} \PYG{n}{compute\PYGZus{}vnew\PYGZus{}code}\PYG{p}{(}\PYG{n}{vnew\PYGZus{}fld}\PYG{p}{,}\PYG{n}{vold\PYGZus{}fld}\PYG{p}{,}
                                            \PYG{o}{.}\PYG{o}{.}\PYG{o}{.}\PYG{p}{)} \PYG{p}{[}\PYG{n}{module\PYGZus{}inline}\PYG{o}{=}\PYG{k+kc}{False}\PYG{p}{]}
            \PYG{l+m+mi}{2}\PYG{p}{:} \PYG{p}{[}\PYG{n}{Profile}\PYG{p}{]}
                \PYG{n}{Schedule}\PYG{p}{[}\PYG{p}{]}
                    \PYG{l+m+mi}{0}\PYG{p}{:} \PYG{n}{Loop}\PYG{p}{[}\PYG{n+nb}{type}\PYG{o}{=}\PYG{l+s+s1}{\PYGZsq{}}\PYG{l+s+s1}{outer}\PYG{l+s+s1}{\PYGZsq{}}\PYG{p}{,}\PYG{n}{field\PYGZus{}space}\PYG{o}{=}\PYG{l+s+s1}{\PYGZsq{}}\PYG{l+s+s1}{go\PYGZus{}ct}\PYG{l+s+s1}{\PYGZsq{}}\PYG{p}{,}
                            \PYG{n}{it\PYGZus{}space}\PYG{o}{=}\PYG{l+s+s1}{\PYGZsq{}}\PYG{l+s+s1}{go\PYGZus{}internal\PYGZus{}pts}\PYG{l+s+s1}{\PYGZsq{}}\PYG{p}{]}
                        \PYG{o}{.}\PYG{o}{.}\PYG{o}{.}
                        \PYG{n}{Schedule}\PYG{p}{[}\PYG{p}{]}
                            \PYG{l+m+mi}{0}\PYG{p}{:} \PYG{n}{Loop}\PYG{p}{[}\PYG{n+nb}{type}\PYG{o}{=}\PYG{l+s+s1}{\PYGZsq{}}\PYG{l+s+s1}{inner}\PYG{l+s+s1}{\PYGZsq{}}\PYG{p}{,}\PYG{n}{field\PYGZus{}space}\PYG{o}{=}\PYG{l+s+s1}{\PYGZsq{}}\PYG{l+s+s1}{go\PYGZus{}ct}\PYG{l+s+s1}{\PYGZsq{}}\PYG{p}{,}
                                    \PYG{n}{it\PYGZus{}space}\PYG{o}{=}\PYG{l+s+s1}{\PYGZsq{}}\PYG{l+s+s1}{go\PYGZus{}internal\PYGZus{}pts}\PYG{l+s+s1}{\PYGZsq{}}\PYG{p}{]}
                                \PYG{o}{.}\PYG{o}{.}\PYG{o}{.}
                                \PYG{n}{Schedule}\PYG{p}{[}\PYG{p}{]}
                                    \PYG{l+m+mi}{0}\PYG{p}{:} \PYG{n}{CodedKern} \PYG{n}{compute\PYGZus{}pnew\PYGZus{}code}\PYG{p}{(}\PYG{n}{pnew\PYGZus{}fld}\PYG{p}{,}\PYG{n}{pold\PYGZus{}fld}\PYG{p}{,}
                                            \PYG{o}{.}\PYG{o}{.}\PYG{o}{.}\PYG{p}{)} \PYG{p}{[}\PYG{n}{module\PYGZus{}inline}\PYG{o}{=}\PYG{k+kc}{False}\PYG{p}{]}
\end{sphinxVerbatim}
\sphinxresetverbatimhllines


\section{Profiling in Scripts \sphinxhyphen{} \sphinxstyleliteralintitle{\sphinxupquote{ProfileTrans}}}
\label{\detokenize{profiling:profiling-in-scripts-profiletrans}}
\sphinxAtStartPar
The greatest flexibility is achieved by using the profiler
transformation explicitly in a transformation script. The script
takes either a single PSyIR Node or a list of PSyIR Nodes as argument,
and will insert a Profile Node into the PSyIR, with the
specified nodes as children. At code creation time the
listed children will all be enclosed in one profile region.
As an example:
\sphinxSetupCaptionForVerbatim{Explicitly adding profiling regions.}
\def\sphinxLiteralBlockLabel{\label{\detokenize{profiling:id7}}}
\fvset{hllines={, 3, 8,}}%
\begin{sphinxVerbatim}[commandchars=\\\{\}]
\PYG{k+kn}{from}\PYG{+w}{ }\PYG{n+nn}{psyclone}\PYG{n+nn}{.}\PYG{n+nn}{psyir}\PYG{n+nn}{.}\PYG{n+nn}{transformations}\PYG{+w}{ }\PYG{k+kn}{import} \PYG{n}{ProfileTrans}

\PYG{n}{p\PYGZus{}trans} \PYG{o}{=} \PYG{n}{ProfileTrans}\PYG{p}{(}\PYG{p}{)}
\PYG{n}{schedule} \PYG{o}{=} \PYG{n}{psy}\PYG{o}{.}\PYG{n}{invokes}\PYG{o}{.}\PYG{n}{get}\PYG{p}{(}\PYG{l+s+s1}{\PYGZsq{}}\PYG{l+s+s1}{invoke\PYGZus{}0}\PYG{l+s+s1}{\PYGZsq{}}\PYG{p}{)}\PYG{o}{.}\PYG{n}{schedule}
\PYG{n+nb}{print}\PYG{p}{(}\PYG{n}{schedule}\PYG{o}{.}\PYG{n}{view}\PYG{p}{(}\PYG{p}{)}\PYG{p}{)}

\PYG{c+c1}{\PYGZsh{} Enclose some children within a single profile region}
\PYG{n}{p\PYGZus{}trans}\PYG{o}{.}\PYG{n}{apply}\PYG{p}{(}\PYG{n}{schedule}\PYG{o}{.}\PYG{n}{children}\PYG{p}{[}\PYG{l+m+mi}{1}\PYG{p}{:}\PYG{l+m+mi}{3}\PYG{p}{]}\PYG{p}{)}
\PYG{n+nb}{print}\PYG{p}{(}\PYG{n}{schedule}\PYG{o}{.}\PYG{n}{view}\PYG{p}{(}\PYG{p}{)}\PYG{p}{)}
\end{sphinxVerbatim}
\sphinxresetverbatimhllines

\sphinxAtStartPar
The profiler transformation also allows the profile name to be set
explicitly, rather than being automatically created (see
{\hyperref[\detokenize{profiling:profile-names}]{\sphinxcrossref{\DUrole{std,std-ref}{Naming Profiling Regions}}}} for details). This allows for potentially
more intuitive names or finer grain control over profiling
(as particular regions could be provided with the same profile
names). For example:
\sphinxSetupCaptionForVerbatim{Setting profile region names.}
\def\sphinxLiteralBlockLabel{\label{\detokenize{profiling:id8}}}
\fvset{hllines={, 5, 8,}}%
\begin{sphinxVerbatim}[commandchars=\\\{\}]
\PYG{n}{invoke} \PYG{o}{=} \PYG{n}{psy}\PYG{o}{.}\PYG{n}{invokes}\PYG{o}{.}\PYG{n}{invoke\PYGZus{}list}\PYG{p}{[}\PYG{l+m+mi}{0}\PYG{p}{]}
\PYG{n}{schedule} \PYG{o}{=} \PYG{n}{invoke}\PYG{o}{.}\PYG{n}{schedule}
\PYG{n}{profile\PYGZus{}trans} \PYG{o}{=} \PYG{n}{ProfileTrans}\PYG{p}{(}\PYG{p}{)}
\PYG{c+c1}{\PYGZsh{} Use the actual PSy\PYGZhy{}layer module and subroutine names.}
\PYG{n}{options} \PYG{o}{=} \PYG{p}{\PYGZob{}}\PYG{l+s+s2}{\PYGZdq{}}\PYG{l+s+s2}{region\PYGZus{}name}\PYG{l+s+s2}{\PYGZdq{}}\PYG{p}{:} \PYG{p}{(}\PYG{n}{psy}\PYG{o}{.}\PYG{n}{name}\PYG{p}{,} \PYG{n}{invoke}\PYG{o}{.}\PYG{n}{name}\PYG{p}{)}\PYG{p}{\PYGZcb{}}
\PYG{n}{profile\PYGZus{}trans}\PYG{o}{.}\PYG{n}{apply}\PYG{p}{(}\PYG{n}{schedule}\PYG{o}{.}\PYG{n}{children}\PYG{p}{,} \PYG{n}{options}\PYG{o}{=}\PYG{n}{options}\PYG{p}{)}
\PYG{c+c1}{\PYGZsh{} Use own names and repeat for different regions to aggregate profile.}
\PYG{n}{options} \PYG{o}{=} \PYG{p}{\PYGZob{}}\PYG{l+s+s2}{\PYGZdq{}}\PYG{l+s+s2}{region\PYGZus{}name}\PYG{l+s+s2}{\PYGZdq{}}\PYG{p}{:} \PYG{p}{(}\PYG{l+s+s2}{\PYGZdq{}}\PYG{l+s+s2}{my\PYGZus{}location}\PYG{l+s+s2}{\PYGZdq{}}\PYG{p}{,} \PYG{l+s+s2}{\PYGZdq{}}\PYG{l+s+s2}{my\PYGZus{}region}\PYG{l+s+s2}{\PYGZdq{}}\PYG{p}{)}\PYG{p}{\PYGZcb{}}
\PYG{n}{profile\PYGZus{}trans}\PYG{o}{.}\PYG{n}{apply}\PYG{p}{(}\PYG{n}{schedule}\PYG{p}{[}\PYG{l+m+mi}{0}\PYG{p}{]}\PYG{o}{.}\PYG{n}{children}\PYG{p}{[}\PYG{l+m+mi}{1}\PYG{p}{:}\PYG{l+m+mi}{2}\PYG{p}{]}\PYG{p}{,} \PYG{n}{options}\PYG{o}{=}\PYG{n}{options}\PYG{p}{)}
\PYG{n}{profile\PYGZus{}trans}\PYG{o}{.}\PYG{n}{apply}\PYG{p}{(}\PYG{n}{schedule}\PYG{p}{[}\PYG{l+m+mi}{0}\PYG{p}{]}\PYG{o}{.}\PYG{n}{children}\PYG{p}{[}\PYG{l+m+mi}{5}\PYG{p}{:}\PYG{l+m+mi}{7}\PYG{p}{]}\PYG{p}{,} \PYG{n}{options}\PYG{o}{=}\PYG{n}{options}\PYG{p}{)}
\end{sphinxVerbatim}
\sphinxresetverbatimhllines

\begin{sphinxadmonition}{warning}{Warning:}
\sphinxAtStartPar
If “region\_name” is misspelt in the options dictionary then the
option will be silently ignored. This is true for all
options. Issue \#613 captures this problem.
\end{sphinxadmonition}

\begin{sphinxadmonition}{warning}{Warning:}
\sphinxAtStartPar
It is the responsibility of the user to make sure that a profile
region is only created inside a multi\sphinxhyphen{}threaded region if the
profiling library used is thread\sphinxhyphen{}safe!
\end{sphinxadmonition}


\section{Naming Profiling Regions}
\label{\detokenize{profiling:naming-profiling-regions}}\label{\detokenize{profiling:profile-names}}
\sphinxAtStartPar
A profile region derives its name from two components:
\begin{description}
\sphinxlineitem{\sphinxcode{\sphinxupquote{module\_name}}}
\sphinxAtStartPar
A string identifying the PSy\sphinxhyphen{}layer (PSyKAl DSL) or module (existing code)
containing this profile node.

\sphinxlineitem{\sphinxcode{\sphinxupquote{region\_name}}}
\sphinxAtStartPar
A string identifying the invoke (PSyKAl DSL) or routine (existing code)
containing this profile node and its location within the invoke/routine
(where necessary).

\end{description}

\sphinxAtStartPar
By default PSyclone will generate appropriate names to uniquely
determine a particular region. Since those names can be
somewhat cryptic, alternative names can be specified by the user
when adding profiling via a transformation script, see
\sphinxhref{https://psyclone-dev.readthedocs.io/en/latest/psy\_data.html\#psy-data-parameters-to-constructor}{Passing Parameters From the User to the Node Constructor}.

\sphinxAtStartPar
The automatic name generation depends on whether you are using a
PSyKAl DSL or only the transformation capabilities of PSyclone. If
you are transforming existing code:
\begin{itemize}
\item {} 
\sphinxAtStartPar
the \sphinxcode{\sphinxupquote{module\_name}} string is set to the module which contains the
current code. If there is no module (e.g. a stand\sphinxhyphen{}alone subroutine),
the subroutine name is used instead. This name is unique as Fortran requires
these names to be unique within a program.

\item {} 
\sphinxAtStartPar
the \sphinxcode{\sphinxupquote{region\_name}} is set to the name of the subroutine, followed by
an \sphinxcode{\sphinxupquote{r}} (standing for region) followed by an integer which uniquely identifies
the profile within the parent function/subroutine/program (based on the
profile node’s position in the PSyIR representation relative to any other
profile nodes). If there is no module name (which means the \sphinxcode{\sphinxupquote{module\_name}}
is already set to the subroutine name), only the \sphinxcode{\sphinxupquote{r}} followed by an
integer number is specified.
\begin{description}
\sphinxlineitem{Example:}\sphinxSetupCaptionForVerbatim{Profiling names used when transforming existing code.}
\def\sphinxLiteralBlockLabel{\label{\detokenize{profiling:id9}}}
\fvset{hllines={, 2, 4,}}%
\begin{sphinxVerbatim}[commandchars=\\\{\}]
! If the subroutine tra\PYGZus{}adv is contained in module tra\PYGZus{}adv\PYGZus{}mod:
CALL profile\PYGZus{}psy\PYGZus{}data \PYGZpc{} PreStart(\PYGZdq{}tra\PYGZus{}adv\PYGZus{}mod\PYGZdq{}, \PYGZdq{}tra\PYGZus{}adv\PYGZhy{}r0\PYGZdq{}, 0, 0)
! If the subroutinetra\PYGZus{}adv is not contained in a module:
CALL profile\PYGZus{}psy\PYGZus{}data \PYGZpc{} PreStart(\PYGZdq{}tra\PYGZus{}adv\PYGZdq{}, \PYGZdq{}r0\PYGZdq{}, 0, 0)
\end{sphinxVerbatim}
\sphinxresetverbatimhllines

\end{description}

\end{itemize}

\sphinxAtStartPar
For the {\hyperref[\detokenize{dynamo0p3:lfric-api}]{\sphinxcrossref{\DUrole{std,std-ref}{LFRic}}}} and
{\hyperref[\detokenize{gocean1p0:gocean-api}]{\sphinxcrossref{\DUrole{std,std-ref}{GOcean}}}} APIs:
\begin{itemize}
\item {} 
\sphinxAtStartPar
the \sphinxcode{\sphinxupquote{module\_name}} string is set to the module name of the generated
PSy\sphinxhyphen{}layer. This name should be unique by design (otherwise module
names would clash when compiling).

\item {} 
\sphinxAtStartPar
the \sphinxcode{\sphinxupquote{region\_name}} is set to the name of the invoke in which it
resides, followed by a \sphinxcode{\sphinxupquote{\sphinxhyphen{}}} and a kernel name if the
profile region contains a single kernel, and is completed by \sphinxcode{\sphinxupquote{\sphinxhyphen{}r}}
(standing for region) followed by an integer which uniquely
identifies the profile within the invoke (based on the profile
node’s position in the PSyIR representation relative to any other
profile nodes). For example:
\begin{quote}
\sphinxSetupCaptionForVerbatim{PSyIR with profiling nodes.}
\def\sphinxLiteralBlockLabel{\label{\detokenize{profiling:id10}}}
\fvset{hllines={, 2, 4, 12,}}%
\begin{sphinxVerbatim}[commandchars=\\\{\}]
\PYG{n}{InvokeSchedule}\PYG{p}{[}\PYG{n}{invoke}\PYG{o}{=}\PYG{l+s+s1}{\PYGZsq{}}\PYG{l+s+s1}{invoke\PYGZus{}0}\PYG{l+s+s1}{\PYGZsq{}}\PYG{p}{,} \PYG{n}{dm}\PYG{o}{=}\PYG{k+kc}{True}\PYG{p}{]}
  \PYG{l+m+mi}{0}\PYG{p}{:} \PYG{n}{Profile}\PYG{p}{[}\PYG{p}{]}
      \PYG{n}{Schedule}\PYG{p}{[}\PYG{p}{]}
          \PYG{l+m+mi}{0}\PYG{p}{:} \PYG{n}{Profile}\PYG{p}{[}\PYG{p}{]}
              \PYG{n}{Schedule}\PYG{p}{[}\PYG{p}{]}
                  \PYG{l+m+mi}{0}\PYG{p}{:} \PYG{n}{HaloExchange}\PYG{p}{[}\PYG{n}{field}\PYG{o}{=}\PYG{l+s+s1}{\PYGZsq{}}\PYG{l+s+s1}{f2}\PYG{l+s+s1}{\PYGZsq{}}\PYG{p}{,} \PYG{n+nb}{type}\PYG{o}{=}\PYG{l+s+s1}{\PYGZsq{}}\PYG{l+s+s1}{region}\PYG{l+s+s1}{\PYGZsq{}}\PYG{p}{,} \PYG{n}{depth}\PYG{o}{=}\PYG{l+m+mi}{1}\PYG{p}{,}
                                  \PYG{n}{check\PYGZus{}dirty}\PYG{o}{=}\PYG{k+kc}{True}\PYG{p}{]}
                  \PYG{l+m+mi}{1}\PYG{p}{:} \PYG{n}{HaloExchange}\PYG{p}{[}\PYG{n}{field}\PYG{o}{=}\PYG{l+s+s1}{\PYGZsq{}}\PYG{l+s+s1}{m1}\PYG{l+s+s1}{\PYGZsq{}}\PYG{p}{,} \PYG{n+nb}{type}\PYG{o}{=}\PYG{l+s+s1}{\PYGZsq{}}\PYG{l+s+s1}{region}\PYG{l+s+s1}{\PYGZsq{}}\PYG{p}{,} \PYG{n}{depth}\PYG{o}{=}\PYG{l+m+mi}{1}\PYG{p}{,}
                                  \PYG{n}{check\PYGZus{}dirty}\PYG{o}{=}\PYG{k+kc}{True}\PYG{p}{]}
                  \PYG{l+m+mi}{2}\PYG{p}{:} \PYG{n}{HaloExchange}\PYG{p}{[}\PYG{n}{field}\PYG{o}{=}\PYG{l+s+s1}{\PYGZsq{}}\PYG{l+s+s1}{m2}\PYG{l+s+s1}{\PYGZsq{}}\PYG{p}{,} \PYG{n+nb}{type}\PYG{o}{=}\PYG{l+s+s1}{\PYGZsq{}}\PYG{l+s+s1}{region}\PYG{l+s+s1}{\PYGZsq{}}\PYG{p}{,} \PYG{n}{depth}\PYG{o}{=}\PYG{l+m+mi}{1}\PYG{p}{,}
                                  \PYG{n}{check\PYGZus{}dirty}\PYG{o}{=}\PYG{k+kc}{True}\PYG{p}{]}
          \PYG{l+m+mi}{1}\PYG{p}{:} \PYG{n}{Profile}\PYG{p}{[}\PYG{p}{]}
              \PYG{n}{Schedule}\PYG{p}{[}\PYG{p}{]}
                  \PYG{l+m+mi}{0}\PYG{p}{:} \PYG{n}{Loop}\PYG{p}{[}\PYG{n+nb}{type}\PYG{o}{=}\PYG{l+s+s1}{\PYGZsq{}}\PYG{l+s+s1}{\PYGZsq{}}\PYG{p}{,} \PYG{n}{field\PYGZus{}space}\PYG{o}{=}\PYG{l+s+s1}{\PYGZsq{}}\PYG{l+s+s1}{w1}\PYG{l+s+s1}{\PYGZsq{}}\PYG{p}{,} \PYG{n}{it\PYGZus{}space}\PYG{o}{=}\PYG{l+s+s1}{\PYGZsq{}}\PYG{l+s+s1}{cells}\PYG{l+s+s1}{\PYGZsq{}}\PYG{p}{,}
                          \PYG{n}{upper\PYGZus{}bound}\PYG{o}{=}\PYG{l+s+s1}{\PYGZsq{}}\PYG{l+s+s1}{cell\PYGZus{}halo(1)}\PYG{l+s+s1}{\PYGZsq{}}\PYG{p}{]}
                      \PYG{n}{Literal}\PYG{p}{[}\PYG{n}{value}\PYG{p}{:}\PYG{l+s+s1}{\PYGZsq{}}\PYG{l+s+s1}{1}\PYG{l+s+s1}{\PYGZsq{}}\PYG{p}{,} \PYG{n}{DataType}\PYG{o}{.}\PYG{n}{INTEGER}\PYG{p}{]}
                      \PYG{n}{Literal}\PYG{p}{[}\PYG{n}{value}\PYG{p}{:}\PYG{l+s+s1}{\PYGZsq{}}\PYG{l+s+s1}{mesh}\PYG{l+s+si}{\PYGZpc{}g}\PYG{l+s+s1}{et\PYGZus{}last\PYGZus{}halo\PYGZus{}cell(1)}\PYG{l+s+s1}{\PYGZsq{}}\PYG{p}{,}
                              \PYG{n}{DataType}\PYG{o}{.}\PYG{n}{INTEGER}\PYG{p}{]}
                      \PYG{n}{Literal}\PYG{p}{[}\PYG{n}{value}\PYG{p}{:}\PYG{l+s+s1}{\PYGZsq{}}\PYG{l+s+s1}{1}\PYG{l+s+s1}{\PYGZsq{}}\PYG{p}{,} \PYG{n}{DataType}\PYG{o}{.}\PYG{n}{INTEGER}\PYG{p}{]}
                      \PYG{n}{Schedule}\PYG{p}{[}\PYG{p}{]}
                          \PYG{l+m+mi}{0}\PYG{p}{:} \PYG{n}{CodedKern} \PYG{n}{testkern\PYGZus{}code}\PYG{p}{(}\PYG{n}{a}\PYG{p}{,}\PYG{n}{f1}\PYG{p}{,}\PYG{n}{f2}\PYG{p}{,}\PYG{n}{m1}\PYG{p}{,}\PYG{n}{m2}\PYG{p}{)}
                             \PYG{p}{[}\PYG{n}{module\PYGZus{}inline}\PYG{o}{=}\PYG{k+kc}{False}\PYG{p}{]}
\end{sphinxVerbatim}
\sphinxresetverbatimhllines

\sphinxAtStartPar
This is the code created for this example:
\sphinxSetupCaptionForVerbatim{Created Fortran source code with profiling regions.}
\def\sphinxLiteralBlockLabel{\label{\detokenize{profiling:id11}}}
\fvset{hllines={, 5, 6, 7, 17, 18, 19, 24, 30,}}%
\begin{sphinxVerbatim}[commandchars=\\\{\}]
 \PYG{n}{MODULE} \PYG{n}{container}
  \PYG{n}{CONTAINS}
  \PYG{n}{SUBROUTINE} \PYG{n}{invoke\PYGZus{}0}\PYG{p}{(}\PYG{n}{a}\PYG{p}{,} \PYG{n}{f1}\PYG{p}{,} \PYG{n}{f2}\PYG{p}{,} \PYG{n}{m1}\PYG{p}{,} \PYG{n}{m2}\PYG{p}{,} \PYG{n}{istp}\PYG{p}{,} \PYG{n}{qr}\PYG{p}{)}
    \PYG{o}{.}\PYG{o}{.}\PYG{o}{.}
    \PYG{n}{CALL} \PYG{n}{psy\PYGZus{}data\PYGZus{}2}\PYG{o}{\PYGZpc{}}\PYG{n}{PreStart}\PYG{p}{(}\PYG{l+s+s2}{\PYGZdq{}}\PYG{l+s+s2}{multi\PYGZus{}functions\PYGZus{}multi\PYGZus{}invokes\PYGZus{}psy}\PYG{l+s+s2}{\PYGZdq{}}\PYG{p}{,} \PYG{l+s+s2}{\PYGZdq{}}\PYG{l+s+s2}{invoke\PYGZus{}0\PYGZhy{}r0}\PYG{l+s+s2}{\PYGZdq{}}\PYG{p}{,} \PYG{o}{\PYGZam{}}
                             \PYG{l+m+mi}{0}\PYG{p}{,} \PYG{l+m+mi}{0}\PYG{p}{)}
    \PYG{n}{CALL} \PYG{n}{psy\PYGZus{}data}\PYG{o}{\PYGZpc{}}\PYG{n}{PreStart}\PYG{p}{(}\PYG{l+s+s2}{\PYGZdq{}}\PYG{l+s+s2}{multi\PYGZus{}functions\PYGZus{}multi\PYGZus{}invokes\PYGZus{}psy}\PYG{l+s+s2}{\PYGZdq{}}\PYG{p}{,} \PYG{l+s+s2}{\PYGZdq{}}\PYG{l+s+s2}{invoke\PYGZus{}0\PYGZhy{}r1}\PYG{l+s+s2}{\PYGZdq{}}\PYG{p}{,} \PYG{l+m+mi}{0}\PYG{p}{,} \PYG{l+m+mi}{0}\PYG{p}{)}
    \PYG{n}{IF} \PYG{p}{(}\PYG{n}{f2\PYGZus{}proxy}\PYG{o}{\PYGZpc{}}\PYG{n}{is\PYGZus{}dirty}\PYG{p}{(}\PYG{n}{depth}\PYG{o}{=}\PYG{l+m+mi}{1}\PYG{p}{)}\PYG{p}{)} \PYG{n}{THEN}
      \PYG{n}{CALL} \PYG{n}{f2\PYGZus{}proxy}\PYG{o}{\PYGZpc{}}\PYG{n}{halo\PYGZus{}exchange}\PYG{p}{(}\PYG{n}{depth}\PYG{o}{=}\PYG{l+m+mi}{1}\PYG{p}{)}
    \PYG{n}{END} \PYG{n}{IF}
    \PYG{n}{IF} \PYG{p}{(}\PYG{n}{m1\PYGZus{}proxy}\PYG{o}{\PYGZpc{}}\PYG{n}{is\PYGZus{}dirty}\PYG{p}{(}\PYG{n}{depth}\PYG{o}{=}\PYG{l+m+mi}{1}\PYG{p}{)}\PYG{p}{)} \PYG{n}{THEN}
      \PYG{n}{CALL} \PYG{n}{m1\PYGZus{}proxy}\PYG{o}{\PYGZpc{}}\PYG{n}{halo\PYGZus{}exchange}\PYG{p}{(}\PYG{n}{depth}\PYG{o}{=}\PYG{l+m+mi}{1}\PYG{p}{)}
    \PYG{n}{END} \PYG{n}{IF}
    \PYG{n}{IF} \PYG{p}{(}\PYG{n}{m2\PYGZus{}proxy}\PYG{o}{\PYGZpc{}}\PYG{n}{is\PYGZus{}dirty}\PYG{p}{(}\PYG{n}{depth}\PYG{o}{=}\PYG{l+m+mi}{1}\PYG{p}{)}\PYG{p}{)} \PYG{n}{THEN}
      \PYG{n}{CALL} \PYG{n}{m2\PYGZus{}proxy}\PYG{o}{\PYGZpc{}}\PYG{n}{halo\PYGZus{}exchange}\PYG{p}{(}\PYG{n}{depth}\PYG{o}{=}\PYG{l+m+mi}{1}\PYG{p}{)}
    \PYG{n}{END} \PYG{n}{IF}
    \PYG{n}{CALL} \PYG{n}{psy\PYGZus{}data}\PYG{o}{\PYGZpc{}}\PYG{n}{PreEnd}\PYG{p}{(}\PYG{p}{)}
    \PYG{n}{CALL} \PYG{n}{psy\PYGZus{}data\PYGZus{}1}\PYG{o}{\PYGZpc{}}\PYG{n}{PreStart}\PYG{p}{(}\PYG{l+s+s2}{\PYGZdq{}}\PYG{l+s+s2}{multi\PYGZus{}functions\PYGZus{}multi\PYGZus{}invokes\PYGZus{}psy}\PYG{l+s+s2}{\PYGZdq{}}\PYG{p}{,} \PYG{l+s+s2}{\PYGZdq{}}\PYG{l+s+s2}{invoke\PYGZus{}0\PYGZhy{}r2}\PYG{l+s+s2}{\PYGZdq{}}\PYG{p}{,} \PYG{o}{\PYGZam{}}
                             \PYG{l+m+mi}{0}\PYG{p}{,} \PYG{l+m+mi}{0}\PYG{p}{)}
    \PYG{n}{DO} \PYG{n}{cell}\PYG{o}{=}\PYG{l+m+mi}{1}\PYG{p}{,}\PYG{n}{mesh}\PYG{o}{\PYGZpc{}}\PYG{n}{get\PYGZus{}last\PYGZus{}halo\PYGZus{}cell}\PYG{p}{(}\PYG{l+m+mi}{1}\PYG{p}{)}
      \PYG{n}{CALL} \PYG{n}{testkern\PYGZus{}code}\PYG{p}{(}\PYG{o}{.}\PYG{o}{.}\PYG{o}{.}\PYG{p}{)}
    \PYG{n}{END} \PYG{n}{DO}

    \PYG{n}{CALL} \PYG{n}{psy\PYGZus{}data\PYGZus{}1}\PYG{o}{\PYGZpc{}}\PYG{n}{PostEnd}\PYG{p}{(}\PYG{p}{)}
    \PYG{o}{.}\PYG{o}{.}\PYG{o}{.}
    \PYG{n}{DO} \PYG{n}{cell}\PYG{o}{=}\PYG{l+m+mi}{1}\PYG{p}{,}\PYG{n}{mesh}\PYG{o}{\PYGZpc{}}\PYG{n}{get\PYGZus{}last\PYGZus{}halo\PYGZus{}cell}\PYG{p}{(}\PYG{l+m+mi}{1}\PYG{p}{)}
      \PYG{n}{CALL} \PYG{n}{testkern\PYGZus{}qr\PYGZus{}code}\PYG{p}{(}\PYG{o}{.}\PYG{o}{.}\PYG{o}{.}\PYG{p}{)}
    \PYG{n}{END} \PYG{n}{DO}
    \PYG{o}{.}\PYG{o}{.}\PYG{o}{.}
    \PYG{n}{CALL} \PYG{n}{psy\PYGZus{}data\PYGZus{}2}\PYG{o}{\PYGZpc{}}\PYG{n}{PostEnd}\PYG{p}{(}\PYG{p}{)}
    \PYG{o}{.}\PYG{o}{.}\PYG{o}{.}
  \PYG{n}{END} \PYG{n}{SUBROUTINE} \PYG{n}{invoke\PYGZus{}0}
\PYG{n}{END} \PYG{n}{MODULE} \PYG{n}{container}
\end{sphinxVerbatim}
\sphinxresetverbatimhllines
\end{quote}

\end{itemize}

\sphinxstepscope


\chapter{PSy Kernel Extractor (PSyKE)}
\label{\detokenize{psyke:psy-kernel-extractor-psyke}}\label{\detokenize{psyke:psyke}}\label{\detokenize{psyke::doc}}

\section{Introduction}
\label{\detokenize{psyke:introduction}}\label{\detokenize{psyke:psyke-intro}}
\sphinxAtStartPar
PSyclone has the ability to define regions of a PSyclone\sphinxhyphen{}conformant code
to be extracted and run as a stand\sphinxhyphen{}alone application. This ability, called
PSyKE (PSy Kernel Extractor), can be useful for benchmarking parts of a
model, such as LFRic, without the need for using its infrastructure.


\subsection{Mechanism}
\label{\detokenize{psyke:mechanism}}\label{\detokenize{psyke:psyke-intro-mechanism}}
\sphinxAtStartPar
The code marked for extraction can be (subject to
{\hyperref[\detokenize{psyke:psyke-intro-restrictions}]{\sphinxcrossref{\DUrole{std,std-ref}{Restrictions}}}}):
\begin{itemize}
\item {} 
\sphinxAtStartPar
One or more Nodes in an Invoke (e.g. Loops containing Kernel or
Built\sphinxhyphen{}In calls, a Directive enclosing one or more Loops) or

\item {} 
\sphinxAtStartPar
The entire Invoke (extraction applied to all Nodes).

\end{itemize}

\sphinxAtStartPar
The basic mechanism of code extraction is through applying the
\sphinxcode{\sphinxupquote{ExtractTrans}} transformation to selected Nodes. This
transformation is further sub\sphinxhyphen{}classed into API\sphinxhyphen{}specific implementations,
\sphinxcode{\sphinxupquote{LFRicExtractTrans}} and \sphinxcode{\sphinxupquote{GOceanExtractTrans}}. Both
sub\sphinxhyphen{}classed transformations insert an instance of the \sphinxcode{\sphinxupquote{ExtractNode}}
object into the Schedule of a specific Invoke. All Nodes marked for
extraction become children of the \sphinxcode{\sphinxupquote{ExtractNode}}.

\sphinxAtStartPar
The \sphinxcode{\sphinxupquote{ExtractNode}} class uses the dependency analysis to detect
which variables are input\sphinxhyphen{}, and which ones are output\sphinxhyphen{}parameters.
The lists of variables are then passed to the \sphinxcode{\sphinxupquote{PSyDataNode}},
which is the base class of any \sphinxcode{\sphinxupquote{ExtractNode}} (details of
the \sphinxcode{\sphinxupquote{PSyDataNode}} can be found in \sphinxhref{https://psyclone-dev.readthedocs.io/en/latest/psy\_data.html\#psy-data}{psy\_data}). This
node then creates the actual code, as in the following LFRic example:

\begin{sphinxVerbatim}[commandchars=\\\{\}]
\PYG{c}{! ExtractStart}
\PYG{c}{!}
\PYG{k}{CALL }\PYG{n}{extract\PYGZus{}psy\PYGZus{}data}\PYG{p}{\PYGZpc{}}\PYG{n}{PreStart}\PYG{p}{(}\PYG{l+s+s2}{\PYGZdq{}testkern\PYGZus{}mod\PYGZdq{}}\PYG{p}{,}\PYG{+w}{ }\PYG{l+s+s2}{\PYGZdq{}testkern\PYGZus{}code\PYGZdq{}}\PYG{p}{,}\PYG{+w}{ }\PYG{l+m+mi}{4}\PYG{p}{,}\PYG{+w}{ }\PYG{l+m+mi}{2}\PYG{p}{)}
\PYG{k}{CALL }\PYG{n}{extract\PYGZus{}psy\PYGZus{}data}\PYG{p}{\PYGZpc{}}\PYG{n}{PreDeclareVariable}\PYG{p}{(}\PYG{l+s+s2}{\PYGZdq{}a\PYGZdq{}}\PYG{p}{,}\PYG{+w}{ }\PYG{n}{a}\PYG{p}{)}
\PYG{k}{CALL }\PYG{n}{extract\PYGZus{}psy\PYGZus{}data}\PYG{p}{\PYGZpc{}}\PYG{n}{PreDeclareVariable}\PYG{p}{(}\PYG{l+s+s2}{\PYGZdq{}f2\PYGZdq{}}\PYG{p}{,}\PYG{+w}{ }\PYG{n}{f2}\PYG{p}{)}
\PYG{k}{CALL }\PYG{n}{extract\PYGZus{}psy\PYGZus{}data}\PYG{p}{\PYGZpc{}}\PYG{n}{PreDeclareVariable}\PYG{p}{(}\PYG{l+s+s2}{\PYGZdq{}m1\PYGZdq{}}\PYG{p}{,}\PYG{+w}{ }\PYG{n}{m1}\PYG{p}{)}
\PYG{k}{CALL }\PYG{n}{extract\PYGZus{}psy\PYGZus{}data}\PYG{p}{\PYGZpc{}}\PYG{n}{PreDeclareVariable}\PYG{p}{(}\PYG{l+s+s2}{\PYGZdq{}m2\PYGZdq{}}\PYG{p}{,}\PYG{+w}{ }\PYG{n}{m2}\PYG{p}{)}
\PYG{k}{CALL }\PYG{n}{extract\PYGZus{}psy\PYGZus{}data}\PYG{p}{\PYGZpc{}}\PYG{n}{PreDeclareVariable}\PYG{p}{(}\PYG{l+s+s2}{\PYGZdq{}map\PYGZus{}w1\PYGZdq{}}\PYG{p}{,}\PYG{+w}{ }\PYG{n}{map\PYGZus{}w1}\PYG{p}{)}
\PYG{p}{.}\PYG{p}{.}\PYG{p}{.}
\PYG{k}{CALL }\PYG{n}{extract\PYGZus{}psy\PYGZus{}data}\PYG{p}{\PYGZpc{}}\PYG{n}{PreDeclareVariable}\PYG{p}{(}\PYG{l+s+s2}{\PYGZdq{}undf\PYGZus{}w3\PYGZdq{}}\PYG{p}{,}\PYG{+w}{ }\PYG{n}{undf\PYGZus{}w3}\PYG{p}{)}
\PYG{k}{CALL }\PYG{n}{extract\PYGZus{}psy\PYGZus{}data}\PYG{p}{\PYGZpc{}}\PYG{n}{PreDeclareVariable}\PYG{p}{(}\PYG{l+s+s2}{\PYGZdq{}f1\PYGZus{}post\PYGZdq{}}\PYG{p}{,}\PYG{+w}{ }\PYG{n}{f1}\PYG{p}{)}
\PYG{k}{CALL }\PYG{n}{extract\PYGZus{}psy\PYGZus{}data}\PYG{p}{\PYGZpc{}}\PYG{n}{PreDeclareVariable}\PYG{p}{(}\PYG{l+s+s2}{\PYGZdq{}cell\PYGZus{}post\PYGZdq{}}\PYG{p}{,}\PYG{+w}{ }\PYG{n}{cell}\PYG{p}{)}
\PYG{k}{CALL }\PYG{n}{extract\PYGZus{}psy\PYGZus{}data}\PYG{p}{\PYGZpc{}}\PYG{n}{PreEndDeclaration}
\PYG{k}{CALL }\PYG{n}{extract\PYGZus{}psy\PYGZus{}data}\PYG{p}{\PYGZpc{}}\PYG{n}{ProvideVariable}\PYG{p}{(}\PYG{l+s+s2}{\PYGZdq{}a\PYGZdq{}}\PYG{p}{,}\PYG{+w}{ }\PYG{n}{a}\PYG{p}{)}
\PYG{k}{CALL }\PYG{n}{extract\PYGZus{}psy\PYGZus{}data}\PYG{p}{\PYGZpc{}}\PYG{n}{ProvideVariable}\PYG{p}{(}\PYG{l+s+s2}{\PYGZdq{}f2\PYGZdq{}}\PYG{p}{,}\PYG{+w}{ }\PYG{n}{f2}\PYG{p}{)}
\PYG{k}{CALL }\PYG{n}{extract\PYGZus{}psy\PYGZus{}data}\PYG{p}{\PYGZpc{}}\PYG{n}{ProvideVariable}\PYG{p}{(}\PYG{l+s+s2}{\PYGZdq{}m1\PYGZdq{}}\PYG{p}{,}\PYG{+w}{ }\PYG{n}{m1}\PYG{p}{)}
\PYG{k}{CALL }\PYG{n}{extract\PYGZus{}psy\PYGZus{}data}\PYG{p}{\PYGZpc{}}\PYG{n}{ProvideVariable}\PYG{p}{(}\PYG{l+s+s2}{\PYGZdq{}m2\PYGZdq{}}\PYG{p}{,}\PYG{+w}{ }\PYG{n}{m2}\PYG{p}{)}
\PYG{k}{CALL }\PYG{n}{extract\PYGZus{}psy\PYGZus{}data}\PYG{p}{\PYGZpc{}}\PYG{n}{ProvideVariable}\PYG{p}{(}\PYG{l+s+s2}{\PYGZdq{}map\PYGZus{}w1\PYGZdq{}}\PYG{p}{,}\PYG{+w}{ }\PYG{n}{map\PYGZus{}w1}\PYG{p}{)}
\PYG{p}{.}\PYG{p}{.}\PYG{p}{.}
\PYG{k}{CALL }\PYG{n}{extract\PYGZus{}psy\PYGZus{}data}\PYG{p}{\PYGZpc{}}\PYG{n}{ProvideVariable}\PYG{p}{(}\PYG{l+s+s2}{\PYGZdq{}undf\PYGZus{}w3\PYGZdq{}}\PYG{p}{,}\PYG{+w}{ }\PYG{n}{undf\PYGZus{}w3}\PYG{p}{)}
\PYG{k}{CALL }\PYG{n}{extract\PYGZus{}psy\PYGZus{}data}\PYG{p}{\PYGZpc{}}\PYG{n}{PreEnd}
\PYG{k}{DO }\PYG{n}{cell}\PYG{o}{=}\PYG{l+m+mi}{1}\PYG{p}{,}\PYG{n}{f1\PYGZus{}proxy}\PYG{p}{\PYGZpc{}}\PYG{n}{vspace}\PYG{p}{\PYGZpc{}}\PYG{n}{get\PYGZus{}ncell}\PYG{p}{(}\PYG{p}{)}
\PYG{+w}{  }\PYG{c}{!}
\PYG{+w}{  }\PYG{k}{CALL }\PYG{n}{testkern\PYGZus{}code}\PYG{p}{(}\PYG{n}{nlayers}\PYG{p}{,}\PYG{+w}{ }\PYG{n}{a}\PYG{p}{,}\PYG{+w}{ }\PYG{n}{f1\PYGZus{}proxy}\PYG{p}{\PYGZpc{}}\PYG{k}{data}\PYG{p}{,}\PYG{+w}{ }\PYG{n}{f2\PYGZus{}proxy}\PYG{p}{\PYGZpc{}}\PYG{k}{data}\PYG{p}{,}\PYG{+w}{  }\PYG{p}{\PYGZam{}}
\PYG{+w}{       }\PYG{n}{m1\PYGZus{}proxy}\PYG{p}{\PYGZpc{}}\PYG{k}{data}\PYG{p}{,}\PYG{+w}{ }\PYG{n}{m2\PYGZus{}proxy}\PYG{p}{\PYGZpc{}}\PYG{k}{data}\PYG{p}{,}\PYG{+w}{ }\PYG{n}{ndf\PYGZus{}w1}\PYG{p}{,}\PYG{+w}{ }\PYG{n}{undf\PYGZus{}w1}\PYG{p}{,}\PYG{+w}{           }\PYG{p}{\PYGZam{}}
\PYG{+w}{       }\PYG{n}{map\PYGZus{}w1}\PYG{p}{(}\PYG{p}{:}\PYG{p}{,}\PYG{n}{cell}\PYG{p}{)}\PYG{p}{,}\PYG{+w}{ }\PYG{n}{ndf\PYGZus{}w2}\PYG{p}{,}\PYG{+w}{ }\PYG{n}{undf\PYGZus{}w2}\PYG{p}{,}\PYG{+w}{ }\PYG{n}{map\PYGZus{}w2}\PYG{p}{(}\PYG{p}{:}\PYG{p}{,}\PYG{n}{cell}\PYG{p}{)}\PYG{p}{,}\PYG{+w}{ }\PYG{n}{ndf\PYGZus{}w3}\PYG{p}{,}\PYG{+w}{ }\PYG{p}{\PYGZam{}}
\PYG{+w}{       }\PYG{n}{undf\PYGZus{}w3}\PYG{p}{,}\PYG{+w}{ }\PYG{n}{map\PYGZus{}w3}\PYG{p}{(}\PYG{p}{:}\PYG{p}{,}\PYG{n}{cell}\PYG{p}{)}\PYG{p}{)}
\PYG{k}{END }\PYG{k}{DO}
\PYG{k}{CALL }\PYG{n}{extract\PYGZus{}psy\PYGZus{}data}\PYG{p}{\PYGZpc{}}\PYG{n}{PostStart}
\PYG{k}{CALL }\PYG{n}{extract\PYGZus{}psy\PYGZus{}data}\PYG{p}{\PYGZpc{}}\PYG{n}{ProvideVariable}\PYG{p}{(}\PYG{l+s+s2}{\PYGZdq{}cell\PYGZus{}post\PYGZdq{}}\PYG{p}{,}\PYG{+w}{ }\PYG{n}{cell}\PYG{p}{)}
\PYG{k}{CALL }\PYG{n}{extract\PYGZus{}psy\PYGZus{}data}\PYG{p}{\PYGZpc{}}\PYG{n}{ProvideVariable}\PYG{p}{(}\PYG{l+s+s2}{\PYGZdq{}f1\PYGZus{}post\PYGZdq{}}\PYG{p}{,}\PYG{+w}{ }\PYG{n}{f1}\PYG{p}{)}
\PYG{k}{CALL }\PYG{n}{extract\PYGZus{}psy\PYGZus{}data}\PYG{p}{\PYGZpc{}}\PYG{n}{PostEnd}
\PYG{c}{!}
\PYG{c}{! ExtractEnd}
\end{sphinxVerbatim}

\sphinxAtStartPar
The {\hyperref[\detokenize{psy_data:psy-data}]{\sphinxcrossref{\DUrole{std,std-ref}{PSyData API}}}} relies on generic Fortran interfaces to
provide the  field\sphinxhyphen{}type\sphinxhyphen{}specific implementations of the \sphinxcode{\sphinxupquote{ProvideVariable}}
for different types. This means that a different version of the external
PSyData library that PSyKE uses must be supplied for each PSyclone API.


\subsection{Restrictions}
\label{\detokenize{psyke:restrictions}}\label{\detokenize{psyke:psyke-intro-restrictions}}
\sphinxAtStartPar
Code extraction can be applied to unoptimised or optimised code. There are
restrictions that check for correctness of optimising transformations when
extraction is applied, as well as restrictions that eliminate dependence on
the specific model infrastructure.


\subsubsection{General}
\label{\detokenize{psyke:general}}\label{\detokenize{psyke:psyke-intro-restrictions-gen}}
\sphinxAtStartPar
This group of restrictions is enforced irrespective of whether optimisations
are used or not.
\begin{itemize}
\item {} 
\sphinxAtStartPar
Extraction can be applied to a single Node or a list of Nodes in a
Schedule. For the latter, Nodes in the list must be consecutive children
of the same parent Schedule.

\item {} 
\sphinxAtStartPar
Extraction cannot be applied to an \sphinxcode{\sphinxupquote{ExtractNode}} or a Node list that
already contains one (otherwise we would have an extract region within
another extract region).

\item {} 
\sphinxAtStartPar
A Kernel or a Built\sphinxhyphen{}In call cannot be extracted without its parent Loop.

\end{itemize}


\subsubsection{Distributed memory}
\label{\detokenize{psyke:distributed-memory}}\label{\detokenize{psyke:psyke-intro-restrictions-dm}}
\sphinxAtStartPar
Kernel extraction for distributed memory is supported in as much as each
process will write its own output file by adding its rank to the output
file name. So each kernel and each rank will produce one file. It is possible
to extract several consecutive kernels, but there must be no halo exchange
calls between the kernels. The extraction transformation will test for this
and raise an exception if this should happen.
The compiled driver program accepts the name of the extracted kernel file as
a command line parameter. If this is not specified, it will use the default
name (\sphinxcode{\sphinxupquote{module\sphinxhyphen{}region}} without a rank).


\subsubsection{Shared memory and API\sphinxhyphen{}specific}
\label{\detokenize{psyke:shared-memory-and-api-specific}}\label{\detokenize{psyke:psyke-intro-restrictions-shared}}
\sphinxAtStartPar
The \sphinxcode{\sphinxupquote{ExtractTrans}} transformation cannot be applied to:
\begin{itemize}
\item {} 
\sphinxAtStartPar
A Loop without its parent Directive,

\item {} 
\sphinxAtStartPar
An orphaned Directive (e.g. \sphinxcode{\sphinxupquote{OMPDoDirective}}, \sphinxcode{\sphinxupquote{ACCLoopDirective}})
without its parent Directive (e.g. ACC or OMP Parallel Directive),

\item {} 
\sphinxAtStartPar
A Loop over cells in a colour without its parent Loop over colours in
the LFRic API,

\item {} 
\sphinxAtStartPar
An inner Loop without its parent outer Loop in the GOcean API.

\item {} 
\sphinxAtStartPar
Kernels that have a halo exchange call between them.

\end{itemize}


\section{Use}
\label{\detokenize{psyke:use}}\label{\detokenize{psyke:psyke-use}}
\sphinxAtStartPar
The code extraction is currently enabled by utilising a transformation
script (see {\hyperref[\detokenize{transformations:sec-transformations-script}]{\sphinxcrossref{\DUrole{std,std-ref}{PSyclone User Scripts}}}} section for more details).

\sphinxAtStartPar
For example, the transformation script which extracts the first Kernel call
in LFRic API test example \sphinxcode{\sphinxupquote{15.1.2\_builtin\_and\_normal\_kernel\_invoke.f90}}
would be written as:

\begin{sphinxVerbatim}[commandchars=\\\{\}]
\PYG{k+kn}{from}\PYG{+w}{ }\PYG{n+nn}{psyclone}\PYG{n+nn}{.}\PYG{n+nn}{domain}\PYG{n+nn}{.}\PYG{n+nn}{lfric}\PYG{n+nn}{.}\PYG{n+nn}{transformations}\PYG{+w}{ }\PYG{k+kn}{import} \PYG{n}{LFRicExtractTrans}

\PYG{c+c1}{\PYGZsh{} Get instance of the ExtractRegionTrans transformation}
\PYG{n}{etrans} \PYG{o}{=} \PYG{n}{LFRicExtractTrans}\PYG{p}{(}\PYG{p}{)}

\PYG{c+c1}{\PYGZsh{} Get Invoke and its Schedule}
\PYG{n}{invoke} \PYG{o}{=} \PYG{n}{psy}\PYG{o}{.}\PYG{n}{invokes}\PYG{o}{.}\PYG{n}{get}\PYG{p}{(}\PYG{l+s+s2}{\PYGZdq{}}\PYG{l+s+s2}{invoke\PYGZus{}0}\PYG{l+s+s2}{\PYGZdq{}}\PYG{p}{)}
\PYG{n}{schedule} \PYG{o}{=} \PYG{n}{invoke}\PYG{o}{.}\PYG{n}{schedule}

\PYG{c+c1}{\PYGZsh{} Apply extract transformation to the selected Node}
\PYG{n}{etrans}\PYG{o}{.}\PYG{n}{apply}\PYG{p}{(}\PYG{n}{schedule}\PYG{o}{.}\PYG{n}{children}\PYG{p}{[}\PYG{l+m+mi}{2}\PYG{p}{]}\PYG{p}{)}
\PYG{n+nb}{print}\PYG{p}{(}\PYG{n}{schedule}\PYG{o}{.}\PYG{n}{view}\PYG{p}{(}\PYG{p}{)}\PYG{p}{)}
\end{sphinxVerbatim}

\sphinxAtStartPar
and called as:

\begin{sphinxVerbatim}[commandchars=\\\{\}]
\PYGZgt{}\PYG{+w}{ }psyclone\PYG{+w}{ }\PYGZhy{}nodm\PYG{+w}{ }\PYGZhy{}s\PYG{+w}{ }./extract\PYGZus{}single\PYGZus{}node.py\PYG{+w}{ }\PYG{l+s+se}{\PYGZbs{}}
\PYG{+w}{    }\PYGZlt{}path\PYGZhy{}to\PYGZhy{}example\PYGZgt{}/15.1.2\PYGZus{}builtin\PYGZus{}and\PYGZus{}normal\PYGZus{}kernel\PYGZus{}invoke.f90
\end{sphinxVerbatim}

\sphinxAtStartPar
PSyclone modifies the Schedule of the selected \sphinxcode{\sphinxupquote{invoke\_0}}:

\begin{sphinxVerbatim}[commandchars=\\\{\}]
\PYG{n}{Schedule}\PYG{p}{[}\PYG{n}{invoke}\PYG{o}{=}\PYG{l+s+s1}{\PYGZsq{}invoke\PYGZus{}0\PYGZsq{}}\PYG{+w}{ }\PYG{n}{dm}\PYG{o}{=}\PYG{n}{False}\PYG{p}{]}
\PYG{+w}{    }\PYG{l+m+mi}{0}\PYG{p}{:}\PYG{+w}{ }\PYG{n}{Loop}\PYG{p}{[}\PYG{k}{type}\PYG{o}{=}\PYG{l+s+s1}{\PYGZsq{}dofs\PYGZsq{}}\PYG{p}{,}\PYG{n}{field\PYGZus{}space}\PYG{o}{=}\PYG{l+s+s1}{\PYGZsq{}any\PYGZus{}space\PYGZus{}1\PYGZsq{}}\PYG{p}{,}\PYG{n}{it\PYGZus{}space}\PYG{o}{=}\PYG{l+s+s1}{\PYGZsq{}dofs\PYGZsq{}}\PYG{p}{,}
\PYG{+w}{            }\PYG{n}{upper\PYGZus{}bound}\PYG{o}{=}\PYG{l+s+s1}{\PYGZsq{}ndofs\PYGZsq{}}\PYG{p}{]}
\PYG{+w}{        }\PYG{n}{Literal}\PYG{p}{[}\PYG{k}{value}\PYG{p}{:}\PYG{l+s+s1}{\PYGZsq{}NOT\PYGZus{}INITIALISED\PYGZsq{}}\PYG{p}{]}
\PYG{+w}{        }\PYG{n}{Literal}\PYG{p}{[}\PYG{k}{value}\PYG{p}{:}\PYG{l+s+s1}{\PYGZsq{}NOT\PYGZus{}INITIALISED\PYGZsq{}}\PYG{p}{]}
\PYG{+w}{        }\PYG{n}{Literal}\PYG{p}{[}\PYG{k}{value}\PYG{p}{:}\PYG{l+s+s1}{\PYGZsq{}1\PYGZsq{}}\PYG{p}{]}
\PYG{+w}{        }\PYG{n}{Schedule}\PYG{p}{[}\PYG{p}{]}
\PYG{+w}{            }\PYG{l+m+mi}{0}\PYG{p}{:}\PYG{+w}{ }\PYG{n}{BuiltIn}\PYG{+w}{ }\PYG{n}{setval\PYGZus{}c}\PYG{p}{(}\PYG{n}{f5}\PYG{p}{,}\PYG{l+m+mf}{0.0}\PYG{p}{)}
\PYG{+w}{    }\PYG{l+m+mi}{1}\PYG{p}{:}\PYG{+w}{ }\PYG{n}{Loop}\PYG{p}{[}\PYG{k}{type}\PYG{o}{=}\PYG{l+s+s1}{\PYGZsq{}dofs\PYGZsq{}}\PYG{p}{,}\PYG{n}{field\PYGZus{}space}\PYG{o}{=}\PYG{l+s+s1}{\PYGZsq{}any\PYGZus{}space\PYGZus{}1\PYGZsq{}}\PYG{p}{,}\PYG{n}{it\PYGZus{}space}\PYG{o}{=}\PYG{l+s+s1}{\PYGZsq{}dofs\PYGZsq{}}\PYG{p}{,}
\PYG{+w}{            }\PYG{n}{upper\PYGZus{}bound}\PYG{o}{=}\PYG{l+s+s1}{\PYGZsq{}ndofs\PYGZsq{}}\PYG{p}{]}
\PYG{+w}{        }\PYG{p}{.}\PYG{p}{.}\PYG{p}{.}
\PYG{+w}{        }\PYG{n}{Schedule}\PYG{p}{[}\PYG{p}{]}
\PYG{+w}{            }\PYG{l+m+mi}{0}\PYG{p}{:}\PYG{+w}{ }\PYG{n}{BuiltIn}\PYG{+w}{ }\PYG{n}{setval\PYGZus{}c}\PYG{p}{(}\PYG{n}{f2}\PYG{p}{,}\PYG{l+m+mf}{0.0}\PYG{p}{)}
\PYG{+w}{    }\PYG{l+m+mi}{2}\PYG{p}{:}\PYG{+w}{ }\PYG{n}{Loop}\PYG{p}{[}\PYG{k}{type}\PYG{o}{=}\PYG{l+s+s1}{\PYGZsq{}\PYGZsq{}}\PYG{p}{,}\PYG{n}{field\PYGZus{}space}\PYG{o}{=}\PYG{l+s+s1}{\PYGZsq{}w2\PYGZsq{}}\PYG{p}{,}\PYG{n}{it\PYGZus{}space}\PYG{o}{=}\PYG{l+s+s1}{\PYGZsq{}cells\PYGZsq{}}\PYG{p}{,}\PYG{+w}{ }\PYG{n}{upper\PYGZus{}bound}\PYG{o}{=}\PYG{l+s+s1}{\PYGZsq{}ncells\PYGZsq{}}\PYG{p}{]}
\PYG{+w}{        }\PYG{p}{.}\PYG{p}{.}\PYG{p}{.}
\PYG{+w}{        }\PYG{n}{Schedule}\PYG{p}{[}\PYG{p}{]}
\PYG{+w}{            }\PYG{l+m+mi}{0}\PYG{p}{:}\PYG{+w}{ }\PYG{n}{CodedKern}\PYG{+w}{ }\PYG{n}{testkern\PYGZus{}code\PYGZus{}w2\PYGZus{}only}\PYG{p}{(}\PYG{n}{f3}\PYG{p}{,}\PYG{n}{f2}\PYG{p}{)}\PYG{+w}{ }\PYG{p}{[}\PYG{n}{module\PYGZus{}inline}\PYG{o}{=}\PYG{n}{False}\PYG{p}{]}
\PYG{+w}{    }\PYG{l+m+mi}{3}\PYG{p}{:}\PYG{+w}{ }\PYG{n}{Loop}\PYG{p}{[}\PYG{k}{type}\PYG{o}{=}\PYG{l+s+s1}{\PYGZsq{}\PYGZsq{}}\PYG{p}{,}\PYG{n}{field\PYGZus{}space}\PYG{o}{=}\PYG{l+s+s1}{\PYGZsq{}wtheta\PYGZsq{}}\PYG{p}{,}\PYG{n}{it\PYGZus{}space}\PYG{o}{=}\PYG{l+s+s1}{\PYGZsq{}cells\PYGZsq{}}\PYG{p}{,}\PYG{+w}{ }\PYG{n}{upper\PYGZus{}bound}\PYG{o}{=}\PYG{l+s+s1}{\PYGZsq{}ncells\PYGZsq{}}\PYG{p}{]}
\PYG{+w}{        }\PYG{p}{.}\PYG{p}{.}\PYG{p}{.}
\PYG{+w}{        }\PYG{n}{Schedule}\PYG{p}{[}\PYG{p}{]}
\PYG{+w}{            }\PYG{l+m+mi}{0}\PYG{p}{:}\PYG{+w}{ }\PYG{n}{CodedKern}\PYG{+w}{ }\PYG{n}{testkern\PYGZus{}wtheta\PYGZus{}code}\PYG{p}{(}\PYG{n}{f4}\PYG{p}{,}\PYG{n}{f5}\PYG{p}{)}\PYG{+w}{ }\PYG{p}{[}\PYG{n}{module\PYGZus{}inline}\PYG{o}{=}\PYG{n}{False}\PYG{p}{]}
\PYG{+w}{    }\PYG{l+m+mi}{4}\PYG{p}{:}\PYG{+w}{ }\PYG{n}{Loop}\PYG{p}{[}\PYG{k}{type}\PYG{o}{=}\PYG{l+s+s1}{\PYGZsq{}\PYGZsq{}}\PYG{p}{,}\PYG{n}{field\PYGZus{}space}\PYG{o}{=}\PYG{l+s+s1}{\PYGZsq{}w1\PYGZsq{}}\PYG{p}{,}\PYG{n}{it\PYGZus{}space}\PYG{o}{=}\PYG{l+s+s1}{\PYGZsq{}cells\PYGZsq{}}\PYG{p}{,}\PYG{+w}{ }\PYG{n}{upper\PYGZus{}bound}\PYG{o}{=}\PYG{l+s+s1}{\PYGZsq{}ncells\PYGZsq{}}\PYG{p}{]}
\PYG{+w}{        }\PYG{p}{.}\PYG{p}{.}\PYG{p}{.}
\PYG{+w}{        }\PYG{n}{Schedule}\PYG{p}{[}\PYG{p}{]}
\PYG{+w}{            }\PYG{l+m+mi}{0}\PYG{p}{:}\PYG{+w}{ }\PYG{n}{CodedKern}\PYG{+w}{ }\PYG{n}{testkern\PYGZus{}code}\PYG{p}{(}\PYG{n}{scalar}\PYG{p}{,}\PYG{n}{f1}\PYG{p}{,}\PYG{n}{f2}\PYG{p}{,}\PYG{n}{f3}\PYG{p}{,}\PYG{n}{f4}\PYG{p}{)}\PYG{+w}{ }\PYG{p}{[}\PYG{n}{module\PYGZus{}inline}\PYG{o}{=}\PYG{n}{False}\PYG{p}{]}
\end{sphinxVerbatim}

\sphinxAtStartPar
to insert the extract region. As shown below, all children of an
\sphinxcode{\sphinxupquote{ExtractNode}} will be part of the region:

\begin{sphinxVerbatim}[commandchars=\\\{\}]
\PYG{n}{Schedule}\PYG{p}{[}\PYG{n}{invoke}\PYG{o}{=}\PYG{l+s+s1}{\PYGZsq{}invoke\PYGZus{}0\PYGZsq{}}\PYG{+w}{ }\PYG{n}{dm}\PYG{o}{=}\PYG{n}{False}\PYG{p}{]}
\PYG{+w}{    }\PYG{l+m+mi}{0}\PYG{p}{:}\PYG{+w}{ }\PYG{n}{Loop}\PYG{p}{[}\PYG{k}{type}\PYG{o}{=}\PYG{l+s+s1}{\PYGZsq{}dofs\PYGZsq{}}\PYG{p}{,}\PYG{n}{field\PYGZus{}space}\PYG{o}{=}\PYG{l+s+s1}{\PYGZsq{}any\PYGZus{}space\PYGZus{}1\PYGZsq{}}\PYG{p}{,}\PYG{n}{it\PYGZus{}space}\PYG{o}{=}\PYG{l+s+s1}{\PYGZsq{}dofs\PYGZsq{}}\PYG{p}{,}
\PYG{+w}{            }\PYG{n}{upper\PYGZus{}bound}\PYG{o}{=}\PYG{l+s+s1}{\PYGZsq{}ndofs\PYGZsq{}}\PYG{p}{]}
\PYG{+w}{        }\PYG{p}{.}\PYG{p}{.}\PYG{p}{.}
\PYG{+w}{        }\PYG{n}{Schedule}\PYG{p}{[}\PYG{p}{]}
\PYG{+w}{            }\PYG{l+m+mi}{0}\PYG{p}{:}\PYG{+w}{ }\PYG{n}{BuiltIn}\PYG{+w}{ }\PYG{n}{setval\PYGZus{}c}\PYG{p}{(}\PYG{n}{f5}\PYG{p}{,}\PYG{l+m+mf}{0.0}\PYG{p}{)}
\PYG{+w}{    }\PYG{l+m+mi}{1}\PYG{p}{:}\PYG{+w}{ }\PYG{n}{Loop}\PYG{p}{[}\PYG{k}{type}\PYG{o}{=}\PYG{l+s+s1}{\PYGZsq{}dofs\PYGZsq{}}\PYG{p}{,}\PYG{n}{field\PYGZus{}space}\PYG{o}{=}\PYG{l+s+s1}{\PYGZsq{}any\PYGZus{}space\PYGZus{}1\PYGZsq{}}\PYG{p}{,}\PYG{n}{it\PYGZus{}space}\PYG{o}{=}\PYG{l+s+s1}{\PYGZsq{}dofs\PYGZsq{}}\PYG{p}{,}
\PYG{+w}{            }\PYG{n}{upper\PYGZus{}bound}\PYG{o}{=}\PYG{l+s+s1}{\PYGZsq{}ndofs\PYGZsq{}}\PYG{p}{]}
\PYG{+w}{        }\PYG{p}{.}\PYG{p}{.}\PYG{p}{.}
\PYG{+w}{        }\PYG{n}{Schedule}\PYG{p}{[}\PYG{p}{]}
\PYG{+w}{            }\PYG{l+m+mi}{0}\PYG{p}{:}\PYG{+w}{ }\PYG{n}{BuiltIn}\PYG{+w}{ }\PYG{n}{setval\PYGZus{}c}\PYG{p}{(}\PYG{n}{f2}\PYG{p}{,}\PYG{l+m+mf}{0.0}\PYG{p}{)}
\PYG{+w}{    }\PYG{l+m+mi}{2}\PYG{p}{:}\PYG{+w}{ }\PYG{n}{Extract}
\PYG{+w}{        }\PYG{n}{Schedule}\PYG{p}{[}\PYG{p}{]}
\PYG{+w}{            }\PYG{l+m+mi}{0}\PYG{p}{:}\PYG{+w}{ }\PYG{n}{Loop}\PYG{p}{[}\PYG{k}{type}\PYG{o}{=}\PYG{l+s+s1}{\PYGZsq{}\PYGZsq{}}\PYG{p}{,}\PYG{n}{field\PYGZus{}space}\PYG{o}{=}\PYG{l+s+s1}{\PYGZsq{}w2\PYGZsq{}}\PYG{p}{,}\PYG{n}{it\PYGZus{}space}\PYG{o}{=}\PYG{l+s+s1}{\PYGZsq{}cells\PYGZsq{}}\PYG{p}{,}\PYG{+w}{ }\PYG{n}{upper\PYGZus{}bound}\PYG{o}{=}\PYG{l+s+s1}{\PYGZsq{}ncells\PYGZsq{}}\PYG{p}{]}
\PYG{+w}{                }\PYG{p}{.}\PYG{p}{.}\PYG{p}{.}
\PYG{+w}{                }\PYG{n}{Schedule}\PYG{p}{[}\PYG{p}{]}
\PYG{+w}{                    }\PYG{l+m+mi}{0}\PYG{p}{:}\PYG{+w}{ }\PYG{n}{CodedKern}\PYG{+w}{ }\PYG{n}{testkern\PYGZus{}code\PYGZus{}w2\PYGZus{}only}\PYG{p}{(}\PYG{n}{f3}\PYG{p}{,}\PYG{n}{f2}\PYG{p}{)}\PYG{+w}{ }\PYG{p}{[}\PYG{n}{module\PYGZus{}inline}\PYG{o}{=}\PYG{n}{False}\PYG{p}{]}
\PYG{+w}{    }\PYG{l+m+mi}{3}\PYG{p}{:}\PYG{+w}{ }\PYG{n}{Loop}\PYG{p}{[}\PYG{k}{type}\PYG{o}{=}\PYG{l+s+s1}{\PYGZsq{}\PYGZsq{}}\PYG{p}{,}\PYG{n}{field\PYGZus{}space}\PYG{o}{=}\PYG{l+s+s1}{\PYGZsq{}wtheta\PYGZsq{}}\PYG{p}{,}\PYG{n}{it\PYGZus{}space}\PYG{o}{=}\PYG{l+s+s1}{\PYGZsq{}cells\PYGZsq{}}\PYG{p}{,}\PYG{+w}{ }\PYG{n}{upper\PYGZus{}bound}\PYG{o}{=}\PYG{l+s+s1}{\PYGZsq{}ncells\PYGZsq{}}\PYG{p}{]}
\PYG{+w}{        }\PYG{p}{.}\PYG{p}{.}\PYG{p}{.}
\PYG{+w}{        }\PYG{n}{Schedule}\PYG{p}{[}\PYG{p}{]}
\PYG{+w}{            }\PYG{l+m+mi}{0}\PYG{p}{:}\PYG{+w}{ }\PYG{n}{CodedKern}\PYG{+w}{ }\PYG{n}{testkern\PYGZus{}wtheta\PYGZus{}code}\PYG{p}{(}\PYG{n}{f4}\PYG{p}{,}\PYG{n}{f5}\PYG{p}{)}\PYG{+w}{ }\PYG{p}{[}\PYG{n}{module\PYGZus{}inline}\PYG{o}{=}\PYG{n}{False}\PYG{p}{]}
\PYG{+w}{    }\PYG{l+m+mi}{4}\PYG{p}{:}\PYG{+w}{ }\PYG{n}{Loop}\PYG{p}{[}\PYG{k}{type}\PYG{o}{=}\PYG{l+s+s1}{\PYGZsq{}\PYGZsq{}}\PYG{p}{,}\PYG{n}{field\PYGZus{}space}\PYG{o}{=}\PYG{l+s+s1}{\PYGZsq{}w1\PYGZsq{}}\PYG{p}{,}\PYG{n}{it\PYGZus{}space}\PYG{o}{=}\PYG{l+s+s1}{\PYGZsq{}cells\PYGZsq{}}\PYG{p}{,}\PYG{+w}{ }\PYG{n}{upper\PYGZus{}bound}\PYG{o}{=}\PYG{l+s+s1}{\PYGZsq{}ncells\PYGZsq{}}\PYG{p}{]}
\PYG{+w}{        }\PYG{p}{.}\PYG{p}{.}\PYG{p}{.}
\PYG{+w}{        }\PYG{n}{Schedule}\PYG{p}{[}\PYG{p}{]}
\PYG{+w}{            }\PYG{l+m+mi}{0}\PYG{p}{:}\PYG{+w}{ }\PYG{n}{CodedKern}\PYG{+w}{ }\PYG{n}{testkern\PYGZus{}code}\PYG{p}{(}\PYG{n}{scalar}\PYG{p}{,}\PYG{n}{f1}\PYG{p}{,}\PYG{n}{f2}\PYG{p}{,}\PYG{n}{f3}\PYG{p}{,}\PYG{n}{f4}\PYG{p}{)}\PYG{+w}{ }\PYG{p}{[}\PYG{n}{module\PYGZus{}inline}\PYG{o}{=}\PYG{n}{False}\PYG{p}{]}
\end{sphinxVerbatim}

\sphinxAtStartPar
To extract multiple Nodes, \sphinxcode{\sphinxupquote{ExtractTrans}} can be applied to the list
of Nodes (subject to {\hyperref[\detokenize{psyke:psyke-intro-restrictions-gen}]{\sphinxcrossref{\DUrole{std,std-ref}{General}}}} restrictions above):

\begin{sphinxVerbatim}[commandchars=\\\{\}]
\PYG{c+c1}{\PYGZsh{} Apply extract transformation to the selected Nodes}
\PYG{n}{etrans}\PYG{o}{.}\PYG{n}{apply}\PYG{p}{(}\PYG{n}{schedule}\PYG{o}{.}\PYG{n}{children}\PYG{p}{[}\PYG{l+m+mi}{1}\PYG{p}{:}\PYG{l+m+mi}{3}\PYG{p}{]}\PYG{p}{)}
\end{sphinxVerbatim}

\sphinxAtStartPar
This modifies the above Schedule as:

\begin{sphinxVerbatim}[commandchars=\\\{\}]
\PYG{p}{.}\PYG{p}{.}\PYG{p}{.}
\PYG{+w}{    }\PYG{n}{Extract}
\PYG{+w}{        }\PYG{n}{Schedule}\PYG{p}{[}\PYG{p}{]}
\PYG{+w}{            }\PYG{l+m+mi}{0}\PYG{p}{:}\PYG{+w}{ }\PYG{n}{Loop}\PYG{p}{[}\PYG{k}{type}\PYG{o}{=}\PYG{l+s+s1}{\PYGZsq{}dofs\PYGZsq{}}\PYG{p}{,}\PYG{n}{field\PYGZus{}space}\PYG{o}{=}\PYG{l+s+s1}{\PYGZsq{}any\PYGZus{}space\PYGZus{}1\PYGZsq{}}\PYG{p}{,}\PYG{n}{it\PYGZus{}space}\PYG{o}{=}\PYG{l+s+s1}{\PYGZsq{}dofs\PYGZsq{}}\PYG{p}{,}
\PYG{+w}{                    }\PYG{n}{upper\PYGZus{}bound}\PYG{o}{=}\PYG{l+s+s1}{\PYGZsq{}ndofs\PYGZsq{}}\PYG{p}{]}
\PYG{+w}{                }\PYG{p}{.}\PYG{p}{.}\PYG{p}{.}
\PYG{+w}{                }\PYG{n}{Schedule}\PYG{p}{[}\PYG{p}{]}
\PYG{+w}{                    }\PYG{l+m+mi}{0}\PYG{p}{:}\PYG{+w}{ }\PYG{n}{BuiltIn}\PYG{+w}{ }\PYG{n}{setval\PYGZus{}c}\PYG{p}{(}\PYG{n}{f2}\PYG{p}{,}\PYG{l+m+mf}{0.0}\PYG{p}{)}
\PYG{+w}{            }\PYG{l+m+mi}{1}\PYG{p}{:}\PYG{+w}{ }\PYG{n}{Loop}\PYG{p}{[}\PYG{k}{type}\PYG{o}{=}\PYG{l+s+s1}{\PYGZsq{}\PYGZsq{}}\PYG{p}{,}\PYG{n}{field\PYGZus{}space}\PYG{o}{=}\PYG{l+s+s1}{\PYGZsq{}w2\PYGZsq{}}\PYG{p}{,}\PYG{n}{it\PYGZus{}space}\PYG{o}{=}\PYG{l+s+s1}{\PYGZsq{}cells\PYGZsq{}}\PYG{p}{,}\PYG{+w}{ }\PYG{n}{upper\PYGZus{}bound}\PYG{o}{=}\PYG{l+s+s1}{\PYGZsq{}ncells\PYGZsq{}}\PYG{p}{]}
\PYG{+w}{                }\PYG{p}{.}\PYG{p}{.}\PYG{p}{.}
\PYG{+w}{                }\PYG{n}{Schedule}\PYG{p}{[}\PYG{p}{]}
\PYG{+w}{                    }\PYG{l+m+mi}{0}\PYG{p}{:}\PYG{+w}{ }\PYG{n}{CodedKern}\PYG{+w}{ }\PYG{n}{testkern\PYGZus{}code\PYGZus{}w2\PYGZus{}only}\PYG{p}{(}\PYG{n}{f3}\PYG{p}{,}\PYG{n}{f2}\PYG{p}{)}\PYG{+w}{ }\PYG{p}{[}\PYG{n}{module\PYGZus{}inline}\PYG{o}{=}\PYG{n}{False}\PYG{p}{]}
\PYG{p}{.}\PYG{p}{.}\PYG{p}{.}
\end{sphinxVerbatim}

\sphinxAtStartPar
As said above, extraction can be performed on optimised code. For example,
the following example transformation script first adds \sphinxcode{\sphinxupquote{!\$OMP PARALLEL DO}}
directive and then extracts the optimised code in LFRic API test
example \sphinxcode{\sphinxupquote{15.1.2\_builtin\_and\_normal\_kernel\_invoke.f90}}:

\begin{sphinxVerbatim}[commandchars=\\\{\}]
\PYG{k+kn}{from}\PYG{+w}{ }\PYG{n+nn}{psyclone}\PYG{n+nn}{.}\PYG{n+nn}{domain}\PYG{n+nn}{.}\PYG{n+nn}{lfric}\PYG{n+nn}{.}\PYG{n+nn}{transformations}\PYG{+w}{ }\PYG{k+kn}{import} \PYG{n}{LFRicExtractTrans}
\PYG{k+kn}{from}\PYG{+w}{ }\PYG{n+nn}{psyclone}\PYG{n+nn}{.}\PYG{n+nn}{transformations}\PYG{+w}{ }\PYG{k+kn}{import} \PYG{n}{DynamoOMPParallelLoopTrans}

\PYG{c+c1}{\PYGZsh{} Get instances of the transformations}
\PYG{n}{etrans} \PYG{o}{=} \PYG{n}{LFRicExtractTrans}\PYG{p}{(}\PYG{p}{)}
\PYG{n}{otrans} \PYG{o}{=} \PYG{n}{DynamoOMPParallelLoopTrans}\PYG{p}{(}\PYG{p}{)}

\PYG{c+c1}{\PYGZsh{} Get Invoke and its Schedule}
\PYG{n}{invoke} \PYG{o}{=} \PYG{n}{psy}\PYG{o}{.}\PYG{n}{invokes}\PYG{o}{.}\PYG{n}{get}\PYG{p}{(}\PYG{l+s+s2}{\PYGZdq{}}\PYG{l+s+s2}{invoke\PYGZus{}0}\PYG{l+s+s2}{\PYGZdq{}}\PYG{p}{)}
\PYG{n}{schedule} \PYG{o}{=} \PYG{n}{invoke}\PYG{o}{.}\PYG{n}{schedule}

\PYG{c+c1}{\PYGZsh{} Add OMP PARALLEL DO directives}
\PYG{n}{otrans}\PYG{o}{.}\PYG{n}{apply}\PYG{p}{(}\PYG{n}{schedule}\PYG{o}{.}\PYG{n}{children}\PYG{p}{[}\PYG{l+m+mi}{1}\PYG{p}{]}\PYG{p}{)}
\PYG{n}{otrans}\PYG{o}{.}\PYG{n}{apply}\PYG{p}{(}\PYG{n}{schedule}\PYG{o}{.}\PYG{n}{children}\PYG{p}{[}\PYG{l+m+mi}{2}\PYG{p}{]}\PYG{p}{)}
\PYG{c+c1}{\PYGZsh{} Apply extract transformation to the selected Nodes}
\PYG{n}{etrans}\PYG{o}{.}\PYG{n}{apply}\PYG{p}{(}\PYG{n}{schedule}\PYG{o}{.}\PYG{n}{children}\PYG{p}{[}\PYG{l+m+mi}{1}\PYG{p}{:}\PYG{l+m+mi}{3}\PYG{p}{]}\PYG{p}{)}
\PYG{n+nb}{print}\PYG{p}{(}\PYG{n}{schedule}\PYG{o}{.}\PYG{n}{view}\PYG{p}{(}\PYG{p}{)}\PYG{p}{)}
\end{sphinxVerbatim}

\sphinxAtStartPar
The generated code is now:

\begin{sphinxVerbatim}[commandchars=\\\{\}]
\PYG{c}{! ExtractStart}
\PYG{k}{CALL }\PYG{n}{extract\PYGZus{}psy\PYGZus{}data}\PYG{p}{\PYGZpc{}}\PYG{n}{PreStart}\PYG{p}{(}\PYG{l+s+s2}{\PYGZdq{}unknown\PYGZhy{}module\PYGZdq{}}\PYG{p}{,}\PYG{+w}{ }\PYG{l+s+s2}{\PYGZdq{}setval\PYGZus{}c\PYGZdq{}}\PYG{p}{,}\PYG{+w}{ }\PYG{l+m+mi}{0}\PYG{p}{,}\PYG{+w}{ }\PYG{l+m+mi}{4}\PYG{p}{)}
\PYG{k}{CALL }\PYG{n}{extract\PYGZus{}psy\PYGZus{}data}\PYG{p}{\PYGZpc{}}\PYG{n}{PreDeclareVariable}\PYG{p}{(}\PYG{l+s+s2}{\PYGZdq{}cell\PYGZus{}post\PYGZdq{}}\PYG{p}{,}\PYG{+w}{ }\PYG{n}{cell}\PYG{p}{)}
\PYG{k}{CALL }\PYG{n}{extract\PYGZus{}psy\PYGZus{}data}\PYG{p}{\PYGZpc{}}\PYG{n}{PreDeclareVariable}\PYG{p}{(}\PYG{l+s+s2}{\PYGZdq{}df\PYGZus{}post\PYGZdq{}}\PYG{p}{,}\PYG{+w}{ }\PYG{n}{df}\PYG{p}{)}
\PYG{k}{CALL }\PYG{n}{extract\PYGZus{}psy\PYGZus{}data}\PYG{p}{\PYGZpc{}}\PYG{n}{PreDeclareVariable}\PYG{p}{(}\PYG{l+s+s2}{\PYGZdq{}f2\PYGZus{}post\PYGZdq{}}\PYG{p}{,}\PYG{+w}{ }\PYG{n}{f2}\PYG{p}{)}
\PYG{k}{CALL }\PYG{n}{extract\PYGZus{}psy\PYGZus{}data}\PYG{p}{\PYGZpc{}}\PYG{n}{PreDeclareVariable}\PYG{p}{(}\PYG{l+s+s2}{\PYGZdq{}f3\PYGZus{}post\PYGZdq{}}\PYG{p}{,}\PYG{+w}{ }\PYG{n}{f3}\PYG{p}{)}
\PYG{p}{.}\PYG{p}{.}\PYG{p}{.}
\PYG{k}{CALL }\PYG{n}{extract\PYGZus{}psy\PYGZus{}data}\PYG{p}{\PYGZpc{}}\PYG{n}{PreEndDeclaration}
\PYG{p}{.}\PYG{p}{.}\PYG{p}{.}
\PYG{k}{CALL }\PYG{n}{extract\PYGZus{}psy\PYGZus{}data}\PYG{p}{\PYGZpc{}}\PYG{n}{PreEnd}
\PYG{c}{!}
\PYG{c}{!\PYGZdl{}omp parallel do default(shared), private(df), schedule(static)}
\PYG{k}{DO }\PYG{n}{df}\PYG{o}{=}\PYG{l+m+mi}{1}\PYG{p}{,}\PYG{n}{undf\PYGZus{}aspc1\PYGZus{}f2}
\PYG{+w}{  }\PYG{n}{f2\PYGZus{}proxy}\PYG{p}{\PYGZpc{}}\PYG{k}{data}\PYG{p}{(}\PYG{n}{df}\PYG{p}{)}\PYG{+w}{ }\PYG{o}{=}\PYG{+w}{ }\PYG{l+m+mf}{0.0}
\PYG{k}{END }\PYG{k}{DO}
\PYG{c}{!\PYGZdl{}omp end parallel do}
\PYG{c}{!\PYGZdl{}omp parallel do default(shared), private(cell), schedule(static)}
\PYG{k}{DO }\PYG{n}{cell}\PYG{o}{=}\PYG{l+m+mi}{1}\PYG{p}{,}\PYG{n}{f3\PYGZus{}proxy}\PYG{p}{\PYGZpc{}}\PYG{n}{vspace}\PYG{p}{\PYGZpc{}}\PYG{n}{get\PYGZus{}ncell}\PYG{p}{(}\PYG{p}{)}
\PYG{+w}{  }\PYG{c}{!}
\PYG{+w}{  }\PYG{k}{CALL }\PYG{n}{testkern\PYGZus{}code\PYGZus{}w2\PYGZus{}only}\PYG{p}{(}\PYG{n}{nlayers}\PYG{p}{,}\PYG{+w}{ }\PYG{n}{f3\PYGZus{}proxy}\PYG{p}{\PYGZpc{}}\PYG{k}{data}\PYG{p}{,}\PYG{+w}{ }\PYG{n}{f2\PYGZus{}proxy}\PYG{p}{\PYGZpc{}}\PYG{k}{data}\PYG{p}{,}\PYG{+w}{ }\PYG{n}{ndf\PYGZus{}w2}\PYG{p}{,}\PYG{+w}{ }\PYG{n}{undf\PYGZus{}w2}\PYG{p}{,}\PYG{+w}{ }\PYG{n}{map\PYGZus{}w2}\PYG{p}{(}\PYG{p}{:}\PYG{p}{,}\PYG{n}{cell}\PYG{p}{)}\PYG{p}{)}
\PYG{k}{END }\PYG{k}{DO}
\PYG{c}{!\PYGZdl{}omp end parallel do}
\PYG{k}{CALL }\PYG{n}{extract\PYGZus{}psy\PYGZus{}data}\PYG{p}{\PYGZpc{}}\PYG{n}{PostStart}
\PYG{k}{CALL }\PYG{n}{extract\PYGZus{}psy\PYGZus{}data}\PYG{p}{\PYGZpc{}}\PYG{n}{ProvideVariable}\PYG{p}{(}\PYG{l+s+s2}{\PYGZdq{}cell\PYGZus{}post\PYGZdq{}}\PYG{p}{,}\PYG{+w}{ }\PYG{n}{cell}\PYG{p}{)}
\PYG{k}{CALL }\PYG{n}{extract\PYGZus{}psy\PYGZus{}data}\PYG{p}{\PYGZpc{}}\PYG{n}{ProvideVariable}\PYG{p}{(}\PYG{l+s+s2}{\PYGZdq{}df\PYGZus{}post\PYGZdq{}}\PYG{p}{,}\PYG{+w}{ }\PYG{n}{df}\PYG{p}{)}
\PYG{k}{CALL }\PYG{n}{extract\PYGZus{}psy\PYGZus{}data}\PYG{p}{\PYGZpc{}}\PYG{n}{ProvideVariable}\PYG{p}{(}\PYG{l+s+s2}{\PYGZdq{}f2\PYGZus{}post\PYGZdq{}}\PYG{p}{,}\PYG{+w}{ }\PYG{n}{f2}\PYG{p}{)}
\PYG{k}{CALL }\PYG{n}{extract\PYGZus{}psy\PYGZus{}data}\PYG{p}{\PYGZpc{}}\PYG{n}{ProvideVariable}\PYG{p}{(}\PYG{l+s+s2}{\PYGZdq{}f3\PYGZus{}post\PYGZdq{}}\PYG{p}{,}\PYG{+w}{ }\PYG{n}{f3}\PYG{p}{)}
\PYG{k}{CALL }\PYG{n}{extract\PYGZus{}psy\PYGZus{}data}\PYG{p}{\PYGZpc{}}\PYG{n}{PostEnd}
\PYG{c}{!}
\PYG{c}{! ExtractEnd}
\end{sphinxVerbatim}

\sphinxAtStartPar
Examples in \sphinxcode{\sphinxupquote{examples/lfric/eg12}} directory demonstrate how to
apply code extraction by utilising PSyclone transformation scripts
(see {\hyperref[\detokenize{tutorials_and_examples:examples}]{\sphinxcrossref{\DUrole{std,std-ref}{Examples}}}} section for more information). The code
in \sphinxcode{\sphinxupquote{examples/lfric/eg17/full\_example\_extract}} can be compiled and
run, and it will create two kernel data files.


\section{Extraction Libraries}
\label{\detokenize{psyke:extraction-libraries}}\label{\detokenize{psyke:id1}}
\sphinxAtStartPar
PSyclone comes with two extraction libraries: one is based on NetCDF
and will create NetCDF files to contain all input\sphinxhyphen{} and output\sphinxhyphen{}parameters.
The second one is a stand\sphinxhyphen{}alone library which uses only standard Fortran
IO to write and read kernel data. The binary files produced using this
library may not be portable between machines and compilers. If you
require such portability then please use the NetCDF extraction library.

\sphinxAtStartPar
The two extraction {\hyperref[\detokenize{libraries:libraries}]{\sphinxcrossref{\DUrole{std,std-ref}{libraries}}}} are in
\sphinxhref{https://github.com/stfc/PSyclone/tree/master/lib/extract/standalone}{lib/extract/standalone}.
and in
\sphinxhref{https://github.com/stfc/PSyclone/tree/master/lib/extract/netcdf}{lib/extract/netcdf}.

\sphinxAtStartPar
All versions of the extraction libraries can be compiled with MPI
support by setting the variable \sphinxcode{\sphinxupquote{MPI=yes}}:

\begin{sphinxVerbatim}[commandchars=\\\{\}]
make\PYG{+w}{ }\PYG{n+nv}{MPI}\PYG{o}{=}yes\PYG{+w}{ }...
\end{sphinxVerbatim}

\sphinxAtStartPar
The only difference is that the output files will now have the process
rank in the name. The compiled driver program accepts the name of the
extracted kernel file as a command line parameter. If this is not specified,
it will use the default name (\sphinxcode{\sphinxupquote{module\sphinxhyphen{}region}} without a rank).


\subsection{Extraction for GOcean}
\label{\detokenize{psyke:extraction-for-gocean}}\label{\detokenize{psyke:id2}}
\sphinxAtStartPar
The extraction libraries in
\sphinxhref{https://github.com/stfc/PSyclone/tree/master/lib/extract/standalone/dl\_esm\_inf}{lib/extract/standalone/dl\_esm\_inf}
and
\sphinxhref{https://github.com/stfc/PSyclone/tree/master/lib/extract/netcdf/dl\_esm\_inf}{lib/extract/netcdf/dl\_esm\_inf}
implement the full PSyData API for use with the
{\hyperref[\detokenize{gocean1p0:gocean-api}]{\sphinxcrossref{\DUrole{std,std-ref}{GOcean}}}} dl\_esm\_inf infrastructure library.
When running the instrumented executable, it will create either a binary or
a NetCDF file for each instrumented
code region. It includes all variables that are read before the code
is executed, and all variables that have been modified. The output
variables have the postfix \sphinxcode{\sphinxupquote{\_post}} attached to the names,
e.g. a variable \sphinxcode{\sphinxupquote{xyz}} that is read and written will be stored
with the name \sphinxcode{\sphinxupquote{xyz}} containing the input values, and the name
\sphinxcode{\sphinxupquote{xyz\_post}} containing the output values. Arrays have their size
explicitly stored (in case of NetCDF as dimensions): again the
variable \sphinxcode{\sphinxupquote{xyz}} will have its
sizes stored as \sphinxcode{\sphinxupquote{xyzdim1}}, \sphinxcode{\sphinxupquote{xyzdim2}} for the input values,
and output arrays use the name \sphinxcode{\sphinxupquote{xyz\_postdim1}}, \sphinxcode{\sphinxupquote{xyz\_postdim2}}.

\begin{sphinxadmonition}{note}{Note:}
\sphinxAtStartPar
The stand\sphinxhyphen{}alone library does not store the names of the
variables in the output file, but these names will be used
as variable names in the created driver.
\end{sphinxadmonition}

\sphinxAtStartPar
The output file contains the values of all variables used in the
subroutine. The \sphinxcode{\sphinxupquote{GOceanExtractTrans}} transformation can automatically
create a driver program which will read the corresponding output file,
call the instrumented region, and compare the results. In order to create
this driver program, the options parameter \sphinxcode{\sphinxupquote{create\_driver}} must
be set to true:

\begin{sphinxVerbatim}[commandchars=\\\{\}]
\PYG{n}{extract} \PYG{o}{=} \PYG{n}{GOceanExtractTrans}\PYG{p}{(}\PYG{p}{)}
\PYG{n}{extract}\PYG{o}{.}\PYG{n}{apply}\PYG{p}{(}\PYG{n}{schedule}\PYG{o}{.}\PYG{n}{children}\PYG{p}{,}
              \PYG{p}{\PYGZob{}}\PYG{l+s+s2}{\PYGZdq{}}\PYG{l+s+s2}{create\PYGZus{}driver}\PYG{l+s+s2}{\PYGZdq{}}\PYG{p}{:} \PYG{k+kc}{True}\PYG{p}{,}
               \PYG{l+s+s2}{\PYGZdq{}}\PYG{l+s+s2}{region\PYGZus{}name}\PYG{l+s+s2}{\PYGZdq{}}\PYG{p}{:} \PYG{p}{(}\PYG{l+s+s2}{\PYGZdq{}}\PYG{l+s+s2}{main}\PYG{l+s+s2}{\PYGZdq{}}\PYG{p}{,} \PYG{l+s+s2}{\PYGZdq{}}\PYG{l+s+s2}{init}\PYG{l+s+s2}{\PYGZdq{}}\PYG{p}{)}\PYG{p}{\PYGZcb{}}\PYG{p}{)}
\end{sphinxVerbatim}

\sphinxAtStartPar
This will create a Fortran file called \sphinxcode{\sphinxupquote{driver\sphinxhyphen{}main\sphinxhyphen{}init.f90}}, which
can then be compiled and executed. This stand\sphinxhyphen{}alone program will read
the output file created during an execution of the actual program, call
the kernel with all required input parameter, and compare the output
variables with the original output variables. This can be used to create
stand\sphinxhyphen{}alone test cases to reproduce a bug, or for performance
optimisation of a stand\sphinxhyphen{}alone kernel.

\begin{sphinxadmonition}{warning}{Warning:}
\sphinxAtStartPar
Care has to be taken that the driver matches the version
of the code that was used to create the output file, otherwise the
driver will likely crash. The stand\sphinxhyphen{}alone driver relies on a
strict ordering of variable values in the output file and e.g.
even renaming one variable can affect this. The NetCDF version
stores the variable names and will not be able to find a variable
if its name has changed.
\end{sphinxadmonition}


\subsection{Extraction for LFRic}
\label{\detokenize{psyke:extraction-for-lfric}}
\sphinxAtStartPar
The libraries in
\sphinxhref{https://github.com/stfc/PSyclone/tree/master/lib/extract/standalone/lfric}{lib/extract/standalone/lfric}
and
\sphinxhref{https://github.com/stfc/PSyclone/tree/master/lib/extract/netcdf/lfric}{lib/extract/netcdf/lfric}
implement the full PSyData API for use with the
{\hyperref[\detokenize{dynamo0p3:lfric-api}]{\sphinxcrossref{\DUrole{std,std-ref}{LFRic}}}} infrastructure library. When running the
code, it will create an output file for each instrumented code region.
The same logic for naming variables (using \sphinxcode{\sphinxupquote{\_post}} for output variables)
used in {\hyperref[\detokenize{psyke:extraction-for-gocean}]{\sphinxcrossref{\DUrole{std,std-ref}{Extraction for GOcean}}}} is used here.

\sphinxAtStartPar
Check {\hyperref[\detokenize{psy_data:integrating-psy-data-lfric}]{\sphinxcrossref{\DUrole{std,std-ref}{Integrating PSyData Libraries into the LFRic Build Environment}}}} for the recommended way of linking
an extraction library to LFRic.

\sphinxAtStartPar
The output file contains the values of all variables used in the
subroutine. The \sphinxcode{\sphinxupquote{LFRicExtractTrans}} transformation can automatically
create a driver program which will read the corresponding output file,
call the instrumented region, and compare the results. In order to create
this driver program, the options parameter \sphinxcode{\sphinxupquote{create\_driver}} must
be set to true:

\begin{sphinxVerbatim}[commandchars=\\\{\}]
\PYG{n}{extract} \PYG{o}{=} \PYG{n}{LFRicExtractTrans}\PYG{p}{(}\PYG{p}{)}
\PYG{n}{extract}\PYG{o}{.}\PYG{n}{apply}\PYG{p}{(}\PYG{n}{schedule}\PYG{o}{.}\PYG{n}{children}\PYG{p}{,}
              \PYG{p}{\PYGZob{}}\PYG{l+s+s2}{\PYGZdq{}}\PYG{l+s+s2}{create\PYGZus{}driver}\PYG{l+s+s2}{\PYGZdq{}}\PYG{p}{:} \PYG{k+kc}{True}\PYG{p}{,}
               \PYG{l+s+s2}{\PYGZdq{}}\PYG{l+s+s2}{region\PYGZus{}name}\PYG{l+s+s2}{\PYGZdq{}}\PYG{p}{:} \PYG{p}{(}\PYG{l+s+s2}{\PYGZdq{}}\PYG{l+s+s2}{main}\PYG{l+s+s2}{\PYGZdq{}}\PYG{p}{,} \PYG{l+s+s2}{\PYGZdq{}}\PYG{l+s+s2}{init}\PYG{l+s+s2}{\PYGZdq{}}\PYG{p}{)}\PYG{p}{\PYGZcb{}}\PYG{p}{)}
\end{sphinxVerbatim}

\sphinxAtStartPar
This will create a Fortran file called \sphinxcode{\sphinxupquote{driver\sphinxhyphen{}main\sphinxhyphen{}init.F90}}, which
can then be compiled and executed. This stand\sphinxhyphen{}alone program will read
the output file created during an execution of the actual program, call
the kernel with all required input parameter, and compare the output
variables with the original output variables. This can be used to create
stand\sphinxhyphen{}alone test cases to reproduce a bug, or for performance
optimisation of a stand\sphinxhyphen{}alone kernel.

\begin{sphinxadmonition}{warning}{Warning:}
\sphinxAtStartPar
Care has to be taken that the driver matches the version
of the code that was used to create the output file, otherwise the
driver will likely crash. The stand\sphinxhyphen{}alone driver relies on a
strict ordering of variable values in the output file and e.g.
even renaming one variable can affect this. The NetCDF version
stores the variable names and will not be able to find a variable
if its name has changed.
\end{sphinxadmonition}

\sphinxAtStartPar
The LFRic kernel driver will inline all required external modules into the
driver. It uses a \sphinxcode{\sphinxupquote{ModuleManager}} to find the required modules, based on the
assumption that a file \sphinxcode{\sphinxupquote{my\_special\_mod.f90}} will define exactly one module
called \sphinxcode{\sphinxupquote{my\_special\_mod}} (the \sphinxcode{\sphinxupquote{\_mod}} is required to be part of the
filename). The driver creator will sort the modules in the appropriate order
and add the source code directly into the driver. As a result, the driver
program is truly stand\sphinxhyphen{}alone and does not need any external dependency (the
only exception being NetCDF if the NetCDF\sphinxhyphen{}based extraction library is used).
The \sphinxcode{\sphinxupquote{ModuleManager}} uses all kernel search paths specified on the
command line (see \sphinxcode{\sphinxupquote{\sphinxhyphen{}d}} option in {\hyperref[\detokenize{psyclone_command:psyclone-command}]{\sphinxcrossref{\DUrole{std,std-ref}{The psyclone command}}}}), and it will
recursively search for all files under each path specified on the command
line.

\sphinxAtStartPar
Therefore, compilation for a created driver, e.g. the one created in
\sphinxcode{\sphinxupquote{examples/lfric/eg17/full\_example\_extract}}, is simple:

\begin{sphinxVerbatim}[commandchars=\\\{\}]
\PYG{g+go}{\PYGZdl{} gfortran \PYGZhy{}g \PYGZhy{}O0 driver\PYGZhy{}main\PYGZhy{}update.F90 \PYGZhy{}o driver\PYGZhy{}main\PYGZhy{}update}
\PYG{g+go}{\PYGZdl{} ./driver\PYGZhy{}main\PYGZhy{}update}
\PYG{g+go}{    Variable      max\PYGZus{}abs      max\PYGZus{}rel      l2\PYGZus{}diff       l2\PYGZus{}cos    identical    \PYGZsh{}rel\PYGZlt{}1E\PYGZhy{}9    \PYGZsh{}rel\PYGZlt{}1E\PYGZhy{}6    \PYGZsh{}rel\PYGZlt{}1E\PYGZhy{}3}
\PYG{g+go}{        cell .0000000E+00 .0000000E+00 .0000000E+00 .1000000E+01 .1000000E+01 .0000000E+00 .0000000E+00 .0000000E+00}
\PYG{g+go}{ field1\PYGZus{}data .0000000E+00 .0000000E+00 .0000000E+00 .1000000E+01 .5390000E+03 .0000000E+00 .0000000E+00 .0000000E+00}
\PYG{g+go}{  dummy\PYGZus{}var1 .0000000E+00 .0000000E+00 .0000000E+00 .1000000E+01 .1000000E+01 .0000000E+00 .0000000E+00 .0000000E+00}
\end{sphinxVerbatim}

\sphinxAtStartPar
(see {\hyperref[\detokenize{psyke:driver-summary-statistics}]{\sphinxcrossref{\DUrole{std,std-ref}{Driver Summary Statistics}}}} for details about the statistics\textasciigrave{}).
Note that the Makefile in the example will actually provide additional include
paths (infrastructure files and extraction library) for the compiler, but
these flags are actually only required for compiling the example program, not
for the driver.


\subsubsection{Restrictions of Kernel Extraction and Driver Creation}
\label{\detokenize{psyke:restrictions-of-kernel-extraction-and-driver-creation}}
\sphinxAtStartPar
A few restrictions still apply to the current implementation of the driver
creation code:
\begin{itemize}
\item {} 
\sphinxAtStartPar
Distributed memory is not yet supported. See \#1992.

\item {} 
\sphinxAtStartPar
The extraction code will now write variables that are used from other
modules to the kernel data file, and the driver will read these values in.
Unfortunately, if a variable is used that is defined as private,
the value cannot be written to the file, and compilation will abort.
The only solution is to modify this file and make all variables public.
This mostly affects \sphinxcode{\sphinxupquote{log\_mod.F90}}, but a few other modules as well.

\item {} 
\sphinxAtStartPar
The new build system FAB will be able to remove \sphinxcode{\sphinxupquote{private}} and
\sphinxcode{\sphinxupquote{protected}} declarations in any source files, meaning no manual
modification of files is required anymore (TODO \#2536).

\end{itemize}


\subsection{Extraction for generic Fortran}
\label{\detokenize{psyke:extraction-for-generic-fortran}}
\sphinxAtStartPar
The libraries in
\sphinxhref{https://github.com/stfc/PSyclone/tree/master/lib/extract/standalone/generic}{lib/extract/standalone/generic}
and
\sphinxhref{https://github.com/stfc/PSyclone/tree/master/lib/extract/netcdf/generic}{lib/extract/netcdf/generic}
implement the full PSyData API for use with generic code transformation.
When running the code, it will create an output file for each instrumented
code region.
The same logic for naming variables used in {\hyperref[\detokenize{psyke:extraction-for-gocean}]{\sphinxcrossref{\DUrole{std,std-ref}{Extraction for GOcean}}}}
is used here.

\begin{sphinxadmonition}{note}{Note:}
\sphinxAtStartPar
Driver creation for generic Fortran is not yet supported, and is
tracked in issue \#2058.
\end{sphinxadmonition}


\section{Driver Summary Statistics}
\label{\detokenize{psyke:driver-summary-statistics}}\label{\detokenize{psyke:id3}}
\sphinxAtStartPar
When a driver is executed, it will print summary statistics at the end
for each variable that was modified, indicating the difference between the
\sphinxtitleref{original} values from when the data file was created, and the \sphinxtitleref{new} ones
computed when executing the kernel. These differences can be caused
by changing the compilation options, or compiler version. Example output:

\begin{sphinxVerbatim}[commandchars=\\\{\}]
\PYG{g+go}{   Variable      max\PYGZus{}abs      max\PYGZus{}rel      l2\PYGZus{}diff       l2\PYGZus{}cos    identical    \PYGZsh{}rel\PYGZlt{}1E\PYGZhy{}9    \PYGZsh{}rel\PYGZlt{}1E\PYGZhy{}6    \PYGZsh{}rel\PYGZlt{}1E\PYGZhy{}3}
\PYG{g+go}{       cell .0000000E+00 .0000000E+00 .0000000E+00 .1000000E+01 .1000000E+01 .0000000E+00 .0000000E+00 .0000000E+00}
\PYG{g+go}{field1\PYGZus{}data .0000000E+00 .0000000E+00 .0000000E+00 .1000000E+01 .5390000E+03 .0000000E+00 .0000000E+00 .0000000E+00}
\PYG{g+go}{ dummy\PYGZus{}var1 .0000000E+00 .0000000E+00 .0000000E+00 .1000000E+01 .1000000E+01 .0000000E+00 .0000000E+00 .0000000E+00}
\end{sphinxVerbatim}

\sphinxAtStartPar
The columns from left to right are:
\begin{itemize}
\item {} 
\sphinxAtStartPar
The variable name.

\item {} 
\sphinxAtStartPar
The maximum absolute error of all elements.

\item {} 
\sphinxAtStartPar
The maximum relative error of all elements. If an element has the value
0, the relative error for this element is considered to be 1.0.

\item {} 
\sphinxAtStartPar
The L2 difference: \(\sqrt{\sum{(original-new)^2}}\).

\item {} 
\sphinxAtStartPar
The cosine of the angle between the two vectors: \(\frac{\sum{original*new}}{\sqrt{\sum{original*original}}*\sqrt{\sum{new*new}}}\).

\item {} 
\sphinxAtStartPar
How many values are identical.

\item {} 
\sphinxAtStartPar
How many values have a relative error of less than 10$^{\text{\sphinxhyphen{}9}}$ but are not identical.

\item {} 
\sphinxAtStartPar
How many values have a relative error of less than 10$^{\text{\sphinxhyphen{}6}}$ but more than 10$^{\text{\sphinxhyphen{}9}}$.

\item {} 
\sphinxAtStartPar
How many values have a relative error of less than 10$^{\text{\sphinxhyphen{}3}}$ but more than 10$^{\text{\sphinxhyphen{}6}}$.

\end{itemize}

\begin{sphinxadmonition}{note}{Note:}
\sphinxAtStartPar
The usefulness of the columns printed is still being evaluated. Early
indications are that the cosine of the angle between the two vectors,
which is commonly used in AI, might not be sensitive enough to give
a good indication of the differences.
\end{sphinxadmonition}

\sphinxstepscope

\begin{sphinxthebibliography}{AFH+19}
\bibitem[AFH+19]{zz_bibliography:id5}
\sphinxAtStartPar
S. V. Adams, R. W. Ford, M. Hambley, J. M. Hobson, I. Kavčič, C. M. Maynard, T. Melvin, E. H. Müller, S. Mullerworth, A. R. Porter, M. Rezny, B. J. Shipway, and R. Wong. LFRic: Meeting the challenges of scalability and performance portability in Weather and Climate models. \sphinxstyleemphasis{Journal of Parallel and Distributed Computing}, 132:383\textendash{}396, 2019. \sphinxhref{https://doi.org/https://doi.org/10.1016/j.jpdc.2019.02.007}{doi:https://doi.org/10.1016/j.jpdc.2019.02.007}.
\end{sphinxthebibliography}



\renewcommand{\indexname}{Index}
\printindex
\end{document}